\documentclass[10pt,pagesize]{article}
\usepackage{fontspec}
\usepackage{tocloft}
\usepackage[paperheight=250mm, paperwidth=139mm, margin=0cm, bottom=1cm, heightrounded, includefoot]{geometry}
%\usepackage{titlesec}
%\titleformat{\section}
%  {\normalfont\fontsize{10}{12}\bfseries}{\thesection}{0.1em}{}
%\titleformat{\subsection}
%  {\normalfont\fontsize{8}{10}\bfseries}{\thesection}{0.1em}{}
%\titlespacing*{\section}
%{0pt}{2ex plus 1ex minus .2ex}{1ex plus .2ex}
%\titlespacing*{\subsection}
%{0pt}{2ex plus 1ex minus .2ex}{1ex plus .2ex}
\usepackage[english]{babel}
\usepackage{tipa}
\usepackage[utf8]{inputenc}
\usepackage{fontspec}
\setmainfont{DejaVu Sans}
\usepackage{hyperref}
\usepackage{multicol}
\providecommand{\tightlist}{%
  \setlength{\itemsep}{0pt}\setlength{\parskip}{0pt}}
\title{pyacqriskwonksuh \\ Pyash English Dictionary.}
\author{htafdwesgina hlutli \\ Authored by Andrii Zvorygin.}
\date{2021}
\begin{document}
\maketitle
\begin{multicols}{3}
\scriptsize
\tableofcontents
\end{multicols}
\begin{multicols}{2}
\scriptsize
\section{Introduction}
This book will cover a basic introduction to the Pyash Language. Pyash is a high
precision language that is easy for homo sapiens to learn and use. 
It was designed as a computer programming language based on linguistic 
  universals.

The vocabulary was generated as follows. A list of around 30,000 English words
were translated into 28 other human languages, the words were then translated to
their phonetic transcription. Each language was given a weight based on the
  number of speakers that it had, or the number of people in the  speech family 
  that it represents.   Then a 3-4 letter word was generated by finding what the most common sounds
  were for that concept. With higher frequency words being easier to pronounce
  than lower frequency words. 

Ambigiuous and excessively borrowed terms were
  blacklisted.  As a net result any of the 30,000 or so English words
  can be mapped to the roughly 7,000 words in Pyash. The dictionary part of this
  book will help you do so.

We will start this book by reviewing the phonetics, then the prosody, then an
  introduction to the grammar. Afterwards will be the dictionary. And we'll end with an
  appendix of the various grammar classes. 


\section{Phonetics}
Mostly Pyash conforms to International Phonetic Alphabet,
except that c is the post-alveolar fricative, where the toung tip is slightly
  more inward than the alveolar ridge when you expel air to make the phoneme.
  In English the `sh' sound. 
The y is palatal approximant, where the middle of toung presses near the palet
  as one expels air. In English the `y' in `you'. The x is a velar fricative, to
  make the sound, is similar to make a k sound where the toung presses into the
  velar position, but instead it is merely near the velar position, allowing the
  air to pass without being stopped. In English
  it is the `ch' is `loch' of Scotish English, a strong `h' like sound. 

The vowels for English readers I'll explain here:
  \begin{itemize}\tightlist
    \item i -- this is the closed front unrounded vowel. In English "eee".
    \item a -- this the open central unrounded vowel. In English "ah". 
    \item u -- this is the back closed rounded vowel. In English "oo" as in "moo". 
    \item o -- this is close mid rounded vowel. In English "oh". 
    \item 6 -- this is the scwha or mid central unrounded vowel. In English "uh".
  \end{itemize}
There are also two tones:
  \begin{itemize}\tightlist
    \item 7 -- this is a high pitch tone
    \item 2 -- this a low pitch tone
  \end{itemize}
unmarked words are in medium tone. 

\subsection{Prosody}

The prosodic cadence is indicated through loudness, with primary stress on initial syllable of a word, with secondary stress on every odd syllable therafter.

\section{Grammar}
  \subsection{Word Order}
  The Pyash Word order is subject-object-verb (SOV) with postpositions 
  which is the same word order that is most natural for homo-sapiens, 
  as exemplified by various deaf groups that made up their own sign language 
  using SOV, and many of the oldest human languages using SOV. Indo-Aryan, 
  Turkic, Japanese, Astro-Asiatic, and a majority of Indigenous langauges all 
  use SOV word order. 

  Subject-Verb-Object (SVO) as in English, and other word 
  orders are typically the result of 
  language atrophy, where the case markings are lost, and people are relegated
  to using adjectives which turn into prepositions or prefixes.  While certainly
  SVO languages are still usable, they are however less optimal in terms of
  homo-sapien brain cohesion, so the SOV word order was used for Pyash. 

  In order to reduce the risk of atrophy, the verb is always at the final
  location. 

  Here is some examples of the word order in Pyash:
  {noun phrase}{verb phrase} \\
  {adjective}{noun} {adverb}{verb}\\
  {stem}{number}{gender}{context}{case} {stem}{tense}{aspect}{mood}\\

\subsection{Grammatical Mood}
Grammatical mood tells a person if the sentence is declaring a fact, a possibility, a command, a question or otherwise. 

The basic grammatical moods are:

  \begin{itemize}\tightlist
    \item li -- realis mood, for declaring a fact. In some ways it is similar to a period. For example: pyanli "love.".
    \item tu -- deontic mood, for expressing a command. For example: pyantu "love!".
    \item si  -- epistemic mood, for expressing a possibility. For example: pyansi "may love.".
    \item psi -- optatative mood, for expressing a request. For example: pyanpsi "please love.".
    \item ri -- interrogative mood, for expressing. For example: pyanri "love?"
  \end{itemize}

  \subsection{Pronouns}

Pronouns are for indicating a person.
The basic pronouns are:

  \begin{itemize}\tightlist
    \item mi -- me/I. For example: mi pyanli "I love."
    \item sa -- you/thee. For example: sa pyanli "You love."
    \item ya -- it.  For example: ya pyanli "It love."
    \item ke -- question word, who/what/where/when depending on context. ke pyanri "what love?"
  \end{itemize}

  \subsection{Grammatical Cases}
Grammatical cases are for indicating how a noun phrase is related to the verb phrase. Brief history of cases, is that in human languages cases used to be verbs or adjectives which became abbreviated over time allowing for more complicated sentences.

The basic cases are 

  \begin{itemize}\tightlist
    \item na -- which is the nominative case or subject. For example: mina pyanli "I love.".
    \item ka -- which is the accusative case or object. For example: saka pyanli "You be loved.".
    \item ti -- which is the genitive case, like "of", or "'s". For example: miti pyanka "my love".
  \end{itemize}

An example that uses both: mina saka pyanli "I love you.".

Lets add several more stems to work with. 

\begin{itemize}\tightlist
  \item kwan -- light
  \item kwin -- give
  \item tlan -- God
  \item cyit -- heart
\end{itemize}

\subsubsection{Motion}


\begin{itemize}\tightlist
  \item so -- source, roughly "from". mina tlanso li "I am from God.".
  \item te -- location, roughly "at". mina kwante li "I am at light.".
  \item ga -- way, through, roughly "via". cyitga pyanli "via heart be love.".
  \item ma -- destination, roughly "to". For example: tlanna sama pyanka kwinli "God, to you, love, gives.".
\end{itemize}

\subsubsection{Context}
Cases can also be used to express relationship within several deixis or contexts, we will cover some of the most common.

Lets add a couple more words
\begin{itemize}\tightlist
  \item tlap -- floor
  \item hrat -- roof
\end{itemize}

\begin{itemize}\tightlist
  \item ce -- under context, "under". For example: mina hratce li "I be under roof.".
  \item nu -- interior context, "in". For example: sana cyitnu keka ri "you, in heart, what?".
  \item pi -- surface context, "on". For example: mina tlappi li "I be on floor.".
  \item se -- time context, indicating something in time. For example: kese pyanri "when love?"
  \item to -- space context, indicating someting in space. For example: keto kwanri "where light?"
  \item pa -- person context, indicating a space/time or time/space person or event. For example: kepa pyanri "Who be love?"
\end{itemize}

contexts and cases can also be combined. as a general rule, start with the context, end with the case. 

Lets add another couple words:

\begin{itemize}\tightlist
  \item hkin -- go
  \item hlaf -- stay
  \item hlam -- come
  \item hcas -- house
\end{itemize}

\begin{itemize}\tightlist
  \item numa -- into. For example: mina hcasnuma hkinsi "I may go into house.".
  \item piso -- out from. For example: mina hcaspiso hlamsi "I may come out from house.".
  \item pama -- for someone. For example: tlanpama hlaftu "stay for God."
  \item pwah -- benefactive case, same as pama but shorter to say. 
  \item paso -- because of. For example: tlanpaso kwanka li "because of God, be light".
  \item ksih -- causal case, same as paso but shorter to say. 
\end{itemize}

\subsection{Number}
\subsection{Quantifiers}
\end{multicols}
\scriptsize
\begin{multicols}{2}
\scriptsize
\section{грамматика}
\subsection{грамматический настроение}
\textbf{li} [\emph{li}] 0x217E: realis mood, for declaring something, the most common kind of sentence.\\
\textbf{tu} [\emph{tu}] 0x295E: deontic mood, to indicate what should be different.\\
\textbf{cu} [\emph{ʃu}] 0x29BE: conditional mood, for if then statements. The if clause comes before the `cu' the then clause comes after.\\
\textbf{si} [\emph{si}] 0x213E: epistemic mood, to indicate that something is possible.\\
\textbf{tyah} [\emph{tjäʰ}] 0x4857: независимый пункт\\
\textbf{la} [\emph{lä}] 0x257E: dependent clause, similar in usage to `that', a sentence that is within another sentence\\
\textbf{ri} [\emph{ri}] 0x21DE: interrogative mood, for questions.\\
\textbf{di} [\emph{di}] 0x227E: директива настроение\\
\textbf{tcih} [\emph{tʃiʰ}] 0x4557: hortative mood, for pleading or begging.\\
\textbf{lo} [\emph{lo}] 0x317E: волевых настроение\\
\textbf{cyih} [\emph{ʃjiʰ}] 0x406F: совещательный настроение\\
\textbf{dwih} [\emph{dwiʰ}] 0x419F: desiderative mood, to indicate what the speaker wants to happen.\\
\textbf{pcih} [\emph{pʃiʰ}] 0x4527: imperative mood, to indicate direct commands.\\
\textbf{psih} [\emph{psiʰ}] 0x4227: optatative mood, for expressing wishes or hopes.\\
\textbf{tseh} [\emph{tseʰ}] 0x5A57: potential mood, something that is probable according to the speaker.\\
\textbf{twuh} [\emph{twuʰ}] 0x5157: dubitative mood, for saying things the speaker is doubtful of, as for expressing sarcasm.\\
\textbf{rwih} [\emph{rwiʰ}] 0x4177: precative mood, for requests.\\
\textbf{jwih} [\emph{ʒwiʰ}] 0x41AF: jussive mood, for remote commands that are to be conveyed to another recipient.\\
\textbf{mwih} [\emph{mwiʰ}] 0x410F: permissive mood, for granting permission to do something, potential response to admonitive mood.\\
\textbf{mu} [\emph{mu}] 0x283E: commissive mood, for making promises, commitments\\
\textbf{nwoh} [\emph{nwoʰ}] 0x6137: Событийный настроение\\
\textbf{lu} [\emph{lu}] 0x297E: speculative mood, for expressing speculation, hypothesis, ``could be''\\
\textbf{byih} [\emph{bjiʰ}] 0x408F: benedictive mood, for expressing blessings.\\
\textbf{kcu7h} [\emph{kʃu˥ʰ}] 0x9517: индуктивный настроение\\
\textbf{dyah} [\emph{djäʰ}] 0x489F: admonitive mood, a request for permission to do something.\\
\textbf{pceh} [\emph{pʃeʰ}] 0x5D27: apprehensive mood, for expressing fears or warnings.\\
\textbf{pruh} [\emph{pruʰ}] 0x5627: imprecative mood, for uttering curses.\\
\textbf{swuh} [\emph{swuʰ}] 0x514F: предположительный настроение\\
\textbf{ru} [\emph{ru}] 0x29DE: prohibitive mood, for expressing something is not allowed, inverse of permissive-mood-, potential response to admonitive mood\\
\textbf{ksuh} [\emph{ksuʰ}] 0x5217: declarative mood, for declaring the begining of a program or essay.\\
\textbf{t6} [\emph{tə}] 0x355E: утвердительный настроение\\
\textbf{zlih} [\emph{zliʰ}] 0x43A7: irrealis настроение\\
\textbf{n6} [\emph{nə}] 0x34DE: сенсорный доказательный настроение\\
\textbf{gyih} [\emph{gjiʰ}] 0x4097: gnomic mood, for expressing an axiom.\\
\textbf{pi7} [\emph{pi˥}] 0x409E: propositive mood, for making proposals or offers, let's\\
\textbf{la7} [\emph{lä˥}] 0x457E: successive clause, for putting several sentences in a sequence, similar to this happened, then that happened\\
\textbf{si2} [\emph{si˩}] 0x613E: necessitative настроение\\
\textbf{syi7h} [\emph{sji˥ʰ}] 0x804F: existential clause, for asserting something that can be imagined.\\
\subsection{грамматический напряженный}
\textbf{ki} [\emph{ki}] 0x205E: мимо напряженный\\
\textbf{ra} [\emph{rä}] 0x25DE: настоящее время напряженный\\
\textbf{bi} [\emph{bi}] 0x223E: будущее напряженный\\
\textbf{tleh} [\emph{tleʰ}] 0x5B57: hesternal tense, past day tense, for things that happened yesterday.\\
\textbf{cleh} [\emph{ʃleʰ}] 0x5B6F: недавний мимо напряженный\\
\textbf{dleh} [\emph{dleʰ}] 0x5B9F: hodiernal tense, for expressing today.\\
\textbf{dxah} [\emph{dxäʰ}] 0x4F9F: дистанционный пульт мимо напряженный\\
\textbf{vrih} [\emph{vriʰ}] 0x46B7: исторический напряженный\\
\textbf{s6} [\emph{sə}] 0x353E: crastinal tense, for expressing tomorrow.\\
\textbf{gweh} [\emph{gweʰ}] 0x5997: скоро будущее напряженный\\
\textbf{tli7h} [\emph{tli˥ʰ}] 0x8357: дистанционный пульт будущее напряженный\\
\subsection{грамматический аспект}
\textbf{fa} [\emph{fä}] 0x259E: perfective aspect, a completed action taken as a whole.\\
\textbf{me} [\emph{me}] 0x2C3E: imperfective aspect, an incomplete action in process.\\
\textbf{tyih} [\emph{tjiʰ}] 0x4057: ретроспективный аспект\\
\textbf{pfih} [\emph{pfiʰ}] 0x4427: progressive aspect, the action is in progress, it is happening actively. \\
\textbf{swah} [\emph{swäʰ}] 0x494F: semelfactive aspect, for expressing a sudden quick event, like a blink or a sneeze.\\
\textbf{rwah} [\emph{rwäʰ}] 0x4977: prospective aspect, the event happens after the other things under consideration. \\
\textbf{go} [\emph{go}] 0x325E: gnomic aspect, for expressing a general truth.\\
\textbf{za} [\emph{zä}] 0x269E: inchoative aspect, for expressing the begining of a state or action\\
\textbf{mweh} [\emph{mweʰ}] 0x590F: заканчивающий аспект\\
\textbf{qa} [\emph{ŋä}] 0x273E: cessative aspect, for referring to the end of an action, when it is ceasing. inchoative-aspect- is begining\\
\textbf{pseh} [\emph{pseʰ}] 0x5A27: аспект\\
\textbf{dweh} [\emph{dweʰ}] 0x599F: delimitive aspect, for referring to events of limited duration, like a limited time offer\\
\textbf{tfeh} [\emph{tfeʰ}] 0x5C57: telic aspect, for actions that have a definite end point. finite or limited.\\
\textbf{lyeh} [\emph{ljeʰ}] 0x585F: atelic aspect, for actions that have no particular end point. infinte or unlimited.\\
\textbf{mreh} [\emph{mreʰ}] 0x5E0F: momentane aspect, for referring to momentary actions, like the unfurling of flower petals in the morning\\
\textbf{xi} [\emph{xi}] 0x235E: habitual aspect, something that happens habitualy.\\
\textbf{ta2} [\emph{tä˩}] 0x655E: continuative aspect, the state is continuing to occur, it is `still' occuring. \\
\textbf{ra2} [\emph{rä˩}] 0x65DE: frequentative aspect, indicates the verb happens multiple times, for example bite frequentative-aspect is equivalent to chewing. \\
\subsection{грамматический доказательный}
\textbf{crih} [\emph{ʃriʰ}] 0x466F: видимый доказательный\\
\textbf{kwih} [\emph{kwiʰ}] 0x4117: каузальный доказательный\\
\textbf{ryih} [\emph{rjiʰ}] 0x4077: доказательный\\
\textbf{proh} [\emph{proʰ}] 0x6627: сообщается доказательный\\
\textbf{troh} [\emph{troʰ}] 0x6657: непосредственный доказательный\\
\textbf{droh} [\emph{droʰ}] 0x669F: слуховой доказательный\\
\textbf{flih} [\emph{fliʰ}] 0x4367: выведенный доказательный\\
\subsection{грамматический дело}
\textbf{na} [\emph{nä}] 0x24DE: именительный дело\\
\textbf{ka} [\emph{kä}] 0x245E: винительный дело\\
\textbf{ta} [\emph{tä}] 0x255E: topic case, the topic about which the sentence is. `regarding topic, some sentence.''. discourse-context way-case\\
\textbf{wu} [\emph{wu}] 0x28BE: vocative case, the intended recipient of the sentence. for example `O' in `O god, hear my prayers.' discourse-context destination-case\\
\textbf{ma} [\emph{mä}] 0x243E: место назначения дело\\
\textbf{yi} [\emph{ji}] 0x207E: dative case, indicates motion towards the noun phrase. destination-case. \\
\textbf{yu} [\emph{ju}] 0x287E: instrumental case, indicates the method or tool used. person-context-way-case.\\
\textbf{mwah} [\emph{mwäʰ}] 0x490F: comitative case, indicates the noun phrase is accompanying the other actors of the sentence. social-context way-case.\\
\textbf{pwah} [\emph{pwäʰ}] 0x4927: benefactive case, indicates the noun phrase is the benficiary of the sentence. social-context destination-case.\\
\textbf{swih} [\emph{swiʰ}] 0x414F: essive case, indicates a duration of time at which the noun phrase occurs. time-context way-case. \\
\textbf{da} [\emph{dä}] 0x267E: adpositional case, indicates preceding stem is the case for this noun phrase.\\
\textbf{te} [\emph{te}] 0x2D5E: место нахождения дело\\
\textbf{ga} [\emph{gä}] 0x265E: путь дело\\
\textbf{so} [\emph{so}] 0x313E: source case, indicates a source noun phrase, from. is used for constructing complex cases or by itself.\\
\textbf{pwih} [\emph{pwiʰ}] 0x4127: ablative case, indicates a source noun phrase, from. is used as a standalone case. space-context source-case\\
\textbf{twah} [\emph{twäʰ}] 0x4957: elative case, indicates motion `out of' some noun phrase. interior-context source-case.\\
\textbf{twih} [\emph{twiʰ}] 0x4157: ilative case, indicates motion `into' some noun phrase. interior-context destination-case.\\
\textbf{myah} [\emph{mjäʰ}] 0x480F: temporal case, indicates a specific time. time-context location-case. \\
\textbf{lwah} [\emph{lwäʰ}] 0x495F: perlative case, indicates motion through space. space-context way-case. \\
\textbf{lwih} [\emph{lwiʰ}] 0x415F: delative case, indicates motion away from noun phrase. surface-context source-case. \\
\textbf{ksih} [\emph{ksiʰ}] 0x4217: causal case, indicates the noun phrase is the cause. person-context source-case. \\
\textbf{tweh} [\emph{tweʰ}] 0x5957: terminative case, indicates temporal destination or goal. time-context destination-case. \\
\textbf{nweh} [\emph{nweʰ}] 0x5937: inessive case, indicates the noun phrase being in at. interior-context location-case. \\
\textbf{lweh} [\emph{lweʰ}] 0x595F: allative case, indicates motion onto a surface. surface-context destination-case. \\
\textbf{we} [\emph{we}] 0x2CBE: vialis case, indicates motion through the noun phrase. used a standalone case. way-case. \\
\textbf{twoh} [\emph{twoʰ}] 0x6157: quotative case, indicates a quote. discource-context location-case.\\
\textbf{de} [\emph{de}] 0x2E7E: distributive case, indicates for-each the noun phrase, or per the noun phrase. iterative for loops.\\
\textbf{nwah} [\emph{nwäʰ}] 0x4937: evidential case, for citing a source as evidence. discource-context source-case\\
\textbf{tloh} [\emph{tloʰ}] 0x6357: multiplicative case, for indicating how many times the event occurs. as for frequentative-aspect or simple repeat loops.\\
\textbf{nyeh} [\emph{njeʰ}] 0x5837: initiative case, indicates the begining of an action. time-context source-case\\
\textbf{sweh} [\emph{sweʰ}] 0x594F: sublative case, indicates motion toward the underside of nounphrase. under-context destination-case.\\
\textbf{txih} [\emph{txiʰ}] 0x4757: excessive case, indicates motion from a state. state-context source-case.\\
\textbf{gvih} [\emph{gviʰ}] 0x4497: pegative case, indicates motion from a person. social-context source-case.\\
\textbf{ryoh} [\emph{rjoʰ}] 0x6077: prosecutive case, indicates motion along a surface. surface-context way-case.\\
\textbf{psuh} [\emph{psuʰ}] 0x5227: superessive case, indicates location on a surface. surface-context location-case.\\
\textbf{ts6h} [\emph{tsəʰ}] 0x6A57: intransitive case, the subject of an intransitive verb.\\
\textbf{pxeh} [\emph{pxeʰ}] 0x5F27: partitive case, a subcase that indicates only partial usage of the noun phrase. lazy argument loading.\\
\textbf{pxoh} [\emph{pxoʰ}] 0x6727: prolative case, indicates motion via an interior. interior-context way-case.\\
\textbf{dxoh} [\emph{dxoʰ}] 0x679F: modal case, indicates the ability or probability of the sentence. state-context location-case. neural weight.\\
\textbf{greh} [\emph{greʰ}] 0x5E97: centric case, indicates the noun phrase being the centre. Such as a star or actor.\\
\textbf{rweh} [\emph{rweʰ}] 0x5977: oblique case, indicates intersection with noun phrase, such as an orbit or street that is being crossed at an oblique angle.\\
\textbf{xwih} [\emph{xwiʰ}] 0x41D7: patient case, the object of a transitive verb.\\
\textbf{kxeh} [\emph{kxeʰ}] 0x5F17: к дело\\
\textbf{lwoh} [\emph{lwoʰ}] 0x615F: locative case, indicates location in space. space-context location-case. \\
\textbf{bvah} [\emph{bväʰ}] 0x4C8F: abessive case, indicating the absence of the noun phrase referred to object. state-context negatory-quantifier\\
\textbf{dvah} [\emph{dväʰ}] 0x4C9F: superessive case, indicates location on a surface. surface-context location-case.\\
\textbf{je} [\emph{ʒe}] 0x2EBE: agentive case, the agent of a transitive verb.\\
\textbf{gveh} [\emph{gveʰ}] 0x5C97: эргатив дело\\
\textbf{vwih} [\emph{vwiʰ}] 0x41B7: evitative case, indicates that the noun phrase is feared or avoided.\\
\textbf{bveh} [\emph{bveʰ}] 0x5C8F: subessive case, indicates location at the underside of nounphrase. under-context location-case.\\
\textbf{tsi7h} [\emph{tsi˥ʰ}] 0x8257: antessive case, indicates the noun phrase precedes the rest of the sentence. space-context retrospective-aspect \\
\textbf{vyah} [\emph{vjäʰ}] 0x48B7: адвербиальный дело\\
\textbf{gvah} [\emph{gväʰ}] 0x4C97: lative case, indicates spatial motion towards the noun phrase. space-context destination-case. \\
\textbf{zreh} [\emph{zreʰ}] 0x5EA7: exocentric case, indicates the noun phrase orbiting the other participants. Such as a moon or audience.\\
\textbf{bu} [\emph{bu}] 0x2A3E: абсолютива дело\\
\textbf{sa7} [\emph{sä˥}] 0x453E: associative case, indicates the noun phrase is an associated group of the agent. social-context location-case.\\
\textbf{si7} [\emph{si˥}] 0x413E: superlative case, indicates motion above a surface. surface-context way-case.\\
\textbf{tsi2h} [\emph{tsi˩ʰ}] 0xC257: послелоги дело\\
\textbf{ta7} [\emph{tä˥}] 0x455E: местный направленный дело\\
\subsection{грамматический притяжательный падеж маркер}
\textbf{ti} [\emph{ti}] 0x215E: родительный дело\\
\textbf{lyah} [\emph{ljäʰ}] 0x485F: отчуждаемый владение\\
\textbf{nyih} [\emph{njiʰ}] 0x4037: неотъемлемое владение\\
\textbf{ye} [\emph{je}] 0x2C7E: притяжательный падеж маркер\\
\textbf{dzih} [\emph{dziʰ}] 0x429F: одержимый дело\\
\subsection{грамматический контекст}
\textbf{nu} [\emph{nu}] 0x28DE: интерьер контекст\\
\textbf{pa} [\emph{pä}] 0x249E: человек контекст\\
\textbf{pi} [\emph{pi}] 0x209E: поверхность контекст\\
\textbf{to} [\emph{to}] 0x315E: пространство контекст\\
\textbf{co} [\emph{ʃo}] 0x31BE: Социальное контекст\\
\textbf{se} [\emph{se}] 0x2D3E: время контекст\\
\textbf{ce} [\emph{ʃe}] 0x2DBE: под контекст\\
\textbf{ro} [\emph{ro}] 0x31DE: государство контекст\\
\textbf{xa} [\emph{xä}] 0x275E: речь контекст\\
\subsection{грамматический Пол}
\textbf{ne} [\emph{ne}] 0x2CDE: мужской Пол\\
\textbf{pyah} [\emph{pjäʰ}] 0x4827: love gender, 4th density life form, learning to forgive and love all, akin to Jesus when he was on Earth, Bodhisattva, Hatonn, Gray aliens\\
\textbf{ji} [\emph{ʒi}] 0x22BE: женский Пол\\
\textbf{nrih} [\emph{nriʰ}] 0x4637: антропный Пол\\
\textbf{bwih} [\emph{bwiʰ}] 0x418F: wise gender, 5th density life form, learning wisdom, akin to Jesus Ascended, Latwii, Arcturians, shape shifting aliens\\
\textbf{zo} [\emph{zo}] 0x329E: zoic gender, 2nd density life form, learning to fullfill desires, akin to animals, basic computers, basic robots\\
\textbf{cloh} [\emph{ʃloʰ}] 0x636F: rational gender, any entity capable of reason, 3rd density life form or higher\\
\textbf{sleh} [\emph{sleʰ}] 0x5B4F: vegetal gender, 2nd density life form, learning to fulfill desires, akin to plant, fungi, basic electronics\\
\textbf{kcuh} [\emph{kʃuʰ}] 0x9517: choice gender, 3rd density life form, learning to make choices, choosing whether to be service to self, service to others, or to remain lost, akin to basic homo sapiens, cetaceans (dolphins/whales), advanced animals\\
\textbf{gjih} [\emph{gʒiʰ}] 0x4597: galaxy gender, 11th density life form, akin to spirit of cosmos, Milky Way, Andromeda\\
\textbf{gxih} [\emph{gxiʰ}] 0x4797: общий Пол\\
\textbf{qe} [\emph{ŋe}] 0x2F3E: star gender, 10th density life form, akin to spirit of a star, Sol, Sirius\\
\textbf{qyih} [\emph{ŋjiʰ}] 0x40CF: оживить Пол\\
\textbf{xleh} [\emph{xleʰ}] 0x5BD7: vehicle gender, 7th density life form, learning about timelines, akin to spirit of a ship/vehicle/spaceship\\
\textbf{vlih} [\emph{vliʰ}] 0x43B7: mineral gender, learning about being aware, 1st density life form\\
\textbf{bxih} [\emph{bxiʰ}] 0x478F: universe gender, 12th density life form, akin to spirit of cosmos, Sophia\\
\textbf{bjeh} [\emph{bʒeʰ}] 0x5D8F: locality gender, 8th density life form, learning to be oversoul, akin to spirit of a place or flock\\
\textbf{nwuh} [\emph{nwuʰ}] 0x5137: research gender, 6th density life form, learning to blend love and wisdom, akin to Ra, Gandhi, Tesla, Prophet Mohammad (pbuh)\\
\textbf{y6} [\emph{jə}] 0x347E: художественный Пол\\
\textbf{sr6h} [\emph{srəʰ}] 0x6E4F: abstract gender, for making new genders, for example tlansr6h would be the gender of the one infinite creator God.\\
\textbf{pi2} [\emph{pi˩}] 0x609E: planetary gender, 9th density life form, akin to spirit of a planet, Gaia, Venus, Luna\\
\textbf{nri7h} [\emph{nri˥ʰ}] 0x8637: неодушевленный Пол\\
\subsection{грамматический номер}
\textbf{pu} [\emph{pu}] 0x289E: множественное число номер\\
\textbf{do} [\emph{do}] 0x327E: номер\\
\textbf{po} [\emph{po}] 0x309E: paucal номер\\
\textbf{fyih} [\emph{fjiʰ}] 0x4067: с плавающей точкой номер\\
\textbf{dluh} [\emph{dluʰ}] 0x539F: двойной номер\\
\textbf{ryuh} [\emph{rjuʰ}] 0x5077: порядковый\\
\textbf{tc6h} [\emph{tʃəʰ}] 0x6D57: целое число\\
\textbf{ve} [\emph{ve}] 0x2EDE: вектор\\
\textbf{pc6h} [\emph{pʃəʰ}] 0x6D27: испытание номер\\
\textbf{tl6h} [\emph{tləʰ}] 0x6B57: multal номер\\
\textbf{gluh} [\emph{gluʰ}] 0x5397: singulative номер\\
\subsection{грамматический местоимение}
\textbf{mi} [\emph{mi}] 0x203E: меня\\
\textbf{sa} [\emph{sä}] 0x253E: вы\\
\textbf{nyah} [\emph{njäʰ}] 0x4837: нас\\
\textbf{ya} [\emph{jä}] 0x247E: Это\\
\textbf{ksah} [\emph{ksäʰ}] 0x4A17: близость\\
\textbf{ke} [\emph{ke}] 0x2C5E: вопрос\\
\textbf{clah} [\emph{ʃläʰ}] 0x4B6F: Другие\\
\textbf{tsuh} [\emph{tsuʰ}] 0x5257: дистальный демонстративный\\
\textbf{fyoh} [\emph{fjoʰ}] 0x6067: возвратный голос\\
\subsection{грамматический эмоции}
\textbf{dyoh} [\emph{djoʰ}] 0x609F: приспособление\\
\textbf{grah} [\emph{gräʰ}] 0x4E97: восхищение\\
\textbf{zwah} [\emph{zwäʰ}] 0x49A7: амбиция\\
\textbf{nroh} [\emph{nroʰ}] 0x6637: гнев\\
\textbf{nyoh} [\emph{njoʰ}] 0x6037: анимационный\\
\textbf{jlih} [\emph{ʒliʰ}] 0x43AF: антиклимакс\\
\textbf{bzih} [\emph{bziʰ}] 0x428F: произвольный\\
\textbf{syoh} [\emph{sjoʰ}] 0x604F: благоговение\\
\textbf{myih} [\emph{mjiʰ}] 0x400F: красивая\\
\textbf{crah} [\emph{ʃräʰ}] 0x4E6F: считают,\\
\textbf{pfah} [\emph{pfäʰ}] 0x4C27: выгодный\\
\textbf{bleh} [\emph{bleʰ}] 0x5B8F: благожелательный\\
\textbf{kceh} [\emph{kʃeʰ}] 0x5D17: лучше\\
\textbf{kruh} [\emph{kruʰ}] 0x5617: горький\\
\textbf{bxah} [\emph{bxäʰ}] 0x4F8F: пресыщенный\\
\textbf{blih} [\emph{bliʰ}] 0x438F: ограниченный\\
\textbf{kroh} [\emph{kroʰ}] 0x6617: очарование\\
\textbf{clih} [\emph{ʃliʰ}] 0x436F: холодно\\
\textbf{kluh} [\emph{kluʰ}] 0x5317: путаница\\
\textbf{xyih} [\emph{xjiʰ}] 0x40D7: раскаяние\\
\textbf{gruh} [\emph{gruʰ}] 0x5697: приседать\\
\textbf{glih} [\emph{gliʰ}] 0x4397: переполненный\\
\textbf{kreh} [\emph{kreʰ}] 0x5E17: жестокий\\
\textbf{krih} [\emph{kriʰ}] 0x4617: любопытный\\
\textbf{kwah} [\emph{kwäʰ}] 0x4917: желание\\
\textbf{psah} [\emph{psäʰ}] 0x4A27: отчаяние\\
\textbf{vrah} [\emph{vräʰ}] 0x4EB7: неудовлетворенность\\
\textbf{cyuh} [\emph{ʃjuʰ}] 0x506F: сомнение\\
\textbf{kxah} [\emph{kxäʰ}] 0x4F17: экстатический\\
\textbf{zyih} [\emph{zjiʰ}] 0x40A7: элита\\
\textbf{dlah} [\emph{dläʰ}] 0x4B9F: наслаждаться\\
\textbf{tcuh} [\emph{tʃuʰ}] 0x5557: энтузиазм\\
\textbf{pxah} [\emph{pxäʰ}] 0x4F27: завлечь\\
\textbf{kxih} [\emph{kxiʰ}] 0x4717: невозмутимость\\
\textbf{zrah} [\emph{zräʰ}] 0x4EA7: истекающий\\
\textbf{croh} [\emph{ʃroʰ}] 0x666F: страх\\
\textbf{gjah} [\emph{gʒäʰ}] 0x4D97: лестный\\
\textbf{mruh} [\emph{mruʰ}] 0x560F: дурак\\
\textbf{txoh} [\emph{txoʰ}] 0x6757: скромный\\
\textbf{xrah} [\emph{xräʰ}] 0x4ED7: жадность\\
\textbf{klih} [\emph{kliʰ}] 0x4317: виновный\\
\textbf{dyih} [\emph{djiʰ}] 0x409F: ненавидеть\\
\textbf{glah} [\emph{gläʰ}] 0x4B97: здоровье\\
\textbf{xruh} [\emph{xruʰ}] 0x56D7: веселый\\
\textbf{cwih} [\emph{ʃwiʰ}] 0x416F: надеяться\\
\textbf{cweh} [\emph{ʃweʰ}] 0x596F: смиренный\\
\textbf{myuh} [\emph{mjuʰ}] 0x500F: юмор\\
\textbf{kcah} [\emph{kʃäʰ}] 0x4D17: голодный\\
\textbf{gwih} [\emph{gwiʰ}] 0x4197: невежественный\\
\textbf{dwa2h} [\emph{dwä˩ʰ}] 0xC99F: равнодушие\\
\textbf{nruh} [\emph{nruʰ}] 0x5637: невинный\\
\textbf{lyih} [\emph{ljiʰ}] 0x405F: интерес\\
\textbf{psoh} [\emph{psoʰ}] 0x6227: непобедимый\\
\textbf{rwuh} [\emph{rwuʰ}] 0x5177: ревнивый\\
\textbf{bzah} [\emph{bzäʰ}] 0x4A8F: шутить\\
\textbf{ryah} [\emph{rjäʰ}] 0x4877: смех\\
\textbf{lwuh} [\emph{lwuʰ}] 0x515F: ленивый\\
\textbf{kloh} [\emph{kloʰ}] 0x6317: Одинокий\\
\textbf{plah} [\emph{pläʰ}] 0x4B27: люблю\\
\textbf{zlah} [\emph{zläʰ}] 0x4BA7: верность\\
\textbf{fwih} [\emph{fwiʰ}] 0x4167: величественный\\
\textbf{mwoh} [\emph{mwoʰ}] 0x610F: меланхолия\\
\textbf{mroh} [\emph{mroʰ}] 0x660F: моральный\\
\textbf{gwah} [\emph{gwäʰ}] 0x4997: нервный\\
\textbf{kseh} [\emph{kseʰ}] 0x5A17: навязчивая идея\\
\textbf{vweh} [\emph{vweʰ}] 0x59B7: перегружены работой\\
\textbf{tcoh} [\emph{tʃoʰ}] 0x6557: боль\\
\textbf{sroh} [\emph{sroʰ}] 0x664F: помилование\\
\textbf{jrih} [\emph{ʒriʰ}] 0x46AF: терпение\\
\textbf{pyih} [\emph{pjiʰ}] 0x4027: мир\\
\textbf{kcih} [\emph{kʃiʰ}] 0x4517: пожалуйста\\
\textbf{qyah} [\emph{ŋjäʰ}] 0x48CF: озабоченность\\
\textbf{krah} [\emph{kräʰ}] 0x4E17: гордость\\
\textbf{pleh} [\emph{pleʰ}] 0x5B27: привилегированный\\
\textbf{kyeh} [\emph{kjeʰ}] 0x5817: пожалел\\
\textbf{tsoh} [\emph{tsoʰ}] 0x6257: уважение\\
\textbf{qrih} [\emph{ŋriʰ}] 0x46CF: мечтательность\\
\textbf{kfah} [\emph{kfäʰ}] 0x4C17: жертва\\
\textbf{srih} [\emph{sriʰ}] 0x464F: грустный\\
\textbf{drah} [\emph{dräʰ}] 0x4E9F: доволен\\
\textbf{xlah} [\emph{xläʰ}] 0x4BD7: сатир\\
\textbf{ksoh} [\emph{ksoʰ}] 0x6217: сексуальное\\
\textbf{cruh} [\emph{ʃruʰ}] 0x566F: позор\\
\textbf{cwuh} [\emph{ʃwuʰ}] 0x516F: шок\\
\textbf{xrih} [\emph{xriʰ}] 0x46D7: застенчивость\\
\textbf{slih} [\emph{sliʰ}] 0x434F: спать\\
\textbf{syah} [\emph{sjäʰ}] 0x484F: улыбка\\
\textbf{syi7h} [\emph{sji˥ʰ}] 0x804F: Социальное\\
\textbf{bwah} [\emph{bwäʰ}] 0x498F: злость\\
\textbf{sloh} [\emph{sloʰ}] 0x634F: успех\\
\textbf{syuh} [\emph{sjuʰ}] 0x504F: супер\\
\textbf{cyo2h} [\emph{ʃjo˩ʰ}] 0xE06F: подозрительный\\
\textbf{djih} [\emph{dʒiʰ}] 0x459F: симпатия\\
\textbf{kyah} [\emph{kjäʰ}] 0x4817: спасибо\\
\textbf{sruh} [\emph{sruʰ}] 0x564F: трепет\\
\textbf{kcoh} [\emph{kʃoʰ}] 0x6517: устала\\
\textbf{brah} [\emph{bräʰ}] 0x4E8F: ловушка\\
\textbf{tyoh} [\emph{tjoʰ}] 0x6057: тривиальный\\
\textbf{slah} [\emph{släʰ}] 0x4B4F: беда\\
\textbf{txuh} [\emph{txuʰ}] 0x5757: приступ боли\\
\textbf{cluh} [\emph{ʃluʰ}] 0x536F: уродливый\\
\textbf{qwah} [\emph{ŋwäʰ}] 0x49CF: неамбициозны\\
\textbf{pcuh} [\emph{pʃuʰ}] 0x5527: неопределенность\\
\textbf{mrah} [\emph{mräʰ}] 0x4E0F: Понимаю\\
\textbf{treh} [\emph{treʰ}] 0x5E57: тщеславный\\
\textbf{pxih} [\emph{pxiʰ}] 0x4727: истребование\\
\textbf{nrah} [\emph{nräʰ}] 0x4E37: тепло\\
\textbf{cyeh} [\emph{ʃjeʰ}] 0x586F: задаваться вопросом\\
\textbf{frah} [\emph{fräʰ}] 0x4E67: ранен\\
\textbf{xwah} [\emph{xwäʰ}] 0x49D7: Вау\\
\textbf{xlih} [\emph{xliʰ}] 0x43D7: несчастность\\
\subsection{грамматический стандарт Международный си единицы}
\textbf{myeh} [\emph{mjeʰ}] 0x580F: ампер\\
\textbf{gyah} [\emph{gjäʰ}] 0x4897: угол\\
\textbf{sl6h} [\emph{sləʰ}] 0x6B4F: цельсию\\
\textbf{drih} [\emph{driʰ}] 0x469F: степень\\
\textbf{xe} [\emph{xe}] 0x2F5E: частота\\
\textbf{jo} [\emph{ʒo}] 0x32BE: высокая температура\\
\textbf{creh} [\emph{ʃreʰ}] 0x5E6F: Генри\\
\textbf{xreh} [\emph{xreʰ}] 0x5ED7: герц\\
\textbf{dxih} [\emph{dxiʰ}] 0x479F: индуктивность\\
\textbf{kyoh} [\emph{kjoʰ}] 0x6017: кило\\
\textbf{kxoh} [\emph{kxoʰ}] 0x6717: килограмм\\
\textbf{m6} [\emph{mə}] 0x343E: метр\\
\textbf{tfoh} [\emph{tfoʰ}] 0x6457: моль\\
\textbf{ry6h} [\emph{rjəʰ}] 0x6877: радиан\\
\textbf{qo} [\emph{ŋo}] 0x333E: радиоактивность\\
\textbf{du} [\emph{du}] 0x2A7E: второй\\
\textbf{truh} [\emph{truʰ}] 0x5657: температура\\
\textbf{vloh} [\emph{vloʰ}] 0x63B7: вольт\\
\textbf{vlah} [\emph{vläʰ}] 0x4BB7: напряжение\\
\textbf{va} [\emph{vä}] 0x26DE: ватт\\
\textbf{blah} [\emph{bläʰ}] 0x4B8F: вес\\
\section{русский -> пяш}
\subsection{ 1 }
10: \textbf{htip} [\emph{ʰtip}] 0x4850 \\
11: \textbf{slen} [\emph{slen}] 0x6B69 \\
13: \textbf{tses} [\emph{tses}] 0x8B4A \\
14: \textbf{hses} [\emph{ʰses}] 0x8B48 \\
15: \textbf{hpet} [\emph{ʰpet}] 0xAB20 \\
18: \textbf{hdap} [\emph{ʰdäp}] 0x4998 \\
19: \textbf{dzin} [\emph{dzin}] 0x6853 \\
\subsection{ 2 }
20: \textbf{syis} [\emph{sjis}] 0x8809 \\
\subsection{ 4 }
4: \textbf{ksas} [\emph{ksäs}] 0x8942 \\
\subsection{ 5 }
5: \textbf{hfak} [\emph{ʰfäk}] 0x2960 \\
\subsection{ 8 }
8: \textbf{hwap} [\emph{ʰwäp}] 0x4928 \\
\subsection{ 9 }
9: \textbf{twun} [\emph{twun}] 0x6A2A \\
\subsection{ A }
Antwerp: \textbf{nw6p} [\emph{nwəp}] 0x4D26 \\
Archy: \textbf{cri7c} [\emph{ʃri˥ʃ}] 0xF0CD \\
Attends: \textbf{txa2t} [\emph{txä˩t}] 0xB9EA \\
\subsection{ B }
Bavaria: \textbf{bves} [\emph{bves}] 0x8B91 \\
\subsection{ C }
Cегодня: \textbf{tren} [\emph{tren}] 0x6BCA \\
\subsection{ D }
DonT: \textbf{hdo2t} [\emph{ʰdo˩t}] 0xBC98 \\
\subsection{ f }
ferous: \textbf{fr6s} [\emph{frəs}] 0x8DCC \\
friborg: \textbf{hfi2p} [\emph{ʰfi˩p}] 0x5860 \\
\subsection{ F }
FY: \textbf{fwif} [\emph{fwif}] 0xC82C \\
\subsection{ I }
IC: \textbf{hji7k} [\emph{ʰʒi˥k}] 0x30A8 \\
IST: \textbf{dxi7s} [\emph{dxi˥s}] 0x90F3 \\
\subsection{ K }
Kathmandu: \textbf{kfu7t} [\emph{kfu˥t}] 0xB282 \\
\subsection{ L }
LY: \textbf{lwi7k} [\emph{lwi˥k}] 0x302B \\
\subsection{ N }
Noman: \textbf{hn62n} [\emph{ʰnə˩n}] 0x7D30 \\
\subsection{ O }
OMA: \textbf{hpo7m} [\emph{ʰpo˥m}] 0x1420 \\
OSE: \textbf{bz6s} [\emph{bzəs}] 0x8D51 \\
\subsection{ P }
Phenicia: \textbf{hf6c} [\emph{ʰfəʃ}] 0xED60 \\
\subsection{ S }
Seychelles: \textbf{sye7s} [\emph{sje˥s}] 0x9309 \\
Sindhi: \textbf{hc67t} [\emph{ʰʃə˥t}] 0xB568 \\
\subsection{ T }
Tinder: \textbf{hte7k} [\emph{ʰte˥k}] 0x3350 \\
Transact: \textbf{tla2k} [\emph{tlä˩k}] 0x396A \\
\subsection{ U }
UAE: \textbf{druc} [\emph{druʃ}] 0xEAD3 \\
UTF-8: \textbf{tfo7f} [\emph{tfo˥f}] 0xD48A \\
\subsection{ W }
Web: \textbf{hwep} [\emph{ʰwep}] 0x4B28 \\
\subsection{ Z }
Z: \textbf{hze7c} [\emph{ʰze˥ʃ}] 0xF3A0 \\
\subsection{ А }
Ааронова: \textbf{xru2n} [\emph{xru˩n}] 0x7ADA \\
Аббасидов: \textbf{bya7s} [\emph{bjä˥s}] 0x9111 \\
аббат: \textbf{hva7t} [\emph{ʰvä˥t}] 0xB1B0 \\
аббатиса: \textbf{hba2s} [\emph{ʰbä˩s}] 0x9988 \\
аббатство: \textbf{hzap} [\emph{ʰzäp}] 0x49A0 \\
аберраций: \textbf{bri7n} [\emph{bri˥n}] 0x70D1 \\
Абиссиния: \textbf{bxa7n} [\emph{bxä˥n}] 0x71F1 \\
аболиционизм: \textbf{bjim} [\emph{bʒim}] 0x08B1 \\
аболиционисты: \textbf{byi7t} [\emph{bji˥t}] 0xB011 \\
абрикос: \textbf{vrip} [\emph{vrip}] 0x48D6 \\
абсент: \textbf{kca7n} [\emph{kʃä˥n}] 0x71A2 \\
Абстрактные: \textbf{srut} [\emph{srut}] 0xAAC9 \\
абсцисса: \textbf{hqa2s} [\emph{ʰŋä˩s}] 0x99C8 \\
абхазия: \textbf{bxi7s} [\emph{bxi˥s}] 0x90F1 \\
авантитул: \textbf{xw6t} [\emph{xwət}] 0xAD3A \\
авария: \textbf{hvak} [\emph{ʰväk}] 0x29B0 \\
Аввакум: \textbf{hxu7k} [\emph{ʰxu˥k}] 0x32D0 \\
август: \textbf{kwas} [\emph{kwäs}] 0x8922 \\
августинец: \textbf{hgo7n} [\emph{ʰgo˥n}] 0x7490 \\
Авестийский: \textbf{vya7t} [\emph{vjä˥t}] 0xB116 \\
авокадо: \textbf{vwok} [\emph{vwok}] 0x2C36 \\
Авраам: \textbf{bra7m} [\emph{brä˥m}] 0x11D1 \\
Австралия: \textbf{vras} [\emph{vräs}] 0x89D6 \\
Австрия: \textbf{vric} [\emph{vriʃ}] 0xE8D6 \\
австро-венгерской: \textbf{txu7n} [\emph{txu˥n}] 0x72EA \\
авто: \textbf{lyot} [\emph{ljot}] 0xAC0B \\
автобиографический: \textbf{tya2k} [\emph{tjä˩k}] 0x390A \\
автобус: \textbf{blos} [\emph{blos}] 0x8C71 \\
автоматизация: \textbf{tfof} [\emph{tfof}] 0xCC8A \\
автоматически: \textbf{txok} [\emph{txok}] 0x2CEA \\
автомотриса: \textbf{xlak} [\emph{xläk}] 0x297A \\
автономный: \textbf{tyi7m} [\emph{tji˥m}] 0x100A \\
Автопортрет: \textbf{tr62t} [\emph{trə˩t}] 0xBDCA \\
автор: \textbf{hlut} [\emph{ʰlut}] 0xAA58 \\
авторизоваться: \textbf{nrun} [\emph{nrun}] 0x6AC6 \\
агглютинативный: \textbf{gxi2t} [\emph{gxi˩t}] 0xB8F2 \\
агент: \textbf{hgen} [\emph{ʰgen}] 0x6B90 \\
агиография: \textbf{hx6f} [\emph{ʰxəf}] 0xCDD0 \\
агностик: \textbf{gwot} [\emph{gwot}] 0xAC32 \\
агорафобия: \textbf{qrop} [\emph{ŋrop}] 0x4CD9 \\
агрегирование: \textbf{txe2n} [\emph{txe˩n}] 0x7BEA \\
агрессивный: \textbf{gris} [\emph{gris}] 0x88D2 \\
агрономия: \textbf{gyom} [\emph{gjom}] 0x0C12 \\
ад: \textbf{crik} [\emph{ʃrik}] 0x28CD \\
адамантин: \textbf{dya7n} [\emph{djä˥n}] 0x7113 \\
адаптер: \textbf{dwap} [\emph{dwäp}] 0x4933 \\
адвентисты: \textbf{dyi2t} [\emph{dji˩t}] 0xB813 \\
адвокат: \textbf{hdac} [\emph{ʰdäʃ}] 0xE998 \\
Аделаида: \textbf{dli7t} [\emph{dli˥t}] 0xB073 \\
адиабатический: \textbf{dja2t} [\emph{dʒä˩t}] 0xB9B3 \\
администрация: \textbf{dran} [\emph{drän}] 0x69D3 \\
адмирал: \textbf{mra2n} [\emph{mrä˩n}] 0x79C1 \\
адреса: \textbf{hdi7s} [\emph{ʰdi˥s}] 0x9098 \\
Азербайджан: \textbf{zr6n} [\emph{zrən}] 0x6DD4 \\
Азия: \textbf{zyas} [\emph{zjäs}] 0x8914 \\
азот: \textbf{nrok} [\emph{nrok}] 0x2CC6 \\
аист: \textbf{xwas} [\emph{xwäs}] 0x893A \\
айва: \textbf{kxif} [\emph{kxif}] 0xC8E2 \\
аймара: \textbf{hya2m} [\emph{ʰjä˩m}] 0x1918 \\
акация: \textbf{kxaf} [\emph{kxäf}] 0xC9E2 \\
акварель: \textbf{hwa7k} [\emph{ʰwä˥k}] 0x3128 \\
аккадское: \textbf{kx67t} [\emph{kxə˥t}] 0xB5E2 \\
акклиматизировались: \textbf{kli2s} [\emph{kli˩s}] 0x9862 \\
аккорд: \textbf{kyot} [\emph{kjot}] 0xAC02 \\
аккордеоны: \textbf{kfo7n} [\emph{kfo˥n}] 0x7482 \\
Аккредитующая: \textbf{kye7n} [\emph{kje˥n}] 0x7302 \\
Аккув: \textbf{kwu7k} [\emph{kwu˥k}] 0x3222 \\
аконит: \textbf{kfi7t} [\emph{kfi˥t}] 0xB082 \\
акробатический: \textbf{kca7k} [\emph{kʃä˥k}] 0x31A2 \\
акров: \textbf{krem} [\emph{krem}] 0x0BC2 \\
акромегалия: \textbf{kfa2k} [\emph{kfä˩k}] 0x3982 \\
акрополь: \textbf{ksi2s} [\emph{ksi˩s}] 0x9842 \\
акростих: \textbf{kso2s} [\emph{kso˩s}] 0x9C42 \\
аксиллярный: \textbf{ksi2p} [\emph{ksi˩p}] 0x5842 \\
аксиоматический: \textbf{ky6t} [\emph{kjət}] 0xAD02 \\
актер: \textbf{bzat} [\emph{bzät}] 0xA951 \\
активизация: \textbf{vri7n} [\emph{vri˥n}] 0x70D6 \\
актиний: \textbf{txan} [\emph{txän}] 0x69EA \\
актуарий: \textbf{tcu7s} [\emph{tʃu˥s}] 0x92AA \\
акула: \textbf{zrak} [\emph{zräk}] 0x29D4 \\
Акупрессура: \textbf{krup} [\emph{krup}] 0x4AC2 \\
акустика: \textbf{ks6k} [\emph{ksək}] 0x2D42 \\
акушерка: \textbf{dvas} [\emph{dväs}] 0x8993 \\
акушерство: \textbf{bwa7k} [\emph{bwä˥k}] 0x3131 \\
акцент: \textbf{kson} [\emph{kson}] 0x6C42 \\
акцепторы: \textbf{syu7t} [\emph{sju˥t}] 0xB209 \\
акции: \textbf{hsoc} [\emph{ʰsoʃ}] 0xEC48 \\
акционер: \textbf{ksum} [\emph{ksum}] 0x0A42 \\
Албания: \textbf{bvap} [\emph{bväp}] 0x4991 \\
алгебра: \textbf{djap} [\emph{dʒäp}] 0x49B3 \\
алгоритм: \textbf{graf} [\emph{gräf}] 0xC9D2 \\
алебастр: \textbf{hla7m} [\emph{ʰlä˥m}] 0x1158 \\
Александр: \textbf{hla7t} [\emph{ʰlä˥t}] 0xB158 \\
алеутская: \textbf{hl62n} [\emph{ʰlə˩n}] 0x7D58 \\
Алжир: \textbf{gjic} [\emph{gʒiʃ}] 0xE8B2 \\
аликвоту: \textbf{lwa7k} [\emph{lwä˥k}] 0x312B \\
алименты: \textbf{lyaf} [\emph{ljäf}] 0xC90B \\
алифатический: \textbf{kfi7k} [\emph{kfi˥k}] 0x3082 \\
алкалоид: \textbf{kla7c} [\emph{klä˥ʃ}] 0xF162 \\
алкил: \textbf{lyi7c} [\emph{lji˥ʃ}] 0xF00B \\
алкоголик: \textbf{hd67k} [\emph{ʰdə˥k}] 0x3598 \\
аллегория: \textbf{hyi2k} [\emph{ʰji˩k}] 0x3818 \\
аллели: \textbf{lwec} [\emph{lweʃ}] 0xEB2B \\
аллергический: \textbf{gri2k} [\emph{gri˩k}] 0x38D2 \\
аллил: \textbf{cli7c} [\emph{ʃli˥ʃ}] 0xF06D \\
аллилуйя: \textbf{xlu2n} [\emph{xlu˩n}] 0x7A7A \\
аллитерация: \textbf{xr6n} [\emph{xrən}] 0x6DDA \\
алмаз: \textbf{dram} [\emph{dräm}] 0x09D3 \\
алоза: \textbf{hca2m} [\emph{ʰʃä˩m}] 0x1968 \\
алтарь: \textbf{blat} [\emph{blät}] 0xA971 \\
алфавит: \textbf{lwam} [\emph{lwäm}] 0x092B \\
алхимик: \textbf{qyas} [\emph{ŋjäs}] 0x8919 \\
алхимия: \textbf{hli2m} [\emph{ʰli˩m}] 0x1858 \\
алый: \textbf{hda2s} [\emph{ʰdä˩s}] 0x9998 \\
альбатрос: \textbf{bra7s} [\emph{brä˥s}] 0x91D1 \\
альбинос: \textbf{lwi2n} [\emph{lwi˩n}] 0x782B \\
альбом: \textbf{hlup} [\emph{ʰlup}] 0x4A58 \\
альбумин: \textbf{bvim} [\emph{bvim}] 0x0891 \\
альдегид: \textbf{dxe2t} [\emph{dxe˩t}] 0xBBF3 \\
альпака: \textbf{hfa2k} [\emph{ʰfä˩k}] 0x3960 \\
альпинизм: \textbf{pram} [\emph{präm}] 0x09C4 \\
Альпы: \textbf{hla7s} [\emph{ʰlä˥s}] 0x9158 \\
альтруизм: \textbf{twus} [\emph{twus}] 0x8A2A \\
альфа: \textbf{flaf} [\emph{fläf}] 0xC96C \\
алюминий: \textbf{lwim} [\emph{lwim}] 0x082B \\
Амазонка: \textbf{dzon} [\emph{dzon}] 0x6C53 \\
амарант: \textbf{mra2p} [\emph{mrä˩p}] 0x59C1 \\
амбиция: \textbf{hgac} [\emph{ʰgäʃ}] 0xE990 \\
амбра: \textbf{mra2s} [\emph{mrä˩s}] 0x99C1 \\
амбразура: \textbf{hba7c} [\emph{ʰbä˥ʃ}] 0xF188 \\
амбулаторный: \textbf{tce7c} [\emph{tʃe˥ʃ}] 0xF3AA \\
амеба: \textbf{vlap} [\emph{vläp}] 0x4976 \\
аменорея: \textbf{mro7c} [\emph{mro˥ʃ}] 0xF4C1 \\
Америка: \textbf{mr6k} [\emph{mrək}] 0x2DC1 \\
амид: \textbf{myi7s} [\emph{mji˥s}] 0x9001 \\
амин: \textbf{qyi7n} [\emph{ŋji˥n}] 0x7019 \\
аммиак: \textbf{hmo7n} [\emph{ʰmo˥n}] 0x7408 \\
амниотический: \textbf{mwi7k} [\emph{mwi˥k}] 0x3021 \\
аморальный: \textbf{nram} [\emph{nräm}] 0x09C6 \\
ампер: \textbf{mrep} [\emph{mrep}] 0x4BC1 \\
амперметр: \textbf{jyap} [\emph{ʒjäp}] 0x4915 \\
Амстердам: \textbf{mra2t} [\emph{mrä˩t}] 0xB9C1 \\
Амур: \textbf{hku2t} [\emph{ʰku˩t}] 0xBA10 \\
амфетамин: \textbf{hfi2m} [\emph{ʰfi˩m}] 0x1860 \\
амфибия: \textbf{fyap} [\emph{fjäp}] 0x490C \\
амфитеатр: \textbf{mya7t} [\emph{mjä˥t}] 0xB101 \\
амфора: \textbf{mwa7k} [\emph{mwä˥k}] 0x3121 \\
амхарский: \textbf{mro7k} [\emph{mro˥k}] 0x34C1 \\
ан: \textbf{hye7n} [\emph{ʰje˥n}] 0x7318 \\
анабаптист: \textbf{hba2p} [\emph{ʰbä˩p}] 0x5988 \\
анаболический: \textbf{nro7k} [\emph{nro˥k}] 0x34C6 \\
анаграмма: \textbf{nri2m} [\emph{nri˩m}] 0x18C6 \\
анаконда: \textbf{nro2n} [\emph{nro˩n}] 0x7CC6 \\
анализировать: \textbf{flis} [\emph{flis}] 0x886C \\
аналитик: \textbf{fles} [\emph{fles}] 0x8B6C \\
аналоговый: \textbf{nyok} [\emph{njok}] 0x2C06 \\
ананас: \textbf{pl6s} [\emph{pləs}] 0x8D64 \\
анархизм: \textbf{hn67s} [\emph{ʰnə˥s}] 0x9530 \\
анастомоз: \textbf{hs67s} [\emph{ʰsə˥s}] 0x9548 \\
анатомия: \textbf{nwom} [\emph{nwom}] 0x0C26 \\
анафилаксия: \textbf{fla2s} [\emph{flä˩s}] 0x996C \\
анахронизм: \textbf{from} [\emph{from}] 0x0CCC \\
анаэробный: \textbf{nwep} [\emph{nwep}] 0x4B26 \\
ангел: \textbf{htet} [\emph{ʰtet}] 0xAB50 \\
ангидрид: \textbf{zra7t} [\emph{zrä˥t}] 0xB1D4 \\
английский: \textbf{qris} [\emph{ŋris}] 0x88D9 \\
англичанка: \textbf{hqe7n} [\emph{ʰŋe˥n}] 0x73C8 \\
англосаксонской: \textbf{kso7n} [\emph{kso˥n}] 0x7442 \\
ангольский: \textbf{qwok} [\emph{ŋwok}] 0x2C39 \\
Анд: \textbf{ny67n} [\emph{njə˥n}] 0x7506 \\
Андалусия: \textbf{dluc} [\emph{dluʃ}] 0xEA73 \\
андеррайтер: \textbf{qr6n} [\emph{ŋrən}] 0x6DD9 \\
андорра: \textbf{qwot} [\emph{ŋwot}] 0xAC39 \\
Андрей: \textbf{nra7t} [\emph{nrä˥t}] 0xB1C6 \\
анекдот: \textbf{hna2k} [\emph{ʰnä˩k}] 0x3930 \\
анемия: \textbf{hxa2m} [\emph{ʰxä˩m}] 0x19D0 \\
анемометр: \textbf{qyam} [\emph{ŋjäm}] 0x0919 \\
анимационный: \textbf{hnem} [\emph{ʰnem}] 0x0B30 \\
анимизм: \textbf{ny6m} [\emph{njəm}] 0x0D06 \\
анклав: \textbf{hne2f} [\emph{ʰne˩f}] 0xDB30 \\
анналы: \textbf{nra2s} [\emph{nrä˩s}] 0x99C6 \\
аннуитет: \textbf{hna2c} [\emph{ʰnä˩ʃ}] 0xF930 \\
анод: \textbf{zlot} [\emph{zlot}] 0xAC74 \\
анонимный: \textbf{nyon} [\emph{njon}] 0x6C06 \\
анорексия: \textbf{hn62s} [\emph{ʰnə˩s}] 0x9D30 \\
антарктический: \textbf{nra2k} [\emph{nrä˩k}] 0x39C6 \\
антацидный: \textbf{tsu7s} [\emph{tsu˥s}] 0x924A \\
антенна: \textbf{ny6n} [\emph{njən}] 0x6D06 \\
антибиотик: \textbf{xyat} [\emph{xjät}] 0xA91A \\
антиген: \textbf{xw6n} [\emph{xwən}] 0x6D3A \\
антидемократические: \textbf{dji7t} [\emph{dʒi˥t}] 0xB0B3 \\
антидот: \textbf{hta7c} [\emph{ʰtä˥ʃ}] 0xF150 \\
антиклимакс: \textbf{tlus} [\emph{tlus}] 0x8A6A \\
антимонопольный: \textbf{nro7s} [\emph{nro˥s}] 0x94C6 \\
антиномия: \textbf{tli2m} [\emph{tli˩m}] 0x186A \\
антиподы: \textbf{tyo7t} [\emph{tjo˥t}] 0xB40A \\
антисептик: \textbf{qwit} [\emph{ŋwit}] 0xA839 \\
антитело: \textbf{hxip} [\emph{ʰxip}] 0x48D0 \\
антифедерализм: \textbf{fre7s} [\emph{fre˥s}] 0x93CC \\
антихрист: \textbf{tru2t} [\emph{tru˩t}] 0xBACA \\
антициклон: \textbf{two7n} [\emph{two˥n}] 0x742A \\
антропология: \textbf{nyoc} [\emph{njoʃ}] 0xEC06 \\
антропоморфизм: \textbf{nr6s} [\emph{nrəs}] 0x8DC6 \\
антропософия: \textbf{nra2f} [\emph{nrä˩f}] 0xD9C6 \\
антропоцентрический: \textbf{nwe7n} [\emph{nwe˥n}] 0x7326 \\
анус: \textbf{gyam} [\emph{gjäm}] 0x0912 \\
анчоус: \textbf{nwi7f} [\emph{nwi˥f}] 0xD026 \\
аорта: \textbf{vrom} [\emph{vrom}] 0x0CD6 \\
апеллянт: \textbf{hp62n} [\emph{ʰpə˩n}] 0x7D20 \\
аплодисменты: \textbf{pwas} [\emph{pwäs}] 0x8924 \\
Аполлон: \textbf{hpo7n} [\emph{ʰpo˥n}] 0x7420 \\
апологетика: \textbf{pli2k} [\emph{pli˩k}] 0x3864 \\
апостольский: \textbf{psi7k} [\emph{psi˥k}] 0x3044 \\
апостроф: \textbf{psa7f} [\emph{psä˥f}] 0xD144 \\
апофеоз: \textbf{pfa2s} [\emph{pfä˩s}] 0x9984 \\
аппаратный: \textbf{rwac} [\emph{rwäʃ}] 0xE92E \\
апперцепция: \textbf{pra7p} [\emph{prä˥p}] 0x51C4 \\
апрель: \textbf{pyip} [\emph{pjip}] 0x4804 \\
аптека: \textbf{fram} [\emph{främ}] 0x09CC \\
арабский: \textbf{hra7p} [\emph{ʰrä˥p}] 0x5170 \\
арак: \textbf{rya7k} [\emph{rjä˥k}] 0x310E \\
арамейский: \textbf{rya7m} [\emph{rjä˥m}] 0x110E \\
арбалет: \textbf{ksu7s} [\emph{ksu˥s}] 0x9242 \\
арбитр: \textbf{bxis} [\emph{bxis}] 0x88F1 \\
арбитраж: \textbf{hra2c} [\emph{ʰrä˩ʃ}] 0xF970 \\
Аргентина: \textbf{gjen} [\emph{gʒen}] 0x6BB2 \\
аргон: \textbf{gron} [\emph{gron}] 0x6CD2 \\
аргумент: \textbf{hrut} [\emph{ʰrut}] 0xAA70 \\
аренда: \textbf{kcet} [\emph{kʃet}] 0xABA2 \\
арендодатель: \textbf{tse2c} [\emph{tse˩ʃ}] 0xFB4A \\
ареометр: \textbf{hdo7m} [\emph{ʰdo˥m}] 0x1498 \\
арестовать: \textbf{grat} [\emph{grät}] 0xA9D2 \\
арии: \textbf{pra7s} [\emph{prä˥s}] 0x91C4 \\
арийский: \textbf{rya7n} [\emph{rjä˥n}] 0x710E \\
Ариохом: \textbf{ryo2k} [\emph{rjo˩k}] 0x3C0E \\
Аристотель: \textbf{hq6t} [\emph{ʰŋət}] 0xADC8 \\
аркебуза: \textbf{ksa7p} [\emph{ksä˥p}] 0x5142 \\
арки: \textbf{kra7n} [\emph{krä˥n}] 0x71C2 \\
арлекин: \textbf{xla2k} [\emph{xlä˩k}] 0x397A \\
армагеддон: \textbf{mye2t} [\emph{mje˩t}] 0xBB01 \\
Армения: \textbf{my6n} [\emph{mjən}] 0x6D01 \\
армия: \textbf{hyim} [\emph{ʰjim}] 0x0818 \\
ароматный: \textbf{hgun} [\emph{ʰgun}] 0x6A90 \\
артериосклероз: \textbf{tli2s} [\emph{tli˩s}] 0x986A \\
артефакт: \textbf{hr6t} [\emph{ʰrət}] 0xAD70 \\
артиллерия: \textbf{hrap} [\emph{ʰräp}] 0x4970 \\
Артишок: \textbf{tc6k} [\emph{tʃək}] 0x2DAA \\
архаизм: \textbf{kci7s} [\emph{kʃi˥s}] 0x90A2 \\
археологический: \textbf{kluk} [\emph{kluk}] 0x2A62 \\
архивариус: \textbf{hra2f} [\emph{ʰrä˩f}] 0xD970 \\
архивируются: \textbf{gva2t} [\emph{gvä˩t}] 0xB992 \\
архиепископ: \textbf{hra7k} [\emph{ʰrä˥k}] 0x3170 \\
архитектура: \textbf{hruk} [\emph{ʰruk}] 0x2A70 \\
архитрав: \textbf{jraf} [\emph{ʒräf}] 0xC9D5 \\
арьергард: \textbf{rwa2t} [\emph{rwä˩t}] 0xB92E \\
асептика: \textbf{pse2p} [\emph{pse˩p}] 0x5B44 \\
асимптотический: \textbf{myi2k} [\emph{mji˩k}] 0x3801 \\
аспект: \textbf{hsek} [\emph{ʰsek}] 0x2B48 \\
Ассамский: \textbf{ksa7m} [\emph{ksä˥m}] 0x1142 \\
астения: \textbf{fwa2n} [\emph{fwä˩n}] 0x792C \\
астигматизм: \textbf{swa7m} [\emph{swä˥m}] 0x1129 \\
астма: \textbf{dzam} [\emph{dzäm}] 0x0953 \\
астральный: \textbf{sra7c} [\emph{srä˥ʃ}] 0xF1C9 \\
астрология: \textbf{tli2c} [\emph{tli˩ʃ}] 0xF86A \\
астронавт: \textbf{xran} [\emph{xrän}] 0x69DA \\
астроном: \textbf{xlam} [\emph{xläm}] 0x097A \\
астрономия: \textbf{tyon} [\emph{tjon}] 0x6C0A \\
Астурийский: \textbf{sru7s} [\emph{sru˥s}] 0x92C9 \\
асфиксия: \textbf{sya2k} [\emph{sjä˩k}] 0x3909 \\
атавизм: \textbf{tsi7m} [\emph{tsi˥m}] 0x104A \\
атака: \textbf{nyat} [\emph{njät}] 0xA906 \\
атеизм: \textbf{tsem} [\emph{tsem}] 0x0B4A \\
Атлантика: \textbf{tl6t} [\emph{tlət}] 0xAD6A \\
атлас: \textbf{tla2s} [\emph{tlä˩s}] 0x996A \\
атолл: \textbf{kca2t} [\emph{kʃä˩t}] 0xB9A2 \\
атом: \textbf{kram} [\emph{kräm}] 0x09C2 \\
атрибут: \textbf{trup} [\emph{trup}] 0x4ACA \\
атропин: \textbf{tra7p} [\emph{trä˥p}] 0x51CA \\
атрофированный: \textbf{tfi7f} [\emph{tfi˥f}] 0xD08A \\
аудио: \textbf{dyim} [\emph{djim}] 0x0813 \\
аукцион: \textbf{glam} [\emph{gläm}] 0x0972 \\
аутизм: \textbf{ht6s} [\emph{ʰtəs}] 0x8D50 \\
афазия: \textbf{hfi7c} [\emph{ʰfi˥ʃ}] 0xF060 \\
Афганистан: \textbf{hf6n} [\emph{ʰfən}] 0x6D60 \\
афелий: \textbf{flo7n} [\emph{flo˥n}] 0x746C \\
Афины: \textbf{hte7n} [\emph{ʰte˥n}] 0x7350 \\
афоризмы: \textbf{fri7m} [\emph{fri˥m}] 0x10CC \\
Африка: \textbf{fruk} [\emph{fruk}] 0x2ACC \\
афроамериканец: \textbf{hri7k} [\emph{ʰri˥k}] 0x3070 \\
афродизиак: \textbf{rwa2k} [\emph{rwä˩k}] 0x392E \\
афферентный: \textbf{fr62n} [\emph{frə˩n}] 0x7DCC \\
аффидавит: \textbf{xyap} [\emph{xjäp}] 0x491A \\
Ахан: \textbf{dja7n} [\emph{dʒä˥n}] 0x71B3 \\
Ахиллес: \textbf{xlic} [\emph{xliʃ}] 0xE87A \\
ахнула: \textbf{hxa7t} [\emph{ʰxä˥t}] 0xB1D0 \\
ахроматический: \textbf{ryos} [\emph{rjos}] 0x8C0E \\
ацетальдегид: \textbf{tla7f} [\emph{tlä˥f}] 0xD16A \\
ацетат: \textbf{sye2t} [\emph{sje˩t}] 0xBB09 \\
ацетил: \textbf{syi7c} [\emph{sji˥ʃ}] 0xF009 \\
ацетилен: \textbf{sri7n} [\emph{sri˥n}] 0x70C9 \\
ацетон: \textbf{hqo7n} [\emph{ʰŋo˥n}] 0x74C8 \\
ацтекский: \textbf{hza2k} [\emph{ʰzä˩k}] 0x39A0 \\
ачехский: \textbf{kca2s} [\emph{kʃä˩s}] 0x99A2 \\
аэрация: \textbf{vyon} [\emph{vjon}] 0x6C16 \\
аэробный: \textbf{byop} [\emph{bjop}] 0x4C11 \\
аэродинамический: \textbf{hro2m} [\emph{ʰro˩m}] 0x1C70 \\
аэрозоль: \textbf{rwos} [\emph{rwos}] 0x8C2E \\
аэропорт: \textbf{fwac} [\emph{fwäʃ}] 0xE92C \\
\subsection{ б }
бабочка: \textbf{brip} [\emph{brip}] 0x48D1 \\
\subsection{ Б }
Багаж: \textbf{blac} [\emph{bläʃ}] 0xE971 \\
багамы: \textbf{bx6m} [\emph{bxəm}] 0x0DF1 \\
Багдад: \textbf{bxa7t} [\emph{bxä˥t}] 0xB1F1 \\
база: \textbf{hbis} [\emph{ʰbis}] 0x8888 \\
базел: \textbf{bza7s} [\emph{bzä˥s}] 0x9151 \\
байпас: \textbf{hba2f} [\emph{ʰbä˩f}] 0xD988 \\
байт: \textbf{bret} [\emph{bret}] 0xABD1 \\
бак: \textbf{tlok} [\emph{tlok}] 0x2C6A \\
бакалавр: \textbf{hba2c} [\emph{ʰbä˩ʃ}] 0xF988 \\
бакелит: \textbf{ble2t} [\emph{ble˩t}] 0xBB71 \\
баклажан: \textbf{bzet} [\emph{bzet}] 0xAB51 \\
бактерии: \textbf{byic} [\emph{bjiʃ}] 0xE811 \\
балеарский: \textbf{blep} [\emph{blep}] 0x4B71 \\
балет: \textbf{blek} [\emph{blek}] 0x2B71 \\
балканский: \textbf{hba7k} [\emph{ʰbä˥k}] 0x3188 \\
балки: \textbf{hbi2m} [\emph{ʰbi˩m}] 0x1888 \\
балкон: \textbf{blak} [\emph{bläk}] 0x2971 \\
балласт: \textbf{bl6t} [\emph{blət}] 0xAD71 \\
баллистический: \textbf{bla7k} [\emph{blä˥k}] 0x3171 \\
баллонах: \textbf{bzot} [\emph{bzot}] 0xAC51 \\
баловать: \textbf{flot} [\emph{flot}] 0xAC6C \\
балтийский: \textbf{bl6k} [\emph{blək}] 0x2D71 \\
бальзам: \textbf{bxam} [\emph{bxäm}] 0x09F1 \\
бальзамировать: \textbf{hc6m} [\emph{ʰʃəm}] 0x0D68 \\
бальный: \textbf{blam} [\emph{bläm}] 0x0971 \\
балюстрада: \textbf{bla2s} [\emph{blä˩s}] 0x9971 \\
бамбук: \textbf{hbup} [\emph{ʰbup}] 0x4A88 \\
бампер: \textbf{hb6p} [\emph{ʰbəp}] 0x4D88 \\
банан: \textbf{hzam} [\emph{ʰzäm}] 0x09A0 \\
Бангладеш: \textbf{blec} [\emph{bleʃ}] 0xEB71 \\
банджо: \textbf{bjoc} [\emph{bʒoʃ}] 0xECB1 \\
бандит: \textbf{xyo7n} [\emph{xjo˥n}] 0x741A \\
банка: \textbf{trim} [\emph{trim}] 0x08CA \\
банка: \textbf{pcun} [\emph{pʃun}] 0x6AA4 \\
банкноты: \textbf{byif} [\emph{bjif}] 0xC811 \\
Банкомат: \textbf{hge7m} [\emph{ʰge˥m}] 0x1390 \\
банкрот: \textbf{jyat} [\emph{ʒjät}] 0xA915 \\
баннер: \textbf{hbes} [\emph{ʰbes}] 0x8B88 \\
барак: \textbf{bra2k} [\emph{brä˩k}] 0x39D1 \\
баратрия: \textbf{bri2k} [\emph{bri˩k}] 0x38D1 \\
барий: \textbf{brim} [\emph{brim}] 0x08D1 \\
баритон: \textbf{bro7n} [\emph{bro˥n}] 0x74D1 \\
бармен: \textbf{bya7n} [\emph{bjä˥n}] 0x7111 \\
барокко: \textbf{bros} [\emph{bros}] 0x8CD1 \\
барометр: \textbf{bxet} [\emph{bxet}] 0xABF1 \\
баронесса: \textbf{bra7n} [\emph{brä˥n}] 0x71D1 \\
баррель: \textbf{brak} [\emph{bräk}] 0x29D1 \\
баррикада: \textbf{bruk} [\emph{bruk}] 0x2AD1 \\
барсук: \textbf{bra7k} [\emph{brä˥k}] 0x31D1 \\
бархат: \textbf{hjet} [\emph{ʰʒet}] 0xABA8 \\
баскский: \textbf{psa7k} [\emph{psä˥k}] 0x3144 \\
бассейн: \textbf{hvun} [\emph{ʰvun}] 0x6AB0 \\
батарея: \textbf{bric} [\emph{briʃ}] 0xE8D1 \\
батик: \textbf{bja7k} [\emph{bʒä˥k}] 0x31B1 \\
батут: \textbf{hr62n} [\emph{ʰrə˩n}] 0x7D70 \\
Бахрейн: \textbf{bre7n} [\emph{bre˥n}] 0x73D1 \\
бацилла: \textbf{bla7s} [\emph{blä˥s}] 0x9171 \\
бачки: \textbf{hqa7c} [\emph{ʰŋä˥ʃ}] 0xF1C8 \\
бдение: \textbf{vyic} [\emph{vjiʃ}] 0xE816 \\
беатификации: \textbf{br6k} [\emph{brək}] 0x2DD1 \\
бег: \textbf{gyok} [\emph{gjok}] 0x2C12 \\
Бег: \textbf{slac} [\emph{släʃ}] 0xE969 \\
бегемот: \textbf{xyom} [\emph{xjom}] 0x0C1A \\
беглый: \textbf{fluc} [\emph{fluʃ}] 0xEA6C \\
беда: \textbf{hmaf} [\emph{ʰmäf}] 0xC908 \\
бедность: \textbf{dj6n} [\emph{dʒən}] 0x6DB3 \\
бедные: \textbf{hyis} [\emph{ʰjis}] 0x8818 \\
беженец: \textbf{fyim} [\emph{fjim}] 0x080C \\
безалкогольный: \textbf{zlok} [\emph{zlok}] 0x2C74 \\
безводный: \textbf{nwi2k} [\emph{nwi˩k}] 0x3826 \\
безвредный: \textbf{nrac} [\emph{nräʃ}] 0xE9C6 \\
безвременно: \textbf{hta2m} [\emph{ʰtä˩m}] 0x1950 \\
безделушка: \textbf{jyok} [\emph{ʒjok}] 0x2C15 \\
бездомный: \textbf{hwom} [\emph{ʰwom}] 0x0C28 \\
бездонный: \textbf{hja2n} [\emph{ʰʒä˩n}] 0x79A8 \\
Беззаботная: \textbf{hxu7n} [\emph{ʰxu˥n}] 0x72D0 \\
безногий: \textbf{zles} [\emph{zles}] 0x8B74 \\
безопасно: \textbf{syam} [\emph{sjäm}] 0x0909 \\
безоружный: \textbf{hno7s} [\emph{ʰno˥s}] 0x9430 \\
безработные: \textbf{byet} [\emph{bjet}] 0xAB11 \\
безрассудный: \textbf{hlem} [\emph{ʰlem}] 0x0B58 \\
безрукавный: \textbf{bvis} [\emph{bvis}] 0x8891 \\
безумие: \textbf{hga7n} [\emph{ʰgä˥n}] 0x7190 \\
безумная: \textbf{dwa7n} [\emph{dwä˥n}] 0x7133 \\
безутешный: \textbf{kli2p} [\emph{kli˩p}] 0x5862 \\
Бейрут: \textbf{bru2t} [\emph{bru˩t}] 0xBAD1 \\
бейсбол: \textbf{hbep} [\emph{ʰbep}] 0x4B88 \\
бекас: \textbf{hsi2p} [\emph{ʰsi˩p}] 0x5848 \\
бекон: \textbf{bjon} [\emph{bʒon}] 0x6CB1 \\
Беларусь: \textbf{dlus} [\emph{dlus}] 0x8A73 \\
Белград: \textbf{ble7t} [\emph{ble˥t}] 0xB371 \\
белка: \textbf{sluk} [\emph{sluk}] 0x2A69 \\
белок: \textbf{pwot} [\emph{pwot}] 0xAC24 \\
белый: \textbf{hpit} [\emph{ʰpit}] 0xA820 \\
Бельгия: \textbf{blem} [\emph{blem}] 0x0B71 \\
белье: \textbf{lyen} [\emph{ljen}] 0x6B0B \\
бельевой: \textbf{lya2s} [\emph{ljä˩s}] 0x990B \\
Бенгалия: \textbf{hba2k} [\emph{ʰbä˩k}] 0x3988 \\
Бенджамен: \textbf{hni7m} [\emph{ʰni˥m}] 0x1030 \\
бенди: \textbf{bxi7t} [\emph{bxi˥t}] 0xB0F1 \\
Бенедикт: \textbf{hne7t} [\emph{ʰne˥t}] 0xB330 \\
бензин: \textbf{glis} [\emph{glis}] 0x8872 \\
бензойный: \textbf{bzok} [\emph{bzok}] 0x2C51 \\
бензол: \textbf{bzen} [\emph{bzen}] 0x6B51 \\
берберский: \textbf{br6p} [\emph{brəp}] 0x4DD1 \\
бергамот: \textbf{hbo7m} [\emph{ʰbo˥m}] 0x1488 \\
берег: \textbf{nwo7t} [\emph{nwo˥t}] 0xB426 \\
берег: \textbf{psic} [\emph{psiʃ}] 0xE844 \\
береза: \textbf{bwut} [\emph{bwut}] 0xAA31 \\
беременна: \textbf{gzit} [\emph{gzit}] 0xA852 \\
беременность: \textbf{hges} [\emph{ʰges}] 0x8B90 \\
бериллий: \textbf{blim} [\emph{blim}] 0x0871 \\
Берлин: \textbf{bli2n} [\emph{bli˩n}] 0x7871 \\
бермудский: \textbf{byu2t} [\emph{bju˩t}] 0xBA11 \\
бесплодие: \textbf{blup} [\emph{blup}] 0x4A71 \\
бесплодный: \textbf{xrun} [\emph{xrun}] 0x6ADA \\
беспозвоночный: \textbf{vrut} [\emph{vrut}] 0xAAD6 \\
беспокоит: \textbf{hxe7t} [\emph{ʰxe˥t}] 0xB3D0 \\
бесполый: \textbf{hqec} [\emph{ʰŋeʃ}] 0xEBC8 \\
бесправных: \textbf{dra2c} [\emph{drä˩ʃ}] 0xF9D3 \\
беспрецедентный: \textbf{bxon} [\emph{bxon}] 0x6CF1 \\
беспризорник: \textbf{cli7t} [\emph{ʃli˥t}] 0xB06D \\
беспричинный: \textbf{ryi7t} [\emph{rji˥t}] 0xB00E \\
беспроводной: \textbf{bris} [\emph{bris}] 0x88D1 \\
беспутный: \textbf{hsu2f} [\emph{ʰsu˩f}] 0xDA48 \\
бессменный: \textbf{tca2s} [\emph{tʃä˩s}] 0x99AA \\
бессмертный: \textbf{hmi7t} [\emph{ʰmi˥t}] 0xB008 \\
бессовестный: \textbf{dvam} [\emph{dväm}] 0x0993 \\
бесспорно: \textbf{htu7p} [\emph{ʰtu˥p}] 0x5250 \\
бесспорный: \textbf{hna7p} [\emph{ʰnä˥p}] 0x5130 \\
бестселлер: \textbf{sya2p} [\emph{sjä˩p}] 0x5909 \\
бесчувственный: \textbf{hye7s} [\emph{ʰje˥s}] 0x9318 \\
бешено: \textbf{kxe2n} [\emph{kxe˩n}] 0x7BE2 \\
библиотека: \textbf{htek} [\emph{ʰtek}] 0x2B50 \\
библиофил: \textbf{twu7k} [\emph{twu˥k}] 0x322A \\
библия: \textbf{byip} [\emph{bjip}] 0x4811 \\
бигуди: \textbf{hki2t} [\emph{ʰki˩t}] 0xB810 \\
бизон: \textbf{bzon} [\emph{bzon}] 0x6C51 \\
бикарбонат: \textbf{qra7n} [\emph{ŋrä˥n}] 0x71D9 \\
бикини: \textbf{bwis} [\emph{bwis}] 0x8831 \\
Билеты: \textbf{txe2t} [\emph{txe˩t}] 0xBBEA \\
билирубин: \textbf{bru7p} [\emph{bru˥p}] 0x52D1 \\
Билли: \textbf{bli7t} [\emph{bli˥t}] 0xB071 \\
било: \textbf{hbi2k} [\emph{ʰbi˩k}] 0x3888 \\
бильярд: \textbf{bla2t} [\emph{blä˩t}] 0xB971 \\
биография: \textbf{hbaf} [\emph{ʰbäf}] 0xC988 \\
биологическая: \textbf{blok} [\emph{blok}] 0x2C71 \\
биом: \textbf{byom} [\emph{bjom}] 0x0C11 \\
биометрия: \textbf{bxe7t} [\emph{bxe˥t}] 0xB3F1 \\
биоразнообразия: \textbf{qyin} [\emph{ŋjin}] 0x6819 \\
биотехнология: \textbf{byok} [\emph{bjok}] 0x2C11 \\
биофизика: \textbf{bli7k} [\emph{bli˥k}] 0x3071 \\
биохимия: \textbf{bwim} [\emph{bwim}] 0x0831 \\
бирюзовый: \textbf{fros} [\emph{fros}] 0x8CCC \\
бисексуальность: \textbf{bwi7t} [\emph{bwi˥t}] 0xB031 \\
бисквит: \textbf{hbuk} [\emph{ʰbuk}] 0x2A88 \\
биты: \textbf{hbi7t} [\emph{ʰbi˥t}] 0xB088 \\
бицепс: \textbf{gxip} [\emph{gxip}] 0x48F2 \\
благоговейный: \textbf{vr6n} [\emph{vrən}] 0x6DD6 \\
благоговение: \textbf{ryot} [\emph{rjot}] 0xAC0E \\
благожелательный: \textbf{hbec} [\emph{ʰbeʃ}] 0xEB88 \\
благоприятствовать: \textbf{pfak} [\emph{pfäk}] 0x2984 \\
благородный: \textbf{hkif} [\emph{ʰkif}] 0xC810 \\
благословить: \textbf{hbik} [\emph{ʰbik}] 0x2888 \\
благочестивый: \textbf{tsi2t} [\emph{tsi˩t}] 0xB84A \\
блаженный: \textbf{hle2t} [\emph{ʰle˩t}] 0xBB58 \\
блаженство: \textbf{bruf} [\emph{bruf}] 0xCAD1 \\
бланшировать: \textbf{bxa2n} [\emph{bxä˩n}] 0x79F1 \\
бледность: \textbf{hpa7p} [\emph{ʰpä˥p}] 0x5120 \\
блестящий: \textbf{blen} [\emph{blen}] 0x6B71 \\
близнец: \textbf{swam} [\emph{swäm}] 0x0929 \\
Близнецы: \textbf{jwis} [\emph{ʒwis}] 0x8835 \\
близорукость: \textbf{hzi7s} [\emph{ʰzi˥s}] 0x90A0 \\
близость: \textbf{hcat} [\emph{ʰʃät}] 0xA968 \\
блин: \textbf{pces} [\emph{pʃes}] 0x8BA4 \\
блог: \textbf{glok} [\emph{glok}] 0x2C72 \\
блок-схема: \textbf{hvot} [\emph{ʰvot}] 0xACB0 \\
блокада: \textbf{hbo7k} [\emph{ʰbo˥k}] 0x3488 \\
блондинка: \textbf{hvaf} [\emph{ʰväf}] 0xC9B0 \\
блюдо: \textbf{pla2t} [\emph{plä˩t}] 0xB964 \\
боб: \textbf{byap} [\emph{bjäp}] 0x4911 \\
бобр: \textbf{bvip} [\emph{bvip}] 0x4891 \\
Бог: \textbf{tlan} [\emph{tlän}] 0x696A \\
богатые: \textbf{hfin} [\emph{ʰfin}] 0x6860 \\
Богемия: \textbf{bxem} [\emph{bxem}] 0x0BF1 \\
богомол: \textbf{gzas} [\emph{gzäs}] 0x8952 \\
божествами: \textbf{hde2t} [\emph{ʰde˩t}] 0xBB98 \\
бозон: \textbf{bzos} [\emph{bzos}] 0x8C51 \\
бойкотировать: \textbf{byik} [\emph{bjik}] 0x2811 \\
бойня: \textbf{bzaf} [\emph{bzäf}] 0xC951 \\
боковая: \textbf{bras} [\emph{bräs}] 0x89D1 \\
боксеров: \textbf{bwos} [\emph{bwos}] 0x8C31 \\
Болгария: \textbf{blif} [\emph{blif}] 0xC871 \\
болезнь: \textbf{hdik} [\emph{ʰdik}] 0x2898 \\
болеутоляющий: \textbf{gya7n} [\emph{gjä˥n}] 0x7112 \\
Боливия: \textbf{blof} [\emph{blof}] 0xCC71 \\
Болонья: \textbf{hlo2n} [\emph{ʰlo˩n}] 0x7C58 \\
болото: \textbf{swot} [\emph{swot}] 0xAC29 \\
болт: \textbf{blot} [\emph{blot}] 0xAC71 \\
боль: \textbf{tcuk} [\emph{tʃuk}] 0x2AAA \\
больверк: \textbf{hve7n} [\emph{ʰve˥n}] 0x73B0 \\
больница: \textbf{fwat} [\emph{fwät}] 0xA92C \\
больной: \textbf{syim} [\emph{sjim}] 0x0809 \\
большинство: \textbf{htos} [\emph{ʰtos}] 0x8C50 \\
большой: \textbf{bya2n} [\emph{bjä˩n}] 0x7911 \\
бомбардировка: \textbf{hbof} [\emph{ʰbof}] 0xCC88 \\
Бомбей: \textbf{hba7p} [\emph{ʰbä˥p}] 0x5188 \\
бондарь: \textbf{xwup} [\emph{xwup}] 0x4A3A \\
бор: \textbf{brop} [\emph{brop}] 0x4CD1 \\
бордель: \textbf{byef} [\emph{bjef}] 0xCB11 \\
борзая: \textbf{hgu7n} [\emph{ʰgu˥n}] 0x7290 \\
бормотание: \textbf{bram} [\emph{bräm}] 0x09D1 \\
борода: \textbf{drit} [\emph{drit}] 0xA8D3 \\
бородавка: \textbf{hyo2t} [\emph{ʰjo˩t}] 0xBC18 \\
бородка: \textbf{hga2c} [\emph{ʰgä˩ʃ}] 0xF990 \\
борозда: \textbf{kroc} [\emph{kroʃ}] 0xECC2 \\
бороться: \textbf{lyut} [\emph{ljut}] 0xAA0B \\
борьба: \textbf{hnep} [\emph{ʰnep}] 0x4B30 \\
босния: \textbf{byon} [\emph{bjon}] 0x6C11 \\
ботаника: \textbf{vwin} [\emph{vwin}] 0x6836 \\
ботик: \textbf{htu2c} [\emph{ʰtu˩ʃ}] 0xFA50 \\
Ботсвана: \textbf{twof} [\emph{twof}] 0xCC2A \\
боулинг: \textbf{bloc} [\emph{bloʃ}] 0xEC71 \\
боцман: \textbf{bwe7n} [\emph{bwe˥n}] 0x7331 \\
боярышник: \textbf{xwa2n} [\emph{xwä˩n}] 0x793A \\
Бразилия: \textbf{bzic} [\emph{bziʃ}] 0xE851 \\
брамин: \textbf{py67t} [\emph{pjə˥t}] 0xB504 \\
брандмауэр: \textbf{fraf} [\emph{fräf}] 0xC9CC \\
бранить: \textbf{tca2n} [\emph{tʃä˩n}] 0x79AA \\
браслет: \textbf{bles} [\emph{bles}] 0x8B71 \\
брат: \textbf{pyak} [\emph{pjäk}] 0x2904 \\
братва: \textbf{hya7t} [\emph{ʰjä˥t}] 0xB118 \\
братоубийство: \textbf{fra2t} [\emph{frä˩t}] 0xB9CC \\
браузер: \textbf{bwas} [\emph{bwäs}] 0x8931 \\
брахман: \textbf{hb6m} [\emph{ʰbəm}] 0x0D88 \\
брачный: \textbf{bxi2n} [\emph{bxi˩n}] 0x78F1 \\
бревна: \textbf{hlo2k} [\emph{ʰlo˩k}] 0x3C58 \\
бремен: \textbf{hre7m} [\emph{ʰre˥m}] 0x1370 \\
бренди: \textbf{hb6t} [\emph{ʰbət}] 0xAD88 \\
бретонский: \textbf{hbe7n} [\emph{ʰbe˥n}] 0x7388 \\
брикет: \textbf{hbe7k} [\emph{ʰbe˥k}] 0x3388 \\
Британия: \textbf{hbi7n} [\emph{ʰbi˥n}] 0x7088 \\
бритва: \textbf{zyat} [\emph{zjät}] 0xA914 \\
бровь: \textbf{brec} [\emph{breʃ}] 0xEBD1 \\
бродить: \textbf{bjan} [\emph{bʒän}] 0x69B1 \\
брокколи: \textbf{bwok} [\emph{bwok}] 0x2C31 \\
бром: \textbf{brom} [\emph{brom}] 0x0CD1 \\
броненосец: \textbf{rwo2t} [\emph{rwo˩t}] 0xBC2E \\
бронх: \textbf{bwo7n} [\emph{bwo˥n}] 0x7431 \\
бронхит: \textbf{bxi7n} [\emph{bxi˥n}] 0x70F1 \\
бросать: \textbf{tsas} [\emph{tsäs}] 0x894A \\
Бруней: \textbf{bl6n} [\emph{blən}] 0x6D71 \\
брюки: \textbf{tlun} [\emph{tlun}] 0x6A6A \\
Брюссель: \textbf{hru7s} [\emph{ʰru˥s}] 0x9270 \\
брюшина: \textbf{hre2m} [\emph{ʰre˩m}] 0x1B70 \\
бубен: \textbf{dvon} [\emph{dvon}] 0x6C93 \\
бубонный: \textbf{byu7n} [\emph{bju˥n}] 0x7211 \\
Будапешт: \textbf{dyus} [\emph{djus}] 0x8A13 \\
буддизм: \textbf{bxum} [\emph{bxum}] 0x0AF1 \\
будуар: \textbf{bwaf} [\emph{bwäf}] 0xC931 \\
будущее: \textbf{plic} [\emph{pliʃ}] 0xE864 \\
буквально: \textbf{sram} [\emph{sräm}] 0x09C9 \\
букет: \textbf{bwuk} [\emph{bwuk}] 0x2A31 \\
букмекер: \textbf{bvok} [\emph{bvok}] 0x2C91 \\
булинь: \textbf{blu2n} [\emph{blu˩n}] 0x7A71 \\
бульдозер: \textbf{blu7t} [\emph{blu˥t}] 0xB271 \\
бульон: \textbf{hbu2n} [\emph{ʰbu˩n}] 0x7A88 \\
бумага: \textbf{kyis} [\emph{kjis}] 0x8802 \\
бумажник: \textbf{blap} [\emph{bläp}] 0x4971 \\
бумеранг: \textbf{bru7n} [\emph{bru˥n}] 0x72D1 \\
бунт: \textbf{bzan} [\emph{bzän}] 0x6951 \\
бунтарь: \textbf{frac} [\emph{fräʃ}] 0xE9CC \\
бур: \textbf{mro2n} [\emph{mro˩n}] 0x7CC1 \\
буравчик: \textbf{gri2t} [\emph{gri˩t}] 0xB8D2 \\
бурдюк: \textbf{hwe2k} [\emph{ʰwe˩k}] 0x3B28 \\
Бурунди: \textbf{br6t} [\emph{brət}] 0xADD1 \\
буря: \textbf{hfup} [\emph{ʰfup}] 0x4A60 \\
Бутан: \textbf{bwun} [\emph{bwun}] 0x6A31 \\
бутил: \textbf{byuc} [\emph{bjuʃ}] 0xEA11 \\
бутон: \textbf{brot} [\emph{brot}] 0xACD1 \\
буфетная: \textbf{bla7t} [\emph{blä˥t}] 0xB171 \\
Бухарест: \textbf{br6s} [\emph{brəs}] 0x8DD1 \\
бухточка: \textbf{hko7f} [\emph{ʰko˥f}] 0xD410 \\
бушель: \textbf{hbu7c} [\emph{ʰbu˥ʃ}] 0xF288 \\
быстро: \textbf{twis} [\emph{twis}] 0x882A \\
бычий: \textbf{bwot} [\emph{bwot}] 0xAC31 \\
Бэзил: \textbf{bjis} [\emph{bʒis}] 0x88B1 \\
Бэтмен: \textbf{hba2m} [\emph{ʰbä˩m}] 0x1988 \\
бюджет: \textbf{bjet} [\emph{bʒet}] 0xABB1 \\
бюрократия: \textbf{jras} [\emph{ʒräs}] 0x89D5 \\
бюстгальтеры: \textbf{bra2s} [\emph{brä˩s}] 0x99D1 \\
\subsection{ в }
в: \textbf{pre7k} [\emph{pre˥k}] 0x33C4 \\
\subsection{ В }
Вавилон: \textbf{bla7n} [\emph{blä˥n}] 0x7171 \\
важность: \textbf{myun} [\emph{mjun}] 0x6A01 \\
вазелин: \textbf{zli7n} [\emph{zli˥n}] 0x7074 \\
вакансия: \textbf{kfis} [\emph{kfis}] 0x8882 \\
вакуоль: \textbf{hvup} [\emph{ʰvup}] 0x4AB0 \\
вакцина: \textbf{byas} [\emph{bjäs}] 0x8911 \\
валентность: \textbf{vlen} [\emph{vlen}] 0x6B76 \\
валериана: \textbf{vra7n} [\emph{vrä˥n}] 0x71D6 \\
валлийский: \textbf{vles} [\emph{vles}] 0x8B76 \\
валлонский: \textbf{vl6n} [\emph{vlən}] 0x6D76 \\
валовой: \textbf{mrut} [\emph{mrut}] 0xAAC1 \\
вальс: \textbf{hwa7s} [\emph{ʰwä˥s}] 0x9128 \\
ванадий: \textbf{hva2m} [\emph{ʰvä˩m}] 0x19B0 \\
ваниль: \textbf{blis} [\emph{blis}] 0x8871 \\
ванна: \textbf{byan} [\emph{bjän}] 0x6911 \\
вареный: \textbf{hbuf} [\emph{ʰbuf}] 0xCA88 \\
варенье: \textbf{hjam} [\emph{ʰʒäm}] 0x09A8 \\
Варшавская: \textbf{hwa2s} [\emph{ʰwä˩s}] 0x9928 \\
ватерлиния: \textbf{vli7n} [\emph{vli˥n}] 0x7076 \\
ватт: \textbf{hwop} [\emph{ʰwop}] 0x4C28 \\
Вау: \textbf{gwas} [\emph{gwäs}] 0x8932 \\
вафельный: \textbf{hw6f} [\emph{ʰwəf}] 0xCD28 \\
Ваххабиты: \textbf{hwa2p} [\emph{ʰwä˩p}] 0x5928 \\
ваш: \textbf{nyut} [\emph{njut}] 0xAA06 \\
вводить: \textbf{hqik} [\emph{ʰŋik}] 0x28C8 \\
вводный: \textbf{tluk} [\emph{tluk}] 0x2A6A \\
вдова: \textbf{hwuf} [\emph{ʰwuf}] 0xCA28 \\
вдохновение: \textbf{nren} [\emph{nren}] 0x6BC6 \\
Веб-сайт: \textbf{hqas} [\emph{ʰŋäs}] 0x89C8 \\
ведомственный: \textbf{pxen} [\emph{pxen}] 0x6BE4 \\
ведьма: \textbf{hwoc} [\emph{ʰwoʃ}] 0xEC28 \\
везде: \textbf{hrac} [\emph{ʰräʃ}] 0xE970 \\
вездесущность: \textbf{zres} [\emph{zres}] 0x8BD4 \\
века: \textbf{zren} [\emph{zren}] 0x6BD4 \\
веко: \textbf{hyep} [\emph{ʰjep}] 0x4B18 \\
вектор: \textbf{dlek} [\emph{dlek}] 0x2B73 \\
великодушие: \textbf{xra7t} [\emph{xrä˥t}] 0xB1DA \\
величественный: \textbf{cwik} [\emph{ʃwik}] 0x282D \\
величина: \textbf{gjat} [\emph{gʒät}] 0xA9B2 \\
велосипед: \textbf{slek} [\emph{slek}] 0x2B69 \\
вена: \textbf{cwam} [\emph{ʃwäm}] 0x092D \\
Венгрия: \textbf{hqif} [\emph{ʰŋif}] 0xC8C8 \\
вендетта: \textbf{hve2t} [\emph{ʰve˩t}] 0xBBB0 \\
Венесуэла: \textbf{hnuf} [\emph{ʰnuf}] 0xCA30 \\
венецианский: \textbf{hves} [\emph{ʰves}] 0x8BB0 \\
венский: \textbf{vy6n} [\emph{vjən}] 0x6D16 \\
верблюд: \textbf{hlum} [\emph{ʰlum}] 0x0A58 \\
вербовщик: \textbf{vrin} [\emph{vrin}] 0x68D6 \\
веревка: \textbf{hrop} [\emph{ʰrop}] 0x4C70 \\
вереск: \textbf{hxef} [\emph{ʰxef}] 0xCBD0 \\
вермишель: \textbf{hvem} [\emph{ʰvem}] 0x0BB0 \\
верность: \textbf{tsot} [\emph{tsot}] 0xAC4A \\
вернуть: \textbf{hwit} [\emph{ʰwit}] 0xA828 \\
вероломство: \textbf{hri2t} [\emph{ʰri˩t}] 0xB870 \\
вероятность: \textbf{pfan} [\emph{pfän}] 0x6984 \\
версии: \textbf{hpe7n} [\emph{ʰpe˥n}] 0x7320 \\
вертикальный: \textbf{rwik} [\emph{rwik}] 0x282E \\
вертолет: \textbf{lwet} [\emph{lwet}] 0xAB2B \\
вертячка: \textbf{txa7k} [\emph{txä˥k}] 0x31EA \\
верфь: \textbf{vyas} [\emph{vjäs}] 0x8916 \\
верховой: \textbf{cwin} [\emph{ʃwin}] 0x682D \\
верхом: \textbf{kxa2t} [\emph{kxä˩t}] 0xB9E2 \\
вершина: \textbf{hvec} [\emph{ʰveʃ}] 0xEBB0 \\
вес: \textbf{psin} [\emph{psin}] 0x6844 \\
веселый: \textbf{clus} [\emph{ʃlus}] 0x8A6D \\
весло: \textbf{zyan} [\emph{zjän}] 0x6914 \\
весна: \textbf{pwan} [\emph{pwän}] 0x6924 \\
веснушчатый: \textbf{hre7k} [\emph{ʰre˥k}] 0x3370 \\
вести: \textbf{tyip} [\emph{tjip}] 0x480A \\
ветреница: \textbf{nwem} [\emph{nwem}] 0x0B26 \\
ветреный: \textbf{hjun} [\emph{ʰʒun}] 0x6AA8 \\
ветчина: \textbf{hgam} [\emph{ʰgäm}] 0x0990 \\
вечерня: \textbf{hve7p} [\emph{ʰve˥p}] 0x53B0 \\
вечнозеленый: \textbf{vren} [\emph{vren}] 0x6BD6 \\
вечность: \textbf{zlon} [\emph{zlon}] 0x6C74 \\
вечный: \textbf{hnen} [\emph{ʰnen}] 0x6B30 \\
вешать: \textbf{dwan} [\emph{dwän}] 0x6933 \\
взаимное: \textbf{hzac} [\emph{ʰzäʃ}] 0xE9A0 \\
взаимодействие: \textbf{zrec} [\emph{zreʃ}] 0xEBD4 \\
взаимопроникновение: \textbf{pr6c} [\emph{prəʃ}] 0xEDC4 \\
взбодрить: \textbf{hbo2n} [\emph{ʰbo˩n}] 0x7C88 \\
взгляд: \textbf{pcas} [\emph{pʃäs}] 0x89A4 \\
вздрогнуть: \textbf{tci7n} [\emph{tʃi˥n}] 0x70AA \\
взимание: \textbf{htu7n} [\emph{ʰtu˥n}] 0x7250 \\
взломать: \textbf{hbop} [\emph{ʰbop}] 0x4C88 \\
взносы: \textbf{hta2s} [\emph{ʰtä˩s}] 0x9950 \\
взрастить: \textbf{xyen} [\emph{xjen}] 0x6B1A \\
взревел: \textbf{nro7t} [\emph{nro˥t}] 0xB4C6 \\
взрыв: \textbf{plas} [\emph{pläs}] 0x8964 \\
вибратор: \textbf{vra7t} [\emph{vrä˥t}] 0xB1D6 \\
виброфон: \textbf{vraf} [\emph{vräf}] 0xC9D6 \\
вигны: \textbf{kwa7p} [\emph{kwä˥p}] 0x5122 \\
видео: \textbf{byot} [\emph{bjot}] 0xAC11 \\
Видео-конференция: \textbf{dxen} [\emph{dxen}] 0x6BF3 \\
видеть: \textbf{hsat} [\emph{ʰsät}] 0xA948 \\
виза: \textbf{bzis} [\emph{bzis}] 0x8851 \\
византийский: \textbf{hza7n} [\emph{ʰzä˥n}] 0x71A0 \\
визг: \textbf{hci7c} [\emph{ʰʃi˥ʃ}] 0xF068 \\
визирь: \textbf{hvi7s} [\emph{ʰvi˥s}] 0x90B0 \\
викторианский: \textbf{kro7n} [\emph{kro˥n}] 0x74C2 \\
викторина: \textbf{kwim} [\emph{kwim}] 0x0822 \\
вилла: \textbf{plif} [\emph{plif}] 0xC864 \\
виллан: \textbf{vla2n} [\emph{vlä˩n}] 0x7976 \\
вилы: \textbf{tfa7k} [\emph{tfä˥k}] 0x318A \\
винил: \textbf{vyi7n} [\emph{vji˥n}] 0x7016 \\
вино: \textbf{pwin} [\emph{pwin}] 0x6824 \\
виновный: \textbf{hdis} [\emph{ʰdis}] 0x8898 \\
виноградарство: \textbf{vli2t} [\emph{vli˩t}] 0xB876 \\
виночерпием: \textbf{px6p} [\emph{pxəp}] 0x4DE4 \\
винтовка: \textbf{tfic} [\emph{tfiʃ}] 0xE88A \\
виолончель: \textbf{vlef} [\emph{vlef}] 0xCB76 \\
Виргиния: \textbf{hri7n} [\emph{ʰri˥n}] 0x7070 \\
виртуальный: \textbf{brut} [\emph{brut}] 0xAAD1 \\
вирус: \textbf{pris} [\emph{pris}] 0x88C4 \\
вирши: \textbf{hgo7s} [\emph{ʰgo˥s}] 0x9490 \\
виски: \textbf{fyik} [\emph{fjik}] 0x280C \\
висмут: \textbf{hzim} [\emph{ʰzim}] 0x08A0 \\
висцеральный: \textbf{vrim} [\emph{vrim}] 0x08D6 \\
витамин: \textbf{hbem} [\emph{ʰbem}] 0x0B88 \\
витой: \textbf{hwot} [\emph{ʰwot}] 0xAC28 \\
вихри: \textbf{dwes} [\emph{dwes}] 0x8B33 \\
вихрь: \textbf{hw6n} [\emph{ʰwən}] 0x6D28 \\
вице: \textbf{hwif} [\emph{ʰwif}] 0xC828 \\
вице-король: \textbf{tra7f} [\emph{trä˥f}] 0xD1CA \\
Вишну: \textbf{hvuc} [\emph{ʰvuʃ}] 0xEAB0 \\
вишня: \textbf{tcef} [\emph{tʃef}] 0xCBAA \\
вкладыш: \textbf{lya2k} [\emph{ljä˩k}] 0x390B \\
вкрутую: \textbf{xru7t} [\emph{xru˥t}] 0xB2DA \\
вкус: \textbf{mwat} [\emph{mwät}] 0xA921 \\
вкусно: \textbf{ksuc} [\emph{ksuʃ}] 0xEA42 \\
владелец: \textbf{slit} [\emph{slit}] 0xA869 \\
владивосток: \textbf{dvok} [\emph{dvok}] 0x2C93 \\
влажный: \textbf{tcup} [\emph{tʃup}] 0x4AAA \\
властители: \textbf{sli7n} [\emph{sli˥n}] 0x7069 \\
властитель: \textbf{qra7t} [\emph{ŋrä˥t}] 0xB1D9 \\
власть: \textbf{cyik} [\emph{ʃjik}] 0x280D \\
влияние: \textbf{hfic} [\emph{ʰfiʃ}] 0xE860 \\
вместе: \textbf{kcat} [\emph{kʃät}] 0xA9A2 \\
вместимость: \textbf{hkot} [\emph{ʰkot}] 0xAC10 \\
вместо: \textbf{psat} [\emph{psät}] 0xA944 \\
вмешиваться: \textbf{kwip} [\emph{kwip}] 0x4822 \\
вмещать: \textbf{jyak} [\emph{ʒjäk}] 0x2915 \\
внеземной: \textbf{tr6t} [\emph{trət}] 0xADCA \\
вниз: \textbf{hcic} [\emph{ʰʃiʃ}] 0xE868 \\
внимательно: \textbf{tcaf} [\emph{tʃäf}] 0xC9AA \\
внутренний: \textbf{brin} [\emph{brin}] 0x68D1 \\
внутренности: \textbf{vya2n} [\emph{vjä˩n}] 0x7916 \\
внутриклеточный: \textbf{tle2c} [\emph{tle˩ʃ}] 0xFB6A \\
внутриутробный: \textbf{nrum} [\emph{nrum}] 0x0AC6 \\
внучата: \textbf{hnu7t} [\emph{ʰnu˥t}] 0xB230 \\
водевиль: \textbf{hvi7f} [\emph{ʰvi˥f}] 0xD0B0 \\
водиться: \textbf{jya7n} [\emph{ʒjä˥n}] 0x7115 \\
водка: \textbf{hvok} [\emph{ʰvok}] 0x2CB0 \\
водобоязнь: \textbf{hga7p} [\emph{ʰgä˥p}] 0x5190 \\
водовод: \textbf{sla2s} [\emph{slä˩s}] 0x9969 \\
водолаз: \textbf{dwaf} [\emph{dwäf}] 0xC933 \\
Водолей: \textbf{kru7s} [\emph{kru˥s}] 0x92C2 \\
водопад: \textbf{crop} [\emph{ʃrop}] 0x4CCD \\
водопровод: \textbf{sl6n} [\emph{slən}] 0x6D69 \\
водород: \textbf{nron} [\emph{nron}] 0x6CC6 \\
воды: \textbf{hwic} [\emph{ʰwiʃ}] 0xE828 \\
водянистый: \textbf{hw6c} [\emph{ʰwəʃ}] 0xED28 \\
вождение: \textbf{hdok} [\emph{ʰdok}] 0x2C98 \\
возврат: \textbf{tfit} [\emph{tfit}] 0xA88A \\
возвратный: \textbf{hfes} [\emph{ʰfes}] 0x8B60 \\
воздух: \textbf{cwat} [\emph{ʃwät}] 0xA92D \\
воздухоплавание: \textbf{kxi2k} [\emph{kxi˩k}] 0x38E2 \\
возмещать: \textbf{rwam} [\emph{rwäm}] 0x092E \\
возможно: \textbf{hyaf} [\emph{ʰjäf}] 0xC918 \\
возможности: \textbf{gxus} [\emph{gxus}] 0x8AF2 \\
возня: \textbf{xlum} [\emph{xlum}] 0x0A7A \\
возражение: \textbf{brat} [\emph{brät}] 0xA9D1 \\
возраст: \textbf{hnis} [\emph{ʰnis}] 0x8830 \\
возчик: \textbf{kra7t} [\emph{krä˥t}] 0xB1C2 \\
воин: \textbf{sron} [\emph{sron}] 0x6CC9 \\
воинственный: \textbf{hy67n} [\emph{ʰjə˥n}] 0x7518 \\
война: \textbf{twat} [\emph{twät}] 0xA92A \\
войти: \textbf{syak} [\emph{sjäk}] 0x2909 \\
вокализации: \textbf{vlas} [\emph{vläs}] 0x8976 \\
волан: \textbf{vlok} [\emph{vlok}] 0x2C76 \\
волдырь: \textbf{hwi7f} [\emph{ʰwi˥f}] 0xD028 \\
волей-неволей: \textbf{pro7k} [\emph{pro˥k}] 0x34C4 \\
волейбол: \textbf{xlip} [\emph{xlip}] 0x487A \\
волк: \textbf{hwek} [\emph{ʰwek}] 0x2B28 \\
волна: \textbf{hzat} [\emph{ʰzät}] 0xA9A0 \\
волнистый: \textbf{dlop} [\emph{dlop}] 0x4C73 \\
волокнистый: \textbf{fro7t} [\emph{fro˥t}] 0xB4CC \\
волоске: \textbf{kxo2t} [\emph{kxo˩t}] 0xBCE2 \\
волосной: \textbf{fre2n} [\emph{fre˩n}] 0x7BCC \\
Волоф: \textbf{hwo7f} [\emph{ʰwo˥f}] 0xD428 \\
волочиться: \textbf{jl6t} [\emph{ʒlət}] 0xAD75 \\
волчица: \textbf{cla7k} [\emph{ʃlä˥k}] 0x316D \\
волынка: \textbf{hgi2p} [\emph{ʰgi˩p}] 0x5890 \\
волынщик: \textbf{hpi2p} [\emph{ʰpi˩p}] 0x5820 \\
вольная: \textbf{hlu7n} [\emph{ʰlu˥n}] 0x7258 \\
вольт: \textbf{fwut} [\emph{fwut}] 0xAA2C \\
вольтметр: \textbf{lyo7t} [\emph{ljo˥t}] 0xB40B \\
вольфрам: \textbf{tsum} [\emph{tsum}] 0x0A4A \\
вонь: \textbf{sruc} [\emph{sruʃ}] 0xEAC9 \\
воображаемый: \textbf{blin} [\emph{blin}] 0x6871 \\
вопить: \textbf{bwus} [\emph{bwus}] 0x8A31 \\
вопрос: \textbf{hwas} [\emph{ʰwäs}] 0x8928 \\
Вопросы-Ответы: \textbf{fwa7t} [\emph{fwä˥t}] 0xB12C \\
вор: \textbf{htof} [\emph{ʰtof}] 0xCC50 \\
ворковать: \textbf{hku7k} [\emph{ʰku˥k}] 0x3210 \\
воробей: \textbf{mrec} [\emph{mreʃ}] 0xEBC1 \\
ворона: \textbf{kraf} [\emph{kräf}] 0xC9C2 \\
воротник: \textbf{kloc} [\emph{kloʃ}] 0xEC62 \\
ворсянка: \textbf{cr6s} [\emph{ʃrəs}] 0x8DCD \\
ворчливый: \textbf{cyup} [\emph{ʃjup}] 0x4A0D \\
воск: \textbf{flam} [\emph{fläm}] 0x096C \\
воскликнул: \textbf{kxum} [\emph{kxum}] 0x0AE2 \\
восклицание: \textbf{kxac} [\emph{kxäʃ}] 0xE9E2 \\
Воскресенье: \textbf{hwip} [\emph{ʰwip}] 0x4828 \\
воскрешение: \textbf{hfuc} [\emph{ʰfuʃ}] 0xEA60 \\
воспалительный: \textbf{fla7t} [\emph{flä˥t}] 0xB16C \\
воспитанным: \textbf{hli7k} [\emph{ʰli˥k}] 0x3058 \\
воспламеняться: \textbf{gzac} [\emph{gzäʃ}] 0xE952 \\
воспроизводимый: \textbf{psu7k} [\emph{psu˥k}] 0x3244 \\
восстановить: \textbf{hrek} [\emph{ʰrek}] 0x2B70 \\
восстановление: \textbf{hfus} [\emph{ʰfus}] 0x8A60 \\
восток: \textbf{tyuk} [\emph{tjuk}] 0x2A0A \\
восхищение: \textbf{crin} [\emph{ʃrin}] 0x68CD \\
Восходящий: \textbf{hqa2n} [\emph{ʰŋä˩n}] 0x79C8 \\
восьмеричный: \textbf{kc6s} [\emph{kʃəs}] 0x8DA2 \\
восьмидесятилетний: \textbf{tle2n} [\emph{tle˩n}] 0x7B6A \\
восьмидесятый: \textbf{tsi2f} [\emph{tsi˩f}] 0xD84A \\
восьмиугольник: \textbf{ty6k} [\emph{tjək}] 0x2D0A \\
Вот: \textbf{hsic} [\emph{ʰsiʃ}] 0xE848 \\
воющий: \textbf{xla2n} [\emph{xlä˩n}] 0x797A \\
вперед: \textbf{hyif} [\emph{ʰjif}] 0xC818 \\
впередсмотрящие: \textbf{lwu2t} [\emph{lwu˩t}] 0xBA2B \\
впечатление: \textbf{bwin} [\emph{bwin}] 0x6831 \\
впитывает: \textbf{hco7p} [\emph{ʰʃo˥p}] 0x5468 \\
враг: \textbf{tcam} [\emph{tʃäm}] 0x09AA \\
враждебность: \textbf{htu2t} [\emph{ʰtu˩t}] 0xBA50 \\
врач: \textbf{tsut} [\emph{tsut}] 0xAA4A \\
вращение: \textbf{prun} [\emph{prun}] 0x6AC4 \\
временный: \textbf{tric} [\emph{triʃ}] 0xE8CA \\
все: \textbf{hsut} [\emph{ʰsut}] 0xAA48 \\
всевидящий: \textbf{hse7k} [\emph{ʰse˥k}] 0x3348 \\
всевозможный: \textbf{xr6m} [\emph{xrəm}] 0x0DDA \\
всегда: \textbf{hsis} [\emph{ʰsis}] 0x8848 \\
всемогущий: \textbf{hwa7n} [\emph{ʰwä˥n}] 0x7128 \\
всеобъемлющая: \textbf{hwo7n} [\emph{ʰwo˥n}] 0x7428 \\
всеядный: \textbf{hm6c} [\emph{ʰməʃ}] 0xED08 \\
вскользь: \textbf{sla2n} [\emph{slä˩n}] 0x7969 \\
всмятку: \textbf{sla7m} [\emph{slä˥m}] 0x1169 \\
вспашка: \textbf{qrac} [\emph{ŋräʃ}] 0xE9D9 \\
всплеснула: \textbf{hla2s} [\emph{ʰlä˩s}] 0x9958 \\
вспыльчивый: \textbf{psa2p} [\emph{psä˩p}] 0x5944 \\
вспышка: \textbf{flac} [\emph{fläʃ}] 0xE96C \\
встряхивают: \textbf{hxi2t} [\emph{ʰxi˩t}] 0xB8D0 \\
вторгшийся: \textbf{txic} [\emph{txiʃ}] 0xE8EA \\
вторник: \textbf{mrac} [\emph{mräʃ}] 0xE9C1 \\
второй: \textbf{tlit} [\emph{tlit}] 0xA86A \\
вуаль: \textbf{pras} [\emph{präs}] 0x89C4 \\
вулкан: \textbf{glak} [\emph{gläk}] 0x2972 \\
вульгарность: \textbf{hva2t} [\emph{ʰvä˩t}] 0xB9B0 \\
вход: \textbf{nrup} [\emph{nrup}] 0x4AC6 \\
входящие: \textbf{dxis} [\emph{dxis}] 0x88F3 \\
вчера: \textbf{hten} [\emph{ʰten}] 0x6B50 \\
вшей: \textbf{hsu7s} [\emph{ʰsu˥s}] 0x9248 \\
вы: \textbf{hnan} [\emph{ʰnän}] 0x6930 \\
выберите: \textbf{kcut} [\emph{kʃut}] 0xAAA2 \\
выбивать: \textbf{kra7f} [\emph{krä˥f}] 0xD1C2 \\
выборы: \textbf{hben} [\emph{ʰben}] 0x6B88 \\
выбрасывать: \textbf{xwa2t} [\emph{xwä˩t}] 0xB93A \\
выбритый: \textbf{mwa2t} [\emph{mwä˩t}] 0xB921 \\
выветривание: \textbf{fr6n} [\emph{frən}] 0x6DCC \\
выводить: \textbf{prac} [\emph{präʃ}] 0xE9C4 \\
выводок: \textbf{jrif} [\emph{ʒrif}] 0xC8D5 \\
выглядывает: \textbf{pwus} [\emph{pwus}] 0x8A24 \\
выгодный: \textbf{hyof} [\emph{ʰjof}] 0xCC18 \\
выгравирован: \textbf{kri7n} [\emph{kri˥n}] 0x70C2 \\
выгрузка: \textbf{dxan} [\emph{dxän}] 0x69F3 \\
выдача: \textbf{kfot} [\emph{kfot}] 0xAC82 \\
выдвижение: \textbf{nyim} [\emph{njim}] 0x0806 \\
выделение: \textbf{plon} [\emph{plon}] 0x6C64 \\
выживаемость: \textbf{sr6f} [\emph{srəf}] 0xCDC9 \\
вызванный: \textbf{txa2n} [\emph{txä˩n}] 0x79EA \\
выздоравливающий: \textbf{kle2t} [\emph{kle˩t}] 0xBB62 \\
выкапывать: \textbf{hxaf} [\emph{ʰxäf}] 0xC9D0 \\
выкидыш: \textbf{gjac} [\emph{gʒäʃ}] 0xE9B2 \\
выкорчевывать: \textbf{dvut} [\emph{dvut}] 0xAA93 \\
вылетающих: \textbf{bxan} [\emph{bxän}] 0x69F1 \\
вымачивать: \textbf{mra7c} [\emph{mrä˥ʃ}] 0xF1C1 \\
вымпел: \textbf{hce2t} [\emph{ʰʃe˩t}] 0xBB68 \\
выпад: \textbf{vrot} [\emph{vrot}] 0xACD6 \\
выполнимый: \textbf{xli7n} [\emph{xli˥n}] 0x707A \\
выполнять: \textbf{pcim} [\emph{pʃim}] 0x08A4 \\
выпрямитель: \textbf{hre7f} [\emph{ʰre˥f}] 0xD370 \\
выпуклость: \textbf{txuc} [\emph{txuʃ}] 0xEAEA \\
выпуклый: \textbf{kxuk} [\emph{kxuk}] 0x2AE2 \\
выпускник: \textbf{ryet} [\emph{rjet}] 0xAB0E \\
выпускники: \textbf{lyu7n} [\emph{lju˥n}] 0x720B \\
выравнивание: \textbf{jren} [\emph{ʒren}] 0x6BD5 \\
выравнивать: \textbf{zric} [\emph{zriʃ}] 0xE8D4 \\
высиживать: \textbf{nyup} [\emph{njup}] 0x4A06 \\
выслушивание: \textbf{kli7s} [\emph{kli˥s}] 0x9062 \\
высокая: \textbf{hkuc} [\emph{ʰkuʃ}] 0xEA10 \\
высокомерный: \textbf{glus} [\emph{glus}] 0x8A72 \\
высокотональный: \textbf{xyi7t} [\emph{xji˥t}] 0xB01A \\
высосанный: \textbf{ksi7t} [\emph{ksi˥t}] 0xB042 \\
высотомер: \textbf{tca7p} [\emph{tʃä˥p}] 0x51AA \\
высший: \textbf{srum} [\emph{srum}] 0x0AC9 \\
высылка: \textbf{ksa7n} [\emph{ksä˥n}] 0x7142 \\
вытеснять: \textbf{srus} [\emph{srus}] 0x8AC9 \\
вычеркивать: \textbf{hja7s} [\emph{ʰʒä˥s}] 0x91A8 \\
вычерпывать: \textbf{hl67t} [\emph{ʰlə˥t}] 0xB558 \\
вычет: \textbf{ksoc} [\emph{ksoʃ}] 0xEC42 \\
вычислимость: \textbf{hku7t} [\emph{ʰku˥t}] 0xB210 \\
вычитаемое: \textbf{cro7n} [\emph{ʃro˥n}] 0x74CD \\
вычитать: \textbf{gret} [\emph{gret}] 0xABD2 \\
вышеупомянутый: \textbf{tri7t} [\emph{tri˥t}] 0xB0CA \\
вышитый: \textbf{bra7t} [\emph{brä˥t}] 0xB1D1 \\
выщелачивание: \textbf{hl6c} [\emph{ʰləʃ}] 0xED58 \\
выявляя: \textbf{fya2n} [\emph{fjä˩n}] 0x790C \\
Вьетнам: \textbf{vyam} [\emph{vjäm}] 0x0916 \\
Вэстон: \textbf{hw6m} [\emph{ʰwəm}] 0x0D28 \\
вяжущий: \textbf{sri2s} [\emph{sri˩s}] 0x98C9 \\
вяз: \textbf{hyem} [\emph{ʰjem}] 0x0B18 \\
вязаный: \textbf{nwet} [\emph{nwet}] 0xAB26 \\
вязкость: \textbf{vwot} [\emph{vwot}] 0xAC36 \\
вялость: \textbf{psa2t} [\emph{psä˩t}] 0xB944 \\
\subsection{ Г }
Г-образный: \textbf{hle2c} [\emph{ʰle˩ʃ}] 0xFB58 \\
Г.Д.: \textbf{gji7t} [\emph{gʒi˥t}] 0xB0B2 \\
Габон: \textbf{gwap} [\emph{gwäp}] 0x4932 \\
Гавана: \textbf{xwa7n} [\emph{xwä˥n}] 0x713A \\
гагачий: \textbf{hdi2t} [\emph{ʰdi˩t}] 0xB898 \\
гадюка: \textbf{ksa2p} [\emph{ksä˩p}] 0x5942 \\
газ: \textbf{hgas} [\emph{ʰgäs}] 0x8990 \\
газета: \textbf{gzat} [\emph{gzät}] 0xA952 \\
газифицировать: \textbf{gzaf} [\emph{gzäf}] 0xC952 \\
газонокосилка: \textbf{tso7m} [\emph{tso˥m}] 0x144A \\
Гаити: \textbf{xyes} [\emph{xjes}] 0x8B1A \\
галактика: \textbf{gjak} [\emph{gʒäk}] 0x29B2 \\
галактоза: \textbf{gjos} [\emph{gʒos}] 0x8CB2 \\
галантерея: \textbf{bwap} [\emph{bwäp}] 0x4931 \\
галерея: \textbf{glan} [\emph{glän}] 0x6972 \\
галисийский: \textbf{gla7c} [\emph{glä˥ʃ}] 0xF172 \\
галлия: \textbf{glim} [\emph{glim}] 0x0872 \\
галлон: \textbf{glon} [\emph{glon}] 0x6C72 \\
галлюциногенный: \textbf{xwuk} [\emph{xwuk}] 0x2A3A \\
гало: \textbf{xla7n} [\emph{xlä˥n}] 0x717A \\
галоген: \textbf{xl6k} [\emph{xlək}] 0x2D7A \\
галопом: \textbf{hg6p} [\emph{ʰgəp}] 0x4D90 \\
галоши: \textbf{gluc} [\emph{gluʃ}] 0xEA72 \\
галстук: \textbf{rwat} [\emph{rwät}] 0xA92E \\
гальванизировать: \textbf{gvun} [\emph{gvun}] 0x6A92 \\
Гамбия: \textbf{gyap} [\emph{gjäp}] 0x4912 \\
гамбургер: \textbf{crep} [\emph{ʃrep}] 0x4BCD \\
гамета: \textbf{gwam} [\emph{gwäm}] 0x0932 \\
гамма: \textbf{hgom} [\emph{ʰgom}] 0x0C90 \\
Гана: \textbf{gx6n} [\emph{gxən}] 0x6DF2 \\
Ганновер: \textbf{hxuf} [\emph{ʰxuf}] 0xCAD0 \\
гантель: \textbf{dyep} [\emph{djep}] 0x4B13 \\
гараж: \textbf{grek} [\emph{grek}] 0x2BD2 \\
гарантия: \textbf{gzan} [\emph{gzän}] 0x6952 \\
гармоника: \textbf{xrop} [\emph{xrop}] 0x4CDA \\
гармония: \textbf{hcom} [\emph{ʰʃom}] 0x0C68 \\
гарпия: \textbf{xra7p} [\emph{xrä˥p}] 0x51DA \\
гаррота: \textbf{gra7s} [\emph{grä˥s}] 0x91D2 \\
гастрит: \textbf{gri2s} [\emph{gri˩s}] 0x98D2 \\
гастроном: \textbf{dxes} [\emph{dxes}] 0x8BF3 \\
гаубица: \textbf{hxo2s} [\emph{ʰxo˩s}] 0x9CD0 \\
гашиш: \textbf{tca2c} [\emph{tʃä˩ʃ}] 0xF9AA \\
гвардия: \textbf{hga2s} [\emph{ʰgä˩s}] 0x9990 \\
Гватемала: \textbf{gxam} [\emph{gxäm}] 0x09F2 \\
Гвиана: \textbf{gyun} [\emph{gjun}] 0x6A12 \\
Гвинея: \textbf{gyif} [\emph{gjif}] 0xC812 \\
гвоздика: \textbf{hki7c} [\emph{ʰki˥ʃ}] 0xF010 \\
где-то: \textbf{hmoc} [\emph{ʰmoʃ}] 0xEC08 \\
гегельянец: \textbf{xle2n} [\emph{xle˩n}] 0x7B7A \\
гедонистический: \textbf{dzi7n} [\emph{dzi˥n}] 0x7053 \\
гей: \textbf{hgek} [\emph{ʰgek}] 0x2B90 \\
гейзер: \textbf{gzis} [\emph{gzis}] 0x8852 \\
гейша: \textbf{hg6c} [\emph{ʰgəʃ}] 0xED90 \\
гекзаметр: \textbf{ksa2m} [\emph{ksä˩m}] 0x1942 \\
гексан: \textbf{hze7n} [\emph{ʰze˥n}] 0x73A0 \\
гелий: \textbf{clum} [\emph{ʃlum}] 0x0A6D \\
гелиоцентрический: \textbf{cyi2t} [\emph{ʃji˩t}] 0xB80D \\
гем: \textbf{hxe2m} [\emph{ʰxe˩m}] 0x1BD0 \\
гематит: \textbf{hxi7t} [\emph{ʰxi˥t}] 0xB0D0 \\
гемоглобин: \textbf{xwop} [\emph{xwop}] 0x4C3A \\
геморрой: \textbf{mro2s} [\emph{mro˩s}] 0x9CC1 \\
ген: \textbf{gyen} [\emph{gjen}] 0x6B12 \\
генералиссимус: \textbf{jlam} [\emph{ʒläm}] 0x0975 \\
генетика: \textbf{nwek} [\emph{nwek}] 0x2B26 \\
гениталии: \textbf{hsi2t} [\emph{ʰsi˩t}] 0xB848 \\
генотипическая: \textbf{jyip} [\emph{ʒjip}] 0x4815 \\
геноцид: \textbf{gxit} [\emph{gxit}] 0xA8F2 \\
Генри: \textbf{clin} [\emph{ʃlin}] 0x686D \\
Гентский: \textbf{hge2t} [\emph{ʰge˩t}] 0xBB90 \\
Генуя: \textbf{jyon} [\emph{ʒjon}] 0x6C15 \\
география: \textbf{trif} [\emph{trif}] 0xC8CA \\
геодезический: \textbf{gyi7k} [\emph{gji˥k}] 0x3012 \\
геолог: \textbf{jlek} [\emph{ʒlek}] 0x2B75 \\
геологических: \textbf{jlik} [\emph{ʒlik}] 0x2875 \\
геометрия: \textbf{gxet} [\emph{gxet}] 0xABF2 \\
георгин: \textbf{dli2k} [\emph{dli˩k}] 0x3873 \\
геофизика: \textbf{gzok} [\emph{gzok}] 0x2C52 \\
гепатит: \textbf{hxas} [\emph{ʰxäs}] 0x89D0 \\
геральдика: \textbf{hr62t} [\emph{ʰrə˩t}] 0xBD70 \\
герань: \textbf{jrem} [\emph{ʒrem}] 0x0BD5 \\
Геркулес: \textbf{xlus} [\emph{xlus}] 0x8A7A \\
геркулесов: \textbf{xle7t} [\emph{xle˥t}] 0xB37A \\
германий: \textbf{grem} [\emph{grem}] 0x0BD2 \\
гермафродит: \textbf{xre7t} [\emph{xre˥t}] 0xB3DA \\
герменевтика: \textbf{xye2k} [\emph{xje˩k}] 0x3B1A \\
герметизация: \textbf{psa7c} [\emph{psä˥ʃ}] 0xF144 \\
герметический: \textbf{kwi2t} [\emph{kwi˩t}] 0xB822 \\
Геродот: \textbf{hxo2t} [\emph{ʰxo˩t}] 0xBCD0 \\
герои: \textbf{xro7n} [\emph{xro˥n}] 0x74DA \\
героин: \textbf{fron} [\emph{fron}] 0x6CCC \\
герц: \textbf{hxes} [\emph{ʰxes}] 0x8BD0 \\
герцеговина: \textbf{hxof} [\emph{ʰxof}] 0xCCD0 \\
герцог: \textbf{hnuk} [\emph{ʰnuk}] 0x2A30 \\
Гессе: \textbf{hxe7s} [\emph{ʰxe˥s}] 0x93D0 \\
гессенский: \textbf{xyep} [\emph{xjep}] 0x4B1A \\
гетеродинный: \textbf{twe7n} [\emph{twe˥n}] 0x732A \\
гетерозиготный: \textbf{tre2k} [\emph{tre˩k}] 0x3BCA \\
гетеросексуальный: \textbf{tro2k} [\emph{tro˩k}] 0x3CCA \\
гиацинт: \textbf{xyi7n} [\emph{xji˥n}] 0x701A \\
Гибралтар: \textbf{bli2t} [\emph{bli˩t}] 0xB871 \\
гибридизация: \textbf{bvac} [\emph{bväʃ}] 0xE991 \\
гигант: \textbf{djan} [\emph{dʒän}] 0x69B3 \\
гигантизм: \textbf{hj6s} [\emph{ʰʒəs}] 0x8DA8 \\
гигрометр: \textbf{gro2t} [\emph{gro˩t}] 0xBCD2 \\
гидравлический: \textbf{xwik} [\emph{xwik}] 0x283A \\
гидразин: \textbf{xru7n} [\emph{xru˥n}] 0x72DA \\
гидрат: \textbf{xre2t} [\emph{xre˩t}] 0xBBDA \\
гидрокси: \textbf{xro7s} [\emph{xro˥s}] 0x94DA \\
гидроксид: \textbf{xrus} [\emph{xrus}] 0x8ADA \\
гидроксил: \textbf{xri7k} [\emph{xri˥k}] 0x30DA \\
гидролиз: \textbf{dles} [\emph{dles}] 0x8B73 \\
гидропонный: \textbf{dr6k} [\emph{drək}] 0x2DD3 \\
гидростатический: \textbf{dri7k} [\emph{dri˥k}] 0x30D3 \\
гидросфера: \textbf{dref} [\emph{dref}] 0xCBD3 \\
гидротермальный: \textbf{dri2m} [\emph{dri˩m}] 0x18D3 \\
гидрофобный: \textbf{dri7p} [\emph{dri˥p}] 0x50D3 \\
гидрохлорид: \textbf{dlo2t} [\emph{dlo˩t}] 0xBC73 \\
гидроэлектрический: \textbf{dwek} [\emph{dwek}] 0x2B33 \\
гиена: \textbf{xre7n} [\emph{xre˥n}] 0x73DA \\
гильдия: \textbf{hgop} [\emph{ʰgop}] 0x4C90 \\
гильотина: \textbf{gli2t} [\emph{gli˩t}] 0xB872 \\
Гималаи: \textbf{xles} [\emph{xles}] 0x8B7A \\
гимн: \textbf{hso7n} [\emph{ʰso˥n}] 0x7448 \\
гимназия: \textbf{mwam} [\emph{mwäm}] 0x0921 \\
гимнастика: \textbf{gyas} [\emph{gjäs}] 0x8912 \\
гинеколог: \textbf{gxi7s} [\emph{gxi˥s}] 0x90F2 \\
гинекология: \textbf{gl6k} [\emph{glək}] 0x2D72 \\
гиперактивный: \textbf{pro2k} [\emph{pro˩k}] 0x3CC4 \\
гипербола: \textbf{xwip} [\emph{xwip}] 0x483A \\
гиперинфляция: \textbf{xl6s} [\emph{xləs}] 0x8D7A \\
гиперссылка: \textbf{xlik} [\emph{xlik}] 0x287A \\
гипертония: \textbf{pra7c} [\emph{prä˥ʃ}] 0xF1C4 \\
гипертрофия: \textbf{pri2f} [\emph{pri˩f}] 0xD8C4 \\
гипнотерапия: \textbf{pfep} [\emph{pfep}] 0x4B84 \\
гипнотизировать: \textbf{hxif} [\emph{ʰxif}] 0xC8D0 \\
гипотония: \textbf{xyec} [\emph{xjeʃ}] 0xEB1A \\
гипофиз: \textbf{pfi7s} [\emph{pfi˥s}] 0x9084 \\
гипохлорит: \textbf{xlu2t} [\emph{xlu˩t}] 0xBA7A \\
гироскоп: \textbf{hjop} [\emph{ʰʒop}] 0x4CA8 \\
гистология: \textbf{sli7c} [\emph{sli˥ʃ}] 0xF069 \\
гитара: \textbf{gyic} [\emph{gjiʃ}] 0xE812 \\
глагол: \textbf{kres} [\emph{kres}] 0x8BC2 \\
глаголический: \textbf{glo7k} [\emph{glo˥k}] 0x3472 \\
гладильный: \textbf{hj6t} [\emph{ʰʒət}] 0xADA8 \\
гладиолус: \textbf{glos} [\emph{glos}] 0x8C72 \\
глаз: \textbf{kyac} [\emph{kjäʃ}] 0xE902 \\
глазурь: \textbf{hgi2s} [\emph{ʰgi˩s}] 0x9890 \\
глаукома: \textbf{lwom} [\emph{lwom}] 0x0C2B \\
глашатаи: \textbf{xla2t} [\emph{xlä˩t}] 0xB97A \\
гликоль: \textbf{gri7k} [\emph{gri˥k}] 0x30D2 \\
глифы: \textbf{gwif} [\emph{gwif}] 0xC832 \\
глицерин: \textbf{gri2n} [\emph{gri˩n}] 0x78D2 \\
глицин: \textbf{gyi7n} [\emph{gji˥n}] 0x7012 \\
глобализация: \textbf{xlan} [\emph{xlän}] 0x697A \\
глотать: \textbf{gwan} [\emph{gwän}] 0x6932 \\
глотка: \textbf{hye2n} [\emph{ʰje˩n}] 0x7B18 \\
глубина: \textbf{hf6k} [\emph{ʰfək}] 0x2D60 \\
Глубинные: \textbf{gruk} [\emph{gruk}] 0x2AD2 \\
глубоко: \textbf{kcin} [\emph{kʃin}] 0x68A2 \\
глухота: \textbf{dro2n} [\emph{dro˩n}] 0x7CD3 \\
глушь: \textbf{bwa7t} [\emph{bwä˥t}] 0xB131 \\
глыбы: \textbf{gla7t} [\emph{glä˥t}] 0xB172 \\
глюкоза: \textbf{gwuk} [\emph{gwuk}] 0x2A32 \\
глютен: \textbf{gvut} [\emph{gvut}] 0xAA92 \\
гм: \textbf{hgu7m} [\emph{ʰgu˥m}] 0x1290 \\
гнев: \textbf{hfun} [\emph{ʰfun}] 0x6A60 \\
гнездо: \textbf{sras} [\emph{sräs}] 0x89C9 \\
гнейс: \textbf{pyes} [\emph{pjes}] 0x8B04 \\
гниль: \textbf{frun} [\emph{frun}] 0x6ACC \\
гной: \textbf{hpo2k} [\emph{ʰpo˩k}] 0x3C20 \\
гностицизм: \textbf{hzi2s} [\emph{ʰzi˩s}] 0x98A0 \\
гностический: \textbf{nyo2t} [\emph{njo˩t}] 0xBC06 \\
говорливость: \textbf{vlit} [\emph{vlit}] 0xA876 \\
говоря: \textbf{klam} [\emph{kläm}] 0x0962 \\
годовалый: \textbf{bri2n} [\emph{bri˩n}] 0x78D1 \\
годовой: \textbf{trak} [\emph{träk}] 0x29CA \\
Голгофа: \textbf{kl6f} [\emph{kləf}] 0xCD62 \\
Голландский: \textbf{dloc} [\emph{dloʃ}] 0xEC73 \\
головастик: \textbf{hxup} [\emph{ʰxup}] 0x4AD0 \\
головокружительный: \textbf{vyi2n} [\emph{vji˩n}] 0x7816 \\
головоломка: \textbf{pwet} [\emph{pwet}] 0xAB24 \\
головорез: \textbf{hco2k} [\emph{ʰʃo˩k}] 0x3C68 \\
голодный: \textbf{hcac} [\emph{ʰʃäʃ}] 0xE968 \\
голубовато-черный: \textbf{hlu7k} [\emph{ʰlu˥k}] 0x3258 \\
голубь: \textbf{pyot} [\emph{pjot}] 0xAC04 \\
гольф: \textbf{glof} [\emph{glof}] 0xCC72 \\
гомеопатия: \textbf{xlop} [\emph{xlop}] 0x4C7A \\
гомосексуализм: \textbf{xwim} [\emph{xwim}] 0x083A \\
гонады: \textbf{gwos} [\emph{gwos}] 0x8C32 \\
гондурасский: \textbf{nru7s} [\emph{nru˥s}] 0x92C6 \\
гонорея: \textbf{grep} [\emph{grep}] 0x4BD2 \\
гончар: \textbf{xroc} [\emph{xroʃ}] 0xECDA \\
горбатый: \textbf{hxo7p} [\emph{ʰxo˥p}] 0x54D0 \\
гордость: \textbf{kcam} [\emph{kʃäm}] 0x09A2 \\
горизонт: \textbf{prin} [\emph{prin}] 0x68C4 \\
гормон: \textbf{crum} [\emph{ʃrum}] 0x0ACD \\
город: \textbf{syit} [\emph{sjit}] 0xA809 \\
Гороскоп: \textbf{ksok} [\emph{ksok}] 0x2C42 \\
гортань: \textbf{qron} [\emph{ŋron}] 0x6CD9 \\
горчица: \textbf{hmom} [\emph{ʰmom}] 0x0C08 \\
горшок: \textbf{pfot} [\emph{pfot}] 0xAC84 \\
горький: \textbf{kcuc} [\emph{kʃuʃ}] 0xEAA2 \\
горьковато-сладкий: \textbf{ks67t} [\emph{ksə˥t}] 0xB542 \\
горючий: \textbf{kx6t} [\emph{kxət}] 0xADE2 \\
господа: \textbf{hsa2n} [\emph{ʰsä˩n}] 0x7948 \\
Госпожа: \textbf{hma7t} [\emph{ʰmä˥t}] 0xB108 \\
гостеприимство: \textbf{hjik} [\emph{ʰʒik}] 0x28A8 \\
Гостиница: \textbf{hnot} [\emph{ʰnot}] 0xAC30 \\
гость: \textbf{hget} [\emph{ʰget}] 0xAB90 \\
готика: \textbf{gxok} [\emph{gxok}] 0x2CF2 \\
готовить: \textbf{kcus} [\emph{kʃus}] 0x8AA2 \\
гравий: \textbf{grik} [\emph{grik}] 0x28D2 \\
градина: \textbf{xlup} [\emph{xlup}] 0x4A7A \\
граждане: \textbf{nri2n} [\emph{nri˩n}] 0x78C6 \\
гражданского: \textbf{swin} [\emph{swin}] 0x6829 \\
грамматика: \textbf{hgaf} [\emph{ʰgäf}] 0xC990 \\
грамотный: \textbf{hs62t} [\emph{ʰsə˩t}] 0xBD48 \\
гранатовый: \textbf{dlaf} [\emph{dläf}] 0xC973 \\
гранит: \textbf{gwen} [\emph{gwen}] 0x6B32 \\
граница: \textbf{bwit} [\emph{bwit}] 0xA831 \\
графит: \textbf{gxif} [\emph{gxif}] 0xC8F2 \\
графия: \textbf{hga7f} [\emph{ʰgä˥f}] 0xD190 \\
графство: \textbf{hca2k} [\emph{ʰʃä˩k}] 0x3968 \\
грациозность: \textbf{rya2n} [\emph{rjä˩n}] 0x790E \\
гребец: \textbf{rwus} [\emph{rwus}] 0x8A2E \\
гребешок: \textbf{gxap} [\emph{gxäp}] 0x49F2 \\
гребни: \textbf{hri2c} [\emph{ʰri˩ʃ}] 0xF870 \\
грейпфрут: \textbf{hruf} [\emph{ʰruf}] 0xCA70 \\
Гренада: \textbf{gy6t} [\emph{gjət}] 0xAD12 \\
Гренландия: \textbf{glen} [\emph{glen}] 0x6B72 \\
грех: \textbf{tsip} [\emph{tsip}] 0x484A \\
Греция: \textbf{gyis} [\emph{gjis}] 0x8812 \\
греческий: \textbf{glif} [\emph{glif}] 0xC872 \\
грешник: \textbf{hki2n} [\emph{ʰki˩n}] 0x7810 \\
григорианский: \textbf{gre2n} [\emph{gre˩n}] 0x7BD2 \\
гриль: \textbf{grip} [\emph{grip}] 0x48D2 \\
гримаса: \textbf{klim} [\emph{klim}] 0x0862 \\
Гринпис: \textbf{gxep} [\emph{gxep}] 0x4BF2 \\
грифон: \textbf{hri7f} [\emph{ʰri˥f}] 0xD070 \\
гроб: \textbf{kras} [\emph{kräs}] 0x89C2 \\
гробница: \textbf{lwi2t} [\emph{lwi˩t}] 0xB82B \\
гром: \textbf{grit} [\emph{grit}] 0xA8D2 \\
грот: \textbf{hma7n} [\emph{ʰmä˥n}] 0x7108 \\
грубый: \textbf{hruc} [\emph{ʰruʃ}] 0xEA70 \\
грудина: \textbf{hsu7m} [\emph{ʰsu˥m}] 0x1248 \\
грудь: \textbf{pfut} [\emph{pfut}] 0xAA84 \\
Грузия: \textbf{groc} [\emph{groʃ}] 0xECD2 \\
груминг: \textbf{gjun} [\emph{gʒun}] 0x6AB2 \\
грунтовый: \textbf{bxic} [\emph{bxiʃ}] 0xE8F1 \\
группы: \textbf{hba2t} [\emph{ʰbä˩t}] 0xB988 \\
грустный: \textbf{syic} [\emph{sjiʃ}] 0xE809 \\
груши: \textbf{pr6p} [\emph{prəp}] 0x4DC4 \\
грыжа: \textbf{gra2n} [\emph{grä˩n}] 0x79D2 \\
грязь: \textbf{nyik} [\emph{njik}] 0x2806 \\
гу: \textbf{hgu7k} [\emph{ʰgu˥k}] 0x3290 \\
гуанин: \textbf{gwun} [\emph{gwun}] 0x6A32 \\
губа: \textbf{tcop} [\emph{tʃop}] 0x4CAA \\
губка: \textbf{pfon} [\emph{pfon}] 0x6C84 \\
губчатый: \textbf{hmi7c} [\emph{ʰmi˥ʃ}] 0xF008 \\
гудок: \textbf{byup} [\emph{bjup}] 0x4A11 \\
гуж: \textbf{hto2k} [\emph{ʰto˩k}] 0x3C50 \\
гумусовый: \textbf{hxu7m} [\emph{ʰxu˥m}] 0x12D0 \\
гуртовщик: \textbf{hra7f} [\emph{ʰrä˥f}] 0xD170 \\
гусар: \textbf{hxa2p} [\emph{ʰxä˩p}] 0x59D0 \\
гусеница: \textbf{mruc} [\emph{mruʃ}] 0xEAC1 \\
гуттаперча: \textbf{gruc} [\emph{gruʃ}] 0xEAD2 \\
гыйба: \textbf{xya7k} [\emph{xjä˥k}] 0x311A \\
\subsection{ Д }
Да: \textbf{dwam} [\emph{dwäm}] 0x0933 \\
Далида: \textbf{hxa2s} [\emph{ʰxä˩s}] 0x99D0 \\
далматинец: \textbf{hd6c} [\emph{ʰdəʃ}] 0xED98 \\
дальтонизм: \textbf{qyan} [\emph{ŋjän}] 0x6919 \\
Дания: \textbf{drem} [\emph{drem}] 0x0BD3 \\
данник: \textbf{twif} [\emph{twif}] 0xC82A \\
данный: \textbf{gvin} [\emph{gvin}] 0x6892 \\
дарохранительница: \textbf{pyi7s} [\emph{pji˥s}] 0x9004 \\
датчанин: \textbf{dra7n} [\emph{drä˥n}] 0x71D3 \\
даты: \textbf{hta2c} [\emph{ʰtä˩ʃ}] 0xF950 \\
дающий: \textbf{kwin} [\emph{kwin}] 0x6822 \\
два: \textbf{tyut} [\emph{tjut}] 0xAA0A \\
дважды: \textbf{twas} [\emph{twäs}] 0x892A \\
двенадцать: \textbf{tfat} [\emph{tfät}] 0xA98A \\
дверь: \textbf{tyan} [\emph{tjän}] 0x690A \\
двигатель: \textbf{mrot} [\emph{mrot}] 0xACC1 \\
движения: \textbf{hdo2n} [\emph{ʰdo˩n}] 0x7C98 \\
движимый: \textbf{kxa7t} [\emph{kxä˥t}] 0xB1E2 \\
двоек: \textbf{tli7s} [\emph{tli˥s}] 0x906A \\
двоеточие: \textbf{jlon} [\emph{ʒlon}] 0x6C75 \\
двоичный: \textbf{byin} [\emph{bjin}] 0x6811 \\
двойники: \textbf{hdo7s} [\emph{ʰdo˥s}] 0x9498 \\
двойственность: \textbf{dvic} [\emph{dviʃ}] 0xE893 \\
дворец: \textbf{plac} [\emph{pläʃ}] 0xE964 \\
дворняжка: \textbf{hja7k} [\emph{ʰʒä˥k}] 0x31A8 \\
двояковыпуклый: \textbf{bvif} [\emph{bvif}] 0xC891 \\
двугранный: \textbf{dzi7t} [\emph{dzi˥t}] 0xB053 \\
двуколка: \textbf{hgi7k} [\emph{ʰgi˥k}] 0x3090 \\
двумужница: \textbf{gxi7t} [\emph{gxi˥t}] 0xB0F2 \\
двуногое: \textbf{hbu2t} [\emph{ʰbu˩t}] 0xBA88 \\
двухгодичный: \textbf{qyac} [\emph{ŋjäʃ}] 0xE919 \\
двухпартийный: \textbf{dxi7n} [\emph{dxi˥n}] 0x70F3 \\
двухпросветный: \textbf{dla7k} [\emph{dlä˥k}] 0x3173 \\
двухсотлетие: \textbf{bva7n} [\emph{bvä˥n}] 0x7191 \\
двухстороннее: \textbf{dva7n} [\emph{dvä˥n}] 0x7193 \\
двухтактный: \textbf{hdo2k} [\emph{ʰdo˩k}] 0x3C98 \\
двухцветный: \textbf{dri7m} [\emph{dri˥m}] 0x10D3 \\
двуязычный: \textbf{dzik} [\emph{dzik}] 0x2853 \\
де: \textbf{dye7t} [\emph{dje˥t}] 0xB313 \\
дебет: \textbf{hdef} [\emph{ʰdef}] 0xCB98 \\
дебушировать: \textbf{hde7k} [\emph{ʰde˥k}] 0x3398 \\
девальвация: \textbf{vlus} [\emph{vlus}] 0x8A76 \\
девонский: \textbf{dv6n} [\emph{dvən}] 0x6D93 \\
девственница: \textbf{kcen} [\emph{kʃen}] 0x6BA2 \\
девушка: \textbf{tcac} [\emph{tʃäʃ}] 0xE9AA \\
девяностый: \textbf{hni2f} [\emph{ʰni˩f}] 0xD830 \\
девятикратный: \textbf{hno7f} [\emph{ʰno˥f}] 0xD430 \\
дегенерат: \textbf{xyet} [\emph{xjet}] 0xAB1A \\
дегустатор: \textbf{tce7s} [\emph{tʃe˥s}] 0x93AA \\
дедовщина: \textbf{xwi7n} [\emph{xwi˥n}] 0x703A \\
дезертир: \textbf{dri2t} [\emph{dri˩t}] 0xB8D3 \\
дезинформация: \textbf{dr6m} [\emph{drəm}] 0x0DD3 \\
дезинформировать: \textbf{myof} [\emph{mjof}] 0xCC01 \\
дезодорант: \textbf{dxoc} [\emph{dxoʃ}] 0xECF3 \\
деизм: \textbf{hdi2m} [\emph{ʰdi˩m}] 0x1898 \\
действительный: \textbf{hyi7k} [\emph{ʰji˥k}] 0x3018 \\
Декабрь: \textbf{tyim} [\emph{tjim}] 0x080A \\
декагон: \textbf{dzek} [\emph{dzek}] 0x2B53 \\
декартовский: \textbf{kx6s} [\emph{kxəs}] 0x8DE2 \\
декорации: \textbf{frin} [\emph{frin}] 0x68CC \\
декстрин: \textbf{kce2t} [\emph{kʃe˩t}] 0xBBA2 \\
дела: \textbf{swan} [\emph{swän}] 0x6929 \\
деликт: \textbf{dli7k} [\emph{dli˥k}] 0x3073 \\
делитель: \textbf{bvas} [\emph{bväs}] 0x8991 \\
дельфийский: \textbf{dlef} [\emph{dlef}] 0xCB73 \\
дельфин: \textbf{dlif} [\emph{dlif}] 0xC873 \\
демагог: \textbf{hdo7k} [\emph{ʰdo˥k}] 0x3498 \\
демагогия: \textbf{dxa7k} [\emph{dxä˥k}] 0x31F3 \\
демерредж: \textbf{tci2c} [\emph{tʃi˩ʃ}] 0xF8AA \\
демократия: \textbf{mris} [\emph{mris}] 0x88C1 \\
демон: \textbf{hde7m} [\emph{ʰde˥m}] 0x1398 \\
демонстративный: \textbf{hmat} [\emph{ʰmät}] 0xA908 \\
демонстрировать: \textbf{dren} [\emph{dren}] 0x6BD3 \\
деморализация: \textbf{dli2n} [\emph{dli˩n}] 0x7873 \\
демпинг: \textbf{hdi7n} [\emph{ʰdi˥n}] 0x7098 \\
денатурировать: \textbf{dyec} [\emph{djeʃ}] 0xEB13 \\
дендрология: \textbf{dri2k} [\emph{dri˩k}] 0x38D3 \\
день: \textbf{tcin} [\emph{tʃin}] 0x68AA \\
Деньги: \textbf{pcan} [\emph{pʃän}] 0x69A4 \\
депиляция: \textbf{ple2c} [\emph{ple˩ʃ}] 0xFB64 \\
деполяризация: \textbf{dla2s} [\emph{dlä˩s}] 0x9973 \\
депутация: \textbf{hpu2t} [\emph{ʰpu˩t}] 0xBA20 \\
дервиш: \textbf{dvec} [\emph{dveʃ}] 0xEB93 \\
деревенщина: \textbf{xyi2k} [\emph{xji˩k}] 0x381A \\
деревня: \textbf{tsun} [\emph{tsun}] 0x6A4A \\
дерево: \textbf{syuc} [\emph{sjuʃ}] 0xEA09 \\
держа: \textbf{pwic} [\emph{pwiʃ}] 0xE824 \\
дерматолог: \textbf{dwoc} [\emph{dwoʃ}] 0xEC33 \\
деспотизм: \textbf{dxem} [\emph{dxem}] 0x0BF3 \\
десть: \textbf{kwa7s} [\emph{kwä˥s}] 0x9122 \\
десятикратный: \textbf{hse2p} [\emph{ʰse˩p}] 0x5B48 \\
десятилетия: \textbf{hda2k} [\emph{ʰdä˩k}] 0x3998 \\
десятилетний: \textbf{sye2n} [\emph{sje˩n}] 0x7B09 \\
десятина: \textbf{dl6k} [\emph{dlək}] 0x2D73 \\
детализировать: \textbf{zlam} [\emph{zläm}] 0x0974 \\
детектор: \textbf{dyik} [\emph{djik}] 0x2813 \\
дети: \textbf{clat} [\emph{ʃlät}] 0xA96D \\
детоксификация: \textbf{kfi2k} [\emph{kfi˩k}] 0x3882 \\
детонатор: \textbf{dwe7t} [\emph{dwe˥t}] 0xB333 \\
детство: \textbf{tlot} [\emph{tlot}] 0xAC6A \\
дефектный: \textbf{hdet} [\emph{ʰdet}] 0xAB98 \\
дефиле: \textbf{dwim} [\emph{dwim}] 0x0833 \\
дефицит: \textbf{kcif} [\emph{kʃif}] 0xC8A2 \\
дефлорация: \textbf{flo2n} [\emph{flo˩n}] 0x7C6C \\
децентрализация: \textbf{dri2s} [\emph{dri˩s}] 0x98D3 \\
децибел: \textbf{dwip} [\emph{dwip}] 0x4833 \\
децимация: \textbf{tci7c} [\emph{tʃi˥ʃ}] 0xF0AA \\
дециметр: \textbf{dxe7t} [\emph{dxe˥t}] 0xB3F3 \\
дешево: \textbf{srac} [\emph{sräʃ}] 0xE9C9 \\
джаз: \textbf{djas} [\emph{dʒäs}] 0x89B3 \\
джайнизм: \textbf{jyam} [\emph{ʒjäm}] 0x0915 \\
Джекоб: \textbf{hya7k} [\emph{ʰjä˥k}] 0x3118 \\
Дженни: \textbf{hji2n} [\emph{ʰʒi˩n}] 0x78A8 \\
Джибути: \textbf{gwut} [\emph{gwut}] 0xAA32 \\
Джимми: \textbf{hci7m} [\emph{ʰʃi˥m}] 0x1068 \\
джинсы: \textbf{gjis} [\emph{gʒis}] 0x88B2 \\
джип: \textbf{gjip} [\emph{gʒip}] 0x48B2 \\
джихад: \textbf{jya2t} [\emph{ʒjä˩t}] 0xB915 \\
джокер: \textbf{gwok} [\emph{gwok}] 0x2C32 \\
Джон: \textbf{jwon} [\emph{ʒwon}] 0x6C35 \\
Джордж: \textbf{qroc} [\emph{ŋroʃ}] 0xECD9 \\
диабет: \textbf{dyes} [\emph{djes}] 0x8B13 \\
диагностировать: \textbf{dyun} [\emph{djun}] 0x6A13 \\
диагональ: \textbf{dwon} [\emph{dwon}] 0x6C33 \\
диаграмма: \textbf{tfap} [\emph{tfäp}] 0x498A \\
диалекты: \textbf{dla2k} [\emph{dlä˩k}] 0x3973 \\
диалог: \textbf{dwac} [\emph{dwäʃ}] 0xE933 \\
диастола: \textbf{dza7t} [\emph{dzä˥t}] 0xB153 \\
диатомовый: \textbf{dya7m} [\emph{djä˥m}] 0x1113 \\
диатонический: \textbf{dya2n} [\emph{djä˩n}] 0x7913 \\
диван: \textbf{dvaf} [\emph{dväf}] 0xC993 \\
дивиденд: \textbf{dwik} [\emph{dwik}] 0x2833 \\
дидактический: \textbf{dxa2t} [\emph{dxä˩t}] 0xB9F3 \\
дизайнер: \textbf{dzic} [\emph{dziʃ}] 0xE853 \\
дизель: \textbf{dzes} [\emph{dzes}] 0x8B53 \\
дизентерия: \textbf{hdi2c} [\emph{ʰdi˩ʃ}] 0xF898 \\
дикари: \textbf{hsa7c} [\emph{ʰsä˥ʃ}] 0xF148 \\
дикий: \textbf{pyin} [\emph{pjin}] 0x6804 \\
диктатор: \textbf{dwot} [\emph{dwot}] 0xAC33 \\
дикция: \textbf{dxi2n} [\emph{dxi˩n}] 0x78F3 \\
дилер: \textbf{dyas} [\emph{djäs}] 0x8913 \\
диморфизм: \textbf{drom} [\emph{drom}] 0x0CD3 \\
динамический: \textbf{dvim} [\emph{dvim}] 0x0893 \\
династия: \textbf{hna2s} [\emph{ʰnä˩s}] 0x9930 \\
динозавр: \textbf{dyos} [\emph{djos}] 0x8C13 \\
диод: \textbf{dy6t} [\emph{djət}] 0xAD13 \\
диоксид: \textbf{dwis} [\emph{dwis}] 0x8833 \\
дионисическое: \textbf{dzi7s} [\emph{dzi˥s}] 0x9053 \\
диплоид: \textbf{dl67t} [\emph{dlə˥t}] 0xB573 \\
дипломат: \textbf{dli2t} [\emph{dli˩t}] 0xB873 \\
дипломатия: \textbf{dlom} [\emph{dlom}] 0x0C73 \\
директива: \textbf{hvit} [\emph{ʰvit}] 0xA8B0 \\
директора: \textbf{hri7s} [\emph{ʰri˥s}] 0x9070 \\
дирижабль: \textbf{hli7p} [\emph{ʰli˥p}] 0x5058 \\
дис: \textbf{dyi7s} [\emph{dji˥s}] 0x9013 \\
диск: \textbf{hdic} [\emph{ʰdiʃ}] 0xE898 \\
дисквалификация: \textbf{kla7k} [\emph{klä˥k}] 0x3162 \\
дискриминация: \textbf{krip} [\emph{krip}] 0x48C2 \\
диспепсия: \textbf{dwep} [\emph{dwep}] 0x4B33 \\
диссертаций: \textbf{fwes} [\emph{fwes}] 0x8B2C \\
дистрактор: \textbf{dr6t} [\emph{drət}] 0xADD3 \\
дисциплина: \textbf{dlin} [\emph{dlin}] 0x6873 \\
дифениламин: \textbf{dla7m} [\emph{dlä˥m}] 0x1173 \\
дифтерия: \textbf{dxi2t} [\emph{dxi˩t}] 0xB8F3 \\
дифтонг: \textbf{dyi2n} [\emph{dji˩n}] 0x7813 \\
диффузный: \textbf{kfus} [\emph{kfus}] 0x8A82 \\
дихотомический: \textbf{hd6m} [\emph{ʰdəm}] 0x0D98 \\
длинноногий: \textbf{hle2k} [\emph{ʰle˩k}] 0x3B58 \\
длинный: \textbf{lyan} [\emph{ljän}] 0x690B \\
дмитрий: \textbf{hd67t} [\emph{ʰdə˥t}] 0xB598 \\
до: \textbf{syat} [\emph{sjät}] 0xA909 \\
добавлений: \textbf{dva2t} [\emph{dvä˩t}] 0xB993 \\
доброволец: \textbf{swon} [\emph{swon}] 0x6C29 \\
доброжелатель: \textbf{hyi7c} [\emph{ʰji˥ʃ}] 0xF018 \\
добыча: \textbf{hjuc} [\emph{ʰʒuʃ}] 0xEAA8 \\
доверчивость: \textbf{hdi2n} [\emph{ʰdi˩n}] 0x7898 \\
довоенный: \textbf{pce2n} [\emph{pʃe˩n}] 0x7BA4 \\
доволен: \textbf{mwac} [\emph{mwäʃ}] 0xE921 \\
дог: \textbf{hma7f} [\emph{ʰmä˥f}] 0xD108 \\
догадка: \textbf{djut} [\emph{dʒut}] 0xAAB3 \\
догадок: \textbf{gjok} [\emph{gʒok}] 0x2CB2 \\
догматизм: \textbf{hda2m} [\emph{ʰdä˩m}] 0x1998 \\
договор: \textbf{tyot} [\emph{tjot}] 0xAC0A \\
додекаэдр: \textbf{dre7t} [\emph{dre˥t}] 0xB3D3 \\
дождемер: \textbf{hle7m} [\emph{ʰle˥m}] 0x1358 \\
дождь: \textbf{myic} [\emph{mjiʃ}] 0xE801 \\
дозировать: \textbf{dros} [\emph{dros}] 0x8CD3 \\
доисторический: \textbf{pxek} [\emph{pxek}] 0x2BE4 \\
доказательный: \textbf{tfak} [\emph{tfäk}] 0x298A \\
доктринер: \textbf{kyi2n} [\emph{kji˩n}] 0x7802 \\
документальный: \textbf{dwut} [\emph{dwut}] 0xAA33 \\
документация: \textbf{dlen} [\emph{dlen}] 0x6B73 \\
долг: \textbf{htup} [\emph{ʰtup}] 0x4A50 \\
долговязый: \textbf{hqa2k} [\emph{ʰŋä˩k}] 0x39C8 \\
должен: \textbf{pcic} [\emph{pʃiʃ}] 0xE8A4 \\
долина: \textbf{klik} [\emph{klik}] 0x2862 \\
доллар: \textbf{dlan} [\emph{dlän}] 0x6973 \\
дом: \textbf{hcas} [\emph{ʰʃäs}] 0x8968 \\
домен: \textbf{tyun} [\emph{tjun}] 0x6A0A \\
Доминика: \textbf{dyok} [\emph{djok}] 0x2C13 \\
доминиканский: \textbf{dxok} [\emph{dxok}] 0x2CF3 \\
домино: \textbf{dxim} [\emph{dxim}] 0x08F3 \\
домовой: \textbf{bvun} [\emph{bvun}] 0x6A91 \\
домогательство: \textbf{gjam} [\emph{gʒäm}] 0x09B2 \\
домой: \textbf{hco2t} [\emph{ʰʃo˩t}] 0xBC68 \\
донор: \textbf{dzot} [\emph{dzot}] 0xAC53 \\
допамин: \textbf{dwo7n} [\emph{dwo˥n}] 0x7433 \\
дополнение: \textbf{plok} [\emph{plok}] 0x2C64 \\
допустимый: \textbf{hyi2s} [\emph{ʰji˩s}] 0x9818 \\
дорогостоящий: \textbf{gxak} [\emph{gxäk}] 0x29F2 \\
дородный: \textbf{hke2t} [\emph{ʰke˩t}] 0xBB10 \\
доска: \textbf{pfat} [\emph{pfät}] 0xA984 \\
достаточно: \textbf{hton} [\emph{ʰton}] 0x6C50 \\
достиг: \textbf{hte2t} [\emph{ʰte˩t}] 0xBB50 \\
достигать: \textbf{hdaf} [\emph{ʰdäf}] 0xC998 \\
доступность: \textbf{slus} [\emph{slus}] 0x8A69 \\
доступный: \textbf{hpot} [\emph{ʰpot}] 0xAC20 \\
дочери: \textbf{nri7n} [\emph{nri˥n}] 0x70C6 \\
дошкольного: \textbf{pso2k} [\emph{pso˩k}] 0x3C44 \\
дравидский: \textbf{dvi7t} [\emph{dvi˥t}] 0xB093 \\
дразнить: \textbf{tfon} [\emph{tfon}] 0x6C8A \\
Дракон: \textbf{dron} [\emph{dron}] 0x6CD3 \\
Дракула: \textbf{hru7k} [\emph{ʰru˥k}] 0x3270 \\
драматично: \textbf{dxit} [\emph{dxit}] 0xA8F3 \\
драматург: \textbf{tsa2k} [\emph{tsä˩k}] 0x394A \\
драпировщик: \textbf{dzap} [\emph{dzäp}] 0x4953 \\
драхма: \textbf{hra2m} [\emph{ʰrä˩m}] 0x1970 \\
древесный: \textbf{hwu7t} [\emph{ʰwu˥t}] 0xB228 \\
дрейфовать: \textbf{hta2f} [\emph{ʰtä˩f}] 0xD950 \\
дремать: \textbf{dzos} [\emph{dzos}] 0x8C53 \\
дробовик: \textbf{cla2n} [\emph{ʃlä˩n}] 0x796D \\
дрова: \textbf{txa2k} [\emph{txä˩k}] 0x39EA \\
дрожжи: \textbf{cyum} [\emph{ʃjum}] 0x0A0D \\
дрок: \textbf{txa7s} [\emph{txä˥s}] 0x91EA \\
друг: \textbf{pyat} [\emph{pjät}] 0xA904 \\
Другие: \textbf{tcan} [\emph{tʃän}] 0x69AA \\
дружелюбие: \textbf{bzim} [\emph{bzim}] 0x0851 \\
дружки: \textbf{kci2n} [\emph{kʃi˩n}] 0x78A2 \\
друид: \textbf{dlu7t} [\emph{dlu˥t}] 0xB273 \\
Дрю: \textbf{drus} [\emph{drus}] 0x8AD3 \\
дуалистический: \textbf{dli2s} [\emph{dli˩s}] 0x9873 \\
дуб: \textbf{clun} [\emph{ʃlun}] 0x6A6D \\
думать: \textbf{hcik} [\emph{ʰʃik}] 0x2868 \\
Дунай: \textbf{hdu7p} [\emph{ʰdu˥p}] 0x5298 \\
дуплет: \textbf{dlo7t} [\emph{dlo˥t}] 0xB473 \\
дурак: \textbf{myuk} [\emph{mjuk}] 0x2A01 \\
дурачество: \textbf{tfi7s} [\emph{tfi˥s}] 0x908A \\
дуться: \textbf{hyo2k} [\emph{ʰjo˩k}] 0x3C18 \\
душа: \textbf{twan} [\emph{twän}] 0x692A \\
душеприказчица: \textbf{jyet} [\emph{ʒjet}] 0xAB15 \\
дуэль: \textbf{hduf} [\emph{ʰduf}] 0xCA98 \\
дым: \textbf{twuc} [\emph{twuʃ}] 0xEA2A \\
дыня: \textbf{kl6n} [\emph{klən}] 0x6D62 \\
дыхание: \textbf{hsaf} [\emph{ʰsäf}] 0xC948 \\
дышло: \textbf{hnu7p} [\emph{ʰnu˥p}] 0x5230 \\
дьявол: \textbf{hdan} [\emph{ʰdän}] 0x6998 \\
дьякон: \textbf{hd6s} [\emph{ʰdəs}] 0x8D98 \\
дюйм: \textbf{nwic} [\emph{nwiʃ}] 0xE826 \\
дядя: \textbf{ksuk} [\emph{ksuk}] 0x2A42 \\
дятел: \textbf{xlok} [\emph{xlok}] 0x2C7A \\
\subsection{ е }
е: \textbf{xli7m} [\emph{xli˥m}] 0x107A \\
евангелистский: \textbf{nya7s} [\emph{njä˥s}] 0x9106 \\
евгеника: \textbf{hyu2n} [\emph{ʰju˩n}] 0x7A18 \\
евклидов: \textbf{fl6t} [\emph{flət}] 0xAD6C \\
евнух: \textbf{nwa2n} [\emph{nwä˩n}] 0x7926 \\
\subsection{ Е }
Евразия: \textbf{vrus} [\emph{vrus}] 0x8AD6 \\
еврейка: \textbf{hyu7n} [\emph{ʰju˥n}] 0x7218 \\
евро: \textbf{vrof} [\emph{vrof}] 0xCCD6 \\
Европа: \textbf{hyup} [\emph{ʰjup}] 0x4A18 \\
европеизация: \textbf{ryuf} [\emph{rjuf}] 0xCA0E \\
Евфрат: \textbf{fra7t} [\emph{frä˥t}] 0xB1CC \\
евхаристия: \textbf{qrus} [\emph{ŋrus}] 0x8AD9 \\
еда: \textbf{kcak} [\emph{kʃäk}] 0x29A2 \\
единодушное: \textbf{jlit} [\emph{ʒlit}] 0xA875 \\
единорог: \textbf{nyi2n} [\emph{nji˩n}] 0x7806 \\
едок: \textbf{sw6k} [\emph{swək}] 0x2D29 \\
еж: \textbf{djek} [\emph{dʒek}] 0x2BB3 \\
ежевика: \textbf{bjip} [\emph{bʒip}] 0x48B1 \\
ежеквартальный: \textbf{rwap} [\emph{rwäp}] 0x492E \\
енот: \textbf{ryuk} [\emph{rjuk}] 0x2A0E \\
епископальный: \textbf{psi2p} [\emph{psi˩p}] 0x5844 \\
епископы: \textbf{hbu7p} [\emph{ʰbu˥p}] 0x5288 \\
ЕС: \textbf{vyun} [\emph{vjun}] 0x6A16 \\
еси: \textbf{hwi7s} [\emph{ʰwi˥s}] 0x9028 \\
Еф: \textbf{hfe7s} [\emph{ʰfe˥s}] 0x9360 \\
Ефрем: \textbf{fre2m} [\emph{fre˩m}] 0x1BCC \\
ЕЦБ: \textbf{byep} [\emph{bjep}] 0x4B11 \\
ЕЭС: \textbf{tce7k} [\emph{tʃe˥k}] 0x33AA \\
\subsection{ ж }
жадность: \textbf{glac} [\emph{gläʃ}] 0xE972 \\
жадностью: \textbf{vrec} [\emph{vreʃ}] 0xEBD6 \\
жажда: \textbf{hguc} [\emph{ʰguʃ}] 0xEA90 \\
жакеты: \textbf{hce7t} [\emph{ʰʃe˥t}] 0xB368 \\
жалоба: \textbf{blan} [\emph{blän}] 0x6971 \\
жаргон: \textbf{jyan} [\emph{ʒjän}] 0x6915 \\
жареные: \textbf{jrot} [\emph{ʒrot}] 0xACD5 \\
жаропонижающий: \textbf{bruc} [\emph{bruʃ}] 0xEAD1 \\
жасмин: \textbf{hmi7m} [\emph{ʰmi˥m}] 0x1008 \\
жатка: \textbf{swo7t} [\emph{swo˥t}] 0xB429 \\
жвачка: \textbf{hgum} [\emph{ʰgum}] 0x0A90 \\
жевание: \textbf{cwun} [\emph{ʃwun}] 0x6A2D \\
желание: \textbf{twin} [\emph{twin}] 0x682A \\
желатин: \textbf{jla2t} [\emph{ʒlä˩t}] 0xB975 \\
железо: \textbf{fyac} [\emph{fjäʃ}] 0xE90C \\
желто-зеленый: \textbf{xla7s} [\emph{xlä˥s}] 0x917A \\
желтовато-коричневый: \textbf{twa7n} [\emph{twä˥n}] 0x712A \\
желток: \textbf{jwik} [\emph{ʒwik}] 0x2835 \\
желтуха: \textbf{xya7n} [\emph{xjä˥n}] 0x711A \\
желтый: \textbf{clas} [\emph{ʃläs}] 0x896D \\
желудок: \textbf{hwet} [\emph{ʰwet}] 0xAB28 \\
желудочек: \textbf{nri7k} [\emph{nri˥k}] 0x30C6 \\
желчь: \textbf{hyo7t} [\emph{ʰjo˥t}] 0xB418 \\
жемчуг: \textbf{jrut} [\emph{ʒrut}] 0xAAD5 \\
жена: \textbf{hcis} [\emph{ʰʃis}] 0x8868 \\
\subsection{ Ж }
Женева: \textbf{hjef} [\emph{ʰʒef}] 0xCBA8 \\
жених: \textbf{hna2f} [\emph{ʰnä˩f}] 0xD930 \\
женоненавистник: \textbf{gxi2n} [\emph{gxi˩n}] 0x78F2 \\
женоненавистничество: \textbf{jli7n} [\emph{ʒli˥n}] 0x7075 \\
женский: \textbf{hjen} [\emph{ʰʒen}] 0x6BA8 \\
жеребенок: \textbf{txoc} [\emph{txoʃ}] 0xECEA \\
жертва: \textbf{cran} [\emph{ʃrän}] 0x69CD \\
жертвовать: \textbf{hdon} [\emph{ʰdon}] 0x6C98 \\
жертвы: \textbf{kri7t} [\emph{kri˥t}] 0xB0C2 \\
жестикуляция: \textbf{syi7n} [\emph{sji˥n}] 0x7009 \\
жестокий: \textbf{krac} [\emph{kräʃ}] 0xE9C2 \\
жестянщик: \textbf{tca7f} [\emph{tʃä˥f}] 0xD1AA \\
живой: \textbf{cwit} [\emph{ʃwit}] 0xA82D \\
живородящий: \textbf{tr6s} [\emph{trəs}] 0x8DCA \\
животное: \textbf{nwam} [\emph{nwäm}] 0x0926 \\
жидкости: \textbf{vwit} [\emph{vwit}] 0xA836 \\
жилет: \textbf{hwi7t} [\emph{ʰwi˥t}] 0xB028 \\
жирафа: \textbf{xrif} [\emph{xrif}] 0xC8DA \\
житница: \textbf{bjes} [\emph{bʒes}] 0x8BB1 \\
жокей: \textbf{hjok} [\emph{ʰʒok}] 0x2CA8 \\
жужжание: \textbf{hb6s} [\emph{ʰbəs}] 0x8D88 \\
жуки: \textbf{bxek} [\emph{bxek}] 0x2BF1 \\
жучок: \textbf{cy6t} [\emph{ʃjət}] 0xAD0D \\
жюри: \textbf{prut} [\emph{prut}] 0xAAC4 \\
\subsection{ з }
за: \textbf{pcin} [\emph{pʃin}] 0x68A4 \\
забавляться: \textbf{hla2k} [\emph{ʰlä˩k}] 0x3958 \\
забастовки: \textbf{hpa7k} [\emph{ʰpä˥k}] 0x3120 \\
забвение: \textbf{hyi2n} [\emph{ʰji˩n}] 0x7818 \\
забита: \textbf{hza7t} [\emph{ʰzä˥t}] 0xB1A0 \\
заблудившись: \textbf{hsu7t} [\emph{ʰsu˥t}] 0xB248 \\
заботливость: \textbf{hli2t} [\emph{ʰli˩t}] 0xB858 \\
заброшенный: \textbf{crit} [\emph{ʃrit}] 0xA8CD \\
забывать: \textbf{hwut} [\emph{ʰwut}] 0xAA28 \\
завернутый: \textbf{pxak} [\emph{pxäk}] 0x29E4 \\
завещание: \textbf{sya7t} [\emph{sjä˥t}] 0xB109 \\
завещать: \textbf{hvi2t} [\emph{ʰvi˩t}] 0xB8B0 \\
завидный: \textbf{hvi2n} [\emph{ʰvi˩n}] 0x78B0 \\
зависеть: \textbf{drat} [\emph{drät}] 0xA9D3 \\
завлечь: \textbf{drap} [\emph{dräp}] 0x49D3 \\
завоевывая: \textbf{gxin} [\emph{gxin}] 0x68F2 \\
завышенный: \textbf{qrek} [\emph{ŋrek}] 0x2BD9 \\
загар: \textbf{vrap} [\emph{vräp}] 0x49D6 \\
заговор: \textbf{hnom} [\emph{ʰnom}] 0x0C30 \\
заголовок: \textbf{hxe2k} [\emph{ʰxe˩k}] 0x3BD0 \\
загрязнение: \textbf{dlun} [\emph{dlun}] 0x6A73 \\
задергать: \textbf{drof} [\emph{drof}] 0xCCD3 \\
задерживаться: \textbf{djik} [\emph{dʒik}] 0x28B3 \\
задумчивый: \textbf{tsa2s} [\emph{tsä˩s}] 0x994A \\
зажимать: \textbf{hzi2t} [\emph{ʰzi˩t}] 0xB8A0 \\
заземление: \textbf{tfef} [\emph{tfef}] 0xCB8A \\
заика: \textbf{hca7m} [\emph{ʰʃä˥m}] 0x1168 \\
заикание: \textbf{kxo7t} [\emph{kxo˥t}] 0xB4E2 \\
\subsection{ З }
Заир: \textbf{hza7s} [\emph{ʰzä˥s}] 0x91A0 \\
зайчонок: \textbf{hve7t} [\emph{ʰve˥t}] 0xB3B0 \\
закал: \textbf{txo7n} [\emph{txo˥n}] 0x74EA \\
закладка: \textbf{bxak} [\emph{bxäk}] 0x29F1 \\
заклепка: \textbf{kli2t} [\emph{kli˩t}] 0xB862 \\
заклинатель: \textbf{sre7s} [\emph{sre˥s}] 0x93C9 \\
заклинать: \textbf{ksa2k} [\emph{ksä˩k}] 0x3942 \\
закон: \textbf{hfaf} [\emph{ʰfäf}] 0xC960 \\
законопослушный: \textbf{sla7k} [\emph{slä˥k}] 0x3169 \\
закупоренный: \textbf{zyot} [\emph{zjot}] 0xAC14 \\
закуска: \textbf{pwis} [\emph{pwis}] 0x8824 \\
зал: \textbf{hcup} [\emph{ʰʃup}] 0x4A68 \\
залогодержатель: \textbf{dxec} [\emph{dxeʃ}] 0xEBF3 \\
заложник: \textbf{jyic} [\emph{ʒjiʃ}] 0xE815 \\
Замбия: \textbf{zyap} [\emph{zjäp}] 0x4914 \\
замена: \textbf{sre2c} [\emph{sre˩ʃ}] 0xFBC9 \\
заместитель: \textbf{kren} [\emph{kren}] 0x6BC2 \\
заметный: \textbf{vlot} [\emph{vlot}] 0xAC76 \\
замечающий: \textbf{txi2t} [\emph{txi˩t}] 0xB8EA \\
замещенный: \textbf{txi7t} [\emph{txi˥t}] 0xB0EA \\
заминированный: \textbf{kxa7s} [\emph{kxä˥s}] 0x91E2 \\
замки: \textbf{hka7k} [\emph{ʰkä˥k}] 0x3110 \\
замок: \textbf{hkon} [\emph{ʰkon}] 0x6C10 \\
замороженные: \textbf{hgot} [\emph{ʰgot}] 0xAC90 \\
замуровывать: \textbf{mye7n} [\emph{mje˥n}] 0x7301 \\
Занзибар: \textbf{hna2p} [\emph{ʰnä˩p}] 0x5930 \\
занос: \textbf{kfi2n} [\emph{kfi˩n}] 0x7882 \\
занято: \textbf{pcut} [\emph{pʃut}] 0xAAA4 \\
заокеанский: \textbf{rye2n} [\emph{rje˩n}] 0x7B0E \\
заостренный: \textbf{tcif} [\emph{tʃif}] 0xC8AA \\
запад: \textbf{pcak} [\emph{pʃäk}] 0x29A4 \\
запаздывание: \textbf{cra7t} [\emph{ʃrä˥t}] 0xB1CD \\
запачканный: \textbf{pwo7t} [\emph{pwo˥t}] 0xB424 \\
запечатывание: \textbf{hm6f} [\emph{ʰməf}] 0xCD08 \\
записывать: \textbf{lyit} [\emph{ljit}] 0xA80B \\
заполнитель: \textbf{txa7f} [\emph{txä˥f}] 0xD1EA \\
заполнить: \textbf{hdin} [\emph{ʰdin}] 0x6898 \\
запомнить: \textbf{myat} [\emph{mjät}] 0xA901 \\
запор: \textbf{hki2m} [\emph{ʰki˩m}] 0x1810 \\
запрет: \textbf{pwon} [\emph{pwon}] 0x6C24 \\
Запрос: \textbf{twam} [\emph{twäm}] 0x092A \\
запятая: \textbf{kwam} [\emph{kwäm}] 0x0922 \\
зарабатывать: \textbf{zlan} [\emph{zlän}] 0x6974 \\
зарплаты: \textbf{cwep} [\emph{ʃwep}] 0x4B2D \\
зарядка: \textbf{tra2k} [\emph{trä˩k}] 0x39CA \\
засахаренный: \textbf{xri2t} [\emph{xri˩t}] 0xB8DA \\
заслонять: \textbf{bla2n} [\emph{blä˩n}] 0x7971 \\
заспать: \textbf{vluf} [\emph{vluf}] 0xCA76 \\
застава: \textbf{tfos} [\emph{tfos}] 0x8C8A \\
застежка: \textbf{hxat} [\emph{ʰxät}] 0xA9D0 \\
застежка-молния: \textbf{hzip} [\emph{ʰzip}] 0x48A0 \\
застенчивость: \textbf{crif} [\emph{ʃrif}] 0xC8CD \\
затвердение: \textbf{dre2n} [\emph{dre˩n}] 0x7BD3 \\
затемнение: \textbf{tw6t} [\emph{twət}] 0xAD2A \\
затопленный: \textbf{fla2n} [\emph{flä˩n}] 0x796C \\
затылок: \textbf{hzep} [\emph{ʰzep}] 0x4BA0 \\
Захария: \textbf{zra7c} [\emph{zrä˥ʃ}] 0xF1D4 \\
захват: \textbf{grap} [\emph{gräp}] 0x49D2 \\
зашикать: \textbf{hco7c} [\emph{ʰʃo˥ʃ}] 0xF468 \\
защелка: \textbf{sla7t} [\emph{slä˥t}] 0xB169 \\
защелками: \textbf{swa7s} [\emph{swä˥s}] 0x9129 \\
защиты: \textbf{tfen} [\emph{tfen}] 0x6B8A \\
звезда: \textbf{srat} [\emph{srät}] 0xA9C9 \\
звенящий: \textbf{hqi2n} [\emph{ʰŋi˩n}] 0x78C8 \\
звукоподражательный: \textbf{hn67p} [\emph{ʰnə˥p}] 0x5530 \\
здание: \textbf{plit} [\emph{plit}] 0xA864 \\
здоровье: \textbf{clac} [\emph{ʃläʃ}] 0xE96D \\
Здравствуйте: \textbf{hcan} [\emph{ʰʃän}] 0x6968 \\
зебра: \textbf{zlap} [\emph{zläp}] 0x4974 \\
Зевадия: \textbf{zya7t} [\emph{zjä˥t}] 0xB114 \\
зевать: \textbf{bwac} [\emph{bwäʃ}] 0xE931 \\
Зеведей: \textbf{zyet} [\emph{zjet}] 0xAB14 \\
зеландии: \textbf{zl6n} [\emph{zlən}] 0x6D74 \\
зеленщик: \textbf{xlos} [\emph{xlos}] 0x8C7A \\
зеленый: \textbf{ksin} [\emph{ksin}] 0x6842 \\
землевладельческий: \textbf{hji7s} [\emph{ʰʒi˥s}] 0x90A8 \\
землекоп: \textbf{bvaf} [\emph{bväf}] 0xC991 \\
землеройка: \textbf{xruc} [\emph{xruʃ}] 0xEADA \\
землетрясение: \textbf{hcem} [\emph{ʰʃem}] 0x0B68 \\
земля: \textbf{vla7t} [\emph{vlä˥t}] 0xB176 \\
земной: \textbf{ts6k} [\emph{tsək}] 0x2D4A \\
зенитная: \textbf{txa7m} [\emph{txä˥m}] 0x11EA \\
зеркала: \textbf{tri7s} [\emph{tri˥s}] 0x90CA \\
зерноядный: \textbf{gvis} [\emph{gvis}] 0x8892 \\
зефир: \textbf{hzif} [\emph{ʰzif}] 0xC8A0 \\
зигота: \textbf{hz6t} [\emph{ʰzət}] 0xADA0 \\
зима: \textbf{myit} [\emph{mjit}] 0xA801 \\
Зимбабве: \textbf{hzup} [\emph{ʰzup}] 0x4AA0 \\
зимородок: \textbf{tre7k} [\emph{tre˥k}] 0x33CA \\
зло: \textbf{cwak} [\emph{ʃwäk}] 0x292D \\
зловещий: \textbf{pcup} [\emph{pʃup}] 0x4AA4 \\
зловонный: \textbf{pxi7t} [\emph{pxi˥t}] 0xB0E4 \\
злокачественный: \textbf{gvan} [\emph{gvän}] 0x6992 \\
злорадствовать: \textbf{cyus} [\emph{ʃjus}] 0x8A0D \\
злость: \textbf{tsik} [\emph{tsik}] 0x284A \\
злоупотребление: \textbf{lyap} [\emph{ljäp}] 0x490B \\
змея: \textbf{hsef} [\emph{ʰsef}] 0xCB48 \\
знак: \textbf{bzin} [\emph{bzin}] 0x6851 \\
знакомства: \textbf{dxek} [\emph{dxek}] 0x2BF3 \\
знакомые: \textbf{flak} [\emph{fläk}] 0x296C \\
знаменатель: \textbf{hd6n} [\emph{ʰdən}] 0x6D98 \\
знаменоносец: \textbf{hxe7p} [\emph{ʰxe˥p}] 0x53D0 \\
знать: \textbf{hkut} [\emph{ʰkut}] 0xAA10 \\
значительность: \textbf{txut} [\emph{txut}] 0xAAEA \\
зоб: \textbf{sro2s} [\emph{sro˩s}] 0x9CC9 \\
зодиак: \textbf{vyac} [\emph{vjäʃ}] 0xE916 \\
золото: \textbf{hwun} [\emph{ʰwun}] 0x6A28 \\
золотуха: \textbf{klo7n} [\emph{klo˥n}] 0x7462 \\
Золушка: \textbf{nre7n} [\emph{nre˥n}] 0x73C6 \\
зонд: \textbf{trop} [\emph{trop}] 0x4CCA \\
зонирование: \textbf{cwi2n} [\emph{ʃwi˩n}] 0x782D \\
зонтик: \textbf{mras} [\emph{mräs}] 0x89C1 \\
зоология: \textbf{gloc} [\emph{gloʃ}] 0xEC72 \\
зоопарк: \textbf{hzuk} [\emph{ʰzuk}] 0x2AA0 \\
зороастрийский: \textbf{hzo7t} [\emph{ʰzo˥t}] 0xB4A0 \\
зубочистка: \textbf{tci2k} [\emph{tʃi˩k}] 0x38AA \\
зубчатый: \textbf{tfa7t} [\emph{tfä˥t}] 0xB18A \\
зубы: \textbf{tlis} [\emph{tlis}] 0x886A \\
зуд: \textbf{pyu7t} [\emph{pju˥t}] 0xB204 \\
зулусский: \textbf{zluk} [\emph{zluk}] 0x2A74 \\
зум: \textbf{zwum} [\emph{zwum}] 0x0A34 \\
зыбь: \textbf{gran} [\emph{grän}] 0x69D2 \\
зычный: \textbf{sre2n} [\emph{sre˩n}] 0x7BC9 \\
зяблик: \textbf{tfi2c} [\emph{tfi˩ʃ}] 0xF88A \\
зяблики: \textbf{pci2c} [\emph{pʃi˩ʃ}] 0xF8A4 \\
\subsection{ и }
иберийский: \textbf{bri7s} [\emph{bri˥s}] 0x90D1 \\
ива: \textbf{hwof} [\emph{ʰwof}] 0xCC28 \\
иволга: \textbf{ryof} [\emph{rjof}] 0xCC0E \\
иврит: \textbf{xrup} [\emph{xrup}] 0x4ADA \\
ивуариец: \textbf{vri2n} [\emph{vri˩n}] 0x78D6 \\
игла: \textbf{hjis} [\emph{ʰʒis}] 0x88A8 \\
иглу: \textbf{gli7k} [\emph{gli˥k}] 0x3072 \\
иго: \textbf{hko7k} [\emph{ʰko˥k}] 0x3410 \\
играть: \textbf{plan} [\emph{plän}] 0x6964 \\
идеально: \textbf{hwaf} [\emph{ʰwäf}] 0xC928 \\
\subsection{ И }
Идеально: \textbf{dlic} [\emph{dliʃ}] 0xE873 \\
идентичность: \textbf{hdif} [\emph{ʰdif}] 0xC898 \\
идеограмма: \textbf{dya7k} [\emph{djä˥k}] 0x3113 \\
идеология: \textbf{dlim} [\emph{dlim}] 0x0873 \\
идиоматический: \textbf{dwa7t} [\emph{dwä˥t}] 0xB133 \\
идиомы: \textbf{hdi7m} [\emph{ʰdi˥m}] 0x1098 \\
идиопатический: \textbf{dxa2k} [\emph{dxä˩k}] 0x39F3 \\
идиосинкразия: \textbf{dlis} [\emph{dlis}] 0x8873 \\
идиофон: \textbf{hdof} [\emph{ʰdof}] 0xCC98 \\
идолопоклонник: \textbf{pca7t} [\emph{pʃä˥t}] 0xB1A4 \\
идти: \textbf{hkin} [\emph{ʰkin}] 0x6810 \\
иды: \textbf{hdi2s} [\emph{ʰdi˩s}] 0x9898 \\
Иегова: \textbf{dxof} [\emph{dxof}] 0xCCF3 \\
Иезекииль: \textbf{hze7k} [\emph{ʰze˥k}] 0x33A0 \\
иезуитский: \textbf{hyi7t} [\emph{ʰji˥t}] 0xB018 \\
иена: \textbf{jyen} [\emph{ʒjen}] 0x6B15 \\
иерархия: \textbf{crus} [\emph{ʃrus}] 0x8ACD \\
Иеремия: \textbf{jrim} [\emph{ʒrim}] 0x08D5 \\
иероглифы: \textbf{xr6s} [\emph{xrəs}] 0x8DDA \\
Иерусалим: \textbf{hjum} [\emph{ʰʒum}] 0x0AA8 \\
избавлены: \textbf{hca7t} [\emph{ʰʃä˥t}] 0xB168 \\
избежать: \textbf{hbit} [\emph{ʰbit}] 0xA888 \\
избиратель: \textbf{bwam} [\emph{bwäm}] 0x0931 \\
избыток: \textbf{klok} [\emph{klok}] 0x2C62 \\
известкование: \textbf{kfa7k} [\emph{kfä˥k}] 0x3182 \\
извиваться: \textbf{hzok} [\emph{ʰzok}] 0x2CA0 \\
извилистый: \textbf{vran} [\emph{vrän}] 0x69D6 \\
извинить: \textbf{hme7t} [\emph{ʰme˥t}] 0xB308 \\
изготовление: \textbf{htis} [\emph{ʰtis}] 0x8850 \\
издавна: \textbf{qya7t} [\emph{ŋjä˥t}] 0xB119 \\
издатель: \textbf{plet} [\emph{plet}] 0xAB64 \\
Изе: \textbf{dji7s} [\emph{dʒi˥s}] 0x90B3 \\
излишество: \textbf{zyut} [\emph{zjut}] 0xAA14 \\
излучающий: \textbf{hm6n} [\emph{ʰmən}] 0x6D08 \\
изменник: \textbf{hma2t} [\emph{ʰmä˩t}] 0xB908 \\
изменчивый: \textbf{rw6n} [\emph{rwən}] 0x6D2E \\
изменяемую: \textbf{flet} [\emph{flet}] 0xAB6C \\
изморось: \textbf{dzis} [\emph{dzis}] 0x8853 \\
изнасилование: \textbf{blaf} [\emph{bläf}] 0xC971 \\
изнашивания: \textbf{hxuk} [\emph{ʰxuk}] 0x2AD0 \\
изнеженность: \textbf{hje2n} [\emph{ʰʒe˩n}] 0x7BA8 \\
изнурение: \textbf{tce7n} [\emph{tʃe˥n}] 0x73AA \\
изобара: \textbf{hq6p} [\emph{ʰŋəp}] 0x4DC8 \\
изобличают: \textbf{bli7n} [\emph{bli˥n}] 0x7071 \\
изображенный: \textbf{dra2s} [\emph{drä˩s}] 0x99D3 \\
изобретенной: \textbf{tci7t} [\emph{tʃi˥t}] 0xB0AA \\
изолятор: \textbf{cl6t} [\emph{ʃlət}] 0xAD6D \\
изоляция: \textbf{klis} [\emph{klis}] 0x8862 \\
изомер: \textbf{zwim} [\emph{zwim}] 0x0834 \\
изометрический: \textbf{swi7k} [\emph{swi˥k}] 0x3029 \\
изоморфизм: \textbf{zrom} [\emph{zrom}] 0x0CD4 \\
изопрен: \textbf{sri2p} [\emph{sri˩p}] 0x58C9 \\
изотерма: \textbf{hqem} [\emph{ʰŋem}] 0x0BC8 \\
изотонический: \textbf{sy67n} [\emph{sjə˥n}] 0x7509 \\
изотоп: \textbf{zwop} [\emph{zwop}] 0x4C34 \\
изотропный: \textbf{sro2k} [\emph{sro˩k}] 0x3CC9 \\
Израиль: \textbf{zris} [\emph{zris}] 0x88D4 \\
израсходовать: \textbf{kfak} [\emph{kfäk}] 0x2982 \\
изумруд: \textbf{zrut} [\emph{zrut}] 0xAAD4 \\
изымать: \textbf{bjot} [\emph{bʒot}] 0xACB1 \\
Иисус: \textbf{hjus} [\emph{ʰʒus}] 0x8AA8 \\
иконоборец: \textbf{kxa7k} [\emph{kxä˥k}] 0x31E2 \\
икота: \textbf{kfok} [\emph{kfok}] 0x2C82 \\
икра: \textbf{kya7c} [\emph{kjä˥ʃ}] 0xF102 \\
Илиада: \textbf{hli7t} [\emph{ʰli˥t}] 0xB058 \\
иллюзия: \textbf{blic} [\emph{bliʃ}] 0xE871 \\
иллюстратор: \textbf{sro2t} [\emph{sro˩t}] 0xBCC9 \\
имитатор: \textbf{myo7t} [\emph{mjo˥t}] 0xB401 \\
иммобилизация: \textbf{bxin} [\emph{bxin}] 0x68F1 \\
иммунный: \textbf{myim} [\emph{mjim}] 0x0801 \\
иммунология: \textbf{myi2c} [\emph{mji˩ʃ}] 0xF801 \\
император: \textbf{krat} [\emph{krät}] 0xA9C2 \\
императрица: \textbf{xra2n} [\emph{xrä˩n}] 0x79DA \\
имперский: \textbf{prim} [\emph{prim}] 0x08C4 \\
Импортировать: \textbf{cyot} [\emph{ʃjot}] 0xAC0D \\
импрессионизм: \textbf{pci2s} [\emph{pʃi˩s}] 0x98A4 \\
импровизация: \textbf{pci2n} [\emph{pʃi˩n}] 0x78A4 \\
импровизировать: \textbf{tru7s} [\emph{tru˥s}] 0x92CA \\
импульсов: \textbf{hmo2k} [\emph{ʰmo˩k}] 0x3C08 \\
имущественный: \textbf{hke7n} [\emph{ʰke˥n}] 0x7310 \\
имя: \textbf{hnim} [\emph{ʰnim}] 0x0830 \\
инбридинг: \textbf{hni7n} [\emph{ʰni˥n}] 0x7030 \\
инвертировать: \textbf{grut} [\emph{grut}] 0xAAD2 \\
инвестиции: \textbf{nwes} [\emph{nwes}] 0x8B26 \\
ИНГ: \textbf{jri7n} [\emph{ʒri˥n}] 0x70D5 \\
ингалятор: \textbf{jlip} [\emph{ʒlip}] 0x4875 \\
ингибирование: \textbf{xyit} [\emph{xjit}] 0xA81A \\
ингибитор: \textbf{ny6t} [\emph{njət}] 0xAD06 \\
ингредиенты: \textbf{xrin} [\emph{xrin}] 0x68DA \\
индеек: \textbf{xyuk} [\emph{xjuk}] 0x2A1A \\
индексация: \textbf{nye7n} [\emph{nje˥n}] 0x7306 \\
индивидуальные: \textbf{qri2t} [\emph{ŋri˩t}] 0xB8D9 \\
индийский: \textbf{nr6t} [\emph{nrət}] 0xADC6 \\
индис: \textbf{nyi2t} [\emph{nji˩t}] 0xB806 \\
индифферентный: \textbf{lya7c} [\emph{ljä˥ʃ}] 0xF10B \\
индол: \textbf{dli7s} [\emph{dli˥s}] 0x9073 \\
Индонезия: \textbf{nyec} [\emph{njeʃ}] 0xEB06 \\
Индостан: \textbf{hnu2s} [\emph{ʰnu˩s}] 0x9A30 \\
индуизм: \textbf{nwum} [\emph{nwum}] 0x0A26 \\
индуктивность: \textbf{dyen} [\emph{djen}] 0x6B13 \\
индуктивный: \textbf{dyi7t} [\emph{dji˥t}] 0xB013 \\
индуктор: \textbf{nyus} [\emph{njus}] 0x8A06 \\
индюков: \textbf{txu2k} [\emph{txu˩k}] 0x3AEA \\
иней: \textbf{swuc} [\emph{swuʃ}] 0xEA29 \\
инертность: \textbf{gxac} [\emph{gxäʃ}] 0xE9F2 \\
инжиниринг: \textbf{jyin} [\emph{ʒjin}] 0x6815 \\
инквизиция: \textbf{hka2c} [\emph{ʰkä˩ʃ}] 0xF910 \\
инкрустация: \textbf{ks6c} [\emph{ksəʃ}] 0xED42 \\
иннервация: \textbf{nri7p} [\emph{nri˥p}] 0x50C6 \\
инновация: \textbf{nwos} [\emph{nwos}] 0x8C26 \\
иносказательный: \textbf{lwa2n} [\emph{lwä˩n}] 0x792B \\
иностранный: \textbf{kwic} [\emph{kwiʃ}] 0xE822 \\
инстилляция: \textbf{tl62n} [\emph{tlə˩n}] 0x7D6A \\
инстинкт: \textbf{srin} [\emph{srin}] 0x68C9 \\
инструмент: \textbf{srit} [\emph{srit}] 0xA8C9 \\
инструменты: \textbf{kluc} [\emph{kluʃ}] 0xEA62 \\
инсулин: \textbf{zlun} [\emph{zlun}] 0x6A74 \\
инталия: \textbf{tyi2k} [\emph{tji˩k}] 0x380A \\
интеллектуальный: \textbf{zli7k} [\emph{zli˥k}] 0x3074 \\
интервалы: \textbf{tri2k} [\emph{tri˩k}] 0x38CA \\
интерес: \textbf{lyas} [\emph{ljäs}] 0x890B \\
интеркостельный: \textbf{nyes} [\emph{njes}] 0x8B06 \\
интернационализация: \textbf{jyac} [\emph{ʒjäʃ}] 0xE915 \\
интерпелляция: \textbf{tl67n} [\emph{tlə˥n}] 0x756A \\
интерполяция: \textbf{tl6s} [\emph{tləs}] 0x8D6A \\
интерпретировать: \textbf{tyif} [\emph{tjif}] 0xC80A \\
интерферометр: \textbf{nro2m} [\emph{nro˩m}] 0x1CC6 \\
интонировать: \textbf{nru7n} [\emph{nru˥n}] 0x72C6 \\
интровертированный: \textbf{tri2c} [\emph{tri˩ʃ}] 0xF8CA \\
инфекционный: \textbf{nruk} [\emph{nruk}] 0x2AC6 \\
инфинитив: \textbf{hnit} [\emph{ʰnit}] 0xA830 \\
инфракрасный: \textbf{grec} [\emph{greʃ}] 0xEBD2 \\
инфраструктура: \textbf{frot} [\emph{frot}] 0xACCC \\
ион: \textbf{lyun} [\emph{ljun}] 0x6A0B \\
Иона: \textbf{hyo2n} [\emph{ʰjo˩n}] 0x7C18 \\
ионосфера: \textbf{nyof} [\emph{njof}] 0xCC06 \\
Иордания: \textbf{jron} [\emph{ʒron}] 0x6CD5 \\
ипостась: \textbf{pla2s} [\emph{plä˩s}] 0x9964 \\
ипохондрик: \textbf{pyi2n} [\emph{pji˩n}] 0x7804 \\
ипохондрия: \textbf{hp62t} [\emph{ʰpə˩t}] 0xBD20 \\
ипподром: \textbf{pri7m} [\emph{pri˥m}] 0x10C4 \\
Ирак: \textbf{hra2k} [\emph{ʰrä˩k}] 0x3970 \\
иридий: \textbf{ryom} [\emph{rjom}] 0x0C0E \\
Ирландия: \textbf{dl6n} [\emph{dlən}] 0x6D73 \\
ирландцы: \textbf{qric} [\emph{ŋriʃ}] 0xE8D9 \\
иррациональный: \textbf{ry6n} [\emph{rjən}] 0x6D0E \\
Исайя: \textbf{hya7s} [\emph{ʰjä˥s}] 0x9118 \\
искатель: \textbf{fret} [\emph{fret}] 0xABCC \\
исключительность: \textbf{klu7s} [\emph{klu˥s}] 0x9262 \\
искусный: \textbf{hwe2n} [\emph{ʰwe˩n}] 0x7B28 \\
искусственный: \textbf{jrat} [\emph{ʒrät}] 0xA9D5 \\
искушение: \textbf{hyoc} [\emph{ʰjoʃ}] 0xEC18 \\
ислам: \textbf{slum} [\emph{slum}] 0x0A69 \\
Исландия: \textbf{pli7t} [\emph{pli˥t}] 0xB064 \\
испанский: \textbf{ps6n} [\emph{psən}] 0x6D44 \\
испарение: \textbf{tx6n} [\emph{txən}] 0x6DEA \\
испачканный: \textbf{dra7k} [\emph{drä˥k}] 0x31D3 \\
испольщик: \textbf{crof} [\emph{ʃrof}] 0xCCCD \\
испорченный: \textbf{vyet} [\emph{vjet}] 0xAB16 \\
испражняться: \textbf{tfa2k} [\emph{tfä˩k}] 0x398A \\
испытывающий: \textbf{cyoc} [\emph{ʃjoʃ}] 0xEC0D \\
исследователь: \textbf{slec} [\emph{sleʃ}] 0xEB69 \\
истекающий: \textbf{srec} [\emph{sreʃ}] 0xEBC9 \\
историки: \textbf{sri7t} [\emph{sri˥t}] 0xB0C9 \\
исторический: \textbf{sris} [\emph{sris}] 0x88C9 \\
история: \textbf{ksis} [\emph{ksis}] 0x8842 \\
источники: \textbf{hso2n} [\emph{ʰso˩n}] 0x7C48 \\
истребование: \textbf{pfik} [\emph{pfik}] 0x2884 \\
Италия: \textbf{xlis} [\emph{xlis}] 0x887A \\
итерация: \textbf{dret} [\emph{dret}] 0xABD3 \\
иттрий: \textbf{xrim} [\emph{xrim}] 0x08DA \\
Иуда: \textbf{hyu2t} [\emph{ʰju˩t}] 0xBA18 \\
иудаизм: \textbf{hdum} [\emph{ʰdum}] 0x0A98 \\
их: \textbf{htic} [\emph{ʰtiʃ}] 0xE850 \\
ихтиология: \textbf{fl6k} [\emph{flək}] 0x2D6C \\
ишиас: \textbf{tci7k} [\emph{tʃi˥k}] 0x30AA \\
Ишмаэль: \textbf{hsi7m} [\emph{ʰsi˥m}] 0x1048 \\
июль: \textbf{clis} [\emph{ʃlis}] 0x886D \\
июнь: \textbf{kwun} [\emph{kwun}] 0x6A22 \\
\subsection{ й }
йеллоунайф: \textbf{lwef} [\emph{lwef}] 0xCB2B \\
\subsection{ Й }
Йемен: \textbf{hy6m} [\emph{ʰjəm}] 0x0D18 \\
йога: \textbf{jwak} [\emph{ʒwäk}] 0x2935 \\
йогурт: \textbf{dxot} [\emph{dxot}] 0xACF3 \\
йод: \textbf{hyom} [\emph{ʰjom}] 0x0C18 \\
йодированной: \textbf{dyo7t} [\emph{djo˥t}] 0xB413 \\
йодль: \textbf{hy67t} [\emph{ʰjə˥t}] 0xB518 \\
йота: \textbf{cla7t} [\emph{ʃlä˥t}] 0xB16D \\
Йоханнесбург: \textbf{hxa7p} [\emph{ʰxä˥p}] 0x51D0 \\
\subsection{ к }
кабала: \textbf{kla7p} [\emph{klä˥p}] 0x5162 \\
кабестан: \textbf{ksu7n} [\emph{ksu˥n}] 0x7242 \\
кабина: \textbf{hken} [\emph{ʰken}] 0x6B10 \\
\subsection{ К }
Кабул: \textbf{hku7p} [\emph{ʰku˥p}] 0x5210 \\
кавалерия: \textbf{kli2n} [\emph{kli˩n}] 0x7862 \\
каверны: \textbf{hvep} [\emph{ʰvep}] 0x4BB0 \\
Кавказ: \textbf{kca7s} [\emph{kʃä˥s}] 0x91A2 \\
кадмий: \textbf{kyem} [\emph{kjem}] 0x0B02 \\
казалось: \textbf{sri2t} [\emph{sri˩t}] 0xB8C9 \\
Казахстан: \textbf{hzos} [\emph{ʰzos}] 0x8CA0 \\
казеин: \textbf{kse7n} [\emph{kse˥n}] 0x7342 \\
казино: \textbf{gzon} [\emph{gzon}] 0x6C52 \\
казуистический: \textbf{ksu7k} [\emph{ksu˥k}] 0x3242 \\
Каир: \textbf{hka7c} [\emph{ʰkä˥ʃ}] 0xF110 \\
как: \textbf{kyan} [\emph{kjän}] 0x6902 \\
как-то: \textbf{htoc} [\emph{ʰtoʃ}] 0xEC50 \\
какао: \textbf{kxok} [\emph{kxok}] 0x2CE2 \\
какие: \textbf{hkam} [\emph{ʰkäm}] 0x0910 \\
какофония: \textbf{kfa2n} [\emph{kfä˩n}] 0x7982 \\
кактус: \textbf{gjas} [\emph{gʒäs}] 0x89B2 \\
каламбур: \textbf{swa7n} [\emph{swä˥n}] 0x7129 \\
калейдоскоп: \textbf{kxo7k} [\emph{kxo˥k}] 0x34E2 \\
календарь: \textbf{jlet} [\emph{ʒlet}] 0xAB75 \\
календы: \textbf{kla2s} [\emph{klä˩s}] 0x9962 \\
калечить: \textbf{hge7n} [\emph{ʰge˥n}] 0x7390 \\
калий: \textbf{plim} [\emph{plim}] 0x0864 \\
Калифорния: \textbf{hla2n} [\emph{ʰlä˩n}] 0x7958 \\
каллиграф: \textbf{kla7f} [\emph{klä˥f}] 0xD162 \\
калория: \textbf{klop} [\emph{klop}] 0x4C62 \\
кальвинист: \textbf{lwa7n} [\emph{lwä˥n}] 0x712B \\
калька: \textbf{vlec} [\emph{vleʃ}] 0xEB76 \\
калькулятор: \textbf{glut} [\emph{glut}] 0xAA72 \\
Калькутта: \textbf{hla7k} [\emph{ʰlä˥k}] 0x3158 \\
кальций: \textbf{klif} [\emph{klif}] 0xC862 \\
кальцит: \textbf{flus} [\emph{flus}] 0x8A6C \\
камбала: \textbf{pwa7t} [\emph{pwä˥t}] 0xB124 \\
камбуз: \textbf{hku2s} [\emph{ʰku˩s}] 0x9A10 \\
камелия: \textbf{kla2m} [\emph{klä˩m}] 0x1962 \\
каменщик: \textbf{hme2n} [\emph{ʰme˩n}] 0x7B08 \\
камень: \textbf{pwac} [\emph{pwäʃ}] 0xE924 \\
Камерун: \textbf{jrun} [\emph{ʒrun}] 0x6AD5 \\
камеры: \textbf{kra2c} [\emph{krä˩ʃ}] 0xF9C2 \\
камуфляж: \textbf{kfa7n} [\emph{kfä˥n}] 0x7182 \\
Канада: \textbf{hx6t} [\emph{ʰxət}] 0xADD0 \\
канал: \textbf{clen} [\emph{ʃlen}] 0x6B6D \\
канарейка: \textbf{hk6c} [\emph{ʰkəʃ}] 0xED10 \\
канату: \textbf{tri7f} [\emph{tri˥f}] 0xD0CA \\
кандидат: \textbf{zrat} [\emph{zrät}] 0xA9D4 \\
каннибал: \textbf{hki2p} [\emph{ʰki˩p}] 0x5810 \\
канон: \textbf{kya7n} [\emph{kjä˥n}] 0x7102 \\
канонизация: \textbf{nwi7s} [\emph{nwi˥s}] 0x9026 \\
каноэ: \textbf{hkos} [\emph{ʰkos}] 0x8C10 \\
кантонский: \textbf{kxo7n} [\emph{kxo˥n}] 0x74E2 \\
канцелярия: \textbf{hca2s} [\emph{ʰʃä˩s}] 0x9968 \\
капать: \textbf{drip} [\emph{drip}] 0x48D3 \\
капающий: \textbf{tsi7k} [\emph{tsi˥k}] 0x304A \\
капер: \textbf{pya7t} [\emph{pjä˥t}] 0xB104 \\
капиллярный: \textbf{plek} [\emph{plek}] 0x2B64 \\
капитан: \textbf{klat} [\emph{klät}] 0xA962 \\
капитуляция: \textbf{txac} [\emph{txäʃ}] 0xE9EA \\
капли: \textbf{gvot} [\emph{gvot}] 0xAC92 \\
каприз: \textbf{kri2k} [\emph{kri˩k}] 0x38C2 \\
капсула: \textbf{kwus} [\emph{kwus}] 0x8A22 \\
капуста: \textbf{hgip} [\emph{ʰgip}] 0x4890 \\
капуцин: \textbf{hku2c} [\emph{ʰku˩ʃ}] 0xFA10 \\
карабин: \textbf{kxa7n} [\emph{kxä˥n}] 0x71E2 \\
каракуль: \textbf{tra7m} [\emph{trä˥m}] 0x11CA \\
карамель: \textbf{zram} [\emph{zräm}] 0x09D4 \\
карандаш: \textbf{pcip} [\emph{pʃip}] 0x48A4 \\
карапуз: \textbf{tfo2t} [\emph{tfo˩t}] 0xBC8A \\
каратэ: \textbf{vr6t} [\emph{vrət}] 0xADD6 \\
карбид: \textbf{kra2p} [\emph{krä˩p}] 0x59C2 \\
карбоксил: \textbf{kwi7k} [\emph{kwi˥k}] 0x3022 \\
карболка: \textbf{kra7k} [\emph{krä˥k}] 0x31C2 \\
карбюратор: \textbf{bra2c} [\emph{brä˩ʃ}] 0xF9D1 \\
кардинальный: \textbf{krin} [\emph{krin}] 0x68C2 \\
кариатида: \textbf{jra2t} [\emph{ʒrä˩t}] 0xB9D5 \\
карибский: \textbf{kr6p} [\emph{krəp}] 0x4DC2 \\
карибу: \textbf{kri2p} [\emph{kri˩p}] 0x58C2 \\
карикатурист: \textbf{rwa7s} [\emph{rwä˥s}] 0x912E \\
Каринтии: \textbf{kr6f} [\emph{krəf}] 0xCDC2 \\
каркать: \textbf{kwa2k} [\emph{kwä˩k}] 0x3922 \\
карликовый: \textbf{hne7n} [\emph{ʰne˥n}] 0x7330 \\
карман: \textbf{clit} [\emph{ʃlit}] 0xA86D \\
карниз: \textbf{gven} [\emph{gven}] 0x6B92 \\
карп: \textbf{kra7p} [\emph{krä˥p}] 0x51C2 \\
карта: \textbf{jrac} [\emph{ʒräʃ}] 0xE9D5 \\
картель: \textbf{kre2t} [\emph{kre˩t}] 0xBBC2 \\
картечь: \textbf{kco2t} [\emph{kʃo˩t}] 0xBCA2 \\
картография: \textbf{tra2f} [\emph{trä˩f}] 0xD9CA \\
картон: \textbf{kxan} [\emph{kxän}] 0x69E2 \\
картошка: \textbf{plut} [\emph{plut}] 0xAA64 \\
картуш: \textbf{hko7c} [\emph{ʰko˥ʃ}] 0xF410 \\
касательный: \textbf{ty6n} [\emph{tjən}] 0x6D0A \\
Каспийская: \textbf{kx6p} [\emph{kxəp}] 0x4DE2 \\
кассета: \textbf{kces} [\emph{kʃes}] 0x8BA2 \\
каста: \textbf{ksa2t} [\emph{ksä˩t}] 0xB942 \\
кастаньеты: \textbf{hke7p} [\emph{ʰke˥p}] 0x5310 \\
кастильский: \textbf{kli7n} [\emph{kli˥n}] 0x7062 \\
кастрировать: \textbf{gjaf} [\emph{gʒäf}] 0xC9B2 \\
катаболизм: \textbf{kla7m} [\emph{klä˥m}] 0x1162 \\
катализ: \textbf{tle7s} [\emph{tle˥s}] 0x936A \\
каталог: \textbf{dlet} [\emph{dlet}] 0xAB73 \\
Каталония: \textbf{tlo7n} [\emph{tlo˥n}] 0x746A \\
катамаран: \textbf{tx6m} [\emph{txəm}] 0x0DEA \\
катаракта: \textbf{pra2t} [\emph{prä˩t}] 0xB9C4 \\
катарсис: \textbf{tfa7s} [\emph{tfä˥s}] 0x918A \\
катафалк: \textbf{hto2s} [\emph{ʰto˩s}] 0x9C50 \\
катехизис: \textbf{tya2s} [\emph{tjä˩s}] 0x990A \\
Катись: \textbf{kra2m} [\emph{krä˩m}] 0x19C2 \\
катод: \textbf{kfet} [\emph{kfet}] 0xAB82 \\
каток: \textbf{gzik} [\emph{gzik}] 0x2852 \\
каучук: \textbf{kca7c} [\emph{kʃä˥ʃ}] 0xF1A2 \\
кафе: \textbf{hkef} [\emph{ʰkef}] 0xCB10 \\
кафтан: \textbf{kfa2t} [\emph{kfä˩t}] 0xB982 \\
качественный: \textbf{klaf} [\emph{kläf}] 0xC962 \\
кашель: \textbf{ksos} [\emph{ksos}] 0x8C42 \\
Кашмир: \textbf{kci7m} [\emph{kʃi˥m}] 0x10A2 \\
каштан: \textbf{kx6n} [\emph{kxən}] 0x6DE2 \\
квадрат: \textbf{kwap} [\emph{kwäp}] 0x4922 \\
квадратура: \textbf{dr6c} [\emph{drəʃ}] 0xEDD3 \\
квадрильон: \textbf{dlo7n} [\emph{dlo˥n}] 0x7473 \\
квазар: \textbf{hka7s} [\emph{ʰkä˥s}] 0x9110 \\
квакер: \textbf{hke2k} [\emph{ʰke˩k}] 0x3B10 \\
квалифицированный: \textbf{hwok} [\emph{ʰwok}] 0x2C28 \\
квантификатором: \textbf{kfa2f} [\emph{kfä˩f}] 0xD982 \\
квантовая: \textbf{kfam} [\emph{kfäm}] 0x0982 \\
кварк: \textbf{kw6k} [\emph{kwək}] 0x2D22 \\
квартет: \textbf{kru7t} [\emph{kru˥t}] 0xB2C2 \\
квартиры: \textbf{vlat} [\emph{vlät}] 0xA976 \\
кварц: \textbf{kros} [\emph{kros}] 0x8CC2 \\
квасцы: \textbf{hla2m} [\emph{ʰlä˩m}] 0x1958 \\
квесты: \textbf{ksu2s} [\emph{ksu˩s}] 0x9A42 \\
квинтет: \textbf{kxe2t} [\emph{kxe˩t}] 0xBBE2 \\
квинтэссенция: \textbf{kse7s} [\emph{kse˥s}] 0x9342 \\
квитанция: \textbf{ryit} [\emph{rjit}] 0xA80E \\
кворум: \textbf{kro7m} [\emph{kro˥m}] 0x14C2 \\
кекс: \textbf{kyek} [\emph{kjek}] 0x2B02 \\
кельтская: \textbf{kli7t} [\emph{kli˥t}] 0xB062 \\
кемпинг: \textbf{kwe2t} [\emph{kwe˩t}] 0xBB22 \\
кенгуру: \textbf{hquk} [\emph{ʰŋuk}] 0x2AC8 \\
Кения: \textbf{xy6n} [\emph{xjən}] 0x6D1A \\
керамика: \textbf{frip} [\emph{frip}] 0x48CC \\
кератин: \textbf{kr67n} [\emph{krə˥n}] 0x75C2 \\
кервель: \textbf{cwef} [\emph{ʃwef}] 0xCB2D \\
кессон: \textbf{kce7n} [\emph{kʃe˥n}] 0x73A2 \\
кетон: \textbf{kxo2n} [\emph{kxo˩n}] 0x7CE2 \\
кетчуп: \textbf{kcep} [\emph{kʃep}] 0x4BA2 \\
кефаль: \textbf{tlo7t} [\emph{tlo˥t}] 0xB46A \\
кеч: \textbf{kfec} [\emph{kfeʃ}] 0xEB82 \\
кибернетический: \textbf{bre2n} [\emph{bre˩n}] 0x7BD1 \\
киберпространство: \textbf{srep} [\emph{srep}] 0x4BC9 \\
кивнул: \textbf{hna2t} [\emph{ʰnä˩t}] 0xB930 \\
Киев: \textbf{hki2f} [\emph{ʰki˩f}] 0xD810 \\
кило: \textbf{kyom} [\emph{kjom}] 0x0C02 \\
киловатт: \textbf{kwif} [\emph{kwif}] 0xC822 \\
килограмм: \textbf{kwom} [\emph{kwom}] 0x0C22 \\
километр: \textbf{klom} [\emph{klom}] 0x0C62 \\
кими: \textbf{bl6p} [\emph{bləp}] 0x4D71 \\
кинематика: \textbf{hn62k} [\emph{ʰnə˩k}] 0x3D30 \\
кинестетическая: \textbf{kfef} [\emph{kfef}] 0xCB82 \\
кинетика: \textbf{gyek} [\emph{gjek}] 0x2B12 \\
кинжал: \textbf{hda7s} [\emph{ʰdä˥s}] 0x9198 \\
кинотеатр: \textbf{vyin} [\emph{vjin}] 0x6816 \\
киоск: \textbf{pxat} [\emph{pxät}] 0xA9E4 \\
Киото: \textbf{kwo7t} [\emph{kwo˥t}] 0xB422 \\
Кипр: \textbf{sr6s} [\emph{srəs}] 0x8DC9 \\
кипятильник: \textbf{tsa2c} [\emph{tsä˩ʃ}] 0xF94A \\
Киргизия: \textbf{ry6s} [\emph{rjəs}] 0x8D0E \\
кириллица: \textbf{cli7k} [\emph{ʃli˥k}] 0x306D \\
киркомотыга: \textbf{hxi2k} [\emph{ʰxi˩k}] 0x38D0 \\
кирпич: \textbf{hto2t} [\emph{ʰto˩t}] 0xBC50 \\
кислород: \textbf{ksen} [\emph{ksen}] 0x6B42 \\
кислота: \textbf{flat} [\emph{flät}] 0xA96C \\
китолов: \textbf{hce7n} [\emph{ʰʃe˥n}] 0x7368 \\
киты: \textbf{cyes} [\emph{ʃjes}] 0x8B0D \\
кишечник: \textbf{bwet} [\emph{bwet}] 0xAB31 \\
клавиатура: \textbf{kwot} [\emph{kwot}] 0xAC22 \\
кладбище: \textbf{mwet} [\emph{mwet}] 0xAB21 \\
кланяясь: \textbf{byun} [\emph{bjun}] 0x6A11 \\
кларнет: \textbf{kw6n} [\emph{kwən}] 0x6D22 \\
классический: \textbf{kwis} [\emph{kwis}] 0x8822 \\
кластер: \textbf{ksus} [\emph{ksus}] 0x8A42 \\
клев: \textbf{nrif} [\emph{nrif}] 0xC8C6 \\
клевание: \textbf{hpe2k} [\emph{ʰpe˩k}] 0x3B20 \\
клевер: \textbf{klof} [\emph{klof}] 0xCC62 \\
клеветнический: \textbf{tya7s} [\emph{tjä˥s}] 0x910A \\
клеенка: \textbf{kli2k} [\emph{kli˩k}] 0x3862 \\
Клемент: \textbf{kxem} [\emph{kxem}] 0x0BE2 \\
клептоман: \textbf{kxa7m} [\emph{kxä˥m}] 0x11E2 \\
клетка: \textbf{hcep} [\emph{ʰʃep}] 0x4B68 \\
клизма: \textbf{hne2m} [\emph{ʰne˩m}] 0x1B30 \\
климат: \textbf{hlim} [\emph{ʰlim}] 0x0858 \\
клиническая: \textbf{lyip} [\emph{ljip}] 0x480B \\
клинометр: \textbf{hlo2m} [\emph{ʰlo˩m}] 0x1C58 \\
клитор: \textbf{kc6t} [\emph{kʃət}] 0xADA2 \\
клонический: \textbf{lyo7n} [\emph{ljo˥n}] 0x740B \\
клуб: \textbf{klup} [\emph{klup}] 0x4A62 \\
клубника: \textbf{trep} [\emph{trep}] 0x4BCA \\
клык: \textbf{fla7n} [\emph{flä˥n}] 0x716C \\
клюква: \textbf{ryep} [\emph{rjep}] 0x4B0E \\
ключ: \textbf{kluf} [\emph{kluf}] 0xCA62 \\
ключица: \textbf{hk67k} [\emph{ʰkə˥k}] 0x3510 \\
кляузничество: \textbf{cli2n} [\emph{ʃli˩n}] 0x786D \\
книга: \textbf{hpuk} [\emph{ʰpuk}] 0x2A20 \\
книголюб: \textbf{blef} [\emph{blef}] 0xCB71 \\
кнопки: \textbf{hbo7n} [\emph{ʰbo˥n}] 0x7488 \\
кнут: \textbf{trip} [\emph{trip}] 0x48CA \\
князья: \textbf{hpa7s} [\emph{ʰpä˥s}] 0x9120 \\
коаксиальный: \textbf{hka2s} [\emph{ʰkä˩s}] 0x9910 \\
кобальт: \textbf{kwok} [\emph{kwok}] 0x2C22 \\
кобура: \textbf{hxu2t} [\emph{ʰxu˩t}] 0xBAD0 \\
кобыла: \textbf{hma7m} [\emph{ʰmä˥m}] 0x1108 \\
ковер: \textbf{klap} [\emph{kläp}] 0x4962 \\
ковочный: \textbf{froc} [\emph{froʃ}] 0xECCC \\
когда: \textbf{hkac} [\emph{ʰkäʃ}] 0xE910 \\
код: \textbf{hkom} [\emph{ʰkom}] 0x0C10 \\
кодирование: \textbf{kfon} [\emph{kfon}] 0x6C82 \\
кодифицированный: \textbf{kfi7n} [\emph{kfi˥n}] 0x7082 \\
кодициль: \textbf{hku7s} [\emph{ʰku˥s}] 0x9210 \\
коечные: \textbf{bre2t} [\emph{bre˩t}] 0xBBD1 \\
кожа: \textbf{hcif} [\emph{ʰʃif}] 0xC868 \\
козел: \textbf{byak} [\emph{bjäk}] 0x2911 \\
Козерог: \textbf{jroc} [\emph{ʒroʃ}] 0xECD5 \\
козыри: \textbf{txa2m} [\emph{txä˩m}] 0x19EA \\
Кока-Кола: \textbf{kwo7k} [\emph{kwo˥k}] 0x3422 \\
кокаин: \textbf{ky6n} [\emph{kjən}] 0x6D02 \\
кокарда: \textbf{kya2t} [\emph{kjä˩t}] 0xB902 \\
кокетка: \textbf{hke7t} [\emph{ʰke˥t}] 0xB310 \\
коклюш: \textbf{pfi7k} [\emph{pfi˥k}] 0x3084 \\
кокос: \textbf{kr6s} [\emph{krəs}] 0x8DC2 \\
коктейль: \textbf{kxec} [\emph{kxeʃ}] 0xEBE2 \\
кола: \textbf{jlok} [\emph{ʒlok}] 0x2C75 \\
колба: \textbf{blop} [\emph{blop}] 0x4C71 \\
колбаса: \textbf{slos} [\emph{slos}] 0x8C69 \\
колготки: \textbf{tf6t} [\emph{tfət}] 0xAD8A \\
колдовство: \textbf{hwu2t} [\emph{ʰwu˩t}] 0xBA28 \\
колено: \textbf{hgik} [\emph{ʰgik}] 0x2890 \\
коленопреклонение: \textbf{jlep} [\emph{ʒlep}] 0x4B75 \\
колесница: \textbf{hru2t} [\emph{ʰru˩t}] 0xBA70 \\
колесный: \textbf{kra7s} [\emph{krä˥s}] 0x91C2 \\
колибри: \textbf{xyop} [\emph{xjop}] 0x4C1A \\
колит: \textbf{cla2s} [\emph{ʃlä˩s}] 0x996D \\
коллаген: \textbf{kl6c} [\emph{kləʃ}] 0xED62 \\
колледж: \textbf{klic} [\emph{kliʃ}] 0xE862 \\
коллектор: \textbf{kyet} [\emph{kjet}] 0xAB02 \\
коллодий: \textbf{klo2t} [\emph{klo˩t}] 0xBC62 \\
коллоквиализм: \textbf{lwum} [\emph{lwum}] 0x0A2B \\
колоколообразный: \textbf{tco7n} [\emph{tʃo˥n}] 0x74AA \\
колония: \textbf{klen} [\emph{klen}] 0x6B62 \\
колоннада: \textbf{kle2n} [\emph{kle˩n}] 0x7B62 \\
Колумбия: \textbf{hk6p} [\emph{ʰkəp}] 0x4D10 \\
Колумбус: \textbf{kl6p} [\emph{kləp}] 0x4D62 \\
Колыбельная: \textbf{lyi7n} [\emph{lji˥n}] 0x700B \\
кольраби: \textbf{kli7p} [\emph{kli˥p}] 0x5062 \\
кольцо: \textbf{clic} [\emph{ʃliʃ}] 0xE86D \\
колючий: \textbf{gxis} [\emph{gxis}] 0x88F2 \\
команда: \textbf{myin} [\emph{mjin}] 0x6801 \\
комар: \textbf{mrok} [\emph{mrok}] 0x2CC1 \\
комбинаторика: \textbf{hbo7t} [\emph{ʰbo˥t}] 0xB488 \\
комета: \textbf{hkem} [\emph{ʰkem}] 0x0B10 \\
комикс: \textbf{mwap} [\emph{mwäp}] 0x4921 \\
комиссии: \textbf{kci7n} [\emph{kʃi˥n}] 0x70A2 \\
Комментарии: \textbf{kyes} [\emph{kjes}] 0x8B02 \\
коммерциализация: \textbf{qyaf} [\emph{ŋjäf}] 0xC919 \\
коммодор: \textbf{hk67t} [\emph{ʰkə˥t}] 0xB510 \\
коммунист: \textbf{kyos} [\emph{kjos}] 0x8C02 \\
коммутатор: \textbf{tfo7t} [\emph{tfo˥t}] 0xB48A \\
комод: \textbf{kco7t} [\emph{kʃo˥t}] 0xB4A2 \\
Коморские: \textbf{jros} [\emph{ʒros}] 0x8CD5 \\
Компания: \textbf{kyin} [\emph{kjin}] 0x6802 \\
компас: \textbf{dxap} [\emph{dxäp}] 0x49F3 \\
компенсатор: \textbf{ps67t} [\emph{psə˥t}] 0xB544 \\
компилировать: \textbf{hk6n} [\emph{ʰkən}] 0x6D10 \\
компостирование: \textbf{kfos} [\emph{kfos}] 0x8C82 \\
компьютер: \textbf{kyut} [\emph{kjut}] 0xAA02 \\
комфорт: \textbf{kwat} [\emph{kwät}] 0xA922 \\
конвекция: \textbf{kfek} [\emph{kfek}] 0x2B82 \\
конвенции: \textbf{hke7k} [\emph{ʰke˥k}] 0x3310 \\
конверт: \textbf{cwaf} [\emph{ʃwäf}] 0xC92D \\
конго: \textbf{kwuk} [\emph{kwuk}] 0x2A22 \\
конгрессменов: \textbf{kwes} [\emph{kwes}] 0x8B22 \\
конденсатор: \textbf{jyot} [\emph{ʒjot}] 0xAC15 \\
кондор: \textbf{kyo7t} [\emph{kjo˥t}] 0xB402 \\
коническая: \textbf{ts6t} [\emph{tsət}] 0xAD4A \\
конкатенация: \textbf{kle7n} [\emph{kle˥n}] 0x7362 \\
конкурс: \textbf{tsos} [\emph{tsos}] 0x8C4A \\
коннотация: \textbf{nyif} [\emph{njif}] 0xC806 \\
консалтинг: \textbf{ksi7n} [\emph{ksi˥n}] 0x7042 \\
консерватория: \textbf{nra2t} [\emph{nrä˩t}] 0xB9C6 \\
консервирование: \textbf{kwa7n} [\emph{kwä˥n}] 0x7122 \\
конспект: \textbf{kse2k} [\emph{kse˩k}] 0x3B42 \\
Константинопольский: \textbf{hsi7p} [\emph{ʰsi˥p}] 0x5048 \\
конституция: \textbf{nwin} [\emph{nwin}] 0x6826 \\
консультант: \textbf{gwac} [\emph{gwäʃ}] 0xE932 \\
контакт: \textbf{hl6n} [\emph{ʰlən}] 0x6D58 \\
контейнеры: \textbf{hven} [\emph{ʰven}] 0x6BB0 \\
континент: \textbf{nrak} [\emph{nräk}] 0x29C6 \\
контр-атака: \textbf{gxat} [\emph{gxät}] 0xA9F2 \\
контрабандист: \textbf{tso7k} [\emph{tso˥k}] 0x344A \\
контральто: \textbf{nro2t} [\emph{nro˩t}] 0xBCC6 \\
контрапункт: \textbf{twup} [\emph{twup}] 0x4A2A \\
контрастные: \textbf{nros} [\emph{nros}] 0x8CC6 \\
Контратенор: \textbf{kru7n} [\emph{kru˥n}] 0x72C2 \\
контрацептив: \textbf{gr6t} [\emph{grət}] 0xADD2 \\
контролируемый: \textbf{drot} [\emph{drot}] 0xACD3 \\
конфабуляция: \textbf{fyuk} [\emph{fjuk}] 0x2A0C \\
конфигурировать: \textbf{hxik} [\emph{ʰxik}] 0x28D0 \\
конфликт: \textbf{htok} [\emph{ʰtok}] 0x2C50 \\
конформист: \textbf{kfi7s} [\emph{kfi˥s}] 0x9082 \\
концентрический: \textbf{nro7n} [\emph{nro˥n}] 0x74C6 \\
концессионер: \textbf{kye2n} [\emph{kje˩n}] 0x7B02 \\
конькобежец: \textbf{kwi7t} [\emph{kwi˥t}] 0xB022 \\
конюшни: \textbf{hmu7c} [\emph{ʰmu˥ʃ}] 0xF208 \\
координировать: \textbf{krop} [\emph{krop}] 0x4CC2 \\
Копенгаген: \textbf{pxo7n} [\emph{pxo˥n}] 0x74E4 \\
Коперник: \textbf{jrok} [\emph{ʒrok}] 0x2CD5 \\
копия: \textbf{hkup} [\emph{ʰkup}] 0x4A10 \\
копра: \textbf{hko2p} [\emph{ʰko˩p}] 0x5C10 \\
копчик: \textbf{kyi7k} [\emph{kji˥k}] 0x3002 \\
копытный: \textbf{glu7t} [\emph{glu˥t}] 0xB272 \\
копье: \textbf{mrak} [\emph{mräk}] 0x29C1 \\
кора: \textbf{pres} [\emph{pres}] 0x8BC4 \\
кораблекрушение: \textbf{cra2n} [\emph{ʃrä˩n}] 0x79CD \\
корабль: \textbf{kcap} [\emph{kʃäp}] 0x49A2 \\
коралловый: \textbf{xrac} [\emph{xräʃ}] 0xE9DA \\
Коран: \textbf{qran} [\emph{ŋrän}] 0x69D9 \\
Корея: \textbf{vrok} [\emph{vrok}] 0x2CD6 \\
корзина: \textbf{bran} [\emph{brän}] 0x69D1 \\
Коринф: \textbf{kri2n} [\emph{kri˩n}] 0x78C2 \\
корица: \textbf{dxik} [\emph{dxik}] 0x28F3 \\
коричневый: \textbf{pyun} [\emph{pjun}] 0x6A04 \\
кормилец: \textbf{qwek} [\emph{ŋwek}] 0x2B39 \\
кормушка: \textbf{mru7n} [\emph{mru˥n}] 0x72C1 \\
корневище: \textbf{kfim} [\emph{kfim}] 0x0882 \\
корнеплоды: \textbf{kr6n} [\emph{krən}] 0x6DC2 \\
Корнуолл: \textbf{kwo7n} [\emph{kwo˥n}] 0x7422 \\
коробейник: \textbf{hta7m} [\emph{ʰtä˥m}] 0x1150 \\
коробка: \textbf{pwak} [\emph{pwäk}] 0x2924 \\
корова: \textbf{hgap} [\emph{ʰgäp}] 0x4990 \\
Королева: \textbf{rwin} [\emph{rwin}] 0x682E \\
король: \textbf{hyap} [\emph{ʰjäp}] 0x4918 \\
корона: \textbf{kran} [\emph{krän}] 0x69C2 \\
коронер: \textbf{kro2n} [\emph{kro˩n}] 0x7CC2 \\
коронование: \textbf{hgu2n} [\emph{ʰgu˩n}] 0x7A90 \\
коротковолновый: \textbf{rwop} [\emph{rwop}] 0x4C2E \\
коррелированный: \textbf{kya2n} [\emph{kjä˩n}] 0x7902 \\
коррупция: \textbf{hrup} [\emph{ʰrup}] 0x4A70 \\
корсаж: \textbf{hko7s} [\emph{ʰko˥s}] 0x9410 \\
Корсика: \textbf{krof} [\emph{krof}] 0xCCC2 \\
кортеж: \textbf{kro2k} [\emph{kro˩k}] 0x3CC2 \\
корунд: \textbf{kru2n} [\emph{kru˩n}] 0x7AC2 \\
корь: \textbf{ksa7s} [\emph{ksä˥s}] 0x9142 \\
коса: \textbf{hke2n} [\emph{ʰke˩n}] 0x7B10 \\
косинус: \textbf{ksup} [\emph{ksup}] 0x4A42 \\
косички: \textbf{hbe7s} [\emph{ʰbe˥s}] 0x9388 \\
косметический: \textbf{ksop} [\emph{ksop}] 0x4C42 \\
космодром: \textbf{pse7p} [\emph{pse˥p}] 0x5344 \\
космология: \textbf{zloc} [\emph{zloʃ}] 0xEC74 \\
космос: \textbf{ksom} [\emph{ksom}] 0x0C42 \\
коснуться: \textbf{tcum} [\emph{tʃum}] 0x0AAA \\
Косово: \textbf{kso7f} [\emph{kso˥f}] 0xD442 \\
косоглазие: \textbf{gwim} [\emph{gwim}] 0x0832 \\
костыль: \textbf{kruk} [\emph{kruk}] 0x2AC2 \\
кость: \textbf{ksot} [\emph{ksot}] 0xAC42 \\
Кот: \textbf{hgak} [\emph{ʰgäk}] 0x2990 \\
котиковый: \textbf{hli7f} [\emph{ʰli˥f}] 0xD058 \\
кофеин: \textbf{kfe7n} [\emph{kfe˥n}] 0x7382 \\
кочевник: \textbf{hzom} [\emph{ʰzom}] 0x0CA0 \\
кошениль: \textbf{tco2n} [\emph{tʃo˩n}] 0x7CAA \\
кошерный: \textbf{hko2c} [\emph{ʰko˩ʃ}] 0xFC10 \\
кошмар: \textbf{hdom} [\emph{ʰdom}] 0x0C98 \\
кощунство: \textbf{qret} [\emph{ŋret}] 0xABD9 \\
коэволюция: \textbf{vloc} [\emph{vloʃ}] 0xEC76 \\
коэффициент: \textbf{hgus} [\emph{ʰgus}] 0x8A90 \\
краб: \textbf{kxek} [\emph{kxek}] 0x2BE2 \\
кран: \textbf{krec} [\emph{kreʃ}] 0xEBC2 \\
крапать: \textbf{bva7t} [\emph{bvä˥t}] 0xB191 \\
крапива: \textbf{xyim} [\emph{xjim}] 0x081A \\
красивая: \textbf{hyus} [\emph{ʰjus}] 0x8A18 \\
красильщик: \textbf{hda7c} [\emph{ʰdä˥ʃ}] 0xF198 \\
красноломкий: \textbf{hs67k} [\emph{ʰsə˥k}] 0x3548 \\
красный: \textbf{hyuk} [\emph{ʰjuk}] 0x2A18 \\
кратеры: \textbf{hke7s} [\emph{ʰke˥s}] 0x9310 \\
краткое: \textbf{hrip} [\emph{ʰrip}] 0x4870 \\
края: \textbf{hda2n} [\emph{ʰdä˩n}] 0x7998 \\
креативность: \textbf{ryaf} [\emph{rjäf}] 0xC90E \\
креационизм: \textbf{twa2m} [\emph{twä˩m}] 0x192A \\
креветка: \textbf{cra7m} [\emph{ʃrä˥m}] 0x11CD \\
кредит: \textbf{hrit} [\emph{ʰrit}] 0xA870 \\
кредитор: \textbf{xwa7t} [\emph{xwä˥t}] 0xB13A \\
крейсерский: \textbf{kl6s} [\emph{kləs}] 0x8D62 \\
крем: \textbf{krim} [\emph{krim}] 0x08C2 \\
крематорий: \textbf{kca7m} [\emph{kʃä˥m}] 0x11A2 \\
кремировать: \textbf{kxa2m} [\emph{kxä˩m}] 0x19E2 \\
кремний: \textbf{zlik} [\emph{zlik}] 0x2874 \\
кренделек: \textbf{hre7p} [\emph{ʰre˥p}] 0x5370 \\
креольский: \textbf{hko2m} [\emph{ʰko˩m}] 0x1C10 \\
кресс: \textbf{gres} [\emph{gres}] 0x8BD2 \\
крестец: \textbf{sri7k} [\emph{sri˥k}] 0x30C9 \\
крестник: \textbf{hgu7t} [\emph{ʰgu˥t}] 0xB290 \\
крестоцветный: \textbf{kfuf} [\emph{kfuf}] 0xCA82 \\
кретинизм: \textbf{kyi7s} [\emph{kji˥s}] 0x9002 \\
крещение: \textbf{hbi7s} [\emph{ʰbi˥s}] 0x9088 \\
кривая: \textbf{kcek} [\emph{kʃek}] 0x2BA2 \\
кризис: \textbf{ksic} [\emph{ksiʃ}] 0xE842 \\
крикет: \textbf{jrik} [\emph{ʒrik}] 0x28D5 \\
крики: \textbf{hxa2t} [\emph{ʰxä˩t}] 0xB9D0 \\
криминолог: \textbf{hmi2c} [\emph{ʰmi˩ʃ}] 0xF808 \\
крипта: \textbf{qrip} [\emph{ŋrip}] 0x48D9 \\
криптограмма: \textbf{kri2m} [\emph{kri˩m}] 0x18C2 \\
криптон: \textbf{kxot} [\emph{kxot}] 0xACE2 \\
кристаллография: \textbf{tl6f} [\emph{tləf}] 0xCD6A \\
кристаллы: \textbf{cw6t} [\emph{ʃwət}] 0xAD2D \\
Крит: \textbf{kre7t} [\emph{kre˥t}] 0xB3C2 \\
критерий: \textbf{rwun} [\emph{rwun}] 0x6A2E \\
критика: \textbf{krik} [\emph{krik}] 0x28C2 \\
кров: \textbf{cric} [\emph{ʃriʃ}] 0xE8CD \\
кровавый: \textbf{brit} [\emph{brit}] 0xA8D1 \\
кровельщик: \textbf{lye2t} [\emph{lje˩t}] 0xBB0B \\
кровожадный: \textbf{gvic} [\emph{gviʃ}] 0xE892 \\
кровосмешение: \textbf{nwi2s} [\emph{nwi˩s}] 0x9826 \\
кровь: \textbf{cyac} [\emph{ʃjäʃ}] 0xE90D \\
крокет: \textbf{hre2k} [\emph{ʰre˩k}] 0x3B70 \\
крокодил: \textbf{kyop} [\emph{kjop}] 0x4C02 \\
кролик: \textbf{cruk} [\emph{ʃruk}] 0x2ACD \\
Кроме: \textbf{hsap} [\emph{ʰsäp}] 0x4948 \\
кропотливый: \textbf{hxem} [\emph{ʰxem}] 0x0BD0 \\
крошечный: \textbf{cyon} [\emph{ʃjon}] 0x6C0D \\
круг: \textbf{trin} [\emph{trin}] 0x68CA \\
круглый: \textbf{hyuc} [\emph{ʰjuʃ}] 0xEA18 \\
кружевной: \textbf{hwe7s} [\emph{ʰwe˥s}] 0x9328 \\
кружка: \textbf{hkop} [\emph{ʰkop}] 0x4C10 \\
круто: \textbf{sri2k} [\emph{sri˩k}] 0x38C9 \\
крутой: \textbf{tcoc} [\emph{tʃoʃ}] 0xECAA \\
крыло: \textbf{gyin} [\emph{gjin}] 0x6812 \\
крыша: \textbf{hrat} [\emph{ʰrät}] 0xA970 \\
крышка: \textbf{kfin} [\emph{kfin}] 0x6882 \\
крюк: \textbf{hduk} [\emph{ʰduk}] 0x2A98 \\
ксенон: \textbf{zyen} [\emph{zjen}] 0x6B14 \\
ксенофобские: \textbf{zwip} [\emph{zwip}] 0x4834 \\
ксилола: \textbf{zli2n} [\emph{zli˩n}] 0x7874 \\
ксилофон: \textbf{zlif} [\emph{zlif}] 0xC874 \\
кто: \textbf{hkik} [\emph{ʰkik}] 0x2810 \\
куб: \textbf{kfup} [\emph{kfup}] 0x4A82 \\
Куба: \textbf{kyup} [\emph{kjup}] 0x4A02 \\
кубизм: \textbf{ksi7m} [\emph{ksi˥m}] 0x1042 \\
Кувейт: \textbf{vwut} [\emph{vwut}] 0xAA36 \\
кувшин: \textbf{xric} [\emph{xriʃ}] 0xE8DA \\
кувыркаться: \textbf{zrot} [\emph{zrot}] 0xACD4 \\
кудахтать: \textbf{klo7k} [\emph{klo˥k}] 0x3462 \\
кузнец: \textbf{xra7n} [\emph{xrä˥n}] 0x71DA \\
кузнечик: \textbf{tsep} [\emph{tsep}] 0x4B4A \\
кузовостроение: \textbf{dwos} [\emph{dwos}] 0x8C33 \\
кукла: \textbf{bwak} [\emph{bwäk}] 0x2931 \\
кукуруза: \textbf{mran} [\emph{mrän}] 0x69C1 \\
Кульминацией: \textbf{hmu2t} [\emph{ʰmu˩t}] 0xBA08 \\
купить: \textbf{kyim} [\emph{kjim}] 0x0802 \\
купол: \textbf{hgup} [\emph{ʰgup}] 0x4A90 \\
купон: \textbf{kcup} [\emph{kʃup}] 0x4AA2 \\
купорос: \textbf{vra2n} [\emph{vrä˩n}] 0x79D6 \\
куратор: \textbf{kxus} [\emph{kxus}] 0x8AE2 \\
курд: \textbf{kr67t} [\emph{krə˥t}] 0xB5C2 \\
Курдистан: \textbf{ksu2n} [\emph{ksu˩n}] 0x7A42 \\
курильщик: \textbf{myu7t} [\emph{mju˥t}] 0xB201 \\
курица: \textbf{kric} [\emph{kriʃ}] 0xE8C2 \\
курносый: \textbf{hnof} [\emph{ʰnof}] 0xCC30 \\
курсор: \textbf{kro7s} [\emph{kro˥s}] 0x94C2 \\
курсы: \textbf{kri2s} [\emph{kri˩s}] 0x98C2 \\
курятник: \textbf{kcop} [\emph{kʃop}] 0x4CA2 \\
кускус: \textbf{ksuf} [\emph{ksuf}] 0xCA42 \\
куст: \textbf{tsuc} [\emph{tsuʃ}] 0xEA4A \\
кухня: \textbf{hcon} [\emph{ʰʃon}] 0x6C68 \\
куча: \textbf{mwak} [\emph{mwäk}] 0x2921 \\
кушетка: \textbf{swup} [\emph{swup}] 0x4A29 \\
кхмерская: \textbf{hke7m} [\emph{ʰke˥m}] 0x1310 \\
\subsection{ л }
лабиринт: \textbf{bren} [\emph{bren}] 0x6BD1 \\
лаборатория: \textbf{hbos} [\emph{ʰbos}] 0x8C88 \\
лавина: \textbf{hv6n} [\emph{ʰvən}] 0x6DB0 \\
лавировали: \textbf{tsa7k} [\emph{tsä˥k}] 0x314A \\
лаг: \textbf{kxa2k} [\emph{kxä˩k}] 0x39E2 \\
лагерь: \textbf{kyam} [\emph{kjäm}] 0x0902 \\
ладан: \textbf{lya7n} [\emph{ljä˥n}] 0x710B \\
лазер: \textbf{lwes} [\emph{lwes}] 0x8B2B \\
лазурный: \textbf{slu2n} [\emph{slu˩n}] 0x7A69 \\
лак: \textbf{vrac} [\emph{vräʃ}] 0xE9D6 \\
лаконично: \textbf{kci7k} [\emph{kʃi˥k}] 0x30A2 \\
лакрица: \textbf{kra7m} [\emph{krä˥m}] 0x11C2 \\
лактоза: \textbf{kxo7s} [\emph{kxo˥s}] 0x94E2 \\
лактон: \textbf{hk62t} [\emph{ʰkə˩t}] 0xBD10 \\
ламантин: \textbf{mra7p} [\emph{mrä˥p}] 0x51C1 \\
ламинария: \textbf{klep} [\emph{klep}] 0x4B62 \\
ламинарный: \textbf{myo2n} [\emph{mjo˩n}] 0x7C01 \\
лантан: \textbf{lw6n} [\emph{lwən}] 0x6D2B \\
ланцетовидный: \textbf{nye2t} [\emph{nje˩t}] 0xBB06 \\
лань: \textbf{hmu7n} [\emph{ʰmu˥n}] 0x7208 \\
лаосский: \textbf{lwoc} [\emph{lwoʃ}] 0xEC2B \\
лапками: \textbf{twe2t} [\emph{twe˩t}] 0xBB2A \\
\subsection{ Л }
Лапландия: \textbf{pla7n} [\emph{plä˥n}] 0x7164 \\
ларингоскоп: \textbf{hro2p} [\emph{ʰro˩p}] 0x5C70 \\
ласка: \textbf{hka7m} [\emph{ʰkä˥m}] 0x1110 \\
лассо: \textbf{lwa2t} [\emph{lwä˩t}] 0xB92B \\
Латвия: \textbf{lwof} [\emph{lwof}] 0xCC2B \\
латинский: \textbf{lyek} [\emph{ljek}] 0x2B0B \\
латунь: \textbf{pcon} [\emph{pʃon}] 0x6CA4 \\
латыш: \textbf{bxep} [\emph{bxep}] 0x4BF1 \\
лауриновая: \textbf{lwi2k} [\emph{lwi˩k}] 0x382B \\
лачуга: \textbf{hguf} [\emph{ʰguf}] 0xCA90 \\
лебедь: \textbf{xwon} [\emph{xwon}] 0x6C3A \\
лев: \textbf{hsif} [\emph{ʰsif}] 0xC848 \\
легальность: \textbf{xlac} [\emph{xläʃ}] 0xE97A \\
легенда: \textbf{lwen} [\emph{lwen}] 0x6B2B \\
легированная: \textbf{xlo2t} [\emph{xlo˩t}] 0xBC7A \\
легкий: \textbf{kwan} [\emph{kwän}] 0x6922 \\
легкое: \textbf{pcep} [\emph{pʃep}] 0x4BA4 \\
лед: \textbf{hbif} [\emph{ʰbif}] 0xC888 \\
леденец: \textbf{hqep} [\emph{ʰŋep}] 0x4BC8 \\
ледник: \textbf{glic} [\emph{gliʃ}] 0xE872 \\
лежачий: \textbf{jwen} [\emph{ʒwen}] 0x6B35 \\
лейкемия: \textbf{hle2m} [\emph{ʰle˩m}] 0x1B58 \\
лейтенант: \textbf{hle2n} [\emph{ʰle˩n}] 0x7B58 \\
лейтмотив: \textbf{hte7f} [\emph{ʰte˥f}] 0xD350 \\
лексикограф: \textbf{hse7f} [\emph{ʰse˥f}] 0xD348 \\
лексикография: \textbf{sre7k} [\emph{sre˥k}] 0x33C9 \\
ленивый: \textbf{lwas} [\emph{lwäs}] 0x892B \\
Ленинград: \textbf{lyi2k} [\emph{lji˩k}] 0x380B \\
лепестки: \textbf{pl6p} [\emph{pləp}] 0x4D64 \\
лепились: \textbf{hla2t} [\emph{ʰlä˩t}] 0xB958 \\
лесоводство: \textbf{lyis} [\emph{ljis}] 0x880B \\
лесок: \textbf{slup} [\emph{slup}] 0x4A69 \\
Лесото: \textbf{lw6t} [\emph{lwət}] 0xAD2B \\
лесс: \textbf{lwus} [\emph{lwus}] 0x8A2B \\
лестница: \textbf{hlet} [\emph{ʰlet}] 0xAB58 \\
лестный: \textbf{cram} [\emph{ʃräm}] 0x09CD \\
летать: \textbf{fran} [\emph{frän}] 0x69CC \\
лето: \textbf{cyam} [\emph{ʃjäm}] 0x090D \\
лецитин: \textbf{lye7n} [\emph{lje˥n}] 0x730B \\
лечение: \textbf{tlim} [\emph{tlim}] 0x086A \\
лещ: \textbf{brem} [\emph{brem}] 0x0BD1 \\
лженаука: \textbf{dxo7n} [\emph{dxo˥n}] 0x74F3 \\
лживый: \textbf{mwep} [\emph{mwep}] 0x4B21 \\
либеральная: \textbf{trap} [\emph{träp}] 0x49CA \\
Либерия: \textbf{lyep} [\emph{ljep}] 0x4B0B \\
либреттист: \textbf{bre7s} [\emph{bre˥s}] 0x93D1 \\
либретто: \textbf{bri2t} [\emph{bri˩t}] 0xB8D1 \\
Ливан: \textbf{bven} [\emph{bven}] 0x6B91 \\
Ливия: \textbf{vyip} [\emph{vjip}] 0x4816 \\
лиги: \textbf{hli2k} [\emph{ʰli˩k}] 0x3858 \\
ликование: \textbf{hga7c} [\emph{ʰgä˥ʃ}] 0xF190 \\
лимон: \textbf{hlom} [\emph{ʰlom}] 0x0C58 \\
лимфа: \textbf{lwip} [\emph{lwip}] 0x482B \\
лингвист: \textbf{lwa7t} [\emph{lwä˥t}] 0xB12B \\
линейный: \textbf{lyec} [\emph{ljeʃ}] 0xEB0B \\
линия: \textbf{hlas} [\emph{ʰläs}] 0x8958 \\
линкор: \textbf{vlip} [\emph{vlip}] 0x4876 \\
линолевая: \textbf{nyo2k} [\emph{njo˩k}] 0x3C06 \\
линька: \textbf{xruk} [\emph{xruk}] 0x2ADA \\
липид: \textbf{djip} [\emph{dʒip}] 0x48B3 \\
липкий: \textbf{lyik} [\emph{ljik}] 0x280B \\
лира: \textbf{lyi2t} [\emph{lji˩t}] 0xB80B \\
лиса: \textbf{flik} [\emph{flik}] 0x286C \\
Лиссабон: \textbf{hlo7p} [\emph{ʰlo˥p}] 0x5458 \\
лист: \textbf{dlit} [\emph{dlit}] 0xA873 \\
листать: \textbf{hfu7n} [\emph{ʰfu˥n}] 0x7260 \\
лиственница: \textbf{lyoc} [\emph{ljoʃ}] 0xEC0B \\
листовка: \textbf{pfet} [\emph{pfet}] 0xAB84 \\
литания: \textbf{lyi2n} [\emph{lji˩n}] 0x780B \\
Литва: \textbf{lyif} [\emph{ljif}] 0xC80B \\
литература: \textbf{lwit} [\emph{lwit}] 0xA82B \\
литий: \textbf{tfim} [\emph{tfim}] 0x088A \\
литография: \textbf{frof} [\emph{frof}] 0xCCCC \\
литология: \textbf{fli2k} [\emph{fli˩k}] 0x386C \\
литосфера: \textbf{hl6f} [\emph{ʰləf}] 0xCD58 \\
литр: \textbf{hl6t} [\emph{ʰlət}] 0xAD58 \\
лифт: \textbf{kfit} [\emph{kfit}] 0xA882 \\
лихорадочный: \textbf{hfa2c} [\emph{ʰfä˩ʃ}] 0xF960 \\
Лихтенштейн: \textbf{tc6s} [\emph{tʃəs}] 0x8DAA \\
лицемер: \textbf{pci7t} [\emph{pʃi˥t}] 0xB0A4 \\
лицензия: \textbf{hlen} [\emph{ʰlen}] 0x6B58 \\
личность: \textbf{kwit} [\emph{kwit}] 0xA822 \\
лишай: \textbf{hl6k} [\emph{ʰlək}] 0x2D58 \\
лоб: \textbf{htot} [\emph{ʰtot}] 0xAC50 \\
лоббист: \textbf{hlo2s} [\emph{ʰlo˩s}] 0x9C58 \\
лобковый: \textbf{hpu2k} [\emph{ʰpu˩k}] 0x3A20 \\
лоботомия: \textbf{bwom} [\emph{bwom}] 0x0C31 \\
ловушка: \textbf{crap} [\emph{ʃräp}] 0x49CD \\
логарифм: \textbf{grim} [\emph{grim}] 0x08D2 \\
логический: \textbf{blu7n} [\emph{blu˥n}] 0x7271 \\
логотипы: \textbf{lwop} [\emph{lwop}] 0x4C2B \\
лодыжка: \textbf{gwak} [\emph{gwäk}] 0x2932 \\
лож: \textbf{hla7c} [\emph{ʰlä˥ʃ}] 0xF158 \\
ложка: \textbf{syuk} [\emph{sjuk}] 0x2A09 \\
ложь: \textbf{cyin} [\emph{ʃjin}] 0x680D \\
лозоискательство: \textbf{dxa7s} [\emph{dxä˥s}] 0x91F3 \\
лозунг: \textbf{lwok} [\emph{lwok}] 0x2C2B \\
локоть: \textbf{kcoc} [\emph{kʃoʃ}] 0xECA2 \\
лом: \textbf{brap} [\emph{bräp}] 0x49D1 \\
ломбард: \textbf{pca2p} [\emph{pʃä˩p}] 0x59A4 \\
Ломбардия: \textbf{hmo2t} [\emph{ʰmo˩t}] 0xBC08 \\
ломкий: \textbf{dxuk} [\emph{dxuk}] 0x2AF3 \\
ломонос: \textbf{kfes} [\emph{kfes}] 0x8B82 \\
ломоть: \textbf{hx67n} [\emph{ʰxə˥n}] 0x75D0 \\
Лондон: \textbf{hlo2t} [\emph{ʰlo˩t}] 0xBC58 \\
лосины: \textbf{hle7s} [\emph{ʰle˥s}] 0x9358 \\
лосось: \textbf{slam} [\emph{släm}] 0x0969 \\
лотерея: \textbf{hli7c} [\emph{ʰli˥ʃ}] 0xF058 \\
лотос: \textbf{hl6s} [\emph{ʰləs}] 0x8D58 \\
лошадь: \textbf{kyas} [\emph{kjäs}] 0x8902 \\
луг: \textbf{crot} [\emph{ʃrot}] 0xACCD \\
лужи: \textbf{hs6s} [\emph{ʰsəs}] 0x8D48 \\
лук: \textbf{cyif} [\emph{ʃjif}] 0xC80D \\
лунатик: \textbf{myi7t} [\emph{mji˥t}] 0xB001 \\
лупа: \textbf{hlu2p} [\emph{ʰlu˩p}] 0x5A58 \\
луч: \textbf{hrem} [\emph{ʰrem}] 0x0B70 \\
лучистый: \textbf{hre7t} [\emph{ʰre˥t}] 0xB370 \\
лучник: \textbf{hru7n} [\emph{ʰru˥n}] 0x7270 \\
лучше: \textbf{pcac} [\emph{pʃäʃ}] 0xE9A4 \\
лущеный: \textbf{xlo2k} [\emph{xlo˩k}] 0x3C7A \\
лысый: \textbf{bxat} [\emph{bxät}] 0xA9F1 \\
любитель: \textbf{myam} [\emph{mjäm}] 0x0901 \\
люблю: \textbf{pyan} [\emph{pjän}] 0x6904 \\
любовно: \textbf{hfe7t} [\emph{ʰfe˥t}] 0xB360 \\
любовный: \textbf{mwa7t} [\emph{mwä˥t}] 0xB121 \\
любопытный: \textbf{kcic} [\emph{kʃiʃ}] 0xE8A2 \\
люк: \textbf{xruf} [\emph{xruf}] 0xCADA \\
Люк: \textbf{hlu2k} [\emph{ʰlu˩k}] 0x3A58 \\
Люксембург: \textbf{hl6p} [\emph{ʰləp}] 0x4D58 \\
лютеранский: \textbf{hlu2n} [\emph{ʰlu˩n}] 0x7A58 \\
лютик: \textbf{brep} [\emph{brep}] 0x4BD1 \\
лютня: \textbf{lwap} [\emph{lwäp}] 0x492B \\
лягушка: \textbf{jrak} [\emph{ʒräk}] 0x29D5 \\
\subsection{ М }
Маврикий: \textbf{mrus} [\emph{mrus}] 0x8AC1 \\
Мавритания: \textbf{vrun} [\emph{vrun}] 0x6AD6 \\
маг: \textbf{hme7c} [\emph{ʰme˥ʃ}] 0xF308 \\
магазин: \textbf{hdak} [\emph{ʰdäk}] 0x2998 \\
магнат: \textbf{nra2n} [\emph{nrä˩n}] 0x79C6 \\
магнезит: \textbf{hge7s} [\emph{ʰge˥s}] 0x9390 \\
магнезия: \textbf{hga7m} [\emph{ʰgä˥m}] 0x1190 \\
магнетит: \textbf{dxos} [\emph{dxos}] 0x8CF3 \\
магний: \textbf{hgem} [\emph{ʰgem}] 0x0B90 \\
магнитный: \textbf{mwik} [\emph{mwik}] 0x2821 \\
магнитометр: \textbf{gyem} [\emph{gjem}] 0x0B12 \\
магнолия: \textbf{gla2n} [\emph{glä˩n}] 0x7972 \\
Магомет: \textbf{hmu7m} [\emph{ʰmu˥m}] 0x1208 \\
Мадагаскар: \textbf{hm6s} [\emph{ʰməs}] 0x8D08 \\
Мадрид: \textbf{mra7t} [\emph{mrä˥t}] 0xB1C1 \\
мажордом: \textbf{dyom} [\emph{djom}] 0x0C13 \\
мазохизм: \textbf{hso2m} [\emph{ʰso˩m}] 0x1C48 \\
Майкл: \textbf{hma7k} [\emph{ʰmä˥k}] 0x3108 \\
майонез: \textbf{myes} [\emph{mjes}] 0x8B01 \\
майоран: \textbf{cyom} [\emph{ʃjom}] 0x0C0D \\
майя: \textbf{mya7f} [\emph{mjä˥f}] 0xD101 \\
мак: \textbf{pcus} [\emph{pʃus}] 0x8AA4 \\
Македония: \textbf{mw6n} [\emph{mwən}] 0x6D21 \\
макраме: \textbf{jram} [\emph{ʒräm}] 0x09D5 \\
макроскопический: \textbf{kso2p} [\emph{kso˩p}] 0x5C42 \\
максимизировать: \textbf{hxit} [\emph{ʰxit}] 0xA8D0 \\
макулатура: \textbf{lye7t} [\emph{lje˥t}] 0xB30B \\
Малави: \textbf{mwif} [\emph{mwif}] 0xC821 \\
Малайзия: \textbf{myec} [\emph{mjeʃ}] 0xEB01 \\
малайский: \textbf{hma7c} [\emph{ʰmä˥ʃ}] 0xF108 \\
Малакка: \textbf{mya7k} [\emph{mjä˥k}] 0x3101 \\
малахит: \textbf{gxim} [\emph{gxim}] 0x08F2 \\
Малахия: \textbf{xla7k} [\emph{xlä˥k}] 0x317A \\
малаялам: \textbf{lya2m} [\emph{ljä˩m}] 0x190B \\
Мальдивы: \textbf{myaf} [\emph{mjäf}] 0xC901 \\
мальтоза: \textbf{mwo2t} [\emph{mwo˩t}] 0xBC21 \\
мальтузианский: \textbf{hm67s} [\emph{ʰmə˥s}] 0x9508 \\
Мама: \textbf{hmam} [\emph{ʰmäm}] 0x0908 \\
маммона: \textbf{mya2n} [\emph{mjä˩n}] 0x7901 \\
мангал: \textbf{hqa2t} [\emph{ʰŋä˩t}] 0xB9C8 \\
мангольд: \textbf{tc6c} [\emph{tʃəʃ}] 0xEDAA \\
мандарин: \textbf{mwon} [\emph{mwon}] 0x6C21 \\
мандолина: \textbf{mwi7n} [\emph{mwi˥n}] 0x7021 \\
манжета: \textbf{kfof} [\emph{kfof}] 0xCC82 \\
маниакальный: \textbf{myu2n} [\emph{mju˩n}] 0x7A01 \\
маникюр: \textbf{nyuc} [\emph{njuʃ}] 0xEA06 \\
манильской: \textbf{hmi7k} [\emph{ʰmi˥k}] 0x3008 \\
манить: \textbf{sra7k} [\emph{srä˥k}] 0x31C9 \\
манифест: \textbf{nya2t} [\emph{njä˩t}] 0xB906 \\
манихей: \textbf{hme2c} [\emph{ʰme˩ʃ}] 0xFB08 \\
манна: \textbf{hmo2n} [\emph{ʰmo˩n}] 0x7C08 \\
маори: \textbf{hmi2m} [\emph{ʰmi˩m}] 0x1808 \\
маразм: \textbf{cra2s} [\emph{ʃrä˩s}] 0x99CD \\
маратхи: \textbf{mra2f} [\emph{mrä˩f}] 0xD9C1 \\
марафон: \textbf{mraf} [\emph{mräf}] 0xC9C1 \\
марганец: \textbf{hzek} [\emph{ʰzek}] 0x2BA0 \\
маргарин: \textbf{jr6n} [\emph{ʒrən}] 0x6DD5 \\
мариновать: \textbf{qri7n} [\emph{ŋri˥n}] 0x70D9 \\
Марки: \textbf{txa2p} [\emph{txä˩p}] 0x59EA \\
маркиз: \textbf{mri2k} [\emph{mri˩k}] 0x38C1 \\
маркированный: \textbf{hye7t} [\emph{ʰje˥t}] 0xB318 \\
марксист: \textbf{mwi7s} [\emph{mwi˥s}] 0x9021 \\
Марокко: \textbf{zrok} [\emph{zrok}] 0x2CD4 \\
марочный: \textbf{hna7c} [\emph{ʰnä˥ʃ}] 0xF130 \\
марсель: \textbf{tsa2f} [\emph{tsä˩f}] 0xD94A \\
Марсель: \textbf{hma2s} [\emph{ʰmä˩s}] 0x9908 \\
марсианин: \textbf{mri7n} [\emph{mri˥n}] 0x70C1 \\
Март: \textbf{myak} [\emph{mjäk}] 0x2901 \\
Мартиника: \textbf{mra7k} [\emph{mrä˥k}] 0x31C1 \\
мартышка: \textbf{jre2t} [\emph{ʒre˩t}] 0xBBD5 \\
маршалльский: \textbf{mya7c} [\emph{mjä˥ʃ}] 0xF101 \\
маршрут: \textbf{jric} [\emph{ʒriʃ}] 0xE8D5 \\
маршрутизатор: \textbf{hxut} [\emph{ʰxut}] 0xAAD0 \\
маршруты: \textbf{hru2n} [\emph{ʰru˩n}] 0x7A70 \\
маска: \textbf{myas} [\emph{mjäs}] 0x8901 \\
масленка: \textbf{zlos} [\emph{zlos}] 0x8C74 \\
масло: \textbf{hyet} [\emph{ʰjet}] 0xAB18 \\
маслобойня: \textbf{kci2m} [\emph{kʃi˩m}] 0x18A2 \\
массаж: \textbf{jlac} [\emph{ʒläʃ}] 0xE975 \\
массажист: \textbf{mya2s} [\emph{mjä˩s}] 0x9901 \\
массив: \textbf{psuf} [\emph{psuf}] 0xCA44 \\
мастер: \textbf{mrap} [\emph{mräp}] 0x49C1 \\
мастит: \textbf{jli7t} [\emph{ʒli˥t}] 0xB075 \\
мастодонт: \textbf{swo7n} [\emph{swo˥n}] 0x7429 \\
мастоидный: \textbf{hjo2t} [\emph{ʰʒo˩t}] 0xBCA8 \\
мастурбация: \textbf{sru7n} [\emph{sru˥n}] 0x72C9 \\
математика: \textbf{gyik} [\emph{gjik}] 0x2812 \\
математики: \textbf{hmi2t} [\emph{ʰmi˩t}] 0xB808 \\
матереубийство: \textbf{mri2m} [\emph{mri˩m}] 0x18C1 \\
материализм: \textbf{mri7s} [\emph{mri˥s}] 0x90C1 \\
материнство: \textbf{mri2n} [\emph{mri˩n}] 0x78C1 \\
матриарх: \textbf{jra7k} [\emph{ʒrä˥k}] 0x31D5 \\
матрицы: \textbf{mre7s} [\emph{mre˥s}] 0x93C1 \\
мафия: \textbf{hma2f} [\emph{ʰmä˩f}] 0xD908 \\
Мафусаил: \textbf{hmu2s} [\emph{ʰmu˩s}] 0x9A08 \\
Махараштра: \textbf{mr6c} [\emph{mrəʃ}] 0xEDC1 \\
машина: \textbf{myik} [\emph{mjik}] 0x2801 \\
маяк: \textbf{gxas} [\emph{gxäs}] 0x89F2 \\
МВФ: \textbf{hxi7m} [\emph{ʰxi˥m}] 0x10D0 \\
мебель: \textbf{hcek} [\emph{ʰʃek}] 0x2B68 \\
меблировать: \textbf{pwa2s} [\emph{pwä˩s}] 0x9924 \\
мегалит: \textbf{glef} [\emph{glef}] 0xCB72 \\
мегаломания: \textbf{gle2n} [\emph{gle˩n}] 0x7B72 \\
мед: \textbf{hfon} [\emph{ʰfon}] 0x6C60 \\
медальон: \textbf{hlo7k} [\emph{ʰlo˥k}] 0x3458 \\
медведь: \textbf{cyas} [\emph{ʃjäs}] 0x890D \\
медиально: \textbf{myi2s} [\emph{mji˩s}] 0x9801 \\
медиана: \textbf{mwa7s} [\emph{mwä˥s}] 0x9121 \\
медлительный: \textbf{dla2t} [\emph{dlä˩t}] 0xB973 \\
медлить: \textbf{tre2t} [\emph{tre˩t}] 0xBBCA \\
медсестра: \textbf{nras} [\emph{nräs}] 0x89C6 \\
медуза: \textbf{hme7s} [\emph{ʰme˥s}] 0x9308 \\
медулла: \textbf{hmu7k} [\emph{ʰmu˥k}] 0x3208 \\
медь: \textbf{hkep} [\emph{ʰkep}] 0x4B10 \\
междисциплинарный: \textbf{hti2n} [\emph{ʰti˩n}] 0x7850 \\
между: \textbf{hkic} [\emph{ʰkiʃ}] 0xE810 \\
Международный: \textbf{nrot} [\emph{nrot}] 0xACC6 \\
межпланетный: \textbf{zlen} [\emph{zlen}] 0x6B74 \\
межсезонье: \textbf{fwam} [\emph{fwäm}] 0x092C \\
мезодерма: \textbf{hme2m} [\emph{ʰme˩m}] 0x1B08 \\
мезолит: \textbf{slo2k} [\emph{slo˩k}] 0x3C69 \\
мейоз: \textbf{mwus} [\emph{mwus}] 0x8A21 \\
Мексика: \textbf{hm6k} [\emph{ʰmək}] 0x2D08 \\
мексиканский: \textbf{hk6k} [\emph{ʰkək}] 0x2D10 \\
мел: \textbf{tcep} [\emph{tʃep}] 0x4BAA \\
Меланезия: \textbf{hle7c} [\emph{ʰle˥ʃ}] 0xF358 \\
меланхолия: \textbf{hloc} [\emph{ʰloʃ}] 0xEC58 \\
мелодичный: \textbf{xref} [\emph{xref}] 0xCBDA \\
мелодия: \textbf{myet} [\emph{mjet}] 0xAB01 \\
мелодрама: \textbf{mra7m} [\emph{mrä˥m}] 0x11C1 \\
мелодраматический: \textbf{mri7k} [\emph{mri˥k}] 0x30C1 \\
мельница: \textbf{hmok} [\emph{ʰmok}] 0x2C08 \\
мениск: \textbf{hn67k} [\emph{ʰnə˥k}] 0x3530 \\
менструация: \textbf{myi2t} [\emph{mji˩t}] 0xB801 \\
ментол: \textbf{hme7f} [\emph{ʰme˥f}] 0xD308 \\
менуэт: \textbf{mre7t} [\emph{mre˥t}] 0xB3C1 \\
меньшинство: \textbf{sros} [\emph{sros}] 0x8CC9 \\
меня: \textbf{hman} [\emph{ʰmän}] 0x6908 \\
меридиан: \textbf{mr6n} [\emph{mrən}] 0x6DC1 \\
Меркурий: \textbf{krek} [\emph{krek}] 0x2BC2 \\
Мероприятия: \textbf{twot} [\emph{twot}] 0xAC2A \\
меры: \textbf{hcos} [\emph{ʰʃos}] 0x8C68 \\
Месопотамия: \textbf{hmo7m} [\emph{ʰmo˥m}] 0x1408 \\
мессианский: \textbf{sye7n} [\emph{sje˥n}] 0x7309 \\
места: \textbf{tcas} [\emph{tʃäs}] 0x89AA \\
месяц: \textbf{mwas} [\emph{mwäs}] 0x8921 \\
метаданные: \textbf{hze2t} [\emph{ʰze˩t}] 0xBBA0 \\
металлоид: \textbf{tles} [\emph{tles}] 0x8B6A \\
металлы: \textbf{hma7s} [\emph{ʰmä˥s}] 0x9108 \\
метаморфоза: \textbf{pxaf} [\emph{pxäf}] 0xC9E4 \\
метан: \textbf{mwes} [\emph{mwes}] 0x8B21 \\
метатеза: \textbf{tfa2s} [\emph{tfä˩s}] 0x998A \\
метафизика: \textbf{hmi2k} [\emph{ʰmi˩k}] 0x3808 \\
метель: \textbf{bzut} [\emph{bzut}] 0xAA51 \\
метемпсихоз: \textbf{hme2s} [\emph{ʰme˩s}] 0x9B08 \\
метеор: \textbf{myos} [\emph{mjos}] 0x8C01 \\
метил: \textbf{hmi7f} [\emph{ʰmi˥f}] 0xD008 \\
метилен: \textbf{myi7c} [\emph{mji˥ʃ}] 0xF001 \\
метловище: \textbf{psu7s} [\emph{psu˥s}] 0x9244 \\
метонимия: \textbf{mwim} [\emph{mwim}] 0x0821 \\
метр: \textbf{mrip} [\emph{mrip}] 0x48C1 \\
метрический: \textbf{mre2k} [\emph{mre˩k}] 0x3BC1 \\
механический: \textbf{mrik} [\emph{mrik}] 0x28C1 \\
меццо-сопрано: \textbf{hzo7n} [\emph{ʰzo˥n}] 0x74A0 \\
меч: \textbf{swat} [\emph{swät}] 0xA929 \\
мечтатель: \textbf{hqa7n} [\emph{ʰŋä˥n}] 0x71C8 \\
мечтательность: \textbf{mwen} [\emph{mwen}] 0x6B21 \\
мешок: \textbf{cyek} [\emph{ʃjek}] 0x2B0D \\
миграция: \textbf{gric} [\emph{griʃ}] 0xE8D2 \\
мигрень: \textbf{xre2n} [\emph{xre˩n}] 0x7BDA \\
миелома: \textbf{slo7m} [\emph{slo˥m}] 0x1469 \\
мизантроп: \textbf{mri7p} [\emph{mri˥p}] 0x50C1 \\
мизантропия: \textbf{mri2p} [\emph{mri˩p}] 0x58C1 \\
мизинец: \textbf{pci7k} [\emph{pʃi˥k}] 0x30A4 \\
микология: \textbf{kli2c} [\emph{kli˩ʃ}] 0xF862 \\
микроавтобус: \textbf{hmi2p} [\emph{ʰmi˩p}] 0x5808 \\
микроб: \textbf{mrup} [\emph{mrup}] 0x4AC1 \\
микробиология: \textbf{mwoc} [\emph{mwoʃ}] 0xEC21 \\
микрометр: \textbf{mri7m} [\emph{mri˥m}] 0x10C1 \\
микрон: \textbf{mro2k} [\emph{mro˩k}] 0x3CC1 \\
Микронезия: \textbf{mra2c} [\emph{mrä˩ʃ}] 0xF9C1 \\
микроорганизмы: \textbf{hg6m} [\emph{ʰgəm}] 0x0D90 \\
микропроцессор: \textbf{mro7s} [\emph{mro˥s}] 0x94C1 \\
микроскопия: \textbf{kri7p} [\emph{kri˥p}] 0x50C2 \\
микрофон: \textbf{mrif} [\emph{mrif}] 0xC8C1 \\
микрочип: \textbf{mr6p} [\emph{mrəp}] 0x4DC1 \\
Милан: \textbf{hmi7n} [\emph{ʰmi˥n}] 0x7008 \\
милая: \textbf{hmic} [\emph{ʰmiʃ}] 0xE808 \\
мили: \textbf{hmif} [\emph{ʰmif}] 0xC808 \\
милитаризм: \textbf{cri7s} [\emph{ʃri˥s}] 0x90CD \\
милиция: \textbf{myip} [\emph{mjip}] 0x4801 \\
миллиард: \textbf{blip} [\emph{blip}] 0x4871 \\
миллиграмм: \textbf{hxam} [\emph{ʰxäm}] 0x09D0 \\
миллилитр: \textbf{vl6t} [\emph{vlət}] 0xAD76 \\
миллиметр: \textbf{myem} [\emph{mjem}] 0x0B01 \\
миллиона: \textbf{mwin} [\emph{mwin}] 0x6821 \\
миллионер: \textbf{lwi7n} [\emph{lwi˥n}] 0x702B \\
мим: \textbf{mwi7m} [\emph{mwi˥m}] 0x1021 \\
мимоза: \textbf{hxo7s} [\emph{ʰxo˥s}] 0x94D0 \\
миндаль: \textbf{bj6t} [\emph{bʒət}] 0xADB1 \\
минеральная: \textbf{mram} [\emph{mräm}] 0x09C1 \\
мини-бар: \textbf{hmi7p} [\emph{ʰmi˥p}] 0x5008 \\
миниатюрами: \textbf{mru2t} [\emph{mru˩t}] 0xBAC1 \\
министерство: \textbf{mret} [\emph{mret}] 0xABC1 \\
министр: \textbf{hmos} [\emph{ʰmos}] 0x8C08 \\
минут: \textbf{hmuk} [\emph{ʰmuk}] 0x2A08 \\
миокард: \textbf{mr67t} [\emph{mrə˥t}] 0xB5C1 \\
миоцен: \textbf{hmo2s} [\emph{ʰmo˩s}] 0x9C08 \\
мир: \textbf{cyan} [\emph{ʃjän}] 0x690D \\
Мир: \textbf{hwac} [\emph{ʰwäʃ}] 0xE928 \\
Мировой: \textbf{dyet} [\emph{djet}] 0xAB13 \\
миссионер: \textbf{xyic} [\emph{xjiʃ}] 0xE81A \\
Миссисипи: \textbf{hmi7s} [\emph{ʰmi˥s}] 0x9008 \\
миссия: \textbf{hmon} [\emph{ʰmon}] 0x6C08 \\
мистифицирован: \textbf{hmi2f} [\emph{ʰmi˩f}] 0xD808 \\
митра: \textbf{hgi2k} [\emph{ʰgi˩k}] 0x3890 \\
мифический: \textbf{swi2k} [\emph{swi˩k}] 0x3829 \\
млекопитающее: \textbf{hjaf} [\emph{ʰʒäf}] 0xC9A8 \\
мнемоника: \textbf{mwi2k} [\emph{mwi˩k}] 0x3821 \\
многогранник: \textbf{lyo2t} [\emph{ljo˩t}] 0xBC0B \\
многозадачных: \textbf{tx67n} [\emph{txə˥n}] 0x75EA \\
многокрасочный: \textbf{kro7k} [\emph{kro˥k}] 0x34C2 \\
многоразовый: \textbf{tfup} [\emph{tfup}] 0x4A8A \\
многосложный: \textbf{pla2p} [\emph{plä˩p}] 0x5964 \\
многочлен: \textbf{xlom} [\emph{xlom}] 0x0C7A \\
Моавитянин: \textbf{mwop} [\emph{mwop}] 0x4C21 \\
мобильность: \textbf{glot} [\emph{glot}] 0xAC72 \\
мог: \textbf{hkas} [\emph{ʰkäs}] 0x8910 \\
могила: \textbf{ksit} [\emph{ksit}] 0xA842 \\
могильный: \textbf{pla7k} [\emph{plä˥k}] 0x3164 \\
мода: \textbf{mwun} [\emph{mwun}] 0x6A21 \\
моделирование: \textbf{dlem} [\emph{dlem}] 0x0B73 \\
модернизировать: \textbf{draf} [\emph{dräf}] 0xC9D3 \\
модификатор: \textbf{pfok} [\emph{pfok}] 0x2C84 \\
модный: \textbf{mro7t} [\emph{mro˥t}] 0xB4C1 \\
модулировать: \textbf{zlut} [\emph{zlut}] 0xAA74 \\
модуль: \textbf{myoc} [\emph{mjoʃ}] 0xEC01 \\
мозаика: \textbf{myok} [\emph{mjok}] 0x2C01 \\
Мозамбик: \textbf{hzop} [\emph{ʰzop}] 0x4CA0 \\
мозжечок: \textbf{hse7c} [\emph{ʰse˥ʃ}] 0xF348 \\
Моисей: \textbf{hmo7s} [\emph{ʰmo˥s}] 0x9408 \\
мой: \textbf{myan} [\emph{mjän}] 0x6901 \\
мокрота: \textbf{pre7n} [\emph{pre˥n}] 0x73C4 \\
мол: \textbf{bxac} [\emph{bxäʃ}] 0xE9F1 \\
Молдова: \textbf{mwof} [\emph{mwof}] 0xCC21 \\
молекулярная: \textbf{myis} [\emph{mjis}] 0x8801 \\
молибден: \textbf{myom} [\emph{mjom}] 0x0C01 \\
молиться: \textbf{pcam} [\emph{pʃäm}] 0x09A4 \\
моллюсками: \textbf{jluk} [\emph{ʒluk}] 0x2A75 \\
молодка: \textbf{plu7t} [\emph{plu˥t}] 0xB264 \\
молодой: \textbf{nwan} [\emph{nwän}] 0x6926 \\
молоко: \textbf{hdut} [\emph{ʰdut}] 0xAA98 \\
молотилка: \textbf{hx6c} [\emph{ʰxəʃ}] 0xEDD0 \\
молоток: \textbf{clut} [\emph{ʃlut}] 0xAA6D \\
молочай: \textbf{hpe7c} [\emph{ʰpe˥ʃ}] 0xF320 \\
молочник: \textbf{gla7n} [\emph{glä˥n}] 0x7172 \\
молочница: \textbf{xra2c} [\emph{xrä˩ʃ}] 0xF9DA \\
моль: \textbf{hmop} [\emph{ʰmop}] 0x4C08 \\
молярный: \textbf{mrof} [\emph{mrof}] 0xCCC1 \\
момент: \textbf{hmac} [\emph{ʰmäʃ}] 0xE908 \\
монада: \textbf{mya7s} [\emph{mjä˥s}] 0x9101 \\
Монако: \textbf{hn6k} [\emph{ʰnək}] 0x2D30 \\
монахи: \textbf{hxa7n} [\emph{ʰxä˥n}] 0x71D0 \\
монахини: \textbf{nya2n} [\emph{njä˩n}] 0x7906 \\
монахиня: \textbf{nwi7n} [\emph{nwi˥n}] 0x7026 \\
монашество: \textbf{nya2s} [\emph{njä˩s}] 0x9906 \\
Монголия: \textbf{mwok} [\emph{mwok}] 0x2C21 \\
монгольский: \textbf{hm62n} [\emph{ʰmə˩n}] 0x7D08 \\
монета: \textbf{djin} [\emph{dʒin}] 0x68B3 \\
монизм: \textbf{mwom} [\emph{mwom}] 0x0C21 \\
монитор: \textbf{mros} [\emph{mros}] 0x8CC1 \\
моногамия: \textbf{byim} [\emph{bjim}] 0x0811 \\
монограмма: \textbf{hno2m} [\emph{ʰno˩m}] 0x1C30 \\
монокль: \textbf{nyo7k} [\emph{njo˥k}] 0x3406 \\
монолит: \textbf{nw6t} [\emph{nwət}] 0xAD26 \\
монолог: \textbf{tli2k} [\emph{tli˩k}] 0x386A \\
монооксид: \textbf{mwi2s} [\emph{mwi˩s}] 0x9821 \\
монополист: \textbf{hge7t} [\emph{ʰge˥t}] 0xB390 \\
монополия: \textbf{hnop} [\emph{ʰnop}] 0x4C30 \\
монорельс: \textbf{mref} [\emph{mref}] 0xCBC1 \\
монохромный: \textbf{nro7m} [\emph{nro˥m}] 0x14C6 \\
монтаж: \textbf{hjoc} [\emph{ʰʒoʃ}] 0xECA8 \\
мопед: \textbf{slo2t} [\emph{slo˩t}] 0xBC69 \\
моральный: \textbf{myut} [\emph{mjut}] 0xAA01 \\
морг: \textbf{mra2k} [\emph{mrä˩k}] 0x39C1 \\
море: \textbf{hmit} [\emph{ʰmit}] 0xA808 \\
морена: \textbf{mre2n} [\emph{mre˩n}] 0x7BC1 \\
мореходный: \textbf{swe2n} [\emph{swe˩n}] 0x7B29 \\
морковь: \textbf{grok} [\emph{grok}] 0x2CD2 \\
мормон: \textbf{mr6m} [\emph{mrəm}] 0x0DC1 \\
морфема: \textbf{mrom} [\emph{mrom}] 0x0CC1 \\
морфий: \textbf{mri7f} [\emph{mri˥f}] 0xD0C1 \\
морфологический: \textbf{cloc} [\emph{ʃloʃ}] 0xEC6D \\
морщина: \textbf{dwun} [\emph{dwun}] 0x6A33 \\
моряк: \textbf{swis} [\emph{swis}] 0x8829 \\
Москва: \textbf{hmo7k} [\emph{ʰmo˥k}] 0x3408 \\
мост: \textbf{hdat} [\emph{ʰdät}] 0xA998 \\
мостильщик: \textbf{pli2f} [\emph{pli˩f}] 0xD864 \\
мотыги: \textbf{hxu2s} [\emph{ʰxu˩s}] 0x9AD0 \\
моча: \textbf{nruc} [\emph{nruʃ}] 0xEAC6 \\
мочевина: \textbf{hyu2s} [\emph{ʰju˩s}] 0x9A18 \\
мочеиспускание: \textbf{pco2n} [\emph{pʃo˩n}] 0x7CA4 \\
мочеполовой: \textbf{hru7c} [\emph{ʰru˥ʃ}] 0xF270 \\
мошонка: \textbf{hk6m} [\emph{ʰkəm}] 0x0D10 \\
мракобесие: \textbf{kyu7t} [\emph{kju˥t}] 0xB202 \\
мститель: \textbf{txe7t} [\emph{txe˥t}] 0xB3EA \\
мудрый: \textbf{hmis} [\emph{ʰmis}] 0x8808 \\
муж: \textbf{hsam} [\emph{ʰsäm}] 0x0948 \\
мужественность: \textbf{tsa7t} [\emph{tsä˥t}] 0xB14A \\
мужской: \textbf{nyan} [\emph{njän}] 0x6906 \\
муза: \textbf{my6s} [\emph{mjəs}] 0x8D01 \\
музей: \textbf{hmem} [\emph{ʰmem}] 0x0B08 \\
Музыка: \textbf{hmik} [\emph{ʰmik}] 0x2808 \\
мука: \textbf{mrin} [\emph{mrin}] 0x68C1 \\
мультимедиа: \textbf{myum} [\emph{mjum}] 0x0A01 \\
муляж: \textbf{zwak} [\emph{zwäk}] 0x2934 \\
муравей: \textbf{mrim} [\emph{mrim}] 0x08C1 \\
мускулатура: \textbf{slu2t} [\emph{slu˩t}] 0xBA69 \\
мускулистый: \textbf{bro2n} [\emph{bro˩n}] 0x7CD1 \\
мускус: \textbf{hso2s} [\emph{ʰso˩s}] 0x9C48 \\
мусульманка: \textbf{zlum} [\emph{zlum}] 0x0A74 \\
мутант: \textbf{xwun} [\emph{xwun}] 0x6A3A \\
муфтий: \textbf{myuf} [\emph{mjuf}] 0xCA01 \\
мученичество: \textbf{cra2t} [\emph{ʃrä˩t}] 0xB9CD \\
муэдзин: \textbf{mwu7n} [\emph{mwu˥n}] 0x7221 \\
мыза: \textbf{gzen} [\emph{gzen}] 0x6B52 \\
мыканица: \textbf{cla2k} [\emph{ʃlä˩k}] 0x396D \\
мыло: \textbf{swun} [\emph{swun}] 0x6A29 \\
мыть: \textbf{dlac} [\emph{dläʃ}] 0xE973 \\
мычание: \textbf{mwum} [\emph{mwum}] 0x0A21 \\
мышеловка: \textbf{hfu7t} [\emph{ʰfu˥t}] 0xB260 \\
мышца: \textbf{mric} [\emph{mriʃ}] 0xE8C1 \\
мышь: \textbf{myus} [\emph{mjus}] 0x8A01 \\
мышьяк: \textbf{ksip} [\emph{ksip}] 0x4842 \\
Мьянма: \textbf{my6m} [\emph{mjəm}] 0x0D01 \\
мэйнфреймов: \textbf{mrem} [\emph{mrem}] 0x0BC1 \\
мэр: \textbf{hmes} [\emph{ʰmes}] 0x8B08 \\
Мэри: \textbf{mri7t} [\emph{mri˥t}] 0xB0C1 \\
Мэтью: \textbf{hmu7t} [\emph{ʰmu˥t}] 0xB208 \\
Мюнхен: \textbf{hmu2n} [\emph{ʰmu˩n}] 0x7A08 \\
мягкий: \textbf{hmun} [\emph{ʰmun}] 0x6A08 \\
мясник: \textbf{hkuf} [\emph{ʰkuf}] 0xCA10 \\
мясо: \textbf{hmot} [\emph{ʰmot}] 0xAC08 \\
мята: \textbf{pxe7n} [\emph{pxe˥n}] 0x73E4 \\
мяукать: \textbf{hmo2m} [\emph{ʰmo˩m}] 0x1C08 \\
мяч: \textbf{hbot} [\emph{ʰbot}] 0xAC88 \\
\subsection{ н }
наблюдение: \textbf{tcen} [\emph{tʃen}] 0x6BAA \\
наветренный: \textbf{nwi2t} [\emph{nwi˩t}] 0xB826 \\
навигация: \textbf{bxen} [\emph{bxen}] 0x6BF1 \\
наводнение: \textbf{hnos} [\emph{ʰnos}] 0x8C30 \\
наволочка: \textbf{pwuk} [\emph{pwuk}] 0x2A24 \\
\subsection{ Н }
Навуходоносор: \textbf{hnu7s} [\emph{ʰnu˥s}] 0x9230 \\
навязанный: \textbf{tli7n} [\emph{tli˥n}] 0x706A \\
наглость: \textbf{pca2t} [\emph{pʃä˩t}] 0xB9A4 \\
нагота: \textbf{hn67t} [\emph{ʰnə˥t}] 0xB530 \\
награда: \textbf{rwan} [\emph{rwän}] 0x692E \\
наградные: \textbf{gra7n} [\emph{grä˥n}] 0x71D2 \\
нагрузка: \textbf{kro7p} [\emph{kro˥p}] 0x54C2 \\
надгортанник: \textbf{vlek} [\emph{vlek}] 0x2B76 \\
надежность: \textbf{cwim} [\emph{ʃwim}] 0x082D \\
наделять: \textbf{nyu7t} [\emph{nju˥t}] 0xB206 \\
надеяться: \textbf{cwan} [\emph{ʃwän}] 0x692D \\
надколенник: \textbf{pfuk} [\emph{pfuk}] 0x2A84 \\
надписи: \textbf{hni2n} [\emph{ʰni˩n}] 0x7830 \\
надпочечная: \textbf{jra2n} [\emph{ʒrä˩n}] 0x79D5 \\
надстройка: \textbf{qruk} [\emph{ŋruk}] 0x2AD9 \\
надуманным: \textbf{hxa7c} [\emph{ʰxä˥ʃ}] 0xF1D0 \\
Нажмите: \textbf{pran} [\emph{prän}] 0x69C4 \\
Назарет: \textbf{nr6f} [\emph{nrəf}] 0xCDC6 \\
назидательный: \textbf{tfa7n} [\emph{tfä˥n}] 0x718A \\
назначенный: \textbf{frat} [\emph{frät}] 0xA9CC \\
наименее: \textbf{ksan} [\emph{ksän}] 0x6942 \\
наихудший: \textbf{psap} [\emph{psäp}] 0x4944 \\
найденный: \textbf{hfan} [\emph{ʰfän}] 0x6960 \\
наказание: \textbf{hjan} [\emph{ʰʒän}] 0x69A8 \\
наковальня: \textbf{nra7s} [\emph{nrä˥s}] 0x91C6 \\
накопление: \textbf{klun} [\emph{klun}] 0x6A62 \\
налогоплательщик: \textbf{nraf} [\emph{nräf}] 0xC9C6 \\
наложить: \textbf{dzan} [\emph{dzän}] 0x6953 \\
наложница: \textbf{hci2p} [\emph{ʰʃi˩p}] 0x5868 \\
наметка: \textbf{bwa2n} [\emph{bwä˩n}] 0x7931 \\
Намибия: \textbf{myep} [\emph{mjep}] 0x4B01 \\
нанка: \textbf{nya7k} [\emph{njä˥k}] 0x3106 \\
нанотехнология: \textbf{hno2c} [\emph{ʰno˩ʃ}] 0xFC30 \\
наоборот: \textbf{hq6s} [\emph{ʰŋəs}] 0x8DC8 \\
наперсток: \textbf{hq6n} [\emph{ʰŋən}] 0x6DC8 \\
написание: \textbf{jlin} [\emph{ʒlin}] 0x6875 \\
напиток: \textbf{hpuc} [\emph{ʰpuʃ}] 0xEA20 \\
наплыв: \textbf{nri7c} [\emph{nri˥ʃ}] 0xF0C6 \\
напоминание: \textbf{rwon} [\emph{rwon}] 0x6C2E \\
направление: \textbf{nyac} [\emph{njäʃ}] 0xE906 \\
направления: \textbf{hr6s} [\emph{ʰrəs}] 0x8D70 \\
направляющие: \textbf{hga7t} [\emph{ʰgä˥t}] 0xB190 \\
напряжение: \textbf{glet} [\emph{glet}] 0xAB72 \\
напыщенность: \textbf{dla7s} [\emph{dlä˥s}] 0x9173 \\
напыщенный: \textbf{gjut} [\emph{gʒut}] 0xAAB2 \\
нарды: \textbf{bvam} [\emph{bväm}] 0x0991 \\
наркоз: \textbf{zwas} [\emph{zwäs}] 0x8934 \\
наркоман: \textbf{hji2k} [\emph{ʰʒi˩k}] 0x38A8 \\
наружу: \textbf{qwas} [\emph{ŋwäs}] 0x8939 \\
наручник: \textbf{xyok} [\emph{xjok}] 0x2C1A \\
нарушать: \textbf{zlak} [\emph{zläk}] 0x2974 \\
нарушитель: \textbf{tru7t} [\emph{tru˥t}] 0xB2CA \\
нарцисс: \textbf{nri2s} [\emph{nri˩s}] 0x98C6 \\
нас: \textbf{hwam} [\emph{ʰwäm}] 0x0928 \\
насаженные: \textbf{pxi7s} [\emph{pxi˥s}] 0x90E4 \\
насекомое: \textbf{ksek} [\emph{ksek}] 0x2B42 \\
Население: \textbf{plak} [\emph{pläk}] 0x2964 \\
наслаждаться: \textbf{hcot} [\emph{ʰʃot}] 0xAC68 \\
наследник: \textbf{cris} [\emph{ʃris}] 0x88CD \\
насос: \textbf{pcop} [\emph{pʃop}] 0x4CA4 \\
наставник: \textbf{psi2t} [\emph{psi˩t}] 0xB844 \\
настоящим: \textbf{crip} [\emph{ʃrip}] 0x48CD \\
насыщенный: \textbf{hcu7t} [\emph{ʰʃu˥t}] 0xB268 \\
натальный: \textbf{hna7m} [\emph{ʰnä˥m}] 0x1130 \\
натирать: \textbf{ryac} [\emph{rjäʃ}] 0xE90E \\
натрий: \textbf{syom} [\emph{sjom}] 0x0C09 \\
наука: \textbf{flan} [\emph{flän}] 0x696C \\
наушники: \textbf{xren} [\emph{xren}] 0x6BDA \\
нафталин: \textbf{fle2n} [\emph{fle˩n}] 0x7B6C \\
нафтол: \textbf{pfa7f} [\emph{pfä˥f}] 0xD184 \\
нацарапал: \textbf{gre7t} [\emph{gre˥t}] 0xB3D2 \\
национальный: \textbf{ksam} [\emph{ksäm}] 0x0942 \\
начальство: \textbf{qras} [\emph{ŋräs}] 0x89D9 \\
начать: \textbf{cwas} [\emph{ʃwäs}] 0x892D \\
начитанный: \textbf{hlo7s} [\emph{ʰlo˥s}] 0x9458 \\
неамбициозны: \textbf{nyos} [\emph{njos}] 0x8C06 \\
Неаполь: \textbf{hna7s} [\emph{ʰnä˥s}] 0x9130 \\
небиблейский: \textbf{sri7c} [\emph{sri˥ʃ}] 0xF0C9 \\
неблагодарный: \textbf{qraf} [\emph{ŋräf}] 0xC9D9 \\
небный: \textbf{pla2c} [\emph{plä˩ʃ}] 0xF964 \\
небо: \textbf{cyak} [\emph{ʃjäk}] 0x290D \\
небоскреб: \textbf{qrap} [\emph{ŋräp}] 0x49D9 \\
небрежность: \textbf{sre7t} [\emph{sre˥t}] 0xB3C9 \\
невежественный: \textbf{nrin} [\emph{nrin}] 0x68C6 \\
невероятна: \textbf{jli2n} [\emph{ʒli˩n}] 0x7875 \\
невеста: \textbf{nrat} [\emph{nrät}] 0xA9C6 \\
невидимый: \textbf{dric} [\emph{driʃ}] 0xE8D3 \\
невинный: \textbf{nris} [\emph{nris}] 0x88C6 \\
невод: \textbf{hse7n} [\emph{ʰse˥n}] 0x7348 \\
невозможно: \textbf{psac} [\emph{psäʃ}] 0xE944 \\
невозмутимость: \textbf{jyit} [\emph{ʒjit}] 0xA815 \\
неволя: \textbf{cwot} [\emph{ʃwot}] 0xAC2D \\
невралгический: \textbf{nyi2k} [\emph{nji˩k}] 0x3806 \\
неврит: \textbf{nru2s} [\emph{nru˩s}] 0x9AC6 \\
неврология: \textbf{nro2k} [\emph{nro˩k}] 0x3CC6 \\
невропатический: \textbf{nri2p} [\emph{nri˩p}] 0x58C6 \\
невротический: \textbf{nru2t} [\emph{nru˩t}] 0xBAC6 \\
невыносимый: \textbf{tra7n} [\emph{trä˥n}] 0x71CA \\
невысказанный: \textbf{nyo7t} [\emph{njo˥t}] 0xB406 \\
негодяй: \textbf{hri7c} [\emph{ʰri˥ʃ}] 0xF070 \\
неграмотность: \textbf{nra2m} [\emph{nrä˩m}] 0x19C6 \\
негритянка: \textbf{nre2s} [\emph{nre˩s}] 0x9BC6 \\
негры: \textbf{hni7k} [\emph{ʰni˥k}] 0x3030 \\
недалекий: \textbf{xyi2t} [\emph{xji˩t}] 0xB81A \\
неделю: \textbf{hsus} [\emph{ʰsus}] 0x8A48 \\
недипломатичный: \textbf{dwom} [\emph{dwom}] 0x0C33 \\
недисциплинированный: \textbf{nya7n} [\emph{njä˥n}] 0x7106 \\
недоделанный: \textbf{dxo2t} [\emph{dxo˩t}] 0xBCF3 \\
недоедание: \textbf{hqu7n} [\emph{ʰŋu˥n}] 0x72C8 \\
недоразумение: \textbf{hli7m} [\emph{ʰli˥m}] 0x1058 \\
недостаток: \textbf{zwac} [\emph{zwäʃ}] 0xE934 \\
недостаточный: \textbf{pret} [\emph{pret}] 0xABC4 \\
недостойный: \textbf{dri7n} [\emph{dri˥n}] 0x70D3 \\
недоступен: \textbf{pxon} [\emph{pxon}] 0x6CE4 \\
недуги: \textbf{hma2p} [\emph{ʰmä˩p}] 0x5908 \\
Неемия: \textbf{hne7m} [\emph{ʰne˥m}] 0x1330 \\
неестественный: \textbf{hpa7c} [\emph{ʰpä˥ʃ}] 0xF120 \\
нежелание: \textbf{glin} [\emph{glin}] 0x6872 \\
незавидный: \textbf{gvip} [\emph{gvip}] 0x4892 \\
незагрязненный: \textbf{drum} [\emph{drum}] 0x0AD3 \\
незадачливый: \textbf{hba7m} [\emph{ʰbä˥m}] 0x1188 \\
незаконное: \textbf{hjif} [\emph{ʰʒif}] 0xC8A8 \\
незаметный: \textbf{nyi7k} [\emph{nji˥k}] 0x3006 \\
незамеченной: \textbf{pfec} [\emph{pfeʃ}] 0xEB84 \\
незатребованный: \textbf{hti7s} [\emph{ʰti˥s}] 0x9050 \\
нездоровый: \textbf{sy6s} [\emph{sjəs}] 0x8D09 \\
незначительно: \textbf{hpe2t} [\emph{ʰpe˩t}] 0xBB20 \\
незрелость: \textbf{txi2k} [\emph{txi˩k}] 0x38EA \\
неизбежный: \textbf{nwip} [\emph{nwip}] 0x4826 \\
неизвестности: \textbf{tya7n} [\emph{tjä˥n}] 0x710A \\
неизвестный: \textbf{klit} [\emph{klit}] 0xA862 \\
неизмеримый: \textbf{hka7p} [\emph{ʰkä˥p}] 0x5110 \\
неинтересный: \textbf{nri7s} [\emph{nri˥s}] 0x90C6 \\
неиспользованный: \textbf{hxok} [\emph{ʰxok}] 0x2CD0 \\
неисторично: \textbf{xro7k} [\emph{xro˥k}] 0x34DA \\
неисчерпаемый: \textbf{hti7p} [\emph{ʰti˥p}] 0x5050 \\
неисчислимый: \textbf{hqus} [\emph{ʰŋus}] 0x8AC8 \\
нейлон: \textbf{nye2n} [\emph{nje˩n}] 0x7B06 \\
нейрон: \textbf{nr6n} [\emph{nrən}] 0x6DC6 \\
нейтрон: \textbf{qrot} [\emph{ŋrot}] 0xACD9 \\
неквалифицированный: \textbf{sra2s} [\emph{srä˩s}] 0x99C9 \\
некоммерческий: \textbf{nri2t} [\emph{nri˩t}] 0xB8C6 \\
неконституционный: \textbf{nwi2n} [\emph{nwi˩n}] 0x7826 \\
неконтролируемый: \textbf{hko7p} [\emph{ʰko˥p}] 0x5410 \\
некритичный: \textbf{kri7k} [\emph{kri˥k}] 0x30C2 \\
некроз: \textbf{nro2s} [\emph{nro˩s}] 0x9CC6 \\
нелинейный: \textbf{hqin} [\emph{ʰŋin}] 0x68C8 \\
нелицензированный: \textbf{hne2s} [\emph{ʰne˩s}] 0x9B30 \\
неловкость: \textbf{hga7k} [\emph{ʰgä˥k}] 0x3190 \\
нелюдимый: \textbf{hqi7p} [\emph{ʰŋi˥p}] 0x50C8 \\
нематода: \textbf{kwum} [\emph{kwum}] 0x0A22 \\
Немецкий: \textbf{hjem} [\emph{ʰʒem}] 0x0BA8 \\
немного: \textbf{hye7k} [\emph{ʰje˥k}] 0x3318 \\
немодный: \textbf{hfo2t} [\emph{ʰfo˩t}] 0xBC60 \\
немытый: \textbf{hsa2c} [\emph{ʰsä˩ʃ}] 0xF948 \\
ненавидеть: \textbf{hcet} [\emph{ʰʃet}] 0xAB68 \\
ненадежный: \textbf{jlas} [\emph{ʒläs}] 0x8975 \\
ненароком: \textbf{nyi7n} [\emph{nji˥n}] 0x7006 \\
ненормальный: \textbf{kxit} [\emph{kxit}] 0xA8E2 \\
нео: \textbf{hnun} [\emph{ʰnun}] 0x6A30 \\
неограниченный: \textbf{kcim} [\emph{kʃim}] 0x08A2 \\
неодим: \textbf{nyom} [\emph{njom}] 0x0C06 \\
неоднозначность: \textbf{kca2n} [\emph{kʃä˩n}] 0x79A2 \\
неоднократно: \textbf{frak} [\emph{fräk}] 0x29CC \\
неоднородный: \textbf{xros} [\emph{xros}] 0x8CDA \\
неодобрительно: \textbf{dvan} [\emph{dvän}] 0x6993 \\
неоклассический: \textbf{ny6s} [\emph{njəs}] 0x8D06 \\
неокончательный: \textbf{tlu2t} [\emph{tlu˩t}] 0xBA6A \\
неолит: \textbf{hno2k} [\emph{ʰno˩k}] 0x3C30 \\
неологизмы: \textbf{jlim} [\emph{ʒlim}] 0x0875 \\
неон: \textbf{hn6n} [\emph{ʰnən}] 0x6D30 \\
неописуемый: \textbf{nra7n} [\emph{nrä˥n}] 0x71C6 \\
неоплаченный: \textbf{hwef} [\emph{ʰwef}] 0xCB28 \\
неопределенность: \textbf{pcet} [\emph{pʃet}] 0xABA4 \\
неопределенные: \textbf{pyif} [\emph{pjif}] 0xC804 \\
неопределенный: \textbf{hmec} [\emph{ʰmeʃ}] 0xEB08 \\
неопытность: \textbf{xyin} [\emph{xjin}] 0x681A \\
неоседланный: \textbf{jrep} [\emph{ʒrep}] 0x4BD5 \\
неосуществимость: \textbf{kca7t} [\emph{kʃä˥t}] 0xB1A2 \\
неотделимый: \textbf{sra7n} [\emph{srä˥n}] 0x71C9 \\
неотобранный: \textbf{hni7t} [\emph{ʰni˥t}] 0xB030 \\
неофициальный: \textbf{fris} [\emph{fris}] 0x88CC \\
неохотно: \textbf{hme7n} [\emph{ʰme˥n}] 0x7308 \\
неоценимый: \textbf{hwup} [\emph{ʰwup}] 0x4A28 \\
Непал: \textbf{hn6p} [\emph{ʰnəp}] 0x4D30 \\
непереходный: \textbf{trof} [\emph{trof}] 0xCCCA \\
неплатежеспособность: \textbf{dli7n} [\emph{dli˥n}] 0x7073 \\
непобедимый: \textbf{pson} [\emph{pson}] 0x6C44 \\
неповрежденный: \textbf{vyat} [\emph{vjät}] 0xA916 \\
неподдельный: \textbf{tfi7t} [\emph{tfi˥t}] 0xB08A \\
неподписанный: \textbf{pfi2t} [\emph{pfi˩t}] 0xB884 \\
неподтвержденный: \textbf{tfe2t} [\emph{tfe˩t}] 0xBB8A \\
неполный: \textbf{clet} [\emph{ʃlet}] 0xAB6D \\
непонятый: \textbf{pri2k} [\emph{pri˩k}] 0x38C4 \\
непослушание: \textbf{pso2n} [\emph{pso˩n}] 0x7C44 \\
непосредственно: \textbf{tsam} [\emph{tsäm}] 0x094A \\
непостижимый: \textbf{ksa2s} [\emph{ksä˩s}] 0x9942 \\
непостоянство: \textbf{vwen} [\emph{vwen}] 0x6B36 \\
неправдоподобно: \textbf{pli7s} [\emph{pli˥s}] 0x9064 \\
непредсказуемость: \textbf{tca2t} [\emph{tʃä˩t}] 0xB9AA \\
непредсказуемый: \textbf{xyip} [\emph{xjip}] 0x481A \\
непреклонный: \textbf{hbi7c} [\emph{ʰbi˥ʃ}] 0xF088 \\
непреложный: \textbf{pla7p} [\emph{plä˥p}] 0x5164 \\
непрерывность: \textbf{nwot} [\emph{nwot}] 0xAC26 \\
непригодность: \textbf{psu2t} [\emph{psu˩t}] 0xBA44 \\
неприемлемо: \textbf{kfip} [\emph{kfip}] 0x4882 \\
непримиримость: \textbf{tso7c} [\emph{tso˥ʃ}] 0xF44A \\
неприятный: \textbf{hpes} [\emph{ʰpes}] 0x8B20 \\
непроверенный: \textbf{tya2t} [\emph{tjä˩t}] 0xB90A \\
непродаваемый: \textbf{pli7p} [\emph{pli˥p}] 0x5064 \\
непродуктивный: \textbf{hsi7f} [\emph{ʰsi˥f}] 0xD048 \\
непрозрачный: \textbf{pxum} [\emph{pxum}] 0x0AE4 \\
непроизвольный: \textbf{vlif} [\emph{vlif}] 0xC876 \\
Непропеченный: \textbf{pxu7k} [\emph{pxu˥k}] 0x32E4 \\
непроходимый: \textbf{hpi7p} [\emph{ʰpi˥p}] 0x5020 \\
непрошеный: \textbf{bvot} [\emph{bvot}] 0xAC91 \\
Нептун: \textbf{xwen} [\emph{xwen}] 0x6B3A \\
неравенства: \textbf{pxa7t} [\emph{pxä˥t}] 0xB1E4 \\
неразборчивость: \textbf{kri2c} [\emph{kri˩ʃ}] 0xF8C2 \\
неразвитой: \textbf{bvit} [\emph{bvit}] 0xA891 \\
неразрезанный: \textbf{nwa7t} [\emph{nwä˥t}] 0xB126 \\
нераспечатанный: \textbf{hnu2n} [\emph{ʰnu˩n}] 0x7A30 \\
нераспространение: \textbf{pfa7s} [\emph{pfä˥s}] 0x9184 \\
нерастворимый: \textbf{xlun} [\emph{xlun}] 0x6A7A \\
нервный: \textbf{hnes} [\emph{ʰnes}] 0x8B30 \\
нерентабельный: \textbf{hno7m} [\emph{ʰno˥m}] 0x1430 \\
Нериум: \textbf{hn6m} [\emph{ʰnəm}] 0x0D30 \\
нерожденный: \textbf{nyo7n} [\emph{njo˥n}] 0x7406 \\
несбалансированность: \textbf{sl6p} [\emph{sləp}] 0x4D69 \\
несведущий: \textbf{pfi7t} [\emph{pfi˥t}] 0xB084 \\
несвязанный: \textbf{hko2k} [\emph{ʰko˩k}] 0x3C10 \\
несклонный: \textbf{hli2c} [\emph{ʰli˩ʃ}] 0xF858 \\
несколько: \textbf{kyak} [\emph{kjäk}] 0x2902 \\
неслышный: \textbf{hqi7n} [\emph{ʰŋi˥n}] 0x70C8 \\
несмешивающийся: \textbf{ps6p} [\emph{psəp}] 0x4D44 \\
несмываемый: \textbf{hki7m} [\emph{ʰki˥m}] 0x1010 \\
несогласный: \textbf{hdo7n} [\emph{ʰdo˥n}] 0x7498 \\
несогласованность: \textbf{hgi7t} [\emph{ʰgi˥t}] 0xB090 \\
несомненно: \textbf{hdi2p} [\emph{ʰdi˩p}] 0x5898 \\
неспортивный: \textbf{pyo7t} [\emph{pjo˥t}] 0xB404 \\
неспособность: \textbf{djat} [\emph{dʒät}] 0xA9B3 \\
несправедливость: \textbf{hni2t} [\emph{ʰni˩t}] 0xB830 \\
несправедливый: \textbf{hno7n} [\emph{ʰno˥n}] 0x7430 \\
несравненный: \textbf{hne7p} [\emph{ʰne˥p}] 0x5330 \\
несчастность: \textbf{kyen} [\emph{kjen}] 0x6B02 \\
Несчастные: \textbf{qro7t} [\emph{ŋro˥t}] 0xB4D9 \\
несъедобный: \textbf{kso7p} [\emph{kso˥p}] 0x5442 \\
нетерпение: \textbf{hpi7n} [\emph{ʰpi˥n}] 0x7020 \\
нетоксичен: \textbf{byos} [\emph{bjos}] 0x8C11 \\
нетрадиционный: \textbf{hxo7n} [\emph{ʰxo˥n}] 0x74D0 \\
неудачник: \textbf{pxa7s} [\emph{pxä˥s}] 0x91E4 \\
неудовлетворенность: \textbf{drun} [\emph{drun}] 0x6AD3 \\
неумело: \textbf{hno2s} [\emph{ʰno˩s}] 0x9C30 \\
неуместность: \textbf{zwen} [\emph{zwen}] 0x6B34 \\
неустойчивый: \textbf{swoc} [\emph{swoʃ}] 0xEC29 \\
неф: \textbf{hna7f} [\emph{ʰnä˥f}] 0xD130 \\
нефрит: \textbf{gyet} [\emph{gjet}] 0xAB12 \\
нефть: \textbf{prem} [\emph{prem}] 0x0BC4 \\
нечитаемый: \textbf{pxe7t} [\emph{pxe˥t}] 0xB3E4 \\
нива: \textbf{kfi2t} [\emph{kfi˩t}] 0xB882 \\
Нигер: \textbf{hn6c} [\emph{ʰnəʃ}] 0xED30 \\
Нигерия: \textbf{nr6c} [\emph{nrəʃ}] 0xEDC6 \\
нигилизм: \textbf{xli7s} [\emph{xli˥s}] 0x907A \\
нижеподписавшиеся: \textbf{xya2t} [\emph{xjä˩t}] 0xB91A \\
нижеприведенный: \textbf{dvun} [\emph{dvun}] 0x6A93 \\
низводить: \textbf{hze7t} [\emph{ʰze˥t}] 0xB3A0 \\
низкооплачиваемый: \textbf{dro2t} [\emph{dro˩t}] 0xBCD3 \\
Никарагуа: \textbf{nr6k} [\emph{nrək}] 0x2DC6 \\
никель: \textbf{nyek} [\emph{njek}] 0x2B06 \\
никогда: \textbf{hkap} [\emph{ʰkäp}] 0x4910 \\
Николай: \textbf{hni2s} [\emph{ʰni˩s}] 0x9830 \\
никотин: \textbf{hn67n} [\emph{ʰnə˥n}] 0x7530 \\
Нил: \textbf{hni2c} [\emph{ʰni˩ʃ}] 0xF830 \\
нимфомания: \textbf{mwe2n} [\emph{mwe˩n}] 0x7B21 \\
ниобий: \textbf{nyop} [\emph{njop}] 0x4C06 \\
нитрат: \textbf{nrem} [\emph{nrem}] 0x0BC6 \\
нитроглицерин: \textbf{tlo2s} [\emph{tlo˩s}] 0x9C6A \\
нить: \textbf{xl6n} [\emph{xlən}] 0x6D7A \\
ничего: \textbf{hmin} [\emph{ʰmin}] 0x6808 \\
нищенство: \textbf{hci7n} [\emph{ʰʃi˥n}] 0x7068 \\
нищий: \textbf{hxi7k} [\emph{ʰxi˥k}] 0x30D0 \\
НМУ: \textbf{ly6m} [\emph{ljəm}] 0x0D0B \\
но: \textbf{hlan} [\emph{ʰlän}] 0x6958 \\
новорожденный: \textbf{nw6n} [\emph{nwən}] 0x6D26 \\
Новости: \textbf{nwas} [\emph{nwäs}] 0x8926 \\
нога: \textbf{txe7k} [\emph{txe˥k}] 0x33EA \\
ногти: \textbf{nwis} [\emph{nwis}] 0x8826 \\
ножницы: \textbf{gjit} [\emph{gʒit}] 0xA8B2 \\
ножовка: \textbf{xwam} [\emph{xwäm}] 0x093A \\
номер: \textbf{hnuc} [\emph{ʰnuʃ}] 0xEA30 \\
номера: \textbf{kfom} [\emph{kfom}] 0x0C82 \\
номинальный: \textbf{nwon} [\emph{nwon}] 0x6C26 \\
Норвегия: \textbf{nrof} [\emph{nrof}] 0xCCC6 \\
норка: \textbf{mwuk} [\emph{mwuk}] 0x2A21 \\
нормандский: \textbf{hno2n} [\emph{ʰno˩n}] 0x7C30 \\
нос: \textbf{hnus} [\emph{ʰnus}] 0x8A30 \\
носить: \textbf{hban} [\emph{ʰbän}] 0x6988 \\
носок: \textbf{bzak} [\emph{bzäk}] 0x2951 \\
носорог: \textbf{nrus} [\emph{nrus}] 0x8AC6 \\
ностальгия: \textbf{xloc} [\emph{xloʃ}] 0xEC7A \\
ноу-хау: \textbf{vwaf} [\emph{vwäf}] 0xC936 \\
ночник: \textbf{cya7k} [\emph{ʃjä˥k}] 0x310D \\
ночь: \textbf{hyam} [\emph{ʰjäm}] 0x0918 \\
ноябрь: \textbf{nwap} [\emph{nwäp}] 0x4926 \\
нравоучительный: \textbf{hne7s} [\emph{ʰne˥s}] 0x9330 \\
нуклеотид: \textbf{klu2t} [\emph{klu˩t}] 0xBA62 \\
нуль: \textbf{zron} [\emph{zron}] 0x6CD4 \\
нумизмат: \textbf{hm62t} [\emph{ʰmə˩t}] 0xBD08 \\
нумизматика: \textbf{hmu2k} [\emph{ʰmu˩k}] 0x3A08 \\
Ньюфаундленд: \textbf{hnu7n} [\emph{ʰnu˥n}] 0x7230 \\
\subsection{ О }
О.Б.: \textbf{hno2p} [\emph{ʰno˩p}] 0x5C30 \\
оазис: \textbf{hvos} [\emph{ʰvos}] 0x8CB0 \\
обагрять: \textbf{hmu2p} [\emph{ʰmu˩p}] 0x5A08 \\
обвинить: \textbf{kcos} [\emph{kʃos}] 0x8CA2 \\
обдирать: \textbf{crem} [\emph{ʃrem}] 0x0BCD \\
обед: \textbf{lwan} [\emph{lwän}] 0x692B \\
обедающий: \textbf{dzac} [\emph{dzäʃ}] 0xE953 \\
обеднять: \textbf{cya7n} [\emph{ʃjä˥n}] 0x710D \\
обезглавить: \textbf{kles} [\emph{kles}] 0x8B62 \\
обезжиренное: \textbf{hsi7t} [\emph{ʰsi˥t}] 0xB048 \\
обезображивание: \textbf{fri7n} [\emph{fri˥n}] 0x70CC \\
обезьяна: \textbf{hzon} [\emph{ʰzon}] 0x6CA0 \\
обелиск: \textbf{bli7s} [\emph{bli˥s}] 0x9071 \\
обескуражена: \textbf{dra7t} [\emph{drä˥t}] 0xB1D3 \\
обескураживать: \textbf{dza2t} [\emph{dzä˩t}] 0xB953 \\
обеспечение: \textbf{hfec} [\emph{ʰfeʃ}] 0xEB60 \\
обесцвечивание: \textbf{plo7n} [\emph{plo˥n}] 0x7464 \\
обещания: \textbf{hci7t} [\emph{ʰʃi˥t}] 0xB068 \\
обжора: \textbf{sla7s} [\emph{slä˥s}] 0x9169 \\
обитать: \textbf{bxik} [\emph{bxik}] 0x28F1 \\
облагаемый: \textbf{kse2p} [\emph{kse˩p}] 0x5B42 \\
облако: \textbf{hyen} [\emph{ʰjen}] 0x6B18 \\
область: \textbf{kcec} [\emph{kʃeʃ}] 0xEBA2 \\
облачный: \textbf{hjut} [\emph{ʰʒut}] 0xAAA8 \\
облегающий: \textbf{cyi7k} [\emph{ʃji˥k}] 0x300D \\
облегчение: \textbf{gli2n} [\emph{gli˩n}] 0x7872 \\
обличительный: \textbf{bza7t} [\emph{bzä˥t}] 0xB151 \\
облучать: \textbf{byam} [\emph{bjäm}] 0x0911 \\
обмен: \textbf{tcec} [\emph{tʃeʃ}] 0xEBAA \\
обморок: \textbf{hwi7c} [\emph{ʰwi˥ʃ}] 0xF028 \\
обмотанный: \textbf{kri2t} [\emph{kri˩t}] 0xB8C2 \\
обналичить: \textbf{kca2c} [\emph{kʃä˩ʃ}] 0xF9A2 \\
обнаружений: \textbf{tx6c} [\emph{txəʃ}] 0xEDEA \\
обнаруживать: \textbf{dvek} [\emph{dvek}] 0x2B93 \\
обновление: \textbf{zrac} [\emph{zräʃ}] 0xE9D4 \\
обозначать: \textbf{dlut} [\emph{dlut}] 0xAA73 \\
обои: \textbf{xlep} [\emph{xlep}] 0x4B7A \\
оборотни: \textbf{rwef} [\emph{rwef}] 0xCB2E \\
обострение: \textbf{zra2n} [\emph{zrä˩n}] 0x79D4 \\
обработка: \textbf{tsen} [\emph{tsen}] 0x6B4A \\
образно: \textbf{pri7k} [\emph{pri˥k}] 0x30C4 \\
образовательный: \textbf{dyic} [\emph{djiʃ}] 0xE813 \\
обращенных: \textbf{kya7t} [\emph{kjä˥t}] 0xB102 \\
обрезываете: \textbf{kfa7t} [\emph{kfä˥t}] 0xB182 \\
обручи: \textbf{xwus} [\emph{xwus}] 0x8A3A \\
обручить: \textbf{brof} [\emph{brof}] 0xCCD1 \\
обстоятельства: \textbf{ksac} [\emph{ksäʃ}] 0xE942 \\
обсуждать: \textbf{dlak} [\emph{dläk}] 0x2973 \\
обувь: \textbf{hbut} [\emph{ʰbut}] 0xAA88 \\
обуздана: \textbf{kso7t} [\emph{kso˥t}] 0xB442 \\
обучение: \textbf{ryin} [\emph{rjin}] 0x680E \\
обхват: \textbf{hgef} [\emph{ʰgef}] 0xCB90 \\
Обще: \textbf{hko7n} [\emph{ʰko˥n}] 0x7410 \\
общеизвестный: \textbf{tyen} [\emph{tjen}] 0x6B0A \\
общественного: \textbf{hsik} [\emph{ʰsik}] 0x2848 \\
общинники: \textbf{hri2m} [\emph{ʰri˩m}] 0x1870 \\
объединение: \textbf{tfi7n} [\emph{tfi˥n}] 0x708A \\
объектив: \textbf{lyet} [\emph{ljet}] 0xAB0B \\
объекты: \textbf{psut} [\emph{psut}] 0xAA44 \\
обывательский: \textbf{fli2s} [\emph{fli˩s}] 0x986C \\
обычай: \textbf{hris} [\emph{ʰris}] 0x8870 \\
обязанный: \textbf{hba7t} [\emph{ʰbä˥t}] 0xB188 \\
обязательный: \textbf{pyic} [\emph{pjiʃ}] 0xE804 \\
обязать: \textbf{hboc} [\emph{ʰboʃ}] 0xEC88 \\
Овадия: \textbf{bya7t} [\emph{bjä˥t}] 0xB111 \\
овальный: \textbf{cwon} [\emph{ʃwon}] 0x6C2D \\
Овен: \textbf{jyis} [\emph{ʒjis}] 0x8815 \\
овес: \textbf{gwon} [\emph{gwon}] 0x6C32 \\
овощной: \textbf{srap} [\emph{sräp}] 0x49C9 \\
овсяница: \textbf{fyum} [\emph{fjum}] 0x0A0C \\
овуляция: \textbf{vla7n} [\emph{vlä˥n}] 0x7176 \\
овца: \textbf{hyec} [\emph{ʰjeʃ}] 0xEB18 \\
оглушительный: \textbf{dy6n} [\emph{djən}] 0x6D13 \\
огнемет: \textbf{pfo7t} [\emph{pfo˥t}] 0xB484 \\
огненный: \textbf{gxic} [\emph{gxiʃ}] 0xE8F2 \\
огнетушитель: \textbf{kxic} [\emph{kxiʃ}] 0xE8E2 \\
Огниво: \textbf{tse7k} [\emph{tse˥k}] 0x334A \\
оголять: \textbf{gyut} [\emph{gjut}] 0xAA12 \\
Огонь: \textbf{hfuk} [\emph{ʰfuk}] 0x2A60 \\
ограниченный: \textbf{blit} [\emph{blit}] 0xA871 \\
огрубевший: \textbf{klu7t} [\emph{klu˥t}] 0xB262 \\
огрубение: \textbf{bzuc} [\emph{bzuʃ}] 0xEA51 \\
огурец: \textbf{kyok} [\emph{kjok}] 0x2C02 \\
одаренный: \textbf{dxas} [\emph{dxäs}] 0x89F3 \\
одетый: \textbf{nyic} [\emph{njiʃ}] 0xE806 \\
одеяло: \textbf{bwat} [\emph{bwät}] 0xA931 \\
один: \textbf{hyik} [\emph{ʰjik}] 0x2818 \\
Одинокий: \textbf{hkok} [\emph{ʰkok}] 0x2C10 \\
одновременно: \textbf{nr67n} [\emph{nrə˥n}] 0x75C6 \\
одновременность: \textbf{hzec} [\emph{ʰzeʃ}] 0xEBA0 \\
однодневка: \textbf{hf6m} [\emph{ʰfəm}] 0x0D60 \\
однозначно: \textbf{mwec} [\emph{mweʃ}] 0xEB21 \\
одноклассник: \textbf{ksec} [\emph{kseʃ}] 0xEB42 \\
одноклеточный: \textbf{nwup} [\emph{nwup}] 0x4A26 \\
однопалатный: \textbf{hni2m} [\emph{ʰni˩m}] 0x1830 \\
однородный: \textbf{gjon} [\emph{gʒon}] 0x6CB2 \\
односложный: \textbf{nyi7p} [\emph{nji˥p}] 0x5006 \\
односторонними: \textbf{hpa2m} [\emph{ʰpä˩m}] 0x1920 \\
одолжил: \textbf{hde7t} [\emph{ʰde˥t}] 0xB398 \\
одометр: \textbf{dxom} [\emph{dxom}] 0x0CF3 \\
одуванчик: \textbf{vlon} [\emph{vlon}] 0x6C76 \\
оживление: \textbf{hri2n} [\emph{ʰri˩n}] 0x7870 \\
оживлять: \textbf{xli2n} [\emph{xli˩n}] 0x787A \\
ожидание: \textbf{tyuc} [\emph{tjuʃ}] 0xEA0A \\
ожирение: \textbf{hfa2p} [\emph{ʰfä˩p}] 0x5960 \\
озабоченность: \textbf{zyac} [\emph{zjäʃ}] 0xE914 \\
озон: \textbf{zyon} [\emph{zjon}] 0x6C14 \\
ой: \textbf{hdoc} [\emph{ʰdoʃ}] 0xEC98 \\
окалина: \textbf{hqaf} [\emph{ʰŋäf}] 0xC9C8 \\
Океания: \textbf{tyo2n} [\emph{tjo˩n}] 0x7C0A \\
океанография: \textbf{xyif} [\emph{xjif}] 0xC81A \\
окислять: \textbf{ksif} [\emph{ksif}] 0xC842 \\
окисляться: \textbf{sla2f} [\emph{slä˩f}] 0xD969 \\
оккультист: \textbf{klo7s} [\emph{klo˥s}] 0x9462 \\
окно: \textbf{twac} [\emph{twäʃ}] 0xE92A \\
около: \textbf{hkip} [\emph{ʰkip}] 0x4810 \\
околополярный: \textbf{hku2p} [\emph{ʰku˩p}] 0x5A10 \\
околоцветник: \textbf{pri7p} [\emph{pri˥p}] 0x50C4 \\
окольный: \textbf{kyi7t} [\emph{kji˥t}] 0xB002 \\
окра: \textbf{bxom} [\emph{bxom}] 0x0CF1 \\
окрестности: \textbf{grin} [\emph{grin}] 0x68D2 \\
оксалат: \textbf{sle2t} [\emph{sle˩t}] 0xBB69 \\
оксихлорид: \textbf{kli7k} [\emph{kli˥k}] 0x3062 \\
оксюморон: \textbf{kf6n} [\emph{kfən}] 0x6D82 \\
октава: \textbf{hka2f} [\emph{ʰkä˩f}] 0xD910 \\
октан: \textbf{kwo2n} [\emph{kwo˩n}] 0x7C22 \\
октет: \textbf{hk6t} [\emph{ʰkət}] 0xAD10 \\
октября: \textbf{kfap} [\emph{kfäp}] 0x4982 \\
окунать: \textbf{pyum} [\emph{pjum}] 0x0A04 \\
олень: \textbf{hlus} [\emph{ʰlus}] 0x8A58 \\
оливковый: \textbf{zlin} [\emph{zlin}] 0x6874 \\
олигополия: \textbf{hwo7p} [\emph{ʰwo˥p}] 0x5428 \\
олигоцен: \textbf{lyi2s} [\emph{lji˩s}] 0x980B \\
Оман: \textbf{txom} [\emph{txom}] 0x0CEA \\
омега: \textbf{hmo7c} [\emph{ʰmo˥ʃ}] 0xF408 \\
омлет: \textbf{zlet} [\emph{zlet}] 0xAB74 \\
омовения: \textbf{byus} [\emph{bjus}] 0x8A11 \\
омоложение: \textbf{jwic} [\emph{ʒwiʃ}] 0xE835 \\
омоним: \textbf{xrom} [\emph{xrom}] 0x0CDA \\
омыляться: \textbf{sra2f} [\emph{srä˩f}] 0xD9C9 \\
ондатра: \textbf{sra2k} [\emph{srä˩k}] 0x39C9 \\
онлайн: \textbf{clon} [\emph{ʃlon}] 0x6C6D \\
ономатопея: \textbf{nr6p} [\emph{nrəp}] 0x4DC6 \\
онтологический: \textbf{tyi7k} [\emph{tji˥k}] 0x300A \\
опадение: \textbf{hbi2c} [\emph{ʰbi˩ʃ}] 0xF888 \\
опал: \textbf{pxop} [\emph{pxop}] 0x4CE4 \\
опасающийся: \textbf{hyos} [\emph{ʰjos}] 0x8C18 \\
опасно: \textbf{pcis} [\emph{pʃis}] 0x88A4 \\
опера: \textbf{hpop} [\emph{ʰpop}] 0x4C20 \\
операционная: \textbf{pros} [\emph{pros}] 0x8CC4 \\
оперетта: \textbf{pro2t} [\emph{pro˩t}] 0xBCC4 \\
оперять: \textbf{lyem} [\emph{ljem}] 0x0B0B \\
опечатка: \textbf{txi2p} [\emph{txi˩p}] 0x58EA \\
опилки: \textbf{tra2m} [\emph{trä˩m}] 0x19CA \\
описательный: \textbf{drut} [\emph{drut}] 0xAAD3 \\
описки: \textbf{rwa7p} [\emph{rwä˥p}] 0x512E \\
опиум: \textbf{fyom} [\emph{fjom}] 0x0C0C \\
опора: \textbf{blas} [\emph{bläs}] 0x8971 \\
опорос: \textbf{frup} [\emph{frup}] 0x4ACC \\
опоссум: \textbf{pso7s} [\emph{pso˥s}] 0x9444 \\
опоясанный: \textbf{gzet} [\emph{gzet}] 0xAB52 \\
оправданный: \textbf{bra2t} [\emph{brä˩t}] 0xB9D1 \\
оправдывать: \textbf{jrit} [\emph{ʒrit}] 0xA8D5 \\
определитель: \textbf{cri2n} [\emph{ʃri˩n}] 0x78CD \\
опреснение: \textbf{dw6n} [\emph{dwən}] 0x6D33 \\
опровержение: \textbf{hfo2n} [\emph{ʰfo˩n}] 0x7C60 \\
опрометчивость: \textbf{dxat} [\emph{dxät}] 0xA9F3 \\
оптимист: \textbf{txi7s} [\emph{txi˥s}] 0x90EA \\
оптический: \textbf{prok} [\emph{prok}] 0x2CC4 \\
оптовая: \textbf{xlof} [\emph{xlof}] 0xCC7A \\
опыление: \textbf{pf6n} [\emph{pfən}] 0x6D84 \\
опыт: \textbf{tyap} [\emph{tjäp}] 0x490A \\
опьяняющий: \textbf{hma2k} [\emph{ʰmä˩k}] 0x3908 \\
оранжад: \textbf{gjot} [\emph{gʒot}] 0xACB2 \\
оранжевый: \textbf{tron} [\emph{tron}] 0x6CCA \\
оратор: \textbf{psik} [\emph{psik}] 0x2844 \\
орбита: \textbf{rwot} [\emph{rwot}] 0xAC2E \\
оргазм: \textbf{gzam} [\emph{gzäm}] 0x0952 \\
органист: \textbf{gxos} [\emph{gxos}] 0x8CF2 \\
органический: \textbf{hgok} [\emph{ʰgok}] 0x2C90 \\
оргия: \textbf{nroc} [\emph{nroʃ}] 0xECC6 \\
орегано: \textbf{hqon} [\emph{ʰŋon}] 0x6CC8 \\
орех-пекан: \textbf{pfek} [\emph{pfek}] 0x2B84 \\
оригинальность: \textbf{kla2t} [\emph{klä˩t}] 0xB962 \\
ориентализм: \textbf{ryo7s} [\emph{rjo˥s}] 0x940E \\
оркестр: \textbf{kset} [\emph{kset}] 0xAB42 \\
оркестровать: \textbf{cre2t} [\emph{ʃre˩t}] 0xBBCD \\
орленок: \textbf{gla7k} [\emph{glä˥k}] 0x3172 \\
орлиный: \textbf{kla7n} [\emph{klä˥n}] 0x7162 \\
орнамент: \textbf{twon} [\emph{twon}] 0x6C2A \\
орнитология: \textbf{vlic} [\emph{vliʃ}] 0xE876 \\
орошение: \textbf{glik} [\emph{glik}] 0x2872 \\
ортопедия: \textbf{hfo7k} [\emph{ʰfo˥k}] 0x3460 \\
оружейник: \textbf{hre7c} [\emph{ʰre˥ʃ}] 0xF370 \\
оружейный: \textbf{tca2k} [\emph{tʃä˩k}] 0x39AA \\
оружие: \textbf{slic} [\emph{sliʃ}] 0xE869 \\
орфографический: \textbf{hqi7f} [\emph{ʰŋi˥f}] 0xD0C8 \\
орфоэпия: \textbf{fyop} [\emph{fjop}] 0x4C0C \\
орхидея: \textbf{ry6t} [\emph{rjət}] 0xAD0E \\
осада: \textbf{hwec} [\emph{ʰweʃ}] 0xEB28 \\
осадок: \textbf{djot} [\emph{dʒot}] 0xACB3 \\
Осанна: \textbf{hx62n} [\emph{ʰxə˩n}] 0x7DD0 \\
осведомленность: \textbf{hyes} [\emph{ʰjes}] 0x8B18 \\
осевой: \textbf{kci2c} [\emph{kʃi˩ʃ}] 0xF8A2 \\
оседлал: \textbf{psa7t} [\emph{psä˥t}] 0xB144 \\
осел: \textbf{lwac} [\emph{lwäʃ}] 0xE92B \\
осень: \textbf{hpoc} [\emph{ʰpoʃ}] 0xEC20 \\
осетия: \textbf{sya2c} [\emph{sjä˩ʃ}] 0xF909 \\
осетр: \textbf{hs62n} [\emph{ʰsə˩n}] 0x7D48 \\
Осия: \textbf{xwos} [\emph{xwos}] 0x8C3A \\
ослепительный: \textbf{hya7n} [\emph{ʰjä˥n}] 0x7118 \\
осматривает: \textbf{hzi7n} [\emph{ʰzi˥n}] 0x70A0 \\
осмос: \textbf{srof} [\emph{srof}] 0xCCC9 \\
особенно: \textbf{hcim} [\emph{ʰʃim}] 0x0868 \\
особы: \textbf{sw6s} [\emph{swəs}] 0x8D29 \\
осознанный: \textbf{syi7s} [\emph{sji˥s}] 0x9009 \\
осока: \textbf{hse2k} [\emph{ʰse˩k}] 0x3B48 \\
оспа: \textbf{vri2t} [\emph{vri˩t}] 0xB8D6 \\
оссифицировать: \textbf{swa2f} [\emph{swä˩f}] 0xD929 \\
оставаться: \textbf{hlaf} [\emph{ʰläf}] 0xC958 \\
оставил: \textbf{kcac} [\emph{kʃäʃ}] 0xE9A2 \\
останец: \textbf{twa2c} [\emph{twä˩ʃ}] 0xF92A \\
остановился: \textbf{tsit} [\emph{tsit}] 0xA84A \\
остатки: \textbf{twic} [\emph{twiʃ}] 0xE82A \\
остеопатия: \textbf{tx6p} [\emph{txəp}] 0x4DEA \\
остров: \textbf{tsap} [\emph{tsäp}] 0x494A \\
островок: \textbf{hzi7t} [\emph{ʰzi˥t}] 0xB0A0 \\
Острота: \textbf{jret} [\emph{ʒret}] 0xABD5 \\
остроугольный: \textbf{jyun} [\emph{ʒjun}] 0x6A15 \\
осциллограф: \textbf{flop} [\emph{flop}] 0x4C6C \\
осьминог: \textbf{tyop} [\emph{tjop}] 0x4C0A \\
отверстие: \textbf{kcun} [\emph{kʃun}] 0x6AA2 \\
отвертка: \textbf{jrap} [\emph{ʒräp}] 0x49D5 \\
отвесно: \textbf{dwec} [\emph{dweʃ}] 0xEB33 \\
ответ: \textbf{hnap} [\emph{ʰnäp}] 0x4930 \\
отвислый: \textbf{txe2s} [\emph{txe˩s}] 0x9BEA \\
отвлекаться: \textbf{djic} [\emph{dʒiʃ}] 0xE8B3 \\
отделимый: \textbf{pfi7n} [\emph{pfi˥n}] 0x7084 \\
отдых: \textbf{fric} [\emph{friʃ}] 0xE8CC \\
отдыхает: \textbf{tsi7t} [\emph{tsi˥t}] 0xB04A \\
отек: \textbf{djem} [\emph{dʒem}] 0x0BB3 \\
отечество: \textbf{htu7k} [\emph{ʰtu˥k}] 0x3250 \\
отиатрия: \textbf{tl67k} [\emph{tlə˥k}] 0x356A \\
отказ: \textbf{tfis} [\emph{tfis}] 0x888A \\
отказываться: \textbf{bvi7t} [\emph{bvi˥t}] 0xB091 \\
откармливать: \textbf{tfa2n} [\emph{tfä˩n}] 0x798A \\
отклонение: \textbf{dwic} [\emph{dwiʃ}] 0xE833 \\
отклоненный: \textbf{hxe7n} [\emph{ʰxe˥n}] 0x73D0 \\
Отключить: \textbf{dluk} [\emph{dluk}] 0x2A73 \\
откровение: \textbf{kcon} [\emph{kʃon}] 0x6CA2 \\
откровенность: \textbf{kra2t} [\emph{krä˩t}] 0xB9C2 \\
открытие: \textbf{plin} [\emph{plin}] 0x6864 \\
открытка: \textbf{psok} [\emph{psok}] 0x2C44 \\
отлаживать: \textbf{dyek} [\emph{djek}] 0x2B13 \\
отмель: \textbf{xlaf} [\emph{xläf}] 0xC97A \\
отмучивать: \textbf{bvek} [\emph{bvek}] 0x2B91 \\
относительно: \textbf{klet} [\emph{klet}] 0xAB62 \\
отозвала: \textbf{hti7t} [\emph{ʰti˥t}] 0xB050 \\
отпереть: \textbf{kla2k} [\emph{klä˩k}] 0x3962 \\
отправитель: \textbf{hp6n} [\emph{ʰpən}] 0x6D20 \\
Отправить: \textbf{hjac} [\emph{ʰʒäʃ}] 0xE9A8 \\
отпуск: \textbf{kfic} [\emph{kfiʃ}] 0xE882 \\
отражать: \textbf{tfik} [\emph{tfik}] 0x288A \\
отродье: \textbf{cwa7s} [\emph{ʃwä˥s}] 0x912D \\
отруби: \textbf{hfuf} [\emph{ʰfuf}] 0xCA60 \\
отрыжка: \textbf{kci7c} [\emph{kʃi˥ʃ}] 0xF0A2 \\
отслаиваться: \textbf{kf6t} [\emph{kfət}] 0xAD82 \\
отставка: \textbf{hfis} [\emph{ʰfis}] 0x8860 \\
отступ: \textbf{syo7t} [\emph{sjo˥t}] 0xB409 \\
отсутствие: \textbf{psaf} [\emph{psäf}] 0xC944 \\
оттуда: \textbf{hdu7t} [\emph{ʰdu˥t}] 0xB298 \\
отцеубийство: \textbf{hpi2f} [\emph{ʰpi˩f}] 0xD820 \\
отчаяние: \textbf{tran} [\emph{trän}] 0x69CA \\
отчество: \textbf{pra2m} [\emph{prä˩m}] 0x19C4 \\
отчим: \textbf{hja2t} [\emph{ʰʒä˩t}] 0xB9A8 \\
отчужденность: \textbf{qwon} [\emph{ŋwon}] 0x6C39 \\
отшельник: \textbf{ryis} [\emph{rjis}] 0x880E \\
офис: \textbf{hpus} [\emph{ʰpus}] 0x8A20 \\
официальный: \textbf{hkaf} [\emph{ʰkäf}] 0xC910 \\
офтальмология: \textbf{fwok} [\emph{fwok}] 0x2C2C \\
офтальмоскоп: \textbf{flo7k} [\emph{flo˥k}] 0x346C \\
охватывать: \textbf{hjak} [\emph{ʰʒäk}] 0x29A8 \\
охваченном: \textbf{hxi2n} [\emph{ʰxi˩n}] 0x78D0 \\
охлократия: \textbf{gr6s} [\emph{grəs}] 0x8DD2 \\
охота: \textbf{clik} [\emph{ʃlik}] 0x286D \\
охотно: \textbf{hxa2n} [\emph{ʰxä˩n}] 0x79D0 \\
оценить: \textbf{pluk} [\emph{pluk}] 0x2A64 \\
оценочный: \textbf{xlat} [\emph{xlät}] 0xA97A \\
оцепенение: \textbf{hma2m} [\emph{ʰmä˩m}] 0x1908 \\
очарование: \textbf{crak} [\emph{ʃräk}] 0x29CD \\
очаровывать: \textbf{hgif} [\emph{ʰgif}] 0xC890 \\
очевидец: \textbf{hci7k} [\emph{ʰʃi˥k}] 0x3068 \\
очевидность: \textbf{hc6n} [\emph{ʰʃən}] 0x6D68 \\
очистка: \textbf{twuk} [\emph{twuk}] 0x2A2A \\
очки: \textbf{hka7n} [\emph{ʰkä˥n}] 0x7110 \\
очный: \textbf{nri7m} [\emph{nri˥m}] 0x10C6 \\
ошибка: \textbf{croc} [\emph{ʃroʃ}] 0xECCD \\
\subsection{ п }
паб: \textbf{psop} [\emph{psop}] 0x4C44 \\
пава: \textbf{mri7c} [\emph{mri˥ʃ}] 0xF0C1 \\
павильон: \textbf{pla2n} [\emph{plä˩n}] 0x7964 \\
пагинация: \textbf{pyo2n} [\emph{pjo˩n}] 0x7C04 \\
пагода: \textbf{pwa2t} [\emph{pwä˩t}] 0xB924 \\
падать: \textbf{hfut} [\emph{ʰfut}] 0xAA60 \\
падение: \textbf{hdup} [\emph{ʰdup}] 0x4A98 \\
пажитник: \textbf{xrep} [\emph{xrep}] 0x4BDA \\
пакет: \textbf{hpek} [\emph{ʰpek}] 0x2B20 \\
\subsection{ П }
Пакистан: \textbf{ks6s} [\emph{ksəs}] 0x8D42 \\
пакля: \textbf{hme2k} [\emph{ʰme˩k}] 0x3B08 \\
палантин: \textbf{tci7p} [\emph{tʃi˥p}] 0x50AA \\
палатки: \textbf{txa7n} [\emph{txä˥n}] 0x71EA \\
палеозойский: \textbf{hpu7s} [\emph{ʰpu˥s}] 0x9220 \\
палеолитический: \textbf{plo7k} [\emph{plo˥k}] 0x3464 \\
палеонтолог: \textbf{gl6t} [\emph{glət}] 0xAD72 \\
палестинец: \textbf{ly6n} [\emph{ljən}] 0x6D0B \\
Палец: \textbf{hjit} [\emph{ʰʒit}] 0xA8A8 \\
палиндром: \textbf{lwem} [\emph{lwem}] 0x0B2B \\
палитра: \textbf{plep} [\emph{plep}] 0x4B64 \\
палладий: \textbf{plop} [\emph{plop}] 0x4C64 \\
паломник: \textbf{hpa7m} [\emph{ʰpä˥m}] 0x1120 \\
Пальма: \textbf{hpem} [\emph{ʰpem}] 0x0B20 \\
Пальто: \textbf{kwut} [\emph{kwut}] 0xAA22 \\
памятный: \textbf{mren} [\emph{mren}] 0x6BC1 \\
панамский: \textbf{pya7m} [\emph{pjä˥m}] 0x1104 \\
панегирик: \textbf{hti2k} [\emph{ʰti˩k}] 0x3850 \\
панегирист: \textbf{kya7s} [\emph{kjä˥s}] 0x9102 \\
панели: \textbf{pxep} [\emph{pxep}] 0x4BE4 \\
панический: \textbf{pwa7k} [\emph{pwä˥k}] 0x3124 \\
панк: \textbf{px6k} [\emph{pxək}] 0x2DE4 \\
панорамный: \textbf{nra7m} [\emph{nrä˥m}] 0x11C6 \\
панталоны: \textbf{xrot} [\emph{xrot}] 0xACDA \\
пантеизм: \textbf{hna2m} [\emph{ʰnä˩m}] 0x1930 \\
пантеон: \textbf{pfa2n} [\emph{pfä˩n}] 0x7984 \\
панцирь: \textbf{kra7c} [\emph{krä˥ʃ}] 0xF1C2 \\
папа: \textbf{hpap} [\emph{ʰpäp}] 0x4920 \\
папайя: \textbf{py6k} [\emph{pjək}] 0x2D04 \\
папирус: \textbf{pra2s} [\emph{prä˩s}] 0x99C4 \\
папоротник-орляк: \textbf{tsi2k} [\emph{tsi˩k}] 0x384A \\
папский: \textbf{pwa2n} [\emph{pwä˩n}] 0x7924 \\
папуасов: \textbf{pyup} [\emph{pjup}] 0x4A04 \\
пар: \textbf{xyon} [\emph{xjon}] 0x6C1A \\
парабола: \textbf{pwop} [\emph{pwop}] 0x4C24 \\
Парагвай: \textbf{pr6k} [\emph{prək}] 0x2DC4 \\
параграф: \textbf{tyaf} [\emph{tjäf}] 0xC90A \\
паразитизм: \textbf{hc6s} [\emph{ʰʃəs}] 0x8D68 \\
парализованный: \textbf{xlas} [\emph{xläs}] 0x897A \\
параллакс: \textbf{dlap} [\emph{dläp}] 0x4973 \\
параллельно: \textbf{pren} [\emph{pren}] 0x6BC4 \\
парамагнитный: \textbf{pwe2n} [\emph{pwe˩n}] 0x7B24 \\
парапланеризм: \textbf{plef} [\emph{plef}] 0xCB64 \\
параплегия: \textbf{pli7c} [\emph{pli˥ʃ}] 0xF064 \\
парашютист: \textbf{pca7p} [\emph{pʃä˥p}] 0x51A4 \\
параюридический: \textbf{pxik} [\emph{pxik}] 0x28E4 \\
паренхима: \textbf{pci2m} [\emph{pʃi˩m}] 0x18A4 \\
Париж: \textbf{pri7s} [\emph{pri˥s}] 0x90C4 \\
парик: \textbf{vrak} [\emph{vräk}] 0x29D6 \\
парикмахер: \textbf{lyes} [\emph{ljes}] 0x8B0B \\
паркет: \textbf{pra2k} [\emph{prä˩k}] 0x39C4 \\
пародия: \textbf{xri7t} [\emph{xri˥t}] 0xB0DA \\
пароксизм: \textbf{psa2s} [\emph{psä˩s}] 0x9944 \\
парообразный: \textbf{hva2s} [\emph{ʰvä˩s}] 0x99B0 \\
партеногенез: \textbf{tw6s} [\emph{twəs}] 0x8D2A \\
партер: \textbf{pref} [\emph{pref}] 0xCBC4 \\
партизанский: \textbf{grif} [\emph{grif}] 0xC8D2 \\
партитивный: \textbf{pfip} [\emph{pfip}] 0x4884 \\
партнерство: \textbf{grot} [\emph{grot}] 0xACD2 \\
парфюмерия: \textbf{pyem} [\emph{pjem}] 0x0B04 \\
паршивый: \textbf{pxi7k} [\emph{pxi˥k}] 0x30E4 \\
пасека: \textbf{qy6n} [\emph{ŋjən}] 0x6D19 \\
пассажиры: \textbf{hsa2k} [\emph{ʰsä˩k}] 0x3948 \\
пассивный: \textbf{prik} [\emph{prik}] 0x28C4 \\
пастораль: \textbf{pxok} [\emph{pxok}] 0x2CE4 \\
пасторат: \textbf{psi7n} [\emph{psi˥n}] 0x7044 \\
пастушьи: \textbf{hma2c} [\emph{ʰmä˩ʃ}] 0xF908 \\
Пасха: \textbf{ps6t} [\emph{psət}] 0xAD44 \\
пасынок: \textbf{tfe7s} [\emph{tfe˥s}] 0x938A \\
патернализм: \textbf{tli7m} [\emph{tli˥m}] 0x106A \\
патология: \textbf{pli7k} [\emph{pli˥k}] 0x3064 \\
патриархальный: \textbf{tra7c} [\emph{trä˥ʃ}] 0xF1CA \\
патриотизм: \textbf{hte7s} [\emph{ʰte˥s}] 0x9350 \\
патриций: \textbf{pri2t} [\emph{pri˩t}] 0xB8C4 \\
патронташ: \textbf{hba7f} [\emph{ʰbä˥f}] 0xD188 \\
патруль: \textbf{pruc} [\emph{pruʃ}] 0xEAC4 \\
патчи: \textbf{hpi2c} [\emph{ʰpi˩ʃ}] 0xF820 \\
паук: \textbf{bzit} [\emph{bzit}] 0xA851 \\
паутина: \textbf{mwa7n} [\emph{mwä˥n}] 0x7121 \\
пацифист: \textbf{pfa7t} [\emph{pfä˥t}] 0xB184 \\
пачули: \textbf{pxuc} [\emph{pxuʃ}] 0xEAE4 \\
паша: \textbf{pca7c} [\emph{pʃä˥ʃ}] 0xF1A4 \\
певцы: \textbf{hga7s} [\emph{ʰgä˥s}] 0x9190 \\
Пегас: \textbf{hpe7s} [\emph{ʰpe˥s}] 0x9320 \\
педагогический: \textbf{pya2k} [\emph{pjä˩k}] 0x3904 \\
педантичность: \textbf{pci2t} [\emph{pʃi˩t}] 0xB8A4 \\
педантка: \textbf{bzun} [\emph{bzun}] 0x6A51 \\
педиатрия: \textbf{hd6k} [\emph{ʰdək}] 0x2D98 \\
педикюр: \textbf{dyuk} [\emph{djuk}] 0x2A13 \\
пекарь: \textbf{hbek} [\emph{ʰbek}] 0x2B88 \\
Пекин: \textbf{hpe7k} [\emph{ʰpe˥k}] 0x3320 \\
пектин: \textbf{pce7n} [\emph{pʃe˥n}] 0x73A4 \\
пелагический: \textbf{lya7k} [\emph{ljä˥k}] 0x310B \\
Пелопоннес: \textbf{hpe2s} [\emph{ʰpe˩s}] 0x9B20 \\
пемза: \textbf{psum} [\emph{psum}] 0x0A44 \\
пена: \textbf{hfum} [\emph{ʰfum}] 0x0A60 \\
Пенджаб: \textbf{hpa2p} [\emph{ʰpä˩p}] 0x5920 \\
пенистый: \textbf{pro7t} [\emph{pro˥t}] 0xB4C4 \\
Пенсильвания: \textbf{pli7n} [\emph{pli˥n}] 0x7064 \\
пенсия: \textbf{psen} [\emph{psen}] 0x6B44 \\
пенсне: \textbf{pye7n} [\emph{pje˥n}] 0x7304 \\
пентан: \textbf{pwe7n} [\emph{pwe˥n}] 0x7324 \\
пентхауз: \textbf{pxa7f} [\emph{pxä˥f}] 0xD1E4 \\
пеньюар: \textbf{glec} [\emph{gleʃ}] 0xEB72 \\
пепельница: \textbf{frek} [\emph{frek}] 0x2BCC \\
первоисточник: \textbf{nru2n} [\emph{nru˩n}] 0x7AC6 \\
первокурсник: \textbf{drim} [\emph{drim}] 0x08D3 \\
первый: \textbf{hpam} [\emph{ʰpäm}] 0x0920 \\
переваливаться: \textbf{hw67t} [\emph{ʰwə˥t}] 0xB528 \\
перевалка: \textbf{tya7m} [\emph{tjä˥m}] 0x110A \\
перевод: \textbf{ryan} [\emph{rjän}] 0x690E \\
переводный: \textbf{tfa7c} [\emph{tfä˥ʃ}] 0xF18A \\
перевооружать: \textbf{sra2m} [\emph{srä˩m}] 0x19C9 \\
переделанный: \textbf{rwi2t} [\emph{rwi˩t}] 0xB82E \\
переделывать: \textbf{rwe2k} [\emph{rwe˩k}] 0x3B2E \\
переедать: \textbf{brok} [\emph{brok}] 0x2CD1 \\
пережить: \textbf{xri7p} [\emph{xri˥p}] 0x50DA \\
перезапуск: \textbf{hxac} [\emph{ʰxäʃ}] 0xE9D0 \\
переименовать: \textbf{kya7m} [\emph{kjä˥m}] 0x1102 \\
перекись: \textbf{pso7t} [\emph{pso˥t}] 0xB444 \\
переключатель: \textbf{gwic} [\emph{gwiʃ}] 0xE832 \\
перекристаллизация: \textbf{kli7c} [\emph{kli˥ʃ}] 0xF062 \\
перекрытие: \textbf{crup} [\emph{ʃrup}] 0x4ACD \\
перекупщик: \textbf{tli7c} [\emph{tli˥ʃ}] 0xF06A \\
переложение: \textbf{txen} [\emph{txen}] 0x6BEA \\
перемежаются: \textbf{nyi7s} [\emph{nji˥s}] 0x9006 \\
переменная: \textbf{prif} [\emph{prif}] 0xC8C4 \\
перемешивает: \textbf{gxa7t} [\emph{gxä˥t}] 0xB1F2 \\
перемирие: \textbf{cwut} [\emph{ʃwut}] 0xAA2D \\
Перенаправление: \textbf{hxen} [\emph{ʰxen}] 0x6BD0 \\
перепев: \textbf{pxa7c} [\emph{pxä˥ʃ}] 0xF1E4 \\
перепел: \textbf{bxap} [\emph{bxäp}] 0x49F1 \\
перепись: \textbf{hjes} [\emph{ʰʒes}] 0x8BA8 \\
переплачивать: \textbf{twep} [\emph{twep}] 0x4B2A \\
переплетенный: \textbf{twa7t} [\emph{twä˥t}] 0xB12A \\
переполненный: \textbf{hgic} [\emph{ʰgiʃ}] 0xE890 \\
переполнять: \textbf{hvut} [\emph{ʰvut}] 0xAAB0 \\
перерабатывать: \textbf{hxak} [\emph{ʰxäk}] 0x29D0 \\
переросший: \textbf{tra2c} [\emph{trä˩ʃ}] 0xF9CA \\
пересекать: \textbf{kyus} [\emph{kjus}] 0x8A02 \\
пересидеть: \textbf{hje7t} [\emph{ʰʒe˥t}] 0xB3A8 \\
перестановка: \textbf{pyef} [\emph{pjef}] 0xCB04 \\
перестраивать: \textbf{ryap} [\emph{rjäp}] 0x490E \\
перестрахование: \textbf{pru7n} [\emph{pru˥n}] 0x72C4 \\
переучет: \textbf{kce7k} [\emph{kʃe˥k}] 0x33A2 \\
перехитрить: \textbf{tyuf} [\emph{tjuf}] 0xCA0A \\
переходный: \textbf{tsom} [\emph{tsom}] 0x0C4A \\
перец: \textbf{hlec} [\emph{ʰleʃ}] 0xEB58 \\
перигей: \textbf{pri7c} [\emph{pri˥ʃ}] 0xF0C4 \\
перигелий: \textbf{ple7n} [\emph{ple˥n}] 0x7364 \\
перикард: \textbf{pra7m} [\emph{prä˥m}] 0x11C4 \\
периодичность: \textbf{pxo2t} [\emph{pxo˩t}] 0xBCE4 \\
перископ: \textbf{psi2k} [\emph{psi˩k}] 0x3844 \\
перистальтика: \textbf{pse2k} [\emph{pse˩k}] 0x3B44 \\
перитонит: \textbf{pye2t} [\emph{pje˩t}] 0xBB04 \\
перламутровый: \textbf{htu7m} [\emph{ʰtu˥m}] 0x1250 \\
перманганат: \textbf{hp67n} [\emph{ʰpə˥n}] 0x7520 \\
пермский: \textbf{pr6m} [\emph{prəm}] 0x0DC4 \\
пернатых: \textbf{pxut} [\emph{pxut}] 0xAAE4 \\
перс: \textbf{pr6s} [\emph{prəs}] 0x8DC4 \\
персик: \textbf{pxec} [\emph{pxeʃ}] 0xEBE4 \\
перспективный: \textbf{pwik} [\emph{pwik}] 0x2824 \\
Перу: \textbf{prup} [\emph{prup}] 0x4AC4 \\
перфоратор: \textbf{psoc} [\emph{psoʃ}] 0xEC44 \\
перфорировать: \textbf{fren} [\emph{fren}] 0x6BCC \\
перчатка: \textbf{zlat} [\emph{zlät}] 0xA974 \\
песенка: \textbf{dri7t} [\emph{dri˥t}] 0xB0D3 \\
пескарь: \textbf{mwi2n} [\emph{mwi˩n}] 0x7821 \\
песок: \textbf{hras} [\emph{ʰräs}] 0x8970 \\
пессарий: \textbf{psi7c} [\emph{psi˥ʃ}] 0xF044 \\
пестик: \textbf{pxus} [\emph{pxus}] 0x8AE4 \\
пестрый: \textbf{xy6t} [\emph{xjət}] 0xAD1A \\
песчаник: \textbf{qyat} [\emph{ŋjät}] 0xA919 \\
петрушка: \textbf{psus} [\emph{psus}] 0x8A44 \\
петушок: \textbf{kxoc} [\emph{kxoʃ}] 0xECE2 \\
пехота: \textbf{hpi2n} [\emph{ʰpi˩n}] 0x7820 \\
печень: \textbf{jwin} [\emph{ʒwin}] 0x6835 \\
пешка: \textbf{pwen} [\emph{pwen}] 0x6B24 \\
пещера: \textbf{hgec} [\emph{ʰgeʃ}] 0xEB90 \\
пианино: \textbf{pyon} [\emph{pjon}] 0x6C04 \\
пиво: \textbf{hbip} [\emph{ʰbip}] 0x4888 \\
пижма: \textbf{tca7c} [\emph{tʃä˥ʃ}] 0xF1AA \\
пикколо: \textbf{hpuf} [\emph{ʰpuf}] 0xCA20 \\
пиксель: \textbf{hxis} [\emph{ʰxis}] 0x88D0 \\
пиломатериалы: \textbf{hma7p} [\emph{ʰmä˥p}] 0x5108 \\
пилот: \textbf{plot} [\emph{plot}] 0xAC64 \\
пинг: \textbf{pli2n} [\emph{pli˩n}] 0x7864 \\
пингвин: \textbf{qwin} [\emph{ŋwin}] 0x6839 \\
пион: \textbf{py6n} [\emph{pjən}] 0x6D04 \\
пионер: \textbf{gzin} [\emph{gzin}] 0x6852 \\
пипетка: \textbf{pwep} [\emph{pwep}] 0x4B24 \\
пирамида: \textbf{crim} [\emph{ʃrim}] 0x08CD \\
пират: \textbf{jlat} [\emph{ʒlät}] 0xA975 \\
пиратство: \textbf{praf} [\emph{präf}] 0xC9C4 \\
Пиренеи: \textbf{pre7s} [\emph{pre˥s}] 0x93C4 \\
пиридин: \textbf{pyi7k} [\emph{pji˥k}] 0x3004 \\
пирит: \textbf{qri7t} [\emph{ŋri˥t}] 0xB0D9 \\
пирогов: \textbf{txuf} [\emph{txuf}] 0xCAEA \\
пироксен: \textbf{pro2s} [\emph{pro˩s}] 0x9CC4 \\
писец: \textbf{sri7s} [\emph{sri˥s}] 0x90C9 \\
письмо: \textbf{lyat} [\emph{ljät}] 0xA90B \\
питаясь: \textbf{xwit} [\emph{xwit}] 0xA83A \\
Питер: \textbf{hpi7t} [\emph{ʰpi˥t}] 0xB020 \\
питомники: \textbf{ky6s} [\emph{kjəs}] 0x8D02 \\
Пифагор: \textbf{hpa2s} [\emph{ʰpä˩s}] 0x9920 \\
пищеварение: \textbf{dxam} [\emph{dxäm}] 0x09F3 \\
плавать: \textbf{hzen} [\emph{ʰzen}] 0x6BA0 \\
плавиться: \textbf{hjec} [\emph{ʰʒeʃ}] 0xEBA8 \\
плавучесть: \textbf{fyun} [\emph{fjun}] 0x6A0C \\
плагиат: \textbf{hpi7c} [\emph{ʰpi˥ʃ}] 0xF020 \\
плазма: \textbf{psem} [\emph{psem}] 0x0B44 \\
плакать: \textbf{hkus} [\emph{ʰkus}] 0x8A10 \\
планер: \textbf{gvet} [\emph{gvet}] 0xAB92 \\
планета: \textbf{gren} [\emph{gren}] 0x6BD2 \\
планируемый: \textbf{hlot} [\emph{ʰlot}] 0xAC58 \\
планка: \textbf{hla7f} [\emph{ʰlä˥f}] 0xD158 \\
планктон: \textbf{pw6n} [\emph{pwən}] 0x6D24 \\
пластика: \textbf{psu2s} [\emph{psu˩s}] 0x9A44 \\
пластичность: \textbf{psuc} [\emph{psuʃ}] 0xEA44 \\
платина: \textbf{pyim} [\emph{pjim}] 0x0804 \\
платить: \textbf{hpis} [\emph{ʰpis}] 0x8820 \\
Платон: \textbf{plo2t} [\emph{plo˩t}] 0xBC64 \\
платонический: \textbf{pxa7n} [\emph{pxä˥n}] 0x71E4 \\
Платформа: \textbf{hpom} [\emph{ʰpom}] 0x0C20 \\
плацента: \textbf{px6n} [\emph{pxən}] 0x6DE4 \\
плащ: \textbf{pxet} [\emph{pxet}] 0xABE4 \\
плевательница: \textbf{pwi2n} [\emph{pwi˩n}] 0x7824 \\
плевра: \textbf{hpe2m} [\emph{ʰpe˩m}] 0x1B20 \\
плеврит: \textbf{pre2n} [\emph{pre˩n}] 0x7BC4 \\
плей-офф: \textbf{pyop} [\emph{pjop}] 0x4C04 \\
плейстоцен: \textbf{pse7s} [\emph{pse˥s}] 0x9344 \\
племя: \textbf{hpok} [\emph{ʰpok}] 0x2C20 \\
племянница: \textbf{bwan} [\emph{bwän}] 0x6931 \\
пленочный: \textbf{mre7n} [\emph{mre˥n}] 0x73C1 \\
пленум: \textbf{ple2n} [\emph{ple˩n}] 0x7B64 \\
плеоназм: \textbf{lyos} [\emph{ljos}] 0x8C0B \\
плечо: \textbf{cwac} [\emph{ʃwäʃ}] 0xE92D \\
Плеяды: \textbf{lya7s} [\emph{ljä˥s}] 0x910B \\
плиоцен: \textbf{pci7n} [\emph{pʃi˥n}] 0x70A4 \\
плита: \textbf{hcop} [\emph{ʰʃop}] 0x4C68 \\
плодоносить: \textbf{fwaf} [\emph{fwäf}] 0xC92C \\
плодородный: \textbf{prop} [\emph{prop}] 0x4CC4 \\
пломбир: \textbf{syu2t} [\emph{sju˩t}] 0xBA09 \\
плотник: \textbf{mrun} [\emph{mrun}] 0x6AC1 \\
плотоядный: \textbf{kxos} [\emph{kxos}] 0x8CE2 \\
плутократия: \textbf{tru2s} [\emph{tru˩s}] 0x9ACA \\
Плутон: \textbf{pwoc} [\emph{pwoʃ}] 0xEC24 \\
плутоний: \textbf{pxun} [\emph{pxun}] 0x6AE4 \\
плюс: \textbf{plus} [\emph{plus}] 0x8A64 \\
пляжный: \textbf{bjen} [\emph{bʒen}] 0x6BB1 \\
пневматический: \textbf{qwik} [\emph{ŋwik}] 0x2839 \\
пневмония: \textbf{myon} [\emph{mjon}] 0x6C01 \\
побег: \textbf{tcap} [\emph{tʃäp}] 0x49AA \\
поведение: \textbf{kfan} [\emph{kfän}] 0x6982 \\
поверхностно: \textbf{py6s} [\emph{pjəs}] 0x8D04 \\
поверхность: \textbf{hpum} [\emph{ʰpum}] 0x0A20 \\
поверь: \textbf{tso7t} [\emph{tso˥t}] 0xB44A \\
поверять: \textbf{hc6f} [\emph{ʰʃəf}] 0xCD68 \\
повесток: \textbf{qwis} [\emph{ŋwis}] 0x8839 \\
поворотливость: \textbf{xyi2n} [\emph{xji˩n}] 0x781A \\
повторять: \textbf{tri2s} [\emph{tri˩s}] 0x98CA \\
поганка: \textbf{tcu2t} [\emph{tʃu˩t}] 0xBAAA \\
погребать: \textbf{mre2t} [\emph{mre˩t}] 0xBBC1 \\
погром: \textbf{pro2m} [\emph{pro˩m}] 0x1CC4 \\
погруженный: \textbf{hto2m} [\emph{ʰto˩m}] 0x1C50 \\
погрузка: \textbf{hlo7t} [\emph{ʰlo˥t}] 0xB458 \\
подавлять: \textbf{hjas} [\emph{ʰʒäs}] 0x89A8 \\
подагра: \textbf{hgo7t} [\emph{ʰgo˥t}] 0xB490 \\
податливый: \textbf{lwi7t} [\emph{lwi˥t}] 0xB02B \\
подача: \textbf{pric} [\emph{priʃ}] 0xE8C4 \\
подбоченясь: \textbf{xyup} [\emph{xjup}] 0x4A1A \\
подвал: \textbf{byac} [\emph{bjäʃ}] 0xE911 \\
подведение: \textbf{syos} [\emph{sjos}] 0x8C09 \\
подвергать: \textbf{pfac} [\emph{pfäʃ}] 0xE984 \\
подветренный: \textbf{pwef} [\emph{pwef}] 0xCB24 \\
подвид: \textbf{bjas} [\emph{bʒäs}] 0x89B1 \\
подводить: \textbf{hvum} [\emph{ʰvum}] 0x0AB0 \\
подголовник: \textbf{hxe2t} [\emph{ʰxe˩t}] 0xBBD0 \\
подгорный: \textbf{dxon} [\emph{dxon}] 0x6CF3 \\
подготовка: \textbf{prip} [\emph{prip}] 0x48C4 \\
подгузник: \textbf{hd6p} [\emph{ʰdəp}] 0x4D98 \\
подделывать: \textbf{flek} [\emph{flek}] 0x2B6C \\
поддерживать: \textbf{qros} [\emph{ŋros}] 0x8CD9 \\
подержанный: \textbf{hr67n} [\emph{ʰrə˥n}] 0x7570 \\
поджилки: \textbf{xwoc} [\emph{xwoʃ}] 0xEC3A \\
подкожный: \textbf{pcem} [\emph{pʃem}] 0x0BA4 \\
подколенный: \textbf{plu2t} [\emph{plu˩t}] 0xBA64 \\
подлокотник: \textbf{mr6s} [\emph{mrəs}] 0x8DC1 \\
подметать: \textbf{swap} [\emph{swäp}] 0x4929 \\
подмножество: \textbf{psec} [\emph{pseʃ}] 0xEB44 \\
подножка: \textbf{pwe7t} [\emph{pwe˥t}] 0xB324 \\
поднятый: \textbf{kxap} [\emph{kxäp}] 0x49E2 \\
подозрительный: \textbf{cyim} [\emph{ʃjim}] 0x080D \\
подотряд: \textbf{brum} [\emph{brum}] 0x0AD1 \\
подписан: \textbf{hqi7k} [\emph{ʰŋi˥k}] 0x30C8 \\
подписываться: \textbf{brik} [\emph{brik}] 0x28D1 \\
подпись: \textbf{srim} [\emph{srim}] 0x08C9 \\
подпочва: \textbf{bxuk} [\emph{bxuk}] 0x2AF1 \\
подразделения: \textbf{dvis} [\emph{dvis}] 0x8893 \\
подруливающее: \textbf{hro7t} [\emph{ʰro˥t}] 0xB470 \\
подрядчик: \textbf{trek} [\emph{trek}] 0x2BCA \\
подсвечник: \textbf{kxi2t} [\emph{kxi˩t}] 0xB8E2 \\
подсемейство: \textbf{pfa7k} [\emph{pfä˥k}] 0x3184 \\
подслушивать: \textbf{zret} [\emph{zret}] 0xABD4 \\
подсознательный: \textbf{byis} [\emph{bjis}] 0x8811 \\
подсолнечник: \textbf{bjic} [\emph{bʒiʃ}] 0xE8B1 \\
подставка: \textbf{kl6t} [\emph{klət}] 0xAD62 \\
подстерегать: \textbf{fli2t} [\emph{fli˩t}] 0xB86C \\
подстрекательстве: \textbf{bya2t} [\emph{bjä˩t}] 0xB911 \\
подстроенный: \textbf{txi7k} [\emph{txi˥k}] 0x30EA \\
подстрочный: \textbf{gri7n} [\emph{gri˥n}] 0x70D2 \\
подсчет: \textbf{hgut} [\emph{ʰgut}] 0xAA90 \\
подтекст: \textbf{htu7t} [\emph{ʰtu˥t}] 0xB250 \\
подушка: \textbf{nyet} [\emph{njet}] 0xAB06 \\
подушки: \textbf{hku2n} [\emph{ʰku˩n}] 0x7A10 \\
подчеркивание: \textbf{hco7n} [\emph{ʰʃo˥n}] 0x7468 \\
подъязычная: \textbf{pco7t} [\emph{pʃo˥t}] 0xB4A4 \\
поезд: \textbf{pyik} [\emph{pjik}] 0x2804 \\
пожалел: \textbf{cyet} [\emph{ʃjet}] 0xAB0D \\
пожалуйста: \textbf{lyic} [\emph{ljiʃ}] 0xE80B \\
пожарник: \textbf{fra2n} [\emph{frä˩n}] 0x79CC \\
пожаротушение: \textbf{hvan} [\emph{ʰvän}] 0x69B0 \\
поза: \textbf{psam} [\emph{psäm}] 0x0944 \\
позвоночник: \textbf{gwit} [\emph{gwit}] 0xA832 \\
позвоночный: \textbf{vron} [\emph{vron}] 0x6CD6 \\
поздно: \textbf{tyam} [\emph{tjäm}] 0x090A \\
поздравлять: \textbf{tfut} [\emph{tfut}] 0xAA8A \\
позитивизм: \textbf{hz6s} [\emph{ʰzəs}] 0x8DA0 \\
позолоченный: \textbf{hgi2t} [\emph{ʰgi˩t}] 0xB890 \\
позор: \textbf{clam} [\emph{ʃläm}] 0x096D \\
позывной: \textbf{zla7n} [\emph{zlä˥n}] 0x7174 \\
поиск: \textbf{tsoc} [\emph{tsoʃ}] 0xEC4A \\
пойло: \textbf{kwa7k} [\emph{kwä˥k}] 0x3122 \\
показной: \textbf{hse2s} [\emph{ʰse˩s}] 0x9B48 \\
покалывание: \textbf{hka2k} [\emph{ʰkä˩k}] 0x3910 \\
покатый: \textbf{hti2c} [\emph{ʰti˩ʃ}] 0xF850 \\
поклонение: \textbf{htep} [\emph{ʰtep}] 0x4B50 \\
поклонник: \textbf{flas} [\emph{fläs}] 0x896C \\
поколение: \textbf{nric} [\emph{nriʃ}] 0xE8C6 \\
покорный: \textbf{txi2n} [\emph{txi˩n}] 0x78EA \\
покровитель: \textbf{proc} [\emph{proʃ}] 0xECC4 \\
покрывало: \textbf{bref} [\emph{bref}] 0xCBD1 \\
покрытосемянное: \textbf{psi7p} [\emph{psi˥p}] 0x5044 \\
покрытый: \textbf{kwak} [\emph{kwäk}] 0x2922 \\
покушаться: \textbf{rwa7k} [\emph{rwä˥k}] 0x312E \\
пол: \textbf{tlap} [\emph{tläp}] 0x496A \\
Пол: \textbf{hjep} [\emph{ʰʒep}] 0x4BA8 \\
полдень: \textbf{tcon} [\emph{tʃon}] 0x6CAA \\
полдюйма: \textbf{pfi2n} [\emph{pfi˩n}] 0x7884 \\
поле: \textbf{jlan} [\emph{ʒlän}] 0x6975 \\
полеводства: \textbf{hta2t} [\emph{ʰtä˩t}] 0xB950 \\
Полегче: \textbf{hsi2s} [\emph{ʰsi˩s}] 0x9848 \\
полемизировать: \textbf{nre7t} [\emph{nre˥t}] 0xB3C6 \\
полемический: \textbf{hg6k} [\emph{ʰgək}] 0x2D90 \\
поли: \textbf{pli2c} [\emph{pli˩ʃ}] 0xF864 \\
полигамия: \textbf{xlim} [\emph{xlim}] 0x087A \\
полиглот: \textbf{xli7k} [\emph{xli˥k}] 0x307A \\
полигональный: \textbf{pxo2n} [\emph{pxo˩n}] 0x7CE4 \\
полимер: \textbf{plem} [\emph{plem}] 0x0B64 \\
полиморфный: \textbf{pxof} [\emph{pxof}] 0xCCE4 \\
Полинезия: \textbf{pyoc} [\emph{pjoʃ}] 0xEC04 \\
полиомиелит: \textbf{pli7f} [\emph{pli˥f}] 0xD064 \\
политеистический: \textbf{hlu2t} [\emph{ʰlu˩t}] 0xBA58 \\
политическое: \textbf{syik} [\emph{sjik}] 0x2809 \\
полиция: \textbf{plis} [\emph{plis}] 0x8864 \\
полка: \textbf{claf} [\emph{ʃläf}] 0xC96D \\
полки: \textbf{hre2n} [\emph{ʰre˩n}] 0x7B70 \\
полковник: \textbf{hko2n} [\emph{ʰko˩n}] 0x7C10 \\
полнота: \textbf{krun} [\emph{krun}] 0x6AC2 \\
полноценный: \textbf{hl67n} [\emph{ʰlə˥n}] 0x7558 \\
полночь: \textbf{hdi7t} [\emph{ʰdi˥t}] 0xB098 \\
половик: \textbf{dvet} [\emph{dvet}] 0xAB93 \\
половина: \textbf{psan} [\emph{psän}] 0x6944 \\
положение: \textbf{hwa2t} [\emph{ʰwä˩t}] 0xB928 \\
полоса: \textbf{traf} [\emph{träf}] 0xC9CA \\
полотенце: \textbf{tluc} [\emph{tluʃ}] 0xEA6A \\
полуавтоматический: \textbf{swo2k} [\emph{swo˩k}] 0x3C29 \\
полубак: \textbf{hfa7p} [\emph{ʰfä˥p}] 0x5160 \\
полубог: \textbf{pce7t} [\emph{pʃe˥t}] 0xB3A4 \\
полугодовой: \textbf{fya7n} [\emph{fjä˥n}] 0x710C \\
полукруг: \textbf{hdi2k} [\emph{ʰdi˩k}] 0x3898 \\
полукруглый: \textbf{druk} [\emph{druk}] 0x2AD3 \\
полупроводник: \textbf{dzet} [\emph{dzet}] 0xAB53 \\
полупрозрачность: \textbf{tlu2n} [\emph{tlu˩n}] 0x7A6A \\
полутень: \textbf{hp6p} [\emph{ʰpəp}] 0x4D20 \\
получатель: \textbf{pses} [\emph{pses}] 0x8B44 \\
полученный: \textbf{hga2t} [\emph{ʰgä˩t}] 0xB990 \\
полый: \textbf{hcoc} [\emph{ʰʃoʃ}] 0xEC68 \\
пользователь: \textbf{hbuc} [\emph{ʰbuʃ}] 0xEA88 \\
Польша: \textbf{pl6n} [\emph{plən}] 0x6D64 \\
поля: \textbf{flit} [\emph{flit}] 0xA86C \\
полярность: \textbf{pr6t} [\emph{prət}] 0xADC4 \\
помадить: \textbf{pyo2t} [\emph{pjo˩t}] 0xBC04 \\
поместья: \textbf{hsa7n} [\emph{ʰsä˥n}] 0x7148 \\
помидор: \textbf{mwot} [\emph{mwot}] 0xAC21 \\
помилование: \textbf{hson} [\emph{ʰson}] 0x6C48 \\
помогать: \textbf{pfa2t} [\emph{pfä˩t}] 0xB984 \\
Помогите: \textbf{hput} [\emph{ʰput}] 0xAA20 \\
понедельник: \textbf{hnoc} [\emph{ʰnoʃ}] 0xEC30 \\
Понимаю: \textbf{pcat} [\emph{pʃät}] 0xA9A4 \\
пономарь: \textbf{kse2t} [\emph{kse˩t}] 0xBB42 \\
понятия: \textbf{kfen} [\emph{kfen}] 0x6B82 \\
поперечный: \textbf{xres} [\emph{xres}] 0x8BDA \\
Попкорн: \textbf{hpof} [\emph{ʰpof}] 0xCC20 \\
поплавок: \textbf{flut} [\emph{flut}] 0xAA6C \\
попугай: \textbf{qrut} [\emph{ŋrut}] 0xAAD9 \\
популярный: \textbf{hluc} [\emph{ʰluʃ}] 0xEA58 \\
попурри: \textbf{hpo2p} [\emph{ʰpo˩p}] 0x5C20 \\
поражать: \textbf{mrek} [\emph{mrek}] 0x2BC1 \\
поражения: \textbf{lyo2n} [\emph{ljo˩n}] 0x7C0B \\
пораженный: \textbf{ht67t} [\emph{ʰtə˥t}] 0xB550 \\
поразительный: \textbf{tri2n} [\emph{tri˩n}] 0x78CA \\
порванный: \textbf{hto7t} [\emph{ʰto˥t}] 0xB450 \\
порка: \textbf{hfa7k} [\emph{ʰfä˥k}] 0x3160 \\
порнографический: \textbf{pcof} [\emph{pʃof}] 0xCCA4 \\
порог: \textbf{frit} [\emph{frit}] 0xA8CC \\
породистый: \textbf{pru7t} [\emph{pru˥t}] 0xB2C4 \\
порох: \textbf{bru7t} [\emph{bru˥t}] 0xB2D1 \\
портной: \textbf{htef} [\emph{ʰtef}] 0xCB50 \\
портрет: \textbf{prot} [\emph{prot}] 0xACC4 \\
портретный: \textbf{kc6k} [\emph{kʃək}] 0x2DA2 \\
Португалия: \textbf{pxuk} [\emph{pxuk}] 0x2AE4 \\
португальский: \textbf{hpi7k} [\emph{ʰpi˥k}] 0x3020 \\
портфель: \textbf{bwek} [\emph{bwek}] 0x2B31 \\
порфир: \textbf{pfof} [\emph{pfof}] 0xCC84 \\
поршень: \textbf{ps6k} [\emph{psək}] 0x2D44 \\
порядковый: \textbf{rwuc} [\emph{rwuʃ}] 0xEA2E \\
посадка: \textbf{trun} [\emph{trun}] 0x6ACA \\
посаженный: \textbf{hpa2t} [\emph{ʰpä˩t}] 0xB920 \\
посвятить: \textbf{drin} [\emph{drin}] 0x68D3 \\
посещение: \textbf{fwan} [\emph{fwän}] 0x692C \\
послал: \textbf{hfat} [\emph{ʰfät}] 0xA960 \\
после: \textbf{hcap} [\emph{ʰʃäp}] 0x4968 \\
последний: \textbf{hcus} [\emph{ʰʃus}] 0x8A68 \\
последовательный: \textbf{rwoc} [\emph{rwoʃ}] 0xEC2E \\
послеобеденный: \textbf{sr67n} [\emph{srə˥n}] 0x75C9 \\
пословица: \textbf{hpe7t} [\emph{ʰpe˥t}] 0xB320 \\
послушание: \textbf{hto2n} [\emph{ʰto˩n}] 0x7C50 \\
посмертный: \textbf{sro7t} [\emph{sro˥t}] 0xB4C9 \\
посмешище: \textbf{hyi7p} [\emph{ʰji˥p}] 0x5018 \\
посмотреть: \textbf{htes} [\emph{ʰtes}] 0x8B50 \\
пособие: \textbf{bjat} [\emph{bʒät}] 0xA9B1 \\
посол: \textbf{tfas} [\emph{tfäs}] 0x898A \\
посольство: \textbf{dwas} [\emph{dwäs}] 0x8933 \\
поставка: \textbf{hsok} [\emph{ʰsok}] 0x2C48 \\
постель: \textbf{pwat} [\emph{pwät}] 0xA924 \\
постепенно: \textbf{hric} [\emph{ʰriʃ}] 0xE870 \\
постоянно: \textbf{htem} [\emph{ʰtem}] 0x0B50 \\
пот: \textbf{gyan} [\emph{gjän}] 0x6912 \\
поташ: \textbf{hpa7f} [\emph{ʰpä˥f}] 0xD120 \\
потворствовать: \textbf{hne2k} [\emph{ʰne˩k}] 0x3B30 \\
потерял: \textbf{hmut} [\emph{ʰmut}] 0xAA08 \\
поток: \textbf{lyam} [\emph{ljäm}] 0x090B \\
потребители: \textbf{bxos} [\emph{bxos}] 0x8CF1 \\
потребление: \textbf{hkof} [\emph{ʰkof}] 0xCC10 \\
потрошить: \textbf{dxep} [\emph{dxep}] 0x4BF3 \\
потряхивание: \textbf{dx6k} [\emph{dxək}] 0x2DF3 \\
потуги: \textbf{hg6s} [\emph{ʰgəs}] 0x8D90 \\
потухший: \textbf{hzik} [\emph{ʰzik}] 0x28A0 \\
похвальный: \textbf{txi7n} [\emph{txi˥n}] 0x70EA \\
похитители: \textbf{pca7s} [\emph{pʃä˥s}] 0x91A4 \\
походить: \textbf{hres} [\emph{ʰres}] 0x8B70 \\
похороны: \textbf{sran} [\emph{srän}] 0x69C9 \\
поцелуй: \textbf{hbus} [\emph{ʰbus}] 0x8A88 \\
почвопокровные: \textbf{hgof} [\emph{ʰgof}] 0xCC90 \\
почетной: \textbf{tyi2n} [\emph{tji˩n}] 0x780A \\
почка: \textbf{gjik} [\emph{gʒik}] 0x28B2 \\
Почта: \textbf{pxos} [\emph{pxos}] 0x8CE4 \\
почтальон: \textbf{vrat} [\emph{vrät}] 0xA9D6 \\
почтмейстер: \textbf{hpo7t} [\emph{ʰpo˥t}] 0xB420 \\
поэзия: \textbf{psis} [\emph{psis}] 0x8844 \\
поэт: \textbf{pwip} [\emph{pwip}] 0x4824 \\
поэтому: \textbf{hwe7k} [\emph{ʰwe˥k}] 0x3328 \\
появляться: \textbf{tcuc} [\emph{tʃuʃ}] 0xEAAA \\
поясница: \textbf{hli7s} [\emph{ʰli˥s}] 0x9058 \\
пр`оклятый: \textbf{vra2t} [\emph{vrä˩t}] 0xB9D6 \\
Пра-пра-бабушка: \textbf{jra7t} [\emph{ʒrä˥t}] 0xB1D5 \\
прабабушка: \textbf{hqa7m} [\emph{ʰŋä˥m}] 0x11C8 \\
правда: \textbf{syan} [\emph{sjän}] 0x6909 \\
правительство: \textbf{syun} [\emph{sjun}] 0x6A09 \\
правнук: \textbf{tsu7n} [\emph{tsu˥n}] 0x724A \\
правопреемник: \textbf{tfe2n} [\emph{tfe˩n}] 0x7B8A \\
Прага: \textbf{pra7k} [\emph{prä˥k}] 0x31C4 \\
прагматика: \textbf{pwi7t} [\emph{pwi˥t}] 0xB024 \\
прадед: \textbf{hsa7f} [\emph{ʰsä˥f}] 0xD148 \\
праздник: \textbf{byit} [\emph{bjit}] 0xA811 \\
праздно: \textbf{hsa7m} [\emph{ʰsä˥m}] 0x1148 \\
праздношатание: \textbf{hza2n} [\emph{ʰzä˩n}] 0x79A0 \\
празеодимий: \textbf{hp6m} [\emph{ʰpəm}] 0x0D20 \\
практическое: \textbf{hrik} [\emph{ʰrik}] 0x2870 \\
прачка: \textbf{hlu7t} [\emph{ʰlu˥t}] 0xB258 \\
превозносить: \textbf{pla7t} [\emph{plä˥t}] 0xB164 \\
превосходить: \textbf{kse2s} [\emph{kse˩s}] 0x9B42 \\
превосходящие: \textbf{cyip} [\emph{ʃjip}] 0x480D \\
превратности: \textbf{hsu2t} [\emph{ʰsu˩t}] 0xBA48 \\
превращение: \textbf{tsu2n} [\emph{tsu˩n}] 0x7A4A \\
превью: \textbf{hmu7s} [\emph{ʰmu˥s}] 0x9208 \\
преданно: \textbf{hvo2t} [\emph{ʰvo˩t}] 0xBCB0 \\
предварительно: \textbf{hpen} [\emph{ʰpen}] 0x6B20 \\
предвидели: \textbf{hyo7n} [\emph{ʰjo˥n}] 0x7418 \\
предвосхищать: \textbf{tye7t} [\emph{tje˥t}] 0xB30A \\
предзнаменование: \textbf{psun} [\emph{psun}] 0x6A44 \\
предикативный: \textbf{hpef} [\emph{ʰpef}] 0xCB20 \\
предлагать: \textbf{crek} [\emph{ʃrek}] 0x2BCD \\
предлог: \textbf{psi2s} [\emph{psi˩s}] 0x9844 \\
преднамеренность: \textbf{pyi2t} [\emph{pji˩t}] 0xB804 \\
предок: \textbf{hsuc} [\emph{ʰsuʃ}] 0xEA48 \\
предоплатой: \textbf{ple2t} [\emph{ple˩t}] 0xBB64 \\
предопределенный: \textbf{pwut} [\emph{pwut}] 0xAA24 \\
предопределять: \textbf{dji2n} [\emph{dʒi˩n}] 0x78B3 \\
предотвращение: \textbf{rwen} [\emph{rwen}] 0x6B2E \\
предохранитель: \textbf{pyus} [\emph{pjus}] 0x8A04 \\
предположим,: \textbf{hnaf} [\emph{ʰnäf}] 0xC930 \\
предпоследний: \textbf{slu7t} [\emph{slu˥t}] 0xB269 \\
предродовой: \textbf{bluc} [\emph{bluʃ}] 0xEA71 \\
представитель: \textbf{pfit} [\emph{pfit}] 0xA884 \\
представление: \textbf{pfic} [\emph{pfiʃ}] 0xE884 \\
предстоящий: \textbf{hvat} [\emph{ʰvät}] 0xA9B0 \\
предчувствие: \textbf{hri7m} [\emph{ʰri˥m}] 0x1070 \\
презерватив: \textbf{kr6m} [\emph{krəm}] 0x0DC2 \\
президент: \textbf{pcen} [\emph{pʃen}] 0x6BA4 \\
презирать: \textbf{dras} [\emph{dräs}] 0x89D3 \\
презрительный: \textbf{hti2s} [\emph{ʰti˩s}] 0x9850 \\
прекращение: \textbf{hzot} [\emph{ʰzot}] 0xACA0 \\
прелюбодей: \textbf{hza7k} [\emph{ʰzä˥k}] 0x31A0 \\
премолярный: \textbf{pyos} [\emph{pjos}] 0x8C04 \\
премьер-министр: \textbf{pxim} [\emph{pxim}] 0x08E4 \\
преобразователь: \textbf{ryuc} [\emph{rjuʃ}] 0xEA0E \\
препятствовать: \textbf{hzaf} [\emph{ʰzäf}] 0xC9A0 \\
прерывистый: \textbf{hqun} [\emph{ʰŋun}] 0x6AC8 \\
пресбиопия: \textbf{ryi7p} [\emph{rji˥p}] 0x500E \\
пресный: \textbf{hla2c} [\emph{ʰlä˩ʃ}] 0xF958 \\
пресс-банда: \textbf{pse7n} [\emph{pse˥n}] 0x7344 \\
преступил: \textbf{tra2p} [\emph{trä˩p}] 0x59CA \\
преступления: \textbf{hra2s} [\emph{ʰrä˩s}] 0x9970 \\
преступники: \textbf{pra2n} [\emph{prä˩n}] 0x79C4 \\
пресыщенный: \textbf{bzas} [\emph{bzäs}] 0x8951 \\
преувеличивать: \textbf{hvo7t} [\emph{ʰvo˥t}] 0xB4B0 \\
преуспевать: \textbf{gxe2t} [\emph{gxe˩t}] 0xBBF2 \\
префикс: \textbf{pfus} [\emph{pfus}] 0x8A84 \\
прецессия: \textbf{pse2s} [\emph{pse˩s}] 0x9B44 \\
прибавлять: \textbf{tla2f} [\emph{tlä˩f}] 0xD96A \\
приблизительный: \textbf{psim} [\emph{psim}] 0x0844 \\
прибой: \textbf{hsof} [\emph{ʰsof}] 0xCC48 \\
прибыль: \textbf{hlif} [\emph{ʰlif}] 0xC858 \\
приведены: \textbf{hyu7t} [\emph{ʰju˥t}] 0xB218 \\
привередливый: \textbf{fyut} [\emph{fjut}] 0xAA0C \\
приветливо: \textbf{hme7p} [\emph{ʰme˥p}] 0x5308 \\
приветствовал: \textbf{hwa7t} [\emph{ʰwä˥t}] 0xB128 \\
прививка: \textbf{kfaf} [\emph{kfäf}] 0xC982 \\
привилегированный: \textbf{pcec} [\emph{pʃeʃ}] 0xEBA4 \\
привлечь: \textbf{cyuk} [\emph{ʃjuk}] 0x2A0D \\
привычный: \textbf{zwat} [\emph{zwät}] 0xA934 \\
привязь: \textbf{cwet} [\emph{ʃwet}] 0xAB2D \\
приглашение: \textbf{hwon} [\emph{ʰwon}] 0x6C28 \\
приглушенный: \textbf{xyoc} [\emph{xjoʃ}] 0xEC1A \\
пригород: \textbf{tyup} [\emph{tjup}] 0x4A0A \\
приготовился: \textbf{hre2t} [\emph{ʰre˩t}] 0xBB70 \\
приданое: \textbf{hdi7k} [\emph{ʰdi˥k}] 0x3098 \\
придатки: \textbf{pwem} [\emph{pwem}] 0x0B24 \\
придворный: \textbf{ksa2n} [\emph{ksä˩n}] 0x7942 \\
придерживалась: \textbf{hb67t} [\emph{ʰbə˥t}] 0xB588 \\
придирчивый: \textbf{ht6k} [\emph{ʰtək}] 0x2D50 \\
придыхательный: \textbf{sri2c} [\emph{sri˩ʃ}] 0xF8C9 \\
приехать: \textbf{hlam} [\emph{ʰläm}] 0x0958 \\
прижаться: \textbf{hbe2n} [\emph{ʰbe˩n}] 0x7B88 \\
приземистый: \textbf{bjut} [\emph{bʒut}] 0xAAB1 \\
призма: \textbf{hpi2m} [\emph{ʰpi˩m}] 0x1820 \\
признание: \textbf{krak} [\emph{kräk}] 0x29C2 \\
призывников: \textbf{jrip} [\emph{ʒrip}] 0x48D5 \\
приключение: \textbf{hdam} [\emph{ʰdäm}] 0x0998 \\
прикрепление: \textbf{vlan} [\emph{vlän}] 0x6976 \\
приличествовать: \textbf{xli2s} [\emph{xli˩s}] 0x987A \\
приличия: \textbf{lya7t} [\emph{ljä˥t}] 0xB10B \\
Приложения: \textbf{pyen} [\emph{pjen}] 0x6B04 \\
прильнуть: \textbf{swok} [\emph{swok}] 0x2C29 \\
примат: \textbf{hpi7m} [\emph{ʰpi˥m}] 0x1020 \\
применимый: \textbf{plik} [\emph{plik}] 0x2864 \\
пример: \textbf{tlam} [\emph{tläm}] 0x096A \\
примириться: \textbf{djac} [\emph{dʒäʃ}] 0xE9B3 \\
примитивный: \textbf{hrif} [\emph{ʰrif}] 0xC870 \\
принадлежит: \textbf{sla2t} [\emph{slä˩t}] 0xB969 \\
принаряжать: \textbf{dxef} [\emph{dxef}] 0xCBF3 \\
принимать: \textbf{ksap} [\emph{ksäp}] 0x4942 \\
приносить: \textbf{tlak} [\emph{tläk}] 0x296A \\
принтер: \textbf{pyet} [\emph{pjet}] 0xAB04 \\
принуждение: \textbf{kwuc} [\emph{kwuʃ}] 0xEA22 \\
принцесса: \textbf{hrus} [\emph{ʰrus}] 0x8A70 \\
принцип: \textbf{psip} [\emph{psip}] 0x4844 \\
приобретение: \textbf{kxi7n} [\emph{kxi˥n}] 0x70E2 \\
приозерный: \textbf{hle2p} [\emph{ʰle˩p}] 0x5B58 \\
приостановить: \textbf{slun} [\emph{slun}] 0x6A69 \\
приоткрытый: \textbf{hda7k} [\emph{ʰdä˥k}] 0x3198 \\
припарка: \textbf{hjof} [\emph{ʰʒof}] 0xCCA8 \\
приписывал: \textbf{hta2p} [\emph{ʰtä˩p}] 0x5950 \\
природный: \textbf{tyic} [\emph{tjiʃ}] 0xE80A \\
приручение: \textbf{dx6n} [\emph{dxən}] 0x6DF3 \\
приседать: \textbf{krus} [\emph{krus}] 0x8AC2 \\
приспособление: \textbf{dyon} [\emph{djon}] 0x6C13 \\
приспособленного: \textbf{hs67t} [\emph{ʰsə˥t}] 0xB548 \\
пристойный: \textbf{dr6s} [\emph{drəs}] 0x8DD3 \\
присутствие: \textbf{psit} [\emph{psit}] 0xA844 \\
притуплять: \textbf{bjam} [\emph{bʒäm}] 0x09B1 \\
прихожанка: \textbf{ryi7n} [\emph{rji˥n}] 0x700E \\
прихорашиваться: \textbf{pro7n} [\emph{pro˥n}] 0x74C4 \\
прицветник: \textbf{bxa7k} [\emph{bxä˥k}] 0x31F1 \\
прицеп: \textbf{tlef} [\emph{tlef}] 0xCB6A \\
причалы: \textbf{hpo2s} [\emph{ʰpo˩s}] 0x9C20 \\
причастие: \textbf{pfas} [\emph{pfäs}] 0x8984 \\
причинный: \textbf{ks6t} [\emph{ksət}] 0xAD42 \\
причуды: \textbf{psa7s} [\emph{psä˥s}] 0x9144 \\
пришелец: \textbf{hco7t} [\emph{ʰʃo˥t}] 0xB468 \\
приятель: \textbf{dyat} [\emph{djät}] 0xA913 \\
приятный: \textbf{ksat} [\emph{ksät}] 0xA942 \\
пробки: \textbf{kro2s} [\emph{kro˩s}] 0x9CC2 \\
пробормотал: \textbf{nru7t} [\emph{nru˥t}] 0xB2C6 \\
пробуждая: \textbf{hvik} [\emph{ʰvik}] 0x28B0 \\
провансальский: \textbf{pwa7n} [\emph{pwä˥n}] 0x7124 \\
проверенный: \textbf{pxi2t} [\emph{pxi˩t}] 0xB8E4 \\
проверка: \textbf{pfen} [\emph{pfen}] 0x6B84 \\
проверяемый: \textbf{ryi7s} [\emph{rji˥s}] 0x900E \\
провидение: \textbf{hpi2t} [\emph{ʰpi˩t}] 0xB820 \\
провинция: \textbf{hpon} [\emph{ʰpon}] 0x6C20 \\
проводимость: \textbf{dvot} [\emph{dvot}] 0xAC93 \\
проводной: \textbf{hwo7t} [\emph{ʰwo˥t}] 0xB428 \\
провокатор: \textbf{hvu7t} [\emph{ʰvu˥t}] 0xB2B0 \\
провоцировать: \textbf{rwok} [\emph{rwok}] 0x2C2E \\
прогнозирование: \textbf{bjin} [\emph{bʒin}] 0x68B1 \\
прогорклый: \textbf{hza2t} [\emph{ʰzä˩t}] 0xB9A0 \\
программа: \textbf{hrom} [\emph{ʰrom}] 0x0C70 \\
программист: \textbf{prom} [\emph{prom}] 0x0CC4 \\
прогресс: \textbf{tyis} [\emph{tjis}] 0x880A \\
продавать: \textbf{hbak} [\emph{ʰbäk}] 0x2988 \\
продавцы: \textbf{hs6n} [\emph{ʰsən}] 0x6D48 \\
продавщица: \textbf{swe7n} [\emph{swe˥n}] 0x7329 \\
продажный: \textbf{vri7t} [\emph{vri˥t}] 0xB0D6 \\
продлить: \textbf{pset} [\emph{pset}] 0xAB44 \\
продувание: \textbf{pfo2n} [\emph{pfo˩n}] 0x7C84 \\
продувка: \textbf{cri7n} [\emph{ʃri˥n}] 0x70CD \\
продуктовый: \textbf{twoc} [\emph{twoʃ}] 0xEC2A \\
проектор: \textbf{pyek} [\emph{pjek}] 0x2B04 \\
проза: \textbf{hra7s} [\emph{ʰrä˥s}] 0x9170 \\
прозвище: \textbf{pxi7n} [\emph{pxi˥n}] 0x70E4 \\
прозелитизм: \textbf{plo7s} [\emph{plo˥s}] 0x9464 \\
прозрачный: \textbf{tfam} [\emph{tfäm}] 0x098A \\
производители: \textbf{hvas} [\emph{ʰväs}] 0x89B0 \\
производитель: \textbf{tfet} [\emph{tfet}] 0xAB8A \\
производить: \textbf{psas} [\emph{psäs}] 0x8944 \\
произвольный: \textbf{hbim} [\emph{ʰbim}] 0x0888 \\
произнесение: \textbf{bvan} [\emph{bvän}] 0x6991 \\
происходящий: \textbf{zwan} [\emph{zwän}] 0x6934 \\
проказа: \textbf{kra2n} [\emph{krä˩n}] 0x79C2 \\
проказливый: \textbf{hc6p} [\emph{ʰʃəp}] 0x4D68 \\
проклинать: \textbf{pse7k} [\emph{pse˥k}] 0x3344 \\
проклятие: \textbf{htop} [\emph{ʰtop}] 0x4C50 \\
прокурор: \textbf{vyik} [\emph{vjik}] 0x2816 \\
пролетариат: \textbf{rya7t} [\emph{rjä˥t}] 0xB10E \\
пролив: \textbf{hxe7c} [\emph{ʰxe˥ʃ}] 0xF3D0 \\
проливной: \textbf{pxef} [\emph{pxef}] 0xCBE4 \\
пролог: \textbf{plum} [\emph{plum}] 0x0A64 \\
проложенный: \textbf{lwos} [\emph{lwos}] 0x8C2B \\
Прометей: \textbf{hpo7s} [\emph{ʰpo˥s}] 0x9420 \\
промокашка: \textbf{cli2t} [\emph{ʃli˩t}] 0xB86D \\
промышленность: \textbf{dyit} [\emph{djit}] 0xA813 \\
проницаемый: \textbf{psep} [\emph{psep}] 0x4B44 \\
проницательно: \textbf{hci2k} [\emph{ʰʃi˩k}] 0x3868 \\
проницательность: \textbf{hsi2m} [\emph{ʰsi˩m}] 0x1848 \\
пропадать: \textbf{hqac} [\emph{ʰŋäʃ}] 0xE9C8 \\
пропан: \textbf{hpo7p} [\emph{ʰpo˥p}] 0x5420 \\
пропаривания: \textbf{ht62n} [\emph{ʰtə˩n}] 0x7D50 \\
проповедовать: \textbf{prak} [\emph{präk}] 0x29C4 \\
пропорциональный: \textbf{ryip} [\emph{rjip}] 0x480E \\
прорастание: \textbf{qyen} [\emph{ŋjen}] 0x6B19 \\
пророк: \textbf{hya2t} [\emph{ʰjä˩t}] 0xB918 \\
просматривал: \textbf{pri2n} [\emph{pri˩n}] 0x78C4 \\
проснулся: \textbf{hket} [\emph{ʰket}] 0xAB10 \\
просодия: \textbf{ps6c} [\emph{psəʃ}] 0xED44 \\
проспать: \textbf{hvop} [\emph{ʰvop}] 0x4CB0 \\
проспект: \textbf{jwan} [\emph{ʒwän}] 0x6935 \\
проститутка: \textbf{cla7n} [\emph{ʃlä˥n}] 0x716D \\
проституция: \textbf{pya2n} [\emph{pjä˩n}] 0x7904 \\
простота: \textbf{glat} [\emph{glät}] 0xA972 \\
пространство-время: \textbf{swa2m} [\emph{swä˩m}] 0x1929 \\
просыпался: \textbf{hca7k} [\emph{ʰʃä˥k}] 0x3168 \\
протестант: \textbf{pca7n} [\emph{pʃä˥n}] 0x71A4 \\
против: \textbf{tyit} [\emph{tjit}] 0xA80A \\
противодействовать: \textbf{hfet} [\emph{ʰfet}] 0xAB60 \\
противоядный: \textbf{nwa2k} [\emph{nwä˩k}] 0x3926 \\
протокол: \textbf{ryok} [\emph{rjok}] 0x2C0E \\
протонов: \textbf{pwos} [\emph{pwos}] 0x8C24 \\
протоплазма: \textbf{hp6s} [\emph{ʰpəs}] 0x8D20 \\
прототипичный: \textbf{pyi7p} [\emph{pji˥p}] 0x5004 \\
профессионал: \textbf{pron} [\emph{pron}] 0x6CC4 \\
профессия: \textbf{hpec} [\emph{ʰpeʃ}] 0xEB20 \\
профессор: \textbf{pfos} [\emph{pfos}] 0x8C84 \\
профилактика: \textbf{pfa2k} [\emph{pfä˩k}] 0x3984 \\
профилирование: \textbf{pfoc} [\emph{pfoʃ}] 0xEC84 \\
проходки: \textbf{px6s} [\emph{pxəs}] 0x8DE4 \\
проходы: \textbf{xwac} [\emph{xwäʃ}] 0xE93A \\
прохождение: \textbf{psak} [\emph{psäk}] 0x2944 \\
прохожие: \textbf{psa2n} [\emph{psä˩n}] 0x7944 \\
процветать: \textbf{hfek} [\emph{ʰfek}] 0x2B60 \\
процессор: \textbf{ples} [\emph{ples}] 0x8B64 \\
процитировал: \textbf{kyit} [\emph{kjit}] 0xA802 \\
прочность: \textbf{htaf} [\emph{ʰtäf}] 0xC950 \\
Пруссия: \textbf{hpi7s} [\emph{ʰpi˥s}] 0x9020 \\
прыгнули: \textbf{tla2t} [\emph{tlä˩t}] 0xB96A \\
прыжок: \textbf{hlok} [\emph{ʰlok}] 0x2C58 \\
прыщ: \textbf{xlo2n} [\emph{xlo˩n}] 0x7C7A \\
псалтырь: \textbf{sla2m} [\emph{slä˩m}] 0x1969 \\
психиатр: \textbf{hl62t} [\emph{ʰlə˩t}] 0xBD58 \\
психиатрия: \textbf{qyic} [\emph{ŋjiʃ}] 0xE819 \\
психоанализ: \textbf{swi2n} [\emph{swi˩n}] 0x7829 \\
психология: \textbf{sloc} [\emph{sloʃ}] 0xEC69 \\
психотерапия: \textbf{hs6p} [\emph{ʰsəp}] 0x4D48 \\
псориаз: \textbf{syi2s} [\emph{sji˩s}] 0x9809 \\
птица: \textbf{pyaf} [\emph{pjäf}] 0xC904 \\
пугало: \textbf{sruf} [\emph{sruf}] 0xCAC9 \\
пудель: \textbf{psu7t} [\emph{psu˥t}] 0xB244 \\
пудинг: \textbf{pwun} [\emph{pwun}] 0x6A24 \\
пузырек: \textbf{hci7p} [\emph{ʰʃi˥p}] 0x5068 \\
пуленепробиваемый: \textbf{bluf} [\emph{bluf}] 0xCA71 \\
пульпа: \textbf{xlap} [\emph{xläp}] 0x497A \\
пульс: \textbf{plos} [\emph{plos}] 0x8C64 \\
пульсировать: \textbf{pw6t} [\emph{pwət}] 0xAD24 \\
пуля: \textbf{gluk} [\emph{gluk}] 0x2A72 \\
пунктирный: \textbf{tyi7t} [\emph{tji˥t}] 0xB00A \\
пунктуальный: \textbf{nwun} [\emph{nwun}] 0x6A26 \\
пунктуация: \textbf{hpu7n} [\emph{ʰpu˥n}] 0x7220 \\
пуриновых: \textbf{pruf} [\emph{pruf}] 0xCAC4 \\
пурпурный: \textbf{tris} [\emph{tris}] 0x88CA \\
пустельга: \textbf{kse7k} [\emph{kse˥k}] 0x3342 \\
пустыня: \textbf{hgat} [\emph{ʰgät}] 0xA990 \\
пустяки: \textbf{hta7f} [\emph{ʰtä˥f}] 0xD150 \\
путаница: \textbf{hbun} [\emph{ʰbun}] 0x6A88 \\
путешественник: \textbf{hyek} [\emph{ʰjek}] 0x2B18 \\
пушок: \textbf{flum} [\emph{flum}] 0x0A6C \\
пчела: \textbf{mwic} [\emph{mwiʃ}] 0xE821 \\
пшеничный: \textbf{hwi7p} [\emph{ʰwi˥p}] 0x5028 \\
пылесос: \textbf{xluf} [\emph{xluf}] 0xCA7A \\
пыли: \textbf{hduc} [\emph{ʰduʃ}] 0xEA98 \\
пылко: \textbf{hji7t} [\emph{ʰʒi˥t}] 0xB0A8 \\
пыльник: \textbf{xyaf} [\emph{xjäf}] 0xC91A \\
пыльца: \textbf{plaf} [\emph{pläf}] 0xC964 \\
пытаясь: \textbf{tsic} [\emph{tsiʃ}] 0xE84A \\
пьяный: \textbf{trik} [\emph{trik}] 0x28CA \\
Пэдди: \textbf{pxa2n} [\emph{pxä˩n}] 0x79E4 \\
пястная: \textbf{tr6p} [\emph{trəp}] 0x4DCA \\
пятерка: \textbf{hfa7f} [\emph{ʰfä˥f}] 0xD160 \\
пятидесятый: \textbf{twi7s} [\emph{twi˥s}] 0x902A \\
Пятикнижие: \textbf{mye7t} [\emph{mje˥t}] 0xB301 \\
пятикратный: \textbf{hwu7p} [\emph{ʰwu˥p}] 0x5228 \\
пятилетие: \textbf{lwu7t} [\emph{lwu˥t}] 0xB22B \\
пятисотый: \textbf{pxo7t} [\emph{pxo˥t}] 0xB4E4 \\
пятиугольный: \textbf{pyo7n} [\emph{pjo˥n}] 0x7404 \\
пятница: \textbf{crut} [\emph{ʃrut}] 0xAACD \\
пятно: \textbf{hdos} [\emph{ʰdos}] 0x8C98 \\
\subsection{ р }
раб: \textbf{hnup} [\emph{ʰnup}] 0x4A30 \\
раболепствовать: \textbf{tlo7k} [\emph{tlo˥k}] 0x346A \\
работник: \textbf{nrek} [\emph{nrek}] 0x2BC6 \\
рабыня: \textbf{dvip} [\emph{dvip}] 0x4893 \\
равнины: \textbf{hma2n} [\emph{ʰmä˩n}] 0x7908 \\
равновесие: \textbf{plen} [\emph{plen}] 0x6B64 \\
равноденственный: \textbf{kcef} [\emph{kʃef}] 0xCBA2 \\
равнодушие: \textbf{dwin} [\emph{dwin}] 0x6833 \\
равносторонний: \textbf{sla7c} [\emph{slä˥ʃ}] 0xF169 \\
равноугольный: \textbf{qyuk} [\emph{ŋjuk}] 0x2A19 \\
радиальный: \textbf{ryif} [\emph{rjif}] 0xC80E \\
радиан: \textbf{xyan} [\emph{xjän}] 0x691A \\
радий: \textbf{ryem} [\emph{rjem}] 0x0B0E \\
радикальный: \textbf{ryo7n} [\emph{rjo˥n}] 0x740E \\
радио: \textbf{dyot} [\emph{djot}] 0xAC13 \\
радиоактивность: \textbf{ryoc} [\emph{rjoʃ}] 0xEC0E \\
радиология: \textbf{dwok} [\emph{dwok}] 0x2C33 \\
радиотелефон: \textbf{dle2n} [\emph{dle˩n}] 0x7B73 \\
радиоуглерод: \textbf{dyop} [\emph{djop}] 0x4C13 \\
радуга: \textbf{rwic} [\emph{rwiʃ}] 0xE82E \\
радужный: \textbf{rwi7n} [\emph{rwi˥n}] 0x702E \\
разборчиво: \textbf{lwep} [\emph{lwep}] 0x4B2B \\
разброс: \textbf{fras} [\emph{fräs}] 0x89CC \\
разбрызгиватель: \textbf{swem} [\emph{swem}] 0x0B29 \\
разведка: \textbf{sya7n} [\emph{sjä˥n}] 0x7109 \\
развертывание: \textbf{hfi7p} [\emph{ʰfi˥p}] 0x5060 \\
развертываться: \textbf{prum} [\emph{prum}] 0x0AC4 \\
развивать: \textbf{kwaf} [\emph{kwäf}] 0xC922 \\
развлекать: \textbf{nrut} [\emph{nrut}] 0xAAC6 \\
развод: \textbf{lwic} [\emph{lwiʃ}] 0xE82B \\
разврат: \textbf{tfi2k} [\emph{tfi˩k}] 0x388A \\
развязывать: \textbf{zla7t} [\emph{zlä˥t}] 0xB174 \\
разглагольствовать: \textbf{xra2t} [\emph{xrä˩t}] 0xB9DA \\
разговорчивый: \textbf{zwit} [\emph{zwit}] 0xA834 \\
разгружать: \textbf{gre7n} [\emph{gre˥n}] 0x73D2 \\
раздаваться: \textbf{hva2n} [\emph{ʰvä˩n}] 0x79B0 \\
раздвоенный: \textbf{kro7t} [\emph{kro˥t}] 0xB4C2 \\
раздевать: \textbf{hcof} [\emph{ʰʃof}] 0xCC68 \\
разделение: \textbf{zran} [\emph{zrän}] 0x69D4 \\
раздосадованный: \textbf{hfa7n} [\emph{ʰfä˥n}] 0x7160 \\
раздражительный: \textbf{dvat} [\emph{dvät}] 0xA993 \\
раздробить: \textbf{tsup} [\emph{tsup}] 0x4A4A \\
разлаживать: \textbf{dla2n} [\emph{dlä˩n}] 0x7973 \\
размежевание: \textbf{xwan} [\emph{xwän}] 0x693A \\
размещение: \textbf{kfun} [\emph{kfun}] 0x6A82 \\
размотанный: \textbf{nwa7n} [\emph{nwä˥n}] 0x7126 \\
разнообразие: \textbf{jwit} [\emph{ʒwit}] 0xA835 \\
разнообразный: \textbf{hbon} [\emph{ʰbon}] 0x6C88 \\
разносить: \textbf{sra7f} [\emph{srä˥f}] 0xD1C9 \\
разносторонний: \textbf{vyit} [\emph{vjit}] 0xA816 \\
разорение: \textbf{hbic} [\emph{ʰbiʃ}] 0xE888 \\
разрабатывать: \textbf{hqam} [\emph{ʰŋäm}] 0x09C8 \\
разреженный: \textbf{href} [\emph{ʰref}] 0xCB70 \\
разрез`ать: \textbf{hb6k} [\emph{ʰbək}] 0x2D88 \\
разрушающий: \textbf{cruc} [\emph{ʃruʃ}] 0xEACD \\
разряд: \textbf{pcuc} [\emph{pʃuʃ}] 0xEAA4 \\
разум: \textbf{hmas} [\emph{ʰmäs}] 0x8908 \\
разъедающий: \textbf{frus} [\emph{frus}] 0x8ACC \\
разъем: \textbf{djak} [\emph{dʒäk}] 0x29B3 \\
\subsection{ Р }
Райнланд: \textbf{ry67n} [\emph{rjə˥n}] 0x750E \\
районы: \textbf{hci2t} [\emph{ʰʃi˩t}] 0xB868 \\
рак: \textbf{kfas} [\emph{kfäs}] 0x8982 \\
ракета: \textbf{ryek} [\emph{rjek}] 0x2B0E \\
раковина: \textbf{swak} [\emph{swäk}] 0x2929 \\
рамадан: \textbf{my6t} [\emph{mjət}] 0xAD01 \\
Рамка: \textbf{tram} [\emph{träm}] 0x09CA \\
ранг: \textbf{tsuk} [\emph{tsuk}] 0x2A4A \\
Рангун: \textbf{rwa7n} [\emph{rwä˥n}] 0x712E \\
ранен: \textbf{fyas} [\emph{fjäs}] 0x890C \\
раскаленный: \textbf{hro7k} [\emph{ʰro˥k}] 0x3470 \\
раскалывать: \textbf{pxa7k} [\emph{pxä˥k}] 0x31E4 \\
раскаяние: \textbf{clip} [\emph{ʃlip}] 0x486D \\
расколачивающий: \textbf{rya7p} [\emph{rjä˥p}] 0x510E \\
раскольник: \textbf{dxi7t} [\emph{dxi˥t}] 0xB0F3 \\
расовой: \textbf{rwis} [\emph{rwis}] 0x882E \\
распаковывать: \textbf{dres} [\emph{dres}] 0x8BD3 \\
Распечатать: \textbf{trat} [\emph{trät}] 0xA9CA \\
расплата: \textbf{psot} [\emph{psot}] 0xAC44 \\
располагать: \textbf{syec} [\emph{sjeʃ}] 0xEB09 \\
распространять: \textbf{syap} [\emph{sjäp}] 0x4909 \\
рассвет: \textbf{hlon} [\emph{ʰlon}] 0x6C58 \\
рассветало: \textbf{dya7t} [\emph{djä˥t}] 0xB113 \\
рассечение: \textbf{drec} [\emph{dreʃ}] 0xEBD3 \\
рассеянно: \textbf{psa7n} [\emph{psä˥n}] 0x7144 \\
рассказчик: \textbf{rwi7t} [\emph{rwi˥t}] 0xB02E \\
расслаблять: \textbf{fwip} [\emph{fwip}] 0x482C \\
расставался: \textbf{hfa7t} [\emph{ʰfä˥t}] 0xB160 \\
расстояние: \textbf{htif} [\emph{ʰtif}] 0xC850 \\
расстройство: \textbf{drak} [\emph{dräk}] 0x29D3 \\
рассудите: \textbf{hdo7t} [\emph{ʰdo˥t}] 0xB498 \\
рассыпалась: \textbf{hpu7t} [\emph{ʰpu˥t}] 0xB220 \\
растворение: \textbf{ts6n} [\emph{tsən}] 0x6D4A \\
растворитель: \textbf{jlot} [\emph{ʒlot}] 0xAC75 \\
расти: \textbf{nwat} [\emph{nwät}] 0xA926 \\
расточительность: \textbf{vwac} [\emph{vwäʃ}] 0xE936 \\
растр: \textbf{ksa7f} [\emph{ksä˥f}] 0xD142 \\
растянутый: \textbf{dzut} [\emph{dzut}] 0xAA53 \\
расхищение: \textbf{dre2t} [\emph{dre˩t}] 0xBBD3 \\
расходы: \textbf{fyin} [\emph{fjin}] 0x680C \\
ратификация: \textbf{pfis} [\emph{pfis}] 0x8884 \\
рафинирование: \textbf{txin} [\emph{txin}] 0x68EA \\
рахит: \textbf{kxa2s} [\emph{kxä˩s}] 0x99E2 \\
рацемический: \textbf{ryi7m} [\emph{rji˥m}] 0x100E \\
рацион: \textbf{hr6n} [\emph{ʰrən}] 0x6D70 \\
рационализировать: \textbf{hvin} [\emph{ʰvin}] 0x68B0 \\
рациональный: \textbf{ryon} [\emph{rjon}] 0x6C0E \\
рвота: \textbf{hbom} [\emph{ʰbom}] 0x0C88 \\
рвотный: \textbf{hte2k} [\emph{ʰte˩k}] 0x3B50 \\
реабилитировать: \textbf{kle7t} [\emph{kle˥t}] 0xB362 \\
реакционер: \textbf{hra7c} [\emph{ʰrä˥ʃ}] 0xF170 \\
реакция: \textbf{ryak} [\emph{rjäk}] 0x290E \\
реализм: \textbf{cles} [\emph{ʃles}] 0x8B6D \\
реалистичный: \textbf{rwet} [\emph{rwet}] 0xAB2E \\
ребро: \textbf{lwuk} [\emph{lwuk}] 0x2A2B \\
ревнивый: \textbf{ryut} [\emph{rjut}] 0xAA0E \\
револьвер: \textbf{vlof} [\emph{vlof}] 0xCC76 \\
революция: \textbf{klon} [\emph{klon}] 0x6C62 \\
регби: \textbf{rwip} [\emph{rwip}] 0x482E \\
регистратор: \textbf{ryes} [\emph{rjes}] 0x8B0E \\
регулирование: \textbf{ryun} [\emph{rjun}] 0x6A0E \\
редакционный: \textbf{srok} [\emph{srok}] 0x2CC9 \\
редкий: \textbf{nran} [\emph{nrän}] 0x69C6 \\
резец: \textbf{vyis} [\emph{vjis}] 0x8816 \\
резидент: \textbf{zwin} [\emph{zwin}] 0x6834 \\
резка: \textbf{krot} [\emph{krot}] 0xACC2 \\
реинкарнация: \textbf{nref} [\emph{nref}] 0xCBC6 \\
рейнджер: \textbf{rwec} [\emph{rweʃ}] 0xEB2E \\
река: \textbf{hyun} [\emph{ʰjun}] 0x6A18 \\
реквизит: \textbf{qram} [\emph{ŋräm}] 0x09D9 \\
реквизиция: \textbf{rwi2n} [\emph{rwi˩n}] 0x782E \\
рекламирующий: \textbf{byum} [\emph{bjum}] 0x0A11 \\
рекламный: \textbf{pxan} [\emph{pxän}] 0x69E4 \\
рекламодатель: \textbf{vyan} [\emph{vjän}] 0x6916 \\
рекомбинация: \textbf{xyos} [\emph{xjos}] 0x8C1A \\
ректор: \textbf{hro2t} [\emph{ʰro˩t}] 0xBC70 \\
рекурсивный: \textbf{hri2f} [\emph{ʰri˩f}] 0xD870 \\
реле: \textbf{glek} [\emph{glek}] 0x2B72 \\
религия: \textbf{tcut} [\emph{tʃut}] 0xAAAA \\
ремесла: \textbf{kcof} [\emph{kʃof}] 0xCCA2 \\
ремесленник: \textbf{klip} [\emph{klip}] 0x4862 \\
ремонт: \textbf{hrum} [\emph{ʰrum}] 0x0A70 \\
Ренессанс: \textbf{hn6s} [\emph{ʰnəs}] 0x8D30 \\
рений: \textbf{hr6m} [\emph{ʰrəm}] 0x0D70 \\
Рентгеновский: \textbf{ksef} [\emph{ksef}] 0xCB42 \\
реостат: \textbf{hra7t} [\emph{ʰrä˥t}] 0xB170 \\
репа: \textbf{tlup} [\emph{tlup}] 0x4A6A \\
репертуар: \textbf{pre2t} [\emph{pre˩t}] 0xBBC4 \\
репортер: \textbf{hros} [\emph{ʰros}] 0x8C70 \\
рептилии: \textbf{pcef} [\emph{pʃef}] 0xCBA4 \\
рептилия: \textbf{xyot} [\emph{xjot}] 0xAC1A \\
ресница: \textbf{zlim} [\emph{zlim}] 0x0874 \\
республиканец: \textbf{pl6k} [\emph{plək}] 0x2D64 \\
реторты: \textbf{hro7p} [\emph{ʰro˥p}] 0x5470 \\
ретроспективный: \textbf{rwek} [\emph{rwek}] 0x2B2E \\
референдум: \textbf{grum} [\emph{grum}] 0x0AD2 \\
реформа: \textbf{hrok} [\emph{ʰrok}] 0x2C70 \\
рецензент: \textbf{hr6k} [\emph{ʰrək}] 0x2D70 \\
рецепт: \textbf{hrep} [\emph{ʰrep}] 0x4B70 \\
рецептор: \textbf{ksep} [\emph{ksep}] 0x4B42 \\
рецессивный: \textbf{hyef} [\emph{ʰjef}] 0xCB18 \\
рецидив: \textbf{hra2p} [\emph{ʰrä˩p}] 0x5970 \\
речь: \textbf{drac} [\emph{dräʃ}] 0xE9D3 \\
решение: \textbf{hden} [\emph{ʰden}] 0x6B98 \\
решительно: \textbf{tli7t} [\emph{tli˥t}] 0xB06A \\
Рея: \textbf{rya2k} [\emph{rjä˩k}] 0x390E \\
ржавчина: \textbf{hgos} [\emph{ʰgos}] 0x8C90 \\
ржанка: \textbf{hp6f} [\emph{ʰpəf}] 0xCD20 \\
ржать: \textbf{hji2t} [\emph{ʰʒi˩t}] 0xB8A8 \\
ризница: \textbf{hvi7t} [\emph{ʰvi˥t}] 0xB0B0 \\
рикошет: \textbf{tce2t} [\emph{tʃe˩t}] 0xBBAA \\
рикша: \textbf{rya2c} [\emph{rjä˩ʃ}] 0xF90E \\
ритм: \textbf{ryim} [\emph{rjim}] 0x080E \\
риторика: \textbf{tsok} [\emph{tsok}] 0x2C4A \\
риф: \textbf{byaf} [\emph{bjäf}] 0xC911 \\
Рияд: \textbf{ry67t} [\emph{rjə˥t}] 0xB50E \\
робот: \textbf{ryop} [\emph{rjop}] 0x4C0E \\
ров: \textbf{hxoc} [\emph{ʰxoʃ}] 0xECD0 \\
роговица: \textbf{grom} [\emph{grom}] 0x0CD2 \\
рогоносец: \textbf{hwo2t} [\emph{ʰwo˩t}] 0xBC28 \\
Роджер: \textbf{hro2c} [\emph{ʰro˩ʃ}] 0xFC70 \\
Родился: \textbf{tfin} [\emph{tfin}] 0x688A \\
родительский: \textbf{hyat} [\emph{ʰjät}] 0xA918 \\
Родос: \textbf{hro2s} [\emph{ʰro˩s}] 0x9C70 \\
родственный: \textbf{hgo2n} [\emph{ʰgo˩n}] 0x7C90 \\
роды: \textbf{cl6n} [\emph{ʃlən}] 0x6D6D \\
рождество: \textbf{hk6s} [\emph{ʰkəs}] 0x8D10 \\
Роза: \textbf{hmap} [\emph{ʰmäp}] 0x4908 \\
Розмари: \textbf{ryec} [\emph{rjeʃ}] 0xEB0E \\
розовощекий: \textbf{brup} [\emph{brup}] 0x4AD1 \\
рококо: \textbf{bxok} [\emph{bxok}] 0x2CF1 \\
роли: \textbf{xyas} [\emph{xjäs}] 0x891A \\
ром: \textbf{hqom} [\emph{ʰŋom}] 0x0CC8 \\
роман: \textbf{rwas} [\emph{rwäs}] 0x892E \\
романс: \textbf{hra7n} [\emph{ʰrä˥n}] 0x7170 \\
романтический: \textbf{hram} [\emph{ʰräm}] 0x0970 \\
ромашка: \textbf{kcem} [\emph{kʃem}] 0x0BA2 \\
ромб: \textbf{hr6p} [\emph{ʰrəp}] 0x4D70 \\
роскошный: \textbf{dlok} [\emph{dlok}] 0x2C73 \\
Россия: \textbf{ryus} [\emph{rjus}] 0x8A0E \\
ростовщик: \textbf{hzu7t} [\emph{ʰzu˥t}] 0xB2A0 \\
рот: \textbf{hmuf} [\emph{ʰmuf}] 0xCA08 \\
ротонда: \textbf{tco2t} [\emph{tʃo˩t}] 0xBCAA \\
Руанда: \textbf{vwat} [\emph{vwät}] 0xA936 \\
Рубашка: \textbf{cras} [\emph{ʃräs}] 0x89CD \\
рубец: \textbf{kxi7s} [\emph{kxi˥s}] 0x90E2 \\
рубидий: \textbf{bwum} [\emph{bwum}] 0x0A31 \\
рубль: \textbf{hru2p} [\emph{ʰru˩p}] 0x5A70 \\
ругань: \textbf{txon} [\emph{txon}] 0x6CEA \\
рудимент: \textbf{qrim} [\emph{ŋrim}] 0x08D9 \\
руды: \textbf{hwi2s} [\emph{ʰwi˩s}] 0x9828 \\
рука: \textbf{hsun} [\emph{ʰsun}] 0x6A48 \\
рукопись: \textbf{srak} [\emph{sräk}] 0x29C9 \\
рукопожатие: \textbf{hxa7m} [\emph{ʰxä˥m}] 0x11D0 \\
рулон: \textbf{hrun} [\emph{ʰrun}] 0x6A70 \\
Румыния: \textbf{hy6s} [\emph{ʰjəs}] 0x8D18 \\
румянец: \textbf{bxo7n} [\emph{bxo˥n}] 0x74F1 \\
руна: \textbf{rwu7n} [\emph{rwu˥n}] 0x722E \\
рустовка: \textbf{hso2c} [\emph{ʰso˩ʃ}] 0xFC48 \\
Рут: \textbf{hru7t} [\emph{ʰru˥t}] 0xB270 \\
рутений: \textbf{ryu2n} [\emph{rju˩n}] 0x7A0E \\
рутил: \textbf{hru2k} [\emph{ʰru˩k}] 0x3A70 \\
ручеек: \textbf{xwic} [\emph{xwiʃ}] 0xE83A \\
рыба: \textbf{kcik} [\emph{kʃik}] 0x28A2 \\
рыба-меч: \textbf{qrif} [\emph{ŋrif}] 0xC8D9 \\
рыбак: \textbf{myif} [\emph{mjif}] 0xC801 \\
рыботорговец: \textbf{clos} [\emph{ʃlos}] 0x8C6D \\
Рыбы: \textbf{hqi2s} [\emph{ʰŋi˩s}] 0x98C8 \\
рыжеватый: \textbf{kxun} [\emph{kxun}] 0x6AE2 \\
рынки: \textbf{hba7s} [\emph{ʰbä˥s}] 0x9188 \\
рысь: \textbf{lyi7s} [\emph{lji˥s}] 0x900B \\
рытвины: \textbf{ryo7t} [\emph{rjo˥t}] 0xB40E \\
рыться: \textbf{nra7c} [\emph{nrä˥ʃ}] 0xF1C6 \\
рыцарский: \textbf{nri7t} [\emph{nri˥t}] 0xB0C6 \\
рычаг: \textbf{lwek} [\emph{lwek}] 0x2B2B \\
рычание: \textbf{gram} [\emph{gräm}] 0x09D2 \\
рюкзак: \textbf{hxap} [\emph{ʰxäp}] 0x49D0 \\
ряды: \textbf{hra2n} [\emph{ʰrä˩n}] 0x7970 \\
\subsection{ с }
саботаж: \textbf{bwoc} [\emph{bwoʃ}] 0xEC31 \\
саванна: \textbf{dven} [\emph{dven}] 0x6B93 \\
сагиттальный: \textbf{swa2s} [\emph{swä˩s}] 0x9929 \\
саго: \textbf{hsu7k} [\emph{ʰsu˥k}] 0x3248 \\
сад: \textbf{fyan} [\emph{fjän}] 0x690C \\
садизм: \textbf{sri2m} [\emph{sri˩m}] 0x18C9 \\
садоводство: \textbf{bvon} [\emph{bvon}] 0x6C91 \\
сажа: \textbf{myi2n} [\emph{mji˩n}] 0x7801 \\
саженцы: \textbf{myup} [\emph{mjup}] 0x4A01 \\
саксонский: \textbf{hsa7s} [\emph{ʰsä˥s}] 0x9148 \\
саксофон: \textbf{ksof} [\emph{ksof}] 0xCC42 \\
салат: \textbf{slut} [\emph{slut}] 0xAA69 \\
\subsection{ С }
Сальвадорских: \textbf{swo2t} [\emph{swo˩t}] 0xBC29 \\
сальный: \textbf{hgi2n} [\emph{ʰgi˩n}] 0x7890 \\
салют: \textbf{fwik} [\emph{fwik}] 0x282C \\
сам: \textbf{hni7c} [\emph{ʰni˥ʃ}] 0xF030 \\
саман: \textbf{hda7p} [\emph{ʰdä˥p}] 0x5198 \\
Самария: \textbf{sra7m} [\emph{srä˥m}] 0x11C9 \\
самбо: \textbf{sla7p} [\emph{slä˥p}] 0x5169 \\
саммит: \textbf{hsem} [\emph{ʰsem}] 0x0B48 \\
Самоанский: \textbf{hso7f} [\emph{ʰso˥f}] 0xD448 \\
самодельный: \textbf{hqi2t} [\emph{ʰŋi˩t}] 0xB8C8 \\
самозванец: \textbf{mwip} [\emph{mwip}] 0x4821 \\
самоиндукции: \textbf{ht67k} [\emph{ʰtə˥k}] 0x3550 \\
самокат: \textbf{swus} [\emph{swus}] 0x8A29 \\
самонадеянным: \textbf{sli2n} [\emph{sli˩n}] 0x7869 \\
самооборона: \textbf{twa2n} [\emph{twä˩n}] 0x792A \\
самоочевидной: \textbf{psi7t} [\emph{psi˥t}] 0xB044 \\
самописец: \textbf{kyi2t} [\emph{kji˩t}] 0xB802 \\
самопоглощение: \textbf{tsi2c} [\emph{tsi˩ʃ}] 0xF84A \\
самореференция: \textbf{tsa7f} [\emph{tsä˥f}] 0xD14A \\
самосвал: \textbf{tli2p} [\emph{tli˩p}] 0x586A \\
самоубийство: \textbf{tsaf} [\emph{tsäf}] 0xC94A \\
самурай: \textbf{sru7t} [\emph{sru˥t}] 0xB2C9 \\
санаторий: \textbf{sya7m} [\emph{sjä˥m}] 0x1109 \\
санитария: \textbf{sw6n} [\emph{swən}] 0x6D29 \\
санскрит: \textbf{nwa2t} [\emph{nwä˩t}] 0xB926 \\
сантиграмм: \textbf{nye2k} [\emph{nje˩k}] 0x3B06 \\
сантиметр: \textbf{nyem} [\emph{njem}] 0x0B06 \\
сапожник: \textbf{hka2p} [\emph{ʰkä˩p}] 0x5910 \\
сапфир: \textbf{pfaf} [\emph{pfäf}] 0xC984 \\
сарай: \textbf{cren} [\emph{ʃren}] 0x6BCD \\
Сардиния: \textbf{zrun} [\emph{zrun}] 0x6AD4 \\
сатанизм: \textbf{sy6m} [\emph{sjəm}] 0x0D09 \\
сатин: \textbf{sya2n} [\emph{sjä˩n}] 0x7909 \\
сатир: \textbf{slep} [\emph{slep}] 0x4B69 \\
сатира: \textbf{hsi7c} [\emph{ʰsi˥ʃ}] 0xF048 \\
Сатурн: \textbf{xyun} [\emph{xjun}] 0x6A1A \\
саудовская: \textbf{sw6t} [\emph{swət}] 0xAD29 \\
сафлоровое: \textbf{xram} [\emph{xräm}] 0x09DA \\
сахар: \textbf{clak} [\emph{ʃläk}] 0x296D \\
сахарин: \textbf{kri7c} [\emph{kri˥ʃ}] 0xF0C2 \\
сбалансированный: \textbf{lwin} [\emph{lwin}] 0x682B \\
сбивчиво: \textbf{xwi7t} [\emph{xwi˥t}] 0xB03A \\
сболтнуть: \textbf{bre7t} [\emph{bre˥t}] 0xB3D1 \\
Свазиленд: \textbf{sy6t} [\emph{sjət}] 0xAD09 \\
свастика: \textbf{hva2k} [\emph{ʰvä˩k}] 0x39B0 \\
свежий: \textbf{syac} [\emph{sjäʃ}] 0xE909 \\
свекла: \textbf{bvut} [\emph{bvut}] 0xAA91 \\
сверкать: \textbf{hr62k} [\emph{ʰrə˩k}] 0x3D70 \\
сверкающих: \textbf{hli7n} [\emph{ʰli˥n}] 0x7058 \\
сверстник: \textbf{hxa2f} [\emph{ʰxä˩f}] 0xD9D0 \\
сверхзвуковой: \textbf{sru7k} [\emph{sru˥k}] 0x32C9 \\
сверхчеловек: \textbf{sra2n} [\emph{srä˩n}] 0x79C9 \\
Светляк: \textbf{nyuk} [\emph{njuk}] 0x2A06 \\
светобоязнь: \textbf{fwop} [\emph{fwop}] 0x4C2C \\
светоч: \textbf{lya2n} [\emph{ljä˩n}] 0x790B \\
светский: \textbf{hdek} [\emph{ʰdek}] 0x2B98 \\
свидетель: \textbf{hgit} [\emph{ʰgit}] 0xA890 \\
свинарник: \textbf{pxi2s} [\emph{pxi˩s}] 0x98E4 \\
свинг: \textbf{gwin} [\emph{gwin}] 0x6832 \\
свинка: \textbf{cya2n} [\emph{ʃjä˩n}] 0x790D \\
свинопас: \textbf{sle2n} [\emph{sle˩n}] 0x7B69 \\
свинья: \textbf{hdus} [\emph{ʰdus}] 0x8A98 \\
свисток: \textbf{slim} [\emph{slim}] 0x0869 \\
свитер: \textbf{swes} [\emph{swes}] 0x8B29 \\
свобода: \textbf{syut} [\emph{sjut}] 0xAA09 \\
свободный: \textbf{hwuk} [\emph{ʰwuk}] 0x2A28 \\
СВЧ: \textbf{mrop} [\emph{mrop}] 0x4CC1 \\
связи: \textbf{hwuc} [\emph{ʰwuʃ}] 0xEA28 \\
связка: \textbf{htap} [\emph{ʰtäp}] 0x4950 \\
связь: \textbf{gyon} [\emph{gjon}] 0x6C12 \\
святой: \textbf{swut} [\emph{swut}] 0xAA29 \\
святоша: \textbf{pri7t} [\emph{pri˥t}] 0xB0C4 \\
святые: \textbf{hse2n} [\emph{ʰse˩n}] 0x7B48 \\
священник: \textbf{pyis} [\emph{pjis}] 0x8804 \\
сгибать: \textbf{flen} [\emph{flen}] 0x6B6C \\
сговариваться: \textbf{klu2n} [\emph{klu˩n}] 0x7A62 \\
сгущенный: \textbf{hki7t} [\emph{ʰki˥t}] 0xB010 \\
сдачи: \textbf{pyap} [\emph{pjäp}] 0x4904 \\
сделанный: \textbf{tcun} [\emph{tʃun}] 0x6AAA \\
сделка: \textbf{myap} [\emph{mjäp}] 0x4901 \\
сдерживаемый: \textbf{pye7t} [\emph{pje˥t}] 0xB304 \\
сдоба: \textbf{pxe2n} [\emph{pxe˩n}] 0x7BE4 \\
север: \textbf{pyut} [\emph{pjut}] 0xAA04 \\
северный: \textbf{bra2n} [\emph{brä˩n}] 0x79D1 \\
северо-восточный: \textbf{zrop} [\emph{zrop}] 0x4CD4 \\
северо-Запад: \textbf{nrep} [\emph{nrep}] 0x4BC6 \\
севилья: \textbf{hsi2f} [\emph{ʰsi˩f}] 0xD848 \\
сёгунат: \textbf{hcu7f} [\emph{ʰʃu˥f}] 0xD268 \\
седан: \textbf{syes} [\emph{sjes}] 0x8B09 \\
седиль: \textbf{syi2f} [\emph{sji˩f}] 0xD809 \\
сейсмограф: \textbf{hz6k} [\emph{ʰzək}] 0x2DA0 \\
сейсмология: \textbf{zlip} [\emph{zlip}] 0x4874 \\
секрет: \textbf{hkim} [\emph{ʰkim}] 0x0810 \\
секретарь: \textbf{kris} [\emph{kris}] 0x88C2 \\
секреция: \textbf{krif} [\emph{krif}] 0xC8C2 \\
сексизм: \textbf{kyi2s} [\emph{kji˩s}] 0x9802 \\
секстан: \textbf{kse2n} [\emph{kse˩n}] 0x7B42 \\
секстет: \textbf{kye2t} [\emph{kje˩t}] 0xBB02 \\
сексуальное: \textbf{hjon} [\emph{ʰʒon}] 0x6CA8 \\
сексуальный: \textbf{hqes} [\emph{ʰŋes}] 0x8BC8 \\
секта: \textbf{dzak} [\emph{dzäk}] 0x2953 \\
секундомер: \textbf{swop} [\emph{swop}] 0x4C29 \\
селезенка: \textbf{pxic} [\emph{pxiʃ}] 0xE8E4 \\
селен: \textbf{clem} [\emph{ʃlem}] 0x0B6D \\
селитра: \textbf{sla2p} [\emph{slä˩p}] 0x5969 \\
сельскохозяйственное: \textbf{ryic} [\emph{rjiʃ}] 0xE80E \\
семантика: \textbf{hm67k} [\emph{ʰmə˥k}] 0x3508 \\
семафор: \textbf{hsa2f} [\emph{ʰsä˩f}] 0xD948 \\
семена: \textbf{srik} [\emph{srik}] 0x28C9 \\
семидесятый: \textbf{tci2s} [\emph{tʃi˩s}] 0x98AA \\
семинар: \textbf{hxan} [\emph{ʰxän}] 0x69D0 \\
семиотический: \textbf{syi2k} [\emph{sji˩k}] 0x3809 \\
семиугольник: \textbf{xye7n} [\emph{xje˥n}] 0x731A \\
семнадцать: \textbf{sret} [\emph{sret}] 0xABC9 \\
семь: \textbf{hsip} [\emph{ʰsip}] 0x4848 \\
семья: \textbf{hfac} [\emph{ʰfäʃ}] 0xE960 \\
семядоля: \textbf{tlo2t} [\emph{tlo˩t}] 0xBC6A \\
Сен-Пьер: \textbf{sye7p} [\emph{sje˥p}] 0x5309 \\
сенат: \textbf{swen} [\emph{swen}] 0x6B29 \\
сеновал: \textbf{xlef} [\emph{xlef}] 0xCB7A \\
сенсорный: \textbf{nrec} [\emph{nreʃ}] 0xEBC6 \\
сентябрь: \textbf{hsep} [\emph{ʰsep}] 0x4B48 \\
сепия: \textbf{hse7p} [\emph{ʰse˥p}] 0x5348 \\
сепсис: \textbf{pce2s} [\emph{pʃe˩s}] 0x9BA4 \\
сера: \textbf{hlof} [\emph{ʰlof}] 0xCC58 \\
сербия: \textbf{sra7p} [\emph{srä˥p}] 0x51C9 \\
сердечно: \textbf{hko2t} [\emph{ʰko˩t}] 0xBC10 \\
сердце: \textbf{cyit} [\emph{ʃjit}] 0xA80D \\
сердцевидный: \textbf{ryi2k} [\emph{rji˩k}] 0x380E \\
сердцеед: \textbf{hwa7m} [\emph{ʰwä˥m}] 0x1128 \\
серебристо-серый: \textbf{gyes} [\emph{gjes}] 0x8B12 \\
Серебряный: \textbf{slin} [\emph{slin}] 0x6869 \\
серенада: \textbf{rye7t} [\emph{rje˥t}] 0xB30E \\
сержант: \textbf{hca2n} [\emph{ʰʃä˩n}] 0x7968 \\
серии: \textbf{slis} [\emph{slis}] 0x8869 \\
сертификат: \textbf{tfot} [\emph{tfot}] 0xAC8A \\
серый: \textbf{kcis} [\emph{kʃis}] 0x88A2 \\
сессия: \textbf{cyen} [\emph{ʃjen}] 0x6B0D \\
сетчатка: \textbf{vwan} [\emph{vwän}] 0x6936 \\
сетчатый: \textbf{tle2t} [\emph{tle˩t}] 0xBB6A \\
Сеульский: \textbf{syu2n} [\emph{sju˩n}] 0x7A09 \\
сжатие: \textbf{mres} [\emph{mres}] 0x8BC1 \\
сжигание: \textbf{ksun} [\emph{ksun}] 0x6A42 \\
сибилянт: \textbf{sl62n} [\emph{slə˩n}] 0x7D69 \\
Сибирь: \textbf{zrip} [\emph{zrip}] 0x48D4 \\
Сивилла: \textbf{hbi7p} [\emph{ʰbi˥p}] 0x5088 \\
сигарета: \textbf{srip} [\emph{srip}] 0x48C9 \\
Сидон: \textbf{syi2n} [\emph{sji˩n}] 0x7809 \\
сидячий: \textbf{dxip} [\emph{dxip}] 0x48F3 \\
сиккатив: \textbf{dj6t} [\emph{dʒət}] 0xADB3 \\
сикх: \textbf{syi7k} [\emph{sji˥k}] 0x3009 \\
силикат: \textbf{sli2k} [\emph{sli˩k}] 0x3869 \\
сильно: \textbf{xla7t} [\emph{xlä˥t}] 0xB17A \\
сильфон: \textbf{fwun} [\emph{fwun}] 0x6A2C \\
симбиоз: \textbf{mwos} [\emph{mwos}] 0x8C21 \\
симметричный: \textbf{my6k} [\emph{mjək}] 0x2D01 \\
симония: \textbf{myo7n} [\emph{mjo˥n}] 0x7401 \\
симпатия: \textbf{tcip} [\emph{tʃip}] 0x48AA \\
симулировать: \textbf{glas} [\emph{gläs}] 0x8972 \\
синагога: \textbf{zyak} [\emph{zjäk}] 0x2914 \\
синапс: \textbf{sr6p} [\emph{srəp}] 0x4DC9 \\
сингалец: \textbf{sli2c} [\emph{sli˩ʃ}] 0xF869 \\
Сингапур: \textbf{hqop} [\emph{ʰŋop}] 0x4CC8 \\
синий: \textbf{plun} [\emph{plun}] 0x6A64 \\
синица: \textbf{twa2s} [\emph{twä˩s}] 0x992A \\
синкопированная: \textbf{hn62t} [\emph{ʰnə˩t}] 0xBD30 \\
синкретизм: \textbf{sri7m} [\emph{sri˥m}] 0x10C9 \\
синонимичный: \textbf{mr6t} [\emph{mrət}] 0xADC1 \\
синоптический: \textbf{nwi7k} [\emph{nwi˥k}] 0x3026 \\
синтаксический: \textbf{nyaf} [\emph{njäf}] 0xC906 \\
синтезировать: \textbf{hs6f} [\emph{ʰsəf}] 0xCD48 \\
синусит: \textbf{nyu2t} [\emph{nju˩t}] 0xBA06 \\
синхронизировать: \textbf{hxin} [\emph{ʰxin}] 0x68D0 \\
сионизм: \textbf{hyi7m} [\emph{ʰji˥m}] 0x1018 \\
Сирия: \textbf{zras} [\emph{zräs}] 0x89D4 \\
сироп: \textbf{xrip} [\emph{xrip}] 0x48DA \\
сирота: \textbf{hya2n} [\emph{ʰjä˩n}] 0x7918 \\
систематическое: \textbf{djit} [\emph{dʒit}] 0xA8B3 \\
системы: \textbf{swim} [\emph{swim}] 0x0829 \\
систолический: \textbf{slo7k} [\emph{slo˥k}] 0x3469 \\
сито: \textbf{slif} [\emph{slif}] 0xC869 \\
сифилис: \textbf{pli2t} [\emph{pli˩t}] 0xB864 \\
сифон: \textbf{syo2n} [\emph{sjo˩n}] 0x7C09 \\
Сицилия: \textbf{zlis} [\emph{zlis}] 0x8874 \\
сказать: \textbf{hsac} [\emph{ʰsäʃ}] 0xE948 \\
сказывают: \textbf{hca7n} [\emph{ʰʃä˥n}] 0x7168 \\
скакать: \textbf{kfo2t} [\emph{kfo˩t}] 0xBC82 \\
скалистый: \textbf{kra2k} [\emph{krä˩k}] 0x39C2 \\
скалы: \textbf{kla7s} [\emph{klä˥s}] 0x9162 \\
скальпель: \textbf{kxup} [\emph{kxup}] 0x4AE2 \\
скаляр: \textbf{kl6m} [\emph{kləm}] 0x0D62 \\
скандий: \textbf{kxat} [\emph{kxät}] 0xA9E2 \\
скандинавский: \textbf{hka7f} [\emph{ʰkä˥f}] 0xD110 \\
скарлатина: \textbf{qri2n} [\emph{ŋri˩n}] 0x78D9 \\
скат: \textbf{dvif} [\emph{dvif}] 0xC893 \\
скатерть: \textbf{tlop} [\emph{tlop}] 0x4C6A \\
скачать: \textbf{dlot} [\emph{dlot}] 0xAC73 \\
скваттер: \textbf{kw6t} [\emph{kwət}] 0xAD22 \\
сквош: \textbf{kwoc} [\emph{kwoʃ}] 0xEC22 \\
скейтборд: \textbf{hxep} [\emph{ʰxep}] 0x4BD0 \\
скидка: \textbf{dyut} [\emph{djut}] 0xAA13 \\
скипидар: \textbf{tref} [\emph{tref}] 0xCBCA \\
сколиоз: \textbf{klo2s} [\emph{klo˩s}] 0x9C62 \\
скольжение: \textbf{slip} [\emph{slip}] 0x4869 \\
скользкий: \textbf{lwi7p} [\emph{lwi˥p}] 0x502B \\
скопа: \textbf{syi2p} [\emph{sji˩p}] 0x5809 \\
скорая: \textbf{blus} [\emph{blus}] 0x8A71 \\
скоропортящийся: \textbf{hpi7f} [\emph{ʰpi˥f}] 0xD020 \\
скорость: \textbf{hsot} [\emph{ʰsot}] 0xAC48 \\
скорпион: \textbf{kr6k} [\emph{krək}] 0x2DC2 \\
скотобойня: \textbf{bwa7n} [\emph{bwä˥n}] 0x7131 \\
скотство: \textbf{sla2c} [\emph{slä˩ʃ}] 0xF969 \\
скотч: \textbf{hso2t} [\emph{ʰso˩t}] 0xBC48 \\
скрепок: \textbf{syep} [\emph{sjep}] 0x4B09 \\
скрещивание: \textbf{hro7n} [\emph{ʰro˥n}] 0x7470 \\
Скриншот: \textbf{qrit} [\emph{ŋrit}] 0xA8D9 \\
скрипач: \textbf{xli7t} [\emph{xli˥t}] 0xB07A \\
скрипка: \textbf{dyin} [\emph{djin}] 0x6813 \\
скромный: \textbf{hfok} [\emph{ʰfok}] 0x2C60 \\
скрывать: \textbf{hgan} [\emph{ʰgän}] 0x6990 \\
скула: \textbf{txu2n} [\emph{txu˩n}] 0x7AEA \\
скульптор: \textbf{hquc} [\emph{ʰŋuʃ}] 0xEAC8 \\
скумбрия: \textbf{krum} [\emph{krum}] 0x0AC2 \\
скунс: \textbf{hsu2k} [\emph{ʰsu˩k}] 0x3A48 \\
скупо: \textbf{kci7t} [\emph{kʃi˥t}] 0xB0A2 \\
слабеть: \textbf{xwut} [\emph{xwut}] 0xAA3A \\
слабый: \textbf{hfip} [\emph{ʰfip}] 0x4860 \\
слава: \textbf{nyam} [\emph{njäm}] 0x0906 \\
славословие: \textbf{sli7k} [\emph{sli˥k}] 0x3069 \\
славянский: \textbf{sla7f} [\emph{slä˥f}] 0xD169 \\
след: \textbf{dla7n} [\emph{dlä˥n}] 0x7173 \\
следовать: \textbf{slan} [\emph{slän}] 0x6969 \\
следствие: \textbf{ksik} [\emph{ksik}] 0x2842 \\
следы: \textbf{cret} [\emph{ʃret}] 0xABCD \\
слеза: \textbf{hzis} [\emph{ʰzis}] 0x88A0 \\
слезящимися: \textbf{xra7s} [\emph{xrä˥s}] 0x91DA \\
слепой: \textbf{hbat} [\emph{ʰbät}] 0xA988 \\
слесарь: \textbf{hla2f} [\emph{ʰlä˩f}] 0xD958 \\
слива: \textbf{hles} [\emph{ʰles}] 0x8B58 \\
Словакия: \textbf{vwak} [\emph{vwäk}] 0x2936 \\
Словарь: \textbf{kwon} [\emph{kwon}] 0x6C22 \\
Словения: \textbf{swof} [\emph{swof}] 0xCC29 \\
словесный: \textbf{cwec} [\emph{ʃweʃ}] 0xEB2D \\
слово: \textbf{tlat} [\emph{tlät}] 0xA96A \\
слоги: \textbf{sli7p} [\emph{sli˥p}] 0x5069 \\
сложить: \textbf{syot} [\emph{sjot}] 0xAC09 \\
сложный: \textbf{hkes} [\emph{ʰkes}] 0x8B10 \\
сломанный: \textbf{psuk} [\emph{psuk}] 0x2A44 \\
слоновость: \textbf{hqi7s} [\emph{ʰŋi˥s}] 0x90C8 \\
слоны: \textbf{hli2n} [\emph{ʰli˩n}] 0x7858 \\
служащий: \textbf{swac} [\emph{swäʃ}] 0xE929 \\
слух: \textbf{hwum} [\emph{ʰwum}] 0x0A28 \\
слуха: \textbf{hnu7c} [\emph{ʰnu˥ʃ}] 0xF230 \\
случайный: \textbf{swec} [\emph{sweʃ}] 0xEB29 \\
слушатель: \textbf{nyot} [\emph{njot}] 0xAC06 \\
слышать: \textbf{hcun} [\emph{ʰʃun}] 0x6A68 \\
слышимый: \textbf{hde7n} [\emph{ʰde˥n}] 0x7398 \\
слюда: \textbf{hy6k} [\emph{ʰjək}] 0x2D18 \\
смазывать: \textbf{jwam} [\emph{ʒwäm}] 0x0935 \\
смартфон: \textbf{srop} [\emph{srop}] 0x4CC9 \\
смачивание: \textbf{hja7n} [\emph{ʰʒä˥n}] 0x71A8 \\
смело: \textbf{hba7n} [\emph{ʰbä˥n}] 0x7188 \\
смерть: \textbf{mwan} [\emph{mwän}] 0x6921 \\
смеситель: \textbf{bx6t} [\emph{bxət}] 0xADF1 \\
сметь: \textbf{kfat} [\emph{kfät}] 0xA982 \\
смех: \textbf{cyaf} [\emph{ʃjäf}] 0xC90D \\
смешанный: \textbf{mrat} [\emph{mrät}] 0xA9C1 \\
смещение: \textbf{prap} [\emph{präp}] 0x49C4 \\
смиренный: \textbf{hmen} [\emph{ʰmen}] 0x6B08 \\
смокинг: \textbf{tsi7f} [\emph{tsi˥f}] 0xD04A \\
смола: \textbf{hre2s} [\emph{ʰre˩s}] 0x9B70 \\
смородина: \textbf{kxes} [\emph{kxes}] 0x8BE2 \\
смущаясь: \textbf{dru2t} [\emph{dru˩t}] 0xBAD3 \\
смягчающий: \textbf{jl6n} [\emph{ʒlən}] 0x6D75 \\
снаряд: \textbf{pyok} [\emph{pjok}] 0x2C04 \\
Снеговик: \textbf{hc67n} [\emph{ʰʃə˥n}] 0x7568 \\
снегоступ: \textbf{hsu7c} [\emph{ʰsu˥ʃ}] 0xF248 \\
снимок: \textbf{syoc} [\emph{sjoʃ}] 0xEC09 \\
снисходить: \textbf{hci2n} [\emph{ʰʃi˩n}] 0x7868 \\
сноска: \textbf{fyot} [\emph{fjot}] 0xAC0C \\
сносно: \textbf{tra2n} [\emph{trä˩n}] 0x79CA \\
собака: \textbf{klak} [\emph{kläk}] 0x2962 \\
собеседник: \textbf{tle7t} [\emph{tle˥t}] 0xB36A \\
собирать: \textbf{kcot} [\emph{kʃot}] 0xACA2 \\
соблазнять: \textbf{hya2s} [\emph{ʰjä˩s}] 0x9918 \\
соболезнование: \textbf{bw6n} [\emph{bwən}] 0x6D31 \\
соболь: \textbf{swep} [\emph{swep}] 0x4B29 \\
сова: \textbf{mwut} [\emph{mwut}] 0xAA21 \\
совершеннолетие: \textbf{sl62t} [\emph{slə˩t}] 0xBD69 \\
совесть: \textbf{hran} [\emph{ʰrän}] 0x6970 \\
совокупление: \textbf{kc6c} [\emph{kʃəʃ}] 0xEDA2 \\
совокупляться: \textbf{klum} [\emph{klum}] 0x0A62 \\
совпадение: \textbf{tyoc} [\emph{tjoʃ}] 0xEC0A \\
согласный: \textbf{kxom} [\emph{kxom}] 0x0CE2 \\
соглядатай: \textbf{xri7s} [\emph{xri˥s}] 0x90DA \\
содержание: \textbf{hnok} [\emph{ʰnok}] 0x2C30 \\
содержащий: \textbf{hzan} [\emph{ʰzän}] 0x69A0 \\
содомия: \textbf{hda7m} [\emph{ʰdä˥m}] 0x1198 \\
содрогнулся: \textbf{hso7t} [\emph{ʰso˥t}] 0xB448 \\
содружество: \textbf{cri2t} [\emph{ʃri˩t}] 0xB8CD \\
соединители: \textbf{hn6t} [\emph{ʰnət}] 0xAD30 \\
сожжен: \textbf{byat} [\emph{bjät}] 0xA911 \\
сожительство: \textbf{hki2c} [\emph{ʰki˩ʃ}] 0xF810 \\
созвездий: \textbf{hko2s} [\emph{ʰko˩s}] 0x9C10 \\
созданный: \textbf{hta7t} [\emph{ʰtä˥t}] 0xB150 \\
созревание: \textbf{sr6n} [\emph{srən}] 0x6DC9 \\
Сократ: \textbf{sra7t} [\emph{srä˥t}] 0xB1C9 \\
сокращать: \textbf{bzac} [\emph{bzäʃ}] 0xE951 \\
сокровище: \textbf{crun} [\emph{ʃrun}] 0x6ACD \\
солецизм: \textbf{sli2m} [\emph{sli˩m}] 0x1869 \\
солипсизм: \textbf{lyo7s} [\emph{ljo˥s}] 0x940B \\
солнечник: \textbf{dro7n} [\emph{dro˥n}] 0x74D3 \\
солнцами: \textbf{xya2n} [\emph{xjä˩n}] 0x791A \\
солнцестояние: \textbf{qyot} [\emph{ŋjot}] 0xAC19 \\
соловей: \textbf{hne2n} [\emph{ʰne˩n}] 0x7B30 \\
солод: \textbf{mya2t} [\emph{mjä˩t}] 0xB901 \\
Соломон: \textbf{slo7n} [\emph{slo˥n}] 0x7469 \\
Сомали: \textbf{swom} [\emph{swom}] 0x0C29 \\
сомнение: \textbf{hcaf} [\emph{ʰʃäf}] 0xC968 \\
сонар: \textbf{zwon} [\emph{zwon}] 0x6C34 \\
соната: \textbf{swa7t} [\emph{swä˥t}] 0xB129 \\
сонный: \textbf{nwus} [\emph{nwus}] 0x8A26 \\
сообщества: \textbf{hmo7t} [\emph{ʰmo˥t}] 0xB408 \\
соответственно: \textbf{rwim} [\emph{rwim}] 0x082E \\
соперник: \textbf{rwit} [\emph{rwit}] 0xA82E \\
сопли: \textbf{hqo2t} [\emph{ʰŋo˩t}] 0xBCC8 \\
сопровождая: \textbf{hka2n} [\emph{ʰkä˩n}] 0x7910 \\
сопротивление: \textbf{trac} [\emph{träʃ}] 0xE9CA \\
сопряженный: \textbf{djok} [\emph{dʒok}] 0x2CB3 \\
сорняками: \textbf{gvas} [\emph{gväs}] 0x8992 \\
сородичи: \textbf{hki7s} [\emph{ʰki˥s}] 0x9010 \\
сороки: \textbf{hgi7p} [\emph{ʰgi˥p}] 0x5090 \\
сороконожка: \textbf{kyep} [\emph{kjep}] 0x4B02 \\
сорокопут: \textbf{cra2k} [\emph{ʃrä˩k}] 0x39CD \\
Сортировать: \textbf{nyit} [\emph{njit}] 0xA806 \\
сосиски: \textbf{hf6f} [\emph{ʰfəf}] 0xCD60 \\
сосна: \textbf{hpos} [\emph{ʰpos}] 0x8C20 \\
сосок: \textbf{txu7t} [\emph{txu˥t}] 0xB2EA \\
состояние: \textbf{cyun} [\emph{ʃjun}] 0x6A0D \\
состоять: \textbf{hfot} [\emph{ʰfot}] 0xAC60 \\
сосулька: \textbf{xluk} [\emph{xluk}] 0x2A7A \\
соте: \textbf{hso7c} [\emph{ʰso˥ʃ}] 0xF448 \\
сотни: \textbf{syop} [\emph{sjop}] 0x4C09 \\
сотрудники: \textbf{hrak} [\emph{ʰräk}] 0x2970 \\
сотрудничать: \textbf{crom} [\emph{ʃrom}] 0x0CCD \\
софистика: \textbf{ksi7s} [\emph{ksi˥s}] 0x9042 \\
София: \textbf{syi7f} [\emph{sji˥f}] 0xD009 \\
Софония: \textbf{fy6n} [\emph{fjən}] 0x6D0C \\
сохранившийся: \textbf{ksi2n} [\emph{ksi˩n}] 0x7842 \\
сохраняет: \textbf{pra7t} [\emph{prä˥t}] 0xB1C4 \\
социально-экономический: \textbf{syi2c} [\emph{sji˩ʃ}] 0xF809 \\
Социальное: \textbf{swic} [\emph{swiʃ}] 0xE829 \\
социология: \textbf{slok} [\emph{slok}] 0x2C69 \\
сочинение: \textbf{nyas} [\emph{njäs}] 0x8906 \\
сочиться: \textbf{hdu7s} [\emph{ʰdu˥s}] 0x9298 \\
сочлененный: \textbf{vlut} [\emph{vlut}] 0xAA76 \\
союз: \textbf{hnon} [\emph{ʰnon}] 0x6C30 \\
союзы: \textbf{djos} [\emph{dʒos}] 0x8CB3 \\
соя: \textbf{syo2t} [\emph{sjo˩t}] 0xBC09 \\
спа: \textbf{gjap} [\emph{gʒäp}] 0x49B2 \\
Спальня: \textbf{hwos} [\emph{ʰwos}] 0x8C28 \\
спам: \textbf{syem} [\emph{sjem}] 0x0B09 \\
спамер: \textbf{pfam} [\emph{pfäm}] 0x0984 \\
спаржа: \textbf{sra7s} [\emph{srä˥s}] 0x91C9 \\
спартанский: \textbf{sra2t} [\emph{srä˩t}] 0xB9C9 \\
Спасатель: \textbf{hqit} [\emph{ʰŋit}] 0xA8C8 \\
спасающий: \textbf{hxuc} [\emph{ʰxuʃ}] 0xEAD0 \\
спасибо: \textbf{kyaf} [\emph{kjäf}] 0xC902 \\
спать: \textbf{hsim} [\emph{ʰsim}] 0x0848 \\
спектр: \textbf{hwem} [\emph{ʰwem}] 0x0B28 \\
спектроскоп: \textbf{kwep} [\emph{kwep}] 0x4B22 \\
спекулировать: \textbf{klut} [\emph{klut}] 0xAA62 \\
спели: \textbf{zyam} [\emph{zjäm}] 0x0914 \\
сперма: \textbf{hka2m} [\emph{ʰkä˩m}] 0x1910 \\
сперматозоиды: \textbf{cra7s} [\emph{ʃrä˥s}] 0x91CD \\
специи: \textbf{clim} [\emph{ʃlim}] 0x086D \\
Спецификация: \textbf{pfif} [\emph{pfif}] 0xC884 \\
СПИД: \textbf{dzit} [\emph{dzit}] 0xA853 \\
спидометр: \textbf{hvom} [\emph{ʰvom}] 0x0CB0 \\
спинодержатель: \textbf{bjap} [\emph{bʒäp}] 0x49B1 \\
спираль: \textbf{clif} [\emph{ʃlif}] 0xC86D \\
спиритизм: \textbf{swum} [\emph{swum}] 0x0A29 \\
список: \textbf{lwat} [\emph{lwät}] 0xA92B \\
сплюснутый: \textbf{blo7t} [\emph{blo˥t}] 0xB471 \\
спонсируемый: \textbf{pso2t} [\emph{pso˩t}] 0xBC44 \\
спорный: \textbf{trus} [\emph{trus}] 0x8ACA \\
Спортивное: \textbf{tyet} [\emph{tjet}] 0xAB0A \\
спорщик: \textbf{pre7t} [\emph{pre˥t}] 0xB3C4 \\
способности: \textbf{kla7t} [\emph{klä˥t}] 0xB162 \\
спотыкаться: \textbf{pcok} [\emph{pʃok}] 0x2CA4 \\
справедливость: \textbf{tsis} [\emph{tsis}] 0x884A \\
справочный: \textbf{hrec} [\emph{ʰreʃ}] 0xEB70 \\
спринтер: \textbf{pc6t} [\emph{pʃət}] 0xADA4 \\
спросил: \textbf{hpun} [\emph{ʰpun}] 0x6A20 \\
спутанный: \textbf{txa7t} [\emph{txä˥t}] 0xB1EA \\
спутник: \textbf{jlic} [\emph{ʒliʃ}] 0xE875 \\
срабатывает: \textbf{txak} [\emph{txäk}] 0x29EA \\
сравнение: \textbf{tras} [\emph{träs}] 0x89CA \\
сращивание: \textbf{plec} [\emph{pleʃ}] 0xEB64 \\
среда: \textbf{brac} [\emph{bräʃ}] 0xE9D1 \\
Средиземье: \textbf{tl67t} [\emph{tlə˥t}] 0xB56A \\
средневековый: \textbf{dyis} [\emph{djis}] 0x8813 \\
средний: \textbf{hmim} [\emph{ʰmim}] 0x0808 \\
сродство: \textbf{fyon} [\emph{fjon}] 0x6C0C \\
сроки: \textbf{syas} [\emph{sjäs}] 0x8909 \\
сросшийся: \textbf{kr62t} [\emph{krə˩t}] 0xBDC2 \\
срывать: \textbf{dwak} [\emph{dwäk}] 0x2933 \\
ссадить: \textbf{xro2t} [\emph{xro˩t}] 0xBCDA \\
ссылаясь: \textbf{txif} [\emph{txif}] 0xC8EA \\
стабильность: \textbf{glit} [\emph{glit}] 0xA872 \\
ставки: \textbf{xlif} [\emph{xlif}] 0xC87A \\
стагфляция: \textbf{tfa2c} [\emph{tfä˩ʃ}] 0xF98A \\
стадо: \textbf{cwen} [\emph{ʃwen}] 0x6B2D \\
стакан: \textbf{klas} [\emph{kläs}] 0x8962 \\
сталкиваться: \textbf{hx6s} [\emph{ʰxəs}] 0x8DD0 \\
сталкиваются: \textbf{mri2t} [\emph{mri˩t}] 0xB8C1 \\
становление: \textbf{hbin} [\emph{ʰbin}] 0x6888 \\
станций: \textbf{tsa2n} [\emph{tsä˩n}] 0x794A \\
старейшина: \textbf{dra2n} [\emph{drä˩n}] 0x79D3 \\
старейшины: \textbf{sla7n} [\emph{slä˥n}] 0x7169 \\
стареющий: \textbf{sya2s} [\emph{sjä˩s}] 0x9909 \\
старикашка: \textbf{gzes} [\emph{gzes}] 0x8B52 \\
старость: \textbf{hluf} [\emph{ʰluf}] 0xCA58 \\
стартапы: \textbf{hsa7p} [\emph{ʰsä˥p}] 0x5148 \\
стартер: \textbf{tr6k} [\emph{trək}] 0x2DCA \\
старый: \textbf{tyat} [\emph{tjät}] 0xA90A \\
статистическая: \textbf{twok} [\emph{twok}] 0x2C2A \\
статуя: \textbf{trut} [\emph{trut}] 0xAACA \\
стафилококк: \textbf{txo2k} [\emph{txo˩k}] 0x3CEA \\
стебель: \textbf{ht6c} [\emph{ʰtəʃ}] 0xED50 \\
стежка: \textbf{gl6n} [\emph{glən}] 0x6D72 \\
стейк: \textbf{tsek} [\emph{tsek}] 0x2B4A \\
стеклярус: \textbf{twi7k} [\emph{twi˥k}] 0x302A \\
стекольщик: \textbf{lyi7k} [\emph{lji˥k}] 0x300B \\
стена: \textbf{plat} [\emph{plät}] 0xA964 \\
стенографистка: \textbf{sref} [\emph{sref}] 0xCBC9 \\
стенография: \textbf{sri7f} [\emph{sri˥f}] 0xD0C9 \\
стенокардия: \textbf{qya7n} [\emph{ŋjä˥n}] 0x7119 \\
стеньга: \textbf{tsa2m} [\emph{tsä˩m}] 0x194A \\
степень: \textbf{tlut} [\emph{tlut}] 0xAA6A \\
стерео: \textbf{tr6n} [\emph{trən}] 0x6DCA \\
стереоскоп: \textbf{tro7k} [\emph{tro˥k}] 0x34CA \\
стержень: \textbf{kxam} [\emph{kxäm}] 0x09E2 \\
стерилизовать: \textbf{vlak} [\emph{vläk}] 0x2976 \\
стетоскоп: \textbf{tx6k} [\emph{txək}] 0x2DEA \\
стимул: \textbf{tsif} [\emph{tsif}] 0xC84A \\
стипендия: \textbf{nrim} [\emph{nrim}] 0x08C6 \\
стирание: \textbf{xras} [\emph{xräs}] 0x89DA \\
стискивать: \textbf{hya7c} [\emph{ʰjä˥ʃ}] 0xF118 \\
стог: \textbf{kxi7k} [\emph{kxi˥k}] 0x30E2 \\
стоимость: \textbf{tlas} [\emph{tläs}] 0x896A \\
стоицизм: \textbf{hso7m} [\emph{ʰso˥m}] 0x1448 \\
сток: \textbf{rwof} [\emph{rwof}] 0xCC2E \\
Стокгольм: \textbf{hto7m} [\emph{ʰto˥m}] 0x1450 \\
стократный: \textbf{hsu7p} [\emph{ʰsu˥p}] 0x5248 \\
столб: \textbf{pcoc} [\emph{pʃoʃ}] 0xECA4 \\
столбняк: \textbf{pso7n} [\emph{pso˥n}] 0x7444 \\
столбчатый: \textbf{slo2n} [\emph{slo˩n}] 0x7C69 \\
столетний: \textbf{sle7n} [\emph{sle˥n}] 0x7369 \\
столица: \textbf{klin} [\emph{klin}] 0x6862 \\
стон: \textbf{nyen} [\emph{njen}] 0x6B06 \\
сторонник: \textbf{syok} [\emph{sjok}] 0x2C09 \\
стороны: \textbf{syaf} [\emph{sjäf}] 0xC909 \\
стоянка: \textbf{prek} [\emph{prek}] 0x2BC4 \\
страна: \textbf{klac} [\emph{kläʃ}] 0xE962 \\
страница: \textbf{hpif} [\emph{ʰpif}] 0xC820 \\
странно: \textbf{drik} [\emph{drik}] 0x28D3 \\
странность: \textbf{sri2n} [\emph{sri˩n}] 0x78C9 \\
странствование: \textbf{pro2n} [\emph{pro˩n}] 0x7CC4 \\
стратегия: \textbf{trec} [\emph{treʃ}] 0xEBCA \\
стратификация: \textbf{tf6k} [\emph{tfək}] 0x2D8A \\
стратосфера: \textbf{ts6f} [\emph{tsəf}] 0xCD4A \\
страусы: \textbf{truf} [\emph{truf}] 0xCACA \\
страх: \textbf{hcuc} [\emph{ʰʃuʃ}] 0xEA68 \\
страшилище: \textbf{byuk} [\emph{bjuk}] 0x2A11 \\
стрекоза: \textbf{dxif} [\emph{dxif}] 0xC8F3 \\
стрелять: \textbf{slot} [\emph{slot}] 0xAC69 \\
стресс: \textbf{tres} [\emph{tres}] 0x8BCA \\
стригальщик: \textbf{sre7c} [\emph{sre˥ʃ}] 0xF3C9 \\
стрижка: \textbf{braf} [\emph{bräf}] 0xC9D1 \\
строб: \textbf{sro7p} [\emph{sro˥p}] 0x54C9 \\
стройный: \textbf{hxe2n} [\emph{ʰxe˩n}] 0x7BD0 \\
стронций: \textbf{sr6m} [\emph{srəm}] 0x0DC9 \\
строп: \textbf{hga2n} [\emph{ʰgä˩n}] 0x7990 \\
стропило: \textbf{bxaf} [\emph{bxäf}] 0xC9F1 \\
структурный: \textbf{rwut} [\emph{rwut}] 0xAA2E \\
стручок: \textbf{plap} [\emph{pläp}] 0x4964 \\
студент: \textbf{hces} [\emph{ʰʃes}] 0x8B68 \\
студия: \textbf{hjot} [\emph{ʰʒot}] 0xACA8 \\
стул: \textbf{cyis} [\emph{ʃjis}] 0x880D \\
ступенька: \textbf{tloc} [\emph{tloʃ}] 0xEC6A \\
стучать: \textbf{truk} [\emph{truk}] 0x2ACA \\
стягивание: \textbf{gvon} [\emph{gvon}] 0x6C92 \\
суб-лейтенант: \textbf{slu7n} [\emph{slu˥n}] 0x7269 \\
субарктический: \textbf{bja7t} [\emph{bʒä˥t}] 0xB1B1 \\
суббота: \textbf{crat} [\emph{ʃrät}] 0xA9CD \\
субдукции: \textbf{bjok} [\emph{bʒok}] 0x2CB1 \\
сублимировать: \textbf{slem} [\emph{slem}] 0x0B69 \\
субсидировать: \textbf{syus} [\emph{sjus}] 0x8A09 \\
суглинок: \textbf{jlut} [\emph{ʒlut}] 0xAA75 \\
суд: \textbf{fyat} [\emph{fjät}] 0xA90C \\
Судан: \textbf{hzus} [\emph{ʰzus}] 0x8AA0 \\
судорога: \textbf{kxuc} [\emph{kxuʃ}] 0xEAE2 \\
судья: \textbf{fwak} [\emph{fwäk}] 0x292C \\
суеверие: \textbf{hmi2n} [\emph{ʰmi˩n}] 0x7808 \\
сука: \textbf{byut} [\emph{bjut}] 0xAA11 \\
Суккот: \textbf{swu7k} [\emph{swu˥k}] 0x3229 \\
султанат: \textbf{sl6k} [\emph{slək}] 0x2D69 \\
сульфид: \textbf{lwuf} [\emph{lwuf}] 0xCA2B \\
суматоха: \textbf{cyos} [\emph{ʃjos}] 0x8C0D \\
Суматра: \textbf{hm67t} [\emph{ʰmə˥t}] 0xB508 \\
сумах: \textbf{hsu2m} [\emph{ʰsu˩m}] 0x1A48 \\
сумерки: \textbf{gxuk} [\emph{gxuk}] 0x2AF2 \\
сумчатый: \textbf{hru7p} [\emph{ʰru˥p}] 0x5270 \\
суннит: \textbf{hzi2n} [\emph{ʰzi˩n}] 0x78A0 \\
суп: \textbf{srup} [\emph{srup}] 0x4AC9 \\
супер: \textbf{syup} [\emph{sjup}] 0x4A09 \\
супер: \textbf{syuh} [\emph{sjuʰ}] 0x504F \\
супердержавы: \textbf{pca2k} [\emph{pʃä˩k}] 0x39A4 \\
супермаркет: \textbf{pruk} [\emph{pruk}] 0x2AC4 \\
супница: \textbf{twi2n} [\emph{twi˩n}] 0x782A \\
суппозиторий: \textbf{pf6t} [\emph{pfət}] 0xAD84 \\
Суринам: \textbf{zrim} [\emph{zrim}] 0x08D4 \\
сурок: \textbf{hwo7k} [\emph{ʰwo˥k}] 0x3428 \\
сурьма: \textbf{tyom} [\emph{tjom}] 0x0C0A \\
сутенер: \textbf{pxup} [\emph{pxup}] 0x4AE4 \\
суточный: \textbf{dxun} [\emph{dxun}] 0x6AF3 \\
суфлер: \textbf{dxut} [\emph{dxut}] 0xAAF3 \\
суффикс: \textbf{hxi2s} [\emph{ʰxi˩s}] 0x98D0 \\
сухой: \textbf{kyuc} [\emph{kjuʃ}] 0xEA02 \\
сфагнум: \textbf{swa2k} [\emph{swä˩k}] 0x3929 \\
сфинкс: \textbf{pxi2n} [\emph{pxi˩n}] 0x78E4 \\
сфинктер: \textbf{fy6t} [\emph{fjət}] 0xAD0C \\
схема: \textbf{rwe2t} [\emph{rwe˩t}] 0xBB2E \\
сходимость: \textbf{blon} [\emph{blon}] 0x6C71 \\
схоластика: \textbf{sya7s} [\emph{sjä˥s}] 0x9109 \\
Счет: \textbf{hkun} [\emph{ʰkun}] 0x6A10 \\
счетчики: \textbf{ksu2t} [\emph{ksu˩t}] 0xBA42 \\
считают,: \textbf{kcan} [\emph{kʃän}] 0x69A2 \\
съеденный: \textbf{hvuk} [\emph{ʰvuk}] 0x2AB0 \\
сыворотка: \textbf{gxa7n} [\emph{gxä˥n}] 0x71F2 \\
сын: \textbf{hpas} [\emph{ʰpäs}] 0x8920 \\
сыпь: \textbf{pl6t} [\emph{plət}] 0xAD64 \\
сыр: \textbf{hgis} [\emph{ʰgis}] 0x8890 \\
сырой: \textbf{syo7n} [\emph{sjo˥n}] 0x7409 \\
сычужок: \textbf{hqe7t} [\emph{ʰŋe˥t}] 0xB3C8 \\
Сью: \textbf{tcus} [\emph{tʃus}] 0x8AAA \\
сэр: \textbf{syen} [\emph{sjen}] 0x6B09 \\
\subsection{ т }
таблетка: \textbf{blet} [\emph{blet}] 0xAB71 \\
таблица: \textbf{pyec} [\emph{pjeʃ}] 0xEB04 \\
\subsection{ Т }
Таблица: \textbf{tyas} [\emph{tjäs}] 0x890A \\
тавтологический: \textbf{tlu7k} [\emph{tlu˥k}] 0x326A \\
Таджикистан: \textbf{jyas} [\emph{ʒjäs}] 0x8915 \\
таз: \textbf{pxes} [\emph{pxes}] 0x8BE4 \\
Таиланд: \textbf{xl6t} [\emph{xlət}] 0xAD7A \\
таитянского: \textbf{txa7c} [\emph{txä˥ʃ}] 0xF1EA \\
Тайвань: \textbf{tw6n} [\emph{twən}] 0x6D2A \\
тайпей: \textbf{tyi7p} [\emph{tji˥p}] 0x500A \\
также: \textbf{hyac} [\emph{ʰjäʃ}] 0xE918 \\
такса: \textbf{dja2n} [\emph{dʒä˩n}] 0x79B3 \\
такси: \textbf{tsus} [\emph{tsus}] 0x8A4A \\
таксометр: \textbf{tsi2m} [\emph{tsi˩m}] 0x184A \\
тактика: \textbf{bluk} [\emph{bluk}] 0x2A71 \\
талант: \textbf{tlen} [\emph{tlen}] 0x6B6A \\
талибы: \textbf{tli7p} [\emph{tli˥p}] 0x506A \\
таллий: \textbf{tla2m} [\emph{tlä˩m}] 0x196A \\
Таллин: \textbf{xl67n} [\emph{xlə˥n}] 0x757A \\
талмудический: \textbf{tlu7t} [\emph{tlu˥t}] 0xB26A \\
тальк: \textbf{tla7k} [\emph{tlä˥k}] 0x316A \\
там: \textbf{klan} [\emph{klän}] 0x6962 \\
тамаринд: \textbf{mra7f} [\emph{mrä˥f}] 0xD1C1 \\
танец: \textbf{hdas} [\emph{ʰdäs}] 0x8998 \\
Танзания: \textbf{ts6c} [\emph{tsəʃ}] 0xED4A \\
танин: \textbf{tlo2n} [\emph{tlo˩n}] 0x7C6A \\
тантал: \textbf{txam} [\emph{txäm}] 0x09EA \\
танцор: \textbf{dyak} [\emph{djäk}] 0x2913 \\
таоизм: \textbf{tca2m} [\emph{tʃä˩m}] 0x19AA \\
тапиока: \textbf{tya7k} [\emph{tjä˥k}] 0x310A \\
тапочка: \textbf{slop} [\emph{slop}] 0x4C69 \\
таракан: \textbf{zlac} [\emph{zläʃ}] 0xE974 \\
тарантул: \textbf{ryu2t} [\emph{rju˩t}] 0xBA0E \\
таро: \textbf{tro2p} [\emph{tro˩p}] 0x5CCA \\
Тасмания: \textbf{hz6n} [\emph{ʰzən}] 0x6DA0 \\
тасование: \textbf{tyof} [\emph{tjof}] 0xCC0A \\
тафта: \textbf{tfe7t} [\emph{tfe˥t}] 0xB38A \\
тахикардия: \textbf{twi2k} [\emph{twi˩k}] 0x382A \\
тащили: \textbf{tra7t} [\emph{trä˥t}] 0xB1CA \\
твердо: \textbf{hret} [\emph{ʰret}] 0xAB70 \\
театр: \textbf{tret} [\emph{tret}] 0xABCA \\
тег: \textbf{tyec} [\emph{tjeʃ}] 0xEB0A \\
Тегеран: \textbf{tx62n} [\emph{txə˩n}] 0x7DEA \\
Тедди: \textbf{hxec} [\emph{ʰxeʃ}] 0xEBD0 \\
текст: \textbf{hwus} [\emph{ʰwus}] 0x8A28 \\
текстиль: \textbf{ksi7p} [\emph{ksi˥p}] 0x5042 \\
тектонический: \textbf{kx6k} [\emph{kxək}] 0x2DE2 \\
телевидение: \textbf{twes} [\emph{twes}] 0x8B2A \\
телеграф: \textbf{txaf} [\emph{txäf}] 0xC9EA \\
тележка: \textbf{trom} [\emph{trom}] 0x0CCA \\
тележки: \textbf{tro7t} [\emph{tro˥t}] 0xB4CA \\
телезритель: \textbf{lwo7t} [\emph{lwo˥t}] 0xB42B \\
телекоммуникация: \textbf{txec} [\emph{txeʃ}] 0xEBEA \\
телеологический: \textbf{tlof} [\emph{tlof}] 0xCC6A \\
телепатия: \textbf{tla7p} [\emph{tlä˥p}] 0x516A \\
телескоп: \textbf{hlop} [\emph{ʰlop}] 0x4C58 \\
телесный: \textbf{craf} [\emph{ʃräf}] 0xC9CD \\
телефонист: \textbf{lwe7n} [\emph{lwe˥n}] 0x732B \\
Телец: \textbf{tro2s} [\emph{tro˩s}] 0x9CCA \\
теллур: \textbf{tlum} [\emph{tlum}] 0x0A6A \\
тело: \textbf{tcic} [\emph{tʃiʃ}] 0xE8AA \\
телята: \textbf{hca7s} [\emph{ʰʃä˥s}] 0x9168 \\
телятина: \textbf{vram} [\emph{vräm}] 0x09D6 \\
тематический: \textbf{fyi2t} [\emph{fji˩t}] 0xB80C \\
тембр: \textbf{hye7p} [\emph{ʰje˥p}] 0x5318 \\
Темза: \textbf{hte2s} [\emph{ʰte˩s}] 0x9B50 \\
темнота: \textbf{kcas} [\emph{kʃäs}] 0x89A2 \\
температура: \textbf{trot} [\emph{trot}] 0xACCA \\
темы: \textbf{hti7c} [\emph{ʰti˥ʃ}] 0xF050 \\
тенор: \textbf{txe7n} [\emph{txe˥n}] 0x73EA \\
тень: \textbf{cyut} [\emph{ʃjut}] 0xAA0D \\
теодолит: \textbf{fli7t} [\emph{fli˥t}] 0xB06C \\
теократия: \textbf{tre2s} [\emph{tre˩s}] 0x9BCA \\
теология: \textbf{flic} [\emph{fliʃ}] 0xE86C \\
теорема: \textbf{frim} [\emph{frim}] 0x08CC \\
теоретизировать: \textbf{fri7s} [\emph{fri˥s}] 0x90CC \\
теоретический: \textbf{frik} [\emph{frik}] 0x28CC \\
теософия: \textbf{hfi2s} [\emph{ʰfi˩s}] 0x9860 \\
тепло: \textbf{nyun} [\emph{njun}] 0x6A06 \\
тепловатый: \textbf{hwu7n} [\emph{ʰwu˥n}] 0x7228 \\
терабайт: \textbf{xret} [\emph{xret}] 0xABDA \\
терзать: \textbf{xra7k} [\emph{xrä˥k}] 0x31DA \\
терка: \textbf{dxac} [\emph{dxäʃ}] 0xE9F3 \\
Терминал: \textbf{ht6n} [\emph{ʰtən}] 0x6D50 \\
термодинамический: \textbf{myek} [\emph{mjek}] 0x2B01 \\
термометр: \textbf{tr6m} [\emph{trəm}] 0x0DCA \\
термостат: \textbf{rw6t} [\emph{rwət}] 0xAD2E \\
термосфера: \textbf{txef} [\emph{txef}] 0xCBEA \\
термоядерный: \textbf{jyek} [\emph{ʒjek}] 0x2B15 \\
терн: \textbf{hlo7m} [\emph{ʰlo˥m}] 0x1458 \\
терпение: \textbf{hjin} [\emph{ʰʒin}] 0x68A8 \\
терракотовый: \textbf{tre7t} [\emph{tre˥t}] 0xB3CA \\
Терри: \textbf{tri7k} [\emph{tri˥k}] 0x30CA \\
терроризм: \textbf{twos} [\emph{twos}] 0x8C2A \\
терьер: \textbf{tre7n} [\emph{tre˥n}] 0x73CA \\
тетралогия: \textbf{tre7c} [\emph{tre˥ʃ}] 0xF3CA \\
тетя: \textbf{hbam} [\emph{ʰbäm}] 0x0988 \\
техасец: \textbf{tse7s} [\emph{tse˥s}] 0x934A \\
технический: \textbf{kfik} [\emph{kfik}] 0x2882 \\
технократия: \textbf{nre2k} [\emph{nre˩k}] 0x3BC6 \\
течка: \textbf{sre2s} [\emph{sre˩s}] 0x9BC9 \\
тигель: \textbf{ks6p} [\emph{ksəp}] 0x4D42 \\
тигр: \textbf{cwic} [\emph{ʃwiʃ}] 0xE82D \\
тик: \textbf{tli7k} [\emph{tli˥k}] 0x306A \\
тик-так: \textbf{tx67t} [\emph{txə˥t}] 0xB5EA \\
тильда: \textbf{tli2t} [\emph{tli˩t}] 0xB86A \\
тимин: \textbf{qyim} [\emph{ŋjim}] 0x0819 \\
Тимоти: \textbf{hti7f} [\emph{ʰti˥f}] 0xD050 \\
тимпан: \textbf{hm67n} [\emph{ʰmə˥n}] 0x7508 \\
тимус: \textbf{tci7m} [\emph{tʃi˥m}] 0x10AA \\
тиофен: \textbf{hfi7f} [\emph{ʰfi˥f}] 0xD060 \\
типологический: \textbf{hp6c} [\emph{ʰpəʃ}] 0xED20 \\
тиранить: \textbf{qra2n} [\emph{ŋrä˩n}] 0x79D9 \\
тире: \textbf{mwe2t} [\emph{mwe˩t}] 0xBB21 \\
Тироль: \textbf{tro7s} [\emph{tro˥s}] 0x94CA \\
тис: \textbf{xyus} [\emph{xjus}] 0x8A1A \\
титан: \textbf{tyem} [\emph{tjem}] 0x0B0A \\
титрование: \textbf{cr6n} [\emph{ʃrən}] 0x6DCD \\
тихо: \textbf{tlac} [\emph{tläʃ}] 0xE96A \\
ткать: \textbf{twen} [\emph{twen}] 0x6B2A \\
ткач: \textbf{hqa7k} [\emph{ʰŋä˥k}] 0x31C8 \\
тмин: \textbf{kref} [\emph{kref}] 0xCBC2 \\
тобогган-одиночка: \textbf{lyuc} [\emph{ljuʃ}] 0xEA0B \\
товарищи: \textbf{hka2t} [\emph{ʰkä˩t}] 0xB910 \\
Тоголезская: \textbf{two7s} [\emph{two˥s}] 0x942A \\
токарь: \textbf{tro2n} [\emph{tro˩n}] 0x7CCA \\
Токийский: \textbf{tyo7k} [\emph{tjo˥k}] 0x340A \\
токсикология: \textbf{sluc} [\emph{sluʃ}] 0xEA69 \\
толстый: \textbf{hmus} [\emph{ʰmus}] 0x8A08 \\
толуол: \textbf{tlu7p} [\emph{tlu˥p}] 0x526A \\
только: \textbf{htun} [\emph{ʰtun}] 0x6A50 \\
Томас: \textbf{hto7s} [\emph{ʰto˥s}] 0x9450 \\
тональный: \textbf{tyu2n} [\emph{tju˩n}] 0x7A0A \\
Тонганский: \textbf{two7k} [\emph{two˥k}] 0x342A \\
тонер: \textbf{tf6n} [\emph{tfən}] 0x6D8A \\
тонзиллит: \textbf{nwa7s} [\emph{nwä˥s}] 0x9126 \\
тоннель: \textbf{tfun} [\emph{tfun}] 0x6A8A \\
топография: \textbf{tr6f} [\emph{trəf}] 0xCDCA \\
топология: \textbf{tl6p} [\emph{tləp}] 0x4D6A \\
тополь: \textbf{pxap} [\emph{pxäp}] 0x49E4 \\
топор: \textbf{kcok} [\emph{kʃok}] 0x2CA2 \\
торжественность: \textbf{hyi7n} [\emph{ʰji˥n}] 0x7018 \\
торий: \textbf{tro7m} [\emph{tro˥m}] 0x14CA \\
тормоз: \textbf{brek} [\emph{brek}] 0x2BD1 \\
торопиться: \textbf{hcok} [\emph{ʰʃok}] 0x2C68 \\
торус: \textbf{hxos} [\emph{ʰxos}] 0x8CD0 \\
торф: \textbf{pxip} [\emph{pxip}] 0x48E4 \\
Тоскана: \textbf{tso7n} [\emph{tso˥n}] 0x744A \\
тост: \textbf{tros} [\emph{tros}] 0x8CCA \\
точность: \textbf{tcok} [\emph{tʃok}] 0x2CAA \\
тошнота: \textbf{hqoc} [\emph{ʰŋoʃ}] 0xECC8 \\
трава: \textbf{kfac} [\emph{kfäʃ}] 0xE982 \\
травля: \textbf{byi7n} [\emph{bji˥n}] 0x7011 \\
травоядное: \textbf{xyof} [\emph{xjof}] 0xCC1A \\
травяной: \textbf{hx6p} [\emph{ʰxəp}] 0x4DD0 \\
трагедия: \textbf{ryat} [\emph{rjät}] 0xA90E \\
традиция: \textbf{hroc} [\emph{ʰroʃ}] 0xEC70 \\
трактирщик: \textbf{twi7t} [\emph{twi˥t}] 0xB02A \\
трактор: \textbf{tl6k} [\emph{tlək}] 0x2D6A \\
тральщик: \textbf{swa7p} [\emph{swä˥p}] 0x5129 \\
транзистор: \textbf{gzot} [\emph{gzot}] 0xAC52 \\
транс: \textbf{gras} [\emph{gräs}] 0x89D2 \\
Трансильвания: \textbf{tle7n} [\emph{tle˥n}] 0x736A \\
транскрибировать: \textbf{nyis} [\emph{njis}] 0x8806 \\
транслокация: \textbf{tw6k} [\emph{twək}] 0x2D2A \\
транспортир: \textbf{qr6t} [\emph{ŋrət}] 0xADD9 \\
транспортируемый: \textbf{hni2k} [\emph{ʰni˩k}] 0x3830 \\
трансформеры: \textbf{tfa7m} [\emph{tfä˥m}] 0x118A \\
трансцендентный: \textbf{tfaf} [\emph{tfäf}] 0xC98A \\
транш: \textbf{hte2c} [\emph{ʰte˩ʃ}] 0xFB50 \\
трапеция: \textbf{tc6m} [\emph{tʃəm}] 0x0DAA \\
тратить: \textbf{hqa7t} [\emph{ʰŋä˥t}] 0xB1C8 \\
траулер: \textbf{tlo7s} [\emph{tlo˥s}] 0x946A \\
трафик: \textbf{hraf} [\emph{ʰräf}] 0xC970 \\
трахеотомия: \textbf{tce7m} [\emph{tʃe˥m}] 0x13AA \\
трезвенник: \textbf{two7t} [\emph{two˥t}] 0xB42A \\
трезубец: \textbf{txi7c} [\emph{txi˥ʃ}] 0xF0EA \\
трель: \textbf{xri7n} [\emph{xri˥n}] 0x70DA \\
трема: \textbf{dra7s} [\emph{drä˥s}] 0x91D3 \\
тренер: \textbf{klot} [\emph{klot}] 0xAC62 \\
трение: \textbf{mroc} [\emph{mroʃ}] 0xECC1 \\
трепел: \textbf{hto7p} [\emph{ʰto˥p}] 0x5450 \\
трепет: \textbf{cron} [\emph{ʃron}] 0x6CCD \\
трепетный: \textbf{htu7s} [\emph{ʰtu˥s}] 0x9250 \\
треска: \textbf{hki7f} [\emph{ʰki˥f}] 0xD010 \\
треугольник: \textbf{tyek} [\emph{tjek}] 0x2B0A \\
трехвалентный: \textbf{gva7t} [\emph{gvä˥t}] 0xB192 \\
трехгодичного: \textbf{tye2n} [\emph{tje˩n}] 0x7B0A \\
трехмерный: \textbf{twem} [\emph{twem}] 0x0B2A \\
трехъязычный: \textbf{tya7c} [\emph{tjä˥ʃ}] 0xF10A \\
трещина: \textbf{rwak} [\emph{rwäk}] 0x292E \\
трещотка: \textbf{rwa2s} [\emph{rwä˩s}] 0x992E \\
три: \textbf{tyin} [\emph{tjin}] 0x680A \\
триангуляция: \textbf{tla7c} [\emph{tlä˥ʃ}] 0xF16A \\
трибуна: \textbf{htu2n} [\emph{ʰtu˩n}] 0x7A50 \\
тривиальный: \textbf{twaf} [\emph{twäf}] 0xC92A \\
тригонометрия: \textbf{tfom} [\emph{tfom}] 0x0C8A \\
тридцатый: \textbf{tse2s} [\emph{tse˩s}] 0x9B4A \\
трижды: \textbf{tsa2t} [\emph{tsä˩t}] 0xB94A \\
триколор: \textbf{tso2n} [\emph{tso˩n}] 0x7C4A \\
триллион: \textbf{tlos} [\emph{tlos}] 0x8C6A \\
трио: \textbf{tro7n} [\emph{tro˥n}] 0x74CA \\
триоксида: \textbf{ts62s} [\emph{tsə˩s}] 0x9D4A \\
трипсин: \textbf{tsu2s} [\emph{tsu˩s}] 0x9A4A \\
трипсы: \textbf{tro7p} [\emph{tro˥p}] 0x54CA \\
триптих: \textbf{txi2c} [\emph{txi˩ʃ}] 0xF8EA \\
трирема: \textbf{tri7m} [\emph{tri˥m}] 0x10CA \\
тритон: \textbf{nwo2t} [\emph{nwo˩t}] 0xBC26 \\
трихинеллез: \textbf{tco2s} [\emph{tʃo˩s}] 0x9CAA \\
Троица: \textbf{ts67t} [\emph{tsə˥t}] 0xB54A \\
Трой: \textbf{tro2t} [\emph{tro˩t}] 0xBCCA \\
тромбоз: \textbf{hr67s} [\emph{ʰrə˥s}] 0x9570 \\
тромбон: \textbf{twop} [\emph{twop}] 0x4C2A \\
тромбоцит: \textbf{hle7t} [\emph{ʰle˥t}] 0xB358 \\
тропический: \textbf{hjip} [\emph{ʰʒip}] 0x48A8 \\
тропосфера: \textbf{hr6f} [\emph{ʰrəf}] 0xCD70 \\
тротуар: \textbf{jwat} [\emph{ʒwät}] 0xA935 \\
троянец: \textbf{hro2n} [\emph{ʰro˩n}] 0x7C70 \\
труба: \textbf{pwap} [\emph{pwäp}] 0x4924 \\
трубач: \textbf{hto2p} [\emph{ʰto˩p}] 0x5C50 \\
трубка: \textbf{fr6k} [\emph{frək}] 0x2DCC \\
трубочка: \textbf{twu7t} [\emph{twu˥t}] 0xB22A \\
трудоголик: \textbf{kxo2k} [\emph{kxo˩k}] 0x3CE2 \\
трудоемкими: \textbf{hto2c} [\emph{ʰto˩ʃ}] 0xFC50 \\
трудолюбиво: \textbf{bri7k} [\emph{bri˥k}] 0x30D1 \\
трудоспособный: \textbf{cla2t} [\emph{ʃlä˩t}] 0xB96D \\
трусики: \textbf{gya7t} [\emph{gjä˥t}] 0xB112 \\
трусливый: \textbf{grac} [\emph{gräʃ}] 0xE9D2 \\
трюфель: \textbf{tfuf} [\emph{tfuf}] 0xCA8A \\
тряпье: \textbf{hro2k} [\emph{ʰro˩k}] 0x3C70 \\
трясти: \textbf{hzak} [\emph{ʰzäk}] 0x29A0 \\
туберкулез: \textbf{tlo2k} [\emph{tlo˩k}] 0x3C6A \\
туземцы: \textbf{hbi2s} [\emph{ʰbi˩s}] 0x9888 \\
тузы: \textbf{hce7s} [\emph{ʰʃe˥s}] 0x9368 \\
туман: \textbf{hzut} [\emph{ʰzut}] 0xAAA0 \\
тумба: \textbf{blom} [\emph{blom}] 0x0C71 \\
тунец: \textbf{htuf} [\emph{ʰtuf}] 0xCA50 \\
Тунис: \textbf{tsuf} [\emph{tsuf}] 0xCA4A \\
Тупик: \textbf{hle7k} [\emph{ʰle˥k}] 0x3358 \\
тупо: \textbf{hqup} [\emph{ʰŋup}] 0x4AC8 \\
тупой: \textbf{bxut} [\emph{bxut}] 0xAAF1 \\
тупость: \textbf{hbo2t} [\emph{ʰbo˩t}] 0xBC88 \\
турбина: \textbf{hti7n} [\emph{ʰti˥n}] 0x7050 \\
туризм: \textbf{brun} [\emph{brun}] 0x6AD1 \\
Турин: \textbf{tru2n} [\emph{tru˩n}] 0x7ACA \\
Туркменистан: \textbf{tx6s} [\emph{txəs}] 0x8DEA \\
турмалин: \textbf{txup} [\emph{txup}] 0x4AEA \\
турчанка: \textbf{tru2k} [\emph{tru˩k}] 0x3ACA \\
тусклый: \textbf{slo7s} [\emph{slo˥s}] 0x9469 \\
тушить: \textbf{txus} [\emph{txus}] 0x8AEA \\
тщеславный: \textbf{cyic} [\emph{ʃjiʃ}] 0xE80D \\
тыква: \textbf{hguk} [\emph{ʰguk}] 0x2A90 \\
тысяча: \textbf{clan} [\emph{ʃlän}] 0x696D \\
тысячелетний: \textbf{cle2n} [\emph{ʃle˩n}] 0x7B6D \\
тьфу: \textbf{htu7f} [\emph{ʰtu˥f}] 0xD250 \\
тюнер: \textbf{tcu2n} [\emph{tʃu˩n}] 0x7AAA \\
тюремщик: \textbf{hyi7s} [\emph{ʰji˥s}] 0x9018 \\
Тюрингия: \textbf{tru7n} [\emph{tru˥n}] 0x72CA \\
тюрьма: \textbf{hgin} [\emph{ʰgin}] 0x6890 \\
тюряга: \textbf{pyi2k} [\emph{pji˩k}] 0x3804 \\
тяготение: \textbf{gvat} [\emph{gvät}] 0xA992 \\
тяжеловес: \textbf{xwe7n} [\emph{xwe˥n}] 0x733A \\
\subsection{ у }
убегай: \textbf{xwap} [\emph{xwäp}] 0x493A \\
убеждать: \textbf{swet} [\emph{swet}] 0xAB29 \\
убийство: \textbf{ksut} [\emph{ksut}] 0xAA42 \\
\subsection{ У }
Ублюдок: \textbf{hwa2n} [\emph{ʰwä˩n}] 0x7928 \\
уборщица: \textbf{hdu7n} [\emph{ʰdu˥n}] 0x7298 \\
убывать: \textbf{hw6k} [\emph{ʰwək}] 0x2D28 \\
Уважаемые: \textbf{syin} [\emph{sjin}] 0x6809 \\
уважаемый: \textbf{rya2t} [\emph{rjä˩t}] 0xB90E \\
уважение: \textbf{tsak} [\emph{tsäk}] 0x294A \\
увеличительный: \textbf{hxet} [\emph{ʰxet}] 0xABD0 \\
увещевание: \textbf{hce2n} [\emph{ʰʃe˩n}] 0x7B68 \\
увещевать: \textbf{tci2n} [\emph{tʃi˩n}] 0x78AA \\
увидел: \textbf{hyic} [\emph{ʰjiʃ}] 0xE818 \\
увлеченный: \textbf{rwa7t} [\emph{rwä˥t}] 0xB12E \\
Увы: \textbf{hxa7s} [\emph{ʰxä˥s}] 0x91D0 \\
увядать: \textbf{fwot} [\emph{fwot}] 0xAC2C \\
Уганда: \textbf{gxut} [\emph{gxut}] 0xAAF2 \\
углевод: \textbf{bxif} [\emph{bxif}] 0xC8F1 \\
углеводород: \textbf{drop} [\emph{drop}] 0x4CD3 \\
углерод: \textbf{kron} [\emph{kron}] 0x6CC2 \\
углубления: \textbf{kcum} [\emph{kʃum}] 0x0AA2 \\
угнетенный: \textbf{dri2n} [\emph{dri˩n}] 0x78D3 \\
уговаривание: \textbf{hxun} [\emph{ʰxun}] 0x6AD0 \\
угол: \textbf{kyun} [\emph{kjun}] 0x6A02 \\
угол: \textbf{hyok} [\emph{ʰjok}] 0x2C18 \\
угрожающий: \textbf{tcet} [\emph{tʃet}] 0xABAA \\
удав: \textbf{hbo7s} [\emph{ʰbo˥s}] 0x9488 \\
удаление: \textbf{hsuf} [\emph{ʰsuf}] 0xCA48 \\
Удалить: \textbf{dlas} [\emph{dläs}] 0x8973 \\
ударив: \textbf{hsa7k} [\emph{ʰsä˥k}] 0x3148 \\
ударов: \textbf{tla7t} [\emph{tlä˥t}] 0xB16A \\
удержан: \textbf{kxi7t} [\emph{kxi˥t}] 0xB0E2 \\
удушение: \textbf{tcom} [\emph{tʃom}] 0x0CAA \\
ужасный: \textbf{kyap} [\emph{kjäp}] 0x4902 \\
уже: \textbf{hlic} [\emph{ʰliʃ}] 0xE858 \\
Узбекистан: \textbf{bxus} [\emph{bxus}] 0x8AF1 \\
уздечка: \textbf{pxam} [\emph{pxäm}] 0x09E4 \\
узкий: \textbf{nyap} [\emph{njäp}] 0x4906 \\
узурпировать: \textbf{zrup} [\emph{zrup}] 0x4AD4 \\
узуфрукт: \textbf{hyu2k} [\emph{ʰju˩k}] 0x3A18 \\
Уилли: \textbf{hwi2n} [\emph{ʰwi˩n}] 0x7828 \\
Уильям: \textbf{hwi7m} [\emph{ʰwi˥m}] 0x1028 \\
укладчик: \textbf{tye7k} [\emph{tje˥k}] 0x330A \\
уклончиво: \textbf{vwef} [\emph{vwef}] 0xCB36 \\
уклончивый: \textbf{kxim} [\emph{kxim}] 0x08E2 \\
Украина: \textbf{kruf} [\emph{kruf}] 0xCAC2 \\
украсть: \textbf{truc} [\emph{truʃ}] 0xEACA \\
украшенный: \textbf{qrun} [\emph{ŋrun}] 0x6AD9 \\
уксус: \textbf{srek} [\emph{srek}] 0x2BC9 \\
укусить: \textbf{krut} [\emph{krut}] 0xAAC2 \\
улей: \textbf{hxa2k} [\emph{ʰxä˩k}] 0x39D0 \\
улицы: \textbf{hsa7t} [\emph{ʰsä˥t}] 0xB148 \\
улучшать: \textbf{mruk} [\emph{mruk}] 0x2AC1 \\
улыбка: \textbf{myac} [\emph{mjäʃ}] 0xE901 \\
ультиматум: \textbf{hti2m} [\emph{ʰti˩m}] 0x1850 \\
ультрафиолетовый: \textbf{twef} [\emph{twef}] 0xCB2A \\
умирающий: \textbf{mro7n} [\emph{mro˥n}] 0x74C1 \\
умник: \textbf{hwa2k} [\emph{ʰwä˩k}] 0x3928 \\
умножить: \textbf{hgoc} [\emph{ʰgoʃ}] 0xEC90 \\
умолять: \textbf{hno2t} [\emph{ʰno˩t}] 0xBC30 \\
умственно: \textbf{mwis} [\emph{mwis}] 0x8821 \\
умышленно: \textbf{gyit} [\emph{gjit}] 0xA812 \\
ун-американских: \textbf{mre7k} [\emph{mre˥k}] 0x33C1 \\
унаследованный: \textbf{vret} [\emph{vret}] 0xABD6 \\
унгвент: \textbf{hjo7n} [\emph{ʰʒo˥n}] 0x74A8 \\
универсально: \textbf{hci7s} [\emph{ʰʃi˥s}] 0x9068 \\
университеты: \textbf{vyut} [\emph{vjut}] 0xAA16 \\
унижать: \textbf{mya7n} [\emph{mjä˥n}] 0x7101 \\
уникальный: \textbf{nwuk} [\emph{nwuk}] 0x2A26 \\
унитарный: \textbf{hnu2t} [\emph{ʰnu˩t}] 0xBA30 \\
уничижительный: \textbf{prof} [\emph{prof}] 0xCCC4 \\
унция: \textbf{hqos} [\emph{ʰŋos}] 0x8CC8 \\
упал: \textbf{hca2t} [\emph{ʰʃä˩t}] 0xB968 \\
упертый: \textbf{hzun} [\emph{ʰzun}] 0x6AA0 \\
упираться: \textbf{hq6k} [\emph{ʰŋək}] 0x2DC8 \\
уполномоченный: \textbf{kc6n} [\emph{kʃən}] 0x6DA2 \\
упорно: \textbf{hne2t} [\emph{ʰne˩t}] 0xBB30 \\
упорство: \textbf{gvit} [\emph{gvit}] 0xA892 \\
управительница: \textbf{dji7n} [\emph{dʒi˥n}] 0x70B3 \\
упреждающий: \textbf{pfef} [\emph{pfef}] 0xCB84 \\
Ура: \textbf{xre2s} [\emph{xre˩s}] 0x9BDA \\
уравнения: \textbf{qyon} [\emph{ŋjon}] 0x6C19 \\
уравнитель: \textbf{hya2c} [\emph{ʰjä˩ʃ}] 0xF918 \\
урал: \textbf{mru7s} [\emph{mru˥s}] 0x92C1 \\
уран: \textbf{ryum} [\emph{rjum}] 0x0A0E \\
урбанизация: \textbf{bzif} [\emph{bzif}] 0xC851 \\
урду: \textbf{zru7t} [\emph{zru˥t}] 0xB2D4 \\
урны: \textbf{grus} [\emph{grus}] 0x8AD2 \\
урод: \textbf{gwi7k} [\emph{gwi˥k}] 0x3032 \\
уродливый: \textbf{cluk} [\emph{ʃluk}] 0x2A6D \\
урожай: \textbf{tcos} [\emph{tʃos}] 0x8CAA \\
уроженцев: \textbf{hte2p} [\emph{ʰte˩p}] 0x5B50 \\
Уругвай: \textbf{rwuk} [\emph{rwuk}] 0x2A2E \\
усик: \textbf{sr6c} [\emph{srəʃ}] 0xEDC9 \\
усилитель: \textbf{hnif} [\emph{ʰnif}] 0xC830 \\
условный: \textbf{gyot} [\emph{gjot}] 0xAC12 \\
успех: \textbf{slak} [\emph{släk}] 0x2969 \\
успешно: \textbf{hsop} [\emph{ʰsop}] 0x4C48 \\
успокойся: \textbf{tya2m} [\emph{tjä˩m}] 0x190A \\
устала: \textbf{hluk} [\emph{ʰluk}] 0x2A58 \\
установить: \textbf{tlip} [\emph{tlip}] 0x486A \\
установка: \textbf{txas} [\emph{txäs}] 0x89EA \\
устрашение: \textbf{dri7s} [\emph{dri˥s}] 0x90D3 \\
устрица: \textbf{glip} [\emph{glip}] 0x4872 \\
утечка: \textbf{clek} [\emph{ʃlek}] 0x2B6D \\
утих: \textbf{hbe7t} [\emph{ʰbe˥t}] 0xB388 \\
утка: \textbf{hbok} [\emph{ʰbok}] 0x2C88 \\
утконос: \textbf{psa7p} [\emph{psä˥p}] 0x5144 \\
утолщение: \textbf{ht67n} [\emph{ʰtə˥n}] 0x7550 \\
утомляет: \textbf{hfi7k} [\emph{ʰfi˥k}] 0x3060 \\
утопия: \textbf{htu2p} [\emph{ʰtu˩p}] 0x5A50 \\
утреня: \textbf{tya2n} [\emph{tjä˩n}] 0x790A \\
утро: \textbf{hsas} [\emph{ʰsäs}] 0x8948 \\
утроить: \textbf{txep} [\emph{txep}] 0x4BEA \\
ухо: \textbf{kwac} [\emph{kwäʃ}] 0xE922 \\
уховертка: \textbf{vwik} [\emph{vwik}] 0x2836 \\
учение: \textbf{hcec} [\emph{ʰʃeʃ}] 0xEB68 \\
ученый: \textbf{gyac} [\emph{gjäʃ}] 0xE912 \\
учетверять: \textbf{kru2k} [\emph{kru˩k}] 0x3AC2 \\
учил: \textbf{hta2n} [\emph{ʰtä˩n}] 0x7950 \\
учить: \textbf{lyac} [\emph{ljäʃ}] 0xE90B \\
учреждения: \textbf{tson} [\emph{tson}] 0x6C4A \\
\subsection{ ф }
фаг: \textbf{fwi7t} [\emph{fwi˥t}] 0xB02C \\
фагот: \textbf{bzap} [\emph{bzäp}] 0x4951 \\
фагоцитоз: \textbf{gwus} [\emph{gwus}] 0x8A32 \\
фазан: \textbf{hf6t} [\emph{ʰfət}] 0xAD60 \\
файл: \textbf{hfas} [\emph{ʰfäs}] 0x8960 \\
факел: \textbf{troc} [\emph{troʃ}] 0xECCA \\
факс: \textbf{hfos} [\emph{ʰfos}] 0x8C60 \\
фактически: \textbf{fwic} [\emph{fwiʃ}] 0xE82C \\
факториал: \textbf{kxin} [\emph{kxin}] 0x68E2 \\
факультативный: \textbf{kli2f} [\emph{kli˩f}] 0xD862 \\
фаланга: \textbf{hfa2s} [\emph{ʰfä˩s}] 0x9960 \\
фаллический: \textbf{fla2k} [\emph{flä˩k}] 0x396C \\
фаллопиевы: \textbf{hlo7n} [\emph{ʰlo˥n}] 0x7458 \\
фальсифицировать: \textbf{ksaf} [\emph{ksäf}] 0xC942 \\
фальцетом: \textbf{fla7s} [\emph{flä˥s}] 0x916C \\
фальшборт: \textbf{xwek} [\emph{xwek}] 0x2B3A \\
фанатик: \textbf{hko7t} [\emph{ʰko˥t}] 0xB410 \\
фанера: \textbf{fwap} [\emph{fwäp}] 0x492C \\
фантасмагория: \textbf{nwa7k} [\emph{nwä˥k}] 0x3126 \\
фарватер: \textbf{xrem} [\emph{xrem}] 0x0BDA \\
\subsection{ Ф }
Фаренгейт: \textbf{fyet} [\emph{fjet}] 0xAB0C \\
Фарерская: \textbf{fro7s} [\emph{fro˥s}] 0x94CC \\
фармацевт: \textbf{fwas} [\emph{fwäs}] 0x892C \\
фартинг: \textbf{hta7p} [\emph{ʰtä˥p}] 0x5150 \\
фартук: \textbf{zrin} [\emph{zrin}] 0x68D4 \\
фарфор: \textbf{psi2n} [\emph{psi˩n}] 0x7844 \\
фарш: \textbf{tfes} [\emph{tfes}] 0x8B8A \\
фаршированный: \textbf{txof} [\emph{txof}] 0xCCEA \\
Фары: \textbf{xla2s} [\emph{xlä˩s}] 0x997A \\
фасоль: \textbf{bwon} [\emph{bwon}] 0x6C31 \\
фат: \textbf{kso7k} [\emph{kso˥k}] 0x3442 \\
фатализм: \textbf{fl6m} [\emph{fləm}] 0x0D6C \\
фашизм: \textbf{hfa7m} [\emph{ʰfä˥m}] 0x1160 \\
фаэтон: \textbf{hfe2n} [\emph{ʰfe˩n}] 0x7B60 \\
февраль: \textbf{frif} [\emph{frif}] 0xC8CC \\
федеральный: \textbf{frap} [\emph{fräp}] 0x49CC \\
феерия: \textbf{tye7n} [\emph{tje˥n}] 0x730A \\
феи: \textbf{pri7n} [\emph{pri˥n}] 0x70C4 \\
фекалии: \textbf{hvis} [\emph{ʰvis}] 0x88B0 \\
феминистка: \textbf{fwis} [\emph{fwis}] 0x882C \\
Феникс: \textbf{hfi2k} [\emph{ʰfi˩k}] 0x3860 \\
фенил: \textbf{fyi7t} [\emph{fji˥t}] 0xB00C \\
фенол: \textbf{hf67n} [\emph{ʰfə˥n}] 0x7560 \\
фенотипический: \textbf{hni7p} [\emph{ʰni˥p}] 0x5030 \\
фенхель: \textbf{fyen} [\emph{fjen}] 0x6B0C \\
феодальный: \textbf{hfa2t} [\emph{ʰfä˩t}] 0xB960 \\
фермер: \textbf{hfim} [\emph{ʰfim}] 0x0860 \\
фетиш: \textbf{hfe2t} [\emph{ʰfe˩t}] 0xBB60 \\
фехтовальщик: \textbf{hfe7n} [\emph{ʰfe˥n}] 0x7360 \\
фибрилляция: \textbf{flec} [\emph{fleʃ}] 0xEB6C \\
Фивы: \textbf{hte7p} [\emph{ʰte˥p}] 0x5350 \\
Фиджи: \textbf{dxic} [\emph{dxiʃ}] 0xE8F3 \\
физика: \textbf{bzik} [\emph{bzik}] 0x2851 \\
физиология: \textbf{vlik} [\emph{vlik}] 0x2876 \\
физиотерапия: \textbf{zyip} [\emph{zjip}] 0x4814 \\
Филадельфия: \textbf{flef} [\emph{flef}] 0xCB6C \\
филантропия: \textbf{hla2p} [\emph{ʰlä˩p}] 0x5958 \\
филе: \textbf{flif} [\emph{flif}] 0xC86C \\
филиал: \textbf{crac} [\emph{ʃräʃ}] 0xE9CD \\
филигрань: \textbf{fli2n} [\emph{fli˩n}] 0x786C \\
Филимон: \textbf{flom} [\emph{flom}] 0x0C6C \\
Филипп: \textbf{fli7p} [\emph{fli˥p}] 0x506C \\
Филиппины: \textbf{flip} [\emph{flip}] 0x486C \\
филогенетическое: \textbf{fwen} [\emph{fwen}] 0x6B2C \\
филология: \textbf{fli7n} [\emph{fli˥n}] 0x706C \\
философ: \textbf{hfif} [\emph{ʰfif}] 0xC860 \\
фильм: \textbf{flim} [\emph{flim}] 0x086C \\
фильтры: \textbf{xli2t} [\emph{xli˩t}] 0xB87A \\
Финляндия: \textbf{fl6n} [\emph{flən}] 0x6D6C \\
финн: \textbf{hfi7n} [\emph{ʰfi˥n}] 0x7060 \\
фисташковый: \textbf{pwaf} [\emph{pwäf}] 0xC924 \\
фишинг: \textbf{hqic} [\emph{ʰŋiʃ}] 0xE8C8 \\
флагшток: \textbf{lya2t} [\emph{ljä˩t}] 0xB90B \\
фламинго: \textbf{flok} [\emph{flok}] 0x2C6C \\
ФЛАНДРИЯ: \textbf{hfa7s} [\emph{ʰfä˥s}] 0x9160 \\
фланель: \textbf{hfa2n} [\emph{ʰfä˩n}] 0x7960 \\
флебит: \textbf{qyes} [\emph{ŋjes}] 0x8B19 \\
Флер-де-Лис: \textbf{fyes} [\emph{fjes}] 0x8B0C \\
флирт: \textbf{fri2n} [\emph{fri˩n}] 0x78CC \\
флорист: \textbf{flos} [\emph{flos}] 0x8C6C \\
флот: \textbf{nwen} [\emph{nwen}] 0x6B26 \\
флоэма: \textbf{flem} [\emph{flem}] 0x0B6C \\
флуоресценция: \textbf{fwon} [\emph{fwon}] 0x6C2C \\
флюид: \textbf{klem} [\emph{klem}] 0x0B62 \\
Фогарма: \textbf{gra7m} [\emph{grä˥m}] 0x11D2 \\
фокус: \textbf{hcuk} [\emph{ʰʃuk}] 0x2A68 \\
Фолклендские: \textbf{fl6s} [\emph{fləs}] 0x8D6C \\
фолликул: \textbf{fluk} [\emph{fluk}] 0x2A6C \\
фольклор: \textbf{floc} [\emph{floʃ}] 0xEC6C \\
фонарь: \textbf{tla7n} [\emph{tlä˥n}] 0x716A \\
фонема: \textbf{hfom} [\emph{ʰfom}] 0x0C60 \\
фонетический: \textbf{fyek} [\emph{fjek}] 0x2B0C \\
фонограф: \textbf{nwof} [\emph{nwof}] 0xCC26 \\
фонтанирование: \textbf{pca2n} [\emph{pʃä˩n}] 0x79A4 \\
форма: \textbf{hfam} [\emph{ʰfäm}] 0x0960 \\
формализм: \textbf{fli7s} [\emph{fli˥s}] 0x906C \\
формалин: \textbf{fro7n} [\emph{fro˥n}] 0x74CC \\
формальдегид: \textbf{fla2t} [\emph{flä˩t}] 0xB96C \\
фоссилизация: \textbf{sl6c} [\emph{sləʃ}] 0xED69 \\
фосфат: \textbf{psof} [\emph{psof}] 0xCC44 \\
фосфид: \textbf{hfi2f} [\emph{ʰfi˩f}] 0xD860 \\
фосфор: \textbf{hfof} [\emph{ʰfof}] 0xCC60 \\
фосфоресцирующий: \textbf{sre7n} [\emph{sre˥n}] 0x73C9 \\
фотогеничный: \textbf{hqo2n} [\emph{ʰŋo˩n}] 0x7CC8 \\
фотография: \textbf{hfoc} [\emph{ʰfoʃ}] 0xEC60 \\
фотон: \textbf{hfo7t} [\emph{ʰfo˥t}] 0xB460 \\
фотосинтез: \textbf{tso2s} [\emph{tso˩s}] 0x9C4A \\
фотоэлектрические: \textbf{fwek} [\emph{fwek}] 0x2B2C \\
фотоэлектрический: \textbf{hfi7t} [\emph{ʰfi˥t}] 0xB060 \\
фрагментация: \textbf{gxon} [\emph{gxon}] 0x6CF2 \\
фраза: \textbf{fyak} [\emph{fjäk}] 0x290C \\
Фракия: \textbf{hre7s} [\emph{ʰre˥s}] 0x9370 \\
фракталы: \textbf{hf6s} [\emph{ʰfəs}] 0x8D60 \\
фрамбезия: \textbf{hya7m} [\emph{ʰjä˥m}] 0x1118 \\
франк: \textbf{fra7n} [\emph{frä˥n}] 0x71CC \\
Французский: \textbf{frec} [\emph{freʃ}] 0xEBCC \\
франшиза: \textbf{dyac} [\emph{djäʃ}] 0xE913 \\
фрейлин: \textbf{dwi7n} [\emph{dwi˥n}] 0x7033 \\
френология: \textbf{fli2c} [\emph{fli˩ʃ}] 0xF86C \\
фригольд: \textbf{txo7t} [\emph{txo˥t}] 0xB4EA \\
фриз: \textbf{xwis} [\emph{xwis}] 0x883A \\
фронтиспис: \textbf{tso7p} [\emph{tso˥p}] 0x544A \\
фталевый: \textbf{fla7k} [\emph{flä˥k}] 0x316C \\
фтор: \textbf{flun} [\emph{flun}] 0x6A6C \\
фторид: \textbf{fru2t} [\emph{fru˩t}] 0xBACC \\
фунгицид: \textbf{qyut} [\emph{ŋjut}] 0xAA19 \\
фундаментализм: \textbf{nya7m} [\emph{njä˥m}] 0x1106 \\
фундук: \textbf{zl6t} [\emph{zlət}] 0xAD74 \\
функции: \textbf{hzit} [\emph{ʰzit}] 0xA8A0 \\
фураж: \textbf{kra2s} [\emph{krä˩s}] 0x99C2 \\
фургон: \textbf{mron} [\emph{mron}] 0x6CC1 \\
фут: \textbf{pyas} [\emph{pjäs}] 0x8904 \\
футбол: \textbf{hfop} [\emph{ʰfop}] 0x4C60 \\
фыркнул: \textbf{syi7t} [\emph{sji˥t}] 0xB009 \\
\subsection{ х }
ха: \textbf{hx6k} [\emph{ʰxək}] 0x2DD0 \\
ха-ха: \textbf{xr6k} [\emph{xrək}] 0x2DDA \\
хадж: \textbf{xla2c} [\emph{xlä˩ʃ}] 0xF97A \\
хакеры: \textbf{hxe7k} [\emph{ʰxe˥k}] 0x33D0 \\
халат: \textbf{hjap} [\emph{ʰʒäp}] 0x49A8 \\
халдейский: \textbf{cli7n} [\emph{ʃli˥n}] 0x706D \\
халиф: \textbf{kli7f} [\emph{kli˥f}] 0xD062 \\
халтура: \textbf{xwi7k} [\emph{xwi˥k}] 0x303A \\
ханжество: \textbf{hp67s} [\emph{ʰpə˥s}] 0x9520 \\
ханство: \textbf{kya2k} [\emph{kjä˩k}] 0x3902 \\
ханука: \textbf{xwa2k} [\emph{xwä˩k}] 0x393A \\
хаотичный: \textbf{xr6t} [\emph{xrət}] 0xADDA \\
характерно: \textbf{tca7s} [\emph{tʃä˥s}] 0x91AA \\
\subsection{ Х }
Хартум: \textbf{xrum} [\emph{xrum}] 0x0ADA \\
хвастовство: \textbf{grun} [\emph{grun}] 0x6AD2 \\
хвост: \textbf{hkek} [\emph{ʰkek}] 0x2B10 \\
Хелен: \textbf{xle7n} [\emph{xle˥n}] 0x737A \\
херувим: \textbf{kre7p} [\emph{kre˥p}] 0x53C2 \\
хеттский: \textbf{hx67t} [\emph{ʰxə˥t}] 0xB5D0 \\
химическая: \textbf{ryik} [\emph{rjik}] 0x280E \\
химус: \textbf{cyem} [\emph{ʃjem}] 0x0B0D \\
хинин: \textbf{kwi7n} [\emph{kwi˥n}] 0x7022 \\
хинолина: \textbf{hk67n} [\emph{ʰkə˥n}] 0x7510 \\
хиппи: \textbf{xyi7p} [\emph{xji˥p}] 0x501A \\
хиромантия: \textbf{swa7k} [\emph{swä˥k}] 0x3129 \\
хиропрактики: \textbf{cr6t} [\emph{ʃrət}] 0xADCD \\
Хиросима: \textbf{xra7m} [\emph{xrä˥m}] 0x11DA \\
хитрость: \textbf{cluc} [\emph{ʃluʃ}] 0xEA6D \\
хихикающий: \textbf{hgi7n} [\emph{ʰgi˥n}] 0x7090 \\
хлеб: \textbf{hpup} [\emph{ʰpup}] 0x4A20 \\
хлопок: \textbf{krap} [\emph{kräp}] 0x49C2 \\
хлопушка: \textbf{xla7p} [\emph{xlä˥p}] 0x517A \\
хлопьевидный: \textbf{flu2n} [\emph{flu˩n}] 0x7A6C \\
хлор: \textbf{klos} [\emph{klos}] 0x8C62 \\
хлорал: \textbf{hlo7c} [\emph{ʰlo˥ʃ}] 0xF458 \\
хлористый: \textbf{lwot} [\emph{lwot}] 0xAC2B \\
хлорофилл: \textbf{zrif} [\emph{zrif}] 0xC8D4 \\
хмуриться: \textbf{srun} [\emph{srun}] 0x6AC9 \\
хобби: \textbf{cyop} [\emph{ʃjop}] 0x4C0D \\
ходули: \textbf{txa2s} [\emph{txä˩s}] 0x99EA \\
хозяйка: \textbf{hxi7n} [\emph{ʰxi˥n}] 0x70D0 \\
хоккей: \textbf{cwek} [\emph{ʃwek}] 0x2B2D \\
холера: \textbf{kla2n} [\emph{klä˩n}] 0x7962 \\
холестерин: \textbf{kxep} [\emph{kxep}] 0x4BE2 \\
холм: \textbf{slat} [\emph{slät}] 0xA969 \\
холодильник: \textbf{pcif} [\emph{pʃif}] 0xC8A4 \\
холодно: \textbf{tlin} [\emph{tlin}] 0x686A \\
хор: \textbf{cros} [\emph{ʃros}] 0x8CCD \\
Хорватия: \textbf{kfoc} [\emph{kfoʃ}] 0xEC82 \\
хорек: \textbf{frop} [\emph{frop}] 0x4CCC \\
хорошо: \textbf{cwoc} [\emph{ʃwoʃ}] 0xEC2D \\
Храбрый: \textbf{pyam} [\emph{pjäm}] 0x0904 \\
храмы: \textbf{myop} [\emph{mjop}] 0x4C01 \\
храп: \textbf{xro2n} [\emph{xro˩n}] 0x7CDA \\
хрен: \textbf{xyac} [\emph{xjäʃ}] 0xE91A \\
хризантема: \textbf{rwaf} [\emph{rwäf}] 0xC92E \\
хриплый: \textbf{sro7k} [\emph{sro˥k}] 0x34C9 \\
христианство: \textbf{hrim} [\emph{ʰrim}] 0x0870 \\
Христос: \textbf{hri7t} [\emph{ʰri˥t}] 0xB070 \\
хром: \textbf{krom} [\emph{krom}] 0x0CC2 \\
хромат: \textbf{rwom} [\emph{rwom}] 0x0C2E \\
хроматин: \textbf{hjo7t} [\emph{ʰʒo˥t}] 0xB4A8 \\
хромосома: \textbf{hjom} [\emph{ʰʒom}] 0x0CA8 \\
хромосфера: \textbf{hro2f} [\emph{ʰro˩f}] 0xDC70 \\
хромота: \textbf{kco7n} [\emph{kʃo˥n}] 0x74A2 \\
хронический: \textbf{zyin} [\emph{zjin}] 0x6814 \\
хронометраж: \textbf{tyi2p} [\emph{tji˩p}] 0x580A \\
хруст: \textbf{hki7k} [\emph{ʰki˥k}] 0x3010 \\
хрящ: \textbf{jruk} [\emph{ʒruk}] 0x2AD5 \\
хунта: \textbf{hcu7n} [\emph{ʰʃu˥n}] 0x7268 \\
хурма: \textbf{ps6m} [\emph{psəm}] 0x0D44 \\
Хусайн: \textbf{hxu2n} [\emph{ʰxu˩n}] 0x7AD0 \\
Хэллоуин: \textbf{xlec} [\emph{xleʃ}] 0xEB7A \\
хэтчбэк: \textbf{qyap} [\emph{ŋjäp}] 0x4919 \\
\subsection{ ц }
цапфа: \textbf{hro7f} [\emph{ʰro˥f}] 0xD470 \\
царапина: \textbf{grak} [\emph{gräk}] 0x29D2 \\
цареубийство: \textbf{gxa2t} [\emph{gxä˩t}] 0xB9F2 \\
царствовал: \textbf{rwa2n} [\emph{rwä˩n}] 0x792E \\
цвет: \textbf{hron} [\emph{ʰron}] 0x6C70 \\
цветных: \textbf{hyo2s} [\emph{ʰjo˩s}] 0x9C18 \\
цветущий: \textbf{bxun} [\emph{bxun}] 0x6AF1 \\
цезий: \textbf{hs6m} [\emph{ʰsəm}] 0x0D48 \\
целебный: \textbf{sya7k} [\emph{sjä˥k}] 0x3109 \\
цели: \textbf{mya2k} [\emph{mjä˩k}] 0x3901 \\
целибат: \textbf{bl6s} [\emph{bləs}] 0x8D71 \\
целлюлоза: \textbf{lyus} [\emph{ljus}] 0x8A0B \\
целостность: \textbf{nres} [\emph{nres}] 0x8BC6 \\
цель: \textbf{hmuc} [\emph{ʰmuʃ}] 0xEA08 \\
цельсию: \textbf{sl6s} [\emph{sləs}] 0x8D69 \\
цемент: \textbf{sy6n} [\emph{sjən}] 0x6D09 \\
цензор: \textbf{sles} [\emph{sles}] 0x8B69 \\
цент: \textbf{hfen} [\emph{ʰfen}] 0x6B60 \\
центр: \textbf{hset} [\emph{ʰset}] 0xAB48 \\
центробежный: \textbf{nyef} [\emph{njef}] 0xCB06 \\
цепи: \textbf{hle7n} [\emph{ʰle˥n}] 0x7358 \\
церебральный: \textbf{sri7p} [\emph{sri˥p}] 0x50C9 \\
церемония: \textbf{sric} [\emph{sriʃ}] 0xE8C9 \\
церий: \textbf{srem} [\emph{srem}] 0x0BC9 \\
церковь: \textbf{kyic} [\emph{kjiʃ}] 0xE802 \\
циан: \textbf{cyi2n} [\emph{ʃji˩n}] 0x780D \\
цианид: \textbf{hsa2m} [\emph{ʰsä˩m}] 0x1948 \\
цианобактерии: \textbf{hna7t} [\emph{ʰnä˥t}] 0xB130 \\
цианоз: \textbf{hva7s} [\emph{ʰvä˥s}] 0x91B0 \\
циветта: \textbf{gvi2t} [\emph{gvi˩t}] 0xB892 \\
цивилизация: \textbf{hwim} [\emph{ʰwim}] 0x0828 \\
циклический: \textbf{gwik} [\emph{gwik}] 0x2832 \\
циклотрон: \textbf{kl67t} [\emph{klə˥t}] 0xB562 \\
цикорий: \textbf{hc6c} [\emph{ʰʃəʃ}] 0xED68 \\
цилиндрический: \textbf{nri2k} [\emph{nri˩k}] 0x38C6 \\
цинк: \textbf{zyik} [\emph{zjik}] 0x2814 \\
цирк: \textbf{sruk} [\emph{sruk}] 0x2AC9 \\
цирконий: \textbf{hzum} [\emph{ʰzum}] 0x0AA0 \\
циркулярно: \textbf{kwi2n} [\emph{kwi˩n}] 0x7822 \\
цистит: \textbf{qwi2t} [\emph{ŋwi˩t}] 0xB839 \\
цитозин: \textbf{gwis} [\emph{gwis}] 0x8832 \\
цитоплазма: \textbf{clom} [\emph{ʃlom}] 0x0C6D \\
цитрат: \textbf{sru2t} [\emph{sru˩t}] 0xBAC9 \\
цифровой: \textbf{djis} [\emph{dʒis}] 0x88B3 \\
цыганский: \textbf{hgi7s} [\emph{ʰgi˥s}] 0x9090 \\
цюрих: \textbf{zruc} [\emph{zruʃ}] 0xEAD4 \\
\subsection{ ч }
ч.л.: \textbf{tse2t} [\emph{tse˩t}] 0xBB4A \\
\subsection{ Ч }
Чад: \textbf{ty6t} [\emph{tjət}] 0xAD0A \\
чайка: \textbf{qwak} [\emph{ŋwäk}] 0x2939 \\
чалый: \textbf{hro7s} [\emph{ʰro˥s}] 0x9470 \\
чаны: \textbf{txo7s} [\emph{txo˥s}] 0x94EA \\
час: \textbf{hcam} [\emph{ʰʃäm}] 0x0968 \\
часовщик: \textbf{hgi2c} [\emph{ʰgi˩ʃ}] 0xF890 \\
частично: \textbf{pfun} [\emph{pfun}] 0x6A84 \\
частный: \textbf{hpaf} [\emph{ʰpäf}] 0xC920 \\
частота: \textbf{fwin} [\emph{fwin}] 0x682C \\
чашелистник: \textbf{hpe2p} [\emph{ʰpe˩p}] 0x5B20 \\
чашечка: \textbf{kl6k} [\emph{klək}] 0x2D62 \\
чеканка: \textbf{tfi2p} [\emph{tfi˩p}] 0x588A \\
челнок: \textbf{cwos} [\emph{ʃwos}] 0x8C2D \\
человек: \textbf{hpat} [\emph{ʰpät}] 0xA920 \\
человечество: \textbf{hjat} [\emph{ʰʒät}] 0xA9A8 \\
человечный: \textbf{djun} [\emph{dʒun}] 0x6AB3 \\
чемодан: \textbf{hxo2n} [\emph{ʰxo˩n}] 0x7CD0 \\
чемпион: \textbf{kcip} [\emph{kʃip}] 0x48A2 \\
червеобразный: \textbf{hvef} [\emph{ʰvef}] 0xCBB0 \\
червивый: \textbf{tri2p} [\emph{tri˩p}] 0x58CA \\
червоточина: \textbf{hwo2n} [\emph{ʰwo˩n}] 0x7C28 \\
через: \textbf{tru7k} [\emph{tru˥k}] 0x32CA \\
череп: \textbf{krok} [\emph{krok}] 0x2CC2 \\
черепаха: \textbf{kruc} [\emph{kruʃ}] 0xEAC2 \\
чернильница: \textbf{hqi7t} [\emph{ʰŋi˥t}] 0xB0C8 \\
черногория: \textbf{hme7k} [\emph{ʰme˥k}] 0x3308 \\
черный: \textbf{hkis} [\emph{ʰkis}] 0x8810 \\
чеснок: \textbf{glun} [\emph{glun}] 0x6A72 \\
честный: \textbf{tcis} [\emph{tʃis}] 0x88AA \\
Четверг: \textbf{cwis} [\emph{ʃwis}] 0x882D \\
четвероногие: \textbf{kre7k} [\emph{kre˥k}] 0x33C2 \\
четвертичный: \textbf{kce2n} [\emph{kʃe˩n}] 0x7BA2 \\
четырехгранный: \textbf{cre7t} [\emph{ʃre˥t}] 0xB3CD \\
четырехмерный: \textbf{rwe2n} [\emph{rwe˩n}] 0x7B2E \\
четырехугольник: \textbf{kre2n} [\emph{kre˩n}] 0x7BC2 \\
Чехословакией: \textbf{clof} [\emph{ʃlof}] 0xCC6D \\
чеченец: \textbf{tce2n} [\emph{tʃe˩n}] 0x7BAA \\
Чешский: \textbf{hc6k} [\emph{ʰʃək}] 0x2D68 \\
чешуйчатый: \textbf{kwa2s} [\emph{kwä˩s}] 0x9922 \\
чизкейк: \textbf{hce2k} [\emph{ʰʃe˩k}] 0x3B68 \\
Чили: \textbf{zlic} [\emph{zliʃ}] 0xE874 \\
чип: \textbf{cwip} [\emph{ʃwip}] 0x482D \\
чипсов: \textbf{kxi7p} [\emph{kxi˥p}] 0x50E2 \\
чириканье: \textbf{tce7t} [\emph{tʃe˥t}] 0xB3AA \\
чистилище: \textbf{hpo2t} [\emph{ʰpo˩t}] 0xBC20 \\
чистописание: \textbf{hni7f} [\emph{ʰni˥f}] 0xD030 \\
читать: \textbf{tyac} [\emph{tjäʃ}] 0xE90A \\
чихать: \textbf{cyec} [\emph{ʃjeʃ}] 0xEB0D \\
член: \textbf{hcen} [\emph{ʰʃen}] 0x6B68 \\
членистоногие: \textbf{fro2t} [\emph{fro˩t}] 0xBCCC \\
чопорный: \textbf{hti7m} [\emph{ʰti˥m}] 0x1050 \\
чревный: \textbf{cl6k} [\emph{ʃlək}] 0x2D6D \\
чревовещатель: \textbf{vli7t} [\emph{vli˥t}] 0xB076 \\
чувствительный: \textbf{hmek} [\emph{ʰmek}] 0x2B08 \\
чудак: \textbf{kx6c} [\emph{kxəʃ}] 0xEDE2 \\
чудотворный: \textbf{hmi2s} [\emph{ʰmi˩s}] 0x9808 \\
чума: \textbf{pyuk} [\emph{pjuk}] 0x2A04 \\
\subsection{ ш }
шаблон: \textbf{tlep} [\emph{tlep}] 0x4B6A \\
шаблонный: \textbf{bli7p} [\emph{bli˥p}] 0x5071 \\
шагов: \textbf{hpa2k} [\emph{ʰpä˩k}] 0x3920 \\
шагомер: \textbf{hye2t} [\emph{ʰje˩t}] 0xBB18 \\
шалот: \textbf{clop} [\emph{ʃlop}] 0x4C6D \\
шампанское: \textbf{cwap} [\emph{ʃwäp}] 0x492D \\
шампунь: \textbf{cwup} [\emph{ʃwup}] 0x4A2D \\
\subsection{ Ш }
Шанхай: \textbf{hca7f} [\emph{ʰʃä˥f}] 0xD168 \\
шапка: \textbf{hmup} [\emph{ʰmup}] 0x4A08 \\
шарик: \textbf{tyi7s} [\emph{tji˥s}] 0x900A \\
шарить: \textbf{hmo7p} [\emph{ʰmo˥p}] 0x5408 \\
шарлатан: \textbf{hca7p} [\emph{ʰʃä˥p}] 0x5168 \\
шарлатанство: \textbf{xri2s} [\emph{xri˩s}] 0x98DA \\
шарф: \textbf{kcaf} [\emph{kʃäf}] 0xC9A2 \\
шасси: \textbf{xra7c} [\emph{xrä˥ʃ}] 0xF1DA \\
шах: \textbf{cya7c} [\emph{ʃjä˥ʃ}] 0xF10D \\
шахматы: \textbf{crec} [\emph{ʃreʃ}] 0xEBCD \\
шашки: \textbf{txem} [\emph{txem}] 0x0BEA \\
швейцар: \textbf{dwen} [\emph{dwen}] 0x6B33 \\
Швейцария: \textbf{swif} [\emph{swif}] 0xC829 \\
Швеция: \textbf{hw6t} [\emph{ʰwət}] 0xAD28 \\
швырнула: \textbf{xle2t} [\emph{xle˩t}] 0xBB7A \\
шедевр: \textbf{hci2s} [\emph{ʰʃi˩s}] 0x9868 \\
шекель: \textbf{hce7k} [\emph{ʰʃe˥k}] 0x3368 \\
шекспировский: \textbf{kse7p} [\emph{kse˥p}] 0x5342 \\
шелест: \textbf{tsa7s} [\emph{tsä˥s}] 0x914A \\
шелк: \textbf{slik} [\emph{slik}] 0x2869 \\
шелковица: \textbf{clup} [\emph{ʃlup}] 0x4A6D \\
шепелявость: \textbf{tsi7p} [\emph{tsi˥p}] 0x504A \\
шериф: \textbf{cref} [\emph{ʃref}] 0xCBCD \\
шерстяной: \textbf{hwu2n} [\emph{ʰwu˩n}] 0x7A28 \\
шершень: \textbf{br6n} [\emph{brən}] 0x6DD1 \\
шестерки: \textbf{tci7s} [\emph{tʃi˥s}] 0x90AA \\
шестидесятый: \textbf{tyi2s} [\emph{tji˩s}] 0x980A \\
шестиугольник: \textbf{gxun} [\emph{gxun}] 0x6AF2 \\
шестнадцатеричный: \textbf{hse2m} [\emph{ʰse˩m}] 0x1B48 \\
шестнадцать: \textbf{hsos} [\emph{ʰsos}] 0x8C48 \\
шесть: \textbf{hlis} [\emph{ʰlis}] 0x8858 \\
шея: \textbf{kyuk} [\emph{kjuk}] 0x2A02 \\
шиит: \textbf{cyi2p} [\emph{ʃji˩p}] 0x580D \\
шип: \textbf{swip} [\emph{swip}] 0x4829 \\
широта: \textbf{kxut} [\emph{kxut}] 0xAAE2 \\
ширь: \textbf{hw67s} [\emph{ʰwə˥s}] 0x9528 \\
шитье: \textbf{slon} [\emph{slon}] 0x6C69 \\
шифер: \textbf{slef} [\emph{slef}] 0xCB69 \\
шифрование: \textbf{krep} [\emph{krep}] 0x4BC2 \\
шишковидный: \textbf{gri7t} [\emph{gri˥t}] 0xB0D2 \\
шкаф: \textbf{nrip} [\emph{nrip}] 0x48C6 \\
школа: \textbf{kcuk} [\emph{kʃuk}] 0x2AA2 \\
шлагбаумом: \textbf{zlop} [\emph{zlop}] 0x4C74 \\
шлепок: \textbf{pya7k} [\emph{pjä˥k}] 0x3104 \\
шлюп: \textbf{lyup} [\emph{ljup}] 0x4A0B \\
шовинист: \textbf{hqis} [\emph{ʰŋis}] 0x88C8 \\
шок: \textbf{cwok} [\emph{ʃwok}] 0x2C2D \\
шоколад: \textbf{clok} [\emph{ʃlok}] 0x2C6D \\
шорты: \textbf{xrat} [\emph{xrät}] 0xA9DA \\
шоры: \textbf{bx6n} [\emph{bxən}] 0x6DF1 \\
шоссе: \textbf{crok} [\emph{ʃrok}] 0x2CCD \\
Шотландия: \textbf{klo2n} [\emph{klo˩n}] 0x7C62 \\
шпалера: \textbf{gles} [\emph{gles}] 0x8B72 \\
шпатлевка: \textbf{pxu7t} [\emph{pxu˥t}] 0xB2E4 \\
шпиль: \textbf{ht6p} [\emph{ʰtəp}] 0x4D50 \\
шпион: \textbf{gyat} [\emph{gjät}] 0xA912 \\
шплинт: \textbf{vyot} [\emph{vjot}] 0xAC16 \\
шпунт: \textbf{sya7p} [\emph{sjä˥p}] 0x5109 \\
шрапнель: \textbf{hca2p} [\emph{ʰʃä˩p}] 0x5968 \\
шрифт: \textbf{txit} [\emph{txit}] 0xA8EA \\
штаб-квартира: \textbf{drok} [\emph{drok}] 0x2CD3 \\
штампами: \textbf{kci2s} [\emph{kʃi˩s}] 0x98A2 \\
штатив: \textbf{tyep} [\emph{tjep}] 0x4B0A \\
штекер: \textbf{tlek} [\emph{tlek}] 0x2B6A \\
штифтик: \textbf{pru2t} [\emph{pru˩t}] 0xBAC4 \\
штифты: \textbf{pyi7n} [\emph{pji˥n}] 0x7004 \\
штопор: \textbf{kr6c} [\emph{krəʃ}] 0xEDC2 \\
шторы: \textbf{pra7n} [\emph{prä˥n}] 0x71C4 \\
штрейкбрехер: \textbf{qrak} [\emph{ŋräk}] 0x29D9 \\
шум: \textbf{hcum} [\emph{ʰʃum}] 0x0A68 \\
шумно: \textbf{hso7s} [\emph{ʰso˥s}] 0x9448 \\
шутить: \textbf{hzas} [\emph{ʰzäs}] 0x89A0 \\
шхуна: \textbf{kxa2n} [\emph{kxä˩n}] 0x79E2 \\
\subsection{ щ }
щавель: \textbf{xraf} [\emph{xräf}] 0xC9DA \\
щегол: \textbf{gxec} [\emph{gxeʃ}] 0xEBF2 \\
щегольство: \textbf{pl6c} [\emph{pləʃ}] 0xED64 \\
щекотание: \textbf{tlu7n} [\emph{tlu˥n}] 0x726A \\
щелочной: \textbf{lwis} [\emph{lwis}] 0x882B \\
щелчок: \textbf{klek} [\emph{klek}] 0x2B62 \\
щенок: \textbf{hpo7k} [\emph{ʰpo˥k}] 0x3420 \\
щетина: \textbf{tri7p} [\emph{tri˥p}] 0x50CA \\
щетка: \textbf{broc} [\emph{broʃ}] 0xECD1 \\
щипцы: \textbf{tcem} [\emph{tʃem}] 0x0BAA \\
щупальце: \textbf{tsu7k} [\emph{tsu˥k}] 0x324A \\
\subsection{ э }
эвкалипт: \textbf{klus} [\emph{klus}] 0x8A62 \\
эволюционный: \textbf{lwon} [\emph{lwon}] 0x6C2B \\
эврика: \textbf{hyu7k} [\emph{ʰju˥k}] 0x3218 \\
эвристический: \textbf{hxu7s} [\emph{ʰxu˥s}] 0x92D0 \\
эвтаназия: \textbf{fra7s} [\emph{frä˥s}] 0x91CC \\
эвтектика: \textbf{tyu2t} [\emph{tju˩t}] 0xBA0A \\
эвфемизм: \textbf{hw6s} [\emph{ʰwəs}] 0x8D28 \\
эгоизм: \textbf{hsi7s} [\emph{ʰsi˥s}] 0x9048 \\
эгоистический: \textbf{hgi7c} [\emph{ʰgi˥ʃ}] 0xF090 \\
\subsection{ Э }
Эдинбург: \textbf{hdi7p} [\emph{ʰdi˥p}] 0x5098 \\
Эквадор: \textbf{kwof} [\emph{kwof}] 0xCC22 \\
экватор: \textbf{xwet} [\emph{xwet}] 0xAB3A \\
экзогамия: \textbf{gxom} [\emph{gxom}] 0x0CF2 \\
экзоскелет: \textbf{zlek} [\emph{zlek}] 0x2B74 \\
экзотерический: \textbf{hge7k} [\emph{ʰge˥k}] 0x3390 \\
экзотермический: \textbf{hzem} [\emph{ʰzem}] 0x0BA0 \\
Экклезиаст: \textbf{hke2s} [\emph{ʰke˩s}] 0x9B10 \\
эклиптика: \textbf{hwe2t} [\emph{ʰwe˩t}] 0xBB28 \\
экология: \textbf{ploc} [\emph{ploʃ}] 0xEC64 \\
эконометрия: \textbf{kyo2t} [\emph{kjo˩t}] 0xBC02 \\
экономика: \textbf{htom} [\emph{ʰtom}] 0x0C50 \\
экономист: \textbf{kcom} [\emph{kʃom}] 0x0CA2 \\
экономка: \textbf{nrap} [\emph{nräp}] 0x49C6 \\
эксгумация: \textbf{hzuc} [\emph{ʰzuʃ}] 0xEAA0 \\
экскурсия: \textbf{kses} [\emph{kses}] 0x8B42 \\
экспедиция: \textbf{hsen} [\emph{ʰsen}] 0x6B48 \\
эксперт: \textbf{tsec} [\emph{tseʃ}] 0xEB4A \\
эксплуатировать: \textbf{kyon} [\emph{kjon}] 0x6C02 \\
экспонат: \textbf{sren} [\emph{sren}] 0x6BC9 \\
экспоненциальный: \textbf{swos} [\emph{swos}] 0x8C29 \\
экспорт: \textbf{sroc} [\emph{sroʃ}] 0xECC9 \\
экспромтом: \textbf{hfo7n} [\emph{ʰfo˥n}] 0x7460 \\
экстатический: \textbf{bjak} [\emph{bʒäk}] 0x29B1 \\
экстернализация: \textbf{tle2s} [\emph{tle˩s}] 0x9B6A \\
экстракт: \textbf{kret} [\emph{kret}] 0xABC2 \\
экстраполяция: \textbf{klef} [\emph{klef}] 0xCB62 \\
экстремист: \textbf{trum} [\emph{trum}] 0x0ACA \\
экстровертность: \textbf{tre2c} [\emph{tre˩ʃ}] 0xFBCA \\
эктодерма: \textbf{hte2m} [\emph{ʰte˩m}] 0x1B50 \\
эктопический: \textbf{twe2k} [\emph{twe˩k}] 0x3B2A \\
эластичный: \textbf{lyak} [\emph{ljäk}] 0x290B \\
элегия: \textbf{hwi7k} [\emph{ʰwi˥k}] 0x3028 \\
электрод: \textbf{gjet} [\emph{gʒet}] 0xABB2 \\
электролиз: \textbf{vres} [\emph{vres}] 0x8BD6 \\
электролит: \textbf{kce7t} [\emph{kʃe˥t}] 0xB3A2 \\
электромагнитный: \textbf{vrek} [\emph{vrek}] 0x2BD6 \\
электронный: \textbf{lyin} [\emph{ljin}] 0x680B \\
электроскоп: \textbf{kre2k} [\emph{kre˩k}] 0x3BC2 \\
элерон: \textbf{gr6n} [\emph{grən}] 0x6DD2 \\
Элизабет: \textbf{hli2p} [\emph{ʰli˩p}] 0x5858 \\
эликсир: \textbf{psi7s} [\emph{psi˥s}] 0x9044 \\
элита: \textbf{lyim} [\emph{ljim}] 0x080B \\
эллипсы: \textbf{hle7p} [\emph{ʰle˥p}] 0x5358 \\
эльзасский: \textbf{hl67s} [\emph{ʰlə˥s}] 0x9558 \\
Эм: \textbf{hye7m} [\emph{ʰje˥m}] 0x1318 \\
эманация: \textbf{hlu2s} [\emph{ʰlu˩s}] 0x9A58 \\
эмбриология: \textbf{bla2k} [\emph{blä˩k}] 0x3971 \\
эмбрион: \textbf{bron} [\emph{bron}] 0x6CD1 \\
эмир: \textbf{myi7m} [\emph{mji˥m}] 0x1001 \\
эмираты: \textbf{mwek} [\emph{mwek}] 0x2B21 \\
эмиссия: \textbf{psif} [\emph{psif}] 0xC844 \\
эмоционально: \textbf{xyek} [\emph{xjek}] 0x2B1A \\
эмпирик: \textbf{qyi7t} [\emph{ŋji˥t}] 0xB019 \\
эму: \textbf{nyum} [\emph{njum}] 0x0A06 \\
эмульсия: \textbf{myuc} [\emph{mjuʃ}] 0xEA01 \\
эмфизема: \textbf{hfi7m} [\emph{ʰfi˥m}] 0x1060 \\
эндокринный: \textbf{dr6n} [\emph{drən}] 0x6DD3 \\
эндотелий: \textbf{hne2p} [\emph{ʰne˩p}] 0x5B30 \\
эндотермический: \textbf{djim} [\emph{dʒim}] 0x08B3 \\
энергия: \textbf{nrit} [\emph{nrit}] 0xA8C6 \\
энтодерма: \textbf{hde2m} [\emph{ʰde˩m}] 0x1B98 \\
энтомолог: \textbf{tl6c} [\emph{tləʃ}] 0xED6A \\
энтомология: \textbf{tle2k} [\emph{tle˩k}] 0x3B6A \\
энтропия: \textbf{drep} [\emph{drep}] 0x4BD3 \\
энтузиазм: \textbf{htec} [\emph{ʰteʃ}] 0xEB50 \\
энциклопедия: \textbf{slo7t} [\emph{slo˥t}] 0xB469 \\
эоцен: \textbf{hs67n} [\emph{ʰsə˥n}] 0x7548 \\
эпигенетическое: \textbf{hpe2n} [\emph{ʰpe˩n}] 0x7B20 \\
эпидермис: \textbf{prep} [\emph{prep}] 0x4BC4 \\
эпизодический: \textbf{psi2c} [\emph{psi˩ʃ}] 0xF844 \\
эпизоотия: \textbf{gzek} [\emph{gzek}] 0x2B52 \\
эпикуреец: \textbf{hpu7k} [\emph{ʰpu˥k}] 0x3220 \\
эпилепсия: \textbf{pli2s} [\emph{pli˩s}] 0x9864 \\
эпистемология: \textbf{sli7m} [\emph{sli˥m}] 0x1069 \\
эпителий: \textbf{pfem} [\emph{pfem}] 0x0B84 \\
эпицентр: \textbf{psef} [\emph{psef}] 0xCB44 \\
эпический: \textbf{hpi2k} [\emph{ʰpi˩k}] 0x3820 \\
эполет: \textbf{cle7t} [\emph{ʃle˥t}] 0xB36D \\
эр: \textbf{hke2p} [\emph{ʰke˩p}] 0x5B10 \\
эргономика: \textbf{hgi7m} [\emph{ʰgi˥m}] 0x1090 \\
эректильная: \textbf{ryi7k} [\emph{rji˥k}] 0x300E \\
эритема: \textbf{xyem} [\emph{xjem}] 0x0B1A \\
Эритрея: \textbf{py6t} [\emph{pjət}] 0xAD04 \\
эркер: \textbf{ryi2n} [\emph{rji˩n}] 0x780E \\
эротика: \textbf{qyek} [\emph{ŋjek}] 0x2B19 \\
эрцгерцог: \textbf{tca7k} [\emph{tʃä˥k}] 0x31AA \\
эскадрилья: \textbf{hse2t} [\emph{ʰse˩t}] 0xBB48 \\
эскалатор: \textbf{kxet} [\emph{kxet}] 0xABE2 \\
эскиз: \textbf{syet} [\emph{sjet}] 0xAB09 \\
эскимос: \textbf{hse7m} [\emph{ʰse˥m}] 0x1348 \\
эскорт: \textbf{srot} [\emph{srot}] 0xACC9 \\
эспланада: \textbf{ple7t} [\emph{ple˥t}] 0xB364 \\
эспрессо: \textbf{vros} [\emph{vros}] 0x8CD6 \\
эстетически: \textbf{hse2c} [\emph{ʰse˩ʃ}] 0xFB48 \\
Эстония: \textbf{qwen} [\emph{ŋwen}] 0x6B39 \\
эстрагон: \textbf{gzak} [\emph{gzäk}] 0x2952 \\
эстроген: \textbf{sr6k} [\emph{srək}] 0x2DC9 \\
эсхатология: \textbf{clo2t} [\emph{ʃlo˩t}] 0xBC6D \\
этан: \textbf{tfe7n} [\emph{tfe˥n}] 0x738A \\
этидия: \textbf{txim} [\emph{txim}] 0x08EA \\
этил: \textbf{tfi7c} [\emph{tfi˥ʃ}] 0xF08A \\
этилен: \textbf{fle7n} [\emph{fle˥n}] 0x736C \\
этимология: \textbf{tli2n} [\emph{tli˩n}] 0x786A \\
этнология: \textbf{hne2c} [\emph{ʰne˩ʃ}] 0xFB30 \\
Это: \textbf{htit} [\emph{ʰtit}] 0xA850 \\
этология: \textbf{tlo2c} [\emph{tlo˩ʃ}] 0xFC6A \\
этос: \textbf{fyos} [\emph{fjos}] 0x8C0C \\
этрусский: \textbf{tre7s} [\emph{tre˥s}] 0x93CA \\
эукариоты: \textbf{kru2t} [\emph{kru˩t}] 0xBAC2 \\
эфемерность: \textbf{hve2n} [\emph{ʰve˩n}] 0x7BB0 \\
Эфиопия: \textbf{hxop} [\emph{ʰxop}] 0x4CD0 \\
эфир: \textbf{bvet} [\emph{bvet}] 0xAB91 \\
эфирный: \textbf{tri7n} [\emph{tri˥n}] 0x70CA \\
эхо: \textbf{tcek} [\emph{tʃek}] 0x2BAA \\
\subsection{ ю }
юбка: \textbf{twet} [\emph{twet}] 0xAB2A \\
юго-восточный: \textbf{swa2n} [\emph{swä˩n}] 0x7929 \\
юго-юго-восток: \textbf{hnu2k} [\emph{ʰnu˩k}] 0x3A30 \\
\subsection{ Ю }
Югославия: \textbf{glaf} [\emph{gläf}] 0xC972 \\
южный: \textbf{hzin} [\emph{ʰzin}] 0x68A0 \\
Юлиан: \textbf{jlun} [\emph{ʒlun}] 0x6A75 \\
юмор: \textbf{hmum} [\emph{ʰmum}] 0x0A08 \\
Юпитер: \textbf{hvip} [\emph{ʰvip}] 0x48B0 \\
юра: \textbf{jrus} [\emph{ʒrus}] 0x8AD5 \\
юриспруденция: \textbf{hfi7s} [\emph{ʰfi˥s}] 0x9060 \\
юрист: \textbf{fri7t} [\emph{fri˥t}] 0xB0CC \\
\subsection{ я }
я: \textbf{jyim} [\emph{ʒjim}] 0x0815 \\
ябеда: \textbf{xlek} [\emph{xlek}] 0x2B7A \\
яблоко: \textbf{hpep} [\emph{ʰpep}] 0x4B20 \\
\subsection{ Я }
Ява: \textbf{jwac} [\emph{ʒwäʃ}] 0xE935 \\
явление: \textbf{gjan} [\emph{gʒän}] 0x69B2 \\
ягодицы: \textbf{nwop} [\emph{nwop}] 0x4C26 \\
яд: \textbf{hzic} [\emph{ʰziʃ}] 0xE8A0 \\
ядра: \textbf{hne7k} [\emph{ʰne˥k}] 0x3330 \\
язва: \textbf{lyon} [\emph{ljon}] 0x6C0B \\
язык: \textbf{pyac} [\emph{pjäʃ}] 0xE904 \\
язычество: \textbf{hpi2s} [\emph{ʰpi˩s}] 0x9820 \\
язычок: \textbf{hyu7s} [\emph{ʰju˥s}] 0x9218 \\
язь: \textbf{by6t} [\emph{bjət}] 0xAD11 \\
яичник: \textbf{dvac} [\emph{dväʃ}] 0xE993 \\
яйценосный: \textbf{hvus} [\emph{ʰvus}] 0x8AB0 \\
яйцо: \textbf{dyan} [\emph{djän}] 0x6913 \\
якорь: \textbf{jlak} [\emph{ʒläk}] 0x2975 \\
яма: \textbf{hgon} [\emph{ʰgon}] 0x6C90 \\
Ямайка: \textbf{xyak} [\emph{xjäk}] 0x291A \\
ямбический: \textbf{hya2p} [\emph{ʰjä˩p}] 0x5918 \\
ямкоделатель: \textbf{pfap} [\emph{pfäp}] 0x4984 \\
январь: \textbf{nwak} [\emph{nwäk}] 0x2926 \\
янки: \textbf{hya2k} [\emph{ʰjä˩k}] 0x3918 \\
Японский: \textbf{hg6n} [\emph{ʰgən}] 0x6D90 \\
яростно: \textbf{nret} [\emph{nret}] 0xABC6 \\
ярь-медянка: \textbf{hvi2k} [\emph{ʰvi˩k}] 0x38B0 \\
ясновидение: \textbf{kyi7n} [\emph{kji˥n}] 0x7002 \\
ястреб: \textbf{blik} [\emph{blik}] 0x2871 \\
яхта: \textbf{hyop} [\emph{ʰjop}] 0x4C18 \\
ящерица: \textbf{kyip} [\emph{kjip}] 0x4802 \\

\section{пяш -> русский}
\subsection{ b }
\textbf{b6} [\emph{bə}] 0x363E: кими\\
\textbf{ba} [\emph{bä}] 0x263E: а также или\\
\textbf{be} [\emph{be}] 0x2E3E: словесный\\
\textbf{bi} [\emph{bi}] 0x223E: будущее напряженный\\
\textbf{bi7} [\emph{bi˥}] 0x423E: ИНГ\\
\textbf{bj6t} [\emph{bʒət}] 0xADB1: миндаль\\
\textbf{bja2n} [\emph{bʒä˩n}] 0x79B1: раздающий милостыню\\
\textbf{bja7k} [\emph{bʒä˥k}] 0x31B1: батик\\
\textbf{bja7t} [\emph{bʒä˥t}] 0xB1B1: субарктический\\
\textbf{bjac} [\emph{bʒäʃ}] 0xE9B1: пивоваренный завод\\
\textbf{bjak} [\emph{bʒäk}] 0x29B1: экстатический\\
\textbf{bjam} [\emph{bʒäm}] 0x09B1: притуплять\\
\textbf{bjan} [\emph{bʒän}] 0x69B1: бродить\\
\textbf{bjap} [\emph{bʒäp}] 0x49B1: спинодержатель\\
\textbf{bjas} [\emph{bʒäs}] 0x89B1: подвид\\
\textbf{bjat} [\emph{bʒät}] 0xA9B1: пособие\\
\textbf{bjeh} [\emph{bʒeʰ}] 0x5D8F: locality gender, 8th density life form, learning to be oversoul, akin to spirit of a place or flock\\
\textbf{bjen} [\emph{bʒen}] 0x6BB1: пляжный\\
\textbf{bjes} [\emph{bʒes}] 0x8BB1: житница\\
\textbf{bjet} [\emph{bʒet}] 0xABB1: бюджет\\
\textbf{bjic} [\emph{bʒiʃ}] 0xE8B1: подсолнечник\\
\textbf{bjih} [\emph{bʒiʰ}] 0x458F: бывший умерший\\
\textbf{bjik} [\emph{bʒik}] 0x28B1: электронная книга\\
\textbf{bjim} [\emph{bʒim}] 0x08B1: аболиционизм\\
\textbf{bjin} [\emph{bʒin}] 0x68B1: прогнозирование\\
\textbf{bjip} [\emph{bʒip}] 0x48B1: ежевика\\
\textbf{bjis} [\emph{bʒis}] 0x88B1: Бэзил\\
\textbf{bjit} [\emph{bʒit}] 0xA8B1: бывший умерший\\
\textbf{bjoc} [\emph{bʒoʃ}] 0xECB1: банджо\\
\textbf{bjok} [\emph{bʒok}] 0x2CB1: субдукции\\
\textbf{bjon} [\emph{bʒon}] 0x6CB1: бекон\\
\textbf{bjot} [\emph{bʒot}] 0xACB1: изымать\\
\textbf{bjun} [\emph{bʒun}] 0x6AB1: продавец книг\\
\textbf{bjut} [\emph{bʒut}] 0xAAB1: приземистый\\
\textbf{bl6k} [\emph{blək}] 0x2D71: балтийский\\
\textbf{bl6n} [\emph{blən}] 0x6D71: Бруней\\
\textbf{bl6p} [\emph{bləp}] 0x4D71: кими\\
\textbf{bl6s} [\emph{bləs}] 0x8D71: целибат\\
\textbf{bl6t} [\emph{blət}] 0xAD71: балласт\\
\textbf{bla2k} [\emph{blä˩k}] 0x3971: эмбриология\\
\textbf{bla2n} [\emph{blä˩n}] 0x7971: заслонять\\
\textbf{bla2s} [\emph{blä˩s}] 0x9971: балюстрада\\
\textbf{bla2t} [\emph{blä˩t}] 0xB971: бильярд\\
\textbf{bla7k} [\emph{blä˥k}] 0x3171: баллистический\\
\textbf{bla7n} [\emph{blä˥n}] 0x7171: Вавилон\\
\textbf{bla7p} [\emph{blä˥p}] 0x5171: болт канат\\
\textbf{bla7s} [\emph{blä˥s}] 0x9171: бацилла\\
\textbf{bla7t} [\emph{blä˥t}] 0xB171: буфетная\\
\textbf{blac} [\emph{bläʃ}] 0xE971: Багаж\\
\textbf{blaf} [\emph{bläf}] 0xC971: изнасилование\\
\textbf{blah} [\emph{bläʰ}] 0x4B8F: вес\\
\textbf{blak} [\emph{bläk}] 0x2971: балкон\\
\textbf{blam} [\emph{bläm}] 0x0971: бальный\\
\textbf{blan} [\emph{blän}] 0x6971: жалоба\\
\textbf{blap} [\emph{bläp}] 0x4971: бумажник\\
\textbf{blas} [\emph{bläs}] 0x8971: опора\\
\textbf{blat} [\emph{blät}] 0xA971: алтарь\\
\textbf{ble2t} [\emph{ble˩t}] 0xBB71: бакелит\\
\textbf{ble7n} [\emph{ble˥n}] 0x7371: китовый ус\\
\textbf{ble7t} [\emph{ble˥t}] 0xB371: Белград\\
\textbf{blec} [\emph{bleʃ}] 0xEB71: Бангладеш\\
\textbf{blef} [\emph{blef}] 0xCB71: книголюб\\
\textbf{bleh} [\emph{bleʰ}] 0x5B8F: благожелательный\\
\textbf{blek} [\emph{blek}] 0x2B71: балет\\
\textbf{blem} [\emph{blem}] 0x0B71: Бельгия\\
\textbf{blen} [\emph{blen}] 0x6B71: блестящий\\
\textbf{blep} [\emph{blep}] 0x4B71: балеарский\\
\textbf{bles} [\emph{bles}] 0x8B71: браслет\\
\textbf{blet} [\emph{blet}] 0xAB71: таблетка\\
\textbf{bli2n} [\emph{bli˩n}] 0x7871: Берлин\\
\textbf{bli2s} [\emph{bli˩s}] 0x9871: черный список\\
\textbf{bli2t} [\emph{bli˩t}] 0xB871: Гибралтар\\
\textbf{bli7k} [\emph{bli˥k}] 0x3071: биофизика\\
\textbf{bli7n} [\emph{bli˥n}] 0x7071: изобличают\\
\textbf{bli7p} [\emph{bli˥p}] 0x5071: шаблонный\\
\textbf{bli7s} [\emph{bli˥s}] 0x9071: обелиск\\
\textbf{bli7t} [\emph{bli˥t}] 0xB071: Билли\\
\textbf{blic} [\emph{bliʃ}] 0xE871: иллюзия\\
\textbf{blif} [\emph{blif}] 0xC871: Болгария\\
\textbf{blih} [\emph{bliʰ}] 0x438F: ограниченный\\
\textbf{blik} [\emph{blik}] 0x2871: ястреб\\
\textbf{blim} [\emph{blim}] 0x0871: бериллий\\
\textbf{blin} [\emph{blin}] 0x6871: воображаемый\\
\textbf{blip} [\emph{blip}] 0x4871: миллиард\\
\textbf{blis} [\emph{blis}] 0x8871: ваниль\\
\textbf{blit} [\emph{blit}] 0xA871: ограниченный\\
\textbf{blo7t} [\emph{blo˥t}] 0xB471: сплюснутый\\
\textbf{bloc} [\emph{bloʃ}] 0xEC71: боулинг\\
\textbf{blof} [\emph{blof}] 0xCC71: Боливия\\
\textbf{blok} [\emph{blok}] 0x2C71: биологическая\\
\textbf{blom} [\emph{blom}] 0x0C71: тумба\\
\textbf{blon} [\emph{blon}] 0x6C71: сходимость\\
\textbf{blop} [\emph{blop}] 0x4C71: колба\\
\textbf{blos} [\emph{blos}] 0x8C71: автобус\\
\textbf{blot} [\emph{blot}] 0xAC71: болт\\
\textbf{blu2n} [\emph{blu˩n}] 0x7A71: булинь\\
\textbf{blu7n} [\emph{blu˥n}] 0x7271: логический\\
\textbf{blu7p} [\emph{blu˥p}] 0x5271: бык щенок\\
\textbf{blu7t} [\emph{blu˥t}] 0xB271: бульдозер\\
\textbf{bluc} [\emph{bluʃ}] 0xEA71: предродовой\\
\textbf{bluf} [\emph{bluf}] 0xCA71: пуленепробиваемый\\
\textbf{bluh} [\emph{bluʰ}] 0x538F: логический\\
\textbf{bluk} [\emph{bluk}] 0x2A71: тактика\\
\textbf{blun} [\emph{blun}] 0x6A71: воздушный шар\\
\textbf{blup} [\emph{blup}] 0x4A71: бесплодие\\
\textbf{blus} [\emph{blus}] 0x8A71: скорая\\
\textbf{blut} [\emph{blut}] 0xAA71: сдвинуть с места\\
\textbf{bo} [\emph{bo}] 0x323E: вероятность\\
\textbf{br6k} [\emph{brək}] 0x2DD1: беатификации\\
\textbf{br6n} [\emph{brən}] 0x6DD1: шершень\\
\textbf{br6p} [\emph{brəp}] 0x4DD1: берберский\\
\textbf{br6s} [\emph{brəs}] 0x8DD1: Бухарест\\
\textbf{br6t} [\emph{brət}] 0xADD1: Бурунди\\
\textbf{bra2c} [\emph{brä˩ʃ}] 0xF9D1: карбюратор\\
\textbf{bra2k} [\emph{brä˩k}] 0x39D1: барак\\
\textbf{bra2n} [\emph{brä˩n}] 0x79D1: северный\\
\textbf{bra2s} [\emph{brä˩s}] 0x99D1: бюстгальтеры\\
\textbf{bra2t} [\emph{brä˩t}] 0xB9D1: оправданный\\
\textbf{bra7k} [\emph{brä˥k}] 0x31D1: барсук\\
\textbf{bra7m} [\emph{brä˥m}] 0x11D1: Авраам\\
\textbf{bra7n} [\emph{brä˥n}] 0x71D1: баронесса\\
\textbf{bra7p} [\emph{brä˥p}] 0x51D1: обломок кирпича\\
\textbf{bra7s} [\emph{brä˥s}] 0x91D1: альбатрос\\
\textbf{bra7t} [\emph{brä˥t}] 0xB1D1: вышитый\\
\textbf{brac} [\emph{bräʃ}] 0xE9D1: среда\\
\textbf{braf} [\emph{bräf}] 0xC9D1: стрижка\\
\textbf{brah} [\emph{bräʰ}] 0x4E8F: ловушка\\
\textbf{brak} [\emph{bräk}] 0x29D1: баррель\\
\textbf{bram} [\emph{bräm}] 0x09D1: бормотание\\
\textbf{bran} [\emph{brän}] 0x69D1: корзина\\
\textbf{brap} [\emph{bräp}] 0x49D1: лом\\
\textbf{bras} [\emph{bräs}] 0x89D1: боковая\\
\textbf{brat} [\emph{brät}] 0xA9D1: возражение\\
\textbf{bre2n} [\emph{bre˩n}] 0x7BD1: кибернетический\\
\textbf{bre2t} [\emph{bre˩t}] 0xBBD1: коечные\\
\textbf{bre7k} [\emph{bre˥k}] 0x33D1: нарушающий закон\\
\textbf{bre7n} [\emph{bre˥n}] 0x73D1: Бахрейн\\
\textbf{bre7s} [\emph{bre˥s}] 0x93D1: либреттист\\
\textbf{bre7t} [\emph{bre˥t}] 0xB3D1: сболтнуть\\
\textbf{brec} [\emph{breʃ}] 0xEBD1: бровь\\
\textbf{bref} [\emph{bref}] 0xCBD1: покрывало\\
\textbf{breh} [\emph{breʰ}] 0x5E8F: изменение из мероприятие\\
\textbf{brek} [\emph{brek}] 0x2BD1: тормоз\\
\textbf{brem} [\emph{brem}] 0x0BD1: лещ\\
\textbf{bren} [\emph{bren}] 0x6BD1: лабиринт\\
\textbf{brep} [\emph{brep}] 0x4BD1: лютик\\
\textbf{bres} [\emph{bres}] 0x8BD1: слово процессор\\
\textbf{bret} [\emph{bret}] 0xABD1: байт\\
\textbf{bri2k} [\emph{bri˩k}] 0x38D1: баратрия\\
\textbf{bri2n} [\emph{bri˩n}] 0x78D1: годовалый\\
\textbf{bri2s} [\emph{bri˩s}] 0x98D1: вылечить все\\
\textbf{bri2t} [\emph{bri˩t}] 0xB8D1: либретто\\
\textbf{bri7k} [\emph{bri˥k}] 0x30D1: трудолюбиво\\
\textbf{bri7m} [\emph{bri˥m}] 0x10D1: шрифт Брайля\\
\textbf{bri7n} [\emph{bri˥n}] 0x70D1: аберраций\\
\textbf{bri7s} [\emph{bri˥s}] 0x90D1: иберийский\\
\textbf{bri7t} [\emph{bri˥t}] 0xB0D1: болотная лихорадка\\
\textbf{bric} [\emph{briʃ}] 0xE8D1: батарея\\
\textbf{brif} [\emph{brif}] 0xC8D1: двоичный значение\\
\textbf{brih} [\emph{briʰ}] 0x468F: база\\
\textbf{brik} [\emph{brik}] 0x28D1: подписываться\\
\textbf{brim} [\emph{brim}] 0x08D1: барий\\
\textbf{brin} [\emph{brin}] 0x68D1: внутренний\\
\textbf{brip} [\emph{brip}] 0x48D1: бабочка\\
\textbf{bris} [\emph{bris}] 0x88D1: беспроводной\\
\textbf{brit} [\emph{brit}] 0xA8D1: кровавый\\
\textbf{bro2n} [\emph{bro˩n}] 0x7CD1: мускулистый\\
\textbf{bro7n} [\emph{bro˥n}] 0x74D1: баритон\\
\textbf{bro7t} [\emph{bro˥t}] 0xB4D1: подпольный акушер\\
\textbf{broc} [\emph{broʃ}] 0xECD1: щетка\\
\textbf{brof} [\emph{brof}] 0xCCD1: обручить\\
\textbf{broh} [\emph{broʰ}] 0x668F: актер роль\\
\textbf{brok} [\emph{brok}] 0x2CD1: переедать\\
\textbf{brom} [\emph{brom}] 0x0CD1: бром\\
\textbf{bron} [\emph{bron}] 0x6CD1: эмбрион\\
\textbf{brop} [\emph{brop}] 0x4CD1: бор\\
\textbf{bros} [\emph{bros}] 0x8CD1: барокко\\
\textbf{brot} [\emph{brot}] 0xACD1: бутон\\
\textbf{bru2t} [\emph{bru˩t}] 0xBAD1: Бейрут\\
\textbf{bru7n} [\emph{bru˥n}] 0x72D1: бумеранг\\
\textbf{bru7p} [\emph{bru˥p}] 0x52D1: билирубин\\
\textbf{bru7t} [\emph{bru˥t}] 0xB2D1: порох\\
\textbf{bruc} [\emph{bruʃ}] 0xEAD1: жаропонижающий\\
\textbf{bruf} [\emph{bruf}] 0xCAD1: блаженство\\
\textbf{bruk} [\emph{bruk}] 0x2AD1: баррикада\\
\textbf{brum} [\emph{brum}] 0x0AD1: подотряд\\
\textbf{brun} [\emph{brun}] 0x6AD1: туризм\\
\textbf{brup} [\emph{brup}] 0x4AD1: розовощекий\\
\textbf{brus} [\emph{brus}] 0x8AD1: зал заседаний совета директоров\\
\textbf{brut} [\emph{brut}] 0xAAD1: виртуальный\\
\textbf{bu} [\emph{bu}] 0x2A3E: абсолютива дело\\
\textbf{bva7n} [\emph{bvä˥n}] 0x7191: двухсотлетие\\
\textbf{bva7t} [\emph{bvä˥t}] 0xB191: крапать\\
\textbf{bvac} [\emph{bväʃ}] 0xE991: гибридизация\\
\textbf{bvaf} [\emph{bväf}] 0xC991: землекоп\\
\textbf{bvah} [\emph{bväʰ}] 0x4C8F: abessive case, indicating the absence of the noun phrase referred to object. state-context negatory-quantifier\\
\textbf{bvak} [\emph{bväk}] 0x2991: abessive дело\\
\textbf{bvam} [\emph{bväm}] 0x0991: нарды\\
\textbf{bvan} [\emph{bvän}] 0x6991: произнесение\\
\textbf{bvap} [\emph{bväp}] 0x4991: Албания\\
\textbf{bvas} [\emph{bväs}] 0x8991: делитель\\
\textbf{bvat} [\emph{bvät}] 0xA991: жестоком обращении\\
\textbf{bveh} [\emph{bveʰ}] 0x5C8F: subessive case, indicates location at the underside of nounphrase. under-context location-case.\\
\textbf{bvek} [\emph{bvek}] 0x2B91: отмучивать\\
\textbf{bven} [\emph{bven}] 0x6B91: Ливан\\
\textbf{bves} [\emph{bves}] 0x8B91: Bavaria\\
\textbf{bvet} [\emph{bvet}] 0xAB91: эфир\\
\textbf{bvi2t} [\emph{bvi˩t}] 0xB891: детский вязаный башмачок\\
\textbf{bvi7h} [\emph{bvi˥ʰ}] 0x848F: дистрибутивный местоимение\\
\textbf{bvi7t} [\emph{bvi˥t}] 0xB091: отказываться\\
\textbf{bvic} [\emph{bviʃ}] 0xE891: все overish\\
\textbf{bvif} [\emph{bvif}] 0xC891: двояковыпуклый\\
\textbf{bvih} [\emph{bviʰ}] 0x448F: двоичный значение\\
\textbf{bvik} [\emph{bvik}] 0x2891: subessive дело\\
\textbf{bvim} [\emph{bvim}] 0x0891: альбумин\\
\textbf{bvin} [\emph{bvin}] 0x6891: вселенная Пол\\
\textbf{bvip} [\emph{bvip}] 0x4891: бобр\\
\textbf{bvis} [\emph{bvis}] 0x8891: безрукавный\\
\textbf{bvit} [\emph{bvit}] 0xA891: неразвитой\\
\textbf{bvok} [\emph{bvok}] 0x2C91: букмекер\\
\textbf{bvon} [\emph{bvon}] 0x6C91: садоводство\\
\textbf{bvot} [\emph{bvot}] 0xAC91: непрошеный\\
\textbf{bvuk} [\emph{bvuk}] 0x2A91: абсолютива дело\\
\textbf{bvun} [\emph{bvun}] 0x6A91: домовой\\
\textbf{bvut} [\emph{bvut}] 0xAA91: свекла\\
\textbf{bw6n} [\emph{bwən}] 0x6D31: соболезнование\\
\textbf{bwa2n} [\emph{bwä˩n}] 0x7931: наметка\\
\textbf{bwa2t} [\emph{bwä˩t}] 0xB931: эстрада для оркестра\\
\textbf{bwa7k} [\emph{bwä˥k}] 0x3131: акушерство\\
\textbf{bwa7n} [\emph{bwä˥n}] 0x7131: скотобойня\\
\textbf{bwa7t} [\emph{bwä˥t}] 0xB131: глушь\\
\textbf{bwac} [\emph{bwäʃ}] 0xE931: зевать\\
\textbf{bwaf} [\emph{bwäf}] 0xC931: будуар\\
\textbf{bwah} [\emph{bwäʰ}] 0x498F: злость\\
\textbf{bwak} [\emph{bwäk}] 0x2931: кукла\\
\textbf{bwam} [\emph{bwäm}] 0x0931: избиратель\\
\textbf{bwan} [\emph{bwän}] 0x6931: племянница\\
\textbf{bwap} [\emph{bwäp}] 0x4931: галантерея\\
\textbf{bwas} [\emph{bwäs}] 0x8931: браузер\\
\textbf{bwat} [\emph{bwät}] 0xA931: одеяло\\
\textbf{bwe7n} [\emph{bwe˥n}] 0x7331: боцман\\
\textbf{bwek} [\emph{bwek}] 0x2B31: портфель\\
\textbf{bwen} [\emph{bwen}] 0x6B31: навязчивая идея\\
\textbf{bwet} [\emph{bwet}] 0xAB31: кишечник\\
\textbf{bwi2n} [\emph{bwi˩n}] 0x7831: би ежегодный\\
\textbf{bwi7t} [\emph{bwi˥t}] 0xB031: бисексуальность\\
\textbf{bwic} [\emph{bwiʃ}] 0xE831: подводная лодка\\
\textbf{bwih} [\emph{bwiʰ}] 0x418F: wise gender, 5th density life form, learning wisdom, akin to Jesus Ascended, Latwii, Arcturians, shape shifting aliens\\
\textbf{bwik} [\emph{bwik}] 0x2831: sublative дело\\
\textbf{bwim} [\emph{bwim}] 0x0831: биохимия\\
\textbf{bwin} [\emph{bwin}] 0x6831: впечатление\\
\textbf{bwip} [\emph{bwip}] 0x4831: БРИК Брач\\
\textbf{bwis} [\emph{bwis}] 0x8831: бикини\\
\textbf{bwit} [\emph{bwit}] 0xA831: граница\\
\textbf{bwo7n} [\emph{bwo˥n}] 0x7431: бронх\\
\textbf{bwoc} [\emph{bwoʃ}] 0xEC31: саботаж\\
\textbf{bwok} [\emph{bwok}] 0x2C31: брокколи\\
\textbf{bwom} [\emph{bwom}] 0x0C31: лоботомия\\
\textbf{bwon} [\emph{bwon}] 0x6C31: фасоль\\
\textbf{bwop} [\emph{bwop}] 0x4C31: Bow Wow\\
\textbf{bwos} [\emph{bwos}] 0x8C31: боксеров\\
\textbf{bwot} [\emph{bwot}] 0xAC31: бычий\\
\textbf{bwuk} [\emph{bwuk}] 0x2A31: букет\\
\textbf{bwum} [\emph{bwum}] 0x0A31: рубидий\\
\textbf{bwun} [\emph{bwun}] 0x6A31: Бутан\\
\textbf{bwus} [\emph{bwus}] 0x8A31: вопить\\
\textbf{bwut} [\emph{bwut}] 0xAA31: береза\\
\textbf{bx6m} [\emph{bxəm}] 0x0DF1: багамы\\
\textbf{bx6n} [\emph{bxən}] 0x6DF1: шоры\\
\textbf{bx6t} [\emph{bxət}] 0xADF1: смеситель\\
\textbf{bxa2n} [\emph{bxä˩n}] 0x79F1: бланшировать\\
\textbf{bxa2t} [\emph{bxä˩t}] 0xB9F1: повязка на глаза\\
\textbf{bxa7k} [\emph{bxä˥k}] 0x31F1: прицветник\\
\textbf{bxa7m} [\emph{bxä˥m}] 0x11F1: белая лошадь\\
\textbf{bxa7n} [\emph{bxä˥n}] 0x71F1: Абиссиния\\
\textbf{bxa7s} [\emph{bxä˥s}] 0x91F1: торговец табачными изделиями\\
\textbf{bxa7t} [\emph{bxä˥t}] 0xB1F1: Багдад\\
\textbf{bxac} [\emph{bxäʃ}] 0xE9F1: мол\\
\textbf{bxaf} [\emph{bxäf}] 0xC9F1: стропило\\
\textbf{bxah} [\emph{bxäʰ}] 0x4F8F: пресыщенный\\
\textbf{bxak} [\emph{bxäk}] 0x29F1: закладка\\
\textbf{bxam} [\emph{bxäm}] 0x09F1: бальзам\\
\textbf{bxan} [\emph{bxän}] 0x69F1: вылетающих\\
\textbf{bxap} [\emph{bxäp}] 0x49F1: перепел\\
\textbf{bxas} [\emph{bxäs}] 0x89F1: к юго-западу\\
\textbf{bxat} [\emph{bxät}] 0xA9F1: лысый\\
\textbf{bxe2n} [\emph{bxe˩n}] 0x7BF1: расположенный ниже\\
\textbf{bxe7t} [\emph{bxe˥t}] 0xB3F1: биометрия\\
\textbf{bxef} [\emph{bxef}] 0xCBF1: Босния герцеговина\\
\textbf{bxek} [\emph{bxek}] 0x2BF1: жуки\\
\textbf{bxem} [\emph{bxem}] 0x0BF1: Богемия\\
\textbf{bxen} [\emph{bxen}] 0x6BF1: навигация\\
\textbf{bxep} [\emph{bxep}] 0x4BF1: латыш\\
\textbf{bxes} [\emph{bxes}] 0x8BF1: удар слева\\
\textbf{bxet} [\emph{bxet}] 0xABF1: барометр\\
\textbf{bxi2n} [\emph{bxi˩n}] 0x78F1: брачный\\
\textbf{bxi2t} [\emph{bxi˩t}] 0xB8F1: концы света\\
\textbf{bxi7n} [\emph{bxi˥n}] 0x70F1: бронхит\\
\textbf{bxi7s} [\emph{bxi˥s}] 0x90F1: абхазия\\
\textbf{bxi7t} [\emph{bxi˥t}] 0xB0F1: бенди\\
\textbf{bxic} [\emph{bxiʃ}] 0xE8F1: грунтовый\\
\textbf{bxif} [\emph{bxif}] 0xC8F1: углевод\\
\textbf{bxih} [\emph{bxiʰ}] 0x478F: universe gender, 12th density life form, akin to spirit of cosmos, Sophia\\
\textbf{bxik} [\emph{bxik}] 0x28F1: обитать\\
\textbf{bxim} [\emph{bxim}] 0x08F1: кипятить на медленном огне\\
\textbf{bxin} [\emph{bxin}] 0x68F1: иммобилизация\\
\textbf{bxip} [\emph{bxip}] 0x48F1: битовая карта\\
\textbf{bxis} [\emph{bxis}] 0x88F1: арбитр\\
\textbf{bxit} [\emph{bxit}] 0xA8F1: Web хозяин\\
\textbf{bxo7n} [\emph{bxo˥n}] 0x74F1: румянец\\
\textbf{bxo7t} [\emph{bxo˥t}] 0xB4F1: закупоривать отверстие\\
\textbf{bxok} [\emph{bxok}] 0x2CF1: рококо\\
\textbf{bxom} [\emph{bxom}] 0x0CF1: окра\\
\textbf{bxon} [\emph{bxon}] 0x6CF1: беспрецедентный\\
\textbf{bxos} [\emph{bxos}] 0x8CF1: потребители\\
\textbf{bxot} [\emph{bxot}] 0xACF1: само контроль\\
\textbf{bxuk} [\emph{bxuk}] 0x2AF1: подпочва\\
\textbf{bxum} [\emph{bxum}] 0x0AF1: буддизм\\
\textbf{bxun} [\emph{bxun}] 0x6AF1: цветущий\\
\textbf{bxus} [\emph{bxus}] 0x8AF1: Узбекистан\\
\textbf{bxut} [\emph{bxut}] 0xAAF1: тупой\\
\textbf{by6n} [\emph{bjən}] 0x6D11: обратное словообразование\\
\textbf{by6t} [\emph{bjət}] 0xAD11: язь\\
\textbf{bya2n} [\emph{bjä˩n}] 0x7911: большой\\
\textbf{bya2t} [\emph{bjä˩t}] 0xB911: подстрекательстве\\
\textbf{bya7n} [\emph{bjä˥n}] 0x7111: бармен\\
\textbf{bya7s} [\emph{bjä˥s}] 0x9111: Аббасидов\\
\textbf{bya7t} [\emph{bjä˥t}] 0xB111: Овадия\\
\textbf{byac} [\emph{bjäʃ}] 0xE911: подвал\\
\textbf{byaf} [\emph{bjäf}] 0xC911: риф\\
\textbf{byah} [\emph{bjäʰ}] 0x488F: язык\\
\textbf{byak} [\emph{bjäk}] 0x2911: козел\\
\textbf{byam} [\emph{bjäm}] 0x0911: облучать\\
\textbf{byan} [\emph{bjän}] 0x6911: ванна\\
\textbf{byap} [\emph{bjäp}] 0x4911: боб\\
\textbf{byas} [\emph{bjäs}] 0x8911: вакцина\\
\textbf{byat} [\emph{bjät}] 0xA911: сожжен\\
\textbf{byef} [\emph{bjef}] 0xCB11: бордель\\
\textbf{byeh} [\emph{bjeʰ}] 0x588F: байт\\
\textbf{byek} [\emph{bjek}] 0x2B11: божья коровка\\
\textbf{byem} [\emph{bjem}] 0x0B11: хлебный мякиш\\
\textbf{byen} [\emph{bjen}] 0x6B11: катание на санях\\
\textbf{byep} [\emph{bjep}] 0x4B11: ЕЦБ\\
\textbf{byes} [\emph{bjes}] 0x8B11: восемьдесят шесть\\
\textbf{byet} [\emph{bjet}] 0xAB11: безработные\\
\textbf{byi2n} [\emph{bji˩n}] 0x7811: большой обратный порядок байт\\
\textbf{byi2t} [\emph{bji˩t}] 0xB811: козел отпущения\\
\textbf{byi7n} [\emph{bji˥n}] 0x7011: травля\\
\textbf{byi7t} [\emph{bji˥t}] 0xB011: аболиционисты\\
\textbf{byic} [\emph{bjiʃ}] 0xE811: бактерии\\
\textbf{byif} [\emph{bjif}] 0xC811: банкноты\\
\textbf{byih} [\emph{bjiʰ}] 0x408F: benedictive mood, for expressing blessings.\\
\textbf{byik} [\emph{bjik}] 0x2811: бойкотировать\\
\textbf{byim} [\emph{bjim}] 0x0811: моногамия\\
\textbf{byin} [\emph{bjin}] 0x6811: двоичный\\
\textbf{byip} [\emph{bjip}] 0x4811: библия\\
\textbf{byis} [\emph{bjis}] 0x8811: подсознательный\\
\textbf{byit} [\emph{bjit}] 0xA811: праздник\\
\textbf{byo7n} [\emph{bjo˥n}] 0x7411: бурсит большого пальца стопы\\
\textbf{byok} [\emph{bjok}] 0x2C11: биотехнология\\
\textbf{byom} [\emph{bjom}] 0x0C11: биом\\
\textbf{byon} [\emph{bjon}] 0x6C11: босния\\
\textbf{byop} [\emph{bjop}] 0x4C11: аэробный\\
\textbf{byos} [\emph{bjos}] 0x8C11: нетоксичен\\
\textbf{byot} [\emph{bjot}] 0xAC11: видео\\
\textbf{byu2t} [\emph{bju˩t}] 0xBA11: бермудский\\
\textbf{byu7n} [\emph{bju˥n}] 0x7211: бубонный\\
\textbf{byu7t} [\emph{bju˥t}] 0xB211: большая ступня\\
\textbf{byuc} [\emph{bjuʃ}] 0xEA11: бутил\\
\textbf{byuk} [\emph{bjuk}] 0x2A11: страшилище\\
\textbf{byum} [\emph{bjum}] 0x0A11: рекламирующий\\
\textbf{byun} [\emph{bjun}] 0x6A11: кланяясь\\
\textbf{byup} [\emph{bjup}] 0x4A11: гудок\\
\textbf{byus} [\emph{bjus}] 0x8A11: омовения\\
\textbf{byut} [\emph{bjut}] 0xAA11: сука\\
\textbf{bz6s} [\emph{bzəs}] 0x8D51: OSE\\
\textbf{bza7s} [\emph{bzä˥s}] 0x9151: базел\\
\textbf{bza7t} [\emph{bzä˥t}] 0xB151: обличительный\\
\textbf{bzac} [\emph{bzäʃ}] 0xE951: сокращать\\
\textbf{bzaf} [\emph{bzäf}] 0xC951: бойня\\
\textbf{bzah} [\emph{bzäʰ}] 0x4A8F: шутить\\
\textbf{bzak} [\emph{bzäk}] 0x2951: носок\\
\textbf{bzam} [\emph{bzäm}] 0x0951: namby pamby\\
\textbf{bzan} [\emph{bzän}] 0x6951: бунт\\
\textbf{bzap} [\emph{bzäp}] 0x4951: фагот\\
\textbf{bzas} [\emph{bzäs}] 0x8951: пресыщенный\\
\textbf{bzat} [\emph{bzät}] 0xA951: актер\\
\textbf{bzek} [\emph{bzek}] 0x2B51: ребенок мозг\\
\textbf{bzen} [\emph{bzen}] 0x6B51: бензол\\
\textbf{bzet} [\emph{bzet}] 0xAB51: баклажан\\
\textbf{bzic} [\emph{bziʃ}] 0xE851: Бразилия\\
\textbf{bzif} [\emph{bzif}] 0xC851: урбанизация\\
\textbf{bzih} [\emph{bziʰ}] 0x428F: произвольный\\
\textbf{bzik} [\emph{bzik}] 0x2851: физика\\
\textbf{bzim} [\emph{bzim}] 0x0851: дружелюбие\\
\textbf{bzin} [\emph{bzin}] 0x6851: знак\\
\textbf{bzis} [\emph{bzis}] 0x8851: виза\\
\textbf{bzit} [\emph{bzit}] 0xA851: паук\\
\textbf{bzoh} [\emph{bzoʰ}] 0x628F: OSE\\
\textbf{bzok} [\emph{bzok}] 0x2C51: бензойный\\
\textbf{bzon} [\emph{bzon}] 0x6C51: бизон\\
\textbf{bzos} [\emph{bzos}] 0x8C51: бозон\\
\textbf{bzot} [\emph{bzot}] 0xAC51: баллонах\\
\textbf{bzuc} [\emph{bzuʃ}] 0xEA51: огрубение\\
\textbf{bzun} [\emph{bzun}] 0x6A51: педантка\\
\textbf{bzut} [\emph{bzut}] 0xAA51: метель\\
\subsection{ c }
\textbf{ca} [\emph{ʃä}] 0x25BE: Здравствуйте\\
\textbf{ce} [\emph{ʃe}] 0x2DBE: под контекст\\
\textbf{ci} [\emph{ʃi}] 0x21BE: имя прилагательное\\
\textbf{ci7} [\emph{ʃi˥}] 0x41BE: светящийся интенсивность\\
\textbf{cl6k} [\emph{ʃlək}] 0x2D6D: чревный\\
\textbf{cl6n} [\emph{ʃlən}] 0x6D6D: роды\\
\textbf{cl6t} [\emph{ʃlət}] 0xAD6D: изолятор\\
\textbf{cla2k} [\emph{ʃlä˩k}] 0x396D: мыканица\\
\textbf{cla2m} [\emph{ʃlä˩m}] 0x196D: родственная душа\\
\textbf{cla2n} [\emph{ʃlä˩n}] 0x796D: дробовик\\
\textbf{cla2s} [\emph{ʃlä˩s}] 0x996D: колит\\
\textbf{cla2t} [\emph{ʃlä˩t}] 0xB96D: трудоспособный\\
\textbf{cla7c} [\emph{ʃlä˥ʃ}] 0xF16D: убить радость\\
\textbf{cla7k} [\emph{ʃlä˥k}] 0x316D: волчица\\
\textbf{cla7n} [\emph{ʃlä˥n}] 0x716D: проститутка\\
\textbf{cla7p} [\emph{ʃlä˥p}] 0x516D: мытье кожи\\
\textbf{cla7t} [\emph{ʃlä˥t}] 0xB16D: йота\\
\textbf{clac} [\emph{ʃläʃ}] 0xE96D: здоровье\\
\textbf{claf} [\emph{ʃläf}] 0xC96D: полка\\
\textbf{clah} [\emph{ʃläʰ}] 0x4B6F: Другие\\
\textbf{clak} [\emph{ʃläk}] 0x296D: сахар\\
\textbf{clam} [\emph{ʃläm}] 0x096D: позор\\
\textbf{clan} [\emph{ʃlän}] 0x696D: тысяча\\
\textbf{clap} [\emph{ʃläp}] 0x496D: привычный аспект\\
\textbf{clas} [\emph{ʃläs}] 0x896D: желтый\\
\textbf{clat} [\emph{ʃlät}] 0xA96D: дети\\
\textbf{cle2n} [\emph{ʃle˩n}] 0x7B6D: тысячелетний\\
\textbf{cle2t} [\emph{ʃle˩t}] 0xBB6D: шнурок для ботинок\\
\textbf{cle7t} [\emph{ʃle˥t}] 0xB36D: эполет\\
\textbf{clec} [\emph{ʃleʃ}] 0xEB6D: выхлопные газы\\
\textbf{clef} [\emph{ʃlef}] 0xCB6D: кресельный подъемник\\
\textbf{cleh} [\emph{ʃleʰ}] 0x5B6F: недавний мимо напряженный\\
\textbf{clek} [\emph{ʃlek}] 0x2B6D: утечка\\
\textbf{clem} [\emph{ʃlem}] 0x0B6D: селен\\
\textbf{clen} [\emph{ʃlen}] 0x6B6D: канал\\
\textbf{clep} [\emph{ʃlep}] 0x4B6D: дождь со снегом\\
\textbf{cles} [\emph{ʃles}] 0x8B6D: реализм\\
\textbf{clet} [\emph{ʃlet}] 0xAB6D: неполный\\
\textbf{cli2n} [\emph{ʃli˩n}] 0x786D: кляузничество\\
\textbf{cli2t} [\emph{ʃli˩t}] 0xB86D: промокашка\\
\textbf{cli7c} [\emph{ʃli˥ʃ}] 0xF06D: аллил\\
\textbf{cli7k} [\emph{ʃli˥k}] 0x306D: кириллица\\
\textbf{cli7n} [\emph{ʃli˥n}] 0x706D: халдейский\\
\textbf{cli7t} [\emph{ʃli˥t}] 0xB06D: беспризорник\\
\textbf{clic} [\emph{ʃliʃ}] 0xE86D: кольцо\\
\textbf{clif} [\emph{ʃlif}] 0xC86D: спираль\\
\textbf{clih} [\emph{ʃliʰ}] 0x436F: холодно\\
\textbf{clik} [\emph{ʃlik}] 0x286D: охота\\
\textbf{clim} [\emph{ʃlim}] 0x086D: специи\\
\textbf{clin} [\emph{ʃlin}] 0x686D: Генри\\
\textbf{clip} [\emph{ʃlip}] 0x486D: раскаяние\\
\textbf{clis} [\emph{ʃlis}] 0x886D: июль\\
\textbf{clit} [\emph{ʃlit}] 0xA86D: карман\\
\textbf{clo2t} [\emph{ʃlo˩t}] 0xBC6D: эсхатология\\
\textbf{cloc} [\emph{ʃloʃ}] 0xEC6D: морфологический\\
\textbf{clof} [\emph{ʃlof}] 0xCC6D: Чехословакией\\
\textbf{cloh} [\emph{ʃloʰ}] 0x636F: rational gender, any entity capable of reason, 3rd density life form or higher\\
\textbf{clok} [\emph{ʃlok}] 0x2C6D: шоколад\\
\textbf{clom} [\emph{ʃlom}] 0x0C6D: цитоплазма\\
\textbf{clon} [\emph{ʃlon}] 0x6C6D: онлайн\\
\textbf{clop} [\emph{ʃlop}] 0x4C6D: шалот\\
\textbf{clos} [\emph{ʃlos}] 0x8C6D: рыботорговец\\
\textbf{clot} [\emph{ʃlot}] 0xAC6D: Новостная рассылка\\
\textbf{clu2t} [\emph{ʃlu˩t}] 0xBA6D: абсолютный ноль\\
\textbf{cluc} [\emph{ʃluʃ}] 0xEA6D: хитрость\\
\textbf{cluh} [\emph{ʃluʰ}] 0x536F: уродливый\\
\textbf{cluk} [\emph{ʃluk}] 0x2A6D: уродливый\\
\textbf{clum} [\emph{ʃlum}] 0x0A6D: гелий\\
\textbf{clun} [\emph{ʃlun}] 0x6A6D: дуб\\
\textbf{clup} [\emph{ʃlup}] 0x4A6D: шелковица\\
\textbf{clus} [\emph{ʃlus}] 0x8A6D: веселый\\
\textbf{clut} [\emph{ʃlut}] 0xAA6D: молоток\\
\textbf{co} [\emph{ʃo}] 0x31BE: Социальное контекст\\
\textbf{cr6n} [\emph{ʃrən}] 0x6DCD: титрование\\
\textbf{cr6s} [\emph{ʃrəs}] 0x8DCD: ворсянка\\
\textbf{cr6t} [\emph{ʃrət}] 0xADCD: хиропрактики\\
\textbf{cra2k} [\emph{ʃrä˩k}] 0x39CD: сорокопут\\
\textbf{cra2n} [\emph{ʃrä˩n}] 0x79CD: кораблекрушение\\
\textbf{cra2s} [\emph{ʃrä˩s}] 0x99CD: маразм\\
\textbf{cra2t} [\emph{ʃrä˩t}] 0xB9CD: мученичество\\
\textbf{cra7k} [\emph{ʃrä˥k}] 0x31CD: львиный зев\\
\textbf{cra7m} [\emph{ʃrä˥m}] 0x11CD: креветка\\
\textbf{cra7n} [\emph{ʃrä˥n}] 0x71CD: детский сад\\
\textbf{cra7s} [\emph{ʃrä˥s}] 0x91CD: сперматозоиды\\
\textbf{cra7t} [\emph{ʃrä˥t}] 0xB1CD: запаздывание\\
\textbf{crac} [\emph{ʃräʃ}] 0xE9CD: филиал\\
\textbf{craf} [\emph{ʃräf}] 0xC9CD: телесный\\
\textbf{crah} [\emph{ʃräʰ}] 0x4E6F: считают,\\
\textbf{crak} [\emph{ʃräk}] 0x29CD: очарование\\
\textbf{cram} [\emph{ʃräm}] 0x09CD: лестный\\
\textbf{cran} [\emph{ʃrän}] 0x69CD: жертва\\
\textbf{crap} [\emph{ʃräp}] 0x49CD: ловушка\\
\textbf{cras} [\emph{ʃräs}] 0x89CD: Рубашка\\
\textbf{crat} [\emph{ʃrät}] 0xA9CD: суббота\\
\textbf{cre2n} [\emph{ʃre˩n}] 0x7BCD: длящийся четыре года\\
\textbf{cre2t} [\emph{ʃre˩t}] 0xBBCD: оркестровать\\
\textbf{cre7n} [\emph{ʃre˥n}] 0x73CD: в соответствии с ним\\
\textbf{cre7t} [\emph{ʃre˥t}] 0xB3CD: четырехгранный\\
\textbf{crec} [\emph{ʃreʃ}] 0xEBCD: шахматы\\
\textbf{cref} [\emph{ʃref}] 0xCBCD: шериф\\
\textbf{creh} [\emph{ʃreʰ}] 0x5E6F: Генри\\
\textbf{crek} [\emph{ʃrek}] 0x2BCD: предлагать\\
\textbf{crem} [\emph{ʃrem}] 0x0BCD: обдирать\\
\textbf{cren} [\emph{ʃren}] 0x6BCD: сарай\\
\textbf{crep} [\emph{ʃrep}] 0x4BCD: гамбургер\\
\textbf{cres} [\emph{ʃres}] 0x8BCD: розничная торговля\\
\textbf{cret} [\emph{ʃret}] 0xABCD: следы\\
\textbf{cri2n} [\emph{ʃri˩n}] 0x78CD: определитель\\
\textbf{cri2t} [\emph{ʃri˩t}] 0xB8CD: содружество\\
\textbf{cri7c} [\emph{ʃri˥ʃ}] 0xF0CD: Archy\\
\textbf{cri7k} [\emph{ʃri˥k}] 0x30CD: машина для отжимания белья\\
\textbf{cri7n} [\emph{ʃri˥n}] 0x70CD: продувка\\
\textbf{cri7s} [\emph{ʃri˥s}] 0x90CD: милитаризм\\
\textbf{cri7t} [\emph{ʃri˥t}] 0xB0CD: замкнутый круг\\
\textbf{cric} [\emph{ʃriʃ}] 0xE8CD: кров\\
\textbf{crif} [\emph{ʃrif}] 0xC8CD: застенчивость\\
\textbf{crih} [\emph{ʃriʰ}] 0x466F: видимый доказательный\\
\textbf{crik} [\emph{ʃrik}] 0x28CD: ад\\
\textbf{crim} [\emph{ʃrim}] 0x08CD: пирамида\\
\textbf{crin} [\emph{ʃrin}] 0x68CD: восхищение\\
\textbf{crip} [\emph{ʃrip}] 0x48CD: настоящим\\
\textbf{cris} [\emph{ʃris}] 0x88CD: наследник\\
\textbf{crit} [\emph{ʃrit}] 0xA8CD: заброшенный\\
\textbf{cro2t} [\emph{ʃro˩t}] 0xBCCD: короткие разрезы\\
\textbf{cro7n} [\emph{ʃro˥n}] 0x74CD: вычитаемое\\
\textbf{cro7t} [\emph{ʃro˥t}] 0xB4CD: сонная артерия\\
\textbf{croc} [\emph{ʃroʃ}] 0xECCD: ошибка\\
\textbf{crof} [\emph{ʃrof}] 0xCCCD: испольщик\\
\textbf{croh} [\emph{ʃroʰ}] 0x666F: страх\\
\textbf{crok} [\emph{ʃrok}] 0x2CCD: шоссе\\
\textbf{crom} [\emph{ʃrom}] 0x0CCD: сотрудничать\\
\textbf{cron} [\emph{ʃron}] 0x6CCD: трепет\\
\textbf{crop} [\emph{ʃrop}] 0x4CCD: водопад\\
\textbf{cros} [\emph{ʃros}] 0x8CCD: хор\\
\textbf{crot} [\emph{ʃrot}] 0xACCD: луг\\
\textbf{cruc} [\emph{ʃruʃ}] 0xEACD: разрушающий\\
\textbf{cruh} [\emph{ʃruʰ}] 0x566F: позор\\
\textbf{cruk} [\emph{ʃruk}] 0x2ACD: кролик\\
\textbf{crum} [\emph{ʃrum}] 0x0ACD: гормон\\
\textbf{crun} [\emph{ʃrun}] 0x6ACD: сокровище\\
\textbf{crup} [\emph{ʃrup}] 0x4ACD: перекрытие\\
\textbf{crus} [\emph{ʃrus}] 0x8ACD: иерархия\\
\textbf{crut} [\emph{ʃrut}] 0xAACD: пятница\\
\textbf{cu} [\emph{ʃu}] 0x29BE: conditional mood, for if then statements. The if clause comes before the `cu' the then clause comes after.\\
\textbf{cw6t} [\emph{ʃwət}] 0xAD2D: кристаллы\\
\textbf{cwa7s} [\emph{ʃwä˥s}] 0x912D: отродье\\
\textbf{cwac} [\emph{ʃwäʃ}] 0xE92D: плечо\\
\textbf{cwaf} [\emph{ʃwäf}] 0xC92D: конверт\\
\textbf{cwah} [\emph{ʃwäʰ}] 0x496F: после\\
\textbf{cwak} [\emph{ʃwäk}] 0x292D: зло\\
\textbf{cwam} [\emph{ʃwäm}] 0x092D: вена\\
\textbf{cwan} [\emph{ʃwän}] 0x692D: надеяться\\
\textbf{cwap} [\emph{ʃwäp}] 0x492D: шампанское\\
\textbf{cwas} [\emph{ʃwäs}] 0x892D: начать\\
\textbf{cwat} [\emph{ʃwät}] 0xA92D: воздух\\
\textbf{cwec} [\emph{ʃweʃ}] 0xEB2D: словесный\\
\textbf{cwef} [\emph{ʃwef}] 0xCB2D: кервель\\
\textbf{cweh} [\emph{ʃweʰ}] 0x596F: смиренный\\
\textbf{cwek} [\emph{ʃwek}] 0x2B2D: хоккей\\
\textbf{cwen} [\emph{ʃwen}] 0x6B2D: стадо\\
\textbf{cwep} [\emph{ʃwep}] 0x4B2D: зарплаты\\
\textbf{cwes} [\emph{ʃwes}] 0x8B2D: Шлезвиг Гольштейн\\
\textbf{cwet} [\emph{ʃwet}] 0xAB2D: привязь\\
\textbf{cwi2n} [\emph{ʃwi˩n}] 0x782D: зонирование\\
\textbf{cwi7n} [\emph{ʃwi˥n}] 0x702D: молодой лебедь\\
\textbf{cwic} [\emph{ʃwiʃ}] 0xE82D: тигр\\
\textbf{cwih} [\emph{ʃwiʰ}] 0x416F: надеяться\\
\textbf{cwik} [\emph{ʃwik}] 0x282D: величественный\\
\textbf{cwim} [\emph{ʃwim}] 0x082D: надежность\\
\textbf{cwin} [\emph{ʃwin}] 0x682D: верховой\\
\textbf{cwip} [\emph{ʃwip}] 0x482D: чип\\
\textbf{cwis} [\emph{ʃwis}] 0x882D: Четверг\\
\textbf{cwit} [\emph{ʃwit}] 0xA82D: живой\\
\textbf{cwoc} [\emph{ʃwoʃ}] 0xEC2D: хорошо\\
\textbf{cwoh} [\emph{ʃwoʰ}] 0x616F: одна и та же мероприятие маркер\\
\textbf{cwok} [\emph{ʃwok}] 0x2C2D: шок\\
\textbf{cwon} [\emph{ʃwon}] 0x6C2D: овальный\\
\textbf{cwos} [\emph{ʃwos}] 0x8C2D: челнок\\
\textbf{cwot} [\emph{ʃwot}] 0xAC2D: неволя\\
\textbf{cwuc} [\emph{ʃwuʃ}] 0xEA2D: Choo Choo\\
\textbf{cwuh} [\emph{ʃwuʰ}] 0x516F: шок\\
\textbf{cwuk} [\emph{ʃwuk}] 0x2A2D: Старый мир\\
\textbf{cwun} [\emph{ʃwun}] 0x6A2D: жевание\\
\textbf{cwup} [\emph{ʃwup}] 0x4A2D: шампунь\\
\textbf{cwut} [\emph{ʃwut}] 0xAA2D: перемирие\\
\textbf{cy6t} [\emph{ʃjət}] 0xAD0D: жучок\\
\textbf{cya2n} [\emph{ʃjä˩n}] 0x790D: свинка\\
\textbf{cya2t} [\emph{ʃjä˩t}] 0xB90D: половина мачты\\
\textbf{cya7c} [\emph{ʃjä˥ʃ}] 0xF10D: шах\\
\textbf{cya7k} [\emph{ʃjä˥k}] 0x310D: ночник\\
\textbf{cya7n} [\emph{ʃjä˥n}] 0x710D: обеднять\\
\textbf{cya7t} [\emph{ʃjä˥t}] 0xB10D: скамеечка для ног\\
\textbf{cyac} [\emph{ʃjäʃ}] 0xE90D: кровь\\
\textbf{cyaf} [\emph{ʃjäf}] 0xC90D: смех\\
\textbf{cyah} [\emph{ʃjäʰ}] 0x486F: человек\\
\textbf{cyak} [\emph{ʃjäk}] 0x290D: небо\\
\textbf{cyam} [\emph{ʃjäm}] 0x090D: лето\\
\textbf{cyan} [\emph{ʃjän}] 0x690D: мир\\
\textbf{cyap} [\emph{ʃjäp}] 0x490D: задаваться вопросом\\
\textbf{cyas} [\emph{ʃjäs}] 0x890D: медведь\\
\textbf{cyat} [\emph{ʃjät}] 0xA90D: видимый доказательный\\
\textbf{cyec} [\emph{ʃjeʃ}] 0xEB0D: чихать\\
\textbf{cyeh} [\emph{ʃjeʰ}] 0x586F: задаваться вопросом\\
\textbf{cyek} [\emph{ʃjek}] 0x2B0D: мешок\\
\textbf{cyem} [\emph{ʃjem}] 0x0B0D: химус\\
\textbf{cyen} [\emph{ʃjen}] 0x6B0D: сессия\\
\textbf{cyep} [\emph{ʃjep}] 0x4B0D: капризный ребенок\\
\textbf{cyes} [\emph{ʃjes}] 0x8B0D: киты\\
\textbf{cyet} [\emph{ʃjet}] 0xAB0D: пожалел\\
\textbf{cyi2n} [\emph{ʃji˩n}] 0x780D: циан\\
\textbf{cyi2p} [\emph{ʃji˩p}] 0x580D: шиит\\
\textbf{cyi2t} [\emph{ʃji˩t}] 0xB80D: гелиоцентрический\\
\textbf{cyi7k} [\emph{ʃji˥k}] 0x300D: облегающий\\
\textbf{cyi7n} [\emph{ʃji˥n}] 0x700D: китайский квартал\\
\textbf{cyi7t} [\emph{ʃji˥t}] 0xB00D: палочки для еды\\
\textbf{cyic} [\emph{ʃjiʃ}] 0xE80D: тщеславный\\
\textbf{cyif} [\emph{ʃjif}] 0xC80D: лук\\
\textbf{cyih} [\emph{ʃjiʰ}] 0x406F: совещательный настроение\\
\textbf{cyik} [\emph{ʃjik}] 0x280D: власть\\
\textbf{cyim} [\emph{ʃjim}] 0x080D: подозрительный\\
\textbf{cyin} [\emph{ʃjin}] 0x680D: ложь\\
\textbf{cyip} [\emph{ʃjip}] 0x480D: превосходящие\\
\textbf{cyis} [\emph{ʃjis}] 0x880D: стул\\
\textbf{cyit} [\emph{ʃjit}] 0xA80D: сердце\\
\textbf{cyo2h} [\emph{ʃjo˩ʰ}] 0xE06F: подозрительный\\
\textbf{cyoc} [\emph{ʃjoʃ}] 0xEC0D: испытывающий\\
\textbf{cyoh} [\emph{ʃjoʰ}] 0x606F: уменьшительный\\
\textbf{cyok} [\emph{ʃjok}] 0x2C0D: пожимание плечами\\
\textbf{cyom} [\emph{ʃjom}] 0x0C0D: майоран\\
\textbf{cyon} [\emph{ʃjon}] 0x6C0D: крошечный\\
\textbf{cyop} [\emph{ʃjop}] 0x4C0D: хобби\\
\textbf{cyos} [\emph{ʃjos}] 0x8C0D: суматоха\\
\textbf{cyot} [\emph{ʃjot}] 0xAC0D: Импортировать\\
\textbf{cyuc} [\emph{ʃjuʃ}] 0xEA0D: высокая температура\\
\textbf{cyuh} [\emph{ʃjuʰ}] 0x506F: сомнение\\
\textbf{cyuk} [\emph{ʃjuk}] 0x2A0D: привлечь\\
\textbf{cyum} [\emph{ʃjum}] 0x0A0D: дрожжи\\
\textbf{cyun} [\emph{ʃjun}] 0x6A0D: состояние\\
\textbf{cyup} [\emph{ʃjup}] 0x4A0D: ворчливый\\
\textbf{cyus} [\emph{ʃjus}] 0x8A0D: злорадствовать\\
\textbf{cyut} [\emph{ʃjut}] 0xAA0D: тень\\
\subsection{ d }
\textbf{d6} [\emph{də}] 0x367E: идиофон\\
\textbf{da} [\emph{dä}] 0x267E: adpositional case, indicates preceding stem is the case for this noun phrase.\\
\textbf{de} [\emph{de}] 0x2E7E: distributive case, indicates for-each the noun phrase, or per the noun phrase. iterative for loops.\\
\textbf{de7} [\emph{de˥}] 0x4E7E: де\\
\textbf{di} [\emph{di}] 0x227E: директива настроение\\
\textbf{di7} [\emph{di˥}] 0x427E: в ближайщем будущем\\
\textbf{dj6n} [\emph{dʒən}] 0x6DB3: бедность\\
\textbf{dj6t} [\emph{dʒət}] 0xADB3: сиккатив\\
\textbf{dja2n} [\emph{dʒä˩n}] 0x79B3: такса\\
\textbf{dja2t} [\emph{dʒä˩t}] 0xB9B3: адиабатический\\
\textbf{dja7n} [\emph{dʒä˥n}] 0x71B3: Ахан\\
\textbf{dja7t} [\emph{dʒä˥t}] 0xB1B3: смертельный удар\\
\textbf{djac} [\emph{dʒäʃ}] 0xE9B3: примириться\\
\textbf{djak} [\emph{dʒäk}] 0x29B3: разъем\\
\textbf{djam} [\emph{dʒäm}] 0x09B3: большой масштаб\\
\textbf{djan} [\emph{dʒän}] 0x69B3: гигант\\
\textbf{djap} [\emph{dʒäp}] 0x49B3: алгебра\\
\textbf{djas} [\emph{dʒäs}] 0x89B3: джаз\\
\textbf{djat} [\emph{dʒät}] 0xA9B3: неспособность\\
\textbf{djek} [\emph{dʒek}] 0x2BB3: еж\\
\textbf{djem} [\emph{dʒem}] 0x0BB3: отек\\
\textbf{djen} [\emph{dʒen}] 0x6BB3: Средиземное море\\
\textbf{djet} [\emph{dʒet}] 0xABB3: косметический ремонт\\
\textbf{dji2n} [\emph{dʒi˩n}] 0x78B3: предопределять\\
\textbf{dji2t} [\emph{dʒi˩t}] 0xB8B3: конец света\\
\textbf{dji7n} [\emph{dʒi˥n}] 0x70B3: управительница\\
\textbf{dji7s} [\emph{dʒi˥s}] 0x90B3: Изе\\
\textbf{dji7t} [\emph{dʒi˥t}] 0xB0B3: антидемократические\\
\textbf{djic} [\emph{dʒiʃ}] 0xE8B3: отвлекаться\\
\textbf{djih} [\emph{dʒiʰ}] 0x459F: симпатия\\
\textbf{djik} [\emph{dʒik}] 0x28B3: задерживаться\\
\textbf{djim} [\emph{dʒim}] 0x08B3: эндотермический\\
\textbf{djin} [\emph{dʒin}] 0x68B3: монета\\
\textbf{djip} [\emph{dʒip}] 0x48B3: липид\\
\textbf{djis} [\emph{dʒis}] 0x88B3: цифровой\\
\textbf{djit} [\emph{dʒit}] 0xA8B3: систематическое\\
\textbf{djok} [\emph{dʒok}] 0x2CB3: сопряженный\\
\textbf{djon} [\emph{dʒon}] 0x6CB3: по умолчанию\\
\textbf{djos} [\emph{dʒos}] 0x8CB3: союзы\\
\textbf{djot} [\emph{dʒot}] 0xACB3: осадок\\
\textbf{djun} [\emph{dʒun}] 0x6AB3: человечный\\
\textbf{djut} [\emph{dʒut}] 0xAAB3: догадка\\
\textbf{dl67t} [\emph{dlə˥t}] 0xB573: диплоид\\
\textbf{dl6k} [\emph{dlək}] 0x2D73: десятина\\
\textbf{dl6n} [\emph{dlən}] 0x6D73: Ирландия\\
\textbf{dl6t} [\emph{dlət}] 0xAD73: должным образом\\
\textbf{dla2k} [\emph{dlä˩k}] 0x3973: диалекты\\
\textbf{dla2n} [\emph{dlä˩n}] 0x7973: разлаживать\\
\textbf{dla2s} [\emph{dlä˩s}] 0x9973: деполяризация\\
\textbf{dla2t} [\emph{dlä˩t}] 0xB973: медлительный\\
\textbf{dla7k} [\emph{dlä˥k}] 0x3173: двухпросветный\\
\textbf{dla7m} [\emph{dlä˥m}] 0x1173: дифениламин\\
\textbf{dla7n} [\emph{dlä˥n}] 0x7173: след\\
\textbf{dla7p} [\emph{dlä˥p}] 0x5173: неровными краями\\
\textbf{dla7s} [\emph{dlä˥s}] 0x9173: напыщенность\\
\textbf{dla7t} [\emph{dlä˥t}] 0xB173: под ногами\\
\textbf{dlac} [\emph{dläʃ}] 0xE973: мыть\\
\textbf{dlaf} [\emph{dläf}] 0xC973: гранатовый\\
\textbf{dlah} [\emph{dläʰ}] 0x4B9F: наслаждаться\\
\textbf{dlak} [\emph{dläk}] 0x2973: обсуждать\\
\textbf{dlam} [\emph{dläm}] 0x0973: отдать честь\\
\textbf{dlan} [\emph{dlän}] 0x6973: доллар\\
\textbf{dlap} [\emph{dläp}] 0x4973: параллакс\\
\textbf{dlas} [\emph{dläs}] 0x8973: Удалить\\
\textbf{dlat} [\emph{dlät}] 0xA973: сегодняшний напряженный\\
\textbf{dle2n} [\emph{dle˩n}] 0x7B73: радиотелефон\\
\textbf{dle2t} [\emph{dle˩t}] 0xBB73: медлить хвостатых\\
\textbf{dle7t} [\emph{dle˥t}] 0xB373: мир устал\\
\textbf{dlec} [\emph{dleʃ}] 0xEB73: ноготь на пальце ноги\\
\textbf{dlef} [\emph{dlef}] 0xCB73: дельфийский\\
\textbf{dleh} [\emph{dleʰ}] 0x5B9F: hodiernal tense, for expressing today.\\
\textbf{dlek} [\emph{dlek}] 0x2B73: вектор\\
\textbf{dlem} [\emph{dlem}] 0x0B73: моделирование\\
\textbf{dlen} [\emph{dlen}] 0x6B73: документация\\
\textbf{dlep} [\emph{dlep}] 0x4B73: дверной звонок\\
\textbf{dles} [\emph{dles}] 0x8B73: гидролиз\\
\textbf{dlet} [\emph{dlet}] 0xAB73: каталог\\
\textbf{dli2k} [\emph{dli˩k}] 0x3873: георгин\\
\textbf{dli2n} [\emph{dli˩n}] 0x7873: деморализация\\
\textbf{dli2s} [\emph{dli˩s}] 0x9873: дуалистический\\
\textbf{dli2t} [\emph{dli˩t}] 0xB873: дипломат\\
\textbf{dli7k} [\emph{dli˥k}] 0x3073: деликт\\
\textbf{dli7n} [\emph{dli˥n}] 0x7073: неплатежеспособность\\
\textbf{dli7s} [\emph{dli˥s}] 0x9073: индол\\
\textbf{dli7t} [\emph{dli˥t}] 0xB073: Аделаида\\
\textbf{dlic} [\emph{dliʃ}] 0xE873: Идеально\\
\textbf{dlif} [\emph{dlif}] 0xC873: дельфин\\
\textbf{dlih} [\emph{dliʰ}] 0x439F: определенный статья\\
\textbf{dlik} [\emph{dlik}] 0x2873: законодательная власть\\
\textbf{dlim} [\emph{dlim}] 0x0873: идеология\\
\textbf{dlin} [\emph{dlin}] 0x6873: дисциплина\\
\textbf{dlip} [\emph{dlip}] 0x4873: сила воли\\
\textbf{dlis} [\emph{dlis}] 0x8873: идиосинкразия\\
\textbf{dlit} [\emph{dlit}] 0xA873: лист\\
\textbf{dlo2t} [\emph{dlo˩t}] 0xBC73: гидрохлорид\\
\textbf{dlo7n} [\emph{dlo˥n}] 0x7473: квадрильон\\
\textbf{dlo7t} [\emph{dlo˥t}] 0xB473: дуплет\\
\textbf{dloc} [\emph{dloʃ}] 0xEC73: Голландский\\
\textbf{dloh} [\emph{dloʰ}] 0x639F: доллар\\
\textbf{dlok} [\emph{dlok}] 0x2C73: роскошный\\
\textbf{dlom} [\emph{dlom}] 0x0C73: дипломатия\\
\textbf{dlon} [\emph{dlon}] 0x6C73: двойной номер\\
\textbf{dlop} [\emph{dlop}] 0x4C73: волнистый\\
\textbf{dlos} [\emph{dlos}] 0x8C73: модальный дело\\
\textbf{dlot} [\emph{dlot}] 0xAC73: скачать\\
\textbf{dlu7t} [\emph{dlu˥t}] 0xB273: друид\\
\textbf{dluc} [\emph{dluʃ}] 0xEA73: Андалусия\\
\textbf{dluh} [\emph{dluʰ}] 0x539F: двойной номер\\
\textbf{dluk} [\emph{dluk}] 0x2A73: Отключить\\
\textbf{dlun} [\emph{dlun}] 0x6A73: загрязнение\\
\textbf{dlus} [\emph{dlus}] 0x8A73: Беларусь\\
\textbf{dlut} [\emph{dlut}] 0xAA73: обозначать\\
\textbf{do} [\emph{do}] 0x327E: номер\\
\textbf{dr6c} [\emph{drəʃ}] 0xEDD3: квадратура\\
\textbf{dr6k} [\emph{drək}] 0x2DD3: гидропонный\\
\textbf{dr6m} [\emph{drəm}] 0x0DD3: дезинформация\\
\textbf{dr6n} [\emph{drən}] 0x6DD3: эндокринный\\
\textbf{dr6s} [\emph{drəs}] 0x8DD3: пристойный\\
\textbf{dr6t} [\emph{drət}] 0xADD3: дистрактор\\
\textbf{dra2c} [\emph{drä˩ʃ}] 0xF9D3: бесправных\\
\textbf{dra2k} [\emph{drä˩k}] 0x39D3: лишать членства\\
\textbf{dra2n} [\emph{drä˩n}] 0x79D3: старейшина\\
\textbf{dra2s} [\emph{drä˩s}] 0x99D3: изображенный\\
\textbf{dra2t} [\emph{drä˩t}] 0xB9D3: подъемный мост\\
\textbf{dra7k} [\emph{drä˥k}] 0x31D3: испачканный\\
\textbf{dra7m} [\emph{drä˥m}] 0x11D3: темная комната\\
\textbf{dra7n} [\emph{drä˥n}] 0x71D3: датчанин\\
\textbf{dra7s} [\emph{drä˥s}] 0x91D3: трема\\
\textbf{dra7t} [\emph{drä˥t}] 0xB1D3: обескуражена\\
\textbf{drac} [\emph{dräʃ}] 0xE9D3: речь\\
\textbf{draf} [\emph{dräf}] 0xC9D3: модернизировать\\
\textbf{drah} [\emph{dräʰ}] 0x4E9F: доволен\\
\textbf{drak} [\emph{dräk}] 0x29D3: расстройство\\
\textbf{dram} [\emph{dräm}] 0x09D3: алмаз\\
\textbf{dran} [\emph{drän}] 0x69D3: администрация\\
\textbf{drap} [\emph{dräp}] 0x49D3: завлечь\\
\textbf{dras} [\emph{dräs}] 0x89D3: презирать\\
\textbf{drat} [\emph{drät}] 0xA9D3: зависеть\\
\textbf{dre2n} [\emph{dre˩n}] 0x7BD3: затвердение\\
\textbf{dre2t} [\emph{dre˩t}] 0xBBD3: расхищение\\
\textbf{dre7n} [\emph{dre˥n}] 0x73D3: безрассудство делают\\
\textbf{dre7t} [\emph{dre˥t}] 0xB3D3: додекаэдр\\
\textbf{drec} [\emph{dreʃ}] 0xEBD3: рассечение\\
\textbf{dref} [\emph{dref}] 0xCBD3: гидросфера\\
\textbf{dreh} [\emph{dreʰ}] 0x5E9F: второй язык\\
\textbf{drek} [\emph{drek}] 0x2BD3: в отдельности\\
\textbf{drem} [\emph{drem}] 0x0BD3: Дания\\
\textbf{dren} [\emph{dren}] 0x6BD3: демонстрировать\\
\textbf{drep} [\emph{drep}] 0x4BD3: энтропия\\
\textbf{dres} [\emph{dres}] 0x8BD3: распаковывать\\
\textbf{dret} [\emph{dret}] 0xABD3: итерация\\
\textbf{dri2k} [\emph{dri˩k}] 0x38D3: дендрология\\
\textbf{dri2m} [\emph{dri˩m}] 0x18D3: гидротермальный\\
\textbf{dri2n} [\emph{dri˩n}] 0x78D3: угнетенный\\
\textbf{dri2s} [\emph{dri˩s}] 0x98D3: децентрализация\\
\textbf{dri2t} [\emph{dri˩t}] 0xB8D3: дезертир\\
\textbf{dri7k} [\emph{dri˥k}] 0x30D3: гидростатический\\
\textbf{dri7m} [\emph{dri˥m}] 0x10D3: двухцветный\\
\textbf{dri7n} [\emph{dri˥n}] 0x70D3: недостойный\\
\textbf{dri7p} [\emph{dri˥p}] 0x50D3: гидрофобный\\
\textbf{dri7s} [\emph{dri˥s}] 0x90D3: устрашение\\
\textbf{dri7t} [\emph{dri˥t}] 0xB0D3: песенка\\
\textbf{dric} [\emph{driʃ}] 0xE8D3: невидимый\\
\textbf{drif} [\emph{drif}] 0xC8D3: подводное крыло\\
\textbf{drih} [\emph{driʰ}] 0x469F: степень\\
\textbf{drik} [\emph{drik}] 0x28D3: странно\\
\textbf{drim} [\emph{drim}] 0x08D3: первокурсник\\
\textbf{drin} [\emph{drin}] 0x68D3: посвятить\\
\textbf{drip} [\emph{drip}] 0x48D3: капать\\
\textbf{dris} [\emph{dris}] 0x88D3: все больше и больше\\
\textbf{drit} [\emph{drit}] 0xA8D3: борода\\
\textbf{dro2n} [\emph{dro˩n}] 0x7CD3: глухота\\
\textbf{dro2t} [\emph{dro˩t}] 0xBCD3: низкооплачиваемый\\
\textbf{dro7k} [\emph{dro˥k}] 0x34D3: могущий быть оцененным\\
\textbf{dro7n} [\emph{dro˥n}] 0x74D3: солнечник\\
\textbf{dro7t} [\emph{dro˥t}] 0xB4D3: любить до безумия\\
\textbf{droc} [\emph{droʃ}] 0xECD3: докторская степень\\
\textbf{drof} [\emph{drof}] 0xCCD3: задергать\\
\textbf{droh} [\emph{droʰ}] 0x669F: слуховой доказательный\\
\textbf{drok} [\emph{drok}] 0x2CD3: штаб-квартира\\
\textbf{drom} [\emph{drom}] 0x0CD3: диморфизм\\
\textbf{dron} [\emph{dron}] 0x6CD3: Дракон\\
\textbf{drop} [\emph{drop}] 0x4CD3: углеводород\\
\textbf{dros} [\emph{dros}] 0x8CD3: дозировать\\
\textbf{drot} [\emph{drot}] 0xACD3: контролируемый\\
\textbf{dru2t} [\emph{dru˩t}] 0xBAD3: смущаясь\\
\textbf{dru7n} [\emph{dru˥n}] 0x72D3: слуховое окно\\
\textbf{dru7t} [\emph{dru˥t}] 0xB2D3: мочегонное средство\\
\textbf{druc} [\emph{druʃ}] 0xEAD3: UAE\\
\textbf{druk} [\emph{druk}] 0x2AD3: полукруглый\\
\textbf{drum} [\emph{drum}] 0x0AD3: незагрязненный\\
\textbf{drun} [\emph{drun}] 0x6AD3: неудовлетворенность\\
\textbf{drup} [\emph{drup}] 0x4AD3: индо Европейский\\
\textbf{drus} [\emph{drus}] 0x8AD3: Дрю\\
\textbf{drut} [\emph{drut}] 0xAAD3: описательный\\
\textbf{du} [\emph{du}] 0x2A7E: второй\\
\textbf{dv6n} [\emph{dvən}] 0x6D93: девонский\\
\textbf{dva2n} [\emph{dvä˩n}] 0x7993: расстройство мочеиспускания\\
\textbf{dva2t} [\emph{dvä˩t}] 0xB993: добавлений\\
\textbf{dva7n} [\emph{dvä˥n}] 0x7193: двухстороннее\\
\textbf{dvac} [\emph{dväʃ}] 0xE993: яичник\\
\textbf{dvaf} [\emph{dväf}] 0xC993: диван\\
\textbf{dvah} [\emph{dväʰ}] 0x4C9F: superessive case, indicates location on a surface. surface-context location-case.\\
\textbf{dvak} [\emph{dväk}] 0x2993: нигиляционными дело\\
\textbf{dvam} [\emph{dväm}] 0x0993: бессовестный\\
\textbf{dvan} [\emph{dvän}] 0x6993: неодобрительно\\
\textbf{dvap} [\emph{dväp}] 0x4993: половина на половину\\
\textbf{dvas} [\emph{dväs}] 0x8993: акушерка\\
\textbf{dvat} [\emph{dvät}] 0xA993: раздражительный\\
\textbf{dvec} [\emph{dveʃ}] 0xEB93: дервиш\\
\textbf{dvek} [\emph{dvek}] 0x2B93: обнаруживать\\
\textbf{dven} [\emph{dven}] 0x6B93: саванна\\
\textbf{dvet} [\emph{dvet}] 0xAB93: половик\\
\textbf{dvi7t} [\emph{dvi˥t}] 0xB093: дравидский\\
\textbf{dvic} [\emph{dviʃ}] 0xE893: двойственность\\
\textbf{dvif} [\emph{dvif}] 0xC893: скат\\
\textbf{dvih} [\emph{dviʰ}] 0x449F: конечный глагол\\
\textbf{dvik} [\emph{dvik}] 0x2893: адессив дело\\
\textbf{dvim} [\emph{dvim}] 0x0893: динамический\\
\textbf{dvin} [\emph{dvin}] 0x6893: домен имя\\
\textbf{dvip} [\emph{dvip}] 0x4893: рабыня\\
\textbf{dvis} [\emph{dvis}] 0x8893: подразделения\\
\textbf{dvit} [\emph{dvit}] 0xA893: чрезмерное продолжительность\\
\textbf{dvok} [\emph{dvok}] 0x2C93: владивосток\\
\textbf{dvon} [\emph{dvon}] 0x6C93: бубен\\
\textbf{dvot} [\emph{dvot}] 0xAC93: проводимость\\
\textbf{dvun} [\emph{dvun}] 0x6A93: нижеприведенный\\
\textbf{dvut} [\emph{dvut}] 0xAA93: выкорчевывать\\
\textbf{dw6n} [\emph{dwən}] 0x6D33: опреснение\\
\textbf{dwa2h} [\emph{dwä˩ʰ}] 0xC99F: равнодушие\\
\textbf{dwa2n} [\emph{dwä˩n}] 0x7933: бог король\\
\textbf{dwa2t} [\emph{dwä˩t}] 0xB933: сандаловое дерево\\
\textbf{dwa7n} [\emph{dwä˥n}] 0x7133: безумная\\
\textbf{dwa7t} [\emph{dwä˥t}] 0xB133: идиоматический\\
\textbf{dwac} [\emph{dwäʃ}] 0xE933: диалог\\
\textbf{dwaf} [\emph{dwäf}] 0xC933: водолаз\\
\textbf{dwah} [\emph{dwäʰ}] 0x499F: Да\\
\textbf{dwak} [\emph{dwäk}] 0x2933: срывать\\
\textbf{dwam} [\emph{dwäm}] 0x0933: Да\\
\textbf{dwan} [\emph{dwän}] 0x6933: вешать\\
\textbf{dwap} [\emph{dwäp}] 0x4933: адаптер\\
\textbf{dwas} [\emph{dwäs}] 0x8933: посольство\\
\textbf{dwat} [\emph{dwät}] 0xA933: лекарственное средство\\
\textbf{dwe7t} [\emph{dwe˥t}] 0xB333: детонатор\\
\textbf{dwec} [\emph{dweʃ}] 0xEB33: отвесно\\
\textbf{dweh} [\emph{dweʰ}] 0x599F: delimitive aspect, for referring to events of limited duration, like a limited time offer\\
\textbf{dwek} [\emph{dwek}] 0x2B33: гидроэлектрический\\
\textbf{dwen} [\emph{dwen}] 0x6B33: швейцар\\
\textbf{dwep} [\emph{dwep}] 0x4B33: диспепсия\\
\textbf{dwes} [\emph{dwes}] 0x8B33: вихри\\
\textbf{dwet} [\emph{dwet}] 0xAB33: со временем\\
\textbf{dwi7n} [\emph{dwi˥n}] 0x7033: фрейлин\\
\textbf{dwi7t} [\emph{dwi˥t}] 0xB033: раздвоенный хвост\\
\textbf{dwic} [\emph{dwiʃ}] 0xE833: отклонение\\
\textbf{dwif} [\emph{dwif}] 0xC833: деревянных духовых инструментов\\
\textbf{dwih} [\emph{dwiʰ}] 0x419F: desiderative mood, to indicate what the speaker wants to happen.\\
\textbf{dwik} [\emph{dwik}] 0x2833: дивиденд\\
\textbf{dwim} [\emph{dwim}] 0x0833: дефиле\\
\textbf{dwin} [\emph{dwin}] 0x6833: равнодушие\\
\textbf{dwip} [\emph{dwip}] 0x4833: децибел\\
\textbf{dwis} [\emph{dwis}] 0x8833: диоксид\\
\textbf{dwit} [\emph{dwit}] 0xA833: слуховой доказательный\\
\textbf{dwo7n} [\emph{dwo˥n}] 0x7433: допамин\\
\textbf{dwoc} [\emph{dwoʃ}] 0xEC33: дерматолог\\
\textbf{dwok} [\emph{dwok}] 0x2C33: радиология\\
\textbf{dwom} [\emph{dwom}] 0x0C33: недипломатичный\\
\textbf{dwon} [\emph{dwon}] 0x6C33: диагональ\\
\textbf{dwos} [\emph{dwos}] 0x8C33: кузовостроение\\
\textbf{dwot} [\emph{dwot}] 0xAC33: диктатор\\
\textbf{dwun} [\emph{dwun}] 0x6A33: морщина\\
\textbf{dwut} [\emph{dwut}] 0xAA33: документальный\\
\textbf{dx6k} [\emph{dxək}] 0x2DF3: потряхивание\\
\textbf{dx6n} [\emph{dxən}] 0x6DF3: приручение\\
\textbf{dx6t} [\emph{dxət}] 0xADF3: не согласен\\
\textbf{dxa2k} [\emph{dxä˩k}] 0x39F3: идиопатический\\
\textbf{dxa2n} [\emph{dxä˩n}] 0x79F3: безопасное поведение\\
\textbf{dxa2t} [\emph{dxä˩t}] 0xB9F3: дидактический\\
\textbf{dxa7k} [\emph{dxä˥k}] 0x31F3: демагогия\\
\textbf{dxa7n} [\emph{dxä˥n}] 0x71F3: середине воздуха\\
\textbf{dxa7s} [\emph{dxä˥s}] 0x91F3: лозоискательство\\
\textbf{dxa7t} [\emph{dxä˥t}] 0xB1F3: правая рука\\
\textbf{dxac} [\emph{dxäʃ}] 0xE9F3: терка\\
\textbf{dxaf} [\emph{dxäf}] 0xC9F3: лесной пожар\\
\textbf{dxah} [\emph{dxäʰ}] 0x4F9F: дистанционный пульт мимо напряженный\\
\textbf{dxak} [\emph{dxäk}] 0x29F3: данные добыча\\
\textbf{dxam} [\emph{dxäm}] 0x09F3: пищеварение\\
\textbf{dxan} [\emph{dxän}] 0x69F3: выгрузка\\
\textbf{dxap} [\emph{dxäp}] 0x49F3: компас\\
\textbf{dxas} [\emph{dxäs}] 0x89F3: одаренный\\
\textbf{dxat} [\emph{dxät}] 0xA9F3: опрометчивость\\
\textbf{dxe2t} [\emph{dxe˩t}] 0xBBF3: альдегид\\
\textbf{dxe7t} [\emph{dxe˥t}] 0xB3F3: дециметр\\
\textbf{dxec} [\emph{dxeʃ}] 0xEBF3: залогодержатель\\
\textbf{dxef} [\emph{dxef}] 0xCBF3: принаряжать\\
\textbf{dxeh} [\emph{dxeʰ}] 0x5F9F: другой мероприятие\\
\textbf{dxek} [\emph{dxek}] 0x2BF3: знакомства\\
\textbf{dxem} [\emph{dxem}] 0x0BF3: деспотизм\\
\textbf{dxen} [\emph{dxen}] 0x6BF3: Видео-конференция\\
\textbf{dxep} [\emph{dxep}] 0x4BF3: потрошить\\
\textbf{dxes} [\emph{dxes}] 0x8BF3: гастроном\\
\textbf{dxet} [\emph{dxet}] 0xABF3: другой мероприятие\\
\textbf{dxi2n} [\emph{dxi˩n}] 0x78F3: дикция\\
\textbf{dxi2t} [\emph{dxi˩t}] 0xB8F3: дифтерия\\
\textbf{dxi7n} [\emph{dxi˥n}] 0x70F3: двухпартийный\\
\textbf{dxi7s} [\emph{dxi˥s}] 0x90F3: IST\\
\textbf{dxi7t} [\emph{dxi˥t}] 0xB0F3: раскольник\\
\textbf{dxic} [\emph{dxiʃ}] 0xE8F3: Фиджи\\
\textbf{dxif} [\emph{dxif}] 0xC8F3: стрекоза\\
\textbf{dxih} [\emph{dxiʰ}] 0x479F: индуктивность\\
\textbf{dxik} [\emph{dxik}] 0x28F3: корица\\
\textbf{dxim} [\emph{dxim}] 0x08F3: домино\\
\textbf{dxin} [\emph{dxin}] 0x68F3: тело человека\\
\textbf{dxip} [\emph{dxip}] 0x48F3: сидячий\\
\textbf{dxis} [\emph{dxis}] 0x88F3: входящие\\
\textbf{dxit} [\emph{dxit}] 0xA8F3: драматично\\
\textbf{dxo2t} [\emph{dxo˩t}] 0xBCF3: недоделанный\\
\textbf{dxo7n} [\emph{dxo˥n}] 0x74F3: лженаука\\
\textbf{dxo7t} [\emph{dxo˥t}] 0xB4F3: без запаха\\
\textbf{dxoc} [\emph{dxoʃ}] 0xECF3: дезодорант\\
\textbf{dxof} [\emph{dxof}] 0xCCF3: Иегова\\
\textbf{dxoh} [\emph{dxoʰ}] 0x679F: modal case, indicates the ability or probability of the sentence. state-context location-case. neural weight.\\
\textbf{dxok} [\emph{dxok}] 0x2CF3: доминиканский\\
\textbf{dxom} [\emph{dxom}] 0x0CF3: одометр\\
\textbf{dxon} [\emph{dxon}] 0x6CF3: подгорный\\
\textbf{dxop} [\emph{dxop}] 0x4CF3: пыльная одежда\\
\textbf{dxos} [\emph{dxos}] 0x8CF3: магнетит\\
\textbf{dxot} [\emph{dxot}] 0xACF3: йогурт\\
\textbf{dxuk} [\emph{dxuk}] 0x2AF3: ломкий\\
\textbf{dxun} [\emph{dxun}] 0x6AF3: суточный\\
\textbf{dxup} [\emph{dxup}] 0x4AF3: дурная слава\\
\textbf{dxus} [\emph{dxus}] 0x8AF3: кукольный дом\\
\textbf{dxut} [\emph{dxut}] 0xAAF3: суфлер\\
\textbf{dy6n} [\emph{djən}] 0x6D13: оглушительный\\
\textbf{dy6t} [\emph{djət}] 0xAD13: диод\\
\textbf{dya2k} [\emph{djä˩k}] 0x3913: два упаковщик\\
\textbf{dya2n} [\emph{djä˩n}] 0x7913: диатонический\\
\textbf{dya2t} [\emph{djä˩t}] 0xB913: сердце раздирающие\\
\textbf{dya7k} [\emph{djä˥k}] 0x3113: идеограмма\\
\textbf{dya7m} [\emph{djä˥m}] 0x1113: диатомовый\\
\textbf{dya7n} [\emph{djä˥n}] 0x7113: адамантин\\
\textbf{dya7t} [\emph{djä˥t}] 0xB113: рассветало\\
\textbf{dyac} [\emph{djäʃ}] 0xE913: франшиза\\
\textbf{dyaf} [\emph{djäf}] 0xC913: отдаленное место\\
\textbf{dyah} [\emph{djäʰ}] 0x489F: admonitive mood, a request for permission to do something.\\
\textbf{dyak} [\emph{djäk}] 0x2913: танцор\\
\textbf{dyam} [\emph{djäm}] 0x0913: родий атом\\
\textbf{dyan} [\emph{djän}] 0x6913: яйцо\\
\textbf{dyap} [\emph{djäp}] 0x4913: второй язык\\
\textbf{dyas} [\emph{djäs}] 0x8913: дилер\\
\textbf{dyat} [\emph{djät}] 0xA913: приятель\\
\textbf{dye7t} [\emph{dje˥t}] 0xB313: де\\
\textbf{dyec} [\emph{djeʃ}] 0xEB13: денатурировать\\
\textbf{dyeh} [\emph{djeʰ}] 0x589F: непосредственный объект\\
\textbf{dyek} [\emph{djek}] 0x2B13: отлаживать\\
\textbf{dyem} [\emph{djem}] 0x0B13: окрашенных в шерсти\\
\textbf{dyen} [\emph{djen}] 0x6B13: индуктивность\\
\textbf{dyep} [\emph{djep}] 0x4B13: гантель\\
\textbf{dyes} [\emph{djes}] 0x8B13: диабет\\
\textbf{dyet} [\emph{djet}] 0xAB13: Мировой\\
\textbf{dyi2n} [\emph{dji˩n}] 0x7813: дифтонг\\
\textbf{dyi2t} [\emph{dji˩t}] 0xB813: адвентисты\\
\textbf{dyi7h} [\emph{dji˥ʰ}] 0x809F: дис\\
\textbf{dyi7n} [\emph{dji˥n}] 0x7013: два обрезной\\
\textbf{dyi7s} [\emph{dji˥s}] 0x9013: дис\\
\textbf{dyi7t} [\emph{dji˥t}] 0xB013: индуктивный\\
\textbf{dyic} [\emph{djiʃ}] 0xE813: образовательный\\
\textbf{dyif} [\emph{djif}] 0xC813: не доставлен\\
\textbf{dyih} [\emph{djiʰ}] 0x409F: ненавидеть\\
\textbf{dyik} [\emph{djik}] 0x2813: детектор\\
\textbf{dyim} [\emph{djim}] 0x0813: аудио\\
\textbf{dyin} [\emph{djin}] 0x6813: скрипка\\
\textbf{dyip} [\emph{djip}] 0x4813: непрямой речь\\
\textbf{dyis} [\emph{djis}] 0x8813: средневековый\\
\textbf{dyit} [\emph{djit}] 0xA813: промышленность\\
\textbf{dyo2n} [\emph{djo˩n}] 0x7C13: середине января\\
\textbf{dyo7n} [\emph{djo˥n}] 0x7413: индо иранское\\
\textbf{dyo7t} [\emph{djo˥t}] 0xB413: йодированной\\
\textbf{dyoc} [\emph{djoʃ}] 0xEC13: пособие по безработице\\
\textbf{dyoh} [\emph{djoʰ}] 0x609F: приспособление\\
\textbf{dyok} [\emph{djok}] 0x2C13: Доминика\\
\textbf{dyom} [\emph{djom}] 0x0C13: мажордом\\
\textbf{dyon} [\emph{djon}] 0x6C13: приспособление\\
\textbf{dyop} [\emph{djop}] 0x4C13: радиоуглерод\\
\textbf{dyos} [\emph{djos}] 0x8C13: динозавр\\
\textbf{dyot} [\emph{djot}] 0xAC13: радио\\
\textbf{dyuk} [\emph{djuk}] 0x2A13: педикюр\\
\textbf{dyun} [\emph{djun}] 0x6A13: диагностировать\\
\textbf{dyus} [\emph{djus}] 0x8A13: Будапешт\\
\textbf{dyut} [\emph{djut}] 0xAA13: скидка\\
\textbf{dza2n} [\emph{dzä˩n}] 0x7953: весь год\\
\textbf{dza2t} [\emph{dzä˩t}] 0xB953: обескураживать\\
\textbf{dza7n} [\emph{dzä˥n}] 0x7153: средний мозг\\
\textbf{dza7t} [\emph{dzä˥t}] 0xB153: диастола\\
\textbf{dzac} [\emph{dzäʃ}] 0xE953: обедающий\\
\textbf{dzak} [\emph{dzäk}] 0x2953: секта\\
\textbf{dzam} [\emph{dzäm}] 0x0953: астма\\
\textbf{dzan} [\emph{dzän}] 0x6953: наложить\\
\textbf{dzap} [\emph{dzäp}] 0x4953: драпировщик\\
\textbf{dzas} [\emph{dzäs}] 0x8953: вызывать недовольство\\
\textbf{dzat} [\emph{dzät}] 0xA953: дистанционный пульт мимо напряженный\\
\textbf{dzek} [\emph{dzek}] 0x2B53: декагон\\
\textbf{dzen} [\emph{dzen}] 0x6B53: двенадцатиперстная кишка\\
\textbf{dzes} [\emph{dzes}] 0x8B53: дизель\\
\textbf{dzet} [\emph{dzet}] 0xAB53: полупроводник\\
\textbf{dzi7n} [\emph{dzi˥n}] 0x7053: гедонистический\\
\textbf{dzi7s} [\emph{dzi˥s}] 0x9053: дионисическое\\
\textbf{dzi7t} [\emph{dzi˥t}] 0xB053: двугранный\\
\textbf{dzic} [\emph{dziʃ}] 0xE853: дизайнер\\
\textbf{dzih} [\emph{dziʰ}] 0x429F: одержимый дело\\
\textbf{dzik} [\emph{dzik}] 0x2853: двуязычный\\
\textbf{dzim} [\emph{dzim}] 0x0853: бегущая дорожка\\
\textbf{dzin} [\emph{dzin}] 0x6853: 19\\
\textbf{dzip} [\emph{dzip}] 0x4853: в середине декабря\\
\textbf{dzis} [\emph{dzis}] 0x8853: изморось\\
\textbf{dzit} [\emph{dzit}] 0xA853: СПИД\\
\textbf{dzoh} [\emph{dzoʰ}] 0x629F: подъязычная\\
\textbf{dzok} [\emph{dzok}] 0x2C53: сухой док\\
\textbf{dzon} [\emph{dzon}] 0x6C53: Амазонка\\
\textbf{dzos} [\emph{dzos}] 0x8C53: дремать\\
\textbf{dzot} [\emph{dzot}] 0xAC53: донор\\
\textbf{dzun} [\emph{dzun}] 0x6A53: Не знаю\\
\textbf{dzut} [\emph{dzut}] 0xAA53: растянутый\\
\subsection{ f }
\textbf{f6} [\emph{fə}] 0x359E: вставной морфема\\
\textbf{fa} [\emph{fä}] 0x259E: perfective aspect, a completed action taken as a whole.\\
\textbf{fe} [\emph{fe}] 0x2D9E: в конце концов\\
\textbf{fi} [\emph{fi}] 0x219E: вернуть\\
\textbf{fi7} [\emph{fi˥}] 0x419E: FY\\
\textbf{fl6k} [\emph{flək}] 0x2D6C: ихтиология\\
\textbf{fl6m} [\emph{fləm}] 0x0D6C: фатализм\\
\textbf{fl6n} [\emph{flən}] 0x6D6C: Финляндия\\
\textbf{fl6s} [\emph{fləs}] 0x8D6C: Фолклендские\\
\textbf{fl6t} [\emph{flət}] 0xAD6C: евклидов\\
\textbf{fla2k} [\emph{flä˩k}] 0x396C: фаллический\\
\textbf{fla2n} [\emph{flä˩n}] 0x796C: затопленный\\
\textbf{fla2s} [\emph{flä˩s}] 0x996C: анафилаксия\\
\textbf{fla2t} [\emph{flä˩t}] 0xB96C: формальдегид\\
\textbf{fla7k} [\emph{flä˥k}] 0x316C: фталевый\\
\textbf{fla7n} [\emph{flä˥n}] 0x716C: клык\\
\textbf{fla7s} [\emph{flä˥s}] 0x916C: фальцетом\\
\textbf{fla7t} [\emph{flä˥t}] 0xB16C: воспалительный\\
\textbf{flac} [\emph{fläʃ}] 0xE96C: вспышка\\
\textbf{flaf} [\emph{fläf}] 0xC96C: альфа\\
\textbf{flah} [\emph{fläʰ}] 0x4B67: контрфактическое условный\\
\textbf{flak} [\emph{fläk}] 0x296C: знакомые\\
\textbf{flam} [\emph{fläm}] 0x096C: воск\\
\textbf{flan} [\emph{flän}] 0x696C: наука\\
\textbf{flap} [\emph{fläp}] 0x496C: полевой шпат\\
\textbf{flas} [\emph{fläs}] 0x896C: поклонник\\
\textbf{flat} [\emph{flät}] 0xA96C: кислота\\
\textbf{fle2n} [\emph{fle˩n}] 0x7B6C: нафталин\\
\textbf{fle7k} [\emph{fle˥k}] 0x336C: передняя нога или лапа\\
\textbf{fle7n} [\emph{fle˥n}] 0x736C: этилен\\
\textbf{flec} [\emph{fleʃ}] 0xEB6C: фибрилляция\\
\textbf{flef} [\emph{flef}] 0xCB6C: Филадельфия\\
\textbf{fleh} [\emph{fleʰ}] 0x5B67: изменяемую\\
\textbf{flek} [\emph{flek}] 0x2B6C: подделывать\\
\textbf{flem} [\emph{flem}] 0x0B6C: флоэма\\
\textbf{flen} [\emph{flen}] 0x6B6C: сгибать\\
\textbf{fles} [\emph{fles}] 0x8B6C: аналитик\\
\textbf{flet} [\emph{flet}] 0xAB6C: изменяемую\\
\textbf{fli2c} [\emph{fli˩ʃ}] 0xF86C: френология\\
\textbf{fli2k} [\emph{fli˩k}] 0x386C: литология\\
\textbf{fli2n} [\emph{fli˩n}] 0x786C: филигрань\\
\textbf{fli2s} [\emph{fli˩s}] 0x986C: обывательский\\
\textbf{fli2t} [\emph{fli˩t}] 0xB86C: подстерегать\\
\textbf{fli7k} [\emph{fli˥k}] 0x306C: Девочка Мальчик\\
\textbf{fli7n} [\emph{fli˥n}] 0x706C: филология\\
\textbf{fli7p} [\emph{fli˥p}] 0x506C: Филипп\\
\textbf{fli7s} [\emph{fli˥s}] 0x906C: формализм\\
\textbf{fli7t} [\emph{fli˥t}] 0xB06C: теодолит\\
\textbf{flic} [\emph{fliʃ}] 0xE86C: теология\\
\textbf{flif} [\emph{flif}] 0xC86C: филе\\
\textbf{flih} [\emph{fliʰ}] 0x4367: выведенный доказательный\\
\textbf{flik} [\emph{flik}] 0x286C: лиса\\
\textbf{flim} [\emph{flim}] 0x086C: фильм\\
\textbf{flin} [\emph{flin}] 0x686C: выведенный доказательный\\
\textbf{flip} [\emph{flip}] 0x486C: Филиппины\\
\textbf{flis} [\emph{flis}] 0x886C: анализировать\\
\textbf{flit} [\emph{flit}] 0xA86C: поля\\
\textbf{flo2n} [\emph{flo˩n}] 0x7C6C: дефлорация\\
\textbf{flo7k} [\emph{flo˥k}] 0x346C: офтальмоскоп\\
\textbf{flo7n} [\emph{flo˥n}] 0x746C: афелий\\
\textbf{floc} [\emph{floʃ}] 0xEC6C: фольклор\\
\textbf{flok} [\emph{flok}] 0x2C6C: фламинго\\
\textbf{flom} [\emph{flom}] 0x0C6C: Филимон\\
\textbf{flon} [\emph{flon}] 0x6C6C: не в сети\\
\textbf{flop} [\emph{flop}] 0x4C6C: осциллограф\\
\textbf{flos} [\emph{flos}] 0x8C6C: флорист\\
\textbf{flot} [\emph{flot}] 0xAC6C: баловать\\
\textbf{flu2n} [\emph{flu˩n}] 0x7A6C: хлопьевидный\\
\textbf{fluc} [\emph{fluʃ}] 0xEA6C: беглый\\
\textbf{fluk} [\emph{fluk}] 0x2A6C: фолликул\\
\textbf{flum} [\emph{flum}] 0x0A6C: пушок\\
\textbf{flun} [\emph{flun}] 0x6A6C: фтор\\
\textbf{flup} [\emph{flup}] 0x4A6C: более полное\\
\textbf{flus} [\emph{flus}] 0x8A6C: кальцит\\
\textbf{flut} [\emph{flut}] 0xAA6C: поплавок\\
\textbf{fo} [\emph{fo}] 0x319E: окончательный частица\\
\textbf{fr62n} [\emph{frə˩n}] 0x7DCC: афферентный\\
\textbf{fr6h} [\emph{frəʰ}] 0x6E67: ferous\\
\textbf{fr6k} [\emph{frək}] 0x2DCC: трубка\\
\textbf{fr6n} [\emph{frən}] 0x6DCC: выветривание\\
\textbf{fr6s} [\emph{frəs}] 0x8DCC: ferous\\
\textbf{fr6t} [\emph{frət}] 0xADCC: одна треть\\
\textbf{fra2n} [\emph{frä˩n}] 0x79CC: пожарник\\
\textbf{fra2p} [\emph{frä˩p}] 0x59CC: импровизированная трибуна\\
\textbf{fra2t} [\emph{frä˩t}] 0xB9CC: братоубийство\\
\textbf{fra7c} [\emph{frä˥ʃ}] 0xF1CC: новая покрышка\\
\textbf{fra7n} [\emph{frä˥n}] 0x71CC: франк\\
\textbf{fra7s} [\emph{frä˥s}] 0x91CC: эвтаназия\\
\textbf{fra7t} [\emph{frä˥t}] 0xB1CC: Евфрат\\
\textbf{frac} [\emph{fräʃ}] 0xE9CC: бунтарь\\
\textbf{fraf} [\emph{fräf}] 0xC9CC: брандмауэр\\
\textbf{frah} [\emph{fräʰ}] 0x4E67: ранен\\
\textbf{frak} [\emph{fräk}] 0x29CC: неоднократно\\
\textbf{fram} [\emph{främ}] 0x09CC: аптека\\
\textbf{fran} [\emph{frän}] 0x69CC: летать\\
\textbf{frap} [\emph{fräp}] 0x49CC: федеральный\\
\textbf{fras} [\emph{fräs}] 0x89CC: разброс\\
\textbf{frat} [\emph{frät}] 0xA9CC: назначенный\\
\textbf{fre2m} [\emph{fre˩m}] 0x1BCC: Ефрем\\
\textbf{fre2n} [\emph{fre˩n}] 0x7BCC: волосной\\
\textbf{fre7s} [\emph{fre˥s}] 0x93CC: антифедерализм\\
\textbf{frec} [\emph{freʃ}] 0xEBCC: Французский\\
\textbf{freh} [\emph{freʰ}] 0x5E67: справочный\\
\textbf{frek} [\emph{frek}] 0x2BCC: пепельница\\
\textbf{frem} [\emph{frem}] 0x0BCC: восходящий поток теплого воздуха\\
\textbf{fren} [\emph{fren}] 0x6BCC: перфорировать\\
\textbf{frep} [\emph{frep}] 0x4BCC: щетка для волос\\
\textbf{fres} [\emph{fres}] 0x8BCC: мобилизовывать средства\\
\textbf{fret} [\emph{fret}] 0xABCC: искатель\\
\textbf{fri2n} [\emph{fri˩n}] 0x78CC: флирт\\
\textbf{fri2t} [\emph{fri˩t}] 0xB8CC: низкое качество\\
\textbf{fri7m} [\emph{fri˥m}] 0x10CC: афоризмы\\
\textbf{fri7n} [\emph{fri˥n}] 0x70CC: обезображивание\\
\textbf{fri7s} [\emph{fri˥s}] 0x90CC: теоретизировать\\
\textbf{fri7t} [\emph{fri˥t}] 0xB0CC: юрист\\
\textbf{fric} [\emph{friʃ}] 0xE8CC: отдых\\
\textbf{frif} [\emph{frif}] 0xC8CC: февраль\\
\textbf{frih} [\emph{friʰ}] 0x4667: знакомые регистр\\
\textbf{frik} [\emph{frik}] 0x28CC: теоретический\\
\textbf{frim} [\emph{frim}] 0x08CC: теорема\\
\textbf{frin} [\emph{frin}] 0x68CC: декорации\\
\textbf{frip} [\emph{frip}] 0x48CC: керамика\\
\textbf{fris} [\emph{fris}] 0x88CC: неофициальный\\
\textbf{frit} [\emph{frit}] 0xA8CC: порог\\
\textbf{fro2t} [\emph{fro˩t}] 0xBCCC: членистоногие\\
\textbf{fro7n} [\emph{fro˥n}] 0x74CC: формалин\\
\textbf{fro7s} [\emph{fro˥s}] 0x94CC: Фарерская\\
\textbf{fro7t} [\emph{fro˥t}] 0xB4CC: волокнистый\\
\textbf{froc} [\emph{froʃ}] 0xECCC: ковочный\\
\textbf{frof} [\emph{frof}] 0xCCCC: литография\\
\textbf{frok} [\emph{frok}] 0x2CCC: крутящий момент\\
\textbf{from} [\emph{from}] 0x0CCC: анахронизм\\
\textbf{fron} [\emph{fron}] 0x6CCC: героин\\
\textbf{frop} [\emph{frop}] 0x4CCC: хорек\\
\textbf{fros} [\emph{fros}] 0x8CCC: бирюзовый\\
\textbf{frot} [\emph{frot}] 0xACCC: инфраструктура\\
\textbf{fru2t} [\emph{fru˩t}] 0xBACC: фторид\\
\textbf{fruk} [\emph{fruk}] 0x2ACC: Африка\\
\textbf{frum} [\emph{frum}] 0x0ACC: фтор атом\\
\textbf{frun} [\emph{frun}] 0x6ACC: гниль\\
\textbf{frup} [\emph{frup}] 0x4ACC: опорос\\
\textbf{frus} [\emph{frus}] 0x8ACC: разъедающий\\
\textbf{frut} [\emph{frut}] 0xAACC: самостоятельно удовлетворяет\\
\textbf{fu} [\emph{fu}] 0x299E: фокус\\
\textbf{fwa2n} [\emph{fwä˩n}] 0x792C: астения\\
\textbf{fwa7t} [\emph{fwä˥t}] 0xB12C: Вопросы-Ответы\\
\textbf{fwac} [\emph{fwäʃ}] 0xE92C: аэропорт\\
\textbf{fwaf} [\emph{fwäf}] 0xC92C: плодоносить\\
\textbf{fwak} [\emph{fwäk}] 0x292C: судья\\
\textbf{fwam} [\emph{fwäm}] 0x092C: межсезонье\\
\textbf{fwan} [\emph{fwän}] 0x692C: посещение\\
\textbf{fwap} [\emph{fwäp}] 0x492C: фанера\\
\textbf{fwas} [\emph{fwäs}] 0x892C: фармацевт\\
\textbf{fwat} [\emph{fwät}] 0xA92C: больница\\
\textbf{fwek} [\emph{fwek}] 0x2B2C: фотоэлектрические\\
\textbf{fwen} [\emph{fwen}] 0x6B2C: филогенетическое\\
\textbf{fwes} [\emph{fwes}] 0x8B2C: диссертаций\\
\textbf{fwet} [\emph{fwet}] 0xAB2C: бесплатное программное обеспечение\\
\textbf{fwi7t} [\emph{fwi˥t}] 0xB02C: фаг\\
\textbf{fwic} [\emph{fwiʃ}] 0xE82C: фактически\\
\textbf{fwif} [\emph{fwif}] 0xC82C: FY\\
\textbf{fwih} [\emph{fwiʰ}] 0x4167: величественный\\
\textbf{fwik} [\emph{fwik}] 0x282C: салют\\
\textbf{fwin} [\emph{fwin}] 0x682C: частота\\
\textbf{fwip} [\emph{fwip}] 0x482C: расслаблять\\
\textbf{fwis} [\emph{fwis}] 0x882C: феминистка\\
\textbf{fwit} [\emph{fwit}] 0xA82C: утвердительный настроение\\
\textbf{fwok} [\emph{fwok}] 0x2C2C: офтальмология\\
\textbf{fwon} [\emph{fwon}] 0x6C2C: флуоресценция\\
\textbf{fwop} [\emph{fwop}] 0x4C2C: светобоязнь\\
\textbf{fwot} [\emph{fwot}] 0xAC2C: увядать\\
\textbf{fwun} [\emph{fwun}] 0x6A2C: сильфон\\
\textbf{fwut} [\emph{fwut}] 0xAA2C: вольт\\
\textbf{fy6n} [\emph{fjən}] 0x6D0C: Софония\\
\textbf{fy6s} [\emph{fjəs}] 0x8D0C: близкий к народу\\
\textbf{fy6t} [\emph{fjət}] 0xAD0C: сфинктер\\
\textbf{fya2n} [\emph{fjä˩n}] 0x790C: выявляя\\
\textbf{fya2t} [\emph{fjä˩t}] 0xB90C: совместное времяпровождение\\
\textbf{fya7n} [\emph{fjä˥n}] 0x710C: полугодовой\\
\textbf{fya7t} [\emph{fjä˥t}] 0xB10C: без даты\\
\textbf{fyac} [\emph{fjäʃ}] 0xE90C: железо\\
\textbf{fyah} [\emph{fjäʰ}] 0x4867: фраза маркер\\
\textbf{fyak} [\emph{fjäk}] 0x290C: фраза\\
\textbf{fyam} [\emph{fjäm}] 0x090C: фраза маркер\\
\textbf{fyan} [\emph{fjän}] 0x690C: сад\\
\textbf{fyap} [\emph{fjäp}] 0x490C: амфибия\\
\textbf{fyas} [\emph{fjäs}] 0x890C: ранен\\
\textbf{fyat} [\emph{fjät}] 0xA90C: суд\\
\textbf{fyek} [\emph{fjek}] 0x2B0C: фонетический\\
\textbf{fyen} [\emph{fjen}] 0x6B0C: фенхель\\
\textbf{fyes} [\emph{fjes}] 0x8B0C: Флер-де-Лис\\
\textbf{fyet} [\emph{fjet}] 0xAB0C: Фаренгейт\\
\textbf{fyi2n} [\emph{fji˩n}] 0x780C: день выплаты жалованья\\
\textbf{fyi2t} [\emph{fji˩t}] 0xB80C: тематический\\
\textbf{fyi7n} [\emph{fji˥n}] 0x700C: пятьдесят один\\
\textbf{fyi7s} [\emph{fji˥s}] 0x900C: не фантастика\\
\textbf{fyi7t} [\emph{fji˥t}] 0xB00C: фенил\\
\textbf{fyic} [\emph{fjiʃ}] 0xE80C: коэффициент полезного действия\\
\textbf{fyif} [\emph{fjif}] 0xC80C: возвратный голос\\
\textbf{fyih} [\emph{fjiʰ}] 0x4067: с плавающей точкой номер\\
\textbf{fyik} [\emph{fjik}] 0x280C: виски\\
\textbf{fyim} [\emph{fjim}] 0x080C: беженец\\
\textbf{fyin} [\emph{fjin}] 0x680C: расходы\\
\textbf{fyip} [\emph{fjip}] 0x480C: копировальное устройство\\
\textbf{fyis} [\emph{fjis}] 0x880C: Очистительный завод\\
\textbf{fyit} [\emph{fjit}] 0xA80C: знакомые регистр\\
\textbf{fyoh} [\emph{fjoʰ}] 0x6067: возвратный голос\\
\textbf{fyok} [\emph{fjok}] 0x2C0C: все органические\\
\textbf{fyom} [\emph{fjom}] 0x0C0C: опиум\\
\textbf{fyon} [\emph{fjon}] 0x6C0C: сродство\\
\textbf{fyop} [\emph{fjop}] 0x4C0C: орфоэпия\\
\textbf{fyos} [\emph{fjos}] 0x8C0C: этос\\
\textbf{fyot} [\emph{fjot}] 0xAC0C: сноска\\
\textbf{fyuk} [\emph{fjuk}] 0x2A0C: конфабуляция\\
\textbf{fyum} [\emph{fjum}] 0x0A0C: овсяница\\
\textbf{fyun} [\emph{fjun}] 0x6A0C: плавучесть\\
\textbf{fyut} [\emph{fjut}] 0xAA0C: привередливый\\
\subsection{ g }
\textbf{ga} [\emph{gä}] 0x265E: путь дело\\
\textbf{ga7} [\emph{gä˥}] 0x465E: графия\\
\textbf{ge} [\emph{ge}] 0x2E5E: агент\\
\textbf{gi} [\emph{gi}] 0x225E: имя\\
\textbf{gja7t} [\emph{gʒä˥t}] 0xB1B2: анти войны\\
\textbf{gjac} [\emph{gʒäʃ}] 0xE9B2: выкидыш\\
\textbf{gjaf} [\emph{gʒäf}] 0xC9B2: кастрировать\\
\textbf{gjah} [\emph{gʒäʰ}] 0x4D97: лестный\\
\textbf{gjak} [\emph{gʒäk}] 0x29B2: галактика\\
\textbf{gjam} [\emph{gʒäm}] 0x09B2: домогательство\\
\textbf{gjan} [\emph{gʒän}] 0x69B2: явление\\
\textbf{gjap} [\emph{gʒäp}] 0x49B2: спа\\
\textbf{gjas} [\emph{gʒäs}] 0x89B2: кактус\\
\textbf{gjat} [\emph{gʒät}] 0xA9B2: величина\\
\textbf{gjek} [\emph{gʒek}] 0x2BB2: еще - если\\
\textbf{gjen} [\emph{gʒen}] 0x6BB2: Аргентина\\
\textbf{gjet} [\emph{gʒet}] 0xABB2: электрод\\
\textbf{gji7t} [\emph{gʒi˥t}] 0xB0B2: Г.Д.\\
\textbf{gjic} [\emph{gʒiʃ}] 0xE8B2: Алжир\\
\textbf{gjih} [\emph{gʒiʰ}] 0x4597: galaxy gender, 11th density life form, akin to spirit of cosmos, Milky Way, Andromeda\\
\textbf{gjik} [\emph{gʒik}] 0x28B2: почка\\
\textbf{gjim} [\emph{gʒim}] 0x08B2: киска кошка\\
\textbf{gjin} [\emph{gʒin}] 0x68B2: день рождения\\
\textbf{gjip} [\emph{gʒip}] 0x48B2: джип\\
\textbf{gjis} [\emph{gʒis}] 0x88B2: джинсы\\
\textbf{gjit} [\emph{gʒit}] 0xA8B2: ножницы\\
\textbf{gjok} [\emph{gʒok}] 0x2CB2: догадок\\
\textbf{gjon} [\emph{gʒon}] 0x6CB2: однородный\\
\textbf{gjos} [\emph{gʒos}] 0x8CB2: галактоза\\
\textbf{gjot} [\emph{gʒot}] 0xACB2: оранжад\\
\textbf{gjun} [\emph{gʒun}] 0x6AB2: груминг\\
\textbf{gjut} [\emph{gʒut}] 0xAAB2: напыщенный\\
\textbf{gl6k} [\emph{glək}] 0x2D72: гинекология\\
\textbf{gl6n} [\emph{glən}] 0x6D72: стежка\\
\textbf{gl6t} [\emph{glət}] 0xAD72: палеонтолог\\
\textbf{gla2n} [\emph{glä˩n}] 0x7972: магнолия\\
\textbf{gla2t} [\emph{glä˩t}] 0xB972: гранд тетя\\
\textbf{gla7c} [\emph{glä˥ʃ}] 0xF172: галисийский\\
\textbf{gla7k} [\emph{glä˥k}] 0x3172: орленок\\
\textbf{gla7n} [\emph{glä˥n}] 0x7172: молочник\\
\textbf{gla7s} [\emph{glä˥s}] 0x9172: желчный камень\\
\textbf{gla7t} [\emph{glä˥t}] 0xB172: глыбы\\
\textbf{glac} [\emph{gläʃ}] 0xE972: жадность\\
\textbf{glaf} [\emph{gläf}] 0xC972: Югославия\\
\textbf{glah} [\emph{gläʰ}] 0x4B97: здоровье\\
\textbf{glak} [\emph{gläk}] 0x2972: вулкан\\
\textbf{glam} [\emph{gläm}] 0x0972: аукцион\\
\textbf{glan} [\emph{glän}] 0x6972: галерея\\
\textbf{glap} [\emph{gläp}] 0x4972: панель инструментов\\
\textbf{glas} [\emph{gläs}] 0x8972: симулировать\\
\textbf{glat} [\emph{glät}] 0xA972: простота\\
\textbf{gle2n} [\emph{gle˩n}] 0x7B72: мегаломания\\
\textbf{gle7t} [\emph{gle˥t}] 0xB372: буравчик глазами\\
\textbf{glec} [\emph{gleʃ}] 0xEB72: пеньюар\\
\textbf{glef} [\emph{glef}] 0xCB72: мегалит\\
\textbf{glek} [\emph{glek}] 0x2B72: реле\\
\textbf{glem} [\emph{glem}] 0x0B72: страдающий манией величия\\
\textbf{glen} [\emph{glen}] 0x6B72: Гренландия\\
\textbf{gles} [\emph{gles}] 0x8B72: шпалера\\
\textbf{glet} [\emph{glet}] 0xAB72: напряжение\\
\textbf{gli2n} [\emph{gli˩n}] 0x7872: облегчение\\
\textbf{gli2t} [\emph{gli˩t}] 0xB872: гильотина\\
\textbf{gli7k} [\emph{gli˥k}] 0x3072: иглу\\
\textbf{gli7n} [\emph{gli˥n}] 0x7072: продолжительность жизни\\
\textbf{gli7t} [\emph{gli˥t}] 0xB072: недвижимое имущество\\
\textbf{glic} [\emph{gliʃ}] 0xE872: ледник\\
\textbf{glif} [\emph{glif}] 0xC872: греческий\\
\textbf{glih} [\emph{gliʰ}] 0x4397: переполненный\\
\textbf{glik} [\emph{glik}] 0x2872: орошение\\
\textbf{glim} [\emph{glim}] 0x0872: галлия\\
\textbf{glin} [\emph{glin}] 0x6872: нежелание\\
\textbf{glip} [\emph{glip}] 0x4872: устрица\\
\textbf{glis} [\emph{glis}] 0x8872: бензин\\
\textbf{glit} [\emph{glit}] 0xA872: стабильность\\
\textbf{glo7k} [\emph{glo˥k}] 0x3472: глаголический\\
\textbf{gloc} [\emph{gloʃ}] 0xEC72: зоология\\
\textbf{glof} [\emph{glof}] 0xCC72: гольф\\
\textbf{glok} [\emph{glok}] 0x2C72: блог\\
\textbf{glom} [\emph{glom}] 0x0C72: Anglo Norman\\
\textbf{glon} [\emph{glon}] 0x6C72: галлон\\
\textbf{glop} [\emph{glop}] 0x4C72: Ванька встанька\\
\textbf{glos} [\emph{glos}] 0x8C72: гладиолус\\
\textbf{glot} [\emph{glot}] 0xAC72: мобильность\\
\textbf{glu7t} [\emph{glu˥t}] 0xB272: копытный\\
\textbf{gluc} [\emph{gluʃ}] 0xEA72: галоши\\
\textbf{gluh} [\emph{gluʰ}] 0x5397: singulative номер\\
\textbf{gluk} [\emph{gluk}] 0x2A72: пуля\\
\textbf{glun} [\emph{glun}] 0x6A72: чеснок\\
\textbf{glus} [\emph{glus}] 0x8A72: высокомерный\\
\textbf{glut} [\emph{glut}] 0xAA72: калькулятор\\
\textbf{go} [\emph{go}] 0x325E: gnomic aspect, for expressing a general truth.\\
\textbf{gr6n} [\emph{grən}] 0x6DD2: элерон\\
\textbf{gr6s} [\emph{grəs}] 0x8DD2: охлократия\\
\textbf{gr6t} [\emph{grət}] 0xADD2: контрацептив\\
\textbf{gra2n} [\emph{grä˩n}] 0x79D2: грыжа\\
\textbf{gra2t} [\emph{grä˩t}] 0xB9D2: лишать места\\
\textbf{gra7k} [\emph{grä˥k}] 0x31D2: Лошадиные силы\\
\textbf{gra7m} [\emph{grä˥m}] 0x11D2: Фогарма\\
\textbf{gra7n} [\emph{grä˥n}] 0x71D2: наградные\\
\textbf{gra7p} [\emph{grä˥p}] 0x51D2: мостовые весы\\
\textbf{gra7s} [\emph{grä˥s}] 0x91D2: гаррота\\
\textbf{gra7t} [\emph{grä˥t}] 0xB1D2: виноградная лоза\\
\textbf{grac} [\emph{gräʃ}] 0xE9D2: трусливый\\
\textbf{graf} [\emph{gräf}] 0xC9D2: алгоритм\\
\textbf{grah} [\emph{gräʰ}] 0x4E97: восхищение\\
\textbf{grak} [\emph{gräk}] 0x29D2: царапина\\
\textbf{gram} [\emph{gräm}] 0x09D2: рычание\\
\textbf{gran} [\emph{grän}] 0x69D2: зыбь\\
\textbf{grap} [\emph{gräp}] 0x49D2: захват\\
\textbf{gras} [\emph{gräs}] 0x89D2: транс\\
\textbf{grat} [\emph{grät}] 0xA9D2: арестовать\\
\textbf{gre2n} [\emph{gre˩n}] 0x7BD2: григорианский\\
\textbf{gre7n} [\emph{gre˥n}] 0x73D2: разгружать\\
\textbf{gre7t} [\emph{gre˥t}] 0xB3D2: нацарапал\\
\textbf{grec} [\emph{greʃ}] 0xEBD2: инфракрасный\\
\textbf{greh} [\emph{greʰ}] 0x5E97: centric case, indicates the noun phrase being the centre. Such as a star or actor.\\
\textbf{grek} [\emph{grek}] 0x2BD2: гараж\\
\textbf{grem} [\emph{grem}] 0x0BD2: германий\\
\textbf{gren} [\emph{gren}] 0x6BD2: планета\\
\textbf{grep} [\emph{grep}] 0x4BD2: гонорея\\
\textbf{gres} [\emph{gres}] 0x8BD2: кресс\\
\textbf{gret} [\emph{gret}] 0xABD2: вычитать\\
\textbf{gri2k} [\emph{gri˩k}] 0x38D2: аллергический\\
\textbf{gri2n} [\emph{gri˩n}] 0x78D2: глицерин\\
\textbf{gri2s} [\emph{gri˩s}] 0x98D2: гастрит\\
\textbf{gri2t} [\emph{gri˩t}] 0xB8D2: буравчик\\
\textbf{gri7k} [\emph{gri˥k}] 0x30D2: гликоль\\
\textbf{gri7n} [\emph{gri˥n}] 0x70D2: подстрочный\\
\textbf{gri7s} [\emph{gri˥s}] 0x90D2: взаимные услуги\\
\textbf{gri7t} [\emph{gri˥t}] 0xB0D2: шишковидный\\
\textbf{gric} [\emph{griʃ}] 0xE8D2: миграция\\
\textbf{grif} [\emph{grif}] 0xC8D2: партизанский\\
\textbf{grih} [\emph{griʰ}] 0x4697: серый\\
\textbf{grik} [\emph{grik}] 0x28D2: гравий\\
\textbf{grim} [\emph{grim}] 0x08D2: логарифм\\
\textbf{grin} [\emph{grin}] 0x68D2: окрестности\\
\textbf{grip} [\emph{grip}] 0x48D2: гриль\\
\textbf{gris} [\emph{gris}] 0x88D2: агрессивный\\
\textbf{grit} [\emph{grit}] 0xA8D2: гром\\
\textbf{gro2t} [\emph{gro˩t}] 0xBCD2: гигрометр\\
\textbf{gro7t} [\emph{gro˥t}] 0xB4D2: меди дном\\
\textbf{groc} [\emph{groʃ}] 0xECD2: Грузия\\
\textbf{groh} [\emph{groʰ}] 0x6697: Цель вызывать\\
\textbf{grok} [\emph{grok}] 0x2CD2: морковь\\
\textbf{grom} [\emph{grom}] 0x0CD2: роговица\\
\textbf{gron} [\emph{gron}] 0x6CD2: аргон\\
\textbf{grop} [\emph{grop}] 0x4CD2: моллюск из семейства брюхоногих\\
\textbf{gros} [\emph{gros}] 0x8CD2: плечевая кость\\
\textbf{grot} [\emph{grot}] 0xACD2: партнерство\\
\textbf{gruc} [\emph{gruʃ}] 0xEAD2: гуттаперча\\
\textbf{gruh} [\emph{gruʰ}] 0x5697: приседать\\
\textbf{gruk} [\emph{gruk}] 0x2AD2: Глубинные\\
\textbf{grum} [\emph{grum}] 0x0AD2: референдум\\
\textbf{grun} [\emph{grun}] 0x6AD2: хвастовство\\
\textbf{grup} [\emph{grup}] 0x4AD2: морской окунь\\
\textbf{grus} [\emph{grus}] 0x8AD2: урны\\
\textbf{grut} [\emph{grut}] 0xAAD2: инвертировать\\
\textbf{gu} [\emph{gu}] 0x2A5E: гм\\
\textbf{gva2t} [\emph{gvä˩t}] 0xB992: архивируются\\
\textbf{gva7t} [\emph{gvä˥t}] 0xB192: трехвалентный\\
\textbf{gvah} [\emph{gväʰ}] 0x4C97: lative case, indicates spatial motion towards the noun phrase. space-context destination-case. \\
\textbf{gvak} [\emph{gväk}] 0x2992: грамматическая случая\\
\textbf{gvan} [\emph{gvän}] 0x6992: злокачественный\\
\textbf{gvap} [\emph{gväp}] 0x4992: бокал для вина\\
\textbf{gvas} [\emph{gväs}] 0x8992: сорняками\\
\textbf{gvat} [\emph{gvät}] 0xA992: тяготение\\
\textbf{gveh} [\emph{gveʰ}] 0x5C97: эргатив дело\\
\textbf{gven} [\emph{gven}] 0x6B92: карниз\\
\textbf{gvet} [\emph{gvet}] 0xAB92: планер\\
\textbf{gvi2t} [\emph{gvi˩t}] 0xB892: циветта\\
\textbf{gvic} [\emph{gviʃ}] 0xE892: кровожадный\\
\textbf{gvih} [\emph{gviʰ}] 0x4497: pegative case, indicates motion from a person. social-context source-case.\\
\textbf{gvik} [\emph{gvik}] 0x2892: эргатив дело\\
\textbf{gvin} [\emph{gvin}] 0x6892: данный\\
\textbf{gvip} [\emph{gvip}] 0x4892: незавидный\\
\textbf{gvis} [\emph{gvis}] 0x8892: зерноядный\\
\textbf{gvit} [\emph{gvit}] 0xA892: упорство\\
\textbf{gvon} [\emph{gvon}] 0x6C92: стягивание\\
\textbf{gvot} [\emph{gvot}] 0xAC92: капли\\
\textbf{gvuk} [\emph{gvuk}] 0x2A92: вогнуто выпуклая\\
\textbf{gvun} [\emph{gvun}] 0x6A92: гальванизировать\\
\textbf{gvut} [\emph{gvut}] 0xAA92: глютен\\
\textbf{gwa7t} [\emph{gwä˥t}] 0xB132: Береговая охрана\\
\textbf{gwac} [\emph{gwäʃ}] 0xE932: консультант\\
\textbf{gwah} [\emph{gwäʰ}] 0x4997: нервный\\
\textbf{gwak} [\emph{gwäk}] 0x2932: лодыжка\\
\textbf{gwam} [\emph{gwäm}] 0x0932: гамета\\
\textbf{gwan} [\emph{gwän}] 0x6932: глотать\\
\textbf{gwap} [\emph{gwäp}] 0x4932: Габон\\
\textbf{gwas} [\emph{gwäs}] 0x8932: Вау\\
\textbf{gwat} [\emph{gwät}] 0xA932: слоновая кость\\
\textbf{gweh} [\emph{gweʰ}] 0x5997: скоро будущее напряженный\\
\textbf{gwen} [\emph{gwen}] 0x6B32: гранит\\
\textbf{gwet} [\emph{gwet}] 0xAB32: скоро будущее напряженный\\
\textbf{gwi7k} [\emph{gwi˥k}] 0x3032: урод\\
\textbf{gwi7n} [\emph{gwi˥n}] 0x7032: глубокая вода\\
\textbf{gwic} [\emph{gwiʃ}] 0xE832: переключатель\\
\textbf{gwif} [\emph{gwif}] 0xC832: глифы\\
\textbf{gwih} [\emph{gwiʰ}] 0x4197: невежественный\\
\textbf{gwik} [\emph{gwik}] 0x2832: циклический\\
\textbf{gwim} [\emph{gwim}] 0x0832: косоглазие\\
\textbf{gwin} [\emph{gwin}] 0x6832: свинг\\
\textbf{gwis} [\emph{gwis}] 0x8832: цитозин\\
\textbf{gwit} [\emph{gwit}] 0xA832: позвоночник\\
\textbf{gwok} [\emph{gwok}] 0x2C32: джокер\\
\textbf{gwon} [\emph{gwon}] 0x6C32: овес\\
\textbf{gwos} [\emph{gwos}] 0x8C32: гонады\\
\textbf{gwot} [\emph{gwot}] 0xAC32: агностик\\
\textbf{gwuk} [\emph{gwuk}] 0x2A32: глюкоза\\
\textbf{gwun} [\emph{gwun}] 0x6A32: гуанин\\
\textbf{gwus} [\emph{gwus}] 0x8A32: фагоцитоз\\
\textbf{gwut} [\emph{gwut}] 0xAA32: Джибути\\
\textbf{gx6n} [\emph{gxən}] 0x6DF2: Гана\\
\textbf{gxa2n} [\emph{gxä˩n}] 0x79F2: отмороженное место\\
\textbf{gxa2t} [\emph{gxä˩t}] 0xB9F2: цареубийство\\
\textbf{gxa7k} [\emph{gxä˥k}] 0x31F2: она попка\\
\textbf{gxa7n} [\emph{gxä˥n}] 0x71F2: сыворотка\\
\textbf{gxa7t} [\emph{gxä˥t}] 0xB1F2: перемешивает\\
\textbf{gxac} [\emph{gxäʃ}] 0xE9F2: инертность\\
\textbf{gxaf} [\emph{gxäf}] 0xC9F2: давление воздуха\\
\textbf{gxah} [\emph{gxäʰ}] 0x4F97: большой\\
\textbf{gxak} [\emph{gxäk}] 0x29F2: дорогостоящий\\
\textbf{gxam} [\emph{gxäm}] 0x09F2: Гватемала\\
\textbf{gxan} [\emph{gxän}] 0x69F2: галактика Пол\\
\textbf{gxap} [\emph{gxäp}] 0x49F2: гребешок\\
\textbf{gxas} [\emph{gxäs}] 0x89F2: маяк\\
\textbf{gxat} [\emph{gxät}] 0xA9F2: контр-атака\\
\textbf{gxe2t} [\emph{gxe˩t}] 0xBBF2: преуспевать\\
\textbf{gxe7t} [\emph{gxe˥t}] 0xB3F2: со-е изд\\
\textbf{gxec} [\emph{gxeʃ}] 0xEBF2: щегол\\
\textbf{gxek} [\emph{gxek}] 0x2BF2: высокие технологии\\
\textbf{gxem} [\emph{gxem}] 0x0BF2: мм хмм\\
\textbf{gxen} [\emph{gxen}] 0x6BF2: неодушевленный Пол\\
\textbf{gxep} [\emph{gxep}] 0x4BF2: Гринпис\\
\textbf{gxes} [\emph{gxes}] 0x8BF2: после чего\\
\textbf{gxet} [\emph{gxet}] 0xABF2: геометрия\\
\textbf{gxi2n} [\emph{gxi˩n}] 0x78F2: женоненавистник\\
\textbf{gxi2s} [\emph{gxi˩s}] 0x98F2: гри гри\\
\textbf{gxi2t} [\emph{gxi˩t}] 0xB8F2: агглютинативный\\
\textbf{gxi7n} [\emph{gxi˥n}] 0x70F2: глубокая тарелка\\
\textbf{gxi7s} [\emph{gxi˥s}] 0x90F2: гинеколог\\
\textbf{gxi7t} [\emph{gxi˥t}] 0xB0F2: двумужница\\
\textbf{gxic} [\emph{gxiʃ}] 0xE8F2: огненный\\
\textbf{gxif} [\emph{gxif}] 0xC8F2: графит\\
\textbf{gxih} [\emph{gxiʰ}] 0x4797: общий Пол\\
\textbf{gxik} [\emph{gxik}] 0x28F2: гражданские права\\
\textbf{gxim} [\emph{gxim}] 0x08F2: малахит\\
\textbf{gxin} [\emph{gxin}] 0x68F2: завоевывая\\
\textbf{gxip} [\emph{gxip}] 0x48F2: бицепс\\
\textbf{gxis} [\emph{gxis}] 0x88F2: колючий\\
\textbf{gxit} [\emph{gxit}] 0xA8F2: геноцид\\
\textbf{gxo7t} [\emph{gxo˥t}] 0xB4F2: совместно неавтоматического\\
\textbf{gxoc} [\emph{gxoʃ}] 0xECF2: горячий воздух\\
\textbf{gxok} [\emph{gxok}] 0x2CF2: готика\\
\textbf{gxom} [\emph{gxom}] 0x0CF2: экзогамия\\
\textbf{gxon} [\emph{gxon}] 0x6CF2: фрагментация\\
\textbf{gxop} [\emph{gxop}] 0x4CF2: коровье дерево\\
\textbf{gxos} [\emph{gxos}] 0x8CF2: органист\\
\textbf{gxot} [\emph{gxot}] 0xACF2: грецкий орех\\
\textbf{gxuk} [\emph{gxuk}] 0x2AF2: сумерки\\
\textbf{gxun} [\emph{gxun}] 0x6AF2: шестиугольник\\
\textbf{gxus} [\emph{gxus}] 0x8AF2: возможности\\
\textbf{gxut} [\emph{gxut}] 0xAAF2: Уганда\\
\textbf{gy6t} [\emph{gjət}] 0xAD12: Гренада\\
\textbf{gya2n} [\emph{gjä˩n}] 0x7912: земляной орех\\
\textbf{gya2t} [\emph{gjä˩t}] 0xB912: стекловидный глазами\\
\textbf{gya7n} [\emph{gjä˥n}] 0x7112: болеутоляющий\\
\textbf{gya7t} [\emph{gjä˥t}] 0xB112: трусики\\
\textbf{gyac} [\emph{gjäʃ}] 0xE912: ученый\\
\textbf{gyaf} [\emph{gjäf}] 0xC912: двор фермы\\
\textbf{gyah} [\emph{gjäʰ}] 0x4897: угол\\
\textbf{gyak} [\emph{gjäk}] 0x2912: средний класс\\
\textbf{gyam} [\emph{gjäm}] 0x0912: анус\\
\textbf{gyan} [\emph{gjän}] 0x6912: пот\\
\textbf{gyap} [\emph{gjäp}] 0x4912: Гамбия\\
\textbf{gyas} [\emph{gjäs}] 0x8912: гимнастика\\
\textbf{gyat} [\emph{gjät}] 0xA912: шпион\\
\textbf{gyek} [\emph{gjek}] 0x2B12: кинетика\\
\textbf{gyem} [\emph{gjem}] 0x0B12: магнитометр\\
\textbf{gyen} [\emph{gjen}] 0x6B12: ген\\
\textbf{gyes} [\emph{gjes}] 0x8B12: серебристо-серый\\
\textbf{gyet} [\emph{gjet}] 0xAB12: нефрит\\
\textbf{gyi7k} [\emph{gji˥k}] 0x3012: геодезический\\
\textbf{gyi7n} [\emph{gji˥n}] 0x7012: глицин\\
\textbf{gyi7t} [\emph{gji˥t}] 0xB012: идти между\\
\textbf{gyic} [\emph{gjiʃ}] 0xE812: гитара\\
\textbf{gyif} [\emph{gjif}] 0xC812: Гвинея\\
\textbf{gyih} [\emph{gjiʰ}] 0x4097: gnomic mood, for expressing an axiom.\\
\textbf{gyik} [\emph{gjik}] 0x2812: математика\\
\textbf{gyim} [\emph{gjim}] 0x0812: бурый уголь\\
\textbf{gyin} [\emph{gjin}] 0x6812: крыло\\
\textbf{gyip} [\emph{gjip}] 0x4812: золотая рыбка\\
\textbf{gyis} [\emph{gjis}] 0x8812: Греция\\
\textbf{gyit} [\emph{gjit}] 0xA812: умышленно\\
\textbf{gyoc} [\emph{gjoʃ}] 0xEC12: часовое дело\\
\textbf{gyok} [\emph{gjok}] 0x2C12: бег\\
\textbf{gyom} [\emph{gjom}] 0x0C12: агрономия\\
\textbf{gyon} [\emph{gjon}] 0x6C12: связь\\
\textbf{gyop} [\emph{gjop}] 0x4C12: спасательный буй или круг\\
\textbf{gyos} [\emph{gjos}] 0x8C12: глухой звук\\
\textbf{gyot} [\emph{gjot}] 0xAC12: условный\\
\textbf{gyun} [\emph{gjun}] 0x6A12: Гвиана\\
\textbf{gyus} [\emph{gjus}] 0x8A12: кучевые облака\\
\textbf{gyut} [\emph{gjut}] 0xAA12: оголять\\
\textbf{gzac} [\emph{gzäʃ}] 0xE952: воспламеняться\\
\textbf{gzaf} [\emph{gzäf}] 0xC952: газифицировать\\
\textbf{gzak} [\emph{gzäk}] 0x2952: эстрагон\\
\textbf{gzam} [\emph{gzäm}] 0x0952: оргазм\\
\textbf{gzan} [\emph{gzän}] 0x6952: гарантия\\
\textbf{gzap} [\emph{gzäp}] 0x4952: за бортом\\
\textbf{gzas} [\emph{gzäs}] 0x8952: богомол\\
\textbf{gzat} [\emph{gzät}] 0xA952: газета\\
\textbf{gzek} [\emph{gzek}] 0x2B52: эпизоотия\\
\textbf{gzen} [\emph{gzen}] 0x6B52: мыза\\
\textbf{gzes} [\emph{gzes}] 0x8B52: старикашка\\
\textbf{gzet} [\emph{gzet}] 0xAB52: опоясанный\\
\textbf{gzif} [\emph{gzif}] 0xC852: географический справочник\\
\textbf{gzih} [\emph{gziʰ}] 0x4297: переходный глагол\\
\textbf{gzik} [\emph{gzik}] 0x2852: каток\\
\textbf{gzin} [\emph{gzin}] 0x6852: пионер\\
\textbf{gzip} [\emph{gzip}] 0x4852: карданов подвес\\
\textbf{gzis} [\emph{gzis}] 0x8852: гейзер\\
\textbf{gzit} [\emph{gzit}] 0xA852: беременна\\
\textbf{gzok} [\emph{gzok}] 0x2C52: геофизика\\
\textbf{gzon} [\emph{gzon}] 0x6C52: казино\\
\textbf{gzot} [\emph{gzot}] 0xAC52: транзистор\\
\textbf{gzut} [\emph{gzut}] 0xAA52: взвести каракуля ду\\
\subsection{ h }
\textbf{hb67t} [\emph{ʰbə˥t}] 0xB588: придерживалась\\
\textbf{hb6k} [\emph{ʰbək}] 0x2D88: разрез`ать\\
\textbf{hb6m} [\emph{ʰbəm}] 0x0D88: брахман\\
\textbf{hb6n} [\emph{ʰbən}] 0x6D88: дистрибутивный местоимение\\
\textbf{hb6p} [\emph{ʰbəp}] 0x4D88: бампер\\
\textbf{hb6s} [\emph{ʰbəs}] 0x8D88: жужжание\\
\textbf{hb6t} [\emph{ʰbət}] 0xAD88: бренди\\
\textbf{hba2c} [\emph{ʰbä˩ʃ}] 0xF988: бакалавр\\
\textbf{hba2f} [\emph{ʰbä˩f}] 0xD988: байпас\\
\textbf{hba2k} [\emph{ʰbä˩k}] 0x3988: Бенгалия\\
\textbf{hba2m} [\emph{ʰbä˩m}] 0x1988: Бэтмен\\
\textbf{hba2n} [\emph{ʰbä˩n}] 0x7988: часто посещаемый\\
\textbf{hba2p} [\emph{ʰbä˩p}] 0x5988: анабаптист\\
\textbf{hba2s} [\emph{ʰbä˩s}] 0x9988: аббатиса\\
\textbf{hba2t} [\emph{ʰbä˩t}] 0xB988: группы\\
\textbf{hba7c} [\emph{ʰbä˥ʃ}] 0xF188: амбразура\\
\textbf{hba7f} [\emph{ʰbä˥f}] 0xD188: патронташ\\
\textbf{hba7k} [\emph{ʰbä˥k}] 0x3188: балканский\\
\textbf{hba7m} [\emph{ʰbä˥m}] 0x1188: незадачливый\\
\textbf{hba7n} [\emph{ʰbä˥n}] 0x7188: смело\\
\textbf{hba7p} [\emph{ʰbä˥p}] 0x5188: Бомбей\\
\textbf{hba7s} [\emph{ʰbä˥s}] 0x9188: рынки\\
\textbf{hba7t} [\emph{ʰbä˥t}] 0xB188: обязанный\\
\textbf{hbac} [\emph{ʰbäʃ}] 0xE988: двоюродная сестра\\
\textbf{hbaf} [\emph{ʰbäf}] 0xC988: биография\\
\textbf{hbak} [\emph{ʰbäk}] 0x2988: продавать\\
\textbf{hbam} [\emph{ʰbäm}] 0x0988: тетя\\
\textbf{hban} [\emph{ʰbän}] 0x6988: носить\\
\textbf{hbap} [\emph{ʰbäp}] 0x4988: прощальный привет\\
\textbf{hbas} [\emph{ʰbäs}] 0x8988: бред какой то\\
\textbf{hbat} [\emph{ʰbät}] 0xA988: слепой\\
\textbf{hbe2n} [\emph{ʰbe˩n}] 0x7B88: прижаться\\
\textbf{hbe2t} [\emph{ʰbe˩t}] 0xBB88: двор замка\\
\textbf{hbe7f} [\emph{ʰbe˥f}] 0xD388: хлебное дерево\\
\textbf{hbe7k} [\emph{ʰbe˥k}] 0x3388: брикет\\
\textbf{hbe7n} [\emph{ʰbe˥n}] 0x7388: бретонский\\
\textbf{hbe7s} [\emph{ʰbe˥s}] 0x9388: косички\\
\textbf{hbe7t} [\emph{ʰbe˥t}] 0xB388: утих\\
\textbf{hbec} [\emph{ʰbeʃ}] 0xEB88: благожелательный\\
\textbf{hbef} [\emph{ʰbef}] 0xCB88: зимняя спячка\\
\textbf{hbek} [\emph{ʰbek}] 0x2B88: пекарь\\
\textbf{hbem} [\emph{ʰbem}] 0x0B88: витамин\\
\textbf{hben} [\emph{ʰben}] 0x6B88: выборы\\
\textbf{hbep} [\emph{ʰbep}] 0x4B88: бейсбол\\
\textbf{hbes} [\emph{ʰbes}] 0x8B88: баннер\\
\textbf{hbet} [\emph{ʰbet}] 0xAB88: делать ставку\\
\textbf{hbi2c} [\emph{ʰbi˩ʃ}] 0xF888: опадение\\
\textbf{hbi2k} [\emph{ʰbi˩k}] 0x3888: било\\
\textbf{hbi2m} [\emph{ʰbi˩m}] 0x1888: балки\\
\textbf{hbi2n} [\emph{ʰbi˩n}] 0x7888: брат в закон\\
\textbf{hbi2s} [\emph{ʰbi˩s}] 0x9888: туземцы\\
\textbf{hbi2t} [\emph{ʰbi˩t}] 0xB888: двадцать два\\
\textbf{hbi7c} [\emph{ʰbi˥ʃ}] 0xF088: непреклонный\\
\textbf{hbi7k} [\emph{ʰbi˥k}] 0x3088: обшивать панелью\\
\textbf{hbi7m} [\emph{ʰbi˥m}] 0x1088: подружка невесты\\
\textbf{hbi7n} [\emph{ʰbi˥n}] 0x7088: Британия\\
\textbf{hbi7p} [\emph{ʰbi˥p}] 0x5088: Сивилла\\
\textbf{hbi7s} [\emph{ʰbi˥s}] 0x9088: крещение\\
\textbf{hbi7t} [\emph{ʰbi˥t}] 0xB088: биты\\
\textbf{hbic} [\emph{ʰbiʃ}] 0xE888: разорение\\
\textbf{hbif} [\emph{ʰbif}] 0xC888: лед\\
\textbf{hbik} [\emph{ʰbik}] 0x2888: благословить\\
\textbf{hbim} [\emph{ʰbim}] 0x0888: произвольный\\
\textbf{hbin} [\emph{ʰbin}] 0x6888: становление\\
\textbf{hbip} [\emph{ʰbip}] 0x4888: пиво\\
\textbf{hbis} [\emph{ʰbis}] 0x8888: база\\
\textbf{hbit} [\emph{ʰbit}] 0xA888: избежать\\
\textbf{hbo2n} [\emph{ʰbo˩n}] 0x7C88: взбодрить\\
\textbf{hbo2t} [\emph{ʰbo˩t}] 0xBC88: тупость\\
\textbf{hbo7k} [\emph{ʰbo˥k}] 0x3488: блокада\\
\textbf{hbo7m} [\emph{ʰbo˥m}] 0x1488: бергамот\\
\textbf{hbo7n} [\emph{ʰbo˥n}] 0x7488: кнопки\\
\textbf{hbo7p} [\emph{ʰbo˥p}] 0x5488: нажать кнопку\\
\textbf{hbo7s} [\emph{ʰbo˥s}] 0x9488: удав\\
\textbf{hbo7t} [\emph{ʰbo˥t}] 0xB488: комбинаторика\\
\textbf{hboc} [\emph{ʰboʃ}] 0xEC88: обязать\\
\textbf{hbof} [\emph{ʰbof}] 0xCC88: бомбардировка\\
\textbf{hbok} [\emph{ʰbok}] 0x2C88: утка\\
\textbf{hbom} [\emph{ʰbom}] 0x0C88: рвота\\
\textbf{hbon} [\emph{ʰbon}] 0x6C88: разнообразный\\
\textbf{hbop} [\emph{ʰbop}] 0x4C88: взломать\\
\textbf{hbos} [\emph{ʰbos}] 0x8C88: лаборатория\\
\textbf{hbot} [\emph{ʰbot}] 0xAC88: мяч\\
\textbf{hbu2n} [\emph{ʰbu˩n}] 0x7A88: бульон\\
\textbf{hbu2t} [\emph{ʰbu˩t}] 0xBA88: двуногое\\
\textbf{hbu7c} [\emph{ʰbu˥ʃ}] 0xF288: бушель\\
\textbf{hbu7n} [\emph{ʰbu˥n}] 0x7288: черное дерево\\
\textbf{hbu7p} [\emph{ʰbu˥p}] 0x5288: епископы\\
\textbf{hbu7t} [\emph{ʰbu˥t}] 0xB288: девяносто два\\
\textbf{hbuc} [\emph{ʰbuʃ}] 0xEA88: пользователь\\
\textbf{hbuf} [\emph{ʰbuf}] 0xCA88: вареный\\
\textbf{hbuk} [\emph{ʰbuk}] 0x2A88: бисквит\\
\textbf{hbum} [\emph{ʰbum}] 0x0A88: рубидий атом\\
\textbf{hbun} [\emph{ʰbun}] 0x6A88: путаница\\
\textbf{hbup} [\emph{ʰbup}] 0x4A88: бамбук\\
\textbf{hbus} [\emph{ʰbus}] 0x8A88: поцелуй\\
\textbf{hbut} [\emph{ʰbut}] 0xAA88: обувь\\
\textbf{hc67n} [\emph{ʰʃə˥n}] 0x7568: Снеговик\\
\textbf{hc67t} [\emph{ʰʃə˥t}] 0xB568: Sindhi\\
\textbf{hc6c} [\emph{ʰʃəʃ}] 0xED68: цикорий\\
\textbf{hc6f} [\emph{ʰʃəf}] 0xCD68: поверять\\
\textbf{hc6k} [\emph{ʰʃək}] 0x2D68: Чешский\\
\textbf{hc6m} [\emph{ʰʃəm}] 0x0D68: бальзамировать\\
\textbf{hc6n} [\emph{ʰʃən}] 0x6D68: очевидность\\
\textbf{hc6p} [\emph{ʰʃəp}] 0x4D68: проказливый\\
\textbf{hc6s} [\emph{ʰʃəs}] 0x8D68: паразитизм\\
\textbf{hc6t} [\emph{ʰʃət}] 0xAD68: светящийся интенсивность\\
\textbf{hca2c} [\emph{ʰʃä˩ʃ}] 0xF968: хорошо советовали\\
\textbf{hca2k} [\emph{ʰʃä˩k}] 0x3968: графство\\
\textbf{hca2m} [\emph{ʰʃä˩m}] 0x1968: алоза\\
\textbf{hca2n} [\emph{ʰʃä˩n}] 0x7968: сержант\\
\textbf{hca2p} [\emph{ʰʃä˩p}] 0x5968: шрапнель\\
\textbf{hca2s} [\emph{ʰʃä˩s}] 0x9968: канцелярия\\
\textbf{hca2t} [\emph{ʰʃä˩t}] 0xB968: упал\\
\textbf{hca7c} [\emph{ʰʃä˥ʃ}] 0xF168: перочинный нож\\
\textbf{hca7f} [\emph{ʰʃä˥f}] 0xD168: Шанхай\\
\textbf{hca7k} [\emph{ʰʃä˥k}] 0x3168: просыпался\\
\textbf{hca7m} [\emph{ʰʃä˥m}] 0x1168: заика\\
\textbf{hca7n} [\emph{ʰʃä˥n}] 0x7168: сказывают\\
\textbf{hca7p} [\emph{ʰʃä˥p}] 0x5168: шарлатан\\
\textbf{hca7s} [\emph{ʰʃä˥s}] 0x9168: телята\\
\textbf{hca7t} [\emph{ʰʃä˥t}] 0xB168: избавлены\\
\textbf{hcac} [\emph{ʰʃäʃ}] 0xE968: голодный\\
\textbf{hcaf} [\emph{ʰʃäf}] 0xC968: сомнение\\
\textbf{hcak} [\emph{ʰʃäk}] 0x2968: под контекст\\
\textbf{hcam} [\emph{ʰʃäm}] 0x0968: час\\
\textbf{hcan} [\emph{ʰʃän}] 0x6968: Здравствуйте\\
\textbf{hcap} [\emph{ʰʃäp}] 0x4968: после\\
\textbf{hcas} [\emph{ʰʃäs}] 0x8968: дом\\
\textbf{hcat} [\emph{ʰʃät}] 0xA968: близость\\
\textbf{hce2k} [\emph{ʰʃe˩k}] 0x3B68: чизкейк\\
\textbf{hce2n} [\emph{ʰʃe˩n}] 0x7B68: увещевание\\
\textbf{hce2t} [\emph{ʰʃe˩t}] 0xBB68: вымпел\\
\textbf{hce7k} [\emph{ʰʃe˥k}] 0x3368: шекель\\
\textbf{hce7n} [\emph{ʰʃe˥n}] 0x7368: китолов\\
\textbf{hce7s} [\emph{ʰʃe˥s}] 0x9368: тузы\\
\textbf{hce7t} [\emph{ʰʃe˥t}] 0xB368: жакеты\\
\textbf{hcec} [\emph{ʰʃeʃ}] 0xEB68: учение\\
\textbf{hcef} [\emph{ʰʃef}] 0xCB68: конечный глагол\\
\textbf{hcek} [\emph{ʰʃek}] 0x2B68: мебель\\
\textbf{hcem} [\emph{ʰʃem}] 0x0B68: землетрясение\\
\textbf{hcen} [\emph{ʰʃen}] 0x6B68: член\\
\textbf{hcep} [\emph{ʰʃep}] 0x4B68: клетка\\
\textbf{hces} [\emph{ʰʃes}] 0x8B68: студент\\
\textbf{hcet} [\emph{ʰʃet}] 0xAB68: ненавидеть\\
\textbf{hci2c} [\emph{ʰʃi˩ʃ}] 0xF868: Wishy водянистый\\
\textbf{hci2k} [\emph{ʰʃi˩k}] 0x3868: проницательно\\
\textbf{hci2n} [\emph{ʰʃi˩n}] 0x7868: снисходить\\
\textbf{hci2p} [\emph{ʰʃi˩p}] 0x5868: наложница\\
\textbf{hci2s} [\emph{ʰʃi˩s}] 0x9868: шедевр\\
\textbf{hci2t} [\emph{ʰʃi˩t}] 0xB868: районы\\
\textbf{hci7c} [\emph{ʰʃi˥ʃ}] 0xF068: визг\\
\textbf{hci7k} [\emph{ʰʃi˥k}] 0x3068: очевидец\\
\textbf{hci7m} [\emph{ʰʃi˥m}] 0x1068: Джимми\\
\textbf{hci7n} [\emph{ʰʃi˥n}] 0x7068: нищенство\\
\textbf{hci7p} [\emph{ʰʃi˥p}] 0x5068: пузырек\\
\textbf{hci7s} [\emph{ʰʃi˥s}] 0x9068: универсально\\
\textbf{hci7t} [\emph{ʰʃi˥t}] 0xB068: обещания\\
\textbf{hcic} [\emph{ʰʃiʃ}] 0xE868: вниз\\
\textbf{hcif} [\emph{ʰʃif}] 0xC868: кожа\\
\textbf{hcik} [\emph{ʰʃik}] 0x2868: думать\\
\textbf{hcim} [\emph{ʰʃim}] 0x0868: особенно\\
\textbf{hcin} [\emph{ʰʃin}] 0x6868: женский Пол\\
\textbf{hcip} [\emph{ʰʃip}] 0x4868: предполагаемый аспект\\
\textbf{hcis} [\emph{ʰʃis}] 0x8868: жена\\
\textbf{hcit} [\emph{ʰʃit}] 0xA868: имя прилагательное\\
\textbf{hco2k} [\emph{ʰʃo˩k}] 0x3C68: головорез\\
\textbf{hco2n} [\emph{ʰʃo˩n}] 0x7C68: до того времени\\
\textbf{hco2t} [\emph{ʰʃo˩t}] 0xBC68: домой\\
\textbf{hco7c} [\emph{ʰʃo˥ʃ}] 0xF468: зашикать\\
\textbf{hco7f} [\emph{ʰʃo˥f}] 0xD468: условное депонирование\\
\textbf{hco7k} [\emph{ʰʃo˥k}] 0x3468: межправительственная океанографическая комиссия\\
\textbf{hco7n} [\emph{ʰʃo˥n}] 0x7468: подчеркивание\\
\textbf{hco7p} [\emph{ʰʃo˥p}] 0x5468: впитывает\\
\textbf{hco7t} [\emph{ʰʃo˥t}] 0xB468: пришелец\\
\textbf{hcoc} [\emph{ʰʃoʃ}] 0xEC68: полый\\
\textbf{hcof} [\emph{ʰʃof}] 0xCC68: раздевать\\
\textbf{hcok} [\emph{ʰʃok}] 0x2C68: торопиться\\
\textbf{hcom} [\emph{ʰʃom}] 0x0C68: гармония\\
\textbf{hcon} [\emph{ʰʃon}] 0x6C68: кухня\\
\textbf{hcop} [\emph{ʰʃop}] 0x4C68: плита\\
\textbf{hcos} [\emph{ʰʃos}] 0x8C68: меры\\
\textbf{hcot} [\emph{ʰʃot}] 0xAC68: наслаждаться\\
\textbf{hcu2n} [\emph{ʰʃu˩n}] 0x7A68: недостойное поведение\\
\textbf{hcu2t} [\emph{ʰʃu˩t}] 0xBA68: у лодки\\
\textbf{hcu7f} [\emph{ʰʃu˥f}] 0xD268: сёгунат\\
\textbf{hcu7k} [\emph{ʰʃu˥k}] 0x3268: большеберцовая кость\\
\textbf{hcu7n} [\emph{ʰʃu˥n}] 0x7268: хунта\\
\textbf{hcu7t} [\emph{ʰʃu˥t}] 0xB268: насыщенный\\
\textbf{hcuc} [\emph{ʰʃuʃ}] 0xEA68: страх\\
\textbf{hcuf} [\emph{ʰʃuf}] 0xCA68: От себя\\
\textbf{hcuk} [\emph{ʰʃuk}] 0x2A68: фокус\\
\textbf{hcum} [\emph{ʰʃum}] 0x0A68: шум\\
\textbf{hcun} [\emph{ʰʃun}] 0x6A68: слышать\\
\textbf{hcup} [\emph{ʰʃup}] 0x4A68: зал\\
\textbf{hcus} [\emph{ʰʃus}] 0x8A68: последний\\
\textbf{hcut} [\emph{ʰʃut}] 0xAA68: условный настроение\\
\textbf{hd67k} [\emph{ʰdə˥k}] 0x3598: алкоголик\\
\textbf{hd67t} [\emph{ʰdə˥t}] 0xB598: дмитрий\\
\textbf{hd6c} [\emph{ʰdəʃ}] 0xED98: далматинец\\
\textbf{hd6k} [\emph{ʰdək}] 0x2D98: педиатрия\\
\textbf{hd6m} [\emph{ʰdəm}] 0x0D98: дихотомический\\
\textbf{hd6n} [\emph{ʰdən}] 0x6D98: знаменатель\\
\textbf{hd6p} [\emph{ʰdəp}] 0x4D98: подгузник\\
\textbf{hd6s} [\emph{ʰdəs}] 0x8D98: дьякон\\
\textbf{hd6t} [\emph{ʰdət}] 0xAD98: середине зимы\\
\textbf{hda2c} [\emph{ʰdä˩ʃ}] 0xF998: в середине марта\\
\textbf{hda2f} [\emph{ʰdä˩f}] 0xD998: крестный отец\\
\textbf{hda2k} [\emph{ʰdä˩k}] 0x3998: десятилетия\\
\textbf{hda2m} [\emph{ʰdä˩m}] 0x1998: догматизм\\
\textbf{hda2n} [\emph{ʰdä˩n}] 0x7998: края\\
\textbf{hda2p} [\emph{ʰdä˩p}] 0x5998: твердый и быстрый\\
\textbf{hda2s} [\emph{ʰdä˩s}] 0x9998: алый\\
\textbf{hda2t} [\emph{ʰdä˩t}] 0xB998: приводной вал\\
\textbf{hda7c} [\emph{ʰdä˥ʃ}] 0xF198: красильщик\\
\textbf{hda7f} [\emph{ʰdä˥f}] 0xD198: два четыре\\
\textbf{hda7k} [\emph{ʰdä˥k}] 0x3198: приоткрытый\\
\textbf{hda7m} [\emph{ʰdä˥m}] 0x1198: содомия\\
\textbf{hda7n} [\emph{ʰdä˥n}] 0x7198: водосточная труба\\
\textbf{hda7p} [\emph{ʰdä˥p}] 0x5198: саман\\
\textbf{hda7s} [\emph{ʰdä˥s}] 0x9198: кинжал\\
\textbf{hda7t} [\emph{ʰdä˥t}] 0xB198: легирующая примесь\\
\textbf{hdac} [\emph{ʰdäʃ}] 0xE998: адвокат\\
\textbf{hdaf} [\emph{ʰdäf}] 0xC998: достигать\\
\textbf{hdak} [\emph{ʰdäk}] 0x2998: магазин\\
\textbf{hdam} [\emph{ʰdäm}] 0x0998: приключение\\
\textbf{hdan} [\emph{ʰdän}] 0x6998: дьявол\\
\textbf{hdap} [\emph{ʰdäp}] 0x4998: 18\\
\textbf{hdas} [\emph{ʰdäs}] 0x8998: танец\\
\textbf{hdat} [\emph{ʰdät}] 0xA998: мост\\
\textbf{hde2m} [\emph{ʰde˩m}] 0x1B98: энтодерма\\
\textbf{hde2n} [\emph{ʰde˩n}] 0x7B98: среднего возраста\\
\textbf{hde2t} [\emph{ʰde˩t}] 0xBB98: божествами\\
\textbf{hde7k} [\emph{ʰde˥k}] 0x3398: дебушировать\\
\textbf{hde7m} [\emph{ʰde˥m}] 0x1398: демон\\
\textbf{hde7n} [\emph{ʰde˥n}] 0x7398: слышимый\\
\textbf{hde7s} [\emph{ʰde˥s}] 0x9398: сестра жены\\
\textbf{hde7t} [\emph{ʰde˥t}] 0xB398: одолжил\\
\textbf{hdec} [\emph{ʰdeʃ}] 0xEB98: К сожалению\\
\textbf{hdef} [\emph{ʰdef}] 0xCB98: дебет\\
\textbf{hdek} [\emph{ʰdek}] 0x2B98: светский\\
\textbf{hdem} [\emph{ʰdem}] 0x0B98: рабочий стол\\
\textbf{hden} [\emph{ʰden}] 0x6B98: решение\\
\textbf{hdep} [\emph{ʰdep}] 0x4B98: delimitative аспект\\
\textbf{hdes} [\emph{ʰdes}] 0x8B98: база данных\\
\textbf{hdet} [\emph{ʰdet}] 0xAB98: дефектный\\
\textbf{hdi2c} [\emph{ʰdi˩ʃ}] 0xF898: дизентерия\\
\textbf{hdi2k} [\emph{ʰdi˩k}] 0x3898: полукруг\\
\textbf{hdi2m} [\emph{ʰdi˩m}] 0x1898: деизм\\
\textbf{hdi2n} [\emph{ʰdi˩n}] 0x7898: доверчивость\\
\textbf{hdi2p} [\emph{ʰdi˩p}] 0x5898: несомненно\\
\textbf{hdi2s} [\emph{ʰdi˩s}] 0x9898: иды\\
\textbf{hdi2t} [\emph{ʰdi˩t}] 0xB898: гагачий\\
\textbf{hdi7c} [\emph{ʰdi˥ʃ}] 0xF098: выводить из заблуждения\\
\textbf{hdi7k} [\emph{ʰdi˥k}] 0x3098: приданое\\
\textbf{hdi7m} [\emph{ʰdi˥m}] 0x1098: идиомы\\
\textbf{hdi7n} [\emph{ʰdi˥n}] 0x7098: демпинг\\
\textbf{hdi7p} [\emph{ʰdi˥p}] 0x5098: Эдинбург\\
\textbf{hdi7s} [\emph{ʰdi˥s}] 0x9098: адреса\\
\textbf{hdi7t} [\emph{ʰdi˥t}] 0xB098: полночь\\
\textbf{hdic} [\emph{ʰdiʃ}] 0xE898: диск\\
\textbf{hdif} [\emph{ʰdif}] 0xC898: идентичность\\
\textbf{hdik} [\emph{ʰdik}] 0x2898: болезнь\\
\textbf{hdim} [\emph{ʰdim}] 0x0898: индуктивный настроение\\
\textbf{hdin} [\emph{ʰdin}] 0x6898: заполнить\\
\textbf{hdip} [\emph{ʰdip}] 0x4898: задерживается императив\\
\textbf{hdis} [\emph{ʰdis}] 0x8898: виновный\\
\textbf{hdit} [\emph{ʰdit}] 0xA898: благословение настроение\\
\textbf{hdo2k} [\emph{ʰdo˩k}] 0x3C98: двухтактный\\
\textbf{hdo2n} [\emph{ʰdo˩n}] 0x7C98: движения\\
\textbf{hdo2t} [\emph{ʰdo˩t}] 0xBC98: DonT\\
\textbf{hdo7k} [\emph{ʰdo˥k}] 0x3498: демагог\\
\textbf{hdo7m} [\emph{ʰdo˥m}] 0x1498: ареометр\\
\textbf{hdo7n} [\emph{ʰdo˥n}] 0x7498: несогласный\\
\textbf{hdo7s} [\emph{ʰdo˥s}] 0x9498: двойники\\
\textbf{hdo7t} [\emph{ʰdo˥t}] 0xB498: рассудите\\
\textbf{hdoc} [\emph{ʰdoʃ}] 0xEC98: ой\\
\textbf{hdof} [\emph{ʰdof}] 0xCC98: идиофон\\
\textbf{hdok} [\emph{ʰdok}] 0x2C98: вождение\\
\textbf{hdom} [\emph{ʰdom}] 0x0C98: кошмар\\
\textbf{hdon} [\emph{ʰdon}] 0x6C98: жертвовать\\
\textbf{hdop} [\emph{ʰdop}] 0x4C98: движение вверх по течению\\
\textbf{hdos} [\emph{ʰdos}] 0x8C98: пятно\\
\textbf{hdot} [\emph{ʰdot}] 0xAC98: дистальный демонстративный\\
\textbf{hdu2n} [\emph{ʰdu˩n}] 0x7A98: середина июня\\
\textbf{hdu2t} [\emph{ʰdu˩t}] 0xBA98: второй курс\\
\textbf{hdu7k} [\emph{ʰdu˥k}] 0x3298: совместное обучение\\
\textbf{hdu7n} [\emph{ʰdu˥n}] 0x7298: уборщица\\
\textbf{hdu7p} [\emph{ʰdu˥p}] 0x5298: Дунай\\
\textbf{hdu7s} [\emph{ʰdu˥s}] 0x9298: сочиться\\
\textbf{hdu7t} [\emph{ʰdu˥t}] 0xB298: оттуда\\
\textbf{hduc} [\emph{ʰduʃ}] 0xEA98: пыли\\
\textbf{hduf} [\emph{ʰduf}] 0xCA98: дуэль\\
\textbf{hduk} [\emph{ʰduk}] 0x2A98: крюк\\
\textbf{hdum} [\emph{ʰdum}] 0x0A98: иудаизм\\
\textbf{hdun} [\emph{ʰdun}] 0x6A98: единственное число\\
\textbf{hdup} [\emph{ʰdup}] 0x4A98: падение\\
\textbf{hdus} [\emph{ʰdus}] 0x8A98: свинья\\
\textbf{hdut} [\emph{ʰdut}] 0xAA98: молоко\\
\textbf{hf67n} [\emph{ʰfə˥n}] 0x7560: фенол\\
\textbf{hf6c} [\emph{ʰfəʃ}] 0xED60: Phenicia\\
\textbf{hf6f} [\emph{ʰfəf}] 0xCD60: сосиски\\
\textbf{hf6k} [\emph{ʰfək}] 0x2D60: глубина\\
\textbf{hf6m} [\emph{ʰfəm}] 0x0D60: однодневка\\
\textbf{hf6n} [\emph{ʰfən}] 0x6D60: Афганистан\\
\textbf{hf6s} [\emph{ʰfəs}] 0x8D60: фракталы\\
\textbf{hf6t} [\emph{ʰfət}] 0xAD60: фазан\\
\textbf{hfa2c} [\emph{ʰfä˩ʃ}] 0xF960: лихорадочный\\
\textbf{hfa2k} [\emph{ʰfä˩k}] 0x3960: альпака\\
\textbf{hfa2n} [\emph{ʰfä˩n}] 0x7960: фланель\\
\textbf{hfa2p} [\emph{ʰfä˩p}] 0x5960: ожирение\\
\textbf{hfa2s} [\emph{ʰfä˩s}] 0x9960: фаланга\\
\textbf{hfa2t} [\emph{ʰfä˩t}] 0xB960: феодальный\\
\textbf{hfa7c} [\emph{ʰfä˥ʃ}] 0xF160: франк tireur\\
\textbf{hfa7f} [\emph{ʰfä˥f}] 0xD160: пятерка\\
\textbf{hfa7k} [\emph{ʰfä˥k}] 0x3160: порка\\
\textbf{hfa7m} [\emph{ʰfä˥m}] 0x1160: фашизм\\
\textbf{hfa7n} [\emph{ʰfä˥n}] 0x7160: раздосадованный\\
\textbf{hfa7p} [\emph{ʰfä˥p}] 0x5160: полубак\\
\textbf{hfa7s} [\emph{ʰfä˥s}] 0x9160: ФЛАНДРИЯ\\
\textbf{hfa7t} [\emph{ʰfä˥t}] 0xB160: расставался\\
\textbf{hfac} [\emph{ʰfäʃ}] 0xE960: семья\\
\textbf{hfaf} [\emph{ʰfäf}] 0xC960: закон\\
\textbf{hfak} [\emph{ʰfäk}] 0x2960: 5\\
\textbf{hfam} [\emph{ʰfäm}] 0x0960: форма\\
\textbf{hfan} [\emph{ʰfän}] 0x6960: найденный\\
\textbf{hfap} [\emph{ʰfäp}] 0x4960: совершенный аспект\\
\textbf{hfas} [\emph{ʰfäs}] 0x8960: файл\\
\textbf{hfat} [\emph{ʰfät}] 0xA960: послал\\
\textbf{hfe2n} [\emph{ʰfe˩n}] 0x7B60: фаэтон\\
\textbf{hfe2s} [\emph{ʰfe˩s}] 0x9B60: франки стрелков\\
\textbf{hfe2t} [\emph{ʰfe˩t}] 0xBB60: фетиш\\
\textbf{hfe7n} [\emph{ʰfe˥n}] 0x7360: фехтовальщик\\
\textbf{hfe7s} [\emph{ʰfe˥s}] 0x9360: Еф\\
\textbf{hfe7t} [\emph{ʰfe˥t}] 0xB360: любовно\\
\textbf{hfec} [\emph{ʰfeʃ}] 0xEB60: обеспечение\\
\textbf{hfef} [\emph{ʰfef}] 0xCB60: взаимное развитие\\
\textbf{hfek} [\emph{ʰfek}] 0x2B60: процветать\\
\textbf{hfem} [\emph{ʰfem}] 0x0B60: вставной морфема\\
\textbf{hfen} [\emph{ʰfen}] 0x6B60: цент\\
\textbf{hfep} [\emph{ʰfep}] 0x4B60: многократный аспект\\
\textbf{hfes} [\emph{ʰfes}] 0x8B60: возвратный\\
\textbf{hfet} [\emph{ʰfet}] 0xAB60: противодействовать\\
\textbf{hfi2f} [\emph{ʰfi˩f}] 0xD860: фосфид\\
\textbf{hfi2k} [\emph{ʰfi˩k}] 0x3860: Феникс\\
\textbf{hfi2m} [\emph{ʰfi˩m}] 0x1860: амфетамин\\
\textbf{hfi2n} [\emph{ʰfi˩n}] 0x7860: безрассудная страсть\\
\textbf{hfi2p} [\emph{ʰfi˩p}] 0x5860: friborg\\
\textbf{hfi2s} [\emph{ʰfi˩s}] 0x9860: теософия\\
\textbf{hfi2t} [\emph{ʰfi˩t}] 0xB860: тридцать два\\
\textbf{hfi7c} [\emph{ʰfi˥ʃ}] 0xF060: афазия\\
\textbf{hfi7f} [\emph{ʰfi˥f}] 0xD060: тиофен\\
\textbf{hfi7k} [\emph{ʰfi˥k}] 0x3060: утомляет\\
\textbf{hfi7m} [\emph{ʰfi˥m}] 0x1060: эмфизема\\
\textbf{hfi7n} [\emph{ʰfi˥n}] 0x7060: финн\\
\textbf{hfi7p} [\emph{ʰfi˥p}] 0x5060: развертывание\\
\textbf{hfi7s} [\emph{ʰfi˥s}] 0x9060: юриспруденция\\
\textbf{hfi7t} [\emph{ʰfi˥t}] 0xB060: фотоэлектрический\\
\textbf{hfic} [\emph{ʰfiʃ}] 0xE860: влияние\\
\textbf{hfif} [\emph{ʰfif}] 0xC860: философ\\
\textbf{hfik} [\emph{ʰfik}] 0x2860: распределительный дело\\
\textbf{hfim} [\emph{ʰfim}] 0x0860: фермер\\
\textbf{hfin} [\emph{ʰfin}] 0x6860: богатые\\
\textbf{hfip} [\emph{ʰfip}] 0x4860: слабый\\
\textbf{hfis} [\emph{ʰfis}] 0x8860: отставка\\
\textbf{hfit} [\emph{ʰfit}] 0xA860: каузальный доказательный\\
\textbf{hfo2n} [\emph{ʰfo˩n}] 0x7C60: опровержение\\
\textbf{hfo2t} [\emph{ʰfo˩t}] 0xBC60: немодный\\
\textbf{hfo7k} [\emph{ʰfo˥k}] 0x3460: ортопедия\\
\textbf{hfo7n} [\emph{ʰfo˥n}] 0x7460: экспромтом\\
\textbf{hfo7s} [\emph{ʰfo˥s}] 0x9460: бесплатно для всех\\
\textbf{hfo7t} [\emph{ʰfo˥t}] 0xB460: фотон\\
\textbf{hfoc} [\emph{ʰfoʃ}] 0xEC60: фотография\\
\textbf{hfof} [\emph{ʰfof}] 0xCC60: фосфор\\
\textbf{hfok} [\emph{ʰfok}] 0x2C60: скромный\\
\textbf{hfom} [\emph{ʰfom}] 0x0C60: фонема\\
\textbf{hfon} [\emph{ʰfon}] 0x6C60: мед\\
\textbf{hfop} [\emph{ʰfop}] 0x4C60: футбол\\
\textbf{hfos} [\emph{ʰfos}] 0x8C60: факс\\
\textbf{hfot} [\emph{ʰfot}] 0xAC60: состоять\\
\textbf{hfu2n} [\emph{ʰfu˩n}] 0x7A60: само важно\\
\textbf{hfu7n} [\emph{ʰfu˥n}] 0x7260: листать\\
\textbf{hfu7t} [\emph{ʰfu˥t}] 0xB260: мышеловка\\
\textbf{hfuc} [\emph{ʰfuʃ}] 0xEA60: воскрешение\\
\textbf{hfuf} [\emph{ʰfuf}] 0xCA60: отруби\\
\textbf{hfuk} [\emph{ʰfuk}] 0x2A60: Огонь\\
\textbf{hfum} [\emph{ʰfum}] 0x0A60: пена\\
\textbf{hfun} [\emph{ʰfun}] 0x6A60: гнев\\
\textbf{hfup} [\emph{ʰfup}] 0x4A60: буря\\
\textbf{hfus} [\emph{ʰfus}] 0x8A60: восстановление\\
\textbf{hfut} [\emph{ʰfut}] 0xAA60: падать\\
\textbf{hg6c} [\emph{ʰgəʃ}] 0xED90: гейша\\
\textbf{hg6k} [\emph{ʰgək}] 0x2D90: полемический\\
\textbf{hg6m} [\emph{ʰgəm}] 0x0D90: микроорганизмы\\
\textbf{hg6n} [\emph{ʰgən}] 0x6D90: Японский\\
\textbf{hg6p} [\emph{ʰgəp}] 0x4D90: галопом\\
\textbf{hg6s} [\emph{ʰgəs}] 0x8D90: потуги\\
\textbf{hg6t} [\emph{ʰgət}] 0xAD90: педаль газа\\
\textbf{hga2c} [\emph{ʰgä˩ʃ}] 0xF990: бородка\\
\textbf{hga2k} [\emph{ʰgä˩k}] 0x3990: оглашение предстоящего бракосочетания\\
\textbf{hga2n} [\emph{ʰgä˩n}] 0x7990: строп\\
\textbf{hga2s} [\emph{ʰgä˩s}] 0x9990: гвардия\\
\textbf{hga2t} [\emph{ʰgä˩t}] 0xB990: полученный\\
\textbf{hga7c} [\emph{ʰgä˥ʃ}] 0xF190: ликование\\
\textbf{hga7f} [\emph{ʰgä˥f}] 0xD190: графия\\
\textbf{hga7k} [\emph{ʰgä˥k}] 0x3190: неловкость\\
\textbf{hga7m} [\emph{ʰgä˥m}] 0x1190: магнезия\\
\textbf{hga7n} [\emph{ʰgä˥n}] 0x7190: безумие\\
\textbf{hga7p} [\emph{ʰgä˥p}] 0x5190: водобоязнь\\
\textbf{hga7s} [\emph{ʰgä˥s}] 0x9190: певцы\\
\textbf{hga7t} [\emph{ʰgä˥t}] 0xB190: направляющие\\
\textbf{hgac} [\emph{ʰgäʃ}] 0xE990: амбиция\\
\textbf{hgaf} [\emph{ʰgäf}] 0xC990: грамматика\\
\textbf{hgak} [\emph{ʰgäk}] 0x2990: Кот\\
\textbf{hgam} [\emph{ʰgäm}] 0x0990: ветчина\\
\textbf{hgan} [\emph{ʰgän}] 0x6990: скрывать\\
\textbf{hgap} [\emph{ʰgäp}] 0x4990: корова\\
\textbf{hgas} [\emph{ʰgäs}] 0x8990: газ\\
\textbf{hgat} [\emph{ʰgät}] 0xA990: пустыня\\
\textbf{hge2n} [\emph{ʰge˩n}] 0x7B90: бог ребенок\\
\textbf{hge2t} [\emph{ʰge˩t}] 0xBB90: Гентский\\
\textbf{hge7k} [\emph{ʰge˥k}] 0x3390: экзотерический\\
\textbf{hge7m} [\emph{ʰge˥m}] 0x1390: Банкомат\\
\textbf{hge7n} [\emph{ʰge˥n}] 0x7390: калечить\\
\textbf{hge7s} [\emph{ʰge˥s}] 0x9390: магнезит\\
\textbf{hge7t} [\emph{ʰge˥t}] 0xB390: монополист\\
\textbf{hgec} [\emph{ʰgeʃ}] 0xEB90: пещера\\
\textbf{hgef} [\emph{ʰgef}] 0xCB90: обхват\\
\textbf{hgek} [\emph{ʰgek}] 0x2B90: гей\\
\textbf{hgem} [\emph{ʰgem}] 0x0B90: магний\\
\textbf{hgen} [\emph{ʰgen}] 0x6B90: агент\\
\textbf{hgep} [\emph{ʰgep}] 0x4B90: Постельное белье\\
\textbf{hges} [\emph{ʰges}] 0x8B90: беременность\\
\textbf{hget} [\emph{ʰget}] 0xAB90: гость\\
\textbf{hgi2c} [\emph{ʰgi˩ʃ}] 0xF890: часовщик\\
\textbf{hgi2k} [\emph{ʰgi˩k}] 0x3890: митра\\
\textbf{hgi2n} [\emph{ʰgi˩n}] 0x7890: сальный\\
\textbf{hgi2p} [\emph{ʰgi˩p}] 0x5890: волынка\\
\textbf{hgi2s} [\emph{ʰgi˩s}] 0x9890: глазурь\\
\textbf{hgi2t} [\emph{ʰgi˩t}] 0xB890: позолоченный\\
\textbf{hgi7c} [\emph{ʰgi˥ʃ}] 0xF090: эгоистический\\
\textbf{hgi7k} [\emph{ʰgi˥k}] 0x3090: двуколка\\
\textbf{hgi7m} [\emph{ʰgi˥m}] 0x1090: эргономика\\
\textbf{hgi7n} [\emph{ʰgi˥n}] 0x7090: хихикающий\\
\textbf{hgi7p} [\emph{ʰgi˥p}] 0x5090: сороки\\
\textbf{hgi7s} [\emph{ʰgi˥s}] 0x9090: цыганский\\
\textbf{hgi7t} [\emph{ʰgi˥t}] 0xB090: несогласованность\\
\textbf{hgic} [\emph{ʰgiʃ}] 0xE890: переполненный\\
\textbf{hgif} [\emph{ʰgif}] 0xC890: очаровывать\\
\textbf{hgik} [\emph{ʰgik}] 0x2890: колено\\
\textbf{hgim} [\emph{ʰgim}] 0x0890: дымовая труба\\
\textbf{hgin} [\emph{ʰgin}] 0x6890: тюрьма\\
\textbf{hgip} [\emph{ʰgip}] 0x4890: капуста\\
\textbf{hgis} [\emph{ʰgis}] 0x8890: сыр\\
\textbf{hgit} [\emph{ʰgit}] 0xA890: свидетель\\
\textbf{hgo2n} [\emph{ʰgo˩n}] 0x7C90: родственный\\
\textbf{hgo2t} [\emph{ʰgo˩t}] 0xBC90: стойки ворот\\
\textbf{hgo7k} [\emph{ʰgo˥k}] 0x3490: идти корзина\\
\textbf{hgo7m} [\emph{ʰgo˥m}] 0x1490: Graeco романский\\
\textbf{hgo7n} [\emph{ʰgo˥n}] 0x7490: августинец\\
\textbf{hgo7s} [\emph{ʰgo˥s}] 0x9490: вирши\\
\textbf{hgo7t} [\emph{ʰgo˥t}] 0xB490: подагра\\
\textbf{hgoc} [\emph{ʰgoʃ}] 0xEC90: умножить\\
\textbf{hgof} [\emph{ʰgof}] 0xCC90: почвопокровные\\
\textbf{hgok} [\emph{ʰgok}] 0x2C90: органический\\
\textbf{hgom} [\emph{ʰgom}] 0x0C90: гамма\\
\textbf{hgon} [\emph{ʰgon}] 0x6C90: яма\\
\textbf{hgop} [\emph{ʰgop}] 0x4C90: гильдия\\
\textbf{hgos} [\emph{ʰgos}] 0x8C90: ржавчина\\
\textbf{hgot} [\emph{ʰgot}] 0xAC90: замороженные\\
\textbf{hgu2n} [\emph{ʰgu˩n}] 0x7A90: коронование\\
\textbf{hgu2t} [\emph{ʰgu˩t}] 0xBA90: Argus глазами\\
\textbf{hgu7k} [\emph{ʰgu˥k}] 0x3290: гу\\
\textbf{hgu7m} [\emph{ʰgu˥m}] 0x1290: гм\\
\textbf{hgu7n} [\emph{ʰgu˥n}] 0x7290: борзая\\
\textbf{hgu7s} [\emph{ʰgu˥s}] 0x9290: Гвинея Бисау\\
\textbf{hgu7t} [\emph{ʰgu˥t}] 0xB290: крестник\\
\textbf{hguc} [\emph{ʰguʃ}] 0xEA90: жажда\\
\textbf{hguf} [\emph{ʰguf}] 0xCA90: лачуга\\
\textbf{hguk} [\emph{ʰguk}] 0x2A90: тыква\\
\textbf{hgum} [\emph{ʰgum}] 0x0A90: жвачка\\
\textbf{hgun} [\emph{ʰgun}] 0x6A90: ароматный\\
\textbf{hgup} [\emph{ʰgup}] 0x4A90: купол\\
\textbf{hgus} [\emph{ʰgus}] 0x8A90: коэффициент\\
\textbf{hgut} [\emph{ʰgut}] 0xAA90: подсчет\\
\textbf{hj6n} [\emph{ʰʒən}] 0x6DA8: не человек\\
\textbf{hj6s} [\emph{ʰʒəs}] 0x8DA8: гигантизм\\
\textbf{hj6t} [\emph{ʰʒət}] 0xADA8: гладильный\\
\textbf{hja2n} [\emph{ʰʒä˩n}] 0x79A8: бездонный\\
\textbf{hja2t} [\emph{ʰʒä˩t}] 0xB9A8: отчим\\
\textbf{hja7k} [\emph{ʰʒä˥k}] 0x31A8: дворняжка\\
\textbf{hja7n} [\emph{ʰʒä˥n}] 0x71A8: смачивание\\
\textbf{hja7s} [\emph{ʰʒä˥s}] 0x91A8: вычеркивать\\
\textbf{hja7t} [\emph{ʰʒä˥t}] 0xB1A8: предел слышимости\\
\textbf{hjac} [\emph{ʰʒäʃ}] 0xE9A8: Отправить\\
\textbf{hjaf} [\emph{ʰʒäf}] 0xC9A8: млекопитающее\\
\textbf{hjak} [\emph{ʰʒäk}] 0x29A8: охватывать\\
\textbf{hjam} [\emph{ʰʒäm}] 0x09A8: варенье\\
\textbf{hjan} [\emph{ʰʒän}] 0x69A8: наказание\\
\textbf{hjap} [\emph{ʰʒäp}] 0x49A8: халат\\
\textbf{hjas} [\emph{ʰʒäs}] 0x89A8: подавлять\\
\textbf{hjat} [\emph{ʰʒät}] 0xA9A8: человечество\\
\textbf{hje2n} [\emph{ʰʒe˩n}] 0x7BA8: изнеженность\\
\textbf{hje7n} [\emph{ʰʒe˥n}] 0x73A8: тощая кишка\\
\textbf{hje7t} [\emph{ʰʒe˥t}] 0xB3A8: пересидеть\\
\textbf{hjec} [\emph{ʰʒeʃ}] 0xEBA8: плавиться\\
\textbf{hjef} [\emph{ʰʒef}] 0xCBA8: Женева\\
\textbf{hjek} [\emph{ʰʒek}] 0x2BA8: агентивного дело\\
\textbf{hjem} [\emph{ʰʒem}] 0x0BA8: Немецкий\\
\textbf{hjen} [\emph{ʰʒen}] 0x6BA8: женский\\
\textbf{hjep} [\emph{ʰʒep}] 0x4BA8: Пол\\
\textbf{hjes} [\emph{ʰʒes}] 0x8BA8: перепись\\
\textbf{hjet} [\emph{ʰʒet}] 0xABA8: бархат\\
\textbf{hji2k} [\emph{ʰʒi˩k}] 0x38A8: наркоман\\
\textbf{hji2n} [\emph{ʰʒi˩n}] 0x78A8: Дженни\\
\textbf{hji2s} [\emph{ʰʒi˩s}] 0x98A8: приносить несчастье\\
\textbf{hji2t} [\emph{ʰʒi˩t}] 0xB8A8: ржать\\
\textbf{hji7k} [\emph{ʰʒi˥k}] 0x30A8: IC\\
\textbf{hji7n} [\emph{ʰʒi˥n}] 0x70A8: Женщины мира\\
\textbf{hji7s} [\emph{ʰʒi˥s}] 0x90A8: землевладельческий\\
\textbf{hji7t} [\emph{ʰʒi˥t}] 0xB0A8: пылко\\
\textbf{hjic} [\emph{ʰʒiʃ}] 0xE8A8: выходить на пенсию\\
\textbf{hjif} [\emph{ʰʒif}] 0xC8A8: незаконное\\
\textbf{hjik} [\emph{ʰʒik}] 0x28A8: гостеприимство\\
\textbf{hjim} [\emph{ʰʒim}] 0x08A8: запретительный настроение\\
\textbf{hjin} [\emph{ʰʒin}] 0x68A8: терпение\\
\textbf{hjip} [\emph{ʰʒip}] 0x48A8: тропический\\
\textbf{hjis} [\emph{ʰʒis}] 0x88A8: игла\\
\textbf{hjit} [\emph{ʰʒit}] 0xA8A8: Палец\\
\textbf{hjo2t} [\emph{ʰʒo˩t}] 0xBCA8: мастоидный\\
\textbf{hjo7n} [\emph{ʰʒo˥n}] 0x74A8: унгвент\\
\textbf{hjo7t} [\emph{ʰʒo˥t}] 0xB4A8: хроматин\\
\textbf{hjoc} [\emph{ʰʒoʃ}] 0xECA8: монтаж\\
\textbf{hjof} [\emph{ʰʒof}] 0xCCA8: припарка\\
\textbf{hjok} [\emph{ʰʒok}] 0x2CA8: жокей\\
\textbf{hjom} [\emph{ʰʒom}] 0x0CA8: хромосома\\
\textbf{hjon} [\emph{ʰʒon}] 0x6CA8: сексуальное\\
\textbf{hjop} [\emph{ʰʒop}] 0x4CA8: гироскоп\\
\textbf{hjos} [\emph{ʰʒos}] 0x8CA8: без присмотра\\
\textbf{hjot} [\emph{ʰʒot}] 0xACA8: студия\\
\textbf{hjuc} [\emph{ʰʒuʃ}] 0xEAA8: добыча\\
\textbf{hjuk} [\emph{ʰʒuk}] 0x2AA8: становиться на колени\\
\textbf{hjum} [\emph{ʰʒum}] 0x0AA8: Иерусалим\\
\textbf{hjun} [\emph{ʰʒun}] 0x6AA8: ветреный\\
\textbf{hjup} [\emph{ʰʒup}] 0x4AA8: песчаная буря\\
\textbf{hjus} [\emph{ʰʒus}] 0x8AA8: Иисус\\
\textbf{hjut} [\emph{ʰʒut}] 0xAAA8: облачный\\
\textbf{hk62t} [\emph{ʰkə˩t}] 0xBD10: лактон\\
\textbf{hk67k} [\emph{ʰkə˥k}] 0x3510: ключица\\
\textbf{hk67n} [\emph{ʰkə˥n}] 0x7510: хинолина\\
\textbf{hk67t} [\emph{ʰkə˥t}] 0xB510: коммодор\\
\textbf{hk6c} [\emph{ʰkəʃ}] 0xED10: канарейка\\
\textbf{hk6f} [\emph{ʰkəf}] 0xCD10: хвойное дерево\\
\textbf{hk6k} [\emph{ʰkək}] 0x2D10: мексиканский\\
\textbf{hk6m} [\emph{ʰkəm}] 0x0D10: мошонка\\
\textbf{hk6n} [\emph{ʰkən}] 0x6D10: компилировать\\
\textbf{hk6p} [\emph{ʰkəp}] 0x4D10: Колумбия\\
\textbf{hk6s} [\emph{ʰkəs}] 0x8D10: рождество\\
\textbf{hk6t} [\emph{ʰkət}] 0xAD10: октет\\
\textbf{hka2c} [\emph{ʰkä˩ʃ}] 0xF910: инквизиция\\
\textbf{hka2f} [\emph{ʰkä˩f}] 0xD910: октава\\
\textbf{hka2k} [\emph{ʰkä˩k}] 0x3910: покалывание\\
\textbf{hka2m} [\emph{ʰkä˩m}] 0x1910: сперма\\
\textbf{hka2n} [\emph{ʰkä˩n}] 0x7910: сопровождая\\
\textbf{hka2p} [\emph{ʰkä˩p}] 0x5910: сапожник\\
\textbf{hka2s} [\emph{ʰkä˩s}] 0x9910: коаксиальный\\
\textbf{hka2t} [\emph{ʰkä˩t}] 0xB910: товарищи\\
\textbf{hka7c} [\emph{ʰkä˥ʃ}] 0xF110: Каир\\
\textbf{hka7f} [\emph{ʰkä˥f}] 0xD110: скандинавский\\
\textbf{hka7k} [\emph{ʰkä˥k}] 0x3110: замки\\
\textbf{hka7m} [\emph{ʰkä˥m}] 0x1110: ласка\\
\textbf{hka7n} [\emph{ʰkä˥n}] 0x7110: очки\\
\textbf{hka7p} [\emph{ʰkä˥p}] 0x5110: неизмеримый\\
\textbf{hka7s} [\emph{ʰkä˥s}] 0x9110: квазар\\
\textbf{hka7t} [\emph{ʰkä˥t}] 0xB110: с благодарностью\\
\textbf{hkac} [\emph{ʰkäʃ}] 0xE910: когда\\
\textbf{hkaf} [\emph{ʰkäf}] 0xC910: официальный\\
\textbf{hkak} [\emph{ʰkäk}] 0x2910: винительный дело\\
\textbf{hkam} [\emph{ʰkäm}] 0x0910: какие\\
\textbf{hkan} [\emph{ʰkän}] 0x6910: пространство контекст\\
\textbf{hkap} [\emph{ʰkäp}] 0x4910: никогда\\
\textbf{hkas} [\emph{ʰkäs}] 0x8910: мог\\
\textbf{hkat} [\emph{ʰkät}] 0xA910: мимо напряженный\\
\textbf{hke2k} [\emph{ʰke˩k}] 0x3B10: квакер\\
\textbf{hke2n} [\emph{ʰke˩n}] 0x7B10: коса\\
\textbf{hke2p} [\emph{ʰke˩p}] 0x5B10: эр\\
\textbf{hke2s} [\emph{ʰke˩s}] 0x9B10: Экклезиаст\\
\textbf{hke2t} [\emph{ʰke˩t}] 0xBB10: дородный\\
\textbf{hke7k} [\emph{ʰke˥k}] 0x3310: конвенции\\
\textbf{hke7m} [\emph{ʰke˥m}] 0x1310: кхмерская\\
\textbf{hke7n} [\emph{ʰke˥n}] 0x7310: имущественный\\
\textbf{hke7p} [\emph{ʰke˥p}] 0x5310: кастаньеты\\
\textbf{hke7s} [\emph{ʰke˥s}] 0x9310: кратеры\\
\textbf{hke7t} [\emph{ʰke˥t}] 0xB310: кокетка\\
\textbf{hkec} [\emph{ʰkeʃ}] 0xEB10: каменный уголь\\
\textbf{hkef} [\emph{ʰkef}] 0xCB10: кафе\\
\textbf{hkek} [\emph{ʰkek}] 0x2B10: хвост\\
\textbf{hkem} [\emph{ʰkem}] 0x0B10: комета\\
\textbf{hken} [\emph{ʰken}] 0x6B10: кабина\\
\textbf{hkep} [\emph{ʰkep}] 0x4B10: медь\\
\textbf{hkes} [\emph{ʰkes}] 0x8B10: сложный\\
\textbf{hket} [\emph{ʰket}] 0xAB10: проснулся\\
\textbf{hki2c} [\emph{ʰki˩ʃ}] 0xF810: сожительство\\
\textbf{hki2f} [\emph{ʰki˩f}] 0xD810: Киев\\
\textbf{hki2k} [\emph{ʰki˩k}] 0x3810: оленья кожа\\
\textbf{hki2m} [\emph{ʰki˩m}] 0x1810: запор\\
\textbf{hki2n} [\emph{ʰki˩n}] 0x7810: грешник\\
\textbf{hki2p} [\emph{ʰki˩p}] 0x5810: каннибал\\
\textbf{hki2s} [\emph{ʰki˩s}] 0x9810: сорок шесть\\
\textbf{hki2t} [\emph{ʰki˩t}] 0xB810: бигуди\\
\textbf{hki7c} [\emph{ʰki˥ʃ}] 0xF010: гвоздика\\
\textbf{hki7f} [\emph{ʰki˥f}] 0xD010: треска\\
\textbf{hki7k} [\emph{ʰki˥k}] 0x3010: хруст\\
\textbf{hki7m} [\emph{ʰki˥m}] 0x1010: несмываемый\\
\textbf{hki7n} [\emph{ʰki˥n}] 0x7010: принимая во внимание\\
\textbf{hki7p} [\emph{ʰki˥p}] 0x5010: валовой внутренний продукт\\
\textbf{hki7s} [\emph{ʰki˥s}] 0x9010: сородичи\\
\textbf{hki7t} [\emph{ʰki˥t}] 0xB010: сгущенный\\
\textbf{hkic} [\emph{ʰkiʃ}] 0xE810: между\\
\textbf{hkif} [\emph{ʰkif}] 0xC810: благородный\\
\textbf{hkik} [\emph{ʰkik}] 0x2810: кто\\
\textbf{hkim} [\emph{ʰkim}] 0x0810: секрет\\
\textbf{hkin} [\emph{ʰkin}] 0x6810: идти\\
\textbf{hkip} [\emph{ʰkip}] 0x4810: около\\
\textbf{hkis} [\emph{ʰkis}] 0x8810: черный\\
\textbf{hkit} [\emph{ʰkit}] 0xA810: Эксклюзивный или\\
\textbf{hko2c} [\emph{ʰko˩ʃ}] 0xFC10: кошерный\\
\textbf{hko2k} [\emph{ʰko˩k}] 0x3C10: несвязанный\\
\textbf{hko2m} [\emph{ʰko˩m}] 0x1C10: креольский\\
\textbf{hko2n} [\emph{ʰko˩n}] 0x7C10: полковник\\
\textbf{hko2p} [\emph{ʰko˩p}] 0x5C10: копра\\
\textbf{hko2s} [\emph{ʰko˩s}] 0x9C10: созвездий\\
\textbf{hko2t} [\emph{ʰko˩t}] 0xBC10: сердечно\\
\textbf{hko7c} [\emph{ʰko˥ʃ}] 0xF410: картуш\\
\textbf{hko7f} [\emph{ʰko˥f}] 0xD410: бухточка\\
\textbf{hko7k} [\emph{ʰko˥k}] 0x3410: иго\\
\textbf{hko7n} [\emph{ʰko˥n}] 0x7410: Обще\\
\textbf{hko7p} [\emph{ʰko˥p}] 0x5410: неконтролируемый\\
\textbf{hko7s} [\emph{ʰko˥s}] 0x9410: корсаж\\
\textbf{hko7t} [\emph{ʰko˥t}] 0xB410: фанатик\\
\textbf{hkoc} [\emph{ʰkoʃ}] 0xEC10: время от времени\\
\textbf{hkof} [\emph{ʰkof}] 0xCC10: потребление\\
\textbf{hkok} [\emph{ʰkok}] 0x2C10: Одинокий\\
\textbf{hkom} [\emph{ʰkom}] 0x0C10: код\\
\textbf{hkon} [\emph{ʰkon}] 0x6C10: замок\\
\textbf{hkop} [\emph{ʰkop}] 0x4C10: кружка\\
\textbf{hkos} [\emph{ʰkos}] 0x8C10: каноэ\\
\textbf{hkot} [\emph{ʰkot}] 0xAC10: вместимость\\
\textbf{hku2c} [\emph{ʰku˩ʃ}] 0xFA10: капуцин\\
\textbf{hku2n} [\emph{ʰku˩n}] 0x7A10: подушки\\
\textbf{hku2p} [\emph{ʰku˩p}] 0x5A10: околополярный\\
\textbf{hku2s} [\emph{ʰku˩s}] 0x9A10: камбуз\\
\textbf{hku2t} [\emph{ʰku˩t}] 0xBA10: Амур\\
\textbf{hku7k} [\emph{ʰku˥k}] 0x3210: ворковать\\
\textbf{hku7n} [\emph{ʰku˥n}] 0x7210: никуда не годится\\
\textbf{hku7p} [\emph{ʰku˥p}] 0x5210: Кабул\\
\textbf{hku7s} [\emph{ʰku˥s}] 0x9210: кодициль\\
\textbf{hku7t} [\emph{ʰku˥t}] 0xB210: вычислимость\\
\textbf{hkuc} [\emph{ʰkuʃ}] 0xEA10: высокая\\
\textbf{hkuf} [\emph{ʰkuf}] 0xCA10: мясник\\
\textbf{hkuk} [\emph{ʰkuk}] 0x2A10: Предложение хвост\\
\textbf{hkum} [\emph{ʰkum}] 0x0A10: спекулятивный настроение\\
\textbf{hkun} [\emph{ʰkun}] 0x6A10: Счет\\
\textbf{hkup} [\emph{ʰkup}] 0x4A10: копия\\
\textbf{hkus} [\emph{ʰkus}] 0x8A10: плакать\\
\textbf{hkut} [\emph{ʰkut}] 0xAA10: знать\\
\textbf{hl62n} [\emph{ʰlə˩n}] 0x7D58: алеутская\\
\textbf{hl62t} [\emph{ʰlə˩t}] 0xBD58: психиатр\\
\textbf{hl67n} [\emph{ʰlə˥n}] 0x7558: полноценный\\
\textbf{hl67s} [\emph{ʰlə˥s}] 0x9558: эльзасский\\
\textbf{hl67t} [\emph{ʰlə˥t}] 0xB558: вычерпывать\\
\textbf{hl6c} [\emph{ʰləʃ}] 0xED58: выщелачивание\\
\textbf{hl6f} [\emph{ʰləf}] 0xCD58: литосфера\\
\textbf{hl6k} [\emph{ʰlək}] 0x2D58: лишай\\
\textbf{hl6m} [\emph{ʰləm}] 0x0D58: euler's номер\\
\textbf{hl6n} [\emph{ʰlən}] 0x6D58: контакт\\
\textbf{hl6p} [\emph{ʰləp}] 0x4D58: Люксембург\\
\textbf{hl6s} [\emph{ʰləs}] 0x8D58: лотос\\
\textbf{hl6t} [\emph{ʰlət}] 0xAD58: литр\\
\textbf{hla2c} [\emph{ʰlä˩ʃ}] 0xF958: пресный\\
\textbf{hla2f} [\emph{ʰlä˩f}] 0xD958: слесарь\\
\textbf{hla2k} [\emph{ʰlä˩k}] 0x3958: забавляться\\
\textbf{hla2m} [\emph{ʰlä˩m}] 0x1958: квасцы\\
\textbf{hla2n} [\emph{ʰlä˩n}] 0x7958: Калифорния\\
\textbf{hla2p} [\emph{ʰlä˩p}] 0x5958: филантропия\\
\textbf{hla2s} [\emph{ʰlä˩s}] 0x9958: всплеснула\\
\textbf{hla2t} [\emph{ʰlä˩t}] 0xB958: лепились\\
\textbf{hla7c} [\emph{ʰlä˥ʃ}] 0xF158: лож\\
\textbf{hla7f} [\emph{ʰlä˥f}] 0xD158: планка\\
\textbf{hla7k} [\emph{ʰlä˥k}] 0x3158: Калькутта\\
\textbf{hla7m} [\emph{ʰlä˥m}] 0x1158: алебастр\\
\textbf{hla7n} [\emph{ʰlä˥n}] 0x7158: медленное высвобождение\\
\textbf{hla7p} [\emph{ʰlä˥p}] 0x5158: плач Иеремии\\
\textbf{hla7s} [\emph{ʰlä˥s}] 0x9158: Альпы\\
\textbf{hla7t} [\emph{ʰlä˥t}] 0xB158: Александр\\
\textbf{hlac} [\emph{ʰläʃ}] 0xE958: отчуждаемый владение\\
\textbf{hlaf} [\emph{ʰläf}] 0xC958: оставаться\\
\textbf{hlak} [\emph{ʰläk}] 0x2958: относительный дело\\
\textbf{hlam} [\emph{ʰläm}] 0x0958: приехать\\
\textbf{hlan} [\emph{ʰlän}] 0x6958: но\\
\textbf{hlap} [\emph{ʰläp}] 0x4958: semelfactive аспект\\
\textbf{hlas} [\emph{ʰläs}] 0x8958: линия\\
\textbf{hlat} [\emph{ʰlät}] 0xA958: факультативный квантификатором\\
\textbf{hle2c} [\emph{ʰle˩ʃ}] 0xFB58: Г-образный\\
\textbf{hle2k} [\emph{ʰle˩k}] 0x3B58: длинноногий\\
\textbf{hle2m} [\emph{ʰle˩m}] 0x1B58: лейкемия\\
\textbf{hle2n} [\emph{ʰle˩n}] 0x7B58: лейтенант\\
\textbf{hle2p} [\emph{ʰle˩p}] 0x5B58: приозерный\\
\textbf{hle2s} [\emph{ʰle˩s}] 0x9B58: в любом случае\\
\textbf{hle2t} [\emph{ʰle˩t}] 0xBB58: блаженный\\
\textbf{hle7c} [\emph{ʰle˥ʃ}] 0xF358: Меланезия\\
\textbf{hle7k} [\emph{ʰle˥k}] 0x3358: Тупик\\
\textbf{hle7m} [\emph{ʰle˥m}] 0x1358: дождемер\\
\textbf{hle7n} [\emph{ʰle˥n}] 0x7358: цепи\\
\textbf{hle7p} [\emph{ʰle˥p}] 0x5358: эллипсы\\
\textbf{hle7s} [\emph{ʰle˥s}] 0x9358: лосины\\
\textbf{hle7t} [\emph{ʰle˥t}] 0xB358: тромбоцит\\
\textbf{hlec} [\emph{ʰleʃ}] 0xEB58: перец\\
\textbf{hlef} [\emph{ʰlef}] 0xCB58: агент вызывать\\
\textbf{hlek} [\emph{ʰlek}] 0x2B58: atelic аспект\\
\textbf{hlem} [\emph{ʰlem}] 0x0B58: безрассудный\\
\textbf{hlen} [\emph{ʰlen}] 0x6B58: лицензия\\
\textbf{hlep} [\emph{ʰlep}] 0x4B58: целенаправленный аспект\\
\textbf{hles} [\emph{ʰles}] 0x8B58: слива\\
\textbf{hlet} [\emph{ʰlet}] 0xAB58: лестница\\
\textbf{hli2c} [\emph{ʰli˩ʃ}] 0xF858: несклонный\\
\textbf{hli2f} [\emph{ʰli˩f}] 0xD858: серебряный сладкоречивый\\
\textbf{hli2k} [\emph{ʰli˩k}] 0x3858: лиги\\
\textbf{hli2m} [\emph{ʰli˩m}] 0x1858: алхимия\\
\textbf{hli2n} [\emph{ʰli˩n}] 0x7858: слоны\\
\textbf{hli2p} [\emph{ʰli˩p}] 0x5858: Элизабет\\
\textbf{hli2s} [\emph{ʰli˩s}] 0x9858: злобный взгляд\\
\textbf{hli2t} [\emph{ʰli˩t}] 0xB858: заботливость\\
\textbf{hli7c} [\emph{ʰli˥ʃ}] 0xF058: лотерея\\
\textbf{hli7f} [\emph{ʰli˥f}] 0xD058: котиковый\\
\textbf{hli7k} [\emph{ʰli˥k}] 0x3058: воспитанным\\
\textbf{hli7m} [\emph{ʰli˥m}] 0x1058: недоразумение\\
\textbf{hli7n} [\emph{ʰli˥n}] 0x7058: сверкающих\\
\textbf{hli7p} [\emph{ʰli˥p}] 0x5058: дирижабль\\
\textbf{hli7s} [\emph{ʰli˥s}] 0x9058: поясница\\
\textbf{hli7t} [\emph{ʰli˥t}] 0xB058: Илиада\\
\textbf{hlic} [\emph{ʰliʃ}] 0xE858: уже\\
\textbf{hlif} [\emph{ʰlif}] 0xC858: прибыль\\
\textbf{hlik} [\emph{ʰlik}] 0x2858: выражающий заключение дело\\
\textbf{hlim} [\emph{ʰlim}] 0x0858: климат\\
\textbf{hlin} [\emph{ʰlin}] 0x6858: с плавающей точкой номер\\
\textbf{hlip} [\emph{ʰlip}] 0x4858: заканчивающий аспект\\
\textbf{hlis} [\emph{ʰlis}] 0x8858: шесть\\
\textbf{hlit} [\emph{ʰlit}] 0xA858: волевых настроение\\
\textbf{hlo2k} [\emph{ʰlo˩k}] 0x3C58: бревна\\
\textbf{hlo2m} [\emph{ʰlo˩m}] 0x1C58: клинометр\\
\textbf{hlo2n} [\emph{ʰlo˩n}] 0x7C58: Болонья\\
\textbf{hlo2p} [\emph{ʰlo˩p}] 0x5C58: кувырок провал\\
\textbf{hlo2s} [\emph{ʰlo˩s}] 0x9C58: лоббист\\
\textbf{hlo2t} [\emph{ʰlo˩t}] 0xBC58: Лондон\\
\textbf{hlo7c} [\emph{ʰlo˥ʃ}] 0xF458: хлорал\\
\textbf{hlo7k} [\emph{ʰlo˥k}] 0x3458: медальон\\
\textbf{hlo7m} [\emph{ʰlo˥m}] 0x1458: терн\\
\textbf{hlo7n} [\emph{ʰlo˥n}] 0x7458: фаллопиевы\\
\textbf{hlo7p} [\emph{ʰlo˥p}] 0x5458: Лиссабон\\
\textbf{hlo7s} [\emph{ʰlo˥s}] 0x9458: начитанный\\
\textbf{hlo7t} [\emph{ʰlo˥t}] 0xB458: погрузка\\
\textbf{hloc} [\emph{ʰloʃ}] 0xEC58: меланхолия\\
\textbf{hlof} [\emph{ʰlof}] 0xCC58: сера\\
\textbf{hlok} [\emph{ʰlok}] 0x2C58: прыжок\\
\textbf{hlom} [\emph{ʰlom}] 0x0C58: лимон\\
\textbf{hlon} [\emph{ʰlon}] 0x6C58: рассвет\\
\textbf{hlop} [\emph{ʰlop}] 0x4C58: телескоп\\
\textbf{hlos} [\emph{ʰlos}] 0x8C58: вверх по лестнице\\
\textbf{hlot} [\emph{ʰlot}] 0xAC58: планируемый\\
\textbf{hlu2k} [\emph{ʰlu˩k}] 0x3A58: Люк\\
\textbf{hlu2n} [\emph{ʰlu˩n}] 0x7A58: лютеранский\\
\textbf{hlu2p} [\emph{ʰlu˩p}] 0x5A58: лупа\\
\textbf{hlu2s} [\emph{ʰlu˩s}] 0x9A58: эманация\\
\textbf{hlu2t} [\emph{ʰlu˩t}] 0xBA58: политеистический\\
\textbf{hlu7c} [\emph{ʰlu˥ʃ}] 0xF258: Loup Garou\\
\textbf{hlu7k} [\emph{ʰlu˥k}] 0x3258: голубовато-черный\\
\textbf{hlu7m} [\emph{ʰlu˥m}] 0x1258: оксид алюминия\\
\textbf{hlu7n} [\emph{ʰlu˥n}] 0x7258: вольная\\
\textbf{hlu7s} [\emph{ʰlu˥s}] 0x9258: толстой кожурой\\
\textbf{hlu7t} [\emph{ʰlu˥t}] 0xB258: прачка\\
\textbf{hluc} [\emph{ʰluʃ}] 0xEA58: популярный\\
\textbf{hluf} [\emph{ʰluf}] 0xCA58: старость\\
\textbf{hluk} [\emph{ʰluk}] 0x2A58: устала\\
\textbf{hlum} [\emph{ʰlum}] 0x0A58: верблюд\\
\textbf{hlun} [\emph{ʰlun}] 0x6A58: включительно человек\\
\textbf{hlup} [\emph{ʰlup}] 0x4A58: альбом\\
\textbf{hlus} [\emph{ʰlus}] 0x8A58: олень\\
\textbf{hlut} [\emph{ʰlut}] 0xAA58: автор\\
\textbf{hm62n} [\emph{ʰmə˩n}] 0x7D08: монгольский\\
\textbf{hm62t} [\emph{ʰmə˩t}] 0xBD08: нумизмат\\
\textbf{hm67k} [\emph{ʰmə˥k}] 0x3508: семантика\\
\textbf{hm67n} [\emph{ʰmə˥n}] 0x7508: тимпан\\
\textbf{hm67s} [\emph{ʰmə˥s}] 0x9508: мальтузианский\\
\textbf{hm67t} [\emph{ʰmə˥t}] 0xB508: Суматра\\
\textbf{hm6c} [\emph{ʰməʃ}] 0xED08: всеядный\\
\textbf{hm6f} [\emph{ʰməf}] 0xCD08: запечатывание\\
\textbf{hm6k} [\emph{ʰmək}] 0x2D08: Мексика\\
\textbf{hm6m} [\emph{ʰməm}] 0x0D08: в середине лета\\
\textbf{hm6n} [\emph{ʰmən}] 0x6D08: излучающий\\
\textbf{hm6p} [\emph{ʰməp}] 0x4D08: материнская плата\\
\textbf{hm6s} [\emph{ʰməs}] 0x8D08: Мадагаскар\\
\textbf{hm6t} [\emph{ʰmət}] 0xAD08: первичная объект\\
\textbf{hma2c} [\emph{ʰmä˩ʃ}] 0xF908: пастушьи\\
\textbf{hma2f} [\emph{ʰmä˩f}] 0xD908: мафия\\
\textbf{hma2k} [\emph{ʰmä˩k}] 0x3908: опьяняющий\\
\textbf{hma2m} [\emph{ʰmä˩m}] 0x1908: оцепенение\\
\textbf{hma2n} [\emph{ʰmä˩n}] 0x7908: равнины\\
\textbf{hma2p} [\emph{ʰmä˩p}] 0x5908: недуги\\
\textbf{hma2s} [\emph{ʰmä˩s}] 0x9908: Марсель\\
\textbf{hma2t} [\emph{ʰmä˩t}] 0xB908: изменник\\
\textbf{hma7c} [\emph{ʰmä˥ʃ}] 0xF108: малайский\\
\textbf{hma7f} [\emph{ʰmä˥f}] 0xD108: дог\\
\textbf{hma7k} [\emph{ʰmä˥k}] 0x3108: Майкл\\
\textbf{hma7m} [\emph{ʰmä˥m}] 0x1108: кобыла\\
\textbf{hma7n} [\emph{ʰmä˥n}] 0x7108: грот\\
\textbf{hma7p} [\emph{ʰmä˥p}] 0x5108: пиломатериалы\\
\textbf{hma7s} [\emph{ʰmä˥s}] 0x9108: металлы\\
\textbf{hma7t} [\emph{ʰmä˥t}] 0xB108: Госпожа\\
\textbf{hmac} [\emph{ʰmäʃ}] 0xE908: момент\\
\textbf{hmaf} [\emph{ʰmäf}] 0xC908: беда\\
\textbf{hmak} [\emph{ʰmäk}] 0x2908: comitative дело\\
\textbf{hmam} [\emph{ʰmäm}] 0x0908: Мама\\
\textbf{hman} [\emph{ʰmän}] 0x6908: меня\\
\textbf{hmap} [\emph{ʰmäp}] 0x4908: Роза\\
\textbf{hmas} [\emph{ʰmäs}] 0x8908: разум\\
\textbf{hmat} [\emph{ʰmät}] 0xA908: демонстративный\\
\textbf{hme2c} [\emph{ʰme˩ʃ}] 0xFB08: манихей\\
\textbf{hme2k} [\emph{ʰme˩k}] 0x3B08: пакля\\
\textbf{hme2m} [\emph{ʰme˩m}] 0x1B08: мезодерма\\
\textbf{hme2n} [\emph{ʰme˩n}] 0x7B08: каменщик\\
\textbf{hme2p} [\emph{ʰme˩p}] 0x5B08: табличка с именем\\
\textbf{hme2s} [\emph{ʰme˩s}] 0x9B08: метемпсихоз\\
\textbf{hme2t} [\emph{ʰme˩t}] 0xBB08: на полпути\\
\textbf{hme7c} [\emph{ʰme˥ʃ}] 0xF308: маг\\
\textbf{hme7f} [\emph{ʰme˥f}] 0xD308: ментол\\
\textbf{hme7k} [\emph{ʰme˥k}] 0x3308: черногория\\
\textbf{hme7m} [\emph{ʰme˥m}] 0x1308: середине мая\\
\textbf{hme7n} [\emph{ʰme˥n}] 0x7308: неохотно\\
\textbf{hme7p} [\emph{ʰme˥p}] 0x5308: приветливо\\
\textbf{hme7s} [\emph{ʰme˥s}] 0x9308: медуза\\
\textbf{hme7t} [\emph{ʰme˥t}] 0xB308: извинить\\
\textbf{hmec} [\emph{ʰmeʃ}] 0xEB08: неопределенный\\
\textbf{hmef} [\emph{ʰmef}] 0xCB08: май месяц\\
\textbf{hmek} [\emph{ʰmek}] 0x2B08: чувствительный\\
\textbf{hmem} [\emph{ʰmem}] 0x0B08: музей\\
\textbf{hmen} [\emph{ʰmen}] 0x6B08: смиренный\\
\textbf{hmep} [\emph{ʰmep}] 0x4B08: momentane аспект\\
\textbf{hmes} [\emph{ʰmes}] 0x8B08: мэр\\
\textbf{hmet} [\emph{ʰmet}] 0xAB08: Эл. адрес\\
\textbf{hmi2c} [\emph{ʰmi˩ʃ}] 0xF808: криминолог\\
\textbf{hmi2f} [\emph{ʰmi˩f}] 0xD808: мистифицирован\\
\textbf{hmi2k} [\emph{ʰmi˩k}] 0x3808: метафизика\\
\textbf{hmi2m} [\emph{ʰmi˩m}] 0x1808: маори\\
\textbf{hmi2n} [\emph{ʰmi˩n}] 0x7808: суеверие\\
\textbf{hmi2p} [\emph{ʰmi˩p}] 0x5808: микроавтобус\\
\textbf{hmi2s} [\emph{ʰmi˩s}] 0x9808: чудотворный\\
\textbf{hmi2t} [\emph{ʰmi˩t}] 0xB808: математики\\
\textbf{hmi7c} [\emph{ʰmi˥ʃ}] 0xF008: губчатый\\
\textbf{hmi7f} [\emph{ʰmi˥f}] 0xD008: метил\\
\textbf{hmi7k} [\emph{ʰmi˥k}] 0x3008: манильской\\
\textbf{hmi7m} [\emph{ʰmi˥m}] 0x1008: жасмин\\
\textbf{hmi7n} [\emph{ʰmi˥n}] 0x7008: Милан\\
\textbf{hmi7p} [\emph{ʰmi˥p}] 0x5008: мини-бар\\
\textbf{hmi7s} [\emph{ʰmi˥s}] 0x9008: Миссисипи\\
\textbf{hmi7t} [\emph{ʰmi˥t}] 0xB008: бессмертный\\
\textbf{hmic} [\emph{ʰmiʃ}] 0xE808: милая\\
\textbf{hmif} [\emph{ʰmif}] 0xC808: мили\\
\textbf{hmik} [\emph{ʰmik}] 0x2808: Музыка\\
\textbf{hmim} [\emph{ʰmim}] 0x0808: средний\\
\textbf{hmin} [\emph{ʰmin}] 0x6808: ничего\\
\textbf{hmip} [\emph{ʰmip}] 0x4808: несовершенный аспект\\
\textbf{hmis} [\emph{ʰmis}] 0x8808: мудрый\\
\textbf{hmit} [\emph{ʰmit}] 0xA808: море\\
\textbf{hmo2k} [\emph{ʰmo˩k}] 0x3C08: импульсов\\
\textbf{hmo2m} [\emph{ʰmo˩m}] 0x1C08: мяукать\\
\textbf{hmo2n} [\emph{ʰmo˩n}] 0x7C08: манна\\
\textbf{hmo2p} [\emph{ʰmo˩p}] 0x5C08: в середине октября\\
\textbf{hmo2s} [\emph{ʰmo˩s}] 0x9C08: миоцен\\
\textbf{hmo2t} [\emph{ʰmo˩t}] 0xBC08: Ломбардия\\
\textbf{hmo7c} [\emph{ʰmo˥ʃ}] 0xF408: омега\\
\textbf{hmo7k} [\emph{ʰmo˥k}] 0x3408: Москва\\
\textbf{hmo7m} [\emph{ʰmo˥m}] 0x1408: Месопотамия\\
\textbf{hmo7n} [\emph{ʰmo˥n}] 0x7408: аммиак\\
\textbf{hmo7p} [\emph{ʰmo˥p}] 0x5408: шарить\\
\textbf{hmo7s} [\emph{ʰmo˥s}] 0x9408: Моисей\\
\textbf{hmo7t} [\emph{ʰmo˥t}] 0xB408: сообщества\\
\textbf{hmoc} [\emph{ʰmoʃ}] 0xEC08: где-то\\
\textbf{hmof} [\emph{ʰmof}] 0xCC08: торговый центр\\
\textbf{hmok} [\emph{ʰmok}] 0x2C08: мельница\\
\textbf{hmom} [\emph{ʰmom}] 0x0C08: горчица\\
\textbf{hmon} [\emph{ʰmon}] 0x6C08: миссия\\
\textbf{hmop} [\emph{ʰmop}] 0x4C08: моль\\
\textbf{hmos} [\emph{ʰmos}] 0x8C08: министр\\
\textbf{hmot} [\emph{ʰmot}] 0xAC08: мясо\\
\textbf{hmu2k} [\emph{ʰmu˩k}] 0x3A08: нумизматика\\
\textbf{hmu2n} [\emph{ʰmu˩n}] 0x7A08: Мюнхен\\
\textbf{hmu2p} [\emph{ʰmu˩p}] 0x5A08: обагрять\\
\textbf{hmu2s} [\emph{ʰmu˩s}] 0x9A08: Мафусаил\\
\textbf{hmu2t} [\emph{ʰmu˩t}] 0xBA08: Кульминацией\\
\textbf{hmu7c} [\emph{ʰmu˥ʃ}] 0xF208: конюшни\\
\textbf{hmu7k} [\emph{ʰmu˥k}] 0x3208: медулла\\
\textbf{hmu7m} [\emph{ʰmu˥m}] 0x1208: Магомет\\
\textbf{hmu7n} [\emph{ʰmu˥n}] 0x7208: лань\\
\textbf{hmu7p} [\emph{ʰmu˥p}] 0x5208: надгробная плита\\
\textbf{hmu7s} [\emph{ʰmu˥s}] 0x9208: превью\\
\textbf{hmu7t} [\emph{ʰmu˥t}] 0xB208: Мэтью\\
\textbf{hmuc} [\emph{ʰmuʃ}] 0xEA08: цель\\
\textbf{hmuf} [\emph{ʰmuf}] 0xCA08: рот\\
\textbf{hmuk} [\emph{ʰmuk}] 0x2A08: минут\\
\textbf{hmum} [\emph{ʰmum}] 0x0A08: юмор\\
\textbf{hmun} [\emph{ʰmun}] 0x6A08: мягкий\\
\textbf{hmup} [\emph{ʰmup}] 0x4A08: шапка\\
\textbf{hmus} [\emph{ʰmus}] 0x8A08: толстый\\
\textbf{hmut} [\emph{ʰmut}] 0xAA08: потерял\\
\textbf{hn62k} [\emph{ʰnə˩k}] 0x3D30: кинематика\\
\textbf{hn62n} [\emph{ʰnə˩n}] 0x7D30: Noman\\
\textbf{hn62s} [\emph{ʰnə˩s}] 0x9D30: анорексия\\
\textbf{hn62t} [\emph{ʰnə˩t}] 0xBD30: синкопированная\\
\textbf{hn67k} [\emph{ʰnə˥k}] 0x3530: мениск\\
\textbf{hn67n} [\emph{ʰnə˥n}] 0x7530: никотин\\
\textbf{hn67p} [\emph{ʰnə˥p}] 0x5530: звукоподражательный\\
\textbf{hn67s} [\emph{ʰnə˥s}] 0x9530: анархизм\\
\textbf{hn67t} [\emph{ʰnə˥t}] 0xB530: нагота\\
\textbf{hn6c} [\emph{ʰnəʃ}] 0xED30: Нигер\\
\textbf{hn6f} [\emph{ʰnəf}] 0xCD30: necessitative настроение\\
\textbf{hn6k} [\emph{ʰnək}] 0x2D30: Монако\\
\textbf{hn6m} [\emph{ʰnəm}] 0x0D30: Нериум\\
\textbf{hn6n} [\emph{ʰnən}] 0x6D30: неон\\
\textbf{hn6p} [\emph{ʰnəp}] 0x4D30: Непал\\
\textbf{hn6s} [\emph{ʰnəs}] 0x8D30: Ренессанс\\
\textbf{hn6t} [\emph{ʰnət}] 0xAD30: соединители\\
\textbf{hna2c} [\emph{ʰnä˩ʃ}] 0xF930: аннуитет\\
\textbf{hna2f} [\emph{ʰnä˩f}] 0xD930: жених\\
\textbf{hna2k} [\emph{ʰnä˩k}] 0x3930: анекдот\\
\textbf{hna2m} [\emph{ʰnä˩m}] 0x1930: пантеизм\\
\textbf{hna2n} [\emph{ʰnä˩n}] 0x7930: лунный свет\\
\textbf{hna2p} [\emph{ʰnä˩p}] 0x5930: Занзибар\\
\textbf{hna2s} [\emph{ʰnä˩s}] 0x9930: династия\\
\textbf{hna2t} [\emph{ʰnä˩t}] 0xB930: кивнул\\
\textbf{hna7c} [\emph{ʰnä˥ʃ}] 0xF130: марочный\\
\textbf{hna7f} [\emph{ʰnä˥f}] 0xD130: неф\\
\textbf{hna7k} [\emph{ʰnä˥k}] 0x3130: сын в закон\\
\textbf{hna7m} [\emph{ʰnä˥m}] 0x1130: натальный\\
\textbf{hna7n} [\emph{ʰnä˥n}] 0x7130: без углерода\\
\textbf{hna7p} [\emph{ʰnä˥p}] 0x5130: бесспорный\\
\textbf{hna7s} [\emph{ʰnä˥s}] 0x9130: Неаполь\\
\textbf{hna7t} [\emph{ʰnä˥t}] 0xB130: цианобактерии\\
\textbf{hnac} [\emph{ʰnäʃ}] 0xE930: путь дело\\
\textbf{hnaf} [\emph{ʰnäf}] 0xC930: предположим,\\
\textbf{hnak} [\emph{ʰnäk}] 0x2930: именительный дело\\
\textbf{hnam} [\emph{ʰnäm}] 0x0930: а также или\\
\textbf{hnan} [\emph{ʰnän}] 0x6930: вы\\
\textbf{hnap} [\emph{ʰnäp}] 0x4930: ответ\\
\textbf{hnas} [\emph{ʰnäs}] 0x8930: не могу\\
\textbf{hnat} [\emph{ʰnät}] 0xA930: настоящее время напряженный\\
\textbf{hne2c} [\emph{ʰne˩ʃ}] 0xFB30: этнология\\
\textbf{hne2f} [\emph{ʰne˩f}] 0xDB30: анклав\\
\textbf{hne2k} [\emph{ʰne˩k}] 0x3B30: потворствовать\\
\textbf{hne2m} [\emph{ʰne˩m}] 0x1B30: клизма\\
\textbf{hne2n} [\emph{ʰne˩n}] 0x7B30: соловей\\
\textbf{hne2p} [\emph{ʰne˩p}] 0x5B30: эндотелий\\
\textbf{hne2s} [\emph{ʰne˩s}] 0x9B30: нелицензированный\\
\textbf{hne2t} [\emph{ʰne˩t}] 0xBB30: упорно\\
\textbf{hne7k} [\emph{ʰne˥k}] 0x3330: ядра\\
\textbf{hne7m} [\emph{ʰne˥m}] 0x1330: Неемия\\
\textbf{hne7n} [\emph{ʰne˥n}] 0x7330: карликовый\\
\textbf{hne7p} [\emph{ʰne˥p}] 0x5330: несравненный\\
\textbf{hne7s} [\emph{ʰne˥s}] 0x9330: нравоучительный\\
\textbf{hne7t} [\emph{ʰne˥t}] 0xB330: Бенедикт\\
\textbf{hnec} [\emph{ʰneʃ}] 0xEB30: Солнечный свет\\
\textbf{hnef} [\emph{ʰnef}] 0xCB30: Блок - масс\\
\textbf{hnek} [\emph{ʰnek}] 0x2B30: инициатива дело\\
\textbf{hnem} [\emph{ʰnem}] 0x0B30: анимационный\\
\textbf{hnen} [\emph{ʰnen}] 0x6B30: вечный\\
\textbf{hnep} [\emph{ʰnep}] 0x4B30: борьба\\
\textbf{hnes} [\emph{ʰnes}] 0x8B30: нервный\\
\textbf{hnet} [\emph{ʰnet}] 0xAB30: замечание настроение\\
\textbf{hni2c} [\emph{ʰni˩ʃ}] 0xF830: Нил\\
\textbf{hni2f} [\emph{ʰni˩f}] 0xD830: девяностый\\
\textbf{hni2k} [\emph{ʰni˩k}] 0x3830: транспортируемый\\
\textbf{hni2m} [\emph{ʰni˩m}] 0x1830: однопалатный\\
\textbf{hni2n} [\emph{ʰni˩n}] 0x7830: надписи\\
\textbf{hni2p} [\emph{ʰni˩p}] 0x5830: ночной колпак\\
\textbf{hni2s} [\emph{ʰni˩s}] 0x9830: Николай\\
\textbf{hni2t} [\emph{ʰni˩t}] 0xB830: несправедливость\\
\textbf{hni7c} [\emph{ʰni˥ʃ}] 0xF030: сам\\
\textbf{hni7f} [\emph{ʰni˥f}] 0xD030: чистописание\\
\textbf{hni7k} [\emph{ʰni˥k}] 0x3030: негры\\
\textbf{hni7m} [\emph{ʰni˥m}] 0x1030: Бенджамен\\
\textbf{hni7n} [\emph{ʰni˥n}] 0x7030: инбридинг\\
\textbf{hni7p} [\emph{ʰni˥p}] 0x5030: фенотипический\\
\textbf{hni7s} [\emph{ʰni˥s}] 0x9030: отказ от ответственности\\
\textbf{hni7t} [\emph{ʰni˥t}] 0xB030: неотобранный\\
\textbf{hnic} [\emph{ʰniʃ}] 0xE830: неотъемлемое владение\\
\textbf{hnif} [\emph{ʰnif}] 0xC830: усилитель\\
\textbf{hnik} [\emph{ʰnik}] 0x2830: родительный дело\\
\textbf{hnim} [\emph{ʰnim}] 0x0830: имя\\
\textbf{hnin} [\emph{ʰnin}] 0x6830: мужской Пол\\
\textbf{hnip} [\emph{ʰnip}] 0x4830: гномический аспект\\
\textbf{hnis} [\emph{ʰnis}] 0x8830: возраст\\
\textbf{hnit} [\emph{ʰnit}] 0xA830: инфинитив\\
\textbf{hno2c} [\emph{ʰno˩ʃ}] 0xFC30: нанотехнология\\
\textbf{hno2f} [\emph{ʰno˩f}] 0xDC30: и более\\
\textbf{hno2k} [\emph{ʰno˩k}] 0x3C30: неолит\\
\textbf{hno2m} [\emph{ʰno˩m}] 0x1C30: монограмма\\
\textbf{hno2n} [\emph{ʰno˩n}] 0x7C30: нормандский\\
\textbf{hno2p} [\emph{ʰno˩p}] 0x5C30: О.Б.\\
\textbf{hno2s} [\emph{ʰno˩s}] 0x9C30: неумело\\
\textbf{hno2t} [\emph{ʰno˩t}] 0xBC30: умолять\\
\textbf{hno7f} [\emph{ʰno˥f}] 0xD430: девятикратный\\
\textbf{hno7k} [\emph{ʰno˥k}] 0x3430: солнечный удар\\
\textbf{hno7m} [\emph{ʰno˥m}] 0x1430: нерентабельный\\
\textbf{hno7n} [\emph{ʰno˥n}] 0x7430: несправедливый\\
\textbf{hno7p} [\emph{ʰno˥p}] 0x5430: Point Blank\\
\textbf{hno7s} [\emph{ʰno˥s}] 0x9430: безоружный\\
\textbf{hno7t} [\emph{ʰno˥t}] 0xB430: к северо-востоку\\
\textbf{hnoc} [\emph{ʰnoʃ}] 0xEC30: понедельник\\
\textbf{hnof} [\emph{ʰnof}] 0xCC30: курносый\\
\textbf{hnok} [\emph{ʰnok}] 0x2C30: содержание\\
\textbf{hnom} [\emph{ʰnom}] 0x0C30: заговор\\
\textbf{hnon} [\emph{ʰnon}] 0x6C30: союз\\
\textbf{hnop} [\emph{ʰnop}] 0x4C30: монополия\\
\textbf{hnos} [\emph{ʰnos}] 0x8C30: наводнение\\
\textbf{hnot} [\emph{ʰnot}] 0xAC30: Гостиница\\
\textbf{hnu2k} [\emph{ʰnu˩k}] 0x3A30: юго-юго-восток\\
\textbf{hnu2n} [\emph{ʰnu˩n}] 0x7A30: нераспечатанный\\
\textbf{hnu2s} [\emph{ʰnu˩s}] 0x9A30: Индостан\\
\textbf{hnu2t} [\emph{ʰnu˩t}] 0xBA30: унитарный\\
\textbf{hnu7c} [\emph{ʰnu˥ʃ}] 0xF230: слуха\\
\textbf{hnu7k} [\emph{ʰnu˥k}] 0x3230: лук ноги\\
\textbf{hnu7n} [\emph{ʰnu˥n}] 0x7230: Ньюфаундленд\\
\textbf{hnu7p} [\emph{ʰnu˥p}] 0x5230: дышло\\
\textbf{hnu7s} [\emph{ʰnu˥s}] 0x9230: Навуходоносор\\
\textbf{hnu7t} [\emph{ʰnu˥t}] 0xB230: внучата\\
\textbf{hnuc} [\emph{ʰnuʃ}] 0xEA30: номер\\
\textbf{hnuf} [\emph{ʰnuf}] 0xCA30: Венесуэла\\
\textbf{hnuk} [\emph{ʰnuk}] 0x2A30: герцог\\
\textbf{hnum} [\emph{ʰnum}] 0x0A30: разрешающий настроение\\
\textbf{hnun} [\emph{ʰnun}] 0x6A30: нео\\
\textbf{hnup} [\emph{ʰnup}] 0x4A30: раб\\
\textbf{hnus} [\emph{ʰnus}] 0x8A30: нос\\
\textbf{hnut} [\emph{ʰnut}] 0xAA30: контрфактическое условный\\
\textbf{hp62n} [\emph{ʰpə˩n}] 0x7D20: апеллянт\\
\textbf{hp62t} [\emph{ʰpə˩t}] 0xBD20: ипохондрия\\
\textbf{hp67n} [\emph{ʰpə˥n}] 0x7520: перманганат\\
\textbf{hp67s} [\emph{ʰpə˥s}] 0x9520: ханжество\\
\textbf{hp67t} [\emph{ʰpə˥t}] 0xB520: Pitter скороговорка\\
\textbf{hp6c} [\emph{ʰpəʃ}] 0xED20: типологический\\
\textbf{hp6f} [\emph{ʰpəf}] 0xCD20: ржанка\\
\textbf{hp6k} [\emph{ʰpək}] 0x2D20: prolative дело\\
\textbf{hp6m} [\emph{ʰpəm}] 0x0D20: празеодимий\\
\textbf{hp6n} [\emph{ʰpən}] 0x6D20: отправитель\\
\textbf{hp6p} [\emph{ʰpəp}] 0x4D20: полутень\\
\textbf{hp6s} [\emph{ʰpəs}] 0x8D20: протоплазма\\
\textbf{hp6t} [\emph{ʰpət}] 0xAD20: изменение из мероприятие\\
\textbf{hpa2k} [\emph{ʰpä˩k}] 0x3920: шагов\\
\textbf{hpa2m} [\emph{ʰpä˩m}] 0x1920: односторонними\\
\textbf{hpa2n} [\emph{ʰpä˩n}] 0x7920: патогенный микроорганизм\\
\textbf{hpa2p} [\emph{ʰpä˩p}] 0x5920: Пенджаб\\
\textbf{hpa2s} [\emph{ʰpä˩s}] 0x9920: Пифагор\\
\textbf{hpa2t} [\emph{ʰpä˩t}] 0xB920: посаженный\\
\textbf{hpa7c} [\emph{ʰpä˥ʃ}] 0xF120: неестественный\\
\textbf{hpa7f} [\emph{ʰpä˥f}] 0xD120: поташ\\
\textbf{hpa7k} [\emph{ʰpä˥k}] 0x3120: забастовки\\
\textbf{hpa7m} [\emph{ʰpä˥m}] 0x1120: паломник\\
\textbf{hpa7n} [\emph{ʰpä˥n}] 0x7120: фунтов стерлингов\\
\textbf{hpa7p} [\emph{ʰpä˥p}] 0x5120: бледность\\
\textbf{hpa7s} [\emph{ʰpä˥s}] 0x9120: князья\\
\textbf{hpa7t} [\emph{ʰpä˥t}] 0xB120: пахотный слой почвы\\
\textbf{hpac} [\emph{ʰpäʃ}] 0xE920: множественное число номер\\
\textbf{hpaf} [\emph{ʰpäf}] 0xC920: частный\\
\textbf{hpak} [\emph{ʰpäk}] 0x2920: благословляющий дело\\
\textbf{hpam} [\emph{ʰpäm}] 0x0920: первый\\
\textbf{hpan} [\emph{ʰpän}] 0x6920: человек контекст\\
\textbf{hpap} [\emph{ʰpäp}] 0x4920: папа\\
\textbf{hpas} [\emph{ʰpäs}] 0x8920: сын\\
\textbf{hpat} [\emph{ʰpät}] 0xA920: человек\\
\textbf{hpe2k} [\emph{ʰpe˩k}] 0x3B20: клевание\\
\textbf{hpe2m} [\emph{ʰpe˩m}] 0x1B20: плевра\\
\textbf{hpe2n} [\emph{ʰpe˩n}] 0x7B20: эпигенетическое\\
\textbf{hpe2p} [\emph{ʰpe˩p}] 0x5B20: чашелистник\\
\textbf{hpe2s} [\emph{ʰpe˩s}] 0x9B20: Пелопоннес\\
\textbf{hpe2t} [\emph{ʰpe˩t}] 0xBB20: незначительно\\
\textbf{hpe7c} [\emph{ʰpe˥ʃ}] 0xF320: молочай\\
\textbf{hpe7k} [\emph{ʰpe˥k}] 0x3320: Пекин\\
\textbf{hpe7m} [\emph{ʰpe˥m}] 0x1320: грушевый сидр\\
\textbf{hpe7n} [\emph{ʰpe˥n}] 0x7320: версии\\
\textbf{hpe7p} [\emph{ʰpe˥p}] 0x5320: северная часть штата\\
\textbf{hpe7s} [\emph{ʰpe˥s}] 0x9320: Пегас\\
\textbf{hpe7t} [\emph{ʰpe˥t}] 0xB320: пословица\\
\textbf{hpec} [\emph{ʰpeʃ}] 0xEB20: профессия\\
\textbf{hpef} [\emph{ʰpef}] 0xCB20: предикативный\\
\textbf{hpek} [\emph{ʰpek}] 0x2B20: пакет\\
\textbf{hpem} [\emph{ʰpem}] 0x0B20: Пальма\\
\textbf{hpen} [\emph{ʰpen}] 0x6B20: предварительно\\
\textbf{hpep} [\emph{ʰpep}] 0x4B20: яблоко\\
\textbf{hpes} [\emph{ʰpes}] 0x8B20: неприятный\\
\textbf{hpet} [\emph{ʰpet}] 0xAB20: 15\\
\textbf{hpi2c} [\emph{ʰpi˩ʃ}] 0xF820: патчи\\
\textbf{hpi2f} [\emph{ʰpi˩f}] 0xD820: отцеубийство\\
\textbf{hpi2k} [\emph{ʰpi˩k}] 0x3820: эпический\\
\textbf{hpi2m} [\emph{ʰpi˩m}] 0x1820: призма\\
\textbf{hpi2n} [\emph{ʰpi˩n}] 0x7820: пехота\\
\textbf{hpi2p} [\emph{ʰpi˩p}] 0x5820: волынщик\\
\textbf{hpi2s} [\emph{ʰpi˩s}] 0x9820: язычество\\
\textbf{hpi2t} [\emph{ʰpi˩t}] 0xB820: провидение\\
\textbf{hpi7c} [\emph{ʰpi˥ʃ}] 0xF020: плагиат\\
\textbf{hpi7f} [\emph{ʰpi˥f}] 0xD020: скоропортящийся\\
\textbf{hpi7k} [\emph{ʰpi˥k}] 0x3020: португальский\\
\textbf{hpi7m} [\emph{ʰpi˥m}] 0x1020: примат\\
\textbf{hpi7n} [\emph{ʰpi˥n}] 0x7020: нетерпение\\
\textbf{hpi7p} [\emph{ʰpi˥p}] 0x5020: непроходимый\\
\textbf{hpi7s} [\emph{ʰpi˥s}] 0x9020: Пруссия\\
\textbf{hpi7t} [\emph{ʰpi˥t}] 0xB020: Питер\\
\textbf{hpic} [\emph{ʰpiʃ}] 0xE820: оптатив настроение\\
\textbf{hpif} [\emph{ʰpif}] 0xC820: страница\\
\textbf{hpik} [\emph{ʰpik}] 0x2820: аблятив дело\\
\textbf{hpim} [\emph{ʰpim}] 0x0820: императив настроение\\
\textbf{hpin} [\emph{ʰpin}] 0x6820: adpositional дело\\
\textbf{hpip} [\emph{ʰpip}] 0x4820: прогрессивный аспект\\
\textbf{hpis} [\emph{ʰpis}] 0x8820: платить\\
\textbf{hpit} [\emph{ʰpit}] 0xA820: белый\\
\textbf{hpo2k} [\emph{ʰpo˩k}] 0x3C20: гной\\
\textbf{hpo2n} [\emph{ʰpo˩n}] 0x7C20: форма волны\\
\textbf{hpo2p} [\emph{ʰpo˩p}] 0x5C20: попурри\\
\textbf{hpo2s} [\emph{ʰpo˩s}] 0x9C20: причалы\\
\textbf{hpo2t} [\emph{ʰpo˩t}] 0xBC20: чистилище\\
\textbf{hpo7k} [\emph{ʰpo˥k}] 0x3420: щенок\\
\textbf{hpo7m} [\emph{ʰpo˥m}] 0x1420: OMA\\
\textbf{hpo7n} [\emph{ʰpo˥n}] 0x7420: Аполлон\\
\textbf{hpo7p} [\emph{ʰpo˥p}] 0x5420: пропан\\
\textbf{hpo7s} [\emph{ʰpo˥s}] 0x9420: Прометей\\
\textbf{hpo7t} [\emph{ʰpo˥t}] 0xB420: почтмейстер\\
\textbf{hpoc} [\emph{ʰpoʃ}] 0xEC20: осень\\
\textbf{hpof} [\emph{ʰpof}] 0xCC20: Попкорн\\
\textbf{hpok} [\emph{ʰpok}] 0x2C20: племя\\
\textbf{hpom} [\emph{ʰpom}] 0x0C20: Платформа\\
\textbf{hpon} [\emph{ʰpon}] 0x6C20: провинция\\
\textbf{hpop} [\emph{ʰpop}] 0x4C20: опера\\
\textbf{hpos} [\emph{ʰpos}] 0x8C20: сосна\\
\textbf{hpot} [\emph{ʰpot}] 0xAC20: доступный\\
\textbf{hpu2k} [\emph{ʰpu˩k}] 0x3A20: лобковый\\
\textbf{hpu2n} [\emph{ʰpu˩n}] 0x7A20: в верхней части города\\
\textbf{hpu2t} [\emph{ʰpu˩t}] 0xBA20: депутация\\
\textbf{hpu7c} [\emph{ʰpu˥ʃ}] 0xF220: голубь косолапый\\
\textbf{hpu7k} [\emph{ʰpu˥k}] 0x3220: эпикуреец\\
\textbf{hpu7n} [\emph{ʰpu˥n}] 0x7220: пунктуация\\
\textbf{hpu7p} [\emph{ʰpu˥p}] 0x5220: Pooh Pooh\\
\textbf{hpu7s} [\emph{ʰpu˥s}] 0x9220: палеозойский\\
\textbf{hpu7t} [\emph{ʰpu˥t}] 0xB220: рассыпалась\\
\textbf{hpuc} [\emph{ʰpuʃ}] 0xEA20: напиток\\
\textbf{hpuf} [\emph{ʰpuf}] 0xCA20: пикколо\\
\textbf{hpuk} [\emph{ʰpuk}] 0x2A20: книга\\
\textbf{hpum} [\emph{ʰpum}] 0x0A20: поверхность\\
\textbf{hpun} [\emph{ʰpun}] 0x6A20: спросил\\
\textbf{hpup} [\emph{ʰpup}] 0x4A20: хлеб\\
\textbf{hpus} [\emph{ʰpus}] 0x8A20: офис\\
\textbf{hput} [\emph{ʰput}] 0xAA20: Помогите\\
\textbf{hq6k} [\emph{ʰŋək}] 0x2DC8: упираться\\
\textbf{hq6m} [\emph{ʰŋəm}] 0x0DC8: болен omened\\
\textbf{hq6n} [\emph{ʰŋən}] 0x6DC8: наперсток\\
\textbf{hq6p} [\emph{ʰŋəp}] 0x4DC8: изобара\\
\textbf{hq6s} [\emph{ʰŋəs}] 0x8DC8: наоборот\\
\textbf{hq6t} [\emph{ʰŋət}] 0xADC8: Аристотель\\
\textbf{hqa2k} [\emph{ʰŋä˩k}] 0x39C8: долговязый\\
\textbf{hqa2n} [\emph{ʰŋä˩n}] 0x79C8: Восходящий\\
\textbf{hqa2p} [\emph{ʰŋä˩p}] 0x59C8: превосходить численно\\
\textbf{hqa2s} [\emph{ʰŋä˩s}] 0x99C8: абсцисса\\
\textbf{hqa2t} [\emph{ʰŋä˩t}] 0xB9C8: мангал\\
\textbf{hqa7c} [\emph{ʰŋä˥ʃ}] 0xF1C8: бачки\\
\textbf{hqa7k} [\emph{ʰŋä˥k}] 0x31C8: ткач\\
\textbf{hqa7m} [\emph{ʰŋä˥m}] 0x11C8: прабабушка\\
\textbf{hqa7n} [\emph{ʰŋä˥n}] 0x71C8: мечтатель\\
\textbf{hqa7p} [\emph{ʰŋä˥p}] 0x51C8: черное и белое\\
\textbf{hqa7s} [\emph{ʰŋä˥s}] 0x91C8: сумасшедший дом\\
\textbf{hqa7t} [\emph{ʰŋä˥t}] 0xB1C8: тратить\\
\textbf{hqac} [\emph{ʰŋäʃ}] 0xE9C8: пропадать\\
\textbf{hqaf} [\emph{ʰŋäf}] 0xC9C8: окалина\\
\textbf{hqak} [\emph{ʰŋäk}] 0x29C8: валять дурака\\
\textbf{hqam} [\emph{ʰŋäm}] 0x09C8: разрабатывать\\
\textbf{hqan} [\emph{ʰŋän}] 0x69C8: длина волны\\
\textbf{hqap} [\emph{ʰŋäp}] 0x49C8: настоящее время причастие\\
\textbf{hqas} [\emph{ʰŋäs}] 0x89C8: Веб-сайт\\
\textbf{hqat} [\emph{ʰŋät}] 0xA9C8: время обеда\\
\textbf{hqe2t} [\emph{ʰŋe˩t}] 0xBBC8: ястреб глазами\\
\textbf{hqe7n} [\emph{ʰŋe˥n}] 0x73C8: англичанка\\
\textbf{hqe7t} [\emph{ʰŋe˥t}] 0xB3C8: сычужок\\
\textbf{hqec} [\emph{ʰŋeʃ}] 0xEBC8: бесполый\\
\textbf{hqef} [\emph{ʰŋef}] 0xCBC8: теория относительности\\
\textbf{hqek} [\emph{ʰŋek}] 0x2BC8: косой дело\\
\textbf{hqem} [\emph{ʰŋem}] 0x0BC8: изотерма\\
\textbf{hqen} [\emph{ʰŋen}] 0x6BC8: звезда Пол\\
\textbf{hqep} [\emph{ʰŋep}] 0x4BC8: леденец\\
\textbf{hqes} [\emph{ʰŋes}] 0x8BC8: сексуальный\\
\textbf{hqet} [\emph{ʰŋet}] 0xABC8: розничная торговля магазин\\
\textbf{hqi2n} [\emph{ʰŋi˩n}] 0x78C8: звенящий\\
\textbf{hqi2s} [\emph{ʰŋi˩s}] 0x98C8: Рыбы\\
\textbf{hqi2t} [\emph{ʰŋi˩t}] 0xB8C8: самодельный\\
\textbf{hqi7f} [\emph{ʰŋi˥f}] 0xD0C8: орфографический\\
\textbf{hqi7k} [\emph{ʰŋi˥k}] 0x30C8: подписан\\
\textbf{hqi7n} [\emph{ʰŋi˥n}] 0x70C8: неслышный\\
\textbf{hqi7p} [\emph{ʰŋi˥p}] 0x50C8: нелюдимый\\
\textbf{hqi7s} [\emph{ʰŋi˥s}] 0x90C8: слоновость\\
\textbf{hqi7t} [\emph{ʰŋi˥t}] 0xB0C8: чернильница\\
\textbf{hqic} [\emph{ʰŋiʃ}] 0xE8C8: фишинг\\
\textbf{hqif} [\emph{ʰŋif}] 0xC8C8: Венгрия\\
\textbf{hqik} [\emph{ʰŋik}] 0x28C8: вводить\\
\textbf{hqim} [\emph{ʰŋim}] 0x08C8: цитата форма\\
\textbf{hqin} [\emph{ʰŋin}] 0x68C8: нелинейный\\
\textbf{hqip} [\emph{ʰŋip}] 0x48C8: облако вычисления\\
\textbf{hqis} [\emph{ʰŋis}] 0x88C8: шовинист\\
\textbf{hqit} [\emph{ʰŋit}] 0xA8C8: Спасатель\\
\textbf{hqo2n} [\emph{ʰŋo˩n}] 0x7CC8: фотогеничный\\
\textbf{hqo2t} [\emph{ʰŋo˩t}] 0xBCC8: сопли\\
\textbf{hqo7k} [\emph{ʰŋo˥k}] 0x34C8: Anglo католическое\\
\textbf{hqo7n} [\emph{ʰŋo˥n}] 0x74C8: ацетон\\
\textbf{hqo7t} [\emph{ʰŋo˥t}] 0xB4C8: Anglo индийский\\
\textbf{hqoc} [\emph{ʰŋoʃ}] 0xECC8: тошнота\\
\textbf{hqok} [\emph{ʰŋok}] 0x2CC8: друг друга\\
\textbf{hqom} [\emph{ʰŋom}] 0x0CC8: ром\\
\textbf{hqon} [\emph{ʰŋon}] 0x6CC8: орегано\\
\textbf{hqop} [\emph{ʰŋop}] 0x4CC8: Сингапур\\
\textbf{hqos} [\emph{ʰŋos}] 0x8CC8: унция\\
\textbf{hqot} [\emph{ʰŋot}] 0xACC8: растворенное вещество\\
\textbf{hqu7n} [\emph{ʰŋu˥n}] 0x72C8: недоедание\\
\textbf{hqu7t} [\emph{ʰŋu˥t}] 0xB2C8: заготовку Doux\\
\textbf{hquc} [\emph{ʰŋuʃ}] 0xEAC8: скульптор\\
\textbf{hquk} [\emph{ʰŋuk}] 0x2AC8: кенгуру\\
\textbf{hqum} [\emph{ʰŋum}] 0x0AC8: имя пользователя\\
\textbf{hqun} [\emph{ʰŋun}] 0x6AC8: прерывистый\\
\textbf{hqup} [\emph{ʰŋup}] 0x4AC8: тупо\\
\textbf{hqus} [\emph{ʰŋus}] 0x8AC8: неисчислимый\\
\textbf{hqut} [\emph{ʰŋut}] 0xAAC8: журнал вне\\
\textbf{hr62k} [\emph{ʰrə˩k}] 0x3D70: сверкать\\
\textbf{hr62n} [\emph{ʰrə˩n}] 0x7D70: батут\\
\textbf{hr62t} [\emph{ʰrə˩t}] 0xBD70: геральдика\\
\textbf{hr67k} [\emph{ʰrə˥k}] 0x3570: ушная сера\\
\textbf{hr67n} [\emph{ʰrə˥n}] 0x7570: подержанный\\
\textbf{hr67s} [\emph{ʰrə˥s}] 0x9570: тромбоз\\
\textbf{hr67t} [\emph{ʰrə˥t}] 0xB570: ушная боль\\
\textbf{hr6c} [\emph{ʰrəʃ}] 0xED70: двадцать семь\\
\textbf{hr6f} [\emph{ʰrəf}] 0xCD70: тропосфера\\
\textbf{hr6k} [\emph{ʰrək}] 0x2D70: рецензент\\
\textbf{hr6m} [\emph{ʰrəm}] 0x0D70: рений\\
\textbf{hr6n} [\emph{ʰrən}] 0x6D70: рацион\\
\textbf{hr6p} [\emph{ʰrəp}] 0x4D70: ромб\\
\textbf{hr6s} [\emph{ʰrəs}] 0x8D70: направления\\
\textbf{hr6t} [\emph{ʰrət}] 0xAD70: артефакт\\
\textbf{hra2c} [\emph{ʰrä˩ʃ}] 0xF970: арбитраж\\
\textbf{hra2f} [\emph{ʰrä˩f}] 0xD970: архивариус\\
\textbf{hra2k} [\emph{ʰrä˩k}] 0x3970: Ирак\\
\textbf{hra2m} [\emph{ʰrä˩m}] 0x1970: драхма\\
\textbf{hra2n} [\emph{ʰrä˩n}] 0x7970: ряды\\
\textbf{hra2p} [\emph{ʰrä˩p}] 0x5970: рецидив\\
\textbf{hra2s} [\emph{ʰrä˩s}] 0x9970: преступления\\
\textbf{hra2t} [\emph{ʰrä˩t}] 0xB970: перекрестная проверка\\
\textbf{hra7c} [\emph{ʰrä˥ʃ}] 0xF170: реакционер\\
\textbf{hra7f} [\emph{ʰrä˥f}] 0xD170: гуртовщик\\
\textbf{hra7k} [\emph{ʰrä˥k}] 0x3170: архиепископ\\
\textbf{hra7m} [\emph{ʰrä˥m}] 0x1170: старое время\\
\textbf{hra7n} [\emph{ʰrä˥n}] 0x7170: романс\\
\textbf{hra7p} [\emph{ʰrä˥p}] 0x5170: арабский\\
\textbf{hra7s} [\emph{ʰrä˥s}] 0x9170: проза\\
\textbf{hra7t} [\emph{ʰrä˥t}] 0xB170: реостат\\
\textbf{hrac} [\emph{ʰräʃ}] 0xE970: везде\\
\textbf{hraf} [\emph{ʰräf}] 0xC970: трафик\\
\textbf{hrak} [\emph{ʰräk}] 0x2970: сотрудники\\
\textbf{hram} [\emph{ʰräm}] 0x0970: романтический\\
\textbf{hran} [\emph{ʰrän}] 0x6970: совесть\\
\textbf{hrap} [\emph{ʰräp}] 0x4970: артиллерия\\
\textbf{hras} [\emph{ʰräs}] 0x8970: песок\\
\textbf{hrat} [\emph{ʰrät}] 0xA970: крыша\\
\textbf{hre2c} [\emph{ʰre˩ʃ}] 0xFB70: в форме полумесяца\\
\textbf{hre2k} [\emph{ʰre˩k}] 0x3B70: крокет\\
\textbf{hre2m} [\emph{ʰre˩m}] 0x1B70: брюшина\\
\textbf{hre2n} [\emph{ʰre˩n}] 0x7B70: полки\\
\textbf{hre2p} [\emph{ʰre˩p}] 0x5B70: повторно заложен\\
\textbf{hre2s} [\emph{ʰre˩s}] 0x9B70: смола\\
\textbf{hre2t} [\emph{ʰre˩t}] 0xBB70: приготовился\\
\textbf{hre7c} [\emph{ʰre˥ʃ}] 0xF370: оружейник\\
\textbf{hre7f} [\emph{ʰre˥f}] 0xD370: выпрямитель\\
\textbf{hre7k} [\emph{ʰre˥k}] 0x3370: веснушчатый\\
\textbf{hre7m} [\emph{ʰre˥m}] 0x1370: бремен\\
\textbf{hre7n} [\emph{ʰre˥n}] 0x7370: передний план\\
\textbf{hre7p} [\emph{ʰre˥p}] 0x5370: кренделек\\
\textbf{hre7s} [\emph{ʰre˥s}] 0x9370: Фракия\\
\textbf{hre7t} [\emph{ʰre˥t}] 0xB370: лучистый\\
\textbf{hrec} [\emph{ʰreʃ}] 0xEB70: справочный\\
\textbf{href} [\emph{ʰref}] 0xCB70: разреженный\\
\textbf{hrek} [\emph{ʰrek}] 0x2B70: восстановить\\
\textbf{hrem} [\emph{ʰrem}] 0x0B70: луч\\
\textbf{hren} [\emph{ʰren}] 0x6B70: речь контекст\\
\textbf{hrep} [\emph{ʰrep}] 0x4B70: рецепт\\
\textbf{hres} [\emph{ʰres}] 0x8B70: походить\\
\textbf{hret} [\emph{ʰret}] 0xAB70: твердо\\
\textbf{hri2c} [\emph{ʰri˩ʃ}] 0xF870: гребни\\
\textbf{hri2f} [\emph{ʰri˩f}] 0xD870: рекурсивный\\
\textbf{hri2k} [\emph{ʰri˩k}] 0x3870: двадцать один\\
\textbf{hri2m} [\emph{ʰri˩m}] 0x1870: общинники\\
\textbf{hri2n} [\emph{ʰri˩n}] 0x7870: оживление\\
\textbf{hri2p} [\emph{ʰri˩p}] 0x5870: прокатка штырь\\
\textbf{hri2s} [\emph{ʰri˩s}] 0x9870: двадцать четыре\\
\textbf{hri2t} [\emph{ʰri˩t}] 0xB870: вероломство\\
\textbf{hri7c} [\emph{ʰri˥ʃ}] 0xF070: негодяй\\
\textbf{hri7f} [\emph{ʰri˥f}] 0xD070: грифон\\
\textbf{hri7k} [\emph{ʰri˥k}] 0x3070: афроамериканец\\
\textbf{hri7m} [\emph{ʰri˥m}] 0x1070: предчувствие\\
\textbf{hri7n} [\emph{ʰri˥n}] 0x7070: Виргиния\\
\textbf{hri7p} [\emph{ʰri˥p}] 0x5070: сорок пять\\
\textbf{hri7s} [\emph{ʰri˥s}] 0x9070: директора\\
\textbf{hri7t} [\emph{ʰri˥t}] 0xB070: Христос\\
\textbf{hric} [\emph{ʰriʃ}] 0xE870: постепенно\\
\textbf{hrif} [\emph{ʰrif}] 0xC870: примитивный\\
\textbf{hrik} [\emph{ʰrik}] 0x2870: практическое\\
\textbf{hrim} [\emph{ʰrim}] 0x0870: христианство\\
\textbf{hrin} [\emph{ʰrin}] 0x6870: рациональный Пол\\
\textbf{hrip} [\emph{ʰrip}] 0x4870: краткое\\
\textbf{hris} [\emph{ʰris}] 0x8870: обычай\\
\textbf{hrit} [\emph{ʰrit}] 0xA870: кредит\\
\textbf{hro2c} [\emph{ʰro˩ʃ}] 0xFC70: Роджер\\
\textbf{hro2f} [\emph{ʰro˩f}] 0xDC70: хромосфера\\
\textbf{hro2k} [\emph{ʰro˩k}] 0x3C70: тряпье\\
\textbf{hro2m} [\emph{ʰro˩m}] 0x1C70: аэродинамический\\
\textbf{hro2n} [\emph{ʰro˩n}] 0x7C70: троянец\\
\textbf{hro2p} [\emph{ʰro˩p}] 0x5C70: ларингоскоп\\
\textbf{hro2s} [\emph{ʰro˩s}] 0x9C70: Родос\\
\textbf{hro2t} [\emph{ʰro˩t}] 0xBC70: ректор\\
\textbf{hro7c} [\emph{ʰro˥ʃ}] 0xF470: конский волос\\
\textbf{hro7f} [\emph{ʰro˥f}] 0xD470: цапфа\\
\textbf{hro7k} [\emph{ʰro˥k}] 0x3470: раскаленный\\
\textbf{hro7m} [\emph{ʰro˥m}] 0x1470: золоченая бронза\\
\textbf{hro7n} [\emph{ʰro˥n}] 0x7470: скрещивание\\
\textbf{hro7p} [\emph{ʰro˥p}] 0x5470: реторты\\
\textbf{hro7s} [\emph{ʰro˥s}] 0x9470: чалый\\
\textbf{hro7t} [\emph{ʰro˥t}] 0xB470: подруливающее\\
\textbf{hroc} [\emph{ʰroʃ}] 0xEC70: традиция\\
\textbf{hrof} [\emph{ʰrof}] 0xCC70: троянец лошадь\\
\textbf{hrok} [\emph{ʰrok}] 0x2C70: реформа\\
\textbf{hrom} [\emph{ʰrom}] 0x0C70: программа\\
\textbf{hron} [\emph{ʰron}] 0x6C70: цвет\\
\textbf{hrop} [\emph{ʰrop}] 0x4C70: веревка\\
\textbf{hros} [\emph{ʰros}] 0x8C70: репортер\\
\textbf{hrot} [\emph{ʰrot}] 0xAC70: недавний мимо напряженный\\
\textbf{hru2k} [\emph{ʰru˩k}] 0x3A70: рутил\\
\textbf{hru2n} [\emph{ʰru˩n}] 0x7A70: маршруты\\
\textbf{hru2p} [\emph{ʰru˩p}] 0x5A70: рубль\\
\textbf{hru2s} [\emph{ʰru˩s}] 0x9A70: заключать в себе\\
\textbf{hru2t} [\emph{ʰru˩t}] 0xBA70: колесница\\
\textbf{hru7c} [\emph{ʰru˥ʃ}] 0xF270: мочеполовой\\
\textbf{hru7k} [\emph{ʰru˥k}] 0x3270: Дракула\\
\textbf{hru7m} [\emph{ʰru˥m}] 0x1270: Каникулы в гареме\\
\textbf{hru7n} [\emph{ʰru˥n}] 0x7270: лучник\\
\textbf{hru7p} [\emph{ʰru˥p}] 0x5270: сумчатый\\
\textbf{hru7s} [\emph{ʰru˥s}] 0x9270: Брюссель\\
\textbf{hru7t} [\emph{ʰru˥t}] 0xB270: Рут\\
\textbf{hruc} [\emph{ʰruʃ}] 0xEA70: грубый\\
\textbf{hruf} [\emph{ʰruf}] 0xCA70: грейпфрут\\
\textbf{hruk} [\emph{ʰruk}] 0x2A70: архитектура\\
\textbf{hrum} [\emph{ʰrum}] 0x0A70: ремонт\\
\textbf{hrun} [\emph{ʰrun}] 0x6A70: рулон\\
\textbf{hrup} [\emph{ʰrup}] 0x4A70: коррупция\\
\textbf{hrus} [\emph{ʰrus}] 0x8A70: принцесса\\
\textbf{hrut} [\emph{ʰrut}] 0xAA70: аргумент\\
\textbf{hs62n} [\emph{ʰsə˩n}] 0x7D48: осетр\\
\textbf{hs62t} [\emph{ʰsə˩t}] 0xBD48: грамотный\\
\textbf{hs67k} [\emph{ʰsə˥k}] 0x3548: красноломкий\\
\textbf{hs67n} [\emph{ʰsə˥n}] 0x7548: эоцен\\
\textbf{hs67s} [\emph{ʰsə˥s}] 0x9548: анастомоз\\
\textbf{hs67t} [\emph{ʰsə˥t}] 0xB548: приспособленного\\
\textbf{hs6c} [\emph{ʰsəʃ}] 0xED48: само уверенность\\
\textbf{hs6f} [\emph{ʰsəf}] 0xCD48: синтезировать\\
\textbf{hs6k} [\emph{ʰsək}] 0x2D48: поиск двигатель\\
\textbf{hs6m} [\emph{ʰsəm}] 0x0D48: цезий\\
\textbf{hs6n} [\emph{ʰsən}] 0x6D48: продавцы\\
\textbf{hs6p} [\emph{ʰsəp}] 0x4D48: психотерапия\\
\textbf{hs6s} [\emph{ʰsəs}] 0x8D48: лужи\\
\textbf{hs6t} [\emph{ʰsət}] 0xAD48: комната отдыха\\
\textbf{hsa2c} [\emph{ʰsä˩ʃ}] 0xF948: немытый\\
\textbf{hsa2f} [\emph{ʰsä˩f}] 0xD948: семафор\\
\textbf{hsa2k} [\emph{ʰsä˩k}] 0x3948: пассажиры\\
\textbf{hsa2m} [\emph{ʰsä˩m}] 0x1948: цианид\\
\textbf{hsa2n} [\emph{ʰsä˩n}] 0x7948: господа\\
\textbf{hsa2p} [\emph{ʰsä˩p}] 0x5948: в средней части судна\\
\textbf{hsa2s} [\emph{ʰsä˩s}] 0x9948: сольный концерт\\
\textbf{hsa2t} [\emph{ʰsä˩t}] 0xB948: над урожаем\\
\textbf{hsa7c} [\emph{ʰsä˥ʃ}] 0xF148: дикари\\
\textbf{hsa7f} [\emph{ʰsä˥f}] 0xD148: прадед\\
\textbf{hsa7k} [\emph{ʰsä˥k}] 0x3148: ударив\\
\textbf{hsa7m} [\emph{ʰsä˥m}] 0x1148: праздно\\
\textbf{hsa7n} [\emph{ʰsä˥n}] 0x7148: поместья\\
\textbf{hsa7p} [\emph{ʰsä˥p}] 0x5148: стартапы\\
\textbf{hsa7s} [\emph{ʰsä˥s}] 0x9148: саксонский\\
\textbf{hsa7t} [\emph{ʰsä˥t}] 0xB148: улицы\\
\textbf{hsac} [\emph{ʰsäʃ}] 0xE948: сказать\\
\textbf{hsaf} [\emph{ʰsäf}] 0xC948: дыхание\\
\textbf{hsak} [\emph{ʰsäk}] 0x2948: время контекст\\
\textbf{hsam} [\emph{ʰsäm}] 0x0948: муж\\
\textbf{hsan} [\emph{ʰsän}] 0x6948: поверхность контекст\\
\textbf{hsap} [\emph{ʰsäp}] 0x4948: Кроме\\
\textbf{hsas} [\emph{ʰsäs}] 0x8948: утро\\
\textbf{hsat} [\emph{ʰsät}] 0xA948: видеть\\
\textbf{hse2c} [\emph{ʰse˩ʃ}] 0xFB48: эстетически\\
\textbf{hse2k} [\emph{ʰse˩k}] 0x3B48: осока\\
\textbf{hse2m} [\emph{ʰse˩m}] 0x1B48: шестнадцатеричный\\
\textbf{hse2n} [\emph{ʰse˩n}] 0x7B48: святые\\
\textbf{hse2p} [\emph{ʰse˩p}] 0x5B48: десятикратный\\
\textbf{hse2s} [\emph{ʰse˩s}] 0x9B48: показной\\
\textbf{hse2t} [\emph{ʰse˩t}] 0xBB48: эскадрилья\\
\textbf{hse7c} [\emph{ʰse˥ʃ}] 0xF348: мозжечок\\
\textbf{hse7f} [\emph{ʰse˥f}] 0xD348: лексикограф\\
\textbf{hse7k} [\emph{ʰse˥k}] 0x3348: всевидящий\\
\textbf{hse7m} [\emph{ʰse˥m}] 0x1348: эскимос\\
\textbf{hse7n} [\emph{ʰse˥n}] 0x7348: невод\\
\textbf{hse7p} [\emph{ʰse˥p}] 0x5348: сепия\\
\textbf{hse7s} [\emph{ʰse˥s}] 0x9348: сорок семь\\
\textbf{hse7t} [\emph{ʰse˥t}] 0xB348: шестьдесят пять\\
\textbf{hsec} [\emph{ʰseʃ}] 0xEB48: программного обеспечения\\
\textbf{hsef} [\emph{ʰsef}] 0xCB48: змея\\
\textbf{hsek} [\emph{ʰsek}] 0x2B48: аспект\\
\textbf{hsem} [\emph{ʰsem}] 0x0B48: саммит\\
\textbf{hsen} [\emph{ʰsen}] 0x6B48: экспедиция\\
\textbf{hsep} [\emph{ʰsep}] 0x4B48: сентябрь\\
\textbf{hses} [\emph{ʰses}] 0x8B48: 14\\
\textbf{hset} [\emph{ʰset}] 0xAB48: центр\\
\textbf{hsi2f} [\emph{ʰsi˩f}] 0xD848: севилья\\
\textbf{hsi2k} [\emph{ʰsi˩k}] 0x3848: пространственно временная\\
\textbf{hsi2m} [\emph{ʰsi˩m}] 0x1848: проницательность\\
\textbf{hsi2n} [\emph{ʰsi˩n}] 0x7848: самый широкий\\
\textbf{hsi2p} [\emph{ʰsi˩p}] 0x5848: бекас\\
\textbf{hsi2s} [\emph{ʰsi˩s}] 0x9848: Полегче\\
\textbf{hsi2t} [\emph{ʰsi˩t}] 0xB848: гениталии\\
\textbf{hsi7c} [\emph{ʰsi˥ʃ}] 0xF048: сатира\\
\textbf{hsi7f} [\emph{ʰsi˥f}] 0xD048: непродуктивный\\
\textbf{hsi7k} [\emph{ʰsi˥k}] 0x3048: двадцать пять\\
\textbf{hsi7m} [\emph{ʰsi˥m}] 0x1048: Ишмаэль\\
\textbf{hsi7n} [\emph{ʰsi˥n}] 0x7048: самостоятельного полива\\
\textbf{hsi7p} [\emph{ʰsi˥p}] 0x5048: Константинопольский\\
\textbf{hsi7s} [\emph{ʰsi˥s}] 0x9048: эгоизм\\
\textbf{hsi7t} [\emph{ʰsi˥t}] 0xB048: обезжиренное\\
\textbf{hsic} [\emph{ʰsiʃ}] 0xE848: Вот\\
\textbf{hsif} [\emph{ʰsif}] 0xC848: лев\\
\textbf{hsik} [\emph{ʰsik}] 0x2848: общественного\\
\textbf{hsim} [\emph{ʰsim}] 0x0848: спать\\
\textbf{hsin} [\emph{ʰsin}] 0x6848: Социальное контекст\\
\textbf{hsip} [\emph{ʰsip}] 0x4848: семь\\
\textbf{hsis} [\emph{ʰsis}] 0x8848: всегда\\
\textbf{hsit} [\emph{ʰsit}] 0xA848: эпистемическое настроение\\
\textbf{hso2c} [\emph{ʰso˩ʃ}] 0xFC48: рустовка\\
\textbf{hso2k} [\emph{ʰso˩k}] 0x3C48: запоздалая мысль\\
\textbf{hso2m} [\emph{ʰso˩m}] 0x1C48: мазохизм\\
\textbf{hso2n} [\emph{ʰso˩n}] 0x7C48: источники\\
\textbf{hso2p} [\emph{ʰso˩p}] 0x5C48: снежный шар\\
\textbf{hso2s} [\emph{ʰso˩s}] 0x9C48: мускус\\
\textbf{hso2t} [\emph{ʰso˩t}] 0xBC48: скотч\\
\textbf{hso7c} [\emph{ʰso˥ʃ}] 0xF448: соте\\
\textbf{hso7f} [\emph{ʰso˥f}] 0xD448: Самоанский\\
\textbf{hso7k} [\emph{ʰso˥k}] 0x3448: вяленое мясо\\
\textbf{hso7m} [\emph{ʰso˥m}] 0x1448: стоицизм\\
\textbf{hso7n} [\emph{ʰso˥n}] 0x7448: гимн\\
\textbf{hso7p} [\emph{ʰso˥p}] 0x5448: левый борт\\
\textbf{hso7s} [\emph{ʰso˥s}] 0x9448: шумно\\
\textbf{hso7t} [\emph{ʰso˥t}] 0xB448: содрогнулся\\
\textbf{hsoc} [\emph{ʰsoʃ}] 0xEC48: акции\\
\textbf{hsof} [\emph{ʰsof}] 0xCC48: прибой\\
\textbf{hsok} [\emph{ʰsok}] 0x2C48: поставка\\
\textbf{hsom} [\emph{ʰsom}] 0x0C48: поваренная соль\\
\textbf{hson} [\emph{ʰson}] 0x6C48: помилование\\
\textbf{hsop} [\emph{ʰsop}] 0x4C48: успешно\\
\textbf{hsos} [\emph{ʰsos}] 0x8C48: шестнадцать\\
\textbf{hsot} [\emph{ʰsot}] 0xAC48: скорость\\
\textbf{hsu2f} [\emph{ʰsu˩f}] 0xDA48: беспутный\\
\textbf{hsu2k} [\emph{ʰsu˩k}] 0x3A48: скунс\\
\textbf{hsu2m} [\emph{ʰsu˩m}] 0x1A48: сумах\\
\textbf{hsu2n} [\emph{ʰsu˩n}] 0x7A48: более чистый\\
\textbf{hsu2p} [\emph{ʰsu˩p}] 0x5A48: шаг гуся\\
\textbf{hsu2s} [\emph{ʰsu˩s}] 0x9A48: Советский Союз\\
\textbf{hsu2t} [\emph{ʰsu˩t}] 0xBA48: превратности\\
\textbf{hsu7c} [\emph{ʰsu˥ʃ}] 0xF248: снегоступ\\
\textbf{hsu7k} [\emph{ʰsu˥k}] 0x3248: саго\\
\textbf{hsu7m} [\emph{ʰsu˥m}] 0x1248: грудина\\
\textbf{hsu7n} [\emph{ʰsu˥n}] 0x7248: двадцать девять\\
\textbf{hsu7p} [\emph{ʰsu˥p}] 0x5248: стократный\\
\textbf{hsu7s} [\emph{ʰsu˥s}] 0x9248: вшей\\
\textbf{hsu7t} [\emph{ʰsu˥t}] 0xB248: заблудившись\\
\textbf{hsuc} [\emph{ʰsuʃ}] 0xEA48: предок\\
\textbf{hsuf} [\emph{ʰsuf}] 0xCA48: удаление\\
\textbf{hsuk} [\emph{ʰsuk}] 0x2A48: источник дело\\
\textbf{hsum} [\emph{ʰsum}] 0x0A48: притяжательный падеж маркер\\
\textbf{hsun} [\emph{ʰsun}] 0x6A48: рука\\
\textbf{hsup} [\emph{ʰsup}] 0x4A48: содержащий окаменелости Пол\\
\textbf{hsus} [\emph{ʰsus}] 0x8A48: неделю\\
\textbf{hsut} [\emph{ʰsut}] 0xAA48: все\\
\textbf{ht62n} [\emph{ʰtə˩n}] 0x7D50: пропаривания\\
\textbf{ht62t} [\emph{ʰtə˩t}] 0xBD50: семьдесят два\\
\textbf{ht67k} [\emph{ʰtə˥k}] 0x3550: самоиндукции\\
\textbf{ht67m} [\emph{ʰtə˥m}] 0x1550: Траляля и Труляля\\
\textbf{ht67n} [\emph{ʰtə˥n}] 0x7550: утолщение\\
\textbf{ht67t} [\emph{ʰtə˥t}] 0xB550: пораженный\\
\textbf{ht6c} [\emph{ʰtəʃ}] 0xED50: стебель\\
\textbf{ht6f} [\emph{ʰtəf}] 0xCD50: Roto румпель\\
\textbf{ht6k} [\emph{ʰtək}] 0x2D50: придирчивый\\
\textbf{ht6m} [\emph{ʰtəm}] 0x0D50: тербий атом\\
\textbf{ht6n} [\emph{ʰtən}] 0x6D50: Терминал\\
\textbf{ht6p} [\emph{ʰtəp}] 0x4D50: шпиль\\
\textbf{ht6s} [\emph{ʰtəs}] 0x8D50: аутизм\\
\textbf{ht6t} [\emph{ʰtət}] 0xAD50: реальный время\\
\textbf{hta2c} [\emph{ʰtä˩ʃ}] 0xF950: даты\\
\textbf{hta2f} [\emph{ʰtä˩f}] 0xD950: дрейфовать\\
\textbf{hta2k} [\emph{ʰtä˩k}] 0x3950: небесный свод\\
\textbf{hta2m} [\emph{ʰtä˩m}] 0x1950: безвременно\\
\textbf{hta2n} [\emph{ʰtä˩n}] 0x7950: учил\\
\textbf{hta2p} [\emph{ʰtä˩p}] 0x5950: приписывал\\
\textbf{hta2s} [\emph{ʰtä˩s}] 0x9950: взносы\\
\textbf{hta2t} [\emph{ʰtä˩t}] 0xB950: полеводства\\
\textbf{hta7c} [\emph{ʰtä˥ʃ}] 0xF150: антидот\\
\textbf{hta7f} [\emph{ʰtä˥f}] 0xD150: пустяки\\
\textbf{hta7k} [\emph{ʰtä˥k}] 0x3150: ультразвуковая эхография\\
\textbf{hta7m} [\emph{ʰtä˥m}] 0x1150: коробейник\\
\textbf{hta7n} [\emph{ʰtä˥n}] 0x7150: вырубают и выжигают\\
\textbf{hta7p} [\emph{ʰtä˥p}] 0x5150: фартинг\\
\textbf{hta7s} [\emph{ʰtä˥s}] 0x9150: формальной одежды\\
\textbf{hta7t} [\emph{ʰtä˥t}] 0xB150: созданный\\
\textbf{htac} [\emph{ʰtäʃ}] 0xE950: место назначения дело\\
\textbf{htaf} [\emph{ʰtäf}] 0xC950: прочность\\
\textbf{htak} [\emph{ʰtäk}] 0x2950: тема дело\\
\textbf{htam} [\emph{ʰtäm}] 0x0950: директива настроение\\
\textbf{htan} [\emph{ʰtän}] 0x6950: интерьер контекст\\
\textbf{htap} [\emph{ʰtäp}] 0x4950: связка\\
\textbf{htas} [\emph{ʰtäs}] 0x8950: еще раз\\
\textbf{htat} [\emph{ʰtät}] 0xA950: бабушка или дедушка\\
\textbf{hte2c} [\emph{ʰte˩ʃ}] 0xFB50: транш\\
\textbf{hte2k} [\emph{ʰte˩k}] 0x3B50: рвотный\\
\textbf{hte2m} [\emph{ʰte˩m}] 0x1B50: эктодерма\\
\textbf{hte2n} [\emph{ʰte˩n}] 0x7B50: в настоящее время\\
\textbf{hte2p} [\emph{ʰte˩p}] 0x5B50: уроженцев\\
\textbf{hte2s} [\emph{ʰte˩s}] 0x9B50: Темза\\
\textbf{hte2t} [\emph{ʰte˩t}] 0xBB50: достиг\\
\textbf{hte7c} [\emph{ʰte˥ʃ}] 0xF350: торговец скобяными изделиями\\
\textbf{hte7f} [\emph{ʰte˥f}] 0xD350: лейтмотив\\
\textbf{hte7k} [\emph{ʰte˥k}] 0x3350: Tinder\\
\textbf{hte7m} [\emph{ʰte˥m}] 0x1350: трамвайная линия\\
\textbf{hte7n} [\emph{ʰte˥n}] 0x7350: Афины\\
\textbf{hte7p} [\emph{ʰte˥p}] 0x5350: Фивы\\
\textbf{hte7s} [\emph{ʰte˥s}] 0x9350: патриотизм\\
\textbf{hte7t} [\emph{ʰte˥t}] 0xB350: и т.д\\
\textbf{htec} [\emph{ʰteʃ}] 0xEB50: энтузиазм\\
\textbf{htef} [\emph{ʰtef}] 0xCB50: портной\\
\textbf{htek} [\emph{ʰtek}] 0x2B50: библиотека\\
\textbf{htem} [\emph{ʰtem}] 0x0B50: постоянно\\
\textbf{hten} [\emph{ʰten}] 0x6B50: вчера\\
\textbf{htep} [\emph{ʰtep}] 0x4B50: поклонение\\
\textbf{htes} [\emph{ʰtes}] 0x8B50: посмотреть\\
\textbf{htet} [\emph{ʰtet}] 0xAB50: ангел\\
\textbf{hti2c} [\emph{ʰti˩ʃ}] 0xF850: покатый\\
\textbf{hti2f} [\emph{ʰti˩f}] 0xD850: сводить на нет\\
\textbf{hti2k} [\emph{ʰti˩k}] 0x3850: панегирик\\
\textbf{hti2m} [\emph{ʰti˩m}] 0x1850: ультиматум\\
\textbf{hti2n} [\emph{ʰti˩n}] 0x7850: междисциплинарный\\
\textbf{hti2p} [\emph{ʰti˩p}] 0x5850: хорошо одетым\\
\textbf{hti2s} [\emph{ʰti˩s}] 0x9850: презрительный\\
\textbf{hti2t} [\emph{ʰti˩t}] 0xB850: сохраняющий нейтралитет\\
\textbf{hti7c} [\emph{ʰti˥ʃ}] 0xF050: темы\\
\textbf{hti7f} [\emph{ʰti˥f}] 0xD050: Тимоти\\
\textbf{hti7m} [\emph{ʰti˥m}] 0x1050: чопорный\\
\textbf{hti7n} [\emph{ʰti˥n}] 0x7050: турбина\\
\textbf{hti7p} [\emph{ʰti˥p}] 0x5050: неисчерпаемый\\
\textbf{hti7s} [\emph{ʰti˥s}] 0x9050: незатребованный\\
\textbf{hti7t} [\emph{ʰti˥t}] 0xB050: отозвала\\
\textbf{htic} [\emph{ʰtiʃ}] 0xE850: их\\
\textbf{htif} [\emph{ʰtif}] 0xC850: расстояние\\
\textbf{htik} [\emph{ʰtik}] 0x2850: дательный дело\\
\textbf{htim} [\emph{ʰtim}] 0x0850: деонтические настроение\\
\textbf{htin} [\emph{ʰtin}] 0x6850: независимый пункт\\
\textbf{htip} [\emph{ʰtip}] 0x4850: 10\\
\textbf{htis} [\emph{ʰtis}] 0x8850: изготовление\\
\textbf{htit} [\emph{ʰtit}] 0xA850: Это\\
\textbf{hto2c} [\emph{ʰto˩ʃ}] 0xFC50: трудоемкими\\
\textbf{hto2k} [\emph{ʰto˩k}] 0x3C50: гуж\\
\textbf{hto2m} [\emph{ʰto˩m}] 0x1C50: погруженный\\
\textbf{hto2n} [\emph{ʰto˩n}] 0x7C50: послушание\\
\textbf{hto2p} [\emph{ʰto˩p}] 0x5C50: трубач\\
\textbf{hto2s} [\emph{ʰto˩s}] 0x9C50: катафалк\\
\textbf{hto2t} [\emph{ʰto˩t}] 0xBC50: кирпич\\
\textbf{hto7c} [\emph{ʰto˥ʃ}] 0xF450: три загнаны\\
\textbf{hto7f} [\emph{ʰto˥f}] 0xD450: легкая победа\\
\textbf{hto7k} [\emph{ʰto˥k}] 0x3450: потребительский фокус\\
\textbf{hto7m} [\emph{ʰto˥m}] 0x1450: Стокгольм\\
\textbf{hto7n} [\emph{ʰto˥n}] 0x7450: штурм мозга\\
\textbf{hto7p} [\emph{ʰto˥p}] 0x5450: трепел\\
\textbf{hto7s} [\emph{ʰto˥s}] 0x9450: Томас\\
\textbf{hto7t} [\emph{ʰto˥t}] 0xB450: порванный\\
\textbf{htoc} [\emph{ʰtoʃ}] 0xEC50: как-то\\
\textbf{htof} [\emph{ʰtof}] 0xCC50: вор\\
\textbf{htok} [\emph{ʰtok}] 0x2C50: конфликт\\
\textbf{htom} [\emph{ʰtom}] 0x0C50: экономика\\
\textbf{hton} [\emph{ʰton}] 0x6C50: достаточно\\
\textbf{htop} [\emph{ʰtop}] 0x4C50: проклятие\\
\textbf{htos} [\emph{ʰtos}] 0x8C50: большинство\\
\textbf{htot} [\emph{ʰtot}] 0xAC50: лоб\\
\textbf{htu2c} [\emph{ʰtu˩ʃ}] 0xFA50: ботик\\
\textbf{htu2k} [\emph{ʰtu˩k}] 0x3A50: чековая книжка\\
\textbf{htu2n} [\emph{ʰtu˩n}] 0x7A50: трибуна\\
\textbf{htu2p} [\emph{ʰtu˩p}] 0x5A50: утопия\\
\textbf{htu2s} [\emph{ʰtu˩s}] 0x9A50: кулачная драка\\
\textbf{htu2t} [\emph{ʰtu˩t}] 0xBA50: враждебность\\
\textbf{htu7f} [\emph{ʰtu˥f}] 0xD250: тьфу\\
\textbf{htu7k} [\emph{ʰtu˥k}] 0x3250: отечество\\
\textbf{htu7m} [\emph{ʰtu˥m}] 0x1250: перламутровый\\
\textbf{htu7n} [\emph{ʰtu˥n}] 0x7250: взимание\\
\textbf{htu7p} [\emph{ʰtu˥p}] 0x5250: бесспорно\\
\textbf{htu7s} [\emph{ʰtu˥s}] 0x9250: трепетный\\
\textbf{htu7t} [\emph{ʰtu˥t}] 0xB250: подтекст\\
\textbf{htuc} [\emph{ʰtuʃ}] 0xEA50: другой предмет\\
\textbf{htuf} [\emph{ʰtuf}] 0xCA50: тунец\\
\textbf{htuk} [\emph{ʰtuk}] 0x2A50: временной дело\\
\textbf{htum} [\emph{ʰtum}] 0x0A50: наставительный настроение\\
\textbf{htun} [\emph{ʰtun}] 0x6A50: только\\
\textbf{htup} [\emph{ʰtup}] 0x4A50: долг\\
\textbf{htus} [\emph{ʰtus}] 0x8A50: государство контекст\\
\textbf{htut} [\emph{ʰtut}] 0xAA50: негативный квантификатором\\
\textbf{hv6n} [\emph{ʰvən}] 0x6DB0: лавина\\
\textbf{hv6t} [\emph{ʰvət}] 0xADB0: вырубка леса\\
\textbf{hva2k} [\emph{ʰvä˩k}] 0x39B0: свастика\\
\textbf{hva2m} [\emph{ʰvä˩m}] 0x19B0: ванадий\\
\textbf{hva2n} [\emph{ʰvä˩n}] 0x79B0: раздаваться\\
\textbf{hva2s} [\emph{ʰvä˩s}] 0x99B0: парообразный\\
\textbf{hva2t} [\emph{ʰvä˩t}] 0xB9B0: вульгарность\\
\textbf{hva7k} [\emph{ʰvä˥k}] 0x31B0: белый воротничок\\
\textbf{hva7n} [\emph{ʰvä˥n}] 0x71B0: люблю делать\\
\textbf{hva7s} [\emph{ʰvä˥s}] 0x91B0: цианоз\\
\textbf{hva7t} [\emph{ʰvä˥t}] 0xB1B0: аббат\\
\textbf{hvac} [\emph{ʰväʃ}] 0xE9B0: посудомоечная машина\\
\textbf{hvaf} [\emph{ʰväf}] 0xC9B0: блондинка\\
\textbf{hvak} [\emph{ʰväk}] 0x29B0: авария\\
\textbf{hvam} [\emph{ʰväm}] 0x09B0: ванадий атом\\
\textbf{hvan} [\emph{ʰvän}] 0x69B0: пожаротушение\\
\textbf{hvap} [\emph{ʰväp}] 0x49B0: свиток бар\\
\textbf{hvas} [\emph{ʰväs}] 0x89B0: производители\\
\textbf{hvat} [\emph{ʰvät}] 0xA9B0: предстоящий\\
\textbf{hve2n} [\emph{ʰve˩n}] 0x7BB0: эфемерность\\
\textbf{hve2t} [\emph{ʰve˩t}] 0xBBB0: вендетта\\
\textbf{hve7n} [\emph{ʰve˥n}] 0x73B0: больверк\\
\textbf{hve7p} [\emph{ʰve˥p}] 0x53B0: вечерня\\
\textbf{hve7s} [\emph{ʰve˥s}] 0x93B0: ветряная оспа\\
\textbf{hve7t} [\emph{ʰve˥t}] 0xB3B0: зайчонок\\
\textbf{hvec} [\emph{ʰveʃ}] 0xEBB0: вершина\\
\textbf{hvef} [\emph{ʰvef}] 0xCBB0: червеобразный\\
\textbf{hvek} [\emph{ʰvek}] 0x2BB0: evitative дело\\
\textbf{hvem} [\emph{ʰvem}] 0x0BB0: вермишель\\
\textbf{hven} [\emph{ʰven}] 0x6BB0: контейнеры\\
\textbf{hvep} [\emph{ʰvep}] 0x4BB0: каверны\\
\textbf{hves} [\emph{ʰves}] 0x8BB0: венецианский\\
\textbf{hvet} [\emph{ʰvet}] 0xABB0: рядом с\\
\textbf{hvi2k} [\emph{ʰvi˩k}] 0x38B0: ярь-медянка\\
\textbf{hvi2n} [\emph{ʰvi˩n}] 0x78B0: завидный\\
\textbf{hvi2s} [\emph{ʰvi˩s}] 0x98B0: наиболее важные части\\
\textbf{hvi2t} [\emph{ʰvi˩t}] 0xB8B0: завещать\\
\textbf{hvi7f} [\emph{ʰvi˥f}] 0xD0B0: водевиль\\
\textbf{hvi7k} [\emph{ʰvi˥k}] 0x30B0: водные виды спорта\\
\textbf{hvi7n} [\emph{ʰvi˥n}] 0x70B0: исполненный по обету\\
\textbf{hvi7s} [\emph{ʰvi˥s}] 0x90B0: визирь\\
\textbf{hvi7t} [\emph{ʰvi˥t}] 0xB0B0: ризница\\
\textbf{hvic} [\emph{ʰviʃ}] 0xE8B0: виртуальный машина\\
\textbf{hvif} [\emph{ʰvif}] 0xC8B0: глагол - суффиксов\\
\textbf{hvik} [\emph{ʰvik}] 0x28B0: пробуждая\\
\textbf{hvim} [\emph{ʰvim}] 0x08B0: время выполнения\\
\textbf{hvin} [\emph{ʰvin}] 0x68B0: рационализировать\\
\textbf{hvip} [\emph{ʰvip}] 0x48B0: Юпитер\\
\textbf{hvis} [\emph{ʰvis}] 0x88B0: фекалии\\
\textbf{hvit} [\emph{ʰvit}] 0xA8B0: директива\\
\textbf{hvo2t} [\emph{ʰvo˩t}] 0xBCB0: преданно\\
\textbf{hvo7t} [\emph{ʰvo˥t}] 0xB4B0: преувеличивать\\
\textbf{hvok} [\emph{ʰvok}] 0x2CB0: водка\\
\textbf{hvom} [\emph{ʰvom}] 0x0CB0: спидометр\\
\textbf{hvon} [\emph{ʰvon}] 0x6CB0: воздушные силы\\
\textbf{hvop} [\emph{ʰvop}] 0x4CB0: проспать\\
\textbf{hvos} [\emph{ʰvos}] 0x8CB0: оазис\\
\textbf{hvot} [\emph{ʰvot}] 0xACB0: блок-схема\\
\textbf{hvu7t} [\emph{ʰvu˥t}] 0xB2B0: провокатор\\
\textbf{hvuc} [\emph{ʰvuʃ}] 0xEAB0: Вишну\\
\textbf{hvuk} [\emph{ʰvuk}] 0x2AB0: съеденный\\
\textbf{hvum} [\emph{ʰvum}] 0x0AB0: подводить\\
\textbf{hvun} [\emph{ʰvun}] 0x6AB0: бассейн\\
\textbf{hvup} [\emph{ʰvup}] 0x4AB0: вакуоль\\
\textbf{hvus} [\emph{ʰvus}] 0x8AB0: яйценосный\\
\textbf{hvut} [\emph{ʰvut}] 0xAAB0: переполнять\\
\textbf{hw67s} [\emph{ʰwə˥s}] 0x9528: ширь\\
\textbf{hw67t} [\emph{ʰwə˥t}] 0xB528: переваливаться\\
\textbf{hw6c} [\emph{ʰwəʃ}] 0xED28: водянистый\\
\textbf{hw6f} [\emph{ʰwəf}] 0xCD28: вафельный\\
\textbf{hw6k} [\emph{ʰwək}] 0x2D28: убывать\\
\textbf{hw6m} [\emph{ʰwəm}] 0x0D28: Вэстон\\
\textbf{hw6n} [\emph{ʰwən}] 0x6D28: вихрь\\
\textbf{hw6s} [\emph{ʰwəs}] 0x8D28: эвфемизм\\
\textbf{hw6t} [\emph{ʰwət}] 0xAD28: Швеция\\
\textbf{hwa2k} [\emph{ʰwä˩k}] 0x3928: умник\\
\textbf{hwa2n} [\emph{ʰwä˩n}] 0x7928: Ублюдок\\
\textbf{hwa2p} [\emph{ʰwä˩p}] 0x5928: Ваххабиты\\
\textbf{hwa2s} [\emph{ʰwä˩s}] 0x9928: Варшавская\\
\textbf{hwa2t} [\emph{ʰwä˩t}] 0xB928: положение\\
\textbf{hwa7c} [\emph{ʰwä˥ʃ}] 0xF128: пятьдесят девять\\
\textbf{hwa7f} [\emph{ʰwä˥f}] 0xD128: письмо совершенным\\
\textbf{hwa7k} [\emph{ʰwä˥k}] 0x3128: акварель\\
\textbf{hwa7m} [\emph{ʰwä˥m}] 0x1128: сердцеед\\
\textbf{hwa7n} [\emph{ʰwä˥n}] 0x7128: всемогущий\\
\textbf{hwa7p} [\emph{ʰwä˥p}] 0x5128: пятьдесят восемь\\
\textbf{hwa7s} [\emph{ʰwä˥s}] 0x9128: вальс\\
\textbf{hwa7t} [\emph{ʰwä˥t}] 0xB128: приветствовал\\
\textbf{hwac} [\emph{ʰwäʃ}] 0xE928: Мир\\
\textbf{hwaf} [\emph{ʰwäf}] 0xC928: идеально\\
\textbf{hwak} [\emph{ʰwäk}] 0x2928: гласная буква\\
\textbf{hwam} [\emph{ʰwäm}] 0x0928: нас\\
\textbf{hwan} [\emph{ʰwän}] 0x6928: а также\\
\textbf{hwap} [\emph{ʰwäp}] 0x4928: 8\\
\textbf{hwas} [\emph{ʰwäs}] 0x8928: вопрос\\
\textbf{hwat} [\emph{ʰwät}] 0xA928: вопрос слово\\
\textbf{hwe2k} [\emph{ʰwe˩k}] 0x3B28: бурдюк\\
\textbf{hwe2n} [\emph{ʰwe˩n}] 0x7B28: искусный\\
\textbf{hwe2s} [\emph{ʰwe˩s}] 0x9B28: снята с охраны\\
\textbf{hwe2t} [\emph{ʰwe˩t}] 0xBB28: эклиптика\\
\textbf{hwe7k} [\emph{ʰwe˥k}] 0x3328: поэтому\\
\textbf{hwe7n} [\emph{ʰwe˥n}] 0x7328: пятьдесят четыре\\
\textbf{hwe7s} [\emph{ʰwe˥s}] 0x9328: кружевной\\
\textbf{hwe7t} [\emph{ʰwe˥t}] 0xB328: друг детства\\
\textbf{hwec} [\emph{ʰweʃ}] 0xEB28: осада\\
\textbf{hwef} [\emph{ʰwef}] 0xCB28: неоплаченный\\
\textbf{hwek} [\emph{ʰwek}] 0x2B28: волк\\
\textbf{hwem} [\emph{ʰwem}] 0x0B28: спектр\\
\textbf{hwen} [\emph{ʰwen}] 0x6B28: растительный Пол\\
\textbf{hwep} [\emph{ʰwep}] 0x4B28: Web\\
\textbf{hwes} [\emph{ʰwes}] 0x8B28: без сознания\\
\textbf{hwet} [\emph{ʰwet}] 0xAB28: желудок\\
\textbf{hwi2c} [\emph{ʰwi˩ʃ}] 0xF828: приговоренный к пожизненному заключению\\
\textbf{hwi2k} [\emph{ʰwi˩k}] 0x3828: волна частиц\\
\textbf{hwi2m} [\emph{ʰwi˩m}] 0x1828: опасный для жизни\\
\textbf{hwi2n} [\emph{ʰwi˩n}] 0x7828: Уилли\\
\textbf{hwi2s} [\emph{ʰwi˩s}] 0x9828: руды\\
\textbf{hwi2t} [\emph{ʰwi˩t}] 0xB828: болотные сапоги\\
\textbf{hwi7c} [\emph{ʰwi˥ʃ}] 0xF028: обморок\\
\textbf{hwi7f} [\emph{ʰwi˥f}] 0xD028: волдырь\\
\textbf{hwi7k} [\emph{ʰwi˥k}] 0x3028: элегия\\
\textbf{hwi7m} [\emph{ʰwi˥m}] 0x1028: Уильям\\
\textbf{hwi7n} [\emph{ʰwi˥n}] 0x7028: пятьдесят пять\\
\textbf{hwi7p} [\emph{ʰwi˥p}] 0x5028: пшеничный\\
\textbf{hwi7s} [\emph{ʰwi˥s}] 0x9028: еси\\
\textbf{hwi7t} [\emph{ʰwi˥t}] 0xB028: жилет\\
\textbf{hwic} [\emph{ʰwiʃ}] 0xE828: воды\\
\textbf{hwif} [\emph{ʰwif}] 0xC828: вице\\
\textbf{hwik} [\emph{ʰwik}] 0x2828: звательный дело\\
\textbf{hwim} [\emph{ʰwim}] 0x0828: цивилизация\\
\textbf{hwin} [\emph{ʰwin}] 0x6828: место нахождения дело\\
\textbf{hwip} [\emph{ʰwip}] 0x4828: Воскресенье\\
\textbf{hwis} [\emph{ʰwis}] 0x8828: Vialis дело\\
\textbf{hwit} [\emph{ʰwit}] 0xA828: вернуть\\
\textbf{hwo2n} [\emph{ʰwo˩n}] 0x7C28: червоточина\\
\textbf{hwo2t} [\emph{ʰwo˩t}] 0xBC28: рогоносец\\
\textbf{hwo7f} [\emph{ʰwo˥f}] 0xD428: Волоф\\
\textbf{hwo7k} [\emph{ʰwo˥k}] 0x3428: сурок\\
\textbf{hwo7n} [\emph{ʰwo˥n}] 0x7428: всеобъемлющая\\
\textbf{hwo7p} [\emph{ʰwo˥p}] 0x5428: олигополия\\
\textbf{hwo7t} [\emph{ʰwo˥t}] 0xB428: проводной\\
\textbf{hwoc} [\emph{ʰwoʃ}] 0xEC28: ведьма\\
\textbf{hwof} [\emph{ʰwof}] 0xCC28: ива\\
\textbf{hwok} [\emph{ʰwok}] 0x2C28: квалифицированный\\
\textbf{hwom} [\emph{ʰwom}] 0x0C28: бездомный\\
\textbf{hwon} [\emph{ʰwon}] 0x6C28: приглашение\\
\textbf{hwop} [\emph{ʰwop}] 0x4C28: ватт\\
\textbf{hwos} [\emph{ʰwos}] 0x8C28: Спальня\\
\textbf{hwot} [\emph{ʰwot}] 0xAC28: витой\\
\textbf{hwu2n} [\emph{ʰwu˩n}] 0x7A28: шерстяной\\
\textbf{hwu2t} [\emph{ʰwu˩t}] 0xBA28: колдовство\\
\textbf{hwu7n} [\emph{ʰwu˥n}] 0x7228: тепловатый\\
\textbf{hwu7p} [\emph{ʰwu˥p}] 0x5228: пятикратный\\
\textbf{hwu7t} [\emph{ʰwu˥t}] 0xB228: древесный\\
\textbf{hwuc} [\emph{ʰwuʃ}] 0xEA28: связи\\
\textbf{hwuf} [\emph{ʰwuf}] 0xCA28: вдова\\
\textbf{hwuk} [\emph{ʰwuk}] 0x2A28: свободный\\
\textbf{hwum} [\emph{ʰwum}] 0x0A28: слух\\
\textbf{hwun} [\emph{ʰwun}] 0x6A28: золото\\
\textbf{hwup} [\emph{ʰwup}] 0x4A28: неоценимый\\
\textbf{hwus} [\emph{ʰwus}] 0x8A28: текст\\
\textbf{hwut} [\emph{ʰwut}] 0xAA28: забывать\\
\textbf{hx62n} [\emph{ʰxə˩n}] 0x7DD0: Осанна\\
\textbf{hx67n} [\emph{ʰxə˥n}] 0x75D0: ломоть\\
\textbf{hx67t} [\emph{ʰxə˥t}] 0xB5D0: хеттский\\
\textbf{hx6c} [\emph{ʰxəʃ}] 0xEDD0: молотилка\\
\textbf{hx6f} [\emph{ʰxəf}] 0xCDD0: агиография\\
\textbf{hx6k} [\emph{ʰxək}] 0x2DD0: ха\\
\textbf{hx6m} [\emph{ʰxəm}] 0x0DD0: почтовый голубь\\
\textbf{hx6n} [\emph{ʰxən}] 0x6DD0: предварительно сегодняшний\\
\textbf{hx6p} [\emph{ʰxəp}] 0x4DD0: травяной\\
\textbf{hx6s} [\emph{ʰxəs}] 0x8DD0: сталкиваться\\
\textbf{hx6t} [\emph{ʰxət}] 0xADD0: Канада\\
\textbf{hxa2c} [\emph{ʰxä˩ʃ}] 0xF9D0: хорошее общение\\
\textbf{hxa2f} [\emph{ʰxä˩f}] 0xD9D0: сверстник\\
\textbf{hxa2k} [\emph{ʰxä˩k}] 0x39D0: улей\\
\textbf{hxa2m} [\emph{ʰxä˩m}] 0x19D0: анемия\\
\textbf{hxa2n} [\emph{ʰxä˩n}] 0x79D0: охотно\\
\textbf{hxa2p} [\emph{ʰxä˩p}] 0x59D0: гусар\\
\textbf{hxa2s} [\emph{ʰxä˩s}] 0x99D0: Далида\\
\textbf{hxa2t} [\emph{ʰxä˩t}] 0xB9D0: крики\\
\textbf{hxa7c} [\emph{ʰxä˥ʃ}] 0xF1D0: надуманным\\
\textbf{hxa7f} [\emph{ʰxä˥f}] 0xD1D0: общее право\\
\textbf{hxa7k} [\emph{ʰxä˥k}] 0x31D0: horsehoe краба\\
\textbf{hxa7m} [\emph{ʰxä˥m}] 0x11D0: рукопожатие\\
\textbf{hxa7n} [\emph{ʰxä˥n}] 0x71D0: монахи\\
\textbf{hxa7p} [\emph{ʰxä˥p}] 0x51D0: Йоханнесбург\\
\textbf{hxa7s} [\emph{ʰxä˥s}] 0x91D0: Увы\\
\textbf{hxa7t} [\emph{ʰxä˥t}] 0xB1D0: ахнула\\
\textbf{hxac} [\emph{ʰxäʃ}] 0xE9D0: перезапуск\\
\textbf{hxaf} [\emph{ʰxäf}] 0xC9D0: выкапывать\\
\textbf{hxak} [\emph{ʰxäk}] 0x29D0: перерабатывать\\
\textbf{hxam} [\emph{ʰxäm}] 0x09D0: миллиграмм\\
\textbf{hxan} [\emph{ʰxän}] 0x69D0: семинар\\
\textbf{hxap} [\emph{ʰxäp}] 0x49D0: рюкзак\\
\textbf{hxas} [\emph{ʰxäs}] 0x89D0: гепатит\\
\textbf{hxat} [\emph{ʰxät}] 0xA9D0: застежка\\
\textbf{hxe2k} [\emph{ʰxe˩k}] 0x3BD0: заголовок\\
\textbf{hxe2m} [\emph{ʰxe˩m}] 0x1BD0: гем\\
\textbf{hxe2n} [\emph{ʰxe˩n}] 0x7BD0: стройный\\
\textbf{hxe2s} [\emph{ʰxe˩s}] 0x9BD0: семьдесят шесть\\
\textbf{hxe2t} [\emph{ʰxe˩t}] 0xBBD0: подголовник\\
\textbf{hxe7c} [\emph{ʰxe˥ʃ}] 0xF3D0: пролив\\
\textbf{hxe7k} [\emph{ʰxe˥k}] 0x33D0: хакеры\\
\textbf{hxe7m} [\emph{ʰxe˥m}] 0x13D0: девственная плева\\
\textbf{hxe7n} [\emph{ʰxe˥n}] 0x73D0: отклоненный\\
\textbf{hxe7p} [\emph{ʰxe˥p}] 0x53D0: знаменоносец\\
\textbf{hxe7s} [\emph{ʰxe˥s}] 0x93D0: Гессе\\
\textbf{hxe7t} [\emph{ʰxe˥t}] 0xB3D0: беспокоит\\
\textbf{hxec} [\emph{ʰxeʃ}] 0xEBD0: Тедди\\
\textbf{hxef} [\emph{ʰxef}] 0xCBD0: вереск\\
\textbf{hxek} [\emph{ʰxek}] 0x2BD0: изменение формы\\
\textbf{hxem} [\emph{ʰxem}] 0x0BD0: кропотливый\\
\textbf{hxen} [\emph{ʰxen}] 0x6BD0: Перенаправление\\
\textbf{hxep} [\emph{ʰxep}] 0x4BD0: скейтборд\\
\textbf{hxes} [\emph{ʰxes}] 0x8BD0: герц\\
\textbf{hxet} [\emph{ʰxet}] 0xABD0: увеличительный\\
\textbf{hxi2k} [\emph{ʰxi˩k}] 0x38D0: киркомотыга\\
\textbf{hxi2m} [\emph{ʰxi˩m}] 0x18D0: Крик осла\\
\textbf{hxi2n} [\emph{ʰxi˩n}] 0x78D0: охваченном\\
\textbf{hxi2p} [\emph{ʰxi˩p}] 0x58D0: делать орфографические ошибки\\
\textbf{hxi2s} [\emph{ʰxi˩s}] 0x98D0: суффикс\\
\textbf{hxi2t} [\emph{ʰxi˩t}] 0xB8D0: встряхивают\\
\textbf{hxi7c} [\emph{ʰxi˥ʃ}] 0xF0D0: семьдесят девять\\
\textbf{hxi7f} [\emph{ʰxi˥f}] 0xD0D0: верный огонь\\
\textbf{hxi7k} [\emph{ʰxi˥k}] 0x30D0: нищий\\
\textbf{hxi7m} [\emph{ʰxi˥m}] 0x10D0: МВФ\\
\textbf{hxi7n} [\emph{ʰxi˥n}] 0x70D0: хозяйка\\
\textbf{hxi7p} [\emph{ʰxi˥p}] 0x50D0: постоялый двор\\
\textbf{hxi7s} [\emph{ʰxi˥s}] 0x90D0: место отдыха\\
\textbf{hxi7t} [\emph{ʰxi˥t}] 0xB0D0: гематит\\
\textbf{hxic} [\emph{ʰxiʃ}] 0xE8D0: сдвиг ключ\\
\textbf{hxif} [\emph{ʰxif}] 0xC8D0: гипнотизировать\\
\textbf{hxik} [\emph{ʰxik}] 0x28D0: конфигурировать\\
\textbf{hxim} [\emph{ʰxim}] 0x08D0: торий атом\\
\textbf{hxin} [\emph{ʰxin}] 0x68D0: синхронизировать\\
\textbf{hxip} [\emph{ʰxip}] 0x48D0: антитело\\
\textbf{hxis} [\emph{ʰxis}] 0x88D0: пиксель\\
\textbf{hxit} [\emph{ʰxit}] 0xA8D0: максимизировать\\
\textbf{hxo2k} [\emph{ʰxo˩k}] 0x3CD0: губная гармоника\\
\textbf{hxo2n} [\emph{ʰxo˩n}] 0x7CD0: чемодан\\
\textbf{hxo2p} [\emph{ʰxo˩p}] 0x5CD0: Ходж мешанина\\
\textbf{hxo2s} [\emph{ʰxo˩s}] 0x9CD0: гаубица\\
\textbf{hxo2t} [\emph{ʰxo˩t}] 0xBCD0: Геродот\\
\textbf{hxo7k} [\emph{ʰxo˥k}] 0x34D0: тоскующий по дому\\
\textbf{hxo7m} [\emph{ʰxo˥m}] 0x14D0: Пожиратель огня\\
\textbf{hxo7n} [\emph{ʰxo˥n}] 0x74D0: нетрадиционный\\
\textbf{hxo7p} [\emph{ʰxo˥p}] 0x54D0: горбатый\\
\textbf{hxo7s} [\emph{ʰxo˥s}] 0x94D0: мимоза\\
\textbf{hxo7t} [\emph{ʰxo˥t}] 0xB4D0: над породой\\
\textbf{hxoc} [\emph{ʰxoʃ}] 0xECD0: ров\\
\textbf{hxof} [\emph{ʰxof}] 0xCCD0: герцеговина\\
\textbf{hxok} [\emph{ʰxok}] 0x2CD0: неиспользованный\\
\textbf{hxom} [\emph{ʰxom}] 0x0CD0: утиль почта\\
\textbf{hxon} [\emph{ʰxon}] 0x6CD0: середина весны\\
\textbf{hxop} [\emph{ʰxop}] 0x4CD0: Эфиопия\\
\textbf{hxos} [\emph{ʰxos}] 0x8CD0: торус\\
\textbf{hxot} [\emph{ʰxot}] 0xACD0: наручные часы\\
\textbf{hxu2n} [\emph{ʰxu˩n}] 0x7AD0: Хусайн\\
\textbf{hxu2s} [\emph{ʰxu˩s}] 0x9AD0: мотыги\\
\textbf{hxu2t} [\emph{ʰxu˩t}] 0xBAD0: кобура\\
\textbf{hxu7c} [\emph{ʰxu˥ʃ}] 0xF2D0: с корпусом, скрытым за горизонтом\\
\textbf{hxu7k} [\emph{ʰxu˥k}] 0x32D0: Аввакум\\
\textbf{hxu7m} [\emph{ʰxu˥m}] 0x12D0: гумусовый\\
\textbf{hxu7n} [\emph{ʰxu˥n}] 0x72D0: Беззаботная\\
\textbf{hxu7s} [\emph{ʰxu˥s}] 0x92D0: эвристический\\
\textbf{hxu7t} [\emph{ʰxu˥t}] 0xB2D0: хорошо делать\\
\textbf{hxuc} [\emph{ʰxuʃ}] 0xEAD0: спасающий\\
\textbf{hxuf} [\emph{ʰxuf}] 0xCAD0: Ганновер\\
\textbf{hxuk} [\emph{ʰxuk}] 0x2AD0: изнашивания\\
\textbf{hxum} [\emph{ʰxum}] 0x0AD0: бедренная кость\\
\textbf{hxun} [\emph{ʰxun}] 0x6AD0: уговаривание\\
\textbf{hxup} [\emph{ʰxup}] 0x4AD0: головастик\\
\textbf{hxus} [\emph{ʰxus}] 0x8AD0: открытый источник\\
\textbf{hxut} [\emph{ʰxut}] 0xAAD0: маршрутизатор\\
\textbf{hy67n} [\emph{ʰjə˥n}] 0x7518: воинственный\\
\textbf{hy67t} [\emph{ʰjə˥t}] 0xB518: йодль\\
\textbf{hy6k} [\emph{ʰjək}] 0x2D18: слюда\\
\textbf{hy6m} [\emph{ʰjəm}] 0x0D18: Йемен\\
\textbf{hy6n} [\emph{ʰjən}] 0x6D18: оживить Пол\\
\textbf{hy6s} [\emph{ʰjəs}] 0x8D18: Румыния\\
\textbf{hy6t} [\emph{ʰjət}] 0xAD18: грамматическая - объект\\
\textbf{hya2c} [\emph{ʰjä˩ʃ}] 0xF918: уравнитель\\
\textbf{hya2k} [\emph{ʰjä˩k}] 0x3918: янки\\
\textbf{hya2m} [\emph{ʰjä˩m}] 0x1918: аймара\\
\textbf{hya2n} [\emph{ʰjä˩n}] 0x7918: сирота\\
\textbf{hya2p} [\emph{ʰjä˩p}] 0x5918: ямбический\\
\textbf{hya2s} [\emph{ʰjä˩s}] 0x9918: соблазнять\\
\textbf{hya2t} [\emph{ʰjä˩t}] 0xB918: пророк\\
\textbf{hya7c} [\emph{ʰjä˥ʃ}] 0xF118: стискивать\\
\textbf{hya7f} [\emph{ʰjä˥f}] 0xD118: один путь\\
\textbf{hya7k} [\emph{ʰjä˥k}] 0x3118: Джекоб\\
\textbf{hya7m} [\emph{ʰjä˥m}] 0x1118: фрамбезия\\
\textbf{hya7n} [\emph{ʰjä˥n}] 0x7118: ослепительный\\
\textbf{hya7p} [\emph{ʰjä˥p}] 0x5118: хорошо обоснован\\
\textbf{hya7s} [\emph{ʰjä˥s}] 0x9118: Исайя\\
\textbf{hya7t} [\emph{ʰjä˥t}] 0xB118: братва\\
\textbf{hyac} [\emph{ʰjäʃ}] 0xE918: также\\
\textbf{hyaf} [\emph{ʰjäf}] 0xC918: возможно\\
\textbf{hyak} [\emph{ʰjäk}] 0x2918: инструментальный дело\\
\textbf{hyam} [\emph{ʰjäm}] 0x0918: ночь\\
\textbf{hyan} [\emph{ʰjän}] 0x6918: зависимый пункт\\
\textbf{hyap} [\emph{ʰjäp}] 0x4918: король\\
\textbf{hyas} [\emph{ʰjäs}] 0x8918: один раз\\
\textbf{hyat} [\emph{ʰjät}] 0xA918: родительский\\
\textbf{hye2k} [\emph{ʰje˩k}] 0x3B18: один к одному\\
\textbf{hye2n} [\emph{ʰje˩n}] 0x7B18: глотка\\
\textbf{hye2t} [\emph{ʰje˩t}] 0xBB18: шагомер\\
\textbf{hye7k} [\emph{ʰje˥k}] 0x3318: немного\\
\textbf{hye7m} [\emph{ʰje˥m}] 0x1318: Эм\\
\textbf{hye7n} [\emph{ʰje˥n}] 0x7318: ан\\
\textbf{hye7p} [\emph{ʰje˥p}] 0x5318: тембр\\
\textbf{hye7s} [\emph{ʰje˥s}] 0x9318: бесчувственный\\
\textbf{hye7t} [\emph{ʰje˥t}] 0xB318: маркированный\\
\textbf{hyec} [\emph{ʰjeʃ}] 0xEB18: овца\\
\textbf{hyef} [\emph{ʰjef}] 0xCB18: рецессивный\\
\textbf{hyek} [\emph{ʰjek}] 0x2B18: путешественник\\
\textbf{hyem} [\emph{ʰjem}] 0x0B18: вяз\\
\textbf{hyen} [\emph{ʰjen}] 0x6B18: облако\\
\textbf{hyep} [\emph{ʰjep}] 0x4B18: веко\\
\textbf{hyes} [\emph{ʰjes}] 0x8B18: осведомленность\\
\textbf{hyet} [\emph{ʰjet}] 0xAB18: масло\\
\textbf{hyi2k} [\emph{ʰji˩k}] 0x3818: аллегория\\
\textbf{hyi2n} [\emph{ʰji˩n}] 0x7818: забвение\\
\textbf{hyi2p} [\emph{ʰji˩p}] 0x5818: ночной клуб\\
\textbf{hyi2s} [\emph{ʰji˩s}] 0x9818: допустимый\\
\textbf{hyi7c} [\emph{ʰji˥ʃ}] 0xF018: доброжелатель\\
\textbf{hyi7k} [\emph{ʰji˥k}] 0x3018: действительный\\
\textbf{hyi7m} [\emph{ʰji˥m}] 0x1018: сионизм\\
\textbf{hyi7n} [\emph{ʰji˥n}] 0x7018: торжественность\\
\textbf{hyi7p} [\emph{ʰji˥p}] 0x5018: посмешище\\
\textbf{hyi7s} [\emph{ʰji˥s}] 0x9018: тюремщик\\
\textbf{hyi7t} [\emph{ʰji˥t}] 0xB018: иезуитский\\
\textbf{hyic} [\emph{ʰjiʃ}] 0xE818: увидел\\
\textbf{hyif} [\emph{ʰjif}] 0xC818: вперед\\
\textbf{hyik} [\emph{ʰjik}] 0x2818: один\\
\textbf{hyim} [\emph{ʰjim}] 0x0818: армия\\
\textbf{hyin} [\emph{ʰjin}] 0x6818: универсальный квантификатором\\
\textbf{hyip} [\emph{ʰjip}] 0x4818: ретроспективный аспект\\
\textbf{hyis} [\emph{ʰjis}] 0x8818: бедные\\
\textbf{hyit} [\emph{ʰjit}] 0xA818: REALiS настроение\\
\textbf{hyo2k} [\emph{ʰjo˩k}] 0x3C18: дуться\\
\textbf{hyo2n} [\emph{ʰjo˩n}] 0x7C18: Иона\\
\textbf{hyo2s} [\emph{ʰjo˩s}] 0x9C18: цветных\\
\textbf{hyo2t} [\emph{ʰjo˩t}] 0xBC18: бородавка\\
\textbf{hyo7k} [\emph{ʰjo˥k}] 0x3418: боль в спине\\
\textbf{hyo7n} [\emph{ʰjo˥n}] 0x7418: предвидели\\
\textbf{hyo7s} [\emph{ʰjo˥s}] 0x9418: самоуправления обладали\\
\textbf{hyo7t} [\emph{ʰjo˥t}] 0xB418: желчь\\
\textbf{hyoc} [\emph{ʰjoʃ}] 0xEC18: искушение\\
\textbf{hyof} [\emph{ʰjof}] 0xCC18: выгодный\\
\textbf{hyok} [\emph{ʰjok}] 0x2C18: угол\\
\textbf{hyom} [\emph{ʰjom}] 0x0C18: йод\\
\textbf{hyon} [\emph{ʰjon}] 0x6C18: тем временем\\
\textbf{hyop} [\emph{ʰjop}] 0x4C18: яхта\\
\textbf{hyos} [\emph{ʰjos}] 0x8C18: опасающийся\\
\textbf{hyot} [\emph{ʰjot}] 0xAC18: опасающийся настроение\\
\textbf{hyu2k} [\emph{ʰju˩k}] 0x3A18: узуфрукт\\
\textbf{hyu2n} [\emph{ʰju˩n}] 0x7A18: евгеника\\
\textbf{hyu2s} [\emph{ʰju˩s}] 0x9A18: мочевина\\
\textbf{hyu2t} [\emph{ʰju˩t}] 0xBA18: Иуда\\
\textbf{hyu7k} [\emph{ʰju˥k}] 0x3218: эврика\\
\textbf{hyu7n} [\emph{ʰju˥n}] 0x7218: еврейка\\
\textbf{hyu7p} [\emph{ʰju˥p}] 0x5218: антилопа гну\\
\textbf{hyu7s} [\emph{ʰju˥s}] 0x9218: язычок\\
\textbf{hyu7t} [\emph{ʰju˥t}] 0xB218: приведены\\
\textbf{hyuc} [\emph{ʰjuʃ}] 0xEA18: круглый\\
\textbf{hyuf} [\emph{ʰjuf}] 0xCA18: мочеиспускательный канал\\
\textbf{hyuk} [\emph{ʰjuk}] 0x2A18: красный\\
\textbf{hyum} [\emph{ʰjum}] 0x0A18: с помощью\\
\textbf{hyun} [\emph{ʰjun}] 0x6A18: река\\
\textbf{hyup} [\emph{ʰjup}] 0x4A18: Европа\\
\textbf{hyus} [\emph{ʰjus}] 0x8A18: красивая\\
\textbf{hyut} [\emph{ʰjut}] 0xAA18: вопросительный настроение\\
\textbf{hz6k} [\emph{ʰzək}] 0x2DA0: сейсмограф\\
\textbf{hz6n} [\emph{ʰzən}] 0x6DA0: Тасмания\\
\textbf{hz6s} [\emph{ʰzəs}] 0x8DA0: позитивизм\\
\textbf{hz6t} [\emph{ʰzət}] 0xADA0: зигота\\
\textbf{hza2k} [\emph{ʰzä˩k}] 0x39A0: ацтекский\\
\textbf{hza2n} [\emph{ʰzä˩n}] 0x79A0: праздношатание\\
\textbf{hza2t} [\emph{ʰzä˩t}] 0xB9A0: прогорклый\\
\textbf{hza7k} [\emph{ʰzä˥k}] 0x31A0: прелюбодей\\
\textbf{hza7n} [\emph{ʰzä˥n}] 0x71A0: византийский\\
\textbf{hza7s} [\emph{ʰzä˥s}] 0x91A0: Заир\\
\textbf{hza7t} [\emph{ʰzä˥t}] 0xB1A0: забита\\
\textbf{hzac} [\emph{ʰzäʃ}] 0xE9A0: взаимное\\
\textbf{hzaf} [\emph{ʰzäf}] 0xC9A0: препятствовать\\
\textbf{hzak} [\emph{ʰzäk}] 0x29A0: трясти\\
\textbf{hzam} [\emph{ʰzäm}] 0x09A0: банан\\
\textbf{hzan} [\emph{ʰzän}] 0x69A0: содержащий\\
\textbf{hzap} [\emph{ʰzäp}] 0x49A0: аббатство\\
\textbf{hzas} [\emph{ʰzäs}] 0x89A0: шутить\\
\textbf{hzat} [\emph{ʰzät}] 0xA9A0: волна\\
\textbf{hze2t} [\emph{ʰze˩t}] 0xBBA0: метаданные\\
\textbf{hze7c} [\emph{ʰze˥ʃ}] 0xF3A0: Z\\
\textbf{hze7k} [\emph{ʰze˥k}] 0x33A0: Иезекииль\\
\textbf{hze7n} [\emph{ʰze˥n}] 0x73A0: гексан\\
\textbf{hze7t} [\emph{ʰze˥t}] 0xB3A0: низводить\\
\textbf{hzec} [\emph{ʰzeʃ}] 0xEBA0: одновременность\\
\textbf{hzek} [\emph{ʰzek}] 0x2BA0: марганец\\
\textbf{hzem} [\emph{ʰzem}] 0x0BA0: экзотермический\\
\textbf{hzen} [\emph{ʰzen}] 0x6BA0: плавать\\
\textbf{hzep} [\emph{ʰzep}] 0x4BA0: затылок\\
\textbf{hzes} [\emph{ʰzes}] 0x8BA0: одержимый дело\\
\textbf{hzet} [\emph{ʰzet}] 0xABA0: морской узел\\
\textbf{hzi2n} [\emph{ʰzi˩n}] 0x78A0: суннит\\
\textbf{hzi2s} [\emph{ʰzi˩s}] 0x98A0: гностицизм\\
\textbf{hzi2t} [\emph{ʰzi˩t}] 0xB8A0: зажимать\\
\textbf{hzi7k} [\emph{ʰzi˥k}] 0x30A0: все американские\\
\textbf{hzi7n} [\emph{ʰzi˥n}] 0x70A0: осматривает\\
\textbf{hzi7p} [\emph{ʰzi˥p}] 0x50A0: живой мертвец\\
\textbf{hzi7s} [\emph{ʰzi˥s}] 0x90A0: близорукость\\
\textbf{hzi7t} [\emph{ʰzi˥t}] 0xB0A0: островок\\
\textbf{hzic} [\emph{ʰziʃ}] 0xE8A0: яд\\
\textbf{hzif} [\emph{ʰzif}] 0xC8A0: зефир\\
\textbf{hzik} [\emph{ʰzik}] 0x28A0: потухший\\
\textbf{hzim} [\emph{ʰzim}] 0x08A0: висмут\\
\textbf{hzin} [\emph{ʰzin}] 0x68A0: южный\\
\textbf{hzip} [\emph{ʰzip}] 0x48A0: застежка-молния\\
\textbf{hzis} [\emph{ʰzis}] 0x88A0: слеза\\
\textbf{hzit} [\emph{ʰzit}] 0xA8A0: функции\\
\textbf{hzo2s} [\emph{ʰzo˩s}] 0x9CA0: вы UNS\\
\textbf{hzo7n} [\emph{ʰzo˥n}] 0x74A0: меццо-сопрано\\
\textbf{hzo7t} [\emph{ʰzo˥t}] 0xB4A0: зороастрийский\\
\textbf{hzoc} [\emph{ʰzoʃ}] 0xECA0: лежать в основе\\
\textbf{hzok} [\emph{ʰzok}] 0x2CA0: извиваться\\
\textbf{hzom} [\emph{ʰzom}] 0x0CA0: кочевник\\
\textbf{hzon} [\emph{ʰzon}] 0x6CA0: обезьяна\\
\textbf{hzop} [\emph{ʰzop}] 0x4CA0: Мозамбик\\
\textbf{hzos} [\emph{ʰzos}] 0x8CA0: Казахстан\\
\textbf{hzot} [\emph{ʰzot}] 0xACA0: прекращение\\
\textbf{hzu2t} [\emph{ʰzu˩t}] 0xBAA0: толстовка с капюшоном\\
\textbf{hzu7t} [\emph{ʰzu˥t}] 0xB2A0: ростовщик\\
\textbf{hzuc} [\emph{ʰzuʃ}] 0xEAA0: эксгумация\\
\textbf{hzuk} [\emph{ʰzuk}] 0x2AA0: зоопарк\\
\textbf{hzum} [\emph{ʰzum}] 0x0AA0: цирконий\\
\textbf{hzun} [\emph{ʰzun}] 0x6AA0: упертый\\
\textbf{hzup} [\emph{ʰzup}] 0x4AA0: Зимбабве\\
\textbf{hzus} [\emph{ʰzus}] 0x8AA0: Судан\\
\textbf{hzut} [\emph{ʰzut}] 0xAAA0: туман\\
\subsection{ j }
\textbf{ja} [\emph{ʒä}] 0x26BE: связка\\
\textbf{je} [\emph{ʒe}] 0x2EBE: agentive case, the agent of a transitive verb.\\
\textbf{ji} [\emph{ʒi}] 0x22BE: женский Пол\\
\textbf{ji7} [\emph{ʒi˥}] 0x42BE: Изе\\
\textbf{jl6n} [\emph{ʒlən}] 0x6D75: смягчающий\\
\textbf{jl6t} [\emph{ʒlət}] 0xAD75: волочиться\\
\textbf{jla2t} [\emph{ʒlä˩t}] 0xB975: желатин\\
\textbf{jla7n} [\emph{ʒlä˥n}] 0x7175: лопатой носатый\\
\textbf{jla7t} [\emph{ʒlä˥t}] 0xB175: красящее вещество\\
\textbf{jlac} [\emph{ʒläʃ}] 0xE975: массаж\\
\textbf{jlaf} [\emph{ʒläf}] 0xC975: золь Ф.А.\\
\textbf{jlak} [\emph{ʒläk}] 0x2975: якорь\\
\textbf{jlam} [\emph{ʒläm}] 0x0975: генералиссимус\\
\textbf{jlan} [\emph{ʒlän}] 0x6975: поле\\
\textbf{jlap} [\emph{ʒläp}] 0x4975: наждачная бумага\\
\textbf{jlas} [\emph{ʒläs}] 0x8975: ненадежный\\
\textbf{jlat} [\emph{ʒlät}] 0xA975: пират\\
\textbf{jlek} [\emph{ʒlek}] 0x2B75: геолог\\
\textbf{jlen} [\emph{ʒlen}] 0x6B75: местонахождение Пол\\
\textbf{jlep} [\emph{ʒlep}] 0x4B75: коленопреклонение\\
\textbf{jles} [\emph{ʒles}] 0x8B75: газетный штамп\\
\textbf{jlet} [\emph{ʒlet}] 0xAB75: календарь\\
\textbf{jli2n} [\emph{ʒli˩n}] 0x7875: невероятна\\
\textbf{jli7n} [\emph{ʒli˥n}] 0x7075: женоненавистничество\\
\textbf{jli7t} [\emph{ʒli˥t}] 0xB075: мастит\\
\textbf{jlic} [\emph{ʒliʃ}] 0xE875: спутник\\
\textbf{jlih} [\emph{ʒliʰ}] 0x43AF: антиклимакс\\
\textbf{jlik} [\emph{ʒlik}] 0x2875: геологических\\
\textbf{jlim} [\emph{ʒlim}] 0x0875: неологизмы\\
\textbf{jlin} [\emph{ʒlin}] 0x6875: написание\\
\textbf{jlip} [\emph{ʒlip}] 0x4875: ингалятор\\
\textbf{jlis} [\emph{ʒlis}] 0x8875: песочные часы\\
\textbf{jlit} [\emph{ʒlit}] 0xA875: единодушное\\
\textbf{jloc} [\emph{ʒloʃ}] 0xEC75: острие иглы\\
\textbf{jlok} [\emph{ʒlok}] 0x2C75: кола\\
\textbf{jlon} [\emph{ʒlon}] 0x6C75: двоеточие\\
\textbf{jlos} [\emph{ʒlos}] 0x8C75: почтовый ящик\\
\textbf{jlot} [\emph{ʒlot}] 0xAC75: растворитель\\
\textbf{jluk} [\emph{ʒluk}] 0x2A75: моллюсками\\
\textbf{jlun} [\emph{ʒlun}] 0x6A75: Юлиан\\
\textbf{jlut} [\emph{ʒlut}] 0xAA75: суглинок\\
\textbf{jo} [\emph{ʒo}] 0x32BE: высокая температура\\
\textbf{jr6n} [\emph{ʒrən}] 0x6DD5: маргарин\\
\textbf{jra2n} [\emph{ʒrä˩n}] 0x79D5: надпочечная\\
\textbf{jra2t} [\emph{ʒrä˩t}] 0xB9D5: кариатида\\
\textbf{jra7k} [\emph{ʒrä˥k}] 0x31D5: матриарх\\
\textbf{jra7n} [\emph{ʒrä˥n}] 0x71D5: кесарево сечение\\
\textbf{jra7t} [\emph{ʒrä˥t}] 0xB1D5: Пра-пра-бабушка\\
\textbf{jrac} [\emph{ʒräʃ}] 0xE9D5: карта\\
\textbf{jraf} [\emph{ʒräf}] 0xC9D5: архитрав\\
\textbf{jrah} [\emph{ʒräʰ}] 0x4EAF: настоящее время причастие\\
\textbf{jrak} [\emph{ʒräk}] 0x29D5: лягушка\\
\textbf{jram} [\emph{ʒräm}] 0x09D5: макраме\\
\textbf{jran} [\emph{ʒrän}] 0x69D5: врач хирург\\
\textbf{jrap} [\emph{ʒräp}] 0x49D5: отвертка\\
\textbf{jras} [\emph{ʒräs}] 0x89D5: бюрократия\\
\textbf{jrat} [\emph{ʒrät}] 0xA9D5: искусственный\\
\textbf{jre2t} [\emph{ʒre˩t}] 0xBBD5: мартышка\\
\textbf{jrek} [\emph{ʒrek}] 0x2BD5: черный как смоль\\
\textbf{jrem} [\emph{ʒrem}] 0x0BD5: герань\\
\textbf{jren} [\emph{ʒren}] 0x6BD5: выравнивание\\
\textbf{jrep} [\emph{ʒrep}] 0x4BD5: неоседланный\\
\textbf{jres} [\emph{ʒres}] 0x8BD5: ночная рубашка\\
\textbf{jret} [\emph{ʒret}] 0xABD5: Острота\\
\textbf{jri7n} [\emph{ʒri˥n}] 0x70D5: ИНГ\\
\textbf{jri7t} [\emph{ʒri˥t}] 0xB0D5: пешеходный мост\\
\textbf{jric} [\emph{ʒriʃ}] 0xE8D5: маршрут\\
\textbf{jrif} [\emph{ʒrif}] 0xC8D5: выводок\\
\textbf{jrih} [\emph{ʒriʰ}] 0x46AF: терпение\\
\textbf{jrik} [\emph{ʒrik}] 0x28D5: крикет\\
\textbf{jrim} [\emph{ʒrim}] 0x08D5: Иеремия\\
\textbf{jrin} [\emph{ʒrin}] 0x68D5: каждый день\\
\textbf{jrip} [\emph{ʒrip}] 0x48D5: призывников\\
\textbf{jris} [\emph{ʒris}] 0x88D5: пассивный причастие\\
\textbf{jrit} [\emph{ʒrit}] 0xA8D5: оправдывать\\
\textbf{jroc} [\emph{ʒroʃ}] 0xECD5: Козерог\\
\textbf{jrok} [\emph{ʒrok}] 0x2CD5: Коперник\\
\textbf{jrom} [\emph{ʒrom}] 0x0CD5: теплым сердцем\\
\textbf{jron} [\emph{ʒron}] 0x6CD5: Иордания\\
\textbf{jros} [\emph{ʒros}] 0x8CD5: Коморские\\
\textbf{jrot} [\emph{ʒrot}] 0xACD5: жареные\\
\textbf{jruk} [\emph{ʒruk}] 0x2AD5: хрящ\\
\textbf{jrun} [\emph{ʒrun}] 0x6AD5: Камерун\\
\textbf{jrus} [\emph{ʒrus}] 0x8AD5: юра\\
\textbf{jrut} [\emph{ʒrut}] 0xAAD5: жемчуг\\
\textbf{jwac} [\emph{ʒwäʃ}] 0xE935: Ява\\
\textbf{jwak} [\emph{ʒwäk}] 0x2935: йога\\
\textbf{jwam} [\emph{ʒwäm}] 0x0935: смазывать\\
\textbf{jwan} [\emph{ʒwän}] 0x6935: проспект\\
\textbf{jwat} [\emph{ʒwät}] 0xA935: тротуар\\
\textbf{jwen} [\emph{ʒwen}] 0x6B35: лежачий\\
\textbf{jwet} [\emph{ʒwet}] 0xAB35: лишать наследства\\
\textbf{jwic} [\emph{ʒwiʃ}] 0xE835: омоложение\\
\textbf{jwih} [\emph{ʒwiʰ}] 0x41AF: jussive mood, for remote commands that are to be conveyed to another recipient.\\
\textbf{jwik} [\emph{ʒwik}] 0x2835: желток\\
\textbf{jwin} [\emph{ʒwin}] 0x6835: печень\\
\textbf{jwis} [\emph{ʒwis}] 0x8835: Близнецы\\
\textbf{jwit} [\emph{ʒwit}] 0xA835: разнообразие\\
\textbf{jwon} [\emph{ʒwon}] 0x6C35: Джон\\
\textbf{jwot} [\emph{ʒwot}] 0xAC35: красное пальто\\
\textbf{jya2n} [\emph{ʒjä˩n}] 0x7915: Инь Янь\\
\textbf{jya2t} [\emph{ʒjä˩t}] 0xB915: джихад\\
\textbf{jya7n} [\emph{ʒjä˥n}] 0x7115: водиться\\
\textbf{jya7t} [\emph{ʒjä˥t}] 0xB115: в натуральную величину\\
\textbf{jyac} [\emph{ʒjäʃ}] 0xE915: интернационализация\\
\textbf{jyak} [\emph{ʒjäk}] 0x2915: вмещать\\
\textbf{jyam} [\emph{ʒjäm}] 0x0915: джайнизм\\
\textbf{jyan} [\emph{ʒjän}] 0x6915: жаргон\\
\textbf{jyap} [\emph{ʒjäp}] 0x4915: амперметр\\
\textbf{jyas} [\emph{ʒjäs}] 0x8915: Таджикистан\\
\textbf{jyat} [\emph{ʒjät}] 0xA915: банкрот\\
\textbf{jyek} [\emph{ʒjek}] 0x2B15: термоядерный\\
\textbf{jyen} [\emph{ʒjen}] 0x6B15: иена\\
\textbf{jyet} [\emph{ʒjet}] 0xAB15: душеприказчица\\
\textbf{jyi7t} [\emph{ʒji˥t}] 0xB015: Последняя минута\\
\textbf{jyic} [\emph{ʒjiʃ}] 0xE815: заложник\\
\textbf{jyih} [\emph{ʒjiʰ}] 0x40AF: не человек\\
\textbf{jyik} [\emph{ʒjik}] 0x2815: безличный глагол\\
\textbf{jyim} [\emph{ʒjim}] 0x0815: я\\
\textbf{jyin} [\emph{ʒjin}] 0x6815: инжиниринг\\
\textbf{jyip} [\emph{ʒjip}] 0x4815: генотипическая\\
\textbf{jyis} [\emph{ʒjis}] 0x8815: Овен\\
\textbf{jyit} [\emph{ʒjit}] 0xA815: невозмутимость\\
\textbf{jyok} [\emph{ʒjok}] 0x2C15: безделушка\\
\textbf{jyon} [\emph{ʒjon}] 0x6C15: Генуя\\
\textbf{jyot} [\emph{ʒjot}] 0xAC15: конденсатор\\
\textbf{jyun} [\emph{ʒjun}] 0x6A15: остроугольный\\
\textbf{jyut} [\emph{ʒjut}] 0xAA15: все цели\\
\subsection{ k }
\textbf{k6} [\emph{kə}] 0x345E: назидательная речь соглашение\\
\textbf{ka} [\emph{kä}] 0x245E: винительный дело\\
\textbf{kc6c} [\emph{kʃəʃ}] 0xEDA2: совокупление\\
\textbf{kc6k} [\emph{kʃək}] 0x2DA2: портретный\\
\textbf{kc6n} [\emph{kʃən}] 0x6DA2: уполномоченный\\
\textbf{kc6s} [\emph{kʃəs}] 0x8DA2: восьмеричный\\
\textbf{kc6t} [\emph{kʃət}] 0xADA2: клитор\\
\textbf{kca2c} [\emph{kʃä˩ʃ}] 0xF9A2: обналичить\\
\textbf{kca2k} [\emph{kʃä˩k}] 0x39A2: черный рынок\\
\textbf{kca2n} [\emph{kʃä˩n}] 0x79A2: неоднозначность\\
\textbf{kca2s} [\emph{kʃä˩s}] 0x99A2: ачехский\\
\textbf{kca2t} [\emph{kʃä˩t}] 0xB9A2: атолл\\
\textbf{kca7c} [\emph{kʃä˥ʃ}] 0xF1A2: каучук\\
\textbf{kca7k} [\emph{kʃä˥k}] 0x31A2: акробатический\\
\textbf{kca7m} [\emph{kʃä˥m}] 0x11A2: крематорий\\
\textbf{kca7n} [\emph{kʃä˥n}] 0x71A2: абсент\\
\textbf{kca7p} [\emph{kʃä˥p}] 0x51A2: довершение пирог\\
\textbf{kca7s} [\emph{kʃä˥s}] 0x91A2: Кавказ\\
\textbf{kca7t} [\emph{kʃä˥t}] 0xB1A2: неосуществимость\\
\textbf{kcac} [\emph{kʃäʃ}] 0xE9A2: оставил\\
\textbf{kcaf} [\emph{kʃäf}] 0xC9A2: шарф\\
\textbf{kcah} [\emph{kʃäʰ}] 0x4D17: голодный\\
\textbf{kcak} [\emph{kʃäk}] 0x29A2: еда\\
\textbf{kcam} [\emph{kʃäm}] 0x09A2: гордость\\
\textbf{kcan} [\emph{kʃän}] 0x69A2: считают,\\
\textbf{kcap} [\emph{kʃäp}] 0x49A2: корабль\\
\textbf{kcas} [\emph{kʃäs}] 0x89A2: темнота\\
\textbf{kcat} [\emph{kʃät}] 0xA9A2: вместе\\
\textbf{kce2n} [\emph{kʃe˩n}] 0x7BA2: четвертичный\\
\textbf{kce2t} [\emph{kʃe˩t}] 0xBBA2: декстрин\\
\textbf{kce7k} [\emph{kʃe˥k}] 0x33A2: переучет\\
\textbf{kce7n} [\emph{kʃe˥n}] 0x73A2: кессон\\
\textbf{kce7t} [\emph{kʃe˥t}] 0xB3A2: электролит\\
\textbf{kcec} [\emph{kʃeʃ}] 0xEBA2: область\\
\textbf{kcef} [\emph{kʃef}] 0xCBA2: равноденственный\\
\textbf{kceh} [\emph{kʃeʰ}] 0x5D17: лучше\\
\textbf{kcek} [\emph{kʃek}] 0x2BA2: кривая\\
\textbf{kcem} [\emph{kʃem}] 0x0BA2: ромашка\\
\textbf{kcen} [\emph{kʃen}] 0x6BA2: девственница\\
\textbf{kcep} [\emph{kʃep}] 0x4BA2: кетчуп\\
\textbf{kces} [\emph{kʃes}] 0x8BA2: кассета\\
\textbf{kcet} [\emph{kʃet}] 0xABA2: аренда\\
\textbf{kci2c} [\emph{kʃi˩ʃ}] 0xF8A2: осевой\\
\textbf{kci2k} [\emph{kʃi˩k}] 0x38A2: решетчатая постройка\\
\textbf{kci2m} [\emph{kʃi˩m}] 0x18A2: маслобойня\\
\textbf{kci2n} [\emph{kʃi˩n}] 0x78A2: дружки\\
\textbf{kci2s} [\emph{kʃi˩s}] 0x98A2: штампами\\
\textbf{kci2t} [\emph{kʃi˩t}] 0xB8A2: с открытым сердцем\\
\textbf{kci7c} [\emph{kʃi˥ʃ}] 0xF0A2: отрыжка\\
\textbf{kci7k} [\emph{kʃi˥k}] 0x30A2: лаконично\\
\textbf{kci7m} [\emph{kʃi˥m}] 0x10A2: Кашмир\\
\textbf{kci7n} [\emph{kʃi˥n}] 0x70A2: комиссии\\
\textbf{kci7s} [\emph{kʃi˥s}] 0x90A2: архаизм\\
\textbf{kci7t} [\emph{kʃi˥t}] 0xB0A2: скупо\\
\textbf{kcic} [\emph{kʃiʃ}] 0xE8A2: любопытный\\
\textbf{kcif} [\emph{kʃif}] 0xC8A2: дефицит\\
\textbf{kcih} [\emph{kʃiʰ}] 0x4517: пожалуйста\\
\textbf{kcik} [\emph{kʃik}] 0x28A2: рыба\\
\textbf{kcim} [\emph{kʃim}] 0x08A2: неограниченный\\
\textbf{kcin} [\emph{kʃin}] 0x68A2: глубоко\\
\textbf{kcip} [\emph{kʃip}] 0x48A2: чемпион\\
\textbf{kcis} [\emph{kʃis}] 0x88A2: серый\\
\textbf{kcit} [\emph{kʃit}] 0xA8A2: в браке\\
\textbf{kco2t} [\emph{kʃo˩t}] 0xBCA2: картечь\\
\textbf{kco7n} [\emph{kʃo˥n}] 0x74A2: хромота\\
\textbf{kco7t} [\emph{kʃo˥t}] 0xB4A2: комод\\
\textbf{kcoc} [\emph{kʃoʃ}] 0xECA2: локоть\\
\textbf{kcof} [\emph{kʃof}] 0xCCA2: ремесла\\
\textbf{kcoh} [\emph{kʃoʰ}] 0x6517: устала\\
\textbf{kcok} [\emph{kʃok}] 0x2CA2: топор\\
\textbf{kcom} [\emph{kʃom}] 0x0CA2: экономист\\
\textbf{kcon} [\emph{kʃon}] 0x6CA2: откровение\\
\textbf{kcop} [\emph{kʃop}] 0x4CA2: курятник\\
\textbf{kcos} [\emph{kʃos}] 0x8CA2: обвинить\\
\textbf{kcot} [\emph{kʃot}] 0xACA2: собирать\\
\textbf{kcu2n} [\emph{kʃu˩n}] 0x7AA2: близко вяжут\\
\textbf{kcu7h} [\emph{kʃu˥ʰ}] 0x9517: индуктивный настроение\\
\textbf{kcu7t} [\emph{kʃu˥t}] 0xB2A2: открытого состава\\
\textbf{kcuc} [\emph{kʃuʃ}] 0xEAA2: горький\\
\textbf{kcuh} [\emph{kʃuʰ}] 0x9517: choice gender, 3rd density life form, learning to make choices, choosing whether to be service to self, service to others, or to remain lost, akin to basic homo sapiens, cetaceans (dolphins/whales), advanced animals\\
\textbf{kcuk} [\emph{kʃuk}] 0x2AA2: школа\\
\textbf{kcum} [\emph{kʃum}] 0x0AA2: углубления\\
\textbf{kcun} [\emph{kʃun}] 0x6AA2: отверстие\\
\textbf{kcup} [\emph{kʃup}] 0x4AA2: купон\\
\textbf{kcus} [\emph{kʃus}] 0x8AA2: готовить\\
\textbf{kcut} [\emph{kʃut}] 0xAAA2: выберите\\
\textbf{ke} [\emph{ke}] 0x2C5E: вопрос\\
\textbf{ke2} [\emph{ke˩}] 0x6C5E: эр\\
\textbf{ke7} [\emph{ke˥}] 0x4C5E: ан\\
\textbf{kf6n} [\emph{kfən}] 0x6D82: оксюморон\\
\textbf{kf6t} [\emph{kfət}] 0xAD82: отслаиваться\\
\textbf{kfa2f} [\emph{kfä˩f}] 0xD982: квантификатором\\
\textbf{kfa2k} [\emph{kfä˩k}] 0x3982: акромегалия\\
\textbf{kfa2n} [\emph{kfä˩n}] 0x7982: какофония\\
\textbf{kfa2t} [\emph{kfä˩t}] 0xB982: кафтан\\
\textbf{kfa7k} [\emph{kfä˥k}] 0x3182: известкование\\
\textbf{kfa7n} [\emph{kfä˥n}] 0x7182: камуфляж\\
\textbf{kfa7t} [\emph{kfä˥t}] 0xB182: обрезываете\\
\textbf{kfac} [\emph{kfäʃ}] 0xE982: трава\\
\textbf{kfaf} [\emph{kfäf}] 0xC982: прививка\\
\textbf{kfah} [\emph{kfäʰ}] 0x4C17: жертва\\
\textbf{kfak} [\emph{kfäk}] 0x2982: израсходовать\\
\textbf{kfam} [\emph{kfäm}] 0x0982: квантовая\\
\textbf{kfan} [\emph{kfän}] 0x6982: поведение\\
\textbf{kfap} [\emph{kfäp}] 0x4982: октября\\
\textbf{kfas} [\emph{kfäs}] 0x8982: рак\\
\textbf{kfat} [\emph{kfät}] 0xA982: сметь\\
\textbf{kfe7n} [\emph{kfe˥n}] 0x7382: кофеин\\
\textbf{kfe7t} [\emph{kfe˥t}] 0xB382: налог бесплатный\\
\textbf{kfec} [\emph{kfeʃ}] 0xEB82: кеч\\
\textbf{kfef} [\emph{kfef}] 0xCB82: кинестетическая\\
\textbf{kfeh} [\emph{kfeʰ}] 0x5C17: ядро благословляющий\\
\textbf{kfek} [\emph{kfek}] 0x2B82: конвекция\\
\textbf{kfen} [\emph{kfen}] 0x6B82: понятия\\
\textbf{kfep} [\emph{kfep}] 0x4B82: отлично справляется\\
\textbf{kfes} [\emph{kfes}] 0x8B82: ломонос\\
\textbf{kfet} [\emph{kfet}] 0xAB82: катод\\
\textbf{kfi2k} [\emph{kfi˩k}] 0x3882: детоксификация\\
\textbf{kfi2n} [\emph{kfi˩n}] 0x7882: занос\\
\textbf{kfi2t} [\emph{kfi˩t}] 0xB882: нива\\
\textbf{kfi7k} [\emph{kfi˥k}] 0x3082: алифатический\\
\textbf{kfi7n} [\emph{kfi˥n}] 0x7082: кодифицированный\\
\textbf{kfi7s} [\emph{kfi˥s}] 0x9082: конформист\\
\textbf{kfi7t} [\emph{kfi˥t}] 0xB082: аконит\\
\textbf{kfic} [\emph{kfiʃ}] 0xE882: отпуск\\
\textbf{kfif} [\emph{kfif}] 0xC882: ядро благословляющий\\
\textbf{kfih} [\emph{kfiʰ}] 0x4417: само уверенность\\
\textbf{kfik} [\emph{kfik}] 0x2882: технический\\
\textbf{kfim} [\emph{kfim}] 0x0882: корневище\\
\textbf{kfin} [\emph{kfin}] 0x6882: крышка\\
\textbf{kfip} [\emph{kfip}] 0x4882: неприемлемо\\
\textbf{kfis} [\emph{kfis}] 0x8882: вакансия\\
\textbf{kfit} [\emph{kfit}] 0xA882: лифт\\
\textbf{kfo2t} [\emph{kfo˩t}] 0xBC82: скакать\\
\textbf{kfo7n} [\emph{kfo˥n}] 0x7482: аккордеоны\\
\textbf{kfo7t} [\emph{kfo˥t}] 0xB482: я делаю\\
\textbf{kfoc} [\emph{kfoʃ}] 0xEC82: Хорватия\\
\textbf{kfof} [\emph{kfof}] 0xCC82: манжета\\
\textbf{kfok} [\emph{kfok}] 0x2C82: икота\\
\textbf{kfom} [\emph{kfom}] 0x0C82: номера\\
\textbf{kfon} [\emph{kfon}] 0x6C82: кодирование\\
\textbf{kfos} [\emph{kfos}] 0x8C82: компостирование\\
\textbf{kfot} [\emph{kfot}] 0xAC82: выдача\\
\textbf{kfu7t} [\emph{kfu˥t}] 0xB282: Kathmandu\\
\textbf{kfuf} [\emph{kfuf}] 0xCA82: крестоцветный\\
\textbf{kfuk} [\emph{kfuk}] 0x2A82: уход за кожей\\
\textbf{kfun} [\emph{kfun}] 0x6A82: размещение\\
\textbf{kfup} [\emph{kfup}] 0x4A82: куб\\
\textbf{kfus} [\emph{kfus}] 0x8A82: диффузный\\
\textbf{kfut} [\emph{kfut}] 0xAA82: комендантский час\\
\textbf{ki} [\emph{ki}] 0x205E: мимо напряженный\\
\textbf{ki2} [\emph{ki˩}] 0x605E: чрезмерное продолжительность\\
\textbf{ki7} [\emph{ki˥}] 0x405E: словесный имя существительное\\
\textbf{kl67t} [\emph{klə˥t}] 0xB562: циклотрон\\
\textbf{kl6c} [\emph{kləʃ}] 0xED62: коллаген\\
\textbf{kl6f} [\emph{kləf}] 0xCD62: Голгофа\\
\textbf{kl6k} [\emph{klək}] 0x2D62: чашечка\\
\textbf{kl6m} [\emph{kləm}] 0x0D62: скаляр\\
\textbf{kl6n} [\emph{klən}] 0x6D62: дыня\\
\textbf{kl6p} [\emph{kləp}] 0x4D62: Колумбус\\
\textbf{kl6s} [\emph{kləs}] 0x8D62: крейсерский\\
\textbf{kl6t} [\emph{klət}] 0xAD62: подставка\\
\textbf{kla2k} [\emph{klä˩k}] 0x3962: отпереть\\
\textbf{kla2m} [\emph{klä˩m}] 0x1962: камелия\\
\textbf{kla2n} [\emph{klä˩n}] 0x7962: холера\\
\textbf{kla2p} [\emph{klä˩p}] 0x5962: спина к спине\\
\textbf{kla2s} [\emph{klä˩s}] 0x9962: календы\\
\textbf{kla2t} [\emph{klä˩t}] 0xB962: оригинальность\\
\textbf{kla7c} [\emph{klä˥ʃ}] 0xF162: алкалоид\\
\textbf{kla7f} [\emph{klä˥f}] 0xD162: каллиграф\\
\textbf{kla7k} [\emph{klä˥k}] 0x3162: дисквалификация\\
\textbf{kla7m} [\emph{klä˥m}] 0x1162: катаболизм\\
\textbf{kla7n} [\emph{klä˥n}] 0x7162: орлиный\\
\textbf{kla7p} [\emph{klä˥p}] 0x5162: кабала\\
\textbf{kla7s} [\emph{klä˥s}] 0x9162: скалы\\
\textbf{kla7t} [\emph{klä˥t}] 0xB162: способности\\
\textbf{klac} [\emph{kläʃ}] 0xE962: страна\\
\textbf{klaf} [\emph{kläf}] 0xC962: качественный\\
\textbf{klah} [\emph{kläʰ}] 0x4B17: там\\
\textbf{klak} [\emph{kläk}] 0x2962: собака\\
\textbf{klam} [\emph{kläm}] 0x0962: говоря\\
\textbf{klan} [\emph{klän}] 0x6962: там\\
\textbf{klap} [\emph{kläp}] 0x4962: ковер\\
\textbf{klas} [\emph{kläs}] 0x8962: стакан\\
\textbf{klat} [\emph{klät}] 0xA962: капитан\\
\textbf{kle2n} [\emph{kle˩n}] 0x7B62: колоннада\\
\textbf{kle2t} [\emph{kle˩t}] 0xBB62: выздоравливающий\\
\textbf{kle7n} [\emph{kle˥n}] 0x7362: конкатенация\\
\textbf{kle7t} [\emph{kle˥t}] 0xB362: реабилитировать\\
\textbf{klec} [\emph{kleʃ}] 0xEB62: собака устала\\
\textbf{klef} [\emph{klef}] 0xCB62: экстраполяция\\
\textbf{kleh} [\emph{kleʰ}] 0x5B17: собака устала\\
\textbf{klek} [\emph{klek}] 0x2B62: щелчок\\
\textbf{klem} [\emph{klem}] 0x0B62: флюид\\
\textbf{klen} [\emph{klen}] 0x6B62: колония\\
\textbf{klep} [\emph{klep}] 0x4B62: ламинария\\
\textbf{kles} [\emph{kles}] 0x8B62: обезглавить\\
\textbf{klet} [\emph{klet}] 0xAB62: относительно\\
\textbf{kli2c} [\emph{kli˩ʃ}] 0xF862: микология\\
\textbf{kli2f} [\emph{kli˩f}] 0xD862: факультативный\\
\textbf{kli2k} [\emph{kli˩k}] 0x3862: клеенка\\
\textbf{kli2n} [\emph{kli˩n}] 0x7862: кавалерия\\
\textbf{kli2p} [\emph{kli˩p}] 0x5862: безутешный\\
\textbf{kli2s} [\emph{kli˩s}] 0x9862: акклиматизировались\\
\textbf{kli2t} [\emph{kli˩t}] 0xB862: заклепка\\
\textbf{kli7c} [\emph{kli˥ʃ}] 0xF062: перекристаллизация\\
\textbf{kli7f} [\emph{kli˥f}] 0xD062: халиф\\
\textbf{kli7k} [\emph{kli˥k}] 0x3062: оксихлорид\\
\textbf{kli7n} [\emph{kli˥n}] 0x7062: кастильский\\
\textbf{kli7p} [\emph{kli˥p}] 0x5062: кольраби\\
\textbf{kli7s} [\emph{kli˥s}] 0x9062: выслушивание\\
\textbf{kli7t} [\emph{kli˥t}] 0xB062: кельтская\\
\textbf{klic} [\emph{kliʃ}] 0xE862: колледж\\
\textbf{klif} [\emph{klif}] 0xC862: кальций\\
\textbf{klih} [\emph{kliʰ}] 0x4317: виновный\\
\textbf{klik} [\emph{klik}] 0x2862: долина\\
\textbf{klim} [\emph{klim}] 0x0862: гримаса\\
\textbf{klin} [\emph{klin}] 0x6862: столица\\
\textbf{klip} [\emph{klip}] 0x4862: ремесленник\\
\textbf{klis} [\emph{klis}] 0x8862: изоляция\\
\textbf{klit} [\emph{klit}] 0xA862: неизвестный\\
\textbf{klo2n} [\emph{klo˩n}] 0x7C62: Шотландия\\
\textbf{klo2s} [\emph{klo˩s}] 0x9C62: сколиоз\\
\textbf{klo2t} [\emph{klo˩t}] 0xBC62: коллодий\\
\textbf{klo7k} [\emph{klo˥k}] 0x3462: кудахтать\\
\textbf{klo7n} [\emph{klo˥n}] 0x7462: золотуха\\
\textbf{klo7p} [\emph{klo˥p}] 0x5462: предварительно колумбийский\\
\textbf{klo7s} [\emph{klo˥s}] 0x9462: оккультист\\
\textbf{klo7t} [\emph{klo˥t}] 0xB462: с открытым ртом\\
\textbf{kloc} [\emph{kloʃ}] 0xEC62: воротник\\
\textbf{klof} [\emph{klof}] 0xCC62: клевер\\
\textbf{kloh} [\emph{kloʰ}] 0x6317: Одинокий\\
\textbf{klok} [\emph{klok}] 0x2C62: избыток\\
\textbf{klom} [\emph{klom}] 0x0C62: километр\\
\textbf{klon} [\emph{klon}] 0x6C62: революция\\
\textbf{klop} [\emph{klop}] 0x4C62: калория\\
\textbf{klos} [\emph{klos}] 0x8C62: хлор\\
\textbf{klot} [\emph{klot}] 0xAC62: тренер\\
\textbf{klu2n} [\emph{klu˩n}] 0x7A62: сговариваться\\
\textbf{klu2t} [\emph{klu˩t}] 0xBA62: нуклеотид\\
\textbf{klu7s} [\emph{klu˥s}] 0x9262: исключительность\\
\textbf{klu7t} [\emph{klu˥t}] 0xB262: огрубевший\\
\textbf{kluc} [\emph{kluʃ}] 0xEA62: инструменты\\
\textbf{kluf} [\emph{kluf}] 0xCA62: ключ\\
\textbf{kluh} [\emph{kluʰ}] 0x5317: путаница\\
\textbf{kluk} [\emph{kluk}] 0x2A62: археологический\\
\textbf{klum} [\emph{klum}] 0x0A62: совокупляться\\
\textbf{klun} [\emph{klun}] 0x6A62: накопление\\
\textbf{klup} [\emph{klup}] 0x4A62: клуб\\
\textbf{klus} [\emph{klus}] 0x8A62: эвкалипт\\
\textbf{klut} [\emph{klut}] 0xAA62: спекулировать\\
\textbf{ko} [\emph{ko}] 0x305E: Предложение хвост\\
\textbf{kr62t} [\emph{krə˩t}] 0xBDC2: сросшийся\\
\textbf{kr67n} [\emph{krə˥n}] 0x75C2: кератин\\
\textbf{kr67t} [\emph{krə˥t}] 0xB5C2: курд\\
\textbf{kr6c} [\emph{krəʃ}] 0xEDC2: штопор\\
\textbf{kr6f} [\emph{krəf}] 0xCDC2: Каринтии\\
\textbf{kr6k} [\emph{krək}] 0x2DC2: скорпион\\
\textbf{kr6m} [\emph{krəm}] 0x0DC2: презерватив\\
\textbf{kr6n} [\emph{krən}] 0x6DC2: корнеплоды\\
\textbf{kr6p} [\emph{krəp}] 0x4DC2: карибский\\
\textbf{kr6s} [\emph{krəs}] 0x8DC2: кокос\\
\textbf{kr6t} [\emph{krət}] 0xADC2: электрический текущий\\
\textbf{kra2c} [\emph{krä˩ʃ}] 0xF9C2: камеры\\
\textbf{kra2k} [\emph{krä˩k}] 0x39C2: скалистый\\
\textbf{kra2m} [\emph{krä˩m}] 0x19C2: Катись\\
\textbf{kra2n} [\emph{krä˩n}] 0x79C2: проказа\\
\textbf{kra2p} [\emph{krä˩p}] 0x59C2: карбид\\
\textbf{kra2s} [\emph{krä˩s}] 0x99C2: фураж\\
\textbf{kra2t} [\emph{krä˩t}] 0xB9C2: откровенность\\
\textbf{kra7c} [\emph{krä˥ʃ}] 0xF1C2: панцирь\\
\textbf{kra7f} [\emph{krä˥f}] 0xD1C2: выбивать\\
\textbf{kra7k} [\emph{krä˥k}] 0x31C2: карболка\\
\textbf{kra7m} [\emph{krä˥m}] 0x11C2: лакрица\\
\textbf{kra7n} [\emph{krä˥n}] 0x71C2: арки\\
\textbf{kra7p} [\emph{krä˥p}] 0x51C2: карп\\
\textbf{kra7s} [\emph{krä˥s}] 0x91C2: колесный\\
\textbf{kra7t} [\emph{krä˥t}] 0xB1C2: возчик\\
\textbf{krac} [\emph{kräʃ}] 0xE9C2: жестокий\\
\textbf{kraf} [\emph{kräf}] 0xC9C2: ворона\\
\textbf{krah} [\emph{kräʰ}] 0x4E17: гордость\\
\textbf{krak} [\emph{kräk}] 0x29C2: признание\\
\textbf{kram} [\emph{kräm}] 0x09C2: атом\\
\textbf{kran} [\emph{krän}] 0x69C2: корона\\
\textbf{krap} [\emph{kräp}] 0x49C2: хлопок\\
\textbf{kras} [\emph{kräs}] 0x89C2: гроб\\
\textbf{krat} [\emph{krät}] 0xA9C2: император\\
\textbf{kre2k} [\emph{kre˩k}] 0x3BC2: электроскоп\\
\textbf{kre2n} [\emph{kre˩n}] 0x7BC2: четырехугольник\\
\textbf{kre2t} [\emph{kre˩t}] 0xBBC2: картель\\
\textbf{kre7k} [\emph{kre˥k}] 0x33C2: четвероногие\\
\textbf{kre7n} [\emph{kre˥n}] 0x73C2: Looker на\\
\textbf{kre7p} [\emph{kre˥p}] 0x53C2: херувим\\
\textbf{kre7s} [\emph{kre˥s}] 0x93C2: азартная игра в кости\\
\textbf{kre7t} [\emph{kre˥t}] 0xB3C2: Крит\\
\textbf{krec} [\emph{kreʃ}] 0xEBC2: кран\\
\textbf{kref} [\emph{kref}] 0xCBC2: тмин\\
\textbf{kreh} [\emph{kreʰ}] 0x5E17: жестокий\\
\textbf{krek} [\emph{krek}] 0x2BC2: Меркурий\\
\textbf{krem} [\emph{krem}] 0x0BC2: акров\\
\textbf{kren} [\emph{kren}] 0x6BC2: заместитель\\
\textbf{krep} [\emph{krep}] 0x4BC2: шифрование\\
\textbf{kres} [\emph{kres}] 0x8BC2: глагол\\
\textbf{kret} [\emph{kret}] 0xABC2: экстракт\\
\textbf{kri2c} [\emph{kri˩ʃ}] 0xF8C2: неразборчивость\\
\textbf{kri2k} [\emph{kri˩k}] 0x38C2: каприз\\
\textbf{kri2m} [\emph{kri˩m}] 0x18C2: криптограмма\\
\textbf{kri2n} [\emph{kri˩n}] 0x78C2: Коринф\\
\textbf{kri2p} [\emph{kri˩p}] 0x58C2: карибу\\
\textbf{kri2s} [\emph{kri˩s}] 0x98C2: курсы\\
\textbf{kri2t} [\emph{kri˩t}] 0xB8C2: обмотанный\\
\textbf{kri7c} [\emph{kri˥ʃ}] 0xF0C2: сахарин\\
\textbf{kri7f} [\emph{kri˥f}] 0xD0C2: выделять курсивом\\
\textbf{kri7k} [\emph{kri˥k}] 0x30C2: некритичный\\
\textbf{kri7n} [\emph{kri˥n}] 0x70C2: выгравирован\\
\textbf{kri7p} [\emph{kri˥p}] 0x50C2: микроскопия\\
\textbf{kri7s} [\emph{kri˥s}] 0x90C2: замковый камень\\
\textbf{kri7t} [\emph{kri˥t}] 0xB0C2: жертвы\\
\textbf{kric} [\emph{kriʃ}] 0xE8C2: курица\\
\textbf{krif} [\emph{krif}] 0xC8C2: секреция\\
\textbf{krih} [\emph{kriʰ}] 0x4617: любопытный\\
\textbf{krik} [\emph{krik}] 0x28C2: критика\\
\textbf{krim} [\emph{krim}] 0x08C2: крем\\
\textbf{krin} [\emph{krin}] 0x68C2: кардинальный\\
\textbf{krip} [\emph{krip}] 0x48C2: дискриминация\\
\textbf{kris} [\emph{kris}] 0x88C2: секретарь\\
\textbf{krit} [\emph{krit}] 0xA8C2: читать корректуру\\
\textbf{kro2k} [\emph{kro˩k}] 0x3CC2: кортеж\\
\textbf{kro2n} [\emph{kro˩n}] 0x7CC2: коронер\\
\textbf{kro2s} [\emph{kro˩s}] 0x9CC2: пробки\\
\textbf{kro2t} [\emph{kro˩t}] 0xBCC2: приусадебный участок\\
\textbf{kro7k} [\emph{kro˥k}] 0x34C2: многокрасочный\\
\textbf{kro7m} [\emph{kro˥m}] 0x14C2: кворум\\
\textbf{kro7n} [\emph{kro˥n}] 0x74C2: викторианский\\
\textbf{kro7p} [\emph{kro˥p}] 0x54C2: нагрузка\\
\textbf{kro7s} [\emph{kro˥s}] 0x94C2: курсор\\
\textbf{kro7t} [\emph{kro˥t}] 0xB4C2: раздвоенный\\
\textbf{kroc} [\emph{kroʃ}] 0xECC2: борозда\\
\textbf{krof} [\emph{krof}] 0xCCC2: Корсика\\
\textbf{kroh} [\emph{kroʰ}] 0x6617: очарование\\
\textbf{krok} [\emph{krok}] 0x2CC2: череп\\
\textbf{krom} [\emph{krom}] 0x0CC2: хром\\
\textbf{kron} [\emph{kron}] 0x6CC2: углерод\\
\textbf{krop} [\emph{krop}] 0x4CC2: координировать\\
\textbf{kros} [\emph{kros}] 0x8CC2: кварц\\
\textbf{krot} [\emph{krot}] 0xACC2: резка\\
\textbf{kru2k} [\emph{kru˩k}] 0x3AC2: учетверять\\
\textbf{kru2n} [\emph{kru˩n}] 0x7AC2: корунд\\
\textbf{kru2t} [\emph{kru˩t}] 0xBAC2: эукариоты\\
\textbf{kru7n} [\emph{kru˥n}] 0x72C2: Контратенор\\
\textbf{kru7s} [\emph{kru˥s}] 0x92C2: Водолей\\
\textbf{kru7t} [\emph{kru˥t}] 0xB2C2: квартет\\
\textbf{kruc} [\emph{kruʃ}] 0xEAC2: черепаха\\
\textbf{kruf} [\emph{kruf}] 0xCAC2: Украина\\
\textbf{kruh} [\emph{kruʰ}] 0x5617: горький\\
\textbf{kruk} [\emph{kruk}] 0x2AC2: костыль\\
\textbf{krum} [\emph{krum}] 0x0AC2: скумбрия\\
\textbf{krun} [\emph{krun}] 0x6AC2: полнота\\
\textbf{krup} [\emph{krup}] 0x4AC2: Акупрессура\\
\textbf{krus} [\emph{krus}] 0x8AC2: приседать\\
\textbf{krut} [\emph{krut}] 0xAAC2: укусить\\
\textbf{ks67t} [\emph{ksə˥t}] 0xB542: горьковато-сладкий\\
\textbf{ks6c} [\emph{ksəʃ}] 0xED42: инкрустация\\
\textbf{ks6h} [\emph{ksəʰ}] 0x6A17: причинный\\
\textbf{ks6k} [\emph{ksək}] 0x2D42: акустика\\
\textbf{ks6n} [\emph{ksən}] 0x6D42: рабочая станция\\
\textbf{ks6p} [\emph{ksəp}] 0x4D42: тигель\\
\textbf{ks6s} [\emph{ksəs}] 0x8D42: Пакистан\\
\textbf{ks6t} [\emph{ksət}] 0xAD42: причинный\\
\textbf{ksa2c} [\emph{ksä˩ʃ}] 0xF942: за кулисы\\
\textbf{ksa2k} [\emph{ksä˩k}] 0x3942: заклинать\\
\textbf{ksa2m} [\emph{ksä˩m}] 0x1942: гекзаметр\\
\textbf{ksa2n} [\emph{ksä˩n}] 0x7942: придворный\\
\textbf{ksa2p} [\emph{ksä˩p}] 0x5942: гадюка\\
\textbf{ksa2s} [\emph{ksä˩s}] 0x9942: непостижимый\\
\textbf{ksa2t} [\emph{ksä˩t}] 0xB942: каста\\
\textbf{ksa7c} [\emph{ksä˥ʃ}] 0xF142: экспертиза главный\\
\textbf{ksa7f} [\emph{ksä˥f}] 0xD142: растр\\
\textbf{ksa7k} [\emph{ksä˥k}] 0x3142: кукурузный крахмал\\
\textbf{ksa7m} [\emph{ksä˥m}] 0x1142: Ассамский\\
\textbf{ksa7n} [\emph{ksä˥n}] 0x7142: высылка\\
\textbf{ksa7p} [\emph{ksä˥p}] 0x5142: аркебуза\\
\textbf{ksa7s} [\emph{ksä˥s}] 0x9142: корь\\
\textbf{ksa7t} [\emph{ksä˥t}] 0xB142: бортовой залп\\
\textbf{ksac} [\emph{ksäʃ}] 0xE942: обстоятельства\\
\textbf{ksaf} [\emph{ksäf}] 0xC942: фальсифицировать\\
\textbf{ksah} [\emph{ksäʰ}] 0x4A17: близость\\
\textbf{ksak} [\emph{ksäk}] 0x2942: причинный дело\\
\textbf{ksam} [\emph{ksäm}] 0x0942: национальный\\
\textbf{ksan} [\emph{ksän}] 0x6942: наименее\\
\textbf{ksap} [\emph{ksäp}] 0x4942: принимать\\
\textbf{ksas} [\emph{ksäs}] 0x8942: 4\\
\textbf{ksat} [\emph{ksät}] 0xA942: приятный\\
\textbf{kse2k} [\emph{kse˩k}] 0x3B42: конспект\\
\textbf{kse2n} [\emph{kse˩n}] 0x7B42: секстан\\
\textbf{kse2p} [\emph{kse˩p}] 0x5B42: облагаемый\\
\textbf{kse2s} [\emph{kse˩s}] 0x9B42: превосходить\\
\textbf{kse2t} [\emph{kse˩t}] 0xBB42: пономарь\\
\textbf{kse7k} [\emph{kse˥k}] 0x3342: пустельга\\
\textbf{kse7n} [\emph{kse˥n}] 0x7342: казеин\\
\textbf{kse7p} [\emph{kse˥p}] 0x5342: шекспировский\\
\textbf{kse7s} [\emph{kse˥s}] 0x9342: квинтэссенция\\
\textbf{kse7t} [\emph{kse˥t}] 0xB342: этих заседаниях\\
\textbf{ksec} [\emph{kseʃ}] 0xEB42: одноклассник\\
\textbf{ksef} [\emph{ksef}] 0xCB42: Рентгеновский\\
\textbf{kseh} [\emph{kseʰ}] 0x5A17: навязчивая идея\\
\textbf{ksek} [\emph{ksek}] 0x2B42: насекомое\\
\textbf{ksem} [\emph{ksem}] 0x0B42: случайный доступ Память\\
\textbf{ksen} [\emph{ksen}] 0x6B42: кислород\\
\textbf{ksep} [\emph{ksep}] 0x4B42: рецептор\\
\textbf{kses} [\emph{kses}] 0x8B42: экскурсия\\
\textbf{kset} [\emph{kset}] 0xAB42: оркестр\\
\textbf{ksi2k} [\emph{ksi˩k}] 0x3842: шестьдесят шесть\\
\textbf{ksi2n} [\emph{ksi˩n}] 0x7842: сохранившийся\\
\textbf{ksi2p} [\emph{ksi˩p}] 0x5842: аксиллярный\\
\textbf{ksi2s} [\emph{ksi˩s}] 0x9842: акрополь\\
\textbf{ksi2t} [\emph{ksi˩t}] 0xB842: тридцать один\\
\textbf{ksi7h} [\emph{ksi˥ʰ}] 0x8217: IST\\
\textbf{ksi7k} [\emph{ksi˥k}] 0x3042: перекрестный допрос\\
\textbf{ksi7m} [\emph{ksi˥m}] 0x1042: кубизм\\
\textbf{ksi7n} [\emph{ksi˥n}] 0x7042: консалтинг\\
\textbf{ksi7p} [\emph{ksi˥p}] 0x5042: текстиль\\
\textbf{ksi7s} [\emph{ksi˥s}] 0x9042: софистика\\
\textbf{ksi7t} [\emph{ksi˥t}] 0xB042: высосанный\\
\textbf{ksic} [\emph{ksiʃ}] 0xE842: кризис\\
\textbf{ksif} [\emph{ksif}] 0xC842: окислять\\
\textbf{ksih} [\emph{ksiʰ}] 0x4217: causal case, indicates the noun phrase is the cause. person-context source-case. \\
\textbf{ksik} [\emph{ksik}] 0x2842: следствие\\
\textbf{ksim} [\emph{ksim}] 0x0842: декларативный настроение\\
\textbf{ksin} [\emph{ksin}] 0x6842: зеленый\\
\textbf{ksip} [\emph{ksip}] 0x4842: мышьяк\\
\textbf{ksis} [\emph{ksis}] 0x8842: история\\
\textbf{ksit} [\emph{ksit}] 0xA842: могила\\
\textbf{kso2k} [\emph{kso˩k}] 0x3C42: Джек в коробке\\
\textbf{kso2n} [\emph{kso˩n}] 0x7C42: с учетом\\
\textbf{kso2p} [\emph{kso˩p}] 0x5C42: макроскопический\\
\textbf{kso2s} [\emph{kso˩s}] 0x9C42: акростих\\
\textbf{kso2t} [\emph{kso˩t}] 0xBC42: украшенный гребнем\\
\textbf{kso7f} [\emph{kso˥f}] 0xD442: Косово\\
\textbf{kso7k} [\emph{kso˥k}] 0x3442: фат\\
\textbf{kso7n} [\emph{kso˥n}] 0x7442: англосаксонской\\
\textbf{kso7p} [\emph{kso˥p}] 0x5442: несъедобный\\
\textbf{kso7s} [\emph{kso˥s}] 0x9442: коробка офис\\
\textbf{kso7t} [\emph{kso˥t}] 0xB442: обуздана\\
\textbf{ksoc} [\emph{ksoʃ}] 0xEC42: вычет\\
\textbf{ksof} [\emph{ksof}] 0xCC42: саксофон\\
\textbf{ksoh} [\emph{ksoʰ}] 0x6217: сексуальное\\
\textbf{ksok} [\emph{ksok}] 0x2C42: Гороскоп\\
\textbf{ksom} [\emph{ksom}] 0x0C42: космос\\
\textbf{kson} [\emph{kson}] 0x6C42: акцент\\
\textbf{ksop} [\emph{ksop}] 0x4C42: косметический\\
\textbf{ksos} [\emph{ksos}] 0x8C42: кашель\\
\textbf{ksot} [\emph{ksot}] 0xAC42: кость\\
\textbf{ksu2n} [\emph{ksu˩n}] 0x7A42: Курдистан\\
\textbf{ksu2s} [\emph{ksu˩s}] 0x9A42: квесты\\
\textbf{ksu2t} [\emph{ksu˩t}] 0xBA42: счетчики\\
\textbf{ksu7k} [\emph{ksu˥k}] 0x3242: казуистический\\
\textbf{ksu7n} [\emph{ksu˥n}] 0x7242: кабестан\\
\textbf{ksu7s} [\emph{ksu˥s}] 0x9242: арбалет\\
\textbf{ksu7t} [\emph{ksu˥t}] 0xB242: шестьдесят один\\
\textbf{ksuc} [\emph{ksuʃ}] 0xEA42: вкусно\\
\textbf{ksuf} [\emph{ksuf}] 0xCA42: кускус\\
\textbf{ksuh} [\emph{ksuʰ}] 0x5217: declarative mood, for declaring the begining of a program or essay.\\
\textbf{ksuk} [\emph{ksuk}] 0x2A42: дядя\\
\textbf{ksum} [\emph{ksum}] 0x0A42: акционер\\
\textbf{ksun} [\emph{ksun}] 0x6A42: сжигание\\
\textbf{ksup} [\emph{ksup}] 0x4A42: косинус\\
\textbf{ksus} [\emph{ksus}] 0x8A42: кластер\\
\textbf{ksut} [\emph{ksut}] 0xAA42: убийство\\
\textbf{ku} [\emph{ku}] 0x285E: Эксклюзивный или\\
\textbf{kw6k} [\emph{kwək}] 0x2D22: кварк\\
\textbf{kw6n} [\emph{kwən}] 0x6D22: кларнет\\
\textbf{kw6t} [\emph{kwət}] 0xAD22: скваттер\\
\textbf{kwa2k} [\emph{kwä˩k}] 0x3922: каркать\\
\textbf{kwa2n} [\emph{kwä˩n}] 0x7922: кругосветное плавание\\
\textbf{kwa2s} [\emph{kwä˩s}] 0x9922: чешуйчатый\\
\textbf{kwa2t} [\emph{kwä˩t}] 0xB922: благоговение ударил\\
\textbf{kwa7k} [\emph{kwä˥k}] 0x3122: пойло\\
\textbf{kwa7n} [\emph{kwä˥n}] 0x7122: консервирование\\
\textbf{kwa7p} [\emph{kwä˥p}] 0x5122: вигны\\
\textbf{kwa7s} [\emph{kwä˥s}] 0x9122: десть\\
\textbf{kwa7t} [\emph{kwä˥t}] 0xB122: один вооруженный\\
\textbf{kwac} [\emph{kwäʃ}] 0xE922: ухо\\
\textbf{kwaf} [\emph{kwäf}] 0xC922: развивать\\
\textbf{kwah} [\emph{kwäʰ}] 0x4917: желание\\
\textbf{kwak} [\emph{kwäk}] 0x2922: покрытый\\
\textbf{kwam} [\emph{kwäm}] 0x0922: запятая\\
\textbf{kwan} [\emph{kwän}] 0x6922: легкий\\
\textbf{kwap} [\emph{kwäp}] 0x4922: квадрат\\
\textbf{kwas} [\emph{kwäs}] 0x8922: август\\
\textbf{kwat} [\emph{kwät}] 0xA922: комфорт\\
\textbf{kwe2t} [\emph{kwe˩t}] 0xBB22: кемпинг\\
\textbf{kwe7t} [\emph{kwe˥t}] 0xB322: собака ушастый\\
\textbf{kweh} [\emph{kweʰ}] 0x5917: эксклюзивный человек\\
\textbf{kwek} [\emph{kwek}] 0x2B22: местный дело\\
\textbf{kwen} [\emph{kwen}] 0x6B22: эксклюзивный человек\\
\textbf{kwep} [\emph{kwep}] 0x4B22: спектроскоп\\
\textbf{kwes} [\emph{kwes}] 0x8B22: конгрессменов\\
\textbf{kwet} [\emph{kwet}] 0xAB22: кататься на коньках\\
\textbf{kwi2k} [\emph{kwi˩k}] 0x3822: калитка хранитель\\
\textbf{kwi2n} [\emph{kwi˩n}] 0x7822: циркулярно\\
\textbf{kwi2t} [\emph{kwi˩t}] 0xB822: герметический\\
\textbf{kwi7k} [\emph{kwi˥k}] 0x3022: карбоксил\\
\textbf{kwi7n} [\emph{kwi˥n}] 0x7022: хинин\\
\textbf{kwi7t} [\emph{kwi˥t}] 0xB022: конькобежец\\
\textbf{kwic} [\emph{kwiʃ}] 0xE822: иностранный\\
\textbf{kwif} [\emph{kwif}] 0xC822: киловатт\\
\textbf{kwih} [\emph{kwiʰ}] 0x4117: каузальный доказательный\\
\textbf{kwik} [\emph{kwik}] 0x2822: квотатив дело\\
\textbf{kwim} [\emph{kwim}] 0x0822: викторина\\
\textbf{kwin} [\emph{kwin}] 0x6822: дающий\\
\textbf{kwip} [\emph{kwip}] 0x4822: вмешиваться\\
\textbf{kwis} [\emph{kwis}] 0x8822: классический\\
\textbf{kwit} [\emph{kwit}] 0xA822: личность\\
\textbf{kwo2n} [\emph{kwo˩n}] 0x7C22: октан\\
\textbf{kwo2t} [\emph{kwo˩t}] 0xBC22: воловий хвост\\
\textbf{kwo7k} [\emph{kwo˥k}] 0x3422: Кока-Кола\\
\textbf{kwo7n} [\emph{kwo˥n}] 0x7422: Корнуолл\\
\textbf{kwo7t} [\emph{kwo˥t}] 0xB422: Киото\\
\textbf{kwoc} [\emph{kwoʃ}] 0xEC22: сквош\\
\textbf{kwof} [\emph{kwof}] 0xCC22: Эквадор\\
\textbf{kwok} [\emph{kwok}] 0x2C22: кобальт\\
\textbf{kwom} [\emph{kwom}] 0x0C22: килограмм\\
\textbf{kwon} [\emph{kwon}] 0x6C22: Словарь\\
\textbf{kwop} [\emph{kwop}] 0x4C22: в гору\\
\textbf{kwos} [\emph{kwos}] 0x8C22: по часовой стрелке\\
\textbf{kwot} [\emph{kwot}] 0xAC22: клавиатура\\
\textbf{kwu7k} [\emph{kwu˥k}] 0x3222: Аккув\\
\textbf{kwuc} [\emph{kwuʃ}] 0xEA22: принуждение\\
\textbf{kwuk} [\emph{kwuk}] 0x2A22: конго\\
\textbf{kwum} [\emph{kwum}] 0x0A22: нематода\\
\textbf{kwun} [\emph{kwun}] 0x6A22: июнь\\
\textbf{kwup} [\emph{kwup}] 0x4A22: высокий уровень\\
\textbf{kwus} [\emph{kwus}] 0x8A22: капсула\\
\textbf{kwut} [\emph{kwut}] 0xAA22: Пальто\\
\textbf{kx67t} [\emph{kxə˥t}] 0xB5E2: аккадское\\
\textbf{kx6c} [\emph{kxəʃ}] 0xEDE2: чудак\\
\textbf{kx6k} [\emph{kxək}] 0x2DE2: тектонический\\
\textbf{kx6n} [\emph{kxən}] 0x6DE2: каштан\\
\textbf{kx6p} [\emph{kxəp}] 0x4DE2: Каспийская\\
\textbf{kx6s} [\emph{kxəs}] 0x8DE2: декартовский\\
\textbf{kx6t} [\emph{kxət}] 0xADE2: горючий\\
\textbf{kxa2k} [\emph{kxä˩k}] 0x39E2: лаг\\
\textbf{kxa2m} [\emph{kxä˩m}] 0x19E2: кремировать\\
\textbf{kxa2n} [\emph{kxä˩n}] 0x79E2: шхуна\\
\textbf{kxa2s} [\emph{kxä˩s}] 0x99E2: рахит\\
\textbf{kxa2t} [\emph{kxä˩t}] 0xB9E2: верхом\\
\textbf{kxa7k} [\emph{kxä˥k}] 0x31E2: иконоборец\\
\textbf{kxa7m} [\emph{kxä˥m}] 0x11E2: клептоман\\
\textbf{kxa7n} [\emph{kxä˥n}] 0x71E2: карабин\\
\textbf{kxa7p} [\emph{kxä˥p}] 0x51E2: взять сумку\\
\textbf{kxa7s} [\emph{kxä˥s}] 0x91E2: заминированный\\
\textbf{kxa7t} [\emph{kxä˥t}] 0xB1E2: движимый\\
\textbf{kxac} [\emph{kxäʃ}] 0xE9E2: восклицание\\
\textbf{kxaf} [\emph{kxäf}] 0xC9E2: акация\\
\textbf{kxah} [\emph{kxäʰ}] 0x4F17: экстатический\\
\textbf{kxak} [\emph{kxäk}] 0x29E2: к дело\\
\textbf{kxam} [\emph{kxäm}] 0x09E2: стержень\\
\textbf{kxan} [\emph{kxän}] 0x69E2: картон\\
\textbf{kxap} [\emph{kxäp}] 0x49E2: поднятый\\
\textbf{kxas} [\emph{kxäs}] 0x89E2: скорость света\\
\textbf{kxat} [\emph{kxät}] 0xA9E2: скандий\\
\textbf{kxe2n} [\emph{kxe˩n}] 0x7BE2: бешено\\
\textbf{kxe2t} [\emph{kxe˩t}] 0xBBE2: квинтет\\
\textbf{kxe7n} [\emph{kxe˥n}] 0x73E2: слабый коленчатый\\
\textbf{kxe7t} [\emph{kxe˥t}] 0xB3E2: снова впадать\\
\textbf{kxec} [\emph{kxeʃ}] 0xEBE2: коктейль\\
\textbf{kxef} [\emph{kxef}] 0xCBE2: коленчатый вал\\
\textbf{kxeh} [\emph{kxeʰ}] 0x5F17: к дело\\
\textbf{kxek} [\emph{kxek}] 0x2BE2: краб\\
\textbf{kxem} [\emph{kxem}] 0x0BE2: Клемент\\
\textbf{kxen} [\emph{kxen}] 0x6BE2: общий Пол\\
\textbf{kxep} [\emph{kxep}] 0x4BE2: холестерин\\
\textbf{kxes} [\emph{kxes}] 0x8BE2: смородина\\
\textbf{kxet} [\emph{kxet}] 0xABE2: эскалатор\\
\textbf{kxi2k} [\emph{kxi˩k}] 0x38E2: воздухоплавание\\
\textbf{kxi2n} [\emph{kxi˩n}] 0x78E2: сорок один\\
\textbf{kxi2s} [\emph{kxi˩s}] 0x98E2: разъездной священник\\
\textbf{kxi2t} [\emph{kxi˩t}] 0xB8E2: подсвечник\\
\textbf{kxi7k} [\emph{kxi˥k}] 0x30E2: стог\\
\textbf{kxi7n} [\emph{kxi˥n}] 0x70E2: приобретение\\
\textbf{kxi7p} [\emph{kxi˥p}] 0x50E2: чипсов\\
\textbf{kxi7s} [\emph{kxi˥s}] 0x90E2: рубец\\
\textbf{kxi7t} [\emph{kxi˥t}] 0xB0E2: удержан\\
\textbf{kxic} [\emph{kxiʃ}] 0xE8E2: огнетушитель\\
\textbf{kxif} [\emph{kxif}] 0xC8E2: айва\\
\textbf{kxih} [\emph{kxiʰ}] 0x4717: невозмутимость\\
\textbf{kxik} [\emph{kxik}] 0x28E2: Народная музыка\\
\textbf{kxim} [\emph{kxim}] 0x08E2: уклончивый\\
\textbf{kxin} [\emph{kxin}] 0x68E2: факториал\\
\textbf{kxip} [\emph{kxip}] 0x48E2: клип доска\\
\textbf{kxis} [\emph{kxis}] 0x88E2: клип Изобразительное искусство\\
\textbf{kxit} [\emph{kxit}] 0xA8E2: ненормальный\\
\textbf{kxo2k} [\emph{kxo˩k}] 0x3CE2: трудоголик\\
\textbf{kxo2n} [\emph{kxo˩n}] 0x7CE2: кетон\\
\textbf{kxo2t} [\emph{kxo˩t}] 0xBCE2: волоске\\
\textbf{kxo7k} [\emph{kxo˥k}] 0x34E2: калейдоскоп\\
\textbf{kxo7n} [\emph{kxo˥n}] 0x74E2: кантонский\\
\textbf{kxo7s} [\emph{kxo˥s}] 0x94E2: лактоза\\
\textbf{kxo7t} [\emph{kxo˥t}] 0xB4E2: заикание\\
\textbf{kxoc} [\emph{kxoʃ}] 0xECE2: петушок\\
\textbf{kxof} [\emph{kxof}] 0xCCE2: гравюра на дереве\\
\textbf{kxoh} [\emph{kxoʰ}] 0x6717: килограмм\\
\textbf{kxok} [\emph{kxok}] 0x2CE2: какао\\
\textbf{kxom} [\emph{kxom}] 0x0CE2: согласный\\
\textbf{kxon} [\emph{kxon}] 0x6CE2: гномический настроение\\
\textbf{kxop} [\emph{kxop}] 0x4CE2: крайняя плоть\\
\textbf{kxos} [\emph{kxos}] 0x8CE2: плотоядный\\
\textbf{kxot} [\emph{kxot}] 0xACE2: криптон\\
\textbf{kxu2t} [\emph{kxu˩t}] 0xBAE2: лицо боль\\
\textbf{kxu7n} [\emph{kxu˥n}] 0x72E2: оптовый торговец\\
\textbf{kxu7t} [\emph{kxu˥t}] 0xB2E2: открытой рукой\\
\textbf{kxuc} [\emph{kxuʃ}] 0xEAE2: судорога\\
\textbf{kxuk} [\emph{kxuk}] 0x2AE2: выпуклый\\
\textbf{kxum} [\emph{kxum}] 0x0AE2: воскликнул\\
\textbf{kxun} [\emph{kxun}] 0x6AE2: рыжеватый\\
\textbf{kxup} [\emph{kxup}] 0x4AE2: скальпель\\
\textbf{kxus} [\emph{kxus}] 0x8AE2: куратор\\
\textbf{kxut} [\emph{kxut}] 0xAAE2: широта\\
\textbf{ky6n} [\emph{kjən}] 0x6D02: кокаин\\
\textbf{ky6s} [\emph{kjəs}] 0x8D02: питомники\\
\textbf{ky6t} [\emph{kjət}] 0xAD02: аксиоматический\\
\textbf{kya2k} [\emph{kjä˩k}] 0x3902: ханство\\
\textbf{kya2n} [\emph{kjä˩n}] 0x7902: коррелированный\\
\textbf{kya2t} [\emph{kjä˩t}] 0xB902: кокарда\\
\textbf{kya7c} [\emph{kjä˥ʃ}] 0xF102: икра\\
\textbf{kya7k} [\emph{kjä˥k}] 0x3102: ручная работа\\
\textbf{kya7m} [\emph{kjä˥m}] 0x1102: переименовать\\
\textbf{kya7n} [\emph{kjä˥n}] 0x7102: канон\\
\textbf{kya7s} [\emph{kjä˥s}] 0x9102: панегирист\\
\textbf{kya7t} [\emph{kjä˥t}] 0xB102: обращенных\\
\textbf{kyac} [\emph{kjäʃ}] 0xE902: глаз\\
\textbf{kyaf} [\emph{kjäf}] 0xC902: спасибо\\
\textbf{kyah} [\emph{kjäʰ}] 0x4817: спасибо\\
\textbf{kyak} [\emph{kjäk}] 0x2902: несколько\\
\textbf{kyam} [\emph{kjäm}] 0x0902: лагерь\\
\textbf{kyan} [\emph{kjän}] 0x6902: как\\
\textbf{kyap} [\emph{kjäp}] 0x4902: ужасный\\
\textbf{kyas} [\emph{kjäs}] 0x8902: лошадь\\
\textbf{kyat} [\emph{kjät}] 0xA902: в самом деле\\
\textbf{kye2n} [\emph{kje˩n}] 0x7B02: концессионер\\
\textbf{kye2t} [\emph{kje˩t}] 0xBB02: секстет\\
\textbf{kye7n} [\emph{kje˥n}] 0x7302: Аккредитующая\\
\textbf{kye7t} [\emph{kje˥t}] 0xB302: х по рейтингу\\
\textbf{kyec} [\emph{kjeʃ}] 0xEB02: рабочее место\\
\textbf{kyeh} [\emph{kjeʰ}] 0x5817: пожалел\\
\textbf{kyek} [\emph{kjek}] 0x2B02: кекс\\
\textbf{kyem} [\emph{kjem}] 0x0B02: кадмий\\
\textbf{kyen} [\emph{kjen}] 0x6B02: несчастность\\
\textbf{kyep} [\emph{kjep}] 0x4B02: сороконожка\\
\textbf{kyes} [\emph{kjes}] 0x8B02: Комментарии\\
\textbf{kyet} [\emph{kjet}] 0xAB02: коллектор\\
\textbf{kyi2n} [\emph{kji˩n}] 0x7802: доктринер\\
\textbf{kyi2s} [\emph{kji˩s}] 0x9802: сексизм\\
\textbf{kyi2t} [\emph{kji˩t}] 0xB802: самописец\\
\textbf{kyi7k} [\emph{kji˥k}] 0x3002: копчик\\
\textbf{kyi7n} [\emph{kji˥n}] 0x7002: ясновидение\\
\textbf{kyi7p} [\emph{kji˥p}] 0x5002: раскол быстро и без хлопот\\
\textbf{kyi7s} [\emph{kji˥s}] 0x9002: кретинизм\\
\textbf{kyi7t} [\emph{kji˥t}] 0xB002: окольный\\
\textbf{kyic} [\emph{kjiʃ}] 0xE802: церковь\\
\textbf{kyif} [\emph{kjif}] 0xC802: аппликативны голос\\
\textbf{kyih} [\emph{kjiʰ}] 0x4017: потому как\\
\textbf{kyik} [\emph{kjik}] 0x2802: потому как\\
\textbf{kyim} [\emph{kjim}] 0x0802: купить\\
\textbf{kyin} [\emph{kjin}] 0x6802: Компания\\
\textbf{kyip} [\emph{kjip}] 0x4802: ящерица\\
\textbf{kyis} [\emph{kjis}] 0x8802: бумага\\
\textbf{kyit} [\emph{kjit}] 0xA802: процитировал\\
\textbf{kyo2t} [\emph{kjo˩t}] 0xBC02: эконометрия\\
\textbf{kyo7n} [\emph{kjo˥n}] 0x7402: комедийная актриса\\
\textbf{kyo7t} [\emph{kjo˥t}] 0xB402: кондор\\
\textbf{kyoc} [\emph{kjoʃ}] 0xEC02: причинная связь\\
\textbf{kyoh} [\emph{kjoʰ}] 0x6017: кило\\
\textbf{kyok} [\emph{kjok}] 0x2C02: огурец\\
\textbf{kyom} [\emph{kjom}] 0x0C02: кило\\
\textbf{kyon} [\emph{kjon}] 0x6C02: эксплуатировать\\
\textbf{kyop} [\emph{kjop}] 0x4C02: крокодил\\
\textbf{kyos} [\emph{kjos}] 0x8C02: коммунист\\
\textbf{kyot} [\emph{kjot}] 0xAC02: аккорд\\
\textbf{kyu7t} [\emph{kju˥t}] 0xB202: мракобесие\\
\textbf{kyuc} [\emph{kjuʃ}] 0xEA02: сухой\\
\textbf{kyuk} [\emph{kjuk}] 0x2A02: шея\\
\textbf{kyum} [\emph{kjum}] 0x0A02: негашеная известь\\
\textbf{kyun} [\emph{kjun}] 0x6A02: угол\\
\textbf{kyup} [\emph{kjup}] 0x4A02: Куба\\
\textbf{kyus} [\emph{kjus}] 0x8A02: пересекать\\
\textbf{kyut} [\emph{kjut}] 0xAA02: компьютер\\
\subsection{ l }
\textbf{l6} [\emph{lə}] 0x357E: агент вызывать\\
\textbf{la} [\emph{lä}] 0x257E: dependent clause, similar in usage to `that', a sentence that is within another sentence\\
\textbf{la7} [\emph{lä˥}] 0x457E: successive clause, for putting several sentences in a sequence, similar to this happened, then that happened\\
\textbf{le} [\emph{le}] 0x2D7E: задерживается императив\\
\textbf{le7} [\emph{le˥}] 0x4D7E: еще - если\\
\textbf{li} [\emph{li}] 0x217E: realis mood, for declaring something, the most common kind of sentence.\\
\textbf{li7} [\emph{li˥}] 0x417E: LY\\
\textbf{lo} [\emph{lo}] 0x317E: волевых настроение\\
\textbf{lu} [\emph{lu}] 0x297E: speculative mood, for expressing speculation, hypothesis, ``could be''\\
\textbf{lw6n} [\emph{lwən}] 0x6D2B: лантан\\
\textbf{lw6s} [\emph{lwəs}] 0x8D2B: долго windedness\\
\textbf{lw6t} [\emph{lwət}] 0xAD2B: Лесото\\
\textbf{lwa2n} [\emph{lwä˩n}] 0x792B: иносказательный\\
\textbf{lwa2t} [\emph{lwä˩t}] 0xB92B: лассо\\
\textbf{lwa7k} [\emph{lwä˥k}] 0x312B: аликвоту\\
\textbf{lwa7n} [\emph{lwä˥n}] 0x712B: кальвинист\\
\textbf{lwa7t} [\emph{lwä˥t}] 0xB12B: лингвист\\
\textbf{lwac} [\emph{lwäʃ}] 0xE92B: осел\\
\textbf{lwaf} [\emph{lwäf}] 0xC92B: вспышка водить машину\\
\textbf{lwah} [\emph{lwäʰ}] 0x495F: perlative case, indicates motion through space. space-context way-case. \\
\textbf{lwak} [\emph{lwäk}] 0x292B: perlative дело\\
\textbf{lwam} [\emph{lwäm}] 0x092B: алфавит\\
\textbf{lwan} [\emph{lwän}] 0x692B: обед\\
\textbf{lwap} [\emph{lwäp}] 0x492B: лютня\\
\textbf{lwas} [\emph{lwäs}] 0x892B: ленивый\\
\textbf{lwat} [\emph{lwät}] 0xA92B: список\\
\textbf{lwe2n} [\emph{lwe˩n}] 0x7B2B: красное Крыло\\
\textbf{lwe7n} [\emph{lwe˥n}] 0x732B: телефонист\\
\textbf{lwec} [\emph{lweʃ}] 0xEB2B: аллели\\
\textbf{lwef} [\emph{lwef}] 0xCB2B: йеллоунайф\\
\textbf{lweh} [\emph{lweʰ}] 0x595F: allative case, indicates motion onto a surface. surface-context destination-case. \\
\textbf{lwek} [\emph{lwek}] 0x2B2B: рычаг\\
\textbf{lwem} [\emph{lwem}] 0x0B2B: палиндром\\
\textbf{lwen} [\emph{lwen}] 0x6B2B: легенда\\
\textbf{lwep} [\emph{lwep}] 0x4B2B: разборчиво\\
\textbf{lwes} [\emph{lwes}] 0x8B2B: лазер\\
\textbf{lwet} [\emph{lwet}] 0xAB2B: вертолет\\
\textbf{lwi2k} [\emph{lwi˩k}] 0x382B: лауриновая\\
\textbf{lwi2n} [\emph{lwi˩n}] 0x782B: альбинос\\
\textbf{lwi2t} [\emph{lwi˩t}] 0xB82B: гробница\\
\textbf{lwi7k} [\emph{lwi˥k}] 0x302B: LY\\
\textbf{lwi7n} [\emph{lwi˥n}] 0x702B: миллионер\\
\textbf{lwi7p} [\emph{lwi˥p}] 0x502B: скользкий\\
\textbf{lwi7s} [\emph{lwi˥s}] 0x902B: Linsey Вулси\\
\textbf{lwi7t} [\emph{lwi˥t}] 0xB02B: податливый\\
\textbf{lwic} [\emph{lwiʃ}] 0xE82B: развод\\
\textbf{lwif} [\emph{lwif}] 0xC82B: местный падеж\\
\textbf{lwih} [\emph{lwiʰ}] 0x415F: delative case, indicates motion away from noun phrase. surface-context source-case. \\
\textbf{lwik} [\emph{lwik}] 0x282B: allative дело\\
\textbf{lwim} [\emph{lwim}] 0x082B: алюминий\\
\textbf{lwin} [\emph{lwin}] 0x682B: сбалансированный\\
\textbf{lwip} [\emph{lwip}] 0x482B: лимфа\\
\textbf{lwis} [\emph{lwis}] 0x882B: щелочной\\
\textbf{lwit} [\emph{lwit}] 0xA82B: литература\\
\textbf{lwo7k} [\emph{lwo˥k}] 0x342B: левое крыло\\
\textbf{lwo7s} [\emph{lwo˥s}] 0x942B: все желающие\\
\textbf{lwo7t} [\emph{lwo˥t}] 0xB42B: телезритель\\
\textbf{lwoc} [\emph{lwoʃ}] 0xEC2B: лаосский\\
\textbf{lwof} [\emph{lwof}] 0xCC2B: Латвия\\
\textbf{lwoh} [\emph{lwoʰ}] 0x615F: locative case, indicates location in space. space-context location-case. \\
\textbf{lwok} [\emph{lwok}] 0x2C2B: лозунг\\
\textbf{lwom} [\emph{lwom}] 0x0C2B: глаукома\\
\textbf{lwon} [\emph{lwon}] 0x6C2B: эволюционный\\
\textbf{lwop} [\emph{lwop}] 0x4C2B: логотипы\\
\textbf{lwos} [\emph{lwos}] 0x8C2B: проложенный\\
\textbf{lwot} [\emph{lwot}] 0xAC2B: хлористый\\
\textbf{lwu2t} [\emph{lwu˩t}] 0xBA2B: впередсмотрящие\\
\textbf{lwu7t} [\emph{lwu˥t}] 0xB22B: пятилетие\\
\textbf{lwuf} [\emph{lwuf}] 0xCA2B: сульфид\\
\textbf{lwuh} [\emph{lwuʰ}] 0x515F: ленивый\\
\textbf{lwuk} [\emph{lwuk}] 0x2A2B: ребро\\
\textbf{lwum} [\emph{lwum}] 0x0A2B: коллоквиализм\\
\textbf{lwun} [\emph{lwun}] 0x6A2B: синие джинсы\\
\textbf{lwus} [\emph{lwus}] 0x8A2B: лесс\\
\textbf{lwut} [\emph{lwut}] 0xAA2B: поток тока\\
\textbf{ly6m} [\emph{ljəm}] 0x0D0B: НМУ\\
\textbf{ly6n} [\emph{ljən}] 0x6D0B: палестинец\\
\textbf{ly6t} [\emph{ljət}] 0xAD0B: назидательная речь соглашение\\
\textbf{lya2c} [\emph{ljä˩ʃ}] 0xF90B: все звезды\\
\textbf{lya2k} [\emph{ljä˩k}] 0x390B: вкладыш\\
\textbf{lya2m} [\emph{ljä˩m}] 0x190B: малаялам\\
\textbf{lya2n} [\emph{ljä˩n}] 0x790B: светоч\\
\textbf{lya2s} [\emph{ljä˩s}] 0x990B: бельевой\\
\textbf{lya2t} [\emph{ljä˩t}] 0xB90B: флагшток\\
\textbf{lya7c} [\emph{ljä˥ʃ}] 0xF10B: индифферентный\\
\textbf{lya7k} [\emph{ljä˥k}] 0x310B: пелагический\\
\textbf{lya7m} [\emph{ljä˥m}] 0x110B: стручковый перец\\
\textbf{lya7n} [\emph{ljä˥n}] 0x710B: ладан\\
\textbf{lya7p} [\emph{ljä˥p}] 0x510B: синий воротник\\
\textbf{lya7s} [\emph{ljä˥s}] 0x910B: Плеяды\\
\textbf{lya7t} [\emph{ljä˥t}] 0xB10B: приличия\\
\textbf{lyac} [\emph{ljäʃ}] 0xE90B: учить\\
\textbf{lyaf} [\emph{ljäf}] 0xC90B: алименты\\
\textbf{lyah} [\emph{ljäʰ}] 0x485F: отчуждаемый владение\\
\textbf{lyak} [\emph{ljäk}] 0x290B: эластичный\\
\textbf{lyam} [\emph{ljäm}] 0x090B: поток\\
\textbf{lyan} [\emph{ljän}] 0x690B: длинный\\
\textbf{lyap} [\emph{ljäp}] 0x490B: злоупотребление\\
\textbf{lyas} [\emph{ljäs}] 0x890B: интерес\\
\textbf{lyat} [\emph{ljät}] 0xA90B: письмо\\
\textbf{lye2t} [\emph{lje˩t}] 0xBB0B: кровельщик\\
\textbf{lye7n} [\emph{lje˥n}] 0x730B: лецитин\\
\textbf{lye7t} [\emph{lje˥t}] 0xB30B: макулатура\\
\textbf{lyec} [\emph{ljeʃ}] 0xEB0B: линейный\\
\textbf{lyeh} [\emph{ljeʰ}] 0x585F: atelic aspect, for actions that have no particular end point. infinte or unlimited.\\
\textbf{lyek} [\emph{ljek}] 0x2B0B: латинский\\
\textbf{lyem} [\emph{ljem}] 0x0B0B: оперять\\
\textbf{lyen} [\emph{ljen}] 0x6B0B: белье\\
\textbf{lyep} [\emph{ljep}] 0x4B0B: Либерия\\
\textbf{lyes} [\emph{ljes}] 0x8B0B: парикмахер\\
\textbf{lyet} [\emph{ljet}] 0xAB0B: объектив\\
\textbf{lyi2k} [\emph{lji˩k}] 0x380B: Ленинград\\
\textbf{lyi2n} [\emph{lji˩n}] 0x780B: литания\\
\textbf{lyi2s} [\emph{lji˩s}] 0x980B: олигоцен\\
\textbf{lyi2t} [\emph{lji˩t}] 0xB80B: лира\\
\textbf{lyi7c} [\emph{lji˥ʃ}] 0xF00B: алкил\\
\textbf{lyi7f} [\emph{lji˥f}] 0xD00B: лилия трусливый\\
\textbf{lyi7k} [\emph{lji˥k}] 0x300B: стекольщик\\
\textbf{lyi7m} [\emph{lji˥m}] 0x100B: подвздошная кость\\
\textbf{lyi7n} [\emph{lji˥n}] 0x700B: Колыбельная\\
\textbf{lyi7s} [\emph{lji˥s}] 0x900B: рысь\\
\textbf{lyi7t} [\emph{lji˥t}] 0xB00B: льняное семя\\
\textbf{lyic} [\emph{ljiʃ}] 0xE80B: пожалуйста\\
\textbf{lyif} [\emph{ljif}] 0xC80B: Литва\\
\textbf{lyih} [\emph{ljiʰ}] 0x405F: интерес\\
\textbf{lyik} [\emph{ljik}] 0x280B: липкий\\
\textbf{lyim} [\emph{ljim}] 0x080B: элита\\
\textbf{lyin} [\emph{ljin}] 0x680B: электронный\\
\textbf{lyip} [\emph{ljip}] 0x480B: клиническая\\
\textbf{lyis} [\emph{ljis}] 0x880B: лесоводство\\
\textbf{lyit} [\emph{ljit}] 0xA80B: записывать\\
\textbf{lyo2n} [\emph{ljo˩n}] 0x7C0B: поражения\\
\textbf{lyo2t} [\emph{ljo˩t}] 0xBC0B: многогранник\\
\textbf{lyo7n} [\emph{ljo˥n}] 0x740B: клонический\\
\textbf{lyo7s} [\emph{ljo˥s}] 0x940B: солипсизм\\
\textbf{lyo7t} [\emph{ljo˥t}] 0xB40B: вольтметр\\
\textbf{lyoc} [\emph{ljoʃ}] 0xEC0B: лиственница\\
\textbf{lyok} [\emph{ljok}] 0x2C0B: Цель вызывать\\
\textbf{lyom} [\emph{ljom}] 0x0C0B: Molly охранник\\
\textbf{lyon} [\emph{ljon}] 0x6C0B: язва\\
\textbf{lyop} [\emph{ljop}] 0x4C0B: Олимпийские игры\\
\textbf{lyos} [\emph{ljos}] 0x8C0B: плеоназм\\
\textbf{lyot} [\emph{ljot}] 0xAC0B: авто\\
\textbf{lyu2t} [\emph{lju˩t}] 0xBA0B: уличный фонарь\\
\textbf{lyu7k} [\emph{lju˥k}] 0x320B: урал Алтайская\\
\textbf{lyu7n} [\emph{lju˥n}] 0x720B: выпускники\\
\textbf{lyu7t} [\emph{lju˥t}] 0xB20B: кормление грудью\\
\textbf{lyuc} [\emph{ljuʃ}] 0xEA0B: тобогган-одиночка\\
\textbf{lyuk} [\emph{ljuk}] 0x2A0B: последующий пункт\\
\textbf{lyum} [\emph{ljum}] 0x0A0B: тулий атом\\
\textbf{lyun} [\emph{ljun}] 0x6A0B: ион\\
\textbf{lyup} [\emph{ljup}] 0x4A0B: шлюп\\
\textbf{lyus} [\emph{ljus}] 0x8A0B: целлюлоза\\
\textbf{lyut} [\emph{ljut}] 0xAA0B: бороться\\
\subsection{ m }
\textbf{m6} [\emph{mə}] 0x343E: метр\\
\textbf{ma} [\emph{mä}] 0x243E: место назначения дело\\
\textbf{ma7} [\emph{mä˥}] 0x443E: количество из вещество\\
\textbf{me} [\emph{me}] 0x2C3E: imperfective aspect, an incomplete action in process.\\
\textbf{mi} [\emph{mi}] 0x203E: меня\\
\textbf{mi7} [\emph{mi˥}] 0x403E: е\\
\textbf{mo} [\emph{mo}] 0x303E: включительно человек\\
\textbf{mr67t} [\emph{mrə˥t}] 0xB5C1: миокард\\
\textbf{mr6c} [\emph{mrəʃ}] 0xEDC1: Махараштра\\
\textbf{mr6k} [\emph{mrək}] 0x2DC1: Америка\\
\textbf{mr6m} [\emph{mrəm}] 0x0DC1: мормон\\
\textbf{mr6n} [\emph{mrən}] 0x6DC1: меридиан\\
\textbf{mr6p} [\emph{mrəp}] 0x4DC1: микрочип\\
\textbf{mr6s} [\emph{mrəs}] 0x8DC1: подлокотник\\
\textbf{mr6t} [\emph{mrət}] 0xADC1: синонимичный\\
\textbf{mra2c} [\emph{mrä˩ʃ}] 0xF9C1: Микронезия\\
\textbf{mra2f} [\emph{mrä˩f}] 0xD9C1: маратхи\\
\textbf{mra2k} [\emph{mrä˩k}] 0x39C1: морг\\
\textbf{mra2m} [\emph{mrä˩m}] 0x19C1: подобный мрамору\\
\textbf{mra2n} [\emph{mrä˩n}] 0x79C1: адмирал\\
\textbf{mra2p} [\emph{mrä˩p}] 0x59C1: амарант\\
\textbf{mra2s} [\emph{mrä˩s}] 0x99C1: амбра\\
\textbf{mra2t} [\emph{mrä˩t}] 0xB9C1: Амстердам\\
\textbf{mra7c} [\emph{mrä˥ʃ}] 0xF1C1: вымачивать\\
\textbf{mra7f} [\emph{mrä˥f}] 0xD1C1: тамаринд\\
\textbf{mra7k} [\emph{mrä˥k}] 0x31C1: Мартиника\\
\textbf{mra7m} [\emph{mrä˥m}] 0x11C1: мелодрама\\
\textbf{mra7n} [\emph{mrä˥n}] 0x71C1: мама в закон\\
\textbf{mra7p} [\emph{mrä˥p}] 0x51C1: ламантин\\
\textbf{mra7s} [\emph{mrä˥s}] 0x91C1: хлопчатобумажная ткань в полоску\\
\textbf{mra7t} [\emph{mrä˥t}] 0xB1C1: Мадрид\\
\textbf{mrac} [\emph{mräʃ}] 0xE9C1: вторник\\
\textbf{mraf} [\emph{mräf}] 0xC9C1: марафон\\
\textbf{mrah} [\emph{mräʰ}] 0x4E0F: Понимаю\\
\textbf{mrak} [\emph{mräk}] 0x29C1: копье\\
\textbf{mram} [\emph{mräm}] 0x09C1: минеральная\\
\textbf{mran} [\emph{mrän}] 0x69C1: кукуруза\\
\textbf{mrap} [\emph{mräp}] 0x49C1: мастер\\
\textbf{mras} [\emph{mräs}] 0x89C1: зонтик\\
\textbf{mrat} [\emph{mrät}] 0xA9C1: смешанный\\
\textbf{mre2k} [\emph{mre˩k}] 0x3BC1: метрический\\
\textbf{mre2n} [\emph{mre˩n}] 0x7BC1: морена\\
\textbf{mre2s} [\emph{mre˩s}] 0x9BC1: памятные предметы\\
\textbf{mre2t} [\emph{mre˩t}] 0xBBC1: погребать\\
\textbf{mre7k} [\emph{mre˥k}] 0x33C1: ун-американских\\
\textbf{mre7n} [\emph{mre˥n}] 0x73C1: пленочный\\
\textbf{mre7s} [\emph{mre˥s}] 0x93C1: матрицы\\
\textbf{mre7t} [\emph{mre˥t}] 0xB3C1: менуэт\\
\textbf{mrec} [\emph{mreʃ}] 0xEBC1: воробей\\
\textbf{mref} [\emph{mref}] 0xCBC1: монорельс\\
\textbf{mreh} [\emph{mreʰ}] 0x5E0F: momentane aspect, for referring to momentary actions, like the unfurling of flower petals in the morning\\
\textbf{mrek} [\emph{mrek}] 0x2BC1: поражать\\
\textbf{mrem} [\emph{mrem}] 0x0BC1: мэйнфреймов\\
\textbf{mren} [\emph{mren}] 0x6BC1: памятный\\
\textbf{mrep} [\emph{mrep}] 0x4BC1: ампер\\
\textbf{mres} [\emph{mres}] 0x8BC1: сжатие\\
\textbf{mret} [\emph{mret}] 0xABC1: министерство\\
\textbf{mri2k} [\emph{mri˩k}] 0x38C1: маркиз\\
\textbf{mri2m} [\emph{mri˩m}] 0x18C1: матереубийство\\
\textbf{mri2n} [\emph{mri˩n}] 0x78C1: материнство\\
\textbf{mri2p} [\emph{mri˩p}] 0x58C1: мизантропия\\
\textbf{mri2s} [\emph{mri˩s}] 0x98C1: , себя\\
\textbf{mri2t} [\emph{mri˩t}] 0xB8C1: сталкиваются\\
\textbf{mri7c} [\emph{mri˥ʃ}] 0xF0C1: пава\\
\textbf{mri7f} [\emph{mri˥f}] 0xD0C1: морфий\\
\textbf{mri7k} [\emph{mri˥k}] 0x30C1: мелодраматический\\
\textbf{mri7m} [\emph{mri˥m}] 0x10C1: микрометр\\
\textbf{mri7n} [\emph{mri˥n}] 0x70C1: марсианин\\
\textbf{mri7p} [\emph{mri˥p}] 0x50C1: мизантроп\\
\textbf{mri7s} [\emph{mri˥s}] 0x90C1: материализм\\
\textbf{mri7t} [\emph{mri˥t}] 0xB0C1: Мэри\\
\textbf{mric} [\emph{mriʃ}] 0xE8C1: мышца\\
\textbf{mrif} [\emph{mrif}] 0xC8C1: микрофон\\
\textbf{mrih} [\emph{mriʰ}] 0x460F: вежливый регистр\\
\textbf{mrik} [\emph{mrik}] 0x28C1: механический\\
\textbf{mrim} [\emph{mrim}] 0x08C1: муравей\\
\textbf{mrin} [\emph{mrin}] 0x68C1: мука\\
\textbf{mrip} [\emph{mrip}] 0x48C1: метр\\
\textbf{mris} [\emph{mris}] 0x88C1: демократия\\
\textbf{mrit} [\emph{mrit}] 0xA8C1: предположительный настроение\\
\textbf{mro2k} [\emph{mro˩k}] 0x3CC1: микрон\\
\textbf{mro2n} [\emph{mro˩n}] 0x7CC1: бур\\
\textbf{mro2s} [\emph{mro˩s}] 0x9CC1: геморрой\\
\textbf{mro2t} [\emph{mro˩t}] 0xBCC1: таращить глаза\\
\textbf{mro7c} [\emph{mro˥ʃ}] 0xF4C1: аменорея\\
\textbf{mro7k} [\emph{mro˥k}] 0x34C1: амхарский\\
\textbf{mro7n} [\emph{mro˥n}] 0x74C1: умирающий\\
\textbf{mro7s} [\emph{mro˥s}] 0x94C1: микропроцессор\\
\textbf{mro7t} [\emph{mro˥t}] 0xB4C1: модный\\
\textbf{mroc} [\emph{mroʃ}] 0xECC1: трение\\
\textbf{mrof} [\emph{mrof}] 0xCCC1: молярный\\
\textbf{mroh} [\emph{mroʰ}] 0x660F: моральный\\
\textbf{mrok} [\emph{mrok}] 0x2CC1: комар\\
\textbf{mrom} [\emph{mrom}] 0x0CC1: морфема\\
\textbf{mron} [\emph{mron}] 0x6CC1: фургон\\
\textbf{mrop} [\emph{mrop}] 0x4CC1: СВЧ\\
\textbf{mros} [\emph{mros}] 0x8CC1: монитор\\
\textbf{mrot} [\emph{mrot}] 0xACC1: двигатель\\
\textbf{mru2t} [\emph{mru˩t}] 0xBAC1: миниатюрами\\
\textbf{mru7n} [\emph{mru˥n}] 0x72C1: кормушка\\
\textbf{mru7s} [\emph{mru˥s}] 0x92C1: урал\\
\textbf{mru7t} [\emph{mru˥t}] 0xB2C1: плохим советом\\
\textbf{mruc} [\emph{mruʃ}] 0xEAC1: гусеница\\
\textbf{mruh} [\emph{mruʰ}] 0x560F: дурак\\
\textbf{mruk} [\emph{mruk}] 0x2AC1: улучшать\\
\textbf{mrum} [\emph{mrum}] 0x0AC1: компьютер программа\\
\textbf{mrun} [\emph{mrun}] 0x6AC1: плотник\\
\textbf{mrup} [\emph{mrup}] 0x4AC1: микроб\\
\textbf{mrus} [\emph{mrus}] 0x8AC1: Маврикий\\
\textbf{mrut} [\emph{mrut}] 0xAAC1: валовой\\
\textbf{mu} [\emph{mu}] 0x283E: commissive mood, for making promises, commitments\\
\textbf{mw6n} [\emph{mwən}] 0x6D21: Македония\\
\textbf{mw6t} [\emph{mwət}] 0xAD21: от двери до двери\\
\textbf{mwa2n} [\emph{mwä˩n}] 0x7921: научная статья\\
\textbf{mwa2t} [\emph{mwä˩t}] 0xB921: выбритый\\
\textbf{mwa7k} [\emph{mwä˥k}] 0x3121: амфора\\
\textbf{mwa7n} [\emph{mwä˥n}] 0x7121: паутина\\
\textbf{mwa7s} [\emph{mwä˥s}] 0x9121: медиана\\
\textbf{mwa7t} [\emph{mwä˥t}] 0xB121: любовный\\
\textbf{mwac} [\emph{mwäʃ}] 0xE921: доволен\\
\textbf{mwaf} [\emph{mwäf}] 0xC921: вредоносных программ\\
\textbf{mwah} [\emph{mwäʰ}] 0x490F: comitative case, indicates the noun phrase is accompanying the other actors of the sentence. social-context way-case.\\
\textbf{mwak} [\emph{mwäk}] 0x2921: куча\\
\textbf{mwam} [\emph{mwäm}] 0x0921: гимназия\\
\textbf{mwan} [\emph{mwän}] 0x6921: смерть\\
\textbf{mwap} [\emph{mwäp}] 0x4921: комикс\\
\textbf{mwas} [\emph{mwäs}] 0x8921: месяц\\
\textbf{mwat} [\emph{mwät}] 0xA921: вкус\\
\textbf{mwe2n} [\emph{mwe˩n}] 0x7B21: нимфомания\\
\textbf{mwe2t} [\emph{mwe˩t}] 0xBB21: тире\\
\textbf{mwec} [\emph{mweʃ}] 0xEB21: однозначно\\
\textbf{mweh} [\emph{mweʰ}] 0x590F: заканчивающий аспект\\
\textbf{mwek} [\emph{mwek}] 0x2B21: эмираты\\
\textbf{mwem} [\emph{mwem}] 0x0B21: забери меня\\
\textbf{mwen} [\emph{mwen}] 0x6B21: мечтательность\\
\textbf{mwep} [\emph{mwep}] 0x4B21: лживый\\
\textbf{mwes} [\emph{mwes}] 0x8B21: метан\\
\textbf{mwet} [\emph{mwet}] 0xAB21: кладбище\\
\textbf{mwi2k} [\emph{mwi˩k}] 0x3821: мнемоника\\
\textbf{mwi2n} [\emph{mwi˩n}] 0x7821: пескарь\\
\textbf{mwi2s} [\emph{mwi˩s}] 0x9821: монооксид\\
\textbf{mwi2t} [\emph{mwi˩t}] 0xB821: сторонник строгой дисциплины\\
\textbf{mwi7k} [\emph{mwi˥k}] 0x3021: амниотический\\
\textbf{mwi7m} [\emph{mwi˥m}] 0x1021: мим\\
\textbf{mwi7n} [\emph{mwi˥n}] 0x7021: мандолина\\
\textbf{mwi7s} [\emph{mwi˥s}] 0x9021: марксист\\
\textbf{mwi7t} [\emph{mwi˥t}] 0xB021: пятьдесят два\\
\textbf{mwic} [\emph{mwiʃ}] 0xE821: пчела\\
\textbf{mwif} [\emph{mwif}] 0xC821: Малави\\
\textbf{mwih} [\emph{mwiʰ}] 0x410F: permissive mood, for granting permission to do something, potential response to admonitive mood.\\
\textbf{mwik} [\emph{mwik}] 0x2821: магнитный\\
\textbf{mwim} [\emph{mwim}] 0x0821: метонимия\\
\textbf{mwin} [\emph{mwin}] 0x6821: миллиона\\
\textbf{mwip} [\emph{mwip}] 0x4821: самозванец\\
\textbf{mwis} [\emph{mwis}] 0x8821: умственно\\
\textbf{mwit} [\emph{mwit}] 0xA821: commissive настроение\\
\textbf{mwo2t} [\emph{mwo˩t}] 0xBC21: мальтоза\\
\textbf{mwoc} [\emph{mwoʃ}] 0xEC21: микробиология\\
\textbf{mwof} [\emph{mwof}] 0xCC21: Молдова\\
\textbf{mwoh} [\emph{mwoʰ}] 0x610F: меланхолия\\
\textbf{mwok} [\emph{mwok}] 0x2C21: Монголия\\
\textbf{mwom} [\emph{mwom}] 0x0C21: монизм\\
\textbf{mwon} [\emph{mwon}] 0x6C21: мандарин\\
\textbf{mwop} [\emph{mwop}] 0x4C21: Моавитянин\\
\textbf{mwos} [\emph{mwos}] 0x8C21: симбиоз\\
\textbf{mwot} [\emph{mwot}] 0xAC21: помидор\\
\textbf{mwu7n} [\emph{mwu˥n}] 0x7221: муэдзин\\
\textbf{mwuc} [\emph{mwuʃ}] 0xEA21: полоскание для рта\\
\textbf{mwuh} [\emph{mwuʰ}] 0x510F: мычание\\
\textbf{mwuk} [\emph{mwuk}] 0x2A21: норка\\
\textbf{mwum} [\emph{mwum}] 0x0A21: мычание\\
\textbf{mwun} [\emph{mwun}] 0x6A21: мода\\
\textbf{mwus} [\emph{mwus}] 0x8A21: мейоз\\
\textbf{mwut} [\emph{mwut}] 0xAA21: сова\\
\textbf{my6k} [\emph{mjək}] 0x2D01: симметричный\\
\textbf{my6m} [\emph{mjəm}] 0x0D01: Мьянма\\
\textbf{my6n} [\emph{mjən}] 0x6D01: Армения\\
\textbf{my6s} [\emph{mjəs}] 0x8D01: муза\\
\textbf{my6t} [\emph{mjət}] 0xAD01: рамадан\\
\textbf{mya2k} [\emph{mjä˩k}] 0x3901: цели\\
\textbf{mya2m} [\emph{mjä˩m}] 0x1901: плоть муха\\
\textbf{mya2n} [\emph{mjä˩n}] 0x7901: маммона\\
\textbf{mya2s} [\emph{mjä˩s}] 0x9901: массажист\\
\textbf{mya2t} [\emph{mjä˩t}] 0xB901: солод\\
\textbf{mya7c} [\emph{mjä˥ʃ}] 0xF101: маршалльский\\
\textbf{mya7f} [\emph{mjä˥f}] 0xD101: майя\\
\textbf{mya7h} [\emph{mjä˥ʰ}] 0x880F: мяукать\\
\textbf{mya7k} [\emph{mjä˥k}] 0x3101: Малакка\\
\textbf{mya7m} [\emph{mjä˥m}] 0x1101: человек людоед\\
\textbf{mya7n} [\emph{mjä˥n}] 0x7101: унижать\\
\textbf{mya7s} [\emph{mjä˥s}] 0x9101: монада\\
\textbf{mya7t} [\emph{mjä˥t}] 0xB101: амфитеатр\\
\textbf{myac} [\emph{mjäʃ}] 0xE901: улыбка\\
\textbf{myaf} [\emph{mjäf}] 0xC901: Мальдивы\\
\textbf{myah} [\emph{mjäʰ}] 0x480F: temporal case, indicates a specific time. time-context location-case. \\
\textbf{myak} [\emph{mjäk}] 0x2901: Март\\
\textbf{myam} [\emph{mjäm}] 0x0901: любитель\\
\textbf{myan} [\emph{mjän}] 0x6901: мой\\
\textbf{myap} [\emph{mjäp}] 0x4901: сделка\\
\textbf{myas} [\emph{mjäs}] 0x8901: маска\\
\textbf{myat} [\emph{mjät}] 0xA901: запомнить\\
\textbf{mye2t} [\emph{mje˩t}] 0xBB01: армагеддон\\
\textbf{mye7n} [\emph{mje˥n}] 0x7301: замуровывать\\
\textbf{mye7t} [\emph{mje˥t}] 0xB301: Пятикнижие\\
\textbf{myec} [\emph{mjeʃ}] 0xEB01: Малайзия\\
\textbf{myeh} [\emph{mjeʰ}] 0x580F: ампер\\
\textbf{myek} [\emph{mjek}] 0x2B01: термодинамический\\
\textbf{myem} [\emph{mjem}] 0x0B01: миллиметр\\
\textbf{myen} [\emph{mjen}] 0x6B01: по аналогии\\
\textbf{myep} [\emph{mjep}] 0x4B01: Намибия\\
\textbf{myes} [\emph{mjes}] 0x8B01: майонез\\
\textbf{myet} [\emph{mjet}] 0xAB01: мелодия\\
\textbf{myi2c} [\emph{mji˩ʃ}] 0xF801: иммунология\\
\textbf{myi2k} [\emph{mji˩k}] 0x3801: асимптотический\\
\textbf{myi2n} [\emph{mji˩n}] 0x7801: сажа\\
\textbf{myi2s} [\emph{mji˩s}] 0x9801: медиально\\
\textbf{myi2t} [\emph{mji˩t}] 0xB801: менструация\\
\textbf{myi7c} [\emph{mji˥ʃ}] 0xF001: метилен\\
\textbf{myi7h} [\emph{mji˥ʰ}] 0x800F: я\\
\textbf{myi7k} [\emph{mji˥k}] 0x3001: середина июля\\
\textbf{myi7m} [\emph{mji˥m}] 0x1001: эмир\\
\textbf{myi7n} [\emph{mji˥n}] 0x7001: неправильное название\\
\textbf{myi7s} [\emph{mji˥s}] 0x9001: амид\\
\textbf{myi7t} [\emph{mji˥t}] 0xB001: лунатик\\
\textbf{myic} [\emph{mjiʃ}] 0xE801: дождь\\
\textbf{myif} [\emph{mjif}] 0xC801: рыбак\\
\textbf{myih} [\emph{mjiʰ}] 0x400F: красивая\\
\textbf{myik} [\emph{mjik}] 0x2801: машина\\
\textbf{myim} [\emph{mjim}] 0x0801: иммунный\\
\textbf{myin} [\emph{mjin}] 0x6801: команда\\
\textbf{myip} [\emph{mjip}] 0x4801: милиция\\
\textbf{myis} [\emph{mjis}] 0x8801: молекулярная\\
\textbf{myit} [\emph{mjit}] 0xA801: зима\\
\textbf{myo2n} [\emph{mjo˩n}] 0x7C01: ламинарный\\
\textbf{myo2t} [\emph{mjo˩t}] 0xBC01: Человек о войне\\
\textbf{myo7n} [\emph{mjo˥n}] 0x7401: симония\\
\textbf{myo7t} [\emph{mjo˥t}] 0xB401: имитатор\\
\textbf{myoc} [\emph{mjoʃ}] 0xEC01: модуль\\
\textbf{myof} [\emph{mjof}] 0xCC01: дезинформировать\\
\textbf{myoh} [\emph{mjoʰ}] 0x600F: средний голос\\
\textbf{myok} [\emph{mjok}] 0x2C01: мозаика\\
\textbf{myom} [\emph{mjom}] 0x0C01: молибден\\
\textbf{myon} [\emph{mjon}] 0x6C01: пневмония\\
\textbf{myop} [\emph{mjop}] 0x4C01: храмы\\
\textbf{myos} [\emph{mjos}] 0x8C01: метеор\\
\textbf{myot} [\emph{mjot}] 0xAC01: печально известный\\
\textbf{myu2n} [\emph{mju˩n}] 0x7A01: маниакальный\\
\textbf{myu7t} [\emph{mju˥t}] 0xB201: курильщик\\
\textbf{myuc} [\emph{mjuʃ}] 0xEA01: эмульсия\\
\textbf{myuf} [\emph{mjuf}] 0xCA01: муфтий\\
\textbf{myuh} [\emph{mjuʰ}] 0x500F: юмор\\
\textbf{myuk} [\emph{mjuk}] 0x2A01: дурак\\
\textbf{myum} [\emph{mjum}] 0x0A01: мультимедиа\\
\textbf{myun} [\emph{mjun}] 0x6A01: важность\\
\textbf{myup} [\emph{mjup}] 0x4A01: саженцы\\
\textbf{myus} [\emph{mjus}] 0x8A01: мышь\\
\textbf{myut} [\emph{mjut}] 0xAA01: моральный\\
\subsection{ n }
\textbf{n6} [\emph{nə}] 0x34DE: сенсорный доказательный настроение\\
\textbf{na} [\emph{nä}] 0x24DE: именительный дело\\
\textbf{na2} [\emph{nä˩}] 0x64DE: нафтол\\
\textbf{ne} [\emph{ne}] 0x2CDE: мужской Пол\\
\textbf{ne2} [\emph{ne˩}] 0x6CDE: энный мощность\\
\textbf{ni} [\emph{ni}] 0x20DE: инфинитив\\
\textbf{ni2} [\emph{ni˩}] 0x60DE: безличный глагол\\
\textbf{ni7} [\emph{ni˥}] 0x40DE: соединительный частица\\
\textbf{no} [\emph{no}] 0x30DE: негативный квантификатором\\
\textbf{no7} [\emph{no˥}] 0x50DE: О.Б.\\
\textbf{nr67n} [\emph{nrə˥n}] 0x75C6: одновременно\\
\textbf{nr6c} [\emph{nrəʃ}] 0xEDC6: Нигерия\\
\textbf{nr6f} [\emph{nrəf}] 0xCDC6: Назарет\\
\textbf{nr6h} [\emph{nrəʰ}] 0x6E37: пациент вызывать\\
\textbf{nr6k} [\emph{nrək}] 0x2DC6: Никарагуа\\
\textbf{nr6m} [\emph{nrəm}] 0x0DC6: патентованное средство\\
\textbf{nr6n} [\emph{nrən}] 0x6DC6: нейрон\\
\textbf{nr6p} [\emph{nrəp}] 0x4DC6: ономатопея\\
\textbf{nr6s} [\emph{nrəs}] 0x8DC6: антропоморфизм\\
\textbf{nr6t} [\emph{nrət}] 0xADC6: индийский\\
\textbf{nra2c} [\emph{nrä˩ʃ}] 0xF9C6: мучение нерва\\
\textbf{nra2f} [\emph{nrä˩f}] 0xD9C6: антропософия\\
\textbf{nra2k} [\emph{nrä˩k}] 0x39C6: антарктический\\
\textbf{nra2m} [\emph{nrä˩m}] 0x19C6: неграмотность\\
\textbf{nra2n} [\emph{nrä˩n}] 0x79C6: магнат\\
\textbf{nra2s} [\emph{nrä˩s}] 0x99C6: анналы\\
\textbf{nra2t} [\emph{nrä˩t}] 0xB9C6: консерватория\\
\textbf{nra7c} [\emph{nrä˥ʃ}] 0xF1C6: рыться\\
\textbf{nra7k} [\emph{nrä˥k}] 0x31C6: дочь в закон\\
\textbf{nra7m} [\emph{nrä˥m}] 0x11C6: панорамный\\
\textbf{nra7n} [\emph{nrä˥n}] 0x71C6: неописуемый\\
\textbf{nra7s} [\emph{nrä˥s}] 0x91C6: наковальня\\
\textbf{nra7t} [\emph{nrä˥t}] 0xB1C6: Андрей\\
\textbf{nrac} [\emph{nräʃ}] 0xE9C6: безвредный\\
\textbf{nraf} [\emph{nräf}] 0xC9C6: налогоплательщик\\
\textbf{nrah} [\emph{nräʰ}] 0x4E37: тепло\\
\textbf{nrak} [\emph{nräk}] 0x29C6: континент\\
\textbf{nram} [\emph{nräm}] 0x09C6: аморальный\\
\textbf{nran} [\emph{nrän}] 0x69C6: редкий\\
\textbf{nrap} [\emph{nräp}] 0x49C6: экономка\\
\textbf{nras} [\emph{nräs}] 0x89C6: медсестра\\
\textbf{nrat} [\emph{nrät}] 0xA9C6: невеста\\
\textbf{nre2k} [\emph{nre˩k}] 0x3BC6: технократия\\
\textbf{nre2n} [\emph{nre˩n}] 0x7BC6: грозовая туча\\
\textbf{nre2s} [\emph{nre˩s}] 0x9BC6: негритянка\\
\textbf{nre2t} [\emph{nre˩t}] 0xBBC6: Волшебная страна\\
\textbf{nre7n} [\emph{nre˥n}] 0x73C6: Золушка\\
\textbf{nre7s} [\emph{nre˥s}] 0x93C6: лестничная площадка\\
\textbf{nre7t} [\emph{nre˥t}] 0xB3C6: полемизировать\\
\textbf{nrec} [\emph{nreʃ}] 0xEBC6: сенсорный\\
\textbf{nref} [\emph{nref}] 0xCBC6: реинкарнация\\
\textbf{nreh} [\emph{nreʰ}] 0x5E37: энергия\\
\textbf{nrek} [\emph{nrek}] 0x2BC6: работник\\
\textbf{nrem} [\emph{nrem}] 0x0BC6: нитрат\\
\textbf{nren} [\emph{nren}] 0x6BC6: вдохновение\\
\textbf{nrep} [\emph{nrep}] 0x4BC6: северо-Запад\\
\textbf{nres} [\emph{nres}] 0x8BC6: целостность\\
\textbf{nret} [\emph{nret}] 0xABC6: яростно\\
\textbf{nri2c} [\emph{nri˩ʃ}] 0xF8C6: приближается к\\
\textbf{nri2k} [\emph{nri˩k}] 0x38C6: цилиндрический\\
\textbf{nri2m} [\emph{nri˩m}] 0x18C6: анаграмма\\
\textbf{nri2n} [\emph{nri˩n}] 0x78C6: граждане\\
\textbf{nri2p} [\emph{nri˩p}] 0x58C6: невропатический\\
\textbf{nri2s} [\emph{nri˩s}] 0x98C6: нарцисс\\
\textbf{nri2t} [\emph{nri˩t}] 0xB8C6: некоммерческий\\
\textbf{nri7c} [\emph{nri˥ʃ}] 0xF0C6: наплыв\\
\textbf{nri7h} [\emph{nri˥ʰ}] 0x8637: неодушевленный Пол\\
\textbf{nri7k} [\emph{nri˥k}] 0x30C6: желудочек\\
\textbf{nri7m} [\emph{nri˥m}] 0x10C6: очный\\
\textbf{nri7n} [\emph{nri˥n}] 0x70C6: дочери\\
\textbf{nri7p} [\emph{nri˥p}] 0x50C6: иннервация\\
\textbf{nri7s} [\emph{nri˥s}] 0x90C6: неинтересный\\
\textbf{nri7t} [\emph{nri˥t}] 0xB0C6: рыцарский\\
\textbf{nric} [\emph{nriʃ}] 0xE8C6: поколение\\
\textbf{nrif} [\emph{nrif}] 0xC8C6: клев\\
\textbf{nrih} [\emph{nriʰ}] 0x4637: антропный Пол\\
\textbf{nrik} [\emph{nrik}] 0x28C6: определенный статья\\
\textbf{nrim} [\emph{nrim}] 0x08C6: стипендия\\
\textbf{nrin} [\emph{nrin}] 0x68C6: невежественный\\
\textbf{nrip} [\emph{nrip}] 0x48C6: шкаф\\
\textbf{nris} [\emph{nris}] 0x88C6: невинный\\
\textbf{nrit} [\emph{nrit}] 0xA8C6: энергия\\
\textbf{nro2k} [\emph{nro˩k}] 0x3CC6: неврология\\
\textbf{nro2m} [\emph{nro˩m}] 0x1CC6: интерферометр\\
\textbf{nro2n} [\emph{nro˩n}] 0x7CC6: анаконда\\
\textbf{nro2s} [\emph{nro˩s}] 0x9CC6: некроз\\
\textbf{nro2t} [\emph{nro˩t}] 0xBCC6: контральто\\
\textbf{nro7k} [\emph{nro˥k}] 0x34C6: анаболический\\
\textbf{nro7m} [\emph{nro˥m}] 0x14C6: монохромный\\
\textbf{nro7n} [\emph{nro˥n}] 0x74C6: концентрический\\
\textbf{nro7s} [\emph{nro˥s}] 0x94C6: антимонопольный\\
\textbf{nro7t} [\emph{nro˥t}] 0xB4C6: взревел\\
\textbf{nroc} [\emph{nroʃ}] 0xECC6: оргия\\
\textbf{nrof} [\emph{nrof}] 0xCCC6: Норвегия\\
\textbf{nroh} [\emph{nroʰ}] 0x6637: гнев\\
\textbf{nrok} [\emph{nrok}] 0x2CC6: азот\\
\textbf{nrom} [\emph{nrom}] 0x0CC6: актер роль\\
\textbf{nron} [\emph{nron}] 0x6CC6: водород\\
\textbf{nrop} [\emph{nrop}] 0x4CC6: начисление заработной платы\\
\textbf{nros} [\emph{nros}] 0x8CC6: контрастные\\
\textbf{nrot} [\emph{nrot}] 0xACC6: Международный\\
\textbf{nru2n} [\emph{nru˩n}] 0x7AC6: первоисточник\\
\textbf{nru2s} [\emph{nru˩s}] 0x9AC6: неврит\\
\textbf{nru2t} [\emph{nru˩t}] 0xBAC6: невротический\\
\textbf{nru7n} [\emph{nru˥n}] 0x72C6: интонировать\\
\textbf{nru7s} [\emph{nru˥s}] 0x92C6: гондурасский\\
\textbf{nru7t} [\emph{nru˥t}] 0xB2C6: пробормотал\\
\textbf{nruc} [\emph{nruʃ}] 0xEAC6: моча\\
\textbf{nruf} [\emph{nruf}] 0xCAC6: отлив прибоя\\
\textbf{nruh} [\emph{nruʰ}] 0x5637: невинный\\
\textbf{nruk} [\emph{nruk}] 0x2AC6: инфекционный\\
\textbf{nrum} [\emph{nrum}] 0x0AC6: внутриутробный\\
\textbf{nrun} [\emph{nrun}] 0x6AC6: авторизоваться\\
\textbf{nrup} [\emph{nrup}] 0x4AC6: вход\\
\textbf{nrus} [\emph{nrus}] 0x8AC6: носорог\\
\textbf{nrut} [\emph{nrut}] 0xAAC6: развлекать\\
\textbf{nu} [\emph{nu}] 0x28DE: интерьер контекст\\
\textbf{nw6n} [\emph{nwən}] 0x6D26: новорожденный\\
\textbf{nw6p} [\emph{nwəp}] 0x4D26: Antwerp\\
\textbf{nw6t} [\emph{nwət}] 0xAD26: монолит\\
\textbf{nwa2k} [\emph{nwä˩k}] 0x3926: противоядный\\
\textbf{nwa2n} [\emph{nwä˩n}] 0x7926: евнух\\
\textbf{nwa2t} [\emph{nwä˩t}] 0xB926: санскрит\\
\textbf{nwa7k} [\emph{nwä˥k}] 0x3126: фантасмагория\\
\textbf{nwa7n} [\emph{nwä˥n}] 0x7126: размотанный\\
\textbf{nwa7s} [\emph{nwä˥s}] 0x9126: тонзиллит\\
\textbf{nwa7t} [\emph{nwä˥t}] 0xB126: неразрезанный\\
\textbf{nwac} [\emph{nwäʃ}] 0xE926: кто угодно\\
\textbf{nwaf} [\emph{nwäf}] 0xC926: пропускная способность\\
\textbf{nwah} [\emph{nwäʰ}] 0x4937: evidential case, for citing a source as evidence. discource-context source-case\\
\textbf{nwak} [\emph{nwäk}] 0x2926: январь\\
\textbf{nwam} [\emph{nwäm}] 0x0926: животное\\
\textbf{nwan} [\emph{nwän}] 0x6926: молодой\\
\textbf{nwap} [\emph{nwäp}] 0x4926: ноябрь\\
\textbf{nwas} [\emph{nwäs}] 0x8926: Новости\\
\textbf{nwat} [\emph{nwät}] 0xA926: расти\\
\textbf{nwe2n} [\emph{nwe˩n}] 0x7B26: один лайнер\\
\textbf{nwe2t} [\emph{nwe˩t}] 0xBB26: женское нижнее белье\\
\textbf{nwe7n} [\emph{nwe˥n}] 0x7326: антропоцентрический\\
\textbf{nwec} [\emph{nweʃ}] 0xEB26: тем не менее\\
\textbf{nweh} [\emph{nweʰ}] 0x5937: inessive case, indicates the noun phrase being in at. interior-context location-case. \\
\textbf{nwek} [\emph{nwek}] 0x2B26: генетика\\
\textbf{nwem} [\emph{nwem}] 0x0B26: ветреница\\
\textbf{nwen} [\emph{nwen}] 0x6B26: флот\\
\textbf{nwep} [\emph{nwep}] 0x4B26: анаэробный\\
\textbf{nwes} [\emph{nwes}] 0x8B26: инвестиции\\
\textbf{nwet} [\emph{nwet}] 0xAB26: вязаный\\
\textbf{nwi2k} [\emph{nwi˩k}] 0x3826: безводный\\
\textbf{nwi2n} [\emph{nwi˩n}] 0x7826: неконституционный\\
\textbf{nwi2s} [\emph{nwi˩s}] 0x9826: кровосмешение\\
\textbf{nwi2t} [\emph{nwi˩t}] 0xB826: наветренный\\
\textbf{nwi7f} [\emph{nwi˥f}] 0xD026: анчоус\\
\textbf{nwi7k} [\emph{nwi˥k}] 0x3026: синоптический\\
\textbf{nwi7n} [\emph{nwi˥n}] 0x7026: монахиня\\
\textbf{nwi7s} [\emph{nwi˥s}] 0x9026: канонизация\\
\textbf{nwi7t} [\emph{nwi˥t}] 0xB026: не введен в эксплуатацию\\
\textbf{nwic} [\emph{nwiʃ}] 0xE826: дюйм\\
\textbf{nwif} [\emph{nwif}] 0xC826: не живой\\
\textbf{nwih} [\emph{nwiʰ}] 0x4137: нео\\
\textbf{nwik} [\emph{nwik}] 0x2826: инессив дело\\
\textbf{nwim} [\emph{nwim}] 0x0826: сенсорный доказательный настроение\\
\textbf{nwin} [\emph{nwin}] 0x6826: конституция\\
\textbf{nwip} [\emph{nwip}] 0x4826: неизбежный\\
\textbf{nwis} [\emph{nwis}] 0x8826: ногти\\
\textbf{nwit} [\emph{nwit}] 0xA826: Событийный настроение\\
\textbf{nwo2t} [\emph{nwo˩t}] 0xBC26: тритон\\
\textbf{nwo7t} [\emph{nwo˥t}] 0xB426: берег\\
\textbf{nwoc} [\emph{nwoʃ}] 0xEC26: с этого времени\\
\textbf{nwof} [\emph{nwof}] 0xCC26: фонограф\\
\textbf{nwoh} [\emph{nwoʰ}] 0x6137: Событийный настроение\\
\textbf{nwok} [\emph{nwok}] 0x2C26: домашнее задание\\
\textbf{nwom} [\emph{nwom}] 0x0C26: анатомия\\
\textbf{nwon} [\emph{nwon}] 0x6C26: номинальный\\
\textbf{nwop} [\emph{nwop}] 0x4C26: ягодицы\\
\textbf{nwos} [\emph{nwos}] 0x8C26: инновация\\
\textbf{nwot} [\emph{nwot}] 0xAC26: непрерывность\\
\textbf{nwu7n} [\emph{nwu˥n}] 0x7226: Не забывай меня\\
\textbf{nwuc} [\emph{nwuʃ}] 0xEA26: движение спуск\\
\textbf{nwuh} [\emph{nwuʰ}] 0x5137: research gender, 6th density life form, learning to blend love and wisdom, akin to Ra, Gandhi, Tesla, Prophet Mohammad (pbuh)\\
\textbf{nwuk} [\emph{nwuk}] 0x2A26: уникальный\\
\textbf{nwum} [\emph{nwum}] 0x0A26: индуизм\\
\textbf{nwun} [\emph{nwun}] 0x6A26: пунктуальный\\
\textbf{nwup} [\emph{nwup}] 0x4A26: одноклеточный\\
\textbf{nwus} [\emph{nwus}] 0x8A26: сонный\\
\textbf{nwut} [\emph{nwut}] 0xAA26: на цыпочках\\
\textbf{ny67n} [\emph{njə˥n}] 0x7506: Анд\\
\textbf{ny6k} [\emph{njək}] 0x2D06: энный мощность\\
\textbf{ny6m} [\emph{njəm}] 0x0D06: анимизм\\
\textbf{ny6n} [\emph{njən}] 0x6D06: антенна\\
\textbf{ny6s} [\emph{njəs}] 0x8D06: неоклассический\\
\textbf{ny6t} [\emph{njət}] 0xAD06: ингибитор\\
\textbf{nya2k} [\emph{njä˩k}] 0x3906: выводить из строя\\
\textbf{nya2n} [\emph{njä˩n}] 0x7906: монахини\\
\textbf{nya2s} [\emph{njä˩s}] 0x9906: монашество\\
\textbf{nya2t} [\emph{njä˩t}] 0xB906: манифест\\
\textbf{nya7c} [\emph{njä˥ʃ}] 0xF106: отрезанный от внешнего мира\\
\textbf{nya7f} [\emph{njä˥f}] 0xD106: Племянница в законе\\
\textbf{nya7k} [\emph{njä˥k}] 0x3106: нанка\\
\textbf{nya7m} [\emph{njä˥m}] 0x1106: фундаментализм\\
\textbf{nya7n} [\emph{njä˥n}] 0x7106: недисциплинированный\\
\textbf{nya7p} [\emph{njä˥p}] 0x5106: односложное слово\\
\textbf{nya7s} [\emph{njä˥s}] 0x9106: евангелистский\\
\textbf{nya7t} [\emph{njä˥t}] 0xB106: сводный брат\\
\textbf{nyac} [\emph{njäʃ}] 0xE906: направление\\
\textbf{nyaf} [\emph{njäf}] 0xC906: синтаксический\\
\textbf{nyah} [\emph{njäʰ}] 0x4837: нас\\
\textbf{nyak} [\emph{njäk}] 0x2906: товарный знак\\
\textbf{nyam} [\emph{njäm}] 0x0906: слава\\
\textbf{nyan} [\emph{njän}] 0x6906: мужской\\
\textbf{nyap} [\emph{njäp}] 0x4906: узкий\\
\textbf{nyas} [\emph{njäs}] 0x8906: сочинение\\
\textbf{nyat} [\emph{njät}] 0xA906: атака\\
\textbf{nye2k} [\emph{nje˩k}] 0x3B06: сантиграмм\\
\textbf{nye2n} [\emph{nje˩n}] 0x7B06: нейлон\\
\textbf{nye2t} [\emph{nje˩t}] 0xBB06: ланцетовидный\\
\textbf{nye7n} [\emph{nje˥n}] 0x7306: индексация\\
\textbf{nye7t} [\emph{nje˥t}] 0xB306: анти семитизм\\
\textbf{nyec} [\emph{njeʃ}] 0xEB06: Индонезия\\
\textbf{nyef} [\emph{njef}] 0xCB06: центробежный\\
\textbf{nyeh} [\emph{njeʰ}] 0x5837: initiative case, indicates the begining of an action. time-context source-case\\
\textbf{nyek} [\emph{njek}] 0x2B06: никель\\
\textbf{nyem} [\emph{njem}] 0x0B06: сантиметр\\
\textbf{nyen} [\emph{njen}] 0x6B06: стон\\
\textbf{nyep} [\emph{njep}] 0x4B06: нюхательный табак\\
\textbf{nyes} [\emph{njes}] 0x8B06: интеркостельный\\
\textbf{nyet} [\emph{njet}] 0xAB06: подушка\\
\textbf{nyi2k} [\emph{nji˩k}] 0x3806: невралгический\\
\textbf{nyi2n} [\emph{nji˩n}] 0x7806: единорог\\
\textbf{nyi2s} [\emph{nji˩s}] 0x9806: самовлюбленный человек\\
\textbf{nyi2t} [\emph{nji˩t}] 0xB806: индис\\
\textbf{nyi7c} [\emph{nji˥ʃ}] 0xF006: вновь ввести\\
\textbf{nyi7k} [\emph{nji˥k}] 0x3006: незаметный\\
\textbf{nyi7n} [\emph{nji˥n}] 0x7006: ненароком\\
\textbf{nyi7p} [\emph{nji˥p}] 0x5006: односложный\\
\textbf{nyi7s} [\emph{nji˥s}] 0x9006: перемежаются\\
\textbf{nyi7t} [\emph{nji˥t}] 0xB006: тридцать восемь\\
\textbf{nyic} [\emph{njiʃ}] 0xE806: одетый\\
\textbf{nyif} [\emph{njif}] 0xC806: коннотация\\
\textbf{nyih} [\emph{njiʰ}] 0x4037: неотъемлемое владение\\
\textbf{nyik} [\emph{njik}] 0x2806: грязь\\
\textbf{nyim} [\emph{njim}] 0x0806: выдвижение\\
\textbf{nyin} [\emph{njin}] 0x6806: антропный Пол\\
\textbf{nyip} [\emph{njip}] 0x4806: в местном масштабе\\
\textbf{nyis} [\emph{njis}] 0x8806: транскрибировать\\
\textbf{nyit} [\emph{njit}] 0xA806: Сортировать\\
\textbf{nyo2k} [\emph{njo˩k}] 0x3C06: линолевая\\
\textbf{nyo2n} [\emph{njo˩n}] 0x7C06: назначаемое лицо\\
\textbf{nyo2t} [\emph{njo˩t}] 0xBC06: гностический\\
\textbf{nyo7k} [\emph{njo˥k}] 0x3406: монокль\\
\textbf{nyo7n} [\emph{njo˥n}] 0x7406: нерожденный\\
\textbf{nyo7t} [\emph{njo˥t}] 0xB406: невысказанный\\
\textbf{nyoc} [\emph{njoʃ}] 0xEC06: антропология\\
\textbf{nyof} [\emph{njof}] 0xCC06: ионосфера\\
\textbf{nyoh} [\emph{njoʰ}] 0x6037: анимационный\\
\textbf{nyok} [\emph{njok}] 0x2C06: аналоговый\\
\textbf{nyom} [\emph{njom}] 0x0C06: неодим\\
\textbf{nyon} [\emph{njon}] 0x6C06: анонимный\\
\textbf{nyop} [\emph{njop}] 0x4C06: ниобий\\
\textbf{nyos} [\emph{njos}] 0x8C06: неамбициозны\\
\textbf{nyot} [\emph{njot}] 0xAC06: слушатель\\
\textbf{nyu2t} [\emph{nju˩t}] 0xBA06: синусит\\
\textbf{nyu7n} [\emph{nju˥n}] 0x7206: нет нет\\
\textbf{nyu7t} [\emph{nju˥t}] 0xB206: наделять\\
\textbf{nyuc} [\emph{njuʃ}] 0xEA06: маникюр\\
\textbf{nyuh} [\emph{njuʰ}] 0x5037: транс цифра\\
\textbf{nyuk} [\emph{njuk}] 0x2A06: Светляк\\
\textbf{nyum} [\emph{njum}] 0x0A06: эму\\
\textbf{nyun} [\emph{njun}] 0x6A06: тепло\\
\textbf{nyup} [\emph{njup}] 0x4A06: высиживать\\
\textbf{nyus} [\emph{njus}] 0x8A06: индуктор\\
\textbf{nyut} [\emph{njut}] 0xAA06: ваш\\
\subsection{ p }
\textbf{p6} [\emph{pə}] 0x349E: сращивание\\
\textbf{pa} [\emph{pä}] 0x249E: человек контекст\\
\textbf{pc6h} [\emph{pʃəʰ}] 0x6D27: испытание номер\\
\textbf{pc6n} [\emph{pʃən}] 0x6DA4: испытание номер\\
\textbf{pc6t} [\emph{pʃət}] 0xADA4: спринтер\\
\textbf{pca2k} [\emph{pʃä˩k}] 0x39A4: супердержавы\\
\textbf{pca2n} [\emph{pʃä˩n}] 0x79A4: фонтанирование\\
\textbf{pca2p} [\emph{pʃä˩p}] 0x59A4: ломбард\\
\textbf{pca2t} [\emph{pʃä˩t}] 0xB9A4: наглость\\
\textbf{pca7c} [\emph{pʃä˥ʃ}] 0xF1A4: паша\\
\textbf{pca7k} [\emph{pʃä˥k}] 0x31A4: магическая фигура\\
\textbf{pca7n} [\emph{pʃä˥n}] 0x71A4: протестант\\
\textbf{pca7p} [\emph{pʃä˥p}] 0x51A4: парашютист\\
\textbf{pca7s} [\emph{pʃä˥s}] 0x91A4: похитители\\
\textbf{pca7t} [\emph{pʃä˥t}] 0xB1A4: идолопоклонник\\
\textbf{pcac} [\emph{pʃäʃ}] 0xE9A4: лучше\\
\textbf{pcaf} [\emph{pʃäf}] 0xC9A4: круглая скобка\\
\textbf{pcah} [\emph{pʃäʰ}] 0x4D27: также\\
\textbf{pcak} [\emph{pʃäk}] 0x29A4: запад\\
\textbf{pcam} [\emph{pʃäm}] 0x09A4: молиться\\
\textbf{pcan} [\emph{pʃän}] 0x69A4: Деньги\\
\textbf{pcap} [\emph{pʃäp}] 0x49A4: в стороне\\
\textbf{pcas} [\emph{pʃäs}] 0x89A4: взгляд\\
\textbf{pcat} [\emph{pʃät}] 0xA9A4: Понимаю\\
\textbf{pce2n} [\emph{pʃe˩n}] 0x7BA4: довоенный\\
\textbf{pce2s} [\emph{pʃe˩s}] 0x9BA4: сепсис\\
\textbf{pce2t} [\emph{pʃe˩t}] 0xBBA4: яблоко щеками\\
\textbf{pce7n} [\emph{pʃe˥n}] 0x73A4: пектин\\
\textbf{pce7t} [\emph{pʃe˥t}] 0xB3A4: полубог\\
\textbf{pcec} [\emph{pʃeʃ}] 0xEBA4: привилегированный\\
\textbf{pcef} [\emph{pʃef}] 0xCBA4: рептилии\\
\textbf{pceh} [\emph{pʃeʰ}] 0x5D27: apprehensive mood, for expressing fears or warnings.\\
\textbf{pcek} [\emph{pʃek}] 0x2BA4: похищение людей\\
\textbf{pcem} [\emph{pʃem}] 0x0BA4: подкожный\\
\textbf{pcen} [\emph{pʃen}] 0x6BA4: президент\\
\textbf{pcep} [\emph{pʃep}] 0x4BA4: легкое\\
\textbf{pces} [\emph{pʃes}] 0x8BA4: блин\\
\textbf{pcet} [\emph{pʃet}] 0xABA4: неопределенность\\
\textbf{pci2c} [\emph{pʃi˩ʃ}] 0xF8A4: зяблики\\
\textbf{pci2m} [\emph{pʃi˩m}] 0x18A4: паренхима\\
\textbf{pci2n} [\emph{pʃi˩n}] 0x78A4: импровизация\\
\textbf{pci2s} [\emph{pʃi˩s}] 0x98A4: импрессионизм\\
\textbf{pci2t} [\emph{pʃi˩t}] 0xB8A4: педантичность\\
\textbf{pci7c} [\emph{pʃi˥ʃ}] 0xF0A4: путеводная звезда\\
\textbf{pci7k} [\emph{pʃi˥k}] 0x30A4: мизинец\\
\textbf{pci7n} [\emph{pʃi˥n}] 0x70A4: плиоцен\\
\textbf{pci7p} [\emph{pʃi˥p}] 0x50A4: пирог столкнулся\\
\textbf{pci7t} [\emph{pʃi˥t}] 0xB0A4: лицемер\\
\textbf{pcic} [\emph{pʃiʃ}] 0xE8A4: должен\\
\textbf{pcif} [\emph{pʃif}] 0xC8A4: холодильник\\
\textbf{pcih} [\emph{pʃiʰ}] 0x4527: imperative mood, to indicate direct commands.\\
\textbf{pcik} [\emph{pʃik}] 0x28A4: по всей видимости\\
\textbf{pcim} [\emph{pʃim}] 0x08A4: выполнять\\
\textbf{pcin} [\emph{pʃin}] 0x68A4: за\\
\textbf{pcip} [\emph{pʃip}] 0x48A4: карандаш\\
\textbf{pcis} [\emph{pʃis}] 0x88A4: опасно\\
\textbf{pcit} [\emph{pʃit}] 0xA8A4: Авторские права\\
\textbf{pco2n} [\emph{pʃo˩n}] 0x7CA4: мочеиспускание\\
\textbf{pco7t} [\emph{pʃo˥t}] 0xB4A4: подъязычная\\
\textbf{pcoc} [\emph{pʃoʃ}] 0xECA4: столб\\
\textbf{pcof} [\emph{pʃof}] 0xCCA4: порнографический\\
\textbf{pcoh} [\emph{pʃoʰ}] 0x6527: заостренный\\
\textbf{pcok} [\emph{pʃok}] 0x2CA4: спотыкаться\\
\textbf{pcom} [\emph{pʃom}] 0x0CA4: операционная система\\
\textbf{pcon} [\emph{pʃon}] 0x6CA4: латунь\\
\textbf{pcop} [\emph{pʃop}] 0x4CA4: насос\\
\textbf{pcos} [\emph{pʃos}] 0x8CA4: виды спорта\\
\textbf{pcot} [\emph{pʃot}] 0xACA4: домашнее животное\\
\textbf{pcuc} [\emph{pʃuʃ}] 0xEAA4: разряд\\
\textbf{pcuh} [\emph{pʃuʰ}] 0x5527: неопределенность\\
\textbf{pcuk} [\emph{pʃuk}] 0x2AA4: послелоги дело\\
\textbf{pcum} [\emph{pʃum}] 0x0AA4: персиковое дерево\\
\textbf{pcun} [\emph{pʃun}] 0x6AA4: банка\\
\textbf{pcup} [\emph{pʃup}] 0x4AA4: зловещий\\
\textbf{pcus} [\emph{pʃus}] 0x8AA4: мак\\
\textbf{pcut} [\emph{pʃut}] 0xAAA4: занято\\
\textbf{pe} [\emph{pe}] 0x2C9E: в стороне\\
\textbf{pe7} [\emph{pe˥}] 0x4C9E: за\\
\textbf{pf6m} [\emph{pfəm}] 0x0D84: в изобилии\\
\textbf{pf6n} [\emph{pfən}] 0x6D84: опыление\\
\textbf{pf6t} [\emph{pfət}] 0xAD84: суппозиторий\\
\textbf{pfa2k} [\emph{pfä˩k}] 0x3984: профилактика\\
\textbf{pfa2n} [\emph{pfä˩n}] 0x7984: пантеон\\
\textbf{pfa2s} [\emph{pfä˩s}] 0x9984: апофеоз\\
\textbf{pfa2t} [\emph{pfä˩t}] 0xB984: помогать\\
\textbf{pfa7f} [\emph{pfä˥f}] 0xD184: нафтол\\
\textbf{pfa7k} [\emph{pfä˥k}] 0x3184: подсемейство\\
\textbf{pfa7n} [\emph{pfä˥n}] 0x7184: восемьдесят пять\\
\textbf{pfa7s} [\emph{pfä˥s}] 0x9184: нераспространение\\
\textbf{pfa7t} [\emph{pfä˥t}] 0xB184: пацифист\\
\textbf{pfac} [\emph{pfäʃ}] 0xE984: подвергать\\
\textbf{pfaf} [\emph{pfäf}] 0xC984: сапфир\\
\textbf{pfah} [\emph{pfäʰ}] 0x4C27: выгодный\\
\textbf{pfak} [\emph{pfäk}] 0x2984: благоприятствовать\\
\textbf{pfam} [\emph{pfäm}] 0x0984: спамер\\
\textbf{pfan} [\emph{pfän}] 0x6984: вероятность\\
\textbf{pfap} [\emph{pfäp}] 0x4984: ямкоделатель\\
\textbf{pfas} [\emph{pfäs}] 0x8984: причастие\\
\textbf{pfat} [\emph{pfät}] 0xA984: доска\\
\textbf{pfe2n} [\emph{pfe˩n}] 0x7B84: полу профессиональный\\
\textbf{pfec} [\emph{pfeʃ}] 0xEB84: незамеченной\\
\textbf{pfef} [\emph{pfef}] 0xCB84: упреждающий\\
\textbf{pfeh} [\emph{pfeʰ}] 0x5C27: чуткими радость\\
\textbf{pfek} [\emph{pfek}] 0x2B84: орех-пекан\\
\textbf{pfem} [\emph{pfem}] 0x0B84: эпителий\\
\textbf{pfen} [\emph{pfen}] 0x6B84: проверка\\
\textbf{pfep} [\emph{pfep}] 0x4B84: гипнотерапия\\
\textbf{pfes} [\emph{pfes}] 0x8B84: бледное лицо\\
\textbf{pfet} [\emph{pfet}] 0xAB84: листовка\\
\textbf{pfi2n} [\emph{pfi˩n}] 0x7884: полдюйма\\
\textbf{pfi2t} [\emph{pfi˩t}] 0xB884: неподписанный\\
\textbf{pfi7k} [\emph{pfi˥k}] 0x3084: коклюш\\
\textbf{pfi7n} [\emph{pfi˥n}] 0x7084: отделимый\\
\textbf{pfi7s} [\emph{pfi˥s}] 0x9084: гипофиз\\
\textbf{pfi7t} [\emph{pfi˥t}] 0xB084: несведущий\\
\textbf{pfic} [\emph{pfiʃ}] 0xE884: представление\\
\textbf{pfif} [\emph{pfif}] 0xC884: Спецификация\\
\textbf{pfih} [\emph{pfiʰ}] 0x4427: progressive aspect, the action is in progress, it is happening actively. \\
\textbf{pfik} [\emph{pfik}] 0x2884: истребование\\
\textbf{pfim} [\emph{pfim}] 0x0884: период полураспада\\
\textbf{pfin} [\emph{pfin}] 0x6884: гулять пешком\\
\textbf{pfip} [\emph{pfip}] 0x4884: партитивный\\
\textbf{pfis} [\emph{pfis}] 0x8884: ратификация\\
\textbf{pfit} [\emph{pfit}] 0xA884: представитель\\
\textbf{pfo2n} [\emph{pfo˩n}] 0x7C84: продувание\\
\textbf{pfo7t} [\emph{pfo˥t}] 0xB484: огнемет\\
\textbf{pfoc} [\emph{pfoʃ}] 0xEC84: профилирование\\
\textbf{pfof} [\emph{pfof}] 0xCC84: порфир\\
\textbf{pfok} [\emph{pfok}] 0x2C84: модификатор\\
\textbf{pfon} [\emph{pfon}] 0x6C84: губка\\
\textbf{pfop} [\emph{pfop}] 0x4C84: поп вверх\\
\textbf{pfos} [\emph{pfos}] 0x8C84: профессор\\
\textbf{pfot} [\emph{pfot}] 0xAC84: горшок\\
\textbf{pfuf} [\emph{pfuf}] 0xCA84: летучие мыши\\
\textbf{pfuh} [\emph{pfuʰ}] 0x5427: прощальный привет\\
\textbf{pfuk} [\emph{pfuk}] 0x2A84: надколенник\\
\textbf{pfun} [\emph{pfun}] 0x6A84: частично\\
\textbf{pfus} [\emph{pfus}] 0x8A84: префикс\\
\textbf{pfut} [\emph{pfut}] 0xAA84: грудь\\
\textbf{pi} [\emph{pi}] 0x209E: поверхность контекст\\
\textbf{pi2} [\emph{pi˩}] 0x609E: planetary gender, 9th density life form, akin to spirit of a planet, Gaia, Venus, Luna\\
\textbf{pi7} [\emph{pi˥}] 0x409E: propositive mood, for making proposals or offers, let's\\
\textbf{pl6c} [\emph{pləʃ}] 0xED64: щегольство\\
\textbf{pl6k} [\emph{plək}] 0x2D64: республиканец\\
\textbf{pl6m} [\emph{pləm}] 0x0D64: привратник желудка\\
\textbf{pl6n} [\emph{plən}] 0x6D64: Польша\\
\textbf{pl6p} [\emph{pləp}] 0x4D64: лепестки\\
\textbf{pl6s} [\emph{pləs}] 0x8D64: ананас\\
\textbf{pl6t} [\emph{plət}] 0xAD64: сыпь\\
\textbf{pla2c} [\emph{plä˩ʃ}] 0xF964: небный\\
\textbf{pla2k} [\emph{plä˩k}] 0x3964: раз в две недели\\
\textbf{pla2n} [\emph{plä˩n}] 0x7964: павильон\\
\textbf{pla2p} [\emph{plä˩p}] 0x5964: многосложный\\
\textbf{pla2s} [\emph{plä˩s}] 0x9964: ипостась\\
\textbf{pla2t} [\emph{plä˩t}] 0xB964: блюдо\\
\textbf{pla7c} [\emph{plä˥ʃ}] 0xF164: яблочный соус\\
\textbf{pla7k} [\emph{plä˥k}] 0x3164: могильный\\
\textbf{pla7n} [\emph{plä˥n}] 0x7164: Лапландия\\
\textbf{pla7p} [\emph{plä˥p}] 0x5164: непреложный\\
\textbf{pla7s} [\emph{plä˥s}] 0x9164: орудийный огонь\\
\textbf{pla7t} [\emph{plä˥t}] 0xB164: превозносить\\
\textbf{plac} [\emph{pläʃ}] 0xE964: дворец\\
\textbf{plaf} [\emph{pläf}] 0xC964: пыльца\\
\textbf{plah} [\emph{pläʰ}] 0x4B27: люблю\\
\textbf{plak} [\emph{pläk}] 0x2964: Население\\
\textbf{plam} [\emph{pläm}] 0x0964: аварийная сигнализация\\
\textbf{plan} [\emph{plän}] 0x6964: играть\\
\textbf{plap} [\emph{pläp}] 0x4964: стручок\\
\textbf{plas} [\emph{pläs}] 0x8964: взрыв\\
\textbf{plat} [\emph{plät}] 0xA964: стена\\
\textbf{ple2c} [\emph{ple˩ʃ}] 0xFB64: депиляция\\
\textbf{ple2n} [\emph{ple˩n}] 0x7B64: пленум\\
\textbf{ple2t} [\emph{ple˩t}] 0xBB64: предоплатой\\
\textbf{ple7k} [\emph{ple˥k}] 0x3364: умереть раньше\\
\textbf{ple7m} [\emph{ple˥m}] 0x1364: Пелл Мелл\\
\textbf{ple7n} [\emph{ple˥n}] 0x7364: перигелий\\
\textbf{ple7t} [\emph{ple˥t}] 0xB364: эспланада\\
\textbf{plec} [\emph{pleʃ}] 0xEB64: сращивание\\
\textbf{plef} [\emph{plef}] 0xCB64: парапланеризм\\
\textbf{pleh} [\emph{pleʰ}] 0x5B27: привилегированный\\
\textbf{plek} [\emph{plek}] 0x2B64: капиллярный\\
\textbf{plem} [\emph{plem}] 0x0B64: полимер\\
\textbf{plen} [\emph{plen}] 0x6B64: равновесие\\
\textbf{plep} [\emph{plep}] 0x4B64: палитра\\
\textbf{ples} [\emph{ples}] 0x8B64: процессор\\
\textbf{plet} [\emph{plet}] 0xAB64: издатель\\
\textbf{pli2c} [\emph{pli˩ʃ}] 0xF864: поли\\
\textbf{pli2f} [\emph{pli˩f}] 0xD864: мостильщик\\
\textbf{pli2k} [\emph{pli˩k}] 0x3864: апологетика\\
\textbf{pli2n} [\emph{pli˩n}] 0x7864: пинг\\
\textbf{pli2s} [\emph{pli˩s}] 0x9864: эпилепсия\\
\textbf{pli2t} [\emph{pli˩t}] 0xB864: сифилис\\
\textbf{pli7c} [\emph{pli˥ʃ}] 0xF064: параплегия\\
\textbf{pli7f} [\emph{pli˥f}] 0xD064: полиомиелит\\
\textbf{pli7k} [\emph{pli˥k}] 0x3064: патология\\
\textbf{pli7n} [\emph{pli˥n}] 0x7064: Пенсильвания\\
\textbf{pli7p} [\emph{pli˥p}] 0x5064: непродаваемый\\
\textbf{pli7s} [\emph{pli˥s}] 0x9064: неправдоподобно\\
\textbf{pli7t} [\emph{pli˥t}] 0xB064: Исландия\\
\textbf{plic} [\emph{pliʃ}] 0xE864: будущее\\
\textbf{plif} [\emph{plif}] 0xC864: вилла\\
\textbf{plih} [\emph{pliʰ}] 0x4327: частный\\
\textbf{plik} [\emph{plik}] 0x2864: применимый\\
\textbf{plim} [\emph{plim}] 0x0864: калий\\
\textbf{plin} [\emph{plin}] 0x6864: открытие\\
\textbf{plip} [\emph{plip}] 0x4864: высказал мнение,\\
\textbf{plis} [\emph{plis}] 0x8864: полиция\\
\textbf{plit} [\emph{plit}] 0xA864: здание\\
\textbf{plo2t} [\emph{plo˩t}] 0xBC64: Платон\\
\textbf{plo7k} [\emph{plo˥k}] 0x3464: палеолитический\\
\textbf{plo7n} [\emph{plo˥n}] 0x7464: обесцвечивание\\
\textbf{plo7s} [\emph{plo˥s}] 0x9464: прозелитизм\\
\textbf{plo7t} [\emph{plo˥t}] 0xB464: орфографический справочник\\
\textbf{ploc} [\emph{ploʃ}] 0xEC64: экология\\
\textbf{plof} [\emph{plof}] 0xCC64: силовая установка\\
\textbf{ploh} [\emph{ploʰ}] 0x6327: за предложение\\
\textbf{plok} [\emph{plok}] 0x2C64: дополнение\\
\textbf{plom} [\emph{plom}] 0x0C64: почтовый штемпель\\
\textbf{plon} [\emph{plon}] 0x6C64: выделение\\
\textbf{plop} [\emph{plop}] 0x4C64: палладий\\
\textbf{plos} [\emph{plos}] 0x8C64: пульс\\
\textbf{plot} [\emph{plot}] 0xAC64: пилот\\
\textbf{plu2t} [\emph{plu˩t}] 0xBA64: подколенный\\
\textbf{plu7t} [\emph{plu˥t}] 0xB264: молодка\\
\textbf{pluc} [\emph{pluʃ}] 0xEA64: paucal номер\\
\textbf{pluh} [\emph{pluʰ}] 0x5327: превосходный степень\\
\textbf{pluk} [\emph{pluk}] 0x2A64: оценить\\
\textbf{plum} [\emph{plum}] 0x0A64: пролог\\
\textbf{plun} [\emph{plun}] 0x6A64: синий\\
\textbf{plup} [\emph{plup}] 0x4A64: тонкой кожурой\\
\textbf{plus} [\emph{plus}] 0x8A64: плюс\\
\textbf{plut} [\emph{plut}] 0xAA64: картошка\\
\textbf{po} [\emph{po}] 0x309E: paucal номер\\
\textbf{pr6c} [\emph{prəʃ}] 0xEDC4: взаимопроникновение\\
\textbf{pr6k} [\emph{prək}] 0x2DC4: Парагвай\\
\textbf{pr6m} [\emph{prəm}] 0x0DC4: пермский\\
\textbf{pr6n} [\emph{prən}] 0x6DC4: Я побежал\\
\textbf{pr6p} [\emph{prəp}] 0x4DC4: груши\\
\textbf{pr6s} [\emph{prəs}] 0x8DC4: перс\\
\textbf{pr6t} [\emph{prət}] 0xADC4: полярность\\
\textbf{pra2c} [\emph{prä˩ʃ}] 0xF9C4: черная работа\\
\textbf{pra2k} [\emph{prä˩k}] 0x39C4: паркет\\
\textbf{pra2m} [\emph{prä˩m}] 0x19C4: отчество\\
\textbf{pra2n} [\emph{prä˩n}] 0x79C4: преступники\\
\textbf{pra2s} [\emph{prä˩s}] 0x99C4: папирус\\
\textbf{pra2t} [\emph{prä˩t}] 0xB9C4: катаракта\\
\textbf{pra7c} [\emph{prä˥ʃ}] 0xF1C4: гипертония\\
\textbf{pra7k} [\emph{prä˥k}] 0x31C4: Прага\\
\textbf{pra7m} [\emph{prä˥m}] 0x11C4: перикард\\
\textbf{pra7n} [\emph{prä˥n}] 0x71C4: шторы\\
\textbf{pra7p} [\emph{prä˥p}] 0x51C4: апперцепция\\
\textbf{pra7s} [\emph{prä˥s}] 0x91C4: арии\\
\textbf{pra7t} [\emph{prä˥t}] 0xB1C4: сохраняет\\
\textbf{prac} [\emph{präʃ}] 0xE9C4: выводить\\
\textbf{praf} [\emph{präf}] 0xC9C4: пиратство\\
\textbf{prah} [\emph{präʰ}] 0x4E27: параграф\\
\textbf{prak} [\emph{präk}] 0x29C4: проповедовать\\
\textbf{pram} [\emph{präm}] 0x09C4: альпинизм\\
\textbf{pran} [\emph{prän}] 0x69C4: Нажмите\\
\textbf{prap} [\emph{präp}] 0x49C4: смещение\\
\textbf{pras} [\emph{präs}] 0x89C4: вуаль\\
\textbf{prat} [\emph{prät}] 0xA9C4: сообщается доказательный\\
\textbf{pre2k} [\emph{pre˩k}] 0x3BC4: грушевидной формы\\
\textbf{pre2n} [\emph{pre˩n}] 0x7BC4: плеврит\\
\textbf{pre2t} [\emph{pre˩t}] 0xBBC4: репертуар\\
\textbf{pre7k} [\emph{pre˥k}] 0x33C4: за\\
\textbf{pre7n} [\emph{pre˥n}] 0x73C4: мокрота\\
\textbf{pre7s} [\emph{pre˥s}] 0x93C4: Пиренеи\\
\textbf{pre7t} [\emph{pre˥t}] 0xB3C4: спорщик\\
\textbf{prec} [\emph{preʃ}] 0xEBC4: вести себя\\
\textbf{pref} [\emph{pref}] 0xCBC4: партер\\
\textbf{preh} [\emph{preʰ}] 0x5E27: предопределенный\\
\textbf{prek} [\emph{prek}] 0x2BC4: стоянка\\
\textbf{prem} [\emph{prem}] 0x0BC4: нефть\\
\textbf{pren} [\emph{pren}] 0x6BC4: параллельно\\
\textbf{prep} [\emph{prep}] 0x4BC4: эпидермис\\
\textbf{pres} [\emph{pres}] 0x8BC4: кора\\
\textbf{pret} [\emph{pret}] 0xABC4: недостаточный\\
\textbf{pri2f} [\emph{pri˩f}] 0xD8C4: гипертрофия\\
\textbf{pri2k} [\emph{pri˩k}] 0x38C4: непонятый\\
\textbf{pri2n} [\emph{pri˩n}] 0x78C4: просматривал\\
\textbf{pri2s} [\emph{pri˩s}] 0x98C4: восемьдесят три\\
\textbf{pri2t} [\emph{pri˩t}] 0xB8C4: патриций\\
\textbf{pri7c} [\emph{pri˥ʃ}] 0xF0C4: перигей\\
\textbf{pri7h} [\emph{pri˥ʰ}] 0x8627: пассивный причастие\\
\textbf{pri7k} [\emph{pri˥k}] 0x30C4: образно\\
\textbf{pri7m} [\emph{pri˥m}] 0x10C4: ипподром\\
\textbf{pri7n} [\emph{pri˥n}] 0x70C4: феи\\
\textbf{pri7p} [\emph{pri˥p}] 0x50C4: околоцветник\\
\textbf{pri7s} [\emph{pri˥s}] 0x90C4: Париж\\
\textbf{pri7t} [\emph{pri˥t}] 0xB0C4: святоша\\
\textbf{pric} [\emph{priʃ}] 0xE8C4: подача\\
\textbf{prif} [\emph{prif}] 0xC8C4: переменная\\
\textbf{prih} [\emph{priʰ}] 0x4627: письмо\\
\textbf{prik} [\emph{prik}] 0x28C4: пассивный\\
\textbf{prim} [\emph{prim}] 0x08C4: имперский\\
\textbf{prin} [\emph{prin}] 0x68C4: горизонт\\
\textbf{prip} [\emph{prip}] 0x48C4: подготовка\\
\textbf{pris} [\emph{pris}] 0x88C4: вирус\\
\textbf{prit} [\emph{prit}] 0xA8C4: проклинающий настроение\\
\textbf{pro2k} [\emph{pro˩k}] 0x3CC4: гиперактивный\\
\textbf{pro2m} [\emph{pro˩m}] 0x1CC4: погром\\
\textbf{pro2n} [\emph{pro˩n}] 0x7CC4: странствование\\
\textbf{pro2s} [\emph{pro˩s}] 0x9CC4: пироксен\\
\textbf{pro2t} [\emph{pro˩t}] 0xBCC4: оперетта\\
\textbf{pro7k} [\emph{pro˥k}] 0x34C4: волей-неволей\\
\textbf{pro7n} [\emph{pro˥n}] 0x74C4: прихорашиваться\\
\textbf{pro7t} [\emph{pro˥t}] 0xB4C4: пенистый\\
\textbf{proc} [\emph{proʃ}] 0xECC4: покровитель\\
\textbf{prof} [\emph{prof}] 0xCCC4: уничижительный\\
\textbf{proh} [\emph{proʰ}] 0x6627: сообщается доказательный\\
\textbf{prok} [\emph{prok}] 0x2CC4: оптический\\
\textbf{prom} [\emph{prom}] 0x0CC4: программист\\
\textbf{pron} [\emph{pron}] 0x6CC4: профессионал\\
\textbf{prop} [\emph{prop}] 0x4CC4: плодородный\\
\textbf{pros} [\emph{pros}] 0x8CC4: операционная\\
\textbf{prot} [\emph{prot}] 0xACC4: портрет\\
\textbf{pru2t} [\emph{pru˩t}] 0xBAC4: штифтик\\
\textbf{pru7n} [\emph{pru˥n}] 0x72C4: перестрахование\\
\textbf{pru7t} [\emph{pru˥t}] 0xB2C4: породистый\\
\textbf{pruc} [\emph{pruʃ}] 0xEAC4: патруль\\
\textbf{pruf} [\emph{pruf}] 0xCAC4: пуриновых\\
\textbf{pruh} [\emph{pruʰ}] 0x5627: imprecative mood, for uttering curses.\\
\textbf{pruk} [\emph{pruk}] 0x2AC4: супермаркет\\
\textbf{prum} [\emph{prum}] 0x0AC4: развертываться\\
\textbf{prun} [\emph{prun}] 0x6AC4: вращение\\
\textbf{prup} [\emph{prup}] 0x4AC4: Перу\\
\textbf{prus} [\emph{prus}] 0x8AC4: Компьютерный вирус\\
\textbf{prut} [\emph{prut}] 0xAAC4: жюри\\
\textbf{ps67t} [\emph{psə˥t}] 0xB544: компенсатор\\
\textbf{ps6c} [\emph{psəʃ}] 0xED44: просодия\\
\textbf{ps6k} [\emph{psək}] 0x2D44: поршень\\
\textbf{ps6m} [\emph{psəm}] 0x0D44: хурма\\
\textbf{ps6n} [\emph{psən}] 0x6D44: испанский\\
\textbf{ps6p} [\emph{psəp}] 0x4D44: несмешивающийся\\
\textbf{ps6s} [\emph{psəs}] 0x8D44: макаронные изделия\\
\textbf{ps6t} [\emph{psət}] 0xAD44: Пасха\\
\textbf{psa2c} [\emph{psä˩ʃ}] 0xF944: восемьдесят семь\\
\textbf{psa2k} [\emph{psä˩k}] 0x3944: полу еженедельно\\
\textbf{psa2n} [\emph{psä˩n}] 0x7944: прохожие\\
\textbf{psa2p} [\emph{psä˩p}] 0x5944: вспыльчивый\\
\textbf{psa2s} [\emph{psä˩s}] 0x9944: пароксизм\\
\textbf{psa2t} [\emph{psä˩t}] 0xB944: вялость\\
\textbf{psa7c} [\emph{psä˥ʃ}] 0xF144: герметизация\\
\textbf{psa7f} [\emph{psä˥f}] 0xD144: апостроф\\
\textbf{psa7h} [\emph{psä˥ʰ}] 0x8A27: изнашивания\\
\textbf{psa7k} [\emph{psä˥k}] 0x3144: баскский\\
\textbf{psa7n} [\emph{psä˥n}] 0x7144: рассеянно\\
\textbf{psa7p} [\emph{psä˥p}] 0x5144: утконос\\
\textbf{psa7s} [\emph{psä˥s}] 0x9144: причуды\\
\textbf{psa7t} [\emph{psä˥t}] 0xB144: оседлал\\
\textbf{psac} [\emph{psäʃ}] 0xE944: невозможно\\
\textbf{psaf} [\emph{psäf}] 0xC944: отсутствие\\
\textbf{psah} [\emph{psäʰ}] 0x4A27: отчаяние\\
\textbf{psak} [\emph{psäk}] 0x2944: прохождение\\
\textbf{psam} [\emph{psäm}] 0x0944: поза\\
\textbf{psan} [\emph{psän}] 0x6944: половина\\
\textbf{psap} [\emph{psäp}] 0x4944: наихудший\\
\textbf{psas} [\emph{psäs}] 0x8944: производить\\
\textbf{psat} [\emph{psät}] 0xA944: вместо\\
\textbf{pse2k} [\emph{pse˩k}] 0x3B44: перистальтика\\
\textbf{pse2n} [\emph{pse˩n}] 0x7B44: точильный камень\\
\textbf{pse2p} [\emph{pse˩p}] 0x5B44: асептика\\
\textbf{pse2s} [\emph{pse˩s}] 0x9B44: прецессия\\
\textbf{pse2t} [\emph{pse˩t}] 0xBB44: шестьдесят семь\\
\textbf{pse7k} [\emph{pse˥k}] 0x3344: проклинать\\
\textbf{pse7n} [\emph{pse˥n}] 0x7344: пресс-банда\\
\textbf{pse7p} [\emph{pse˥p}] 0x5344: космодром\\
\textbf{pse7s} [\emph{pse˥s}] 0x9344: плейстоцен\\
\textbf{psec} [\emph{pseʃ}] 0xEB44: подмножество\\
\textbf{psef} [\emph{psef}] 0xCB44: эпицентр\\
\textbf{pseh} [\emph{pseʰ}] 0x5A27: аспект\\
\textbf{psek} [\emph{psek}] 0x2B44: superessive дело\\
\textbf{psem} [\emph{psem}] 0x0B44: плазма\\
\textbf{psen} [\emph{psen}] 0x6B44: пенсия\\
\textbf{psep} [\emph{psep}] 0x4B44: проницаемый\\
\textbf{pses} [\emph{pses}] 0x8B44: получатель\\
\textbf{pset} [\emph{pset}] 0xAB44: продлить\\
\textbf{psi2c} [\emph{psi˩ʃ}] 0xF844: эпизодический\\
\textbf{psi2h} [\emph{psi˩ʰ}] 0xC227: пластика\\
\textbf{psi2k} [\emph{psi˩k}] 0x3844: перископ\\
\textbf{psi2n} [\emph{psi˩n}] 0x7844: фарфор\\
\textbf{psi2p} [\emph{psi˩p}] 0x5844: епископальный\\
\textbf{psi2s} [\emph{psi˩s}] 0x9844: предлог\\
\textbf{psi2t} [\emph{psi˩t}] 0xB844: наставник\\
\textbf{psi7c} [\emph{psi˥ʃ}] 0xF044: пессарий\\
\textbf{psi7k} [\emph{psi˥k}] 0x3044: апостольский\\
\textbf{psi7n} [\emph{psi˥n}] 0x7044: пасторат\\
\textbf{psi7p} [\emph{psi˥p}] 0x5044: покрытосемянное\\
\textbf{psi7s} [\emph{psi˥s}] 0x9044: эликсир\\
\textbf{psi7t} [\emph{psi˥t}] 0xB044: самоочевидной\\
\textbf{psic} [\emph{psiʃ}] 0xE844: берег\\
\textbf{psif} [\emph{psif}] 0xC844: эмиссия\\
\textbf{psih} [\emph{psiʰ}] 0x4227: optatative mood, for expressing wishes or hopes.\\
\textbf{psik} [\emph{psik}] 0x2844: оратор\\
\textbf{psim} [\emph{psim}] 0x0844: приблизительный\\
\textbf{psin} [\emph{psin}] 0x6844: вес\\
\textbf{psip} [\emph{psip}] 0x4844: принцип\\
\textbf{psis} [\emph{psis}] 0x8844: поэзия\\
\textbf{psit} [\emph{psit}] 0xA844: присутствие\\
\textbf{pso2k} [\emph{pso˩k}] 0x3C44: дошкольного\\
\textbf{pso2n} [\emph{pso˩n}] 0x7C44: непослушание\\
\textbf{pso2t} [\emph{pso˩t}] 0xBC44: спонсируемый\\
\textbf{pso7k} [\emph{pso˥k}] 0x3444: поймать все\\
\textbf{pso7n} [\emph{pso˥n}] 0x7444: столбняк\\
\textbf{pso7s} [\emph{pso˥s}] 0x9444: опоссум\\
\textbf{pso7t} [\emph{pso˥t}] 0xB444: перекись\\
\textbf{psoc} [\emph{psoʃ}] 0xEC44: перфоратор\\
\textbf{psof} [\emph{psof}] 0xCC44: фосфат\\
\textbf{psoh} [\emph{psoʰ}] 0x6227: непобедимый\\
\textbf{psok} [\emph{psok}] 0x2C44: открытка\\
\textbf{psom} [\emph{psom}] 0x0C44: читать только\\
\textbf{pson} [\emph{pson}] 0x6C44: непобедимый\\
\textbf{psop} [\emph{psop}] 0x4C44: паб\\
\textbf{psos} [\emph{psos}] 0x8C44: заграничный пасспорт\\
\textbf{psot} [\emph{psot}] 0xAC44: расплата\\
\textbf{psu2s} [\emph{psu˩s}] 0x9A44: пластика\\
\textbf{psu2t} [\emph{psu˩t}] 0xBA44: непригодность\\
\textbf{psu7k} [\emph{psu˥k}] 0x3244: воспроизводимый\\
\textbf{psu7n} [\emph{psu˥n}] 0x7244: плати как сможешь\\
\textbf{psu7s} [\emph{psu˥s}] 0x9244: метловище\\
\textbf{psu7t} [\emph{psu˥t}] 0xB244: пудель\\
\textbf{psuc} [\emph{psuʃ}] 0xEA44: пластичность\\
\textbf{psuf} [\emph{psuf}] 0xCA44: массив\\
\textbf{psuh} [\emph{psuʰ}] 0x5227: superessive case, indicates location on a surface. surface-context location-case.\\
\textbf{psuk} [\emph{psuk}] 0x2A44: сломанный\\
\textbf{psum} [\emph{psum}] 0x0A44: пемза\\
\textbf{psun} [\emph{psun}] 0x6A44: предзнаменование\\
\textbf{psup} [\emph{psup}] 0x4A44: губная помада\\
\textbf{psus} [\emph{psus}] 0x8A44: петрушка\\
\textbf{psut} [\emph{psut}] 0xAA44: объекты\\
\textbf{pu} [\emph{pu}] 0x289E: множественное число номер\\
\textbf{pw6n} [\emph{pwən}] 0x6D24: планктон\\
\textbf{pw6t} [\emph{pwət}] 0xAD24: пульсировать\\
\textbf{pwa2n} [\emph{pwä˩n}] 0x7924: папский\\
\textbf{pwa2s} [\emph{pwä˩s}] 0x9924: меблировать\\
\textbf{pwa2t} [\emph{pwä˩t}] 0xB924: пагода\\
\textbf{pwa7k} [\emph{pwä˥k}] 0x3124: панический\\
\textbf{pwa7n} [\emph{pwä˥n}] 0x7124: провансальский\\
\textbf{pwa7t} [\emph{pwä˥t}] 0xB124: камбала\\
\textbf{pwac} [\emph{pwäʃ}] 0xE924: камень\\
\textbf{pwaf} [\emph{pwäf}] 0xC924: фисташковый\\
\textbf{pwah} [\emph{pwäʰ}] 0x4927: benefactive case, indicates the noun phrase is the benficiary of the sentence. social-context destination-case.\\
\textbf{pwak} [\emph{pwäk}] 0x2924: коробка\\
\textbf{pwam} [\emph{pwäm}] 0x0924: шпионских программ\\
\textbf{pwan} [\emph{pwän}] 0x6924: весна\\
\textbf{pwap} [\emph{pwäp}] 0x4924: труба\\
\textbf{pwas} [\emph{pwäs}] 0x8924: аплодисменты\\
\textbf{pwat} [\emph{pwät}] 0xA924: постель\\
\textbf{pwe2n} [\emph{pwe˩n}] 0x7B24: парамагнитный\\
\textbf{pwe7n} [\emph{pwe˥n}] 0x7324: пентан\\
\textbf{pwe7t} [\emph{pwe˥t}] 0xB324: подножка\\
\textbf{pwec} [\emph{pweʃ}] 0xEB24: партитивный дело\\
\textbf{pwef} [\emph{pwef}] 0xCB24: подветренный\\
\textbf{pwek} [\emph{pwek}] 0x2B24: pegative дело\\
\textbf{pwem} [\emph{pwem}] 0x0B24: придатки\\
\textbf{pwen} [\emph{pwen}] 0x6B24: пешка\\
\textbf{pwep} [\emph{pwep}] 0x4B24: пипетка\\
\textbf{pwes} [\emph{pwes}] 0x8B24: восемьдесят четыре\\
\textbf{pwet} [\emph{pwet}] 0xAB24: головоломка\\
\textbf{pwi2n} [\emph{pwi˩n}] 0x7824: плевательница\\
\textbf{pwi2t} [\emph{pwi˩t}] 0xB824: предварительно adamite\\
\textbf{pwi7k} [\emph{pwi˥k}] 0x3024: дело всей жизни\\
\textbf{pwi7n} [\emph{pwi˥n}] 0x7024: восемьдесят девять\\
\textbf{pwi7t} [\emph{pwi˥t}] 0xB024: прагматика\\
\textbf{pwic} [\emph{pwiʃ}] 0xE824: держа\\
\textbf{pwif} [\emph{pwif}] 0xC824: предварительный просмотр\\
\textbf{pwih} [\emph{pwiʰ}] 0x4127: ablative case, indicates a source noun phrase, from. is used as a standalone case. space-context source-case\\
\textbf{pwik} [\emph{pwik}] 0x2824: перспективный\\
\textbf{pwim} [\emph{pwim}] 0x0824: неполная занятость\\
\textbf{pwin} [\emph{pwin}] 0x6824: вино\\
\textbf{pwip} [\emph{pwip}] 0x4824: поэт\\
\textbf{pwis} [\emph{pwis}] 0x8824: закуска\\
\textbf{pwit} [\emph{pwit}] 0xA824: просительный настроение\\
\textbf{pwo7t} [\emph{pwo˥t}] 0xB424: запачканный\\
\textbf{pwoc} [\emph{pwoʃ}] 0xEC24: Плутон\\
\textbf{pwok} [\emph{pwok}] 0x2C24: обвинительный дело\\
\textbf{pwom} [\emph{pwom}] 0x0C24: задняя дверь\\
\textbf{pwon} [\emph{pwon}] 0x6C24: запрет\\
\textbf{pwop} [\emph{pwop}] 0x4C24: парабола\\
\textbf{pwos} [\emph{pwos}] 0x8C24: протонов\\
\textbf{pwot} [\emph{pwot}] 0xAC24: белок\\
\textbf{pwu7n} [\emph{pwu˥n}] 0x7224: предварительная оплата\\
\textbf{pwuk} [\emph{pwuk}] 0x2A24: наволочка\\
\textbf{pwun} [\emph{pwun}] 0x6A24: пудинг\\
\textbf{pwus} [\emph{pwus}] 0x8A24: выглядывает\\
\textbf{pwut} [\emph{pwut}] 0xAA24: предопределенный\\
\textbf{px6k} [\emph{pxək}] 0x2DE4: панк\\
\textbf{px6n} [\emph{pxən}] 0x6DE4: плацента\\
\textbf{px6p} [\emph{pxəp}] 0x4DE4: виночерпием\\
\textbf{px6s} [\emph{pxəs}] 0x8DE4: проходки\\
\textbf{px6t} [\emph{pxət}] 0xADE4: лампа накаливания\\
\textbf{pxa2k} [\emph{pxä˩k}] 0x39E4: града молодец\\
\textbf{pxa2n} [\emph{pxä˩n}] 0x79E4: Пэдди\\
\textbf{pxa2s} [\emph{pxä˩s}] 0x99E4: судья на линии\\
\textbf{pxa2t} [\emph{pxä˩t}] 0xB9E4: топ мачты\\
\textbf{pxa7c} [\emph{pxä˥ʃ}] 0xF1E4: перепев\\
\textbf{pxa7f} [\emph{pxä˥f}] 0xD1E4: пентхауз\\
\textbf{pxa7k} [\emph{pxä˥k}] 0x31E4: раскалывать\\
\textbf{pxa7n} [\emph{pxä˥n}] 0x71E4: платонический\\
\textbf{pxa7p} [\emph{pxä˥p}] 0x51E4: после того, как хлопают\\
\textbf{pxa7s} [\emph{pxä˥s}] 0x91E4: неудачник\\
\textbf{pxa7t} [\emph{pxä˥t}] 0xB1E4: неравенства\\
\textbf{pxac} [\emph{pxäʃ}] 0xE9E4: чуткими радость\\
\textbf{pxaf} [\emph{pxäf}] 0xC9E4: метаморфоза\\
\textbf{pxah} [\emph{pxäʰ}] 0x4F27: завлечь\\
\textbf{pxak} [\emph{pxäk}] 0x29E4: завернутый\\
\textbf{pxam} [\emph{pxäm}] 0x09E4: уздечка\\
\textbf{pxan} [\emph{pxän}] 0x69E4: рекламный\\
\textbf{pxap} [\emph{pxäp}] 0x49E4: тополь\\
\textbf{pxas} [\emph{pxäs}] 0x89E4: колпачки замок\\
\textbf{pxat} [\emph{pxät}] 0xA9E4: киоск\\
\textbf{pxe2n} [\emph{pxe˩n}] 0x7BE4: сдоба\\
\textbf{pxe2t} [\emph{pxe˩t}] 0xBBE4: наборная машина\\
\textbf{pxe7n} [\emph{pxe˥n}] 0x73E4: мята\\
\textbf{pxe7t} [\emph{pxe˥t}] 0xB3E4: нечитаемый\\
\textbf{pxec} [\emph{pxeʃ}] 0xEBE4: персик\\
\textbf{pxef} [\emph{pxef}] 0xCBE4: проливной\\
\textbf{pxeh} [\emph{pxeʰ}] 0x5F27: partitive case, a subcase that indicates only partial usage of the noun phrase. lazy argument loading.\\
\textbf{pxek} [\emph{pxek}] 0x2BE4: доисторический\\
\textbf{pxem} [\emph{pxem}] 0x0BE4: раскаленный добела\\
\textbf{pxen} [\emph{pxen}] 0x6BE4: ведомственный\\
\textbf{pxep} [\emph{pxep}] 0x4BE4: панели\\
\textbf{pxes} [\emph{pxes}] 0x8BE4: таз\\
\textbf{pxet} [\emph{pxet}] 0xABE4: плащ\\
\textbf{pxi2k} [\emph{pxi˩k}] 0x38E4: носовое кровотечение\\
\textbf{pxi2n} [\emph{pxi˩n}] 0x78E4: сфинкс\\
\textbf{pxi2s} [\emph{pxi˩s}] 0x98E4: свинарник\\
\textbf{pxi2t} [\emph{pxi˩t}] 0xB8E4: проверенный\\
\textbf{pxi7k} [\emph{pxi˥k}] 0x30E4: паршивый\\
\textbf{pxi7n} [\emph{pxi˥n}] 0x70E4: прозвище\\
\textbf{pxi7s} [\emph{pxi˥s}] 0x90E4: насаженные\\
\textbf{pxi7t} [\emph{pxi˥t}] 0xB0E4: зловонный\\
\textbf{pxic} [\emph{pxiʃ}] 0xE8E4: селезенка\\
\textbf{pxif} [\emph{pxif}] 0xC8E4: полиция сотрудник\\
\textbf{pxih} [\emph{pxiʰ}] 0x4727: истребование\\
\textbf{pxik} [\emph{pxik}] 0x28E4: параюридический\\
\textbf{pxim} [\emph{pxim}] 0x08E4: премьер-министр\\
\textbf{pxin} [\emph{pxin}] 0x68E4: клянусь в\\
\textbf{pxip} [\emph{pxip}] 0x48E4: торф\\
\textbf{pxis} [\emph{pxis}] 0x88E4: пассивный голос\\
\textbf{pxit} [\emph{pxit}] 0xA8E4: IP адрес\\
\textbf{pxo2n} [\emph{pxo˩n}] 0x7CE4: полигональный\\
\textbf{pxo2t} [\emph{pxo˩t}] 0xBCE4: периодичность\\
\textbf{pxo7n} [\emph{pxo˥n}] 0x74E4: Копенгаген\\
\textbf{pxo7t} [\emph{pxo˥t}] 0xB4E4: пятисотый\\
\textbf{pxoc} [\emph{pxoʃ}] 0xECE4: за предложение\\
\textbf{pxof} [\emph{pxof}] 0xCCE4: полиморфный\\
\textbf{pxoh} [\emph{pxoʰ}] 0x6727: prolative case, indicates motion via an interior. interior-context way-case.\\
\textbf{pxok} [\emph{pxok}] 0x2CE4: пастораль\\
\textbf{pxom} [\emph{pxom}] 0x0CE4: propositive настроение\\
\textbf{pxon} [\emph{pxon}] 0x6CE4: недоступен\\
\textbf{pxop} [\emph{pxop}] 0x4CE4: опал\\
\textbf{pxos} [\emph{pxos}] 0x8CE4: Почта\\
\textbf{pxot} [\emph{pxot}] 0xACE4: из первых рук\\
\textbf{pxu7k} [\emph{pxu˥k}] 0x32E4: Непропеченный\\
\textbf{pxu7t} [\emph{pxu˥t}] 0xB2E4: шпатлевка\\
\textbf{pxuc} [\emph{pxuʃ}] 0xEAE4: пачули\\
\textbf{pxuk} [\emph{pxuk}] 0x2AE4: Португалия\\
\textbf{pxum} [\emph{pxum}] 0x0AE4: непрозрачный\\
\textbf{pxun} [\emph{pxun}] 0x6AE4: плутоний\\
\textbf{pxup} [\emph{pxup}] 0x4AE4: сутенер\\
\textbf{pxus} [\emph{pxus}] 0x8AE4: пестик\\
\textbf{pxut} [\emph{pxut}] 0xAAE4: пернатых\\
\textbf{py67t} [\emph{pjə˥t}] 0xB504: брамин\\
\textbf{py6k} [\emph{pjək}] 0x2D04: папайя\\
\textbf{py6n} [\emph{pjən}] 0x6D04: пион\\
\textbf{py6s} [\emph{pjəs}] 0x8D04: поверхностно\\
\textbf{py6t} [\emph{pjət}] 0xAD04: Эритрея\\
\textbf{pya2k} [\emph{pjä˩k}] 0x3904: педагогический\\
\textbf{pya2n} [\emph{pjä˩n}] 0x7904: проституция\\
\textbf{pya2p} [\emph{pjä˩p}] 0x5904: назад ткань\\
\textbf{pya2t} [\emph{pjä˩t}] 0xB904: торговец шляпами\\
\textbf{pya7k} [\emph{pjä˥k}] 0x3104: шлепок\\
\textbf{pya7m} [\emph{pjä˥m}] 0x1104: панамский\\
\textbf{pya7n} [\emph{pjä˥n}] 0x7104: восемьдесят один\\
\textbf{pya7t} [\emph{pjä˥t}] 0xB104: капер\\
\textbf{pyac} [\emph{pjäʃ}] 0xE904: язык\\
\textbf{pyaf} [\emph{pjäf}] 0xC904: птица\\
\textbf{pyah} [\emph{pjäʰ}] 0x4827: love gender, 4th density life form, learning to forgive and love all, akin to Jesus when he was on Earth, Bodhisattva, Hatonn, Gray aliens\\
\textbf{pyak} [\emph{pjäk}] 0x2904: брат\\
\textbf{pyam} [\emph{pjäm}] 0x0904: Храбрый\\
\textbf{pyan} [\emph{pjän}] 0x6904: люблю\\
\textbf{pyap} [\emph{pjäp}] 0x4904: сдачи\\
\textbf{pyas} [\emph{pjäs}] 0x8904: фут\\
\textbf{pyat} [\emph{pjät}] 0xA904: друг\\
\textbf{pye2t} [\emph{pje˩t}] 0xBB04: перитонит\\
\textbf{pye7n} [\emph{pje˥n}] 0x7304: пенсне\\
\textbf{pye7t} [\emph{pje˥t}] 0xB304: сдерживаемый\\
\textbf{pyec} [\emph{pjeʃ}] 0xEB04: таблица\\
\textbf{pyef} [\emph{pjef}] 0xCB04: перестановка\\
\textbf{pyeh} [\emph{pjeʰ}] 0x5827: предикативный\\
\textbf{pyek} [\emph{pjek}] 0x2B04: проектор\\
\textbf{pyem} [\emph{pjem}] 0x0B04: парфюмерия\\
\textbf{pyen} [\emph{pjen}] 0x6B04: Приложения\\
\textbf{pyep} [\emph{pjep}] 0x4B04: детский манеж\\
\textbf{pyes} [\emph{pjes}] 0x8B04: гнейс\\
\textbf{pyet} [\emph{pjet}] 0xAB04: принтер\\
\textbf{pyi2k} [\emph{pji˩k}] 0x3804: тюряга\\
\textbf{pyi2n} [\emph{pji˩n}] 0x7804: ипохондрик\\
\textbf{pyi2t} [\emph{pji˩t}] 0xB804: преднамеренность\\
\textbf{pyi7f} [\emph{pji˥f}] 0xD004: пирогов глазами\\
\textbf{pyi7h} [\emph{pji˥ʰ}] 0x8027: аппликативны голос\\
\textbf{pyi7k} [\emph{pji˥k}] 0x3004: пиридин\\
\textbf{pyi7n} [\emph{pji˥n}] 0x7004: штифты\\
\textbf{pyi7p} [\emph{pji˥p}] 0x5004: прототипичный\\
\textbf{pyi7s} [\emph{pji˥s}] 0x9004: дарохранительница\\
\textbf{pyi7t} [\emph{pji˥t}] 0xB004: пол пинты\\
\textbf{pyic} [\emph{pjiʃ}] 0xE804: обязательный\\
\textbf{pyif} [\emph{pjif}] 0xC804: неопределенные\\
\textbf{pyih} [\emph{pjiʰ}] 0x4027: мир\\
\textbf{pyik} [\emph{pjik}] 0x2804: поезд\\
\textbf{pyim} [\emph{pjim}] 0x0804: платина\\
\textbf{pyin} [\emph{pjin}] 0x6804: дикий\\
\textbf{pyip} [\emph{pjip}] 0x4804: апрель\\
\textbf{pyis} [\emph{pjis}] 0x8804: священник\\
\textbf{pyit} [\emph{pjit}] 0xA804: вежливый регистр\\
\textbf{pyo2n} [\emph{pjo˩n}] 0x7C04: пагинация\\
\textbf{pyo2t} [\emph{pjo˩t}] 0xBC04: помадить\\
\textbf{pyo7n} [\emph{pjo˥n}] 0x7404: пятиугольный\\
\textbf{pyo7t} [\emph{pjo˥t}] 0xB404: неспортивный\\
\textbf{pyoc} [\emph{pjoʃ}] 0xEC04: Полинезия\\
\textbf{pyoh} [\emph{pjoʰ}] 0x6027: пассивный голос\\
\textbf{pyok} [\emph{pjok}] 0x2C04: снаряд\\
\textbf{pyom} [\emph{pjom}] 0x0C04: полоний атом\\
\textbf{pyon} [\emph{pjon}] 0x6C04: пианино\\
\textbf{pyop} [\emph{pjop}] 0x4C04: плей-офф\\
\textbf{pyos} [\emph{pjos}] 0x8C04: премолярный\\
\textbf{pyot} [\emph{pjot}] 0xAC04: голубь\\
\textbf{pyu7t} [\emph{pju˥t}] 0xB204: зуд\\
\textbf{pyuc} [\emph{pjuʃ}] 0xEA04: истребление населения\\
\textbf{pyuh} [\emph{pjuʰ}] 0x5027: но\\
\textbf{pyuk} [\emph{pjuk}] 0x2A04: чума\\
\textbf{pyum} [\emph{pjum}] 0x0A04: окунать\\
\textbf{pyun} [\emph{pjun}] 0x6A04: коричневый\\
\textbf{pyup} [\emph{pjup}] 0x4A04: папуасов\\
\textbf{pyus} [\emph{pjus}] 0x8A04: предохранитель\\
\textbf{pyut} [\emph{pjut}] 0xAA04: север\\
\subsection{ q }
\textbf{qa} [\emph{ŋä}] 0x273E: cessative aspect, for referring to the end of an action, when it is ceasing. inchoative-aspect- is begining\\
\textbf{qe} [\emph{ŋe}] 0x2F3E: star gender, 10th density life form, akin to spirit of a star, Sol, Sirius\\
\textbf{qi} [\emph{ŋi}] 0x233E: переходный\\
\textbf{qo} [\emph{ŋo}] 0x333E: радиоактивность\\
\textbf{qr6n} [\emph{ŋrən}] 0x6DD9: андеррайтер\\
\textbf{qr6t} [\emph{ŋrət}] 0xADD9: транспортир\\
\textbf{qra2n} [\emph{ŋrä˩n}] 0x79D9: тиранить\\
\textbf{qra2t} [\emph{ŋrä˩t}] 0xB9D9: расположить к себе\\
\textbf{qra7k} [\emph{ŋrä˥k}] 0x31D9: гнездо взломщик\\
\textbf{qra7n} [\emph{ŋrä˥n}] 0x71D9: бикарбонат\\
\textbf{qra7t} [\emph{ŋrä˥t}] 0xB1D9: властитель\\
\textbf{qrac} [\emph{ŋräʃ}] 0xE9D9: вспашка\\
\textbf{qraf} [\emph{ŋräf}] 0xC9D9: неблагодарный\\
\textbf{qrah} [\emph{ŋräʰ}] 0x4ECF: текущий актуальность маркер\\
\textbf{qrak} [\emph{ŋräk}] 0x29D9: штрейкбрехер\\
\textbf{qram} [\emph{ŋräm}] 0x09D9: реквизит\\
\textbf{qran} [\emph{ŋrän}] 0x69D9: Коран\\
\textbf{qrap} [\emph{ŋräp}] 0x49D9: небоскреб\\
\textbf{qras} [\emph{ŋräs}] 0x89D9: начальство\\
\textbf{qrat} [\emph{ŋrät}] 0xA9D9: вверх ногами\\
\textbf{qrec} [\emph{ŋreʃ}] 0xEBD9: поджелудочная железа\\
\textbf{qrek} [\emph{ŋrek}] 0x2BD9: завышенный\\
\textbf{qren} [\emph{ŋren}] 0x6BD9: планетарный Пол\\
\textbf{qres} [\emph{ŋres}] 0x8BD9: внушающих благоговение\\
\textbf{qret} [\emph{ŋret}] 0xABD9: кощунство\\
\textbf{qri2n} [\emph{ŋri˩n}] 0x78D9: скарлатина\\
\textbf{qri2s} [\emph{ŋri˩s}] 0x98D9: цирроз печени\\
\textbf{qri2t} [\emph{ŋri˩t}] 0xB8D9: индивидуальные\\
\textbf{qri7n} [\emph{ŋri˥n}] 0x70D9: мариновать\\
\textbf{qri7t} [\emph{ŋri˥t}] 0xB0D9: пирит\\
\textbf{qric} [\emph{ŋriʃ}] 0xE8D9: ирландцы\\
\textbf{qrif} [\emph{ŋrif}] 0xC8D9: рыба-меч\\
\textbf{qrih} [\emph{ŋriʰ}] 0x46CF: мечтательность\\
\textbf{qrik} [\emph{ŋrik}] 0x28D9: центральный дело\\
\textbf{qrim} [\emph{ŋrim}] 0x08D9: рудимент\\
\textbf{qrin} [\emph{ŋrin}] 0x68D9: мнимая единица -\\
\textbf{qrip} [\emph{ŋrip}] 0x48D9: крипта\\
\textbf{qris} [\emph{ŋris}] 0x88D9: английский\\
\textbf{qrit} [\emph{ŋrit}] 0xA8D9: Скриншот\\
\textbf{qro7t} [\emph{ŋro˥t}] 0xB4D9: Несчастные\\
\textbf{qroc} [\emph{ŋroʃ}] 0xECD9: Джордж\\
\textbf{qrok} [\emph{ŋrok}] 0x2CD9: нет до\\
\textbf{qron} [\emph{ŋron}] 0x6CD9: гортань\\
\textbf{qrop} [\emph{ŋrop}] 0x4CD9: агорафобия\\
\textbf{qros} [\emph{ŋros}] 0x8CD9: поддерживать\\
\textbf{qrot} [\emph{ŋrot}] 0xACD9: нейтрон\\
\textbf{qruk} [\emph{ŋruk}] 0x2AD9: надстройка\\
\textbf{qrun} [\emph{ŋrun}] 0x6AD9: украшенный\\
\textbf{qrus} [\emph{ŋrus}] 0x8AD9: евхаристия\\
\textbf{qrut} [\emph{ŋrut}] 0xAAD9: попугай\\
\textbf{qwac} [\emph{ŋwäʃ}] 0xE939: Бог оставил\\
\textbf{qwah} [\emph{ŋwäʰ}] 0x49CF: неамбициозны\\
\textbf{qwak} [\emph{ŋwäk}] 0x2939: чайка\\
\textbf{qwam} [\emph{ŋwäm}] 0x0939: празднование новоселья\\
\textbf{qwan} [\emph{ŋwän}] 0x6939: часто спросил вопросов\\
\textbf{qwap} [\emph{ŋwäp}] 0x4939: искусственное освещение\\
\textbf{qwas} [\emph{ŋwäs}] 0x8939: наружу\\
\textbf{qwat} [\emph{ŋwät}] 0xA939: количество из вещество\\
\textbf{qwek} [\emph{ŋwek}] 0x2B39: кормилец\\
\textbf{qwen} [\emph{ŋwen}] 0x6B39: Эстония\\
\textbf{qwes} [\emph{ŋwes}] 0x8B39: Ви Ви\\
\textbf{qwet} [\emph{ŋwet}] 0xAB39: конский хвост\\
\textbf{qwi2t} [\emph{ŋwi˩t}] 0xB839: цистит\\
\textbf{qwik} [\emph{ŋwik}] 0x2839: пневматический\\
\textbf{qwin} [\emph{ŋwin}] 0x6839: пингвин\\
\textbf{qwis} [\emph{ŋwis}] 0x8839: повесток\\
\textbf{qwit} [\emph{ŋwit}] 0xA839: антисептик\\
\textbf{qwok} [\emph{ŋwok}] 0x2C39: ангольский\\
\textbf{qwon} [\emph{ŋwon}] 0x6C39: отчужденность\\
\textbf{qwot} [\emph{ŋwot}] 0xAC39: андорра\\
\textbf{qwut} [\emph{ŋwut}] 0xAA39: площади сфальсифицированы\\
\textbf{qy6n} [\emph{ŋjən}] 0x6D19: пасека\\
\textbf{qya2t} [\emph{ŋjä˩t}] 0xB919: воздушная почта\\
\textbf{qya7n} [\emph{ŋjä˥n}] 0x7119: стенокардия\\
\textbf{qya7t} [\emph{ŋjä˥t}] 0xB119: издавна\\
\textbf{qyac} [\emph{ŋjäʃ}] 0xE919: двухгодичный\\
\textbf{qyaf} [\emph{ŋjäf}] 0xC919: коммерциализация\\
\textbf{qyah} [\emph{ŋjäʰ}] 0x48CF: озабоченность\\
\textbf{qyak} [\emph{ŋjäk}] 0x2919: Один действие глагол\\
\textbf{qyam} [\emph{ŋjäm}] 0x0919: анемометр\\
\textbf{qyan} [\emph{ŋjän}] 0x6919: дальтонизм\\
\textbf{qyap} [\emph{ŋjäp}] 0x4919: хэтчбэк\\
\textbf{qyas} [\emph{ŋjäs}] 0x8919: алхимик\\
\textbf{qyat} [\emph{ŋjät}] 0xA919: песчаник\\
\textbf{qyeh} [\emph{ŋjeʰ}] 0x58CF: Один действие глагол\\
\textbf{qyek} [\emph{ŋjek}] 0x2B19: эротика\\
\textbf{qyen} [\emph{ŋjen}] 0x6B19: прорастание\\
\textbf{qyes} [\emph{ŋjes}] 0x8B19: флебит\\
\textbf{qyet} [\emph{ŋjet}] 0xAB19: подвергать опасности\\
\textbf{qyi7n} [\emph{ŋji˥n}] 0x7019: амин\\
\textbf{qyi7t} [\emph{ŋji˥t}] 0xB019: эмпирик\\
\textbf{qyic} [\emph{ŋjiʃ}] 0xE819: психиатрия\\
\textbf{qyih} [\emph{ŋjiʰ}] 0x40CF: оживить Пол\\
\textbf{qyik} [\emph{ŋjik}] 0x2819: соединительный частица\\
\textbf{qyim} [\emph{ŋjim}] 0x0819: тимин\\
\textbf{qyin} [\emph{ŋjin}] 0x6819: биоразнообразия\\
\textbf{qyip} [\emph{ŋjip}] 0x4819: Покойся с миром\\
\textbf{qyis} [\emph{ŋjis}] 0x8819: бизнес человек\\
\textbf{qyit} [\emph{ŋjit}] 0xA819: критическая точка\\
\textbf{qyon} [\emph{ŋjon}] 0x6C19: уравнения\\
\textbf{qyot} [\emph{ŋjot}] 0xAC19: солнцестояние\\
\textbf{qyuk} [\emph{ŋjuk}] 0x2A19: равноугольный\\
\textbf{qyut} [\emph{ŋjut}] 0xAA19: фунгицид\\
\subsection{ r }
\textbf{r6} [\emph{rə}] 0x35DE: причастие\\
\textbf{ra} [\emph{rä}] 0x25DE: настоящее время напряженный\\
\textbf{ra2} [\emph{rä˩}] 0x65DE: frequentative aspect, indicates the verb happens multiple times, for example bite frequentative-aspect is equivalent to chewing. \\
\textbf{re} [\emph{re}] 0x2DDE: другой предмет\\
\textbf{ri} [\emph{ri}] 0x21DE: interrogative mood, for questions.\\
\textbf{ro} [\emph{ro}] 0x31DE: государство контекст\\
\textbf{ru} [\emph{ru}] 0x29DE: prohibitive mood, for expressing something is not allowed, inverse of permissive-mood-, potential response to admonitive mood\\
\textbf{rw6k} [\emph{rwək}] 0x2D2E: барвинок малый\\
\textbf{rw6n} [\emph{rwən}] 0x6D2E: изменчивый\\
\textbf{rw6t} [\emph{rwət}] 0xAD2E: термостат\\
\textbf{rwa2k} [\emph{rwä˩k}] 0x392E: афродизиак\\
\textbf{rwa2n} [\emph{rwä˩n}] 0x792E: царствовал\\
\textbf{rwa2s} [\emph{rwä˩s}] 0x992E: трещотка\\
\textbf{rwa2t} [\emph{rwä˩t}] 0xB92E: арьергард\\
\textbf{rwa7k} [\emph{rwä˥k}] 0x312E: покушаться\\
\textbf{rwa7n} [\emph{rwä˥n}] 0x712E: Рангун\\
\textbf{rwa7p} [\emph{rwä˥p}] 0x512E: описки\\
\textbf{rwa7s} [\emph{rwä˥s}] 0x912E: карикатурист\\
\textbf{rwa7t} [\emph{rwä˥t}] 0xB12E: увлеченный\\
\textbf{rwac} [\emph{rwäʃ}] 0xE92E: аппаратный\\
\textbf{rwaf} [\emph{rwäf}] 0xC92E: хризантема\\
\textbf{rwah} [\emph{rwäʰ}] 0x4977: prospective aspect, the event happens after the other things under consideration. \\
\textbf{rwak} [\emph{rwäk}] 0x292E: трещина\\
\textbf{rwam} [\emph{rwäm}] 0x092E: возмещать\\
\textbf{rwan} [\emph{rwän}] 0x692E: награда\\
\textbf{rwap} [\emph{rwäp}] 0x492E: ежеквартальный\\
\textbf{rwas} [\emph{rwäs}] 0x892E: роман\\
\textbf{rwat} [\emph{rwät}] 0xA92E: галстук\\
\textbf{rwe2k} [\emph{rwe˩k}] 0x3B2E: переделывать\\
\textbf{rwe2n} [\emph{rwe˩n}] 0x7B2E: четырехмерный\\
\textbf{rwe2t} [\emph{rwe˩t}] 0xBB2E: схема\\
\textbf{rwe7t} [\emph{rwe˥t}] 0xB32E: сосед по комнате\\
\textbf{rwec} [\emph{rweʃ}] 0xEB2E: рейнджер\\
\textbf{rwef} [\emph{rwef}] 0xCB2E: оборотни\\
\textbf{rweh} [\emph{rweʰ}] 0x5977: oblique case, indicates intersection with noun phrase, such as an orbit or street that is being crossed at an oblique angle.\\
\textbf{rwek} [\emph{rwek}] 0x2B2E: ретроспективный\\
\textbf{rwem} [\emph{rwem}] 0x0B2E: красное дерево\\
\textbf{rwen} [\emph{rwen}] 0x6B2E: предотвращение\\
\textbf{rwep} [\emph{rwep}] 0x4B2E: сорок восемь\\
\textbf{rwes} [\emph{rwes}] 0x8B2E: изотопный индикатор\\
\textbf{rwet} [\emph{rwet}] 0xAB2E: реалистичный\\
\textbf{rwi2n} [\emph{rwi˩n}] 0x782E: реквизиция\\
\textbf{rwi2t} [\emph{rwi˩t}] 0xB82E: переделанный\\
\textbf{rwi7n} [\emph{rwi˥n}] 0x702E: радужный\\
\textbf{rwi7t} [\emph{rwi˥t}] 0xB02E: рассказчик\\
\textbf{rwic} [\emph{rwiʃ}] 0xE82E: радуга\\
\textbf{rwif} [\emph{rwif}] 0xC82E: сорок четыре\\
\textbf{rwih} [\emph{rwiʰ}] 0x4177: precative mood, for requests.\\
\textbf{rwik} [\emph{rwik}] 0x282E: вертикальный\\
\textbf{rwim} [\emph{rwim}] 0x082E: соответственно\\
\textbf{rwin} [\emph{rwin}] 0x682E: Королева\\
\textbf{rwip} [\emph{rwip}] 0x482E: регби\\
\textbf{rwis} [\emph{rwis}] 0x882E: расовой\\
\textbf{rwit} [\emph{rwit}] 0xA82E: соперник\\
\textbf{rwo2n} [\emph{rwo˩n}] 0x7C2E: ретороманский диалект\\
\textbf{rwo2t} [\emph{rwo˩t}] 0xBC2E: броненосец\\
\textbf{rwo7t} [\emph{rwo˥t}] 0xB42E: красный верх\\
\textbf{rwoc} [\emph{rwoʃ}] 0xEC2E: последовательный\\
\textbf{rwof} [\emph{rwof}] 0xCC2E: сток\\
\textbf{rwoh} [\emph{rwoʰ}] 0x6177: движение в гору\\
\textbf{rwok} [\emph{rwok}] 0x2C2E: провоцировать\\
\textbf{rwom} [\emph{rwom}] 0x0C2E: хромат\\
\textbf{rwon} [\emph{rwon}] 0x6C2E: напоминание\\
\textbf{rwop} [\emph{rwop}] 0x4C2E: коротковолновый\\
\textbf{rwos} [\emph{rwos}] 0x8C2E: аэрозоль\\
\textbf{rwot} [\emph{rwot}] 0xAC2E: орбита\\
\textbf{rwu7n} [\emph{rwu˥n}] 0x722E: руна\\
\textbf{rwuc} [\emph{rwuʃ}] 0xEA2E: порядковый\\
\textbf{rwuh} [\emph{rwuʰ}] 0x5177: ревнивый\\
\textbf{rwuk} [\emph{rwuk}] 0x2A2E: Уругвай\\
\textbf{rwum} [\emph{rwum}] 0x0A2E: розовое дерево\\
\textbf{rwun} [\emph{rwun}] 0x6A2E: критерий\\
\textbf{rwus} [\emph{rwus}] 0x8A2E: гребец\\
\textbf{rwut} [\emph{rwut}] 0xAA2E: структурный\\
\textbf{ry67n} [\emph{rjə˥n}] 0x750E: Райнланд\\
\textbf{ry67t} [\emph{rjə˥t}] 0xB50E: Рияд\\
\textbf{ry6h} [\emph{rjəʰ}] 0x6877: радиан\\
\textbf{ry6n} [\emph{rjən}] 0x6D0E: иррациональный\\
\textbf{ry6s} [\emph{rjəs}] 0x8D0E: Киргизия\\
\textbf{ry6t} [\emph{rjət}] 0xAD0E: орхидея\\
\textbf{rya2c} [\emph{rjä˩ʃ}] 0xF90E: рикша\\
\textbf{rya2k} [\emph{rjä˩k}] 0x390E: Рея\\
\textbf{rya2n} [\emph{rjä˩n}] 0x790E: грациозность\\
\textbf{rya2s} [\emph{rjä˩s}] 0x990E: Королевская битва\\
\textbf{rya2t} [\emph{rjä˩t}] 0xB90E: уважаемый\\
\textbf{rya7k} [\emph{rjä˥k}] 0x310E: арак\\
\textbf{rya7m} [\emph{rjä˥m}] 0x110E: арамейский\\
\textbf{rya7n} [\emph{rjä˥n}] 0x710E: арийский\\
\textbf{rya7p} [\emph{rjä˥p}] 0x510E: расколачивающий\\
\textbf{rya7s} [\emph{rjä˥s}] 0x910E: английская лапта\\
\textbf{rya7t} [\emph{rjä˥t}] 0xB10E: пролетариат\\
\textbf{ryac} [\emph{rjäʃ}] 0xE90E: натирать\\
\textbf{ryaf} [\emph{rjäf}] 0xC90E: креативность\\
\textbf{ryah} [\emph{rjäʰ}] 0x4877: смех\\
\textbf{ryak} [\emph{rjäk}] 0x290E: реакция\\
\textbf{ryam} [\emph{rjäm}] 0x090E: местоимение маркер\\
\textbf{ryan} [\emph{rjän}] 0x690E: перевод\\
\textbf{ryap} [\emph{rjäp}] 0x490E: перестраивать\\
\textbf{ryas} [\emph{rjäs}] 0x890E: коробка передач\\
\textbf{ryat} [\emph{rjät}] 0xA90E: трагедия\\
\textbf{rye2n} [\emph{rje˩n}] 0x7B0E: заокеанский\\
\textbf{rye2t} [\emph{rje˩t}] 0xBB0E: повторного входа\\
\textbf{rye7n} [\emph{rje˥n}] 0x730E: зайца крикливый\\
\textbf{rye7t} [\emph{rje˥t}] 0xB30E: серенада\\
\textbf{ryec} [\emph{rjeʃ}] 0xEB0E: Розмари\\
\textbf{ryeh} [\emph{rjeʰ}] 0x5877: ретроспективный\\
\textbf{ryek} [\emph{rjek}] 0x2B0E: ракета\\
\textbf{ryem} [\emph{rjem}] 0x0B0E: радий\\
\textbf{ryen} [\emph{rjen}] 0x6B0E: несмотря на\\
\textbf{ryep} [\emph{rjep}] 0x4B0E: клюква\\
\textbf{ryes} [\emph{rjes}] 0x8B0E: регистратор\\
\textbf{ryet} [\emph{rjet}] 0xAB0E: выпускник\\
\textbf{ryi2k} [\emph{rji˩k}] 0x380E: сердцевидный\\
\textbf{ryi2n} [\emph{rji˩n}] 0x780E: эркер\\
\textbf{ryi2t} [\emph{rji˩t}] 0xB80E: сорок два\\
\textbf{ryi7k} [\emph{rji˥k}] 0x300E: эректильная\\
\textbf{ryi7m} [\emph{rji˥m}] 0x100E: рацемический\\
\textbf{ryi7n} [\emph{rji˥n}] 0x700E: прихожанка\\
\textbf{ryi7p} [\emph{rji˥p}] 0x500E: пресбиопия\\
\textbf{ryi7s} [\emph{rji˥s}] 0x900E: проверяемый\\
\textbf{ryi7t} [\emph{rji˥t}] 0xB00E: беспричинный\\
\textbf{ryic} [\emph{rjiʃ}] 0xE80E: сельскохозяйственное\\
\textbf{ryif} [\emph{rjif}] 0xC80E: радиальный\\
\textbf{ryih} [\emph{rjiʰ}] 0x4077: доказательный\\
\textbf{ryik} [\emph{rjik}] 0x280E: химическая\\
\textbf{ryim} [\emph{rjim}] 0x080E: ритм\\
\textbf{ryin} [\emph{rjin}] 0x680E: обучение\\
\textbf{ryip} [\emph{rjip}] 0x480E: пропорциональный\\
\textbf{ryis} [\emph{rjis}] 0x880E: отшельник\\
\textbf{ryit} [\emph{rjit}] 0xA80E: квитанция\\
\textbf{ryo2k} [\emph{rjo˩k}] 0x3C0E: Ариохом\\
\textbf{ryo2t} [\emph{rjo˩t}] 0xBC0E: женское общество\\
\textbf{ryo7n} [\emph{rjo˥n}] 0x740E: радикальный\\
\textbf{ryo7s} [\emph{rjo˥s}] 0x940E: ориентализм\\
\textbf{ryo7t} [\emph{rjo˥t}] 0xB40E: рытвины\\
\textbf{ryoc} [\emph{rjoʃ}] 0xEC0E: радиоактивность\\
\textbf{ryof} [\emph{rjof}] 0xCC0E: иволга\\
\textbf{ryoh} [\emph{rjoʰ}] 0x6077: prosecutive case, indicates motion along a surface. surface-context way-case.\\
\textbf{ryok} [\emph{rjok}] 0x2C0E: протокол\\
\textbf{ryom} [\emph{rjom}] 0x0C0E: иридий\\
\textbf{ryon} [\emph{rjon}] 0x6C0E: рациональный\\
\textbf{ryop} [\emph{rjop}] 0x4C0E: робот\\
\textbf{ryos} [\emph{rjos}] 0x8C0E: ахроматический\\
\textbf{ryot} [\emph{rjot}] 0xAC0E: благоговение\\
\textbf{ryu2n} [\emph{rju˩n}] 0x7A0E: рутений\\
\textbf{ryu2t} [\emph{rju˩t}] 0xBA0E: тарантул\\
\textbf{ryu7n} [\emph{rju˥n}] 0x720E: Веселая карусель\\
\textbf{ryuc} [\emph{rjuʃ}] 0xEA0E: преобразователь\\
\textbf{ryuf} [\emph{rjuf}] 0xCA0E: европеизация\\
\textbf{ryuh} [\emph{rjuʰ}] 0x5077: порядковый\\
\textbf{ryuk} [\emph{rjuk}] 0x2A0E: енот\\
\textbf{ryum} [\emph{rjum}] 0x0A0E: уран\\
\textbf{ryun} [\emph{rjun}] 0x6A0E: регулирование\\
\textbf{ryup} [\emph{rjup}] 0x4A0E: повторное использование\\
\textbf{ryus} [\emph{rjus}] 0x8A0E: Россия\\
\textbf{ryut} [\emph{rjut}] 0xAA0E: ревнивый\\
\subsection{ s }
\textbf{s6} [\emph{sə}] 0x353E: crastinal tense, for expressing tomorrow.\\
\textbf{sa} [\emph{sä}] 0x253E: вы\\
\textbf{sa7} [\emph{sä˥}] 0x453E: associative case, indicates the noun phrase is an associated group of the agent. social-context location-case.\\
\textbf{se} [\emph{se}] 0x2D3E: время контекст\\
\textbf{si} [\emph{si}] 0x213E: epistemic mood, to indicate that something is possible.\\
\textbf{si2} [\emph{si˩}] 0x613E: necessitative настроение\\
\textbf{si7} [\emph{si˥}] 0x413E: superlative case, indicates motion above a surface. surface-context way-case.\\
\textbf{sl62n} [\emph{slə˩n}] 0x7D69: сибилянт\\
\textbf{sl62t} [\emph{slə˩t}] 0xBD69: совершеннолетие\\
\textbf{sl6c} [\emph{sləʃ}] 0xED69: фоссилизация\\
\textbf{sl6h} [\emph{sləʰ}] 0x6B4F: цельсию\\
\textbf{sl6k} [\emph{slək}] 0x2D69: султанат\\
\textbf{sl6m} [\emph{sləm}] 0x0D69: манная крупа\\
\textbf{sl6n} [\emph{slən}] 0x6D69: водопровод\\
\textbf{sl6p} [\emph{sləp}] 0x4D69: несбалансированность\\
\textbf{sl6s} [\emph{sləs}] 0x8D69: цельсию\\
\textbf{sl6t} [\emph{slət}] 0xAD69: crastinal напряженный\\
\textbf{sla2c} [\emph{slä˩ʃ}] 0xF969: скотство\\
\textbf{sla2f} [\emph{slä˩f}] 0xD969: окисляться\\
\textbf{sla2k} [\emph{slä˩k}] 0x3969: высшее общество\\
\textbf{sla2m} [\emph{slä˩m}] 0x1969: псалтырь\\
\textbf{sla2n} [\emph{slä˩n}] 0x7969: вскользь\\
\textbf{sla2p} [\emph{slä˩p}] 0x5969: селитра\\
\textbf{sla2s} [\emph{slä˩s}] 0x9969: водовод\\
\textbf{sla2t} [\emph{slä˩t}] 0xB969: принадлежит\\
\textbf{sla7c} [\emph{slä˥ʃ}] 0xF169: равносторонний\\
\textbf{sla7f} [\emph{slä˥f}] 0xD169: славянский\\
\textbf{sla7k} [\emph{slä˥k}] 0x3169: законопослушный\\
\textbf{sla7m} [\emph{slä˥m}] 0x1169: всмятку\\
\textbf{sla7n} [\emph{slä˥n}] 0x7169: старейшины\\
\textbf{sla7p} [\emph{slä˥p}] 0x5169: самбо\\
\textbf{sla7s} [\emph{slä˥s}] 0x9169: обжора\\
\textbf{sla7t} [\emph{slä˥t}] 0xB169: защелка\\
\textbf{slac} [\emph{släʃ}] 0xE969: Бег\\
\textbf{slaf} [\emph{släf}] 0xC969: Социальное сетей\\
\textbf{slah} [\emph{släʰ}] 0x4B4F: беда\\
\textbf{slak} [\emph{släk}] 0x2969: успех\\
\textbf{slam} [\emph{släm}] 0x0969: лосось\\
\textbf{slan} [\emph{slän}] 0x6969: следовать\\
\textbf{slap} [\emph{släp}] 0x4969: потерпеть неудачу\\
\textbf{slas} [\emph{släs}] 0x8969: Изобразительное искусство\\
\textbf{slat} [\emph{slät}] 0xA969: холм\\
\textbf{sle2n} [\emph{sle˩n}] 0x7B69: свинопас\\
\textbf{sle2t} [\emph{sle˩t}] 0xBB69: оксалат\\
\textbf{sle7n} [\emph{sle˥n}] 0x7369: столетний\\
\textbf{sle7t} [\emph{sle˥t}] 0xB369: омела белая\\
\textbf{slec} [\emph{sleʃ}] 0xEB69: исследователь\\
\textbf{slef} [\emph{slef}] 0xCB69: шифер\\
\textbf{sleh} [\emph{sleʰ}] 0x5B4F: vegetal gender, 2nd density life form, learning to fulfill desires, akin to plant, fungi, basic electronics\\
\textbf{slek} [\emph{slek}] 0x2B69: велосипед\\
\textbf{slem} [\emph{slem}] 0x0B69: сублимировать\\
\textbf{slen} [\emph{slen}] 0x6B69: 11\\
\textbf{slep} [\emph{slep}] 0x4B69: сатир\\
\textbf{sles} [\emph{sles}] 0x8B69: цензор\\
\textbf{slet} [\emph{slet}] 0xAB69: живая изгородь\\
\textbf{sli2c} [\emph{sli˩ʃ}] 0xF869: сингалец\\
\textbf{sli2k} [\emph{sli˩k}] 0x3869: силикат\\
\textbf{sli2m} [\emph{sli˩m}] 0x1869: солецизм\\
\textbf{sli2n} [\emph{sli˩n}] 0x7869: самонадеянным\\
\textbf{sli2p} [\emph{sli˩p}] 0x5869: несколько миллиардов\\
\textbf{sli2s} [\emph{sli˩s}] 0x9869: двадцать шесть\\
\textbf{sli7c} [\emph{sli˥ʃ}] 0xF069: гистология\\
\textbf{sli7f} [\emph{sli˥f}] 0xD069: пан славянство\\
\textbf{sli7k} [\emph{sli˥k}] 0x3069: славословие\\
\textbf{sli7m} [\emph{sli˥m}] 0x1069: эпистемология\\
\textbf{sli7n} [\emph{sli˥n}] 0x7069: властители\\
\textbf{sli7p} [\emph{sli˥p}] 0x5069: слоги\\
\textbf{sli7s} [\emph{sli˥s}] 0x9069: пятьдесят шесть\\
\textbf{sli7t} [\emph{sli˥t}] 0xB069: тридцать шесть\\
\textbf{slic} [\emph{sliʃ}] 0xE869: оружие\\
\textbf{slif} [\emph{slif}] 0xC869: сито\\
\textbf{slih} [\emph{sliʰ}] 0x434F: спать\\
\textbf{slik} [\emph{slik}] 0x2869: шелк\\
\textbf{slim} [\emph{slim}] 0x0869: свисток\\
\textbf{slin} [\emph{slin}] 0x6869: Серебряный\\
\textbf{slip} [\emph{slip}] 0x4869: скольжение\\
\textbf{slis} [\emph{slis}] 0x8869: серии\\
\textbf{slit} [\emph{slit}] 0xA869: владелец\\
\textbf{slo2k} [\emph{slo˩k}] 0x3C69: мезолит\\
\textbf{slo2n} [\emph{slo˩n}] 0x7C69: столбчатый\\
\textbf{slo2s} [\emph{slo˩s}] 0x9C69: все отводы\\
\textbf{slo2t} [\emph{slo˩t}] 0xBC69: мопед\\
\textbf{slo7k} [\emph{slo˥k}] 0x3469: систолический\\
\textbf{slo7m} [\emph{slo˥m}] 0x1469: миелома\\
\textbf{slo7n} [\emph{slo˥n}] 0x7469: Соломон\\
\textbf{slo7s} [\emph{slo˥s}] 0x9469: тусклый\\
\textbf{slo7t} [\emph{slo˥t}] 0xB469: энциклопедия\\
\textbf{sloc} [\emph{sloʃ}] 0xEC69: психология\\
\textbf{slof} [\emph{slof}] 0xCC69: морское дно\\
\textbf{sloh} [\emph{sloʰ}] 0x634F: успех\\
\textbf{slok} [\emph{slok}] 0x2C69: социология\\
\textbf{slom} [\emph{slom}] 0x0C69: читать только Память\\
\textbf{slon} [\emph{slon}] 0x6C69: шитье\\
\textbf{slop} [\emph{slop}] 0x4C69: тапочка\\
\textbf{slos} [\emph{slos}] 0x8C69: колбаса\\
\textbf{slot} [\emph{slot}] 0xAC69: стрелять\\
\textbf{slu2n} [\emph{slu˩n}] 0x7A69: лазурный\\
\textbf{slu2t} [\emph{slu˩t}] 0xBA69: мускулатура\\
\textbf{slu7n} [\emph{slu˥n}] 0x7269: суб-лейтенант\\
\textbf{slu7t} [\emph{slu˥t}] 0xB269: предпоследний\\
\textbf{sluc} [\emph{sluʃ}] 0xEA69: токсикология\\
\textbf{sluf} [\emph{sluf}] 0xCA69: singulative номер\\
\textbf{sluh} [\emph{sluʰ}] 0x534F: самостоятельно удовлетворяет\\
\textbf{sluk} [\emph{sluk}] 0x2A69: белка\\
\textbf{slum} [\emph{slum}] 0x0A69: ислам\\
\textbf{slun} [\emph{slun}] 0x6A69: приостановить\\
\textbf{slup} [\emph{slup}] 0x4A69: лесок\\
\textbf{slus} [\emph{slus}] 0x8A69: доступность\\
\textbf{slut} [\emph{slut}] 0xAA69: салат\\
\textbf{so} [\emph{so}] 0x313E: source case, indicates a source noun phrase, from. is used for constructing complex cases or by itself.\\
\textbf{sr67n} [\emph{srə˥n}] 0x75C9: послеобеденный\\
\textbf{sr67t} [\emph{srə˥t}] 0xB5C9: прямой край\\
\textbf{sr6c} [\emph{srəʃ}] 0xEDC9: усик\\
\textbf{sr6f} [\emph{srəf}] 0xCDC9: выживаемость\\
\textbf{sr6h} [\emph{srəʰ}] 0x6E4F: abstract gender, for making new genders, for example tlansr6h would be the gender of the one infinite creator God.\\
\textbf{sr6k} [\emph{srək}] 0x2DC9: эстроген\\
\textbf{sr6m} [\emph{srəm}] 0x0DC9: стронций\\
\textbf{sr6n} [\emph{srən}] 0x6DC9: созревание\\
\textbf{sr6p} [\emph{srəp}] 0x4DC9: синапс\\
\textbf{sr6s} [\emph{srəs}] 0x8DC9: Кипр\\
\textbf{sr6t} [\emph{srət}] 0xADC9: сверхкритической жидкости\\
\textbf{sra2f} [\emph{srä˩f}] 0xD9C9: омыляться\\
\textbf{sra2k} [\emph{srä˩k}] 0x39C9: ондатра\\
\textbf{sra2m} [\emph{srä˩m}] 0x19C9: перевооружать\\
\textbf{sra2n} [\emph{srä˩n}] 0x79C9: сверхчеловек\\
\textbf{sra2p} [\emph{srä˩p}] 0x59C9: вызов в суд\\
\textbf{sra2s} [\emph{srä˩s}] 0x99C9: неквалифицированный\\
\textbf{sra2t} [\emph{srä˩t}] 0xB9C9: спартанский\\
\textbf{sra7c} [\emph{srä˥ʃ}] 0xF1C9: астральный\\
\textbf{sra7f} [\emph{srä˥f}] 0xD1C9: разносить\\
\textbf{sra7k} [\emph{srä˥k}] 0x31C9: манить\\
\textbf{sra7m} [\emph{srä˥m}] 0x11C9: Самария\\
\textbf{sra7n} [\emph{srä˥n}] 0x71C9: неотделимый\\
\textbf{sra7p} [\emph{srä˥p}] 0x51C9: сербия\\
\textbf{sra7s} [\emph{srä˥s}] 0x91C9: спаржа\\
\textbf{sra7t} [\emph{srä˥t}] 0xB1C9: Сократ\\
\textbf{srac} [\emph{sräʃ}] 0xE9C9: дешево\\
\textbf{sraf} [\emph{sräf}] 0xC9C9: на открытом воздухе\\
\textbf{srah} [\emph{sräʰ}] 0x4E4F: до\\
\textbf{srak} [\emph{sräk}] 0x29C9: рукопись\\
\textbf{sram} [\emph{sräm}] 0x09C9: буквально\\
\textbf{sran} [\emph{srän}] 0x69C9: похороны\\
\textbf{srap} [\emph{sräp}] 0x49C9: овощной\\
\textbf{sras} [\emph{sräs}] 0x89C9: гнездо\\
\textbf{srat} [\emph{srät}] 0xA9C9: звезда\\
\textbf{sre2c} [\emph{sre˩ʃ}] 0xFBC9: замена\\
\textbf{sre2k} [\emph{sre˩k}] 0x3BC9: сэр благоговением\\
\textbf{sre2n} [\emph{sre˩n}] 0x7BC9: зычный\\
\textbf{sre2s} [\emph{sre˩s}] 0x9BC9: течка\\
\textbf{sre2t} [\emph{sre˩t}] 0xBBC9: опытным путем\\
\textbf{sre7c} [\emph{sre˥ʃ}] 0xF3C9: стригальщик\\
\textbf{sre7k} [\emph{sre˥k}] 0x33C9: лексикография\\
\textbf{sre7n} [\emph{sre˥n}] 0x73C9: фосфоресцирующий\\
\textbf{sre7s} [\emph{sre˥s}] 0x93C9: заклинатель\\
\textbf{sre7t} [\emph{sre˥t}] 0xB3C9: небрежность\\
\textbf{srec} [\emph{sreʃ}] 0xEBC9: истекающий\\
\textbf{sref} [\emph{sref}] 0xCBC9: стенографистка\\
\textbf{sreh} [\emph{sreʰ}] 0x5E4F: серии\\
\textbf{srek} [\emph{srek}] 0x2BC9: уксус\\
\textbf{srem} [\emph{srem}] 0x0BC9: церий\\
\textbf{sren} [\emph{sren}] 0x6BC9: экспонат\\
\textbf{srep} [\emph{srep}] 0x4BC9: киберпространство\\
\textbf{sres} [\emph{sres}] 0x8BC9: С уважением\\
\textbf{sret} [\emph{sret}] 0xABC9: семнадцать\\
\textbf{sri2c} [\emph{sri˩ʃ}] 0xF8C9: придыхательный\\
\textbf{sri2k} [\emph{sri˩k}] 0x38C9: круто\\
\textbf{sri2m} [\emph{sri˩m}] 0x18C9: садизм\\
\textbf{sri2n} [\emph{sri˩n}] 0x78C9: странность\\
\textbf{sri2p} [\emph{sri˩p}] 0x58C9: изопрен\\
\textbf{sri2s} [\emph{sri˩s}] 0x98C9: вяжущий\\
\textbf{sri2t} [\emph{sri˩t}] 0xB8C9: казалось\\
\textbf{sri7c} [\emph{sri˥ʃ}] 0xF0C9: небиблейский\\
\textbf{sri7f} [\emph{sri˥f}] 0xD0C9: стенография\\
\textbf{sri7h} [\emph{sri˥ʰ}] 0x864F: взаимная голос\\
\textbf{sri7k} [\emph{sri˥k}] 0x30C9: крестец\\
\textbf{sri7m} [\emph{sri˥m}] 0x10C9: синкретизм\\
\textbf{sri7n} [\emph{sri˥n}] 0x70C9: ацетилен\\
\textbf{sri7p} [\emph{sri˥p}] 0x50C9: церебральный\\
\textbf{sri7s} [\emph{sri˥s}] 0x90C9: писец\\
\textbf{sri7t} [\emph{sri˥t}] 0xB0C9: историки\\
\textbf{sric} [\emph{sriʃ}] 0xE8C9: церемония\\
\textbf{srif} [\emph{srif}] 0xC8C9: исторический напряженный\\
\textbf{srih} [\emph{sriʰ}] 0x464F: грустный\\
\textbf{srik} [\emph{srik}] 0x28C9: семена\\
\textbf{srim} [\emph{srim}] 0x08C9: подпись\\
\textbf{srin} [\emph{srin}] 0x68C9: инстинкт\\
\textbf{srip} [\emph{srip}] 0x48C9: сигарета\\
\textbf{sris} [\emph{sris}] 0x88C9: исторический\\
\textbf{srit} [\emph{srit}] 0xA8C9: инструмент\\
\textbf{sro2k} [\emph{sro˩k}] 0x3CC9: изотропный\\
\textbf{sro2n} [\emph{sro˩n}] 0x7CC9: правый борт\\
\textbf{sro2s} [\emph{sro˩s}] 0x9CC9: зоб\\
\textbf{sro2t} [\emph{sro˩t}] 0xBCC9: иллюстратор\\
\textbf{sro7c} [\emph{sro˥ʃ}] 0xF4C9: сербско хорватской\\
\textbf{sro7k} [\emph{sro˥k}] 0x34C9: хриплый\\
\textbf{sro7n} [\emph{sro˥n}] 0x74C9: все важные\\
\textbf{sro7p} [\emph{sro˥p}] 0x54C9: строб\\
\textbf{sro7s} [\emph{sro˥s}] 0x94C9: слишком любящий свою жену\\
\textbf{sro7t} [\emph{sro˥t}] 0xB4C9: посмертный\\
\textbf{sroc} [\emph{sroʃ}] 0xECC9: экспорт\\
\textbf{srof} [\emph{srof}] 0xCCC9: осмос\\
\textbf{sroh} [\emph{sroʰ}] 0x664F: помилование\\
\textbf{srok} [\emph{srok}] 0x2CC9: редакционный\\
\textbf{srom} [\emph{srom}] 0x0CC9: выставочный зал\\
\textbf{sron} [\emph{sron}] 0x6CC9: воин\\
\textbf{srop} [\emph{srop}] 0x4CC9: смартфон\\
\textbf{sros} [\emph{sros}] 0x8CC9: меньшинство\\
\textbf{srot} [\emph{srot}] 0xACC9: эскорт\\
\textbf{sru2n} [\emph{sru˩n}] 0x7AC9: распятие на кресте\\
\textbf{sru2t} [\emph{sru˩t}] 0xBAC9: цитрат\\
\textbf{sru7k} [\emph{sru˥k}] 0x32C9: сверхзвуковой\\
\textbf{sru7n} [\emph{sru˥n}] 0x72C9: мастурбация\\
\textbf{sru7s} [\emph{sru˥s}] 0x92C9: Астурийский\\
\textbf{sru7t} [\emph{sru˥t}] 0xB2C9: самурай\\
\textbf{sruc} [\emph{sruʃ}] 0xEAC9: вонь\\
\textbf{sruf} [\emph{sruf}] 0xCAC9: пугало\\
\textbf{sruh} [\emph{sruʰ}] 0x564F: трепет\\
\textbf{sruk} [\emph{sruk}] 0x2AC9: цирк\\
\textbf{srum} [\emph{srum}] 0x0AC9: высший\\
\textbf{srun} [\emph{srun}] 0x6AC9: хмуриться\\
\textbf{srup} [\emph{srup}] 0x4AC9: суп\\
\textbf{srus} [\emph{srus}] 0x8AC9: вытеснять\\
\textbf{srut} [\emph{srut}] 0xAAC9: Абстрактные\\
\textbf{su} [\emph{su}] 0x293E: Кроме\\
\textbf{sw6k} [\emph{swək}] 0x2D29: едок\\
\textbf{sw6n} [\emph{swən}] 0x6D29: санитария\\
\textbf{sw6s} [\emph{swəs}] 0x8D29: особы\\
\textbf{sw6t} [\emph{swət}] 0xAD29: саудовская\\
\textbf{swa2f} [\emph{swä˩f}] 0xD929: оссифицировать\\
\textbf{swa2k} [\emph{swä˩k}] 0x3929: сфагнум\\
\textbf{swa2m} [\emph{swä˩m}] 0x1929: пространство-время\\
\textbf{swa2n} [\emph{swä˩n}] 0x7929: юго-восточный\\
\textbf{swa2s} [\emph{swä˩s}] 0x9929: сагиттальный\\
\textbf{swa2t} [\emph{swä˩t}] 0xB929: восемьдесят два\\
\textbf{swa7c} [\emph{swä˥ʃ}] 0xF129: образчик ткани\\
\textbf{swa7k} [\emph{swä˥k}] 0x3129: хиромантия\\
\textbf{swa7m} [\emph{swä˥m}] 0x1129: астигматизм\\
\textbf{swa7n} [\emph{swä˥n}] 0x7129: каламбур\\
\textbf{swa7p} [\emph{swä˥p}] 0x5129: тральщик\\
\textbf{swa7s} [\emph{swä˥s}] 0x9129: защелками\\
\textbf{swa7t} [\emph{swä˥t}] 0xB129: соната\\
\textbf{swac} [\emph{swäʃ}] 0xE929: служащий\\
\textbf{swaf} [\emph{swäf}] 0xC929: ассоциативный дело\\
\textbf{swah} [\emph{swäʰ}] 0x494F: semelfactive aspect, for expressing a sudden quick event, like a blink or a sneeze.\\
\textbf{swak} [\emph{swäk}] 0x2929: раковина\\
\textbf{swam} [\emph{swäm}] 0x0929: близнец\\
\textbf{swan} [\emph{swän}] 0x6929: дела\\
\textbf{swap} [\emph{swäp}] 0x4929: подметать\\
\textbf{swas} [\emph{swäs}] 0x8929: оказание услуг\\
\textbf{swat} [\emph{swät}] 0xA929: меч\\
\textbf{swe2n} [\emph{swe˩n}] 0x7B29: мореходный\\
\textbf{swe2t} [\emph{swe˩t}] 0xBB29: так и так\\
\textbf{swe7n} [\emph{swe˥n}] 0x7329: продавщица\\
\textbf{swec} [\emph{sweʃ}] 0xEB29: случайный\\
\textbf{sweh} [\emph{sweʰ}] 0x594F: sublative case, indicates motion toward the underside of nounphrase. under-context destination-case.\\
\textbf{swek} [\emph{swek}] 0x2B29: перегружены работой\\
\textbf{swem} [\emph{swem}] 0x0B29: разбрызгиватель\\
\textbf{swen} [\emph{swen}] 0x6B29: сенат\\
\textbf{swep} [\emph{swep}] 0x4B29: соболь\\
\textbf{swes} [\emph{swes}] 0x8B29: свитер\\
\textbf{swet} [\emph{swet}] 0xAB29: убеждать\\
\textbf{swi2k} [\emph{swi˩k}] 0x3829: мифический\\
\textbf{swi2n} [\emph{swi˩n}] 0x7829: психоанализ\\
\textbf{swi2t} [\emph{swi˩t}] 0xB829: тридцать семь\\
\textbf{swi7c} [\emph{swi˥ʃ}] 0xF029: относящийся к пищеводу\\
\textbf{swi7k} [\emph{swi˥k}] 0x3029: изометрический\\
\textbf{swi7n} [\emph{swi˥n}] 0x7029: сорок девять\\
\textbf{swi7s} [\emph{swi˥s}] 0x9029: в сторону моря\\
\textbf{swi7t} [\emph{swi˥t}] 0xB029: пятьдесят семь\\
\textbf{swic} [\emph{swiʃ}] 0xE829: Социальное\\
\textbf{swif} [\emph{swif}] 0xC829: Швейцария\\
\textbf{swih} [\emph{swiʰ}] 0x414F: essive case, indicates a duration of time at which the noun phrase occurs. time-context way-case. \\
\textbf{swik} [\emph{swik}] 0x2829: щий дело\\
\textbf{swim} [\emph{swim}] 0x0829: системы\\
\textbf{swin} [\emph{swin}] 0x6829: гражданского\\
\textbf{swip} [\emph{swip}] 0x4829: шип\\
\textbf{swis} [\emph{swis}] 0x8829: моряк\\
\textbf{swit} [\emph{swit}] 0xA829: повелительный настроение\\
\textbf{swo2k} [\emph{swo˩k}] 0x3C29: полуавтоматический\\
\textbf{swo2t} [\emph{swo˩t}] 0xBC29: Сальвадорских\\
\textbf{swo7n} [\emph{swo˥n}] 0x7429: мастодонт\\
\textbf{swo7s} [\emph{swo˥s}] 0x9429: сказать так\\
\textbf{swo7t} [\emph{swo˥t}] 0xB429: жатка\\
\textbf{swoc} [\emph{swoʃ}] 0xEC29: неустойчивый\\
\textbf{swof} [\emph{swof}] 0xCC29: Словения\\
\textbf{swoh} [\emph{swoʰ}] 0x614F: последовательный\\
\textbf{swok} [\emph{swok}] 0x2C29: прильнуть\\
\textbf{swom} [\emph{swom}] 0x0C29: Сомали\\
\textbf{swon} [\emph{swon}] 0x6C29: доброволец\\
\textbf{swop} [\emph{swop}] 0x4C29: секундомер\\
\textbf{swos} [\emph{swos}] 0x8C29: экспоненциальный\\
\textbf{swot} [\emph{swot}] 0xAC29: болото\\
\textbf{swu7k} [\emph{swu˥k}] 0x3229: Суккот\\
\textbf{swuc} [\emph{swuʃ}] 0xEA29: иней\\
\textbf{swuh} [\emph{swuʰ}] 0x514F: предположительный настроение\\
\textbf{swuk} [\emph{swuk}] 0x2A29: превосходный дело\\
\textbf{swum} [\emph{swum}] 0x0A29: спиритизм\\
\textbf{swun} [\emph{swun}] 0x6A29: мыло\\
\textbf{swup} [\emph{swup}] 0x4A29: кушетка\\
\textbf{swus} [\emph{swus}] 0x8A29: самокат\\
\textbf{swut} [\emph{swut}] 0xAA29: святой\\
\textbf{sy67n} [\emph{sjə˥n}] 0x7509: изотонический\\
\textbf{sy6k} [\emph{sjək}] 0x2D09: экзистенциальный пункт\\
\textbf{sy6m} [\emph{sjəm}] 0x0D09: сатанизм\\
\textbf{sy6n} [\emph{sjən}] 0x6D09: цемент\\
\textbf{sy6s} [\emph{sjəs}] 0x8D09: нездоровый\\
\textbf{sy6t} [\emph{sjət}] 0xAD09: Свазиленд\\
\textbf{sya2c} [\emph{sjä˩ʃ}] 0xF909: осетия\\
\textbf{sya2k} [\emph{sjä˩k}] 0x3909: асфиксия\\
\textbf{sya2n} [\emph{sjä˩n}] 0x7909: сатин\\
\textbf{sya2p} [\emph{sjä˩p}] 0x5909: бестселлер\\
\textbf{sya2s} [\emph{sjä˩s}] 0x9909: стареющий\\
\textbf{sya2t} [\emph{sjä˩t}] 0xB909: шестьдесят восемь\\
\textbf{sya7c} [\emph{sjä˥ʃ}] 0xF109: половина дядя\\
\textbf{sya7k} [\emph{sjä˥k}] 0x3109: целебный\\
\textbf{sya7m} [\emph{sjä˥m}] 0x1109: санаторий\\
\textbf{sya7n} [\emph{sjä˥n}] 0x7109: разведка\\
\textbf{sya7p} [\emph{sjä˥p}] 0x5109: шпунт\\
\textbf{sya7s} [\emph{sjä˥s}] 0x9109: схоластика\\
\textbf{sya7t} [\emph{sjä˥t}] 0xB109: завещание\\
\textbf{syac} [\emph{sjäʃ}] 0xE909: свежий\\
\textbf{syaf} [\emph{sjäf}] 0xC909: стороны\\
\textbf{syah} [\emph{sjäʰ}] 0x484F: улыбка\\
\textbf{syak} [\emph{sjäk}] 0x2909: войти\\
\textbf{syam} [\emph{sjäm}] 0x0909: безопасно\\
\textbf{syan} [\emph{sjän}] 0x6909: правда\\
\textbf{syap} [\emph{sjäp}] 0x4909: распространять\\
\textbf{syas} [\emph{sjäs}] 0x8909: сроки\\
\textbf{syat} [\emph{sjät}] 0xA909: до\\
\textbf{sye2n} [\emph{sje˩n}] 0x7B09: десятилетний\\
\textbf{sye2t} [\emph{sje˩t}] 0xBB09: ацетат\\
\textbf{sye7n} [\emph{sje˥n}] 0x7309: мессианский\\
\textbf{sye7p} [\emph{sje˥p}] 0x5309: Сен-Пьер\\
\textbf{sye7s} [\emph{sje˥s}] 0x9309: Seychelles\\
\textbf{sye7t} [\emph{sje˥t}] 0xB309: шестьдесят девять\\
\textbf{syec} [\emph{sjeʃ}] 0xEB09: располагать\\
\textbf{syef} [\emph{sjef}] 0xCB09: Scot бесплатно\\
\textbf{syeh} [\emph{sjeʰ}] 0x584F: с этого времени\\
\textbf{syek} [\emph{sjek}] 0x2B09: класть не на место\\
\textbf{syem} [\emph{sjem}] 0x0B09: спам\\
\textbf{syen} [\emph{sjen}] 0x6B09: сэр\\
\textbf{syep} [\emph{sjep}] 0x4B09: скрепок\\
\textbf{syes} [\emph{sjes}] 0x8B09: седан\\
\textbf{syet} [\emph{sjet}] 0xAB09: эскиз\\
\textbf{syi2c} [\emph{sji˩ʃ}] 0xF809: социально-экономический\\
\textbf{syi2f} [\emph{sji˩f}] 0xD809: седиль\\
\textbf{syi2k} [\emph{sji˩k}] 0x3809: семиотический\\
\textbf{syi2n} [\emph{sji˩n}] 0x7809: Сидон\\
\textbf{syi2p} [\emph{sji˩p}] 0x5809: скопа\\
\textbf{syi2s} [\emph{sji˩s}] 0x9809: псориаз\\
\textbf{syi2t} [\emph{sji˩t}] 0xB809: сослагательное наклонение\\
\textbf{syi7c} [\emph{sji˥ʃ}] 0xF009: ацетил\\
\textbf{syi7f} [\emph{sji˥f}] 0xD009: София\\
\textbf{syi7h} [\emph{sji˥ʰ}] 0x804F: existential clause, for asserting something that can be imagined.\\
\textbf{syi7h} [\emph{sji˥ʰ}] 0x804F: existential clause, for asserting something that can be imagined.\\
\textbf{syi7k} [\emph{sji˥k}] 0x3009: сикх\\
\textbf{syi7n} [\emph{sji˥n}] 0x7009: жестикуляция\\
\textbf{syi7s} [\emph{sji˥s}] 0x9009: осознанный\\
\textbf{syi7t} [\emph{sji˥t}] 0xB009: фыркнул\\
\textbf{syic} [\emph{sjiʃ}] 0xE809: грустный\\
\textbf{syif} [\emph{sjif}] 0xC809: си единицы\\
\textbf{syih} [\emph{sjiʰ}] 0x404F: в равной степени\\
\textbf{syik} [\emph{sjik}] 0x2809: политическое\\
\textbf{syim} [\emph{sjim}] 0x0809: больной\\
\textbf{syin} [\emph{sjin}] 0x6809: Уважаемые\\
\textbf{syip} [\emph{sjip}] 0x4809: в равной степени\\
\textbf{syis} [\emph{sjis}] 0x8809: 20\\
\textbf{syit} [\emph{sjit}] 0xA809: город\\
\textbf{syo2n} [\emph{sjo˩n}] 0x7C09: сифон\\
\textbf{syo2t} [\emph{sjo˩t}] 0xBC09: соя\\
\textbf{syo7n} [\emph{sjo˥n}] 0x7409: сырой\\
\textbf{syo7t} [\emph{sjo˥t}] 0xB409: отступ\\
\textbf{syoc} [\emph{sjoʃ}] 0xEC09: снимок\\
\textbf{syof} [\emph{sjof}] 0xCC09: Исполнительный директор\\
\textbf{syoh} [\emph{sjoʰ}] 0x604F: благоговение\\
\textbf{syok} [\emph{sjok}] 0x2C09: сторонник\\
\textbf{syom} [\emph{sjom}] 0x0C09: натрий\\
\textbf{syon} [\emph{sjon}] 0x6C09: тяжелое дыхание\\
\textbf{syop} [\emph{sjop}] 0x4C09: сотни\\
\textbf{syos} [\emph{sjos}] 0x8C09: подведение\\
\textbf{syot} [\emph{sjot}] 0xAC09: сложить\\
\textbf{syu2n} [\emph{sju˩n}] 0x7A09: Сеульский\\
\textbf{syu2t} [\emph{sju˩t}] 0xBA09: пломбир\\
\textbf{syu7n} [\emph{sju˥n}] 0x7209: укус ядовитой змеи\\
\textbf{syu7t} [\emph{sju˥t}] 0xB209: акцепторы\\
\textbf{syuc} [\emph{sjuʃ}] 0xEA09: дерево\\
\textbf{syuh} [\emph{sjuʰ}] 0x504F: супер\\
\textbf{syuk} [\emph{sjuk}] 0x2A09: ложка\\
\textbf{syum} [\emph{sjum}] 0x0A09: стронций атом\\
\textbf{syun} [\emph{sjun}] 0x6A09: правительство\\
\textbf{syup} [\emph{sjup}] 0x4A09: супер\\
\textbf{syus} [\emph{sjus}] 0x8A09: субсидировать\\
\textbf{syut} [\emph{sjut}] 0xAA09: свобода\\
\subsection{ t }
\textbf{t6} [\emph{tə}] 0x355E: утвердительный настроение\\
\textbf{ta} [\emph{tä}] 0x255E: topic case, the topic about which the sentence is. `regarding topic, some sentence.''. discourse-context way-case\\
\textbf{ta2} [\emph{tä˩}] 0x655E: continuative aspect, the state is continuing to occur, it is `still' occuring. \\
\textbf{ta7} [\emph{tä˥}] 0x455E: местный направленный дело\\
\textbf{tc67t} [\emph{tʃə˥t}] 0xB5AA: том Тит\\
\textbf{tc6c} [\emph{tʃəʃ}] 0xEDAA: мангольд\\
\textbf{tc6h} [\emph{tʃəʰ}] 0x6D57: целое число\\
\textbf{tc6k} [\emph{tʃək}] 0x2DAA: Артишок\\
\textbf{tc6m} [\emph{tʃəm}] 0x0DAA: трапеция\\
\textbf{tc6n} [\emph{tʃən}] 0x6DAA: художественный Пол\\
\textbf{tc6s} [\emph{tʃəs}] 0x8DAA: Лихтенштейн\\
\textbf{tc6t} [\emph{tʃət}] 0xADAA: непосредственный объект\\
\textbf{tca2c} [\emph{tʃä˩ʃ}] 0xF9AA: гашиш\\
\textbf{tca2f} [\emph{tʃä˩f}] 0xD9AA: каменная стена\\
\textbf{tca2k} [\emph{tʃä˩k}] 0x39AA: оружейный\\
\textbf{tca2m} [\emph{tʃä˩m}] 0x19AA: таоизм\\
\textbf{tca2n} [\emph{tʃä˩n}] 0x79AA: бранить\\
\textbf{tca2p} [\emph{tʃä˩p}] 0x59AA: магазин игрушек\\
\textbf{tca2s} [\emph{tʃä˩s}] 0x99AA: бессменный\\
\textbf{tca2t} [\emph{tʃä˩t}] 0xB9AA: непредсказуемость\\
\textbf{tca7c} [\emph{tʃä˥ʃ}] 0xF1AA: пижма\\
\textbf{tca7f} [\emph{tʃä˥f}] 0xD1AA: жестянщик\\
\textbf{tca7k} [\emph{tʃä˥k}] 0x31AA: эрцгерцог\\
\textbf{tca7m} [\emph{tʃä˥m}] 0x11AA: меткий стрелок\\
\textbf{tca7n} [\emph{tʃä˥n}] 0x71AA: тридцать девять\\
\textbf{tca7p} [\emph{tʃä˥p}] 0x51AA: высотомер\\
\textbf{tca7s} [\emph{tʃä˥s}] 0x91AA: характерно\\
\textbf{tca7t} [\emph{tʃä˥t}] 0xB1AA: право по рождению\\
\textbf{tcac} [\emph{tʃäʃ}] 0xE9AA: девушка\\
\textbf{tcaf} [\emph{tʃäf}] 0xC9AA: внимательно\\
\textbf{tcak} [\emph{tʃäk}] 0x29AA: одна и та же предмет маркер\\
\textbf{tcam} [\emph{tʃäm}] 0x09AA: враг\\
\textbf{tcan} [\emph{tʃän}] 0x69AA: Другие\\
\textbf{tcap} [\emph{tʃäp}] 0x49AA: побег\\
\textbf{tcas} [\emph{tʃäs}] 0x89AA: места\\
\textbf{tcat} [\emph{tʃät}] 0xA9AA: будущее напряженный\\
\textbf{tce2n} [\emph{tʃe˩n}] 0x7BAA: чеченец\\
\textbf{tce2t} [\emph{tʃe˩t}] 0xBBAA: рикошет\\
\textbf{tce7c} [\emph{tʃe˥ʃ}] 0xF3AA: амбулаторный\\
\textbf{tce7k} [\emph{tʃe˥k}] 0x33AA: ЕЭС\\
\textbf{tce7m} [\emph{tʃe˥m}] 0x13AA: трахеотомия\\
\textbf{tce7n} [\emph{tʃe˥n}] 0x73AA: изнурение\\
\textbf{tce7s} [\emph{tʃe˥s}] 0x93AA: дегустатор\\
\textbf{tce7t} [\emph{tʃe˥t}] 0xB3AA: чириканье\\
\textbf{tcec} [\emph{tʃeʃ}] 0xEBAA: обмен\\
\textbf{tcef} [\emph{tʃef}] 0xCBAA: вишня\\
\textbf{tceh} [\emph{tʃeʰ}] 0x5D57: одна и та же предмет маркер\\
\textbf{tcek} [\emph{tʃek}] 0x2BAA: эхо\\
\textbf{tcem} [\emph{tʃem}] 0x0BAA: щипцы\\
\textbf{tcen} [\emph{tʃen}] 0x6BAA: наблюдение\\
\textbf{tcep} [\emph{tʃep}] 0x4BAA: мел\\
\textbf{tces} [\emph{tʃes}] 0x8BAA: большой теннис\\
\textbf{tcet} [\emph{tʃet}] 0xABAA: угрожающий\\
\textbf{tci2c} [\emph{tʃi˩ʃ}] 0xF8AA: демерредж\\
\textbf{tci2k} [\emph{tʃi˩k}] 0x38AA: зубочистка\\
\textbf{tci2m} [\emph{tʃi˩m}] 0x18AA: лимит времени\\
\textbf{tci2n} [\emph{tʃi˩n}] 0x78AA: увещевать\\
\textbf{tci2p} [\emph{tʃi˩p}] 0x58AA: растирать в порошок\\
\textbf{tci2s} [\emph{tʃi˩s}] 0x98AA: семидесятый\\
\textbf{tci2t} [\emph{tʃi˩t}] 0xB8AA: к тому же\\
\textbf{tci7c} [\emph{tʃi˥ʃ}] 0xF0AA: децимация\\
\textbf{tci7h} [\emph{tʃi˥ʰ}] 0x8557: Archy\\
\textbf{tci7k} [\emph{tʃi˥k}] 0x30AA: ишиас\\
\textbf{tci7m} [\emph{tʃi˥m}] 0x10AA: тимус\\
\textbf{tci7n} [\emph{tʃi˥n}] 0x70AA: вздрогнуть\\
\textbf{tci7p} [\emph{tʃi˥p}] 0x50AA: палантин\\
\textbf{tci7s} [\emph{tʃi˥s}] 0x90AA: шестерки\\
\textbf{tci7t} [\emph{tʃi˥t}] 0xB0AA: изобретенной\\
\textbf{tcic} [\emph{tʃiʃ}] 0xE8AA: тело\\
\textbf{tcif} [\emph{tʃif}] 0xC8AA: заостренный\\
\textbf{tcih} [\emph{tʃiʰ}] 0x4557: hortative mood, for pleading or begging.\\
\textbf{tcik} [\emph{tʃik}] 0x28AA: доказательный дело\\
\textbf{tcim} [\emph{tʃim}] 0x08AA: одна и та же мероприятие маркер\\
\textbf{tcin} [\emph{tʃin}] 0x68AA: день\\
\textbf{tcip} [\emph{tʃip}] 0x48AA: симпатия\\
\textbf{tcis} [\emph{tʃis}] 0x88AA: честный\\
\textbf{tcit} [\emph{tʃit}] 0xA8AA: совещательный настроение\\
\textbf{tco2n} [\emph{tʃo˩n}] 0x7CAA: кошениль\\
\textbf{tco2s} [\emph{tʃo˩s}] 0x9CAA: трихинеллез\\
\textbf{tco2t} [\emph{tʃo˩t}] 0xBCAA: ротонда\\
\textbf{tco7n} [\emph{tʃo˥n}] 0x74AA: колоколообразный\\
\textbf{tco7s} [\emph{tʃo˥s}] 0x94AA: пятно на солнце\\
\textbf{tco7t} [\emph{tʃo˥t}] 0xB4AA: смертное ложе\\
\textbf{tcoc} [\emph{tʃoʃ}] 0xECAA: крутой\\
\textbf{tcof} [\emph{tʃof}] 0xCCAA: стрельба из лука\\
\textbf{tcoh} [\emph{tʃoʰ}] 0x6557: боль\\
\textbf{tcok} [\emph{tʃok}] 0x2CAA: точность\\
\textbf{tcom} [\emph{tʃom}] 0x0CAA: удушение\\
\textbf{tcon} [\emph{tʃon}] 0x6CAA: полдень\\
\textbf{tcop} [\emph{tʃop}] 0x4CAA: губа\\
\textbf{tcos} [\emph{tʃos}] 0x8CAA: урожай\\
\textbf{tcot} [\emph{tʃot}] 0xACAA: твердое вещество\\
\textbf{tcu2n} [\emph{tʃu˩n}] 0x7AAA: тюнер\\
\textbf{tcu2t} [\emph{tʃu˩t}] 0xBAAA: поганка\\
\textbf{tcu7s} [\emph{tʃu˥s}] 0x92AA: актуарий\\
\textbf{tcu7t} [\emph{tʃu˥t}] 0xB2AA: плотно прижимистый\\
\textbf{tcuc} [\emph{tʃuʃ}] 0xEAAA: появляться\\
\textbf{tcuh} [\emph{tʃuʰ}] 0x5557: энтузиазм\\
\textbf{tcuk} [\emph{tʃuk}] 0x2AAA: боль\\
\textbf{tcum} [\emph{tʃum}] 0x0AAA: коснуться\\
\textbf{tcun} [\emph{tʃun}] 0x6AAA: сделанный\\
\textbf{tcup} [\emph{tʃup}] 0x4AAA: влажный\\
\textbf{tcus} [\emph{tʃus}] 0x8AAA: Сью\\
\textbf{tcut} [\emph{tʃut}] 0xAAAA: религия\\
\textbf{te} [\emph{te}] 0x2D5E: место нахождения дело\\
\textbf{tf6k} [\emph{tfək}] 0x2D8A: стратификация\\
\textbf{tf6n} [\emph{tfən}] 0x6D8A: тонер\\
\textbf{tf6s} [\emph{tfəs}] 0x8D8A: две трети\\
\textbf{tf6t} [\emph{tfət}] 0xAD8A: колготки\\
\textbf{tfa2c} [\emph{tfä˩ʃ}] 0xF98A: стагфляция\\
\textbf{tfa2k} [\emph{tfä˩k}] 0x398A: испражняться\\
\textbf{tfa2n} [\emph{tfä˩n}] 0x798A: откармливать\\
\textbf{tfa2s} [\emph{tfä˩s}] 0x998A: метатеза\\
\textbf{tfa2t} [\emph{tfä˩t}] 0xB98A: брюшной тиф\\
\textbf{tfa7c} [\emph{tfä˥ʃ}] 0xF18A: переводный\\
\textbf{tfa7k} [\emph{tfä˥k}] 0x318A: вилы\\
\textbf{tfa7m} [\emph{tfä˥m}] 0x118A: трансформеры\\
\textbf{tfa7n} [\emph{tfä˥n}] 0x718A: назидательный\\
\textbf{tfa7s} [\emph{tfä˥s}] 0x918A: катарсис\\
\textbf{tfa7t} [\emph{tfä˥t}] 0xB18A: зубчатый\\
\textbf{tfac} [\emph{tfäʃ}] 0xE98A: место действия\\
\textbf{tfaf} [\emph{tfäf}] 0xC98A: трансцендентный\\
\textbf{tfah} [\emph{tfäʰ}] 0x4C57: надлежащий имя существительное\\
\textbf{tfak} [\emph{tfäk}] 0x298A: доказательный\\
\textbf{tfam} [\emph{tfäm}] 0x098A: прозрачный\\
\textbf{tfan} [\emph{tfän}] 0x698A: вдруг, внезапно\\
\textbf{tfap} [\emph{tfäp}] 0x498A: диаграмма\\
\textbf{tfas} [\emph{tfäs}] 0x898A: посол\\
\textbf{tfat} [\emph{tfät}] 0xA98A: двенадцать\\
\textbf{tfe2n} [\emph{tfe˩n}] 0x7B8A: правопреемник\\
\textbf{tfe2t} [\emph{tfe˩t}] 0xBB8A: неподтвержденный\\
\textbf{tfe7n} [\emph{tfe˥n}] 0x738A: этан\\
\textbf{tfe7s} [\emph{tfe˥s}] 0x938A: пасынок\\
\textbf{tfe7t} [\emph{tfe˥t}] 0xB38A: тафта\\
\textbf{tfec} [\emph{tfeʃ}] 0xEB8A: канцелярские товары\\
\textbf{tfef} [\emph{tfef}] 0xCB8A: заземление\\
\textbf{tfeh} [\emph{tfeʰ}] 0x5C57: telic aspect, for actions that have a definite end point. finite or limited.\\
\textbf{tfek} [\emph{tfek}] 0x2B8A: местный направленный дело\\
\textbf{tfem} [\emph{tfem}] 0x0B8A: сводная сестра\\
\textbf{tfen} [\emph{tfen}] 0x6B8A: защиты\\
\textbf{tfep} [\emph{tfep}] 0x4B8A: чайная чашка\\
\textbf{tfes} [\emph{tfes}] 0x8B8A: фарш\\
\textbf{tfet} [\emph{tfet}] 0xAB8A: производитель\\
\textbf{tfi2c} [\emph{tfi˩ʃ}] 0xF88A: зяблик\\
\textbf{tfi2k} [\emph{tfi˩k}] 0x388A: разврат\\
\textbf{tfi2n} [\emph{tfi˩n}] 0x788A: рассеянность внимания\\
\textbf{tfi2p} [\emph{tfi˩p}] 0x588A: чеканка\\
\textbf{tfi2s} [\emph{tfi˩s}] 0x988A: тридцать три\\
\textbf{tfi2t} [\emph{tfi˩t}] 0xB88A: тридцать четыре\\
\textbf{tfi7c} [\emph{tfi˥ʃ}] 0xF08A: этил\\
\textbf{tfi7f} [\emph{tfi˥f}] 0xD08A: атрофированный\\
\textbf{tfi7k} [\emph{tfi˥k}] 0x308A: превращение в стекло\\
\textbf{tfi7n} [\emph{tfi˥n}] 0x708A: объединение\\
\textbf{tfi7s} [\emph{tfi˥s}] 0x908A: дурачество\\
\textbf{tfi7t} [\emph{tfi˥t}] 0xB08A: неподдельный\\
\textbf{tfic} [\emph{tfiʃ}] 0xE88A: винтовка\\
\textbf{tfif} [\emph{tfif}] 0xC88A: переходный глагол\\
\textbf{tfih} [\emph{tfiʰ}] 0x4457: усилитель\\
\textbf{tfik} [\emph{tfik}] 0x288A: отражать\\
\textbf{tfim} [\emph{tfim}] 0x088A: литий\\
\textbf{tfin} [\emph{tfin}] 0x688A: Родился\\
\textbf{tfip} [\emph{tfip}] 0x488A: continuative аспект\\
\textbf{tfis} [\emph{tfis}] 0x888A: отказ\\
\textbf{tfit} [\emph{tfit}] 0xA88A: возврат\\
\textbf{tfo2n} [\emph{tfo˩n}] 0x7C8A: фиксация азота\\
\textbf{tfo2t} [\emph{tfo˩t}] 0xBC8A: карапуз\\
\textbf{tfo7f} [\emph{tfo˥f}] 0xD48A: UTF-8\\
\textbf{tfo7n} [\emph{tfo˥n}] 0x748A: буксировка цветные\\
\textbf{tfo7t} [\emph{tfo˥t}] 0xB48A: коммутатор\\
\textbf{tfoc} [\emph{tfoʃ}] 0xEC8A: приступ боли\\
\textbf{tfof} [\emph{tfof}] 0xCC8A: автоматизация\\
\textbf{tfoh} [\emph{tfoʰ}] 0x6457: моль\\
\textbf{tfok} [\emph{tfok}] 0x2C8A: мускатный орех\\
\textbf{tfom} [\emph{tfom}] 0x0C8A: тригонометрия\\
\textbf{tfon} [\emph{tfon}] 0x6C8A: дразнить\\
\textbf{tfop} [\emph{tfop}] 0x4C8A: зубная щетка\\
\textbf{tfos} [\emph{tfos}] 0x8C8A: застава\\
\textbf{tfot} [\emph{tfot}] 0xAC8A: сертификат\\
\textbf{tfu2n} [\emph{tfu˩n}] 0x7A8A: бой быков\\
\textbf{tfu2t} [\emph{tfu˩t}] 0xBA8A: тутти фрутти\\
\textbf{tfu7t} [\emph{tfu˥t}] 0xB28A: сыщик гончая\\
\textbf{tfuc} [\emph{tfuʃ}] 0xEA8A: ткацкий станок\\
\textbf{tfuf} [\emph{tfuf}] 0xCA8A: трюфель\\
\textbf{tfuh} [\emph{tfuʰ}] 0x5457: цитата форма\\
\textbf{tfuk} [\emph{tfuk}] 0x2A8A: домашний скот\\
\textbf{tfun} [\emph{tfun}] 0x6A8A: тоннель\\
\textbf{tfup} [\emph{tfup}] 0x4A8A: многоразовый\\
\textbf{tfus} [\emph{tfus}] 0x8A8A: зубная паста\\
\textbf{tfut} [\emph{tfut}] 0xAA8A: поздравлять\\
\textbf{ti} [\emph{ti}] 0x215E: родительный дело\\
\textbf{ti2} [\emph{ti˩}] 0x615E: электрический текущий\\
\textbf{ti7} [\emph{ti˥}] 0x415E: мимо совершенный\\
\textbf{tl62n} [\emph{tlə˩n}] 0x7D6A: инстилляция\\
\textbf{tl67k} [\emph{tlə˥k}] 0x356A: отиатрия\\
\textbf{tl67n} [\emph{tlə˥n}] 0x756A: интерпелляция\\
\textbf{tl67t} [\emph{tlə˥t}] 0xB56A: Средиземье\\
\textbf{tl6c} [\emph{tləʃ}] 0xED6A: энтомолог\\
\textbf{tl6f} [\emph{tləf}] 0xCD6A: кристаллография\\
\textbf{tl6h} [\emph{tləʰ}] 0x6B57: multal номер\\
\textbf{tl6k} [\emph{tlək}] 0x2D6A: трактор\\
\textbf{tl6m} [\emph{tləm}] 0x0D6A: Твидл ди\\
\textbf{tl6n} [\emph{tlən}] 0x6D6A: центральный обработка Ед. изм\\
\textbf{tl6p} [\emph{tləp}] 0x4D6A: топология\\
\textbf{tl6s} [\emph{tləs}] 0x8D6A: интерполяция\\
\textbf{tl6t} [\emph{tlət}] 0xAD6A: Атлантика\\
\textbf{tla2c} [\emph{tlä˩ʃ}] 0xF96A: желтофиоль садовая\\
\textbf{tla2f} [\emph{tlä˩f}] 0xD96A: прибавлять\\
\textbf{tla2k} [\emph{tlä˩k}] 0x396A: Transact\\
\textbf{tla2m} [\emph{tlä˩m}] 0x196A: таллий\\
\textbf{tla2n} [\emph{tlä˩n}] 0x796A: младший офицер\\
\textbf{tla2p} [\emph{tlä˩p}] 0x596A: боевой топор\\
\textbf{tla2s} [\emph{tlä˩s}] 0x996A: атлас\\
\textbf{tla2t} [\emph{tlä˩t}] 0xB96A: прыгнули\\
\textbf{tla7c} [\emph{tlä˥ʃ}] 0xF16A: триангуляция\\
\textbf{tla7f} [\emph{tlä˥f}] 0xD16A: ацетальдегид\\
\textbf{tla7k} [\emph{tlä˥k}] 0x316A: тальк\\
\textbf{tla7m} [\emph{tlä˥m}] 0x116A: звон колоколов\\
\textbf{tla7n} [\emph{tlä˥n}] 0x716A: фонарь\\
\textbf{tla7p} [\emph{tlä˥p}] 0x516A: телепатия\\
\textbf{tla7s} [\emph{tlä˥s}] 0x916A: сорок три\\
\textbf{tla7t} [\emph{tlä˥t}] 0xB16A: ударов\\
\textbf{tlac} [\emph{tläʃ}] 0xE96A: тихо\\
\textbf{tlaf} [\emph{tläf}] 0xC96A: принимать участие\\
\textbf{tlah} [\emph{tläʰ}] 0x4B57: факультативный квантификатором\\
\textbf{tlak} [\emph{tläk}] 0x296A: приносить\\
\textbf{tlam} [\emph{tläm}] 0x096A: пример\\
\textbf{tlan} [\emph{tlän}] 0x696A: Бог\\
\textbf{tlap} [\emph{tläp}] 0x496A: пол\\
\textbf{tlas} [\emph{tläs}] 0x896A: стоимость\\
\textbf{tlat} [\emph{tlät}] 0xA96A: слово\\
\textbf{tle2c} [\emph{tle˩ʃ}] 0xFB6A: внутриклеточный\\
\textbf{tle2k} [\emph{tle˩k}] 0x3B6A: энтомология\\
\textbf{tle2n} [\emph{tle˩n}] 0x7B6A: восьмидесятилетний\\
\textbf{tle2s} [\emph{tle˩s}] 0x9B6A: экстернализация\\
\textbf{tle2t} [\emph{tle˩t}] 0xBB6A: сетчатый\\
\textbf{tle7k} [\emph{tle˥k}] 0x336A: электронной коммерции\\
\textbf{tle7n} [\emph{tle˥n}] 0x736A: Трансильвания\\
\textbf{tle7s} [\emph{tle˥s}] 0x936A: катализ\\
\textbf{tle7t} [\emph{tle˥t}] 0xB36A: собеседник\\
\textbf{tlec} [\emph{tleʃ}] 0xEB6A: поставить галочку\\
\textbf{tlef} [\emph{tlef}] 0xCB6A: прицеп\\
\textbf{tleh} [\emph{tleʰ}] 0x5B57: hesternal tense, past day tense, for things that happened yesterday.\\
\textbf{tlek} [\emph{tlek}] 0x2B6A: штекер\\
\textbf{tlem} [\emph{tlem}] 0x0B6A: текущий актуальность маркер\\
\textbf{tlen} [\emph{tlen}] 0x6B6A: талант\\
\textbf{tlep} [\emph{tlep}] 0x4B6A: шаблон\\
\textbf{tles} [\emph{tles}] 0x8B6A: металлоид\\
\textbf{tlet} [\emph{tlet}] 0xAB6A: hesternal напряженный\\
\textbf{tli2c} [\emph{tli˩ʃ}] 0xF86A: астрология\\
\textbf{tli2k} [\emph{tli˩k}] 0x386A: монолог\\
\textbf{tli2m} [\emph{tli˩m}] 0x186A: антиномия\\
\textbf{tli2n} [\emph{tli˩n}] 0x786A: этимология\\
\textbf{tli2p} [\emph{tli˩p}] 0x586A: самосвал\\
\textbf{tli2s} [\emph{tli˩s}] 0x986A: артериосклероз\\
\textbf{tli2t} [\emph{tli˩t}] 0xB86A: тильда\\
\textbf{tli7c} [\emph{tli˥ʃ}] 0xF06A: перекупщик\\
\textbf{tli7h} [\emph{tli˥ʰ}] 0x8357: дистанционный пульт будущее напряженный\\
\textbf{tli7k} [\emph{tli˥k}] 0x306A: тик\\
\textbf{tli7m} [\emph{tli˥m}] 0x106A: патернализм\\
\textbf{tli7n} [\emph{tli˥n}] 0x706A: навязанный\\
\textbf{tli7p} [\emph{tli˥p}] 0x506A: талибы\\
\textbf{tli7s} [\emph{tli˥s}] 0x906A: двоек\\
\textbf{tli7t} [\emph{tli˥t}] 0xB06A: решительно\\
\textbf{tlic} [\emph{tliʃ}] 0xE86A: в конце концов\\
\textbf{tlif} [\emph{tlif}] 0xC86A: кафельная плитка\\
\textbf{tlih} [\emph{tliʰ}] 0x4357: утвердительный квантификатором\\
\textbf{tlik} [\emph{tlik}] 0x286A: мультипликативный дело\\
\textbf{tlim} [\emph{tlim}] 0x086A: лечение\\
\textbf{tlin} [\emph{tlin}] 0x686A: холодно\\
\textbf{tlip} [\emph{tlip}] 0x486A: установить\\
\textbf{tlis} [\emph{tlis}] 0x886A: зубы\\
\textbf{tlit} [\emph{tlit}] 0xA86A: второй\\
\textbf{tlo2c} [\emph{tlo˩ʃ}] 0xFC6A: этология\\
\textbf{tlo2k} [\emph{tlo˩k}] 0x3C6A: туберкулез\\
\textbf{tlo2n} [\emph{tlo˩n}] 0x7C6A: танин\\
\textbf{tlo2s} [\emph{tlo˩s}] 0x9C6A: нитроглицерин\\
\textbf{tlo2t} [\emph{tlo˩t}] 0xBC6A: семядоля\\
\textbf{tlo7k} [\emph{tlo˥k}] 0x346A: раболепствовать\\
\textbf{tlo7n} [\emph{tlo˥n}] 0x746A: Каталония\\
\textbf{tlo7s} [\emph{tlo˥s}] 0x946A: траулер\\
\textbf{tlo7t} [\emph{tlo˥t}] 0xB46A: кефаль\\
\textbf{tloc} [\emph{tloʃ}] 0xEC6A: ступенька\\
\textbf{tlof} [\emph{tlof}] 0xCC6A: телеологический\\
\textbf{tloh} [\emph{tloʰ}] 0x6357: multiplicative case, for indicating how many times the event occurs. as for frequentative-aspect or simple repeat loops.\\
\textbf{tlok} [\emph{tlok}] 0x2C6A: бак\\
\textbf{tlom} [\emph{tlom}] 0x0C6A: multal номер\\
\textbf{tlon} [\emph{tlon}] 0x6C6A: средство передвижения\\
\textbf{tlop} [\emph{tlop}] 0x4C6A: скатерть\\
\textbf{tlos} [\emph{tlos}] 0x8C6A: триллион\\
\textbf{tlot} [\emph{tlot}] 0xAC6A: детство\\
\textbf{tlu2n} [\emph{tlu˩n}] 0x7A6A: полупрозрачность\\
\textbf{tlu2t} [\emph{tlu˩t}] 0xBA6A: неокончательный\\
\textbf{tlu7k} [\emph{tlu˥k}] 0x326A: тавтологический\\
\textbf{tlu7n} [\emph{tlu˥n}] 0x726A: щекотание\\
\textbf{tlu7p} [\emph{tlu˥p}] 0x526A: толуол\\
\textbf{tlu7t} [\emph{tlu˥t}] 0xB26A: талмудический\\
\textbf{tluc} [\emph{tluʃ}] 0xEA6A: полотенце\\
\textbf{tluf} [\emph{tluf}] 0xCA6A: капелька болтовня\\
\textbf{tluh} [\emph{tluʰ}] 0x5357: единственное число\\
\textbf{tluk} [\emph{tluk}] 0x2A6A: вводный\\
\textbf{tlum} [\emph{tlum}] 0x0A6A: теллур\\
\textbf{tlun} [\emph{tlun}] 0x6A6A: брюки\\
\textbf{tlup} [\emph{tlup}] 0x4A6A: репа\\
\textbf{tlus} [\emph{tlus}] 0x8A6A: антиклимакс\\
\textbf{tlut} [\emph{tlut}] 0xAA6A: степень\\
\textbf{to} [\emph{to}] 0x315E: пространство контекст\\
\textbf{to7} [\emph{to˥}] 0x515E: OMA\\
\textbf{tr62t} [\emph{trə˩t}] 0xBDCA: Автопортрет\\
\textbf{tr67n} [\emph{trə˥n}] 0x75CA: каровое озеро\\
\textbf{tr67t} [\emph{trə˥t}] 0xB5CA: звездная пыль\\
\textbf{tr6c} [\emph{trəʃ}] 0xEDCA: щитовидная железа\\
\textbf{tr6f} [\emph{trəf}] 0xCDCA: топография\\
\textbf{tr6h} [\emph{trəʰ}] 0x6E57: движение вверх по течению\\
\textbf{tr6k} [\emph{trək}] 0x2DCA: стартер\\
\textbf{tr6m} [\emph{trəm}] 0x0DCA: термометр\\
\textbf{tr6n} [\emph{trən}] 0x6DCA: стерео\\
\textbf{tr6p} [\emph{trəp}] 0x4DCA: пястная\\
\textbf{tr6s} [\emph{trəs}] 0x8DCA: живородящий\\
\textbf{tr6t} [\emph{trət}] 0xADCA: внеземной\\
\textbf{tra2c} [\emph{trä˩ʃ}] 0xF9CA: переросший\\
\textbf{tra2f} [\emph{trä˩f}] 0xD9CA: картография\\
\textbf{tra2k} [\emph{trä˩k}] 0x39CA: зарядка\\
\textbf{tra2m} [\emph{trä˩m}] 0x19CA: опилки\\
\textbf{tra2n} [\emph{trä˩n}] 0x79CA: сносно\\
\textbf{tra2p} [\emph{trä˩p}] 0x59CA: преступил\\
\textbf{tra2s} [\emph{trä˩s}] 0x99CA: устало тащиться\\
\textbf{tra7c} [\emph{trä˥ʃ}] 0xF1CA: патриархальный\\
\textbf{tra7f} [\emph{trä˥f}] 0xD1CA: вице-король\\
\textbf{tra7h} [\emph{trä˥ʰ}] 0x8E57: природный\\
\textbf{tra7k} [\emph{trä˥k}] 0x31CA: Звездный свет\\
\textbf{tra7m} [\emph{trä˥m}] 0x11CA: каракуль\\
\textbf{tra7n} [\emph{trä˥n}] 0x71CA: невыносимый\\
\textbf{tra7p} [\emph{trä˥p}] 0x51CA: атропин\\
\textbf{tra7s} [\emph{trä˥s}] 0x91CA: три четверти\\
\textbf{tra7t} [\emph{trä˥t}] 0xB1CA: тащили\\
\textbf{trac} [\emph{träʃ}] 0xE9CA: сопротивление\\
\textbf{traf} [\emph{träf}] 0xC9CA: полоса\\
\textbf{trah} [\emph{träʰ}] 0x4E57: сравнительная степень\\
\textbf{trak} [\emph{träk}] 0x29CA: годовой\\
\textbf{tram} [\emph{träm}] 0x09CA: Рамка\\
\textbf{tran} [\emph{trän}] 0x69CA: отчаяние\\
\textbf{trap} [\emph{träp}] 0x49CA: либеральная\\
\textbf{tras} [\emph{träs}] 0x89CA: сравнение\\
\textbf{trat} [\emph{trät}] 0xA9CA: Распечатать\\
\textbf{tre2c} [\emph{tre˩ʃ}] 0xFBCA: экстровертность\\
\textbf{tre2k} [\emph{tre˩k}] 0x3BCA: гетерозиготный\\
\textbf{tre2n} [\emph{tre˩n}] 0x7BCA: токарный станок\\
\textbf{tre2s} [\emph{tre˩s}] 0x9BCA: теократия\\
\textbf{tre2t} [\emph{tre˩t}] 0xBBCA: медлить\\
\textbf{tre7c} [\emph{tre˥ʃ}] 0xF3CA: тетралогия\\
\textbf{tre7k} [\emph{tre˥k}] 0x33CA: зимородок\\
\textbf{tre7n} [\emph{tre˥n}] 0x73CA: терьер\\
\textbf{tre7s} [\emph{tre˥s}] 0x93CA: этрусский\\
\textbf{tre7t} [\emph{tre˥t}] 0xB3CA: терракотовый\\
\textbf{trec} [\emph{treʃ}] 0xEBCA: стратегия\\
\textbf{tref} [\emph{tref}] 0xCBCA: скипидар\\
\textbf{treh} [\emph{treʰ}] 0x5E57: тщеславный\\
\textbf{trek} [\emph{trek}] 0x2BCA: подрядчик\\
\textbf{trem} [\emph{trem}] 0x0BCA: личное снаряжение\\
\textbf{tren} [\emph{tren}] 0x6BCA: Cегодня\\
\textbf{trep} [\emph{trep}] 0x4BCA: клубника\\
\textbf{tres} [\emph{tres}] 0x8BCA: стресс\\
\textbf{tret} [\emph{tret}] 0xABCA: театр\\
\textbf{tri2c} [\emph{tri˩ʃ}] 0xF8CA: интровертированный\\
\textbf{tri2k} [\emph{tri˩k}] 0x38CA: интервалы\\
\textbf{tri2m} [\emph{tri˩m}] 0x18CA: с электроприводом\\
\textbf{tri2n} [\emph{tri˩n}] 0x78CA: поразительный\\
\textbf{tri2p} [\emph{tri˩p}] 0x58CA: червивый\\
\textbf{tri2s} [\emph{tri˩s}] 0x98CA: повторять\\
\textbf{tri7c} [\emph{tri˥ʃ}] 0xF0CA: печатная машинка\\
\textbf{tri7f} [\emph{tri˥f}] 0xD0CA: канату\\
\textbf{tri7h} [\emph{tri˥ʰ}] 0x8657: мимо пассивный причастие\\
\textbf{tri7k} [\emph{tri˥k}] 0x30CA: Терри\\
\textbf{tri7m} [\emph{tri˥m}] 0x10CA: трирема\\
\textbf{tri7n} [\emph{tri˥n}] 0x70CA: эфирный\\
\textbf{tri7p} [\emph{tri˥p}] 0x50CA: щетина\\
\textbf{tri7s} [\emph{tri˥s}] 0x90CA: зеркала\\
\textbf{tri7t} [\emph{tri˥t}] 0xB0CA: вышеупомянутый\\
\textbf{tric} [\emph{triʃ}] 0xE8CA: временный\\
\textbf{trif} [\emph{trif}] 0xC8CA: география\\
\textbf{trih} [\emph{triʰ}] 0x4657: Сортировать\\
\textbf{trik} [\emph{trik}] 0x28CA: пьяный\\
\textbf{trim} [\emph{trim}] 0x08CA: банка\\
\textbf{trin} [\emph{trin}] 0x68CA: круг\\
\textbf{trip} [\emph{trip}] 0x48CA: кнут\\
\textbf{tris} [\emph{tris}] 0x88CA: пурпурный\\
\textbf{trit} [\emph{trit}] 0xA8CA: непосредственный доказательный\\
\textbf{tro2k} [\emph{tro˩k}] 0x3CCA: гетеросексуальный\\
\textbf{tro2n} [\emph{tro˩n}] 0x7CCA: токарь\\
\textbf{tro2p} [\emph{tro˩p}] 0x5CCA: таро\\
\textbf{tro2s} [\emph{tro˩s}] 0x9CCA: Телец\\
\textbf{tro2t} [\emph{tro˩t}] 0xBCCA: Трой\\
\textbf{tro7k} [\emph{tro˥k}] 0x34CA: стереоскоп\\
\textbf{tro7m} [\emph{tro˥m}] 0x14CA: торий\\
\textbf{tro7n} [\emph{tro˥n}] 0x74CA: трио\\
\textbf{tro7p} [\emph{tro˥p}] 0x54CA: трипсы\\
\textbf{tro7s} [\emph{tro˥s}] 0x94CA: Тироль\\
\textbf{tro7t} [\emph{tro˥t}] 0xB4CA: тележки\\
\textbf{troc} [\emph{troʃ}] 0xECCA: факел\\
\textbf{trof} [\emph{trof}] 0xCCCA: непереходный\\
\textbf{troh} [\emph{troʰ}] 0x6657: непосредственный доказательный\\
\textbf{trok} [\emph{trok}] 0x2CCA: Головная боль\\
\textbf{trom} [\emph{trom}] 0x0CCA: тележка\\
\textbf{tron} [\emph{tron}] 0x6CCA: оранжевый\\
\textbf{trop} [\emph{trop}] 0x4CCA: зонд\\
\textbf{tros} [\emph{tros}] 0x8CCA: тост\\
\textbf{trot} [\emph{trot}] 0xACCA: температура\\
\textbf{tru2k} [\emph{tru˩k}] 0x3ACA: турчанка\\
\textbf{tru2n} [\emph{tru˩n}] 0x7ACA: Турин\\
\textbf{tru2s} [\emph{tru˩s}] 0x9ACA: плутократия\\
\textbf{tru2t} [\emph{tru˩t}] 0xBACA: антихрист\\
\textbf{tru7k} [\emph{tru˥k}] 0x32CA: через\\
\textbf{tru7n} [\emph{tru˥n}] 0x72CA: Тюрингия\\
\textbf{tru7s} [\emph{tru˥s}] 0x92CA: импровизировать\\
\textbf{tru7t} [\emph{tru˥t}] 0xB2CA: нарушитель\\
\textbf{truc} [\emph{truʃ}] 0xEACA: украсть\\
\textbf{truf} [\emph{truf}] 0xCACA: страусы\\
\textbf{truh} [\emph{truʰ}] 0x5657: температура\\
\textbf{truk} [\emph{truk}] 0x2ACA: стучать\\
\textbf{trum} [\emph{trum}] 0x0ACA: экстремист\\
\textbf{trun} [\emph{trun}] 0x6ACA: посадка\\
\textbf{trup} [\emph{trup}] 0x4ACA: атрибут\\
\textbf{trus} [\emph{trus}] 0x8ACA: спорный\\
\textbf{trut} [\emph{trut}] 0xAACA: статуя\\
\textbf{ts62s} [\emph{tsə˩s}] 0x9D4A: триоксида\\
\textbf{ts62t} [\emph{tsə˩t}] 0xBD4A: семьдесят семь\\
\textbf{ts67n} [\emph{tsə˥n}] 0x754A: да сэр\\
\textbf{ts67t} [\emph{tsə˥t}] 0xB54A: Троица\\
\textbf{ts6c} [\emph{tsəʃ}] 0xED4A: Танзания\\
\textbf{ts6f} [\emph{tsəf}] 0xCD4A: стратосфера\\
\textbf{ts6h} [\emph{tsəʰ}] 0x6A57: intransitive case, the subject of an intransitive verb.\\
\textbf{ts6k} [\emph{tsək}] 0x2D4A: земной\\
\textbf{ts6m} [\emph{tsəm}] 0x0D4A: набивка чучел\\
\textbf{ts6n} [\emph{tsən}] 0x6D4A: растворение\\
\textbf{ts6p} [\emph{tsəp}] 0x4D4A: горячая точка\\
\textbf{ts6s} [\emph{tsəs}] 0x8D4A: транс цифра\\
\textbf{ts6t} [\emph{tsət}] 0xAD4A: коническая\\
\textbf{tsa2c} [\emph{tsä˩ʃ}] 0xF94A: кипятильник\\
\textbf{tsa2f} [\emph{tsä˩f}] 0xD94A: марсель\\
\textbf{tsa2k} [\emph{tsä˩k}] 0x394A: драматург\\
\textbf{tsa2m} [\emph{tsä˩m}] 0x194A: стеньга\\
\textbf{tsa2n} [\emph{tsä˩n}] 0x794A: станций\\
\textbf{tsa2p} [\emph{tsä˩p}] 0x594A: мусорный ящик\\
\textbf{tsa2s} [\emph{tsä˩s}] 0x994A: задумчивый\\
\textbf{tsa2t} [\emph{tsä˩t}] 0xB94A: трижды\\
\textbf{tsa7c} [\emph{tsä˥ʃ}] 0xF14A: девяносто семь\\
\textbf{tsa7f} [\emph{tsä˥f}] 0xD14A: самореференция\\
\textbf{tsa7h} [\emph{tsä˥ʰ}] 0x8A57: antipassive голос\\
\textbf{tsa7k} [\emph{tsä˥k}] 0x314A: лавировали\\
\textbf{tsa7m} [\emph{tsä˥m}] 0x114A: человек в руках\\
\textbf{tsa7n} [\emph{tsä˥n}] 0x714A: государственные мужи\\
\textbf{tsa7p} [\emph{tsä˥p}] 0x514A: девяносто восемь\\
\textbf{tsa7s} [\emph{tsä˥s}] 0x914A: шелест\\
\textbf{tsa7t} [\emph{tsä˥t}] 0xB14A: мужественность\\
\textbf{tsac} [\emph{tsäʃ}] 0xE94A: довольно часто\\
\textbf{tsaf} [\emph{tsäf}] 0xC94A: самоубийство\\
\textbf{tsah} [\emph{tsäʰ}] 0x4A57: демонстративный\\
\textbf{tsak} [\emph{tsäk}] 0x294A: уважение\\
\textbf{tsam} [\emph{tsäm}] 0x094A: непосредственно\\
\textbf{tsan} [\emph{tsän}] 0x694A: надлежащий имя существительное\\
\textbf{tsap} [\emph{tsäp}] 0x494A: остров\\
\textbf{tsas} [\emph{tsäs}] 0x894A: бросать\\
\textbf{tsat} [\emph{tsät}] 0xA94A: потенциал настроение\\
\textbf{tse2c} [\emph{tse˩ʃ}] 0xFB4A: арендодатель\\
\textbf{tse2k} [\emph{tse˩k}] 0x3B4A: постельный клоп\\
\textbf{tse2n} [\emph{tse˩n}] 0x7B4A: двадцать секунд\\
\textbf{tse2p} [\emph{tse˩p}] 0x5B4A: в середине сентября\\
\textbf{tse2s} [\emph{tse˩s}] 0x9B4A: тридцатый\\
\textbf{tse2t} [\emph{tse˩t}] 0xBB4A: ч.л.\\
\textbf{tse7c} [\emph{tse˥ʃ}] 0xF34A: проверенный временем\\
\textbf{tse7f} [\emph{tse˥f}] 0xD34A: сохранение времени\\
\textbf{tse7k} [\emph{tse˥k}] 0x334A: Огниво\\
\textbf{tse7m} [\emph{tse˥m}] 0x134A: время чая\\
\textbf{tse7n} [\emph{tse˥n}] 0x734A: девяносто пять\\
\textbf{tse7p} [\emph{tse˥p}] 0x534A: девяносто шесть\\
\textbf{tse7s} [\emph{tse˥s}] 0x934A: техасец\\
\textbf{tse7t} [\emph{tse˥t}] 0xB34A: время отхода ко сну\\
\textbf{tsec} [\emph{tseʃ}] 0xEB4A: эксперт\\
\textbf{tsef} [\emph{tsef}] 0xCB4A: antessive дело\\
\textbf{tseh} [\emph{tseʰ}] 0x5A57: potential mood, something that is probable according to the speaker.\\
\textbf{tsek} [\emph{tsek}] 0x2B4A: стейк\\
\textbf{tsem} [\emph{tsem}] 0x0B4A: атеизм\\
\textbf{tsen} [\emph{tsen}] 0x6B4A: обработка\\
\textbf{tsep} [\emph{tsep}] 0x4B4A: кузнечик\\
\textbf{tses} [\emph{tses}] 0x8B4A: 13\\
\textbf{tset} [\emph{tset}] 0xAB4A: жажда мести\\
\textbf{tsi2c} [\emph{tsi˩ʃ}] 0xF84A: самопоглощение\\
\textbf{tsi2f} [\emph{tsi˩f}] 0xD84A: восьмидесятый\\
\textbf{tsi2h} [\emph{tsi˩ʰ}] 0xC257: послелоги дело\\
\textbf{tsi2k} [\emph{tsi˩k}] 0x384A: папоротник-орляк\\
\textbf{tsi2m} [\emph{tsi˩m}] 0x184A: таксометр\\
\textbf{tsi2p} [\emph{tsi˩p}] 0x584A: семьдесят один\\
\textbf{tsi2t} [\emph{tsi˩t}] 0xB84A: благочестивый\\
\textbf{tsi7c} [\emph{tsi˥ʃ}] 0xF04A: девяносто девять\\
\textbf{tsi7f} [\emph{tsi˥f}] 0xD04A: смокинг\\
\textbf{tsi7h} [\emph{tsi˥ʰ}] 0x8257: antessive case, indicates the noun phrase precedes the rest of the sentence. space-context retrospective-aspect \\
\textbf{tsi7k} [\emph{tsi˥k}] 0x304A: капающий\\
\textbf{tsi7m} [\emph{tsi˥m}] 0x104A: атавизм\\
\textbf{tsi7n} [\emph{tsi˥n}] 0x704A: семьдесят пять\\
\textbf{tsi7p} [\emph{tsi˥p}] 0x504A: шепелявость\\
\textbf{tsi7s} [\emph{tsi˥s}] 0x904A: двадцать три\\
\textbf{tsi7t} [\emph{tsi˥t}] 0xB04A: отдыхает\\
\textbf{tsic} [\emph{tsiʃ}] 0xE84A: пытаясь\\
\textbf{tsif} [\emph{tsif}] 0xC84A: стимул\\
\textbf{tsih} [\emph{tsiʰ}] 0x4257: Вот\\
\textbf{tsik} [\emph{tsik}] 0x284A: злость\\
\textbf{tsim} [\emph{tsim}] 0x084A: время года\\
\textbf{tsip} [\emph{tsip}] 0x484A: грех\\
\textbf{tsis} [\emph{tsis}] 0x884A: справедливость\\
\textbf{tsit} [\emph{tsit}] 0xA84A: остановился\\
\textbf{tso2k} [\emph{tso˩k}] 0x3C4A: в середине августа\\
\textbf{tso2n} [\emph{tso˩n}] 0x7C4A: триколор\\
\textbf{tso2s} [\emph{tso˩s}] 0x9C4A: фотосинтез\\
\textbf{tso2t} [\emph{tso˩t}] 0xBC4A: шестьдесят четыре\\
\textbf{tso7c} [\emph{tso˥ʃ}] 0xF44A: непримиримость\\
\textbf{tso7k} [\emph{tso˥k}] 0x344A: контрабандист\\
\textbf{tso7m} [\emph{tso˥m}] 0x144A: газонокосилка\\
\textbf{tso7n} [\emph{tso˥n}] 0x744A: Тоскана\\
\textbf{tso7p} [\emph{tso˥p}] 0x544A: фронтиспис\\
\textbf{tso7s} [\emph{tso˥s}] 0x944A: девяносто четыре\\
\textbf{tso7t} [\emph{tso˥t}] 0xB44A: поверь\\
\textbf{tsoc} [\emph{tsoʃ}] 0xEC4A: поиск\\
\textbf{tsof} [\emph{tsof}] 0xCC4A: antipassive голос\\
\textbf{tsoh} [\emph{tsoʰ}] 0x6257: уважение\\
\textbf{tsok} [\emph{tsok}] 0x2C4A: риторика\\
\textbf{tsom} [\emph{tsom}] 0x0C4A: переходный\\
\textbf{tson} [\emph{tson}] 0x6C4A: учреждения\\
\textbf{tsop} [\emph{tsop}] 0x4C4A: движение в гору\\
\textbf{tsos} [\emph{tsos}] 0x8C4A: конкурс\\
\textbf{tsot} [\emph{tsot}] 0xAC4A: верность\\
\textbf{tsu2n} [\emph{tsu˩n}] 0x7A4A: превращение\\
\textbf{tsu2s} [\emph{tsu˩s}] 0x9A4A: трипсин\\
\textbf{tsu2t} [\emph{tsu˩t}] 0xBA4A: шестьдесят два\\
\textbf{tsu7k} [\emph{tsu˥k}] 0x324A: щупальце\\
\textbf{tsu7n} [\emph{tsu˥n}] 0x724A: правнук\\
\textbf{tsu7s} [\emph{tsu˥s}] 0x924A: антацидный\\
\textbf{tsu7t} [\emph{tsu˥t}] 0xB24A: шестьдесят три\\
\textbf{tsuc} [\emph{tsuʃ}] 0xEA4A: куст\\
\textbf{tsuf} [\emph{tsuf}] 0xCA4A: Тунис\\
\textbf{tsuh} [\emph{tsuʰ}] 0x5257: дистальный демонстративный\\
\textbf{tsuk} [\emph{tsuk}] 0x2A4A: ранг\\
\textbf{tsum} [\emph{tsum}] 0x0A4A: вольфрам\\
\textbf{tsun} [\emph{tsun}] 0x6A4A: деревня\\
\textbf{tsup} [\emph{tsup}] 0x4A4A: раздробить\\
\textbf{tsus} [\emph{tsus}] 0x8A4A: такси\\
\textbf{tsut} [\emph{tsut}] 0xAA4A: врач\\
\textbf{tu} [\emph{tu}] 0x295E: deontic mood, to indicate what should be different.\\
\textbf{tw6k} [\emph{twək}] 0x2D2A: транслокация\\
\textbf{tw6n} [\emph{twən}] 0x6D2A: Тайвань\\
\textbf{tw6s} [\emph{twəs}] 0x8D2A: партеногенез\\
\textbf{tw6t} [\emph{twət}] 0xAD2A: затемнение\\
\textbf{twa2c} [\emph{twä˩ʃ}] 0xF92A: останец\\
\textbf{twa2k} [\emph{twä˩k}] 0x392A: командная работа\\
\textbf{twa2m} [\emph{twä˩m}] 0x192A: креационизм\\
\textbf{twa2n} [\emph{twä˩n}] 0x792A: самооборона\\
\textbf{twa2s} [\emph{twä˩s}] 0x992A: синица\\
\textbf{twa2t} [\emph{twä˩t}] 0xB92A: паническое бегство\\
\textbf{twa7k} [\emph{twä˥k}] 0x312A: собиратель древностей\\
\textbf{twa7n} [\emph{twä˥n}] 0x712A: желтовато-коричневый\\
\textbf{twa7s} [\emph{twä˥s}] 0x912A: девяносто три\\
\textbf{twa7t} [\emph{twä˥t}] 0xB12A: переплетенный\\
\textbf{twac} [\emph{twäʃ}] 0xE92A: окно\\
\textbf{twaf} [\emph{twäf}] 0xC92A: тривиальный\\
\textbf{twah} [\emph{twäʰ}] 0x4957: elative case, indicates motion `out of' some noun phrase. interior-context source-case.\\
\textbf{twak} [\emph{twäk}] 0x292A: delative дело\\
\textbf{twam} [\emph{twäm}] 0x092A: Запрос\\
\textbf{twan} [\emph{twän}] 0x692A: душа\\
\textbf{twap} [\emph{twäp}] 0x492A: прекращение аспект\\
\textbf{twas} [\emph{twäs}] 0x892A: дважды\\
\textbf{twat} [\emph{twät}] 0xA92A: война\\
\textbf{twe2k} [\emph{twe˩k}] 0x3B2A: эктопический\\
\textbf{twe2t} [\emph{twe˩t}] 0xBB2A: лапками\\
\textbf{twe7k} [\emph{twe˥k}] 0x332A: т квадрат\\
\textbf{twe7n} [\emph{twe˥n}] 0x732A: гетеродинный\\
\textbf{twe7t} [\emph{twe˥t}] 0xB32A: е текст\\
\textbf{twec} [\emph{tweʃ}] 0xEB2A: дистанционный пульт будущее напряженный\\
\textbf{twef} [\emph{twef}] 0xCB2A: ультрафиолетовый\\
\textbf{tweh} [\emph{tweʰ}] 0x5957: terminative case, indicates temporal destination or goal. time-context destination-case. \\
\textbf{twek} [\emph{twek}] 0x2B2A: чрезмерное дело\\
\textbf{twem} [\emph{twem}] 0x0B2A: трехмерный\\
\textbf{twen} [\emph{twen}] 0x6B2A: ткать\\
\textbf{twep} [\emph{twep}] 0x4B2A: переплачивать\\
\textbf{twes} [\emph{twes}] 0x8B2A: телевидение\\
\textbf{twet} [\emph{twet}] 0xAB2A: юбка\\
\textbf{twi2k} [\emph{twi˩k}] 0x382A: тахикардия\\
\textbf{twi2n} [\emph{twi˩n}] 0x782A: супница\\
\textbf{twi2t} [\emph{twi˩t}] 0xB82A: опросный лист\\
\textbf{twi7k} [\emph{twi˥k}] 0x302A: стеклярус\\
\textbf{twi7n} [\emph{twi˥n}] 0x702A: пятьдесят три\\
\textbf{twi7p} [\emph{twi˥p}] 0x502A: лакомый кусок\\
\textbf{twi7s} [\emph{twi˥s}] 0x902A: пятидесятый\\
\textbf{twi7t} [\emph{twi˥t}] 0xB02A: трактирщик\\
\textbf{twic} [\emph{twiʃ}] 0xE82A: остатки\\
\textbf{twif} [\emph{twif}] 0xC82A: данник\\
\textbf{twih} [\emph{twiʰ}] 0x4157: ilative case, indicates motion `into' some noun phrase. interior-context destination-case.\\
\textbf{twik} [\emph{twik}] 0x282A: terminative дело\\
\textbf{twim} [\emph{twim}] 0x082A: дезидеративный настроение\\
\textbf{twin} [\emph{twin}] 0x682A: желание\\
\textbf{twip} [\emph{twip}] 0x482A: начинательный аспект\\
\textbf{twis} [\emph{twis}] 0x882A: быстро\\
\textbf{twit} [\emph{twit}] 0xA82A: утвердительный квантификатором\\
\textbf{two2n} [\emph{two˩n}] 0x7C2A: потрясающее зрелище\\
\textbf{two7k} [\emph{two˥k}] 0x342A: Тонганский\\
\textbf{two7n} [\emph{two˥n}] 0x742A: антициклон\\
\textbf{two7s} [\emph{two˥s}] 0x942A: Тоголезская\\
\textbf{two7t} [\emph{two˥t}] 0xB42A: трезвенник\\
\textbf{twoc} [\emph{twoʃ}] 0xEC2A: продуктовый\\
\textbf{twof} [\emph{twof}] 0xCC2A: Ботсвана\\
\textbf{twoh} [\emph{twoʰ}] 0x6157: quotative case, indicates a quote. discource-context location-case.\\
\textbf{twok} [\emph{twok}] 0x2C2A: статистическая\\
\textbf{twom} [\emph{twom}] 0x0C2A: внутренняя связь\\
\textbf{twon} [\emph{twon}] 0x6C2A: орнамент\\
\textbf{twop} [\emph{twop}] 0x4C2A: тромбон\\
\textbf{twos} [\emph{twos}] 0x8C2A: терроризм\\
\textbf{twot} [\emph{twot}] 0xAC2A: Мероприятия\\
\textbf{twu2k} [\emph{twu˩k}] 0x3A2A: середина недели\\
\textbf{twu7k} [\emph{twu˥k}] 0x322A: библиофил\\
\textbf{twu7n} [\emph{twu˥n}] 0x722A: мучнистая горлышком\\
\textbf{twu7t} [\emph{twu˥t}] 0xB22A: трубочка\\
\textbf{twuc} [\emph{twuʃ}] 0xEA2A: дым\\
\textbf{twuh} [\emph{twuʰ}] 0x5157: dubitative mood, for saying things the speaker is doubtful of, as for expressing sarcasm.\\
\textbf{twuk} [\emph{twuk}] 0x2A2A: очистка\\
\textbf{twun} [\emph{twun}] 0x6A2A: 9\\
\textbf{twup} [\emph{twup}] 0x4A2A: контрапункт\\
\textbf{twus} [\emph{twus}] 0x8A2A: альтруизм\\
\textbf{twut} [\emph{twut}] 0xAA2A: dubitative настроение\\
\textbf{tx62n} [\emph{txə˩n}] 0x7DEA: Тегеран\\
\textbf{tx62t} [\emph{txə˩t}] 0xBDEA: держатель билетов\\
\textbf{tx67n} [\emph{txə˥n}] 0x75EA: многозадачных\\
\textbf{tx67t} [\emph{txə˥t}] 0xB5EA: тик-так\\
\textbf{tx6c} [\emph{txəʃ}] 0xEDEA: обнаружений\\
\textbf{tx6k} [\emph{txək}] 0x2DEA: стетоскоп\\
\textbf{tx6m} [\emph{txəm}] 0x0DEA: катамаран\\
\textbf{tx6n} [\emph{txən}] 0x6DEA: испарение\\
\textbf{tx6p} [\emph{txəp}] 0x4DEA: остеопатия\\
\textbf{tx6s} [\emph{txəs}] 0x8DEA: Туркменистан\\
\textbf{tx6t} [\emph{txət}] 0xADEA: текст редактор\\
\textbf{txa2c} [\emph{txä˩ʃ}] 0xF9EA: анти герой\\
\textbf{txa2f} [\emph{txä˩f}] 0xD9EA: перетягивание каната\\
\textbf{txa2k} [\emph{txä˩k}] 0x39EA: дрова\\
\textbf{txa2m} [\emph{txä˩m}] 0x19EA: козыри\\
\textbf{txa2n} [\emph{txä˩n}] 0x79EA: вызванный\\
\textbf{txa2p} [\emph{txä˩p}] 0x59EA: Марки\\
\textbf{txa2s} [\emph{txä˩s}] 0x99EA: ходули\\
\textbf{txa2t} [\emph{txä˩t}] 0xB9EA: Attends\\
\textbf{txa7c} [\emph{txä˥ʃ}] 0xF1EA: таитянского\\
\textbf{txa7f} [\emph{txä˥f}] 0xD1EA: заполнитель\\
\textbf{txa7k} [\emph{txä˥k}] 0x31EA: вертячка\\
\textbf{txa7m} [\emph{txä˥m}] 0x11EA: зенитная\\
\textbf{txa7n} [\emph{txä˥n}] 0x71EA: палатки\\
\textbf{txa7p} [\emph{txä˥p}] 0x51EA: хорошо носится\\
\textbf{txa7s} [\emph{txä˥s}] 0x91EA: дрок\\
\textbf{txa7t} [\emph{txä˥t}] 0xB1EA: спутанный\\
\textbf{txac} [\emph{txäʃ}] 0xE9EA: капитуляция\\
\textbf{txaf} [\emph{txäf}] 0xC9EA: телеграф\\
\textbf{txah} [\emph{txäʰ}] 0x4F57: местоимение маркер\\
\textbf{txak} [\emph{txäk}] 0x29EA: срабатывает\\
\textbf{txam} [\emph{txäm}] 0x09EA: тантал\\
\textbf{txan} [\emph{txän}] 0x69EA: актиний\\
\textbf{txap} [\emph{txäp}] 0x49EA: положение дел бар\\
\textbf{txas} [\emph{txäs}] 0x89EA: установка\\
\textbf{txat} [\emph{txät}] 0xA9EA: агент по продаже недвижимости\\
\textbf{txe2k} [\emph{txe˩k}] 0x3BEA: коленный рефлекс\\
\textbf{txe2n} [\emph{txe˩n}] 0x7BEA: агрегирование\\
\textbf{txe2s} [\emph{txe˩s}] 0x9BEA: отвислый\\
\textbf{txe2t} [\emph{txe˩t}] 0xBBEA: Билеты\\
\textbf{txe7k} [\emph{txe˥k}] 0x33EA: нога\\
\textbf{txe7n} [\emph{txe˥n}] 0x73EA: тенор\\
\textbf{txe7s} [\emph{txe˥s}] 0x93EA: поймать как улова может\\
\textbf{txe7t} [\emph{txe˥t}] 0xB3EA: мститель\\
\textbf{txec} [\emph{txeʃ}] 0xEBEA: телекоммуникация\\
\textbf{txef} [\emph{txef}] 0xCBEA: термосфера\\
\textbf{txeh} [\emph{txeʰ}] 0x5F57: увеличительный\\
\textbf{txek} [\emph{txek}] 0x2BEA: п- корень\\
\textbf{txem} [\emph{txem}] 0x0BEA: шашки\\
\textbf{txen} [\emph{txen}] 0x6BEA: переложение\\
\textbf{txep} [\emph{txep}] 0x4BEA: утроить\\
\textbf{txes} [\emph{txes}] 0x8BEA: коллектив имя существительное\\
\textbf{txet} [\emph{txet}] 0xABEA: в чате\\
\textbf{txi2c} [\emph{txi˩ʃ}] 0xF8EA: триптих\\
\textbf{txi2k} [\emph{txi˩k}] 0x38EA: незрелость\\
\textbf{txi2n} [\emph{txi˩n}] 0x78EA: покорный\\
\textbf{txi2p} [\emph{txi˩p}] 0x58EA: опечатка\\
\textbf{txi2s} [\emph{txi˩s}] 0x98EA: семьдесят три\\
\textbf{txi2t} [\emph{txi˩t}] 0xB8EA: замечающий\\
\textbf{txi7c} [\emph{txi˥ʃ}] 0xF0EA: трезубец\\
\textbf{txi7h} [\emph{txi˥ʰ}] 0x8757: коллектив имя существительное\\
\textbf{txi7k} [\emph{txi˥k}] 0x30EA: подстроенный\\
\textbf{txi7n} [\emph{txi˥n}] 0x70EA: похвальный\\
\textbf{txi7p} [\emph{txi˥p}] 0x50EA: лежал в\\
\textbf{txi7s} [\emph{txi˥s}] 0x90EA: оптимист\\
\textbf{txi7t} [\emph{txi˥t}] 0xB0EA: замещенный\\
\textbf{txic} [\emph{txiʃ}] 0xE8EA: вторгшийся\\
\textbf{txif} [\emph{txif}] 0xC8EA: ссылаясь\\
\textbf{txih} [\emph{txiʰ}] 0x4757: excessive case, indicates motion from a state. state-context source-case.\\
\textbf{txik} [\emph{txik}] 0x28EA: указательный палец\\
\textbf{txim} [\emph{txim}] 0x08EA: этидия\\
\textbf{txin} [\emph{txin}] 0x68EA: рафинирование\\
\textbf{txip} [\emph{txip}] 0x48EA: тройной точкой\\
\textbf{txis} [\emph{txis}] 0x88EA: точка матрица\\
\textbf{txit} [\emph{txit}] 0xA8EA: шрифт\\
\textbf{txo2k} [\emph{txo˩k}] 0x3CEA: стафилококк\\
\textbf{txo2n} [\emph{txo˩n}] 0x7CEA: низкий скатных\\
\textbf{txo2t} [\emph{txo˩t}] 0xBCEA: право аренды\\
\textbf{txo7k} [\emph{txo˥k}] 0x34EA: замочная скважина\\
\textbf{txo7m} [\emph{txo˥m}] 0x14EA: Tom Tom\\
\textbf{txo7n} [\emph{txo˥n}] 0x74EA: закал\\
\textbf{txo7s} [\emph{txo˥s}] 0x94EA: чаны\\
\textbf{txo7t} [\emph{txo˥t}] 0xB4EA: фригольд\\
\textbf{txoc} [\emph{txoʃ}] 0xECEA: жеребенок\\
\textbf{txof} [\emph{txof}] 0xCCEA: фаршированный\\
\textbf{txoh} [\emph{txoʰ}] 0x6757: скромный\\
\textbf{txok} [\emph{txok}] 0x2CEA: автоматически\\
\textbf{txom} [\emph{txom}] 0x0CEA: Оман\\
\textbf{txon} [\emph{txon}] 0x6CEA: ругань\\
\textbf{txop} [\emph{txop}] 0x4CEA: шахматная доска\\
\textbf{txos} [\emph{txos}] 0x8CEA: средний голос\\
\textbf{txot} [\emph{txot}] 0xACEA: в середине осени\\
\textbf{txu2k} [\emph{txu˩k}] 0x3AEA: индюков\\
\textbf{txu2n} [\emph{txu˩n}] 0x7AEA: скула\\
\textbf{txu2s} [\emph{txu˩s}] 0x9AEA: очень большой размер\\
\textbf{txu2t} [\emph{txu˩t}] 0xBAEA: трудно носатый\\
\textbf{txu7n} [\emph{txu˥n}] 0x72EA: австро-венгерской\\
\textbf{txu7t} [\emph{txu˥t}] 0xB2EA: сосок\\
\textbf{txuc} [\emph{txuʃ}] 0xEAEA: выпуклость\\
\textbf{txuf} [\emph{txuf}] 0xCAEA: пирогов\\
\textbf{txuh} [\emph{txuʰ}] 0x5757: приступ боли\\
\textbf{txuk} [\emph{txuk}] 0x2AEA: короткое замыкание\\
\textbf{txum} [\emph{txum}] 0x0AEA: локтевая кость\\
\textbf{txun} [\emph{txun}] 0x6AEA: исследование Пол\\
\textbf{txup} [\emph{txup}] 0x4AEA: турмалин\\
\textbf{txus} [\emph{txus}] 0x8AEA: тушить\\
\textbf{txut} [\emph{txut}] 0xAAEA: значительность\\
\textbf{ty6k} [\emph{tjək}] 0x2D0A: восьмиугольник\\
\textbf{ty6n} [\emph{tjən}] 0x6D0A: касательный\\
\textbf{ty6p} [\emph{tjəp}] 0x4D0A: в середине апреля\\
\textbf{ty6s} [\emph{tjəs}] 0x8D0A: подземные толчки\\
\textbf{ty6t} [\emph{tjət}] 0xAD0A: Чад\\
\textbf{tya2c} [\emph{tjä˩ʃ}] 0xF90A: анте Никейский\\
\textbf{tya2k} [\emph{tjä˩k}] 0x390A: автобиографический\\
\textbf{tya2m} [\emph{tjä˩m}] 0x190A: успокойся\\
\textbf{tya2n} [\emph{tjä˩n}] 0x790A: утреня\\
\textbf{tya2p} [\emph{tjä˩p}] 0x590A: банка открывашка\\
\textbf{tya2s} [\emph{tjä˩s}] 0x990A: катехизис\\
\textbf{tya2t} [\emph{tjä˩t}] 0xB90A: непроверенный\\
\textbf{tya7c} [\emph{tjä˥ʃ}] 0xF10A: трехъязычный\\
\textbf{tya7k} [\emph{tjä˥k}] 0x310A: тапиока\\
\textbf{tya7m} [\emph{tjä˥m}] 0x110A: перевалка\\
\textbf{tya7n} [\emph{tjä˥n}] 0x710A: неизвестности\\
\textbf{tya7p} [\emph{tjä˥p}] 0x510A: летать по ночам\\
\textbf{tya7s} [\emph{tjä˥s}] 0x910A: клеветнический\\
\textbf{tya7t} [\emph{tjä˥t}] 0xB10A: левая рука\\
\textbf{tyac} [\emph{tjäʃ}] 0xE90A: читать\\
\textbf{tyaf} [\emph{tjäf}] 0xC90A: параграф\\
\textbf{tyah} [\emph{tjäʰ}] 0x4857: независимый пункт\\
\textbf{tyak} [\emph{tjäk}] 0x290A: таким образом\\
\textbf{tyam} [\emph{tjäm}] 0x090A: поздно\\
\textbf{tyan} [\emph{tjän}] 0x690A: дверь\\
\textbf{tyap} [\emph{tjäp}] 0x490A: опыт\\
\textbf{tyas} [\emph{tjäs}] 0x890A: Таблица\\
\textbf{tyat} [\emph{tjät}] 0xA90A: старый\\
\textbf{tye2n} [\emph{tje˩n}] 0x7B0A: трехгодичного\\
\textbf{tye2t} [\emph{tje˩t}] 0xBB0A: семьдесят восемь\\
\textbf{tye7k} [\emph{tje˥k}] 0x330A: укладчик\\
\textbf{tye7n} [\emph{tje˥n}] 0x730A: феерия\\
\textbf{tye7t} [\emph{tje˥t}] 0xB30A: предвосхищать\\
\textbf{tyec} [\emph{tjeʃ}] 0xEB0A: тег\\
\textbf{tyef} [\emph{tjef}] 0xCB0A: моток пряжи\\
\textbf{tyeh} [\emph{tjeʰ}] 0x5857: текст\\
\textbf{tyek} [\emph{tjek}] 0x2B0A: треугольник\\
\textbf{tyem} [\emph{tjem}] 0x0B0A: титан\\
\textbf{tyen} [\emph{tjen}] 0x6B0A: общеизвестный\\
\textbf{tyep} [\emph{tjep}] 0x4B0A: штатив\\
\textbf{tyes} [\emph{tjes}] 0x8B0A: поставщик провизии\\
\textbf{tyet} [\emph{tjet}] 0xAB0A: Спортивное\\
\textbf{tyi2k} [\emph{tji˩k}] 0x380A: инталия\\
\textbf{tyi2m} [\emph{tji˩m}] 0x180A: метка времени\\
\textbf{tyi2n} [\emph{tji˩n}] 0x780A: почетной\\
\textbf{tyi2p} [\emph{tji˩p}] 0x580A: хронометраж\\
\textbf{tyi2s} [\emph{tji˩s}] 0x980A: шестидесятый\\
\textbf{tyi2t} [\emph{tji˩t}] 0xB80A: двадцать восемь\\
\textbf{tyi7c} [\emph{tji˥ʃ}] 0xF00A: быть принятым в высшее учебное заведение\\
\textbf{tyi7k} [\emph{tji˥k}] 0x300A: онтологический\\
\textbf{tyi7m} [\emph{tji˥m}] 0x100A: автономный\\
\textbf{tyi7n} [\emph{tji˥n}] 0x700A: диоксид титана\\
\textbf{tyi7p} [\emph{tji˥p}] 0x500A: тайпей\\
\textbf{tyi7s} [\emph{tji˥s}] 0x900A: шарик\\
\textbf{tyi7t} [\emph{tji˥t}] 0xB00A: пунктирный\\
\textbf{tyic} [\emph{tjiʃ}] 0xE80A: природный\\
\textbf{tyif} [\emph{tjif}] 0xC80A: интерпретировать\\
\textbf{tyih} [\emph{tjiʰ}] 0x4057: ретроспективный аспект\\
\textbf{tyik} [\emph{tjik}] 0x280A: дать согласие\\
\textbf{tyim} [\emph{tjim}] 0x080A: Декабрь\\
\textbf{tyin} [\emph{tjin}] 0x680A: три\\
\textbf{tyip} [\emph{tjip}] 0x480A: вести\\
\textbf{tyis} [\emph{tjis}] 0x880A: прогресс\\
\textbf{tyit} [\emph{tjit}] 0xA80A: против\\
\textbf{tyo2n} [\emph{tjo˩n}] 0x7C0A: Океания\\
\textbf{tyo2t} [\emph{tjo˩t}] 0xBC0A: эфирное масло\\
\textbf{tyo7k} [\emph{tjo˥k}] 0x340A: Токийский\\
\textbf{tyo7n} [\emph{tjo˥n}] 0x740A: девяносто один\\
\textbf{tyo7t} [\emph{tjo˥t}] 0xB40A: антиподы\\
\textbf{tyoc} [\emph{tjoʃ}] 0xEC0A: совпадение\\
\textbf{tyof} [\emph{tjof}] 0xCC0A: тасование\\
\textbf{tyoh} [\emph{tjoʰ}] 0x6057: тривиальный\\
\textbf{tyok} [\emph{tjok}] 0x2C0A: грузовая машина\\
\textbf{tyom} [\emph{tjom}] 0x0C0A: сурьма\\
\textbf{tyon} [\emph{tjon}] 0x6C0A: астрономия\\
\textbf{tyop} [\emph{tjop}] 0x4C0A: осьминог\\
\textbf{tyos} [\emph{tjos}] 0x8C0A: доверенное лицо\\
\textbf{tyot} [\emph{tjot}] 0xAC0A: договор\\
\textbf{tyu2n} [\emph{tju˩n}] 0x7A0A: тональный\\
\textbf{tyu2t} [\emph{tju˩t}] 0xBA0A: эвтектика\\
\textbf{tyu7k} [\emph{tju˥k}] 0x320A: тупые под углом\\
\textbf{tyu7n} [\emph{tju˥n}] 0x720A: отнятие ребенка от груди\\
\textbf{tyu7t} [\emph{tju˥t}] 0xB20A: тат тат крыса\\
\textbf{tyuc} [\emph{tjuʃ}] 0xEA0A: ожидание\\
\textbf{tyuf} [\emph{tjuf}] 0xCA0A: перехитрить\\
\textbf{tyuh} [\emph{tjuʰ}] 0x5057: только\\
\textbf{tyuk} [\emph{tjuk}] 0x2A0A: восток\\
\textbf{tyum} [\emph{tjum}] 0x0A0A: рутений атом\\
\textbf{tyun} [\emph{tjun}] 0x6A0A: домен\\
\textbf{tyup} [\emph{tjup}] 0x4A0A: пригород\\
\textbf{tyus} [\emph{tjus}] 0x8A0A: целое число\\
\textbf{tyut} [\emph{tjut}] 0xAA0A: два\\
\subsection{ v }
\textbf{va} [\emph{vä}] 0x26DE: ватт\\
\textbf{ve} [\emph{ve}] 0x2EDE: вектор\\
\textbf{vi} [\emph{vi}] 0x22DE: глагол - суффиксов\\
\textbf{vl6n} [\emph{vlən}] 0x6D76: валлонский\\
\textbf{vl6t} [\emph{vlət}] 0xAD76: миллилитр\\
\textbf{vla2n} [\emph{vlä˩n}] 0x7976: виллан\\
\textbf{vla7n} [\emph{vlä˥n}] 0x7176: овуляция\\
\textbf{vla7t} [\emph{vlä˥t}] 0xB176: земля\\
\textbf{vlac} [\emph{vläʃ}] 0xE976: морская звезда\\
\textbf{vlaf} [\emph{vläf}] 0xC976: пис`ать стихи\\
\textbf{vlah} [\emph{vläʰ}] 0x4BB7: напряжение\\
\textbf{vlak} [\emph{vläk}] 0x2976: стерилизовать\\
\textbf{vlam} [\emph{vläm}] 0x0976: имя файла\\
\textbf{vlan} [\emph{vlän}] 0x6976: прикрепление\\
\textbf{vlap} [\emph{vläp}] 0x4976: амеба\\
\textbf{vlas} [\emph{vläs}] 0x8976: вокализации\\
\textbf{vlat} [\emph{vlät}] 0xA976: квартиры\\
\textbf{vlec} [\emph{vleʃ}] 0xEB76: калька\\
\textbf{vlef} [\emph{vlef}] 0xCB76: виолончель\\
\textbf{vlek} [\emph{vlek}] 0x2B76: надгортанник\\
\textbf{vlen} [\emph{vlen}] 0x6B76: валентность\\
\textbf{vles} [\emph{vles}] 0x8B76: валлийский\\
\textbf{vlet} [\emph{vlet}] 0xAB76: сварной шов\\
\textbf{vli2n} [\emph{vli˩n}] 0x7876: прямая линия\\
\textbf{vli2t} [\emph{vli˩t}] 0xB876: виноградарство\\
\textbf{vli7n} [\emph{vli˥n}] 0x7076: ватерлиния\\
\textbf{vli7t} [\emph{vli˥t}] 0xB076: чревовещатель\\
\textbf{vlic} [\emph{vliʃ}] 0xE876: орнитология\\
\textbf{vlif} [\emph{vlif}] 0xC876: непроизвольный\\
\textbf{vlih} [\emph{vliʰ}] 0x43B7: mineral gender, learning about being aware, 1st density life form\\
\textbf{vlik} [\emph{vlik}] 0x2876: физиология\\
\textbf{vlim} [\emph{vlim}] 0x0876: виртуальный Память\\
\textbf{vlin} [\emph{vlin}] 0x6876: средство передвижения Пол\\
\textbf{vlip} [\emph{vlip}] 0x4876: линкор\\
\textbf{vlis} [\emph{vlis}] 0x8876: виртуальный реальность\\
\textbf{vlit} [\emph{vlit}] 0xA876: говорливость\\
\textbf{vloc} [\emph{vloʃ}] 0xEC76: коэволюция\\
\textbf{vlof} [\emph{vlof}] 0xCC76: револьвер\\
\textbf{vloh} [\emph{vloʰ}] 0x63B7: вольт\\
\textbf{vlok} [\emph{vlok}] 0x2C76: волан\\
\textbf{vlon} [\emph{vlon}] 0x6C76: одуванчик\\
\textbf{vlot} [\emph{vlot}] 0xAC76: заметный\\
\textbf{vluf} [\emph{vluf}] 0xCA76: заспать\\
\textbf{vluk} [\emph{vluk}] 0x2A76: вольта лицо\\
\textbf{vlun} [\emph{vlun}] 0x6A76: линия электропередачи\\
\textbf{vlus} [\emph{vlus}] 0x8A76: девальвация\\
\textbf{vlut} [\emph{vlut}] 0xAA76: сочлененный\\
\textbf{vr6n} [\emph{vrən}] 0x6DD6: благоговейный\\
\textbf{vr6t} [\emph{vrət}] 0xADD6: каратэ\\
\textbf{vra2n} [\emph{vrä˩n}] 0x79D6: купорос\\
\textbf{vra2t} [\emph{vrä˩t}] 0xB9D6: пр`оклятый\\
\textbf{vra7n} [\emph{vrä˥n}] 0x71D6: валериана\\
\textbf{vra7t} [\emph{vrä˥t}] 0xB1D6: вибратор\\
\textbf{vrac} [\emph{vräʃ}] 0xE9D6: лак\\
\textbf{vraf} [\emph{vräf}] 0xC9D6: виброфон\\
\textbf{vrah} [\emph{vräʰ}] 0x4EB7: неудовлетворенность\\
\textbf{vrak} [\emph{vräk}] 0x29D6: парик\\
\textbf{vram} [\emph{vräm}] 0x09D6: телятина\\
\textbf{vran} [\emph{vrän}] 0x69D6: извилистый\\
\textbf{vrap} [\emph{vräp}] 0x49D6: загар\\
\textbf{vras} [\emph{vräs}] 0x89D6: Австралия\\
\textbf{vrat} [\emph{vrät}] 0xA9D6: почтальон\\
\textbf{vrec} [\emph{vreʃ}] 0xEBD6: жадностью\\
\textbf{vrek} [\emph{vrek}] 0x2BD6: электромагнитный\\
\textbf{vren} [\emph{vren}] 0x6BD6: вечнозеленый\\
\textbf{vres} [\emph{vres}] 0x8BD6: электролиз\\
\textbf{vret} [\emph{vret}] 0xABD6: унаследованный\\
\textbf{vri2n} [\emph{vri˩n}] 0x78D6: ивуариец\\
\textbf{vri2t} [\emph{vri˩t}] 0xB8D6: оспа\\
\textbf{vri7n} [\emph{vri˥n}] 0x70D6: активизация\\
\textbf{vri7t} [\emph{vri˥t}] 0xB0D6: продажный\\
\textbf{vric} [\emph{vriʃ}] 0xE8D6: Австрия\\
\textbf{vrif} [\emph{vrif}] 0xC8D6: заднего вида\\
\textbf{vrih} [\emph{vriʰ}] 0x46B7: исторический напряженный\\
\textbf{vrik} [\emph{vrik}] 0x28D6: окончательный частица\\
\textbf{vrim} [\emph{vrim}] 0x08D6: висцеральный\\
\textbf{vrin} [\emph{vrin}] 0x68D6: вербовщик\\
\textbf{vrip} [\emph{vrip}] 0x48D6: абрикос\\
\textbf{vris} [\emph{vris}] 0x88D6: взаимная голос\\
\textbf{vrit} [\emph{vrit}] 0xA8D6: литературный раб\\
\textbf{vrof} [\emph{vrof}] 0xCCD6: евро\\
\textbf{vrok} [\emph{vrok}] 0x2CD6: Корея\\
\textbf{vrom} [\emph{vrom}] 0x0CD6: аорта\\
\textbf{vron} [\emph{vron}] 0x6CD6: позвоночный\\
\textbf{vros} [\emph{vros}] 0x8CD6: эспрессо\\
\textbf{vrot} [\emph{vrot}] 0xACD6: выпад\\
\textbf{vrun} [\emph{vrun}] 0x6AD6: Мавритания\\
\textbf{vrus} [\emph{vrus}] 0x8AD6: Евразия\\
\textbf{vrut} [\emph{vrut}] 0xAAD6: беспозвоночный\\
\textbf{vwac} [\emph{vwäʃ}] 0xE936: расточительность\\
\textbf{vwaf} [\emph{vwäf}] 0xC936: ноу-хау\\
\textbf{vwah} [\emph{vwäʰ}] 0x49B7: движение спуск\\
\textbf{vwak} [\emph{vwäk}] 0x2936: Словакия\\
\textbf{vwan} [\emph{vwän}] 0x6936: сетчатка\\
\textbf{vwas} [\emph{vwäs}] 0x8936: винный камень\\
\textbf{vwat} [\emph{vwät}] 0xA936: Руанда\\
\textbf{vwef} [\emph{vwef}] 0xCB36: уклончиво\\
\textbf{vweh} [\emph{vweʰ}] 0x59B7: перегружены работой\\
\textbf{vwek} [\emph{vwek}] 0x2B36: Авант курьер\\
\textbf{vwen} [\emph{vwen}] 0x6B36: непостоянство\\
\textbf{vwet} [\emph{vwet}] 0xAB36: мировой класс\\
\textbf{vwih} [\emph{vwiʰ}] 0x41B7: evitative case, indicates that the noun phrase is feared or avoided.\\
\textbf{vwik} [\emph{vwik}] 0x2836: уховертка\\
\textbf{vwin} [\emph{vwin}] 0x6836: ботаника\\
\textbf{vwis} [\emph{vwis}] 0x8836: выщипывать пинцетом\\
\textbf{vwit} [\emph{vwit}] 0xA836: жидкости\\
\textbf{vwok} [\emph{vwok}] 0x2C36: авокадо\\
\textbf{vwot} [\emph{vwot}] 0xAC36: вязкость\\
\textbf{vwut} [\emph{vwut}] 0xAA36: Кувейт\\
\textbf{vy6n} [\emph{vjən}] 0x6D16: венский\\
\textbf{vya2n} [\emph{vjä˩n}] 0x7916: внутренности\\
\textbf{vya7n} [\emph{vjä˥n}] 0x7116: руки меня вниз\\
\textbf{vya7t} [\emph{vjä˥t}] 0xB116: Авестийский\\
\textbf{vyac} [\emph{vjäʃ}] 0xE916: зодиак\\
\textbf{vyah} [\emph{vjäʰ}] 0x48B7: адвербиальный дело\\
\textbf{vyak} [\emph{vjäk}] 0x2916: адвербиальный дело\\
\textbf{vyam} [\emph{vjäm}] 0x0916: Вьетнам\\
\textbf{vyan} [\emph{vjän}] 0x6916: рекламодатель\\
\textbf{vyap} [\emph{vjäp}] 0x4916: восемьдесят восемь\\
\textbf{vyas} [\emph{vjäs}] 0x8916: верфь\\
\textbf{vyat} [\emph{vjät}] 0xA916: неповрежденный\\
\textbf{vyen} [\emph{vjen}] 0x6B16: ярмарочная площадь\\
\textbf{vyet} [\emph{vjet}] 0xAB16: испорченный\\
\textbf{vyi2n} [\emph{vji˩n}] 0x7816: головокружительный\\
\textbf{vyi7n} [\emph{vji˥n}] 0x7016: винил\\
\textbf{vyic} [\emph{vjiʃ}] 0xE816: бдение\\
\textbf{vyih} [\emph{vjiʰ}] 0x40B7: непрямой речь\\
\textbf{vyik} [\emph{vjik}] 0x2816: прокурор\\
\textbf{vyin} [\emph{vjin}] 0x6816: кинотеатр\\
\textbf{vyip} [\emph{vjip}] 0x4816: Ливия\\
\textbf{vyis} [\emph{vjis}] 0x8816: резец\\
\textbf{vyit} [\emph{vjit}] 0xA816: разносторонний\\
\textbf{vyon} [\emph{vjon}] 0x6C16: аэрация\\
\textbf{vyot} [\emph{vjot}] 0xAC16: шплинт\\
\textbf{vyun} [\emph{vjun}] 0x6A16: ЕС\\
\textbf{vyut} [\emph{vjut}] 0xAA16: университеты\\
\subsection{ w }
\textbf{w6} [\emph{wə}] 0x34BE: НМУ\\
\textbf{wa} [\emph{wä}] 0x24BE: а также\\
\textbf{we} [\emph{we}] 0x2CBE: vialis case, indicates motion through the noun phrase. used a standalone case. way-case. \\
\textbf{wi} [\emph{wi}] 0x20BE: универсальный квантификатором\\
\textbf{wi7} [\emph{wi˥}] 0x40BE: IC\\
\textbf{wo} [\emph{wo}] 0x30BE: слово\\
\textbf{wu} [\emph{wu}] 0x28BE: vocative case, the intended recipient of the sentence. for example `O' in `O god, hear my prayers.' discourse-context destination-case\\
\subsection{ x }
\textbf{xa} [\emph{xä}] 0x275E: речь контекст\\
\textbf{xe} [\emph{xe}] 0x2F5E: частота\\
\textbf{xe7} [\emph{xe˥}] 0x4F5E: Эм\\
\textbf{xi} [\emph{xi}] 0x235E: habitual aspect, something that happens habitualy.\\
\textbf{xl67n} [\emph{xlə˥n}] 0x757A: Таллин\\
\textbf{xl6k} [\emph{xlək}] 0x2D7A: галоген\\
\textbf{xl6n} [\emph{xlən}] 0x6D7A: нить\\
\textbf{xl6p} [\emph{xləp}] 0x4D7A: заячья губа\\
\textbf{xl6s} [\emph{xləs}] 0x8D7A: гиперинфляция\\
\textbf{xl6t} [\emph{xlət}] 0xAD7A: Таиланд\\
\textbf{xla2c} [\emph{xlä˩ʃ}] 0xF97A: хадж\\
\textbf{xla2k} [\emph{xlä˩k}] 0x397A: арлекин\\
\textbf{xla2n} [\emph{xlä˩n}] 0x797A: воющий\\
\textbf{xla2s} [\emph{xlä˩s}] 0x997A: Фары\\
\textbf{xla2t} [\emph{xlä˩t}] 0xB97A: глашатаи\\
\textbf{xla7c} [\emph{xlä˥ʃ}] 0xF17A: специалист по генеалогии\\
\textbf{xla7k} [\emph{xlä˥k}] 0x317A: Малахия\\
\textbf{xla7n} [\emph{xlä˥n}] 0x717A: гало\\
\textbf{xla7p} [\emph{xlä˥p}] 0x517A: хлопушка\\
\textbf{xla7s} [\emph{xlä˥s}] 0x917A: желто-зеленый\\
\textbf{xla7t} [\emph{xlä˥t}] 0xB17A: сильно\\
\textbf{xlac} [\emph{xläʃ}] 0xE97A: легальность\\
\textbf{xlaf} [\emph{xläf}] 0xC97A: отмель\\
\textbf{xlah} [\emph{xläʰ}] 0x4BD7: сатир\\
\textbf{xlak} [\emph{xläk}] 0x297A: автомотриса\\
\textbf{xlam} [\emph{xläm}] 0x097A: астроном\\
\textbf{xlan} [\emph{xlän}] 0x697A: глобализация\\
\textbf{xlap} [\emph{xläp}] 0x497A: пульпа\\
\textbf{xlas} [\emph{xläs}] 0x897A: парализованный\\
\textbf{xlat} [\emph{xlät}] 0xA97A: оценочный\\
\textbf{xle2n} [\emph{xle˩n}] 0x7B7A: гегельянец\\
\textbf{xle2t} [\emph{xle˩t}] 0xBB7A: швырнула\\
\textbf{xle7n} [\emph{xle˥n}] 0x737A: Хелен\\
\textbf{xle7t} [\emph{xle˥t}] 0xB37A: геркулесов\\
\textbf{xlec} [\emph{xleʃ}] 0xEB7A: Хэллоуин\\
\textbf{xlef} [\emph{xlef}] 0xCB7A: сеновал\\
\textbf{xleh} [\emph{xleʰ}] 0x5BD7: vehicle gender, 7th density life form, learning about timelines, akin to spirit of a ship/vehicle/spaceship\\
\textbf{xlek} [\emph{xlek}] 0x2B7A: ябеда\\
\textbf{xlem} [\emph{xlem}] 0x0B7A: бархатный червь\\
\textbf{xlen} [\emph{xlen}] 0x6B7A: словесный имя существительное\\
\textbf{xlep} [\emph{xlep}] 0x4B7A: обои\\
\textbf{xles} [\emph{xles}] 0x8B7A: Гималаи\\
\textbf{xlet} [\emph{xlet}] 0xAB7A: превосходный степень\\
\textbf{xli2f} [\emph{xli˩f}] 0xD87A: воздушные перевозки\\
\textbf{xli2k} [\emph{xli˩k}] 0x387A: штокроза розовая\\
\textbf{xli2n} [\emph{xli˩n}] 0x787A: оживлять\\
\textbf{xli2s} [\emph{xli˩s}] 0x987A: приличествовать\\
\textbf{xli2t} [\emph{xli˩t}] 0xB87A: фильтры\\
\textbf{xli7k} [\emph{xli˥k}] 0x307A: полиглот\\
\textbf{xli7m} [\emph{xli˥m}] 0x107A: е\\
\textbf{xli7n} [\emph{xli˥n}] 0x707A: выполнимый\\
\textbf{xli7s} [\emph{xli˥s}] 0x907A: нигилизм\\
\textbf{xli7t} [\emph{xli˥t}] 0xB07A: скрипач\\
\textbf{xlic} [\emph{xliʃ}] 0xE87A: Ахиллес\\
\textbf{xlif} [\emph{xlif}] 0xC87A: ставки\\
\textbf{xlih} [\emph{xliʰ}] 0x43D7: несчастность\\
\textbf{xlik} [\emph{xlik}] 0x287A: гиперссылка\\
\textbf{xlim} [\emph{xlim}] 0x087A: полигамия\\
\textbf{xlin} [\emph{xlin}] 0x687A: горячая линия\\
\textbf{xlip} [\emph{xlip}] 0x487A: волейбол\\
\textbf{xlis} [\emph{xlis}] 0x887A: Италия\\
\textbf{xlit} [\emph{xlit}] 0xA87A: инвалидная коляска\\
\textbf{xlo2k} [\emph{xlo˩k}] 0x3C7A: лущеный\\
\textbf{xlo2n} [\emph{xlo˩n}] 0x7C7A: прыщ\\
\textbf{xlo2t} [\emph{xlo˩t}] 0xBC7A: легированная\\
\textbf{xlo7n} [\emph{xlo˥n}] 0x747A: с воздушным охлаждением\\
\textbf{xlo7t} [\emph{xlo˥t}] 0xB47A: колесный мастер\\
\textbf{xloc} [\emph{xloʃ}] 0xEC7A: ностальгия\\
\textbf{xlof} [\emph{xlof}] 0xCC7A: оптовая\\
\textbf{xlok} [\emph{xlok}] 0x2C7A: дятел\\
\textbf{xlom} [\emph{xlom}] 0x0C7A: многочлен\\
\textbf{xlon} [\emph{xlon}] 0x6C7A: мобильный телефон\\
\textbf{xlop} [\emph{xlop}] 0x4C7A: гомеопатия\\
\textbf{xlos} [\emph{xlos}] 0x8C7A: зеленщик\\
\textbf{xlot} [\emph{xlot}] 0xAC7A: закат солнца\\
\textbf{xlu2n} [\emph{xlu˩n}] 0x7A7A: аллилуйя\\
\textbf{xlu2t} [\emph{xlu˩t}] 0xBA7A: гипохлорит\\
\textbf{xluc} [\emph{xluʃ}] 0xEA7A: ремень кнута\\
\textbf{xluf} [\emph{xluf}] 0xCA7A: пылесос\\
\textbf{xluk} [\emph{xluk}] 0x2A7A: сосулька\\
\textbf{xlum} [\emph{xlum}] 0x0A7A: возня\\
\textbf{xlun} [\emph{xlun}] 0x6A7A: нерастворимый\\
\textbf{xlup} [\emph{xlup}] 0x4A7A: градина\\
\textbf{xlus} [\emph{xlus}] 0x8A7A: Геркулес\\
\textbf{xlut} [\emph{xlut}] 0xAA7A: домашнее растение\\
\textbf{xo} [\emph{xo}] 0x335E: предварительно сегодняшний\\
\textbf{xr6k} [\emph{xrək}] 0x2DDA: ха-ха\\
\textbf{xr6m} [\emph{xrəm}] 0x0DDA: всевозможный\\
\textbf{xr6n} [\emph{xrən}] 0x6DDA: аллитерация\\
\textbf{xr6s} [\emph{xrəs}] 0x8DDA: иероглифы\\
\textbf{xr6t} [\emph{xrət}] 0xADDA: хаотичный\\
\textbf{xra2c} [\emph{xrä˩ʃ}] 0xF9DA: молочница\\
\textbf{xra2k} [\emph{xrä˩k}] 0x39DA: вес тары\\
\textbf{xra2n} [\emph{xrä˩n}] 0x79DA: императрица\\
\textbf{xra2s} [\emph{xrä˩s}] 0x99DA: торговец трикотажными изделиями\\
\textbf{xra2t} [\emph{xrä˩t}] 0xB9DA: разглагольствовать\\
\textbf{xra7c} [\emph{xrä˥ʃ}] 0xF1DA: шасси\\
\textbf{xra7k} [\emph{xrä˥k}] 0x31DA: терзать\\
\textbf{xra7m} [\emph{xrä˥m}] 0x11DA: Хиросима\\
\textbf{xra7n} [\emph{xrä˥n}] 0x71DA: кузнец\\
\textbf{xra7p} [\emph{xrä˥p}] 0x51DA: гарпия\\
\textbf{xra7s} [\emph{xrä˥s}] 0x91DA: слезящимися\\
\textbf{xra7t} [\emph{xrä˥t}] 0xB1DA: великодушие\\
\textbf{xrac} [\emph{xräʃ}] 0xE9DA: коралловый\\
\textbf{xraf} [\emph{xräf}] 0xC9DA: щавель\\
\textbf{xrah} [\emph{xräʰ}] 0x4ED7: жадность\\
\textbf{xrak} [\emph{xräk}] 0x29DA: как ни странно\\
\textbf{xram} [\emph{xräm}] 0x09DA: сафлоровое\\
\textbf{xran} [\emph{xrän}] 0x69DA: астронавт\\
\textbf{xrap} [\emph{xräp}] 0x49DA: носовой платок\\
\textbf{xras} [\emph{xräs}] 0x89DA: стирание\\
\textbf{xrat} [\emph{xrät}] 0xA9DA: шорты\\
\textbf{xre2n} [\emph{xre˩n}] 0x7BDA: мигрень\\
\textbf{xre2s} [\emph{xre˩s}] 0x9BDA: Ура\\
\textbf{xre2t} [\emph{xre˩t}] 0xBBDA: гидрат\\
\textbf{xre7n} [\emph{xre˥n}] 0x73DA: гиена\\
\textbf{xre7s} [\emph{xre˥s}] 0x93DA: о лице\\
\textbf{xre7t} [\emph{xre˥t}] 0xB3DA: гермафродит\\
\textbf{xrec} [\emph{xreʃ}] 0xEBDA: фамильная ценность\\
\textbf{xref} [\emph{xref}] 0xCBDA: мелодичный\\
\textbf{xreh} [\emph{xreʰ}] 0x5ED7: герц\\
\textbf{xrek} [\emph{xrek}] 0x2BDA: пациент вызывать\\
\textbf{xrem} [\emph{xrem}] 0x0BDA: фарватер\\
\textbf{xren} [\emph{xren}] 0x6BDA: наушники\\
\textbf{xrep} [\emph{xrep}] 0x4BDA: пажитник\\
\textbf{xres} [\emph{xres}] 0x8BDA: поперечный\\
\textbf{xret} [\emph{xret}] 0xABDA: терабайт\\
\textbf{xri2k} [\emph{xri˩k}] 0x38DA: Hara Kiri\\
\textbf{xri2n} [\emph{xri˩n}] 0x78DA: служить примером\\
\textbf{xri2s} [\emph{xri˩s}] 0x98DA: шарлатанство\\
\textbf{xri2t} [\emph{xri˩t}] 0xB8DA: засахаренный\\
\textbf{xri7k} [\emph{xri˥k}] 0x30DA: гидроксил\\
\textbf{xri7n} [\emph{xri˥n}] 0x70DA: трель\\
\textbf{xri7p} [\emph{xri˥p}] 0x50DA: пережить\\
\textbf{xri7s} [\emph{xri˥s}] 0x90DA: соглядатай\\
\textbf{xri7t} [\emph{xri˥t}] 0xB0DA: пародия\\
\textbf{xric} [\emph{xriʃ}] 0xE8DA: кувшин\\
\textbf{xrif} [\emph{xrif}] 0xC8DA: жирафа\\
\textbf{xrih} [\emph{xriʰ}] 0x46D7: застенчивость\\
\textbf{xrik} [\emph{xrik}] 0x28DA: пациент дело\\
\textbf{xrim} [\emph{xrim}] 0x08DA: иттрий\\
\textbf{xrin} [\emph{xrin}] 0x68DA: ингредиенты\\
\textbf{xrip} [\emph{xrip}] 0x48DA: сироп\\
\textbf{xris} [\emph{xris}] 0x88DA: пользователь интерфейс\\
\textbf{xrit} [\emph{xrit}] 0xA8DA: Главная страница\\
\textbf{xro2n} [\emph{xro˩n}] 0x7CDA: храп\\
\textbf{xro2t} [\emph{xro˩t}] 0xBCDA: ссадить\\
\textbf{xro7k} [\emph{xro˥k}] 0x34DA: неисторично\\
\textbf{xro7n} [\emph{xro˥n}] 0x74DA: герои\\
\textbf{xro7s} [\emph{xro˥s}] 0x94DA: гидрокси\\
\textbf{xro7t} [\emph{xro˥t}] 0xB4DA: белая жесть\\
\textbf{xroc} [\emph{xroʃ}] 0xECDA: гончар\\
\textbf{xrof} [\emph{xrof}] 0xCCDA: корабль на воздушной подушке\\
\textbf{xroh} [\emph{xroʰ}] 0x66D7: коллега по работе\\
\textbf{xrok} [\emph{xrok}] 0x2CDA: коллега по работе\\
\textbf{xrom} [\emph{xrom}] 0x0CDA: омоним\\
\textbf{xron} [\emph{xron}] 0x6CDA: Абстрактные Пол\\
\textbf{xrop} [\emph{xrop}] 0x4CDA: гармоника\\
\textbf{xros} [\emph{xros}] 0x8CDA: неоднородный\\
\textbf{xrot} [\emph{xrot}] 0xACDA: панталоны\\
\textbf{xru2n} [\emph{xru˩n}] 0x7ADA: Ааронова\\
\textbf{xru2t} [\emph{xru˩t}] 0xBADA: казнить на электрическом стуле\\
\textbf{xru7n} [\emph{xru˥n}] 0x72DA: гидразин\\
\textbf{xru7t} [\emph{xru˥t}] 0xB2DA: вкрутую\\
\textbf{xruc} [\emph{xruʃ}] 0xEADA: землеройка\\
\textbf{xruf} [\emph{xruf}] 0xCADA: люк\\
\textbf{xruh} [\emph{xruʰ}] 0x56D7: веселый\\
\textbf{xruk} [\emph{xruk}] 0x2ADA: линька\\
\textbf{xrum} [\emph{xrum}] 0x0ADA: Хартум\\
\textbf{xrun} [\emph{xrun}] 0x6ADA: бесплодный\\
\textbf{xrup} [\emph{xrup}] 0x4ADA: иврит\\
\textbf{xrus} [\emph{xrus}] 0x8ADA: гидроксид\\
\textbf{xrut} [\emph{xrut}] 0xAADA: детская площадка\\
\textbf{xu} [\emph{xu}] 0x2B5E: первичная объект\\
\textbf{xw6n} [\emph{xwən}] 0x6D3A: антиген\\
\textbf{xw6t} [\emph{xwət}] 0xAD3A: авантитул\\
\textbf{xwa2k} [\emph{xwä˩k}] 0x393A: ханука\\
\textbf{xwa2n} [\emph{xwä˩n}] 0x793A: боярышник\\
\textbf{xwa2t} [\emph{xwä˩t}] 0xB93A: выбрасывать\\
\textbf{xwa7m} [\emph{xwä˥m}] 0x113A: он мужчина\\
\textbf{xwa7n} [\emph{xwä˥n}] 0x713A: Гавана\\
\textbf{xwa7t} [\emph{xwä˥t}] 0xB13A: кредитор\\
\textbf{xwac} [\emph{xwäʃ}] 0xE93A: проходы\\
\textbf{xwaf} [\emph{xwäf}] 0xC93A: водоносный горизонт\\
\textbf{xwah} [\emph{xwäʰ}] 0x49D7: Вау\\
\textbf{xwak} [\emph{xwäk}] 0x293A: мочевой пузырь\\
\textbf{xwam} [\emph{xwäm}] 0x093A: ножовка\\
\textbf{xwan} [\emph{xwän}] 0x693A: размежевание\\
\textbf{xwap} [\emph{xwäp}] 0x493A: убегай\\
\textbf{xwas} [\emph{xwäs}] 0x893A: аист\\
\textbf{xwat} [\emph{xwät}] 0xA93A: мимо совершенный\\
\textbf{xwe7n} [\emph{xwe˥n}] 0x733A: тяжеловес\\
\textbf{xwek} [\emph{xwek}] 0x2B3A: фальшборт\\
\textbf{xwem} [\emph{xwem}] 0x0B3A: не член\\
\textbf{xwen} [\emph{xwen}] 0x6B3A: Нептун\\
\textbf{xwes} [\emph{xwes}] 0x8B3A: семьдесят четыре\\
\textbf{xwet} [\emph{xwet}] 0xAB3A: экватор\\
\textbf{xwi2n} [\emph{xwi˩n}] 0x783A: лечь в дрейф\\
\textbf{xwi2t} [\emph{xwi˩t}] 0xB83A: дыра в стене\\
\textbf{xwi7k} [\emph{xwi˥k}] 0x303A: халтура\\
\textbf{xwi7n} [\emph{xwi˥n}] 0x703A: дедовщина\\
\textbf{xwi7t} [\emph{xwi˥t}] 0xB03A: сбивчиво\\
\textbf{xwic} [\emph{xwiʃ}] 0xE83A: ручеек\\
\textbf{xwif} [\emph{xwif}] 0xC83A: встречный ветер\\
\textbf{xwih} [\emph{xwiʰ}] 0x41D7: patient case, the object of a transitive verb.\\
\textbf{xwik} [\emph{xwik}] 0x283A: гидравлический\\
\textbf{xwim} [\emph{xwim}] 0x083A: гомосексуализм\\
\textbf{xwin} [\emph{xwin}] 0x683A: минеральная Пол\\
\textbf{xwip} [\emph{xwip}] 0x483A: гипербола\\
\textbf{xwis} [\emph{xwis}] 0x883A: фриз\\
\textbf{xwit} [\emph{xwit}] 0xA83A: питаясь\\
\textbf{xwoc} [\emph{xwoʃ}] 0xEC3A: поджилки\\
\textbf{xwok} [\emph{xwok}] 0x2C3A: гео наука\\
\textbf{xwom} [\emph{xwom}] 0x0C3A: вся погода\\
\textbf{xwon} [\emph{xwon}] 0x6C3A: лебедь\\
\textbf{xwop} [\emph{xwop}] 0x4C3A: гемоглобин\\
\textbf{xwos} [\emph{xwos}] 0x8C3A: Осия\\
\textbf{xwot} [\emph{xwot}] 0xAC3A: так называемые\\
\textbf{xwuk} [\emph{xwuk}] 0x2A3A: галлюциногенный\\
\textbf{xwun} [\emph{xwun}] 0x6A3A: мутант\\
\textbf{xwup} [\emph{xwup}] 0x4A3A: бондарь\\
\textbf{xwus} [\emph{xwus}] 0x8A3A: обручи\\
\textbf{xwut} [\emph{xwut}] 0xAA3A: слабеть\\
\textbf{xy6n} [\emph{xjən}] 0x6D1A: Кения\\
\textbf{xy6t} [\emph{xjət}] 0xAD1A: пестрый\\
\textbf{xya2n} [\emph{xjä˩n}] 0x791A: солнцами\\
\textbf{xya2t} [\emph{xjä˩t}] 0xB91A: нижеподписавшиеся\\
\textbf{xya7k} [\emph{xjä˥k}] 0x311A: гыйба\\
\textbf{xya7n} [\emph{xjä˥n}] 0x711A: желтуха\\
\textbf{xya7s} [\emph{xjä˥s}] 0x911A: самостоятельно напористой\\
\textbf{xya7t} [\emph{xjä˥t}] 0xB11A: язык связан\\
\textbf{xyac} [\emph{xjäʃ}] 0xE91A: хрен\\
\textbf{xyaf} [\emph{xjäf}] 0xC91A: пыльник\\
\textbf{xyah} [\emph{xjäʰ}] 0x48D7: Мировой\\
\textbf{xyak} [\emph{xjäk}] 0x291A: Ямайка\\
\textbf{xyam} [\emph{xjäm}] 0x091A: таллий атом\\
\textbf{xyan} [\emph{xjän}] 0x691A: радиан\\
\textbf{xyap} [\emph{xjäp}] 0x491A: аффидавит\\
\textbf{xyas} [\emph{xjäs}] 0x891A: роли\\
\textbf{xyat} [\emph{xjät}] 0xA91A: антибиотик\\
\textbf{xye2k} [\emph{xje˩k}] 0x3B1A: герменевтика\\
\textbf{xye2n} [\emph{xje˩n}] 0x7B1A: красный глаз\\
\textbf{xye7n} [\emph{xje˥n}] 0x731A: семиугольник\\
\textbf{xye7t} [\emph{xje˥t}] 0xB31A: один ноги\\
\textbf{xyec} [\emph{xjeʃ}] 0xEB1A: гипотония\\
\textbf{xyek} [\emph{xjek}] 0x2B1A: эмоционально\\
\textbf{xyem} [\emph{xjem}] 0x0B1A: эритема\\
\textbf{xyen} [\emph{xjen}] 0x6B1A: взрастить\\
\textbf{xyep} [\emph{xjep}] 0x4B1A: гессенский\\
\textbf{xyes} [\emph{xjes}] 0x8B1A: Гаити\\
\textbf{xyet} [\emph{xjet}] 0xAB1A: дегенерат\\
\textbf{xyi2k} [\emph{xji˩k}] 0x381A: деревенщина\\
\textbf{xyi2n} [\emph{xji˩n}] 0x781A: поворотливость\\
\textbf{xyi2t} [\emph{xji˩t}] 0xB81A: недалекий\\
\textbf{xyi7n} [\emph{xji˥n}] 0x701A: гиацинт\\
\textbf{xyi7p} [\emph{xji˥p}] 0x501A: хиппи\\
\textbf{xyi7t} [\emph{xji˥t}] 0xB01A: высокотональный\\
\textbf{xyic} [\emph{xjiʃ}] 0xE81A: миссионер\\
\textbf{xyif} [\emph{xjif}] 0xC81A: океанография\\
\textbf{xyih} [\emph{xjiʰ}] 0x40D7: раскаяние\\
\textbf{xyik} [\emph{xjik}] 0x281A: дискета диск\\
\textbf{xyim} [\emph{xjim}] 0x081A: крапива\\
\textbf{xyin} [\emph{xjin}] 0x681A: неопытность\\
\textbf{xyip} [\emph{xjip}] 0x481A: непредсказуемый\\
\textbf{xyis} [\emph{xjis}] 0x881A: мимо пассивный причастие\\
\textbf{xyit} [\emph{xjit}] 0xA81A: ингибирование\\
\textbf{xyo7n} [\emph{xjo˥n}] 0x741A: бандит\\
\textbf{xyo7t} [\emph{xjo˥t}] 0xB41A: hoity toity\\
\textbf{xyoc} [\emph{xjoʃ}] 0xEC1A: приглушенный\\
\textbf{xyof} [\emph{xjof}] 0xCC1A: травоядное\\
\textbf{xyok} [\emph{xjok}] 0x2C1A: наручник\\
\textbf{xyom} [\emph{xjom}] 0x0C1A: бегемот\\
\textbf{xyon} [\emph{xjon}] 0x6C1A: пар\\
\textbf{xyop} [\emph{xjop}] 0x4C1A: колибри\\
\textbf{xyos} [\emph{xjos}] 0x8C1A: рекомбинация\\
\textbf{xyot} [\emph{xjot}] 0xAC1A: рептилия\\
\textbf{xyu7t} [\emph{xju˥t}] 0xB21A: здор`ово сделать\\
\textbf{xyuc} [\emph{xjuʃ}] 0xEA1A: удобство использования\\
\textbf{xyuk} [\emph{xjuk}] 0x2A1A: индеек\\
\textbf{xyum} [\emph{xjum}] 0x0A1A: Медовый месяц\\
\textbf{xyun} [\emph{xjun}] 0x6A1A: Сатурн\\
\textbf{xyup} [\emph{xjup}] 0x4A1A: подбоченясь\\
\textbf{xyus} [\emph{xjus}] 0x8A1A: тис\\
\textbf{xyut} [\emph{xjut}] 0xAA1A: в ближайщем будущем\\
\subsection{ y }
\textbf{y6} [\emph{jə}] 0x347E: художественный Пол\\
\textbf{ya} [\emph{jä}] 0x247E: Это\\
\textbf{ye} [\emph{je}] 0x2C7E: притяжательный падеж маркер\\
\textbf{yi} [\emph{ji}] 0x207E: dative case, indicates motion towards the noun phrase. destination-case. \\
\textbf{yi7} [\emph{ji˥}] 0x407E: язь\\
\textbf{yo} [\emph{jo}] 0x307E: союз\\
\textbf{yu} [\emph{ju}] 0x287E: instrumental case, indicates the method or tool used. person-context-way-case.\\
\subsection{ z }
\textbf{za} [\emph{zä}] 0x269E: inchoative aspect, for expressing the begining of a state or action\\
\textbf{ze} [\emph{ze}] 0x2E9E: Z\\
\textbf{zi} [\emph{zi}] 0x229E: процитировал\\
\textbf{zl6n} [\emph{zlən}] 0x6D74: зеландии\\
\textbf{zl6s} [\emph{zləs}] 0x8D74: кутерьмы на ближний свет\\
\textbf{zl6t} [\emph{zlət}] 0xAD74: фундук\\
\textbf{zla2n} [\emph{zlä˩n}] 0x7974: симулировать болезнь\\
\textbf{zla2t} [\emph{zlä˩t}] 0xB974: долг нагруженные\\
\textbf{zla7n} [\emph{zlä˥n}] 0x7174: позывной\\
\textbf{zla7t} [\emph{zlä˥t}] 0xB174: развязывать\\
\textbf{zlac} [\emph{zläʃ}] 0xE974: таракан\\
\textbf{zlah} [\emph{zläʰ}] 0x4BA7: верность\\
\textbf{zlak} [\emph{zläk}] 0x2974: нарушать\\
\textbf{zlam} [\emph{zläm}] 0x0974: детализировать\\
\textbf{zlan} [\emph{zlän}] 0x6974: зарабатывать\\
\textbf{zlap} [\emph{zläp}] 0x4974: зебра\\
\textbf{zlas} [\emph{zläs}] 0x8974: бег с препятствиями\\
\textbf{zlat} [\emph{zlät}] 0xA974: перчатка\\
\textbf{zlek} [\emph{zlek}] 0x2B74: экзоскелет\\
\textbf{zlen} [\emph{zlen}] 0x6B74: межпланетный\\
\textbf{zles} [\emph{zles}] 0x8B74: безногий\\
\textbf{zlet} [\emph{zlet}] 0xAB74: омлет\\
\textbf{zli2n} [\emph{zli˩n}] 0x7874: ксилола\\
\textbf{zli7k} [\emph{zli˥k}] 0x3074: интеллектуальный\\
\textbf{zli7n} [\emph{zli˥n}] 0x7074: вазелин\\
\textbf{zli7t} [\emph{zli˥t}] 0xB074: немного обратный порядок байт\\
\textbf{zlic} [\emph{zliʃ}] 0xE874: Чили\\
\textbf{zlif} [\emph{zlif}] 0xC874: ксилофон\\
\textbf{zlih} [\emph{zliʰ}] 0x43A7: irrealis настроение\\
\textbf{zlik} [\emph{zlik}] 0x2874: кремний\\
\textbf{zlim} [\emph{zlim}] 0x0874: ресница\\
\textbf{zlin} [\emph{zlin}] 0x6874: оливковый\\
\textbf{zlip} [\emph{zlip}] 0x4874: сейсмология\\
\textbf{zlis} [\emph{zlis}] 0x8874: Сицилия\\
\textbf{zlit} [\emph{zlit}] 0xA874: irrealis настроение\\
\textbf{zloc} [\emph{zloʃ}] 0xEC74: космология\\
\textbf{zlok} [\emph{zlok}] 0x2C74: безалкогольный\\
\textbf{zlon} [\emph{zlon}] 0x6C74: вечность\\
\textbf{zlop} [\emph{zlop}] 0x4C74: шлагбаумом\\
\textbf{zlos} [\emph{zlos}] 0x8C74: масленка\\
\textbf{zlot} [\emph{zlot}] 0xAC74: анод\\
\textbf{zluk} [\emph{zluk}] 0x2A74: зулусский\\
\textbf{zlum} [\emph{zlum}] 0x0A74: мусульманка\\
\textbf{zlun} [\emph{zlun}] 0x6A74: инсулин\\
\textbf{zlut} [\emph{zlut}] 0xAA74: модулировать\\
\textbf{zo} [\emph{zo}] 0x329E: zoic gender, 2nd density life form, learning to fullfill desires, akin to animals, basic computers, basic robots\\
\textbf{zr6n} [\emph{zrən}] 0x6DD4: Азербайджан\\
\textbf{zra2n} [\emph{zrä˩n}] 0x79D4: обострение\\
\textbf{zra2t} [\emph{zrä˩t}] 0xB9D4: разветвленная далеко\\
\textbf{zra7c} [\emph{zrä˥ʃ}] 0xF1D4: Захария\\
\textbf{zra7n} [\emph{zrä˥n}] 0x71D4: сад роз\\
\textbf{zra7t} [\emph{zrä˥t}] 0xB1D4: ангидрид\\
\textbf{zrac} [\emph{zräʃ}] 0xE9D4: обновление\\
\textbf{zraf} [\emph{zräf}] 0xC9D4: нулевое воздействие\\
\textbf{zrah} [\emph{zräʰ}] 0x4EA7: истекающий\\
\textbf{zrak} [\emph{zräk}] 0x29D4: акула\\
\textbf{zram} [\emph{zräm}] 0x09D4: карамель\\
\textbf{zran} [\emph{zrän}] 0x69D4: разделение\\
\textbf{zrap} [\emph{zräp}] 0x49D4: песчаный карьер\\
\textbf{zras} [\emph{zräs}] 0x89D4: Сирия\\
\textbf{zrat} [\emph{zrät}] 0xA9D4: кандидат\\
\textbf{zrec} [\emph{zreʃ}] 0xEBD4: взаимодействие\\
\textbf{zreh} [\emph{zreʰ}] 0x5EA7: exocentric case, indicates the noun phrase orbiting the other participants. Such as a moon or audience.\\
\textbf{zrek} [\emph{zrek}] 0x2BD4: экзоцентричный дело\\
\textbf{zren} [\emph{zren}] 0x6BD4: века\\
\textbf{zres} [\emph{zres}] 0x8BD4: вездесущность\\
\textbf{zret} [\emph{zret}] 0xABD4: подслушивать\\
\textbf{zri2n} [\emph{zri˩n}] 0x78D4: мелодичный перезвон\\
\textbf{zri7n} [\emph{zri˥n}] 0x70D4: петух малиновка\\
\textbf{zri7t} [\emph{zri˥t}] 0xB0D4: самостоятельно осуждающий\\
\textbf{zric} [\emph{zriʃ}] 0xE8D4: выравнивать\\
\textbf{zrif} [\emph{zrif}] 0xC8D4: хлорофилл\\
\textbf{zrih} [\emph{zriʰ}] 0x46A7: формальный регистр\\
\textbf{zrik} [\emph{zrik}] 0x28D4: поставить в известность\\
\textbf{zrim} [\emph{zrim}] 0x08D4: Суринам\\
\textbf{zrin} [\emph{zrin}] 0x68D4: фартук\\
\textbf{zrip} [\emph{zrip}] 0x48D4: Сибирь\\
\textbf{zris} [\emph{zris}] 0x88D4: Израиль\\
\textbf{zrit} [\emph{zrit}] 0xA8D4: формальный регистр\\
\textbf{zrok} [\emph{zrok}] 0x2CD4: Марокко\\
\textbf{zrom} [\emph{zrom}] 0x0CD4: изоморфизм\\
\textbf{zron} [\emph{zron}] 0x6CD4: нуль\\
\textbf{zrop} [\emph{zrop}] 0x4CD4: северо-восточный\\
\textbf{zros} [\emph{zros}] 0x8CD4: Азорские острова\\
\textbf{zrot} [\emph{zrot}] 0xACD4: кувыркаться\\
\textbf{zru7t} [\emph{zru˥t}] 0xB2D4: урду\\
\textbf{zruc} [\emph{zruʃ}] 0xEAD4: цюрих\\
\textbf{zruk} [\emph{zruk}] 0x2AD4: стержневой корень\\
\textbf{zrun} [\emph{zrun}] 0x6AD4: Сардиния\\
\textbf{zrup} [\emph{zrup}] 0x4AD4: узурпировать\\
\textbf{zrut} [\emph{zrut}] 0xAAD4: изумруд\\
\textbf{zwa7n} [\emph{zwä˥n}] 0x7134: вступить в повторный брак\\
\textbf{zwac} [\emph{zwäʃ}] 0xE934: недостаток\\
\textbf{zwah} [\emph{zwäʰ}] 0x49A7: амбиция\\
\textbf{zwak} [\emph{zwäk}] 0x2934: муляж\\
\textbf{zwan} [\emph{zwän}] 0x6934: происходящий\\
\textbf{zwas} [\emph{zwäs}] 0x8934: наркоз\\
\textbf{zwat} [\emph{zwät}] 0xA934: привычный\\
\textbf{zwen} [\emph{zwen}] 0x6B34: неуместность\\
\textbf{zwet} [\emph{zwet}] 0xAB34: второй полусредний вес\\
\textbf{zwik} [\emph{zwik}] 0x2834: горнолыжный спорт\\
\textbf{zwim} [\emph{zwim}] 0x0834: изомер\\
\textbf{zwin} [\emph{zwin}] 0x6834: резидент\\
\textbf{zwip} [\emph{zwip}] 0x4834: ксенофобские\\
\textbf{zwit} [\emph{zwit}] 0xA834: разговорчивый\\
\textbf{zwon} [\emph{zwon}] 0x6C34: сонар\\
\textbf{zwop} [\emph{zwop}] 0x4C34: изотоп\\
\textbf{zwum} [\emph{zwum}] 0x0A34: зум\\
\textbf{zya7t} [\emph{zjä˥t}] 0xB114: Зевадия\\
\textbf{zyac} [\emph{zjäʃ}] 0xE914: озабоченность\\
\textbf{zyaf} [\emph{zjäf}] 0xC914: затяжные прыжки с парашютом\\
\textbf{zyak} [\emph{zjäk}] 0x2914: синагога\\
\textbf{zyam} [\emph{zjäm}] 0x0914: спели\\
\textbf{zyan} [\emph{zjän}] 0x6914: весло\\
\textbf{zyap} [\emph{zjäp}] 0x4914: Замбия\\
\textbf{zyas} [\emph{zjäs}] 0x8914: Азия\\
\textbf{zyat} [\emph{zjät}] 0xA914: бритва\\
\textbf{zyen} [\emph{zjen}] 0x6B14: ксенон\\
\textbf{zyet} [\emph{zjet}] 0xAB14: Зеведей\\
\textbf{zyic} [\emph{zjiʃ}] 0xE814: крайний срок\\
\textbf{zyih} [\emph{zjiʰ}] 0x40A7: элита\\
\textbf{zyik} [\emph{zjik}] 0x2814: цинк\\
\textbf{zyin} [\emph{zjin}] 0x6814: хронический\\
\textbf{zyip} [\emph{zjip}] 0x4814: физиотерапия\\
\textbf{zyis} [\emph{zjis}] 0x8814: солнечные часы\\
\textbf{zyit} [\emph{zjit}] 0xA814: должностное лицо\\
\textbf{zyon} [\emph{zjon}] 0x6C14: озон\\
\textbf{zyot} [\emph{zjot}] 0xAC14: закупоренный\\
\textbf{zyut} [\emph{zjut}] 0xAA14: излишество\\

\section{пяш Numbers, Rhyming dictionary}
\subsection{ 0 }
0x0801 \textbf{myim} [\emph{mjim}]: иммунный\\
0x0802 \textbf{kyim} [\emph{kjim}]: купить\\
0x0804 \textbf{pyim} [\emph{pjim}]: платина\\
0x0806 \textbf{nyim} [\emph{njim}]: выдвижение\\
0x0808 \textbf{hmim} [\emph{ʰmim}]: средний\\
0x0809 \textbf{syim} [\emph{sjim}]: больной\\
0x080A \textbf{tyim} [\emph{tjim}]: Декабрь\\
0x080B \textbf{lyim} [\emph{ljim}]: элита\\
0x080C \textbf{fyim} [\emph{fjim}]: беженец\\
0x080D \textbf{cyim} [\emph{ʃjim}]: подозрительный\\
0x080E \textbf{ryim} [\emph{rjim}]: ритм\\
0x0810 \textbf{hkim} [\emph{ʰkim}]: секрет\\
0x0811 \textbf{byim} [\emph{bjim}]: моногамия\\
0x0812 \textbf{gyim} [\emph{gjim}]: бурый уголь\\
0x0813 \textbf{dyim} [\emph{djim}]: аудио\\
0x0815 \textbf{jyim} [\emph{ʒjim}]: я\\
0x0818 \textbf{hyim} [\emph{ʰjim}]: армия\\
0x0819 \textbf{qyim} [\emph{ŋjim}]: тимин\\
0x081A \textbf{xyim} [\emph{xjim}]: крапива\\
0x0820 \textbf{hpim} [\emph{ʰpim}]: императив настроение\\
0x0821 \textbf{mwim} [\emph{mwim}]: метонимия\\
0x0822 \textbf{kwim} [\emph{kwim}]: викторина\\
0x0824 \textbf{pwim} [\emph{pwim}]: неполная занятость\\
0x0826 \textbf{nwim} [\emph{nwim}]: сенсорный доказательный настроение\\
0x0828 \textbf{hwim} [\emph{ʰwim}]: цивилизация\\
0x0829 \textbf{swim} [\emph{swim}]: системы\\
0x082A \textbf{twim} [\emph{twim}]: дезидеративный настроение\\
0x082B \textbf{lwim} [\emph{lwim}]: алюминий\\
0x082D \textbf{cwim} [\emph{ʃwim}]: надежность\\
0x082E \textbf{rwim} [\emph{rwim}]: соответственно\\
0x0830 \textbf{hnim} [\emph{ʰnim}]: имя\\
0x0831 \textbf{bwim} [\emph{bwim}]: биохимия\\
0x0832 \textbf{gwim} [\emph{gwim}]: косоглазие\\
0x0833 \textbf{dwim} [\emph{dwim}]: дефиле\\
0x0834 \textbf{zwim} [\emph{zwim}]: изомер\\
0x083A \textbf{xwim} [\emph{xwim}]: гомосексуализм\\
0x0842 \textbf{ksim} [\emph{ksim}]: декларативный настроение\\
0x0844 \textbf{psim} [\emph{psim}]: приблизительный\\
0x0848 \textbf{hsim} [\emph{ʰsim}]: спать\\
0x084A \textbf{tsim} [\emph{tsim}]: время года\\
0x0850 \textbf{htim} [\emph{ʰtim}]: деонтические настроение\\
0x0851 \textbf{bzim} [\emph{bzim}]: дружелюбие\\
0x0853 \textbf{dzim} [\emph{dzim}]: бегущая дорожка\\
0x0858 \textbf{hlim} [\emph{ʰlim}]: климат\\
0x0860 \textbf{hfim} [\emph{ʰfim}]: фермер\\
0x0862 \textbf{klim} [\emph{klim}]: гримаса\\
0x0864 \textbf{plim} [\emph{plim}]: калий\\
0x0868 \textbf{hcim} [\emph{ʰʃim}]: особенно\\
0x0869 \textbf{slim} [\emph{slim}]: свисток\\
0x086A \textbf{tlim} [\emph{tlim}]: лечение\\
0x086C \textbf{flim} [\emph{flim}]: фильм\\
0x086D \textbf{clim} [\emph{ʃlim}]: специи\\
0x0870 \textbf{hrim} [\emph{ʰrim}]: христианство\\
0x0871 \textbf{blim} [\emph{blim}]: бериллий\\
0x0872 \textbf{glim} [\emph{glim}]: галлия\\
0x0873 \textbf{dlim} [\emph{dlim}]: идеология\\
0x0874 \textbf{zlim} [\emph{zlim}]: ресница\\
0x0875 \textbf{jlim} [\emph{ʒlim}]: неологизмы\\
0x0876 \textbf{vlim} [\emph{vlim}]: виртуальный Память\\
0x087A \textbf{xlim} [\emph{xlim}]: полигамия\\
0x0882 \textbf{kfim} [\emph{kfim}]: корневище\\
0x0884 \textbf{pfim} [\emph{pfim}]: период полураспада\\
0x0888 \textbf{hbim} [\emph{ʰbim}]: произвольный\\
0x088A \textbf{tfim} [\emph{tfim}]: литий\\
0x0890 \textbf{hgim} [\emph{ʰgim}]: дымовая труба\\
0x0891 \textbf{bvim} [\emph{bvim}]: альбумин\\
0x0893 \textbf{dvim} [\emph{dvim}]: динамический\\
0x0898 \textbf{hdim} [\emph{ʰdim}]: индуктивный настроение\\
0x08A0 \textbf{hzim} [\emph{ʰzim}]: висмут\\
0x08A2 \textbf{kcim} [\emph{kʃim}]: неограниченный\\
0x08A4 \textbf{pcim} [\emph{pʃim}]: выполнять\\
0x08A8 \textbf{hjim} [\emph{ʰʒim}]: запретительный настроение\\
0x08AA \textbf{tcim} [\emph{tʃim}]: одна и та же мероприятие маркер\\
0x08B0 \textbf{hvim} [\emph{ʰvim}]: время выполнения\\
0x08B1 \textbf{bjim} [\emph{bʒim}]: аболиционизм\\
0x08B2 \textbf{gjim} [\emph{gʒim}]: киска кошка\\
0x08B3 \textbf{djim} [\emph{dʒim}]: эндотермический\\
0x08C1 \textbf{mrim} [\emph{mrim}]: муравей\\
0x08C2 \textbf{krim} [\emph{krim}]: крем\\
0x08C4 \textbf{prim} [\emph{prim}]: имперский\\
0x08C6 \textbf{nrim} [\emph{nrim}]: стипендия\\
0x08C8 \textbf{hqim} [\emph{ʰŋim}]: цитата форма\\
0x08C9 \textbf{srim} [\emph{srim}]: подпись\\
0x08CA \textbf{trim} [\emph{trim}]: банка\\
0x08CC \textbf{frim} [\emph{frim}]: теорема\\
0x08CD \textbf{crim} [\emph{ʃrim}]: пирамида\\
0x08D0 \textbf{hxim} [\emph{ʰxim}]: торий атом\\
0x08D1 \textbf{brim} [\emph{brim}]: барий\\
0x08D2 \textbf{grim} [\emph{grim}]: логарифм\\
0x08D3 \textbf{drim} [\emph{drim}]: первокурсник\\
0x08D4 \textbf{zrim} [\emph{zrim}]: Суринам\\
0x08D5 \textbf{jrim} [\emph{ʒrim}]: Иеремия\\
0x08D6 \textbf{vrim} [\emph{vrim}]: висцеральный\\
0x08D9 \textbf{qrim} [\emph{ŋrim}]: рудимент\\
0x08DA \textbf{xrim} [\emph{xrim}]: иттрий\\
0x08E2 \textbf{kxim} [\emph{kxim}]: уклончивый\\
0x08E4 \textbf{pxim} [\emph{pxim}]: премьер-министр\\
0x08EA \textbf{txim} [\emph{txim}]: этидия\\
0x08F1 \textbf{bxim} [\emph{bxim}]: кипятить на медленном огне\\
0x08F2 \textbf{gxim} [\emph{gxim}]: малахит\\
0x08F3 \textbf{dxim} [\emph{dxim}]: домино\\
0x0901 \textbf{myam} [\emph{mjäm}]: любитель\\
0x0902 \textbf{kyam} [\emph{kjäm}]: лагерь\\
0x0904 \textbf{pyam} [\emph{pjäm}]: Храбрый\\
0x0906 \textbf{nyam} [\emph{njäm}]: слава\\
0x0908 \textbf{hmam} [\emph{ʰmäm}]: Мама\\
0x0909 \textbf{syam} [\emph{sjäm}]: безопасно\\
0x090A \textbf{tyam} [\emph{tjäm}]: поздно\\
0x090B \textbf{lyam} [\emph{ljäm}]: поток\\
0x090C \textbf{fyam} [\emph{fjäm}]: фраза маркер\\
0x090D \textbf{cyam} [\emph{ʃjäm}]: лето\\
0x090E \textbf{ryam} [\emph{rjäm}]: местоимение маркер\\
0x0910 \textbf{hkam} [\emph{ʰkäm}]: какие\\
0x0911 \textbf{byam} [\emph{bjäm}]: облучать\\
0x0912 \textbf{gyam} [\emph{gjäm}]: анус\\
0x0913 \textbf{dyam} [\emph{djäm}]: родий атом\\
0x0914 \textbf{zyam} [\emph{zjäm}]: спели\\
0x0915 \textbf{jyam} [\emph{ʒjäm}]: джайнизм\\
0x0916 \textbf{vyam} [\emph{vjäm}]: Вьетнам\\
0x0918 \textbf{hyam} [\emph{ʰjäm}]: ночь\\
0x0919 \textbf{qyam} [\emph{ŋjäm}]: анемометр\\
0x091A \textbf{xyam} [\emph{xjäm}]: таллий атом\\
0x0920 \textbf{hpam} [\emph{ʰpäm}]: первый\\
0x0921 \textbf{mwam} [\emph{mwäm}]: гимназия\\
0x0922 \textbf{kwam} [\emph{kwäm}]: запятая\\
0x0924 \textbf{pwam} [\emph{pwäm}]: шпионских программ\\
0x0926 \textbf{nwam} [\emph{nwäm}]: животное\\
0x0928 \textbf{hwam} [\emph{ʰwäm}]: нас\\
0x0929 \textbf{swam} [\emph{swäm}]: близнец\\
0x092A \textbf{twam} [\emph{twäm}]: Запрос\\
0x092B \textbf{lwam} [\emph{lwäm}]: алфавит\\
0x092C \textbf{fwam} [\emph{fwäm}]: межсезонье\\
0x092D \textbf{cwam} [\emph{ʃwäm}]: вена\\
0x092E \textbf{rwam} [\emph{rwäm}]: возмещать\\
0x0930 \textbf{hnam} [\emph{ʰnäm}]: а также или\\
0x0931 \textbf{bwam} [\emph{bwäm}]: избиратель\\
0x0932 \textbf{gwam} [\emph{gwäm}]: гамета\\
0x0933 \textbf{dwam} [\emph{dwäm}]: Да\\
0x0935 \textbf{jwam} [\emph{ʒwäm}]: смазывать\\
0x0939 \textbf{qwam} [\emph{ŋwäm}]: празднование новоселья\\
0x093A \textbf{xwam} [\emph{xwäm}]: ножовка\\
0x0942 \textbf{ksam} [\emph{ksäm}]: национальный\\
0x0944 \textbf{psam} [\emph{psäm}]: поза\\
0x0948 \textbf{hsam} [\emph{ʰsäm}]: муж\\
0x094A \textbf{tsam} [\emph{tsäm}]: непосредственно\\
0x0950 \textbf{htam} [\emph{ʰtäm}]: директива настроение\\
0x0951 \textbf{bzam} [\emph{bzäm}]: namby pamby\\
0x0952 \textbf{gzam} [\emph{gzäm}]: оргазм\\
0x0953 \textbf{dzam} [\emph{dzäm}]: астма\\
0x0958 \textbf{hlam} [\emph{ʰläm}]: приехать\\
0x0960 \textbf{hfam} [\emph{ʰfäm}]: форма\\
0x0962 \textbf{klam} [\emph{kläm}]: говоря\\
0x0964 \textbf{plam} [\emph{pläm}]: аварийная сигнализация\\
0x0968 \textbf{hcam} [\emph{ʰʃäm}]: час\\
0x0969 \textbf{slam} [\emph{släm}]: лосось\\
0x096A \textbf{tlam} [\emph{tläm}]: пример\\
0x096C \textbf{flam} [\emph{fläm}]: воск\\
0x096D \textbf{clam} [\emph{ʃläm}]: позор\\
0x0970 \textbf{hram} [\emph{ʰräm}]: романтический\\
0x0971 \textbf{blam} [\emph{bläm}]: бальный\\
0x0972 \textbf{glam} [\emph{gläm}]: аукцион\\
0x0973 \textbf{dlam} [\emph{dläm}]: отдать честь\\
0x0974 \textbf{zlam} [\emph{zläm}]: детализировать\\
0x0975 \textbf{jlam} [\emph{ʒläm}]: генералиссимус\\
0x0976 \textbf{vlam} [\emph{vläm}]: имя файла\\
0x097A \textbf{xlam} [\emph{xläm}]: астроном\\
0x0982 \textbf{kfam} [\emph{kfäm}]: квантовая\\
0x0984 \textbf{pfam} [\emph{pfäm}]: спамер\\
0x0988 \textbf{hbam} [\emph{ʰbäm}]: тетя\\
0x098A \textbf{tfam} [\emph{tfäm}]: прозрачный\\
0x0990 \textbf{hgam} [\emph{ʰgäm}]: ветчина\\
0x0991 \textbf{bvam} [\emph{bväm}]: нарды\\
0x0993 \textbf{dvam} [\emph{dväm}]: бессовестный\\
0x0998 \textbf{hdam} [\emph{ʰdäm}]: приключение\\
0x09A0 \textbf{hzam} [\emph{ʰzäm}]: банан\\
0x09A2 \textbf{kcam} [\emph{kʃäm}]: гордость\\
0x09A4 \textbf{pcam} [\emph{pʃäm}]: молиться\\
0x09A8 \textbf{hjam} [\emph{ʰʒäm}]: варенье\\
0x09AA \textbf{tcam} [\emph{tʃäm}]: враг\\
0x09B0 \textbf{hvam} [\emph{ʰväm}]: ванадий атом\\
0x09B1 \textbf{bjam} [\emph{bʒäm}]: притуплять\\
0x09B2 \textbf{gjam} [\emph{gʒäm}]: домогательство\\
0x09B3 \textbf{djam} [\emph{dʒäm}]: большой масштаб\\
0x09C1 \textbf{mram} [\emph{mräm}]: минеральная\\
0x09C2 \textbf{kram} [\emph{kräm}]: атом\\
0x09C4 \textbf{pram} [\emph{präm}]: альпинизм\\
0x09C6 \textbf{nram} [\emph{nräm}]: аморальный\\
0x09C8 \textbf{hqam} [\emph{ʰŋäm}]: разрабатывать\\
0x09C9 \textbf{sram} [\emph{sräm}]: буквально\\
0x09CA \textbf{tram} [\emph{träm}]: Рамка\\
0x09CC \textbf{fram} [\emph{främ}]: аптека\\
0x09CD \textbf{cram} [\emph{ʃräm}]: лестный\\
0x09D0 \textbf{hxam} [\emph{ʰxäm}]: миллиграмм\\
0x09D1 \textbf{bram} [\emph{bräm}]: бормотание\\
0x09D2 \textbf{gram} [\emph{gräm}]: рычание\\
0x09D3 \textbf{dram} [\emph{dräm}]: алмаз\\
0x09D4 \textbf{zram} [\emph{zräm}]: карамель\\
0x09D5 \textbf{jram} [\emph{ʒräm}]: макраме\\
0x09D6 \textbf{vram} [\emph{vräm}]: телятина\\
0x09D9 \textbf{qram} [\emph{ŋräm}]: реквизит\\
0x09DA \textbf{xram} [\emph{xräm}]: сафлоровое\\
0x09E2 \textbf{kxam} [\emph{kxäm}]: стержень\\
0x09E4 \textbf{pxam} [\emph{pxäm}]: уздечка\\
0x09EA \textbf{txam} [\emph{txäm}]: тантал\\
0x09F1 \textbf{bxam} [\emph{bxäm}]: бальзам\\
0x09F2 \textbf{gxam} [\emph{gxäm}]: Гватемала\\
0x09F3 \textbf{dxam} [\emph{dxäm}]: пищеварение\\
0x0A01 \textbf{myum} [\emph{mjum}]: мультимедиа\\
0x0A02 \textbf{kyum} [\emph{kjum}]: негашеная известь\\
0x0A04 \textbf{pyum} [\emph{pjum}]: окунать\\
0x0A06 \textbf{nyum} [\emph{njum}]: эму\\
0x0A08 \textbf{hmum} [\emph{ʰmum}]: юмор\\
0x0A09 \textbf{syum} [\emph{sjum}]: стронций атом\\
0x0A0A \textbf{tyum} [\emph{tjum}]: рутений атом\\
0x0A0B \textbf{lyum} [\emph{ljum}]: тулий атом\\
0x0A0C \textbf{fyum} [\emph{fjum}]: овсяница\\
0x0A0D \textbf{cyum} [\emph{ʃjum}]: дрожжи\\
0x0A0E \textbf{ryum} [\emph{rjum}]: уран\\
0x0A10 \textbf{hkum} [\emph{ʰkum}]: спекулятивный настроение\\
0x0A11 \textbf{byum} [\emph{bjum}]: рекламирующий\\
0x0A18 \textbf{hyum} [\emph{ʰjum}]: с помощью\\
0x0A1A \textbf{xyum} [\emph{xjum}]: Медовый месяц\\
0x0A20 \textbf{hpum} [\emph{ʰpum}]: поверхность\\
0x0A21 \textbf{mwum} [\emph{mwum}]: мычание\\
0x0A22 \textbf{kwum} [\emph{kwum}]: нематода\\
0x0A26 \textbf{nwum} [\emph{nwum}]: индуизм\\
0x0A28 \textbf{hwum} [\emph{ʰwum}]: слух\\
0x0A29 \textbf{swum} [\emph{swum}]: спиритизм\\
0x0A2B \textbf{lwum} [\emph{lwum}]: коллоквиализм\\
0x0A2E \textbf{rwum} [\emph{rwum}]: розовое дерево\\
0x0A30 \textbf{hnum} [\emph{ʰnum}]: разрешающий настроение\\
0x0A31 \textbf{bwum} [\emph{bwum}]: рубидий\\
0x0A34 \textbf{zwum} [\emph{zwum}]: зум\\
0x0A42 \textbf{ksum} [\emph{ksum}]: акционер\\
0x0A44 \textbf{psum} [\emph{psum}]: пемза\\
0x0A48 \textbf{hsum} [\emph{ʰsum}]: притяжательный падеж маркер\\
0x0A4A \textbf{tsum} [\emph{tsum}]: вольфрам\\
0x0A50 \textbf{htum} [\emph{ʰtum}]: наставительный настроение\\
0x0A58 \textbf{hlum} [\emph{ʰlum}]: верблюд\\
0x0A60 \textbf{hfum} [\emph{ʰfum}]: пена\\
0x0A62 \textbf{klum} [\emph{klum}]: совокупляться\\
0x0A64 \textbf{plum} [\emph{plum}]: пролог\\
0x0A68 \textbf{hcum} [\emph{ʰʃum}]: шум\\
0x0A69 \textbf{slum} [\emph{slum}]: ислам\\
0x0A6A \textbf{tlum} [\emph{tlum}]: теллур\\
0x0A6C \textbf{flum} [\emph{flum}]: пушок\\
0x0A6D \textbf{clum} [\emph{ʃlum}]: гелий\\
0x0A70 \textbf{hrum} [\emph{ʰrum}]: ремонт\\
0x0A74 \textbf{zlum} [\emph{zlum}]: мусульманка\\
0x0A7A \textbf{xlum} [\emph{xlum}]: возня\\
0x0A88 \textbf{hbum} [\emph{ʰbum}]: рубидий атом\\
0x0A90 \textbf{hgum} [\emph{ʰgum}]: жвачка\\
0x0A98 \textbf{hdum} [\emph{ʰdum}]: иудаизм\\
0x0AA0 \textbf{hzum} [\emph{ʰzum}]: цирконий\\
0x0AA2 \textbf{kcum} [\emph{kʃum}]: углубления\\
0x0AA4 \textbf{pcum} [\emph{pʃum}]: персиковое дерево\\
0x0AA8 \textbf{hjum} [\emph{ʰʒum}]: Иерусалим\\
0x0AAA \textbf{tcum} [\emph{tʃum}]: коснуться\\
0x0AB0 \textbf{hvum} [\emph{ʰvum}]: подводить\\
0x0AC1 \textbf{mrum} [\emph{mrum}]: компьютер программа\\
0x0AC2 \textbf{krum} [\emph{krum}]: скумбрия\\
0x0AC4 \textbf{prum} [\emph{prum}]: развертываться\\
0x0AC6 \textbf{nrum} [\emph{nrum}]: внутриутробный\\
0x0AC8 \textbf{hqum} [\emph{ʰŋum}]: имя пользователя\\
0x0AC9 \textbf{srum} [\emph{srum}]: высший\\
0x0ACA \textbf{trum} [\emph{trum}]: экстремист\\
0x0ACC \textbf{frum} [\emph{frum}]: фтор атом\\
0x0ACD \textbf{crum} [\emph{ʃrum}]: гормон\\
0x0AD0 \textbf{hxum} [\emph{ʰxum}]: бедренная кость\\
0x0AD1 \textbf{brum} [\emph{brum}]: подотряд\\
0x0AD2 \textbf{grum} [\emph{grum}]: референдум\\
0x0AD3 \textbf{drum} [\emph{drum}]: незагрязненный\\
0x0ADA \textbf{xrum} [\emph{xrum}]: Хартум\\
0x0AE2 \textbf{kxum} [\emph{kxum}]: воскликнул\\
0x0AE4 \textbf{pxum} [\emph{pxum}]: непрозрачный\\
0x0AEA \textbf{txum} [\emph{txum}]: локтевая кость\\
0x0AF1 \textbf{bxum} [\emph{bxum}]: буддизм\\
0x0B01 \textbf{myem} [\emph{mjem}]: миллиметр\\
0x0B02 \textbf{kyem} [\emph{kjem}]: кадмий\\
0x0B04 \textbf{pyem} [\emph{pjem}]: парфюмерия\\
0x0B06 \textbf{nyem} [\emph{njem}]: сантиметр\\
0x0B08 \textbf{hmem} [\emph{ʰmem}]: музей\\
0x0B09 \textbf{syem} [\emph{sjem}]: спам\\
0x0B0A \textbf{tyem} [\emph{tjem}]: титан\\
0x0B0B \textbf{lyem} [\emph{ljem}]: оперять\\
0x0B0D \textbf{cyem} [\emph{ʃjem}]: химус\\
0x0B0E \textbf{ryem} [\emph{rjem}]: радий\\
0x0B10 \textbf{hkem} [\emph{ʰkem}]: комета\\
0x0B11 \textbf{byem} [\emph{bjem}]: хлебный мякиш\\
0x0B12 \textbf{gyem} [\emph{gjem}]: магнитометр\\
0x0B13 \textbf{dyem} [\emph{djem}]: окрашенных в шерсти\\
0x0B18 \textbf{hyem} [\emph{ʰjem}]: вяз\\
0x0B1A \textbf{xyem} [\emph{xjem}]: эритема\\
0x0B20 \textbf{hpem} [\emph{ʰpem}]: Пальма\\
0x0B21 \textbf{mwem} [\emph{mwem}]: забери меня\\
0x0B24 \textbf{pwem} [\emph{pwem}]: придатки\\
0x0B26 \textbf{nwem} [\emph{nwem}]: ветреница\\
0x0B28 \textbf{hwem} [\emph{ʰwem}]: спектр\\
0x0B29 \textbf{swem} [\emph{swem}]: разбрызгиватель\\
0x0B2A \textbf{twem} [\emph{twem}]: трехмерный\\
0x0B2B \textbf{lwem} [\emph{lwem}]: палиндром\\
0x0B2E \textbf{rwem} [\emph{rwem}]: красное дерево\\
0x0B30 \textbf{hnem} [\emph{ʰnem}]: анимационный\\
0x0B3A \textbf{xwem} [\emph{xwem}]: не член\\
0x0B42 \textbf{ksem} [\emph{ksem}]: случайный доступ Память\\
0x0B44 \textbf{psem} [\emph{psem}]: плазма\\
0x0B48 \textbf{hsem} [\emph{ʰsem}]: саммит\\
0x0B4A \textbf{tsem} [\emph{tsem}]: атеизм\\
0x0B50 \textbf{htem} [\emph{ʰtem}]: постоянно\\
0x0B58 \textbf{hlem} [\emph{ʰlem}]: безрассудный\\
0x0B60 \textbf{hfem} [\emph{ʰfem}]: вставной морфема\\
0x0B62 \textbf{klem} [\emph{klem}]: флюид\\
0x0B64 \textbf{plem} [\emph{plem}]: полимер\\
0x0B68 \textbf{hcem} [\emph{ʰʃem}]: землетрясение\\
0x0B69 \textbf{slem} [\emph{slem}]: сублимировать\\
0x0B6A \textbf{tlem} [\emph{tlem}]: текущий актуальность маркер\\
0x0B6C \textbf{flem} [\emph{flem}]: флоэма\\
0x0B6D \textbf{clem} [\emph{ʃlem}]: селен\\
0x0B70 \textbf{hrem} [\emph{ʰrem}]: луч\\
0x0B71 \textbf{blem} [\emph{blem}]: Бельгия\\
0x0B72 \textbf{glem} [\emph{glem}]: страдающий манией величия\\
0x0B73 \textbf{dlem} [\emph{dlem}]: моделирование\\
0x0B7A \textbf{xlem} [\emph{xlem}]: бархатный червь\\
0x0B84 \textbf{pfem} [\emph{pfem}]: эпителий\\
0x0B88 \textbf{hbem} [\emph{ʰbem}]: витамин\\
0x0B8A \textbf{tfem} [\emph{tfem}]: сводная сестра\\
0x0B90 \textbf{hgem} [\emph{ʰgem}]: магний\\
0x0B98 \textbf{hdem} [\emph{ʰdem}]: рабочий стол\\
0x0BA0 \textbf{hzem} [\emph{ʰzem}]: экзотермический\\
0x0BA2 \textbf{kcem} [\emph{kʃem}]: ромашка\\
0x0BA4 \textbf{pcem} [\emph{pʃem}]: подкожный\\
0x0BA8 \textbf{hjem} [\emph{ʰʒem}]: Немецкий\\
0x0BAA \textbf{tcem} [\emph{tʃem}]: щипцы\\
0x0BB0 \textbf{hvem} [\emph{ʰvem}]: вермишель\\
0x0BB3 \textbf{djem} [\emph{dʒem}]: отек\\
0x0BC1 \textbf{mrem} [\emph{mrem}]: мэйнфреймов\\
0x0BC2 \textbf{krem} [\emph{krem}]: акров\\
0x0BC4 \textbf{prem} [\emph{prem}]: нефть\\
0x0BC6 \textbf{nrem} [\emph{nrem}]: нитрат\\
0x0BC8 \textbf{hqem} [\emph{ʰŋem}]: изотерма\\
0x0BC9 \textbf{srem} [\emph{srem}]: церий\\
0x0BCA \textbf{trem} [\emph{trem}]: личное снаряжение\\
0x0BCC \textbf{frem} [\emph{frem}]: восходящий поток теплого воздуха\\
0x0BCD \textbf{crem} [\emph{ʃrem}]: обдирать\\
0x0BD0 \textbf{hxem} [\emph{ʰxem}]: кропотливый\\
0x0BD1 \textbf{brem} [\emph{brem}]: лещ\\
0x0BD2 \textbf{grem} [\emph{grem}]: германий\\
0x0BD3 \textbf{drem} [\emph{drem}]: Дания\\
0x0BD5 \textbf{jrem} [\emph{ʒrem}]: герань\\
0x0BDA \textbf{xrem} [\emph{xrem}]: фарватер\\
0x0BE2 \textbf{kxem} [\emph{kxem}]: Клемент\\
0x0BE4 \textbf{pxem} [\emph{pxem}]: раскаленный добела\\
0x0BEA \textbf{txem} [\emph{txem}]: шашки\\
0x0BF1 \textbf{bxem} [\emph{bxem}]: Богемия\\
0x0BF2 \textbf{gxem} [\emph{gxem}]: мм хмм\\
0x0BF3 \textbf{dxem} [\emph{dxem}]: деспотизм\\
0x0C01 \textbf{myom} [\emph{mjom}]: молибден\\
0x0C02 \textbf{kyom} [\emph{kjom}]: кило\\
0x0C04 \textbf{pyom} [\emph{pjom}]: полоний атом\\
0x0C06 \textbf{nyom} [\emph{njom}]: неодим\\
0x0C08 \textbf{hmom} [\emph{ʰmom}]: горчица\\
0x0C09 \textbf{syom} [\emph{sjom}]: натрий\\
0x0C0A \textbf{tyom} [\emph{tjom}]: сурьма\\
0x0C0B \textbf{lyom} [\emph{ljom}]: Molly охранник\\
0x0C0C \textbf{fyom} [\emph{fjom}]: опиум\\
0x0C0D \textbf{cyom} [\emph{ʃjom}]: майоран\\
0x0C0E \textbf{ryom} [\emph{rjom}]: иридий\\
0x0C10 \textbf{hkom} [\emph{ʰkom}]: код\\
0x0C11 \textbf{byom} [\emph{bjom}]: биом\\
0x0C12 \textbf{gyom} [\emph{gjom}]: агрономия\\
0x0C13 \textbf{dyom} [\emph{djom}]: мажордом\\
0x0C18 \textbf{hyom} [\emph{ʰjom}]: йод\\
0x0C1A \textbf{xyom} [\emph{xjom}]: бегемот\\
0x0C20 \textbf{hpom} [\emph{ʰpom}]: Платформа\\
0x0C21 \textbf{mwom} [\emph{mwom}]: монизм\\
0x0C22 \textbf{kwom} [\emph{kwom}]: килограмм\\
0x0C24 \textbf{pwom} [\emph{pwom}]: задняя дверь\\
0x0C26 \textbf{nwom} [\emph{nwom}]: анатомия\\
0x0C28 \textbf{hwom} [\emph{ʰwom}]: бездомный\\
0x0C29 \textbf{swom} [\emph{swom}]: Сомали\\
0x0C2A \textbf{twom} [\emph{twom}]: внутренняя связь\\
0x0C2B \textbf{lwom} [\emph{lwom}]: глаукома\\
0x0C2E \textbf{rwom} [\emph{rwom}]: хромат\\
0x0C30 \textbf{hnom} [\emph{ʰnom}]: заговор\\
0x0C31 \textbf{bwom} [\emph{bwom}]: лоботомия\\
0x0C33 \textbf{dwom} [\emph{dwom}]: недипломатичный\\
0x0C3A \textbf{xwom} [\emph{xwom}]: вся погода\\
0x0C42 \textbf{ksom} [\emph{ksom}]: космос\\
0x0C44 \textbf{psom} [\emph{psom}]: читать только\\
0x0C48 \textbf{hsom} [\emph{ʰsom}]: поваренная соль\\
0x0C4A \textbf{tsom} [\emph{tsom}]: переходный\\
0x0C50 \textbf{htom} [\emph{ʰtom}]: экономика\\
0x0C58 \textbf{hlom} [\emph{ʰlom}]: лимон\\
0x0C60 \textbf{hfom} [\emph{ʰfom}]: фонема\\
0x0C62 \textbf{klom} [\emph{klom}]: километр\\
0x0C64 \textbf{plom} [\emph{plom}]: почтовый штемпель\\
0x0C68 \textbf{hcom} [\emph{ʰʃom}]: гармония\\
0x0C69 \textbf{slom} [\emph{slom}]: читать только Память\\
0x0C6A \textbf{tlom} [\emph{tlom}]: multal номер\\
0x0C6C \textbf{flom} [\emph{flom}]: Филимон\\
0x0C6D \textbf{clom} [\emph{ʃlom}]: цитоплазма\\
0x0C70 \textbf{hrom} [\emph{ʰrom}]: программа\\
0x0C71 \textbf{blom} [\emph{blom}]: тумба\\
0x0C72 \textbf{glom} [\emph{glom}]: Anglo Norman\\
0x0C73 \textbf{dlom} [\emph{dlom}]: дипломатия\\
0x0C7A \textbf{xlom} [\emph{xlom}]: многочлен\\
0x0C82 \textbf{kfom} [\emph{kfom}]: номера\\
0x0C88 \textbf{hbom} [\emph{ʰbom}]: рвота\\
0x0C8A \textbf{tfom} [\emph{tfom}]: тригонометрия\\
0x0C90 \textbf{hgom} [\emph{ʰgom}]: гамма\\
0x0C98 \textbf{hdom} [\emph{ʰdom}]: кошмар\\
0x0CA0 \textbf{hzom} [\emph{ʰzom}]: кочевник\\
0x0CA2 \textbf{kcom} [\emph{kʃom}]: экономист\\
0x0CA4 \textbf{pcom} [\emph{pʃom}]: операционная система\\
0x0CA8 \textbf{hjom} [\emph{ʰʒom}]: хромосома\\
0x0CAA \textbf{tcom} [\emph{tʃom}]: удушение\\
0x0CB0 \textbf{hvom} [\emph{ʰvom}]: спидометр\\
0x0CC1 \textbf{mrom} [\emph{mrom}]: морфема\\
0x0CC2 \textbf{krom} [\emph{krom}]: хром\\
0x0CC4 \textbf{prom} [\emph{prom}]: программист\\
0x0CC6 \textbf{nrom} [\emph{nrom}]: актер роль\\
0x0CC8 \textbf{hqom} [\emph{ʰŋom}]: ром\\
0x0CC9 \textbf{srom} [\emph{srom}]: выставочный зал\\
0x0CCA \textbf{trom} [\emph{trom}]: тележка\\
0x0CCC \textbf{from} [\emph{from}]: анахронизм\\
0x0CCD \textbf{crom} [\emph{ʃrom}]: сотрудничать\\
0x0CD0 \textbf{hxom} [\emph{ʰxom}]: утиль почта\\
0x0CD1 \textbf{brom} [\emph{brom}]: бром\\
0x0CD2 \textbf{grom} [\emph{grom}]: роговица\\
0x0CD3 \textbf{drom} [\emph{drom}]: диморфизм\\
0x0CD4 \textbf{zrom} [\emph{zrom}]: изоморфизм\\
0x0CD5 \textbf{jrom} [\emph{ʒrom}]: теплым сердцем\\
0x0CD6 \textbf{vrom} [\emph{vrom}]: аорта\\
0x0CDA \textbf{xrom} [\emph{xrom}]: омоним\\
0x0CE2 \textbf{kxom} [\emph{kxom}]: согласный\\
0x0CE4 \textbf{pxom} [\emph{pxom}]: propositive настроение\\
0x0CEA \textbf{txom} [\emph{txom}]: Оман\\
0x0CF1 \textbf{bxom} [\emph{bxom}]: окра\\
0x0CF2 \textbf{gxom} [\emph{gxom}]: экзогамия\\
0x0CF3 \textbf{dxom} [\emph{dxom}]: одометр\\
0x0D01 \textbf{my6m} [\emph{mjəm}]: Мьянма\\
0x0D06 \textbf{ny6m} [\emph{njəm}]: анимизм\\
0x0D08 \textbf{hm6m} [\emph{ʰməm}]: в середине лета\\
0x0D09 \textbf{sy6m} [\emph{sjəm}]: сатанизм\\
0x0D0B \textbf{ly6m} [\emph{ljəm}]: НМУ\\
0x0D10 \textbf{hk6m} [\emph{ʰkəm}]: мошонка\\
0x0D18 \textbf{hy6m} [\emph{ʰjəm}]: Йемен\\
0x0D20 \textbf{hp6m} [\emph{ʰpəm}]: празеодимий\\
0x0D28 \textbf{hw6m} [\emph{ʰwəm}]: Вэстон\\
0x0D30 \textbf{hn6m} [\emph{ʰnəm}]: Нериум\\
0x0D44 \textbf{ps6m} [\emph{psəm}]: хурма\\
0x0D48 \textbf{hs6m} [\emph{ʰsəm}]: цезий\\
0x0D4A \textbf{ts6m} [\emph{tsəm}]: набивка чучел\\
0x0D50 \textbf{ht6m} [\emph{ʰtəm}]: тербий атом\\
0x0D58 \textbf{hl6m} [\emph{ʰləm}]: euler's номер\\
0x0D60 \textbf{hf6m} [\emph{ʰfəm}]: однодневка\\
0x0D62 \textbf{kl6m} [\emph{kləm}]: скаляр\\
0x0D64 \textbf{pl6m} [\emph{pləm}]: привратник желудка\\
0x0D68 \textbf{hc6m} [\emph{ʰʃəm}]: бальзамировать\\
0x0D69 \textbf{sl6m} [\emph{sləm}]: манная крупа\\
0x0D6A \textbf{tl6m} [\emph{tləm}]: Твидл ди\\
0x0D6C \textbf{fl6m} [\emph{fləm}]: фатализм\\
0x0D70 \textbf{hr6m} [\emph{ʰrəm}]: рений\\
0x0D84 \textbf{pf6m} [\emph{pfəm}]: в изобилии\\
0x0D88 \textbf{hb6m} [\emph{ʰbəm}]: брахман\\
0x0D90 \textbf{hg6m} [\emph{ʰgəm}]: микроорганизмы\\
0x0D98 \textbf{hd6m} [\emph{ʰdəm}]: дихотомический\\
0x0DAA \textbf{tc6m} [\emph{tʃəm}]: трапеция\\
0x0DC1 \textbf{mr6m} [\emph{mrəm}]: мормон\\
0x0DC2 \textbf{kr6m} [\emph{krəm}]: презерватив\\
0x0DC4 \textbf{pr6m} [\emph{prəm}]: пермский\\
0x0DC6 \textbf{nr6m} [\emph{nrəm}]: патентованное средство\\
0x0DC8 \textbf{hq6m} [\emph{ʰŋəm}]: болен omened\\
0x0DC9 \textbf{sr6m} [\emph{srəm}]: стронций\\
0x0DCA \textbf{tr6m} [\emph{trəm}]: термометр\\
0x0DD0 \textbf{hx6m} [\emph{ʰxəm}]: почтовый голубь\\
0x0DD3 \textbf{dr6m} [\emph{drəm}]: дезинформация\\
0x0DDA \textbf{xr6m} [\emph{xrəm}]: всевозможный\\
0x0DEA \textbf{tx6m} [\emph{txəm}]: катамаран\\
0x0DF1 \textbf{bx6m} [\emph{bxəm}]: багамы\\
\subsection{ 1 }
0x1001 \textbf{myi7m} [\emph{mji˥m}]: эмир\\
0x1008 \textbf{hmi7m} [\emph{ʰmi˥m}]: жасмин\\
0x100A \textbf{tyi7m} [\emph{tji˥m}]: автономный\\
0x100B \textbf{lyi7m} [\emph{lji˥m}]: подвздошная кость\\
0x100E \textbf{ryi7m} [\emph{rji˥m}]: рацемический\\
0x1010 \textbf{hki7m} [\emph{ʰki˥m}]: несмываемый\\
0x1018 \textbf{hyi7m} [\emph{ʰji˥m}]: сионизм\\
0x1020 \textbf{hpi7m} [\emph{ʰpi˥m}]: примат\\
0x1021 \textbf{mwi7m} [\emph{mwi˥m}]: мим\\
0x1028 \textbf{hwi7m} [\emph{ʰwi˥m}]: Уильям\\
0x1030 \textbf{hni7m} [\emph{ʰni˥m}]: Бенджамен\\
0x1042 \textbf{ksi7m} [\emph{ksi˥m}]: кубизм\\
0x1048 \textbf{hsi7m} [\emph{ʰsi˥m}]: Ишмаэль\\
0x104A \textbf{tsi7m} [\emph{tsi˥m}]: атавизм\\
0x1050 \textbf{hti7m} [\emph{ʰti˥m}]: чопорный\\
0x1058 \textbf{hli7m} [\emph{ʰli˥m}]: недоразумение\\
0x1060 \textbf{hfi7m} [\emph{ʰfi˥m}]: эмфизема\\
0x1068 \textbf{hci7m} [\emph{ʰʃi˥m}]: Джимми\\
0x1069 \textbf{sli7m} [\emph{sli˥m}]: эпистемология\\
0x106A \textbf{tli7m} [\emph{tli˥m}]: патернализм\\
0x1070 \textbf{hri7m} [\emph{ʰri˥m}]: предчувствие\\
0x107A \textbf{xli7m} [\emph{xli˥m}]: е\\
0x1088 \textbf{hbi7m} [\emph{ʰbi˥m}]: подружка невесты\\
0x1090 \textbf{hgi7m} [\emph{ʰgi˥m}]: эргономика\\
0x1098 \textbf{hdi7m} [\emph{ʰdi˥m}]: идиомы\\
0x10A2 \textbf{kci7m} [\emph{kʃi˥m}]: Кашмир\\
0x10AA \textbf{tci7m} [\emph{tʃi˥m}]: тимус\\
0x10C1 \textbf{mri7m} [\emph{mri˥m}]: микрометр\\
0x10C4 \textbf{pri7m} [\emph{pri˥m}]: ипподром\\
0x10C6 \textbf{nri7m} [\emph{nri˥m}]: очный\\
0x10C9 \textbf{sri7m} [\emph{sri˥m}]: синкретизм\\
0x10CA \textbf{tri7m} [\emph{tri˥m}]: трирема\\
0x10CC \textbf{fri7m} [\emph{fri˥m}]: афоризмы\\
0x10D0 \textbf{hxi7m} [\emph{ʰxi˥m}]: МВФ\\
0x10D1 \textbf{bri7m} [\emph{bri˥m}]: шрифт Брайля\\
0x10D3 \textbf{dri7m} [\emph{dri˥m}]: двухцветный\\
0x1101 \textbf{mya7m} [\emph{mjä˥m}]: человек людоед\\
0x1102 \textbf{kya7m} [\emph{kjä˥m}]: переименовать\\
0x1104 \textbf{pya7m} [\emph{pjä˥m}]: панамский\\
0x1106 \textbf{nya7m} [\emph{njä˥m}]: фундаментализм\\
0x1108 \textbf{hma7m} [\emph{ʰmä˥m}]: кобыла\\
0x1109 \textbf{sya7m} [\emph{sjä˥m}]: санаторий\\
0x110A \textbf{tya7m} [\emph{tjä˥m}]: перевалка\\
0x110B \textbf{lya7m} [\emph{ljä˥m}]: стручковый перец\\
0x110E \textbf{rya7m} [\emph{rjä˥m}]: арамейский\\
0x1110 \textbf{hka7m} [\emph{ʰkä˥m}]: ласка\\
0x1113 \textbf{dya7m} [\emph{djä˥m}]: диатомовый\\
0x1118 \textbf{hya7m} [\emph{ʰjä˥m}]: фрамбезия\\
0x1120 \textbf{hpa7m} [\emph{ʰpä˥m}]: паломник\\
0x1128 \textbf{hwa7m} [\emph{ʰwä˥m}]: сердцеед\\
0x1129 \textbf{swa7m} [\emph{swä˥m}]: астигматизм\\
0x1130 \textbf{hna7m} [\emph{ʰnä˥m}]: натальный\\
0x113A \textbf{xwa7m} [\emph{xwä˥m}]: он мужчина\\
0x1142 \textbf{ksa7m} [\emph{ksä˥m}]: Ассамский\\
0x1148 \textbf{hsa7m} [\emph{ʰsä˥m}]: праздно\\
0x114A \textbf{tsa7m} [\emph{tsä˥m}]: человек в руках\\
0x1150 \textbf{hta7m} [\emph{ʰtä˥m}]: коробейник\\
0x1158 \textbf{hla7m} [\emph{ʰlä˥m}]: алебастр\\
0x1160 \textbf{hfa7m} [\emph{ʰfä˥m}]: фашизм\\
0x1162 \textbf{kla7m} [\emph{klä˥m}]: катаболизм\\
0x1168 \textbf{hca7m} [\emph{ʰʃä˥m}]: заика\\
0x1169 \textbf{sla7m} [\emph{slä˥m}]: всмятку\\
0x116A \textbf{tla7m} [\emph{tlä˥m}]: звон колоколов\\
0x1170 \textbf{hra7m} [\emph{ʰrä˥m}]: старое время\\
0x1173 \textbf{dla7m} [\emph{dlä˥m}]: дифениламин\\
0x1188 \textbf{hba7m} [\emph{ʰbä˥m}]: незадачливый\\
0x118A \textbf{tfa7m} [\emph{tfä˥m}]: трансформеры\\
0x1190 \textbf{hga7m} [\emph{ʰgä˥m}]: магнезия\\
0x1198 \textbf{hda7m} [\emph{ʰdä˥m}]: содомия\\
0x11A2 \textbf{kca7m} [\emph{kʃä˥m}]: крематорий\\
0x11AA \textbf{tca7m} [\emph{tʃä˥m}]: меткий стрелок\\
0x11C1 \textbf{mra7m} [\emph{mrä˥m}]: мелодрама\\
0x11C2 \textbf{kra7m} [\emph{krä˥m}]: лакрица\\
0x11C4 \textbf{pra7m} [\emph{prä˥m}]: перикард\\
0x11C6 \textbf{nra7m} [\emph{nrä˥m}]: панорамный\\
0x11C8 \textbf{hqa7m} [\emph{ʰŋä˥m}]: прабабушка\\
0x11C9 \textbf{sra7m} [\emph{srä˥m}]: Самария\\
0x11CA \textbf{tra7m} [\emph{trä˥m}]: каракуль\\
0x11CD \textbf{cra7m} [\emph{ʃrä˥m}]: креветка\\
0x11D0 \textbf{hxa7m} [\emph{ʰxä˥m}]: рукопожатие\\
0x11D1 \textbf{bra7m} [\emph{brä˥m}]: Авраам\\
0x11D2 \textbf{gra7m} [\emph{grä˥m}]: Фогарма\\
0x11D3 \textbf{dra7m} [\emph{drä˥m}]: темная комната\\
0x11DA \textbf{xra7m} [\emph{xrä˥m}]: Хиросима\\
0x11E2 \textbf{kxa7m} [\emph{kxä˥m}]: клептоман\\
0x11EA \textbf{txa7m} [\emph{txä˥m}]: зенитная\\
0x11F1 \textbf{bxa7m} [\emph{bxä˥m}]: белая лошадь\\
0x1208 \textbf{hmu7m} [\emph{ʰmu˥m}]: Магомет\\
0x1248 \textbf{hsu7m} [\emph{ʰsu˥m}]: грудина\\
0x1250 \textbf{htu7m} [\emph{ʰtu˥m}]: перламутровый\\
0x1258 \textbf{hlu7m} [\emph{ʰlu˥m}]: оксид алюминия\\
0x1270 \textbf{hru7m} [\emph{ʰru˥m}]: Каникулы в гареме\\
0x1290 \textbf{hgu7m} [\emph{ʰgu˥m}]: гм\\
0x12D0 \textbf{hxu7m} [\emph{ʰxu˥m}]: гумусовый\\
0x1308 \textbf{hme7m} [\emph{ʰme˥m}]: середине мая\\
0x1310 \textbf{hke7m} [\emph{ʰke˥m}]: кхмерская\\
0x1318 \textbf{hye7m} [\emph{ʰje˥m}]: Эм\\
0x1320 \textbf{hpe7m} [\emph{ʰpe˥m}]: грушевый сидр\\
0x1330 \textbf{hne7m} [\emph{ʰne˥m}]: Неемия\\
0x1348 \textbf{hse7m} [\emph{ʰse˥m}]: эскимос\\
0x134A \textbf{tse7m} [\emph{tse˥m}]: время чая\\
0x1350 \textbf{hte7m} [\emph{ʰte˥m}]: трамвайная линия\\
0x1358 \textbf{hle7m} [\emph{ʰle˥m}]: дождемер\\
0x1364 \textbf{ple7m} [\emph{ple˥m}]: Пелл Мелл\\
0x1370 \textbf{hre7m} [\emph{ʰre˥m}]: бремен\\
0x1390 \textbf{hge7m} [\emph{ʰge˥m}]: Банкомат\\
0x1398 \textbf{hde7m} [\emph{ʰde˥m}]: демон\\
0x13AA \textbf{tce7m} [\emph{tʃe˥m}]: трахеотомия\\
0x13D0 \textbf{hxe7m} [\emph{ʰxe˥m}]: девственная плева\\
0x1408 \textbf{hmo7m} [\emph{ʰmo˥m}]: Месопотамия\\
0x1420 \textbf{hpo7m} [\emph{ʰpo˥m}]: OMA\\
0x1430 \textbf{hno7m} [\emph{ʰno˥m}]: нерентабельный\\
0x1448 \textbf{hso7m} [\emph{ʰso˥m}]: стоицизм\\
0x144A \textbf{tso7m} [\emph{tso˥m}]: газонокосилка\\
0x1450 \textbf{hto7m} [\emph{ʰto˥m}]: Стокгольм\\
0x1458 \textbf{hlo7m} [\emph{ʰlo˥m}]: терн\\
0x1469 \textbf{slo7m} [\emph{slo˥m}]: миелома\\
0x1470 \textbf{hro7m} [\emph{ʰro˥m}]: золоченая бронза\\
0x1488 \textbf{hbo7m} [\emph{ʰbo˥m}]: бергамот\\
0x1490 \textbf{hgo7m} [\emph{ʰgo˥m}]: Graeco романский\\
0x1498 \textbf{hdo7m} [\emph{ʰdo˥m}]: ареометр\\
0x14C2 \textbf{kro7m} [\emph{kro˥m}]: кворум\\
0x14C6 \textbf{nro7m} [\emph{nro˥m}]: монохромный\\
0x14CA \textbf{tro7m} [\emph{tro˥m}]: торий\\
0x14D0 \textbf{hxo7m} [\emph{ʰxo˥m}]: Пожиратель огня\\
0x14EA \textbf{txo7m} [\emph{txo˥m}]: Tom Tom\\
0x1550 \textbf{ht67m} [\emph{ʰtə˥m}]: Траляля и Труляля\\
0x1808 \textbf{hmi2m} [\emph{ʰmi˩m}]: маори\\
0x180A \textbf{tyi2m} [\emph{tji˩m}]: метка времени\\
0x1810 \textbf{hki2m} [\emph{ʰki˩m}]: запор\\
0x1820 \textbf{hpi2m} [\emph{ʰpi˩m}]: призма\\
0x1828 \textbf{hwi2m} [\emph{ʰwi˩m}]: опасный для жизни\\
0x1830 \textbf{hni2m} [\emph{ʰni˩m}]: однопалатный\\
0x1848 \textbf{hsi2m} [\emph{ʰsi˩m}]: проницательность\\
0x184A \textbf{tsi2m} [\emph{tsi˩m}]: таксометр\\
0x1850 \textbf{hti2m} [\emph{ʰti˩m}]: ультиматум\\
0x1858 \textbf{hli2m} [\emph{ʰli˩m}]: алхимия\\
0x1860 \textbf{hfi2m} [\emph{ʰfi˩m}]: амфетамин\\
0x1869 \textbf{sli2m} [\emph{sli˩m}]: солецизм\\
0x186A \textbf{tli2m} [\emph{tli˩m}]: антиномия\\
0x1870 \textbf{hri2m} [\emph{ʰri˩m}]: общинники\\
0x1888 \textbf{hbi2m} [\emph{ʰbi˩m}]: балки\\
0x1898 \textbf{hdi2m} [\emph{ʰdi˩m}]: деизм\\
0x18A2 \textbf{kci2m} [\emph{kʃi˩m}]: маслобойня\\
0x18A4 \textbf{pci2m} [\emph{pʃi˩m}]: паренхима\\
0x18AA \textbf{tci2m} [\emph{tʃi˩m}]: лимит времени\\
0x18C1 \textbf{mri2m} [\emph{mri˩m}]: матереубийство\\
0x18C2 \textbf{kri2m} [\emph{kri˩m}]: криптограмма\\
0x18C6 \textbf{nri2m} [\emph{nri˩m}]: анаграмма\\
0x18C9 \textbf{sri2m} [\emph{sri˩m}]: садизм\\
0x18CA \textbf{tri2m} [\emph{tri˩m}]: с электроприводом\\
0x18D0 \textbf{hxi2m} [\emph{ʰxi˩m}]: Крик осла\\
0x18D3 \textbf{dri2m} [\emph{dri˩m}]: гидротермальный\\
0x1901 \textbf{mya2m} [\emph{mjä˩m}]: плоть муха\\
0x1908 \textbf{hma2m} [\emph{ʰmä˩m}]: оцепенение\\
0x190A \textbf{tya2m} [\emph{tjä˩m}]: успокойся\\
0x190B \textbf{lya2m} [\emph{ljä˩m}]: малаялам\\
0x1910 \textbf{hka2m} [\emph{ʰkä˩m}]: сперма\\
0x1918 \textbf{hya2m} [\emph{ʰjä˩m}]: аймара\\
0x1920 \textbf{hpa2m} [\emph{ʰpä˩m}]: односторонними\\
0x1929 \textbf{swa2m} [\emph{swä˩m}]: пространство-время\\
0x192A \textbf{twa2m} [\emph{twä˩m}]: креационизм\\
0x1930 \textbf{hna2m} [\emph{ʰnä˩m}]: пантеизм\\
0x1942 \textbf{ksa2m} [\emph{ksä˩m}]: гекзаметр\\
0x1948 \textbf{hsa2m} [\emph{ʰsä˩m}]: цианид\\
0x194A \textbf{tsa2m} [\emph{tsä˩m}]: стеньга\\
0x1950 \textbf{hta2m} [\emph{ʰtä˩m}]: безвременно\\
0x1958 \textbf{hla2m} [\emph{ʰlä˩m}]: квасцы\\
0x1962 \textbf{kla2m} [\emph{klä˩m}]: камелия\\
0x1968 \textbf{hca2m} [\emph{ʰʃä˩m}]: алоза\\
0x1969 \textbf{sla2m} [\emph{slä˩m}]: псалтырь\\
0x196A \textbf{tla2m} [\emph{tlä˩m}]: таллий\\
0x196D \textbf{cla2m} [\emph{ʃlä˩m}]: родственная душа\\
0x1970 \textbf{hra2m} [\emph{ʰrä˩m}]: драхма\\
0x1988 \textbf{hba2m} [\emph{ʰbä˩m}]: Бэтмен\\
0x1998 \textbf{hda2m} [\emph{ʰdä˩m}]: догматизм\\
0x19AA \textbf{tca2m} [\emph{tʃä˩m}]: таоизм\\
0x19B0 \textbf{hva2m} [\emph{ʰvä˩m}]: ванадий\\
0x19C1 \textbf{mra2m} [\emph{mrä˩m}]: подобный мрамору\\
0x19C2 \textbf{kra2m} [\emph{krä˩m}]: Катись\\
0x19C4 \textbf{pra2m} [\emph{prä˩m}]: отчество\\
0x19C6 \textbf{nra2m} [\emph{nrä˩m}]: неграмотность\\
0x19C9 \textbf{sra2m} [\emph{srä˩m}]: перевооружать\\
0x19CA \textbf{tra2m} [\emph{trä˩m}]: опилки\\
0x19D0 \textbf{hxa2m} [\emph{ʰxä˩m}]: анемия\\
0x19E2 \textbf{kxa2m} [\emph{kxä˩m}]: кремировать\\
0x19EA \textbf{txa2m} [\emph{txä˩m}]: козыри\\
0x1A48 \textbf{hsu2m} [\emph{ʰsu˩m}]: сумах\\
0x1B08 \textbf{hme2m} [\emph{ʰme˩m}]: мезодерма\\
0x1B20 \textbf{hpe2m} [\emph{ʰpe˩m}]: плевра\\
0x1B30 \textbf{hne2m} [\emph{ʰne˩m}]: клизма\\
0x1B48 \textbf{hse2m} [\emph{ʰse˩m}]: шестнадцатеричный\\
0x1B50 \textbf{hte2m} [\emph{ʰte˩m}]: эктодерма\\
0x1B58 \textbf{hle2m} [\emph{ʰle˩m}]: лейкемия\\
0x1B70 \textbf{hre2m} [\emph{ʰre˩m}]: брюшина\\
0x1B98 \textbf{hde2m} [\emph{ʰde˩m}]: энтодерма\\
0x1BCC \textbf{fre2m} [\emph{fre˩m}]: Ефрем\\
0x1BD0 \textbf{hxe2m} [\emph{ʰxe˩m}]: гем\\
0x1C08 \textbf{hmo2m} [\emph{ʰmo˩m}]: мяукать\\
0x1C10 \textbf{hko2m} [\emph{ʰko˩m}]: креольский\\
0x1C30 \textbf{hno2m} [\emph{ʰno˩m}]: монограмма\\
0x1C48 \textbf{hso2m} [\emph{ʰso˩m}]: мазохизм\\
0x1C50 \textbf{hto2m} [\emph{ʰto˩m}]: погруженный\\
0x1C58 \textbf{hlo2m} [\emph{ʰlo˩m}]: клинометр\\
0x1C70 \textbf{hro2m} [\emph{ʰro˩m}]: аэродинамический\\
0x1CC4 \textbf{pro2m} [\emph{pro˩m}]: погром\\
0x1CC6 \textbf{nro2m} [\emph{nro˩m}]: интерферометр\\
\subsection{ 2 }
0x203E \textbf{mi} [\emph{mi}]: меня\\
0x205E \textbf{ki} [\emph{ki}]: мимо напряженный\\
0x207E \textbf{yi} [\emph{ji}]: dative case, indicates motion towards the noun phrase. destination-case. \\
0x209E \textbf{pi} [\emph{pi}]: поверхность контекст\\
0x20BE \textbf{wi} [\emph{wi}]: универсальный квантификатором\\
0x20DE \textbf{ni} [\emph{ni}]: инфинитив\\
0x213E \textbf{si} [\emph{si}]: epistemic mood, to indicate that something is possible.\\
0x215E \textbf{ti} [\emph{ti}]: родительный дело\\
0x217E \textbf{li} [\emph{li}]: realis mood, for declaring something, the most common kind of sentence.\\
0x219E \textbf{fi} [\emph{fi}]: вернуть\\
0x21BE \textbf{ci} [\emph{ʃi}]: имя прилагательное\\
0x21DE \textbf{ri} [\emph{ri}]: interrogative mood, for questions.\\
0x223E \textbf{bi} [\emph{bi}]: будущее напряженный\\
0x225E \textbf{gi} [\emph{gi}]: имя\\
0x227E \textbf{di} [\emph{di}]: директива настроение\\
0x229E \textbf{zi} [\emph{zi}]: процитировал\\
0x22BE \textbf{ji} [\emph{ʒi}]: женский Пол\\
0x22DE \textbf{vi} [\emph{vi}]: глагол - суффиксов\\
0x233E \textbf{qi} [\emph{ŋi}]: переходный\\
0x235E \textbf{xi} [\emph{xi}]: habitual aspect, something that happens habitualy.\\
0x243E \textbf{ma} [\emph{mä}]: место назначения дело\\
0x245E \textbf{ka} [\emph{kä}]: винительный дело\\
0x247E \textbf{ya} [\emph{jä}]: Это\\
0x249E \textbf{pa} [\emph{pä}]: человек контекст\\
0x24BE \textbf{wa} [\emph{wä}]: а также\\
0x24DE \textbf{na} [\emph{nä}]: именительный дело\\
0x253E \textbf{sa} [\emph{sä}]: вы\\
0x255E \textbf{ta} [\emph{tä}]: topic case, the topic about which the sentence is. `regarding topic, some sentence.''. discourse-context way-case\\
0x257E \textbf{la} [\emph{lä}]: dependent clause, similar in usage to `that', a sentence that is within another sentence\\
0x259E \textbf{fa} [\emph{fä}]: perfective aspect, a completed action taken as a whole.\\
0x25BE \textbf{ca} [\emph{ʃä}]: Здравствуйте\\
0x25DE \textbf{ra} [\emph{rä}]: настоящее время напряженный\\
0x263E \textbf{ba} [\emph{bä}]: а также или\\
0x265E \textbf{ga} [\emph{gä}]: путь дело\\
0x267E \textbf{da} [\emph{dä}]: adpositional case, indicates preceding stem is the case for this noun phrase.\\
0x269E \textbf{za} [\emph{zä}]: inchoative aspect, for expressing the begining of a state or action\\
0x26BE \textbf{ja} [\emph{ʒä}]: связка\\
0x26DE \textbf{va} [\emph{vä}]: ватт\\
0x273E \textbf{qa} [\emph{ŋä}]: cessative aspect, for referring to the end of an action, when it is ceasing. inchoative-aspect- is begining\\
0x275E \textbf{xa} [\emph{xä}]: речь контекст\\
0x2801 \textbf{myik} [\emph{mjik}]: машина\\
0x2802 \textbf{kyik} [\emph{kjik}]: потому как\\
0x2804 \textbf{pyik} [\emph{pjik}]: поезд\\
0x2806 \textbf{nyik} [\emph{njik}]: грязь\\
0x2808 \textbf{hmik} [\emph{ʰmik}]: Музыка\\
0x2809 \textbf{syik} [\emph{sjik}]: политическое\\
0x280A \textbf{tyik} [\emph{tjik}]: дать согласие\\
0x280B \textbf{lyik} [\emph{ljik}]: липкий\\
0x280C \textbf{fyik} [\emph{fjik}]: виски\\
0x280D \textbf{cyik} [\emph{ʃjik}]: власть\\
0x280E \textbf{ryik} [\emph{rjik}]: химическая\\
0x2810 \textbf{hkik} [\emph{ʰkik}]: кто\\
0x2811 \textbf{byik} [\emph{bjik}]: бойкотировать\\
0x2812 \textbf{gyik} [\emph{gjik}]: математика\\
0x2813 \textbf{dyik} [\emph{djik}]: детектор\\
0x2814 \textbf{zyik} [\emph{zjik}]: цинк\\
0x2815 \textbf{jyik} [\emph{ʒjik}]: безличный глагол\\
0x2816 \textbf{vyik} [\emph{vjik}]: прокурор\\
0x2818 \textbf{hyik} [\emph{ʰjik}]: один\\
0x2819 \textbf{qyik} [\emph{ŋjik}]: соединительный частица\\
0x281A \textbf{xyik} [\emph{xjik}]: дискета диск\\
0x2820 \textbf{hpik} [\emph{ʰpik}]: аблятив дело\\
0x2821 \textbf{mwik} [\emph{mwik}]: магнитный\\
0x2822 \textbf{kwik} [\emph{kwik}]: квотатив дело\\
0x2824 \textbf{pwik} [\emph{pwik}]: перспективный\\
0x2826 \textbf{nwik} [\emph{nwik}]: инессив дело\\
0x2828 \textbf{hwik} [\emph{ʰwik}]: звательный дело\\
0x2829 \textbf{swik} [\emph{swik}]: щий дело\\
0x282A \textbf{twik} [\emph{twik}]: terminative дело\\
0x282B \textbf{lwik} [\emph{lwik}]: allative дело\\
0x282C \textbf{fwik} [\emph{fwik}]: салют\\
0x282D \textbf{cwik} [\emph{ʃwik}]: величественный\\
0x282E \textbf{rwik} [\emph{rwik}]: вертикальный\\
0x2830 \textbf{hnik} [\emph{ʰnik}]: родительный дело\\
0x2831 \textbf{bwik} [\emph{bwik}]: sublative дело\\
0x2832 \textbf{gwik} [\emph{gwik}]: циклический\\
0x2833 \textbf{dwik} [\emph{dwik}]: дивиденд\\
0x2834 \textbf{zwik} [\emph{zwik}]: горнолыжный спорт\\
0x2835 \textbf{jwik} [\emph{ʒwik}]: желток\\
0x2836 \textbf{vwik} [\emph{vwik}]: уховертка\\
0x2839 \textbf{qwik} [\emph{ŋwik}]: пневматический\\
0x283A \textbf{xwik} [\emph{xwik}]: гидравлический\\
0x283E \textbf{mu} [\emph{mu}]: commissive mood, for making promises, commitments\\
0x2842 \textbf{ksik} [\emph{ksik}]: следствие\\
0x2844 \textbf{psik} [\emph{psik}]: оратор\\
0x2848 \textbf{hsik} [\emph{ʰsik}]: общественного\\
0x284A \textbf{tsik} [\emph{tsik}]: злость\\
0x2850 \textbf{htik} [\emph{ʰtik}]: дательный дело\\
0x2851 \textbf{bzik} [\emph{bzik}]: физика\\
0x2852 \textbf{gzik} [\emph{gzik}]: каток\\
0x2853 \textbf{dzik} [\emph{dzik}]: двуязычный\\
0x2858 \textbf{hlik} [\emph{ʰlik}]: выражающий заключение дело\\
0x285E \textbf{ku} [\emph{ku}]: Эксклюзивный или\\
0x2860 \textbf{hfik} [\emph{ʰfik}]: распределительный дело\\
0x2862 \textbf{klik} [\emph{klik}]: долина\\
0x2864 \textbf{plik} [\emph{plik}]: применимый\\
0x2868 \textbf{hcik} [\emph{ʰʃik}]: думать\\
0x2869 \textbf{slik} [\emph{slik}]: шелк\\
0x286A \textbf{tlik} [\emph{tlik}]: мультипликативный дело\\
0x286C \textbf{flik} [\emph{flik}]: лиса\\
0x286D \textbf{clik} [\emph{ʃlik}]: охота\\
0x2870 \textbf{hrik} [\emph{ʰrik}]: практическое\\
0x2871 \textbf{blik} [\emph{blik}]: ястреб\\
0x2872 \textbf{glik} [\emph{glik}]: орошение\\
0x2873 \textbf{dlik} [\emph{dlik}]: законодательная власть\\
0x2874 \textbf{zlik} [\emph{zlik}]: кремний\\
0x2875 \textbf{jlik} [\emph{ʒlik}]: геологических\\
0x2876 \textbf{vlik} [\emph{vlik}]: физиология\\
0x287A \textbf{xlik} [\emph{xlik}]: гиперссылка\\
0x287E \textbf{yu} [\emph{ju}]: instrumental case, indicates the method or tool used. person-context-way-case.\\
0x2882 \textbf{kfik} [\emph{kfik}]: технический\\
0x2884 \textbf{pfik} [\emph{pfik}]: истребование\\
0x2888 \textbf{hbik} [\emph{ʰbik}]: благословить\\
0x288A \textbf{tfik} [\emph{tfik}]: отражать\\
0x2890 \textbf{hgik} [\emph{ʰgik}]: колено\\
0x2891 \textbf{bvik} [\emph{bvik}]: subessive дело\\
0x2892 \textbf{gvik} [\emph{gvik}]: эргатив дело\\
0x2893 \textbf{dvik} [\emph{dvik}]: адессив дело\\
0x2898 \textbf{hdik} [\emph{ʰdik}]: болезнь\\
0x289E \textbf{pu} [\emph{pu}]: множественное число номер\\
0x28A0 \textbf{hzik} [\emph{ʰzik}]: потухший\\
0x28A2 \textbf{kcik} [\emph{kʃik}]: рыба\\
0x28A4 \textbf{pcik} [\emph{pʃik}]: по всей видимости\\
0x28A8 \textbf{hjik} [\emph{ʰʒik}]: гостеприимство\\
0x28AA \textbf{tcik} [\emph{tʃik}]: доказательный дело\\
0x28B0 \textbf{hvik} [\emph{ʰvik}]: пробуждая\\
0x28B1 \textbf{bjik} [\emph{bʒik}]: электронная книга\\
0x28B2 \textbf{gjik} [\emph{gʒik}]: почка\\
0x28B3 \textbf{djik} [\emph{dʒik}]: задерживаться\\
0x28BE \textbf{wu} [\emph{wu}]: vocative case, the intended recipient of the sentence. for example `O' in `O god, hear my prayers.' discourse-context destination-case\\
0x28C1 \textbf{mrik} [\emph{mrik}]: механический\\
0x28C2 \textbf{krik} [\emph{krik}]: критика\\
0x28C4 \textbf{prik} [\emph{prik}]: пассивный\\
0x28C6 \textbf{nrik} [\emph{nrik}]: определенный статья\\
0x28C8 \textbf{hqik} [\emph{ʰŋik}]: вводить\\
0x28C9 \textbf{srik} [\emph{srik}]: семена\\
0x28CA \textbf{trik} [\emph{trik}]: пьяный\\
0x28CC \textbf{frik} [\emph{frik}]: теоретический\\
0x28CD \textbf{crik} [\emph{ʃrik}]: ад\\
0x28D0 \textbf{hxik} [\emph{ʰxik}]: конфигурировать\\
0x28D1 \textbf{brik} [\emph{brik}]: подписываться\\
0x28D2 \textbf{grik} [\emph{grik}]: гравий\\
0x28D3 \textbf{drik} [\emph{drik}]: странно\\
0x28D4 \textbf{zrik} [\emph{zrik}]: поставить в известность\\
0x28D5 \textbf{jrik} [\emph{ʒrik}]: крикет\\
0x28D6 \textbf{vrik} [\emph{vrik}]: окончательный частица\\
0x28D9 \textbf{qrik} [\emph{ŋrik}]: центральный дело\\
0x28DA \textbf{xrik} [\emph{xrik}]: пациент дело\\
0x28DE \textbf{nu} [\emph{nu}]: интерьер контекст\\
0x28E2 \textbf{kxik} [\emph{kxik}]: Народная музыка\\
0x28E4 \textbf{pxik} [\emph{pxik}]: параюридический\\
0x28EA \textbf{txik} [\emph{txik}]: указательный палец\\
0x28F1 \textbf{bxik} [\emph{bxik}]: обитать\\
0x28F2 \textbf{gxik} [\emph{gxik}]: гражданские права\\
0x28F3 \textbf{dxik} [\emph{dxik}]: корица\\
0x2901 \textbf{myak} [\emph{mjäk}]: Март\\
0x2902 \textbf{kyak} [\emph{kjäk}]: несколько\\
0x2904 \textbf{pyak} [\emph{pjäk}]: брат\\
0x2906 \textbf{nyak} [\emph{njäk}]: товарный знак\\
0x2908 \textbf{hmak} [\emph{ʰmäk}]: comitative дело\\
0x2909 \textbf{syak} [\emph{sjäk}]: войти\\
0x290A \textbf{tyak} [\emph{tjäk}]: таким образом\\
0x290B \textbf{lyak} [\emph{ljäk}]: эластичный\\
0x290C \textbf{fyak} [\emph{fjäk}]: фраза\\
0x290D \textbf{cyak} [\emph{ʃjäk}]: небо\\
0x290E \textbf{ryak} [\emph{rjäk}]: реакция\\
0x2910 \textbf{hkak} [\emph{ʰkäk}]: винительный дело\\
0x2911 \textbf{byak} [\emph{bjäk}]: козел\\
0x2912 \textbf{gyak} [\emph{gjäk}]: средний класс\\
0x2913 \textbf{dyak} [\emph{djäk}]: танцор\\
0x2914 \textbf{zyak} [\emph{zjäk}]: синагога\\
0x2915 \textbf{jyak} [\emph{ʒjäk}]: вмещать\\
0x2916 \textbf{vyak} [\emph{vjäk}]: адвербиальный дело\\
0x2918 \textbf{hyak} [\emph{ʰjäk}]: инструментальный дело\\
0x2919 \textbf{qyak} [\emph{ŋjäk}]: Один действие глагол\\
0x291A \textbf{xyak} [\emph{xjäk}]: Ямайка\\
0x2920 \textbf{hpak} [\emph{ʰpäk}]: благословляющий дело\\
0x2921 \textbf{mwak} [\emph{mwäk}]: куча\\
0x2922 \textbf{kwak} [\emph{kwäk}]: покрытый\\
0x2924 \textbf{pwak} [\emph{pwäk}]: коробка\\
0x2926 \textbf{nwak} [\emph{nwäk}]: январь\\
0x2928 \textbf{hwak} [\emph{ʰwäk}]: гласная буква\\
0x2929 \textbf{swak} [\emph{swäk}]: раковина\\
0x292A \textbf{twak} [\emph{twäk}]: delative дело\\
0x292B \textbf{lwak} [\emph{lwäk}]: perlative дело\\
0x292C \textbf{fwak} [\emph{fwäk}]: судья\\
0x292D \textbf{cwak} [\emph{ʃwäk}]: зло\\
0x292E \textbf{rwak} [\emph{rwäk}]: трещина\\
0x2930 \textbf{hnak} [\emph{ʰnäk}]: именительный дело\\
0x2931 \textbf{bwak} [\emph{bwäk}]: кукла\\
0x2932 \textbf{gwak} [\emph{gwäk}]: лодыжка\\
0x2933 \textbf{dwak} [\emph{dwäk}]: срывать\\
0x2934 \textbf{zwak} [\emph{zwäk}]: муляж\\
0x2935 \textbf{jwak} [\emph{ʒwäk}]: йога\\
0x2936 \textbf{vwak} [\emph{vwäk}]: Словакия\\
0x2939 \textbf{qwak} [\emph{ŋwäk}]: чайка\\
0x293A \textbf{xwak} [\emph{xwäk}]: мочевой пузырь\\
0x293E \textbf{su} [\emph{su}]: Кроме\\
0x2942 \textbf{ksak} [\emph{ksäk}]: причинный дело\\
0x2944 \textbf{psak} [\emph{psäk}]: прохождение\\
0x2948 \textbf{hsak} [\emph{ʰsäk}]: время контекст\\
0x294A \textbf{tsak} [\emph{tsäk}]: уважение\\
0x2950 \textbf{htak} [\emph{ʰtäk}]: тема дело\\
0x2951 \textbf{bzak} [\emph{bzäk}]: носок\\
0x2952 \textbf{gzak} [\emph{gzäk}]: эстрагон\\
0x2953 \textbf{dzak} [\emph{dzäk}]: секта\\
0x2958 \textbf{hlak} [\emph{ʰläk}]: относительный дело\\
0x295E \textbf{tu} [\emph{tu}]: deontic mood, to indicate what should be different.\\
0x2960 \textbf{hfak} [\emph{ʰfäk}]: 5\\
0x2962 \textbf{klak} [\emph{kläk}]: собака\\
0x2964 \textbf{plak} [\emph{pläk}]: Население\\
0x2968 \textbf{hcak} [\emph{ʰʃäk}]: под контекст\\
0x2969 \textbf{slak} [\emph{släk}]: успех\\
0x296A \textbf{tlak} [\emph{tläk}]: приносить\\
0x296C \textbf{flak} [\emph{fläk}]: знакомые\\
0x296D \textbf{clak} [\emph{ʃläk}]: сахар\\
0x2970 \textbf{hrak} [\emph{ʰräk}]: сотрудники\\
0x2971 \textbf{blak} [\emph{bläk}]: балкон\\
0x2972 \textbf{glak} [\emph{gläk}]: вулкан\\
0x2973 \textbf{dlak} [\emph{dläk}]: обсуждать\\
0x2974 \textbf{zlak} [\emph{zläk}]: нарушать\\
0x2975 \textbf{jlak} [\emph{ʒläk}]: якорь\\
0x2976 \textbf{vlak} [\emph{vläk}]: стерилизовать\\
0x297A \textbf{xlak} [\emph{xläk}]: автомотриса\\
0x297E \textbf{lu} [\emph{lu}]: speculative mood, for expressing speculation, hypothesis, ``could be''\\
0x2982 \textbf{kfak} [\emph{kfäk}]: израсходовать\\
0x2984 \textbf{pfak} [\emph{pfäk}]: благоприятствовать\\
0x2988 \textbf{hbak} [\emph{ʰbäk}]: продавать\\
0x298A \textbf{tfak} [\emph{tfäk}]: доказательный\\
0x2990 \textbf{hgak} [\emph{ʰgäk}]: Кот\\
0x2991 \textbf{bvak} [\emph{bväk}]: abessive дело\\
0x2992 \textbf{gvak} [\emph{gväk}]: грамматическая случая\\
0x2993 \textbf{dvak} [\emph{dväk}]: нигиляционными дело\\
0x2998 \textbf{hdak} [\emph{ʰdäk}]: магазин\\
0x299E \textbf{fu} [\emph{fu}]: фокус\\
0x29A0 \textbf{hzak} [\emph{ʰzäk}]: трясти\\
0x29A2 \textbf{kcak} [\emph{kʃäk}]: еда\\
0x29A4 \textbf{pcak} [\emph{pʃäk}]: запад\\
0x29A8 \textbf{hjak} [\emph{ʰʒäk}]: охватывать\\
0x29AA \textbf{tcak} [\emph{tʃäk}]: одна и та же предмет маркер\\
0x29B0 \textbf{hvak} [\emph{ʰväk}]: авария\\
0x29B1 \textbf{bjak} [\emph{bʒäk}]: экстатический\\
0x29B2 \textbf{gjak} [\emph{gʒäk}]: галактика\\
0x29B3 \textbf{djak} [\emph{dʒäk}]: разъем\\
0x29BE \textbf{cu} [\emph{ʃu}]: conditional mood, for if then statements. The if clause comes before the `cu' the then clause comes after.\\
0x29C1 \textbf{mrak} [\emph{mräk}]: копье\\
0x29C2 \textbf{krak} [\emph{kräk}]: признание\\
0x29C4 \textbf{prak} [\emph{präk}]: проповедовать\\
0x29C6 \textbf{nrak} [\emph{nräk}]: континент\\
0x29C8 \textbf{hqak} [\emph{ʰŋäk}]: валять дурака\\
0x29C9 \textbf{srak} [\emph{sräk}]: рукопись\\
0x29CA \textbf{trak} [\emph{träk}]: годовой\\
0x29CC \textbf{frak} [\emph{fräk}]: неоднократно\\
0x29CD \textbf{crak} [\emph{ʃräk}]: очарование\\
0x29D0 \textbf{hxak} [\emph{ʰxäk}]: перерабатывать\\
0x29D1 \textbf{brak} [\emph{bräk}]: баррель\\
0x29D2 \textbf{grak} [\emph{gräk}]: царапина\\
0x29D3 \textbf{drak} [\emph{dräk}]: расстройство\\
0x29D4 \textbf{zrak} [\emph{zräk}]: акула\\
0x29D5 \textbf{jrak} [\emph{ʒräk}]: лягушка\\
0x29D6 \textbf{vrak} [\emph{vräk}]: парик\\
0x29D9 \textbf{qrak} [\emph{ŋräk}]: штрейкбрехер\\
0x29DA \textbf{xrak} [\emph{xräk}]: как ни странно\\
0x29DE \textbf{ru} [\emph{ru}]: prohibitive mood, for expressing something is not allowed, inverse of permissive-mood-, potential response to admonitive mood\\
0x29E2 \textbf{kxak} [\emph{kxäk}]: к дело\\
0x29E4 \textbf{pxak} [\emph{pxäk}]: завернутый\\
0x29EA \textbf{txak} [\emph{txäk}]: срабатывает\\
0x29F1 \textbf{bxak} [\emph{bxäk}]: закладка\\
0x29F2 \textbf{gxak} [\emph{gxäk}]: дорогостоящий\\
0x29F3 \textbf{dxak} [\emph{dxäk}]: данные добыча\\
0x2A01 \textbf{myuk} [\emph{mjuk}]: дурак\\
0x2A02 \textbf{kyuk} [\emph{kjuk}]: шея\\
0x2A04 \textbf{pyuk} [\emph{pjuk}]: чума\\
0x2A06 \textbf{nyuk} [\emph{njuk}]: Светляк\\
0x2A08 \textbf{hmuk} [\emph{ʰmuk}]: минут\\
0x2A09 \textbf{syuk} [\emph{sjuk}]: ложка\\
0x2A0A \textbf{tyuk} [\emph{tjuk}]: восток\\
0x2A0B \textbf{lyuk} [\emph{ljuk}]: последующий пункт\\
0x2A0C \textbf{fyuk} [\emph{fjuk}]: конфабуляция\\
0x2A0D \textbf{cyuk} [\emph{ʃjuk}]: привлечь\\
0x2A0E \textbf{ryuk} [\emph{rjuk}]: енот\\
0x2A10 \textbf{hkuk} [\emph{ʰkuk}]: Предложение хвост\\
0x2A11 \textbf{byuk} [\emph{bjuk}]: страшилище\\
0x2A13 \textbf{dyuk} [\emph{djuk}]: педикюр\\
0x2A18 \textbf{hyuk} [\emph{ʰjuk}]: красный\\
0x2A19 \textbf{qyuk} [\emph{ŋjuk}]: равноугольный\\
0x2A1A \textbf{xyuk} [\emph{xjuk}]: индеек\\
0x2A20 \textbf{hpuk} [\emph{ʰpuk}]: книга\\
0x2A21 \textbf{mwuk} [\emph{mwuk}]: норка\\
0x2A22 \textbf{kwuk} [\emph{kwuk}]: конго\\
0x2A24 \textbf{pwuk} [\emph{pwuk}]: наволочка\\
0x2A26 \textbf{nwuk} [\emph{nwuk}]: уникальный\\
0x2A28 \textbf{hwuk} [\emph{ʰwuk}]: свободный\\
0x2A29 \textbf{swuk} [\emph{swuk}]: превосходный дело\\
0x2A2A \textbf{twuk} [\emph{twuk}]: очистка\\
0x2A2B \textbf{lwuk} [\emph{lwuk}]: ребро\\
0x2A2D \textbf{cwuk} [\emph{ʃwuk}]: Старый мир\\
0x2A2E \textbf{rwuk} [\emph{rwuk}]: Уругвай\\
0x2A30 \textbf{hnuk} [\emph{ʰnuk}]: герцог\\
0x2A31 \textbf{bwuk} [\emph{bwuk}]: букет\\
0x2A32 \textbf{gwuk} [\emph{gwuk}]: глюкоза\\
0x2A3A \textbf{xwuk} [\emph{xwuk}]: галлюциногенный\\
0x2A3E \textbf{bu} [\emph{bu}]: абсолютива дело\\
0x2A42 \textbf{ksuk} [\emph{ksuk}]: дядя\\
0x2A44 \textbf{psuk} [\emph{psuk}]: сломанный\\
0x2A48 \textbf{hsuk} [\emph{ʰsuk}]: источник дело\\
0x2A4A \textbf{tsuk} [\emph{tsuk}]: ранг\\
0x2A50 \textbf{htuk} [\emph{ʰtuk}]: временной дело\\
0x2A58 \textbf{hluk} [\emph{ʰluk}]: устала\\
0x2A5E \textbf{gu} [\emph{gu}]: гм\\
0x2A60 \textbf{hfuk} [\emph{ʰfuk}]: Огонь\\
0x2A62 \textbf{kluk} [\emph{kluk}]: археологический\\
0x2A64 \textbf{pluk} [\emph{pluk}]: оценить\\
0x2A68 \textbf{hcuk} [\emph{ʰʃuk}]: фокус\\
0x2A69 \textbf{sluk} [\emph{sluk}]: белка\\
0x2A6A \textbf{tluk} [\emph{tluk}]: вводный\\
0x2A6C \textbf{fluk} [\emph{fluk}]: фолликул\\
0x2A6D \textbf{cluk} [\emph{ʃluk}]: уродливый\\
0x2A70 \textbf{hruk} [\emph{ʰruk}]: архитектура\\
0x2A71 \textbf{bluk} [\emph{bluk}]: тактика\\
0x2A72 \textbf{gluk} [\emph{gluk}]: пуля\\
0x2A73 \textbf{dluk} [\emph{dluk}]: Отключить\\
0x2A74 \textbf{zluk} [\emph{zluk}]: зулусский\\
0x2A75 \textbf{jluk} [\emph{ʒluk}]: моллюсками\\
0x2A76 \textbf{vluk} [\emph{vluk}]: вольта лицо\\
0x2A7A \textbf{xluk} [\emph{xluk}]: сосулька\\
0x2A7E \textbf{du} [\emph{du}]: второй\\
0x2A82 \textbf{kfuk} [\emph{kfuk}]: уход за кожей\\
0x2A84 \textbf{pfuk} [\emph{pfuk}]: надколенник\\
0x2A88 \textbf{hbuk} [\emph{ʰbuk}]: бисквит\\
0x2A8A \textbf{tfuk} [\emph{tfuk}]: домашний скот\\
0x2A90 \textbf{hguk} [\emph{ʰguk}]: тыква\\
0x2A91 \textbf{bvuk} [\emph{bvuk}]: абсолютива дело\\
0x2A92 \textbf{gvuk} [\emph{gvuk}]: вогнуто выпуклая\\
0x2A98 \textbf{hduk} [\emph{ʰduk}]: крюк\\
0x2AA0 \textbf{hzuk} [\emph{ʰzuk}]: зоопарк\\
0x2AA2 \textbf{kcuk} [\emph{kʃuk}]: школа\\
0x2AA4 \textbf{pcuk} [\emph{pʃuk}]: послелоги дело\\
0x2AA8 \textbf{hjuk} [\emph{ʰʒuk}]: становиться на колени\\
0x2AAA \textbf{tcuk} [\emph{tʃuk}]: боль\\
0x2AB0 \textbf{hvuk} [\emph{ʰvuk}]: съеденный\\
0x2AC1 \textbf{mruk} [\emph{mruk}]: улучшать\\
0x2AC2 \textbf{kruk} [\emph{kruk}]: костыль\\
0x2AC4 \textbf{pruk} [\emph{pruk}]: супермаркет\\
0x2AC6 \textbf{nruk} [\emph{nruk}]: инфекционный\\
0x2AC8 \textbf{hquk} [\emph{ʰŋuk}]: кенгуру\\
0x2AC9 \textbf{sruk} [\emph{sruk}]: цирк\\
0x2ACA \textbf{truk} [\emph{truk}]: стучать\\
0x2ACC \textbf{fruk} [\emph{fruk}]: Африка\\
0x2ACD \textbf{cruk} [\emph{ʃruk}]: кролик\\
0x2AD0 \textbf{hxuk} [\emph{ʰxuk}]: изнашивания\\
0x2AD1 \textbf{bruk} [\emph{bruk}]: баррикада\\
0x2AD2 \textbf{gruk} [\emph{gruk}]: Глубинные\\
0x2AD3 \textbf{druk} [\emph{druk}]: полукруглый\\
0x2AD4 \textbf{zruk} [\emph{zruk}]: стержневой корень\\
0x2AD5 \textbf{jruk} [\emph{ʒruk}]: хрящ\\
0x2AD9 \textbf{qruk} [\emph{ŋruk}]: надстройка\\
0x2ADA \textbf{xruk} [\emph{xruk}]: линька\\
0x2AE2 \textbf{kxuk} [\emph{kxuk}]: выпуклый\\
0x2AE4 \textbf{pxuk} [\emph{pxuk}]: Португалия\\
0x2AEA \textbf{txuk} [\emph{txuk}]: короткое замыкание\\
0x2AF1 \textbf{bxuk} [\emph{bxuk}]: подпочва\\
0x2AF2 \textbf{gxuk} [\emph{gxuk}]: сумерки\\
0x2AF3 \textbf{dxuk} [\emph{dxuk}]: ломкий\\
0x2B01 \textbf{myek} [\emph{mjek}]: термодинамический\\
0x2B02 \textbf{kyek} [\emph{kjek}]: кекс\\
0x2B04 \textbf{pyek} [\emph{pjek}]: проектор\\
0x2B06 \textbf{nyek} [\emph{njek}]: никель\\
0x2B08 \textbf{hmek} [\emph{ʰmek}]: чувствительный\\
0x2B09 \textbf{syek} [\emph{sjek}]: класть не на место\\
0x2B0A \textbf{tyek} [\emph{tjek}]: треугольник\\
0x2B0B \textbf{lyek} [\emph{ljek}]: латинский\\
0x2B0C \textbf{fyek} [\emph{fjek}]: фонетический\\
0x2B0D \textbf{cyek} [\emph{ʃjek}]: мешок\\
0x2B0E \textbf{ryek} [\emph{rjek}]: ракета\\
0x2B10 \textbf{hkek} [\emph{ʰkek}]: хвост\\
0x2B11 \textbf{byek} [\emph{bjek}]: божья коровка\\
0x2B12 \textbf{gyek} [\emph{gjek}]: кинетика\\
0x2B13 \textbf{dyek} [\emph{djek}]: отлаживать\\
0x2B15 \textbf{jyek} [\emph{ʒjek}]: термоядерный\\
0x2B18 \textbf{hyek} [\emph{ʰjek}]: путешественник\\
0x2B19 \textbf{qyek} [\emph{ŋjek}]: эротика\\
0x2B1A \textbf{xyek} [\emph{xjek}]: эмоционально\\
0x2B20 \textbf{hpek} [\emph{ʰpek}]: пакет\\
0x2B21 \textbf{mwek} [\emph{mwek}]: эмираты\\
0x2B22 \textbf{kwek} [\emph{kwek}]: местный дело\\
0x2B24 \textbf{pwek} [\emph{pwek}]: pegative дело\\
0x2B26 \textbf{nwek} [\emph{nwek}]: генетика\\
0x2B28 \textbf{hwek} [\emph{ʰwek}]: волк\\
0x2B29 \textbf{swek} [\emph{swek}]: перегружены работой\\
0x2B2A \textbf{twek} [\emph{twek}]: чрезмерное дело\\
0x2B2B \textbf{lwek} [\emph{lwek}]: рычаг\\
0x2B2C \textbf{fwek} [\emph{fwek}]: фотоэлектрические\\
0x2B2D \textbf{cwek} [\emph{ʃwek}]: хоккей\\
0x2B2E \textbf{rwek} [\emph{rwek}]: ретроспективный\\
0x2B30 \textbf{hnek} [\emph{ʰnek}]: инициатива дело\\
0x2B31 \textbf{bwek} [\emph{bwek}]: портфель\\
0x2B33 \textbf{dwek} [\emph{dwek}]: гидроэлектрический\\
0x2B36 \textbf{vwek} [\emph{vwek}]: Авант курьер\\
0x2B39 \textbf{qwek} [\emph{ŋwek}]: кормилец\\
0x2B3A \textbf{xwek} [\emph{xwek}]: фальшборт\\
0x2B42 \textbf{ksek} [\emph{ksek}]: насекомое\\
0x2B44 \textbf{psek} [\emph{psek}]: superessive дело\\
0x2B48 \textbf{hsek} [\emph{ʰsek}]: аспект\\
0x2B4A \textbf{tsek} [\emph{tsek}]: стейк\\
0x2B50 \textbf{htek} [\emph{ʰtek}]: библиотека\\
0x2B51 \textbf{bzek} [\emph{bzek}]: ребенок мозг\\
0x2B52 \textbf{gzek} [\emph{gzek}]: эпизоотия\\
0x2B53 \textbf{dzek} [\emph{dzek}]: декагон\\
0x2B58 \textbf{hlek} [\emph{ʰlek}]: atelic аспект\\
0x2B5E \textbf{xu} [\emph{xu}]: первичная объект\\
0x2B60 \textbf{hfek} [\emph{ʰfek}]: процветать\\
0x2B62 \textbf{klek} [\emph{klek}]: щелчок\\
0x2B64 \textbf{plek} [\emph{plek}]: капиллярный\\
0x2B68 \textbf{hcek} [\emph{ʰʃek}]: мебель\\
0x2B69 \textbf{slek} [\emph{slek}]: велосипед\\
0x2B6A \textbf{tlek} [\emph{tlek}]: штекер\\
0x2B6C \textbf{flek} [\emph{flek}]: подделывать\\
0x2B6D \textbf{clek} [\emph{ʃlek}]: утечка\\
0x2B70 \textbf{hrek} [\emph{ʰrek}]: восстановить\\
0x2B71 \textbf{blek} [\emph{blek}]: балет\\
0x2B72 \textbf{glek} [\emph{glek}]: реле\\
0x2B73 \textbf{dlek} [\emph{dlek}]: вектор\\
0x2B74 \textbf{zlek} [\emph{zlek}]: экзоскелет\\
0x2B75 \textbf{jlek} [\emph{ʒlek}]: геолог\\
0x2B76 \textbf{vlek} [\emph{vlek}]: надгортанник\\
0x2B7A \textbf{xlek} [\emph{xlek}]: ябеда\\
0x2B82 \textbf{kfek} [\emph{kfek}]: конвекция\\
0x2B84 \textbf{pfek} [\emph{pfek}]: орех-пекан\\
0x2B88 \textbf{hbek} [\emph{ʰbek}]: пекарь\\
0x2B8A \textbf{tfek} [\emph{tfek}]: местный направленный дело\\
0x2B90 \textbf{hgek} [\emph{ʰgek}]: гей\\
0x2B91 \textbf{bvek} [\emph{bvek}]: отмучивать\\
0x2B93 \textbf{dvek} [\emph{dvek}]: обнаруживать\\
0x2B98 \textbf{hdek} [\emph{ʰdek}]: светский\\
0x2BA0 \textbf{hzek} [\emph{ʰzek}]: марганец\\
0x2BA2 \textbf{kcek} [\emph{kʃek}]: кривая\\
0x2BA4 \textbf{pcek} [\emph{pʃek}]: похищение людей\\
0x2BA8 \textbf{hjek} [\emph{ʰʒek}]: агентивного дело\\
0x2BAA \textbf{tcek} [\emph{tʃek}]: эхо\\
0x2BB0 \textbf{hvek} [\emph{ʰvek}]: evitative дело\\
0x2BB2 \textbf{gjek} [\emph{gʒek}]: еще - если\\
0x2BB3 \textbf{djek} [\emph{dʒek}]: еж\\
0x2BC1 \textbf{mrek} [\emph{mrek}]: поражать\\
0x2BC2 \textbf{krek} [\emph{krek}]: Меркурий\\
0x2BC4 \textbf{prek} [\emph{prek}]: стоянка\\
0x2BC6 \textbf{nrek} [\emph{nrek}]: работник\\
0x2BC8 \textbf{hqek} [\emph{ʰŋek}]: косой дело\\
0x2BC9 \textbf{srek} [\emph{srek}]: уксус\\
0x2BCA \textbf{trek} [\emph{trek}]: подрядчик\\
0x2BCC \textbf{frek} [\emph{frek}]: пепельница\\
0x2BCD \textbf{crek} [\emph{ʃrek}]: предлагать\\
0x2BD0 \textbf{hxek} [\emph{ʰxek}]: изменение формы\\
0x2BD1 \textbf{brek} [\emph{brek}]: тормоз\\
0x2BD2 \textbf{grek} [\emph{grek}]: гараж\\
0x2BD3 \textbf{drek} [\emph{drek}]: в отдельности\\
0x2BD4 \textbf{zrek} [\emph{zrek}]: экзоцентричный дело\\
0x2BD5 \textbf{jrek} [\emph{ʒrek}]: черный как смоль\\
0x2BD6 \textbf{vrek} [\emph{vrek}]: электромагнитный\\
0x2BD9 \textbf{qrek} [\emph{ŋrek}]: завышенный\\
0x2BDA \textbf{xrek} [\emph{xrek}]: пациент вызывать\\
0x2BE2 \textbf{kxek} [\emph{kxek}]: краб\\
0x2BE4 \textbf{pxek} [\emph{pxek}]: доисторический\\
0x2BEA \textbf{txek} [\emph{txek}]: п- корень\\
0x2BF1 \textbf{bxek} [\emph{bxek}]: жуки\\
0x2BF2 \textbf{gxek} [\emph{gxek}]: высокие технологии\\
0x2BF3 \textbf{dxek} [\emph{dxek}]: знакомства\\
0x2C01 \textbf{myok} [\emph{mjok}]: мозаика\\
0x2C02 \textbf{kyok} [\emph{kjok}]: огурец\\
0x2C04 \textbf{pyok} [\emph{pjok}]: снаряд\\
0x2C06 \textbf{nyok} [\emph{njok}]: аналоговый\\
0x2C08 \textbf{hmok} [\emph{ʰmok}]: мельница\\
0x2C09 \textbf{syok} [\emph{sjok}]: сторонник\\
0x2C0A \textbf{tyok} [\emph{tjok}]: грузовая машина\\
0x2C0B \textbf{lyok} [\emph{ljok}]: Цель вызывать\\
0x2C0C \textbf{fyok} [\emph{fjok}]: все органические\\
0x2C0D \textbf{cyok} [\emph{ʃjok}]: пожимание плечами\\
0x2C0E \textbf{ryok} [\emph{rjok}]: протокол\\
0x2C10 \textbf{hkok} [\emph{ʰkok}]: Одинокий\\
0x2C11 \textbf{byok} [\emph{bjok}]: биотехнология\\
0x2C12 \textbf{gyok} [\emph{gjok}]: бег\\
0x2C13 \textbf{dyok} [\emph{djok}]: Доминика\\
0x2C15 \textbf{jyok} [\emph{ʒjok}]: безделушка\\
0x2C18 \textbf{hyok} [\emph{ʰjok}]: угол\\
0x2C1A \textbf{xyok} [\emph{xjok}]: наручник\\
0x2C20 \textbf{hpok} [\emph{ʰpok}]: племя\\
0x2C21 \textbf{mwok} [\emph{mwok}]: Монголия\\
0x2C22 \textbf{kwok} [\emph{kwok}]: кобальт\\
0x2C24 \textbf{pwok} [\emph{pwok}]: обвинительный дело\\
0x2C26 \textbf{nwok} [\emph{nwok}]: домашнее задание\\
0x2C28 \textbf{hwok} [\emph{ʰwok}]: квалифицированный\\
0x2C29 \textbf{swok} [\emph{swok}]: прильнуть\\
0x2C2A \textbf{twok} [\emph{twok}]: статистическая\\
0x2C2B \textbf{lwok} [\emph{lwok}]: лозунг\\
0x2C2C \textbf{fwok} [\emph{fwok}]: офтальмология\\
0x2C2D \textbf{cwok} [\emph{ʃwok}]: шок\\
0x2C2E \textbf{rwok} [\emph{rwok}]: провоцировать\\
0x2C30 \textbf{hnok} [\emph{ʰnok}]: содержание\\
0x2C31 \textbf{bwok} [\emph{bwok}]: брокколи\\
0x2C32 \textbf{gwok} [\emph{gwok}]: джокер\\
0x2C33 \textbf{dwok} [\emph{dwok}]: радиология\\
0x2C36 \textbf{vwok} [\emph{vwok}]: авокадо\\
0x2C39 \textbf{qwok} [\emph{ŋwok}]: ангольский\\
0x2C3A \textbf{xwok} [\emph{xwok}]: гео наука\\
0x2C3E \textbf{me} [\emph{me}]: imperfective aspect, an incomplete action in process.\\
0x2C42 \textbf{ksok} [\emph{ksok}]: Гороскоп\\
0x2C44 \textbf{psok} [\emph{psok}]: открытка\\
0x2C48 \textbf{hsok} [\emph{ʰsok}]: поставка\\
0x2C4A \textbf{tsok} [\emph{tsok}]: риторика\\
0x2C50 \textbf{htok} [\emph{ʰtok}]: конфликт\\
0x2C51 \textbf{bzok} [\emph{bzok}]: бензойный\\
0x2C52 \textbf{gzok} [\emph{gzok}]: геофизика\\
0x2C53 \textbf{dzok} [\emph{dzok}]: сухой док\\
0x2C58 \textbf{hlok} [\emph{ʰlok}]: прыжок\\
0x2C5E \textbf{ke} [\emph{ke}]: вопрос\\
0x2C60 \textbf{hfok} [\emph{ʰfok}]: скромный\\
0x2C62 \textbf{klok} [\emph{klok}]: избыток\\
0x2C64 \textbf{plok} [\emph{plok}]: дополнение\\
0x2C68 \textbf{hcok} [\emph{ʰʃok}]: торопиться\\
0x2C69 \textbf{slok} [\emph{slok}]: социология\\
0x2C6A \textbf{tlok} [\emph{tlok}]: бак\\
0x2C6C \textbf{flok} [\emph{flok}]: фламинго\\
0x2C6D \textbf{clok} [\emph{ʃlok}]: шоколад\\
0x2C70 \textbf{hrok} [\emph{ʰrok}]: реформа\\
0x2C71 \textbf{blok} [\emph{blok}]: биологическая\\
0x2C72 \textbf{glok} [\emph{glok}]: блог\\
0x2C73 \textbf{dlok} [\emph{dlok}]: роскошный\\
0x2C74 \textbf{zlok} [\emph{zlok}]: безалкогольный\\
0x2C75 \textbf{jlok} [\emph{ʒlok}]: кола\\
0x2C76 \textbf{vlok} [\emph{vlok}]: волан\\
0x2C7A \textbf{xlok} [\emph{xlok}]: дятел\\
0x2C7E \textbf{ye} [\emph{je}]: притяжательный падеж маркер\\
0x2C82 \textbf{kfok} [\emph{kfok}]: икота\\
0x2C84 \textbf{pfok} [\emph{pfok}]: модификатор\\
0x2C88 \textbf{hbok} [\emph{ʰbok}]: утка\\
0x2C8A \textbf{tfok} [\emph{tfok}]: мускатный орех\\
0x2C90 \textbf{hgok} [\emph{ʰgok}]: органический\\
0x2C91 \textbf{bvok} [\emph{bvok}]: букмекер\\
0x2C93 \textbf{dvok} [\emph{dvok}]: владивосток\\
0x2C98 \textbf{hdok} [\emph{ʰdok}]: вождение\\
0x2C9E \textbf{pe} [\emph{pe}]: в стороне\\
0x2CA0 \textbf{hzok} [\emph{ʰzok}]: извиваться\\
0x2CA2 \textbf{kcok} [\emph{kʃok}]: топор\\
0x2CA4 \textbf{pcok} [\emph{pʃok}]: спотыкаться\\
0x2CA8 \textbf{hjok} [\emph{ʰʒok}]: жокей\\
0x2CAA \textbf{tcok} [\emph{tʃok}]: точность\\
0x2CB0 \textbf{hvok} [\emph{ʰvok}]: водка\\
0x2CB1 \textbf{bjok} [\emph{bʒok}]: субдукции\\
0x2CB2 \textbf{gjok} [\emph{gʒok}]: догадок\\
0x2CB3 \textbf{djok} [\emph{dʒok}]: сопряженный\\
0x2CBE \textbf{we} [\emph{we}]: vialis case, indicates motion through the noun phrase. used a standalone case. way-case. \\
0x2CC1 \textbf{mrok} [\emph{mrok}]: комар\\
0x2CC2 \textbf{krok} [\emph{krok}]: череп\\
0x2CC4 \textbf{prok} [\emph{prok}]: оптический\\
0x2CC6 \textbf{nrok} [\emph{nrok}]: азот\\
0x2CC8 \textbf{hqok} [\emph{ʰŋok}]: друг друга\\
0x2CC9 \textbf{srok} [\emph{srok}]: редакционный\\
0x2CCA \textbf{trok} [\emph{trok}]: Головная боль\\
0x2CCC \textbf{frok} [\emph{frok}]: крутящий момент\\
0x2CCD \textbf{crok} [\emph{ʃrok}]: шоссе\\
0x2CD0 \textbf{hxok} [\emph{ʰxok}]: неиспользованный\\
0x2CD1 \textbf{brok} [\emph{brok}]: переедать\\
0x2CD2 \textbf{grok} [\emph{grok}]: морковь\\
0x2CD3 \textbf{drok} [\emph{drok}]: штаб-квартира\\
0x2CD4 \textbf{zrok} [\emph{zrok}]: Марокко\\
0x2CD5 \textbf{jrok} [\emph{ʒrok}]: Коперник\\
0x2CD6 \textbf{vrok} [\emph{vrok}]: Корея\\
0x2CD9 \textbf{qrok} [\emph{ŋrok}]: нет до\\
0x2CDA \textbf{xrok} [\emph{xrok}]: коллега по работе\\
0x2CDE \textbf{ne} [\emph{ne}]: мужской Пол\\
0x2CE2 \textbf{kxok} [\emph{kxok}]: какао\\
0x2CE4 \textbf{pxok} [\emph{pxok}]: пастораль\\
0x2CEA \textbf{txok} [\emph{txok}]: автоматически\\
0x2CF1 \textbf{bxok} [\emph{bxok}]: рококо\\
0x2CF2 \textbf{gxok} [\emph{gxok}]: готика\\
0x2CF3 \textbf{dxok} [\emph{dxok}]: доминиканский\\
0x2D01 \textbf{my6k} [\emph{mjək}]: симметричный\\
0x2D04 \textbf{py6k} [\emph{pjək}]: папайя\\
0x2D06 \textbf{ny6k} [\emph{njək}]: энный мощность\\
0x2D08 \textbf{hm6k} [\emph{ʰmək}]: Мексика\\
0x2D09 \textbf{sy6k} [\emph{sjək}]: экзистенциальный пункт\\
0x2D0A \textbf{ty6k} [\emph{tjək}]: восьмиугольник\\
0x2D10 \textbf{hk6k} [\emph{ʰkək}]: мексиканский\\
0x2D18 \textbf{hy6k} [\emph{ʰjək}]: слюда\\
0x2D20 \textbf{hp6k} [\emph{ʰpək}]: prolative дело\\
0x2D22 \textbf{kw6k} [\emph{kwək}]: кварк\\
0x2D28 \textbf{hw6k} [\emph{ʰwək}]: убывать\\
0x2D29 \textbf{sw6k} [\emph{swək}]: едок\\
0x2D2A \textbf{tw6k} [\emph{twək}]: транслокация\\
0x2D2E \textbf{rw6k} [\emph{rwək}]: барвинок малый\\
0x2D30 \textbf{hn6k} [\emph{ʰnək}]: Монако\\
0x2D3E \textbf{se} [\emph{se}]: время контекст\\
0x2D42 \textbf{ks6k} [\emph{ksək}]: акустика\\
0x2D44 \textbf{ps6k} [\emph{psək}]: поршень\\
0x2D48 \textbf{hs6k} [\emph{ʰsək}]: поиск двигатель\\
0x2D4A \textbf{ts6k} [\emph{tsək}]: земной\\
0x2D50 \textbf{ht6k} [\emph{ʰtək}]: придирчивый\\
0x2D58 \textbf{hl6k} [\emph{ʰlək}]: лишай\\
0x2D5E \textbf{te} [\emph{te}]: место нахождения дело\\
0x2D60 \textbf{hf6k} [\emph{ʰfək}]: глубина\\
0x2D62 \textbf{kl6k} [\emph{klək}]: чашечка\\
0x2D64 \textbf{pl6k} [\emph{plək}]: республиканец\\
0x2D68 \textbf{hc6k} [\emph{ʰʃək}]: Чешский\\
0x2D69 \textbf{sl6k} [\emph{slək}]: султанат\\
0x2D6A \textbf{tl6k} [\emph{tlək}]: трактор\\
0x2D6C \textbf{fl6k} [\emph{flək}]: ихтиология\\
0x2D6D \textbf{cl6k} [\emph{ʃlək}]: чревный\\
0x2D70 \textbf{hr6k} [\emph{ʰrək}]: рецензент\\
0x2D71 \textbf{bl6k} [\emph{blək}]: балтийский\\
0x2D72 \textbf{gl6k} [\emph{glək}]: гинекология\\
0x2D73 \textbf{dl6k} [\emph{dlək}]: десятина\\
0x2D7A \textbf{xl6k} [\emph{xlək}]: галоген\\
0x2D7E \textbf{le} [\emph{le}]: задерживается императив\\
0x2D88 \textbf{hb6k} [\emph{ʰbək}]: разрез`ать\\
0x2D8A \textbf{tf6k} [\emph{tfək}]: стратификация\\
0x2D90 \textbf{hg6k} [\emph{ʰgək}]: полемический\\
0x2D98 \textbf{hd6k} [\emph{ʰdək}]: педиатрия\\
0x2D9E \textbf{fe} [\emph{fe}]: в конце концов\\
0x2DA0 \textbf{hz6k} [\emph{ʰzək}]: сейсмограф\\
0x2DA2 \textbf{kc6k} [\emph{kʃək}]: портретный\\
0x2DAA \textbf{tc6k} [\emph{tʃək}]: Артишок\\
0x2DBE \textbf{ce} [\emph{ʃe}]: под контекст\\
0x2DC1 \textbf{mr6k} [\emph{mrək}]: Америка\\
0x2DC2 \textbf{kr6k} [\emph{krək}]: скорпион\\
0x2DC4 \textbf{pr6k} [\emph{prək}]: Парагвай\\
0x2DC6 \textbf{nr6k} [\emph{nrək}]: Никарагуа\\
0x2DC8 \textbf{hq6k} [\emph{ʰŋək}]: упираться\\
0x2DC9 \textbf{sr6k} [\emph{srək}]: эстроген\\
0x2DCA \textbf{tr6k} [\emph{trək}]: стартер\\
0x2DCC \textbf{fr6k} [\emph{frək}]: трубка\\
0x2DD0 \textbf{hx6k} [\emph{ʰxək}]: ха\\
0x2DD1 \textbf{br6k} [\emph{brək}]: беатификации\\
0x2DD3 \textbf{dr6k} [\emph{drək}]: гидропонный\\
0x2DDA \textbf{xr6k} [\emph{xrək}]: ха-ха\\
0x2DDE \textbf{re} [\emph{re}]: другой предмет\\
0x2DE2 \textbf{kx6k} [\emph{kxək}]: тектонический\\
0x2DE4 \textbf{px6k} [\emph{pxək}]: панк\\
0x2DEA \textbf{tx6k} [\emph{txək}]: стетоскоп\\
0x2DF3 \textbf{dx6k} [\emph{dxək}]: потряхивание\\
0x2E3E \textbf{be} [\emph{be}]: словесный\\
0x2E5E \textbf{ge} [\emph{ge}]: агент\\
0x2E7E \textbf{de} [\emph{de}]: distributive case, indicates for-each the noun phrase, or per the noun phrase. iterative for loops.\\
0x2E9E \textbf{ze} [\emph{ze}]: Z\\
0x2EBE \textbf{je} [\emph{ʒe}]: agentive case, the agent of a transitive verb.\\
0x2EDE \textbf{ve} [\emph{ve}]: вектор\\
0x2F3E \textbf{qe} [\emph{ŋe}]: star gender, 10th density life form, akin to spirit of a star, Sol, Sirius\\
0x2F5E \textbf{xe} [\emph{xe}]: частота\\
\subsection{ 3 }
0x3001 \textbf{myi7k} [\emph{mji˥k}]: середина июля\\
0x3002 \textbf{kyi7k} [\emph{kji˥k}]: копчик\\
0x3004 \textbf{pyi7k} [\emph{pji˥k}]: пиридин\\
0x3006 \textbf{nyi7k} [\emph{nji˥k}]: незаметный\\
0x3008 \textbf{hmi7k} [\emph{ʰmi˥k}]: манильской\\
0x3009 \textbf{syi7k} [\emph{sji˥k}]: сикх\\
0x300A \textbf{tyi7k} [\emph{tji˥k}]: онтологический\\
0x300B \textbf{lyi7k} [\emph{lji˥k}]: стекольщик\\
0x300D \textbf{cyi7k} [\emph{ʃji˥k}]: облегающий\\
0x300E \textbf{ryi7k} [\emph{rji˥k}]: эректильная\\
0x3010 \textbf{hki7k} [\emph{ʰki˥k}]: хруст\\
0x3012 \textbf{gyi7k} [\emph{gji˥k}]: геодезический\\
0x3018 \textbf{hyi7k} [\emph{ʰji˥k}]: действительный\\
0x3020 \textbf{hpi7k} [\emph{ʰpi˥k}]: португальский\\
0x3021 \textbf{mwi7k} [\emph{mwi˥k}]: амниотический\\
0x3022 \textbf{kwi7k} [\emph{kwi˥k}]: карбоксил\\
0x3024 \textbf{pwi7k} [\emph{pwi˥k}]: дело всей жизни\\
0x3026 \textbf{nwi7k} [\emph{nwi˥k}]: синоптический\\
0x3028 \textbf{hwi7k} [\emph{ʰwi˥k}]: элегия\\
0x3029 \textbf{swi7k} [\emph{swi˥k}]: изометрический\\
0x302A \textbf{twi7k} [\emph{twi˥k}]: стеклярус\\
0x302B \textbf{lwi7k} [\emph{lwi˥k}]: LY\\
0x3030 \textbf{hni7k} [\emph{ʰni˥k}]: негры\\
0x3032 \textbf{gwi7k} [\emph{gwi˥k}]: урод\\
0x303A \textbf{xwi7k} [\emph{xwi˥k}]: халтура\\
0x303E \textbf{mo} [\emph{mo}]: включительно человек\\
0x3042 \textbf{ksi7k} [\emph{ksi˥k}]: перекрестный допрос\\
0x3044 \textbf{psi7k} [\emph{psi˥k}]: апостольский\\
0x3048 \textbf{hsi7k} [\emph{ʰsi˥k}]: двадцать пять\\
0x304A \textbf{tsi7k} [\emph{tsi˥k}]: капающий\\
0x3058 \textbf{hli7k} [\emph{ʰli˥k}]: воспитанным\\
0x305E \textbf{ko} [\emph{ko}]: Предложение хвост\\
0x3060 \textbf{hfi7k} [\emph{ʰfi˥k}]: утомляет\\
0x3062 \textbf{kli7k} [\emph{kli˥k}]: оксихлорид\\
0x3064 \textbf{pli7k} [\emph{pli˥k}]: патология\\
0x3068 \textbf{hci7k} [\emph{ʰʃi˥k}]: очевидец\\
0x3069 \textbf{sli7k} [\emph{sli˥k}]: славословие\\
0x306A \textbf{tli7k} [\emph{tli˥k}]: тик\\
0x306C \textbf{fli7k} [\emph{fli˥k}]: Девочка Мальчик\\
0x306D \textbf{cli7k} [\emph{ʃli˥k}]: кириллица\\
0x3070 \textbf{hri7k} [\emph{ʰri˥k}]: афроамериканец\\
0x3071 \textbf{bli7k} [\emph{bli˥k}]: биофизика\\
0x3072 \textbf{gli7k} [\emph{gli˥k}]: иглу\\
0x3073 \textbf{dli7k} [\emph{dli˥k}]: деликт\\
0x3074 \textbf{zli7k} [\emph{zli˥k}]: интеллектуальный\\
0x307A \textbf{xli7k} [\emph{xli˥k}]: полиглот\\
0x307E \textbf{yo} [\emph{jo}]: союз\\
0x3082 \textbf{kfi7k} [\emph{kfi˥k}]: алифатический\\
0x3084 \textbf{pfi7k} [\emph{pfi˥k}]: коклюш\\
0x3088 \textbf{hbi7k} [\emph{ʰbi˥k}]: обшивать панелью\\
0x308A \textbf{tfi7k} [\emph{tfi˥k}]: превращение в стекло\\
0x3090 \textbf{hgi7k} [\emph{ʰgi˥k}]: двуколка\\
0x3098 \textbf{hdi7k} [\emph{ʰdi˥k}]: приданое\\
0x309E \textbf{po} [\emph{po}]: paucal номер\\
0x30A0 \textbf{hzi7k} [\emph{ʰzi˥k}]: все американские\\
0x30A2 \textbf{kci7k} [\emph{kʃi˥k}]: лаконично\\
0x30A4 \textbf{pci7k} [\emph{pʃi˥k}]: мизинец\\
0x30A8 \textbf{hji7k} [\emph{ʰʒi˥k}]: IC\\
0x30AA \textbf{tci7k} [\emph{tʃi˥k}]: ишиас\\
0x30B0 \textbf{hvi7k} [\emph{ʰvi˥k}]: водные виды спорта\\
0x30BE \textbf{wo} [\emph{wo}]: слово\\
0x30C1 \textbf{mri7k} [\emph{mri˥k}]: мелодраматический\\
0x30C2 \textbf{kri7k} [\emph{kri˥k}]: некритичный\\
0x30C4 \textbf{pri7k} [\emph{pri˥k}]: образно\\
0x30C6 \textbf{nri7k} [\emph{nri˥k}]: желудочек\\
0x30C8 \textbf{hqi7k} [\emph{ʰŋi˥k}]: подписан\\
0x30C9 \textbf{sri7k} [\emph{sri˥k}]: крестец\\
0x30CA \textbf{tri7k} [\emph{tri˥k}]: Терри\\
0x30CD \textbf{cri7k} [\emph{ʃri˥k}]: машина для отжимания белья\\
0x30D0 \textbf{hxi7k} [\emph{ʰxi˥k}]: нищий\\
0x30D1 \textbf{bri7k} [\emph{bri˥k}]: трудолюбиво\\
0x30D2 \textbf{gri7k} [\emph{gri˥k}]: гликоль\\
0x30D3 \textbf{dri7k} [\emph{dri˥k}]: гидростатический\\
0x30DA \textbf{xri7k} [\emph{xri˥k}]: гидроксил\\
0x30DE \textbf{no} [\emph{no}]: негативный квантификатором\\
0x30E2 \textbf{kxi7k} [\emph{kxi˥k}]: стог\\
0x30E4 \textbf{pxi7k} [\emph{pxi˥k}]: паршивый\\
0x30EA \textbf{txi7k} [\emph{txi˥k}]: подстроенный\\
0x3101 \textbf{mya7k} [\emph{mjä˥k}]: Малакка\\
0x3102 \textbf{kya7k} [\emph{kjä˥k}]: ручная работа\\
0x3104 \textbf{pya7k} [\emph{pjä˥k}]: шлепок\\
0x3106 \textbf{nya7k} [\emph{njä˥k}]: нанка\\
0x3108 \textbf{hma7k} [\emph{ʰmä˥k}]: Майкл\\
0x3109 \textbf{sya7k} [\emph{sjä˥k}]: целебный\\
0x310A \textbf{tya7k} [\emph{tjä˥k}]: тапиока\\
0x310B \textbf{lya7k} [\emph{ljä˥k}]: пелагический\\
0x310D \textbf{cya7k} [\emph{ʃjä˥k}]: ночник\\
0x310E \textbf{rya7k} [\emph{rjä˥k}]: арак\\
0x3110 \textbf{hka7k} [\emph{ʰkä˥k}]: замки\\
0x3113 \textbf{dya7k} [\emph{djä˥k}]: идеограмма\\
0x3118 \textbf{hya7k} [\emph{ʰjä˥k}]: Джекоб\\
0x311A \textbf{xya7k} [\emph{xjä˥k}]: гыйба\\
0x3120 \textbf{hpa7k} [\emph{ʰpä˥k}]: забастовки\\
0x3121 \textbf{mwa7k} [\emph{mwä˥k}]: амфора\\
0x3122 \textbf{kwa7k} [\emph{kwä˥k}]: пойло\\
0x3124 \textbf{pwa7k} [\emph{pwä˥k}]: панический\\
0x3126 \textbf{nwa7k} [\emph{nwä˥k}]: фантасмагория\\
0x3128 \textbf{hwa7k} [\emph{ʰwä˥k}]: акварель\\
0x3129 \textbf{swa7k} [\emph{swä˥k}]: хиромантия\\
0x312A \textbf{twa7k} [\emph{twä˥k}]: собиратель древностей\\
0x312B \textbf{lwa7k} [\emph{lwä˥k}]: аликвоту\\
0x312E \textbf{rwa7k} [\emph{rwä˥k}]: покушаться\\
0x3130 \textbf{hna7k} [\emph{ʰnä˥k}]: сын в закон\\
0x3131 \textbf{bwa7k} [\emph{bwä˥k}]: акушерство\\
0x313E \textbf{so} [\emph{so}]: source case, indicates a source noun phrase, from. is used for constructing complex cases or by itself.\\
0x3142 \textbf{ksa7k} [\emph{ksä˥k}]: кукурузный крахмал\\
0x3144 \textbf{psa7k} [\emph{psä˥k}]: баскский\\
0x3148 \textbf{hsa7k} [\emph{ʰsä˥k}]: ударив\\
0x314A \textbf{tsa7k} [\emph{tsä˥k}]: лавировали\\
0x3150 \textbf{hta7k} [\emph{ʰtä˥k}]: ультразвуковая эхография\\
0x3158 \textbf{hla7k} [\emph{ʰlä˥k}]: Калькутта\\
0x315E \textbf{to} [\emph{to}]: пространство контекст\\
0x3160 \textbf{hfa7k} [\emph{ʰfä˥k}]: порка\\
0x3162 \textbf{kla7k} [\emph{klä˥k}]: дисквалификация\\
0x3164 \textbf{pla7k} [\emph{plä˥k}]: могильный\\
0x3168 \textbf{hca7k} [\emph{ʰʃä˥k}]: просыпался\\
0x3169 \textbf{sla7k} [\emph{slä˥k}]: законопослушный\\
0x316A \textbf{tla7k} [\emph{tlä˥k}]: тальк\\
0x316C \textbf{fla7k} [\emph{flä˥k}]: фталевый\\
0x316D \textbf{cla7k} [\emph{ʃlä˥k}]: волчица\\
0x3170 \textbf{hra7k} [\emph{ʰrä˥k}]: архиепископ\\
0x3171 \textbf{bla7k} [\emph{blä˥k}]: баллистический\\
0x3172 \textbf{gla7k} [\emph{glä˥k}]: орленок\\
0x3173 \textbf{dla7k} [\emph{dlä˥k}]: двухпросветный\\
0x317A \textbf{xla7k} [\emph{xlä˥k}]: Малахия\\
0x317E \textbf{lo} [\emph{lo}]: волевых настроение\\
0x3182 \textbf{kfa7k} [\emph{kfä˥k}]: известкование\\
0x3184 \textbf{pfa7k} [\emph{pfä˥k}]: подсемейство\\
0x3188 \textbf{hba7k} [\emph{ʰbä˥k}]: балканский\\
0x318A \textbf{tfa7k} [\emph{tfä˥k}]: вилы\\
0x3190 \textbf{hga7k} [\emph{ʰgä˥k}]: неловкость\\
0x3198 \textbf{hda7k} [\emph{ʰdä˥k}]: приоткрытый\\
0x319E \textbf{fo} [\emph{fo}]: окончательный частица\\
0x31A0 \textbf{hza7k} [\emph{ʰzä˥k}]: прелюбодей\\
0x31A2 \textbf{kca7k} [\emph{kʃä˥k}]: акробатический\\
0x31A4 \textbf{pca7k} [\emph{pʃä˥k}]: магическая фигура\\
0x31A8 \textbf{hja7k} [\emph{ʰʒä˥k}]: дворняжка\\
0x31AA \textbf{tca7k} [\emph{tʃä˥k}]: эрцгерцог\\
0x31B0 \textbf{hva7k} [\emph{ʰvä˥k}]: белый воротничок\\
0x31B1 \textbf{bja7k} [\emph{bʒä˥k}]: батик\\
0x31BE \textbf{co} [\emph{ʃo}]: Социальное контекст\\
0x31C1 \textbf{mra7k} [\emph{mrä˥k}]: Мартиника\\
0x31C2 \textbf{kra7k} [\emph{krä˥k}]: карболка\\
0x31C4 \textbf{pra7k} [\emph{prä˥k}]: Прага\\
0x31C6 \textbf{nra7k} [\emph{nrä˥k}]: дочь в закон\\
0x31C8 \textbf{hqa7k} [\emph{ʰŋä˥k}]: ткач\\
0x31C9 \textbf{sra7k} [\emph{srä˥k}]: манить\\
0x31CA \textbf{tra7k} [\emph{trä˥k}]: Звездный свет\\
0x31CD \textbf{cra7k} [\emph{ʃrä˥k}]: львиный зев\\
0x31D0 \textbf{hxa7k} [\emph{ʰxä˥k}]: horsehoe краба\\
0x31D1 \textbf{bra7k} [\emph{brä˥k}]: барсук\\
0x31D2 \textbf{gra7k} [\emph{grä˥k}]: Лошадиные силы\\
0x31D3 \textbf{dra7k} [\emph{drä˥k}]: испачканный\\
0x31D5 \textbf{jra7k} [\emph{ʒrä˥k}]: матриарх\\
0x31D9 \textbf{qra7k} [\emph{ŋrä˥k}]: гнездо взломщик\\
0x31DA \textbf{xra7k} [\emph{xrä˥k}]: терзать\\
0x31DE \textbf{ro} [\emph{ro}]: государство контекст\\
0x31E2 \textbf{kxa7k} [\emph{kxä˥k}]: иконоборец\\
0x31E4 \textbf{pxa7k} [\emph{pxä˥k}]: раскалывать\\
0x31EA \textbf{txa7k} [\emph{txä˥k}]: вертячка\\
0x31F1 \textbf{bxa7k} [\emph{bxä˥k}]: прицветник\\
0x31F2 \textbf{gxa7k} [\emph{gxä˥k}]: она попка\\
0x31F3 \textbf{dxa7k} [\emph{dxä˥k}]: демагогия\\
0x3208 \textbf{hmu7k} [\emph{ʰmu˥k}]: медулла\\
0x320A \textbf{tyu7k} [\emph{tju˥k}]: тупые под углом\\
0x320B \textbf{lyu7k} [\emph{lju˥k}]: урал Алтайская\\
0x3210 \textbf{hku7k} [\emph{ʰku˥k}]: ворковать\\
0x3218 \textbf{hyu7k} [\emph{ʰju˥k}]: эврика\\
0x3220 \textbf{hpu7k} [\emph{ʰpu˥k}]: эпикуреец\\
0x3222 \textbf{kwu7k} [\emph{kwu˥k}]: Аккув\\
0x3229 \textbf{swu7k} [\emph{swu˥k}]: Суккот\\
0x322A \textbf{twu7k} [\emph{twu˥k}]: библиофил\\
0x3230 \textbf{hnu7k} [\emph{ʰnu˥k}]: лук ноги\\
0x323E \textbf{bo} [\emph{bo}]: вероятность\\
0x3242 \textbf{ksu7k} [\emph{ksu˥k}]: казуистический\\
0x3244 \textbf{psu7k} [\emph{psu˥k}]: воспроизводимый\\
0x3248 \textbf{hsu7k} [\emph{ʰsu˥k}]: саго\\
0x324A \textbf{tsu7k} [\emph{tsu˥k}]: щупальце\\
0x3250 \textbf{htu7k} [\emph{ʰtu˥k}]: отечество\\
0x3258 \textbf{hlu7k} [\emph{ʰlu˥k}]: голубовато-черный\\
0x325E \textbf{go} [\emph{go}]: gnomic aspect, for expressing a general truth.\\
0x3268 \textbf{hcu7k} [\emph{ʰʃu˥k}]: большеберцовая кость\\
0x326A \textbf{tlu7k} [\emph{tlu˥k}]: тавтологический\\
0x3270 \textbf{hru7k} [\emph{ʰru˥k}]: Дракула\\
0x327E \textbf{do} [\emph{do}]: номер\\
0x3290 \textbf{hgu7k} [\emph{ʰgu˥k}]: гу\\
0x3298 \textbf{hdu7k} [\emph{ʰdu˥k}]: совместное обучение\\
0x329E \textbf{zo} [\emph{zo}]: zoic gender, 2nd density life form, learning to fullfill desires, akin to animals, basic computers, basic robots\\
0x32BE \textbf{jo} [\emph{ʒo}]: высокая температура\\
0x32C9 \textbf{sru7k} [\emph{sru˥k}]: сверхзвуковой\\
0x32CA \textbf{tru7k} [\emph{tru˥k}]: через\\
0x32D0 \textbf{hxu7k} [\emph{ʰxu˥k}]: Аввакум\\
0x32E4 \textbf{pxu7k} [\emph{pxu˥k}]: Непропеченный\\
0x3308 \textbf{hme7k} [\emph{ʰme˥k}]: черногория\\
0x330A \textbf{tye7k} [\emph{tje˥k}]: укладчик\\
0x3310 \textbf{hke7k} [\emph{ʰke˥k}]: конвенции\\
0x3318 \textbf{hye7k} [\emph{ʰje˥k}]: немного\\
0x3320 \textbf{hpe7k} [\emph{ʰpe˥k}]: Пекин\\
0x3328 \textbf{hwe7k} [\emph{ʰwe˥k}]: поэтому\\
0x332A \textbf{twe7k} [\emph{twe˥k}]: т квадрат\\
0x3330 \textbf{hne7k} [\emph{ʰne˥k}]: ядра\\
0x333E \textbf{qo} [\emph{ŋo}]: радиоактивность\\
0x3342 \textbf{kse7k} [\emph{kse˥k}]: пустельга\\
0x3344 \textbf{pse7k} [\emph{pse˥k}]: проклинать\\
0x3348 \textbf{hse7k} [\emph{ʰse˥k}]: всевидящий\\
0x334A \textbf{tse7k} [\emph{tse˥k}]: Огниво\\
0x3350 \textbf{hte7k} [\emph{ʰte˥k}]: Tinder\\
0x3358 \textbf{hle7k} [\emph{ʰle˥k}]: Тупик\\
0x335E \textbf{xo} [\emph{xo}]: предварительно сегодняшний\\
0x3364 \textbf{ple7k} [\emph{ple˥k}]: умереть раньше\\
0x3368 \textbf{hce7k} [\emph{ʰʃe˥k}]: шекель\\
0x336A \textbf{tle7k} [\emph{tle˥k}]: электронной коммерции\\
0x336C \textbf{fle7k} [\emph{fle˥k}]: передняя нога или лапа\\
0x3370 \textbf{hre7k} [\emph{ʰre˥k}]: веснушчатый\\
0x3388 \textbf{hbe7k} [\emph{ʰbe˥k}]: брикет\\
0x3390 \textbf{hge7k} [\emph{ʰge˥k}]: экзотерический\\
0x3398 \textbf{hde7k} [\emph{ʰde˥k}]: дебушировать\\
0x33A0 \textbf{hze7k} [\emph{ʰze˥k}]: Иезекииль\\
0x33A2 \textbf{kce7k} [\emph{kʃe˥k}]: переучет\\
0x33AA \textbf{tce7k} [\emph{tʃe˥k}]: ЕЭС\\
0x33C1 \textbf{mre7k} [\emph{mre˥k}]: ун-американских\\
0x33C2 \textbf{kre7k} [\emph{kre˥k}]: четвероногие\\
0x33C4 \textbf{pre7k} [\emph{pre˥k}]: за\\
0x33C9 \textbf{sre7k} [\emph{sre˥k}]: лексикография\\
0x33CA \textbf{tre7k} [\emph{tre˥k}]: зимородок\\
0x33D0 \textbf{hxe7k} [\emph{ʰxe˥k}]: хакеры\\
0x33D1 \textbf{bre7k} [\emph{bre˥k}]: нарушающий закон\\
0x33EA \textbf{txe7k} [\emph{txe˥k}]: нога\\
0x3406 \textbf{nyo7k} [\emph{njo˥k}]: монокль\\
0x3408 \textbf{hmo7k} [\emph{ʰmo˥k}]: Москва\\
0x340A \textbf{tyo7k} [\emph{tjo˥k}]: Токийский\\
0x3410 \textbf{hko7k} [\emph{ʰko˥k}]: иго\\
0x3418 \textbf{hyo7k} [\emph{ʰjo˥k}]: боль в спине\\
0x3420 \textbf{hpo7k} [\emph{ʰpo˥k}]: щенок\\
0x3422 \textbf{kwo7k} [\emph{kwo˥k}]: Кока-Кола\\
0x3428 \textbf{hwo7k} [\emph{ʰwo˥k}]: сурок\\
0x342A \textbf{two7k} [\emph{two˥k}]: Тонганский\\
0x342B \textbf{lwo7k} [\emph{lwo˥k}]: левое крыло\\
0x3430 \textbf{hno7k} [\emph{ʰno˥k}]: солнечный удар\\
0x343E \textbf{m6} [\emph{mə}]: метр\\
0x3442 \textbf{kso7k} [\emph{kso˥k}]: фат\\
0x3444 \textbf{pso7k} [\emph{pso˥k}]: поймать все\\
0x3448 \textbf{hso7k} [\emph{ʰso˥k}]: вяленое мясо\\
0x344A \textbf{tso7k} [\emph{tso˥k}]: контрабандист\\
0x3450 \textbf{hto7k} [\emph{ʰto˥k}]: потребительский фокус\\
0x3458 \textbf{hlo7k} [\emph{ʰlo˥k}]: медальон\\
0x345E \textbf{k6} [\emph{kə}]: назидательная речь соглашение\\
0x3460 \textbf{hfo7k} [\emph{ʰfo˥k}]: ортопедия\\
0x3462 \textbf{klo7k} [\emph{klo˥k}]: кудахтать\\
0x3464 \textbf{plo7k} [\emph{plo˥k}]: палеолитический\\
0x3468 \textbf{hco7k} [\emph{ʰʃo˥k}]: межправительственная океанографическая комиссия\\
0x3469 \textbf{slo7k} [\emph{slo˥k}]: систолический\\
0x346A \textbf{tlo7k} [\emph{tlo˥k}]: раболепствовать\\
0x346C \textbf{flo7k} [\emph{flo˥k}]: офтальмоскоп\\
0x3470 \textbf{hro7k} [\emph{ʰro˥k}]: раскаленный\\
0x3472 \textbf{glo7k} [\emph{glo˥k}]: глаголический\\
0x347E \textbf{y6} [\emph{jə}]: художественный Пол\\
0x3488 \textbf{hbo7k} [\emph{ʰbo˥k}]: блокада\\
0x3490 \textbf{hgo7k} [\emph{ʰgo˥k}]: идти корзина\\
0x3498 \textbf{hdo7k} [\emph{ʰdo˥k}]: демагог\\
0x349E \textbf{p6} [\emph{pə}]: сращивание\\
0x34BE \textbf{w6} [\emph{wə}]: НМУ\\
0x34C1 \textbf{mro7k} [\emph{mro˥k}]: амхарский\\
0x34C2 \textbf{kro7k} [\emph{kro˥k}]: многокрасочный\\
0x34C4 \textbf{pro7k} [\emph{pro˥k}]: волей-неволей\\
0x34C6 \textbf{nro7k} [\emph{nro˥k}]: анаболический\\
0x34C8 \textbf{hqo7k} [\emph{ʰŋo˥k}]: Anglo католическое\\
0x34C9 \textbf{sro7k} [\emph{sro˥k}]: хриплый\\
0x34CA \textbf{tro7k} [\emph{tro˥k}]: стереоскоп\\
0x34D0 \textbf{hxo7k} [\emph{ʰxo˥k}]: тоскующий по дому\\
0x34D3 \textbf{dro7k} [\emph{dro˥k}]: могущий быть оцененным\\
0x34DA \textbf{xro7k} [\emph{xro˥k}]: неисторично\\
0x34DE \textbf{n6} [\emph{nə}]: сенсорный доказательный настроение\\
0x34E2 \textbf{kxo7k} [\emph{kxo˥k}]: калейдоскоп\\
0x34EA \textbf{txo7k} [\emph{txo˥k}]: замочная скважина\\
0x3508 \textbf{hm67k} [\emph{ʰmə˥k}]: семантика\\
0x3510 \textbf{hk67k} [\emph{ʰkə˥k}]: ключица\\
0x3530 \textbf{hn67k} [\emph{ʰnə˥k}]: мениск\\
0x353E \textbf{s6} [\emph{sə}]: crastinal tense, for expressing tomorrow.\\
0x3548 \textbf{hs67k} [\emph{ʰsə˥k}]: красноломкий\\
0x3550 \textbf{ht67k} [\emph{ʰtə˥k}]: самоиндукции\\
0x355E \textbf{t6} [\emph{tə}]: утвердительный настроение\\
0x356A \textbf{tl67k} [\emph{tlə˥k}]: отиатрия\\
0x3570 \textbf{hr67k} [\emph{ʰrə˥k}]: ушная сера\\
0x357E \textbf{l6} [\emph{lə}]: агент вызывать\\
0x3598 \textbf{hd67k} [\emph{ʰdə˥k}]: алкоголик\\
0x359E \textbf{f6} [\emph{fə}]: вставной морфема\\
0x35DE \textbf{r6} [\emph{rə}]: причастие\\
0x363E \textbf{b6} [\emph{bə}]: кими\\
0x367E \textbf{d6} [\emph{də}]: идиофон\\
0x3801 \textbf{myi2k} [\emph{mji˩k}]: асимптотический\\
0x3804 \textbf{pyi2k} [\emph{pji˩k}]: тюряга\\
0x3806 \textbf{nyi2k} [\emph{nji˩k}]: невралгический\\
0x3808 \textbf{hmi2k} [\emph{ʰmi˩k}]: метафизика\\
0x3809 \textbf{syi2k} [\emph{sji˩k}]: семиотический\\
0x380A \textbf{tyi2k} [\emph{tji˩k}]: инталия\\
0x380B \textbf{lyi2k} [\emph{lji˩k}]: Ленинград\\
0x380E \textbf{ryi2k} [\emph{rji˩k}]: сердцевидный\\
0x3810 \textbf{hki2k} [\emph{ʰki˩k}]: оленья кожа\\
0x3818 \textbf{hyi2k} [\emph{ʰji˩k}]: аллегория\\
0x381A \textbf{xyi2k} [\emph{xji˩k}]: деревенщина\\
0x3820 \textbf{hpi2k} [\emph{ʰpi˩k}]: эпический\\
0x3821 \textbf{mwi2k} [\emph{mwi˩k}]: мнемоника\\
0x3822 \textbf{kwi2k} [\emph{kwi˩k}]: калитка хранитель\\
0x3826 \textbf{nwi2k} [\emph{nwi˩k}]: безводный\\
0x3828 \textbf{hwi2k} [\emph{ʰwi˩k}]: волна частиц\\
0x3829 \textbf{swi2k} [\emph{swi˩k}]: мифический\\
0x382A \textbf{twi2k} [\emph{twi˩k}]: тахикардия\\
0x382B \textbf{lwi2k} [\emph{lwi˩k}]: лауриновая\\
0x3830 \textbf{hni2k} [\emph{ʰni˩k}]: транспортируемый\\
0x3842 \textbf{ksi2k} [\emph{ksi˩k}]: шестьдесят шесть\\
0x3844 \textbf{psi2k} [\emph{psi˩k}]: перископ\\
0x3848 \textbf{hsi2k} [\emph{ʰsi˩k}]: пространственно временная\\
0x384A \textbf{tsi2k} [\emph{tsi˩k}]: папоротник-орляк\\
0x3850 \textbf{hti2k} [\emph{ʰti˩k}]: панегирик\\
0x3858 \textbf{hli2k} [\emph{ʰli˩k}]: лиги\\
0x3860 \textbf{hfi2k} [\emph{ʰfi˩k}]: Феникс\\
0x3862 \textbf{kli2k} [\emph{kli˩k}]: клеенка\\
0x3864 \textbf{pli2k} [\emph{pli˩k}]: апологетика\\
0x3868 \textbf{hci2k} [\emph{ʰʃi˩k}]: проницательно\\
0x3869 \textbf{sli2k} [\emph{sli˩k}]: силикат\\
0x386A \textbf{tli2k} [\emph{tli˩k}]: монолог\\
0x386C \textbf{fli2k} [\emph{fli˩k}]: литология\\
0x3870 \textbf{hri2k} [\emph{ʰri˩k}]: двадцать один\\
0x3873 \textbf{dli2k} [\emph{dli˩k}]: георгин\\
0x387A \textbf{xli2k} [\emph{xli˩k}]: штокроза розовая\\
0x3882 \textbf{kfi2k} [\emph{kfi˩k}]: детоксификация\\
0x3888 \textbf{hbi2k} [\emph{ʰbi˩k}]: било\\
0x388A \textbf{tfi2k} [\emph{tfi˩k}]: разврат\\
0x3890 \textbf{hgi2k} [\emph{ʰgi˩k}]: митра\\
0x3898 \textbf{hdi2k} [\emph{ʰdi˩k}]: полукруг\\
0x38A2 \textbf{kci2k} [\emph{kʃi˩k}]: решетчатая постройка\\
0x38A8 \textbf{hji2k} [\emph{ʰʒi˩k}]: наркоман\\
0x38AA \textbf{tci2k} [\emph{tʃi˩k}]: зубочистка\\
0x38B0 \textbf{hvi2k} [\emph{ʰvi˩k}]: ярь-медянка\\
0x38C1 \textbf{mri2k} [\emph{mri˩k}]: маркиз\\
0x38C2 \textbf{kri2k} [\emph{kri˩k}]: каприз\\
0x38C4 \textbf{pri2k} [\emph{pri˩k}]: непонятый\\
0x38C6 \textbf{nri2k} [\emph{nri˩k}]: цилиндрический\\
0x38C9 \textbf{sri2k} [\emph{sri˩k}]: круто\\
0x38CA \textbf{tri2k} [\emph{tri˩k}]: интервалы\\
0x38D0 \textbf{hxi2k} [\emph{ʰxi˩k}]: киркомотыга\\
0x38D1 \textbf{bri2k} [\emph{bri˩k}]: баратрия\\
0x38D2 \textbf{gri2k} [\emph{gri˩k}]: аллергический\\
0x38D3 \textbf{dri2k} [\emph{dri˩k}]: дендрология\\
0x38DA \textbf{xri2k} [\emph{xri˩k}]: Hara Kiri\\
0x38E2 \textbf{kxi2k} [\emph{kxi˩k}]: воздухоплавание\\
0x38E4 \textbf{pxi2k} [\emph{pxi˩k}]: носовое кровотечение\\
0x38EA \textbf{txi2k} [\emph{txi˩k}]: незрелость\\
0x3901 \textbf{mya2k} [\emph{mjä˩k}]: цели\\
0x3902 \textbf{kya2k} [\emph{kjä˩k}]: ханство\\
0x3904 \textbf{pya2k} [\emph{pjä˩k}]: педагогический\\
0x3906 \textbf{nya2k} [\emph{njä˩k}]: выводить из строя\\
0x3908 \textbf{hma2k} [\emph{ʰmä˩k}]: опьяняющий\\
0x3909 \textbf{sya2k} [\emph{sjä˩k}]: асфиксия\\
0x390A \textbf{tya2k} [\emph{tjä˩k}]: автобиографический\\
0x390B \textbf{lya2k} [\emph{ljä˩k}]: вкладыш\\
0x390E \textbf{rya2k} [\emph{rjä˩k}]: Рея\\
0x3910 \textbf{hka2k} [\emph{ʰkä˩k}]: покалывание\\
0x3913 \textbf{dya2k} [\emph{djä˩k}]: два упаковщик\\
0x3918 \textbf{hya2k} [\emph{ʰjä˩k}]: янки\\
0x3920 \textbf{hpa2k} [\emph{ʰpä˩k}]: шагов\\
0x3922 \textbf{kwa2k} [\emph{kwä˩k}]: каркать\\
0x3926 \textbf{nwa2k} [\emph{nwä˩k}]: противоядный\\
0x3928 \textbf{hwa2k} [\emph{ʰwä˩k}]: умник\\
0x3929 \textbf{swa2k} [\emph{swä˩k}]: сфагнум\\
0x392A \textbf{twa2k} [\emph{twä˩k}]: командная работа\\
0x392E \textbf{rwa2k} [\emph{rwä˩k}]: афродизиак\\
0x3930 \textbf{hna2k} [\emph{ʰnä˩k}]: анекдот\\
0x393A \textbf{xwa2k} [\emph{xwä˩k}]: ханука\\
0x3942 \textbf{ksa2k} [\emph{ksä˩k}]: заклинать\\
0x3944 \textbf{psa2k} [\emph{psä˩k}]: полу еженедельно\\
0x3948 \textbf{hsa2k} [\emph{ʰsä˩k}]: пассажиры\\
0x394A \textbf{tsa2k} [\emph{tsä˩k}]: драматург\\
0x3950 \textbf{hta2k} [\emph{ʰtä˩k}]: небесный свод\\
0x3958 \textbf{hla2k} [\emph{ʰlä˩k}]: забавляться\\
0x3960 \textbf{hfa2k} [\emph{ʰfä˩k}]: альпака\\
0x3962 \textbf{kla2k} [\emph{klä˩k}]: отпереть\\
0x3964 \textbf{pla2k} [\emph{plä˩k}]: раз в две недели\\
0x3968 \textbf{hca2k} [\emph{ʰʃä˩k}]: графство\\
0x3969 \textbf{sla2k} [\emph{slä˩k}]: высшее общество\\
0x396A \textbf{tla2k} [\emph{tlä˩k}]: Transact\\
0x396C \textbf{fla2k} [\emph{flä˩k}]: фаллический\\
0x396D \textbf{cla2k} [\emph{ʃlä˩k}]: мыканица\\
0x3970 \textbf{hra2k} [\emph{ʰrä˩k}]: Ирак\\
0x3971 \textbf{bla2k} [\emph{blä˩k}]: эмбриология\\
0x3973 \textbf{dla2k} [\emph{dlä˩k}]: диалекты\\
0x397A \textbf{xla2k} [\emph{xlä˩k}]: арлекин\\
0x3982 \textbf{kfa2k} [\emph{kfä˩k}]: акромегалия\\
0x3984 \textbf{pfa2k} [\emph{pfä˩k}]: профилактика\\
0x3988 \textbf{hba2k} [\emph{ʰbä˩k}]: Бенгалия\\
0x398A \textbf{tfa2k} [\emph{tfä˩k}]: испражняться\\
0x3990 \textbf{hga2k} [\emph{ʰgä˩k}]: оглашение предстоящего бракосочетания\\
0x3998 \textbf{hda2k} [\emph{ʰdä˩k}]: десятилетия\\
0x39A0 \textbf{hza2k} [\emph{ʰzä˩k}]: ацтекский\\
0x39A2 \textbf{kca2k} [\emph{kʃä˩k}]: черный рынок\\
0x39A4 \textbf{pca2k} [\emph{pʃä˩k}]: супердержавы\\
0x39AA \textbf{tca2k} [\emph{tʃä˩k}]: оружейный\\
0x39B0 \textbf{hva2k} [\emph{ʰvä˩k}]: свастика\\
0x39C1 \textbf{mra2k} [\emph{mrä˩k}]: морг\\
0x39C2 \textbf{kra2k} [\emph{krä˩k}]: скалистый\\
0x39C4 \textbf{pra2k} [\emph{prä˩k}]: паркет\\
0x39C6 \textbf{nra2k} [\emph{nrä˩k}]: антарктический\\
0x39C8 \textbf{hqa2k} [\emph{ʰŋä˩k}]: долговязый\\
0x39C9 \textbf{sra2k} [\emph{srä˩k}]: ондатра\\
0x39CA \textbf{tra2k} [\emph{trä˩k}]: зарядка\\
0x39CD \textbf{cra2k} [\emph{ʃrä˩k}]: сорокопут\\
0x39D0 \textbf{hxa2k} [\emph{ʰxä˩k}]: улей\\
0x39D1 \textbf{bra2k} [\emph{brä˩k}]: барак\\
0x39D3 \textbf{dra2k} [\emph{drä˩k}]: лишать членства\\
0x39DA \textbf{xra2k} [\emph{xrä˩k}]: вес тары\\
0x39E2 \textbf{kxa2k} [\emph{kxä˩k}]: лаг\\
0x39E4 \textbf{pxa2k} [\emph{pxä˩k}]: града молодец\\
0x39EA \textbf{txa2k} [\emph{txä˩k}]: дрова\\
0x39F3 \textbf{dxa2k} [\emph{dxä˩k}]: идиопатический\\
0x3A08 \textbf{hmu2k} [\emph{ʰmu˩k}]: нумизматика\\
0x3A18 \textbf{hyu2k} [\emph{ʰju˩k}]: узуфрукт\\
0x3A20 \textbf{hpu2k} [\emph{ʰpu˩k}]: лобковый\\
0x3A2A \textbf{twu2k} [\emph{twu˩k}]: середина недели\\
0x3A30 \textbf{hnu2k} [\emph{ʰnu˩k}]: юго-юго-восток\\
0x3A48 \textbf{hsu2k} [\emph{ʰsu˩k}]: скунс\\
0x3A50 \textbf{htu2k} [\emph{ʰtu˩k}]: чековая книжка\\
0x3A58 \textbf{hlu2k} [\emph{ʰlu˩k}]: Люк\\
0x3A70 \textbf{hru2k} [\emph{ʰru˩k}]: рутил\\
0x3AC2 \textbf{kru2k} [\emph{kru˩k}]: учетверять\\
0x3ACA \textbf{tru2k} [\emph{tru˩k}]: турчанка\\
0x3AEA \textbf{txu2k} [\emph{txu˩k}]: индюков\\
0x3B06 \textbf{nye2k} [\emph{nje˩k}]: сантиграмм\\
0x3B08 \textbf{hme2k} [\emph{ʰme˩k}]: пакля\\
0x3B10 \textbf{hke2k} [\emph{ʰke˩k}]: квакер\\
0x3B18 \textbf{hye2k} [\emph{ʰje˩k}]: один к одному\\
0x3B1A \textbf{xye2k} [\emph{xje˩k}]: герменевтика\\
0x3B20 \textbf{hpe2k} [\emph{ʰpe˩k}]: клевание\\
0x3B28 \textbf{hwe2k} [\emph{ʰwe˩k}]: бурдюк\\
0x3B2A \textbf{twe2k} [\emph{twe˩k}]: эктопический\\
0x3B2E \textbf{rwe2k} [\emph{rwe˩k}]: переделывать\\
0x3B30 \textbf{hne2k} [\emph{ʰne˩k}]: потворствовать\\
0x3B42 \textbf{kse2k} [\emph{kse˩k}]: конспект\\
0x3B44 \textbf{pse2k} [\emph{pse˩k}]: перистальтика\\
0x3B48 \textbf{hse2k} [\emph{ʰse˩k}]: осока\\
0x3B4A \textbf{tse2k} [\emph{tse˩k}]: постельный клоп\\
0x3B50 \textbf{hte2k} [\emph{ʰte˩k}]: рвотный\\
0x3B58 \textbf{hle2k} [\emph{ʰle˩k}]: длинноногий\\
0x3B68 \textbf{hce2k} [\emph{ʰʃe˩k}]: чизкейк\\
0x3B6A \textbf{tle2k} [\emph{tle˩k}]: энтомология\\
0x3B70 \textbf{hre2k} [\emph{ʰre˩k}]: крокет\\
0x3BC1 \textbf{mre2k} [\emph{mre˩k}]: метрический\\
0x3BC2 \textbf{kre2k} [\emph{kre˩k}]: электроскоп\\
0x3BC4 \textbf{pre2k} [\emph{pre˩k}]: грушевидной формы\\
0x3BC6 \textbf{nre2k} [\emph{nre˩k}]: технократия\\
0x3BC9 \textbf{sre2k} [\emph{sre˩k}]: сэр благоговением\\
0x3BCA \textbf{tre2k} [\emph{tre˩k}]: гетерозиготный\\
0x3BD0 \textbf{hxe2k} [\emph{ʰxe˩k}]: заголовок\\
0x3BEA \textbf{txe2k} [\emph{txe˩k}]: коленный рефлекс\\
0x3C06 \textbf{nyo2k} [\emph{njo˩k}]: линолевая\\
0x3C08 \textbf{hmo2k} [\emph{ʰmo˩k}]: импульсов\\
0x3C0E \textbf{ryo2k} [\emph{rjo˩k}]: Ариохом\\
0x3C10 \textbf{hko2k} [\emph{ʰko˩k}]: несвязанный\\
0x3C18 \textbf{hyo2k} [\emph{ʰjo˩k}]: дуться\\
0x3C20 \textbf{hpo2k} [\emph{ʰpo˩k}]: гной\\
0x3C29 \textbf{swo2k} [\emph{swo˩k}]: полуавтоматический\\
0x3C30 \textbf{hno2k} [\emph{ʰno˩k}]: неолит\\
0x3C42 \textbf{kso2k} [\emph{kso˩k}]: Джек в коробке\\
0x3C44 \textbf{pso2k} [\emph{pso˩k}]: дошкольного\\
0x3C48 \textbf{hso2k} [\emph{ʰso˩k}]: запоздалая мысль\\
0x3C4A \textbf{tso2k} [\emph{tso˩k}]: в середине августа\\
0x3C50 \textbf{hto2k} [\emph{ʰto˩k}]: гуж\\
0x3C58 \textbf{hlo2k} [\emph{ʰlo˩k}]: бревна\\
0x3C68 \textbf{hco2k} [\emph{ʰʃo˩k}]: головорез\\
0x3C69 \textbf{slo2k} [\emph{slo˩k}]: мезолит\\
0x3C6A \textbf{tlo2k} [\emph{tlo˩k}]: туберкулез\\
0x3C70 \textbf{hro2k} [\emph{ʰro˩k}]: тряпье\\
0x3C7A \textbf{xlo2k} [\emph{xlo˩k}]: лущеный\\
0x3C98 \textbf{hdo2k} [\emph{ʰdo˩k}]: двухтактный\\
0x3CC1 \textbf{mro2k} [\emph{mro˩k}]: микрон\\
0x3CC2 \textbf{kro2k} [\emph{kro˩k}]: кортеж\\
0x3CC4 \textbf{pro2k} [\emph{pro˩k}]: гиперактивный\\
0x3CC6 \textbf{nro2k} [\emph{nro˩k}]: неврология\\
0x3CC9 \textbf{sro2k} [\emph{sro˩k}]: изотропный\\
0x3CCA \textbf{tro2k} [\emph{tro˩k}]: гетеросексуальный\\
0x3CD0 \textbf{hxo2k} [\emph{ʰxo˩k}]: губная гармоника\\
0x3CE2 \textbf{kxo2k} [\emph{kxo˩k}]: трудоголик\\
0x3CEA \textbf{txo2k} [\emph{txo˩k}]: стафилококк\\
0x3D30 \textbf{hn62k} [\emph{ʰnə˩k}]: кинематика\\
0x3D70 \textbf{hr62k} [\emph{ʰrə˩k}]: сверкать\\
\subsection{ 4 }
0x400F \textbf{myih} [\emph{mjiʰ}]: красивая\\
0x4017 \textbf{kyih} [\emph{kjiʰ}]: потому как\\
0x4027 \textbf{pyih} [\emph{pjiʰ}]: мир\\
0x4037 \textbf{nyih} [\emph{njiʰ}]: неотъемлемое владение\\
0x403E \textbf{mi7} [\emph{mi˥}]: е\\
0x404F \textbf{syih} [\emph{sjiʰ}]: в равной степени\\
0x4057 \textbf{tyih} [\emph{tjiʰ}]: ретроспективный аспект\\
0x405E \textbf{ki7} [\emph{ki˥}]: словесный имя существительное\\
0x405F \textbf{lyih} [\emph{ljiʰ}]: интерес\\
0x4067 \textbf{fyih} [\emph{fjiʰ}]: с плавающей точкой номер\\
0x406F \textbf{cyih} [\emph{ʃjiʰ}]: совещательный настроение\\
0x4077 \textbf{ryih} [\emph{rjiʰ}]: доказательный\\
0x407E \textbf{yi7} [\emph{ji˥}]: язь\\
0x408F \textbf{byih} [\emph{bjiʰ}]: benedictive mood, for expressing blessings.\\
0x4097 \textbf{gyih} [\emph{gjiʰ}]: gnomic mood, for expressing an axiom.\\
0x409E \textbf{pi7} [\emph{pi˥}]: propositive mood, for making proposals or offers, let's\\
0x409F \textbf{dyih} [\emph{djiʰ}]: ненавидеть\\
0x40A7 \textbf{zyih} [\emph{zjiʰ}]: элита\\
0x40AF \textbf{jyih} [\emph{ʒjiʰ}]: не человек\\
0x40B7 \textbf{vyih} [\emph{vjiʰ}]: непрямой речь\\
0x40BE \textbf{wi7} [\emph{wi˥}]: IC\\
0x40CF \textbf{qyih} [\emph{ŋjiʰ}]: оживить Пол\\
0x40D7 \textbf{xyih} [\emph{xjiʰ}]: раскаяние\\
0x40DE \textbf{ni7} [\emph{ni˥}]: соединительный частица\\
0x410F \textbf{mwih} [\emph{mwiʰ}]: permissive mood, for granting permission to do something, potential response to admonitive mood.\\
0x4117 \textbf{kwih} [\emph{kwiʰ}]: каузальный доказательный\\
0x4127 \textbf{pwih} [\emph{pwiʰ}]: ablative case, indicates a source noun phrase, from. is used as a standalone case. space-context source-case\\
0x4137 \textbf{nwih} [\emph{nwiʰ}]: нео\\
0x413E \textbf{si7} [\emph{si˥}]: superlative case, indicates motion above a surface. surface-context way-case.\\
0x414F \textbf{swih} [\emph{swiʰ}]: essive case, indicates a duration of time at which the noun phrase occurs. time-context way-case. \\
0x4157 \textbf{twih} [\emph{twiʰ}]: ilative case, indicates motion `into' some noun phrase. interior-context destination-case.\\
0x415E \textbf{ti7} [\emph{ti˥}]: мимо совершенный\\
0x415F \textbf{lwih} [\emph{lwiʰ}]: delative case, indicates motion away from noun phrase. surface-context source-case. \\
0x4167 \textbf{fwih} [\emph{fwiʰ}]: величественный\\
0x416F \textbf{cwih} [\emph{ʃwiʰ}]: надеяться\\
0x4177 \textbf{rwih} [\emph{rwiʰ}]: precative mood, for requests.\\
0x417E \textbf{li7} [\emph{li˥}]: LY\\
0x418F \textbf{bwih} [\emph{bwiʰ}]: wise gender, 5th density life form, learning wisdom, akin to Jesus Ascended, Latwii, Arcturians, shape shifting aliens\\
0x4197 \textbf{gwih} [\emph{gwiʰ}]: невежественный\\
0x419E \textbf{fi7} [\emph{fi˥}]: FY\\
0x419F \textbf{dwih} [\emph{dwiʰ}]: desiderative mood, to indicate what the speaker wants to happen.\\
0x41AF \textbf{jwih} [\emph{ʒwiʰ}]: jussive mood, for remote commands that are to be conveyed to another recipient.\\
0x41B7 \textbf{vwih} [\emph{vwiʰ}]: evitative case, indicates that the noun phrase is feared or avoided.\\
0x41BE \textbf{ci7} [\emph{ʃi˥}]: светящийся интенсивность\\
0x41D7 \textbf{xwih} [\emph{xwiʰ}]: patient case, the object of a transitive verb.\\
0x4217 \textbf{ksih} [\emph{ksiʰ}]: causal case, indicates the noun phrase is the cause. person-context source-case. \\
0x4227 \textbf{psih} [\emph{psiʰ}]: optatative mood, for expressing wishes or hopes.\\
0x423E \textbf{bi7} [\emph{bi˥}]: ИНГ\\
0x4257 \textbf{tsih} [\emph{tsiʰ}]: Вот\\
0x427E \textbf{di7} [\emph{di˥}]: в ближайщем будущем\\
0x428F \textbf{bzih} [\emph{bziʰ}]: произвольный\\
0x4297 \textbf{gzih} [\emph{gziʰ}]: переходный глагол\\
0x429F \textbf{dzih} [\emph{dziʰ}]: одержимый дело\\
0x42BE \textbf{ji7} [\emph{ʒi˥}]: Изе\\
0x4317 \textbf{klih} [\emph{kliʰ}]: виновный\\
0x4327 \textbf{plih} [\emph{pliʰ}]: частный\\
0x434F \textbf{slih} [\emph{sliʰ}]: спать\\
0x4357 \textbf{tlih} [\emph{tliʰ}]: утвердительный квантификатором\\
0x4367 \textbf{flih} [\emph{fliʰ}]: выведенный доказательный\\
0x436F \textbf{clih} [\emph{ʃliʰ}]: холодно\\
0x438F \textbf{blih} [\emph{bliʰ}]: ограниченный\\
0x4397 \textbf{glih} [\emph{gliʰ}]: переполненный\\
0x439F \textbf{dlih} [\emph{dliʰ}]: определенный статья\\
0x43A7 \textbf{zlih} [\emph{zliʰ}]: irrealis настроение\\
0x43AF \textbf{jlih} [\emph{ʒliʰ}]: антиклимакс\\
0x43B7 \textbf{vlih} [\emph{vliʰ}]: mineral gender, learning about being aware, 1st density life form\\
0x43D7 \textbf{xlih} [\emph{xliʰ}]: несчастность\\
0x4417 \textbf{kfih} [\emph{kfiʰ}]: само уверенность\\
0x4427 \textbf{pfih} [\emph{pfiʰ}]: progressive aspect, the action is in progress, it is happening actively. \\
0x443E \textbf{ma7} [\emph{mä˥}]: количество из вещество\\
0x4457 \textbf{tfih} [\emph{tfiʰ}]: усилитель\\
0x448F \textbf{bvih} [\emph{bviʰ}]: двоичный значение\\
0x4497 \textbf{gvih} [\emph{gviʰ}]: pegative case, indicates motion from a person. social-context source-case.\\
0x449F \textbf{dvih} [\emph{dviʰ}]: конечный глагол\\
0x4517 \textbf{kcih} [\emph{kʃiʰ}]: пожалуйста\\
0x4527 \textbf{pcih} [\emph{pʃiʰ}]: imperative mood, to indicate direct commands.\\
0x453E \textbf{sa7} [\emph{sä˥}]: associative case, indicates the noun phrase is an associated group of the agent. social-context location-case.\\
0x4557 \textbf{tcih} [\emph{tʃiʰ}]: hortative mood, for pleading or begging.\\
0x455E \textbf{ta7} [\emph{tä˥}]: местный направленный дело\\
0x457E \textbf{la7} [\emph{lä˥}]: successive clause, for putting several sentences in a sequence, similar to this happened, then that happened\\
0x458F \textbf{bjih} [\emph{bʒiʰ}]: бывший умерший\\
0x4597 \textbf{gjih} [\emph{gʒiʰ}]: galaxy gender, 11th density life form, akin to spirit of cosmos, Milky Way, Andromeda\\
0x459F \textbf{djih} [\emph{dʒiʰ}]: симпатия\\
0x460F \textbf{mrih} [\emph{mriʰ}]: вежливый регистр\\
0x4617 \textbf{krih} [\emph{kriʰ}]: любопытный\\
0x4627 \textbf{prih} [\emph{priʰ}]: письмо\\
0x4637 \textbf{nrih} [\emph{nriʰ}]: антропный Пол\\
0x464F \textbf{srih} [\emph{sriʰ}]: грустный\\
0x4657 \textbf{trih} [\emph{triʰ}]: Сортировать\\
0x465E \textbf{ga7} [\emph{gä˥}]: графия\\
0x4667 \textbf{frih} [\emph{friʰ}]: знакомые регистр\\
0x466F \textbf{crih} [\emph{ʃriʰ}]: видимый доказательный\\
0x468F \textbf{brih} [\emph{briʰ}]: база\\
0x4697 \textbf{grih} [\emph{griʰ}]: серый\\
0x469F \textbf{drih} [\emph{driʰ}]: степень\\
0x46A7 \textbf{zrih} [\emph{zriʰ}]: формальный регистр\\
0x46AF \textbf{jrih} [\emph{ʒriʰ}]: терпение\\
0x46B7 \textbf{vrih} [\emph{vriʰ}]: исторический напряженный\\
0x46CF \textbf{qrih} [\emph{ŋriʰ}]: мечтательность\\
0x46D7 \textbf{xrih} [\emph{xriʰ}]: застенчивость\\
0x4717 \textbf{kxih} [\emph{kxiʰ}]: невозмутимость\\
0x4727 \textbf{pxih} [\emph{pxiʰ}]: истребование\\
0x4757 \textbf{txih} [\emph{txiʰ}]: excessive case, indicates motion from a state. state-context source-case.\\
0x478F \textbf{bxih} [\emph{bxiʰ}]: universe gender, 12th density life form, akin to spirit of cosmos, Sophia\\
0x4797 \textbf{gxih} [\emph{gxiʰ}]: общий Пол\\
0x479F \textbf{dxih} [\emph{dxiʰ}]: индуктивность\\
0x4801 \textbf{myip} [\emph{mjip}]: милиция\\
0x4802 \textbf{kyip} [\emph{kjip}]: ящерица\\
0x4804 \textbf{pyip} [\emph{pjip}]: апрель\\
0x4806 \textbf{nyip} [\emph{njip}]: в местном масштабе\\
0x4808 \textbf{hmip} [\emph{ʰmip}]: несовершенный аспект\\
0x4809 \textbf{syip} [\emph{sjip}]: в равной степени\\
0x480A \textbf{tyip} [\emph{tjip}]: вести\\
0x480B \textbf{lyip} [\emph{ljip}]: клиническая\\
0x480C \textbf{fyip} [\emph{fjip}]: копировальное устройство\\
0x480D \textbf{cyip} [\emph{ʃjip}]: превосходящие\\
0x480E \textbf{ryip} [\emph{rjip}]: пропорциональный\\
0x480F \textbf{myah} [\emph{mjäʰ}]: temporal case, indicates a specific time. time-context location-case. \\
0x4810 \textbf{hkip} [\emph{ʰkip}]: около\\
0x4811 \textbf{byip} [\emph{bjip}]: библия\\
0x4812 \textbf{gyip} [\emph{gjip}]: золотая рыбка\\
0x4813 \textbf{dyip} [\emph{djip}]: непрямой речь\\
0x4814 \textbf{zyip} [\emph{zjip}]: физиотерапия\\
0x4815 \textbf{jyip} [\emph{ʒjip}]: генотипическая\\
0x4816 \textbf{vyip} [\emph{vjip}]: Ливия\\
0x4817 \textbf{kyah} [\emph{kjäʰ}]: спасибо\\
0x4818 \textbf{hyip} [\emph{ʰjip}]: ретроспективный аспект\\
0x4819 \textbf{qyip} [\emph{ŋjip}]: Покойся с миром\\
0x481A \textbf{xyip} [\emph{xjip}]: непредсказуемый\\
0x4820 \textbf{hpip} [\emph{ʰpip}]: прогрессивный аспект\\
0x4821 \textbf{mwip} [\emph{mwip}]: самозванец\\
0x4822 \textbf{kwip} [\emph{kwip}]: вмешиваться\\
0x4824 \textbf{pwip} [\emph{pwip}]: поэт\\
0x4826 \textbf{nwip} [\emph{nwip}]: неизбежный\\
0x4827 \textbf{pyah} [\emph{pjäʰ}]: love gender, 4th density life form, learning to forgive and love all, akin to Jesus when he was on Earth, Bodhisattva, Hatonn, Gray aliens\\
0x4828 \textbf{hwip} [\emph{ʰwip}]: Воскресенье\\
0x4829 \textbf{swip} [\emph{swip}]: шип\\
0x482A \textbf{twip} [\emph{twip}]: начинательный аспект\\
0x482B \textbf{lwip} [\emph{lwip}]: лимфа\\
0x482C \textbf{fwip} [\emph{fwip}]: расслаблять\\
0x482D \textbf{cwip} [\emph{ʃwip}]: чип\\
0x482E \textbf{rwip} [\emph{rwip}]: регби\\
0x4830 \textbf{hnip} [\emph{ʰnip}]: гномический аспект\\
0x4831 \textbf{bwip} [\emph{bwip}]: БРИК Брач\\
0x4833 \textbf{dwip} [\emph{dwip}]: децибел\\
0x4834 \textbf{zwip} [\emph{zwip}]: ксенофобские\\
0x4837 \textbf{nyah} [\emph{njäʰ}]: нас\\
0x483A \textbf{xwip} [\emph{xwip}]: гипербола\\
0x4842 \textbf{ksip} [\emph{ksip}]: мышьяк\\
0x4844 \textbf{psip} [\emph{psip}]: принцип\\
0x4848 \textbf{hsip} [\emph{ʰsip}]: семь\\
0x484A \textbf{tsip} [\emph{tsip}]: грех\\
0x484F \textbf{syah} [\emph{sjäʰ}]: улыбка\\
0x4850 \textbf{htip} [\emph{ʰtip}]: 10\\
0x4852 \textbf{gzip} [\emph{gzip}]: карданов подвес\\
0x4853 \textbf{dzip} [\emph{dzip}]: в середине декабря\\
0x4857 \textbf{tyah} [\emph{tjäʰ}]: независимый пункт\\
0x4858 \textbf{hlip} [\emph{ʰlip}]: заканчивающий аспект\\
0x485F \textbf{lyah} [\emph{ljäʰ}]: отчуждаемый владение\\
0x4860 \textbf{hfip} [\emph{ʰfip}]: слабый\\
0x4862 \textbf{klip} [\emph{klip}]: ремесленник\\
0x4864 \textbf{plip} [\emph{plip}]: высказал мнение,\\
0x4867 \textbf{fyah} [\emph{fjäʰ}]: фраза маркер\\
0x4868 \textbf{hcip} [\emph{ʰʃip}]: предполагаемый аспект\\
0x4869 \textbf{slip} [\emph{slip}]: скольжение\\
0x486A \textbf{tlip} [\emph{tlip}]: установить\\
0x486C \textbf{flip} [\emph{flip}]: Филиппины\\
0x486D \textbf{clip} [\emph{ʃlip}]: раскаяние\\
0x486F \textbf{cyah} [\emph{ʃjäʰ}]: человек\\
0x4870 \textbf{hrip} [\emph{ʰrip}]: краткое\\
0x4871 \textbf{blip} [\emph{blip}]: миллиард\\
0x4872 \textbf{glip} [\emph{glip}]: устрица\\
0x4873 \textbf{dlip} [\emph{dlip}]: сила воли\\
0x4874 \textbf{zlip} [\emph{zlip}]: сейсмология\\
0x4875 \textbf{jlip} [\emph{ʒlip}]: ингалятор\\
0x4876 \textbf{vlip} [\emph{vlip}]: линкор\\
0x4877 \textbf{ryah} [\emph{rjäʰ}]: смех\\
0x487A \textbf{xlip} [\emph{xlip}]: волейбол\\
0x4882 \textbf{kfip} [\emph{kfip}]: неприемлемо\\
0x4884 \textbf{pfip} [\emph{pfip}]: партитивный\\
0x4888 \textbf{hbip} [\emph{ʰbip}]: пиво\\
0x488A \textbf{tfip} [\emph{tfip}]: continuative аспект\\
0x488F \textbf{byah} [\emph{bjäʰ}]: язык\\
0x4890 \textbf{hgip} [\emph{ʰgip}]: капуста\\
0x4891 \textbf{bvip} [\emph{bvip}]: бобр\\
0x4892 \textbf{gvip} [\emph{gvip}]: незавидный\\
0x4893 \textbf{dvip} [\emph{dvip}]: рабыня\\
0x4897 \textbf{gyah} [\emph{gjäʰ}]: угол\\
0x4898 \textbf{hdip} [\emph{ʰdip}]: задерживается императив\\
0x489F \textbf{dyah} [\emph{djäʰ}]: admonitive mood, a request for permission to do something.\\
0x48A0 \textbf{hzip} [\emph{ʰzip}]: застежка-молния\\
0x48A2 \textbf{kcip} [\emph{kʃip}]: чемпион\\
0x48A4 \textbf{pcip} [\emph{pʃip}]: карандаш\\
0x48A8 \textbf{hjip} [\emph{ʰʒip}]: тропический\\
0x48AA \textbf{tcip} [\emph{tʃip}]: симпатия\\
0x48B0 \textbf{hvip} [\emph{ʰvip}]: Юпитер\\
0x48B1 \textbf{bjip} [\emph{bʒip}]: ежевика\\
0x48B2 \textbf{gjip} [\emph{gʒip}]: джип\\
0x48B3 \textbf{djip} [\emph{dʒip}]: липид\\
0x48B7 \textbf{vyah} [\emph{vjäʰ}]: адвербиальный дело\\
0x48C1 \textbf{mrip} [\emph{mrip}]: метр\\
0x48C2 \textbf{krip} [\emph{krip}]: дискриминация\\
0x48C4 \textbf{prip} [\emph{prip}]: подготовка\\
0x48C6 \textbf{nrip} [\emph{nrip}]: шкаф\\
0x48C8 \textbf{hqip} [\emph{ʰŋip}]: облако вычисления\\
0x48C9 \textbf{srip} [\emph{srip}]: сигарета\\
0x48CA \textbf{trip} [\emph{trip}]: кнут\\
0x48CC \textbf{frip} [\emph{frip}]: керамика\\
0x48CD \textbf{crip} [\emph{ʃrip}]: настоящим\\
0x48CF \textbf{qyah} [\emph{ŋjäʰ}]: озабоченность\\
0x48D0 \textbf{hxip} [\emph{ʰxip}]: антитело\\
0x48D1 \textbf{brip} [\emph{brip}]: бабочка\\
0x48D2 \textbf{grip} [\emph{grip}]: гриль\\
0x48D3 \textbf{drip} [\emph{drip}]: капать\\
0x48D4 \textbf{zrip} [\emph{zrip}]: Сибирь\\
0x48D5 \textbf{jrip} [\emph{ʒrip}]: призывников\\
0x48D6 \textbf{vrip} [\emph{vrip}]: абрикос\\
0x48D7 \textbf{xyah} [\emph{xjäʰ}]: Мировой\\
0x48D9 \textbf{qrip} [\emph{ŋrip}]: крипта\\
0x48DA \textbf{xrip} [\emph{xrip}]: сироп\\
0x48E2 \textbf{kxip} [\emph{kxip}]: клип доска\\
0x48E4 \textbf{pxip} [\emph{pxip}]: торф\\
0x48EA \textbf{txip} [\emph{txip}]: тройной точкой\\
0x48F1 \textbf{bxip} [\emph{bxip}]: битовая карта\\
0x48F2 \textbf{gxip} [\emph{gxip}]: бицепс\\
0x48F3 \textbf{dxip} [\emph{dxip}]: сидячий\\
0x4901 \textbf{myap} [\emph{mjäp}]: сделка\\
0x4902 \textbf{kyap} [\emph{kjäp}]: ужасный\\
0x4904 \textbf{pyap} [\emph{pjäp}]: сдачи\\
0x4906 \textbf{nyap} [\emph{njäp}]: узкий\\
0x4908 \textbf{hmap} [\emph{ʰmäp}]: Роза\\
0x4909 \textbf{syap} [\emph{sjäp}]: распространять\\
0x490A \textbf{tyap} [\emph{tjäp}]: опыт\\
0x490B \textbf{lyap} [\emph{ljäp}]: злоупотребление\\
0x490C \textbf{fyap} [\emph{fjäp}]: амфибия\\
0x490D \textbf{cyap} [\emph{ʃjäp}]: задаваться вопросом\\
0x490E \textbf{ryap} [\emph{rjäp}]: перестраивать\\
0x490F \textbf{mwah} [\emph{mwäʰ}]: comitative case, indicates the noun phrase is accompanying the other actors of the sentence. social-context way-case.\\
0x4910 \textbf{hkap} [\emph{ʰkäp}]: никогда\\
0x4911 \textbf{byap} [\emph{bjäp}]: боб\\
0x4912 \textbf{gyap} [\emph{gjäp}]: Гамбия\\
0x4913 \textbf{dyap} [\emph{djäp}]: второй язык\\
0x4914 \textbf{zyap} [\emph{zjäp}]: Замбия\\
0x4915 \textbf{jyap} [\emph{ʒjäp}]: амперметр\\
0x4916 \textbf{vyap} [\emph{vjäp}]: восемьдесят восемь\\
0x4917 \textbf{kwah} [\emph{kwäʰ}]: желание\\
0x4918 \textbf{hyap} [\emph{ʰjäp}]: король\\
0x4919 \textbf{qyap} [\emph{ŋjäp}]: хэтчбэк\\
0x491A \textbf{xyap} [\emph{xjäp}]: аффидавит\\
0x4920 \textbf{hpap} [\emph{ʰpäp}]: папа\\
0x4921 \textbf{mwap} [\emph{mwäp}]: комикс\\
0x4922 \textbf{kwap} [\emph{kwäp}]: квадрат\\
0x4924 \textbf{pwap} [\emph{pwäp}]: труба\\
0x4926 \textbf{nwap} [\emph{nwäp}]: ноябрь\\
0x4927 \textbf{pwah} [\emph{pwäʰ}]: benefactive case, indicates the noun phrase is the benficiary of the sentence. social-context destination-case.\\
0x4928 \textbf{hwap} [\emph{ʰwäp}]: 8\\
0x4929 \textbf{swap} [\emph{swäp}]: подметать\\
0x492A \textbf{twap} [\emph{twäp}]: прекращение аспект\\
0x492B \textbf{lwap} [\emph{lwäp}]: лютня\\
0x492C \textbf{fwap} [\emph{fwäp}]: фанера\\
0x492D \textbf{cwap} [\emph{ʃwäp}]: шампанское\\
0x492E \textbf{rwap} [\emph{rwäp}]: ежеквартальный\\
0x4930 \textbf{hnap} [\emph{ʰnäp}]: ответ\\
0x4931 \textbf{bwap} [\emph{bwäp}]: галантерея\\
0x4932 \textbf{gwap} [\emph{gwäp}]: Габон\\
0x4933 \textbf{dwap} [\emph{dwäp}]: адаптер\\
0x4937 \textbf{nwah} [\emph{nwäʰ}]: evidential case, for citing a source as evidence. discource-context source-case\\
0x4939 \textbf{qwap} [\emph{ŋwäp}]: искусственное освещение\\
0x493A \textbf{xwap} [\emph{xwäp}]: убегай\\
0x4942 \textbf{ksap} [\emph{ksäp}]: принимать\\
0x4944 \textbf{psap} [\emph{psäp}]: наихудший\\
0x4948 \textbf{hsap} [\emph{ʰsäp}]: Кроме\\
0x494A \textbf{tsap} [\emph{tsäp}]: остров\\
0x494F \textbf{swah} [\emph{swäʰ}]: semelfactive aspect, for expressing a sudden quick event, like a blink or a sneeze.\\
0x4950 \textbf{htap} [\emph{ʰtäp}]: связка\\
0x4951 \textbf{bzap} [\emph{bzäp}]: фагот\\
0x4952 \textbf{gzap} [\emph{gzäp}]: за бортом\\
0x4953 \textbf{dzap} [\emph{dzäp}]: драпировщик\\
0x4957 \textbf{twah} [\emph{twäʰ}]: elative case, indicates motion `out of' some noun phrase. interior-context source-case.\\
0x4958 \textbf{hlap} [\emph{ʰläp}]: semelfactive аспект\\
0x495F \textbf{lwah} [\emph{lwäʰ}]: perlative case, indicates motion through space. space-context way-case. \\
0x4960 \textbf{hfap} [\emph{ʰfäp}]: совершенный аспект\\
0x4962 \textbf{klap} [\emph{kläp}]: ковер\\
0x4964 \textbf{plap} [\emph{pläp}]: стручок\\
0x4968 \textbf{hcap} [\emph{ʰʃäp}]: после\\
0x4969 \textbf{slap} [\emph{släp}]: потерпеть неудачу\\
0x496A \textbf{tlap} [\emph{tläp}]: пол\\
0x496C \textbf{flap} [\emph{fläp}]: полевой шпат\\
0x496D \textbf{clap} [\emph{ʃläp}]: привычный аспект\\
0x496F \textbf{cwah} [\emph{ʃwäʰ}]: после\\
0x4970 \textbf{hrap} [\emph{ʰräp}]: артиллерия\\
0x4971 \textbf{blap} [\emph{bläp}]: бумажник\\
0x4972 \textbf{glap} [\emph{gläp}]: панель инструментов\\
0x4973 \textbf{dlap} [\emph{dläp}]: параллакс\\
0x4974 \textbf{zlap} [\emph{zläp}]: зебра\\
0x4975 \textbf{jlap} [\emph{ʒläp}]: наждачная бумага\\
0x4976 \textbf{vlap} [\emph{vläp}]: амеба\\
0x4977 \textbf{rwah} [\emph{rwäʰ}]: prospective aspect, the event happens after the other things under consideration. \\
0x497A \textbf{xlap} [\emph{xläp}]: пульпа\\
0x4982 \textbf{kfap} [\emph{kfäp}]: октября\\
0x4984 \textbf{pfap} [\emph{pfäp}]: ямкоделатель\\
0x4988 \textbf{hbap} [\emph{ʰbäp}]: прощальный привет\\
0x498A \textbf{tfap} [\emph{tfäp}]: диаграмма\\
0x498F \textbf{bwah} [\emph{bwäʰ}]: злость\\
0x4990 \textbf{hgap} [\emph{ʰgäp}]: корова\\
0x4991 \textbf{bvap} [\emph{bväp}]: Албания\\
0x4992 \textbf{gvap} [\emph{gväp}]: бокал для вина\\
0x4993 \textbf{dvap} [\emph{dväp}]: половина на половину\\
0x4997 \textbf{gwah} [\emph{gwäʰ}]: нервный\\
0x4998 \textbf{hdap} [\emph{ʰdäp}]: 18\\
0x499F \textbf{dwah} [\emph{dwäʰ}]: Да\\
0x49A0 \textbf{hzap} [\emph{ʰzäp}]: аббатство\\
0x49A2 \textbf{kcap} [\emph{kʃäp}]: корабль\\
0x49A4 \textbf{pcap} [\emph{pʃäp}]: в стороне\\
0x49A7 \textbf{zwah} [\emph{zwäʰ}]: амбиция\\
0x49A8 \textbf{hjap} [\emph{ʰʒäp}]: халат\\
0x49AA \textbf{tcap} [\emph{tʃäp}]: побег\\
0x49B0 \textbf{hvap} [\emph{ʰväp}]: свиток бар\\
0x49B1 \textbf{bjap} [\emph{bʒäp}]: спинодержатель\\
0x49B2 \textbf{gjap} [\emph{gʒäp}]: спа\\
0x49B3 \textbf{djap} [\emph{dʒäp}]: алгебра\\
0x49B7 \textbf{vwah} [\emph{vwäʰ}]: движение спуск\\
0x49C1 \textbf{mrap} [\emph{mräp}]: мастер\\
0x49C2 \textbf{krap} [\emph{kräp}]: хлопок\\
0x49C4 \textbf{prap} [\emph{präp}]: смещение\\
0x49C6 \textbf{nrap} [\emph{nräp}]: экономка\\
0x49C8 \textbf{hqap} [\emph{ʰŋäp}]: настоящее время причастие\\
0x49C9 \textbf{srap} [\emph{sräp}]: овощной\\
0x49CA \textbf{trap} [\emph{träp}]: либеральная\\
0x49CC \textbf{frap} [\emph{fräp}]: федеральный\\
0x49CD \textbf{crap} [\emph{ʃräp}]: ловушка\\
0x49CF \textbf{qwah} [\emph{ŋwäʰ}]: неамбициозны\\
0x49D0 \textbf{hxap} [\emph{ʰxäp}]: рюкзак\\
0x49D1 \textbf{brap} [\emph{bräp}]: лом\\
0x49D2 \textbf{grap} [\emph{gräp}]: захват\\
0x49D3 \textbf{drap} [\emph{dräp}]: завлечь\\
0x49D4 \textbf{zrap} [\emph{zräp}]: песчаный карьер\\
0x49D5 \textbf{jrap} [\emph{ʒräp}]: отвертка\\
0x49D6 \textbf{vrap} [\emph{vräp}]: загар\\
0x49D7 \textbf{xwah} [\emph{xwäʰ}]: Вау\\
0x49D9 \textbf{qrap} [\emph{ŋräp}]: небоскреб\\
0x49DA \textbf{xrap} [\emph{xräp}]: носовой платок\\
0x49E2 \textbf{kxap} [\emph{kxäp}]: поднятый\\
0x49E4 \textbf{pxap} [\emph{pxäp}]: тополь\\
0x49EA \textbf{txap} [\emph{txäp}]: положение дел бар\\
0x49F1 \textbf{bxap} [\emph{bxäp}]: перепел\\
0x49F2 \textbf{gxap} [\emph{gxäp}]: гребешок\\
0x49F3 \textbf{dxap} [\emph{dxäp}]: компас\\
0x4A01 \textbf{myup} [\emph{mjup}]: саженцы\\
0x4A02 \textbf{kyup} [\emph{kjup}]: Куба\\
0x4A04 \textbf{pyup} [\emph{pjup}]: папуасов\\
0x4A06 \textbf{nyup} [\emph{njup}]: высиживать\\
0x4A08 \textbf{hmup} [\emph{ʰmup}]: шапка\\
0x4A09 \textbf{syup} [\emph{sjup}]: супер\\
0x4A0A \textbf{tyup} [\emph{tjup}]: пригород\\
0x4A0B \textbf{lyup} [\emph{ljup}]: шлюп\\
0x4A0D \textbf{cyup} [\emph{ʃjup}]: ворчливый\\
0x4A0E \textbf{ryup} [\emph{rjup}]: повторное использование\\
0x4A10 \textbf{hkup} [\emph{ʰkup}]: копия\\
0x4A11 \textbf{byup} [\emph{bjup}]: гудок\\
0x4A17 \textbf{ksah} [\emph{ksäʰ}]: близость\\
0x4A18 \textbf{hyup} [\emph{ʰjup}]: Европа\\
0x4A1A \textbf{xyup} [\emph{xjup}]: подбоченясь\\
0x4A20 \textbf{hpup} [\emph{ʰpup}]: хлеб\\
0x4A22 \textbf{kwup} [\emph{kwup}]: высокий уровень\\
0x4A26 \textbf{nwup} [\emph{nwup}]: одноклеточный\\
0x4A27 \textbf{psah} [\emph{psäʰ}]: отчаяние\\
0x4A28 \textbf{hwup} [\emph{ʰwup}]: неоценимый\\
0x4A29 \textbf{swup} [\emph{swup}]: кушетка\\
0x4A2A \textbf{twup} [\emph{twup}]: контрапункт\\
0x4A2D \textbf{cwup} [\emph{ʃwup}]: шампунь\\
0x4A30 \textbf{hnup} [\emph{ʰnup}]: раб\\
0x4A3A \textbf{xwup} [\emph{xwup}]: бондарь\\
0x4A42 \textbf{ksup} [\emph{ksup}]: косинус\\
0x4A44 \textbf{psup} [\emph{psup}]: губная помада\\
0x4A48 \textbf{hsup} [\emph{ʰsup}]: содержащий окаменелости Пол\\
0x4A4A \textbf{tsup} [\emph{tsup}]: раздробить\\
0x4A50 \textbf{htup} [\emph{ʰtup}]: долг\\
0x4A57 \textbf{tsah} [\emph{tsäʰ}]: демонстративный\\
0x4A58 \textbf{hlup} [\emph{ʰlup}]: альбом\\
0x4A60 \textbf{hfup} [\emph{ʰfup}]: буря\\
0x4A62 \textbf{klup} [\emph{klup}]: клуб\\
0x4A64 \textbf{plup} [\emph{plup}]: тонкой кожурой\\
0x4A68 \textbf{hcup} [\emph{ʰʃup}]: зал\\
0x4A69 \textbf{slup} [\emph{slup}]: лесок\\
0x4A6A \textbf{tlup} [\emph{tlup}]: репа\\
0x4A6C \textbf{flup} [\emph{flup}]: более полное\\
0x4A6D \textbf{clup} [\emph{ʃlup}]: шелковица\\
0x4A70 \textbf{hrup} [\emph{ʰrup}]: коррупция\\
0x4A71 \textbf{blup} [\emph{blup}]: бесплодие\\
0x4A7A \textbf{xlup} [\emph{xlup}]: градина\\
0x4A82 \textbf{kfup} [\emph{kfup}]: куб\\
0x4A88 \textbf{hbup} [\emph{ʰbup}]: бамбук\\
0x4A8A \textbf{tfup} [\emph{tfup}]: многоразовый\\
0x4A8F \textbf{bzah} [\emph{bzäʰ}]: шутить\\
0x4A90 \textbf{hgup} [\emph{ʰgup}]: купол\\
0x4A98 \textbf{hdup} [\emph{ʰdup}]: падение\\
0x4AA0 \textbf{hzup} [\emph{ʰzup}]: Зимбабве\\
0x4AA2 \textbf{kcup} [\emph{kʃup}]: купон\\
0x4AA4 \textbf{pcup} [\emph{pʃup}]: зловещий\\
0x4AA8 \textbf{hjup} [\emph{ʰʒup}]: песчаная буря\\
0x4AAA \textbf{tcup} [\emph{tʃup}]: влажный\\
0x4AB0 \textbf{hvup} [\emph{ʰvup}]: вакуоль\\
0x4AC1 \textbf{mrup} [\emph{mrup}]: микроб\\
0x4AC2 \textbf{krup} [\emph{krup}]: Акупрессура\\
0x4AC4 \textbf{prup} [\emph{prup}]: Перу\\
0x4AC6 \textbf{nrup} [\emph{nrup}]: вход\\
0x4AC8 \textbf{hqup} [\emph{ʰŋup}]: тупо\\
0x4AC9 \textbf{srup} [\emph{srup}]: суп\\
0x4ACA \textbf{trup} [\emph{trup}]: атрибут\\
0x4ACC \textbf{frup} [\emph{frup}]: опорос\\
0x4ACD \textbf{crup} [\emph{ʃrup}]: перекрытие\\
0x4AD0 \textbf{hxup} [\emph{ʰxup}]: головастик\\
0x4AD1 \textbf{brup} [\emph{brup}]: розовощекий\\
0x4AD2 \textbf{grup} [\emph{grup}]: морской окунь\\
0x4AD3 \textbf{drup} [\emph{drup}]: индо Европейский\\
0x4AD4 \textbf{zrup} [\emph{zrup}]: узурпировать\\
0x4ADA \textbf{xrup} [\emph{xrup}]: иврит\\
0x4AE2 \textbf{kxup} [\emph{kxup}]: скальпель\\
0x4AE4 \textbf{pxup} [\emph{pxup}]: сутенер\\
0x4AEA \textbf{txup} [\emph{txup}]: турмалин\\
0x4AF3 \textbf{dxup} [\emph{dxup}]: дурная слава\\
0x4B01 \textbf{myep} [\emph{mjep}]: Намибия\\
0x4B02 \textbf{kyep} [\emph{kjep}]: сороконожка\\
0x4B04 \textbf{pyep} [\emph{pjep}]: детский манеж\\
0x4B06 \textbf{nyep} [\emph{njep}]: нюхательный табак\\
0x4B08 \textbf{hmep} [\emph{ʰmep}]: momentane аспект\\
0x4B09 \textbf{syep} [\emph{sjep}]: скрепок\\
0x4B0A \textbf{tyep} [\emph{tjep}]: штатив\\
0x4B0B \textbf{lyep} [\emph{ljep}]: Либерия\\
0x4B0D \textbf{cyep} [\emph{ʃjep}]: капризный ребенок\\
0x4B0E \textbf{ryep} [\emph{rjep}]: клюква\\
0x4B10 \textbf{hkep} [\emph{ʰkep}]: медь\\
0x4B11 \textbf{byep} [\emph{bjep}]: ЕЦБ\\
0x4B13 \textbf{dyep} [\emph{djep}]: гантель\\
0x4B17 \textbf{klah} [\emph{kläʰ}]: там\\
0x4B18 \textbf{hyep} [\emph{ʰjep}]: веко\\
0x4B1A \textbf{xyep} [\emph{xjep}]: гессенский\\
0x4B20 \textbf{hpep} [\emph{ʰpep}]: яблоко\\
0x4B21 \textbf{mwep} [\emph{mwep}]: лживый\\
0x4B22 \textbf{kwep} [\emph{kwep}]: спектроскоп\\
0x4B24 \textbf{pwep} [\emph{pwep}]: пипетка\\
0x4B26 \textbf{nwep} [\emph{nwep}]: анаэробный\\
0x4B27 \textbf{plah} [\emph{pläʰ}]: люблю\\
0x4B28 \textbf{hwep} [\emph{ʰwep}]: Web\\
0x4B29 \textbf{swep} [\emph{swep}]: соболь\\
0x4B2A \textbf{twep} [\emph{twep}]: переплачивать\\
0x4B2B \textbf{lwep} [\emph{lwep}]: разборчиво\\
0x4B2D \textbf{cwep} [\emph{ʃwep}]: зарплаты\\
0x4B2E \textbf{rwep} [\emph{rwep}]: сорок восемь\\
0x4B30 \textbf{hnep} [\emph{ʰnep}]: борьба\\
0x4B33 \textbf{dwep} [\emph{dwep}]: диспепсия\\
0x4B42 \textbf{ksep} [\emph{ksep}]: рецептор\\
0x4B44 \textbf{psep} [\emph{psep}]: проницаемый\\
0x4B48 \textbf{hsep} [\emph{ʰsep}]: сентябрь\\
0x4B4A \textbf{tsep} [\emph{tsep}]: кузнечик\\
0x4B4F \textbf{slah} [\emph{släʰ}]: беда\\
0x4B50 \textbf{htep} [\emph{ʰtep}]: поклонение\\
0x4B57 \textbf{tlah} [\emph{tläʰ}]: факультативный квантификатором\\
0x4B58 \textbf{hlep} [\emph{ʰlep}]: целенаправленный аспект\\
0x4B60 \textbf{hfep} [\emph{ʰfep}]: многократный аспект\\
0x4B62 \textbf{klep} [\emph{klep}]: ламинария\\
0x4B64 \textbf{plep} [\emph{plep}]: палитра\\
0x4B67 \textbf{flah} [\emph{fläʰ}]: контрфактическое условный\\
0x4B68 \textbf{hcep} [\emph{ʰʃep}]: клетка\\
0x4B69 \textbf{slep} [\emph{slep}]: сатир\\
0x4B6A \textbf{tlep} [\emph{tlep}]: шаблон\\
0x4B6D \textbf{clep} [\emph{ʃlep}]: дождь со снегом\\
0x4B6F \textbf{clah} [\emph{ʃläʰ}]: Другие\\
0x4B70 \textbf{hrep} [\emph{ʰrep}]: рецепт\\
0x4B71 \textbf{blep} [\emph{blep}]: балеарский\\
0x4B73 \textbf{dlep} [\emph{dlep}]: дверной звонок\\
0x4B75 \textbf{jlep} [\emph{ʒlep}]: коленопреклонение\\
0x4B7A \textbf{xlep} [\emph{xlep}]: обои\\
0x4B82 \textbf{kfep} [\emph{kfep}]: отлично справляется\\
0x4B84 \textbf{pfep} [\emph{pfep}]: гипнотерапия\\
0x4B88 \textbf{hbep} [\emph{ʰbep}]: бейсбол\\
0x4B8A \textbf{tfep} [\emph{tfep}]: чайная чашка\\
0x4B8F \textbf{blah} [\emph{bläʰ}]: вес\\
0x4B90 \textbf{hgep} [\emph{ʰgep}]: Постельное белье\\
0x4B97 \textbf{glah} [\emph{gläʰ}]: здоровье\\
0x4B98 \textbf{hdep} [\emph{ʰdep}]: delimitative аспект\\
0x4B9F \textbf{dlah} [\emph{dläʰ}]: наслаждаться\\
0x4BA0 \textbf{hzep} [\emph{ʰzep}]: затылок\\
0x4BA2 \textbf{kcep} [\emph{kʃep}]: кетчуп\\
0x4BA4 \textbf{pcep} [\emph{pʃep}]: легкое\\
0x4BA7 \textbf{zlah} [\emph{zläʰ}]: верность\\
0x4BA8 \textbf{hjep} [\emph{ʰʒep}]: Пол\\
0x4BAA \textbf{tcep} [\emph{tʃep}]: мел\\
0x4BB0 \textbf{hvep} [\emph{ʰvep}]: каверны\\
0x4BB7 \textbf{vlah} [\emph{vläʰ}]: напряжение\\
0x4BC1 \textbf{mrep} [\emph{mrep}]: ампер\\
0x4BC2 \textbf{krep} [\emph{krep}]: шифрование\\
0x4BC4 \textbf{prep} [\emph{prep}]: эпидермис\\
0x4BC6 \textbf{nrep} [\emph{nrep}]: северо-Запад\\
0x4BC8 \textbf{hqep} [\emph{ʰŋep}]: леденец\\
0x4BC9 \textbf{srep} [\emph{srep}]: киберпространство\\
0x4BCA \textbf{trep} [\emph{trep}]: клубника\\
0x4BCC \textbf{frep} [\emph{frep}]: щетка для волос\\
0x4BCD \textbf{crep} [\emph{ʃrep}]: гамбургер\\
0x4BD0 \textbf{hxep} [\emph{ʰxep}]: скейтборд\\
0x4BD1 \textbf{brep} [\emph{brep}]: лютик\\
0x4BD2 \textbf{grep} [\emph{grep}]: гонорея\\
0x4BD3 \textbf{drep} [\emph{drep}]: энтропия\\
0x4BD5 \textbf{jrep} [\emph{ʒrep}]: неоседланный\\
0x4BD7 \textbf{xlah} [\emph{xläʰ}]: сатир\\
0x4BDA \textbf{xrep} [\emph{xrep}]: пажитник\\
0x4BE2 \textbf{kxep} [\emph{kxep}]: холестерин\\
0x4BE4 \textbf{pxep} [\emph{pxep}]: панели\\
0x4BEA \textbf{txep} [\emph{txep}]: утроить\\
0x4BF1 \textbf{bxep} [\emph{bxep}]: латыш\\
0x4BF2 \textbf{gxep} [\emph{gxep}]: Гринпис\\
0x4BF3 \textbf{dxep} [\emph{dxep}]: потрошить\\
0x4C01 \textbf{myop} [\emph{mjop}]: храмы\\
0x4C02 \textbf{kyop} [\emph{kjop}]: крокодил\\
0x4C04 \textbf{pyop} [\emph{pjop}]: плей-офф\\
0x4C06 \textbf{nyop} [\emph{njop}]: ниобий\\
0x4C08 \textbf{hmop} [\emph{ʰmop}]: моль\\
0x4C09 \textbf{syop} [\emph{sjop}]: сотни\\
0x4C0A \textbf{tyop} [\emph{tjop}]: осьминог\\
0x4C0B \textbf{lyop} [\emph{ljop}]: Олимпийские игры\\
0x4C0C \textbf{fyop} [\emph{fjop}]: орфоэпия\\
0x4C0D \textbf{cyop} [\emph{ʃjop}]: хобби\\
0x4C0E \textbf{ryop} [\emph{rjop}]: робот\\
0x4C10 \textbf{hkop} [\emph{ʰkop}]: кружка\\
0x4C11 \textbf{byop} [\emph{bjop}]: аэробный\\
0x4C12 \textbf{gyop} [\emph{gjop}]: спасательный буй или круг\\
0x4C13 \textbf{dyop} [\emph{djop}]: радиоуглерод\\
0x4C17 \textbf{kfah} [\emph{kfäʰ}]: жертва\\
0x4C18 \textbf{hyop} [\emph{ʰjop}]: яхта\\
0x4C1A \textbf{xyop} [\emph{xjop}]: колибри\\
0x4C20 \textbf{hpop} [\emph{ʰpop}]: опера\\
0x4C21 \textbf{mwop} [\emph{mwop}]: Моавитянин\\
0x4C22 \textbf{kwop} [\emph{kwop}]: в гору\\
0x4C24 \textbf{pwop} [\emph{pwop}]: парабола\\
0x4C26 \textbf{nwop} [\emph{nwop}]: ягодицы\\
0x4C27 \textbf{pfah} [\emph{pfäʰ}]: выгодный\\
0x4C28 \textbf{hwop} [\emph{ʰwop}]: ватт\\
0x4C29 \textbf{swop} [\emph{swop}]: секундомер\\
0x4C2A \textbf{twop} [\emph{twop}]: тромбон\\
0x4C2B \textbf{lwop} [\emph{lwop}]: логотипы\\
0x4C2C \textbf{fwop} [\emph{fwop}]: светобоязнь\\
0x4C2E \textbf{rwop} [\emph{rwop}]: коротковолновый\\
0x4C30 \textbf{hnop} [\emph{ʰnop}]: монополия\\
0x4C31 \textbf{bwop} [\emph{bwop}]: Bow Wow\\
0x4C34 \textbf{zwop} [\emph{zwop}]: изотоп\\
0x4C3A \textbf{xwop} [\emph{xwop}]: гемоглобин\\
0x4C42 \textbf{ksop} [\emph{ksop}]: косметический\\
0x4C44 \textbf{psop} [\emph{psop}]: паб\\
0x4C48 \textbf{hsop} [\emph{ʰsop}]: успешно\\
0x4C4A \textbf{tsop} [\emph{tsop}]: движение в гору\\
0x4C50 \textbf{htop} [\emph{ʰtop}]: проклятие\\
0x4C57 \textbf{tfah} [\emph{tfäʰ}]: надлежащий имя существительное\\
0x4C58 \textbf{hlop} [\emph{ʰlop}]: телескоп\\
0x4C5E \textbf{ke7} [\emph{ke˥}]: ан\\
0x4C60 \textbf{hfop} [\emph{ʰfop}]: футбол\\
0x4C62 \textbf{klop} [\emph{klop}]: калория\\
0x4C64 \textbf{plop} [\emph{plop}]: палладий\\
0x4C68 \textbf{hcop} [\emph{ʰʃop}]: плита\\
0x4C69 \textbf{slop} [\emph{slop}]: тапочка\\
0x4C6A \textbf{tlop} [\emph{tlop}]: скатерть\\
0x4C6C \textbf{flop} [\emph{flop}]: осциллограф\\
0x4C6D \textbf{clop} [\emph{ʃlop}]: шалот\\
0x4C70 \textbf{hrop} [\emph{ʰrop}]: веревка\\
0x4C71 \textbf{blop} [\emph{blop}]: колба\\
0x4C72 \textbf{glop} [\emph{glop}]: Ванька встанька\\
0x4C73 \textbf{dlop} [\emph{dlop}]: волнистый\\
0x4C74 \textbf{zlop} [\emph{zlop}]: шлагбаумом\\
0x4C7A \textbf{xlop} [\emph{xlop}]: гомеопатия\\
0x4C84 \textbf{pfop} [\emph{pfop}]: поп вверх\\
0x4C88 \textbf{hbop} [\emph{ʰbop}]: взломать\\
0x4C8A \textbf{tfop} [\emph{tfop}]: зубная щетка\\
0x4C8F \textbf{bvah} [\emph{bväʰ}]: abessive case, indicating the absence of the noun phrase referred to object. state-context negatory-quantifier\\
0x4C90 \textbf{hgop} [\emph{ʰgop}]: гильдия\\
0x4C97 \textbf{gvah} [\emph{gväʰ}]: lative case, indicates spatial motion towards the noun phrase. space-context destination-case. \\
0x4C98 \textbf{hdop} [\emph{ʰdop}]: движение вверх по течению\\
0x4C9E \textbf{pe7} [\emph{pe˥}]: за\\
0x4C9F \textbf{dvah} [\emph{dväʰ}]: superessive case, indicates location on a surface. surface-context location-case.\\
0x4CA0 \textbf{hzop} [\emph{ʰzop}]: Мозамбик\\
0x4CA2 \textbf{kcop} [\emph{kʃop}]: курятник\\
0x4CA4 \textbf{pcop} [\emph{pʃop}]: насос\\
0x4CA8 \textbf{hjop} [\emph{ʰʒop}]: гироскоп\\
0x4CAA \textbf{tcop} [\emph{tʃop}]: губа\\
0x4CB0 \textbf{hvop} [\emph{ʰvop}]: проспать\\
0x4CC1 \textbf{mrop} [\emph{mrop}]: СВЧ\\
0x4CC2 \textbf{krop} [\emph{krop}]: координировать\\
0x4CC4 \textbf{prop} [\emph{prop}]: плодородный\\
0x4CC6 \textbf{nrop} [\emph{nrop}]: начисление заработной платы\\
0x4CC8 \textbf{hqop} [\emph{ʰŋop}]: Сингапур\\
0x4CC9 \textbf{srop} [\emph{srop}]: смартфон\\
0x4CCA \textbf{trop} [\emph{trop}]: зонд\\
0x4CCC \textbf{frop} [\emph{frop}]: хорек\\
0x4CCD \textbf{crop} [\emph{ʃrop}]: водопад\\
0x4CD0 \textbf{hxop} [\emph{ʰxop}]: Эфиопия\\
0x4CD1 \textbf{brop} [\emph{brop}]: бор\\
0x4CD2 \textbf{grop} [\emph{grop}]: моллюск из семейства брюхоногих\\
0x4CD3 \textbf{drop} [\emph{drop}]: углеводород\\
0x4CD4 \textbf{zrop} [\emph{zrop}]: северо-восточный\\
0x4CD9 \textbf{qrop} [\emph{ŋrop}]: агорафобия\\
0x4CDA \textbf{xrop} [\emph{xrop}]: гармоника\\
0x4CE2 \textbf{kxop} [\emph{kxop}]: крайняя плоть\\
0x4CE4 \textbf{pxop} [\emph{pxop}]: опал\\
0x4CEA \textbf{txop} [\emph{txop}]: шахматная доска\\
0x4CF2 \textbf{gxop} [\emph{gxop}]: коровье дерево\\
0x4CF3 \textbf{dxop} [\emph{dxop}]: пыльная одежда\\
0x4D08 \textbf{hm6p} [\emph{ʰməp}]: материнская плата\\
0x4D0A \textbf{ty6p} [\emph{tjəp}]: в середине апреля\\
0x4D10 \textbf{hk6p} [\emph{ʰkəp}]: Колумбия\\
0x4D17 \textbf{kcah} [\emph{kʃäʰ}]: голодный\\
0x4D20 \textbf{hp6p} [\emph{ʰpəp}]: полутень\\
0x4D26 \textbf{nw6p} [\emph{nwəp}]: Antwerp\\
0x4D27 \textbf{pcah} [\emph{pʃäʰ}]: также\\
0x4D30 \textbf{hn6p} [\emph{ʰnəp}]: Непал\\
0x4D42 \textbf{ks6p} [\emph{ksəp}]: тигель\\
0x4D44 \textbf{ps6p} [\emph{psəp}]: несмешивающийся\\
0x4D48 \textbf{hs6p} [\emph{ʰsəp}]: психотерапия\\
0x4D4A \textbf{ts6p} [\emph{tsəp}]: горячая точка\\
0x4D50 \textbf{ht6p} [\emph{ʰtəp}]: шпиль\\
0x4D58 \textbf{hl6p} [\emph{ʰləp}]: Люксембург\\
0x4D62 \textbf{kl6p} [\emph{kləp}]: Колумбус\\
0x4D64 \textbf{pl6p} [\emph{pləp}]: лепестки\\
0x4D68 \textbf{hc6p} [\emph{ʰʃəp}]: проказливый\\
0x4D69 \textbf{sl6p} [\emph{sləp}]: несбалансированность\\
0x4D6A \textbf{tl6p} [\emph{tləp}]: топология\\
0x4D70 \textbf{hr6p} [\emph{ʰrəp}]: ромб\\
0x4D71 \textbf{bl6p} [\emph{bləp}]: кими\\
0x4D7A \textbf{xl6p} [\emph{xləp}]: заячья губа\\
0x4D7E \textbf{le7} [\emph{le˥}]: еще - если\\
0x4D88 \textbf{hb6p} [\emph{ʰbəp}]: бампер\\
0x4D90 \textbf{hg6p} [\emph{ʰgəp}]: галопом\\
0x4D97 \textbf{gjah} [\emph{gʒäʰ}]: лестный\\
0x4D98 \textbf{hd6p} [\emph{ʰdəp}]: подгузник\\
0x4DC1 \textbf{mr6p} [\emph{mrəp}]: микрочип\\
0x4DC2 \textbf{kr6p} [\emph{krəp}]: карибский\\
0x4DC4 \textbf{pr6p} [\emph{prəp}]: груши\\
0x4DC6 \textbf{nr6p} [\emph{nrəp}]: ономатопея\\
0x4DC8 \textbf{hq6p} [\emph{ʰŋəp}]: изобара\\
0x4DC9 \textbf{sr6p} [\emph{srəp}]: синапс\\
0x4DCA \textbf{tr6p} [\emph{trəp}]: пястная\\
0x4DD0 \textbf{hx6p} [\emph{ʰxəp}]: травяной\\
0x4DD1 \textbf{br6p} [\emph{brəp}]: берберский\\
0x4DE2 \textbf{kx6p} [\emph{kxəp}]: Каспийская\\
0x4DE4 \textbf{px6p} [\emph{pxəp}]: виночерпием\\
0x4DEA \textbf{tx6p} [\emph{txəp}]: остеопатия\\
0x4E0F \textbf{mrah} [\emph{mräʰ}]: Понимаю\\
0x4E17 \textbf{krah} [\emph{kräʰ}]: гордость\\
0x4E27 \textbf{prah} [\emph{präʰ}]: параграф\\
0x4E37 \textbf{nrah} [\emph{nräʰ}]: тепло\\
0x4E4F \textbf{srah} [\emph{sräʰ}]: до\\
0x4E57 \textbf{trah} [\emph{träʰ}]: сравнительная степень\\
0x4E67 \textbf{frah} [\emph{fräʰ}]: ранен\\
0x4E6F \textbf{crah} [\emph{ʃräʰ}]: считают,\\
0x4E7E \textbf{de7} [\emph{de˥}]: де\\
0x4E8F \textbf{brah} [\emph{bräʰ}]: ловушка\\
0x4E97 \textbf{grah} [\emph{gräʰ}]: восхищение\\
0x4E9F \textbf{drah} [\emph{dräʰ}]: доволен\\
0x4EA7 \textbf{zrah} [\emph{zräʰ}]: истекающий\\
0x4EAF \textbf{jrah} [\emph{ʒräʰ}]: настоящее время причастие\\
0x4EB7 \textbf{vrah} [\emph{vräʰ}]: неудовлетворенность\\
0x4ECF \textbf{qrah} [\emph{ŋräʰ}]: текущий актуальность маркер\\
0x4ED7 \textbf{xrah} [\emph{xräʰ}]: жадность\\
0x4F17 \textbf{kxah} [\emph{kxäʰ}]: экстатический\\
0x4F27 \textbf{pxah} [\emph{pxäʰ}]: завлечь\\
0x4F57 \textbf{txah} [\emph{txäʰ}]: местоимение маркер\\
0x4F5E \textbf{xe7} [\emph{xe˥}]: Эм\\
0x4F8F \textbf{bxah} [\emph{bxäʰ}]: пресыщенный\\
0x4F97 \textbf{gxah} [\emph{gxäʰ}]: большой\\
0x4F9F \textbf{dxah} [\emph{dxäʰ}]: дистанционный пульт мимо напряженный\\
\subsection{ 5 }
0x5002 \textbf{kyi7p} [\emph{kji˥p}]: раскол быстро и без хлопот\\
0x5004 \textbf{pyi7p} [\emph{pji˥p}]: прототипичный\\
0x5006 \textbf{nyi7p} [\emph{nji˥p}]: односложный\\
0x5008 \textbf{hmi7p} [\emph{ʰmi˥p}]: мини-бар\\
0x500A \textbf{tyi7p} [\emph{tji˥p}]: тайпей\\
0x500E \textbf{ryi7p} [\emph{rji˥p}]: пресбиопия\\
0x500F \textbf{myuh} [\emph{mjuʰ}]: юмор\\
0x5010 \textbf{hki7p} [\emph{ʰki˥p}]: валовой внутренний продукт\\
0x5018 \textbf{hyi7p} [\emph{ʰji˥p}]: посмешище\\
0x501A \textbf{xyi7p} [\emph{xji˥p}]: хиппи\\
0x5020 \textbf{hpi7p} [\emph{ʰpi˥p}]: непроходимый\\
0x5027 \textbf{pyuh} [\emph{pjuʰ}]: но\\
0x5028 \textbf{hwi7p} [\emph{ʰwi˥p}]: пшеничный\\
0x502A \textbf{twi7p} [\emph{twi˥p}]: лакомый кусок\\
0x502B \textbf{lwi7p} [\emph{lwi˥p}]: скользкий\\
0x5030 \textbf{hni7p} [\emph{ʰni˥p}]: фенотипический\\
0x5037 \textbf{nyuh} [\emph{njuʰ}]: транс цифра\\
0x5042 \textbf{ksi7p} [\emph{ksi˥p}]: текстиль\\
0x5044 \textbf{psi7p} [\emph{psi˥p}]: покрытосемянное\\
0x5048 \textbf{hsi7p} [\emph{ʰsi˥p}]: Константинопольский\\
0x504A \textbf{tsi7p} [\emph{tsi˥p}]: шепелявость\\
0x504F \textbf{syuh} [\emph{sjuʰ}]: супер\\
0x5050 \textbf{hti7p} [\emph{ʰti˥p}]: неисчерпаемый\\
0x5057 \textbf{tyuh} [\emph{tjuʰ}]: только\\
0x5058 \textbf{hli7p} [\emph{ʰli˥p}]: дирижабль\\
0x5060 \textbf{hfi7p} [\emph{ʰfi˥p}]: развертывание\\
0x5062 \textbf{kli7p} [\emph{kli˥p}]: кольраби\\
0x5064 \textbf{pli7p} [\emph{pli˥p}]: непродаваемый\\
0x5068 \textbf{hci7p} [\emph{ʰʃi˥p}]: пузырек\\
0x5069 \textbf{sli7p} [\emph{sli˥p}]: слоги\\
0x506A \textbf{tli7p} [\emph{tli˥p}]: талибы\\
0x506C \textbf{fli7p} [\emph{fli˥p}]: Филипп\\
0x506F \textbf{cyuh} [\emph{ʃjuʰ}]: сомнение\\
0x5070 \textbf{hri7p} [\emph{ʰri˥p}]: сорок пять\\
0x5071 \textbf{bli7p} [\emph{bli˥p}]: шаблонный\\
0x5077 \textbf{ryuh} [\emph{rjuʰ}]: порядковый\\
0x5088 \textbf{hbi7p} [\emph{ʰbi˥p}]: Сивилла\\
0x5090 \textbf{hgi7p} [\emph{ʰgi˥p}]: сороки\\
0x5098 \textbf{hdi7p} [\emph{ʰdi˥p}]: Эдинбург\\
0x50A0 \textbf{hzi7p} [\emph{ʰzi˥p}]: живой мертвец\\
0x50A4 \textbf{pci7p} [\emph{pʃi˥p}]: пирог столкнулся\\
0x50AA \textbf{tci7p} [\emph{tʃi˥p}]: палантин\\
0x50C1 \textbf{mri7p} [\emph{mri˥p}]: мизантроп\\
0x50C2 \textbf{kri7p} [\emph{kri˥p}]: микроскопия\\
0x50C4 \textbf{pri7p} [\emph{pri˥p}]: околоцветник\\
0x50C6 \textbf{nri7p} [\emph{nri˥p}]: иннервация\\
0x50C8 \textbf{hqi7p} [\emph{ʰŋi˥p}]: нелюдимый\\
0x50C9 \textbf{sri7p} [\emph{sri˥p}]: церебральный\\
0x50CA \textbf{tri7p} [\emph{tri˥p}]: щетина\\
0x50D0 \textbf{hxi7p} [\emph{ʰxi˥p}]: постоялый двор\\
0x50D3 \textbf{dri7p} [\emph{dri˥p}]: гидрофобный\\
0x50DA \textbf{xri7p} [\emph{xri˥p}]: пережить\\
0x50DE \textbf{no7} [\emph{no˥}]: О.Б.\\
0x50E2 \textbf{kxi7p} [\emph{kxi˥p}]: чипсов\\
0x50EA \textbf{txi7p} [\emph{txi˥p}]: лежал в\\
0x5106 \textbf{nya7p} [\emph{njä˥p}]: односложное слово\\
0x5108 \textbf{hma7p} [\emph{ʰmä˥p}]: пиломатериалы\\
0x5109 \textbf{sya7p} [\emph{sjä˥p}]: шпунт\\
0x510A \textbf{tya7p} [\emph{tjä˥p}]: летать по ночам\\
0x510B \textbf{lya7p} [\emph{ljä˥p}]: синий воротник\\
0x510E \textbf{rya7p} [\emph{rjä˥p}]: расколачивающий\\
0x510F \textbf{mwuh} [\emph{mwuʰ}]: мычание\\
0x5110 \textbf{hka7p} [\emph{ʰkä˥p}]: неизмеримый\\
0x5118 \textbf{hya7p} [\emph{ʰjä˥p}]: хорошо обоснован\\
0x5120 \textbf{hpa7p} [\emph{ʰpä˥p}]: бледность\\
0x5122 \textbf{kwa7p} [\emph{kwä˥p}]: вигны\\
0x5128 \textbf{hwa7p} [\emph{ʰwä˥p}]: пятьдесят восемь\\
0x5129 \textbf{swa7p} [\emph{swä˥p}]: тральщик\\
0x512E \textbf{rwa7p} [\emph{rwä˥p}]: описки\\
0x5130 \textbf{hna7p} [\emph{ʰnä˥p}]: бесспорный\\
0x5137 \textbf{nwuh} [\emph{nwuʰ}]: research gender, 6th density life form, learning to blend love and wisdom, akin to Ra, Gandhi, Tesla, Prophet Mohammad (pbuh)\\
0x5142 \textbf{ksa7p} [\emph{ksä˥p}]: аркебуза\\
0x5144 \textbf{psa7p} [\emph{psä˥p}]: утконос\\
0x5148 \textbf{hsa7p} [\emph{ʰsä˥p}]: стартапы\\
0x514A \textbf{tsa7p} [\emph{tsä˥p}]: девяносто восемь\\
0x514F \textbf{swuh} [\emph{swuʰ}]: предположительный настроение\\
0x5150 \textbf{hta7p} [\emph{ʰtä˥p}]: фартинг\\
0x5157 \textbf{twuh} [\emph{twuʰ}]: dubitative mood, for saying things the speaker is doubtful of, as for expressing sarcasm.\\
0x5158 \textbf{hla7p} [\emph{ʰlä˥p}]: плач Иеремии\\
0x515E \textbf{to7} [\emph{to˥}]: OMA\\
0x515F \textbf{lwuh} [\emph{lwuʰ}]: ленивый\\
0x5160 \textbf{hfa7p} [\emph{ʰfä˥p}]: полубак\\
0x5162 \textbf{kla7p} [\emph{klä˥p}]: кабала\\
0x5164 \textbf{pla7p} [\emph{plä˥p}]: непреложный\\
0x5168 \textbf{hca7p} [\emph{ʰʃä˥p}]: шарлатан\\
0x5169 \textbf{sla7p} [\emph{slä˥p}]: самбо\\
0x516A \textbf{tla7p} [\emph{tlä˥p}]: телепатия\\
0x516D \textbf{cla7p} [\emph{ʃlä˥p}]: мытье кожи\\
0x516F \textbf{cwuh} [\emph{ʃwuʰ}]: шок\\
0x5170 \textbf{hra7p} [\emph{ʰrä˥p}]: арабский\\
0x5171 \textbf{bla7p} [\emph{blä˥p}]: болт канат\\
0x5173 \textbf{dla7p} [\emph{dlä˥p}]: неровными краями\\
0x5177 \textbf{rwuh} [\emph{rwuʰ}]: ревнивый\\
0x517A \textbf{xla7p} [\emph{xlä˥p}]: хлопушка\\
0x5188 \textbf{hba7p} [\emph{ʰbä˥p}]: Бомбей\\
0x5190 \textbf{hga7p} [\emph{ʰgä˥p}]: водобоязнь\\
0x5198 \textbf{hda7p} [\emph{ʰdä˥p}]: саман\\
0x51A2 \textbf{kca7p} [\emph{kʃä˥p}]: довершение пирог\\
0x51A4 \textbf{pca7p} [\emph{pʃä˥p}]: парашютист\\
0x51AA \textbf{tca7p} [\emph{tʃä˥p}]: высотомер\\
0x51C1 \textbf{mra7p} [\emph{mrä˥p}]: ламантин\\
0x51C2 \textbf{kra7p} [\emph{krä˥p}]: карп\\
0x51C4 \textbf{pra7p} [\emph{prä˥p}]: апперцепция\\
0x51C8 \textbf{hqa7p} [\emph{ʰŋä˥p}]: черное и белое\\
0x51C9 \textbf{sra7p} [\emph{srä˥p}]: сербия\\
0x51CA \textbf{tra7p} [\emph{trä˥p}]: атропин\\
0x51D0 \textbf{hxa7p} [\emph{ʰxä˥p}]: Йоханнесбург\\
0x51D1 \textbf{bra7p} [\emph{brä˥p}]: обломок кирпича\\
0x51D2 \textbf{gra7p} [\emph{grä˥p}]: мостовые весы\\
0x51DA \textbf{xra7p} [\emph{xrä˥p}]: гарпия\\
0x51E2 \textbf{kxa7p} [\emph{kxä˥p}]: взять сумку\\
0x51E4 \textbf{pxa7p} [\emph{pxä˥p}]: после того, как хлопают\\
0x51EA \textbf{txa7p} [\emph{txä˥p}]: хорошо носится\\
0x5208 \textbf{hmu7p} [\emph{ʰmu˥p}]: надгробная плита\\
0x5210 \textbf{hku7p} [\emph{ʰku˥p}]: Кабул\\
0x5217 \textbf{ksuh} [\emph{ksuʰ}]: declarative mood, for declaring the begining of a program or essay.\\
0x5218 \textbf{hyu7p} [\emph{ʰju˥p}]: антилопа гну\\
0x5220 \textbf{hpu7p} [\emph{ʰpu˥p}]: Pooh Pooh\\
0x5227 \textbf{psuh} [\emph{psuʰ}]: superessive case, indicates location on a surface. surface-context location-case.\\
0x5228 \textbf{hwu7p} [\emph{ʰwu˥p}]: пятикратный\\
0x5230 \textbf{hnu7p} [\emph{ʰnu˥p}]: дышло\\
0x5248 \textbf{hsu7p} [\emph{ʰsu˥p}]: стократный\\
0x5250 \textbf{htu7p} [\emph{ʰtu˥p}]: бесспорно\\
0x5257 \textbf{tsuh} [\emph{tsuʰ}]: дистальный демонстративный\\
0x526A \textbf{tlu7p} [\emph{tlu˥p}]: толуол\\
0x5270 \textbf{hru7p} [\emph{ʰru˥p}]: сумчатый\\
0x5271 \textbf{blu7p} [\emph{blu˥p}]: бык щенок\\
0x5288 \textbf{hbu7p} [\emph{ʰbu˥p}]: епископы\\
0x5298 \textbf{hdu7p} [\emph{ʰdu˥p}]: Дунай\\
0x52D1 \textbf{bru7p} [\emph{bru˥p}]: билирубин\\
0x5308 \textbf{hme7p} [\emph{ʰme˥p}]: приветливо\\
0x5309 \textbf{sye7p} [\emph{sje˥p}]: Сен-Пьер\\
0x5310 \textbf{hke7p} [\emph{ʰke˥p}]: кастаньеты\\
0x5317 \textbf{kluh} [\emph{kluʰ}]: путаница\\
0x5318 \textbf{hye7p} [\emph{ʰje˥p}]: тембр\\
0x5320 \textbf{hpe7p} [\emph{ʰpe˥p}]: северная часть штата\\
0x5327 \textbf{pluh} [\emph{pluʰ}]: превосходный степень\\
0x5330 \textbf{hne7p} [\emph{ʰne˥p}]: несравненный\\
0x5342 \textbf{kse7p} [\emph{kse˥p}]: шекспировский\\
0x5344 \textbf{pse7p} [\emph{pse˥p}]: космодром\\
0x5348 \textbf{hse7p} [\emph{ʰse˥p}]: сепия\\
0x534A \textbf{tse7p} [\emph{tse˥p}]: девяносто шесть\\
0x534F \textbf{sluh} [\emph{sluʰ}]: самостоятельно удовлетворяет\\
0x5350 \textbf{hte7p} [\emph{ʰte˥p}]: Фивы\\
0x5357 \textbf{tluh} [\emph{tluʰ}]: единственное число\\
0x5358 \textbf{hle7p} [\emph{ʰle˥p}]: эллипсы\\
0x536F \textbf{cluh} [\emph{ʃluʰ}]: уродливый\\
0x5370 \textbf{hre7p} [\emph{ʰre˥p}]: кренделек\\
0x538F \textbf{bluh} [\emph{bluʰ}]: логический\\
0x5397 \textbf{gluh} [\emph{gluʰ}]: singulative номер\\
0x539F \textbf{dluh} [\emph{dluʰ}]: двойной номер\\
0x53B0 \textbf{hve7p} [\emph{ʰve˥p}]: вечерня\\
0x53C2 \textbf{kre7p} [\emph{kre˥p}]: херувим\\
0x53D0 \textbf{hxe7p} [\emph{ʰxe˥p}]: знаменоносец\\
0x5408 \textbf{hmo7p} [\emph{ʰmo˥p}]: шарить\\
0x5410 \textbf{hko7p} [\emph{ʰko˥p}]: неконтролируемый\\
0x5420 \textbf{hpo7p} [\emph{ʰpo˥p}]: пропан\\
0x5427 \textbf{pfuh} [\emph{pfuʰ}]: прощальный привет\\
0x5428 \textbf{hwo7p} [\emph{ʰwo˥p}]: олигополия\\
0x5430 \textbf{hno7p} [\emph{ʰno˥p}]: Point Blank\\
0x5442 \textbf{kso7p} [\emph{kso˥p}]: несъедобный\\
0x5448 \textbf{hso7p} [\emph{ʰso˥p}]: левый борт\\
0x544A \textbf{tso7p} [\emph{tso˥p}]: фронтиспис\\
0x5450 \textbf{hto7p} [\emph{ʰto˥p}]: трепел\\
0x5457 \textbf{tfuh} [\emph{tfuʰ}]: цитата форма\\
0x5458 \textbf{hlo7p} [\emph{ʰlo˥p}]: Лиссабон\\
0x5462 \textbf{klo7p} [\emph{klo˥p}]: предварительно колумбийский\\
0x5468 \textbf{hco7p} [\emph{ʰʃo˥p}]: впитывает\\
0x5470 \textbf{hro7p} [\emph{ʰro˥p}]: реторты\\
0x5488 \textbf{hbo7p} [\emph{ʰbo˥p}]: нажать кнопку\\
0x54C2 \textbf{kro7p} [\emph{kro˥p}]: нагрузка\\
0x54C9 \textbf{sro7p} [\emph{sro˥p}]: строб\\
0x54CA \textbf{tro7p} [\emph{tro˥p}]: трипсы\\
0x54D0 \textbf{hxo7p} [\emph{ʰxo˥p}]: горбатый\\
0x5527 \textbf{pcuh} [\emph{pʃuʰ}]: неопределенность\\
0x5530 \textbf{hn67p} [\emph{ʰnə˥p}]: звукоподражательный\\
0x5557 \textbf{tcuh} [\emph{tʃuʰ}]: энтузиазм\\
0x560F \textbf{mruh} [\emph{mruʰ}]: дурак\\
0x5617 \textbf{kruh} [\emph{kruʰ}]: горький\\
0x5627 \textbf{pruh} [\emph{pruʰ}]: imprecative mood, for uttering curses.\\
0x5637 \textbf{nruh} [\emph{nruʰ}]: невинный\\
0x564F \textbf{sruh} [\emph{sruʰ}]: трепет\\
0x5657 \textbf{truh} [\emph{truʰ}]: температура\\
0x566F \textbf{cruh} [\emph{ʃruʰ}]: позор\\
0x5697 \textbf{gruh} [\emph{gruʰ}]: приседать\\
0x56D7 \textbf{xruh} [\emph{xruʰ}]: веселый\\
0x5757 \textbf{txuh} [\emph{txuʰ}]: приступ боли\\
0x5808 \textbf{hmi2p} [\emph{ʰmi˩p}]: микроавтобус\\
0x5809 \textbf{syi2p} [\emph{sji˩p}]: скопа\\
0x580A \textbf{tyi2p} [\emph{tji˩p}]: хронометраж\\
0x580D \textbf{cyi2p} [\emph{ʃji˩p}]: шиит\\
0x580F \textbf{myeh} [\emph{mjeʰ}]: ампер\\
0x5810 \textbf{hki2p} [\emph{ʰki˩p}]: каннибал\\
0x5817 \textbf{kyeh} [\emph{kjeʰ}]: пожалел\\
0x5818 \textbf{hyi2p} [\emph{ʰji˩p}]: ночной клуб\\
0x5820 \textbf{hpi2p} [\emph{ʰpi˩p}]: волынщик\\
0x5827 \textbf{pyeh} [\emph{pjeʰ}]: предикативный\\
0x5830 \textbf{hni2p} [\emph{ʰni˩p}]: ночной колпак\\
0x5837 \textbf{nyeh} [\emph{njeʰ}]: initiative case, indicates the begining of an action. time-context source-case\\
0x5842 \textbf{ksi2p} [\emph{ksi˩p}]: аксиллярный\\
0x5844 \textbf{psi2p} [\emph{psi˩p}]: епископальный\\
0x5848 \textbf{hsi2p} [\emph{ʰsi˩p}]: бекас\\
0x584A \textbf{tsi2p} [\emph{tsi˩p}]: семьдесят один\\
0x584F \textbf{syeh} [\emph{sjeʰ}]: с этого времени\\
0x5850 \textbf{hti2p} [\emph{ʰti˩p}]: хорошо одетым\\
0x5857 \textbf{tyeh} [\emph{tjeʰ}]: текст\\
0x5858 \textbf{hli2p} [\emph{ʰli˩p}]: Элизабет\\
0x585F \textbf{lyeh} [\emph{ljeʰ}]: atelic aspect, for actions that have no particular end point. infinte or unlimited.\\
0x5860 \textbf{hfi2p} [\emph{ʰfi˩p}]: friborg\\
0x5862 \textbf{kli2p} [\emph{kli˩p}]: безутешный\\
0x5868 \textbf{hci2p} [\emph{ʰʃi˩p}]: наложница\\
0x5869 \textbf{sli2p} [\emph{sli˩p}]: несколько миллиардов\\
0x586A \textbf{tli2p} [\emph{tli˩p}]: самосвал\\
0x586F \textbf{cyeh} [\emph{ʃjeʰ}]: задаваться вопросом\\
0x5870 \textbf{hri2p} [\emph{ʰri˩p}]: прокатка штырь\\
0x5877 \textbf{ryeh} [\emph{rjeʰ}]: ретроспективный\\
0x588A \textbf{tfi2p} [\emph{tfi˩p}]: чеканка\\
0x588F \textbf{byeh} [\emph{bjeʰ}]: байт\\
0x5890 \textbf{hgi2p} [\emph{ʰgi˩p}]: волынка\\
0x5898 \textbf{hdi2p} [\emph{ʰdi˩p}]: несомненно\\
0x589F \textbf{dyeh} [\emph{djeʰ}]: непосредственный объект\\
0x58AA \textbf{tci2p} [\emph{tʃi˩p}]: растирать в порошок\\
0x58C1 \textbf{mri2p} [\emph{mri˩p}]: мизантропия\\
0x58C2 \textbf{kri2p} [\emph{kri˩p}]: карибу\\
0x58C6 \textbf{nri2p} [\emph{nri˩p}]: невропатический\\
0x58C9 \textbf{sri2p} [\emph{sri˩p}]: изопрен\\
0x58CA \textbf{tri2p} [\emph{tri˩p}]: червивый\\
0x58CF \textbf{qyeh} [\emph{ŋjeʰ}]: Один действие глагол\\
0x58D0 \textbf{hxi2p} [\emph{ʰxi˩p}]: делать орфографические ошибки\\
0x58EA \textbf{txi2p} [\emph{txi˩p}]: опечатка\\
0x5904 \textbf{pya2p} [\emph{pjä˩p}]: назад ткань\\
0x5908 \textbf{hma2p} [\emph{ʰmä˩p}]: недуги\\
0x5909 \textbf{sya2p} [\emph{sjä˩p}]: бестселлер\\
0x590A \textbf{tya2p} [\emph{tjä˩p}]: банка открывашка\\
0x590F \textbf{mweh} [\emph{mweʰ}]: заканчивающий аспект\\
0x5910 \textbf{hka2p} [\emph{ʰkä˩p}]: сапожник\\
0x5917 \textbf{kweh} [\emph{kweʰ}]: эксклюзивный человек\\
0x5918 \textbf{hya2p} [\emph{ʰjä˩p}]: ямбический\\
0x5920 \textbf{hpa2p} [\emph{ʰpä˩p}]: Пенджаб\\
0x5928 \textbf{hwa2p} [\emph{ʰwä˩p}]: Ваххабиты\\
0x5930 \textbf{hna2p} [\emph{ʰnä˩p}]: Занзибар\\
0x5937 \textbf{nweh} [\emph{nweʰ}]: inessive case, indicates the noun phrase being in at. interior-context location-case. \\
0x5942 \textbf{ksa2p} [\emph{ksä˩p}]: гадюка\\
0x5944 \textbf{psa2p} [\emph{psä˩p}]: вспыльчивый\\
0x5948 \textbf{hsa2p} [\emph{ʰsä˩p}]: в средней части судна\\
0x594A \textbf{tsa2p} [\emph{tsä˩p}]: мусорный ящик\\
0x594F \textbf{sweh} [\emph{sweʰ}]: sublative case, indicates motion toward the underside of nounphrase. under-context destination-case.\\
0x5950 \textbf{hta2p} [\emph{ʰtä˩p}]: приписывал\\
0x5957 \textbf{tweh} [\emph{tweʰ}]: terminative case, indicates temporal destination or goal. time-context destination-case. \\
0x5958 \textbf{hla2p} [\emph{ʰlä˩p}]: филантропия\\
0x595F \textbf{lweh} [\emph{lweʰ}]: allative case, indicates motion onto a surface. surface-context destination-case. \\
0x5960 \textbf{hfa2p} [\emph{ʰfä˩p}]: ожирение\\
0x5962 \textbf{kla2p} [\emph{klä˩p}]: спина к спине\\
0x5964 \textbf{pla2p} [\emph{plä˩p}]: многосложный\\
0x5968 \textbf{hca2p} [\emph{ʰʃä˩p}]: шрапнель\\
0x5969 \textbf{sla2p} [\emph{slä˩p}]: селитра\\
0x596A \textbf{tla2p} [\emph{tlä˩p}]: боевой топор\\
0x596F \textbf{cweh} [\emph{ʃweʰ}]: смиренный\\
0x5970 \textbf{hra2p} [\emph{ʰrä˩p}]: рецидив\\
0x5977 \textbf{rweh} [\emph{rweʰ}]: oblique case, indicates intersection with noun phrase, such as an orbit or street that is being crossed at an oblique angle.\\
0x5988 \textbf{hba2p} [\emph{ʰbä˩p}]: анабаптист\\
0x5997 \textbf{gweh} [\emph{gweʰ}]: скоро будущее напряженный\\
0x5998 \textbf{hda2p} [\emph{ʰdä˩p}]: твердый и быстрый\\
0x599F \textbf{dweh} [\emph{dweʰ}]: delimitive aspect, for referring to events of limited duration, like a limited time offer\\
0x59A4 \textbf{pca2p} [\emph{pʃä˩p}]: ломбард\\
0x59AA \textbf{tca2p} [\emph{tʃä˩p}]: магазин игрушек\\
0x59B7 \textbf{vweh} [\emph{vweʰ}]: перегружены работой\\
0x59C1 \textbf{mra2p} [\emph{mrä˩p}]: амарант\\
0x59C2 \textbf{kra2p} [\emph{krä˩p}]: карбид\\
0x59C8 \textbf{hqa2p} [\emph{ʰŋä˩p}]: превосходить численно\\
0x59C9 \textbf{sra2p} [\emph{srä˩p}]: вызов в суд\\
0x59CA \textbf{tra2p} [\emph{trä˩p}]: преступил\\
0x59CC \textbf{fra2p} [\emph{frä˩p}]: импровизированная трибуна\\
0x59D0 \textbf{hxa2p} [\emph{ʰxä˩p}]: гусар\\
0x59EA \textbf{txa2p} [\emph{txä˩p}]: Марки\\
0x5A08 \textbf{hmu2p} [\emph{ʰmu˩p}]: обагрять\\
0x5A10 \textbf{hku2p} [\emph{ʰku˩p}]: околополярный\\
0x5A17 \textbf{kseh} [\emph{kseʰ}]: навязчивая идея\\
0x5A27 \textbf{pseh} [\emph{pseʰ}]: аспект\\
0x5A48 \textbf{hsu2p} [\emph{ʰsu˩p}]: шаг гуся\\
0x5A50 \textbf{htu2p} [\emph{ʰtu˩p}]: утопия\\
0x5A57 \textbf{tseh} [\emph{tseʰ}]: potential mood, something that is probable according to the speaker.\\
0x5A58 \textbf{hlu2p} [\emph{ʰlu˩p}]: лупа\\
0x5A70 \textbf{hru2p} [\emph{ʰru˩p}]: рубль\\
0x5B08 \textbf{hme2p} [\emph{ʰme˩p}]: табличка с именем\\
0x5B10 \textbf{hke2p} [\emph{ʰke˩p}]: эр\\
0x5B17 \textbf{kleh} [\emph{kleʰ}]: собака устала\\
0x5B20 \textbf{hpe2p} [\emph{ʰpe˩p}]: чашелистник\\
0x5B27 \textbf{pleh} [\emph{pleʰ}]: привилегированный\\
0x5B30 \textbf{hne2p} [\emph{ʰne˩p}]: эндотелий\\
0x5B42 \textbf{kse2p} [\emph{kse˩p}]: облагаемый\\
0x5B44 \textbf{pse2p} [\emph{pse˩p}]: асептика\\
0x5B48 \textbf{hse2p} [\emph{ʰse˩p}]: десятикратный\\
0x5B4A \textbf{tse2p} [\emph{tse˩p}]: в середине сентября\\
0x5B4F \textbf{sleh} [\emph{sleʰ}]: vegetal gender, 2nd density life form, learning to fulfill desires, akin to plant, fungi, basic electronics\\
0x5B50 \textbf{hte2p} [\emph{ʰte˩p}]: уроженцев\\
0x5B57 \textbf{tleh} [\emph{tleʰ}]: hesternal tense, past day tense, for things that happened yesterday.\\
0x5B58 \textbf{hle2p} [\emph{ʰle˩p}]: приозерный\\
0x5B67 \textbf{fleh} [\emph{fleʰ}]: изменяемую\\
0x5B6F \textbf{cleh} [\emph{ʃleʰ}]: недавний мимо напряженный\\
0x5B70 \textbf{hre2p} [\emph{ʰre˩p}]: повторно заложен\\
0x5B8F \textbf{bleh} [\emph{bleʰ}]: благожелательный\\
0x5B9F \textbf{dleh} [\emph{dleʰ}]: hodiernal tense, for expressing today.\\
0x5BD7 \textbf{xleh} [\emph{xleʰ}]: vehicle gender, 7th density life form, learning about timelines, akin to spirit of a ship/vehicle/spaceship\\
0x5C08 \textbf{hmo2p} [\emph{ʰmo˩p}]: в середине октября\\
0x5C10 \textbf{hko2p} [\emph{ʰko˩p}]: копра\\
0x5C17 \textbf{kfeh} [\emph{kfeʰ}]: ядро благословляющий\\
0x5C20 \textbf{hpo2p} [\emph{ʰpo˩p}]: попурри\\
0x5C27 \textbf{pfeh} [\emph{pfeʰ}]: чуткими радость\\
0x5C30 \textbf{hno2p} [\emph{ʰno˩p}]: О.Б.\\
0x5C42 \textbf{kso2p} [\emph{kso˩p}]: макроскопический\\
0x5C48 \textbf{hso2p} [\emph{ʰso˩p}]: снежный шар\\
0x5C50 \textbf{hto2p} [\emph{ʰto˩p}]: трубач\\
0x5C57 \textbf{tfeh} [\emph{tfeʰ}]: telic aspect, for actions that have a definite end point. finite or limited.\\
0x5C58 \textbf{hlo2p} [\emph{ʰlo˩p}]: кувырок провал\\
0x5C70 \textbf{hro2p} [\emph{ʰro˩p}]: ларингоскоп\\
0x5C8F \textbf{bveh} [\emph{bveʰ}]: subessive case, indicates location at the underside of nounphrase. under-context location-case.\\
0x5C97 \textbf{gveh} [\emph{gveʰ}]: эргатив дело\\
0x5CCA \textbf{tro2p} [\emph{tro˩p}]: таро\\
0x5CD0 \textbf{hxo2p} [\emph{ʰxo˩p}]: Ходж мешанина\\
0x5D17 \textbf{kceh} [\emph{kʃeʰ}]: лучше\\
0x5D27 \textbf{pceh} [\emph{pʃeʰ}]: apprehensive mood, for expressing fears or warnings.\\
0x5D57 \textbf{tceh} [\emph{tʃeʰ}]: одна и та же предмет маркер\\
0x5D8F \textbf{bjeh} [\emph{bʒeʰ}]: locality gender, 8th density life form, learning to be oversoul, akin to spirit of a place or flock\\
0x5E0F \textbf{mreh} [\emph{mreʰ}]: momentane aspect, for referring to momentary actions, like the unfurling of flower petals in the morning\\
0x5E17 \textbf{kreh} [\emph{kreʰ}]: жестокий\\
0x5E27 \textbf{preh} [\emph{preʰ}]: предопределенный\\
0x5E37 \textbf{nreh} [\emph{nreʰ}]: энергия\\
0x5E4F \textbf{sreh} [\emph{sreʰ}]: серии\\
0x5E57 \textbf{treh} [\emph{treʰ}]: тщеславный\\
0x5E67 \textbf{freh} [\emph{freʰ}]: справочный\\
0x5E6F \textbf{creh} [\emph{ʃreʰ}]: Генри\\
0x5E8F \textbf{breh} [\emph{breʰ}]: изменение из мероприятие\\
0x5E97 \textbf{greh} [\emph{greʰ}]: centric case, indicates the noun phrase being the centre. Such as a star or actor.\\
0x5E9F \textbf{dreh} [\emph{dreʰ}]: второй язык\\
0x5EA7 \textbf{zreh} [\emph{zreʰ}]: exocentric case, indicates the noun phrase orbiting the other participants. Such as a moon or audience.\\
0x5ED7 \textbf{xreh} [\emph{xreʰ}]: герц\\
0x5F17 \textbf{kxeh} [\emph{kxeʰ}]: к дело\\
0x5F27 \textbf{pxeh} [\emph{pxeʰ}]: partitive case, a subcase that indicates only partial usage of the noun phrase. lazy argument loading.\\
0x5F57 \textbf{txeh} [\emph{txeʰ}]: увеличительный\\
0x5F9F \textbf{dxeh} [\emph{dxeʰ}]: другой мероприятие\\
\subsection{ 6 }
0x600F \textbf{myoh} [\emph{mjoʰ}]: средний голос\\
0x6017 \textbf{kyoh} [\emph{kjoʰ}]: кило\\
0x6027 \textbf{pyoh} [\emph{pjoʰ}]: пассивный голос\\
0x6037 \textbf{nyoh} [\emph{njoʰ}]: анимационный\\
0x604F \textbf{syoh} [\emph{sjoʰ}]: благоговение\\
0x6057 \textbf{tyoh} [\emph{tjoʰ}]: тривиальный\\
0x605E \textbf{ki2} [\emph{ki˩}]: чрезмерное продолжительность\\
0x6067 \textbf{fyoh} [\emph{fjoʰ}]: возвратный голос\\
0x606F \textbf{cyoh} [\emph{ʃjoʰ}]: уменьшительный\\
0x6077 \textbf{ryoh} [\emph{rjoʰ}]: prosecutive case, indicates motion along a surface. surface-context way-case.\\
0x609E \textbf{pi2} [\emph{pi˩}]: planetary gender, 9th density life form, akin to spirit of a planet, Gaia, Venus, Luna\\
0x609F \textbf{dyoh} [\emph{djoʰ}]: приспособление\\
0x60DE \textbf{ni2} [\emph{ni˩}]: безличный глагол\\
0x610F \textbf{mwoh} [\emph{mwoʰ}]: меланхолия\\
0x6137 \textbf{nwoh} [\emph{nwoʰ}]: Событийный настроение\\
0x613E \textbf{si2} [\emph{si˩}]: necessitative настроение\\
0x614F \textbf{swoh} [\emph{swoʰ}]: последовательный\\
0x6157 \textbf{twoh} [\emph{twoʰ}]: quotative case, indicates a quote. discource-context location-case.\\
0x615E \textbf{ti2} [\emph{ti˩}]: электрический текущий\\
0x615F \textbf{lwoh} [\emph{lwoʰ}]: locative case, indicates location in space. space-context location-case. \\
0x616F \textbf{cwoh} [\emph{ʃwoʰ}]: одна и та же мероприятие маркер\\
0x6177 \textbf{rwoh} [\emph{rwoʰ}]: движение в гору\\
0x6217 \textbf{ksoh} [\emph{ksoʰ}]: сексуальное\\
0x6227 \textbf{psoh} [\emph{psoʰ}]: непобедимый\\
0x6257 \textbf{tsoh} [\emph{tsoʰ}]: уважение\\
0x628F \textbf{bzoh} [\emph{bzoʰ}]: OSE\\
0x629F \textbf{dzoh} [\emph{dzoʰ}]: подъязычная\\
0x6317 \textbf{kloh} [\emph{kloʰ}]: Одинокий\\
0x6327 \textbf{ploh} [\emph{ploʰ}]: за предложение\\
0x634F \textbf{sloh} [\emph{sloʰ}]: успех\\
0x6357 \textbf{tloh} [\emph{tloʰ}]: multiplicative case, for indicating how many times the event occurs. as for frequentative-aspect or simple repeat loops.\\
0x636F \textbf{cloh} [\emph{ʃloʰ}]: rational gender, any entity capable of reason, 3rd density life form or higher\\
0x639F \textbf{dloh} [\emph{dloʰ}]: доллар\\
0x63B7 \textbf{vloh} [\emph{vloʰ}]: вольт\\
0x6457 \textbf{tfoh} [\emph{tfoʰ}]: моль\\
0x64DE \textbf{na2} [\emph{nä˩}]: нафтол\\
0x6517 \textbf{kcoh} [\emph{kʃoʰ}]: устала\\
0x6527 \textbf{pcoh} [\emph{pʃoʰ}]: заостренный\\
0x6557 \textbf{tcoh} [\emph{tʃoʰ}]: боль\\
0x655E \textbf{ta2} [\emph{tä˩}]: continuative aspect, the state is continuing to occur, it is `still' occuring. \\
0x65DE \textbf{ra2} [\emph{rä˩}]: frequentative aspect, indicates the verb happens multiple times, for example bite frequentative-aspect is equivalent to chewing. \\
0x660F \textbf{mroh} [\emph{mroʰ}]: моральный\\
0x6617 \textbf{kroh} [\emph{kroʰ}]: очарование\\
0x6627 \textbf{proh} [\emph{proʰ}]: сообщается доказательный\\
0x6637 \textbf{nroh} [\emph{nroʰ}]: гнев\\
0x664F \textbf{sroh} [\emph{sroʰ}]: помилование\\
0x6657 \textbf{troh} [\emph{troʰ}]: непосредственный доказательный\\
0x666F \textbf{croh} [\emph{ʃroʰ}]: страх\\
0x668F \textbf{broh} [\emph{broʰ}]: актер роль\\
0x6697 \textbf{groh} [\emph{groʰ}]: Цель вызывать\\
0x669F \textbf{droh} [\emph{droʰ}]: слуховой доказательный\\
0x66D7 \textbf{xroh} [\emph{xroʰ}]: коллега по работе\\
0x6717 \textbf{kxoh} [\emph{kxoʰ}]: килограмм\\
0x6727 \textbf{pxoh} [\emph{pxoʰ}]: prolative case, indicates motion via an interior. interior-context way-case.\\
0x6757 \textbf{txoh} [\emph{txoʰ}]: скромный\\
0x679F \textbf{dxoh} [\emph{dxoʰ}]: modal case, indicates the ability or probability of the sentence. state-context location-case. neural weight.\\
0x6801 \textbf{myin} [\emph{mjin}]: команда\\
0x6802 \textbf{kyin} [\emph{kjin}]: Компания\\
0x6804 \textbf{pyin} [\emph{pjin}]: дикий\\
0x6806 \textbf{nyin} [\emph{njin}]: антропный Пол\\
0x6808 \textbf{hmin} [\emph{ʰmin}]: ничего\\
0x6809 \textbf{syin} [\emph{sjin}]: Уважаемые\\
0x680A \textbf{tyin} [\emph{tjin}]: три\\
0x680B \textbf{lyin} [\emph{ljin}]: электронный\\
0x680C \textbf{fyin} [\emph{fjin}]: расходы\\
0x680D \textbf{cyin} [\emph{ʃjin}]: ложь\\
0x680E \textbf{ryin} [\emph{rjin}]: обучение\\
0x6810 \textbf{hkin} [\emph{ʰkin}]: идти\\
0x6811 \textbf{byin} [\emph{bjin}]: двоичный\\
0x6812 \textbf{gyin} [\emph{gjin}]: крыло\\
0x6813 \textbf{dyin} [\emph{djin}]: скрипка\\
0x6814 \textbf{zyin} [\emph{zjin}]: хронический\\
0x6815 \textbf{jyin} [\emph{ʒjin}]: инжиниринг\\
0x6816 \textbf{vyin} [\emph{vjin}]: кинотеатр\\
0x6818 \textbf{hyin} [\emph{ʰjin}]: универсальный квантификатором\\
0x6819 \textbf{qyin} [\emph{ŋjin}]: биоразнообразия\\
0x681A \textbf{xyin} [\emph{xjin}]: неопытность\\
0x6820 \textbf{hpin} [\emph{ʰpin}]: adpositional дело\\
0x6821 \textbf{mwin} [\emph{mwin}]: миллиона\\
0x6822 \textbf{kwin} [\emph{kwin}]: дающий\\
0x6824 \textbf{pwin} [\emph{pwin}]: вино\\
0x6826 \textbf{nwin} [\emph{nwin}]: конституция\\
0x6828 \textbf{hwin} [\emph{ʰwin}]: место нахождения дело\\
0x6829 \textbf{swin} [\emph{swin}]: гражданского\\
0x682A \textbf{twin} [\emph{twin}]: желание\\
0x682B \textbf{lwin} [\emph{lwin}]: сбалансированный\\
0x682C \textbf{fwin} [\emph{fwin}]: частота\\
0x682D \textbf{cwin} [\emph{ʃwin}]: верховой\\
0x682E \textbf{rwin} [\emph{rwin}]: Королева\\
0x6830 \textbf{hnin} [\emph{ʰnin}]: мужской Пол\\
0x6831 \textbf{bwin} [\emph{bwin}]: впечатление\\
0x6832 \textbf{gwin} [\emph{gwin}]: свинг\\
0x6833 \textbf{dwin} [\emph{dwin}]: равнодушие\\
0x6834 \textbf{zwin} [\emph{zwin}]: резидент\\
0x6835 \textbf{jwin} [\emph{ʒwin}]: печень\\
0x6836 \textbf{vwin} [\emph{vwin}]: ботаника\\
0x6839 \textbf{qwin} [\emph{ŋwin}]: пингвин\\
0x683A \textbf{xwin} [\emph{xwin}]: минеральная Пол\\
0x6842 \textbf{ksin} [\emph{ksin}]: зеленый\\
0x6844 \textbf{psin} [\emph{psin}]: вес\\
0x6848 \textbf{hsin} [\emph{ʰsin}]: Социальное контекст\\
0x6850 \textbf{htin} [\emph{ʰtin}]: независимый пункт\\
0x6851 \textbf{bzin} [\emph{bzin}]: знак\\
0x6852 \textbf{gzin} [\emph{gzin}]: пионер\\
0x6853 \textbf{dzin} [\emph{dzin}]: 19\\
0x6858 \textbf{hlin} [\emph{ʰlin}]: с плавающей точкой номер\\
0x6860 \textbf{hfin} [\emph{ʰfin}]: богатые\\
0x6862 \textbf{klin} [\emph{klin}]: столица\\
0x6864 \textbf{plin} [\emph{plin}]: открытие\\
0x6868 \textbf{hcin} [\emph{ʰʃin}]: женский Пол\\
0x6869 \textbf{slin} [\emph{slin}]: Серебряный\\
0x686A \textbf{tlin} [\emph{tlin}]: холодно\\
0x686C \textbf{flin} [\emph{flin}]: выведенный доказательный\\
0x686D \textbf{clin} [\emph{ʃlin}]: Генри\\
0x6870 \textbf{hrin} [\emph{ʰrin}]: рациональный Пол\\
0x6871 \textbf{blin} [\emph{blin}]: воображаемый\\
0x6872 \textbf{glin} [\emph{glin}]: нежелание\\
0x6873 \textbf{dlin} [\emph{dlin}]: дисциплина\\
0x6874 \textbf{zlin} [\emph{zlin}]: оливковый\\
0x6875 \textbf{jlin} [\emph{ʒlin}]: написание\\
0x6876 \textbf{vlin} [\emph{vlin}]: средство передвижения Пол\\
0x6877 \textbf{ry6h} [\emph{rjəʰ}]: радиан\\
0x687A \textbf{xlin} [\emph{xlin}]: горячая линия\\
0x6882 \textbf{kfin} [\emph{kfin}]: крышка\\
0x6884 \textbf{pfin} [\emph{pfin}]: гулять пешком\\
0x6888 \textbf{hbin} [\emph{ʰbin}]: становление\\
0x688A \textbf{tfin} [\emph{tfin}]: Родился\\
0x6890 \textbf{hgin} [\emph{ʰgin}]: тюрьма\\
0x6891 \textbf{bvin} [\emph{bvin}]: вселенная Пол\\
0x6892 \textbf{gvin} [\emph{gvin}]: данный\\
0x6893 \textbf{dvin} [\emph{dvin}]: домен имя\\
0x6898 \textbf{hdin} [\emph{ʰdin}]: заполнить\\
0x68A0 \textbf{hzin} [\emph{ʰzin}]: южный\\
0x68A2 \textbf{kcin} [\emph{kʃin}]: глубоко\\
0x68A4 \textbf{pcin} [\emph{pʃin}]: за\\
0x68A8 \textbf{hjin} [\emph{ʰʒin}]: терпение\\
0x68AA \textbf{tcin} [\emph{tʃin}]: день\\
0x68B0 \textbf{hvin} [\emph{ʰvin}]: рационализировать\\
0x68B1 \textbf{bjin} [\emph{bʒin}]: прогнозирование\\
0x68B2 \textbf{gjin} [\emph{gʒin}]: день рождения\\
0x68B3 \textbf{djin} [\emph{dʒin}]: монета\\
0x68C1 \textbf{mrin} [\emph{mrin}]: мука\\
0x68C2 \textbf{krin} [\emph{krin}]: кардинальный\\
0x68C4 \textbf{prin} [\emph{prin}]: горизонт\\
0x68C6 \textbf{nrin} [\emph{nrin}]: невежественный\\
0x68C8 \textbf{hqin} [\emph{ʰŋin}]: нелинейный\\
0x68C9 \textbf{srin} [\emph{srin}]: инстинкт\\
0x68CA \textbf{trin} [\emph{trin}]: круг\\
0x68CC \textbf{frin} [\emph{frin}]: декорации\\
0x68CD \textbf{crin} [\emph{ʃrin}]: восхищение\\
0x68D0 \textbf{hxin} [\emph{ʰxin}]: синхронизировать\\
0x68D1 \textbf{brin} [\emph{brin}]: внутренний\\
0x68D2 \textbf{grin} [\emph{grin}]: окрестности\\
0x68D3 \textbf{drin} [\emph{drin}]: посвятить\\
0x68D4 \textbf{zrin} [\emph{zrin}]: фартук\\
0x68D5 \textbf{jrin} [\emph{ʒrin}]: каждый день\\
0x68D6 \textbf{vrin} [\emph{vrin}]: вербовщик\\
0x68D9 \textbf{qrin} [\emph{ŋrin}]: мнимая единица -\\
0x68DA \textbf{xrin} [\emph{xrin}]: ингредиенты\\
0x68E2 \textbf{kxin} [\emph{kxin}]: факториал\\
0x68E4 \textbf{pxin} [\emph{pxin}]: клянусь в\\
0x68EA \textbf{txin} [\emph{txin}]: рафинирование\\
0x68F1 \textbf{bxin} [\emph{bxin}]: иммобилизация\\
0x68F2 \textbf{gxin} [\emph{gxin}]: завоевывая\\
0x68F3 \textbf{dxin} [\emph{dxin}]: тело человека\\
0x6901 \textbf{myan} [\emph{mjän}]: мой\\
0x6902 \textbf{kyan} [\emph{kjän}]: как\\
0x6904 \textbf{pyan} [\emph{pjän}]: люблю\\
0x6906 \textbf{nyan} [\emph{njän}]: мужской\\
0x6908 \textbf{hman} [\emph{ʰmän}]: меня\\
0x6909 \textbf{syan} [\emph{sjän}]: правда\\
0x690A \textbf{tyan} [\emph{tjän}]: дверь\\
0x690B \textbf{lyan} [\emph{ljän}]: длинный\\
0x690C \textbf{fyan} [\emph{fjän}]: сад\\
0x690D \textbf{cyan} [\emph{ʃjän}]: мир\\
0x690E \textbf{ryan} [\emph{rjän}]: перевод\\
0x6910 \textbf{hkan} [\emph{ʰkän}]: пространство контекст\\
0x6911 \textbf{byan} [\emph{bjän}]: ванна\\
0x6912 \textbf{gyan} [\emph{gjän}]: пот\\
0x6913 \textbf{dyan} [\emph{djän}]: яйцо\\
0x6914 \textbf{zyan} [\emph{zjän}]: весло\\
0x6915 \textbf{jyan} [\emph{ʒjän}]: жаргон\\
0x6916 \textbf{vyan} [\emph{vjän}]: рекламодатель\\
0x6918 \textbf{hyan} [\emph{ʰjän}]: зависимый пункт\\
0x6919 \textbf{qyan} [\emph{ŋjän}]: дальтонизм\\
0x691A \textbf{xyan} [\emph{xjän}]: радиан\\
0x6920 \textbf{hpan} [\emph{ʰpän}]: человек контекст\\
0x6921 \textbf{mwan} [\emph{mwän}]: смерть\\
0x6922 \textbf{kwan} [\emph{kwän}]: легкий\\
0x6924 \textbf{pwan} [\emph{pwän}]: весна\\
0x6926 \textbf{nwan} [\emph{nwän}]: молодой\\
0x6928 \textbf{hwan} [\emph{ʰwän}]: а также\\
0x6929 \textbf{swan} [\emph{swän}]: дела\\
0x692A \textbf{twan} [\emph{twän}]: душа\\
0x692B \textbf{lwan} [\emph{lwän}]: обед\\
0x692C \textbf{fwan} [\emph{fwän}]: посещение\\
0x692D \textbf{cwan} [\emph{ʃwän}]: надеяться\\
0x692E \textbf{rwan} [\emph{rwän}]: награда\\
0x6930 \textbf{hnan} [\emph{ʰnän}]: вы\\
0x6931 \textbf{bwan} [\emph{bwän}]: племянница\\
0x6932 \textbf{gwan} [\emph{gwän}]: глотать\\
0x6933 \textbf{dwan} [\emph{dwän}]: вешать\\
0x6934 \textbf{zwan} [\emph{zwän}]: происходящий\\
0x6935 \textbf{jwan} [\emph{ʒwän}]: проспект\\
0x6936 \textbf{vwan} [\emph{vwän}]: сетчатка\\
0x6939 \textbf{qwan} [\emph{ŋwän}]: часто спросил вопросов\\
0x693A \textbf{xwan} [\emph{xwän}]: размежевание\\
0x6942 \textbf{ksan} [\emph{ksän}]: наименее\\
0x6944 \textbf{psan} [\emph{psän}]: половина\\
0x6948 \textbf{hsan} [\emph{ʰsän}]: поверхность контекст\\
0x694A \textbf{tsan} [\emph{tsän}]: надлежащий имя существительное\\
0x6950 \textbf{htan} [\emph{ʰtän}]: интерьер контекст\\
0x6951 \textbf{bzan} [\emph{bzän}]: бунт\\
0x6952 \textbf{gzan} [\emph{gzän}]: гарантия\\
0x6953 \textbf{dzan} [\emph{dzän}]: наложить\\
0x6958 \textbf{hlan} [\emph{ʰlän}]: но\\
0x6960 \textbf{hfan} [\emph{ʰfän}]: найденный\\
0x6962 \textbf{klan} [\emph{klän}]: там\\
0x6964 \textbf{plan} [\emph{plän}]: играть\\
0x6968 \textbf{hcan} [\emph{ʰʃän}]: Здравствуйте\\
0x6969 \textbf{slan} [\emph{slän}]: следовать\\
0x696A \textbf{tlan} [\emph{tlän}]: Бог\\
0x696C \textbf{flan} [\emph{flän}]: наука\\
0x696D \textbf{clan} [\emph{ʃlän}]: тысяча\\
0x6970 \textbf{hran} [\emph{ʰrän}]: совесть\\
0x6971 \textbf{blan} [\emph{blän}]: жалоба\\
0x6972 \textbf{glan} [\emph{glän}]: галерея\\
0x6973 \textbf{dlan} [\emph{dlän}]: доллар\\
0x6974 \textbf{zlan} [\emph{zlän}]: зарабатывать\\
0x6975 \textbf{jlan} [\emph{ʒlän}]: поле\\
0x6976 \textbf{vlan} [\emph{vlän}]: прикрепление\\
0x697A \textbf{xlan} [\emph{xlän}]: глобализация\\
0x6982 \textbf{kfan} [\emph{kfän}]: поведение\\
0x6984 \textbf{pfan} [\emph{pfän}]: вероятность\\
0x6988 \textbf{hban} [\emph{ʰbän}]: носить\\
0x698A \textbf{tfan} [\emph{tfän}]: вдруг, внезапно\\
0x6990 \textbf{hgan} [\emph{ʰgän}]: скрывать\\
0x6991 \textbf{bvan} [\emph{bvän}]: произнесение\\
0x6992 \textbf{gvan} [\emph{gvän}]: злокачественный\\
0x6993 \textbf{dvan} [\emph{dvän}]: неодобрительно\\
0x6998 \textbf{hdan} [\emph{ʰdän}]: дьявол\\
0x69A0 \textbf{hzan} [\emph{ʰzän}]: содержащий\\
0x69A2 \textbf{kcan} [\emph{kʃän}]: считают,\\
0x69A4 \textbf{pcan} [\emph{pʃän}]: Деньги\\
0x69A8 \textbf{hjan} [\emph{ʰʒän}]: наказание\\
0x69AA \textbf{tcan} [\emph{tʃän}]: Другие\\
0x69B0 \textbf{hvan} [\emph{ʰvän}]: пожаротушение\\
0x69B1 \textbf{bjan} [\emph{bʒän}]: бродить\\
0x69B2 \textbf{gjan} [\emph{gʒän}]: явление\\
0x69B3 \textbf{djan} [\emph{dʒän}]: гигант\\
0x69C1 \textbf{mran} [\emph{mrän}]: кукуруза\\
0x69C2 \textbf{kran} [\emph{krän}]: корона\\
0x69C4 \textbf{pran} [\emph{prän}]: Нажмите\\
0x69C6 \textbf{nran} [\emph{nrän}]: редкий\\
0x69C8 \textbf{hqan} [\emph{ʰŋän}]: длина волны\\
0x69C9 \textbf{sran} [\emph{srän}]: похороны\\
0x69CA \textbf{tran} [\emph{trän}]: отчаяние\\
0x69CC \textbf{fran} [\emph{frän}]: летать\\
0x69CD \textbf{cran} [\emph{ʃrän}]: жертва\\
0x69D0 \textbf{hxan} [\emph{ʰxän}]: семинар\\
0x69D1 \textbf{bran} [\emph{brän}]: корзина\\
0x69D2 \textbf{gran} [\emph{grän}]: зыбь\\
0x69D3 \textbf{dran} [\emph{drän}]: администрация\\
0x69D4 \textbf{zran} [\emph{zrän}]: разделение\\
0x69D5 \textbf{jran} [\emph{ʒrän}]: врач хирург\\
0x69D6 \textbf{vran} [\emph{vrän}]: извилистый\\
0x69D9 \textbf{qran} [\emph{ŋrän}]: Коран\\
0x69DA \textbf{xran} [\emph{xrän}]: астронавт\\
0x69E2 \textbf{kxan} [\emph{kxän}]: картон\\
0x69E4 \textbf{pxan} [\emph{pxän}]: рекламный\\
0x69EA \textbf{txan} [\emph{txän}]: актиний\\
0x69F1 \textbf{bxan} [\emph{bxän}]: вылетающих\\
0x69F2 \textbf{gxan} [\emph{gxän}]: галактика Пол\\
0x69F3 \textbf{dxan} [\emph{dxän}]: выгрузка\\
0x6A01 \textbf{myun} [\emph{mjun}]: важность\\
0x6A02 \textbf{kyun} [\emph{kjun}]: угол\\
0x6A04 \textbf{pyun} [\emph{pjun}]: коричневый\\
0x6A06 \textbf{nyun} [\emph{njun}]: тепло\\
0x6A08 \textbf{hmun} [\emph{ʰmun}]: мягкий\\
0x6A09 \textbf{syun} [\emph{sjun}]: правительство\\
0x6A0A \textbf{tyun} [\emph{tjun}]: домен\\
0x6A0B \textbf{lyun} [\emph{ljun}]: ион\\
0x6A0C \textbf{fyun} [\emph{fjun}]: плавучесть\\
0x6A0D \textbf{cyun} [\emph{ʃjun}]: состояние\\
0x6A0E \textbf{ryun} [\emph{rjun}]: регулирование\\
0x6A10 \textbf{hkun} [\emph{ʰkun}]: Счет\\
0x6A11 \textbf{byun} [\emph{bjun}]: кланяясь\\
0x6A12 \textbf{gyun} [\emph{gjun}]: Гвиана\\
0x6A13 \textbf{dyun} [\emph{djun}]: диагностировать\\
0x6A15 \textbf{jyun} [\emph{ʒjun}]: остроугольный\\
0x6A16 \textbf{vyun} [\emph{vjun}]: ЕС\\
0x6A17 \textbf{ks6h} [\emph{ksəʰ}]: причинный\\
0x6A18 \textbf{hyun} [\emph{ʰjun}]: река\\
0x6A1A \textbf{xyun} [\emph{xjun}]: Сатурн\\
0x6A20 \textbf{hpun} [\emph{ʰpun}]: спросил\\
0x6A21 \textbf{mwun} [\emph{mwun}]: мода\\
0x6A22 \textbf{kwun} [\emph{kwun}]: июнь\\
0x6A24 \textbf{pwun} [\emph{pwun}]: пудинг\\
0x6A26 \textbf{nwun} [\emph{nwun}]: пунктуальный\\
0x6A28 \textbf{hwun} [\emph{ʰwun}]: золото\\
0x6A29 \textbf{swun} [\emph{swun}]: мыло\\
0x6A2A \textbf{twun} [\emph{twun}]: 9\\
0x6A2B \textbf{lwun} [\emph{lwun}]: синие джинсы\\
0x6A2C \textbf{fwun} [\emph{fwun}]: сильфон\\
0x6A2D \textbf{cwun} [\emph{ʃwun}]: жевание\\
0x6A2E \textbf{rwun} [\emph{rwun}]: критерий\\
0x6A30 \textbf{hnun} [\emph{ʰnun}]: нео\\
0x6A31 \textbf{bwun} [\emph{bwun}]: Бутан\\
0x6A32 \textbf{gwun} [\emph{gwun}]: гуанин\\
0x6A33 \textbf{dwun} [\emph{dwun}]: морщина\\
0x6A3A \textbf{xwun} [\emph{xwun}]: мутант\\
0x6A42 \textbf{ksun} [\emph{ksun}]: сжигание\\
0x6A44 \textbf{psun} [\emph{psun}]: предзнаменование\\
0x6A48 \textbf{hsun} [\emph{ʰsun}]: рука\\
0x6A4A \textbf{tsun} [\emph{tsun}]: деревня\\
0x6A50 \textbf{htun} [\emph{ʰtun}]: только\\
0x6A51 \textbf{bzun} [\emph{bzun}]: педантка\\
0x6A53 \textbf{dzun} [\emph{dzun}]: Не знаю\\
0x6A57 \textbf{ts6h} [\emph{tsəʰ}]: intransitive case, the subject of an intransitive verb.\\
0x6A58 \textbf{hlun} [\emph{ʰlun}]: включительно человек\\
0x6A60 \textbf{hfun} [\emph{ʰfun}]: гнев\\
0x6A62 \textbf{klun} [\emph{klun}]: накопление\\
0x6A64 \textbf{plun} [\emph{plun}]: синий\\
0x6A68 \textbf{hcun} [\emph{ʰʃun}]: слышать\\
0x6A69 \textbf{slun} [\emph{slun}]: приостановить\\
0x6A6A \textbf{tlun} [\emph{tlun}]: брюки\\
0x6A6C \textbf{flun} [\emph{flun}]: фтор\\
0x6A6D \textbf{clun} [\emph{ʃlun}]: дуб\\
0x6A70 \textbf{hrun} [\emph{ʰrun}]: рулон\\
0x6A71 \textbf{blun} [\emph{blun}]: воздушный шар\\
0x6A72 \textbf{glun} [\emph{glun}]: чеснок\\
0x6A73 \textbf{dlun} [\emph{dlun}]: загрязнение\\
0x6A74 \textbf{zlun} [\emph{zlun}]: инсулин\\
0x6A75 \textbf{jlun} [\emph{ʒlun}]: Юлиан\\
0x6A76 \textbf{vlun} [\emph{vlun}]: линия электропередачи\\
0x6A7A \textbf{xlun} [\emph{xlun}]: нерастворимый\\
0x6A82 \textbf{kfun} [\emph{kfun}]: размещение\\
0x6A84 \textbf{pfun} [\emph{pfun}]: частично\\
0x6A88 \textbf{hbun} [\emph{ʰbun}]: путаница\\
0x6A8A \textbf{tfun} [\emph{tfun}]: тоннель\\
0x6A90 \textbf{hgun} [\emph{ʰgun}]: ароматный\\
0x6A91 \textbf{bvun} [\emph{bvun}]: домовой\\
0x6A92 \textbf{gvun} [\emph{gvun}]: гальванизировать\\
0x6A93 \textbf{dvun} [\emph{dvun}]: нижеприведенный\\
0x6A98 \textbf{hdun} [\emph{ʰdun}]: единственное число\\
0x6AA0 \textbf{hzun} [\emph{ʰzun}]: упертый\\
0x6AA2 \textbf{kcun} [\emph{kʃun}]: отверстие\\
0x6AA4 \textbf{pcun} [\emph{pʃun}]: банка\\
0x6AA8 \textbf{hjun} [\emph{ʰʒun}]: ветреный\\
0x6AAA \textbf{tcun} [\emph{tʃun}]: сделанный\\
0x6AB0 \textbf{hvun} [\emph{ʰvun}]: бассейн\\
0x6AB1 \textbf{bjun} [\emph{bʒun}]: продавец книг\\
0x6AB2 \textbf{gjun} [\emph{gʒun}]: груминг\\
0x6AB3 \textbf{djun} [\emph{dʒun}]: человечный\\
0x6AC1 \textbf{mrun} [\emph{mrun}]: плотник\\
0x6AC2 \textbf{krun} [\emph{krun}]: полнота\\
0x6AC4 \textbf{prun} [\emph{prun}]: вращение\\
0x6AC6 \textbf{nrun} [\emph{nrun}]: авторизоваться\\
0x6AC8 \textbf{hqun} [\emph{ʰŋun}]: прерывистый\\
0x6AC9 \textbf{srun} [\emph{srun}]: хмуриться\\
0x6ACA \textbf{trun} [\emph{trun}]: посадка\\
0x6ACC \textbf{frun} [\emph{frun}]: гниль\\
0x6ACD \textbf{crun} [\emph{ʃrun}]: сокровище\\
0x6AD0 \textbf{hxun} [\emph{ʰxun}]: уговаривание\\
0x6AD1 \textbf{brun} [\emph{brun}]: туризм\\
0x6AD2 \textbf{grun} [\emph{grun}]: хвастовство\\
0x6AD3 \textbf{drun} [\emph{drun}]: неудовлетворенность\\
0x6AD4 \textbf{zrun} [\emph{zrun}]: Сардиния\\
0x6AD5 \textbf{jrun} [\emph{ʒrun}]: Камерун\\
0x6AD6 \textbf{vrun} [\emph{vrun}]: Мавритания\\
0x6AD9 \textbf{qrun} [\emph{ŋrun}]: украшенный\\
0x6ADA \textbf{xrun} [\emph{xrun}]: бесплодный\\
0x6AE2 \textbf{kxun} [\emph{kxun}]: рыжеватый\\
0x6AE4 \textbf{pxun} [\emph{pxun}]: плутоний\\
0x6AEA \textbf{txun} [\emph{txun}]: исследование Пол\\
0x6AF1 \textbf{bxun} [\emph{bxun}]: цветущий\\
0x6AF2 \textbf{gxun} [\emph{gxun}]: шестиугольник\\
0x6AF3 \textbf{dxun} [\emph{dxun}]: суточный\\
0x6B01 \textbf{myen} [\emph{mjen}]: по аналогии\\
0x6B02 \textbf{kyen} [\emph{kjen}]: несчастность\\
0x6B04 \textbf{pyen} [\emph{pjen}]: Приложения\\
0x6B06 \textbf{nyen} [\emph{njen}]: стон\\
0x6B08 \textbf{hmen} [\emph{ʰmen}]: смиренный\\
0x6B09 \textbf{syen} [\emph{sjen}]: сэр\\
0x6B0A \textbf{tyen} [\emph{tjen}]: общеизвестный\\
0x6B0B \textbf{lyen} [\emph{ljen}]: белье\\
0x6B0C \textbf{fyen} [\emph{fjen}]: фенхель\\
0x6B0D \textbf{cyen} [\emph{ʃjen}]: сессия\\
0x6B0E \textbf{ryen} [\emph{rjen}]: несмотря на\\
0x6B10 \textbf{hken} [\emph{ʰken}]: кабина\\
0x6B11 \textbf{byen} [\emph{bjen}]: катание на санях\\
0x6B12 \textbf{gyen} [\emph{gjen}]: ген\\
0x6B13 \textbf{dyen} [\emph{djen}]: индуктивность\\
0x6B14 \textbf{zyen} [\emph{zjen}]: ксенон\\
0x6B15 \textbf{jyen} [\emph{ʒjen}]: иена\\
0x6B16 \textbf{vyen} [\emph{vjen}]: ярмарочная площадь\\
0x6B18 \textbf{hyen} [\emph{ʰjen}]: облако\\
0x6B19 \textbf{qyen} [\emph{ŋjen}]: прорастание\\
0x6B1A \textbf{xyen} [\emph{xjen}]: взрастить\\
0x6B20 \textbf{hpen} [\emph{ʰpen}]: предварительно\\
0x6B21 \textbf{mwen} [\emph{mwen}]: мечтательность\\
0x6B22 \textbf{kwen} [\emph{kwen}]: эксклюзивный человек\\
0x6B24 \textbf{pwen} [\emph{pwen}]: пешка\\
0x6B26 \textbf{nwen} [\emph{nwen}]: флот\\
0x6B28 \textbf{hwen} [\emph{ʰwen}]: растительный Пол\\
0x6B29 \textbf{swen} [\emph{swen}]: сенат\\
0x6B2A \textbf{twen} [\emph{twen}]: ткать\\
0x6B2B \textbf{lwen} [\emph{lwen}]: легенда\\
0x6B2C \textbf{fwen} [\emph{fwen}]: филогенетическое\\
0x6B2D \textbf{cwen} [\emph{ʃwen}]: стадо\\
0x6B2E \textbf{rwen} [\emph{rwen}]: предотвращение\\
0x6B30 \textbf{hnen} [\emph{ʰnen}]: вечный\\
0x6B31 \textbf{bwen} [\emph{bwen}]: навязчивая идея\\
0x6B32 \textbf{gwen} [\emph{gwen}]: гранит\\
0x6B33 \textbf{dwen} [\emph{dwen}]: швейцар\\
0x6B34 \textbf{zwen} [\emph{zwen}]: неуместность\\
0x6B35 \textbf{jwen} [\emph{ʒwen}]: лежачий\\
0x6B36 \textbf{vwen} [\emph{vwen}]: непостоянство\\
0x6B39 \textbf{qwen} [\emph{ŋwen}]: Эстония\\
0x6B3A \textbf{xwen} [\emph{xwen}]: Нептун\\
0x6B42 \textbf{ksen} [\emph{ksen}]: кислород\\
0x6B44 \textbf{psen} [\emph{psen}]: пенсия\\
0x6B48 \textbf{hsen} [\emph{ʰsen}]: экспедиция\\
0x6B4A \textbf{tsen} [\emph{tsen}]: обработка\\
0x6B4F \textbf{sl6h} [\emph{sləʰ}]: цельсию\\
0x6B50 \textbf{hten} [\emph{ʰten}]: вчера\\
0x6B51 \textbf{bzen} [\emph{bzen}]: бензол\\
0x6B52 \textbf{gzen} [\emph{gzen}]: мыза\\
0x6B53 \textbf{dzen} [\emph{dzen}]: двенадцатиперстная кишка\\
0x6B57 \textbf{tl6h} [\emph{tləʰ}]: multal номер\\
0x6B58 \textbf{hlen} [\emph{ʰlen}]: лицензия\\
0x6B60 \textbf{hfen} [\emph{ʰfen}]: цент\\
0x6B62 \textbf{klen} [\emph{klen}]: колония\\
0x6B64 \textbf{plen} [\emph{plen}]: равновесие\\
0x6B68 \textbf{hcen} [\emph{ʰʃen}]: член\\
0x6B69 \textbf{slen} [\emph{slen}]: 11\\
0x6B6A \textbf{tlen} [\emph{tlen}]: талант\\
0x6B6C \textbf{flen} [\emph{flen}]: сгибать\\
0x6B6D \textbf{clen} [\emph{ʃlen}]: канал\\
0x6B70 \textbf{hren} [\emph{ʰren}]: речь контекст\\
0x6B71 \textbf{blen} [\emph{blen}]: блестящий\\
0x6B72 \textbf{glen} [\emph{glen}]: Гренландия\\
0x6B73 \textbf{dlen} [\emph{dlen}]: документация\\
0x6B74 \textbf{zlen} [\emph{zlen}]: межпланетный\\
0x6B75 \textbf{jlen} [\emph{ʒlen}]: местонахождение Пол\\
0x6B76 \textbf{vlen} [\emph{vlen}]: валентность\\
0x6B7A \textbf{xlen} [\emph{xlen}]: словесный имя существительное\\
0x6B82 \textbf{kfen} [\emph{kfen}]: понятия\\
0x6B84 \textbf{pfen} [\emph{pfen}]: проверка\\
0x6B88 \textbf{hben} [\emph{ʰben}]: выборы\\
0x6B8A \textbf{tfen} [\emph{tfen}]: защиты\\
0x6B90 \textbf{hgen} [\emph{ʰgen}]: агент\\
0x6B91 \textbf{bven} [\emph{bven}]: Ливан\\
0x6B92 \textbf{gven} [\emph{gven}]: карниз\\
0x6B93 \textbf{dven} [\emph{dven}]: саванна\\
0x6B98 \textbf{hden} [\emph{ʰden}]: решение\\
0x6BA0 \textbf{hzen} [\emph{ʰzen}]: плавать\\
0x6BA2 \textbf{kcen} [\emph{kʃen}]: девственница\\
0x6BA4 \textbf{pcen} [\emph{pʃen}]: президент\\
0x6BA8 \textbf{hjen} [\emph{ʰʒen}]: женский\\
0x6BAA \textbf{tcen} [\emph{tʃen}]: наблюдение\\
0x6BB0 \textbf{hven} [\emph{ʰven}]: контейнеры\\
0x6BB1 \textbf{bjen} [\emph{bʒen}]: пляжный\\
0x6BB2 \textbf{gjen} [\emph{gʒen}]: Аргентина\\
0x6BB3 \textbf{djen} [\emph{dʒen}]: Средиземное море\\
0x6BC1 \textbf{mren} [\emph{mren}]: памятный\\
0x6BC2 \textbf{kren} [\emph{kren}]: заместитель\\
0x6BC4 \textbf{pren} [\emph{pren}]: параллельно\\
0x6BC6 \textbf{nren} [\emph{nren}]: вдохновение\\
0x6BC8 \textbf{hqen} [\emph{ʰŋen}]: звезда Пол\\
0x6BC9 \textbf{sren} [\emph{sren}]: экспонат\\
0x6BCA \textbf{tren} [\emph{tren}]: Cегодня\\
0x6BCC \textbf{fren} [\emph{fren}]: перфорировать\\
0x6BCD \textbf{cren} [\emph{ʃren}]: сарай\\
0x6BD0 \textbf{hxen} [\emph{ʰxen}]: Перенаправление\\
0x6BD1 \textbf{bren} [\emph{bren}]: лабиринт\\
0x6BD2 \textbf{gren} [\emph{gren}]: планета\\
0x6BD3 \textbf{dren} [\emph{dren}]: демонстрировать\\
0x6BD4 \textbf{zren} [\emph{zren}]: века\\
0x6BD5 \textbf{jren} [\emph{ʒren}]: выравнивание\\
0x6BD6 \textbf{vren} [\emph{vren}]: вечнозеленый\\
0x6BD9 \textbf{qren} [\emph{ŋren}]: планетарный Пол\\
0x6BDA \textbf{xren} [\emph{xren}]: наушники\\
0x6BE2 \textbf{kxen} [\emph{kxen}]: общий Пол\\
0x6BE4 \textbf{pxen} [\emph{pxen}]: ведомственный\\
0x6BEA \textbf{txen} [\emph{txen}]: переложение\\
0x6BF1 \textbf{bxen} [\emph{bxen}]: навигация\\
0x6BF2 \textbf{gxen} [\emph{gxen}]: неодушевленный Пол\\
0x6BF3 \textbf{dxen} [\emph{dxen}]: Видео-конференция\\
0x6C01 \textbf{myon} [\emph{mjon}]: пневмония\\
0x6C02 \textbf{kyon} [\emph{kjon}]: эксплуатировать\\
0x6C04 \textbf{pyon} [\emph{pjon}]: пианино\\
0x6C06 \textbf{nyon} [\emph{njon}]: анонимный\\
0x6C08 \textbf{hmon} [\emph{ʰmon}]: миссия\\
0x6C09 \textbf{syon} [\emph{sjon}]: тяжелое дыхание\\
0x6C0A \textbf{tyon} [\emph{tjon}]: астрономия\\
0x6C0B \textbf{lyon} [\emph{ljon}]: язва\\
0x6C0C \textbf{fyon} [\emph{fjon}]: сродство\\
0x6C0D \textbf{cyon} [\emph{ʃjon}]: крошечный\\
0x6C0E \textbf{ryon} [\emph{rjon}]: рациональный\\
0x6C10 \textbf{hkon} [\emph{ʰkon}]: замок\\
0x6C11 \textbf{byon} [\emph{bjon}]: босния\\
0x6C12 \textbf{gyon} [\emph{gjon}]: связь\\
0x6C13 \textbf{dyon} [\emph{djon}]: приспособление\\
0x6C14 \textbf{zyon} [\emph{zjon}]: озон\\
0x6C15 \textbf{jyon} [\emph{ʒjon}]: Генуя\\
0x6C16 \textbf{vyon} [\emph{vjon}]: аэрация\\
0x6C18 \textbf{hyon} [\emph{ʰjon}]: тем временем\\
0x6C19 \textbf{qyon} [\emph{ŋjon}]: уравнения\\
0x6C1A \textbf{xyon} [\emph{xjon}]: пар\\
0x6C20 \textbf{hpon} [\emph{ʰpon}]: провинция\\
0x6C21 \textbf{mwon} [\emph{mwon}]: мандарин\\
0x6C22 \textbf{kwon} [\emph{kwon}]: Словарь\\
0x6C24 \textbf{pwon} [\emph{pwon}]: запрет\\
0x6C26 \textbf{nwon} [\emph{nwon}]: номинальный\\
0x6C28 \textbf{hwon} [\emph{ʰwon}]: приглашение\\
0x6C29 \textbf{swon} [\emph{swon}]: доброволец\\
0x6C2A \textbf{twon} [\emph{twon}]: орнамент\\
0x6C2B \textbf{lwon} [\emph{lwon}]: эволюционный\\
0x6C2C \textbf{fwon} [\emph{fwon}]: флуоресценция\\
0x6C2D \textbf{cwon} [\emph{ʃwon}]: овальный\\
0x6C2E \textbf{rwon} [\emph{rwon}]: напоминание\\
0x6C30 \textbf{hnon} [\emph{ʰnon}]: союз\\
0x6C31 \textbf{bwon} [\emph{bwon}]: фасоль\\
0x6C32 \textbf{gwon} [\emph{gwon}]: овес\\
0x6C33 \textbf{dwon} [\emph{dwon}]: диагональ\\
0x6C34 \textbf{zwon} [\emph{zwon}]: сонар\\
0x6C35 \textbf{jwon} [\emph{ʒwon}]: Джон\\
0x6C39 \textbf{qwon} [\emph{ŋwon}]: отчужденность\\
0x6C3A \textbf{xwon} [\emph{xwon}]: лебедь\\
0x6C42 \textbf{kson} [\emph{kson}]: акцент\\
0x6C44 \textbf{pson} [\emph{pson}]: непобедимый\\
0x6C48 \textbf{hson} [\emph{ʰson}]: помилование\\
0x6C4A \textbf{tson} [\emph{tson}]: учреждения\\
0x6C50 \textbf{hton} [\emph{ʰton}]: достаточно\\
0x6C51 \textbf{bzon} [\emph{bzon}]: бизон\\
0x6C52 \textbf{gzon} [\emph{gzon}]: казино\\
0x6C53 \textbf{dzon} [\emph{dzon}]: Амазонка\\
0x6C58 \textbf{hlon} [\emph{ʰlon}]: рассвет\\
0x6C5E \textbf{ke2} [\emph{ke˩}]: эр\\
0x6C60 \textbf{hfon} [\emph{ʰfon}]: мед\\
0x6C62 \textbf{klon} [\emph{klon}]: революция\\
0x6C64 \textbf{plon} [\emph{plon}]: выделение\\
0x6C68 \textbf{hcon} [\emph{ʰʃon}]: кухня\\
0x6C69 \textbf{slon} [\emph{slon}]: шитье\\
0x6C6A \textbf{tlon} [\emph{tlon}]: средство передвижения\\
0x6C6C \textbf{flon} [\emph{flon}]: не в сети\\
0x6C6D \textbf{clon} [\emph{ʃlon}]: онлайн\\
0x6C70 \textbf{hron} [\emph{ʰron}]: цвет\\
0x6C71 \textbf{blon} [\emph{blon}]: сходимость\\
0x6C72 \textbf{glon} [\emph{glon}]: галлон\\
0x6C73 \textbf{dlon} [\emph{dlon}]: двойной номер\\
0x6C74 \textbf{zlon} [\emph{zlon}]: вечность\\
0x6C75 \textbf{jlon} [\emph{ʒlon}]: двоеточие\\
0x6C76 \textbf{vlon} [\emph{vlon}]: одуванчик\\
0x6C7A \textbf{xlon} [\emph{xlon}]: мобильный телефон\\
0x6C82 \textbf{kfon} [\emph{kfon}]: кодирование\\
0x6C84 \textbf{pfon} [\emph{pfon}]: губка\\
0x6C88 \textbf{hbon} [\emph{ʰbon}]: разнообразный\\
0x6C8A \textbf{tfon} [\emph{tfon}]: дразнить\\
0x6C90 \textbf{hgon} [\emph{ʰgon}]: яма\\
0x6C91 \textbf{bvon} [\emph{bvon}]: садоводство\\
0x6C92 \textbf{gvon} [\emph{gvon}]: стягивание\\
0x6C93 \textbf{dvon} [\emph{dvon}]: бубен\\
0x6C98 \textbf{hdon} [\emph{ʰdon}]: жертвовать\\
0x6CA0 \textbf{hzon} [\emph{ʰzon}]: обезьяна\\
0x6CA2 \textbf{kcon} [\emph{kʃon}]: откровение\\
0x6CA4 \textbf{pcon} [\emph{pʃon}]: латунь\\
0x6CA8 \textbf{hjon} [\emph{ʰʒon}]: сексуальное\\
0x6CAA \textbf{tcon} [\emph{tʃon}]: полдень\\
0x6CB0 \textbf{hvon} [\emph{ʰvon}]: воздушные силы\\
0x6CB1 \textbf{bjon} [\emph{bʒon}]: бекон\\
0x6CB2 \textbf{gjon} [\emph{gʒon}]: однородный\\
0x6CB3 \textbf{djon} [\emph{dʒon}]: по умолчанию\\
0x6CC1 \textbf{mron} [\emph{mron}]: фургон\\
0x6CC2 \textbf{kron} [\emph{kron}]: углерод\\
0x6CC4 \textbf{pron} [\emph{pron}]: профессионал\\
0x6CC6 \textbf{nron} [\emph{nron}]: водород\\
0x6CC8 \textbf{hqon} [\emph{ʰŋon}]: орегано\\
0x6CC9 \textbf{sron} [\emph{sron}]: воин\\
0x6CCA \textbf{tron} [\emph{tron}]: оранжевый\\
0x6CCC \textbf{fron} [\emph{fron}]: героин\\
0x6CCD \textbf{cron} [\emph{ʃron}]: трепет\\
0x6CD0 \textbf{hxon} [\emph{ʰxon}]: середина весны\\
0x6CD1 \textbf{bron} [\emph{bron}]: эмбрион\\
0x6CD2 \textbf{gron} [\emph{gron}]: аргон\\
0x6CD3 \textbf{dron} [\emph{dron}]: Дракон\\
0x6CD4 \textbf{zron} [\emph{zron}]: нуль\\
0x6CD5 \textbf{jron} [\emph{ʒron}]: Иордания\\
0x6CD6 \textbf{vron} [\emph{vron}]: позвоночный\\
0x6CD9 \textbf{qron} [\emph{ŋron}]: гортань\\
0x6CDA \textbf{xron} [\emph{xron}]: Абстрактные Пол\\
0x6CDE \textbf{ne2} [\emph{ne˩}]: энный мощность\\
0x6CE2 \textbf{kxon} [\emph{kxon}]: гномический настроение\\
0x6CE4 \textbf{pxon} [\emph{pxon}]: недоступен\\
0x6CEA \textbf{txon} [\emph{txon}]: ругань\\
0x6CF1 \textbf{bxon} [\emph{bxon}]: беспрецедентный\\
0x6CF2 \textbf{gxon} [\emph{gxon}]: фрагментация\\
0x6CF3 \textbf{dxon} [\emph{dxon}]: подгорный\\
0x6D01 \textbf{my6n} [\emph{mjən}]: Армения\\
0x6D02 \textbf{ky6n} [\emph{kjən}]: кокаин\\
0x6D04 \textbf{py6n} [\emph{pjən}]: пион\\
0x6D06 \textbf{ny6n} [\emph{njən}]: антенна\\
0x6D08 \textbf{hm6n} [\emph{ʰmən}]: излучающий\\
0x6D09 \textbf{sy6n} [\emph{sjən}]: цемент\\
0x6D0A \textbf{ty6n} [\emph{tjən}]: касательный\\
0x6D0B \textbf{ly6n} [\emph{ljən}]: палестинец\\
0x6D0C \textbf{fy6n} [\emph{fjən}]: Софония\\
0x6D0E \textbf{ry6n} [\emph{rjən}]: иррациональный\\
0x6D10 \textbf{hk6n} [\emph{ʰkən}]: компилировать\\
0x6D11 \textbf{by6n} [\emph{bjən}]: обратное словообразование\\
0x6D13 \textbf{dy6n} [\emph{djən}]: оглушительный\\
0x6D16 \textbf{vy6n} [\emph{vjən}]: венский\\
0x6D18 \textbf{hy6n} [\emph{ʰjən}]: оживить Пол\\
0x6D19 \textbf{qy6n} [\emph{ŋjən}]: пасека\\
0x6D1A \textbf{xy6n} [\emph{xjən}]: Кения\\
0x6D20 \textbf{hp6n} [\emph{ʰpən}]: отправитель\\
0x6D21 \textbf{mw6n} [\emph{mwən}]: Македония\\
0x6D22 \textbf{kw6n} [\emph{kwən}]: кларнет\\
0x6D24 \textbf{pw6n} [\emph{pwən}]: планктон\\
0x6D26 \textbf{nw6n} [\emph{nwən}]: новорожденный\\
0x6D27 \textbf{pc6h} [\emph{pʃəʰ}]: испытание номер\\
0x6D28 \textbf{hw6n} [\emph{ʰwən}]: вихрь\\
0x6D29 \textbf{sw6n} [\emph{swən}]: санитария\\
0x6D2A \textbf{tw6n} [\emph{twən}]: Тайвань\\
0x6D2B \textbf{lw6n} [\emph{lwən}]: лантан\\
0x6D2E \textbf{rw6n} [\emph{rwən}]: изменчивый\\
0x6D30 \textbf{hn6n} [\emph{ʰnən}]: неон\\
0x6D31 \textbf{bw6n} [\emph{bwən}]: соболезнование\\
0x6D33 \textbf{dw6n} [\emph{dwən}]: опреснение\\
0x6D3A \textbf{xw6n} [\emph{xwən}]: антиген\\
0x6D42 \textbf{ks6n} [\emph{ksən}]: рабочая станция\\
0x6D44 \textbf{ps6n} [\emph{psən}]: испанский\\
0x6D48 \textbf{hs6n} [\emph{ʰsən}]: продавцы\\
0x6D4A \textbf{ts6n} [\emph{tsən}]: растворение\\
0x6D50 \textbf{ht6n} [\emph{ʰtən}]: Терминал\\
0x6D57 \textbf{tc6h} [\emph{tʃəʰ}]: целое число\\
0x6D58 \textbf{hl6n} [\emph{ʰlən}]: контакт\\
0x6D60 \textbf{hf6n} [\emph{ʰfən}]: Афганистан\\
0x6D62 \textbf{kl6n} [\emph{klən}]: дыня\\
0x6D64 \textbf{pl6n} [\emph{plən}]: Польша\\
0x6D68 \textbf{hc6n} [\emph{ʰʃən}]: очевидность\\
0x6D69 \textbf{sl6n} [\emph{slən}]: водопровод\\
0x6D6A \textbf{tl6n} [\emph{tlən}]: центральный обработка Ед. изм\\
0x6D6C \textbf{fl6n} [\emph{flən}]: Финляндия\\
0x6D6D \textbf{cl6n} [\emph{ʃlən}]: роды\\
0x6D70 \textbf{hr6n} [\emph{ʰrən}]: рацион\\
0x6D71 \textbf{bl6n} [\emph{blən}]: Бруней\\
0x6D72 \textbf{gl6n} [\emph{glən}]: стежка\\
0x6D73 \textbf{dl6n} [\emph{dlən}]: Ирландия\\
0x6D74 \textbf{zl6n} [\emph{zlən}]: зеландии\\
0x6D75 \textbf{jl6n} [\emph{ʒlən}]: смягчающий\\
0x6D76 \textbf{vl6n} [\emph{vlən}]: валлонский\\
0x6D7A \textbf{xl6n} [\emph{xlən}]: нить\\
0x6D82 \textbf{kf6n} [\emph{kfən}]: оксюморон\\
0x6D84 \textbf{pf6n} [\emph{pfən}]: опыление\\
0x6D88 \textbf{hb6n} [\emph{ʰbən}]: дистрибутивный местоимение\\
0x6D8A \textbf{tf6n} [\emph{tfən}]: тонер\\
0x6D90 \textbf{hg6n} [\emph{ʰgən}]: Японский\\
0x6D93 \textbf{dv6n} [\emph{dvən}]: девонский\\
0x6D98 \textbf{hd6n} [\emph{ʰdən}]: знаменатель\\
0x6DA0 \textbf{hz6n} [\emph{ʰzən}]: Тасмания\\
0x6DA2 \textbf{kc6n} [\emph{kʃən}]: уполномоченный\\
0x6DA4 \textbf{pc6n} [\emph{pʃən}]: испытание номер\\
0x6DA8 \textbf{hj6n} [\emph{ʰʒən}]: не человек\\
0x6DAA \textbf{tc6n} [\emph{tʃən}]: художественный Пол\\
0x6DB0 \textbf{hv6n} [\emph{ʰvən}]: лавина\\
0x6DB3 \textbf{dj6n} [\emph{dʒən}]: бедность\\
0x6DC1 \textbf{mr6n} [\emph{mrən}]: меридиан\\
0x6DC2 \textbf{kr6n} [\emph{krən}]: корнеплоды\\
0x6DC4 \textbf{pr6n} [\emph{prən}]: Я побежал\\
0x6DC6 \textbf{nr6n} [\emph{nrən}]: нейрон\\
0x6DC8 \textbf{hq6n} [\emph{ʰŋən}]: наперсток\\
0x6DC9 \textbf{sr6n} [\emph{srən}]: созревание\\
0x6DCA \textbf{tr6n} [\emph{trən}]: стерео\\
0x6DCC \textbf{fr6n} [\emph{frən}]: выветривание\\
0x6DCD \textbf{cr6n} [\emph{ʃrən}]: титрование\\
0x6DD0 \textbf{hx6n} [\emph{ʰxən}]: предварительно сегодняшний\\
0x6DD1 \textbf{br6n} [\emph{brən}]: шершень\\
0x6DD2 \textbf{gr6n} [\emph{grən}]: элерон\\
0x6DD3 \textbf{dr6n} [\emph{drən}]: эндокринный\\
0x6DD4 \textbf{zr6n} [\emph{zrən}]: Азербайджан\\
0x6DD5 \textbf{jr6n} [\emph{ʒrən}]: маргарин\\
0x6DD6 \textbf{vr6n} [\emph{vrən}]: благоговейный\\
0x6DD9 \textbf{qr6n} [\emph{ŋrən}]: андеррайтер\\
0x6DDA \textbf{xr6n} [\emph{xrən}]: аллитерация\\
0x6DE2 \textbf{kx6n} [\emph{kxən}]: каштан\\
0x6DE4 \textbf{px6n} [\emph{pxən}]: плацента\\
0x6DEA \textbf{tx6n} [\emph{txən}]: испарение\\
0x6DF1 \textbf{bx6n} [\emph{bxən}]: шоры\\
0x6DF2 \textbf{gx6n} [\emph{gxən}]: Гана\\
0x6DF3 \textbf{dx6n} [\emph{dxən}]: приручение\\
0x6E37 \textbf{nr6h} [\emph{nrəʰ}]: пациент вызывать\\
0x6E4F \textbf{sr6h} [\emph{srəʰ}]: abstract gender, for making new genders, for example tlansr6h would be the gender of the one infinite creator God.\\
0x6E57 \textbf{tr6h} [\emph{trəʰ}]: движение вверх по течению\\
0x6E67 \textbf{fr6h} [\emph{frəʰ}]: ferous\\
\subsection{ 7 }
0x7001 \textbf{myi7n} [\emph{mji˥n}]: неправильное название\\
0x7002 \textbf{kyi7n} [\emph{kji˥n}]: ясновидение\\
0x7004 \textbf{pyi7n} [\emph{pji˥n}]: штифты\\
0x7006 \textbf{nyi7n} [\emph{nji˥n}]: ненароком\\
0x7008 \textbf{hmi7n} [\emph{ʰmi˥n}]: Милан\\
0x7009 \textbf{syi7n} [\emph{sji˥n}]: жестикуляция\\
0x700A \textbf{tyi7n} [\emph{tji˥n}]: диоксид титана\\
0x700B \textbf{lyi7n} [\emph{lji˥n}]: Колыбельная\\
0x700C \textbf{fyi7n} [\emph{fji˥n}]: пятьдесят один\\
0x700D \textbf{cyi7n} [\emph{ʃji˥n}]: китайский квартал\\
0x700E \textbf{ryi7n} [\emph{rji˥n}]: прихожанка\\
0x7010 \textbf{hki7n} [\emph{ʰki˥n}]: принимая во внимание\\
0x7011 \textbf{byi7n} [\emph{bji˥n}]: травля\\
0x7012 \textbf{gyi7n} [\emph{gji˥n}]: глицин\\
0x7013 \textbf{dyi7n} [\emph{dji˥n}]: два обрезной\\
0x7016 \textbf{vyi7n} [\emph{vji˥n}]: винил\\
0x7018 \textbf{hyi7n} [\emph{ʰji˥n}]: торжественность\\
0x7019 \textbf{qyi7n} [\emph{ŋji˥n}]: амин\\
0x701A \textbf{xyi7n} [\emph{xji˥n}]: гиацинт\\
0x7020 \textbf{hpi7n} [\emph{ʰpi˥n}]: нетерпение\\
0x7021 \textbf{mwi7n} [\emph{mwi˥n}]: мандолина\\
0x7022 \textbf{kwi7n} [\emph{kwi˥n}]: хинин\\
0x7024 \textbf{pwi7n} [\emph{pwi˥n}]: восемьдесят девять\\
0x7026 \textbf{nwi7n} [\emph{nwi˥n}]: монахиня\\
0x7028 \textbf{hwi7n} [\emph{ʰwi˥n}]: пятьдесят пять\\
0x7029 \textbf{swi7n} [\emph{swi˥n}]: сорок девять\\
0x702A \textbf{twi7n} [\emph{twi˥n}]: пятьдесят три\\
0x702B \textbf{lwi7n} [\emph{lwi˥n}]: миллионер\\
0x702D \textbf{cwi7n} [\emph{ʃwi˥n}]: молодой лебедь\\
0x702E \textbf{rwi7n} [\emph{rwi˥n}]: радужный\\
0x7030 \textbf{hni7n} [\emph{ʰni˥n}]: инбридинг\\
0x7032 \textbf{gwi7n} [\emph{gwi˥n}]: глубокая вода\\
0x7033 \textbf{dwi7n} [\emph{dwi˥n}]: фрейлин\\
0x703A \textbf{xwi7n} [\emph{xwi˥n}]: дедовщина\\
0x7042 \textbf{ksi7n} [\emph{ksi˥n}]: консалтинг\\
0x7044 \textbf{psi7n} [\emph{psi˥n}]: пасторат\\
0x7048 \textbf{hsi7n} [\emph{ʰsi˥n}]: самостоятельного полива\\
0x704A \textbf{tsi7n} [\emph{tsi˥n}]: семьдесят пять\\
0x7050 \textbf{hti7n} [\emph{ʰti˥n}]: турбина\\
0x7053 \textbf{dzi7n} [\emph{dzi˥n}]: гедонистический\\
0x7058 \textbf{hli7n} [\emph{ʰli˥n}]: сверкающих\\
0x7060 \textbf{hfi7n} [\emph{ʰfi˥n}]: финн\\
0x7062 \textbf{kli7n} [\emph{kli˥n}]: кастильский\\
0x7064 \textbf{pli7n} [\emph{pli˥n}]: Пенсильвания\\
0x7068 \textbf{hci7n} [\emph{ʰʃi˥n}]: нищенство\\
0x7069 \textbf{sli7n} [\emph{sli˥n}]: властители\\
0x706A \textbf{tli7n} [\emph{tli˥n}]: навязанный\\
0x706C \textbf{fli7n} [\emph{fli˥n}]: филология\\
0x706D \textbf{cli7n} [\emph{ʃli˥n}]: халдейский\\
0x7070 \textbf{hri7n} [\emph{ʰri˥n}]: Виргиния\\
0x7071 \textbf{bli7n} [\emph{bli˥n}]: изобличают\\
0x7072 \textbf{gli7n} [\emph{gli˥n}]: продолжительность жизни\\
0x7073 \textbf{dli7n} [\emph{dli˥n}]: неплатежеспособность\\
0x7074 \textbf{zli7n} [\emph{zli˥n}]: вазелин\\
0x7075 \textbf{jli7n} [\emph{ʒli˥n}]: женоненавистничество\\
0x7076 \textbf{vli7n} [\emph{vli˥n}]: ватерлиния\\
0x707A \textbf{xli7n} [\emph{xli˥n}]: выполнимый\\
0x7082 \textbf{kfi7n} [\emph{kfi˥n}]: кодифицированный\\
0x7084 \textbf{pfi7n} [\emph{pfi˥n}]: отделимый\\
0x7088 \textbf{hbi7n} [\emph{ʰbi˥n}]: Британия\\
0x708A \textbf{tfi7n} [\emph{tfi˥n}]: объединение\\
0x7090 \textbf{hgi7n} [\emph{ʰgi˥n}]: хихикающий\\
0x7098 \textbf{hdi7n} [\emph{ʰdi˥n}]: демпинг\\
0x70A0 \textbf{hzi7n} [\emph{ʰzi˥n}]: осматривает\\
0x70A2 \textbf{kci7n} [\emph{kʃi˥n}]: комиссии\\
0x70A4 \textbf{pci7n} [\emph{pʃi˥n}]: плиоцен\\
0x70A8 \textbf{hji7n} [\emph{ʰʒi˥n}]: Женщины мира\\
0x70AA \textbf{tci7n} [\emph{tʃi˥n}]: вздрогнуть\\
0x70B0 \textbf{hvi7n} [\emph{ʰvi˥n}]: исполненный по обету\\
0x70B3 \textbf{dji7n} [\emph{dʒi˥n}]: управительница\\
0x70C1 \textbf{mri7n} [\emph{mri˥n}]: марсианин\\
0x70C2 \textbf{kri7n} [\emph{kri˥n}]: выгравирован\\
0x70C4 \textbf{pri7n} [\emph{pri˥n}]: феи\\
0x70C6 \textbf{nri7n} [\emph{nri˥n}]: дочери\\
0x70C8 \textbf{hqi7n} [\emph{ʰŋi˥n}]: неслышный\\
0x70C9 \textbf{sri7n} [\emph{sri˥n}]: ацетилен\\
0x70CA \textbf{tri7n} [\emph{tri˥n}]: эфирный\\
0x70CC \textbf{fri7n} [\emph{fri˥n}]: обезображивание\\
0x70CD \textbf{cri7n} [\emph{ʃri˥n}]: продувка\\
0x70D0 \textbf{hxi7n} [\emph{ʰxi˥n}]: хозяйка\\
0x70D1 \textbf{bri7n} [\emph{bri˥n}]: аберраций\\
0x70D2 \textbf{gri7n} [\emph{gri˥n}]: подстрочный\\
0x70D3 \textbf{dri7n} [\emph{dri˥n}]: недостойный\\
0x70D4 \textbf{zri7n} [\emph{zri˥n}]: петух малиновка\\
0x70D5 \textbf{jri7n} [\emph{ʒri˥n}]: ИНГ\\
0x70D6 \textbf{vri7n} [\emph{vri˥n}]: активизация\\
0x70D9 \textbf{qri7n} [\emph{ŋri˥n}]: мариновать\\
0x70DA \textbf{xri7n} [\emph{xri˥n}]: трель\\
0x70E2 \textbf{kxi7n} [\emph{kxi˥n}]: приобретение\\
0x70E4 \textbf{pxi7n} [\emph{pxi˥n}]: прозвище\\
0x70EA \textbf{txi7n} [\emph{txi˥n}]: похвальный\\
0x70F1 \textbf{bxi7n} [\emph{bxi˥n}]: бронхит\\
0x70F2 \textbf{gxi7n} [\emph{gxi˥n}]: глубокая тарелка\\
0x70F3 \textbf{dxi7n} [\emph{dxi˥n}]: двухпартийный\\
0x7101 \textbf{mya7n} [\emph{mjä˥n}]: унижать\\
0x7102 \textbf{kya7n} [\emph{kjä˥n}]: канон\\
0x7104 \textbf{pya7n} [\emph{pjä˥n}]: восемьдесят один\\
0x7106 \textbf{nya7n} [\emph{njä˥n}]: недисциплинированный\\
0x7108 \textbf{hma7n} [\emph{ʰmä˥n}]: грот\\
0x7109 \textbf{sya7n} [\emph{sjä˥n}]: разведка\\
0x710A \textbf{tya7n} [\emph{tjä˥n}]: неизвестности\\
0x710B \textbf{lya7n} [\emph{ljä˥n}]: ладан\\
0x710C \textbf{fya7n} [\emph{fjä˥n}]: полугодовой\\
0x710D \textbf{cya7n} [\emph{ʃjä˥n}]: обеднять\\
0x710E \textbf{rya7n} [\emph{rjä˥n}]: арийский\\
0x7110 \textbf{hka7n} [\emph{ʰkä˥n}]: очки\\
0x7111 \textbf{bya7n} [\emph{bjä˥n}]: бармен\\
0x7112 \textbf{gya7n} [\emph{gjä˥n}]: болеутоляющий\\
0x7113 \textbf{dya7n} [\emph{djä˥n}]: адамантин\\
0x7115 \textbf{jya7n} [\emph{ʒjä˥n}]: водиться\\
0x7116 \textbf{vya7n} [\emph{vjä˥n}]: руки меня вниз\\
0x7118 \textbf{hya7n} [\emph{ʰjä˥n}]: ослепительный\\
0x7119 \textbf{qya7n} [\emph{ŋjä˥n}]: стенокардия\\
0x711A \textbf{xya7n} [\emph{xjä˥n}]: желтуха\\
0x7120 \textbf{hpa7n} [\emph{ʰpä˥n}]: фунтов стерлингов\\
0x7121 \textbf{mwa7n} [\emph{mwä˥n}]: паутина\\
0x7122 \textbf{kwa7n} [\emph{kwä˥n}]: консервирование\\
0x7124 \textbf{pwa7n} [\emph{pwä˥n}]: провансальский\\
0x7126 \textbf{nwa7n} [\emph{nwä˥n}]: размотанный\\
0x7128 \textbf{hwa7n} [\emph{ʰwä˥n}]: всемогущий\\
0x7129 \textbf{swa7n} [\emph{swä˥n}]: каламбур\\
0x712A \textbf{twa7n} [\emph{twä˥n}]: желтовато-коричневый\\
0x712B \textbf{lwa7n} [\emph{lwä˥n}]: кальвинист\\
0x712E \textbf{rwa7n} [\emph{rwä˥n}]: Рангун\\
0x7130 \textbf{hna7n} [\emph{ʰnä˥n}]: без углерода\\
0x7131 \textbf{bwa7n} [\emph{bwä˥n}]: скотобойня\\
0x7133 \textbf{dwa7n} [\emph{dwä˥n}]: безумная\\
0x7134 \textbf{zwa7n} [\emph{zwä˥n}]: вступить в повторный брак\\
0x713A \textbf{xwa7n} [\emph{xwä˥n}]: Гавана\\
0x7142 \textbf{ksa7n} [\emph{ksä˥n}]: высылка\\
0x7144 \textbf{psa7n} [\emph{psä˥n}]: рассеянно\\
0x7148 \textbf{hsa7n} [\emph{ʰsä˥n}]: поместья\\
0x714A \textbf{tsa7n} [\emph{tsä˥n}]: государственные мужи\\
0x7150 \textbf{hta7n} [\emph{ʰtä˥n}]: вырубают и выжигают\\
0x7153 \textbf{dza7n} [\emph{dzä˥n}]: средний мозг\\
0x7158 \textbf{hla7n} [\emph{ʰlä˥n}]: медленное высвобождение\\
0x7160 \textbf{hfa7n} [\emph{ʰfä˥n}]: раздосадованный\\
0x7162 \textbf{kla7n} [\emph{klä˥n}]: орлиный\\
0x7164 \textbf{pla7n} [\emph{plä˥n}]: Лапландия\\
0x7168 \textbf{hca7n} [\emph{ʰʃä˥n}]: сказывают\\
0x7169 \textbf{sla7n} [\emph{slä˥n}]: старейшины\\
0x716A \textbf{tla7n} [\emph{tlä˥n}]: фонарь\\
0x716C \textbf{fla7n} [\emph{flä˥n}]: клык\\
0x716D \textbf{cla7n} [\emph{ʃlä˥n}]: проститутка\\
0x7170 \textbf{hra7n} [\emph{ʰrä˥n}]: романс\\
0x7171 \textbf{bla7n} [\emph{blä˥n}]: Вавилон\\
0x7172 \textbf{gla7n} [\emph{glä˥n}]: молочник\\
0x7173 \textbf{dla7n} [\emph{dlä˥n}]: след\\
0x7174 \textbf{zla7n} [\emph{zlä˥n}]: позывной\\
0x7175 \textbf{jla7n} [\emph{ʒlä˥n}]: лопатой носатый\\
0x7176 \textbf{vla7n} [\emph{vlä˥n}]: овуляция\\
0x717A \textbf{xla7n} [\emph{xlä˥n}]: гало\\
0x7182 \textbf{kfa7n} [\emph{kfä˥n}]: камуфляж\\
0x7184 \textbf{pfa7n} [\emph{pfä˥n}]: восемьдесят пять\\
0x7188 \textbf{hba7n} [\emph{ʰbä˥n}]: смело\\
0x718A \textbf{tfa7n} [\emph{tfä˥n}]: назидательный\\
0x7190 \textbf{hga7n} [\emph{ʰgä˥n}]: безумие\\
0x7191 \textbf{bva7n} [\emph{bvä˥n}]: двухсотлетие\\
0x7193 \textbf{dva7n} [\emph{dvä˥n}]: двухстороннее\\
0x7198 \textbf{hda7n} [\emph{ʰdä˥n}]: водосточная труба\\
0x71A0 \textbf{hza7n} [\emph{ʰzä˥n}]: византийский\\
0x71A2 \textbf{kca7n} [\emph{kʃä˥n}]: абсент\\
0x71A4 \textbf{pca7n} [\emph{pʃä˥n}]: протестант\\
0x71A8 \textbf{hja7n} [\emph{ʰʒä˥n}]: смачивание\\
0x71AA \textbf{tca7n} [\emph{tʃä˥n}]: тридцать девять\\
0x71B0 \textbf{hva7n} [\emph{ʰvä˥n}]: люблю делать\\
0x71B3 \textbf{dja7n} [\emph{dʒä˥n}]: Ахан\\
0x71C1 \textbf{mra7n} [\emph{mrä˥n}]: мама в закон\\
0x71C2 \textbf{kra7n} [\emph{krä˥n}]: арки\\
0x71C4 \textbf{pra7n} [\emph{prä˥n}]: шторы\\
0x71C6 \textbf{nra7n} [\emph{nrä˥n}]: неописуемый\\
0x71C8 \textbf{hqa7n} [\emph{ʰŋä˥n}]: мечтатель\\
0x71C9 \textbf{sra7n} [\emph{srä˥n}]: неотделимый\\
0x71CA \textbf{tra7n} [\emph{trä˥n}]: невыносимый\\
0x71CC \textbf{fra7n} [\emph{frä˥n}]: франк\\
0x71CD \textbf{cra7n} [\emph{ʃrä˥n}]: детский сад\\
0x71D0 \textbf{hxa7n} [\emph{ʰxä˥n}]: монахи\\
0x71D1 \textbf{bra7n} [\emph{brä˥n}]: баронесса\\
0x71D2 \textbf{gra7n} [\emph{grä˥n}]: наградные\\
0x71D3 \textbf{dra7n} [\emph{drä˥n}]: датчанин\\
0x71D4 \textbf{zra7n} [\emph{zrä˥n}]: сад роз\\
0x71D5 \textbf{jra7n} [\emph{ʒrä˥n}]: кесарево сечение\\
0x71D6 \textbf{vra7n} [\emph{vrä˥n}]: валериана\\
0x71D9 \textbf{qra7n} [\emph{ŋrä˥n}]: бикарбонат\\
0x71DA \textbf{xra7n} [\emph{xrä˥n}]: кузнец\\
0x71E2 \textbf{kxa7n} [\emph{kxä˥n}]: карабин\\
0x71E4 \textbf{pxa7n} [\emph{pxä˥n}]: платонический\\
0x71EA \textbf{txa7n} [\emph{txä˥n}]: палатки\\
0x71F1 \textbf{bxa7n} [\emph{bxä˥n}]: Абиссиния\\
0x71F2 \textbf{gxa7n} [\emph{gxä˥n}]: сыворотка\\
0x71F3 \textbf{dxa7n} [\emph{dxä˥n}]: середине воздуха\\
0x7206 \textbf{nyu7n} [\emph{nju˥n}]: нет нет\\
0x7208 \textbf{hmu7n} [\emph{ʰmu˥n}]: лань\\
0x7209 \textbf{syu7n} [\emph{sju˥n}]: укус ядовитой змеи\\
0x720A \textbf{tyu7n} [\emph{tju˥n}]: отнятие ребенка от груди\\
0x720B \textbf{lyu7n} [\emph{lju˥n}]: выпускники\\
0x720E \textbf{ryu7n} [\emph{rju˥n}]: Веселая карусель\\
0x7210 \textbf{hku7n} [\emph{ʰku˥n}]: никуда не годится\\
0x7211 \textbf{byu7n} [\emph{bju˥n}]: бубонный\\
0x7218 \textbf{hyu7n} [\emph{ʰju˥n}]: еврейка\\
0x7220 \textbf{hpu7n} [\emph{ʰpu˥n}]: пунктуация\\
0x7221 \textbf{mwu7n} [\emph{mwu˥n}]: муэдзин\\
0x7224 \textbf{pwu7n} [\emph{pwu˥n}]: предварительная оплата\\
0x7226 \textbf{nwu7n} [\emph{nwu˥n}]: Не забывай меня\\
0x7228 \textbf{hwu7n} [\emph{ʰwu˥n}]: тепловатый\\
0x722A \textbf{twu7n} [\emph{twu˥n}]: мучнистая горлышком\\
0x722E \textbf{rwu7n} [\emph{rwu˥n}]: руна\\
0x7230 \textbf{hnu7n} [\emph{ʰnu˥n}]: Ньюфаундленд\\
0x7242 \textbf{ksu7n} [\emph{ksu˥n}]: кабестан\\
0x7244 \textbf{psu7n} [\emph{psu˥n}]: плати как сможешь\\
0x7248 \textbf{hsu7n} [\emph{ʰsu˥n}]: двадцать девять\\
0x724A \textbf{tsu7n} [\emph{tsu˥n}]: правнук\\
0x7250 \textbf{htu7n} [\emph{ʰtu˥n}]: взимание\\
0x7258 \textbf{hlu7n} [\emph{ʰlu˥n}]: вольная\\
0x7260 \textbf{hfu7n} [\emph{ʰfu˥n}]: листать\\
0x7268 \textbf{hcu7n} [\emph{ʰʃu˥n}]: хунта\\
0x7269 \textbf{slu7n} [\emph{slu˥n}]: суб-лейтенант\\
0x726A \textbf{tlu7n} [\emph{tlu˥n}]: щекотание\\
0x7270 \textbf{hru7n} [\emph{ʰru˥n}]: лучник\\
0x7271 \textbf{blu7n} [\emph{blu˥n}]: логический\\
0x7288 \textbf{hbu7n} [\emph{ʰbu˥n}]: черное дерево\\
0x7290 \textbf{hgu7n} [\emph{ʰgu˥n}]: борзая\\
0x7298 \textbf{hdu7n} [\emph{ʰdu˥n}]: уборщица\\
0x72C1 \textbf{mru7n} [\emph{mru˥n}]: кормушка\\
0x72C2 \textbf{kru7n} [\emph{kru˥n}]: Контратенор\\
0x72C4 \textbf{pru7n} [\emph{pru˥n}]: перестрахование\\
0x72C6 \textbf{nru7n} [\emph{nru˥n}]: интонировать\\
0x72C8 \textbf{hqu7n} [\emph{ʰŋu˥n}]: недоедание\\
0x72C9 \textbf{sru7n} [\emph{sru˥n}]: мастурбация\\
0x72CA \textbf{tru7n} [\emph{tru˥n}]: Тюрингия\\
0x72D0 \textbf{hxu7n} [\emph{ʰxu˥n}]: Беззаботная\\
0x72D1 \textbf{bru7n} [\emph{bru˥n}]: бумеранг\\
0x72D3 \textbf{dru7n} [\emph{dru˥n}]: слуховое окно\\
0x72DA \textbf{xru7n} [\emph{xru˥n}]: гидразин\\
0x72E2 \textbf{kxu7n} [\emph{kxu˥n}]: оптовый торговец\\
0x72EA \textbf{txu7n} [\emph{txu˥n}]: австро-венгерской\\
0x7301 \textbf{mye7n} [\emph{mje˥n}]: замуровывать\\
0x7302 \textbf{kye7n} [\emph{kje˥n}]: Аккредитующая\\
0x7304 \textbf{pye7n} [\emph{pje˥n}]: пенсне\\
0x7306 \textbf{nye7n} [\emph{nje˥n}]: индексация\\
0x7308 \textbf{hme7n} [\emph{ʰme˥n}]: неохотно\\
0x7309 \textbf{sye7n} [\emph{sje˥n}]: мессианский\\
0x730A \textbf{tye7n} [\emph{tje˥n}]: феерия\\
0x730B \textbf{lye7n} [\emph{lje˥n}]: лецитин\\
0x730E \textbf{rye7n} [\emph{rje˥n}]: зайца крикливый\\
0x7310 \textbf{hke7n} [\emph{ʰke˥n}]: имущественный\\
0x7318 \textbf{hye7n} [\emph{ʰje˥n}]: ан\\
0x731A \textbf{xye7n} [\emph{xje˥n}]: семиугольник\\
0x7320 \textbf{hpe7n} [\emph{ʰpe˥n}]: версии\\
0x7324 \textbf{pwe7n} [\emph{pwe˥n}]: пентан\\
0x7326 \textbf{nwe7n} [\emph{nwe˥n}]: антропоцентрический\\
0x7328 \textbf{hwe7n} [\emph{ʰwe˥n}]: пятьдесят четыре\\
0x7329 \textbf{swe7n} [\emph{swe˥n}]: продавщица\\
0x732A \textbf{twe7n} [\emph{twe˥n}]: гетеродинный\\
0x732B \textbf{lwe7n} [\emph{lwe˥n}]: телефонист\\
0x7330 \textbf{hne7n} [\emph{ʰne˥n}]: карликовый\\
0x7331 \textbf{bwe7n} [\emph{bwe˥n}]: боцман\\
0x733A \textbf{xwe7n} [\emph{xwe˥n}]: тяжеловес\\
0x7342 \textbf{kse7n} [\emph{kse˥n}]: казеин\\
0x7344 \textbf{pse7n} [\emph{pse˥n}]: пресс-банда\\
0x7348 \textbf{hse7n} [\emph{ʰse˥n}]: невод\\
0x734A \textbf{tse7n} [\emph{tse˥n}]: девяносто пять\\
0x7350 \textbf{hte7n} [\emph{ʰte˥n}]: Афины\\
0x7358 \textbf{hle7n} [\emph{ʰle˥n}]: цепи\\
0x7360 \textbf{hfe7n} [\emph{ʰfe˥n}]: фехтовальщик\\
0x7362 \textbf{kle7n} [\emph{kle˥n}]: конкатенация\\
0x7364 \textbf{ple7n} [\emph{ple˥n}]: перигелий\\
0x7368 \textbf{hce7n} [\emph{ʰʃe˥n}]: китолов\\
0x7369 \textbf{sle7n} [\emph{sle˥n}]: столетний\\
0x736A \textbf{tle7n} [\emph{tle˥n}]: Трансильвания\\
0x736C \textbf{fle7n} [\emph{fle˥n}]: этилен\\
0x7370 \textbf{hre7n} [\emph{ʰre˥n}]: передний план\\
0x7371 \textbf{ble7n} [\emph{ble˥n}]: китовый ус\\
0x737A \textbf{xle7n} [\emph{xle˥n}]: Хелен\\
0x7382 \textbf{kfe7n} [\emph{kfe˥n}]: кофеин\\
0x7388 \textbf{hbe7n} [\emph{ʰbe˥n}]: бретонский\\
0x738A \textbf{tfe7n} [\emph{tfe˥n}]: этан\\
0x7390 \textbf{hge7n} [\emph{ʰge˥n}]: калечить\\
0x7398 \textbf{hde7n} [\emph{ʰde˥n}]: слышимый\\
0x73A0 \textbf{hze7n} [\emph{ʰze˥n}]: гексан\\
0x73A2 \textbf{kce7n} [\emph{kʃe˥n}]: кессон\\
0x73A4 \textbf{pce7n} [\emph{pʃe˥n}]: пектин\\
0x73A8 \textbf{hje7n} [\emph{ʰʒe˥n}]: тощая кишка\\
0x73AA \textbf{tce7n} [\emph{tʃe˥n}]: изнурение\\
0x73B0 \textbf{hve7n} [\emph{ʰve˥n}]: больверк\\
0x73C1 \textbf{mre7n} [\emph{mre˥n}]: пленочный\\
0x73C2 \textbf{kre7n} [\emph{kre˥n}]: Looker на\\
0x73C4 \textbf{pre7n} [\emph{pre˥n}]: мокрота\\
0x73C6 \textbf{nre7n} [\emph{nre˥n}]: Золушка\\
0x73C8 \textbf{hqe7n} [\emph{ʰŋe˥n}]: англичанка\\
0x73C9 \textbf{sre7n} [\emph{sre˥n}]: фосфоресцирующий\\
0x73CA \textbf{tre7n} [\emph{tre˥n}]: терьер\\
0x73CD \textbf{cre7n} [\emph{ʃre˥n}]: в соответствии с ним\\
0x73D0 \textbf{hxe7n} [\emph{ʰxe˥n}]: отклоненный\\
0x73D1 \textbf{bre7n} [\emph{bre˥n}]: Бахрейн\\
0x73D2 \textbf{gre7n} [\emph{gre˥n}]: разгружать\\
0x73D3 \textbf{dre7n} [\emph{dre˥n}]: безрассудство делают\\
0x73DA \textbf{xre7n} [\emph{xre˥n}]: гиена\\
0x73E2 \textbf{kxe7n} [\emph{kxe˥n}]: слабый коленчатый\\
0x73E4 \textbf{pxe7n} [\emph{pxe˥n}]: мята\\
0x73EA \textbf{txe7n} [\emph{txe˥n}]: тенор\\
0x7401 \textbf{myo7n} [\emph{mjo˥n}]: симония\\
0x7402 \textbf{kyo7n} [\emph{kjo˥n}]: комедийная актриса\\
0x7404 \textbf{pyo7n} [\emph{pjo˥n}]: пятиугольный\\
0x7406 \textbf{nyo7n} [\emph{njo˥n}]: нерожденный\\
0x7408 \textbf{hmo7n} [\emph{ʰmo˥n}]: аммиак\\
0x7409 \textbf{syo7n} [\emph{sjo˥n}]: сырой\\
0x740A \textbf{tyo7n} [\emph{tjo˥n}]: девяносто один\\
0x740B \textbf{lyo7n} [\emph{ljo˥n}]: клонический\\
0x740E \textbf{ryo7n} [\emph{rjo˥n}]: радикальный\\
0x7410 \textbf{hko7n} [\emph{ʰko˥n}]: Обще\\
0x7411 \textbf{byo7n} [\emph{bjo˥n}]: бурсит большого пальца стопы\\
0x7413 \textbf{dyo7n} [\emph{djo˥n}]: индо иранское\\
0x7418 \textbf{hyo7n} [\emph{ʰjo˥n}]: предвидели\\
0x741A \textbf{xyo7n} [\emph{xjo˥n}]: бандит\\
0x7420 \textbf{hpo7n} [\emph{ʰpo˥n}]: Аполлон\\
0x7422 \textbf{kwo7n} [\emph{kwo˥n}]: Корнуолл\\
0x7428 \textbf{hwo7n} [\emph{ʰwo˥n}]: всеобъемлющая\\
0x7429 \textbf{swo7n} [\emph{swo˥n}]: мастодонт\\
0x742A \textbf{two7n} [\emph{two˥n}]: антициклон\\
0x7430 \textbf{hno7n} [\emph{ʰno˥n}]: несправедливый\\
0x7431 \textbf{bwo7n} [\emph{bwo˥n}]: бронх\\
0x7433 \textbf{dwo7n} [\emph{dwo˥n}]: допамин\\
0x7442 \textbf{kso7n} [\emph{kso˥n}]: англосаксонской\\
0x7444 \textbf{pso7n} [\emph{pso˥n}]: столбняк\\
0x7448 \textbf{hso7n} [\emph{ʰso˥n}]: гимн\\
0x744A \textbf{tso7n} [\emph{tso˥n}]: Тоскана\\
0x7450 \textbf{hto7n} [\emph{ʰto˥n}]: штурм мозга\\
0x7458 \textbf{hlo7n} [\emph{ʰlo˥n}]: фаллопиевы\\
0x7460 \textbf{hfo7n} [\emph{ʰfo˥n}]: экспромтом\\
0x7462 \textbf{klo7n} [\emph{klo˥n}]: золотуха\\
0x7464 \textbf{plo7n} [\emph{plo˥n}]: обесцвечивание\\
0x7468 \textbf{hco7n} [\emph{ʰʃo˥n}]: подчеркивание\\
0x7469 \textbf{slo7n} [\emph{slo˥n}]: Соломон\\
0x746A \textbf{tlo7n} [\emph{tlo˥n}]: Каталония\\
0x746C \textbf{flo7n} [\emph{flo˥n}]: афелий\\
0x7470 \textbf{hro7n} [\emph{ʰro˥n}]: скрещивание\\
0x7473 \textbf{dlo7n} [\emph{dlo˥n}]: квадрильон\\
0x747A \textbf{xlo7n} [\emph{xlo˥n}]: с воздушным охлаждением\\
0x7482 \textbf{kfo7n} [\emph{kfo˥n}]: аккордеоны\\
0x7488 \textbf{hbo7n} [\emph{ʰbo˥n}]: кнопки\\
0x748A \textbf{tfo7n} [\emph{tfo˥n}]: буксировка цветные\\
0x7490 \textbf{hgo7n} [\emph{ʰgo˥n}]: августинец\\
0x7498 \textbf{hdo7n} [\emph{ʰdo˥n}]: несогласный\\
0x74A0 \textbf{hzo7n} [\emph{ʰzo˥n}]: меццо-сопрано\\
0x74A2 \textbf{kco7n} [\emph{kʃo˥n}]: хромота\\
0x74A8 \textbf{hjo7n} [\emph{ʰʒo˥n}]: унгвент\\
0x74AA \textbf{tco7n} [\emph{tʃo˥n}]: колоколообразный\\
0x74C1 \textbf{mro7n} [\emph{mro˥n}]: умирающий\\
0x74C2 \textbf{kro7n} [\emph{kro˥n}]: викторианский\\
0x74C4 \textbf{pro7n} [\emph{pro˥n}]: прихорашиваться\\
0x74C6 \textbf{nro7n} [\emph{nro˥n}]: концентрический\\
0x74C8 \textbf{hqo7n} [\emph{ʰŋo˥n}]: ацетон\\
0x74C9 \textbf{sro7n} [\emph{sro˥n}]: все важные\\
0x74CA \textbf{tro7n} [\emph{tro˥n}]: трио\\
0x74CC \textbf{fro7n} [\emph{fro˥n}]: формалин\\
0x74CD \textbf{cro7n} [\emph{ʃro˥n}]: вычитаемое\\
0x74D0 \textbf{hxo7n} [\emph{ʰxo˥n}]: нетрадиционный\\
0x74D1 \textbf{bro7n} [\emph{bro˥n}]: баритон\\
0x74D3 \textbf{dro7n} [\emph{dro˥n}]: солнечник\\
0x74DA \textbf{xro7n} [\emph{xro˥n}]: герои\\
0x74E2 \textbf{kxo7n} [\emph{kxo˥n}]: кантонский\\
0x74E4 \textbf{pxo7n} [\emph{pxo˥n}]: Копенгаген\\
0x74EA \textbf{txo7n} [\emph{txo˥n}]: закал\\
0x74F1 \textbf{bxo7n} [\emph{bxo˥n}]: румянец\\
0x74F3 \textbf{dxo7n} [\emph{dxo˥n}]: лженаука\\
0x7506 \textbf{ny67n} [\emph{njə˥n}]: Анд\\
0x7508 \textbf{hm67n} [\emph{ʰmə˥n}]: тимпан\\
0x7509 \textbf{sy67n} [\emph{sjə˥n}]: изотонический\\
0x750E \textbf{ry67n} [\emph{rjə˥n}]: Райнланд\\
0x7510 \textbf{hk67n} [\emph{ʰkə˥n}]: хинолина\\
0x7518 \textbf{hy67n} [\emph{ʰjə˥n}]: воинственный\\
0x7520 \textbf{hp67n} [\emph{ʰpə˥n}]: перманганат\\
0x7530 \textbf{hn67n} [\emph{ʰnə˥n}]: никотин\\
0x7548 \textbf{hs67n} [\emph{ʰsə˥n}]: эоцен\\
0x754A \textbf{ts67n} [\emph{tsə˥n}]: да сэр\\
0x7550 \textbf{ht67n} [\emph{ʰtə˥n}]: утолщение\\
0x7558 \textbf{hl67n} [\emph{ʰlə˥n}]: полноценный\\
0x7560 \textbf{hf67n} [\emph{ʰfə˥n}]: фенол\\
0x7568 \textbf{hc67n} [\emph{ʰʃə˥n}]: Снеговик\\
0x756A \textbf{tl67n} [\emph{tlə˥n}]: интерпелляция\\
0x7570 \textbf{hr67n} [\emph{ʰrə˥n}]: подержанный\\
0x757A \textbf{xl67n} [\emph{xlə˥n}]: Таллин\\
0x75C2 \textbf{kr67n} [\emph{krə˥n}]: кератин\\
0x75C6 \textbf{nr67n} [\emph{nrə˥n}]: одновременно\\
0x75C9 \textbf{sr67n} [\emph{srə˥n}]: послеобеденный\\
0x75CA \textbf{tr67n} [\emph{trə˥n}]: каровое озеро\\
0x75D0 \textbf{hx67n} [\emph{ʰxə˥n}]: ломоть\\
0x75EA \textbf{tx67n} [\emph{txə˥n}]: многозадачных\\
0x7801 \textbf{myi2n} [\emph{mji˩n}]: сажа\\
0x7802 \textbf{kyi2n} [\emph{kji˩n}]: доктринер\\
0x7804 \textbf{pyi2n} [\emph{pji˩n}]: ипохондрик\\
0x7806 \textbf{nyi2n} [\emph{nji˩n}]: единорог\\
0x7808 \textbf{hmi2n} [\emph{ʰmi˩n}]: суеверие\\
0x7809 \textbf{syi2n} [\emph{sji˩n}]: Сидон\\
0x780A \textbf{tyi2n} [\emph{tji˩n}]: почетной\\
0x780B \textbf{lyi2n} [\emph{lji˩n}]: литания\\
0x780C \textbf{fyi2n} [\emph{fji˩n}]: день выплаты жалованья\\
0x780D \textbf{cyi2n} [\emph{ʃji˩n}]: циан\\
0x780E \textbf{ryi2n} [\emph{rji˩n}]: эркер\\
0x7810 \textbf{hki2n} [\emph{ʰki˩n}]: грешник\\
0x7811 \textbf{byi2n} [\emph{bji˩n}]: большой обратный порядок байт\\
0x7813 \textbf{dyi2n} [\emph{dji˩n}]: дифтонг\\
0x7816 \textbf{vyi2n} [\emph{vji˩n}]: головокружительный\\
0x7818 \textbf{hyi2n} [\emph{ʰji˩n}]: забвение\\
0x781A \textbf{xyi2n} [\emph{xji˩n}]: поворотливость\\
0x7820 \textbf{hpi2n} [\emph{ʰpi˩n}]: пехота\\
0x7821 \textbf{mwi2n} [\emph{mwi˩n}]: пескарь\\
0x7822 \textbf{kwi2n} [\emph{kwi˩n}]: циркулярно\\
0x7824 \textbf{pwi2n} [\emph{pwi˩n}]: плевательница\\
0x7826 \textbf{nwi2n} [\emph{nwi˩n}]: неконституционный\\
0x7828 \textbf{hwi2n} [\emph{ʰwi˩n}]: Уилли\\
0x7829 \textbf{swi2n} [\emph{swi˩n}]: психоанализ\\
0x782A \textbf{twi2n} [\emph{twi˩n}]: супница\\
0x782B \textbf{lwi2n} [\emph{lwi˩n}]: альбинос\\
0x782D \textbf{cwi2n} [\emph{ʃwi˩n}]: зонирование\\
0x782E \textbf{rwi2n} [\emph{rwi˩n}]: реквизиция\\
0x7830 \textbf{hni2n} [\emph{ʰni˩n}]: надписи\\
0x7831 \textbf{bwi2n} [\emph{bwi˩n}]: би ежегодный\\
0x783A \textbf{xwi2n} [\emph{xwi˩n}]: лечь в дрейф\\
0x7842 \textbf{ksi2n} [\emph{ksi˩n}]: сохранившийся\\
0x7844 \textbf{psi2n} [\emph{psi˩n}]: фарфор\\
0x7848 \textbf{hsi2n} [\emph{ʰsi˩n}]: самый широкий\\
0x7850 \textbf{hti2n} [\emph{ʰti˩n}]: междисциплинарный\\
0x7858 \textbf{hli2n} [\emph{ʰli˩n}]: слоны\\
0x7860 \textbf{hfi2n} [\emph{ʰfi˩n}]: безрассудная страсть\\
0x7862 \textbf{kli2n} [\emph{kli˩n}]: кавалерия\\
0x7864 \textbf{pli2n} [\emph{pli˩n}]: пинг\\
0x7868 \textbf{hci2n} [\emph{ʰʃi˩n}]: снисходить\\
0x7869 \textbf{sli2n} [\emph{sli˩n}]: самонадеянным\\
0x786A \textbf{tli2n} [\emph{tli˩n}]: этимология\\
0x786C \textbf{fli2n} [\emph{fli˩n}]: филигрань\\
0x786D \textbf{cli2n} [\emph{ʃli˩n}]: кляузничество\\
0x7870 \textbf{hri2n} [\emph{ʰri˩n}]: оживление\\
0x7871 \textbf{bli2n} [\emph{bli˩n}]: Берлин\\
0x7872 \textbf{gli2n} [\emph{gli˩n}]: облегчение\\
0x7873 \textbf{dli2n} [\emph{dli˩n}]: деморализация\\
0x7874 \textbf{zli2n} [\emph{zli˩n}]: ксилола\\
0x7875 \textbf{jli2n} [\emph{ʒli˩n}]: невероятна\\
0x7876 \textbf{vli2n} [\emph{vli˩n}]: прямая линия\\
0x787A \textbf{xli2n} [\emph{xli˩n}]: оживлять\\
0x7882 \textbf{kfi2n} [\emph{kfi˩n}]: занос\\
0x7884 \textbf{pfi2n} [\emph{pfi˩n}]: полдюйма\\
0x7888 \textbf{hbi2n} [\emph{ʰbi˩n}]: брат в закон\\
0x788A \textbf{tfi2n} [\emph{tfi˩n}]: рассеянность внимания\\
0x7890 \textbf{hgi2n} [\emph{ʰgi˩n}]: сальный\\
0x7898 \textbf{hdi2n} [\emph{ʰdi˩n}]: доверчивость\\
0x78A0 \textbf{hzi2n} [\emph{ʰzi˩n}]: суннит\\
0x78A2 \textbf{kci2n} [\emph{kʃi˩n}]: дружки\\
0x78A4 \textbf{pci2n} [\emph{pʃi˩n}]: импровизация\\
0x78A8 \textbf{hji2n} [\emph{ʰʒi˩n}]: Дженни\\
0x78AA \textbf{tci2n} [\emph{tʃi˩n}]: увещевать\\
0x78B0 \textbf{hvi2n} [\emph{ʰvi˩n}]: завидный\\
0x78B3 \textbf{dji2n} [\emph{dʒi˩n}]: предопределять\\
0x78C1 \textbf{mri2n} [\emph{mri˩n}]: материнство\\
0x78C2 \textbf{kri2n} [\emph{kri˩n}]: Коринф\\
0x78C4 \textbf{pri2n} [\emph{pri˩n}]: просматривал\\
0x78C6 \textbf{nri2n} [\emph{nri˩n}]: граждане\\
0x78C8 \textbf{hqi2n} [\emph{ʰŋi˩n}]: звенящий\\
0x78C9 \textbf{sri2n} [\emph{sri˩n}]: странность\\
0x78CA \textbf{tri2n} [\emph{tri˩n}]: поразительный\\
0x78CC \textbf{fri2n} [\emph{fri˩n}]: флирт\\
0x78CD \textbf{cri2n} [\emph{ʃri˩n}]: определитель\\
0x78D0 \textbf{hxi2n} [\emph{ʰxi˩n}]: охваченном\\
0x78D1 \textbf{bri2n} [\emph{bri˩n}]: годовалый\\
0x78D2 \textbf{gri2n} [\emph{gri˩n}]: глицерин\\
0x78D3 \textbf{dri2n} [\emph{dri˩n}]: угнетенный\\
0x78D4 \textbf{zri2n} [\emph{zri˩n}]: мелодичный перезвон\\
0x78D6 \textbf{vri2n} [\emph{vri˩n}]: ивуариец\\
0x78D9 \textbf{qri2n} [\emph{ŋri˩n}]: скарлатина\\
0x78DA \textbf{xri2n} [\emph{xri˩n}]: служить примером\\
0x78E2 \textbf{kxi2n} [\emph{kxi˩n}]: сорок один\\
0x78E4 \textbf{pxi2n} [\emph{pxi˩n}]: сфинкс\\
0x78EA \textbf{txi2n} [\emph{txi˩n}]: покорный\\
0x78F1 \textbf{bxi2n} [\emph{bxi˩n}]: брачный\\
0x78F2 \textbf{gxi2n} [\emph{gxi˩n}]: женоненавистник\\
0x78F3 \textbf{dxi2n} [\emph{dxi˩n}]: дикция\\
0x7901 \textbf{mya2n} [\emph{mjä˩n}]: маммона\\
0x7902 \textbf{kya2n} [\emph{kjä˩n}]: коррелированный\\
0x7904 \textbf{pya2n} [\emph{pjä˩n}]: проституция\\
0x7906 \textbf{nya2n} [\emph{njä˩n}]: монахини\\
0x7908 \textbf{hma2n} [\emph{ʰmä˩n}]: равнины\\
0x7909 \textbf{sya2n} [\emph{sjä˩n}]: сатин\\
0x790A \textbf{tya2n} [\emph{tjä˩n}]: утреня\\
0x790B \textbf{lya2n} [\emph{ljä˩n}]: светоч\\
0x790C \textbf{fya2n} [\emph{fjä˩n}]: выявляя\\
0x790D \textbf{cya2n} [\emph{ʃjä˩n}]: свинка\\
0x790E \textbf{rya2n} [\emph{rjä˩n}]: грациозность\\
0x7910 \textbf{hka2n} [\emph{ʰkä˩n}]: сопровождая\\
0x7911 \textbf{bya2n} [\emph{bjä˩n}]: большой\\
0x7912 \textbf{gya2n} [\emph{gjä˩n}]: земляной орех\\
0x7913 \textbf{dya2n} [\emph{djä˩n}]: диатонический\\
0x7915 \textbf{jya2n} [\emph{ʒjä˩n}]: Инь Янь\\
0x7916 \textbf{vya2n} [\emph{vjä˩n}]: внутренности\\
0x7918 \textbf{hya2n} [\emph{ʰjä˩n}]: сирота\\
0x791A \textbf{xya2n} [\emph{xjä˩n}]: солнцами\\
0x7920 \textbf{hpa2n} [\emph{ʰpä˩n}]: патогенный микроорганизм\\
0x7921 \textbf{mwa2n} [\emph{mwä˩n}]: научная статья\\
0x7922 \textbf{kwa2n} [\emph{kwä˩n}]: кругосветное плавание\\
0x7924 \textbf{pwa2n} [\emph{pwä˩n}]: папский\\
0x7926 \textbf{nwa2n} [\emph{nwä˩n}]: евнух\\
0x7928 \textbf{hwa2n} [\emph{ʰwä˩n}]: Ублюдок\\
0x7929 \textbf{swa2n} [\emph{swä˩n}]: юго-восточный\\
0x792A \textbf{twa2n} [\emph{twä˩n}]: самооборона\\
0x792B \textbf{lwa2n} [\emph{lwä˩n}]: иносказательный\\
0x792C \textbf{fwa2n} [\emph{fwä˩n}]: астения\\
0x792E \textbf{rwa2n} [\emph{rwä˩n}]: царствовал\\
0x7930 \textbf{hna2n} [\emph{ʰnä˩n}]: лунный свет\\
0x7931 \textbf{bwa2n} [\emph{bwä˩n}]: наметка\\
0x7933 \textbf{dwa2n} [\emph{dwä˩n}]: бог король\\
0x793A \textbf{xwa2n} [\emph{xwä˩n}]: боярышник\\
0x7942 \textbf{ksa2n} [\emph{ksä˩n}]: придворный\\
0x7944 \textbf{psa2n} [\emph{psä˩n}]: прохожие\\
0x7948 \textbf{hsa2n} [\emph{ʰsä˩n}]: господа\\
0x794A \textbf{tsa2n} [\emph{tsä˩n}]: станций\\
0x7950 \textbf{hta2n} [\emph{ʰtä˩n}]: учил\\
0x7953 \textbf{dza2n} [\emph{dzä˩n}]: весь год\\
0x7958 \textbf{hla2n} [\emph{ʰlä˩n}]: Калифорния\\
0x7960 \textbf{hfa2n} [\emph{ʰfä˩n}]: фланель\\
0x7962 \textbf{kla2n} [\emph{klä˩n}]: холера\\
0x7964 \textbf{pla2n} [\emph{plä˩n}]: павильон\\
0x7968 \textbf{hca2n} [\emph{ʰʃä˩n}]: сержант\\
0x7969 \textbf{sla2n} [\emph{slä˩n}]: вскользь\\
0x796A \textbf{tla2n} [\emph{tlä˩n}]: младший офицер\\
0x796C \textbf{fla2n} [\emph{flä˩n}]: затопленный\\
0x796D \textbf{cla2n} [\emph{ʃlä˩n}]: дробовик\\
0x7970 \textbf{hra2n} [\emph{ʰrä˩n}]: ряды\\
0x7971 \textbf{bla2n} [\emph{blä˩n}]: заслонять\\
0x7972 \textbf{gla2n} [\emph{glä˩n}]: магнолия\\
0x7973 \textbf{dla2n} [\emph{dlä˩n}]: разлаживать\\
0x7974 \textbf{zla2n} [\emph{zlä˩n}]: симулировать болезнь\\
0x7976 \textbf{vla2n} [\emph{vlä˩n}]: виллан\\
0x797A \textbf{xla2n} [\emph{xlä˩n}]: воющий\\
0x7982 \textbf{kfa2n} [\emph{kfä˩n}]: какофония\\
0x7984 \textbf{pfa2n} [\emph{pfä˩n}]: пантеон\\
0x7988 \textbf{hba2n} [\emph{ʰbä˩n}]: часто посещаемый\\
0x798A \textbf{tfa2n} [\emph{tfä˩n}]: откармливать\\
0x7990 \textbf{hga2n} [\emph{ʰgä˩n}]: строп\\
0x7993 \textbf{dva2n} [\emph{dvä˩n}]: расстройство мочеиспускания\\
0x7998 \textbf{hda2n} [\emph{ʰdä˩n}]: края\\
0x79A0 \textbf{hza2n} [\emph{ʰzä˩n}]: праздношатание\\
0x79A2 \textbf{kca2n} [\emph{kʃä˩n}]: неоднозначность\\
0x79A4 \textbf{pca2n} [\emph{pʃä˩n}]: фонтанирование\\
0x79A8 \textbf{hja2n} [\emph{ʰʒä˩n}]: бездонный\\
0x79AA \textbf{tca2n} [\emph{tʃä˩n}]: бранить\\
0x79B0 \textbf{hva2n} [\emph{ʰvä˩n}]: раздаваться\\
0x79B1 \textbf{bja2n} [\emph{bʒä˩n}]: раздающий милостыню\\
0x79B3 \textbf{dja2n} [\emph{dʒä˩n}]: такса\\
0x79C1 \textbf{mra2n} [\emph{mrä˩n}]: адмирал\\
0x79C2 \textbf{kra2n} [\emph{krä˩n}]: проказа\\
0x79C4 \textbf{pra2n} [\emph{prä˩n}]: преступники\\
0x79C6 \textbf{nra2n} [\emph{nrä˩n}]: магнат\\
0x79C8 \textbf{hqa2n} [\emph{ʰŋä˩n}]: Восходящий\\
0x79C9 \textbf{sra2n} [\emph{srä˩n}]: сверхчеловек\\
0x79CA \textbf{tra2n} [\emph{trä˩n}]: сносно\\
0x79CC \textbf{fra2n} [\emph{frä˩n}]: пожарник\\
0x79CD \textbf{cra2n} [\emph{ʃrä˩n}]: кораблекрушение\\
0x79D0 \textbf{hxa2n} [\emph{ʰxä˩n}]: охотно\\
0x79D1 \textbf{bra2n} [\emph{brä˩n}]: северный\\
0x79D2 \textbf{gra2n} [\emph{grä˩n}]: грыжа\\
0x79D3 \textbf{dra2n} [\emph{drä˩n}]: старейшина\\
0x79D4 \textbf{zra2n} [\emph{zrä˩n}]: обострение\\
0x79D5 \textbf{jra2n} [\emph{ʒrä˩n}]: надпочечная\\
0x79D6 \textbf{vra2n} [\emph{vrä˩n}]: купорос\\
0x79D9 \textbf{qra2n} [\emph{ŋrä˩n}]: тиранить\\
0x79DA \textbf{xra2n} [\emph{xrä˩n}]: императрица\\
0x79E2 \textbf{kxa2n} [\emph{kxä˩n}]: шхуна\\
0x79E4 \textbf{pxa2n} [\emph{pxä˩n}]: Пэдди\\
0x79EA \textbf{txa2n} [\emph{txä˩n}]: вызванный\\
0x79F1 \textbf{bxa2n} [\emph{bxä˩n}]: бланшировать\\
0x79F2 \textbf{gxa2n} [\emph{gxä˩n}]: отмороженное место\\
0x79F3 \textbf{dxa2n} [\emph{dxä˩n}]: безопасное поведение\\
0x7A01 \textbf{myu2n} [\emph{mju˩n}]: маниакальный\\
0x7A08 \textbf{hmu2n} [\emph{ʰmu˩n}]: Мюнхен\\
0x7A09 \textbf{syu2n} [\emph{sju˩n}]: Сеульский\\
0x7A0A \textbf{tyu2n} [\emph{tju˩n}]: тональный\\
0x7A0E \textbf{ryu2n} [\emph{rju˩n}]: рутений\\
0x7A10 \textbf{hku2n} [\emph{ʰku˩n}]: подушки\\
0x7A18 \textbf{hyu2n} [\emph{ʰju˩n}]: евгеника\\
0x7A20 \textbf{hpu2n} [\emph{ʰpu˩n}]: в верхней части города\\
0x7A28 \textbf{hwu2n} [\emph{ʰwu˩n}]: шерстяной\\
0x7A30 \textbf{hnu2n} [\emph{ʰnu˩n}]: нераспечатанный\\
0x7A42 \textbf{ksu2n} [\emph{ksu˩n}]: Курдистан\\
0x7A48 \textbf{hsu2n} [\emph{ʰsu˩n}]: более чистый\\
0x7A4A \textbf{tsu2n} [\emph{tsu˩n}]: превращение\\
0x7A50 \textbf{htu2n} [\emph{ʰtu˩n}]: трибуна\\
0x7A58 \textbf{hlu2n} [\emph{ʰlu˩n}]: лютеранский\\
0x7A60 \textbf{hfu2n} [\emph{ʰfu˩n}]: само важно\\
0x7A62 \textbf{klu2n} [\emph{klu˩n}]: сговариваться\\
0x7A68 \textbf{hcu2n} [\emph{ʰʃu˩n}]: недостойное поведение\\
0x7A69 \textbf{slu2n} [\emph{slu˩n}]: лазурный\\
0x7A6A \textbf{tlu2n} [\emph{tlu˩n}]: полупрозрачность\\
0x7A6C \textbf{flu2n} [\emph{flu˩n}]: хлопьевидный\\
0x7A70 \textbf{hru2n} [\emph{ʰru˩n}]: маршруты\\
0x7A71 \textbf{blu2n} [\emph{blu˩n}]: булинь\\
0x7A7A \textbf{xlu2n} [\emph{xlu˩n}]: аллилуйя\\
0x7A88 \textbf{hbu2n} [\emph{ʰbu˩n}]: бульон\\
0x7A8A \textbf{tfu2n} [\emph{tfu˩n}]: бой быков\\
0x7A90 \textbf{hgu2n} [\emph{ʰgu˩n}]: коронование\\
0x7A98 \textbf{hdu2n} [\emph{ʰdu˩n}]: середина июня\\
0x7AA2 \textbf{kcu2n} [\emph{kʃu˩n}]: близко вяжут\\
0x7AAA \textbf{tcu2n} [\emph{tʃu˩n}]: тюнер\\
0x7AC2 \textbf{kru2n} [\emph{kru˩n}]: корунд\\
0x7AC6 \textbf{nru2n} [\emph{nru˩n}]: первоисточник\\
0x7AC9 \textbf{sru2n} [\emph{sru˩n}]: распятие на кресте\\
0x7ACA \textbf{tru2n} [\emph{tru˩n}]: Турин\\
0x7AD0 \textbf{hxu2n} [\emph{ʰxu˩n}]: Хусайн\\
0x7ADA \textbf{xru2n} [\emph{xru˩n}]: Ааронова\\
0x7AEA \textbf{txu2n} [\emph{txu˩n}]: скула\\
0x7B02 \textbf{kye2n} [\emph{kje˩n}]: концессионер\\
0x7B06 \textbf{nye2n} [\emph{nje˩n}]: нейлон\\
0x7B08 \textbf{hme2n} [\emph{ʰme˩n}]: каменщик\\
0x7B09 \textbf{sye2n} [\emph{sje˩n}]: десятилетний\\
0x7B0A \textbf{tye2n} [\emph{tje˩n}]: трехгодичного\\
0x7B0E \textbf{rye2n} [\emph{rje˩n}]: заокеанский\\
0x7B10 \textbf{hke2n} [\emph{ʰke˩n}]: коса\\
0x7B18 \textbf{hye2n} [\emph{ʰje˩n}]: глотка\\
0x7B1A \textbf{xye2n} [\emph{xje˩n}]: красный глаз\\
0x7B20 \textbf{hpe2n} [\emph{ʰpe˩n}]: эпигенетическое\\
0x7B21 \textbf{mwe2n} [\emph{mwe˩n}]: нимфомания\\
0x7B24 \textbf{pwe2n} [\emph{pwe˩n}]: парамагнитный\\
0x7B26 \textbf{nwe2n} [\emph{nwe˩n}]: один лайнер\\
0x7B28 \textbf{hwe2n} [\emph{ʰwe˩n}]: искусный\\
0x7B29 \textbf{swe2n} [\emph{swe˩n}]: мореходный\\
0x7B2B \textbf{lwe2n} [\emph{lwe˩n}]: красное Крыло\\
0x7B2E \textbf{rwe2n} [\emph{rwe˩n}]: четырехмерный\\
0x7B30 \textbf{hne2n} [\emph{ʰne˩n}]: соловей\\
0x7B42 \textbf{kse2n} [\emph{kse˩n}]: секстан\\
0x7B44 \textbf{pse2n} [\emph{pse˩n}]: точильный камень\\
0x7B48 \textbf{hse2n} [\emph{ʰse˩n}]: святые\\
0x7B4A \textbf{tse2n} [\emph{tse˩n}]: двадцать секунд\\
0x7B50 \textbf{hte2n} [\emph{ʰte˩n}]: в настоящее время\\
0x7B58 \textbf{hle2n} [\emph{ʰle˩n}]: лейтенант\\
0x7B60 \textbf{hfe2n} [\emph{ʰfe˩n}]: фаэтон\\
0x7B62 \textbf{kle2n} [\emph{kle˩n}]: колоннада\\
0x7B64 \textbf{ple2n} [\emph{ple˩n}]: пленум\\
0x7B68 \textbf{hce2n} [\emph{ʰʃe˩n}]: увещевание\\
0x7B69 \textbf{sle2n} [\emph{sle˩n}]: свинопас\\
0x7B6A \textbf{tle2n} [\emph{tle˩n}]: восьмидесятилетний\\
0x7B6C \textbf{fle2n} [\emph{fle˩n}]: нафталин\\
0x7B6D \textbf{cle2n} [\emph{ʃle˩n}]: тысячелетний\\
0x7B70 \textbf{hre2n} [\emph{ʰre˩n}]: полки\\
0x7B72 \textbf{gle2n} [\emph{gle˩n}]: мегаломания\\
0x7B73 \textbf{dle2n} [\emph{dle˩n}]: радиотелефон\\
0x7B7A \textbf{xle2n} [\emph{xle˩n}]: гегельянец\\
0x7B84 \textbf{pfe2n} [\emph{pfe˩n}]: полу профессиональный\\
0x7B88 \textbf{hbe2n} [\emph{ʰbe˩n}]: прижаться\\
0x7B8A \textbf{tfe2n} [\emph{tfe˩n}]: правопреемник\\
0x7B90 \textbf{hge2n} [\emph{ʰge˩n}]: бог ребенок\\
0x7B98 \textbf{hde2n} [\emph{ʰde˩n}]: среднего возраста\\
0x7BA2 \textbf{kce2n} [\emph{kʃe˩n}]: четвертичный\\
0x7BA4 \textbf{pce2n} [\emph{pʃe˩n}]: довоенный\\
0x7BA8 \textbf{hje2n} [\emph{ʰʒe˩n}]: изнеженность\\
0x7BAA \textbf{tce2n} [\emph{tʃe˩n}]: чеченец\\
0x7BB0 \textbf{hve2n} [\emph{ʰve˩n}]: эфемерность\\
0x7BC1 \textbf{mre2n} [\emph{mre˩n}]: морена\\
0x7BC2 \textbf{kre2n} [\emph{kre˩n}]: четырехугольник\\
0x7BC4 \textbf{pre2n} [\emph{pre˩n}]: плеврит\\
0x7BC6 \textbf{nre2n} [\emph{nre˩n}]: грозовая туча\\
0x7BC9 \textbf{sre2n} [\emph{sre˩n}]: зычный\\
0x7BCA \textbf{tre2n} [\emph{tre˩n}]: токарный станок\\
0x7BCC \textbf{fre2n} [\emph{fre˩n}]: волосной\\
0x7BCD \textbf{cre2n} [\emph{ʃre˩n}]: длящийся четыре года\\
0x7BD0 \textbf{hxe2n} [\emph{ʰxe˩n}]: стройный\\
0x7BD1 \textbf{bre2n} [\emph{bre˩n}]: кибернетический\\
0x7BD2 \textbf{gre2n} [\emph{gre˩n}]: григорианский\\
0x7BD3 \textbf{dre2n} [\emph{dre˩n}]: затвердение\\
0x7BDA \textbf{xre2n} [\emph{xre˩n}]: мигрень\\
0x7BE2 \textbf{kxe2n} [\emph{kxe˩n}]: бешено\\
0x7BE4 \textbf{pxe2n} [\emph{pxe˩n}]: сдоба\\
0x7BEA \textbf{txe2n} [\emph{txe˩n}]: агрегирование\\
0x7BF1 \textbf{bxe2n} [\emph{bxe˩n}]: расположенный ниже\\
0x7C01 \textbf{myo2n} [\emph{mjo˩n}]: ламинарный\\
0x7C04 \textbf{pyo2n} [\emph{pjo˩n}]: пагинация\\
0x7C06 \textbf{nyo2n} [\emph{njo˩n}]: назначаемое лицо\\
0x7C08 \textbf{hmo2n} [\emph{ʰmo˩n}]: манна\\
0x7C09 \textbf{syo2n} [\emph{sjo˩n}]: сифон\\
0x7C0A \textbf{tyo2n} [\emph{tjo˩n}]: Океания\\
0x7C0B \textbf{lyo2n} [\emph{ljo˩n}]: поражения\\
0x7C10 \textbf{hko2n} [\emph{ʰko˩n}]: полковник\\
0x7C13 \textbf{dyo2n} [\emph{djo˩n}]: середине января\\
0x7C18 \textbf{hyo2n} [\emph{ʰjo˩n}]: Иона\\
0x7C20 \textbf{hpo2n} [\emph{ʰpo˩n}]: форма волны\\
0x7C22 \textbf{kwo2n} [\emph{kwo˩n}]: октан\\
0x7C28 \textbf{hwo2n} [\emph{ʰwo˩n}]: червоточина\\
0x7C2A \textbf{two2n} [\emph{two˩n}]: потрясающее зрелище\\
0x7C2E \textbf{rwo2n} [\emph{rwo˩n}]: ретороманский диалект\\
0x7C30 \textbf{hno2n} [\emph{ʰno˩n}]: нормандский\\
0x7C42 \textbf{kso2n} [\emph{kso˩n}]: с учетом\\
0x7C44 \textbf{pso2n} [\emph{pso˩n}]: непослушание\\
0x7C48 \textbf{hso2n} [\emph{ʰso˩n}]: источники\\
0x7C4A \textbf{tso2n} [\emph{tso˩n}]: триколор\\
0x7C50 \textbf{hto2n} [\emph{ʰto˩n}]: послушание\\
0x7C58 \textbf{hlo2n} [\emph{ʰlo˩n}]: Болонья\\
0x7C60 \textbf{hfo2n} [\emph{ʰfo˩n}]: опровержение\\
0x7C62 \textbf{klo2n} [\emph{klo˩n}]: Шотландия\\
0x7C68 \textbf{hco2n} [\emph{ʰʃo˩n}]: до того времени\\
0x7C69 \textbf{slo2n} [\emph{slo˩n}]: столбчатый\\
0x7C6A \textbf{tlo2n} [\emph{tlo˩n}]: танин\\
0x7C6C \textbf{flo2n} [\emph{flo˩n}]: дефлорация\\
0x7C70 \textbf{hro2n} [\emph{ʰro˩n}]: троянец\\
0x7C7A \textbf{xlo2n} [\emph{xlo˩n}]: прыщ\\
0x7C84 \textbf{pfo2n} [\emph{pfo˩n}]: продувание\\
0x7C88 \textbf{hbo2n} [\emph{ʰbo˩n}]: взбодрить\\
0x7C8A \textbf{tfo2n} [\emph{tfo˩n}]: фиксация азота\\
0x7C90 \textbf{hgo2n} [\emph{ʰgo˩n}]: родственный\\
0x7C98 \textbf{hdo2n} [\emph{ʰdo˩n}]: движения\\
0x7CA4 \textbf{pco2n} [\emph{pʃo˩n}]: мочеиспускание\\
0x7CAA \textbf{tco2n} [\emph{tʃo˩n}]: кошениль\\
0x7CC1 \textbf{mro2n} [\emph{mro˩n}]: бур\\
0x7CC2 \textbf{kro2n} [\emph{kro˩n}]: коронер\\
0x7CC4 \textbf{pro2n} [\emph{pro˩n}]: странствование\\
0x7CC6 \textbf{nro2n} [\emph{nro˩n}]: анаконда\\
0x7CC8 \textbf{hqo2n} [\emph{ʰŋo˩n}]: фотогеничный\\
0x7CC9 \textbf{sro2n} [\emph{sro˩n}]: правый борт\\
0x7CCA \textbf{tro2n} [\emph{tro˩n}]: токарь\\
0x7CD0 \textbf{hxo2n} [\emph{ʰxo˩n}]: чемодан\\
0x7CD1 \textbf{bro2n} [\emph{bro˩n}]: мускулистый\\
0x7CD3 \textbf{dro2n} [\emph{dro˩n}]: глухота\\
0x7CDA \textbf{xro2n} [\emph{xro˩n}]: храп\\
0x7CE2 \textbf{kxo2n} [\emph{kxo˩n}]: кетон\\
0x7CE4 \textbf{pxo2n} [\emph{pxo˩n}]: полигональный\\
0x7CEA \textbf{txo2n} [\emph{txo˩n}]: низкий скатных\\
0x7D08 \textbf{hm62n} [\emph{ʰmə˩n}]: монгольский\\
0x7D20 \textbf{hp62n} [\emph{ʰpə˩n}]: апеллянт\\
0x7D30 \textbf{hn62n} [\emph{ʰnə˩n}]: Noman\\
0x7D48 \textbf{hs62n} [\emph{ʰsə˩n}]: осетр\\
0x7D50 \textbf{ht62n} [\emph{ʰtə˩n}]: пропаривания\\
0x7D58 \textbf{hl62n} [\emph{ʰlə˩n}]: алеутская\\
0x7D69 \textbf{sl62n} [\emph{slə˩n}]: сибилянт\\
0x7D6A \textbf{tl62n} [\emph{tlə˩n}]: инстилляция\\
0x7D70 \textbf{hr62n} [\emph{ʰrə˩n}]: батут\\
0x7DCC \textbf{fr62n} [\emph{frə˩n}]: афферентный\\
0x7DD0 \textbf{hx62n} [\emph{ʰxə˩n}]: Осанна\\
0x7DEA \textbf{tx62n} [\emph{txə˩n}]: Тегеран\\
\subsection{ 8 }
0x800F \textbf{myi7h} [\emph{mji˥ʰ}]: я\\
0x8027 \textbf{pyi7h} [\emph{pji˥ʰ}]: аппликативны голос\\
0x804F \textbf{syi7h} [\emph{sji˥ʰ}]: existential clause, for asserting something that can be imagined.\\
0x804F \textbf{syi7h} [\emph{sji˥ʰ}]: existential clause, for asserting something that can be imagined.\\
0x809F \textbf{dyi7h} [\emph{dji˥ʰ}]: дис\\
0x8217 \textbf{ksi7h} [\emph{ksi˥ʰ}]: IST\\
0x8257 \textbf{tsi7h} [\emph{tsi˥ʰ}]: antessive case, indicates the noun phrase precedes the rest of the sentence. space-context retrospective-aspect \\
0x8357 \textbf{tli7h} [\emph{tli˥ʰ}]: дистанционный пульт будущее напряженный\\
0x848F \textbf{bvi7h} [\emph{bvi˥ʰ}]: дистрибутивный местоимение\\
0x8557 \textbf{tci7h} [\emph{tʃi˥ʰ}]: Archy\\
0x8627 \textbf{pri7h} [\emph{pri˥ʰ}]: пассивный причастие\\
0x8637 \textbf{nri7h} [\emph{nri˥ʰ}]: неодушевленный Пол\\
0x864F \textbf{sri7h} [\emph{sri˥ʰ}]: взаимная голос\\
0x8657 \textbf{tri7h} [\emph{tri˥ʰ}]: мимо пассивный причастие\\
0x8757 \textbf{txi7h} [\emph{txi˥ʰ}]: коллектив имя существительное\\
0x8801 \textbf{myis} [\emph{mjis}]: молекулярная\\
0x8802 \textbf{kyis} [\emph{kjis}]: бумага\\
0x8804 \textbf{pyis} [\emph{pjis}]: священник\\
0x8806 \textbf{nyis} [\emph{njis}]: транскрибировать\\
0x8808 \textbf{hmis} [\emph{ʰmis}]: мудрый\\
0x8809 \textbf{syis} [\emph{sjis}]: 20\\
0x880A \textbf{tyis} [\emph{tjis}]: прогресс\\
0x880B \textbf{lyis} [\emph{ljis}]: лесоводство\\
0x880C \textbf{fyis} [\emph{fjis}]: Очистительный завод\\
0x880D \textbf{cyis} [\emph{ʃjis}]: стул\\
0x880E \textbf{ryis} [\emph{rjis}]: отшельник\\
0x880F \textbf{mya7h} [\emph{mjä˥ʰ}]: мяукать\\
0x8810 \textbf{hkis} [\emph{ʰkis}]: черный\\
0x8811 \textbf{byis} [\emph{bjis}]: подсознательный\\
0x8812 \textbf{gyis} [\emph{gjis}]: Греция\\
0x8813 \textbf{dyis} [\emph{djis}]: средневековый\\
0x8814 \textbf{zyis} [\emph{zjis}]: солнечные часы\\
0x8815 \textbf{jyis} [\emph{ʒjis}]: Овен\\
0x8816 \textbf{vyis} [\emph{vjis}]: резец\\
0x8818 \textbf{hyis} [\emph{ʰjis}]: бедные\\
0x8819 \textbf{qyis} [\emph{ŋjis}]: бизнес человек\\
0x881A \textbf{xyis} [\emph{xjis}]: мимо пассивный причастие\\
0x8820 \textbf{hpis} [\emph{ʰpis}]: платить\\
0x8821 \textbf{mwis} [\emph{mwis}]: умственно\\
0x8822 \textbf{kwis} [\emph{kwis}]: классический\\
0x8824 \textbf{pwis} [\emph{pwis}]: закуска\\
0x8826 \textbf{nwis} [\emph{nwis}]: ногти\\
0x8828 \textbf{hwis} [\emph{ʰwis}]: Vialis дело\\
0x8829 \textbf{swis} [\emph{swis}]: моряк\\
0x882A \textbf{twis} [\emph{twis}]: быстро\\
0x882B \textbf{lwis} [\emph{lwis}]: щелочной\\
0x882C \textbf{fwis} [\emph{fwis}]: феминистка\\
0x882D \textbf{cwis} [\emph{ʃwis}]: Четверг\\
0x882E \textbf{rwis} [\emph{rwis}]: расовой\\
0x8830 \textbf{hnis} [\emph{ʰnis}]: возраст\\
0x8831 \textbf{bwis} [\emph{bwis}]: бикини\\
0x8832 \textbf{gwis} [\emph{gwis}]: цитозин\\
0x8833 \textbf{dwis} [\emph{dwis}]: диоксид\\
0x8835 \textbf{jwis} [\emph{ʒwis}]: Близнецы\\
0x8836 \textbf{vwis} [\emph{vwis}]: выщипывать пинцетом\\
0x8839 \textbf{qwis} [\emph{ŋwis}]: повесток\\
0x883A \textbf{xwis} [\emph{xwis}]: фриз\\
0x8842 \textbf{ksis} [\emph{ksis}]: история\\
0x8844 \textbf{psis} [\emph{psis}]: поэзия\\
0x8848 \textbf{hsis} [\emph{ʰsis}]: всегда\\
0x884A \textbf{tsis} [\emph{tsis}]: справедливость\\
0x8850 \textbf{htis} [\emph{ʰtis}]: изготовление\\
0x8851 \textbf{bzis} [\emph{bzis}]: виза\\
0x8852 \textbf{gzis} [\emph{gzis}]: гейзер\\
0x8853 \textbf{dzis} [\emph{dzis}]: изморось\\
0x8858 \textbf{hlis} [\emph{ʰlis}]: шесть\\
0x8860 \textbf{hfis} [\emph{ʰfis}]: отставка\\
0x8862 \textbf{klis} [\emph{klis}]: изоляция\\
0x8864 \textbf{plis} [\emph{plis}]: полиция\\
0x8868 \textbf{hcis} [\emph{ʰʃis}]: жена\\
0x8869 \textbf{slis} [\emph{slis}]: серии\\
0x886A \textbf{tlis} [\emph{tlis}]: зубы\\
0x886C \textbf{flis} [\emph{flis}]: анализировать\\
0x886D \textbf{clis} [\emph{ʃlis}]: июль\\
0x8870 \textbf{hris} [\emph{ʰris}]: обычай\\
0x8871 \textbf{blis} [\emph{blis}]: ваниль\\
0x8872 \textbf{glis} [\emph{glis}]: бензин\\
0x8873 \textbf{dlis} [\emph{dlis}]: идиосинкразия\\
0x8874 \textbf{zlis} [\emph{zlis}]: Сицилия\\
0x8875 \textbf{jlis} [\emph{ʒlis}]: песочные часы\\
0x8876 \textbf{vlis} [\emph{vlis}]: виртуальный реальность\\
0x887A \textbf{xlis} [\emph{xlis}]: Италия\\
0x8882 \textbf{kfis} [\emph{kfis}]: вакансия\\
0x8884 \textbf{pfis} [\emph{pfis}]: ратификация\\
0x8888 \textbf{hbis} [\emph{ʰbis}]: база\\
0x888A \textbf{tfis} [\emph{tfis}]: отказ\\
0x8890 \textbf{hgis} [\emph{ʰgis}]: сыр\\
0x8891 \textbf{bvis} [\emph{bvis}]: безрукавный\\
0x8892 \textbf{gvis} [\emph{gvis}]: зерноядный\\
0x8893 \textbf{dvis} [\emph{dvis}]: подразделения\\
0x8898 \textbf{hdis} [\emph{ʰdis}]: виновный\\
0x88A0 \textbf{hzis} [\emph{ʰzis}]: слеза\\
0x88A2 \textbf{kcis} [\emph{kʃis}]: серый\\
0x88A4 \textbf{pcis} [\emph{pʃis}]: опасно\\
0x88A8 \textbf{hjis} [\emph{ʰʒis}]: игла\\
0x88AA \textbf{tcis} [\emph{tʃis}]: честный\\
0x88B0 \textbf{hvis} [\emph{ʰvis}]: фекалии\\
0x88B1 \textbf{bjis} [\emph{bʒis}]: Бэзил\\
0x88B2 \textbf{gjis} [\emph{gʒis}]: джинсы\\
0x88B3 \textbf{djis} [\emph{dʒis}]: цифровой\\
0x88C1 \textbf{mris} [\emph{mris}]: демократия\\
0x88C2 \textbf{kris} [\emph{kris}]: секретарь\\
0x88C4 \textbf{pris} [\emph{pris}]: вирус\\
0x88C6 \textbf{nris} [\emph{nris}]: невинный\\
0x88C8 \textbf{hqis} [\emph{ʰŋis}]: шовинист\\
0x88C9 \textbf{sris} [\emph{sris}]: исторический\\
0x88CA \textbf{tris} [\emph{tris}]: пурпурный\\
0x88CC \textbf{fris} [\emph{fris}]: неофициальный\\
0x88CD \textbf{cris} [\emph{ʃris}]: наследник\\
0x88D0 \textbf{hxis} [\emph{ʰxis}]: пиксель\\
0x88D1 \textbf{bris} [\emph{bris}]: беспроводной\\
0x88D2 \textbf{gris} [\emph{gris}]: агрессивный\\
0x88D3 \textbf{dris} [\emph{dris}]: все больше и больше\\
0x88D4 \textbf{zris} [\emph{zris}]: Израиль\\
0x88D5 \textbf{jris} [\emph{ʒris}]: пассивный причастие\\
0x88D6 \textbf{vris} [\emph{vris}]: взаимная голос\\
0x88D9 \textbf{qris} [\emph{ŋris}]: английский\\
0x88DA \textbf{xris} [\emph{xris}]: пользователь интерфейс\\
0x88E2 \textbf{kxis} [\emph{kxis}]: клип Изобразительное искусство\\
0x88E4 \textbf{pxis} [\emph{pxis}]: пассивный голос\\
0x88EA \textbf{txis} [\emph{txis}]: точка матрица\\
0x88F1 \textbf{bxis} [\emph{bxis}]: арбитр\\
0x88F2 \textbf{gxis} [\emph{gxis}]: колючий\\
0x88F3 \textbf{dxis} [\emph{dxis}]: входящие\\
0x8901 \textbf{myas} [\emph{mjäs}]: маска\\
0x8902 \textbf{kyas} [\emph{kjäs}]: лошадь\\
0x8904 \textbf{pyas} [\emph{pjäs}]: фут\\
0x8906 \textbf{nyas} [\emph{njäs}]: сочинение\\
0x8908 \textbf{hmas} [\emph{ʰmäs}]: разум\\
0x8909 \textbf{syas} [\emph{sjäs}]: сроки\\
0x890A \textbf{tyas} [\emph{tjäs}]: Таблица\\
0x890B \textbf{lyas} [\emph{ljäs}]: интерес\\
0x890C \textbf{fyas} [\emph{fjäs}]: ранен\\
0x890D \textbf{cyas} [\emph{ʃjäs}]: медведь\\
0x890E \textbf{ryas} [\emph{rjäs}]: коробка передач\\
0x8910 \textbf{hkas} [\emph{ʰkäs}]: мог\\
0x8911 \textbf{byas} [\emph{bjäs}]: вакцина\\
0x8912 \textbf{gyas} [\emph{gjäs}]: гимнастика\\
0x8913 \textbf{dyas} [\emph{djäs}]: дилер\\
0x8914 \textbf{zyas} [\emph{zjäs}]: Азия\\
0x8915 \textbf{jyas} [\emph{ʒjäs}]: Таджикистан\\
0x8916 \textbf{vyas} [\emph{vjäs}]: верфь\\
0x8918 \textbf{hyas} [\emph{ʰjäs}]: один раз\\
0x8919 \textbf{qyas} [\emph{ŋjäs}]: алхимик\\
0x891A \textbf{xyas} [\emph{xjäs}]: роли\\
0x8920 \textbf{hpas} [\emph{ʰpäs}]: сын\\
0x8921 \textbf{mwas} [\emph{mwäs}]: месяц\\
0x8922 \textbf{kwas} [\emph{kwäs}]: август\\
0x8924 \textbf{pwas} [\emph{pwäs}]: аплодисменты\\
0x8926 \textbf{nwas} [\emph{nwäs}]: Новости\\
0x8928 \textbf{hwas} [\emph{ʰwäs}]: вопрос\\
0x8929 \textbf{swas} [\emph{swäs}]: оказание услуг\\
0x892A \textbf{twas} [\emph{twäs}]: дважды\\
0x892B \textbf{lwas} [\emph{lwäs}]: ленивый\\
0x892C \textbf{fwas} [\emph{fwäs}]: фармацевт\\
0x892D \textbf{cwas} [\emph{ʃwäs}]: начать\\
0x892E \textbf{rwas} [\emph{rwäs}]: роман\\
0x8930 \textbf{hnas} [\emph{ʰnäs}]: не могу\\
0x8931 \textbf{bwas} [\emph{bwäs}]: браузер\\
0x8932 \textbf{gwas} [\emph{gwäs}]: Вау\\
0x8933 \textbf{dwas} [\emph{dwäs}]: посольство\\
0x8934 \textbf{zwas} [\emph{zwäs}]: наркоз\\
0x8936 \textbf{vwas} [\emph{vwäs}]: винный камень\\
0x8939 \textbf{qwas} [\emph{ŋwäs}]: наружу\\
0x893A \textbf{xwas} [\emph{xwäs}]: аист\\
0x8942 \textbf{ksas} [\emph{ksäs}]: 4\\
0x8944 \textbf{psas} [\emph{psäs}]: производить\\
0x8948 \textbf{hsas} [\emph{ʰsäs}]: утро\\
0x894A \textbf{tsas} [\emph{tsäs}]: бросать\\
0x8950 \textbf{htas} [\emph{ʰtäs}]: еще раз\\
0x8951 \textbf{bzas} [\emph{bzäs}]: пресыщенный\\
0x8952 \textbf{gzas} [\emph{gzäs}]: богомол\\
0x8953 \textbf{dzas} [\emph{dzäs}]: вызывать недовольство\\
0x8958 \textbf{hlas} [\emph{ʰläs}]: линия\\
0x8960 \textbf{hfas} [\emph{ʰfäs}]: файл\\
0x8962 \textbf{klas} [\emph{kläs}]: стакан\\
0x8964 \textbf{plas} [\emph{pläs}]: взрыв\\
0x8968 \textbf{hcas} [\emph{ʰʃäs}]: дом\\
0x8969 \textbf{slas} [\emph{släs}]: Изобразительное искусство\\
0x896A \textbf{tlas} [\emph{tläs}]: стоимость\\
0x896C \textbf{flas} [\emph{fläs}]: поклонник\\
0x896D \textbf{clas} [\emph{ʃläs}]: желтый\\
0x8970 \textbf{hras} [\emph{ʰräs}]: песок\\
0x8971 \textbf{blas} [\emph{bläs}]: опора\\
0x8972 \textbf{glas} [\emph{gläs}]: симулировать\\
0x8973 \textbf{dlas} [\emph{dläs}]: Удалить\\
0x8974 \textbf{zlas} [\emph{zläs}]: бег с препятствиями\\
0x8975 \textbf{jlas} [\emph{ʒläs}]: ненадежный\\
0x8976 \textbf{vlas} [\emph{vläs}]: вокализации\\
0x897A \textbf{xlas} [\emph{xläs}]: парализованный\\
0x8982 \textbf{kfas} [\emph{kfäs}]: рак\\
0x8984 \textbf{pfas} [\emph{pfäs}]: причастие\\
0x8988 \textbf{hbas} [\emph{ʰbäs}]: бред какой то\\
0x898A \textbf{tfas} [\emph{tfäs}]: посол\\
0x8990 \textbf{hgas} [\emph{ʰgäs}]: газ\\
0x8991 \textbf{bvas} [\emph{bväs}]: делитель\\
0x8992 \textbf{gvas} [\emph{gväs}]: сорняками\\
0x8993 \textbf{dvas} [\emph{dväs}]: акушерка\\
0x8998 \textbf{hdas} [\emph{ʰdäs}]: танец\\
0x89A0 \textbf{hzas} [\emph{ʰzäs}]: шутить\\
0x89A2 \textbf{kcas} [\emph{kʃäs}]: темнота\\
0x89A4 \textbf{pcas} [\emph{pʃäs}]: взгляд\\
0x89A8 \textbf{hjas} [\emph{ʰʒäs}]: подавлять\\
0x89AA \textbf{tcas} [\emph{tʃäs}]: места\\
0x89B0 \textbf{hvas} [\emph{ʰväs}]: производители\\
0x89B1 \textbf{bjas} [\emph{bʒäs}]: подвид\\
0x89B2 \textbf{gjas} [\emph{gʒäs}]: кактус\\
0x89B3 \textbf{djas} [\emph{dʒäs}]: джаз\\
0x89C1 \textbf{mras} [\emph{mräs}]: зонтик\\
0x89C2 \textbf{kras} [\emph{kräs}]: гроб\\
0x89C4 \textbf{pras} [\emph{präs}]: вуаль\\
0x89C6 \textbf{nras} [\emph{nräs}]: медсестра\\
0x89C8 \textbf{hqas} [\emph{ʰŋäs}]: Веб-сайт\\
0x89C9 \textbf{sras} [\emph{sräs}]: гнездо\\
0x89CA \textbf{tras} [\emph{träs}]: сравнение\\
0x89CC \textbf{fras} [\emph{fräs}]: разброс\\
0x89CD \textbf{cras} [\emph{ʃräs}]: Рубашка\\
0x89D0 \textbf{hxas} [\emph{ʰxäs}]: гепатит\\
0x89D1 \textbf{bras} [\emph{bräs}]: боковая\\
0x89D2 \textbf{gras} [\emph{gräs}]: транс\\
0x89D3 \textbf{dras} [\emph{dräs}]: презирать\\
0x89D4 \textbf{zras} [\emph{zräs}]: Сирия\\
0x89D5 \textbf{jras} [\emph{ʒräs}]: бюрократия\\
0x89D6 \textbf{vras} [\emph{vräs}]: Австралия\\
0x89D9 \textbf{qras} [\emph{ŋräs}]: начальство\\
0x89DA \textbf{xras} [\emph{xräs}]: стирание\\
0x89E2 \textbf{kxas} [\emph{kxäs}]: скорость света\\
0x89E4 \textbf{pxas} [\emph{pxäs}]: колпачки замок\\
0x89EA \textbf{txas} [\emph{txäs}]: установка\\
0x89F1 \textbf{bxas} [\emph{bxäs}]: к юго-западу\\
0x89F2 \textbf{gxas} [\emph{gxäs}]: маяк\\
0x89F3 \textbf{dxas} [\emph{dxäs}]: одаренный\\
0x8A01 \textbf{myus} [\emph{mjus}]: мышь\\
0x8A02 \textbf{kyus} [\emph{kjus}]: пересекать\\
0x8A04 \textbf{pyus} [\emph{pjus}]: предохранитель\\
0x8A06 \textbf{nyus} [\emph{njus}]: индуктор\\
0x8A08 \textbf{hmus} [\emph{ʰmus}]: толстый\\
0x8A09 \textbf{syus} [\emph{sjus}]: субсидировать\\
0x8A0A \textbf{tyus} [\emph{tjus}]: целое число\\
0x8A0B \textbf{lyus} [\emph{ljus}]: целлюлоза\\
0x8A0D \textbf{cyus} [\emph{ʃjus}]: злорадствовать\\
0x8A0E \textbf{ryus} [\emph{rjus}]: Россия\\
0x8A10 \textbf{hkus} [\emph{ʰkus}]: плакать\\
0x8A11 \textbf{byus} [\emph{bjus}]: омовения\\
0x8A12 \textbf{gyus} [\emph{gjus}]: кучевые облака\\
0x8A13 \textbf{dyus} [\emph{djus}]: Будапешт\\
0x8A18 \textbf{hyus} [\emph{ʰjus}]: красивая\\
0x8A1A \textbf{xyus} [\emph{xjus}]: тис\\
0x8A20 \textbf{hpus} [\emph{ʰpus}]: офис\\
0x8A21 \textbf{mwus} [\emph{mwus}]: мейоз\\
0x8A22 \textbf{kwus} [\emph{kwus}]: капсула\\
0x8A24 \textbf{pwus} [\emph{pwus}]: выглядывает\\
0x8A26 \textbf{nwus} [\emph{nwus}]: сонный\\
0x8A27 \textbf{psa7h} [\emph{psä˥ʰ}]: изнашивания\\
0x8A28 \textbf{hwus} [\emph{ʰwus}]: текст\\
0x8A29 \textbf{swus} [\emph{swus}]: самокат\\
0x8A2A \textbf{twus} [\emph{twus}]: альтруизм\\
0x8A2B \textbf{lwus} [\emph{lwus}]: лесс\\
0x8A2E \textbf{rwus} [\emph{rwus}]: гребец\\
0x8A30 \textbf{hnus} [\emph{ʰnus}]: нос\\
0x8A31 \textbf{bwus} [\emph{bwus}]: вопить\\
0x8A32 \textbf{gwus} [\emph{gwus}]: фагоцитоз\\
0x8A3A \textbf{xwus} [\emph{xwus}]: обручи\\
0x8A42 \textbf{ksus} [\emph{ksus}]: кластер\\
0x8A44 \textbf{psus} [\emph{psus}]: петрушка\\
0x8A48 \textbf{hsus} [\emph{ʰsus}]: неделю\\
0x8A4A \textbf{tsus} [\emph{tsus}]: такси\\
0x8A50 \textbf{htus} [\emph{ʰtus}]: государство контекст\\
0x8A57 \textbf{tsa7h} [\emph{tsä˥ʰ}]: antipassive голос\\
0x8A58 \textbf{hlus} [\emph{ʰlus}]: олень\\
0x8A60 \textbf{hfus} [\emph{ʰfus}]: восстановление\\
0x8A62 \textbf{klus} [\emph{klus}]: эвкалипт\\
0x8A64 \textbf{plus} [\emph{plus}]: плюс\\
0x8A68 \textbf{hcus} [\emph{ʰʃus}]: последний\\
0x8A69 \textbf{slus} [\emph{slus}]: доступность\\
0x8A6A \textbf{tlus} [\emph{tlus}]: антиклимакс\\
0x8A6C \textbf{flus} [\emph{flus}]: кальцит\\
0x8A6D \textbf{clus} [\emph{ʃlus}]: веселый\\
0x8A70 \textbf{hrus} [\emph{ʰrus}]: принцесса\\
0x8A71 \textbf{blus} [\emph{blus}]: скорая\\
0x8A72 \textbf{glus} [\emph{glus}]: высокомерный\\
0x8A73 \textbf{dlus} [\emph{dlus}]: Беларусь\\
0x8A76 \textbf{vlus} [\emph{vlus}]: девальвация\\
0x8A7A \textbf{xlus} [\emph{xlus}]: Геркулес\\
0x8A82 \textbf{kfus} [\emph{kfus}]: диффузный\\
0x8A84 \textbf{pfus} [\emph{pfus}]: префикс\\
0x8A88 \textbf{hbus} [\emph{ʰbus}]: поцелуй\\
0x8A8A \textbf{tfus} [\emph{tfus}]: зубная паста\\
0x8A90 \textbf{hgus} [\emph{ʰgus}]: коэффициент\\
0x8A98 \textbf{hdus} [\emph{ʰdus}]: свинья\\
0x8AA0 \textbf{hzus} [\emph{ʰzus}]: Судан\\
0x8AA2 \textbf{kcus} [\emph{kʃus}]: готовить\\
0x8AA4 \textbf{pcus} [\emph{pʃus}]: мак\\
0x8AA8 \textbf{hjus} [\emph{ʰʒus}]: Иисус\\
0x8AAA \textbf{tcus} [\emph{tʃus}]: Сью\\
0x8AB0 \textbf{hvus} [\emph{ʰvus}]: яйценосный\\
0x8AC1 \textbf{mrus} [\emph{mrus}]: Маврикий\\
0x8AC2 \textbf{krus} [\emph{krus}]: приседать\\
0x8AC4 \textbf{prus} [\emph{prus}]: Компьютерный вирус\\
0x8AC6 \textbf{nrus} [\emph{nrus}]: носорог\\
0x8AC8 \textbf{hqus} [\emph{ʰŋus}]: неисчислимый\\
0x8AC9 \textbf{srus} [\emph{srus}]: вытеснять\\
0x8ACA \textbf{trus} [\emph{trus}]: спорный\\
0x8ACC \textbf{frus} [\emph{frus}]: разъедающий\\
0x8ACD \textbf{crus} [\emph{ʃrus}]: иерархия\\
0x8AD0 \textbf{hxus} [\emph{ʰxus}]: открытый источник\\
0x8AD1 \textbf{brus} [\emph{brus}]: зал заседаний совета директоров\\
0x8AD2 \textbf{grus} [\emph{grus}]: урны\\
0x8AD3 \textbf{drus} [\emph{drus}]: Дрю\\
0x8AD5 \textbf{jrus} [\emph{ʒrus}]: юра\\
0x8AD6 \textbf{vrus} [\emph{vrus}]: Евразия\\
0x8AD9 \textbf{qrus} [\emph{ŋrus}]: евхаристия\\
0x8ADA \textbf{xrus} [\emph{xrus}]: гидроксид\\
0x8AE2 \textbf{kxus} [\emph{kxus}]: куратор\\
0x8AE4 \textbf{pxus} [\emph{pxus}]: пестик\\
0x8AEA \textbf{txus} [\emph{txus}]: тушить\\
0x8AF1 \textbf{bxus} [\emph{bxus}]: Узбекистан\\
0x8AF2 \textbf{gxus} [\emph{gxus}]: возможности\\
0x8AF3 \textbf{dxus} [\emph{dxus}]: кукольный дом\\
0x8B01 \textbf{myes} [\emph{mjes}]: майонез\\
0x8B02 \textbf{kyes} [\emph{kjes}]: Комментарии\\
0x8B04 \textbf{pyes} [\emph{pjes}]: гнейс\\
0x8B06 \textbf{nyes} [\emph{njes}]: интеркостельный\\
0x8B08 \textbf{hmes} [\emph{ʰmes}]: мэр\\
0x8B09 \textbf{syes} [\emph{sjes}]: седан\\
0x8B0A \textbf{tyes} [\emph{tjes}]: поставщик провизии\\
0x8B0B \textbf{lyes} [\emph{ljes}]: парикмахер\\
0x8B0C \textbf{fyes} [\emph{fjes}]: Флер-де-Лис\\
0x8B0D \textbf{cyes} [\emph{ʃjes}]: киты\\
0x8B0E \textbf{ryes} [\emph{rjes}]: регистратор\\
0x8B10 \textbf{hkes} [\emph{ʰkes}]: сложный\\
0x8B11 \textbf{byes} [\emph{bjes}]: восемьдесят шесть\\
0x8B12 \textbf{gyes} [\emph{gjes}]: серебристо-серый\\
0x8B13 \textbf{dyes} [\emph{djes}]: диабет\\
0x8B18 \textbf{hyes} [\emph{ʰjes}]: осведомленность\\
0x8B19 \textbf{qyes} [\emph{ŋjes}]: флебит\\
0x8B1A \textbf{xyes} [\emph{xjes}]: Гаити\\
0x8B20 \textbf{hpes} [\emph{ʰpes}]: неприятный\\
0x8B21 \textbf{mwes} [\emph{mwes}]: метан\\
0x8B22 \textbf{kwes} [\emph{kwes}]: конгрессменов\\
0x8B24 \textbf{pwes} [\emph{pwes}]: восемьдесят четыре\\
0x8B26 \textbf{nwes} [\emph{nwes}]: инвестиции\\
0x8B28 \textbf{hwes} [\emph{ʰwes}]: без сознания\\
0x8B29 \textbf{swes} [\emph{swes}]: свитер\\
0x8B2A \textbf{twes} [\emph{twes}]: телевидение\\
0x8B2B \textbf{lwes} [\emph{lwes}]: лазер\\
0x8B2C \textbf{fwes} [\emph{fwes}]: диссертаций\\
0x8B2D \textbf{cwes} [\emph{ʃwes}]: Шлезвиг Гольштейн\\
0x8B2E \textbf{rwes} [\emph{rwes}]: изотопный индикатор\\
0x8B30 \textbf{hnes} [\emph{ʰnes}]: нервный\\
0x8B33 \textbf{dwes} [\emph{dwes}]: вихри\\
0x8B39 \textbf{qwes} [\emph{ŋwes}]: Ви Ви\\
0x8B3A \textbf{xwes} [\emph{xwes}]: семьдесят четыре\\
0x8B42 \textbf{kses} [\emph{kses}]: экскурсия\\
0x8B44 \textbf{pses} [\emph{pses}]: получатель\\
0x8B48 \textbf{hses} [\emph{ʰses}]: 14\\
0x8B4A \textbf{tses} [\emph{tses}]: 13\\
0x8B50 \textbf{htes} [\emph{ʰtes}]: посмотреть\\
0x8B52 \textbf{gzes} [\emph{gzes}]: старикашка\\
0x8B53 \textbf{dzes} [\emph{dzes}]: дизель\\
0x8B58 \textbf{hles} [\emph{ʰles}]: слива\\
0x8B60 \textbf{hfes} [\emph{ʰfes}]: возвратный\\
0x8B62 \textbf{kles} [\emph{kles}]: обезглавить\\
0x8B64 \textbf{ples} [\emph{ples}]: процессор\\
0x8B68 \textbf{hces} [\emph{ʰʃes}]: студент\\
0x8B69 \textbf{sles} [\emph{sles}]: цензор\\
0x8B6A \textbf{tles} [\emph{tles}]: металлоид\\
0x8B6C \textbf{fles} [\emph{fles}]: аналитик\\
0x8B6D \textbf{cles} [\emph{ʃles}]: реализм\\
0x8B70 \textbf{hres} [\emph{ʰres}]: походить\\
0x8B71 \textbf{bles} [\emph{bles}]: браслет\\
0x8B72 \textbf{gles} [\emph{gles}]: шпалера\\
0x8B73 \textbf{dles} [\emph{dles}]: гидролиз\\
0x8B74 \textbf{zles} [\emph{zles}]: безногий\\
0x8B75 \textbf{jles} [\emph{ʒles}]: газетный штамп\\
0x8B76 \textbf{vles} [\emph{vles}]: валлийский\\
0x8B7A \textbf{xles} [\emph{xles}]: Гималаи\\
0x8B82 \textbf{kfes} [\emph{kfes}]: ломонос\\
0x8B84 \textbf{pfes} [\emph{pfes}]: бледное лицо\\
0x8B88 \textbf{hbes} [\emph{ʰbes}]: баннер\\
0x8B8A \textbf{tfes} [\emph{tfes}]: фарш\\
0x8B90 \textbf{hges} [\emph{ʰges}]: беременность\\
0x8B91 \textbf{bves} [\emph{bves}]: Bavaria\\
0x8B98 \textbf{hdes} [\emph{ʰdes}]: база данных\\
0x8BA0 \textbf{hzes} [\emph{ʰzes}]: одержимый дело\\
0x8BA2 \textbf{kces} [\emph{kʃes}]: кассета\\
0x8BA4 \textbf{pces} [\emph{pʃes}]: блин\\
0x8BA8 \textbf{hjes} [\emph{ʰʒes}]: перепись\\
0x8BAA \textbf{tces} [\emph{tʃes}]: большой теннис\\
0x8BB0 \textbf{hves} [\emph{ʰves}]: венецианский\\
0x8BB1 \textbf{bjes} [\emph{bʒes}]: житница\\
0x8BC1 \textbf{mres} [\emph{mres}]: сжатие\\
0x8BC2 \textbf{kres} [\emph{kres}]: глагол\\
0x8BC4 \textbf{pres} [\emph{pres}]: кора\\
0x8BC6 \textbf{nres} [\emph{nres}]: целостность\\
0x8BC8 \textbf{hqes} [\emph{ʰŋes}]: сексуальный\\
0x8BC9 \textbf{sres} [\emph{sres}]: С уважением\\
0x8BCA \textbf{tres} [\emph{tres}]: стресс\\
0x8BCC \textbf{fres} [\emph{fres}]: мобилизовывать средства\\
0x8BCD \textbf{cres} [\emph{ʃres}]: розничная торговля\\
0x8BD0 \textbf{hxes} [\emph{ʰxes}]: герц\\
0x8BD1 \textbf{bres} [\emph{bres}]: слово процессор\\
0x8BD2 \textbf{gres} [\emph{gres}]: кресс\\
0x8BD3 \textbf{dres} [\emph{dres}]: распаковывать\\
0x8BD4 \textbf{zres} [\emph{zres}]: вездесущность\\
0x8BD5 \textbf{jres} [\emph{ʒres}]: ночная рубашка\\
0x8BD6 \textbf{vres} [\emph{vres}]: электролиз\\
0x8BD9 \textbf{qres} [\emph{ŋres}]: внушающих благоговение\\
0x8BDA \textbf{xres} [\emph{xres}]: поперечный\\
0x8BE2 \textbf{kxes} [\emph{kxes}]: смородина\\
0x8BE4 \textbf{pxes} [\emph{pxes}]: таз\\
0x8BEA \textbf{txes} [\emph{txes}]: коллектив имя существительное\\
0x8BF1 \textbf{bxes} [\emph{bxes}]: удар слева\\
0x8BF2 \textbf{gxes} [\emph{gxes}]: после чего\\
0x8BF3 \textbf{dxes} [\emph{dxes}]: гастроном\\
0x8C01 \textbf{myos} [\emph{mjos}]: метеор\\
0x8C02 \textbf{kyos} [\emph{kjos}]: коммунист\\
0x8C04 \textbf{pyos} [\emph{pjos}]: премолярный\\
0x8C06 \textbf{nyos} [\emph{njos}]: неамбициозны\\
0x8C08 \textbf{hmos} [\emph{ʰmos}]: министр\\
0x8C09 \textbf{syos} [\emph{sjos}]: подведение\\
0x8C0A \textbf{tyos} [\emph{tjos}]: доверенное лицо\\
0x8C0B \textbf{lyos} [\emph{ljos}]: плеоназм\\
0x8C0C \textbf{fyos} [\emph{fjos}]: этос\\
0x8C0D \textbf{cyos} [\emph{ʃjos}]: суматоха\\
0x8C0E \textbf{ryos} [\emph{rjos}]: ахроматический\\
0x8C10 \textbf{hkos} [\emph{ʰkos}]: каноэ\\
0x8C11 \textbf{byos} [\emph{bjos}]: нетоксичен\\
0x8C12 \textbf{gyos} [\emph{gjos}]: глухой звук\\
0x8C13 \textbf{dyos} [\emph{djos}]: динозавр\\
0x8C18 \textbf{hyos} [\emph{ʰjos}]: опасающийся\\
0x8C1A \textbf{xyos} [\emph{xjos}]: рекомбинация\\
0x8C20 \textbf{hpos} [\emph{ʰpos}]: сосна\\
0x8C21 \textbf{mwos} [\emph{mwos}]: симбиоз\\
0x8C22 \textbf{kwos} [\emph{kwos}]: по часовой стрелке\\
0x8C24 \textbf{pwos} [\emph{pwos}]: протонов\\
0x8C26 \textbf{nwos} [\emph{nwos}]: инновация\\
0x8C28 \textbf{hwos} [\emph{ʰwos}]: Спальня\\
0x8C29 \textbf{swos} [\emph{swos}]: экспоненциальный\\
0x8C2A \textbf{twos} [\emph{twos}]: терроризм\\
0x8C2B \textbf{lwos} [\emph{lwos}]: проложенный\\
0x8C2D \textbf{cwos} [\emph{ʃwos}]: челнок\\
0x8C2E \textbf{rwos} [\emph{rwos}]: аэрозоль\\
0x8C30 \textbf{hnos} [\emph{ʰnos}]: наводнение\\
0x8C31 \textbf{bwos} [\emph{bwos}]: боксеров\\
0x8C32 \textbf{gwos} [\emph{gwos}]: гонады\\
0x8C33 \textbf{dwos} [\emph{dwos}]: кузовостроение\\
0x8C3A \textbf{xwos} [\emph{xwos}]: Осия\\
0x8C42 \textbf{ksos} [\emph{ksos}]: кашель\\
0x8C44 \textbf{psos} [\emph{psos}]: заграничный пасспорт\\
0x8C48 \textbf{hsos} [\emph{ʰsos}]: шестнадцать\\
0x8C4A \textbf{tsos} [\emph{tsos}]: конкурс\\
0x8C50 \textbf{htos} [\emph{ʰtos}]: большинство\\
0x8C51 \textbf{bzos} [\emph{bzos}]: бозон\\
0x8C53 \textbf{dzos} [\emph{dzos}]: дремать\\
0x8C58 \textbf{hlos} [\emph{ʰlos}]: вверх по лестнице\\
0x8C60 \textbf{hfos} [\emph{ʰfos}]: факс\\
0x8C62 \textbf{klos} [\emph{klos}]: хлор\\
0x8C64 \textbf{plos} [\emph{plos}]: пульс\\
0x8C68 \textbf{hcos} [\emph{ʰʃos}]: меры\\
0x8C69 \textbf{slos} [\emph{slos}]: колбаса\\
0x8C6A \textbf{tlos} [\emph{tlos}]: триллион\\
0x8C6C \textbf{flos} [\emph{flos}]: флорист\\
0x8C6D \textbf{clos} [\emph{ʃlos}]: рыботорговец\\
0x8C70 \textbf{hros} [\emph{ʰros}]: репортер\\
0x8C71 \textbf{blos} [\emph{blos}]: автобус\\
0x8C72 \textbf{glos} [\emph{glos}]: гладиолус\\
0x8C73 \textbf{dlos} [\emph{dlos}]: модальный дело\\
0x8C74 \textbf{zlos} [\emph{zlos}]: масленка\\
0x8C75 \textbf{jlos} [\emph{ʒlos}]: почтовый ящик\\
0x8C7A \textbf{xlos} [\emph{xlos}]: зеленщик\\
0x8C82 \textbf{kfos} [\emph{kfos}]: компостирование\\
0x8C84 \textbf{pfos} [\emph{pfos}]: профессор\\
0x8C88 \textbf{hbos} [\emph{ʰbos}]: лаборатория\\
0x8C8A \textbf{tfos} [\emph{tfos}]: застава\\
0x8C90 \textbf{hgos} [\emph{ʰgos}]: ржавчина\\
0x8C98 \textbf{hdos} [\emph{ʰdos}]: пятно\\
0x8CA0 \textbf{hzos} [\emph{ʰzos}]: Казахстан\\
0x8CA2 \textbf{kcos} [\emph{kʃos}]: обвинить\\
0x8CA4 \textbf{pcos} [\emph{pʃos}]: виды спорта\\
0x8CA8 \textbf{hjos} [\emph{ʰʒos}]: без присмотра\\
0x8CAA \textbf{tcos} [\emph{tʃos}]: урожай\\
0x8CB0 \textbf{hvos} [\emph{ʰvos}]: оазис\\
0x8CB2 \textbf{gjos} [\emph{gʒos}]: галактоза\\
0x8CB3 \textbf{djos} [\emph{dʒos}]: союзы\\
0x8CC1 \textbf{mros} [\emph{mros}]: монитор\\
0x8CC2 \textbf{kros} [\emph{kros}]: кварц\\
0x8CC4 \textbf{pros} [\emph{pros}]: операционная\\
0x8CC6 \textbf{nros} [\emph{nros}]: контрастные\\
0x8CC8 \textbf{hqos} [\emph{ʰŋos}]: унция\\
0x8CC9 \textbf{sros} [\emph{sros}]: меньшинство\\
0x8CCA \textbf{tros} [\emph{tros}]: тост\\
0x8CCC \textbf{fros} [\emph{fros}]: бирюзовый\\
0x8CCD \textbf{cros} [\emph{ʃros}]: хор\\
0x8CD0 \textbf{hxos} [\emph{ʰxos}]: торус\\
0x8CD1 \textbf{bros} [\emph{bros}]: барокко\\
0x8CD2 \textbf{gros} [\emph{gros}]: плечевая кость\\
0x8CD3 \textbf{dros} [\emph{dros}]: дозировать\\
0x8CD4 \textbf{zros} [\emph{zros}]: Азорские острова\\
0x8CD5 \textbf{jros} [\emph{ʒros}]: Коморские\\
0x8CD6 \textbf{vros} [\emph{vros}]: эспрессо\\
0x8CD9 \textbf{qros} [\emph{ŋros}]: поддерживать\\
0x8CDA \textbf{xros} [\emph{xros}]: неоднородный\\
0x8CE2 \textbf{kxos} [\emph{kxos}]: плотоядный\\
0x8CE4 \textbf{pxos} [\emph{pxos}]: Почта\\
0x8CEA \textbf{txos} [\emph{txos}]: средний голос\\
0x8CF1 \textbf{bxos} [\emph{bxos}]: потребители\\
0x8CF2 \textbf{gxos} [\emph{gxos}]: органист\\
0x8CF3 \textbf{dxos} [\emph{dxos}]: магнетит\\
0x8D01 \textbf{my6s} [\emph{mjəs}]: муза\\
0x8D02 \textbf{ky6s} [\emph{kjəs}]: питомники\\
0x8D04 \textbf{py6s} [\emph{pjəs}]: поверхностно\\
0x8D06 \textbf{ny6s} [\emph{njəs}]: неоклассический\\
0x8D08 \textbf{hm6s} [\emph{ʰməs}]: Мадагаскар\\
0x8D09 \textbf{sy6s} [\emph{sjəs}]: нездоровый\\
0x8D0A \textbf{ty6s} [\emph{tjəs}]: подземные толчки\\
0x8D0C \textbf{fy6s} [\emph{fjəs}]: близкий к народу\\
0x8D0E \textbf{ry6s} [\emph{rjəs}]: Киргизия\\
0x8D10 \textbf{hk6s} [\emph{ʰkəs}]: рождество\\
0x8D18 \textbf{hy6s} [\emph{ʰjəs}]: Румыния\\
0x8D20 \textbf{hp6s} [\emph{ʰpəs}]: протоплазма\\
0x8D28 \textbf{hw6s} [\emph{ʰwəs}]: эвфемизм\\
0x8D29 \textbf{sw6s} [\emph{swəs}]: особы\\
0x8D2A \textbf{tw6s} [\emph{twəs}]: партеногенез\\
0x8D2B \textbf{lw6s} [\emph{lwəs}]: долго windedness\\
0x8D30 \textbf{hn6s} [\emph{ʰnəs}]: Ренессанс\\
0x8D42 \textbf{ks6s} [\emph{ksəs}]: Пакистан\\
0x8D44 \textbf{ps6s} [\emph{psəs}]: макаронные изделия\\
0x8D48 \textbf{hs6s} [\emph{ʰsəs}]: лужи\\
0x8D4A \textbf{ts6s} [\emph{tsəs}]: транс цифра\\
0x8D50 \textbf{ht6s} [\emph{ʰtəs}]: аутизм\\
0x8D51 \textbf{bz6s} [\emph{bzəs}]: OSE\\
0x8D58 \textbf{hl6s} [\emph{ʰləs}]: лотос\\
0x8D60 \textbf{hf6s} [\emph{ʰfəs}]: фракталы\\
0x8D62 \textbf{kl6s} [\emph{kləs}]: крейсерский\\
0x8D64 \textbf{pl6s} [\emph{pləs}]: ананас\\
0x8D68 \textbf{hc6s} [\emph{ʰʃəs}]: паразитизм\\
0x8D69 \textbf{sl6s} [\emph{sləs}]: цельсию\\
0x8D6A \textbf{tl6s} [\emph{tləs}]: интерполяция\\
0x8D6C \textbf{fl6s} [\emph{fləs}]: Фолклендские\\
0x8D70 \textbf{hr6s} [\emph{ʰrəs}]: направления\\
0x8D71 \textbf{bl6s} [\emph{bləs}]: целибат\\
0x8D74 \textbf{zl6s} [\emph{zləs}]: кутерьмы на ближний свет\\
0x8D7A \textbf{xl6s} [\emph{xləs}]: гиперинфляция\\
0x8D88 \textbf{hb6s} [\emph{ʰbəs}]: жужжание\\
0x8D8A \textbf{tf6s} [\emph{tfəs}]: две трети\\
0x8D90 \textbf{hg6s} [\emph{ʰgəs}]: потуги\\
0x8D98 \textbf{hd6s} [\emph{ʰdəs}]: дьякон\\
0x8DA0 \textbf{hz6s} [\emph{ʰzəs}]: позитивизм\\
0x8DA2 \textbf{kc6s} [\emph{kʃəs}]: восьмеричный\\
0x8DA8 \textbf{hj6s} [\emph{ʰʒəs}]: гигантизм\\
0x8DAA \textbf{tc6s} [\emph{tʃəs}]: Лихтенштейн\\
0x8DC1 \textbf{mr6s} [\emph{mrəs}]: подлокотник\\
0x8DC2 \textbf{kr6s} [\emph{krəs}]: кокос\\
0x8DC4 \textbf{pr6s} [\emph{prəs}]: перс\\
0x8DC6 \textbf{nr6s} [\emph{nrəs}]: антропоморфизм\\
0x8DC8 \textbf{hq6s} [\emph{ʰŋəs}]: наоборот\\
0x8DC9 \textbf{sr6s} [\emph{srəs}]: Кипр\\
0x8DCA \textbf{tr6s} [\emph{trəs}]: живородящий\\
0x8DCC \textbf{fr6s} [\emph{frəs}]: ferous\\
0x8DCD \textbf{cr6s} [\emph{ʃrəs}]: ворсянка\\
0x8DD0 \textbf{hx6s} [\emph{ʰxəs}]: сталкиваться\\
0x8DD1 \textbf{br6s} [\emph{brəs}]: Бухарест\\
0x8DD2 \textbf{gr6s} [\emph{grəs}]: охлократия\\
0x8DD3 \textbf{dr6s} [\emph{drəs}]: пристойный\\
0x8DDA \textbf{xr6s} [\emph{xrəs}]: иероглифы\\
0x8DE2 \textbf{kx6s} [\emph{kxəs}]: декартовский\\
0x8DE4 \textbf{px6s} [\emph{pxəs}]: проходки\\
0x8DEA \textbf{tx6s} [\emph{txəs}]: Туркменистан\\
0x8E57 \textbf{tra7h} [\emph{trä˥ʰ}]: природный\\
\subsection{ 9 }
0x9001 \textbf{myi7s} [\emph{mji˥s}]: амид\\
0x9002 \textbf{kyi7s} [\emph{kji˥s}]: кретинизм\\
0x9004 \textbf{pyi7s} [\emph{pji˥s}]: дарохранительница\\
0x9006 \textbf{nyi7s} [\emph{nji˥s}]: перемежаются\\
0x9008 \textbf{hmi7s} [\emph{ʰmi˥s}]: Миссисипи\\
0x9009 \textbf{syi7s} [\emph{sji˥s}]: осознанный\\
0x900A \textbf{tyi7s} [\emph{tji˥s}]: шарик\\
0x900B \textbf{lyi7s} [\emph{lji˥s}]: рысь\\
0x900C \textbf{fyi7s} [\emph{fji˥s}]: не фантастика\\
0x900E \textbf{ryi7s} [\emph{rji˥s}]: проверяемый\\
0x9010 \textbf{hki7s} [\emph{ʰki˥s}]: сородичи\\
0x9013 \textbf{dyi7s} [\emph{dji˥s}]: дис\\
0x9018 \textbf{hyi7s} [\emph{ʰji˥s}]: тюремщик\\
0x9020 \textbf{hpi7s} [\emph{ʰpi˥s}]: Пруссия\\
0x9021 \textbf{mwi7s} [\emph{mwi˥s}]: марксист\\
0x9026 \textbf{nwi7s} [\emph{nwi˥s}]: канонизация\\
0x9028 \textbf{hwi7s} [\emph{ʰwi˥s}]: еси\\
0x9029 \textbf{swi7s} [\emph{swi˥s}]: в сторону моря\\
0x902A \textbf{twi7s} [\emph{twi˥s}]: пятидесятый\\
0x902B \textbf{lwi7s} [\emph{lwi˥s}]: Linsey Вулси\\
0x9030 \textbf{hni7s} [\emph{ʰni˥s}]: отказ от ответственности\\
0x9042 \textbf{ksi7s} [\emph{ksi˥s}]: софистика\\
0x9044 \textbf{psi7s} [\emph{psi˥s}]: эликсир\\
0x9048 \textbf{hsi7s} [\emph{ʰsi˥s}]: эгоизм\\
0x904A \textbf{tsi7s} [\emph{tsi˥s}]: двадцать три\\
0x9050 \textbf{hti7s} [\emph{ʰti˥s}]: незатребованный\\
0x9053 \textbf{dzi7s} [\emph{dzi˥s}]: дионисическое\\
0x9058 \textbf{hli7s} [\emph{ʰli˥s}]: поясница\\
0x9060 \textbf{hfi7s} [\emph{ʰfi˥s}]: юриспруденция\\
0x9062 \textbf{kli7s} [\emph{kli˥s}]: выслушивание\\
0x9064 \textbf{pli7s} [\emph{pli˥s}]: неправдоподобно\\
0x9068 \textbf{hci7s} [\emph{ʰʃi˥s}]: универсально\\
0x9069 \textbf{sli7s} [\emph{sli˥s}]: пятьдесят шесть\\
0x906A \textbf{tli7s} [\emph{tli˥s}]: двоек\\
0x906C \textbf{fli7s} [\emph{fli˥s}]: формализм\\
0x9070 \textbf{hri7s} [\emph{ʰri˥s}]: директора\\
0x9071 \textbf{bli7s} [\emph{bli˥s}]: обелиск\\
0x9073 \textbf{dli7s} [\emph{dli˥s}]: индол\\
0x907A \textbf{xli7s} [\emph{xli˥s}]: нигилизм\\
0x9082 \textbf{kfi7s} [\emph{kfi˥s}]: конформист\\
0x9084 \textbf{pfi7s} [\emph{pfi˥s}]: гипофиз\\
0x9088 \textbf{hbi7s} [\emph{ʰbi˥s}]: крещение\\
0x908A \textbf{tfi7s} [\emph{tfi˥s}]: дурачество\\
0x9090 \textbf{hgi7s} [\emph{ʰgi˥s}]: цыганский\\
0x9098 \textbf{hdi7s} [\emph{ʰdi˥s}]: адреса\\
0x90A0 \textbf{hzi7s} [\emph{ʰzi˥s}]: близорукость\\
0x90A2 \textbf{kci7s} [\emph{kʃi˥s}]: архаизм\\
0x90A8 \textbf{hji7s} [\emph{ʰʒi˥s}]: землевладельческий\\
0x90AA \textbf{tci7s} [\emph{tʃi˥s}]: шестерки\\
0x90B0 \textbf{hvi7s} [\emph{ʰvi˥s}]: визирь\\
0x90B3 \textbf{dji7s} [\emph{dʒi˥s}]: Изе\\
0x90C1 \textbf{mri7s} [\emph{mri˥s}]: материализм\\
0x90C2 \textbf{kri7s} [\emph{kri˥s}]: замковый камень\\
0x90C4 \textbf{pri7s} [\emph{pri˥s}]: Париж\\
0x90C6 \textbf{nri7s} [\emph{nri˥s}]: неинтересный\\
0x90C8 \textbf{hqi7s} [\emph{ʰŋi˥s}]: слоновость\\
0x90C9 \textbf{sri7s} [\emph{sri˥s}]: писец\\
0x90CA \textbf{tri7s} [\emph{tri˥s}]: зеркала\\
0x90CC \textbf{fri7s} [\emph{fri˥s}]: теоретизировать\\
0x90CD \textbf{cri7s} [\emph{ʃri˥s}]: милитаризм\\
0x90D0 \textbf{hxi7s} [\emph{ʰxi˥s}]: место отдыха\\
0x90D1 \textbf{bri7s} [\emph{bri˥s}]: иберийский\\
0x90D2 \textbf{gri7s} [\emph{gri˥s}]: взаимные услуги\\
0x90D3 \textbf{dri7s} [\emph{dri˥s}]: устрашение\\
0x90DA \textbf{xri7s} [\emph{xri˥s}]: соглядатай\\
0x90E2 \textbf{kxi7s} [\emph{kxi˥s}]: рубец\\
0x90E4 \textbf{pxi7s} [\emph{pxi˥s}]: насаженные\\
0x90EA \textbf{txi7s} [\emph{txi˥s}]: оптимист\\
0x90F1 \textbf{bxi7s} [\emph{bxi˥s}]: абхазия\\
0x90F2 \textbf{gxi7s} [\emph{gxi˥s}]: гинеколог\\
0x90F3 \textbf{dxi7s} [\emph{dxi˥s}]: IST\\
0x9101 \textbf{mya7s} [\emph{mjä˥s}]: монада\\
0x9102 \textbf{kya7s} [\emph{kjä˥s}]: панегирист\\
0x9106 \textbf{nya7s} [\emph{njä˥s}]: евангелистский\\
0x9108 \textbf{hma7s} [\emph{ʰmä˥s}]: металлы\\
0x9109 \textbf{sya7s} [\emph{sjä˥s}]: схоластика\\
0x910A \textbf{tya7s} [\emph{tjä˥s}]: клеветнический\\
0x910B \textbf{lya7s} [\emph{ljä˥s}]: Плеяды\\
0x910E \textbf{rya7s} [\emph{rjä˥s}]: английская лапта\\
0x9110 \textbf{hka7s} [\emph{ʰkä˥s}]: квазар\\
0x9111 \textbf{bya7s} [\emph{bjä˥s}]: Аббасидов\\
0x9118 \textbf{hya7s} [\emph{ʰjä˥s}]: Исайя\\
0x911A \textbf{xya7s} [\emph{xjä˥s}]: самостоятельно напористой\\
0x9120 \textbf{hpa7s} [\emph{ʰpä˥s}]: князья\\
0x9121 \textbf{mwa7s} [\emph{mwä˥s}]: медиана\\
0x9122 \textbf{kwa7s} [\emph{kwä˥s}]: десть\\
0x9126 \textbf{nwa7s} [\emph{nwä˥s}]: тонзиллит\\
0x9128 \textbf{hwa7s} [\emph{ʰwä˥s}]: вальс\\
0x9129 \textbf{swa7s} [\emph{swä˥s}]: защелками\\
0x912A \textbf{twa7s} [\emph{twä˥s}]: девяносто три\\
0x912D \textbf{cwa7s} [\emph{ʃwä˥s}]: отродье\\
0x912E \textbf{rwa7s} [\emph{rwä˥s}]: карикатурист\\
0x9130 \textbf{hna7s} [\emph{ʰnä˥s}]: Неаполь\\
0x9142 \textbf{ksa7s} [\emph{ksä˥s}]: корь\\
0x9144 \textbf{psa7s} [\emph{psä˥s}]: причуды\\
0x9148 \textbf{hsa7s} [\emph{ʰsä˥s}]: саксонский\\
0x914A \textbf{tsa7s} [\emph{tsä˥s}]: шелест\\
0x9150 \textbf{hta7s} [\emph{ʰtä˥s}]: формальной одежды\\
0x9151 \textbf{bza7s} [\emph{bzä˥s}]: базел\\
0x9158 \textbf{hla7s} [\emph{ʰlä˥s}]: Альпы\\
0x9160 \textbf{hfa7s} [\emph{ʰfä˥s}]: ФЛАНДРИЯ\\
0x9162 \textbf{kla7s} [\emph{klä˥s}]: скалы\\
0x9164 \textbf{pla7s} [\emph{plä˥s}]: орудийный огонь\\
0x9168 \textbf{hca7s} [\emph{ʰʃä˥s}]: телята\\
0x9169 \textbf{sla7s} [\emph{slä˥s}]: обжора\\
0x916A \textbf{tla7s} [\emph{tlä˥s}]: сорок три\\
0x916C \textbf{fla7s} [\emph{flä˥s}]: фальцетом\\
0x9170 \textbf{hra7s} [\emph{ʰrä˥s}]: проза\\
0x9171 \textbf{bla7s} [\emph{blä˥s}]: бацилла\\
0x9172 \textbf{gla7s} [\emph{glä˥s}]: желчный камень\\
0x9173 \textbf{dla7s} [\emph{dlä˥s}]: напыщенность\\
0x917A \textbf{xla7s} [\emph{xlä˥s}]: желто-зеленый\\
0x9184 \textbf{pfa7s} [\emph{pfä˥s}]: нераспространение\\
0x9188 \textbf{hba7s} [\emph{ʰbä˥s}]: рынки\\
0x918A \textbf{tfa7s} [\emph{tfä˥s}]: катарсис\\
0x9190 \textbf{hga7s} [\emph{ʰgä˥s}]: певцы\\
0x9198 \textbf{hda7s} [\emph{ʰdä˥s}]: кинжал\\
0x91A0 \textbf{hza7s} [\emph{ʰzä˥s}]: Заир\\
0x91A2 \textbf{kca7s} [\emph{kʃä˥s}]: Кавказ\\
0x91A4 \textbf{pca7s} [\emph{pʃä˥s}]: похитители\\
0x91A8 \textbf{hja7s} [\emph{ʰʒä˥s}]: вычеркивать\\
0x91AA \textbf{tca7s} [\emph{tʃä˥s}]: характерно\\
0x91B0 \textbf{hva7s} [\emph{ʰvä˥s}]: цианоз\\
0x91C1 \textbf{mra7s} [\emph{mrä˥s}]: хлопчатобумажная ткань в полоску\\
0x91C2 \textbf{kra7s} [\emph{krä˥s}]: колесный\\
0x91C4 \textbf{pra7s} [\emph{prä˥s}]: арии\\
0x91C6 \textbf{nra7s} [\emph{nrä˥s}]: наковальня\\
0x91C8 \textbf{hqa7s} [\emph{ʰŋä˥s}]: сумасшедший дом\\
0x91C9 \textbf{sra7s} [\emph{srä˥s}]: спаржа\\
0x91CA \textbf{tra7s} [\emph{trä˥s}]: три четверти\\
0x91CC \textbf{fra7s} [\emph{frä˥s}]: эвтаназия\\
0x91CD \textbf{cra7s} [\emph{ʃrä˥s}]: сперматозоиды\\
0x91D0 \textbf{hxa7s} [\emph{ʰxä˥s}]: Увы\\
0x91D1 \textbf{bra7s} [\emph{brä˥s}]: альбатрос\\
0x91D2 \textbf{gra7s} [\emph{grä˥s}]: гаррота\\
0x91D3 \textbf{dra7s} [\emph{drä˥s}]: трема\\
0x91DA \textbf{xra7s} [\emph{xrä˥s}]: слезящимися\\
0x91E2 \textbf{kxa7s} [\emph{kxä˥s}]: заминированный\\
0x91E4 \textbf{pxa7s} [\emph{pxä˥s}]: неудачник\\
0x91EA \textbf{txa7s} [\emph{txä˥s}]: дрок\\
0x91F1 \textbf{bxa7s} [\emph{bxä˥s}]: торговец табачными изделиями\\
0x91F3 \textbf{dxa7s} [\emph{dxä˥s}]: лозоискательство\\
0x9208 \textbf{hmu7s} [\emph{ʰmu˥s}]: превью\\
0x9210 \textbf{hku7s} [\emph{ʰku˥s}]: кодициль\\
0x9218 \textbf{hyu7s} [\emph{ʰju˥s}]: язычок\\
0x9220 \textbf{hpu7s} [\emph{ʰpu˥s}]: палеозойский\\
0x9230 \textbf{hnu7s} [\emph{ʰnu˥s}]: Навуходоносор\\
0x9242 \textbf{ksu7s} [\emph{ksu˥s}]: арбалет\\
0x9244 \textbf{psu7s} [\emph{psu˥s}]: метловище\\
0x9248 \textbf{hsu7s} [\emph{ʰsu˥s}]: вшей\\
0x924A \textbf{tsu7s} [\emph{tsu˥s}]: антацидный\\
0x9250 \textbf{htu7s} [\emph{ʰtu˥s}]: трепетный\\
0x9258 \textbf{hlu7s} [\emph{ʰlu˥s}]: толстой кожурой\\
0x9262 \textbf{klu7s} [\emph{klu˥s}]: исключительность\\
0x9270 \textbf{hru7s} [\emph{ʰru˥s}]: Брюссель\\
0x9290 \textbf{hgu7s} [\emph{ʰgu˥s}]: Гвинея Бисау\\
0x9298 \textbf{hdu7s} [\emph{ʰdu˥s}]: сочиться\\
0x92AA \textbf{tcu7s} [\emph{tʃu˥s}]: актуарий\\
0x92C1 \textbf{mru7s} [\emph{mru˥s}]: урал\\
0x92C2 \textbf{kru7s} [\emph{kru˥s}]: Водолей\\
0x92C6 \textbf{nru7s} [\emph{nru˥s}]: гондурасский\\
0x92C9 \textbf{sru7s} [\emph{sru˥s}]: Астурийский\\
0x92CA \textbf{tru7s} [\emph{tru˥s}]: импровизировать\\
0x92D0 \textbf{hxu7s} [\emph{ʰxu˥s}]: эвристический\\
0x9308 \textbf{hme7s} [\emph{ʰme˥s}]: медуза\\
0x9309 \textbf{sye7s} [\emph{sje˥s}]: Seychelles\\
0x9310 \textbf{hke7s} [\emph{ʰke˥s}]: кратеры\\
0x9318 \textbf{hye7s} [\emph{ʰje˥s}]: бесчувственный\\
0x9320 \textbf{hpe7s} [\emph{ʰpe˥s}]: Пегас\\
0x9328 \textbf{hwe7s} [\emph{ʰwe˥s}]: кружевной\\
0x9330 \textbf{hne7s} [\emph{ʰne˥s}]: нравоучительный\\
0x9342 \textbf{kse7s} [\emph{kse˥s}]: квинтэссенция\\
0x9344 \textbf{pse7s} [\emph{pse˥s}]: плейстоцен\\
0x9348 \textbf{hse7s} [\emph{ʰse˥s}]: сорок семь\\
0x934A \textbf{tse7s} [\emph{tse˥s}]: техасец\\
0x9350 \textbf{hte7s} [\emph{ʰte˥s}]: патриотизм\\
0x9358 \textbf{hle7s} [\emph{ʰle˥s}]: лосины\\
0x9360 \textbf{hfe7s} [\emph{ʰfe˥s}]: Еф\\
0x9368 \textbf{hce7s} [\emph{ʰʃe˥s}]: тузы\\
0x936A \textbf{tle7s} [\emph{tle˥s}]: катализ\\
0x9370 \textbf{hre7s} [\emph{ʰre˥s}]: Фракия\\
0x9388 \textbf{hbe7s} [\emph{ʰbe˥s}]: косички\\
0x938A \textbf{tfe7s} [\emph{tfe˥s}]: пасынок\\
0x9390 \textbf{hge7s} [\emph{ʰge˥s}]: магнезит\\
0x9398 \textbf{hde7s} [\emph{ʰde˥s}]: сестра жены\\
0x93AA \textbf{tce7s} [\emph{tʃe˥s}]: дегустатор\\
0x93B0 \textbf{hve7s} [\emph{ʰve˥s}]: ветряная оспа\\
0x93C1 \textbf{mre7s} [\emph{mre˥s}]: матрицы\\
0x93C2 \textbf{kre7s} [\emph{kre˥s}]: азартная игра в кости\\
0x93C4 \textbf{pre7s} [\emph{pre˥s}]: Пиренеи\\
0x93C6 \textbf{nre7s} [\emph{nre˥s}]: лестничная площадка\\
0x93C9 \textbf{sre7s} [\emph{sre˥s}]: заклинатель\\
0x93CA \textbf{tre7s} [\emph{tre˥s}]: этрусский\\
0x93CC \textbf{fre7s} [\emph{fre˥s}]: антифедерализм\\
0x93D0 \textbf{hxe7s} [\emph{ʰxe˥s}]: Гессе\\
0x93D1 \textbf{bre7s} [\emph{bre˥s}]: либреттист\\
0x93DA \textbf{xre7s} [\emph{xre˥s}]: о лице\\
0x93EA \textbf{txe7s} [\emph{txe˥s}]: поймать как улова может\\
0x9408 \textbf{hmo7s} [\emph{ʰmo˥s}]: Моисей\\
0x940B \textbf{lyo7s} [\emph{ljo˥s}]: солипсизм\\
0x940E \textbf{ryo7s} [\emph{rjo˥s}]: ориентализм\\
0x9410 \textbf{hko7s} [\emph{ʰko˥s}]: корсаж\\
0x9418 \textbf{hyo7s} [\emph{ʰjo˥s}]: самоуправления обладали\\
0x9420 \textbf{hpo7s} [\emph{ʰpo˥s}]: Прометей\\
0x9429 \textbf{swo7s} [\emph{swo˥s}]: сказать так\\
0x942A \textbf{two7s} [\emph{two˥s}]: Тоголезская\\
0x942B \textbf{lwo7s} [\emph{lwo˥s}]: все желающие\\
0x9430 \textbf{hno7s} [\emph{ʰno˥s}]: безоружный\\
0x9442 \textbf{kso7s} [\emph{kso˥s}]: коробка офис\\
0x9444 \textbf{pso7s} [\emph{pso˥s}]: опоссум\\
0x9448 \textbf{hso7s} [\emph{ʰso˥s}]: шумно\\
0x944A \textbf{tso7s} [\emph{tso˥s}]: девяносто четыре\\
0x9450 \textbf{hto7s} [\emph{ʰto˥s}]: Томас\\
0x9458 \textbf{hlo7s} [\emph{ʰlo˥s}]: начитанный\\
0x9460 \textbf{hfo7s} [\emph{ʰfo˥s}]: бесплатно для всех\\
0x9462 \textbf{klo7s} [\emph{klo˥s}]: оккультист\\
0x9464 \textbf{plo7s} [\emph{plo˥s}]: прозелитизм\\
0x9469 \textbf{slo7s} [\emph{slo˥s}]: тусклый\\
0x946A \textbf{tlo7s} [\emph{tlo˥s}]: траулер\\
0x9470 \textbf{hro7s} [\emph{ʰro˥s}]: чалый\\
0x9488 \textbf{hbo7s} [\emph{ʰbo˥s}]: удав\\
0x9490 \textbf{hgo7s} [\emph{ʰgo˥s}]: вирши\\
0x9498 \textbf{hdo7s} [\emph{ʰdo˥s}]: двойники\\
0x94AA \textbf{tco7s} [\emph{tʃo˥s}]: пятно на солнце\\
0x94C1 \textbf{mro7s} [\emph{mro˥s}]: микропроцессор\\
0x94C2 \textbf{kro7s} [\emph{kro˥s}]: курсор\\
0x94C6 \textbf{nro7s} [\emph{nro˥s}]: антимонопольный\\
0x94C9 \textbf{sro7s} [\emph{sro˥s}]: слишком любящий свою жену\\
0x94CA \textbf{tro7s} [\emph{tro˥s}]: Тироль\\
0x94CC \textbf{fro7s} [\emph{fro˥s}]: Фарерская\\
0x94D0 \textbf{hxo7s} [\emph{ʰxo˥s}]: мимоза\\
0x94DA \textbf{xro7s} [\emph{xro˥s}]: гидрокси\\
0x94E2 \textbf{kxo7s} [\emph{kxo˥s}]: лактоза\\
0x94EA \textbf{txo7s} [\emph{txo˥s}]: чаны\\
0x9508 \textbf{hm67s} [\emph{ʰmə˥s}]: мальтузианский\\
0x9517 \textbf{kcu7h} [\emph{kʃu˥ʰ}]: индуктивный настроение\\
0x9517 \textbf{kcuh} [\emph{kʃuʰ}]: choice gender, 3rd density life form, learning to make choices, choosing whether to be service to self, service to others, or to remain lost, akin to basic homo sapiens, cetaceans (dolphins/whales), advanced animals\\
0x9520 \textbf{hp67s} [\emph{ʰpə˥s}]: ханжество\\
0x9528 \textbf{hw67s} [\emph{ʰwə˥s}]: ширь\\
0x9530 \textbf{hn67s} [\emph{ʰnə˥s}]: анархизм\\
0x9548 \textbf{hs67s} [\emph{ʰsə˥s}]: анастомоз\\
0x9558 \textbf{hl67s} [\emph{ʰlə˥s}]: эльзасский\\
0x9570 \textbf{hr67s} [\emph{ʰrə˥s}]: тромбоз\\
0x9801 \textbf{myi2s} [\emph{mji˩s}]: медиально\\
0x9802 \textbf{kyi2s} [\emph{kji˩s}]: сексизм\\
0x9806 \textbf{nyi2s} [\emph{nji˩s}]: самовлюбленный человек\\
0x9808 \textbf{hmi2s} [\emph{ʰmi˩s}]: чудотворный\\
0x9809 \textbf{syi2s} [\emph{sji˩s}]: псориаз\\
0x980A \textbf{tyi2s} [\emph{tji˩s}]: шестидесятый\\
0x980B \textbf{lyi2s} [\emph{lji˩s}]: олигоцен\\
0x9810 \textbf{hki2s} [\emph{ʰki˩s}]: сорок шесть\\
0x9818 \textbf{hyi2s} [\emph{ʰji˩s}]: допустимый\\
0x9820 \textbf{hpi2s} [\emph{ʰpi˩s}]: язычество\\
0x9821 \textbf{mwi2s} [\emph{mwi˩s}]: монооксид\\
0x9826 \textbf{nwi2s} [\emph{nwi˩s}]: кровосмешение\\
0x9828 \textbf{hwi2s} [\emph{ʰwi˩s}]: руды\\
0x9830 \textbf{hni2s} [\emph{ʰni˩s}]: Николай\\
0x9842 \textbf{ksi2s} [\emph{ksi˩s}]: акрополь\\
0x9844 \textbf{psi2s} [\emph{psi˩s}]: предлог\\
0x9848 \textbf{hsi2s} [\emph{ʰsi˩s}]: Полегче\\
0x9850 \textbf{hti2s} [\emph{ʰti˩s}]: презрительный\\
0x9858 \textbf{hli2s} [\emph{ʰli˩s}]: злобный взгляд\\
0x9860 \textbf{hfi2s} [\emph{ʰfi˩s}]: теософия\\
0x9862 \textbf{kli2s} [\emph{kli˩s}]: акклиматизировались\\
0x9864 \textbf{pli2s} [\emph{pli˩s}]: эпилепсия\\
0x9868 \textbf{hci2s} [\emph{ʰʃi˩s}]: шедевр\\
0x9869 \textbf{sli2s} [\emph{sli˩s}]: двадцать шесть\\
0x986A \textbf{tli2s} [\emph{tli˩s}]: артериосклероз\\
0x986C \textbf{fli2s} [\emph{fli˩s}]: обывательский\\
0x9870 \textbf{hri2s} [\emph{ʰri˩s}]: двадцать четыре\\
0x9871 \textbf{bli2s} [\emph{bli˩s}]: черный список\\
0x9873 \textbf{dli2s} [\emph{dli˩s}]: дуалистический\\
0x987A \textbf{xli2s} [\emph{xli˩s}]: приличествовать\\
0x9888 \textbf{hbi2s} [\emph{ʰbi˩s}]: туземцы\\
0x988A \textbf{tfi2s} [\emph{tfi˩s}]: тридцать три\\
0x9890 \textbf{hgi2s} [\emph{ʰgi˩s}]: глазурь\\
0x9898 \textbf{hdi2s} [\emph{ʰdi˩s}]: иды\\
0x98A0 \textbf{hzi2s} [\emph{ʰzi˩s}]: гностицизм\\
0x98A2 \textbf{kci2s} [\emph{kʃi˩s}]: штампами\\
0x98A4 \textbf{pci2s} [\emph{pʃi˩s}]: импрессионизм\\
0x98A8 \textbf{hji2s} [\emph{ʰʒi˩s}]: приносить несчастье\\
0x98AA \textbf{tci2s} [\emph{tʃi˩s}]: семидесятый\\
0x98B0 \textbf{hvi2s} [\emph{ʰvi˩s}]: наиболее важные части\\
0x98C1 \textbf{mri2s} [\emph{mri˩s}]: , себя\\
0x98C2 \textbf{kri2s} [\emph{kri˩s}]: курсы\\
0x98C4 \textbf{pri2s} [\emph{pri˩s}]: восемьдесят три\\
0x98C6 \textbf{nri2s} [\emph{nri˩s}]: нарцисс\\
0x98C8 \textbf{hqi2s} [\emph{ʰŋi˩s}]: Рыбы\\
0x98C9 \textbf{sri2s} [\emph{sri˩s}]: вяжущий\\
0x98CA \textbf{tri2s} [\emph{tri˩s}]: повторять\\
0x98D0 \textbf{hxi2s} [\emph{ʰxi˩s}]: суффикс\\
0x98D1 \textbf{bri2s} [\emph{bri˩s}]: вылечить все\\
0x98D2 \textbf{gri2s} [\emph{gri˩s}]: гастрит\\
0x98D3 \textbf{dri2s} [\emph{dri˩s}]: децентрализация\\
0x98D9 \textbf{qri2s} [\emph{ŋri˩s}]: цирроз печени\\
0x98DA \textbf{xri2s} [\emph{xri˩s}]: шарлатанство\\
0x98E2 \textbf{kxi2s} [\emph{kxi˩s}]: разъездной священник\\
0x98E4 \textbf{pxi2s} [\emph{pxi˩s}]: свинарник\\
0x98EA \textbf{txi2s} [\emph{txi˩s}]: семьдесят три\\
0x98F2 \textbf{gxi2s} [\emph{gxi˩s}]: гри гри\\
0x9901 \textbf{mya2s} [\emph{mjä˩s}]: массажист\\
0x9906 \textbf{nya2s} [\emph{njä˩s}]: монашество\\
0x9908 \textbf{hma2s} [\emph{ʰmä˩s}]: Марсель\\
0x9909 \textbf{sya2s} [\emph{sjä˩s}]: стареющий\\
0x990A \textbf{tya2s} [\emph{tjä˩s}]: катехизис\\
0x990B \textbf{lya2s} [\emph{ljä˩s}]: бельевой\\
0x990E \textbf{rya2s} [\emph{rjä˩s}]: Королевская битва\\
0x9910 \textbf{hka2s} [\emph{ʰkä˩s}]: коаксиальный\\
0x9918 \textbf{hya2s} [\emph{ʰjä˩s}]: соблазнять\\
0x9920 \textbf{hpa2s} [\emph{ʰpä˩s}]: Пифагор\\
0x9922 \textbf{kwa2s} [\emph{kwä˩s}]: чешуйчатый\\
0x9924 \textbf{pwa2s} [\emph{pwä˩s}]: меблировать\\
0x9928 \textbf{hwa2s} [\emph{ʰwä˩s}]: Варшавская\\
0x9929 \textbf{swa2s} [\emph{swä˩s}]: сагиттальный\\
0x992A \textbf{twa2s} [\emph{twä˩s}]: синица\\
0x992E \textbf{rwa2s} [\emph{rwä˩s}]: трещотка\\
0x9930 \textbf{hna2s} [\emph{ʰnä˩s}]: династия\\
0x9942 \textbf{ksa2s} [\emph{ksä˩s}]: непостижимый\\
0x9944 \textbf{psa2s} [\emph{psä˩s}]: пароксизм\\
0x9948 \textbf{hsa2s} [\emph{ʰsä˩s}]: сольный концерт\\
0x994A \textbf{tsa2s} [\emph{tsä˩s}]: задумчивый\\
0x9950 \textbf{hta2s} [\emph{ʰtä˩s}]: взносы\\
0x9958 \textbf{hla2s} [\emph{ʰlä˩s}]: всплеснула\\
0x9960 \textbf{hfa2s} [\emph{ʰfä˩s}]: фаланга\\
0x9962 \textbf{kla2s} [\emph{klä˩s}]: календы\\
0x9964 \textbf{pla2s} [\emph{plä˩s}]: ипостась\\
0x9968 \textbf{hca2s} [\emph{ʰʃä˩s}]: канцелярия\\
0x9969 \textbf{sla2s} [\emph{slä˩s}]: водовод\\
0x996A \textbf{tla2s} [\emph{tlä˩s}]: атлас\\
0x996C \textbf{fla2s} [\emph{flä˩s}]: анафилаксия\\
0x996D \textbf{cla2s} [\emph{ʃlä˩s}]: колит\\
0x9970 \textbf{hra2s} [\emph{ʰrä˩s}]: преступления\\
0x9971 \textbf{bla2s} [\emph{blä˩s}]: балюстрада\\
0x9973 \textbf{dla2s} [\emph{dlä˩s}]: деполяризация\\
0x997A \textbf{xla2s} [\emph{xlä˩s}]: Фары\\
0x9984 \textbf{pfa2s} [\emph{pfä˩s}]: апофеоз\\
0x9988 \textbf{hba2s} [\emph{ʰbä˩s}]: аббатиса\\
0x998A \textbf{tfa2s} [\emph{tfä˩s}]: метатеза\\
0x9990 \textbf{hga2s} [\emph{ʰgä˩s}]: гвардия\\
0x9998 \textbf{hda2s} [\emph{ʰdä˩s}]: алый\\
0x99A2 \textbf{kca2s} [\emph{kʃä˩s}]: ачехский\\
0x99AA \textbf{tca2s} [\emph{tʃä˩s}]: бессменный\\
0x99B0 \textbf{hva2s} [\emph{ʰvä˩s}]: парообразный\\
0x99C1 \textbf{mra2s} [\emph{mrä˩s}]: амбра\\
0x99C2 \textbf{kra2s} [\emph{krä˩s}]: фураж\\
0x99C4 \textbf{pra2s} [\emph{prä˩s}]: папирус\\
0x99C6 \textbf{nra2s} [\emph{nrä˩s}]: анналы\\
0x99C8 \textbf{hqa2s} [\emph{ʰŋä˩s}]: абсцисса\\
0x99C9 \textbf{sra2s} [\emph{srä˩s}]: неквалифицированный\\
0x99CA \textbf{tra2s} [\emph{trä˩s}]: устало тащиться\\
0x99CD \textbf{cra2s} [\emph{ʃrä˩s}]: маразм\\
0x99D0 \textbf{hxa2s} [\emph{ʰxä˩s}]: Далида\\
0x99D1 \textbf{bra2s} [\emph{brä˩s}]: бюстгальтеры\\
0x99D3 \textbf{dra2s} [\emph{drä˩s}]: изображенный\\
0x99DA \textbf{xra2s} [\emph{xrä˩s}]: торговец трикотажными изделиями\\
0x99E2 \textbf{kxa2s} [\emph{kxä˩s}]: рахит\\
0x99E4 \textbf{pxa2s} [\emph{pxä˩s}]: судья на линии\\
0x99EA \textbf{txa2s} [\emph{txä˩s}]: ходули\\
0x9A08 \textbf{hmu2s} [\emph{ʰmu˩s}]: Мафусаил\\
0x9A10 \textbf{hku2s} [\emph{ʰku˩s}]: камбуз\\
0x9A18 \textbf{hyu2s} [\emph{ʰju˩s}]: мочевина\\
0x9A30 \textbf{hnu2s} [\emph{ʰnu˩s}]: Индостан\\
0x9A42 \textbf{ksu2s} [\emph{ksu˩s}]: квесты\\
0x9A44 \textbf{psu2s} [\emph{psu˩s}]: пластика\\
0x9A48 \textbf{hsu2s} [\emph{ʰsu˩s}]: Советский Союз\\
0x9A4A \textbf{tsu2s} [\emph{tsu˩s}]: трипсин\\
0x9A50 \textbf{htu2s} [\emph{ʰtu˩s}]: кулачная драка\\
0x9A58 \textbf{hlu2s} [\emph{ʰlu˩s}]: эманация\\
0x9A70 \textbf{hru2s} [\emph{ʰru˩s}]: заключать в себе\\
0x9AC6 \textbf{nru2s} [\emph{nru˩s}]: неврит\\
0x9ACA \textbf{tru2s} [\emph{tru˩s}]: плутократия\\
0x9AD0 \textbf{hxu2s} [\emph{ʰxu˩s}]: мотыги\\
0x9AEA \textbf{txu2s} [\emph{txu˩s}]: очень большой размер\\
0x9B08 \textbf{hme2s} [\emph{ʰme˩s}]: метемпсихоз\\
0x9B10 \textbf{hke2s} [\emph{ʰke˩s}]: Экклезиаст\\
0x9B20 \textbf{hpe2s} [\emph{ʰpe˩s}]: Пелопоннес\\
0x9B28 \textbf{hwe2s} [\emph{ʰwe˩s}]: снята с охраны\\
0x9B30 \textbf{hne2s} [\emph{ʰne˩s}]: нелицензированный\\
0x9B42 \textbf{kse2s} [\emph{kse˩s}]: превосходить\\
0x9B44 \textbf{pse2s} [\emph{pse˩s}]: прецессия\\
0x9B48 \textbf{hse2s} [\emph{ʰse˩s}]: показной\\
0x9B4A \textbf{tse2s} [\emph{tse˩s}]: тридцатый\\
0x9B50 \textbf{hte2s} [\emph{ʰte˩s}]: Темза\\
0x9B58 \textbf{hle2s} [\emph{ʰle˩s}]: в любом случае\\
0x9B60 \textbf{hfe2s} [\emph{ʰfe˩s}]: франки стрелков\\
0x9B6A \textbf{tle2s} [\emph{tle˩s}]: экстернализация\\
0x9B70 \textbf{hre2s} [\emph{ʰre˩s}]: смола\\
0x9BA4 \textbf{pce2s} [\emph{pʃe˩s}]: сепсис\\
0x9BC1 \textbf{mre2s} [\emph{mre˩s}]: памятные предметы\\
0x9BC6 \textbf{nre2s} [\emph{nre˩s}]: негритянка\\
0x9BC9 \textbf{sre2s} [\emph{sre˩s}]: течка\\
0x9BCA \textbf{tre2s} [\emph{tre˩s}]: теократия\\
0x9BD0 \textbf{hxe2s} [\emph{ʰxe˩s}]: семьдесят шесть\\
0x9BDA \textbf{xre2s} [\emph{xre˩s}]: Ура\\
0x9BEA \textbf{txe2s} [\emph{txe˩s}]: отвислый\\
0x9C08 \textbf{hmo2s} [\emph{ʰmo˩s}]: миоцен\\
0x9C10 \textbf{hko2s} [\emph{ʰko˩s}]: созвездий\\
0x9C18 \textbf{hyo2s} [\emph{ʰjo˩s}]: цветных\\
0x9C20 \textbf{hpo2s} [\emph{ʰpo˩s}]: причалы\\
0x9C30 \textbf{hno2s} [\emph{ʰno˩s}]: неумело\\
0x9C42 \textbf{kso2s} [\emph{kso˩s}]: акростих\\
0x9C48 \textbf{hso2s} [\emph{ʰso˩s}]: мускус\\
0x9C4A \textbf{tso2s} [\emph{tso˩s}]: фотосинтез\\
0x9C50 \textbf{hto2s} [\emph{ʰto˩s}]: катафалк\\
0x9C58 \textbf{hlo2s} [\emph{ʰlo˩s}]: лоббист\\
0x9C62 \textbf{klo2s} [\emph{klo˩s}]: сколиоз\\
0x9C69 \textbf{slo2s} [\emph{slo˩s}]: все отводы\\
0x9C6A \textbf{tlo2s} [\emph{tlo˩s}]: нитроглицерин\\
0x9C70 \textbf{hro2s} [\emph{ʰro˩s}]: Родос\\
0x9CA0 \textbf{hzo2s} [\emph{ʰzo˩s}]: вы UNS\\
0x9CAA \textbf{tco2s} [\emph{tʃo˩s}]: трихинеллез\\
0x9CC1 \textbf{mro2s} [\emph{mro˩s}]: геморрой\\
0x9CC2 \textbf{kro2s} [\emph{kro˩s}]: пробки\\
0x9CC4 \textbf{pro2s} [\emph{pro˩s}]: пироксен\\
0x9CC6 \textbf{nro2s} [\emph{nro˩s}]: некроз\\
0x9CC9 \textbf{sro2s} [\emph{sro˩s}]: зоб\\
0x9CCA \textbf{tro2s} [\emph{tro˩s}]: Телец\\
0x9CD0 \textbf{hxo2s} [\emph{ʰxo˩s}]: гаубица\\
0x9D30 \textbf{hn62s} [\emph{ʰnə˩s}]: анорексия\\
0x9D4A \textbf{ts62s} [\emph{tsə˩s}]: триоксида\\
\subsection{ A }
0xA801 \textbf{myit} [\emph{mjit}]: зима\\
0xA802 \textbf{kyit} [\emph{kjit}]: процитировал\\
0xA804 \textbf{pyit} [\emph{pjit}]: вежливый регистр\\
0xA806 \textbf{nyit} [\emph{njit}]: Сортировать\\
0xA808 \textbf{hmit} [\emph{ʰmit}]: море\\
0xA809 \textbf{syit} [\emph{sjit}]: город\\
0xA80A \textbf{tyit} [\emph{tjit}]: против\\
0xA80B \textbf{lyit} [\emph{ljit}]: записывать\\
0xA80C \textbf{fyit} [\emph{fjit}]: знакомые регистр\\
0xA80D \textbf{cyit} [\emph{ʃjit}]: сердце\\
0xA80E \textbf{ryit} [\emph{rjit}]: квитанция\\
0xA810 \textbf{hkit} [\emph{ʰkit}]: Эксклюзивный или\\
0xA811 \textbf{byit} [\emph{bjit}]: праздник\\
0xA812 \textbf{gyit} [\emph{gjit}]: умышленно\\
0xA813 \textbf{dyit} [\emph{djit}]: промышленность\\
0xA814 \textbf{zyit} [\emph{zjit}]: должностное лицо\\
0xA815 \textbf{jyit} [\emph{ʒjit}]: невозмутимость\\
0xA816 \textbf{vyit} [\emph{vjit}]: разносторонний\\
0xA818 \textbf{hyit} [\emph{ʰjit}]: REALiS настроение\\
0xA819 \textbf{qyit} [\emph{ŋjit}]: критическая точка\\
0xA81A \textbf{xyit} [\emph{xjit}]: ингибирование\\
0xA820 \textbf{hpit} [\emph{ʰpit}]: белый\\
0xA821 \textbf{mwit} [\emph{mwit}]: commissive настроение\\
0xA822 \textbf{kwit} [\emph{kwit}]: личность\\
0xA824 \textbf{pwit} [\emph{pwit}]: просительный настроение\\
0xA826 \textbf{nwit} [\emph{nwit}]: Событийный настроение\\
0xA828 \textbf{hwit} [\emph{ʰwit}]: вернуть\\
0xA829 \textbf{swit} [\emph{swit}]: повелительный настроение\\
0xA82A \textbf{twit} [\emph{twit}]: утвердительный квантификатором\\
0xA82B \textbf{lwit} [\emph{lwit}]: литература\\
0xA82C \textbf{fwit} [\emph{fwit}]: утвердительный настроение\\
0xA82D \textbf{cwit} [\emph{ʃwit}]: живой\\
0xA82E \textbf{rwit} [\emph{rwit}]: соперник\\
0xA830 \textbf{hnit} [\emph{ʰnit}]: инфинитив\\
0xA831 \textbf{bwit} [\emph{bwit}]: граница\\
0xA832 \textbf{gwit} [\emph{gwit}]: позвоночник\\
0xA833 \textbf{dwit} [\emph{dwit}]: слуховой доказательный\\
0xA834 \textbf{zwit} [\emph{zwit}]: разговорчивый\\
0xA835 \textbf{jwit} [\emph{ʒwit}]: разнообразие\\
0xA836 \textbf{vwit} [\emph{vwit}]: жидкости\\
0xA839 \textbf{qwit} [\emph{ŋwit}]: антисептик\\
0xA83A \textbf{xwit} [\emph{xwit}]: питаясь\\
0xA842 \textbf{ksit} [\emph{ksit}]: могила\\
0xA844 \textbf{psit} [\emph{psit}]: присутствие\\
0xA848 \textbf{hsit} [\emph{ʰsit}]: эпистемическое настроение\\
0xA84A \textbf{tsit} [\emph{tsit}]: остановился\\
0xA850 \textbf{htit} [\emph{ʰtit}]: Это\\
0xA851 \textbf{bzit} [\emph{bzit}]: паук\\
0xA852 \textbf{gzit} [\emph{gzit}]: беременна\\
0xA853 \textbf{dzit} [\emph{dzit}]: СПИД\\
0xA858 \textbf{hlit} [\emph{ʰlit}]: волевых настроение\\
0xA860 \textbf{hfit} [\emph{ʰfit}]: каузальный доказательный\\
0xA862 \textbf{klit} [\emph{klit}]: неизвестный\\
0xA864 \textbf{plit} [\emph{plit}]: здание\\
0xA868 \textbf{hcit} [\emph{ʰʃit}]: имя прилагательное\\
0xA869 \textbf{slit} [\emph{slit}]: владелец\\
0xA86A \textbf{tlit} [\emph{tlit}]: второй\\
0xA86C \textbf{flit} [\emph{flit}]: поля\\
0xA86D \textbf{clit} [\emph{ʃlit}]: карман\\
0xA870 \textbf{hrit} [\emph{ʰrit}]: кредит\\
0xA871 \textbf{blit} [\emph{blit}]: ограниченный\\
0xA872 \textbf{glit} [\emph{glit}]: стабильность\\
0xA873 \textbf{dlit} [\emph{dlit}]: лист\\
0xA874 \textbf{zlit} [\emph{zlit}]: irrealis настроение\\
0xA875 \textbf{jlit} [\emph{ʒlit}]: единодушное\\
0xA876 \textbf{vlit} [\emph{vlit}]: говорливость\\
0xA87A \textbf{xlit} [\emph{xlit}]: инвалидная коляска\\
0xA882 \textbf{kfit} [\emph{kfit}]: лифт\\
0xA884 \textbf{pfit} [\emph{pfit}]: представитель\\
0xA888 \textbf{hbit} [\emph{ʰbit}]: избежать\\
0xA88A \textbf{tfit} [\emph{tfit}]: возврат\\
0xA890 \textbf{hgit} [\emph{ʰgit}]: свидетель\\
0xA891 \textbf{bvit} [\emph{bvit}]: неразвитой\\
0xA892 \textbf{gvit} [\emph{gvit}]: упорство\\
0xA893 \textbf{dvit} [\emph{dvit}]: чрезмерное продолжительность\\
0xA898 \textbf{hdit} [\emph{ʰdit}]: благословение настроение\\
0xA8A0 \textbf{hzit} [\emph{ʰzit}]: функции\\
0xA8A2 \textbf{kcit} [\emph{kʃit}]: в браке\\
0xA8A4 \textbf{pcit} [\emph{pʃit}]: Авторские права\\
0xA8A8 \textbf{hjit} [\emph{ʰʒit}]: Палец\\
0xA8AA \textbf{tcit} [\emph{tʃit}]: совещательный настроение\\
0xA8B0 \textbf{hvit} [\emph{ʰvit}]: директива\\
0xA8B1 \textbf{bjit} [\emph{bʒit}]: бывший умерший\\
0xA8B2 \textbf{gjit} [\emph{gʒit}]: ножницы\\
0xA8B3 \textbf{djit} [\emph{dʒit}]: систематическое\\
0xA8C1 \textbf{mrit} [\emph{mrit}]: предположительный настроение\\
0xA8C2 \textbf{krit} [\emph{krit}]: читать корректуру\\
0xA8C4 \textbf{prit} [\emph{prit}]: проклинающий настроение\\
0xA8C6 \textbf{nrit} [\emph{nrit}]: энергия\\
0xA8C8 \textbf{hqit} [\emph{ʰŋit}]: Спасатель\\
0xA8C9 \textbf{srit} [\emph{srit}]: инструмент\\
0xA8CA \textbf{trit} [\emph{trit}]: непосредственный доказательный\\
0xA8CC \textbf{frit} [\emph{frit}]: порог\\
0xA8CD \textbf{crit} [\emph{ʃrit}]: заброшенный\\
0xA8D0 \textbf{hxit} [\emph{ʰxit}]: максимизировать\\
0xA8D1 \textbf{brit} [\emph{brit}]: кровавый\\
0xA8D2 \textbf{grit} [\emph{grit}]: гром\\
0xA8D3 \textbf{drit} [\emph{drit}]: борода\\
0xA8D4 \textbf{zrit} [\emph{zrit}]: формальный регистр\\
0xA8D5 \textbf{jrit} [\emph{ʒrit}]: оправдывать\\
0xA8D6 \textbf{vrit} [\emph{vrit}]: литературный раб\\
0xA8D9 \textbf{qrit} [\emph{ŋrit}]: Скриншот\\
0xA8DA \textbf{xrit} [\emph{xrit}]: Главная страница\\
0xA8E2 \textbf{kxit} [\emph{kxit}]: ненормальный\\
0xA8E4 \textbf{pxit} [\emph{pxit}]: IP адрес\\
0xA8EA \textbf{txit} [\emph{txit}]: шрифт\\
0xA8F1 \textbf{bxit} [\emph{bxit}]: Web хозяин\\
0xA8F2 \textbf{gxit} [\emph{gxit}]: геноцид\\
0xA8F3 \textbf{dxit} [\emph{dxit}]: драматично\\
0xA901 \textbf{myat} [\emph{mjät}]: запомнить\\
0xA902 \textbf{kyat} [\emph{kjät}]: в самом деле\\
0xA904 \textbf{pyat} [\emph{pjät}]: друг\\
0xA906 \textbf{nyat} [\emph{njät}]: атака\\
0xA908 \textbf{hmat} [\emph{ʰmät}]: демонстративный\\
0xA909 \textbf{syat} [\emph{sjät}]: до\\
0xA90A \textbf{tyat} [\emph{tjät}]: старый\\
0xA90B \textbf{lyat} [\emph{ljät}]: письмо\\
0xA90C \textbf{fyat} [\emph{fjät}]: суд\\
0xA90D \textbf{cyat} [\emph{ʃjät}]: видимый доказательный\\
0xA90E \textbf{ryat} [\emph{rjät}]: трагедия\\
0xA910 \textbf{hkat} [\emph{ʰkät}]: мимо напряженный\\
0xA911 \textbf{byat} [\emph{bjät}]: сожжен\\
0xA912 \textbf{gyat} [\emph{gjät}]: шпион\\
0xA913 \textbf{dyat} [\emph{djät}]: приятель\\
0xA914 \textbf{zyat} [\emph{zjät}]: бритва\\
0xA915 \textbf{jyat} [\emph{ʒjät}]: банкрот\\
0xA916 \textbf{vyat} [\emph{vjät}]: неповрежденный\\
0xA918 \textbf{hyat} [\emph{ʰjät}]: родительский\\
0xA919 \textbf{qyat} [\emph{ŋjät}]: песчаник\\
0xA91A \textbf{xyat} [\emph{xjät}]: антибиотик\\
0xA920 \textbf{hpat} [\emph{ʰpät}]: человек\\
0xA921 \textbf{mwat} [\emph{mwät}]: вкус\\
0xA922 \textbf{kwat} [\emph{kwät}]: комфорт\\
0xA924 \textbf{pwat} [\emph{pwät}]: постель\\
0xA926 \textbf{nwat} [\emph{nwät}]: расти\\
0xA928 \textbf{hwat} [\emph{ʰwät}]: вопрос слово\\
0xA929 \textbf{swat} [\emph{swät}]: меч\\
0xA92A \textbf{twat} [\emph{twät}]: война\\
0xA92B \textbf{lwat} [\emph{lwät}]: список\\
0xA92C \textbf{fwat} [\emph{fwät}]: больница\\
0xA92D \textbf{cwat} [\emph{ʃwät}]: воздух\\
0xA92E \textbf{rwat} [\emph{rwät}]: галстук\\
0xA930 \textbf{hnat} [\emph{ʰnät}]: настоящее время напряженный\\
0xA931 \textbf{bwat} [\emph{bwät}]: одеяло\\
0xA932 \textbf{gwat} [\emph{gwät}]: слоновая кость\\
0xA933 \textbf{dwat} [\emph{dwät}]: лекарственное средство\\
0xA934 \textbf{zwat} [\emph{zwät}]: привычный\\
0xA935 \textbf{jwat} [\emph{ʒwät}]: тротуар\\
0xA936 \textbf{vwat} [\emph{vwät}]: Руанда\\
0xA939 \textbf{qwat} [\emph{ŋwät}]: количество из вещество\\
0xA93A \textbf{xwat} [\emph{xwät}]: мимо совершенный\\
0xA942 \textbf{ksat} [\emph{ksät}]: приятный\\
0xA944 \textbf{psat} [\emph{psät}]: вместо\\
0xA948 \textbf{hsat} [\emph{ʰsät}]: видеть\\
0xA94A \textbf{tsat} [\emph{tsät}]: потенциал настроение\\
0xA950 \textbf{htat} [\emph{ʰtät}]: бабушка или дедушка\\
0xA951 \textbf{bzat} [\emph{bzät}]: актер\\
0xA952 \textbf{gzat} [\emph{gzät}]: газета\\
0xA953 \textbf{dzat} [\emph{dzät}]: дистанционный пульт мимо напряженный\\
0xA958 \textbf{hlat} [\emph{ʰlät}]: факультативный квантификатором\\
0xA960 \textbf{hfat} [\emph{ʰfät}]: послал\\
0xA962 \textbf{klat} [\emph{klät}]: капитан\\
0xA964 \textbf{plat} [\emph{plät}]: стена\\
0xA968 \textbf{hcat} [\emph{ʰʃät}]: близость\\
0xA969 \textbf{slat} [\emph{slät}]: холм\\
0xA96A \textbf{tlat} [\emph{tlät}]: слово\\
0xA96C \textbf{flat} [\emph{flät}]: кислота\\
0xA96D \textbf{clat} [\emph{ʃlät}]: дети\\
0xA970 \textbf{hrat} [\emph{ʰrät}]: крыша\\
0xA971 \textbf{blat} [\emph{blät}]: алтарь\\
0xA972 \textbf{glat} [\emph{glät}]: простота\\
0xA973 \textbf{dlat} [\emph{dlät}]: сегодняшний напряженный\\
0xA974 \textbf{zlat} [\emph{zlät}]: перчатка\\
0xA975 \textbf{jlat} [\emph{ʒlät}]: пират\\
0xA976 \textbf{vlat} [\emph{vlät}]: квартиры\\
0xA97A \textbf{xlat} [\emph{xlät}]: оценочный\\
0xA982 \textbf{kfat} [\emph{kfät}]: сметь\\
0xA984 \textbf{pfat} [\emph{pfät}]: доска\\
0xA988 \textbf{hbat} [\emph{ʰbät}]: слепой\\
0xA98A \textbf{tfat} [\emph{tfät}]: двенадцать\\
0xA990 \textbf{hgat} [\emph{ʰgät}]: пустыня\\
0xA991 \textbf{bvat} [\emph{bvät}]: жестоком обращении\\
0xA992 \textbf{gvat} [\emph{gvät}]: тяготение\\
0xA993 \textbf{dvat} [\emph{dvät}]: раздражительный\\
0xA998 \textbf{hdat} [\emph{ʰdät}]: мост\\
0xA9A0 \textbf{hzat} [\emph{ʰzät}]: волна\\
0xA9A2 \textbf{kcat} [\emph{kʃät}]: вместе\\
0xA9A4 \textbf{pcat} [\emph{pʃät}]: Понимаю\\
0xA9A8 \textbf{hjat} [\emph{ʰʒät}]: человечество\\
0xA9AA \textbf{tcat} [\emph{tʃät}]: будущее напряженный\\
0xA9B0 \textbf{hvat} [\emph{ʰvät}]: предстоящий\\
0xA9B1 \textbf{bjat} [\emph{bʒät}]: пособие\\
0xA9B2 \textbf{gjat} [\emph{gʒät}]: величина\\
0xA9B3 \textbf{djat} [\emph{dʒät}]: неспособность\\
0xA9C1 \textbf{mrat} [\emph{mrät}]: смешанный\\
0xA9C2 \textbf{krat} [\emph{krät}]: император\\
0xA9C4 \textbf{prat} [\emph{prät}]: сообщается доказательный\\
0xA9C6 \textbf{nrat} [\emph{nrät}]: невеста\\
0xA9C8 \textbf{hqat} [\emph{ʰŋät}]: время обеда\\
0xA9C9 \textbf{srat} [\emph{srät}]: звезда\\
0xA9CA \textbf{trat} [\emph{trät}]: Распечатать\\
0xA9CC \textbf{frat} [\emph{frät}]: назначенный\\
0xA9CD \textbf{crat} [\emph{ʃrät}]: суббота\\
0xA9D0 \textbf{hxat} [\emph{ʰxät}]: застежка\\
0xA9D1 \textbf{brat} [\emph{brät}]: возражение\\
0xA9D2 \textbf{grat} [\emph{grät}]: арестовать\\
0xA9D3 \textbf{drat} [\emph{drät}]: зависеть\\
0xA9D4 \textbf{zrat} [\emph{zrät}]: кандидат\\
0xA9D5 \textbf{jrat} [\emph{ʒrät}]: искусственный\\
0xA9D6 \textbf{vrat} [\emph{vrät}]: почтальон\\
0xA9D9 \textbf{qrat} [\emph{ŋrät}]: вверх ногами\\
0xA9DA \textbf{xrat} [\emph{xrät}]: шорты\\
0xA9E2 \textbf{kxat} [\emph{kxät}]: скандий\\
0xA9E4 \textbf{pxat} [\emph{pxät}]: киоск\\
0xA9EA \textbf{txat} [\emph{txät}]: агент по продаже недвижимости\\
0xA9F1 \textbf{bxat} [\emph{bxät}]: лысый\\
0xA9F2 \textbf{gxat} [\emph{gxät}]: контр-атака\\
0xA9F3 \textbf{dxat} [\emph{dxät}]: опрометчивость\\
0xAA01 \textbf{myut} [\emph{mjut}]: моральный\\
0xAA02 \textbf{kyut} [\emph{kjut}]: компьютер\\
0xAA04 \textbf{pyut} [\emph{pjut}]: север\\
0xAA06 \textbf{nyut} [\emph{njut}]: ваш\\
0xAA08 \textbf{hmut} [\emph{ʰmut}]: потерял\\
0xAA09 \textbf{syut} [\emph{sjut}]: свобода\\
0xAA0A \textbf{tyut} [\emph{tjut}]: два\\
0xAA0B \textbf{lyut} [\emph{ljut}]: бороться\\
0xAA0C \textbf{fyut} [\emph{fjut}]: привередливый\\
0xAA0D \textbf{cyut} [\emph{ʃjut}]: тень\\
0xAA0E \textbf{ryut} [\emph{rjut}]: ревнивый\\
0xAA10 \textbf{hkut} [\emph{ʰkut}]: знать\\
0xAA11 \textbf{byut} [\emph{bjut}]: сука\\
0xAA12 \textbf{gyut} [\emph{gjut}]: оголять\\
0xAA13 \textbf{dyut} [\emph{djut}]: скидка\\
0xAA14 \textbf{zyut} [\emph{zjut}]: излишество\\
0xAA15 \textbf{jyut} [\emph{ʒjut}]: все цели\\
0xAA16 \textbf{vyut} [\emph{vjut}]: университеты\\
0xAA18 \textbf{hyut} [\emph{ʰjut}]: вопросительный настроение\\
0xAA19 \textbf{qyut} [\emph{ŋjut}]: фунгицид\\
0xAA1A \textbf{xyut} [\emph{xjut}]: в ближайщем будущем\\
0xAA20 \textbf{hput} [\emph{ʰput}]: Помогите\\
0xAA21 \textbf{mwut} [\emph{mwut}]: сова\\
0xAA22 \textbf{kwut} [\emph{kwut}]: Пальто\\
0xAA24 \textbf{pwut} [\emph{pwut}]: предопределенный\\
0xAA26 \textbf{nwut} [\emph{nwut}]: на цыпочках\\
0xAA28 \textbf{hwut} [\emph{ʰwut}]: забывать\\
0xAA29 \textbf{swut} [\emph{swut}]: святой\\
0xAA2A \textbf{twut} [\emph{twut}]: dubitative настроение\\
0xAA2B \textbf{lwut} [\emph{lwut}]: поток тока\\
0xAA2C \textbf{fwut} [\emph{fwut}]: вольт\\
0xAA2D \textbf{cwut} [\emph{ʃwut}]: перемирие\\
0xAA2E \textbf{rwut} [\emph{rwut}]: структурный\\
0xAA30 \textbf{hnut} [\emph{ʰnut}]: контрфактическое условный\\
0xAA31 \textbf{bwut} [\emph{bwut}]: береза\\
0xAA32 \textbf{gwut} [\emph{gwut}]: Джибути\\
0xAA33 \textbf{dwut} [\emph{dwut}]: документальный\\
0xAA36 \textbf{vwut} [\emph{vwut}]: Кувейт\\
0xAA39 \textbf{qwut} [\emph{ŋwut}]: площади сфальсифицированы\\
0xAA3A \textbf{xwut} [\emph{xwut}]: слабеть\\
0xAA42 \textbf{ksut} [\emph{ksut}]: убийство\\
0xAA44 \textbf{psut} [\emph{psut}]: объекты\\
0xAA48 \textbf{hsut} [\emph{ʰsut}]: все\\
0xAA4A \textbf{tsut} [\emph{tsut}]: врач\\
0xAA50 \textbf{htut} [\emph{ʰtut}]: негативный квантификатором\\
0xAA51 \textbf{bzut} [\emph{bzut}]: метель\\
0xAA52 \textbf{gzut} [\emph{gzut}]: взвести каракуля ду\\
0xAA53 \textbf{dzut} [\emph{dzut}]: растянутый\\
0xAA58 \textbf{hlut} [\emph{ʰlut}]: автор\\
0xAA60 \textbf{hfut} [\emph{ʰfut}]: падать\\
0xAA62 \textbf{klut} [\emph{klut}]: спекулировать\\
0xAA64 \textbf{plut} [\emph{plut}]: картошка\\
0xAA68 \textbf{hcut} [\emph{ʰʃut}]: условный настроение\\
0xAA69 \textbf{slut} [\emph{slut}]: салат\\
0xAA6A \textbf{tlut} [\emph{tlut}]: степень\\
0xAA6C \textbf{flut} [\emph{flut}]: поплавок\\
0xAA6D \textbf{clut} [\emph{ʃlut}]: молоток\\
0xAA70 \textbf{hrut} [\emph{ʰrut}]: аргумент\\
0xAA71 \textbf{blut} [\emph{blut}]: сдвинуть с места\\
0xAA72 \textbf{glut} [\emph{glut}]: калькулятор\\
0xAA73 \textbf{dlut} [\emph{dlut}]: обозначать\\
0xAA74 \textbf{zlut} [\emph{zlut}]: модулировать\\
0xAA75 \textbf{jlut} [\emph{ʒlut}]: суглинок\\
0xAA76 \textbf{vlut} [\emph{vlut}]: сочлененный\\
0xAA7A \textbf{xlut} [\emph{xlut}]: домашнее растение\\
0xAA82 \textbf{kfut} [\emph{kfut}]: комендантский час\\
0xAA84 \textbf{pfut} [\emph{pfut}]: грудь\\
0xAA88 \textbf{hbut} [\emph{ʰbut}]: обувь\\
0xAA8A \textbf{tfut} [\emph{tfut}]: поздравлять\\
0xAA90 \textbf{hgut} [\emph{ʰgut}]: подсчет\\
0xAA91 \textbf{bvut} [\emph{bvut}]: свекла\\
0xAA92 \textbf{gvut} [\emph{gvut}]: глютен\\
0xAA93 \textbf{dvut} [\emph{dvut}]: выкорчевывать\\
0xAA98 \textbf{hdut} [\emph{ʰdut}]: молоко\\
0xAAA0 \textbf{hzut} [\emph{ʰzut}]: туман\\
0xAAA2 \textbf{kcut} [\emph{kʃut}]: выберите\\
0xAAA4 \textbf{pcut} [\emph{pʃut}]: занято\\
0xAAA8 \textbf{hjut} [\emph{ʰʒut}]: облачный\\
0xAAAA \textbf{tcut} [\emph{tʃut}]: религия\\
0xAAB0 \textbf{hvut} [\emph{ʰvut}]: переполнять\\
0xAAB1 \textbf{bjut} [\emph{bʒut}]: приземистый\\
0xAAB2 \textbf{gjut} [\emph{gʒut}]: напыщенный\\
0xAAB3 \textbf{djut} [\emph{dʒut}]: догадка\\
0xAAC1 \textbf{mrut} [\emph{mrut}]: валовой\\
0xAAC2 \textbf{krut} [\emph{krut}]: укусить\\
0xAAC4 \textbf{prut} [\emph{prut}]: жюри\\
0xAAC6 \textbf{nrut} [\emph{nrut}]: развлекать\\
0xAAC8 \textbf{hqut} [\emph{ʰŋut}]: журнал вне\\
0xAAC9 \textbf{srut} [\emph{srut}]: Абстрактные\\
0xAACA \textbf{trut} [\emph{trut}]: статуя\\
0xAACC \textbf{frut} [\emph{frut}]: самостоятельно удовлетворяет\\
0xAACD \textbf{crut} [\emph{ʃrut}]: пятница\\
0xAAD0 \textbf{hxut} [\emph{ʰxut}]: маршрутизатор\\
0xAAD1 \textbf{brut} [\emph{brut}]: виртуальный\\
0xAAD2 \textbf{grut} [\emph{grut}]: инвертировать\\
0xAAD3 \textbf{drut} [\emph{drut}]: описательный\\
0xAAD4 \textbf{zrut} [\emph{zrut}]: изумруд\\
0xAAD5 \textbf{jrut} [\emph{ʒrut}]: жемчуг\\
0xAAD6 \textbf{vrut} [\emph{vrut}]: беспозвоночный\\
0xAAD9 \textbf{qrut} [\emph{ŋrut}]: попугай\\
0xAADA \textbf{xrut} [\emph{xrut}]: детская площадка\\
0xAAE2 \textbf{kxut} [\emph{kxut}]: широта\\
0xAAE4 \textbf{pxut} [\emph{pxut}]: пернатых\\
0xAAEA \textbf{txut} [\emph{txut}]: значительность\\
0xAAF1 \textbf{bxut} [\emph{bxut}]: тупой\\
0xAAF2 \textbf{gxut} [\emph{gxut}]: Уганда\\
0xAAF3 \textbf{dxut} [\emph{dxut}]: суфлер\\
0xAB01 \textbf{myet} [\emph{mjet}]: мелодия\\
0xAB02 \textbf{kyet} [\emph{kjet}]: коллектор\\
0xAB04 \textbf{pyet} [\emph{pjet}]: принтер\\
0xAB06 \textbf{nyet} [\emph{njet}]: подушка\\
0xAB08 \textbf{hmet} [\emph{ʰmet}]: Эл. адрес\\
0xAB09 \textbf{syet} [\emph{sjet}]: эскиз\\
0xAB0A \textbf{tyet} [\emph{tjet}]: Спортивное\\
0xAB0B \textbf{lyet} [\emph{ljet}]: объектив\\
0xAB0C \textbf{fyet} [\emph{fjet}]: Фаренгейт\\
0xAB0D \textbf{cyet} [\emph{ʃjet}]: пожалел\\
0xAB0E \textbf{ryet} [\emph{rjet}]: выпускник\\
0xAB10 \textbf{hket} [\emph{ʰket}]: проснулся\\
0xAB11 \textbf{byet} [\emph{bjet}]: безработные\\
0xAB12 \textbf{gyet} [\emph{gjet}]: нефрит\\
0xAB13 \textbf{dyet} [\emph{djet}]: Мировой\\
0xAB14 \textbf{zyet} [\emph{zjet}]: Зеведей\\
0xAB15 \textbf{jyet} [\emph{ʒjet}]: душеприказчица\\
0xAB16 \textbf{vyet} [\emph{vjet}]: испорченный\\
0xAB18 \textbf{hyet} [\emph{ʰjet}]: масло\\
0xAB19 \textbf{qyet} [\emph{ŋjet}]: подвергать опасности\\
0xAB1A \textbf{xyet} [\emph{xjet}]: дегенерат\\
0xAB20 \textbf{hpet} [\emph{ʰpet}]: 15\\
0xAB21 \textbf{mwet} [\emph{mwet}]: кладбище\\
0xAB22 \textbf{kwet} [\emph{kwet}]: кататься на коньках\\
0xAB24 \textbf{pwet} [\emph{pwet}]: головоломка\\
0xAB26 \textbf{nwet} [\emph{nwet}]: вязаный\\
0xAB28 \textbf{hwet} [\emph{ʰwet}]: желудок\\
0xAB29 \textbf{swet} [\emph{swet}]: убеждать\\
0xAB2A \textbf{twet} [\emph{twet}]: юбка\\
0xAB2B \textbf{lwet} [\emph{lwet}]: вертолет\\
0xAB2C \textbf{fwet} [\emph{fwet}]: бесплатное программное обеспечение\\
0xAB2D \textbf{cwet} [\emph{ʃwet}]: привязь\\
0xAB2E \textbf{rwet} [\emph{rwet}]: реалистичный\\
0xAB30 \textbf{hnet} [\emph{ʰnet}]: замечание настроение\\
0xAB31 \textbf{bwet} [\emph{bwet}]: кишечник\\
0xAB32 \textbf{gwet} [\emph{gwet}]: скоро будущее напряженный\\
0xAB33 \textbf{dwet} [\emph{dwet}]: со временем\\
0xAB34 \textbf{zwet} [\emph{zwet}]: второй полусредний вес\\
0xAB35 \textbf{jwet} [\emph{ʒwet}]: лишать наследства\\
0xAB36 \textbf{vwet} [\emph{vwet}]: мировой класс\\
0xAB39 \textbf{qwet} [\emph{ŋwet}]: конский хвост\\
0xAB3A \textbf{xwet} [\emph{xwet}]: экватор\\
0xAB42 \textbf{kset} [\emph{kset}]: оркестр\\
0xAB44 \textbf{pset} [\emph{pset}]: продлить\\
0xAB48 \textbf{hset} [\emph{ʰset}]: центр\\
0xAB4A \textbf{tset} [\emph{tset}]: жажда мести\\
0xAB50 \textbf{htet} [\emph{ʰtet}]: ангел\\
0xAB51 \textbf{bzet} [\emph{bzet}]: баклажан\\
0xAB52 \textbf{gzet} [\emph{gzet}]: опоясанный\\
0xAB53 \textbf{dzet} [\emph{dzet}]: полупроводник\\
0xAB58 \textbf{hlet} [\emph{ʰlet}]: лестница\\
0xAB60 \textbf{hfet} [\emph{ʰfet}]: противодействовать\\
0xAB62 \textbf{klet} [\emph{klet}]: относительно\\
0xAB64 \textbf{plet} [\emph{plet}]: издатель\\
0xAB68 \textbf{hcet} [\emph{ʰʃet}]: ненавидеть\\
0xAB69 \textbf{slet} [\emph{slet}]: живая изгородь\\
0xAB6A \textbf{tlet} [\emph{tlet}]: hesternal напряженный\\
0xAB6C \textbf{flet} [\emph{flet}]: изменяемую\\
0xAB6D \textbf{clet} [\emph{ʃlet}]: неполный\\
0xAB70 \textbf{hret} [\emph{ʰret}]: твердо\\
0xAB71 \textbf{blet} [\emph{blet}]: таблетка\\
0xAB72 \textbf{glet} [\emph{glet}]: напряжение\\
0xAB73 \textbf{dlet} [\emph{dlet}]: каталог\\
0xAB74 \textbf{zlet} [\emph{zlet}]: омлет\\
0xAB75 \textbf{jlet} [\emph{ʒlet}]: календарь\\
0xAB76 \textbf{vlet} [\emph{vlet}]: сварной шов\\
0xAB7A \textbf{xlet} [\emph{xlet}]: превосходный степень\\
0xAB82 \textbf{kfet} [\emph{kfet}]: катод\\
0xAB84 \textbf{pfet} [\emph{pfet}]: листовка\\
0xAB88 \textbf{hbet} [\emph{ʰbet}]: делать ставку\\
0xAB8A \textbf{tfet} [\emph{tfet}]: производитель\\
0xAB90 \textbf{hget} [\emph{ʰget}]: гость\\
0xAB91 \textbf{bvet} [\emph{bvet}]: эфир\\
0xAB92 \textbf{gvet} [\emph{gvet}]: планер\\
0xAB93 \textbf{dvet} [\emph{dvet}]: половик\\
0xAB98 \textbf{hdet} [\emph{ʰdet}]: дефектный\\
0xABA0 \textbf{hzet} [\emph{ʰzet}]: морской узел\\
0xABA2 \textbf{kcet} [\emph{kʃet}]: аренда\\
0xABA4 \textbf{pcet} [\emph{pʃet}]: неопределенность\\
0xABA8 \textbf{hjet} [\emph{ʰʒet}]: бархат\\
0xABAA \textbf{tcet} [\emph{tʃet}]: угрожающий\\
0xABB0 \textbf{hvet} [\emph{ʰvet}]: рядом с\\
0xABB1 \textbf{bjet} [\emph{bʒet}]: бюджет\\
0xABB2 \textbf{gjet} [\emph{gʒet}]: электрод\\
0xABB3 \textbf{djet} [\emph{dʒet}]: косметический ремонт\\
0xABC1 \textbf{mret} [\emph{mret}]: министерство\\
0xABC2 \textbf{kret} [\emph{kret}]: экстракт\\
0xABC4 \textbf{pret} [\emph{pret}]: недостаточный\\
0xABC6 \textbf{nret} [\emph{nret}]: яростно\\
0xABC8 \textbf{hqet} [\emph{ʰŋet}]: розничная торговля магазин\\
0xABC9 \textbf{sret} [\emph{sret}]: семнадцать\\
0xABCA \textbf{tret} [\emph{tret}]: театр\\
0xABCC \textbf{fret} [\emph{fret}]: искатель\\
0xABCD \textbf{cret} [\emph{ʃret}]: следы\\
0xABD0 \textbf{hxet} [\emph{ʰxet}]: увеличительный\\
0xABD1 \textbf{bret} [\emph{bret}]: байт\\
0xABD2 \textbf{gret} [\emph{gret}]: вычитать\\
0xABD3 \textbf{dret} [\emph{dret}]: итерация\\
0xABD4 \textbf{zret} [\emph{zret}]: подслушивать\\
0xABD5 \textbf{jret} [\emph{ʒret}]: Острота\\
0xABD6 \textbf{vret} [\emph{vret}]: унаследованный\\
0xABD9 \textbf{qret} [\emph{ŋret}]: кощунство\\
0xABDA \textbf{xret} [\emph{xret}]: терабайт\\
0xABE2 \textbf{kxet} [\emph{kxet}]: эскалатор\\
0xABE4 \textbf{pxet} [\emph{pxet}]: плащ\\
0xABEA \textbf{txet} [\emph{txet}]: в чате\\
0xABF1 \textbf{bxet} [\emph{bxet}]: барометр\\
0xABF2 \textbf{gxet} [\emph{gxet}]: геометрия\\
0xABF3 \textbf{dxet} [\emph{dxet}]: другой мероприятие\\
0xAC01 \textbf{myot} [\emph{mjot}]: печально известный\\
0xAC02 \textbf{kyot} [\emph{kjot}]: аккорд\\
0xAC04 \textbf{pyot} [\emph{pjot}]: голубь\\
0xAC06 \textbf{nyot} [\emph{njot}]: слушатель\\
0xAC08 \textbf{hmot} [\emph{ʰmot}]: мясо\\
0xAC09 \textbf{syot} [\emph{sjot}]: сложить\\
0xAC0A \textbf{tyot} [\emph{tjot}]: договор\\
0xAC0B \textbf{lyot} [\emph{ljot}]: авто\\
0xAC0C \textbf{fyot} [\emph{fjot}]: сноска\\
0xAC0D \textbf{cyot} [\emph{ʃjot}]: Импортировать\\
0xAC0E \textbf{ryot} [\emph{rjot}]: благоговение\\
0xAC10 \textbf{hkot} [\emph{ʰkot}]: вместимость\\
0xAC11 \textbf{byot} [\emph{bjot}]: видео\\
0xAC12 \textbf{gyot} [\emph{gjot}]: условный\\
0xAC13 \textbf{dyot} [\emph{djot}]: радио\\
0xAC14 \textbf{zyot} [\emph{zjot}]: закупоренный\\
0xAC15 \textbf{jyot} [\emph{ʒjot}]: конденсатор\\
0xAC16 \textbf{vyot} [\emph{vjot}]: шплинт\\
0xAC18 \textbf{hyot} [\emph{ʰjot}]: опасающийся настроение\\
0xAC19 \textbf{qyot} [\emph{ŋjot}]: солнцестояние\\
0xAC1A \textbf{xyot} [\emph{xjot}]: рептилия\\
0xAC20 \textbf{hpot} [\emph{ʰpot}]: доступный\\
0xAC21 \textbf{mwot} [\emph{mwot}]: помидор\\
0xAC22 \textbf{kwot} [\emph{kwot}]: клавиатура\\
0xAC24 \textbf{pwot} [\emph{pwot}]: белок\\
0xAC26 \textbf{nwot} [\emph{nwot}]: непрерывность\\
0xAC28 \textbf{hwot} [\emph{ʰwot}]: витой\\
0xAC29 \textbf{swot} [\emph{swot}]: болото\\
0xAC2A \textbf{twot} [\emph{twot}]: Мероприятия\\
0xAC2B \textbf{lwot} [\emph{lwot}]: хлористый\\
0xAC2C \textbf{fwot} [\emph{fwot}]: увядать\\
0xAC2D \textbf{cwot} [\emph{ʃwot}]: неволя\\
0xAC2E \textbf{rwot} [\emph{rwot}]: орбита\\
0xAC30 \textbf{hnot} [\emph{ʰnot}]: Гостиница\\
0xAC31 \textbf{bwot} [\emph{bwot}]: бычий\\
0xAC32 \textbf{gwot} [\emph{gwot}]: агностик\\
0xAC33 \textbf{dwot} [\emph{dwot}]: диктатор\\
0xAC35 \textbf{jwot} [\emph{ʒwot}]: красное пальто\\
0xAC36 \textbf{vwot} [\emph{vwot}]: вязкость\\
0xAC39 \textbf{qwot} [\emph{ŋwot}]: андорра\\
0xAC3A \textbf{xwot} [\emph{xwot}]: так называемые\\
0xAC42 \textbf{ksot} [\emph{ksot}]: кость\\
0xAC44 \textbf{psot} [\emph{psot}]: расплата\\
0xAC48 \textbf{hsot} [\emph{ʰsot}]: скорость\\
0xAC4A \textbf{tsot} [\emph{tsot}]: верность\\
0xAC50 \textbf{htot} [\emph{ʰtot}]: лоб\\
0xAC51 \textbf{bzot} [\emph{bzot}]: баллонах\\
0xAC52 \textbf{gzot} [\emph{gzot}]: транзистор\\
0xAC53 \textbf{dzot} [\emph{dzot}]: донор\\
0xAC58 \textbf{hlot} [\emph{ʰlot}]: планируемый\\
0xAC60 \textbf{hfot} [\emph{ʰfot}]: состоять\\
0xAC62 \textbf{klot} [\emph{klot}]: тренер\\
0xAC64 \textbf{plot} [\emph{plot}]: пилот\\
0xAC68 \textbf{hcot} [\emph{ʰʃot}]: наслаждаться\\
0xAC69 \textbf{slot} [\emph{slot}]: стрелять\\
0xAC6A \textbf{tlot} [\emph{tlot}]: детство\\
0xAC6C \textbf{flot} [\emph{flot}]: баловать\\
0xAC6D \textbf{clot} [\emph{ʃlot}]: Новостная рассылка\\
0xAC70 \textbf{hrot} [\emph{ʰrot}]: недавний мимо напряженный\\
0xAC71 \textbf{blot} [\emph{blot}]: болт\\
0xAC72 \textbf{glot} [\emph{glot}]: мобильность\\
0xAC73 \textbf{dlot} [\emph{dlot}]: скачать\\
0xAC74 \textbf{zlot} [\emph{zlot}]: анод\\
0xAC75 \textbf{jlot} [\emph{ʒlot}]: растворитель\\
0xAC76 \textbf{vlot} [\emph{vlot}]: заметный\\
0xAC7A \textbf{xlot} [\emph{xlot}]: закат солнца\\
0xAC82 \textbf{kfot} [\emph{kfot}]: выдача\\
0xAC84 \textbf{pfot} [\emph{pfot}]: горшок\\
0xAC88 \textbf{hbot} [\emph{ʰbot}]: мяч\\
0xAC8A \textbf{tfot} [\emph{tfot}]: сертификат\\
0xAC90 \textbf{hgot} [\emph{ʰgot}]: замороженные\\
0xAC91 \textbf{bvot} [\emph{bvot}]: непрошеный\\
0xAC92 \textbf{gvot} [\emph{gvot}]: капли\\
0xAC93 \textbf{dvot} [\emph{dvot}]: проводимость\\
0xAC98 \textbf{hdot} [\emph{ʰdot}]: дистальный демонстративный\\
0xACA0 \textbf{hzot} [\emph{ʰzot}]: прекращение\\
0xACA2 \textbf{kcot} [\emph{kʃot}]: собирать\\
0xACA4 \textbf{pcot} [\emph{pʃot}]: домашнее животное\\
0xACA8 \textbf{hjot} [\emph{ʰʒot}]: студия\\
0xACAA \textbf{tcot} [\emph{tʃot}]: твердое вещество\\
0xACB0 \textbf{hvot} [\emph{ʰvot}]: блок-схема\\
0xACB1 \textbf{bjot} [\emph{bʒot}]: изымать\\
0xACB2 \textbf{gjot} [\emph{gʒot}]: оранжад\\
0xACB3 \textbf{djot} [\emph{dʒot}]: осадок\\
0xACC1 \textbf{mrot} [\emph{mrot}]: двигатель\\
0xACC2 \textbf{krot} [\emph{krot}]: резка\\
0xACC4 \textbf{prot} [\emph{prot}]: портрет\\
0xACC6 \textbf{nrot} [\emph{nrot}]: Международный\\
0xACC8 \textbf{hqot} [\emph{ʰŋot}]: растворенное вещество\\
0xACC9 \textbf{srot} [\emph{srot}]: эскорт\\
0xACCA \textbf{trot} [\emph{trot}]: температура\\
0xACCC \textbf{frot} [\emph{frot}]: инфраструктура\\
0xACCD \textbf{crot} [\emph{ʃrot}]: луг\\
0xACD0 \textbf{hxot} [\emph{ʰxot}]: наручные часы\\
0xACD1 \textbf{brot} [\emph{brot}]: бутон\\
0xACD2 \textbf{grot} [\emph{grot}]: партнерство\\
0xACD3 \textbf{drot} [\emph{drot}]: контролируемый\\
0xACD4 \textbf{zrot} [\emph{zrot}]: кувыркаться\\
0xACD5 \textbf{jrot} [\emph{ʒrot}]: жареные\\
0xACD6 \textbf{vrot} [\emph{vrot}]: выпад\\
0xACD9 \textbf{qrot} [\emph{ŋrot}]: нейтрон\\
0xACDA \textbf{xrot} [\emph{xrot}]: панталоны\\
0xACE2 \textbf{kxot} [\emph{kxot}]: криптон\\
0xACE4 \textbf{pxot} [\emph{pxot}]: из первых рук\\
0xACEA \textbf{txot} [\emph{txot}]: в середине осени\\
0xACF1 \textbf{bxot} [\emph{bxot}]: само контроль\\
0xACF2 \textbf{gxot} [\emph{gxot}]: грецкий орех\\
0xACF3 \textbf{dxot} [\emph{dxot}]: йогурт\\
0xAD01 \textbf{my6t} [\emph{mjət}]: рамадан\\
0xAD02 \textbf{ky6t} [\emph{kjət}]: аксиоматический\\
0xAD04 \textbf{py6t} [\emph{pjət}]: Эритрея\\
0xAD06 \textbf{ny6t} [\emph{njət}]: ингибитор\\
0xAD08 \textbf{hm6t} [\emph{ʰmət}]: первичная объект\\
0xAD09 \textbf{sy6t} [\emph{sjət}]: Свазиленд\\
0xAD0A \textbf{ty6t} [\emph{tjət}]: Чад\\
0xAD0B \textbf{ly6t} [\emph{ljət}]: назидательная речь соглашение\\
0xAD0C \textbf{fy6t} [\emph{fjət}]: сфинктер\\
0xAD0D \textbf{cy6t} [\emph{ʃjət}]: жучок\\
0xAD0E \textbf{ry6t} [\emph{rjət}]: орхидея\\
0xAD10 \textbf{hk6t} [\emph{ʰkət}]: октет\\
0xAD11 \textbf{by6t} [\emph{bjət}]: язь\\
0xAD12 \textbf{gy6t} [\emph{gjət}]: Гренада\\
0xAD13 \textbf{dy6t} [\emph{djət}]: диод\\
0xAD18 \textbf{hy6t} [\emph{ʰjət}]: грамматическая - объект\\
0xAD1A \textbf{xy6t} [\emph{xjət}]: пестрый\\
0xAD20 \textbf{hp6t} [\emph{ʰpət}]: изменение из мероприятие\\
0xAD21 \textbf{mw6t} [\emph{mwət}]: от двери до двери\\
0xAD22 \textbf{kw6t} [\emph{kwət}]: скваттер\\
0xAD24 \textbf{pw6t} [\emph{pwət}]: пульсировать\\
0xAD26 \textbf{nw6t} [\emph{nwət}]: монолит\\
0xAD28 \textbf{hw6t} [\emph{ʰwət}]: Швеция\\
0xAD29 \textbf{sw6t} [\emph{swət}]: саудовская\\
0xAD2A \textbf{tw6t} [\emph{twət}]: затемнение\\
0xAD2B \textbf{lw6t} [\emph{lwət}]: Лесото\\
0xAD2D \textbf{cw6t} [\emph{ʃwət}]: кристаллы\\
0xAD2E \textbf{rw6t} [\emph{rwət}]: термостат\\
0xAD30 \textbf{hn6t} [\emph{ʰnət}]: соединители\\
0xAD3A \textbf{xw6t} [\emph{xwət}]: авантитул\\
0xAD42 \textbf{ks6t} [\emph{ksət}]: причинный\\
0xAD44 \textbf{ps6t} [\emph{psət}]: Пасха\\
0xAD48 \textbf{hs6t} [\emph{ʰsət}]: комната отдыха\\
0xAD4A \textbf{ts6t} [\emph{tsət}]: коническая\\
0xAD50 \textbf{ht6t} [\emph{ʰtət}]: реальный время\\
0xAD58 \textbf{hl6t} [\emph{ʰlət}]: литр\\
0xAD60 \textbf{hf6t} [\emph{ʰfət}]: фазан\\
0xAD62 \textbf{kl6t} [\emph{klət}]: подставка\\
0xAD64 \textbf{pl6t} [\emph{plət}]: сыпь\\
0xAD68 \textbf{hc6t} [\emph{ʰʃət}]: светящийся интенсивность\\
0xAD69 \textbf{sl6t} [\emph{slət}]: crastinal напряженный\\
0xAD6A \textbf{tl6t} [\emph{tlət}]: Атлантика\\
0xAD6C \textbf{fl6t} [\emph{flət}]: евклидов\\
0xAD6D \textbf{cl6t} [\emph{ʃlət}]: изолятор\\
0xAD70 \textbf{hr6t} [\emph{ʰrət}]: артефакт\\
0xAD71 \textbf{bl6t} [\emph{blət}]: балласт\\
0xAD72 \textbf{gl6t} [\emph{glət}]: палеонтолог\\
0xAD73 \textbf{dl6t} [\emph{dlət}]: должным образом\\
0xAD74 \textbf{zl6t} [\emph{zlət}]: фундук\\
0xAD75 \textbf{jl6t} [\emph{ʒlət}]: волочиться\\
0xAD76 \textbf{vl6t} [\emph{vlət}]: миллилитр\\
0xAD7A \textbf{xl6t} [\emph{xlət}]: Таиланд\\
0xAD82 \textbf{kf6t} [\emph{kfət}]: отслаиваться\\
0xAD84 \textbf{pf6t} [\emph{pfət}]: суппозиторий\\
0xAD88 \textbf{hb6t} [\emph{ʰbət}]: бренди\\
0xAD8A \textbf{tf6t} [\emph{tfət}]: колготки\\
0xAD90 \textbf{hg6t} [\emph{ʰgət}]: педаль газа\\
0xAD98 \textbf{hd6t} [\emph{ʰdət}]: середине зимы\\
0xADA0 \textbf{hz6t} [\emph{ʰzət}]: зигота\\
0xADA2 \textbf{kc6t} [\emph{kʃət}]: клитор\\
0xADA4 \textbf{pc6t} [\emph{pʃət}]: спринтер\\
0xADA8 \textbf{hj6t} [\emph{ʰʒət}]: гладильный\\
0xADAA \textbf{tc6t} [\emph{tʃət}]: непосредственный объект\\
0xADB0 \textbf{hv6t} [\emph{ʰvət}]: вырубка леса\\
0xADB1 \textbf{bj6t} [\emph{bʒət}]: миндаль\\
0xADB3 \textbf{dj6t} [\emph{dʒət}]: сиккатив\\
0xADC1 \textbf{mr6t} [\emph{mrət}]: синонимичный\\
0xADC2 \textbf{kr6t} [\emph{krət}]: электрический текущий\\
0xADC4 \textbf{pr6t} [\emph{prət}]: полярность\\
0xADC6 \textbf{nr6t} [\emph{nrət}]: индийский\\
0xADC8 \textbf{hq6t} [\emph{ʰŋət}]: Аристотель\\
0xADC9 \textbf{sr6t} [\emph{srət}]: сверхкритической жидкости\\
0xADCA \textbf{tr6t} [\emph{trət}]: внеземной\\
0xADCC \textbf{fr6t} [\emph{frət}]: одна треть\\
0xADCD \textbf{cr6t} [\emph{ʃrət}]: хиропрактики\\
0xADD0 \textbf{hx6t} [\emph{ʰxət}]: Канада\\
0xADD1 \textbf{br6t} [\emph{brət}]: Бурунди\\
0xADD2 \textbf{gr6t} [\emph{grət}]: контрацептив\\
0xADD3 \textbf{dr6t} [\emph{drət}]: дистрактор\\
0xADD6 \textbf{vr6t} [\emph{vrət}]: каратэ\\
0xADD9 \textbf{qr6t} [\emph{ŋrət}]: транспортир\\
0xADDA \textbf{xr6t} [\emph{xrət}]: хаотичный\\
0xADE2 \textbf{kx6t} [\emph{kxət}]: горючий\\
0xADE4 \textbf{px6t} [\emph{pxət}]: лампа накаливания\\
0xADEA \textbf{tx6t} [\emph{txət}]: текст редактор\\
0xADF1 \textbf{bx6t} [\emph{bxət}]: смеситель\\
0xADF3 \textbf{dx6t} [\emph{dxət}]: не согласен\\
\subsection{ B }
0xB001 \textbf{myi7t} [\emph{mji˥t}]: лунатик\\
0xB002 \textbf{kyi7t} [\emph{kji˥t}]: окольный\\
0xB004 \textbf{pyi7t} [\emph{pji˥t}]: пол пинты\\
0xB006 \textbf{nyi7t} [\emph{nji˥t}]: тридцать восемь\\
0xB008 \textbf{hmi7t} [\emph{ʰmi˥t}]: бессмертный\\
0xB009 \textbf{syi7t} [\emph{sji˥t}]: фыркнул\\
0xB00A \textbf{tyi7t} [\emph{tji˥t}]: пунктирный\\
0xB00B \textbf{lyi7t} [\emph{lji˥t}]: льняное семя\\
0xB00C \textbf{fyi7t} [\emph{fji˥t}]: фенил\\
0xB00D \textbf{cyi7t} [\emph{ʃji˥t}]: палочки для еды\\
0xB00E \textbf{ryi7t} [\emph{rji˥t}]: беспричинный\\
0xB010 \textbf{hki7t} [\emph{ʰki˥t}]: сгущенный\\
0xB011 \textbf{byi7t} [\emph{bji˥t}]: аболиционисты\\
0xB012 \textbf{gyi7t} [\emph{gji˥t}]: идти между\\
0xB013 \textbf{dyi7t} [\emph{dji˥t}]: индуктивный\\
0xB015 \textbf{jyi7t} [\emph{ʒji˥t}]: Последняя минута\\
0xB018 \textbf{hyi7t} [\emph{ʰji˥t}]: иезуитский\\
0xB019 \textbf{qyi7t} [\emph{ŋji˥t}]: эмпирик\\
0xB01A \textbf{xyi7t} [\emph{xji˥t}]: высокотональный\\
0xB020 \textbf{hpi7t} [\emph{ʰpi˥t}]: Питер\\
0xB021 \textbf{mwi7t} [\emph{mwi˥t}]: пятьдесят два\\
0xB022 \textbf{kwi7t} [\emph{kwi˥t}]: конькобежец\\
0xB024 \textbf{pwi7t} [\emph{pwi˥t}]: прагматика\\
0xB026 \textbf{nwi7t} [\emph{nwi˥t}]: не введен в эксплуатацию\\
0xB028 \textbf{hwi7t} [\emph{ʰwi˥t}]: жилет\\
0xB029 \textbf{swi7t} [\emph{swi˥t}]: пятьдесят семь\\
0xB02A \textbf{twi7t} [\emph{twi˥t}]: трактирщик\\
0xB02B \textbf{lwi7t} [\emph{lwi˥t}]: податливый\\
0xB02C \textbf{fwi7t} [\emph{fwi˥t}]: фаг\\
0xB02E \textbf{rwi7t} [\emph{rwi˥t}]: рассказчик\\
0xB030 \textbf{hni7t} [\emph{ʰni˥t}]: неотобранный\\
0xB031 \textbf{bwi7t} [\emph{bwi˥t}]: бисексуальность\\
0xB033 \textbf{dwi7t} [\emph{dwi˥t}]: раздвоенный хвост\\
0xB03A \textbf{xwi7t} [\emph{xwi˥t}]: сбивчиво\\
0xB042 \textbf{ksi7t} [\emph{ksi˥t}]: высосанный\\
0xB044 \textbf{psi7t} [\emph{psi˥t}]: самоочевидной\\
0xB048 \textbf{hsi7t} [\emph{ʰsi˥t}]: обезжиренное\\
0xB04A \textbf{tsi7t} [\emph{tsi˥t}]: отдыхает\\
0xB050 \textbf{hti7t} [\emph{ʰti˥t}]: отозвала\\
0xB053 \textbf{dzi7t} [\emph{dzi˥t}]: двугранный\\
0xB058 \textbf{hli7t} [\emph{ʰli˥t}]: Илиада\\
0xB060 \textbf{hfi7t} [\emph{ʰfi˥t}]: фотоэлектрический\\
0xB062 \textbf{kli7t} [\emph{kli˥t}]: кельтская\\
0xB064 \textbf{pli7t} [\emph{pli˥t}]: Исландия\\
0xB068 \textbf{hci7t} [\emph{ʰʃi˥t}]: обещания\\
0xB069 \textbf{sli7t} [\emph{sli˥t}]: тридцать шесть\\
0xB06A \textbf{tli7t} [\emph{tli˥t}]: решительно\\
0xB06C \textbf{fli7t} [\emph{fli˥t}]: теодолит\\
0xB06D \textbf{cli7t} [\emph{ʃli˥t}]: беспризорник\\
0xB070 \textbf{hri7t} [\emph{ʰri˥t}]: Христос\\
0xB071 \textbf{bli7t} [\emph{bli˥t}]: Билли\\
0xB072 \textbf{gli7t} [\emph{gli˥t}]: недвижимое имущество\\
0xB073 \textbf{dli7t} [\emph{dli˥t}]: Аделаида\\
0xB074 \textbf{zli7t} [\emph{zli˥t}]: немного обратный порядок байт\\
0xB075 \textbf{jli7t} [\emph{ʒli˥t}]: мастит\\
0xB076 \textbf{vli7t} [\emph{vli˥t}]: чревовещатель\\
0xB07A \textbf{xli7t} [\emph{xli˥t}]: скрипач\\
0xB082 \textbf{kfi7t} [\emph{kfi˥t}]: аконит\\
0xB084 \textbf{pfi7t} [\emph{pfi˥t}]: несведущий\\
0xB088 \textbf{hbi7t} [\emph{ʰbi˥t}]: биты\\
0xB08A \textbf{tfi7t} [\emph{tfi˥t}]: неподдельный\\
0xB090 \textbf{hgi7t} [\emph{ʰgi˥t}]: несогласованность\\
0xB091 \textbf{bvi7t} [\emph{bvi˥t}]: отказываться\\
0xB093 \textbf{dvi7t} [\emph{dvi˥t}]: дравидский\\
0xB098 \textbf{hdi7t} [\emph{ʰdi˥t}]: полночь\\
0xB0A0 \textbf{hzi7t} [\emph{ʰzi˥t}]: островок\\
0xB0A2 \textbf{kci7t} [\emph{kʃi˥t}]: скупо\\
0xB0A4 \textbf{pci7t} [\emph{pʃi˥t}]: лицемер\\
0xB0A8 \textbf{hji7t} [\emph{ʰʒi˥t}]: пылко\\
0xB0AA \textbf{tci7t} [\emph{tʃi˥t}]: изобретенной\\
0xB0B0 \textbf{hvi7t} [\emph{ʰvi˥t}]: ризница\\
0xB0B2 \textbf{gji7t} [\emph{gʒi˥t}]: Г.Д.\\
0xB0B3 \textbf{dji7t} [\emph{dʒi˥t}]: антидемократические\\
0xB0C1 \textbf{mri7t} [\emph{mri˥t}]: Мэри\\
0xB0C2 \textbf{kri7t} [\emph{kri˥t}]: жертвы\\
0xB0C4 \textbf{pri7t} [\emph{pri˥t}]: святоша\\
0xB0C6 \textbf{nri7t} [\emph{nri˥t}]: рыцарский\\
0xB0C8 \textbf{hqi7t} [\emph{ʰŋi˥t}]: чернильница\\
0xB0C9 \textbf{sri7t} [\emph{sri˥t}]: историки\\
0xB0CA \textbf{tri7t} [\emph{tri˥t}]: вышеупомянутый\\
0xB0CC \textbf{fri7t} [\emph{fri˥t}]: юрист\\
0xB0CD \textbf{cri7t} [\emph{ʃri˥t}]: замкнутый круг\\
0xB0D0 \textbf{hxi7t} [\emph{ʰxi˥t}]: гематит\\
0xB0D1 \textbf{bri7t} [\emph{bri˥t}]: болотная лихорадка\\
0xB0D2 \textbf{gri7t} [\emph{gri˥t}]: шишковидный\\
0xB0D3 \textbf{dri7t} [\emph{dri˥t}]: песенка\\
0xB0D4 \textbf{zri7t} [\emph{zri˥t}]: самостоятельно осуждающий\\
0xB0D5 \textbf{jri7t} [\emph{ʒri˥t}]: пешеходный мост\\
0xB0D6 \textbf{vri7t} [\emph{vri˥t}]: продажный\\
0xB0D9 \textbf{qri7t} [\emph{ŋri˥t}]: пирит\\
0xB0DA \textbf{xri7t} [\emph{xri˥t}]: пародия\\
0xB0E2 \textbf{kxi7t} [\emph{kxi˥t}]: удержан\\
0xB0E4 \textbf{pxi7t} [\emph{pxi˥t}]: зловонный\\
0xB0EA \textbf{txi7t} [\emph{txi˥t}]: замещенный\\
0xB0F1 \textbf{bxi7t} [\emph{bxi˥t}]: бенди\\
0xB0F2 \textbf{gxi7t} [\emph{gxi˥t}]: двумужница\\
0xB0F3 \textbf{dxi7t} [\emph{dxi˥t}]: раскольник\\
0xB101 \textbf{mya7t} [\emph{mjä˥t}]: амфитеатр\\
0xB102 \textbf{kya7t} [\emph{kjä˥t}]: обращенных\\
0xB104 \textbf{pya7t} [\emph{pjä˥t}]: капер\\
0xB106 \textbf{nya7t} [\emph{njä˥t}]: сводный брат\\
0xB108 \textbf{hma7t} [\emph{ʰmä˥t}]: Госпожа\\
0xB109 \textbf{sya7t} [\emph{sjä˥t}]: завещание\\
0xB10A \textbf{tya7t} [\emph{tjä˥t}]: левая рука\\
0xB10B \textbf{lya7t} [\emph{ljä˥t}]: приличия\\
0xB10C \textbf{fya7t} [\emph{fjä˥t}]: без даты\\
0xB10D \textbf{cya7t} [\emph{ʃjä˥t}]: скамеечка для ног\\
0xB10E \textbf{rya7t} [\emph{rjä˥t}]: пролетариат\\
0xB110 \textbf{hka7t} [\emph{ʰkä˥t}]: с благодарностью\\
0xB111 \textbf{bya7t} [\emph{bjä˥t}]: Овадия\\
0xB112 \textbf{gya7t} [\emph{gjä˥t}]: трусики\\
0xB113 \textbf{dya7t} [\emph{djä˥t}]: рассветало\\
0xB114 \textbf{zya7t} [\emph{zjä˥t}]: Зевадия\\
0xB115 \textbf{jya7t} [\emph{ʒjä˥t}]: в натуральную величину\\
0xB116 \textbf{vya7t} [\emph{vjä˥t}]: Авестийский\\
0xB118 \textbf{hya7t} [\emph{ʰjä˥t}]: братва\\
0xB119 \textbf{qya7t} [\emph{ŋjä˥t}]: издавна\\
0xB11A \textbf{xya7t} [\emph{xjä˥t}]: язык связан\\
0xB120 \textbf{hpa7t} [\emph{ʰpä˥t}]: пахотный слой почвы\\
0xB121 \textbf{mwa7t} [\emph{mwä˥t}]: любовный\\
0xB122 \textbf{kwa7t} [\emph{kwä˥t}]: один вооруженный\\
0xB124 \textbf{pwa7t} [\emph{pwä˥t}]: камбала\\
0xB126 \textbf{nwa7t} [\emph{nwä˥t}]: неразрезанный\\
0xB128 \textbf{hwa7t} [\emph{ʰwä˥t}]: приветствовал\\
0xB129 \textbf{swa7t} [\emph{swä˥t}]: соната\\
0xB12A \textbf{twa7t} [\emph{twä˥t}]: переплетенный\\
0xB12B \textbf{lwa7t} [\emph{lwä˥t}]: лингвист\\
0xB12C \textbf{fwa7t} [\emph{fwä˥t}]: Вопросы-Ответы\\
0xB12E \textbf{rwa7t} [\emph{rwä˥t}]: увлеченный\\
0xB130 \textbf{hna7t} [\emph{ʰnä˥t}]: цианобактерии\\
0xB131 \textbf{bwa7t} [\emph{bwä˥t}]: глушь\\
0xB132 \textbf{gwa7t} [\emph{gwä˥t}]: Береговая охрана\\
0xB133 \textbf{dwa7t} [\emph{dwä˥t}]: идиоматический\\
0xB13A \textbf{xwa7t} [\emph{xwä˥t}]: кредитор\\
0xB142 \textbf{ksa7t} [\emph{ksä˥t}]: бортовой залп\\
0xB144 \textbf{psa7t} [\emph{psä˥t}]: оседлал\\
0xB148 \textbf{hsa7t} [\emph{ʰsä˥t}]: улицы\\
0xB14A \textbf{tsa7t} [\emph{tsä˥t}]: мужественность\\
0xB150 \textbf{hta7t} [\emph{ʰtä˥t}]: созданный\\
0xB151 \textbf{bza7t} [\emph{bzä˥t}]: обличительный\\
0xB153 \textbf{dza7t} [\emph{dzä˥t}]: диастола\\
0xB158 \textbf{hla7t} [\emph{ʰlä˥t}]: Александр\\
0xB160 \textbf{hfa7t} [\emph{ʰfä˥t}]: расставался\\
0xB162 \textbf{kla7t} [\emph{klä˥t}]: способности\\
0xB164 \textbf{pla7t} [\emph{plä˥t}]: превозносить\\
0xB168 \textbf{hca7t} [\emph{ʰʃä˥t}]: избавлены\\
0xB169 \textbf{sla7t} [\emph{slä˥t}]: защелка\\
0xB16A \textbf{tla7t} [\emph{tlä˥t}]: ударов\\
0xB16C \textbf{fla7t} [\emph{flä˥t}]: воспалительный\\
0xB16D \textbf{cla7t} [\emph{ʃlä˥t}]: йота\\
0xB170 \textbf{hra7t} [\emph{ʰrä˥t}]: реостат\\
0xB171 \textbf{bla7t} [\emph{blä˥t}]: буфетная\\
0xB172 \textbf{gla7t} [\emph{glä˥t}]: глыбы\\
0xB173 \textbf{dla7t} [\emph{dlä˥t}]: под ногами\\
0xB174 \textbf{zla7t} [\emph{zlä˥t}]: развязывать\\
0xB175 \textbf{jla7t} [\emph{ʒlä˥t}]: красящее вещество\\
0xB176 \textbf{vla7t} [\emph{vlä˥t}]: земля\\
0xB17A \textbf{xla7t} [\emph{xlä˥t}]: сильно\\
0xB182 \textbf{kfa7t} [\emph{kfä˥t}]: обрезываете\\
0xB184 \textbf{pfa7t} [\emph{pfä˥t}]: пацифист\\
0xB188 \textbf{hba7t} [\emph{ʰbä˥t}]: обязанный\\
0xB18A \textbf{tfa7t} [\emph{tfä˥t}]: зубчатый\\
0xB190 \textbf{hga7t} [\emph{ʰgä˥t}]: направляющие\\
0xB191 \textbf{bva7t} [\emph{bvä˥t}]: крапать\\
0xB192 \textbf{gva7t} [\emph{gvä˥t}]: трехвалентный\\
0xB198 \textbf{hda7t} [\emph{ʰdä˥t}]: легирующая примесь\\
0xB1A0 \textbf{hza7t} [\emph{ʰzä˥t}]: забита\\
0xB1A2 \textbf{kca7t} [\emph{kʃä˥t}]: неосуществимость\\
0xB1A4 \textbf{pca7t} [\emph{pʃä˥t}]: идолопоклонник\\
0xB1A8 \textbf{hja7t} [\emph{ʰʒä˥t}]: предел слышимости\\
0xB1AA \textbf{tca7t} [\emph{tʃä˥t}]: право по рождению\\
0xB1B0 \textbf{hva7t} [\emph{ʰvä˥t}]: аббат\\
0xB1B1 \textbf{bja7t} [\emph{bʒä˥t}]: субарктический\\
0xB1B2 \textbf{gja7t} [\emph{gʒä˥t}]: анти войны\\
0xB1B3 \textbf{dja7t} [\emph{dʒä˥t}]: смертельный удар\\
0xB1C1 \textbf{mra7t} [\emph{mrä˥t}]: Мадрид\\
0xB1C2 \textbf{kra7t} [\emph{krä˥t}]: возчик\\
0xB1C4 \textbf{pra7t} [\emph{prä˥t}]: сохраняет\\
0xB1C6 \textbf{nra7t} [\emph{nrä˥t}]: Андрей\\
0xB1C8 \textbf{hqa7t} [\emph{ʰŋä˥t}]: тратить\\
0xB1C9 \textbf{sra7t} [\emph{srä˥t}]: Сократ\\
0xB1CA \textbf{tra7t} [\emph{trä˥t}]: тащили\\
0xB1CC \textbf{fra7t} [\emph{frä˥t}]: Евфрат\\
0xB1CD \textbf{cra7t} [\emph{ʃrä˥t}]: запаздывание\\
0xB1D0 \textbf{hxa7t} [\emph{ʰxä˥t}]: ахнула\\
0xB1D1 \textbf{bra7t} [\emph{brä˥t}]: вышитый\\
0xB1D2 \textbf{gra7t} [\emph{grä˥t}]: виноградная лоза\\
0xB1D3 \textbf{dra7t} [\emph{drä˥t}]: обескуражена\\
0xB1D4 \textbf{zra7t} [\emph{zrä˥t}]: ангидрид\\
0xB1D5 \textbf{jra7t} [\emph{ʒrä˥t}]: Пра-пра-бабушка\\
0xB1D6 \textbf{vra7t} [\emph{vrä˥t}]: вибратор\\
0xB1D9 \textbf{qra7t} [\emph{ŋrä˥t}]: властитель\\
0xB1DA \textbf{xra7t} [\emph{xrä˥t}]: великодушие\\
0xB1E2 \textbf{kxa7t} [\emph{kxä˥t}]: движимый\\
0xB1E4 \textbf{pxa7t} [\emph{pxä˥t}]: неравенства\\
0xB1EA \textbf{txa7t} [\emph{txä˥t}]: спутанный\\
0xB1F1 \textbf{bxa7t} [\emph{bxä˥t}]: Багдад\\
0xB1F2 \textbf{gxa7t} [\emph{gxä˥t}]: перемешивает\\
0xB1F3 \textbf{dxa7t} [\emph{dxä˥t}]: правая рука\\
0xB201 \textbf{myu7t} [\emph{mju˥t}]: курильщик\\
0xB202 \textbf{kyu7t} [\emph{kju˥t}]: мракобесие\\
0xB204 \textbf{pyu7t} [\emph{pju˥t}]: зуд\\
0xB206 \textbf{nyu7t} [\emph{nju˥t}]: наделять\\
0xB208 \textbf{hmu7t} [\emph{ʰmu˥t}]: Мэтью\\
0xB209 \textbf{syu7t} [\emph{sju˥t}]: акцепторы\\
0xB20A \textbf{tyu7t} [\emph{tju˥t}]: тат тат крыса\\
0xB20B \textbf{lyu7t} [\emph{lju˥t}]: кормление грудью\\
0xB210 \textbf{hku7t} [\emph{ʰku˥t}]: вычислимость\\
0xB211 \textbf{byu7t} [\emph{bju˥t}]: большая ступня\\
0xB218 \textbf{hyu7t} [\emph{ʰju˥t}]: приведены\\
0xB21A \textbf{xyu7t} [\emph{xju˥t}]: здор`ово сделать\\
0xB220 \textbf{hpu7t} [\emph{ʰpu˥t}]: рассыпалась\\
0xB228 \textbf{hwu7t} [\emph{ʰwu˥t}]: древесный\\
0xB22A \textbf{twu7t} [\emph{twu˥t}]: трубочка\\
0xB22B \textbf{lwu7t} [\emph{lwu˥t}]: пятилетие\\
0xB230 \textbf{hnu7t} [\emph{ʰnu˥t}]: внучата\\
0xB242 \textbf{ksu7t} [\emph{ksu˥t}]: шестьдесят один\\
0xB244 \textbf{psu7t} [\emph{psu˥t}]: пудель\\
0xB248 \textbf{hsu7t} [\emph{ʰsu˥t}]: заблудившись\\
0xB24A \textbf{tsu7t} [\emph{tsu˥t}]: шестьдесят три\\
0xB250 \textbf{htu7t} [\emph{ʰtu˥t}]: подтекст\\
0xB258 \textbf{hlu7t} [\emph{ʰlu˥t}]: прачка\\
0xB260 \textbf{hfu7t} [\emph{ʰfu˥t}]: мышеловка\\
0xB262 \textbf{klu7t} [\emph{klu˥t}]: огрубевший\\
0xB264 \textbf{plu7t} [\emph{plu˥t}]: молодка\\
0xB268 \textbf{hcu7t} [\emph{ʰʃu˥t}]: насыщенный\\
0xB269 \textbf{slu7t} [\emph{slu˥t}]: предпоследний\\
0xB26A \textbf{tlu7t} [\emph{tlu˥t}]: талмудический\\
0xB270 \textbf{hru7t} [\emph{ʰru˥t}]: Рут\\
0xB271 \textbf{blu7t} [\emph{blu˥t}]: бульдозер\\
0xB272 \textbf{glu7t} [\emph{glu˥t}]: копытный\\
0xB273 \textbf{dlu7t} [\emph{dlu˥t}]: друид\\
0xB282 \textbf{kfu7t} [\emph{kfu˥t}]: Kathmandu\\
0xB288 \textbf{hbu7t} [\emph{ʰbu˥t}]: девяносто два\\
0xB28A \textbf{tfu7t} [\emph{tfu˥t}]: сыщик гончая\\
0xB290 \textbf{hgu7t} [\emph{ʰgu˥t}]: крестник\\
0xB298 \textbf{hdu7t} [\emph{ʰdu˥t}]: оттуда\\
0xB2A0 \textbf{hzu7t} [\emph{ʰzu˥t}]: ростовщик\\
0xB2A2 \textbf{kcu7t} [\emph{kʃu˥t}]: открытого состава\\
0xB2AA \textbf{tcu7t} [\emph{tʃu˥t}]: плотно прижимистый\\
0xB2B0 \textbf{hvu7t} [\emph{ʰvu˥t}]: провокатор\\
0xB2C1 \textbf{mru7t} [\emph{mru˥t}]: плохим советом\\
0xB2C2 \textbf{kru7t} [\emph{kru˥t}]: квартет\\
0xB2C4 \textbf{pru7t} [\emph{pru˥t}]: породистый\\
0xB2C6 \textbf{nru7t} [\emph{nru˥t}]: пробормотал\\
0xB2C8 \textbf{hqu7t} [\emph{ʰŋu˥t}]: заготовку Doux\\
0xB2C9 \textbf{sru7t} [\emph{sru˥t}]: самурай\\
0xB2CA \textbf{tru7t} [\emph{tru˥t}]: нарушитель\\
0xB2D0 \textbf{hxu7t} [\emph{ʰxu˥t}]: хорошо делать\\
0xB2D1 \textbf{bru7t} [\emph{bru˥t}]: порох\\
0xB2D3 \textbf{dru7t} [\emph{dru˥t}]: мочегонное средство\\
0xB2D4 \textbf{zru7t} [\emph{zru˥t}]: урду\\
0xB2DA \textbf{xru7t} [\emph{xru˥t}]: вкрутую\\
0xB2E2 \textbf{kxu7t} [\emph{kxu˥t}]: открытой рукой\\
0xB2E4 \textbf{pxu7t} [\emph{pxu˥t}]: шпатлевка\\
0xB2EA \textbf{txu7t} [\emph{txu˥t}]: сосок\\
0xB301 \textbf{mye7t} [\emph{mje˥t}]: Пятикнижие\\
0xB302 \textbf{kye7t} [\emph{kje˥t}]: х по рейтингу\\
0xB304 \textbf{pye7t} [\emph{pje˥t}]: сдерживаемый\\
0xB306 \textbf{nye7t} [\emph{nje˥t}]: анти семитизм\\
0xB308 \textbf{hme7t} [\emph{ʰme˥t}]: извинить\\
0xB309 \textbf{sye7t} [\emph{sje˥t}]: шестьдесят девять\\
0xB30A \textbf{tye7t} [\emph{tje˥t}]: предвосхищать\\
0xB30B \textbf{lye7t} [\emph{lje˥t}]: макулатура\\
0xB30E \textbf{rye7t} [\emph{rje˥t}]: серенада\\
0xB310 \textbf{hke7t} [\emph{ʰke˥t}]: кокетка\\
0xB313 \textbf{dye7t} [\emph{dje˥t}]: де\\
0xB318 \textbf{hye7t} [\emph{ʰje˥t}]: маркированный\\
0xB31A \textbf{xye7t} [\emph{xje˥t}]: один ноги\\
0xB320 \textbf{hpe7t} [\emph{ʰpe˥t}]: пословица\\
0xB322 \textbf{kwe7t} [\emph{kwe˥t}]: собака ушастый\\
0xB324 \textbf{pwe7t} [\emph{pwe˥t}]: подножка\\
0xB328 \textbf{hwe7t} [\emph{ʰwe˥t}]: друг детства\\
0xB32A \textbf{twe7t} [\emph{twe˥t}]: е текст\\
0xB32E \textbf{rwe7t} [\emph{rwe˥t}]: сосед по комнате\\
0xB330 \textbf{hne7t} [\emph{ʰne˥t}]: Бенедикт\\
0xB333 \textbf{dwe7t} [\emph{dwe˥t}]: детонатор\\
0xB342 \textbf{kse7t} [\emph{kse˥t}]: этих заседаниях\\
0xB348 \textbf{hse7t} [\emph{ʰse˥t}]: шестьдесят пять\\
0xB34A \textbf{tse7t} [\emph{tse˥t}]: время отхода ко сну\\
0xB350 \textbf{hte7t} [\emph{ʰte˥t}]: и т.д\\
0xB358 \textbf{hle7t} [\emph{ʰle˥t}]: тромбоцит\\
0xB360 \textbf{hfe7t} [\emph{ʰfe˥t}]: любовно\\
0xB362 \textbf{kle7t} [\emph{kle˥t}]: реабилитировать\\
0xB364 \textbf{ple7t} [\emph{ple˥t}]: эспланада\\
0xB368 \textbf{hce7t} [\emph{ʰʃe˥t}]: жакеты\\
0xB369 \textbf{sle7t} [\emph{sle˥t}]: омела белая\\
0xB36A \textbf{tle7t} [\emph{tle˥t}]: собеседник\\
0xB36D \textbf{cle7t} [\emph{ʃle˥t}]: эполет\\
0xB370 \textbf{hre7t} [\emph{ʰre˥t}]: лучистый\\
0xB371 \textbf{ble7t} [\emph{ble˥t}]: Белград\\
0xB372 \textbf{gle7t} [\emph{gle˥t}]: буравчик глазами\\
0xB373 \textbf{dle7t} [\emph{dle˥t}]: мир устал\\
0xB37A \textbf{xle7t} [\emph{xle˥t}]: геркулесов\\
0xB382 \textbf{kfe7t} [\emph{kfe˥t}]: налог бесплатный\\
0xB388 \textbf{hbe7t} [\emph{ʰbe˥t}]: утих\\
0xB38A \textbf{tfe7t} [\emph{tfe˥t}]: тафта\\
0xB390 \textbf{hge7t} [\emph{ʰge˥t}]: монополист\\
0xB398 \textbf{hde7t} [\emph{ʰde˥t}]: одолжил\\
0xB3A0 \textbf{hze7t} [\emph{ʰze˥t}]: низводить\\
0xB3A2 \textbf{kce7t} [\emph{kʃe˥t}]: электролит\\
0xB3A4 \textbf{pce7t} [\emph{pʃe˥t}]: полубог\\
0xB3A8 \textbf{hje7t} [\emph{ʰʒe˥t}]: пересидеть\\
0xB3AA \textbf{tce7t} [\emph{tʃe˥t}]: чириканье\\
0xB3B0 \textbf{hve7t} [\emph{ʰve˥t}]: зайчонок\\
0xB3C1 \textbf{mre7t} [\emph{mre˥t}]: менуэт\\
0xB3C2 \textbf{kre7t} [\emph{kre˥t}]: Крит\\
0xB3C4 \textbf{pre7t} [\emph{pre˥t}]: спорщик\\
0xB3C6 \textbf{nre7t} [\emph{nre˥t}]: полемизировать\\
0xB3C8 \textbf{hqe7t} [\emph{ʰŋe˥t}]: сычужок\\
0xB3C9 \textbf{sre7t} [\emph{sre˥t}]: небрежность\\
0xB3CA \textbf{tre7t} [\emph{tre˥t}]: терракотовый\\
0xB3CD \textbf{cre7t} [\emph{ʃre˥t}]: четырехгранный\\
0xB3D0 \textbf{hxe7t} [\emph{ʰxe˥t}]: беспокоит\\
0xB3D1 \textbf{bre7t} [\emph{bre˥t}]: сболтнуть\\
0xB3D2 \textbf{gre7t} [\emph{gre˥t}]: нацарапал\\
0xB3D3 \textbf{dre7t} [\emph{dre˥t}]: додекаэдр\\
0xB3DA \textbf{xre7t} [\emph{xre˥t}]: гермафродит\\
0xB3E2 \textbf{kxe7t} [\emph{kxe˥t}]: снова впадать\\
0xB3E4 \textbf{pxe7t} [\emph{pxe˥t}]: нечитаемый\\
0xB3EA \textbf{txe7t} [\emph{txe˥t}]: мститель\\
0xB3F1 \textbf{bxe7t} [\emph{bxe˥t}]: биометрия\\
0xB3F2 \textbf{gxe7t} [\emph{gxe˥t}]: со-е изд\\
0xB3F3 \textbf{dxe7t} [\emph{dxe˥t}]: дециметр\\
0xB401 \textbf{myo7t} [\emph{mjo˥t}]: имитатор\\
0xB402 \textbf{kyo7t} [\emph{kjo˥t}]: кондор\\
0xB404 \textbf{pyo7t} [\emph{pjo˥t}]: неспортивный\\
0xB406 \textbf{nyo7t} [\emph{njo˥t}]: невысказанный\\
0xB408 \textbf{hmo7t} [\emph{ʰmo˥t}]: сообщества\\
0xB409 \textbf{syo7t} [\emph{sjo˥t}]: отступ\\
0xB40A \textbf{tyo7t} [\emph{tjo˥t}]: антиподы\\
0xB40B \textbf{lyo7t} [\emph{ljo˥t}]: вольтметр\\
0xB40E \textbf{ryo7t} [\emph{rjo˥t}]: рытвины\\
0xB410 \textbf{hko7t} [\emph{ʰko˥t}]: фанатик\\
0xB413 \textbf{dyo7t} [\emph{djo˥t}]: йодированной\\
0xB418 \textbf{hyo7t} [\emph{ʰjo˥t}]: желчь\\
0xB41A \textbf{xyo7t} [\emph{xjo˥t}]: hoity toity\\
0xB420 \textbf{hpo7t} [\emph{ʰpo˥t}]: почтмейстер\\
0xB422 \textbf{kwo7t} [\emph{kwo˥t}]: Киото\\
0xB424 \textbf{pwo7t} [\emph{pwo˥t}]: запачканный\\
0xB426 \textbf{nwo7t} [\emph{nwo˥t}]: берег\\
0xB428 \textbf{hwo7t} [\emph{ʰwo˥t}]: проводной\\
0xB429 \textbf{swo7t} [\emph{swo˥t}]: жатка\\
0xB42A \textbf{two7t} [\emph{two˥t}]: трезвенник\\
0xB42B \textbf{lwo7t} [\emph{lwo˥t}]: телезритель\\
0xB42E \textbf{rwo7t} [\emph{rwo˥t}]: красный верх\\
0xB430 \textbf{hno7t} [\emph{ʰno˥t}]: к северо-востоку\\
0xB442 \textbf{kso7t} [\emph{kso˥t}]: обуздана\\
0xB444 \textbf{pso7t} [\emph{pso˥t}]: перекись\\
0xB448 \textbf{hso7t} [\emph{ʰso˥t}]: содрогнулся\\
0xB44A \textbf{tso7t} [\emph{tso˥t}]: поверь\\
0xB450 \textbf{hto7t} [\emph{ʰto˥t}]: порванный\\
0xB458 \textbf{hlo7t} [\emph{ʰlo˥t}]: погрузка\\
0xB460 \textbf{hfo7t} [\emph{ʰfo˥t}]: фотон\\
0xB462 \textbf{klo7t} [\emph{klo˥t}]: с открытым ртом\\
0xB464 \textbf{plo7t} [\emph{plo˥t}]: орфографический справочник\\
0xB468 \textbf{hco7t} [\emph{ʰʃo˥t}]: пришелец\\
0xB469 \textbf{slo7t} [\emph{slo˥t}]: энциклопедия\\
0xB46A \textbf{tlo7t} [\emph{tlo˥t}]: кефаль\\
0xB470 \textbf{hro7t} [\emph{ʰro˥t}]: подруливающее\\
0xB471 \textbf{blo7t} [\emph{blo˥t}]: сплюснутый\\
0xB473 \textbf{dlo7t} [\emph{dlo˥t}]: дуплет\\
0xB47A \textbf{xlo7t} [\emph{xlo˥t}]: колесный мастер\\
0xB482 \textbf{kfo7t} [\emph{kfo˥t}]: я делаю\\
0xB484 \textbf{pfo7t} [\emph{pfo˥t}]: огнемет\\
0xB488 \textbf{hbo7t} [\emph{ʰbo˥t}]: комбинаторика\\
0xB48A \textbf{tfo7t} [\emph{tfo˥t}]: коммутатор\\
0xB490 \textbf{hgo7t} [\emph{ʰgo˥t}]: подагра\\
0xB498 \textbf{hdo7t} [\emph{ʰdo˥t}]: рассудите\\
0xB4A0 \textbf{hzo7t} [\emph{ʰzo˥t}]: зороастрийский\\
0xB4A2 \textbf{kco7t} [\emph{kʃo˥t}]: комод\\
0xB4A4 \textbf{pco7t} [\emph{pʃo˥t}]: подъязычная\\
0xB4A8 \textbf{hjo7t} [\emph{ʰʒo˥t}]: хроматин\\
0xB4AA \textbf{tco7t} [\emph{tʃo˥t}]: смертное ложе\\
0xB4B0 \textbf{hvo7t} [\emph{ʰvo˥t}]: преувеличивать\\
0xB4C1 \textbf{mro7t} [\emph{mro˥t}]: модный\\
0xB4C2 \textbf{kro7t} [\emph{kro˥t}]: раздвоенный\\
0xB4C4 \textbf{pro7t} [\emph{pro˥t}]: пенистый\\
0xB4C6 \textbf{nro7t} [\emph{nro˥t}]: взревел\\
0xB4C8 \textbf{hqo7t} [\emph{ʰŋo˥t}]: Anglo индийский\\
0xB4C9 \textbf{sro7t} [\emph{sro˥t}]: посмертный\\
0xB4CA \textbf{tro7t} [\emph{tro˥t}]: тележки\\
0xB4CC \textbf{fro7t} [\emph{fro˥t}]: волокнистый\\
0xB4CD \textbf{cro7t} [\emph{ʃro˥t}]: сонная артерия\\
0xB4D0 \textbf{hxo7t} [\emph{ʰxo˥t}]: над породой\\
0xB4D1 \textbf{bro7t} [\emph{bro˥t}]: подпольный акушер\\
0xB4D2 \textbf{gro7t} [\emph{gro˥t}]: меди дном\\
0xB4D3 \textbf{dro7t} [\emph{dro˥t}]: любить до безумия\\
0xB4D9 \textbf{qro7t} [\emph{ŋro˥t}]: Несчастные\\
0xB4DA \textbf{xro7t} [\emph{xro˥t}]: белая жесть\\
0xB4E2 \textbf{kxo7t} [\emph{kxo˥t}]: заикание\\
0xB4E4 \textbf{pxo7t} [\emph{pxo˥t}]: пятисотый\\
0xB4EA \textbf{txo7t} [\emph{txo˥t}]: фригольд\\
0xB4F1 \textbf{bxo7t} [\emph{bxo˥t}]: закупоривать отверстие\\
0xB4F2 \textbf{gxo7t} [\emph{gxo˥t}]: совместно неавтоматического\\
0xB4F3 \textbf{dxo7t} [\emph{dxo˥t}]: без запаха\\
0xB504 \textbf{py67t} [\emph{pjə˥t}]: брамин\\
0xB508 \textbf{hm67t} [\emph{ʰmə˥t}]: Суматра\\
0xB50E \textbf{ry67t} [\emph{rjə˥t}]: Рияд\\
0xB510 \textbf{hk67t} [\emph{ʰkə˥t}]: коммодор\\
0xB518 \textbf{hy67t} [\emph{ʰjə˥t}]: йодль\\
0xB520 \textbf{hp67t} [\emph{ʰpə˥t}]: Pitter скороговорка\\
0xB528 \textbf{hw67t} [\emph{ʰwə˥t}]: переваливаться\\
0xB530 \textbf{hn67t} [\emph{ʰnə˥t}]: нагота\\
0xB542 \textbf{ks67t} [\emph{ksə˥t}]: горьковато-сладкий\\
0xB544 \textbf{ps67t} [\emph{psə˥t}]: компенсатор\\
0xB548 \textbf{hs67t} [\emph{ʰsə˥t}]: приспособленного\\
0xB54A \textbf{ts67t} [\emph{tsə˥t}]: Троица\\
0xB550 \textbf{ht67t} [\emph{ʰtə˥t}]: пораженный\\
0xB558 \textbf{hl67t} [\emph{ʰlə˥t}]: вычерпывать\\
0xB562 \textbf{kl67t} [\emph{klə˥t}]: циклотрон\\
0xB568 \textbf{hc67t} [\emph{ʰʃə˥t}]: Sindhi\\
0xB56A \textbf{tl67t} [\emph{tlə˥t}]: Средиземье\\
0xB570 \textbf{hr67t} [\emph{ʰrə˥t}]: ушная боль\\
0xB573 \textbf{dl67t} [\emph{dlə˥t}]: диплоид\\
0xB588 \textbf{hb67t} [\emph{ʰbə˥t}]: придерживалась\\
0xB598 \textbf{hd67t} [\emph{ʰdə˥t}]: дмитрий\\
0xB5AA \textbf{tc67t} [\emph{tʃə˥t}]: том Тит\\
0xB5C1 \textbf{mr67t} [\emph{mrə˥t}]: миокард\\
0xB5C2 \textbf{kr67t} [\emph{krə˥t}]: курд\\
0xB5C9 \textbf{sr67t} [\emph{srə˥t}]: прямой край\\
0xB5CA \textbf{tr67t} [\emph{trə˥t}]: звездная пыль\\
0xB5D0 \textbf{hx67t} [\emph{ʰxə˥t}]: хеттский\\
0xB5E2 \textbf{kx67t} [\emph{kxə˥t}]: аккадское\\
0xB5EA \textbf{tx67t} [\emph{txə˥t}]: тик-так\\
0xB801 \textbf{myi2t} [\emph{mji˩t}]: менструация\\
0xB802 \textbf{kyi2t} [\emph{kji˩t}]: самописец\\
0xB804 \textbf{pyi2t} [\emph{pji˩t}]: преднамеренность\\
0xB806 \textbf{nyi2t} [\emph{nji˩t}]: индис\\
0xB808 \textbf{hmi2t} [\emph{ʰmi˩t}]: математики\\
0xB809 \textbf{syi2t} [\emph{sji˩t}]: сослагательное наклонение\\
0xB80A \textbf{tyi2t} [\emph{tji˩t}]: двадцать восемь\\
0xB80B \textbf{lyi2t} [\emph{lji˩t}]: лира\\
0xB80C \textbf{fyi2t} [\emph{fji˩t}]: тематический\\
0xB80D \textbf{cyi2t} [\emph{ʃji˩t}]: гелиоцентрический\\
0xB80E \textbf{ryi2t} [\emph{rji˩t}]: сорок два\\
0xB810 \textbf{hki2t} [\emph{ʰki˩t}]: бигуди\\
0xB811 \textbf{byi2t} [\emph{bji˩t}]: козел отпущения\\
0xB813 \textbf{dyi2t} [\emph{dji˩t}]: адвентисты\\
0xB81A \textbf{xyi2t} [\emph{xji˩t}]: недалекий\\
0xB820 \textbf{hpi2t} [\emph{ʰpi˩t}]: провидение\\
0xB821 \textbf{mwi2t} [\emph{mwi˩t}]: сторонник строгой дисциплины\\
0xB822 \textbf{kwi2t} [\emph{kwi˩t}]: герметический\\
0xB824 \textbf{pwi2t} [\emph{pwi˩t}]: предварительно adamite\\
0xB826 \textbf{nwi2t} [\emph{nwi˩t}]: наветренный\\
0xB828 \textbf{hwi2t} [\emph{ʰwi˩t}]: болотные сапоги\\
0xB829 \textbf{swi2t} [\emph{swi˩t}]: тридцать семь\\
0xB82A \textbf{twi2t} [\emph{twi˩t}]: опросный лист\\
0xB82B \textbf{lwi2t} [\emph{lwi˩t}]: гробница\\
0xB82E \textbf{rwi2t} [\emph{rwi˩t}]: переделанный\\
0xB830 \textbf{hni2t} [\emph{ʰni˩t}]: несправедливость\\
0xB839 \textbf{qwi2t} [\emph{ŋwi˩t}]: цистит\\
0xB83A \textbf{xwi2t} [\emph{xwi˩t}]: дыра в стене\\
0xB842 \textbf{ksi2t} [\emph{ksi˩t}]: тридцать один\\
0xB844 \textbf{psi2t} [\emph{psi˩t}]: наставник\\
0xB848 \textbf{hsi2t} [\emph{ʰsi˩t}]: гениталии\\
0xB84A \textbf{tsi2t} [\emph{tsi˩t}]: благочестивый\\
0xB850 \textbf{hti2t} [\emph{ʰti˩t}]: сохраняющий нейтралитет\\
0xB858 \textbf{hli2t} [\emph{ʰli˩t}]: заботливость\\
0xB860 \textbf{hfi2t} [\emph{ʰfi˩t}]: тридцать два\\
0xB862 \textbf{kli2t} [\emph{kli˩t}]: заклепка\\
0xB864 \textbf{pli2t} [\emph{pli˩t}]: сифилис\\
0xB868 \textbf{hci2t} [\emph{ʰʃi˩t}]: районы\\
0xB86A \textbf{tli2t} [\emph{tli˩t}]: тильда\\
0xB86C \textbf{fli2t} [\emph{fli˩t}]: подстерегать\\
0xB86D \textbf{cli2t} [\emph{ʃli˩t}]: промокашка\\
0xB870 \textbf{hri2t} [\emph{ʰri˩t}]: вероломство\\
0xB871 \textbf{bli2t} [\emph{bli˩t}]: Гибралтар\\
0xB872 \textbf{gli2t} [\emph{gli˩t}]: гильотина\\
0xB873 \textbf{dli2t} [\emph{dli˩t}]: дипломат\\
0xB876 \textbf{vli2t} [\emph{vli˩t}]: виноградарство\\
0xB87A \textbf{xli2t} [\emph{xli˩t}]: фильтры\\
0xB882 \textbf{kfi2t} [\emph{kfi˩t}]: нива\\
0xB884 \textbf{pfi2t} [\emph{pfi˩t}]: неподписанный\\
0xB888 \textbf{hbi2t} [\emph{ʰbi˩t}]: двадцать два\\
0xB88A \textbf{tfi2t} [\emph{tfi˩t}]: тридцать четыре\\
0xB890 \textbf{hgi2t} [\emph{ʰgi˩t}]: позолоченный\\
0xB891 \textbf{bvi2t} [\emph{bvi˩t}]: детский вязаный башмачок\\
0xB892 \textbf{gvi2t} [\emph{gvi˩t}]: циветта\\
0xB898 \textbf{hdi2t} [\emph{ʰdi˩t}]: гагачий\\
0xB8A0 \textbf{hzi2t} [\emph{ʰzi˩t}]: зажимать\\
0xB8A2 \textbf{kci2t} [\emph{kʃi˩t}]: с открытым сердцем\\
0xB8A4 \textbf{pci2t} [\emph{pʃi˩t}]: педантичность\\
0xB8A8 \textbf{hji2t} [\emph{ʰʒi˩t}]: ржать\\
0xB8AA \textbf{tci2t} [\emph{tʃi˩t}]: к тому же\\
0xB8B0 \textbf{hvi2t} [\emph{ʰvi˩t}]: завещать\\
0xB8B3 \textbf{dji2t} [\emph{dʒi˩t}]: конец света\\
0xB8C1 \textbf{mri2t} [\emph{mri˩t}]: сталкиваются\\
0xB8C2 \textbf{kri2t} [\emph{kri˩t}]: обмотанный\\
0xB8C4 \textbf{pri2t} [\emph{pri˩t}]: патриций\\
0xB8C6 \textbf{nri2t} [\emph{nri˩t}]: некоммерческий\\
0xB8C8 \textbf{hqi2t} [\emph{ʰŋi˩t}]: самодельный\\
0xB8C9 \textbf{sri2t} [\emph{sri˩t}]: казалось\\
0xB8CC \textbf{fri2t} [\emph{fri˩t}]: низкое качество\\
0xB8CD \textbf{cri2t} [\emph{ʃri˩t}]: содружество\\
0xB8D0 \textbf{hxi2t} [\emph{ʰxi˩t}]: встряхивают\\
0xB8D1 \textbf{bri2t} [\emph{bri˩t}]: либретто\\
0xB8D2 \textbf{gri2t} [\emph{gri˩t}]: буравчик\\
0xB8D3 \textbf{dri2t} [\emph{dri˩t}]: дезертир\\
0xB8D6 \textbf{vri2t} [\emph{vri˩t}]: оспа\\
0xB8D9 \textbf{qri2t} [\emph{ŋri˩t}]: индивидуальные\\
0xB8DA \textbf{xri2t} [\emph{xri˩t}]: засахаренный\\
0xB8E2 \textbf{kxi2t} [\emph{kxi˩t}]: подсвечник\\
0xB8E4 \textbf{pxi2t} [\emph{pxi˩t}]: проверенный\\
0xB8EA \textbf{txi2t} [\emph{txi˩t}]: замечающий\\
0xB8F1 \textbf{bxi2t} [\emph{bxi˩t}]: концы света\\
0xB8F2 \textbf{gxi2t} [\emph{gxi˩t}]: агглютинативный\\
0xB8F3 \textbf{dxi2t} [\emph{dxi˩t}]: дифтерия\\
0xB901 \textbf{mya2t} [\emph{mjä˩t}]: солод\\
0xB902 \textbf{kya2t} [\emph{kjä˩t}]: кокарда\\
0xB904 \textbf{pya2t} [\emph{pjä˩t}]: торговец шляпами\\
0xB906 \textbf{nya2t} [\emph{njä˩t}]: манифест\\
0xB908 \textbf{hma2t} [\emph{ʰmä˩t}]: изменник\\
0xB909 \textbf{sya2t} [\emph{sjä˩t}]: шестьдесят восемь\\
0xB90A \textbf{tya2t} [\emph{tjä˩t}]: непроверенный\\
0xB90B \textbf{lya2t} [\emph{ljä˩t}]: флагшток\\
0xB90C \textbf{fya2t} [\emph{fjä˩t}]: совместное времяпровождение\\
0xB90D \textbf{cya2t} [\emph{ʃjä˩t}]: половина мачты\\
0xB90E \textbf{rya2t} [\emph{rjä˩t}]: уважаемый\\
0xB910 \textbf{hka2t} [\emph{ʰkä˩t}]: товарищи\\
0xB911 \textbf{bya2t} [\emph{bjä˩t}]: подстрекательстве\\
0xB912 \textbf{gya2t} [\emph{gjä˩t}]: стекловидный глазами\\
0xB913 \textbf{dya2t} [\emph{djä˩t}]: сердце раздирающие\\
0xB915 \textbf{jya2t} [\emph{ʒjä˩t}]: джихад\\
0xB918 \textbf{hya2t} [\emph{ʰjä˩t}]: пророк\\
0xB919 \textbf{qya2t} [\emph{ŋjä˩t}]: воздушная почта\\
0xB91A \textbf{xya2t} [\emph{xjä˩t}]: нижеподписавшиеся\\
0xB920 \textbf{hpa2t} [\emph{ʰpä˩t}]: посаженный\\
0xB921 \textbf{mwa2t} [\emph{mwä˩t}]: выбритый\\
0xB922 \textbf{kwa2t} [\emph{kwä˩t}]: благоговение ударил\\
0xB924 \textbf{pwa2t} [\emph{pwä˩t}]: пагода\\
0xB926 \textbf{nwa2t} [\emph{nwä˩t}]: санскрит\\
0xB928 \textbf{hwa2t} [\emph{ʰwä˩t}]: положение\\
0xB929 \textbf{swa2t} [\emph{swä˩t}]: восемьдесят два\\
0xB92A \textbf{twa2t} [\emph{twä˩t}]: паническое бегство\\
0xB92B \textbf{lwa2t} [\emph{lwä˩t}]: лассо\\
0xB92E \textbf{rwa2t} [\emph{rwä˩t}]: арьергард\\
0xB930 \textbf{hna2t} [\emph{ʰnä˩t}]: кивнул\\
0xB931 \textbf{bwa2t} [\emph{bwä˩t}]: эстрада для оркестра\\
0xB933 \textbf{dwa2t} [\emph{dwä˩t}]: сандаловое дерево\\
0xB93A \textbf{xwa2t} [\emph{xwä˩t}]: выбрасывать\\
0xB942 \textbf{ksa2t} [\emph{ksä˩t}]: каста\\
0xB944 \textbf{psa2t} [\emph{psä˩t}]: вялость\\
0xB948 \textbf{hsa2t} [\emph{ʰsä˩t}]: над урожаем\\
0xB94A \textbf{tsa2t} [\emph{tsä˩t}]: трижды\\
0xB950 \textbf{hta2t} [\emph{ʰtä˩t}]: полеводства\\
0xB953 \textbf{dza2t} [\emph{dzä˩t}]: обескураживать\\
0xB958 \textbf{hla2t} [\emph{ʰlä˩t}]: лепились\\
0xB960 \textbf{hfa2t} [\emph{ʰfä˩t}]: феодальный\\
0xB962 \textbf{kla2t} [\emph{klä˩t}]: оригинальность\\
0xB964 \textbf{pla2t} [\emph{plä˩t}]: блюдо\\
0xB968 \textbf{hca2t} [\emph{ʰʃä˩t}]: упал\\
0xB969 \textbf{sla2t} [\emph{slä˩t}]: принадлежит\\
0xB96A \textbf{tla2t} [\emph{tlä˩t}]: прыгнули\\
0xB96C \textbf{fla2t} [\emph{flä˩t}]: формальдегид\\
0xB96D \textbf{cla2t} [\emph{ʃlä˩t}]: трудоспособный\\
0xB970 \textbf{hra2t} [\emph{ʰrä˩t}]: перекрестная проверка\\
0xB971 \textbf{bla2t} [\emph{blä˩t}]: бильярд\\
0xB972 \textbf{gla2t} [\emph{glä˩t}]: гранд тетя\\
0xB973 \textbf{dla2t} [\emph{dlä˩t}]: медлительный\\
0xB974 \textbf{zla2t} [\emph{zlä˩t}]: долг нагруженные\\
0xB975 \textbf{jla2t} [\emph{ʒlä˩t}]: желатин\\
0xB97A \textbf{xla2t} [\emph{xlä˩t}]: глашатаи\\
0xB982 \textbf{kfa2t} [\emph{kfä˩t}]: кафтан\\
0xB984 \textbf{pfa2t} [\emph{pfä˩t}]: помогать\\
0xB988 \textbf{hba2t} [\emph{ʰbä˩t}]: группы\\
0xB98A \textbf{tfa2t} [\emph{tfä˩t}]: брюшной тиф\\
0xB990 \textbf{hga2t} [\emph{ʰgä˩t}]: полученный\\
0xB992 \textbf{gva2t} [\emph{gvä˩t}]: архивируются\\
0xB993 \textbf{dva2t} [\emph{dvä˩t}]: добавлений\\
0xB998 \textbf{hda2t} [\emph{ʰdä˩t}]: приводной вал\\
0xB9A0 \textbf{hza2t} [\emph{ʰzä˩t}]: прогорклый\\
0xB9A2 \textbf{kca2t} [\emph{kʃä˩t}]: атолл\\
0xB9A4 \textbf{pca2t} [\emph{pʃä˩t}]: наглость\\
0xB9A8 \textbf{hja2t} [\emph{ʰʒä˩t}]: отчим\\
0xB9AA \textbf{tca2t} [\emph{tʃä˩t}]: непредсказуемость\\
0xB9B0 \textbf{hva2t} [\emph{ʰvä˩t}]: вульгарность\\
0xB9B3 \textbf{dja2t} [\emph{dʒä˩t}]: адиабатический\\
0xB9C1 \textbf{mra2t} [\emph{mrä˩t}]: Амстердам\\
0xB9C2 \textbf{kra2t} [\emph{krä˩t}]: откровенность\\
0xB9C4 \textbf{pra2t} [\emph{prä˩t}]: катаракта\\
0xB9C6 \textbf{nra2t} [\emph{nrä˩t}]: консерватория\\
0xB9C8 \textbf{hqa2t} [\emph{ʰŋä˩t}]: мангал\\
0xB9C9 \textbf{sra2t} [\emph{srä˩t}]: спартанский\\
0xB9CC \textbf{fra2t} [\emph{frä˩t}]: братоубийство\\
0xB9CD \textbf{cra2t} [\emph{ʃrä˩t}]: мученичество\\
0xB9D0 \textbf{hxa2t} [\emph{ʰxä˩t}]: крики\\
0xB9D1 \textbf{bra2t} [\emph{brä˩t}]: оправданный\\
0xB9D2 \textbf{gra2t} [\emph{grä˩t}]: лишать места\\
0xB9D3 \textbf{dra2t} [\emph{drä˩t}]: подъемный мост\\
0xB9D4 \textbf{zra2t} [\emph{zrä˩t}]: разветвленная далеко\\
0xB9D5 \textbf{jra2t} [\emph{ʒrä˩t}]: кариатида\\
0xB9D6 \textbf{vra2t} [\emph{vrä˩t}]: пр`оклятый\\
0xB9D9 \textbf{qra2t} [\emph{ŋrä˩t}]: расположить к себе\\
0xB9DA \textbf{xra2t} [\emph{xrä˩t}]: разглагольствовать\\
0xB9E2 \textbf{kxa2t} [\emph{kxä˩t}]: верхом\\
0xB9E4 \textbf{pxa2t} [\emph{pxä˩t}]: топ мачты\\
0xB9EA \textbf{txa2t} [\emph{txä˩t}]: Attends\\
0xB9F1 \textbf{bxa2t} [\emph{bxä˩t}]: повязка на глаза\\
0xB9F2 \textbf{gxa2t} [\emph{gxä˩t}]: цареубийство\\
0xB9F3 \textbf{dxa2t} [\emph{dxä˩t}]: дидактический\\
0xBA06 \textbf{nyu2t} [\emph{nju˩t}]: синусит\\
0xBA08 \textbf{hmu2t} [\emph{ʰmu˩t}]: Кульминацией\\
0xBA09 \textbf{syu2t} [\emph{sju˩t}]: пломбир\\
0xBA0A \textbf{tyu2t} [\emph{tju˩t}]: эвтектика\\
0xBA0B \textbf{lyu2t} [\emph{lju˩t}]: уличный фонарь\\
0xBA0E \textbf{ryu2t} [\emph{rju˩t}]: тарантул\\
0xBA10 \textbf{hku2t} [\emph{ʰku˩t}]: Амур\\
0xBA11 \textbf{byu2t} [\emph{bju˩t}]: бермудский\\
0xBA18 \textbf{hyu2t} [\emph{ʰju˩t}]: Иуда\\
0xBA20 \textbf{hpu2t} [\emph{ʰpu˩t}]: депутация\\
0xBA28 \textbf{hwu2t} [\emph{ʰwu˩t}]: колдовство\\
0xBA2B \textbf{lwu2t} [\emph{lwu˩t}]: впередсмотрящие\\
0xBA30 \textbf{hnu2t} [\emph{ʰnu˩t}]: унитарный\\
0xBA42 \textbf{ksu2t} [\emph{ksu˩t}]: счетчики\\
0xBA44 \textbf{psu2t} [\emph{psu˩t}]: непригодность\\
0xBA48 \textbf{hsu2t} [\emph{ʰsu˩t}]: превратности\\
0xBA4A \textbf{tsu2t} [\emph{tsu˩t}]: шестьдесят два\\
0xBA50 \textbf{htu2t} [\emph{ʰtu˩t}]: враждебность\\
0xBA58 \textbf{hlu2t} [\emph{ʰlu˩t}]: политеистический\\
0xBA62 \textbf{klu2t} [\emph{klu˩t}]: нуклеотид\\
0xBA64 \textbf{plu2t} [\emph{plu˩t}]: подколенный\\
0xBA68 \textbf{hcu2t} [\emph{ʰʃu˩t}]: у лодки\\
0xBA69 \textbf{slu2t} [\emph{slu˩t}]: мускулатура\\
0xBA6A \textbf{tlu2t} [\emph{tlu˩t}]: неокончательный\\
0xBA6D \textbf{clu2t} [\emph{ʃlu˩t}]: абсолютный ноль\\
0xBA70 \textbf{hru2t} [\emph{ʰru˩t}]: колесница\\
0xBA7A \textbf{xlu2t} [\emph{xlu˩t}]: гипохлорит\\
0xBA88 \textbf{hbu2t} [\emph{ʰbu˩t}]: двуногое\\
0xBA8A \textbf{tfu2t} [\emph{tfu˩t}]: тутти фрутти\\
0xBA90 \textbf{hgu2t} [\emph{ʰgu˩t}]: Argus глазами\\
0xBA98 \textbf{hdu2t} [\emph{ʰdu˩t}]: второй курс\\
0xBAA0 \textbf{hzu2t} [\emph{ʰzu˩t}]: толстовка с капюшоном\\
0xBAAA \textbf{tcu2t} [\emph{tʃu˩t}]: поганка\\
0xBAC1 \textbf{mru2t} [\emph{mru˩t}]: миниатюрами\\
0xBAC2 \textbf{kru2t} [\emph{kru˩t}]: эукариоты\\
0xBAC4 \textbf{pru2t} [\emph{pru˩t}]: штифтик\\
0xBAC6 \textbf{nru2t} [\emph{nru˩t}]: невротический\\
0xBAC9 \textbf{sru2t} [\emph{sru˩t}]: цитрат\\
0xBACA \textbf{tru2t} [\emph{tru˩t}]: антихрист\\
0xBACC \textbf{fru2t} [\emph{fru˩t}]: фторид\\
0xBAD0 \textbf{hxu2t} [\emph{ʰxu˩t}]: кобура\\
0xBAD1 \textbf{bru2t} [\emph{bru˩t}]: Бейрут\\
0xBAD3 \textbf{dru2t} [\emph{dru˩t}]: смущаясь\\
0xBADA \textbf{xru2t} [\emph{xru˩t}]: казнить на электрическом стуле\\
0xBAE2 \textbf{kxu2t} [\emph{kxu˩t}]: лицо боль\\
0xBAEA \textbf{txu2t} [\emph{txu˩t}]: трудно носатый\\
0xBB01 \textbf{mye2t} [\emph{mje˩t}]: армагеддон\\
0xBB02 \textbf{kye2t} [\emph{kje˩t}]: секстет\\
0xBB04 \textbf{pye2t} [\emph{pje˩t}]: перитонит\\
0xBB06 \textbf{nye2t} [\emph{nje˩t}]: ланцетовидный\\
0xBB08 \textbf{hme2t} [\emph{ʰme˩t}]: на полпути\\
0xBB09 \textbf{sye2t} [\emph{sje˩t}]: ацетат\\
0xBB0A \textbf{tye2t} [\emph{tje˩t}]: семьдесят восемь\\
0xBB0B \textbf{lye2t} [\emph{lje˩t}]: кровельщик\\
0xBB0E \textbf{rye2t} [\emph{rje˩t}]: повторного входа\\
0xBB10 \textbf{hke2t} [\emph{ʰke˩t}]: дородный\\
0xBB18 \textbf{hye2t} [\emph{ʰje˩t}]: шагомер\\
0xBB20 \textbf{hpe2t} [\emph{ʰpe˩t}]: незначительно\\
0xBB21 \textbf{mwe2t} [\emph{mwe˩t}]: тире\\
0xBB22 \textbf{kwe2t} [\emph{kwe˩t}]: кемпинг\\
0xBB26 \textbf{nwe2t} [\emph{nwe˩t}]: женское нижнее белье\\
0xBB28 \textbf{hwe2t} [\emph{ʰwe˩t}]: эклиптика\\
0xBB29 \textbf{swe2t} [\emph{swe˩t}]: так и так\\
0xBB2A \textbf{twe2t} [\emph{twe˩t}]: лапками\\
0xBB2E \textbf{rwe2t} [\emph{rwe˩t}]: схема\\
0xBB30 \textbf{hne2t} [\emph{ʰne˩t}]: упорно\\
0xBB42 \textbf{kse2t} [\emph{kse˩t}]: пономарь\\
0xBB44 \textbf{pse2t} [\emph{pse˩t}]: шестьдесят семь\\
0xBB48 \textbf{hse2t} [\emph{ʰse˩t}]: эскадрилья\\
0xBB4A \textbf{tse2t} [\emph{tse˩t}]: ч.л.\\
0xBB50 \textbf{hte2t} [\emph{ʰte˩t}]: достиг\\
0xBB58 \textbf{hle2t} [\emph{ʰle˩t}]: блаженный\\
0xBB60 \textbf{hfe2t} [\emph{ʰfe˩t}]: фетиш\\
0xBB62 \textbf{kle2t} [\emph{kle˩t}]: выздоравливающий\\
0xBB64 \textbf{ple2t} [\emph{ple˩t}]: предоплатой\\
0xBB68 \textbf{hce2t} [\emph{ʰʃe˩t}]: вымпел\\
0xBB69 \textbf{sle2t} [\emph{sle˩t}]: оксалат\\
0xBB6A \textbf{tle2t} [\emph{tle˩t}]: сетчатый\\
0xBB6D \textbf{cle2t} [\emph{ʃle˩t}]: шнурок для ботинок\\
0xBB70 \textbf{hre2t} [\emph{ʰre˩t}]: приготовился\\
0xBB71 \textbf{ble2t} [\emph{ble˩t}]: бакелит\\
0xBB73 \textbf{dle2t} [\emph{dle˩t}]: медлить хвостатых\\
0xBB7A \textbf{xle2t} [\emph{xle˩t}]: швырнула\\
0xBB88 \textbf{hbe2t} [\emph{ʰbe˩t}]: двор замка\\
0xBB8A \textbf{tfe2t} [\emph{tfe˩t}]: неподтвержденный\\
0xBB90 \textbf{hge2t} [\emph{ʰge˩t}]: Гентский\\
0xBB98 \textbf{hde2t} [\emph{ʰde˩t}]: божествами\\
0xBBA0 \textbf{hze2t} [\emph{ʰze˩t}]: метаданные\\
0xBBA2 \textbf{kce2t} [\emph{kʃe˩t}]: декстрин\\
0xBBA4 \textbf{pce2t} [\emph{pʃe˩t}]: яблоко щеками\\
0xBBAA \textbf{tce2t} [\emph{tʃe˩t}]: рикошет\\
0xBBB0 \textbf{hve2t} [\emph{ʰve˩t}]: вендетта\\
0xBBC1 \textbf{mre2t} [\emph{mre˩t}]: погребать\\
0xBBC2 \textbf{kre2t} [\emph{kre˩t}]: картель\\
0xBBC4 \textbf{pre2t} [\emph{pre˩t}]: репертуар\\
0xBBC6 \textbf{nre2t} [\emph{nre˩t}]: Волшебная страна\\
0xBBC8 \textbf{hqe2t} [\emph{ʰŋe˩t}]: ястреб глазами\\
0xBBC9 \textbf{sre2t} [\emph{sre˩t}]: опытным путем\\
0xBBCA \textbf{tre2t} [\emph{tre˩t}]: медлить\\
0xBBCD \textbf{cre2t} [\emph{ʃre˩t}]: оркестровать\\
0xBBD0 \textbf{hxe2t} [\emph{ʰxe˩t}]: подголовник\\
0xBBD1 \textbf{bre2t} [\emph{bre˩t}]: коечные\\
0xBBD3 \textbf{dre2t} [\emph{dre˩t}]: расхищение\\
0xBBD5 \textbf{jre2t} [\emph{ʒre˩t}]: мартышка\\
0xBBDA \textbf{xre2t} [\emph{xre˩t}]: гидрат\\
0xBBE2 \textbf{kxe2t} [\emph{kxe˩t}]: квинтет\\
0xBBE4 \textbf{pxe2t} [\emph{pxe˩t}]: наборная машина\\
0xBBEA \textbf{txe2t} [\emph{txe˩t}]: Билеты\\
0xBBF2 \textbf{gxe2t} [\emph{gxe˩t}]: преуспевать\\
0xBBF3 \textbf{dxe2t} [\emph{dxe˩t}]: альдегид\\
0xBC01 \textbf{myo2t} [\emph{mjo˩t}]: Человек о войне\\
0xBC02 \textbf{kyo2t} [\emph{kjo˩t}]: эконометрия\\
0xBC04 \textbf{pyo2t} [\emph{pjo˩t}]: помадить\\
0xBC06 \textbf{nyo2t} [\emph{njo˩t}]: гностический\\
0xBC08 \textbf{hmo2t} [\emph{ʰmo˩t}]: Ломбардия\\
0xBC09 \textbf{syo2t} [\emph{sjo˩t}]: соя\\
0xBC0A \textbf{tyo2t} [\emph{tjo˩t}]: эфирное масло\\
0xBC0B \textbf{lyo2t} [\emph{ljo˩t}]: многогранник\\
0xBC0E \textbf{ryo2t} [\emph{rjo˩t}]: женское общество\\
0xBC10 \textbf{hko2t} [\emph{ʰko˩t}]: сердечно\\
0xBC18 \textbf{hyo2t} [\emph{ʰjo˩t}]: бородавка\\
0xBC20 \textbf{hpo2t} [\emph{ʰpo˩t}]: чистилище\\
0xBC21 \textbf{mwo2t} [\emph{mwo˩t}]: мальтоза\\
0xBC22 \textbf{kwo2t} [\emph{kwo˩t}]: воловий хвост\\
0xBC26 \textbf{nwo2t} [\emph{nwo˩t}]: тритон\\
0xBC28 \textbf{hwo2t} [\emph{ʰwo˩t}]: рогоносец\\
0xBC29 \textbf{swo2t} [\emph{swo˩t}]: Сальвадорских\\
0xBC2E \textbf{rwo2t} [\emph{rwo˩t}]: броненосец\\
0xBC30 \textbf{hno2t} [\emph{ʰno˩t}]: умолять\\
0xBC42 \textbf{kso2t} [\emph{kso˩t}]: украшенный гребнем\\
0xBC44 \textbf{pso2t} [\emph{pso˩t}]: спонсируемый\\
0xBC48 \textbf{hso2t} [\emph{ʰso˩t}]: скотч\\
0xBC4A \textbf{tso2t} [\emph{tso˩t}]: шестьдесят четыре\\
0xBC50 \textbf{hto2t} [\emph{ʰto˩t}]: кирпич\\
0xBC58 \textbf{hlo2t} [\emph{ʰlo˩t}]: Лондон\\
0xBC60 \textbf{hfo2t} [\emph{ʰfo˩t}]: немодный\\
0xBC62 \textbf{klo2t} [\emph{klo˩t}]: коллодий\\
0xBC64 \textbf{plo2t} [\emph{plo˩t}]: Платон\\
0xBC68 \textbf{hco2t} [\emph{ʰʃo˩t}]: домой\\
0xBC69 \textbf{slo2t} [\emph{slo˩t}]: мопед\\
0xBC6A \textbf{tlo2t} [\emph{tlo˩t}]: семядоля\\
0xBC6D \textbf{clo2t} [\emph{ʃlo˩t}]: эсхатология\\
0xBC70 \textbf{hro2t} [\emph{ʰro˩t}]: ректор\\
0xBC73 \textbf{dlo2t} [\emph{dlo˩t}]: гидрохлорид\\
0xBC7A \textbf{xlo2t} [\emph{xlo˩t}]: легированная\\
0xBC82 \textbf{kfo2t} [\emph{kfo˩t}]: скакать\\
0xBC88 \textbf{hbo2t} [\emph{ʰbo˩t}]: тупость\\
0xBC8A \textbf{tfo2t} [\emph{tfo˩t}]: карапуз\\
0xBC90 \textbf{hgo2t} [\emph{ʰgo˩t}]: стойки ворот\\
0xBC98 \textbf{hdo2t} [\emph{ʰdo˩t}]: DonT\\
0xBCA2 \textbf{kco2t} [\emph{kʃo˩t}]: картечь\\
0xBCA8 \textbf{hjo2t} [\emph{ʰʒo˩t}]: мастоидный\\
0xBCAA \textbf{tco2t} [\emph{tʃo˩t}]: ротонда\\
0xBCB0 \textbf{hvo2t} [\emph{ʰvo˩t}]: преданно\\
0xBCC1 \textbf{mro2t} [\emph{mro˩t}]: таращить глаза\\
0xBCC2 \textbf{kro2t} [\emph{kro˩t}]: приусадебный участок\\
0xBCC4 \textbf{pro2t} [\emph{pro˩t}]: оперетта\\
0xBCC6 \textbf{nro2t} [\emph{nro˩t}]: контральто\\
0xBCC8 \textbf{hqo2t} [\emph{ʰŋo˩t}]: сопли\\
0xBCC9 \textbf{sro2t} [\emph{sro˩t}]: иллюстратор\\
0xBCCA \textbf{tro2t} [\emph{tro˩t}]: Трой\\
0xBCCC \textbf{fro2t} [\emph{fro˩t}]: членистоногие\\
0xBCCD \textbf{cro2t} [\emph{ʃro˩t}]: короткие разрезы\\
0xBCD0 \textbf{hxo2t} [\emph{ʰxo˩t}]: Геродот\\
0xBCD2 \textbf{gro2t} [\emph{gro˩t}]: гигрометр\\
0xBCD3 \textbf{dro2t} [\emph{dro˩t}]: низкооплачиваемый\\
0xBCDA \textbf{xro2t} [\emph{xro˩t}]: ссадить\\
0xBCE2 \textbf{kxo2t} [\emph{kxo˩t}]: волоске\\
0xBCE4 \textbf{pxo2t} [\emph{pxo˩t}]: периодичность\\
0xBCEA \textbf{txo2t} [\emph{txo˩t}]: право аренды\\
0xBCF3 \textbf{dxo2t} [\emph{dxo˩t}]: недоделанный\\
0xBD08 \textbf{hm62t} [\emph{ʰmə˩t}]: нумизмат\\
0xBD10 \textbf{hk62t} [\emph{ʰkə˩t}]: лактон\\
0xBD20 \textbf{hp62t} [\emph{ʰpə˩t}]: ипохондрия\\
0xBD30 \textbf{hn62t} [\emph{ʰnə˩t}]: синкопированная\\
0xBD48 \textbf{hs62t} [\emph{ʰsə˩t}]: грамотный\\
0xBD4A \textbf{ts62t} [\emph{tsə˩t}]: семьдесят семь\\
0xBD50 \textbf{ht62t} [\emph{ʰtə˩t}]: семьдесят два\\
0xBD58 \textbf{hl62t} [\emph{ʰlə˩t}]: психиатр\\
0xBD69 \textbf{sl62t} [\emph{slə˩t}]: совершеннолетие\\
0xBD70 \textbf{hr62t} [\emph{ʰrə˩t}]: геральдика\\
0xBDC2 \textbf{kr62t} [\emph{krə˩t}]: сросшийся\\
0xBDCA \textbf{tr62t} [\emph{trə˩t}]: Автопортрет\\
0xBDEA \textbf{tx62t} [\emph{txə˩t}]: держатель билетов\\
\subsection{ C }
0xC227 \textbf{psi2h} [\emph{psi˩ʰ}]: пластика\\
0xC257 \textbf{tsi2h} [\emph{tsi˩ʰ}]: послелоги дело\\
0xC801 \textbf{myif} [\emph{mjif}]: рыбак\\
0xC802 \textbf{kyif} [\emph{kjif}]: аппликативны голос\\
0xC804 \textbf{pyif} [\emph{pjif}]: неопределенные\\
0xC806 \textbf{nyif} [\emph{njif}]: коннотация\\
0xC808 \textbf{hmif} [\emph{ʰmif}]: мили\\
0xC809 \textbf{syif} [\emph{sjif}]: си единицы\\
0xC80A \textbf{tyif} [\emph{tjif}]: интерпретировать\\
0xC80B \textbf{lyif} [\emph{ljif}]: Литва\\
0xC80C \textbf{fyif} [\emph{fjif}]: возвратный голос\\
0xC80D \textbf{cyif} [\emph{ʃjif}]: лук\\
0xC80E \textbf{ryif} [\emph{rjif}]: радиальный\\
0xC810 \textbf{hkif} [\emph{ʰkif}]: благородный\\
0xC811 \textbf{byif} [\emph{bjif}]: банкноты\\
0xC812 \textbf{gyif} [\emph{gjif}]: Гвинея\\
0xC813 \textbf{dyif} [\emph{djif}]: не доставлен\\
0xC818 \textbf{hyif} [\emph{ʰjif}]: вперед\\
0xC81A \textbf{xyif} [\emph{xjif}]: океанография\\
0xC820 \textbf{hpif} [\emph{ʰpif}]: страница\\
0xC821 \textbf{mwif} [\emph{mwif}]: Малави\\
0xC822 \textbf{kwif} [\emph{kwif}]: киловатт\\
0xC824 \textbf{pwif} [\emph{pwif}]: предварительный просмотр\\
0xC826 \textbf{nwif} [\emph{nwif}]: не живой\\
0xC828 \textbf{hwif} [\emph{ʰwif}]: вице\\
0xC829 \textbf{swif} [\emph{swif}]: Швейцария\\
0xC82A \textbf{twif} [\emph{twif}]: данник\\
0xC82B \textbf{lwif} [\emph{lwif}]: местный падеж\\
0xC82C \textbf{fwif} [\emph{fwif}]: FY\\
0xC82E \textbf{rwif} [\emph{rwif}]: сорок четыре\\
0xC830 \textbf{hnif} [\emph{ʰnif}]: усилитель\\
0xC832 \textbf{gwif} [\emph{gwif}]: глифы\\
0xC833 \textbf{dwif} [\emph{dwif}]: деревянных духовых инструментов\\
0xC83A \textbf{xwif} [\emph{xwif}]: встречный ветер\\
0xC842 \textbf{ksif} [\emph{ksif}]: окислять\\
0xC844 \textbf{psif} [\emph{psif}]: эмиссия\\
0xC848 \textbf{hsif} [\emph{ʰsif}]: лев\\
0xC84A \textbf{tsif} [\emph{tsif}]: стимул\\
0xC850 \textbf{htif} [\emph{ʰtif}]: расстояние\\
0xC851 \textbf{bzif} [\emph{bzif}]: урбанизация\\
0xC852 \textbf{gzif} [\emph{gzif}]: географический справочник\\
0xC858 \textbf{hlif} [\emph{ʰlif}]: прибыль\\
0xC860 \textbf{hfif} [\emph{ʰfif}]: философ\\
0xC862 \textbf{klif} [\emph{klif}]: кальций\\
0xC864 \textbf{plif} [\emph{plif}]: вилла\\
0xC868 \textbf{hcif} [\emph{ʰʃif}]: кожа\\
0xC869 \textbf{slif} [\emph{slif}]: сито\\
0xC86A \textbf{tlif} [\emph{tlif}]: кафельная плитка\\
0xC86C \textbf{flif} [\emph{flif}]: филе\\
0xC86D \textbf{clif} [\emph{ʃlif}]: спираль\\
0xC870 \textbf{hrif} [\emph{ʰrif}]: примитивный\\
0xC871 \textbf{blif} [\emph{blif}]: Болгария\\
0xC872 \textbf{glif} [\emph{glif}]: греческий\\
0xC873 \textbf{dlif} [\emph{dlif}]: дельфин\\
0xC874 \textbf{zlif} [\emph{zlif}]: ксилофон\\
0xC876 \textbf{vlif} [\emph{vlif}]: непроизвольный\\
0xC87A \textbf{xlif} [\emph{xlif}]: ставки\\
0xC882 \textbf{kfif} [\emph{kfif}]: ядро благословляющий\\
0xC884 \textbf{pfif} [\emph{pfif}]: Спецификация\\
0xC888 \textbf{hbif} [\emph{ʰbif}]: лед\\
0xC88A \textbf{tfif} [\emph{tfif}]: переходный глагол\\
0xC890 \textbf{hgif} [\emph{ʰgif}]: очаровывать\\
0xC891 \textbf{bvif} [\emph{bvif}]: двояковыпуклый\\
0xC893 \textbf{dvif} [\emph{dvif}]: скат\\
0xC898 \textbf{hdif} [\emph{ʰdif}]: идентичность\\
0xC8A0 \textbf{hzif} [\emph{ʰzif}]: зефир\\
0xC8A2 \textbf{kcif} [\emph{kʃif}]: дефицит\\
0xC8A4 \textbf{pcif} [\emph{pʃif}]: холодильник\\
0xC8A8 \textbf{hjif} [\emph{ʰʒif}]: незаконное\\
0xC8AA \textbf{tcif} [\emph{tʃif}]: заостренный\\
0xC8B0 \textbf{hvif} [\emph{ʰvif}]: глагол - суффиксов\\
0xC8C1 \textbf{mrif} [\emph{mrif}]: микрофон\\
0xC8C2 \textbf{krif} [\emph{krif}]: секреция\\
0xC8C4 \textbf{prif} [\emph{prif}]: переменная\\
0xC8C6 \textbf{nrif} [\emph{nrif}]: клев\\
0xC8C8 \textbf{hqif} [\emph{ʰŋif}]: Венгрия\\
0xC8C9 \textbf{srif} [\emph{srif}]: исторический напряженный\\
0xC8CA \textbf{trif} [\emph{trif}]: география\\
0xC8CC \textbf{frif} [\emph{frif}]: февраль\\
0xC8CD \textbf{crif} [\emph{ʃrif}]: застенчивость\\
0xC8D0 \textbf{hxif} [\emph{ʰxif}]: гипнотизировать\\
0xC8D1 \textbf{brif} [\emph{brif}]: двоичный значение\\
0xC8D2 \textbf{grif} [\emph{grif}]: партизанский\\
0xC8D3 \textbf{drif} [\emph{drif}]: подводное крыло\\
0xC8D4 \textbf{zrif} [\emph{zrif}]: хлорофилл\\
0xC8D5 \textbf{jrif} [\emph{ʒrif}]: выводок\\
0xC8D6 \textbf{vrif} [\emph{vrif}]: заднего вида\\
0xC8D9 \textbf{qrif} [\emph{ŋrif}]: рыба-меч\\
0xC8DA \textbf{xrif} [\emph{xrif}]: жирафа\\
0xC8E2 \textbf{kxif} [\emph{kxif}]: айва\\
0xC8E4 \textbf{pxif} [\emph{pxif}]: полиция сотрудник\\
0xC8EA \textbf{txif} [\emph{txif}]: ссылаясь\\
0xC8F1 \textbf{bxif} [\emph{bxif}]: углевод\\
0xC8F2 \textbf{gxif} [\emph{gxif}]: графит\\
0xC8F3 \textbf{dxif} [\emph{dxif}]: стрекоза\\
0xC901 \textbf{myaf} [\emph{mjäf}]: Мальдивы\\
0xC902 \textbf{kyaf} [\emph{kjäf}]: спасибо\\
0xC904 \textbf{pyaf} [\emph{pjäf}]: птица\\
0xC906 \textbf{nyaf} [\emph{njäf}]: синтаксический\\
0xC908 \textbf{hmaf} [\emph{ʰmäf}]: беда\\
0xC909 \textbf{syaf} [\emph{sjäf}]: стороны\\
0xC90A \textbf{tyaf} [\emph{tjäf}]: параграф\\
0xC90B \textbf{lyaf} [\emph{ljäf}]: алименты\\
0xC90D \textbf{cyaf} [\emph{ʃjäf}]: смех\\
0xC90E \textbf{ryaf} [\emph{rjäf}]: креативность\\
0xC910 \textbf{hkaf} [\emph{ʰkäf}]: официальный\\
0xC911 \textbf{byaf} [\emph{bjäf}]: риф\\
0xC912 \textbf{gyaf} [\emph{gjäf}]: двор фермы\\
0xC913 \textbf{dyaf} [\emph{djäf}]: отдаленное место\\
0xC914 \textbf{zyaf} [\emph{zjäf}]: затяжные прыжки с парашютом\\
0xC918 \textbf{hyaf} [\emph{ʰjäf}]: возможно\\
0xC919 \textbf{qyaf} [\emph{ŋjäf}]: коммерциализация\\
0xC91A \textbf{xyaf} [\emph{xjäf}]: пыльник\\
0xC920 \textbf{hpaf} [\emph{ʰpäf}]: частный\\
0xC921 \textbf{mwaf} [\emph{mwäf}]: вредоносных программ\\
0xC922 \textbf{kwaf} [\emph{kwäf}]: развивать\\
0xC924 \textbf{pwaf} [\emph{pwäf}]: фисташковый\\
0xC926 \textbf{nwaf} [\emph{nwäf}]: пропускная способность\\
0xC928 \textbf{hwaf} [\emph{ʰwäf}]: идеально\\
0xC929 \textbf{swaf} [\emph{swäf}]: ассоциативный дело\\
0xC92A \textbf{twaf} [\emph{twäf}]: тривиальный\\
0xC92B \textbf{lwaf} [\emph{lwäf}]: вспышка водить машину\\
0xC92C \textbf{fwaf} [\emph{fwäf}]: плодоносить\\
0xC92D \textbf{cwaf} [\emph{ʃwäf}]: конверт\\
0xC92E \textbf{rwaf} [\emph{rwäf}]: хризантема\\
0xC930 \textbf{hnaf} [\emph{ʰnäf}]: предположим,\\
0xC931 \textbf{bwaf} [\emph{bwäf}]: будуар\\
0xC933 \textbf{dwaf} [\emph{dwäf}]: водолаз\\
0xC936 \textbf{vwaf} [\emph{vwäf}]: ноу-хау\\
0xC93A \textbf{xwaf} [\emph{xwäf}]: водоносный горизонт\\
0xC942 \textbf{ksaf} [\emph{ksäf}]: фальсифицировать\\
0xC944 \textbf{psaf} [\emph{psäf}]: отсутствие\\
0xC948 \textbf{hsaf} [\emph{ʰsäf}]: дыхание\\
0xC94A \textbf{tsaf} [\emph{tsäf}]: самоубийство\\
0xC950 \textbf{htaf} [\emph{ʰtäf}]: прочность\\
0xC951 \textbf{bzaf} [\emph{bzäf}]: бойня\\
0xC952 \textbf{gzaf} [\emph{gzäf}]: газифицировать\\
0xC958 \textbf{hlaf} [\emph{ʰläf}]: оставаться\\
0xC960 \textbf{hfaf} [\emph{ʰfäf}]: закон\\
0xC962 \textbf{klaf} [\emph{kläf}]: качественный\\
0xC964 \textbf{plaf} [\emph{pläf}]: пыльца\\
0xC968 \textbf{hcaf} [\emph{ʰʃäf}]: сомнение\\
0xC969 \textbf{slaf} [\emph{släf}]: Социальное сетей\\
0xC96A \textbf{tlaf} [\emph{tläf}]: принимать участие\\
0xC96C \textbf{flaf} [\emph{fläf}]: альфа\\
0xC96D \textbf{claf} [\emph{ʃläf}]: полка\\
0xC970 \textbf{hraf} [\emph{ʰräf}]: трафик\\
0xC971 \textbf{blaf} [\emph{bläf}]: изнасилование\\
0xC972 \textbf{glaf} [\emph{gläf}]: Югославия\\
0xC973 \textbf{dlaf} [\emph{dläf}]: гранатовый\\
0xC975 \textbf{jlaf} [\emph{ʒläf}]: золь Ф.А.\\
0xC976 \textbf{vlaf} [\emph{vläf}]: пис`ать стихи\\
0xC97A \textbf{xlaf} [\emph{xläf}]: отмель\\
0xC982 \textbf{kfaf} [\emph{kfäf}]: прививка\\
0xC984 \textbf{pfaf} [\emph{pfäf}]: сапфир\\
0xC988 \textbf{hbaf} [\emph{ʰbäf}]: биография\\
0xC98A \textbf{tfaf} [\emph{tfäf}]: трансцендентный\\
0xC990 \textbf{hgaf} [\emph{ʰgäf}]: грамматика\\
0xC991 \textbf{bvaf} [\emph{bväf}]: землекоп\\
0xC993 \textbf{dvaf} [\emph{dväf}]: диван\\
0xC998 \textbf{hdaf} [\emph{ʰdäf}]: достигать\\
0xC99F \textbf{dwa2h} [\emph{dwä˩ʰ}]: равнодушие\\
0xC9A0 \textbf{hzaf} [\emph{ʰzäf}]: препятствовать\\
0xC9A2 \textbf{kcaf} [\emph{kʃäf}]: шарф\\
0xC9A4 \textbf{pcaf} [\emph{pʃäf}]: круглая скобка\\
0xC9A8 \textbf{hjaf} [\emph{ʰʒäf}]: млекопитающее\\
0xC9AA \textbf{tcaf} [\emph{tʃäf}]: внимательно\\
0xC9B0 \textbf{hvaf} [\emph{ʰväf}]: блондинка\\
0xC9B2 \textbf{gjaf} [\emph{gʒäf}]: кастрировать\\
0xC9C1 \textbf{mraf} [\emph{mräf}]: марафон\\
0xC9C2 \textbf{kraf} [\emph{kräf}]: ворона\\
0xC9C4 \textbf{praf} [\emph{präf}]: пиратство\\
0xC9C6 \textbf{nraf} [\emph{nräf}]: налогоплательщик\\
0xC9C8 \textbf{hqaf} [\emph{ʰŋäf}]: окалина\\
0xC9C9 \textbf{sraf} [\emph{sräf}]: на открытом воздухе\\
0xC9CA \textbf{traf} [\emph{träf}]: полоса\\
0xC9CC \textbf{fraf} [\emph{fräf}]: брандмауэр\\
0xC9CD \textbf{craf} [\emph{ʃräf}]: телесный\\
0xC9D0 \textbf{hxaf} [\emph{ʰxäf}]: выкапывать\\
0xC9D1 \textbf{braf} [\emph{bräf}]: стрижка\\
0xC9D2 \textbf{graf} [\emph{gräf}]: алгоритм\\
0xC9D3 \textbf{draf} [\emph{dräf}]: модернизировать\\
0xC9D4 \textbf{zraf} [\emph{zräf}]: нулевое воздействие\\
0xC9D5 \textbf{jraf} [\emph{ʒräf}]: архитрав\\
0xC9D6 \textbf{vraf} [\emph{vräf}]: виброфон\\
0xC9D9 \textbf{qraf} [\emph{ŋräf}]: неблагодарный\\
0xC9DA \textbf{xraf} [\emph{xräf}]: щавель\\
0xC9E2 \textbf{kxaf} [\emph{kxäf}]: акация\\
0xC9E4 \textbf{pxaf} [\emph{pxäf}]: метаморфоза\\
0xC9EA \textbf{txaf} [\emph{txäf}]: телеграф\\
0xC9F1 \textbf{bxaf} [\emph{bxäf}]: стропило\\
0xC9F2 \textbf{gxaf} [\emph{gxäf}]: давление воздуха\\
0xC9F3 \textbf{dxaf} [\emph{dxäf}]: лесной пожар\\
0xCA01 \textbf{myuf} [\emph{mjuf}]: муфтий\\
0xCA08 \textbf{hmuf} [\emph{ʰmuf}]: рот\\
0xCA0A \textbf{tyuf} [\emph{tjuf}]: перехитрить\\
0xCA0E \textbf{ryuf} [\emph{rjuf}]: европеизация\\
0xCA10 \textbf{hkuf} [\emph{ʰkuf}]: мясник\\
0xCA18 \textbf{hyuf} [\emph{ʰjuf}]: мочеиспускательный канал\\
0xCA20 \textbf{hpuf} [\emph{ʰpuf}]: пикколо\\
0xCA28 \textbf{hwuf} [\emph{ʰwuf}]: вдова\\
0xCA2B \textbf{lwuf} [\emph{lwuf}]: сульфид\\
0xCA30 \textbf{hnuf} [\emph{ʰnuf}]: Венесуэла\\
0xCA42 \textbf{ksuf} [\emph{ksuf}]: кускус\\
0xCA44 \textbf{psuf} [\emph{psuf}]: массив\\
0xCA48 \textbf{hsuf} [\emph{ʰsuf}]: удаление\\
0xCA4A \textbf{tsuf} [\emph{tsuf}]: Тунис\\
0xCA50 \textbf{htuf} [\emph{ʰtuf}]: тунец\\
0xCA58 \textbf{hluf} [\emph{ʰluf}]: старость\\
0xCA60 \textbf{hfuf} [\emph{ʰfuf}]: отруби\\
0xCA62 \textbf{kluf} [\emph{kluf}]: ключ\\
0xCA68 \textbf{hcuf} [\emph{ʰʃuf}]: От себя\\
0xCA69 \textbf{sluf} [\emph{sluf}]: singulative номер\\
0xCA6A \textbf{tluf} [\emph{tluf}]: капелька болтовня\\
0xCA70 \textbf{hruf} [\emph{ʰruf}]: грейпфрут\\
0xCA71 \textbf{bluf} [\emph{bluf}]: пуленепробиваемый\\
0xCA76 \textbf{vluf} [\emph{vluf}]: заспать\\
0xCA7A \textbf{xluf} [\emph{xluf}]: пылесос\\
0xCA82 \textbf{kfuf} [\emph{kfuf}]: крестоцветный\\
0xCA84 \textbf{pfuf} [\emph{pfuf}]: летучие мыши\\
0xCA88 \textbf{hbuf} [\emph{ʰbuf}]: вареный\\
0xCA8A \textbf{tfuf} [\emph{tfuf}]: трюфель\\
0xCA90 \textbf{hguf} [\emph{ʰguf}]: лачуга\\
0xCA98 \textbf{hduf} [\emph{ʰduf}]: дуэль\\
0xCAC2 \textbf{kruf} [\emph{kruf}]: Украина\\
0xCAC4 \textbf{pruf} [\emph{pruf}]: пуриновых\\
0xCAC6 \textbf{nruf} [\emph{nruf}]: отлив прибоя\\
0xCAC9 \textbf{sruf} [\emph{sruf}]: пугало\\
0xCACA \textbf{truf} [\emph{truf}]: страусы\\
0xCAD0 \textbf{hxuf} [\emph{ʰxuf}]: Ганновер\\
0xCAD1 \textbf{bruf} [\emph{bruf}]: блаженство\\
0xCADA \textbf{xruf} [\emph{xruf}]: люк\\
0xCAEA \textbf{txuf} [\emph{txuf}]: пирогов\\
0xCB04 \textbf{pyef} [\emph{pjef}]: перестановка\\
0xCB06 \textbf{nyef} [\emph{njef}]: центробежный\\
0xCB08 \textbf{hmef} [\emph{ʰmef}]: май месяц\\
0xCB09 \textbf{syef} [\emph{sjef}]: Scot бесплатно\\
0xCB0A \textbf{tyef} [\emph{tjef}]: моток пряжи\\
0xCB10 \textbf{hkef} [\emph{ʰkef}]: кафе\\
0xCB11 \textbf{byef} [\emph{bjef}]: бордель\\
0xCB18 \textbf{hyef} [\emph{ʰjef}]: рецессивный\\
0xCB20 \textbf{hpef} [\emph{ʰpef}]: предикативный\\
0xCB24 \textbf{pwef} [\emph{pwef}]: подветренный\\
0xCB28 \textbf{hwef} [\emph{ʰwef}]: неоплаченный\\
0xCB2A \textbf{twef} [\emph{twef}]: ультрафиолетовый\\
0xCB2B \textbf{lwef} [\emph{lwef}]: йеллоунайф\\
0xCB2D \textbf{cwef} [\emph{ʃwef}]: кервель\\
0xCB2E \textbf{rwef} [\emph{rwef}]: оборотни\\
0xCB30 \textbf{hnef} [\emph{ʰnef}]: Блок - масс\\
0xCB36 \textbf{vwef} [\emph{vwef}]: уклончиво\\
0xCB42 \textbf{ksef} [\emph{ksef}]: Рентгеновский\\
0xCB44 \textbf{psef} [\emph{psef}]: эпицентр\\
0xCB48 \textbf{hsef} [\emph{ʰsef}]: змея\\
0xCB4A \textbf{tsef} [\emph{tsef}]: antessive дело\\
0xCB50 \textbf{htef} [\emph{ʰtef}]: портной\\
0xCB58 \textbf{hlef} [\emph{ʰlef}]: агент вызывать\\
0xCB60 \textbf{hfef} [\emph{ʰfef}]: взаимное развитие\\
0xCB62 \textbf{klef} [\emph{klef}]: экстраполяция\\
0xCB64 \textbf{plef} [\emph{plef}]: парапланеризм\\
0xCB68 \textbf{hcef} [\emph{ʰʃef}]: конечный глагол\\
0xCB69 \textbf{slef} [\emph{slef}]: шифер\\
0xCB6A \textbf{tlef} [\emph{tlef}]: прицеп\\
0xCB6C \textbf{flef} [\emph{flef}]: Филадельфия\\
0xCB6D \textbf{clef} [\emph{ʃlef}]: кресельный подъемник\\
0xCB70 \textbf{href} [\emph{ʰref}]: разреженный\\
0xCB71 \textbf{blef} [\emph{blef}]: книголюб\\
0xCB72 \textbf{glef} [\emph{glef}]: мегалит\\
0xCB73 \textbf{dlef} [\emph{dlef}]: дельфийский\\
0xCB76 \textbf{vlef} [\emph{vlef}]: виолончель\\
0xCB7A \textbf{xlef} [\emph{xlef}]: сеновал\\
0xCB82 \textbf{kfef} [\emph{kfef}]: кинестетическая\\
0xCB84 \textbf{pfef} [\emph{pfef}]: упреждающий\\
0xCB88 \textbf{hbef} [\emph{ʰbef}]: зимняя спячка\\
0xCB8A \textbf{tfef} [\emph{tfef}]: заземление\\
0xCB90 \textbf{hgef} [\emph{ʰgef}]: обхват\\
0xCB98 \textbf{hdef} [\emph{ʰdef}]: дебет\\
0xCBA2 \textbf{kcef} [\emph{kʃef}]: равноденственный\\
0xCBA4 \textbf{pcef} [\emph{pʃef}]: рептилии\\
0xCBA8 \textbf{hjef} [\emph{ʰʒef}]: Женева\\
0xCBAA \textbf{tcef} [\emph{tʃef}]: вишня\\
0xCBB0 \textbf{hvef} [\emph{ʰvef}]: червеобразный\\
0xCBC1 \textbf{mref} [\emph{mref}]: монорельс\\
0xCBC2 \textbf{kref} [\emph{kref}]: тмин\\
0xCBC4 \textbf{pref} [\emph{pref}]: партер\\
0xCBC6 \textbf{nref} [\emph{nref}]: реинкарнация\\
0xCBC8 \textbf{hqef} [\emph{ʰŋef}]: теория относительности\\
0xCBC9 \textbf{sref} [\emph{sref}]: стенографистка\\
0xCBCA \textbf{tref} [\emph{tref}]: скипидар\\
0xCBCD \textbf{cref} [\emph{ʃref}]: шериф\\
0xCBD0 \textbf{hxef} [\emph{ʰxef}]: вереск\\
0xCBD1 \textbf{bref} [\emph{bref}]: покрывало\\
0xCBD3 \textbf{dref} [\emph{dref}]: гидросфера\\
0xCBDA \textbf{xref} [\emph{xref}]: мелодичный\\
0xCBE2 \textbf{kxef} [\emph{kxef}]: коленчатый вал\\
0xCBE4 \textbf{pxef} [\emph{pxef}]: проливной\\
0xCBEA \textbf{txef} [\emph{txef}]: термосфера\\
0xCBF1 \textbf{bxef} [\emph{bxef}]: Босния герцеговина\\
0xCBF3 \textbf{dxef} [\emph{dxef}]: принаряжать\\
0xCC01 \textbf{myof} [\emph{mjof}]: дезинформировать\\
0xCC06 \textbf{nyof} [\emph{njof}]: ионосфера\\
0xCC08 \textbf{hmof} [\emph{ʰmof}]: торговый центр\\
0xCC09 \textbf{syof} [\emph{sjof}]: Исполнительный директор\\
0xCC0A \textbf{tyof} [\emph{tjof}]: тасование\\
0xCC0E \textbf{ryof} [\emph{rjof}]: иволга\\
0xCC10 \textbf{hkof} [\emph{ʰkof}]: потребление\\
0xCC18 \textbf{hyof} [\emph{ʰjof}]: выгодный\\
0xCC1A \textbf{xyof} [\emph{xjof}]: травоядное\\
0xCC20 \textbf{hpof} [\emph{ʰpof}]: Попкорн\\
0xCC21 \textbf{mwof} [\emph{mwof}]: Молдова\\
0xCC22 \textbf{kwof} [\emph{kwof}]: Эквадор\\
0xCC26 \textbf{nwof} [\emph{nwof}]: фонограф\\
0xCC28 \textbf{hwof} [\emph{ʰwof}]: ива\\
0xCC29 \textbf{swof} [\emph{swof}]: Словения\\
0xCC2A \textbf{twof} [\emph{twof}]: Ботсвана\\
0xCC2B \textbf{lwof} [\emph{lwof}]: Латвия\\
0xCC2E \textbf{rwof} [\emph{rwof}]: сток\\
0xCC30 \textbf{hnof} [\emph{ʰnof}]: курносый\\
0xCC42 \textbf{ksof} [\emph{ksof}]: саксофон\\
0xCC44 \textbf{psof} [\emph{psof}]: фосфат\\
0xCC48 \textbf{hsof} [\emph{ʰsof}]: прибой\\
0xCC4A \textbf{tsof} [\emph{tsof}]: antipassive голос\\
0xCC50 \textbf{htof} [\emph{ʰtof}]: вор\\
0xCC58 \textbf{hlof} [\emph{ʰlof}]: сера\\
0xCC60 \textbf{hfof} [\emph{ʰfof}]: фосфор\\
0xCC62 \textbf{klof} [\emph{klof}]: клевер\\
0xCC64 \textbf{plof} [\emph{plof}]: силовая установка\\
0xCC68 \textbf{hcof} [\emph{ʰʃof}]: раздевать\\
0xCC69 \textbf{slof} [\emph{slof}]: морское дно\\
0xCC6A \textbf{tlof} [\emph{tlof}]: телеологический\\
0xCC6D \textbf{clof} [\emph{ʃlof}]: Чехословакией\\
0xCC70 \textbf{hrof} [\emph{ʰrof}]: троянец лошадь\\
0xCC71 \textbf{blof} [\emph{blof}]: Боливия\\
0xCC72 \textbf{glof} [\emph{glof}]: гольф\\
0xCC76 \textbf{vlof} [\emph{vlof}]: револьвер\\
0xCC7A \textbf{xlof} [\emph{xlof}]: оптовая\\
0xCC82 \textbf{kfof} [\emph{kfof}]: манжета\\
0xCC84 \textbf{pfof} [\emph{pfof}]: порфир\\
0xCC88 \textbf{hbof} [\emph{ʰbof}]: бомбардировка\\
0xCC8A \textbf{tfof} [\emph{tfof}]: автоматизация\\
0xCC90 \textbf{hgof} [\emph{ʰgof}]: почвопокровные\\
0xCC98 \textbf{hdof} [\emph{ʰdof}]: идиофон\\
0xCCA2 \textbf{kcof} [\emph{kʃof}]: ремесла\\
0xCCA4 \textbf{pcof} [\emph{pʃof}]: порнографический\\
0xCCA8 \textbf{hjof} [\emph{ʰʒof}]: припарка\\
0xCCAA \textbf{tcof} [\emph{tʃof}]: стрельба из лука\\
0xCCC1 \textbf{mrof} [\emph{mrof}]: молярный\\
0xCCC2 \textbf{krof} [\emph{krof}]: Корсика\\
0xCCC4 \textbf{prof} [\emph{prof}]: уничижительный\\
0xCCC6 \textbf{nrof} [\emph{nrof}]: Норвегия\\
0xCCC9 \textbf{srof} [\emph{srof}]: осмос\\
0xCCCA \textbf{trof} [\emph{trof}]: непереходный\\
0xCCCC \textbf{frof} [\emph{frof}]: литография\\
0xCCCD \textbf{crof} [\emph{ʃrof}]: испольщик\\
0xCCD0 \textbf{hxof} [\emph{ʰxof}]: герцеговина\\
0xCCD1 \textbf{brof} [\emph{brof}]: обручить\\
0xCCD3 \textbf{drof} [\emph{drof}]: задергать\\
0xCCD6 \textbf{vrof} [\emph{vrof}]: евро\\
0xCCDA \textbf{xrof} [\emph{xrof}]: корабль на воздушной подушке\\
0xCCE2 \textbf{kxof} [\emph{kxof}]: гравюра на дереве\\
0xCCE4 \textbf{pxof} [\emph{pxof}]: полиморфный\\
0xCCEA \textbf{txof} [\emph{txof}]: фаршированный\\
0xCCF3 \textbf{dxof} [\emph{dxof}]: Иегова\\
0xCD08 \textbf{hm6f} [\emph{ʰməf}]: запечатывание\\
0xCD10 \textbf{hk6f} [\emph{ʰkəf}]: хвойное дерево\\
0xCD20 \textbf{hp6f} [\emph{ʰpəf}]: ржанка\\
0xCD28 \textbf{hw6f} [\emph{ʰwəf}]: вафельный\\
0xCD30 \textbf{hn6f} [\emph{ʰnəf}]: necessitative настроение\\
0xCD48 \textbf{hs6f} [\emph{ʰsəf}]: синтезировать\\
0xCD4A \textbf{ts6f} [\emph{tsəf}]: стратосфера\\
0xCD50 \textbf{ht6f} [\emph{ʰtəf}]: Roto румпель\\
0xCD58 \textbf{hl6f} [\emph{ʰləf}]: литосфера\\
0xCD60 \textbf{hf6f} [\emph{ʰfəf}]: сосиски\\
0xCD62 \textbf{kl6f} [\emph{kləf}]: Голгофа\\
0xCD68 \textbf{hc6f} [\emph{ʰʃəf}]: поверять\\
0xCD6A \textbf{tl6f} [\emph{tləf}]: кристаллография\\
0xCD70 \textbf{hr6f} [\emph{ʰrəf}]: тропосфера\\
0xCDC2 \textbf{kr6f} [\emph{krəf}]: Каринтии\\
0xCDC6 \textbf{nr6f} [\emph{nrəf}]: Назарет\\
0xCDC9 \textbf{sr6f} [\emph{srəf}]: выживаемость\\
0xCDCA \textbf{tr6f} [\emph{trəf}]: топография\\
0xCDD0 \textbf{hx6f} [\emph{ʰxəf}]: агиография\\
\subsection{ D }
0xD004 \textbf{pyi7f} [\emph{pji˥f}]: пирогов глазами\\
0xD008 \textbf{hmi7f} [\emph{ʰmi˥f}]: метил\\
0xD009 \textbf{syi7f} [\emph{sji˥f}]: София\\
0xD00B \textbf{lyi7f} [\emph{lji˥f}]: лилия трусливый\\
0xD010 \textbf{hki7f} [\emph{ʰki˥f}]: треска\\
0xD020 \textbf{hpi7f} [\emph{ʰpi˥f}]: скоропортящийся\\
0xD026 \textbf{nwi7f} [\emph{nwi˥f}]: анчоус\\
0xD028 \textbf{hwi7f} [\emph{ʰwi˥f}]: волдырь\\
0xD030 \textbf{hni7f} [\emph{ʰni˥f}]: чистописание\\
0xD048 \textbf{hsi7f} [\emph{ʰsi˥f}]: непродуктивный\\
0xD04A \textbf{tsi7f} [\emph{tsi˥f}]: смокинг\\
0xD050 \textbf{hti7f} [\emph{ʰti˥f}]: Тимоти\\
0xD058 \textbf{hli7f} [\emph{ʰli˥f}]: котиковый\\
0xD060 \textbf{hfi7f} [\emph{ʰfi˥f}]: тиофен\\
0xD062 \textbf{kli7f} [\emph{kli˥f}]: халиф\\
0xD064 \textbf{pli7f} [\emph{pli˥f}]: полиомиелит\\
0xD069 \textbf{sli7f} [\emph{sli˥f}]: пан славянство\\
0xD070 \textbf{hri7f} [\emph{ʰri˥f}]: грифон\\
0xD08A \textbf{tfi7f} [\emph{tfi˥f}]: атрофированный\\
0xD0B0 \textbf{hvi7f} [\emph{ʰvi˥f}]: водевиль\\
0xD0C1 \textbf{mri7f} [\emph{mri˥f}]: морфий\\
0xD0C2 \textbf{kri7f} [\emph{kri˥f}]: выделять курсивом\\
0xD0C8 \textbf{hqi7f} [\emph{ʰŋi˥f}]: орфографический\\
0xD0C9 \textbf{sri7f} [\emph{sri˥f}]: стенография\\
0xD0CA \textbf{tri7f} [\emph{tri˥f}]: канату\\
0xD0D0 \textbf{hxi7f} [\emph{ʰxi˥f}]: верный огонь\\
0xD101 \textbf{mya7f} [\emph{mjä˥f}]: майя\\
0xD106 \textbf{nya7f} [\emph{njä˥f}]: Племянница в законе\\
0xD108 \textbf{hma7f} [\emph{ʰmä˥f}]: дог\\
0xD110 \textbf{hka7f} [\emph{ʰkä˥f}]: скандинавский\\
0xD118 \textbf{hya7f} [\emph{ʰjä˥f}]: один путь\\
0xD120 \textbf{hpa7f} [\emph{ʰpä˥f}]: поташ\\
0xD128 \textbf{hwa7f} [\emph{ʰwä˥f}]: письмо совершенным\\
0xD130 \textbf{hna7f} [\emph{ʰnä˥f}]: неф\\
0xD142 \textbf{ksa7f} [\emph{ksä˥f}]: растр\\
0xD144 \textbf{psa7f} [\emph{psä˥f}]: апостроф\\
0xD148 \textbf{hsa7f} [\emph{ʰsä˥f}]: прадед\\
0xD14A \textbf{tsa7f} [\emph{tsä˥f}]: самореференция\\
0xD150 \textbf{hta7f} [\emph{ʰtä˥f}]: пустяки\\
0xD158 \textbf{hla7f} [\emph{ʰlä˥f}]: планка\\
0xD160 \textbf{hfa7f} [\emph{ʰfä˥f}]: пятерка\\
0xD162 \textbf{kla7f} [\emph{klä˥f}]: каллиграф\\
0xD168 \textbf{hca7f} [\emph{ʰʃä˥f}]: Шанхай\\
0xD169 \textbf{sla7f} [\emph{slä˥f}]: славянский\\
0xD16A \textbf{tla7f} [\emph{tlä˥f}]: ацетальдегид\\
0xD170 \textbf{hra7f} [\emph{ʰrä˥f}]: гуртовщик\\
0xD184 \textbf{pfa7f} [\emph{pfä˥f}]: нафтол\\
0xD188 \textbf{hba7f} [\emph{ʰbä˥f}]: патронташ\\
0xD190 \textbf{hga7f} [\emph{ʰgä˥f}]: графия\\
0xD198 \textbf{hda7f} [\emph{ʰdä˥f}]: два четыре\\
0xD1AA \textbf{tca7f} [\emph{tʃä˥f}]: жестянщик\\
0xD1C1 \textbf{mra7f} [\emph{mrä˥f}]: тамаринд\\
0xD1C2 \textbf{kra7f} [\emph{krä˥f}]: выбивать\\
0xD1C9 \textbf{sra7f} [\emph{srä˥f}]: разносить\\
0xD1CA \textbf{tra7f} [\emph{trä˥f}]: вице-король\\
0xD1D0 \textbf{hxa7f} [\emph{ʰxä˥f}]: общее право\\
0xD1E4 \textbf{pxa7f} [\emph{pxä˥f}]: пентхауз\\
0xD1EA \textbf{txa7f} [\emph{txä˥f}]: заполнитель\\
0xD250 \textbf{htu7f} [\emph{ʰtu˥f}]: тьфу\\
0xD268 \textbf{hcu7f} [\emph{ʰʃu˥f}]: сёгунат\\
0xD308 \textbf{hme7f} [\emph{ʰme˥f}]: ментол\\
0xD348 \textbf{hse7f} [\emph{ʰse˥f}]: лексикограф\\
0xD34A \textbf{tse7f} [\emph{tse˥f}]: сохранение времени\\
0xD350 \textbf{hte7f} [\emph{ʰte˥f}]: лейтмотив\\
0xD370 \textbf{hre7f} [\emph{ʰre˥f}]: выпрямитель\\
0xD388 \textbf{hbe7f} [\emph{ʰbe˥f}]: хлебное дерево\\
0xD410 \textbf{hko7f} [\emph{ʰko˥f}]: бухточка\\
0xD428 \textbf{hwo7f} [\emph{ʰwo˥f}]: Волоф\\
0xD430 \textbf{hno7f} [\emph{ʰno˥f}]: девятикратный\\
0xD442 \textbf{kso7f} [\emph{kso˥f}]: Косово\\
0xD448 \textbf{hso7f} [\emph{ʰso˥f}]: Самоанский\\
0xD450 \textbf{hto7f} [\emph{ʰto˥f}]: легкая победа\\
0xD468 \textbf{hco7f} [\emph{ʰʃo˥f}]: условное депонирование\\
0xD470 \textbf{hro7f} [\emph{ʰro˥f}]: цапфа\\
0xD48A \textbf{tfo7f} [\emph{tfo˥f}]: UTF-8\\
0xD808 \textbf{hmi2f} [\emph{ʰmi˩f}]: мистифицирован\\
0xD809 \textbf{syi2f} [\emph{sji˩f}]: седиль\\
0xD810 \textbf{hki2f} [\emph{ʰki˩f}]: Киев\\
0xD820 \textbf{hpi2f} [\emph{ʰpi˩f}]: отцеубийство\\
0xD830 \textbf{hni2f} [\emph{ʰni˩f}]: девяностый\\
0xD848 \textbf{hsi2f} [\emph{ʰsi˩f}]: севилья\\
0xD84A \textbf{tsi2f} [\emph{tsi˩f}]: восьмидесятый\\
0xD850 \textbf{hti2f} [\emph{ʰti˩f}]: сводить на нет\\
0xD858 \textbf{hli2f} [\emph{ʰli˩f}]: серебряный сладкоречивый\\
0xD860 \textbf{hfi2f} [\emph{ʰfi˩f}]: фосфид\\
0xD862 \textbf{kli2f} [\emph{kli˩f}]: факультативный\\
0xD864 \textbf{pli2f} [\emph{pli˩f}]: мостильщик\\
0xD870 \textbf{hri2f} [\emph{ʰri˩f}]: рекурсивный\\
0xD87A \textbf{xli2f} [\emph{xli˩f}]: воздушные перевозки\\
0xD8C4 \textbf{pri2f} [\emph{pri˩f}]: гипертрофия\\
0xD908 \textbf{hma2f} [\emph{ʰmä˩f}]: мафия\\
0xD910 \textbf{hka2f} [\emph{ʰkä˩f}]: октава\\
0xD929 \textbf{swa2f} [\emph{swä˩f}]: оссифицировать\\
0xD930 \textbf{hna2f} [\emph{ʰnä˩f}]: жених\\
0xD948 \textbf{hsa2f} [\emph{ʰsä˩f}]: семафор\\
0xD94A \textbf{tsa2f} [\emph{tsä˩f}]: марсель\\
0xD950 \textbf{hta2f} [\emph{ʰtä˩f}]: дрейфовать\\
0xD958 \textbf{hla2f} [\emph{ʰlä˩f}]: слесарь\\
0xD969 \textbf{sla2f} [\emph{slä˩f}]: окисляться\\
0xD96A \textbf{tla2f} [\emph{tlä˩f}]: прибавлять\\
0xD970 \textbf{hra2f} [\emph{ʰrä˩f}]: архивариус\\
0xD982 \textbf{kfa2f} [\emph{kfä˩f}]: квантификатором\\
0xD988 \textbf{hba2f} [\emph{ʰbä˩f}]: байпас\\
0xD998 \textbf{hda2f} [\emph{ʰdä˩f}]: крестный отец\\
0xD9AA \textbf{tca2f} [\emph{tʃä˩f}]: каменная стена\\
0xD9C1 \textbf{mra2f} [\emph{mrä˩f}]: маратхи\\
0xD9C6 \textbf{nra2f} [\emph{nrä˩f}]: антропософия\\
0xD9C9 \textbf{sra2f} [\emph{srä˩f}]: омыляться\\
0xD9CA \textbf{tra2f} [\emph{trä˩f}]: картография\\
0xD9D0 \textbf{hxa2f} [\emph{ʰxä˩f}]: сверстник\\
0xD9EA \textbf{txa2f} [\emph{txä˩f}]: перетягивание каната\\
0xDA48 \textbf{hsu2f} [\emph{ʰsu˩f}]: беспутный\\
0xDB30 \textbf{hne2f} [\emph{ʰne˩f}]: анклав\\
0xDC30 \textbf{hno2f} [\emph{ʰno˩f}]: и более\\
0xDC70 \textbf{hro2f} [\emph{ʰro˩f}]: хромосфера\\
\subsection{ E }
0xE06F \textbf{cyo2h} [\emph{ʃjo˩ʰ}]: подозрительный\\
0xE801 \textbf{myic} [\emph{mjiʃ}]: дождь\\
0xE802 \textbf{kyic} [\emph{kjiʃ}]: церковь\\
0xE804 \textbf{pyic} [\emph{pjiʃ}]: обязательный\\
0xE806 \textbf{nyic} [\emph{njiʃ}]: одетый\\
0xE808 \textbf{hmic} [\emph{ʰmiʃ}]: милая\\
0xE809 \textbf{syic} [\emph{sjiʃ}]: грустный\\
0xE80A \textbf{tyic} [\emph{tjiʃ}]: природный\\
0xE80B \textbf{lyic} [\emph{ljiʃ}]: пожалуйста\\
0xE80C \textbf{fyic} [\emph{fjiʃ}]: коэффициент полезного действия\\
0xE80D \textbf{cyic} [\emph{ʃjiʃ}]: тщеславный\\
0xE80E \textbf{ryic} [\emph{rjiʃ}]: сельскохозяйственное\\
0xE810 \textbf{hkic} [\emph{ʰkiʃ}]: между\\
0xE811 \textbf{byic} [\emph{bjiʃ}]: бактерии\\
0xE812 \textbf{gyic} [\emph{gjiʃ}]: гитара\\
0xE813 \textbf{dyic} [\emph{djiʃ}]: образовательный\\
0xE814 \textbf{zyic} [\emph{zjiʃ}]: крайний срок\\
0xE815 \textbf{jyic} [\emph{ʒjiʃ}]: заложник\\
0xE816 \textbf{vyic} [\emph{vjiʃ}]: бдение\\
0xE818 \textbf{hyic} [\emph{ʰjiʃ}]: увидел\\
0xE819 \textbf{qyic} [\emph{ŋjiʃ}]: психиатрия\\
0xE81A \textbf{xyic} [\emph{xjiʃ}]: миссионер\\
0xE820 \textbf{hpic} [\emph{ʰpiʃ}]: оптатив настроение\\
0xE821 \textbf{mwic} [\emph{mwiʃ}]: пчела\\
0xE822 \textbf{kwic} [\emph{kwiʃ}]: иностранный\\
0xE824 \textbf{pwic} [\emph{pwiʃ}]: держа\\
0xE826 \textbf{nwic} [\emph{nwiʃ}]: дюйм\\
0xE828 \textbf{hwic} [\emph{ʰwiʃ}]: воды\\
0xE829 \textbf{swic} [\emph{swiʃ}]: Социальное\\
0xE82A \textbf{twic} [\emph{twiʃ}]: остатки\\
0xE82B \textbf{lwic} [\emph{lwiʃ}]: развод\\
0xE82C \textbf{fwic} [\emph{fwiʃ}]: фактически\\
0xE82D \textbf{cwic} [\emph{ʃwiʃ}]: тигр\\
0xE82E \textbf{rwic} [\emph{rwiʃ}]: радуга\\
0xE830 \textbf{hnic} [\emph{ʰniʃ}]: неотъемлемое владение\\
0xE831 \textbf{bwic} [\emph{bwiʃ}]: подводная лодка\\
0xE832 \textbf{gwic} [\emph{gwiʃ}]: переключатель\\
0xE833 \textbf{dwic} [\emph{dwiʃ}]: отклонение\\
0xE835 \textbf{jwic} [\emph{ʒwiʃ}]: омоложение\\
0xE83A \textbf{xwic} [\emph{xwiʃ}]: ручеек\\
0xE842 \textbf{ksic} [\emph{ksiʃ}]: кризис\\
0xE844 \textbf{psic} [\emph{psiʃ}]: берег\\
0xE848 \textbf{hsic} [\emph{ʰsiʃ}]: Вот\\
0xE84A \textbf{tsic} [\emph{tsiʃ}]: пытаясь\\
0xE850 \textbf{htic} [\emph{ʰtiʃ}]: их\\
0xE851 \textbf{bzic} [\emph{bziʃ}]: Бразилия\\
0xE853 \textbf{dzic} [\emph{dziʃ}]: дизайнер\\
0xE858 \textbf{hlic} [\emph{ʰliʃ}]: уже\\
0xE860 \textbf{hfic} [\emph{ʰfiʃ}]: влияние\\
0xE862 \textbf{klic} [\emph{kliʃ}]: колледж\\
0xE864 \textbf{plic} [\emph{pliʃ}]: будущее\\
0xE868 \textbf{hcic} [\emph{ʰʃiʃ}]: вниз\\
0xE869 \textbf{slic} [\emph{sliʃ}]: оружие\\
0xE86A \textbf{tlic} [\emph{tliʃ}]: в конце концов\\
0xE86C \textbf{flic} [\emph{fliʃ}]: теология\\
0xE86D \textbf{clic} [\emph{ʃliʃ}]: кольцо\\
0xE870 \textbf{hric} [\emph{ʰriʃ}]: постепенно\\
0xE871 \textbf{blic} [\emph{bliʃ}]: иллюзия\\
0xE872 \textbf{glic} [\emph{gliʃ}]: ледник\\
0xE873 \textbf{dlic} [\emph{dliʃ}]: Идеально\\
0xE874 \textbf{zlic} [\emph{zliʃ}]: Чили\\
0xE875 \textbf{jlic} [\emph{ʒliʃ}]: спутник\\
0xE876 \textbf{vlic} [\emph{vliʃ}]: орнитология\\
0xE87A \textbf{xlic} [\emph{xliʃ}]: Ахиллес\\
0xE882 \textbf{kfic} [\emph{kfiʃ}]: отпуск\\
0xE884 \textbf{pfic} [\emph{pfiʃ}]: представление\\
0xE888 \textbf{hbic} [\emph{ʰbiʃ}]: разорение\\
0xE88A \textbf{tfic} [\emph{tfiʃ}]: винтовка\\
0xE890 \textbf{hgic} [\emph{ʰgiʃ}]: переполненный\\
0xE891 \textbf{bvic} [\emph{bviʃ}]: все overish\\
0xE892 \textbf{gvic} [\emph{gviʃ}]: кровожадный\\
0xE893 \textbf{dvic} [\emph{dviʃ}]: двойственность\\
0xE898 \textbf{hdic} [\emph{ʰdiʃ}]: диск\\
0xE8A0 \textbf{hzic} [\emph{ʰziʃ}]: яд\\
0xE8A2 \textbf{kcic} [\emph{kʃiʃ}]: любопытный\\
0xE8A4 \textbf{pcic} [\emph{pʃiʃ}]: должен\\
0xE8A8 \textbf{hjic} [\emph{ʰʒiʃ}]: выходить на пенсию\\
0xE8AA \textbf{tcic} [\emph{tʃiʃ}]: тело\\
0xE8B0 \textbf{hvic} [\emph{ʰviʃ}]: виртуальный машина\\
0xE8B1 \textbf{bjic} [\emph{bʒiʃ}]: подсолнечник\\
0xE8B2 \textbf{gjic} [\emph{gʒiʃ}]: Алжир\\
0xE8B3 \textbf{djic} [\emph{dʒiʃ}]: отвлекаться\\
0xE8C1 \textbf{mric} [\emph{mriʃ}]: мышца\\
0xE8C2 \textbf{kric} [\emph{kriʃ}]: курица\\
0xE8C4 \textbf{pric} [\emph{priʃ}]: подача\\
0xE8C6 \textbf{nric} [\emph{nriʃ}]: поколение\\
0xE8C8 \textbf{hqic} [\emph{ʰŋiʃ}]: фишинг\\
0xE8C9 \textbf{sric} [\emph{sriʃ}]: церемония\\
0xE8CA \textbf{tric} [\emph{triʃ}]: временный\\
0xE8CC \textbf{fric} [\emph{friʃ}]: отдых\\
0xE8CD \textbf{cric} [\emph{ʃriʃ}]: кров\\
0xE8D0 \textbf{hxic} [\emph{ʰxiʃ}]: сдвиг ключ\\
0xE8D1 \textbf{bric} [\emph{briʃ}]: батарея\\
0xE8D2 \textbf{gric} [\emph{griʃ}]: миграция\\
0xE8D3 \textbf{dric} [\emph{driʃ}]: невидимый\\
0xE8D4 \textbf{zric} [\emph{zriʃ}]: выравнивать\\
0xE8D5 \textbf{jric} [\emph{ʒriʃ}]: маршрут\\
0xE8D6 \textbf{vric} [\emph{vriʃ}]: Австрия\\
0xE8D9 \textbf{qric} [\emph{ŋriʃ}]: ирландцы\\
0xE8DA \textbf{xric} [\emph{xriʃ}]: кувшин\\
0xE8E2 \textbf{kxic} [\emph{kxiʃ}]: огнетушитель\\
0xE8E4 \textbf{pxic} [\emph{pxiʃ}]: селезенка\\
0xE8EA \textbf{txic} [\emph{txiʃ}]: вторгшийся\\
0xE8F1 \textbf{bxic} [\emph{bxiʃ}]: грунтовый\\
0xE8F2 \textbf{gxic} [\emph{gxiʃ}]: огненный\\
0xE8F3 \textbf{dxic} [\emph{dxiʃ}]: Фиджи\\
0xE901 \textbf{myac} [\emph{mjäʃ}]: улыбка\\
0xE902 \textbf{kyac} [\emph{kjäʃ}]: глаз\\
0xE904 \textbf{pyac} [\emph{pjäʃ}]: язык\\
0xE906 \textbf{nyac} [\emph{njäʃ}]: направление\\
0xE908 \textbf{hmac} [\emph{ʰmäʃ}]: момент\\
0xE909 \textbf{syac} [\emph{sjäʃ}]: свежий\\
0xE90A \textbf{tyac} [\emph{tjäʃ}]: читать\\
0xE90B \textbf{lyac} [\emph{ljäʃ}]: учить\\
0xE90C \textbf{fyac} [\emph{fjäʃ}]: железо\\
0xE90D \textbf{cyac} [\emph{ʃjäʃ}]: кровь\\
0xE90E \textbf{ryac} [\emph{rjäʃ}]: натирать\\
0xE910 \textbf{hkac} [\emph{ʰkäʃ}]: когда\\
0xE911 \textbf{byac} [\emph{bjäʃ}]: подвал\\
0xE912 \textbf{gyac} [\emph{gjäʃ}]: ученый\\
0xE913 \textbf{dyac} [\emph{djäʃ}]: франшиза\\
0xE914 \textbf{zyac} [\emph{zjäʃ}]: озабоченность\\
0xE915 \textbf{jyac} [\emph{ʒjäʃ}]: интернационализация\\
0xE916 \textbf{vyac} [\emph{vjäʃ}]: зодиак\\
0xE918 \textbf{hyac} [\emph{ʰjäʃ}]: также\\
0xE919 \textbf{qyac} [\emph{ŋjäʃ}]: двухгодичный\\
0xE91A \textbf{xyac} [\emph{xjäʃ}]: хрен\\
0xE920 \textbf{hpac} [\emph{ʰpäʃ}]: множественное число номер\\
0xE921 \textbf{mwac} [\emph{mwäʃ}]: доволен\\
0xE922 \textbf{kwac} [\emph{kwäʃ}]: ухо\\
0xE924 \textbf{pwac} [\emph{pwäʃ}]: камень\\
0xE926 \textbf{nwac} [\emph{nwäʃ}]: кто угодно\\
0xE928 \textbf{hwac} [\emph{ʰwäʃ}]: Мир\\
0xE929 \textbf{swac} [\emph{swäʃ}]: служащий\\
0xE92A \textbf{twac} [\emph{twäʃ}]: окно\\
0xE92B \textbf{lwac} [\emph{lwäʃ}]: осел\\
0xE92C \textbf{fwac} [\emph{fwäʃ}]: аэропорт\\
0xE92D \textbf{cwac} [\emph{ʃwäʃ}]: плечо\\
0xE92E \textbf{rwac} [\emph{rwäʃ}]: аппаратный\\
0xE930 \textbf{hnac} [\emph{ʰnäʃ}]: путь дело\\
0xE931 \textbf{bwac} [\emph{bwäʃ}]: зевать\\
0xE932 \textbf{gwac} [\emph{gwäʃ}]: консультант\\
0xE933 \textbf{dwac} [\emph{dwäʃ}]: диалог\\
0xE934 \textbf{zwac} [\emph{zwäʃ}]: недостаток\\
0xE935 \textbf{jwac} [\emph{ʒwäʃ}]: Ява\\
0xE936 \textbf{vwac} [\emph{vwäʃ}]: расточительность\\
0xE939 \textbf{qwac} [\emph{ŋwäʃ}]: Бог оставил\\
0xE93A \textbf{xwac} [\emph{xwäʃ}]: проходы\\
0xE942 \textbf{ksac} [\emph{ksäʃ}]: обстоятельства\\
0xE944 \textbf{psac} [\emph{psäʃ}]: невозможно\\
0xE948 \textbf{hsac} [\emph{ʰsäʃ}]: сказать\\
0xE94A \textbf{tsac} [\emph{tsäʃ}]: довольно часто\\
0xE950 \textbf{htac} [\emph{ʰtäʃ}]: место назначения дело\\
0xE951 \textbf{bzac} [\emph{bzäʃ}]: сокращать\\
0xE952 \textbf{gzac} [\emph{gzäʃ}]: воспламеняться\\
0xE953 \textbf{dzac} [\emph{dzäʃ}]: обедающий\\
0xE958 \textbf{hlac} [\emph{ʰläʃ}]: отчуждаемый владение\\
0xE960 \textbf{hfac} [\emph{ʰfäʃ}]: семья\\
0xE962 \textbf{klac} [\emph{kläʃ}]: страна\\
0xE964 \textbf{plac} [\emph{pläʃ}]: дворец\\
0xE968 \textbf{hcac} [\emph{ʰʃäʃ}]: голодный\\
0xE969 \textbf{slac} [\emph{släʃ}]: Бег\\
0xE96A \textbf{tlac} [\emph{tläʃ}]: тихо\\
0xE96C \textbf{flac} [\emph{fläʃ}]: вспышка\\
0xE96D \textbf{clac} [\emph{ʃläʃ}]: здоровье\\
0xE970 \textbf{hrac} [\emph{ʰräʃ}]: везде\\
0xE971 \textbf{blac} [\emph{bläʃ}]: Багаж\\
0xE972 \textbf{glac} [\emph{gläʃ}]: жадность\\
0xE973 \textbf{dlac} [\emph{dläʃ}]: мыть\\
0xE974 \textbf{zlac} [\emph{zläʃ}]: таракан\\
0xE975 \textbf{jlac} [\emph{ʒläʃ}]: массаж\\
0xE976 \textbf{vlac} [\emph{vläʃ}]: морская звезда\\
0xE97A \textbf{xlac} [\emph{xläʃ}]: легальность\\
0xE982 \textbf{kfac} [\emph{kfäʃ}]: трава\\
0xE984 \textbf{pfac} [\emph{pfäʃ}]: подвергать\\
0xE988 \textbf{hbac} [\emph{ʰbäʃ}]: двоюродная сестра\\
0xE98A \textbf{tfac} [\emph{tfäʃ}]: место действия\\
0xE990 \textbf{hgac} [\emph{ʰgäʃ}]: амбиция\\
0xE991 \textbf{bvac} [\emph{bväʃ}]: гибридизация\\
0xE993 \textbf{dvac} [\emph{dväʃ}]: яичник\\
0xE998 \textbf{hdac} [\emph{ʰdäʃ}]: адвокат\\
0xE9A0 \textbf{hzac} [\emph{ʰzäʃ}]: взаимное\\
0xE9A2 \textbf{kcac} [\emph{kʃäʃ}]: оставил\\
0xE9A4 \textbf{pcac} [\emph{pʃäʃ}]: лучше\\
0xE9A8 \textbf{hjac} [\emph{ʰʒäʃ}]: Отправить\\
0xE9AA \textbf{tcac} [\emph{tʃäʃ}]: девушка\\
0xE9B0 \textbf{hvac} [\emph{ʰväʃ}]: посудомоечная машина\\
0xE9B1 \textbf{bjac} [\emph{bʒäʃ}]: пивоваренный завод\\
0xE9B2 \textbf{gjac} [\emph{gʒäʃ}]: выкидыш\\
0xE9B3 \textbf{djac} [\emph{dʒäʃ}]: примириться\\
0xE9C1 \textbf{mrac} [\emph{mräʃ}]: вторник\\
0xE9C2 \textbf{krac} [\emph{kräʃ}]: жестокий\\
0xE9C4 \textbf{prac} [\emph{präʃ}]: выводить\\
0xE9C6 \textbf{nrac} [\emph{nräʃ}]: безвредный\\
0xE9C8 \textbf{hqac} [\emph{ʰŋäʃ}]: пропадать\\
0xE9C9 \textbf{srac} [\emph{sräʃ}]: дешево\\
0xE9CA \textbf{trac} [\emph{träʃ}]: сопротивление\\
0xE9CC \textbf{frac} [\emph{fräʃ}]: бунтарь\\
0xE9CD \textbf{crac} [\emph{ʃräʃ}]: филиал\\
0xE9D0 \textbf{hxac} [\emph{ʰxäʃ}]: перезапуск\\
0xE9D1 \textbf{brac} [\emph{bräʃ}]: среда\\
0xE9D2 \textbf{grac} [\emph{gräʃ}]: трусливый\\
0xE9D3 \textbf{drac} [\emph{dräʃ}]: речь\\
0xE9D4 \textbf{zrac} [\emph{zräʃ}]: обновление\\
0xE9D5 \textbf{jrac} [\emph{ʒräʃ}]: карта\\
0xE9D6 \textbf{vrac} [\emph{vräʃ}]: лак\\
0xE9D9 \textbf{qrac} [\emph{ŋräʃ}]: вспашка\\
0xE9DA \textbf{xrac} [\emph{xräʃ}]: коралловый\\
0xE9E2 \textbf{kxac} [\emph{kxäʃ}]: восклицание\\
0xE9E4 \textbf{pxac} [\emph{pxäʃ}]: чуткими радость\\
0xE9EA \textbf{txac} [\emph{txäʃ}]: капитуляция\\
0xE9F1 \textbf{bxac} [\emph{bxäʃ}]: мол\\
0xE9F2 \textbf{gxac} [\emph{gxäʃ}]: инертность\\
0xE9F3 \textbf{dxac} [\emph{dxäʃ}]: терка\\
0xEA01 \textbf{myuc} [\emph{mjuʃ}]: эмульсия\\
0xEA02 \textbf{kyuc} [\emph{kjuʃ}]: сухой\\
0xEA04 \textbf{pyuc} [\emph{pjuʃ}]: истребление населения\\
0xEA06 \textbf{nyuc} [\emph{njuʃ}]: маникюр\\
0xEA08 \textbf{hmuc} [\emph{ʰmuʃ}]: цель\\
0xEA09 \textbf{syuc} [\emph{sjuʃ}]: дерево\\
0xEA0A \textbf{tyuc} [\emph{tjuʃ}]: ожидание\\
0xEA0B \textbf{lyuc} [\emph{ljuʃ}]: тобогган-одиночка\\
0xEA0D \textbf{cyuc} [\emph{ʃjuʃ}]: высокая температура\\
0xEA0E \textbf{ryuc} [\emph{rjuʃ}]: преобразователь\\
0xEA10 \textbf{hkuc} [\emph{ʰkuʃ}]: высокая\\
0xEA11 \textbf{byuc} [\emph{bjuʃ}]: бутил\\
0xEA18 \textbf{hyuc} [\emph{ʰjuʃ}]: круглый\\
0xEA1A \textbf{xyuc} [\emph{xjuʃ}]: удобство использования\\
0xEA20 \textbf{hpuc} [\emph{ʰpuʃ}]: напиток\\
0xEA21 \textbf{mwuc} [\emph{mwuʃ}]: полоскание для рта\\
0xEA22 \textbf{kwuc} [\emph{kwuʃ}]: принуждение\\
0xEA26 \textbf{nwuc} [\emph{nwuʃ}]: движение спуск\\
0xEA28 \textbf{hwuc} [\emph{ʰwuʃ}]: связи\\
0xEA29 \textbf{swuc} [\emph{swuʃ}]: иней\\
0xEA2A \textbf{twuc} [\emph{twuʃ}]: дым\\
0xEA2D \textbf{cwuc} [\emph{ʃwuʃ}]: Choo Choo\\
0xEA2E \textbf{rwuc} [\emph{rwuʃ}]: порядковый\\
0xEA30 \textbf{hnuc} [\emph{ʰnuʃ}]: номер\\
0xEA42 \textbf{ksuc} [\emph{ksuʃ}]: вкусно\\
0xEA44 \textbf{psuc} [\emph{psuʃ}]: пластичность\\
0xEA48 \textbf{hsuc} [\emph{ʰsuʃ}]: предок\\
0xEA4A \textbf{tsuc} [\emph{tsuʃ}]: куст\\
0xEA50 \textbf{htuc} [\emph{ʰtuʃ}]: другой предмет\\
0xEA51 \textbf{bzuc} [\emph{bzuʃ}]: огрубение\\
0xEA58 \textbf{hluc} [\emph{ʰluʃ}]: популярный\\
0xEA60 \textbf{hfuc} [\emph{ʰfuʃ}]: воскрешение\\
0xEA62 \textbf{kluc} [\emph{kluʃ}]: инструменты\\
0xEA64 \textbf{pluc} [\emph{pluʃ}]: paucal номер\\
0xEA68 \textbf{hcuc} [\emph{ʰʃuʃ}]: страх\\
0xEA69 \textbf{sluc} [\emph{sluʃ}]: токсикология\\
0xEA6A \textbf{tluc} [\emph{tluʃ}]: полотенце\\
0xEA6C \textbf{fluc} [\emph{fluʃ}]: беглый\\
0xEA6D \textbf{cluc} [\emph{ʃluʃ}]: хитрость\\
0xEA70 \textbf{hruc} [\emph{ʰruʃ}]: грубый\\
0xEA71 \textbf{bluc} [\emph{bluʃ}]: предродовой\\
0xEA72 \textbf{gluc} [\emph{gluʃ}]: галоши\\
0xEA73 \textbf{dluc} [\emph{dluʃ}]: Андалусия\\
0xEA7A \textbf{xluc} [\emph{xluʃ}]: ремень кнута\\
0xEA88 \textbf{hbuc} [\emph{ʰbuʃ}]: пользователь\\
0xEA8A \textbf{tfuc} [\emph{tfuʃ}]: ткацкий станок\\
0xEA90 \textbf{hguc} [\emph{ʰguʃ}]: жажда\\
0xEA98 \textbf{hduc} [\emph{ʰduʃ}]: пыли\\
0xEAA0 \textbf{hzuc} [\emph{ʰzuʃ}]: эксгумация\\
0xEAA2 \textbf{kcuc} [\emph{kʃuʃ}]: горький\\
0xEAA4 \textbf{pcuc} [\emph{pʃuʃ}]: разряд\\
0xEAA8 \textbf{hjuc} [\emph{ʰʒuʃ}]: добыча\\
0xEAAA \textbf{tcuc} [\emph{tʃuʃ}]: появляться\\
0xEAB0 \textbf{hvuc} [\emph{ʰvuʃ}]: Вишну\\
0xEAC1 \textbf{mruc} [\emph{mruʃ}]: гусеница\\
0xEAC2 \textbf{kruc} [\emph{kruʃ}]: черепаха\\
0xEAC4 \textbf{pruc} [\emph{pruʃ}]: патруль\\
0xEAC6 \textbf{nruc} [\emph{nruʃ}]: моча\\
0xEAC8 \textbf{hquc} [\emph{ʰŋuʃ}]: скульптор\\
0xEAC9 \textbf{sruc} [\emph{sruʃ}]: вонь\\
0xEACA \textbf{truc} [\emph{truʃ}]: украсть\\
0xEACD \textbf{cruc} [\emph{ʃruʃ}]: разрушающий\\
0xEAD0 \textbf{hxuc} [\emph{ʰxuʃ}]: спасающий\\
0xEAD1 \textbf{bruc} [\emph{bruʃ}]: жаропонижающий\\
0xEAD2 \textbf{gruc} [\emph{gruʃ}]: гуттаперча\\
0xEAD3 \textbf{druc} [\emph{druʃ}]: UAE\\
0xEAD4 \textbf{zruc} [\emph{zruʃ}]: цюрих\\
0xEADA \textbf{xruc} [\emph{xruʃ}]: землеройка\\
0xEAE2 \textbf{kxuc} [\emph{kxuʃ}]: судорога\\
0xEAE4 \textbf{pxuc} [\emph{pxuʃ}]: пачули\\
0xEAEA \textbf{txuc} [\emph{txuʃ}]: выпуклость\\
0xEB01 \textbf{myec} [\emph{mjeʃ}]: Малайзия\\
0xEB02 \textbf{kyec} [\emph{kjeʃ}]: рабочее место\\
0xEB04 \textbf{pyec} [\emph{pjeʃ}]: таблица\\
0xEB06 \textbf{nyec} [\emph{njeʃ}]: Индонезия\\
0xEB08 \textbf{hmec} [\emph{ʰmeʃ}]: неопределенный\\
0xEB09 \textbf{syec} [\emph{sjeʃ}]: располагать\\
0xEB0A \textbf{tyec} [\emph{tjeʃ}]: тег\\
0xEB0B \textbf{lyec} [\emph{ljeʃ}]: линейный\\
0xEB0D \textbf{cyec} [\emph{ʃjeʃ}]: чихать\\
0xEB0E \textbf{ryec} [\emph{rjeʃ}]: Розмари\\
0xEB10 \textbf{hkec} [\emph{ʰkeʃ}]: каменный уголь\\
0xEB13 \textbf{dyec} [\emph{djeʃ}]: денатурировать\\
0xEB18 \textbf{hyec} [\emph{ʰjeʃ}]: овца\\
0xEB1A \textbf{xyec} [\emph{xjeʃ}]: гипотония\\
0xEB20 \textbf{hpec} [\emph{ʰpeʃ}]: профессия\\
0xEB21 \textbf{mwec} [\emph{mweʃ}]: однозначно\\
0xEB24 \textbf{pwec} [\emph{pweʃ}]: партитивный дело\\
0xEB26 \textbf{nwec} [\emph{nweʃ}]: тем не менее\\
0xEB28 \textbf{hwec} [\emph{ʰweʃ}]: осада\\
0xEB29 \textbf{swec} [\emph{sweʃ}]: случайный\\
0xEB2A \textbf{twec} [\emph{tweʃ}]: дистанционный пульт будущее напряженный\\
0xEB2B \textbf{lwec} [\emph{lweʃ}]: аллели\\
0xEB2D \textbf{cwec} [\emph{ʃweʃ}]: словесный\\
0xEB2E \textbf{rwec} [\emph{rweʃ}]: рейнджер\\
0xEB30 \textbf{hnec} [\emph{ʰneʃ}]: Солнечный свет\\
0xEB33 \textbf{dwec} [\emph{dweʃ}]: отвесно\\
0xEB42 \textbf{ksec} [\emph{kseʃ}]: одноклассник\\
0xEB44 \textbf{psec} [\emph{pseʃ}]: подмножество\\
0xEB48 \textbf{hsec} [\emph{ʰseʃ}]: программного обеспечения\\
0xEB4A \textbf{tsec} [\emph{tseʃ}]: эксперт\\
0xEB50 \textbf{htec} [\emph{ʰteʃ}]: энтузиазм\\
0xEB58 \textbf{hlec} [\emph{ʰleʃ}]: перец\\
0xEB60 \textbf{hfec} [\emph{ʰfeʃ}]: обеспечение\\
0xEB62 \textbf{klec} [\emph{kleʃ}]: собака устала\\
0xEB64 \textbf{plec} [\emph{pleʃ}]: сращивание\\
0xEB68 \textbf{hcec} [\emph{ʰʃeʃ}]: учение\\
0xEB69 \textbf{slec} [\emph{sleʃ}]: исследователь\\
0xEB6A \textbf{tlec} [\emph{tleʃ}]: поставить галочку\\
0xEB6C \textbf{flec} [\emph{fleʃ}]: фибрилляция\\
0xEB6D \textbf{clec} [\emph{ʃleʃ}]: выхлопные газы\\
0xEB70 \textbf{hrec} [\emph{ʰreʃ}]: справочный\\
0xEB71 \textbf{blec} [\emph{bleʃ}]: Бангладеш\\
0xEB72 \textbf{glec} [\emph{gleʃ}]: пеньюар\\
0xEB73 \textbf{dlec} [\emph{dleʃ}]: ноготь на пальце ноги\\
0xEB76 \textbf{vlec} [\emph{vleʃ}]: калька\\
0xEB7A \textbf{xlec} [\emph{xleʃ}]: Хэллоуин\\
0xEB82 \textbf{kfec} [\emph{kfeʃ}]: кеч\\
0xEB84 \textbf{pfec} [\emph{pfeʃ}]: незамеченной\\
0xEB88 \textbf{hbec} [\emph{ʰbeʃ}]: благожелательный\\
0xEB8A \textbf{tfec} [\emph{tfeʃ}]: канцелярские товары\\
0xEB90 \textbf{hgec} [\emph{ʰgeʃ}]: пещера\\
0xEB93 \textbf{dvec} [\emph{dveʃ}]: дервиш\\
0xEB98 \textbf{hdec} [\emph{ʰdeʃ}]: К сожалению\\
0xEBA0 \textbf{hzec} [\emph{ʰzeʃ}]: одновременность\\
0xEBA2 \textbf{kcec} [\emph{kʃeʃ}]: область\\
0xEBA4 \textbf{pcec} [\emph{pʃeʃ}]: привилегированный\\
0xEBA8 \textbf{hjec} [\emph{ʰʒeʃ}]: плавиться\\
0xEBAA \textbf{tcec} [\emph{tʃeʃ}]: обмен\\
0xEBB0 \textbf{hvec} [\emph{ʰveʃ}]: вершина\\
0xEBC1 \textbf{mrec} [\emph{mreʃ}]: воробей\\
0xEBC2 \textbf{krec} [\emph{kreʃ}]: кран\\
0xEBC4 \textbf{prec} [\emph{preʃ}]: вести себя\\
0xEBC6 \textbf{nrec} [\emph{nreʃ}]: сенсорный\\
0xEBC8 \textbf{hqec} [\emph{ʰŋeʃ}]: бесполый\\
0xEBC9 \textbf{srec} [\emph{sreʃ}]: истекающий\\
0xEBCA \textbf{trec} [\emph{treʃ}]: стратегия\\
0xEBCC \textbf{frec} [\emph{freʃ}]: Французский\\
0xEBCD \textbf{crec} [\emph{ʃreʃ}]: шахматы\\
0xEBD0 \textbf{hxec} [\emph{ʰxeʃ}]: Тедди\\
0xEBD1 \textbf{brec} [\emph{breʃ}]: бровь\\
0xEBD2 \textbf{grec} [\emph{greʃ}]: инфракрасный\\
0xEBD3 \textbf{drec} [\emph{dreʃ}]: рассечение\\
0xEBD4 \textbf{zrec} [\emph{zreʃ}]: взаимодействие\\
0xEBD6 \textbf{vrec} [\emph{vreʃ}]: жадностью\\
0xEBD9 \textbf{qrec} [\emph{ŋreʃ}]: поджелудочная железа\\
0xEBDA \textbf{xrec} [\emph{xreʃ}]: фамильная ценность\\
0xEBE2 \textbf{kxec} [\emph{kxeʃ}]: коктейль\\
0xEBE4 \textbf{pxec} [\emph{pxeʃ}]: персик\\
0xEBEA \textbf{txec} [\emph{txeʃ}]: телекоммуникация\\
0xEBF2 \textbf{gxec} [\emph{gxeʃ}]: щегол\\
0xEBF3 \textbf{dxec} [\emph{dxeʃ}]: залогодержатель\\
0xEC01 \textbf{myoc} [\emph{mjoʃ}]: модуль\\
0xEC02 \textbf{kyoc} [\emph{kjoʃ}]: причинная связь\\
0xEC04 \textbf{pyoc} [\emph{pjoʃ}]: Полинезия\\
0xEC06 \textbf{nyoc} [\emph{njoʃ}]: антропология\\
0xEC08 \textbf{hmoc} [\emph{ʰmoʃ}]: где-то\\
0xEC09 \textbf{syoc} [\emph{sjoʃ}]: снимок\\
0xEC0A \textbf{tyoc} [\emph{tjoʃ}]: совпадение\\
0xEC0B \textbf{lyoc} [\emph{ljoʃ}]: лиственница\\
0xEC0D \textbf{cyoc} [\emph{ʃjoʃ}]: испытывающий\\
0xEC0E \textbf{ryoc} [\emph{rjoʃ}]: радиоактивность\\
0xEC10 \textbf{hkoc} [\emph{ʰkoʃ}]: время от времени\\
0xEC12 \textbf{gyoc} [\emph{gjoʃ}]: часовое дело\\
0xEC13 \textbf{dyoc} [\emph{djoʃ}]: пособие по безработице\\
0xEC18 \textbf{hyoc} [\emph{ʰjoʃ}]: искушение\\
0xEC1A \textbf{xyoc} [\emph{xjoʃ}]: приглушенный\\
0xEC20 \textbf{hpoc} [\emph{ʰpoʃ}]: осень\\
0xEC21 \textbf{mwoc} [\emph{mwoʃ}]: микробиология\\
0xEC22 \textbf{kwoc} [\emph{kwoʃ}]: сквош\\
0xEC24 \textbf{pwoc} [\emph{pwoʃ}]: Плутон\\
0xEC26 \textbf{nwoc} [\emph{nwoʃ}]: с этого времени\\
0xEC28 \textbf{hwoc} [\emph{ʰwoʃ}]: ведьма\\
0xEC29 \textbf{swoc} [\emph{swoʃ}]: неустойчивый\\
0xEC2A \textbf{twoc} [\emph{twoʃ}]: продуктовый\\
0xEC2B \textbf{lwoc} [\emph{lwoʃ}]: лаосский\\
0xEC2D \textbf{cwoc} [\emph{ʃwoʃ}]: хорошо\\
0xEC2E \textbf{rwoc} [\emph{rwoʃ}]: последовательный\\
0xEC30 \textbf{hnoc} [\emph{ʰnoʃ}]: понедельник\\
0xEC31 \textbf{bwoc} [\emph{bwoʃ}]: саботаж\\
0xEC33 \textbf{dwoc} [\emph{dwoʃ}]: дерматолог\\
0xEC3A \textbf{xwoc} [\emph{xwoʃ}]: поджилки\\
0xEC42 \textbf{ksoc} [\emph{ksoʃ}]: вычет\\
0xEC44 \textbf{psoc} [\emph{psoʃ}]: перфоратор\\
0xEC48 \textbf{hsoc} [\emph{ʰsoʃ}]: акции\\
0xEC4A \textbf{tsoc} [\emph{tsoʃ}]: поиск\\
0xEC50 \textbf{htoc} [\emph{ʰtoʃ}]: как-то\\
0xEC58 \textbf{hloc} [\emph{ʰloʃ}]: меланхолия\\
0xEC60 \textbf{hfoc} [\emph{ʰfoʃ}]: фотография\\
0xEC62 \textbf{kloc} [\emph{kloʃ}]: воротник\\
0xEC64 \textbf{ploc} [\emph{ploʃ}]: экология\\
0xEC68 \textbf{hcoc} [\emph{ʰʃoʃ}]: полый\\
0xEC69 \textbf{sloc} [\emph{sloʃ}]: психология\\
0xEC6A \textbf{tloc} [\emph{tloʃ}]: ступенька\\
0xEC6C \textbf{floc} [\emph{floʃ}]: фольклор\\
0xEC6D \textbf{cloc} [\emph{ʃloʃ}]: морфологический\\
0xEC70 \textbf{hroc} [\emph{ʰroʃ}]: традиция\\
0xEC71 \textbf{bloc} [\emph{bloʃ}]: боулинг\\
0xEC72 \textbf{gloc} [\emph{gloʃ}]: зоология\\
0xEC73 \textbf{dloc} [\emph{dloʃ}]: Голландский\\
0xEC74 \textbf{zloc} [\emph{zloʃ}]: космология\\
0xEC75 \textbf{jloc} [\emph{ʒloʃ}]: острие иглы\\
0xEC76 \textbf{vloc} [\emph{vloʃ}]: коэволюция\\
0xEC7A \textbf{xloc} [\emph{xloʃ}]: ностальгия\\
0xEC82 \textbf{kfoc} [\emph{kfoʃ}]: Хорватия\\
0xEC84 \textbf{pfoc} [\emph{pfoʃ}]: профилирование\\
0xEC88 \textbf{hboc} [\emph{ʰboʃ}]: обязать\\
0xEC8A \textbf{tfoc} [\emph{tfoʃ}]: приступ боли\\
0xEC90 \textbf{hgoc} [\emph{ʰgoʃ}]: умножить\\
0xEC98 \textbf{hdoc} [\emph{ʰdoʃ}]: ой\\
0xECA0 \textbf{hzoc} [\emph{ʰzoʃ}]: лежать в основе\\
0xECA2 \textbf{kcoc} [\emph{kʃoʃ}]: локоть\\
0xECA4 \textbf{pcoc} [\emph{pʃoʃ}]: столб\\
0xECA8 \textbf{hjoc} [\emph{ʰʒoʃ}]: монтаж\\
0xECAA \textbf{tcoc} [\emph{tʃoʃ}]: крутой\\
0xECB1 \textbf{bjoc} [\emph{bʒoʃ}]: банджо\\
0xECC1 \textbf{mroc} [\emph{mroʃ}]: трение\\
0xECC2 \textbf{kroc} [\emph{kroʃ}]: борозда\\
0xECC4 \textbf{proc} [\emph{proʃ}]: покровитель\\
0xECC6 \textbf{nroc} [\emph{nroʃ}]: оргия\\
0xECC8 \textbf{hqoc} [\emph{ʰŋoʃ}]: тошнота\\
0xECC9 \textbf{sroc} [\emph{sroʃ}]: экспорт\\
0xECCA \textbf{troc} [\emph{troʃ}]: факел\\
0xECCC \textbf{froc} [\emph{froʃ}]: ковочный\\
0xECCD \textbf{croc} [\emph{ʃroʃ}]: ошибка\\
0xECD0 \textbf{hxoc} [\emph{ʰxoʃ}]: ров\\
0xECD1 \textbf{broc} [\emph{broʃ}]: щетка\\
0xECD2 \textbf{groc} [\emph{groʃ}]: Грузия\\
0xECD3 \textbf{droc} [\emph{droʃ}]: докторская степень\\
0xECD5 \textbf{jroc} [\emph{ʒroʃ}]: Козерог\\
0xECD9 \textbf{qroc} [\emph{ŋroʃ}]: Джордж\\
0xECDA \textbf{xroc} [\emph{xroʃ}]: гончар\\
0xECE2 \textbf{kxoc} [\emph{kxoʃ}]: петушок\\
0xECE4 \textbf{pxoc} [\emph{pxoʃ}]: за предложение\\
0xECEA \textbf{txoc} [\emph{txoʃ}]: жеребенок\\
0xECF2 \textbf{gxoc} [\emph{gxoʃ}]: горячий воздух\\
0xECF3 \textbf{dxoc} [\emph{dxoʃ}]: дезодорант\\
0xED08 \textbf{hm6c} [\emph{ʰməʃ}]: всеядный\\
0xED10 \textbf{hk6c} [\emph{ʰkəʃ}]: канарейка\\
0xED20 \textbf{hp6c} [\emph{ʰpəʃ}]: типологический\\
0xED28 \textbf{hw6c} [\emph{ʰwəʃ}]: водянистый\\
0xED30 \textbf{hn6c} [\emph{ʰnəʃ}]: Нигер\\
0xED42 \textbf{ks6c} [\emph{ksəʃ}]: инкрустация\\
0xED44 \textbf{ps6c} [\emph{psəʃ}]: просодия\\
0xED48 \textbf{hs6c} [\emph{ʰsəʃ}]: само уверенность\\
0xED4A \textbf{ts6c} [\emph{tsəʃ}]: Танзания\\
0xED50 \textbf{ht6c} [\emph{ʰtəʃ}]: стебель\\
0xED58 \textbf{hl6c} [\emph{ʰləʃ}]: выщелачивание\\
0xED60 \textbf{hf6c} [\emph{ʰfəʃ}]: Phenicia\\
0xED62 \textbf{kl6c} [\emph{kləʃ}]: коллаген\\
0xED64 \textbf{pl6c} [\emph{pləʃ}]: щегольство\\
0xED68 \textbf{hc6c} [\emph{ʰʃəʃ}]: цикорий\\
0xED69 \textbf{sl6c} [\emph{sləʃ}]: фоссилизация\\
0xED6A \textbf{tl6c} [\emph{tləʃ}]: энтомолог\\
0xED70 \textbf{hr6c} [\emph{ʰrəʃ}]: двадцать семь\\
0xED90 \textbf{hg6c} [\emph{ʰgəʃ}]: гейша\\
0xED98 \textbf{hd6c} [\emph{ʰdəʃ}]: далматинец\\
0xEDA2 \textbf{kc6c} [\emph{kʃəʃ}]: совокупление\\
0xEDAA \textbf{tc6c} [\emph{tʃəʃ}]: мангольд\\
0xEDC1 \textbf{mr6c} [\emph{mrəʃ}]: Махараштра\\
0xEDC2 \textbf{kr6c} [\emph{krəʃ}]: штопор\\
0xEDC4 \textbf{pr6c} [\emph{prəʃ}]: взаимопроникновение\\
0xEDC6 \textbf{nr6c} [\emph{nrəʃ}]: Нигерия\\
0xEDC9 \textbf{sr6c} [\emph{srəʃ}]: усик\\
0xEDCA \textbf{tr6c} [\emph{trəʃ}]: щитовидная железа\\
0xEDD0 \textbf{hx6c} [\emph{ʰxəʃ}]: молотилка\\
0xEDD3 \textbf{dr6c} [\emph{drəʃ}]: квадратура\\
0xEDE2 \textbf{kx6c} [\emph{kxəʃ}]: чудак\\
0xEDEA \textbf{tx6c} [\emph{txəʃ}]: обнаружений\\
\subsection{ F }
0xF001 \textbf{myi7c} [\emph{mji˥ʃ}]: метилен\\
0xF006 \textbf{nyi7c} [\emph{nji˥ʃ}]: вновь ввести\\
0xF008 \textbf{hmi7c} [\emph{ʰmi˥ʃ}]: губчатый\\
0xF009 \textbf{syi7c} [\emph{sji˥ʃ}]: ацетил\\
0xF00A \textbf{tyi7c} [\emph{tji˥ʃ}]: быть принятым в высшее учебное заведение\\
0xF00B \textbf{lyi7c} [\emph{lji˥ʃ}]: алкил\\
0xF010 \textbf{hki7c} [\emph{ʰki˥ʃ}]: гвоздика\\
0xF018 \textbf{hyi7c} [\emph{ʰji˥ʃ}]: доброжелатель\\
0xF020 \textbf{hpi7c} [\emph{ʰpi˥ʃ}]: плагиат\\
0xF028 \textbf{hwi7c} [\emph{ʰwi˥ʃ}]: обморок\\
0xF029 \textbf{swi7c} [\emph{swi˥ʃ}]: относящийся к пищеводу\\
0xF030 \textbf{hni7c} [\emph{ʰni˥ʃ}]: сам\\
0xF044 \textbf{psi7c} [\emph{psi˥ʃ}]: пессарий\\
0xF048 \textbf{hsi7c} [\emph{ʰsi˥ʃ}]: сатира\\
0xF04A \textbf{tsi7c} [\emph{tsi˥ʃ}]: девяносто девять\\
0xF050 \textbf{hti7c} [\emph{ʰti˥ʃ}]: темы\\
0xF058 \textbf{hli7c} [\emph{ʰli˥ʃ}]: лотерея\\
0xF060 \textbf{hfi7c} [\emph{ʰfi˥ʃ}]: афазия\\
0xF062 \textbf{kli7c} [\emph{kli˥ʃ}]: перекристаллизация\\
0xF064 \textbf{pli7c} [\emph{pli˥ʃ}]: параплегия\\
0xF068 \textbf{hci7c} [\emph{ʰʃi˥ʃ}]: визг\\
0xF069 \textbf{sli7c} [\emph{sli˥ʃ}]: гистология\\
0xF06A \textbf{tli7c} [\emph{tli˥ʃ}]: перекупщик\\
0xF06D \textbf{cli7c} [\emph{ʃli˥ʃ}]: аллил\\
0xF070 \textbf{hri7c} [\emph{ʰri˥ʃ}]: негодяй\\
0xF088 \textbf{hbi7c} [\emph{ʰbi˥ʃ}]: непреклонный\\
0xF08A \textbf{tfi7c} [\emph{tfi˥ʃ}]: этил\\
0xF090 \textbf{hgi7c} [\emph{ʰgi˥ʃ}]: эгоистический\\
0xF098 \textbf{hdi7c} [\emph{ʰdi˥ʃ}]: выводить из заблуждения\\
0xF0A2 \textbf{kci7c} [\emph{kʃi˥ʃ}]: отрыжка\\
0xF0A4 \textbf{pci7c} [\emph{pʃi˥ʃ}]: путеводная звезда\\
0xF0AA \textbf{tci7c} [\emph{tʃi˥ʃ}]: децимация\\
0xF0C1 \textbf{mri7c} [\emph{mri˥ʃ}]: пава\\
0xF0C2 \textbf{kri7c} [\emph{kri˥ʃ}]: сахарин\\
0xF0C4 \textbf{pri7c} [\emph{pri˥ʃ}]: перигей\\
0xF0C6 \textbf{nri7c} [\emph{nri˥ʃ}]: наплыв\\
0xF0C9 \textbf{sri7c} [\emph{sri˥ʃ}]: небиблейский\\
0xF0CA \textbf{tri7c} [\emph{tri˥ʃ}]: печатная машинка\\
0xF0CD \textbf{cri7c} [\emph{ʃri˥ʃ}]: Archy\\
0xF0D0 \textbf{hxi7c} [\emph{ʰxi˥ʃ}]: семьдесят девять\\
0xF0EA \textbf{txi7c} [\emph{txi˥ʃ}]: трезубец\\
0xF101 \textbf{mya7c} [\emph{mjä˥ʃ}]: маршалльский\\
0xF102 \textbf{kya7c} [\emph{kjä˥ʃ}]: икра\\
0xF106 \textbf{nya7c} [\emph{njä˥ʃ}]: отрезанный от внешнего мира\\
0xF108 \textbf{hma7c} [\emph{ʰmä˥ʃ}]: малайский\\
0xF109 \textbf{sya7c} [\emph{sjä˥ʃ}]: половина дядя\\
0xF10A \textbf{tya7c} [\emph{tjä˥ʃ}]: трехъязычный\\
0xF10B \textbf{lya7c} [\emph{ljä˥ʃ}]: индифферентный\\
0xF10D \textbf{cya7c} [\emph{ʃjä˥ʃ}]: шах\\
0xF110 \textbf{hka7c} [\emph{ʰkä˥ʃ}]: Каир\\
0xF118 \textbf{hya7c} [\emph{ʰjä˥ʃ}]: стискивать\\
0xF120 \textbf{hpa7c} [\emph{ʰpä˥ʃ}]: неестественный\\
0xF128 \textbf{hwa7c} [\emph{ʰwä˥ʃ}]: пятьдесят девять\\
0xF129 \textbf{swa7c} [\emph{swä˥ʃ}]: образчик ткани\\
0xF130 \textbf{hna7c} [\emph{ʰnä˥ʃ}]: марочный\\
0xF142 \textbf{ksa7c} [\emph{ksä˥ʃ}]: экспертиза главный\\
0xF144 \textbf{psa7c} [\emph{psä˥ʃ}]: герметизация\\
0xF148 \textbf{hsa7c} [\emph{ʰsä˥ʃ}]: дикари\\
0xF14A \textbf{tsa7c} [\emph{tsä˥ʃ}]: девяносто семь\\
0xF150 \textbf{hta7c} [\emph{ʰtä˥ʃ}]: антидот\\
0xF158 \textbf{hla7c} [\emph{ʰlä˥ʃ}]: лож\\
0xF160 \textbf{hfa7c} [\emph{ʰfä˥ʃ}]: франк tireur\\
0xF162 \textbf{kla7c} [\emph{klä˥ʃ}]: алкалоид\\
0xF164 \textbf{pla7c} [\emph{plä˥ʃ}]: яблочный соус\\
0xF168 \textbf{hca7c} [\emph{ʰʃä˥ʃ}]: перочинный нож\\
0xF169 \textbf{sla7c} [\emph{slä˥ʃ}]: равносторонний\\
0xF16A \textbf{tla7c} [\emph{tlä˥ʃ}]: триангуляция\\
0xF16D \textbf{cla7c} [\emph{ʃlä˥ʃ}]: убить радость\\
0xF170 \textbf{hra7c} [\emph{ʰrä˥ʃ}]: реакционер\\
0xF172 \textbf{gla7c} [\emph{glä˥ʃ}]: галисийский\\
0xF17A \textbf{xla7c} [\emph{xlä˥ʃ}]: специалист по генеалогии\\
0xF188 \textbf{hba7c} [\emph{ʰbä˥ʃ}]: амбразура\\
0xF18A \textbf{tfa7c} [\emph{tfä˥ʃ}]: переводный\\
0xF190 \textbf{hga7c} [\emph{ʰgä˥ʃ}]: ликование\\
0xF198 \textbf{hda7c} [\emph{ʰdä˥ʃ}]: красильщик\\
0xF1A2 \textbf{kca7c} [\emph{kʃä˥ʃ}]: каучук\\
0xF1A4 \textbf{pca7c} [\emph{pʃä˥ʃ}]: паша\\
0xF1AA \textbf{tca7c} [\emph{tʃä˥ʃ}]: пижма\\
0xF1C1 \textbf{mra7c} [\emph{mrä˥ʃ}]: вымачивать\\
0xF1C2 \textbf{kra7c} [\emph{krä˥ʃ}]: панцирь\\
0xF1C4 \textbf{pra7c} [\emph{prä˥ʃ}]: гипертония\\
0xF1C6 \textbf{nra7c} [\emph{nrä˥ʃ}]: рыться\\
0xF1C8 \textbf{hqa7c} [\emph{ʰŋä˥ʃ}]: бачки\\
0xF1C9 \textbf{sra7c} [\emph{srä˥ʃ}]: астральный\\
0xF1CA \textbf{tra7c} [\emph{trä˥ʃ}]: патриархальный\\
0xF1CC \textbf{fra7c} [\emph{frä˥ʃ}]: новая покрышка\\
0xF1D0 \textbf{hxa7c} [\emph{ʰxä˥ʃ}]: надуманным\\
0xF1D4 \textbf{zra7c} [\emph{zrä˥ʃ}]: Захария\\
0xF1DA \textbf{xra7c} [\emph{xrä˥ʃ}]: шасси\\
0xF1E4 \textbf{pxa7c} [\emph{pxä˥ʃ}]: перепев\\
0xF1EA \textbf{txa7c} [\emph{txä˥ʃ}]: таитянского\\
0xF208 \textbf{hmu7c} [\emph{ʰmu˥ʃ}]: конюшни\\
0xF220 \textbf{hpu7c} [\emph{ʰpu˥ʃ}]: голубь косолапый\\
0xF230 \textbf{hnu7c} [\emph{ʰnu˥ʃ}]: слуха\\
0xF248 \textbf{hsu7c} [\emph{ʰsu˥ʃ}]: снегоступ\\
0xF258 \textbf{hlu7c} [\emph{ʰlu˥ʃ}]: Loup Garou\\
0xF270 \textbf{hru7c} [\emph{ʰru˥ʃ}]: мочеполовой\\
0xF288 \textbf{hbu7c} [\emph{ʰbu˥ʃ}]: бушель\\
0xF2D0 \textbf{hxu7c} [\emph{ʰxu˥ʃ}]: с корпусом, скрытым за горизонтом\\
0xF308 \textbf{hme7c} [\emph{ʰme˥ʃ}]: маг\\
0xF320 \textbf{hpe7c} [\emph{ʰpe˥ʃ}]: молочай\\
0xF348 \textbf{hse7c} [\emph{ʰse˥ʃ}]: мозжечок\\
0xF34A \textbf{tse7c} [\emph{tse˥ʃ}]: проверенный временем\\
0xF350 \textbf{hte7c} [\emph{ʰte˥ʃ}]: торговец скобяными изделиями\\
0xF358 \textbf{hle7c} [\emph{ʰle˥ʃ}]: Меланезия\\
0xF370 \textbf{hre7c} [\emph{ʰre˥ʃ}]: оружейник\\
0xF3A0 \textbf{hze7c} [\emph{ʰze˥ʃ}]: Z\\
0xF3AA \textbf{tce7c} [\emph{tʃe˥ʃ}]: амбулаторный\\
0xF3C9 \textbf{sre7c} [\emph{sre˥ʃ}]: стригальщик\\
0xF3CA \textbf{tre7c} [\emph{tre˥ʃ}]: тетралогия\\
0xF3D0 \textbf{hxe7c} [\emph{ʰxe˥ʃ}]: пролив\\
0xF408 \textbf{hmo7c} [\emph{ʰmo˥ʃ}]: омега\\
0xF410 \textbf{hko7c} [\emph{ʰko˥ʃ}]: картуш\\
0xF448 \textbf{hso7c} [\emph{ʰso˥ʃ}]: соте\\
0xF44A \textbf{tso7c} [\emph{tso˥ʃ}]: непримиримость\\
0xF450 \textbf{hto7c} [\emph{ʰto˥ʃ}]: три загнаны\\
0xF458 \textbf{hlo7c} [\emph{ʰlo˥ʃ}]: хлорал\\
0xF468 \textbf{hco7c} [\emph{ʰʃo˥ʃ}]: зашикать\\
0xF470 \textbf{hro7c} [\emph{ʰro˥ʃ}]: конский волос\\
0xF4C1 \textbf{mro7c} [\emph{mro˥ʃ}]: аменорея\\
0xF4C9 \textbf{sro7c} [\emph{sro˥ʃ}]: сербско хорватской\\
0xF801 \textbf{myi2c} [\emph{mji˩ʃ}]: иммунология\\
0xF808 \textbf{hmi2c} [\emph{ʰmi˩ʃ}]: криминолог\\
0xF809 \textbf{syi2c} [\emph{sji˩ʃ}]: социально-экономический\\
0xF810 \textbf{hki2c} [\emph{ʰki˩ʃ}]: сожительство\\
0xF820 \textbf{hpi2c} [\emph{ʰpi˩ʃ}]: патчи\\
0xF828 \textbf{hwi2c} [\emph{ʰwi˩ʃ}]: приговоренный к пожизненному заключению\\
0xF830 \textbf{hni2c} [\emph{ʰni˩ʃ}]: Нил\\
0xF844 \textbf{psi2c} [\emph{psi˩ʃ}]: эпизодический\\
0xF84A \textbf{tsi2c} [\emph{tsi˩ʃ}]: самопоглощение\\
0xF850 \textbf{hti2c} [\emph{ʰti˩ʃ}]: покатый\\
0xF858 \textbf{hli2c} [\emph{ʰli˩ʃ}]: несклонный\\
0xF862 \textbf{kli2c} [\emph{kli˩ʃ}]: микология\\
0xF864 \textbf{pli2c} [\emph{pli˩ʃ}]: поли\\
0xF868 \textbf{hci2c} [\emph{ʰʃi˩ʃ}]: Wishy водянистый\\
0xF869 \textbf{sli2c} [\emph{sli˩ʃ}]: сингалец\\
0xF86A \textbf{tli2c} [\emph{tli˩ʃ}]: астрология\\
0xF86C \textbf{fli2c} [\emph{fli˩ʃ}]: френология\\
0xF870 \textbf{hri2c} [\emph{ʰri˩ʃ}]: гребни\\
0xF888 \textbf{hbi2c} [\emph{ʰbi˩ʃ}]: опадение\\
0xF88A \textbf{tfi2c} [\emph{tfi˩ʃ}]: зяблик\\
0xF890 \textbf{hgi2c} [\emph{ʰgi˩ʃ}]: часовщик\\
0xF898 \textbf{hdi2c} [\emph{ʰdi˩ʃ}]: дизентерия\\
0xF8A2 \textbf{kci2c} [\emph{kʃi˩ʃ}]: осевой\\
0xF8A4 \textbf{pci2c} [\emph{pʃi˩ʃ}]: зяблики\\
0xF8AA \textbf{tci2c} [\emph{tʃi˩ʃ}]: демерредж\\
0xF8C2 \textbf{kri2c} [\emph{kri˩ʃ}]: неразборчивость\\
0xF8C6 \textbf{nri2c} [\emph{nri˩ʃ}]: приближается к\\
0xF8C9 \textbf{sri2c} [\emph{sri˩ʃ}]: придыхательный\\
0xF8CA \textbf{tri2c} [\emph{tri˩ʃ}]: интровертированный\\
0xF8EA \textbf{txi2c} [\emph{txi˩ʃ}]: триптих\\
0xF908 \textbf{hma2c} [\emph{ʰmä˩ʃ}]: пастушьи\\
0xF909 \textbf{sya2c} [\emph{sjä˩ʃ}]: осетия\\
0xF90A \textbf{tya2c} [\emph{tjä˩ʃ}]: анте Никейский\\
0xF90B \textbf{lya2c} [\emph{ljä˩ʃ}]: все звезды\\
0xF90E \textbf{rya2c} [\emph{rjä˩ʃ}]: рикша\\
0xF910 \textbf{hka2c} [\emph{ʰkä˩ʃ}]: инквизиция\\
0xF918 \textbf{hya2c} [\emph{ʰjä˩ʃ}]: уравнитель\\
0xF92A \textbf{twa2c} [\emph{twä˩ʃ}]: останец\\
0xF930 \textbf{hna2c} [\emph{ʰnä˩ʃ}]: аннуитет\\
0xF942 \textbf{ksa2c} [\emph{ksä˩ʃ}]: за кулисы\\
0xF944 \textbf{psa2c} [\emph{psä˩ʃ}]: восемьдесят семь\\
0xF948 \textbf{hsa2c} [\emph{ʰsä˩ʃ}]: немытый\\
0xF94A \textbf{tsa2c} [\emph{tsä˩ʃ}]: кипятильник\\
0xF950 \textbf{hta2c} [\emph{ʰtä˩ʃ}]: даты\\
0xF958 \textbf{hla2c} [\emph{ʰlä˩ʃ}]: пресный\\
0xF960 \textbf{hfa2c} [\emph{ʰfä˩ʃ}]: лихорадочный\\
0xF964 \textbf{pla2c} [\emph{plä˩ʃ}]: небный\\
0xF968 \textbf{hca2c} [\emph{ʰʃä˩ʃ}]: хорошо советовали\\
0xF969 \textbf{sla2c} [\emph{slä˩ʃ}]: скотство\\
0xF96A \textbf{tla2c} [\emph{tlä˩ʃ}]: желтофиоль садовая\\
0xF970 \textbf{hra2c} [\emph{ʰrä˩ʃ}]: арбитраж\\
0xF97A \textbf{xla2c} [\emph{xlä˩ʃ}]: хадж\\
0xF988 \textbf{hba2c} [\emph{ʰbä˩ʃ}]: бакалавр\\
0xF98A \textbf{tfa2c} [\emph{tfä˩ʃ}]: стагфляция\\
0xF990 \textbf{hga2c} [\emph{ʰgä˩ʃ}]: бородка\\
0xF998 \textbf{hda2c} [\emph{ʰdä˩ʃ}]: в середине марта\\
0xF9A2 \textbf{kca2c} [\emph{kʃä˩ʃ}]: обналичить\\
0xF9AA \textbf{tca2c} [\emph{tʃä˩ʃ}]: гашиш\\
0xF9C1 \textbf{mra2c} [\emph{mrä˩ʃ}]: Микронезия\\
0xF9C2 \textbf{kra2c} [\emph{krä˩ʃ}]: камеры\\
0xF9C4 \textbf{pra2c} [\emph{prä˩ʃ}]: черная работа\\
0xF9C6 \textbf{nra2c} [\emph{nrä˩ʃ}]: мучение нерва\\
0xF9CA \textbf{tra2c} [\emph{trä˩ʃ}]: переросший\\
0xF9D0 \textbf{hxa2c} [\emph{ʰxä˩ʃ}]: хорошее общение\\
0xF9D1 \textbf{bra2c} [\emph{brä˩ʃ}]: карбюратор\\
0xF9D3 \textbf{dra2c} [\emph{drä˩ʃ}]: бесправных\\
0xF9DA \textbf{xra2c} [\emph{xrä˩ʃ}]: молочница\\
0xF9EA \textbf{txa2c} [\emph{txä˩ʃ}]: анти герой\\
0xFA10 \textbf{hku2c} [\emph{ʰku˩ʃ}]: капуцин\\
0xFA50 \textbf{htu2c} [\emph{ʰtu˩ʃ}]: ботик\\
0xFB08 \textbf{hme2c} [\emph{ʰme˩ʃ}]: манихей\\
0xFB30 \textbf{hne2c} [\emph{ʰne˩ʃ}]: этнология\\
0xFB48 \textbf{hse2c} [\emph{ʰse˩ʃ}]: эстетически\\
0xFB4A \textbf{tse2c} [\emph{tse˩ʃ}]: арендодатель\\
0xFB50 \textbf{hte2c} [\emph{ʰte˩ʃ}]: транш\\
0xFB58 \textbf{hle2c} [\emph{ʰle˩ʃ}]: Г-образный\\
0xFB64 \textbf{ple2c} [\emph{ple˩ʃ}]: депиляция\\
0xFB6A \textbf{tle2c} [\emph{tle˩ʃ}]: внутриклеточный\\
0xFB70 \textbf{hre2c} [\emph{ʰre˩ʃ}]: в форме полумесяца\\
0xFBC9 \textbf{sre2c} [\emph{sre˩ʃ}]: замена\\
0xFBCA \textbf{tre2c} [\emph{tre˩ʃ}]: экстровертность\\
0xFC10 \textbf{hko2c} [\emph{ʰko˩ʃ}]: кошерный\\
0xFC30 \textbf{hno2c} [\emph{ʰno˩ʃ}]: нанотехнология\\
0xFC48 \textbf{hso2c} [\emph{ʰso˩ʃ}]: рустовка\\
0xFC50 \textbf{hto2c} [\emph{ʰto˩ʃ}]: трудоемкими\\
0xFC6A \textbf{tlo2c} [\emph{tlo˩ʃ}]: этология\\
0xFC70 \textbf{hro2c} [\emph{ʰro˩ʃ}]: Роджер\\

\end{multicols}

\end{document}

